\documentclass[11pt,oneside]{article}
\usepackage[margin=1in]{geometry}
\usepackage{amsmath,amssymb,amsthm,mathtools,booktabs}
\usepackage[colorlinks=true,linkcolor=blue,urlcolor=blue]{hyperref}
\newtheorem{theorem}{Theorem}[section]
\newtheorem{proposition}[theorem]{Proposition}
\newtheorem{record}[theorem]{Record}
\newtheorem{firewall}[theorem]{Claim firewall}
\DeclareMathOperator{\SU}{SU}
\DeclareMathOperator{\Tr}{Tr}
\title{Renewal Yang--Mills Mass-Gap Research Notes\\
Post-Fifteenth-Archive Compact Lineage, Version V89}
\author{Continuation with a scope-checked result ledger}
\date{Research continuation}
\begin{document}
\maketitle
\tableofcontents

\section{Context, restart, and the unchanged objective}
\label{f16v1:context}

The objective is a rigorous nontrivial interacting four-dimensional
continuum Yang--Mills theory with a positive physical Hamiltonian mass
gap. That objective remains unproved. We currently work with an exact
compact \(\SU(2)\) Wilson transfer, its gauge constraints, and particular
weak-field limits. Auxiliary Gaussian separation, a lattice transfer
ratio, and a calibrated physical continuum mass are different claims.

The hundred-version fifteenth lineage is preserved in
\begin{center}\small
\path{Renewal_Yang_Mills_Mass_Gap_Research_Notes_V100_f15.tex}.
\end{center}
This is a compact restart, not a replacement of that archive.
References such as f15 V99 specify sections of that file.
Prior hypotheses and corrections remain binding. The summary below
is a result map, not a fresh certification of every historical claim.

The required V100 checkpoint reread SM continuum V72, NS stability V26,
shell-mixing V16, parallel YM V5, SM source-selection V31,
regenerative stationarity V34 and exact-action Einstein bridge V24.
Their current versions and hashes were unchanged. No companion
provides the missing volume-uniform native YM spectral comparison.
The useful cross-program lesson is to control densities or local
quantities rather than demand a volume-independent tail for an
extensive count. The detailed checkpoint is in f15 V100.
The next companion review is this lineage's V5; its next normal
archive/restart is V100.

The recent upstream change is to retain the normalized magnetic
image of the cube carrier rather than freeze its internal variables
or project its image straight back to the old ground band.
The present note continues by computing the exact overlap of those
conditional vectors across a native bridge kinetic step.
Its weak-field response is not an invariant-subspace theorem:
even the resulting Gaussian compressed family fails the semigroup
law at finite time. We keep that discrepancy explicit.

\section{Major interfaces retained from f15}
\label{f16v1:ledger}

\begin{center}\small
\begin{tabular}{@{}p{0.15\textwidth}p{0.79\textwidth}@{}}
\toprule
Range & Established objects and scope limits\\
\midrule
V1--V11 &
Full compact source/entropy estimates, literal blocking, growing-depth
curvature comparisons and scale-correct pinned-fibre interfaces.
They do not give physical continuum coercivity.\\[1mm]
V12--V37 &
Local chart and source-response comparisons with normalization,
exceptional-set and entropy costs retained; obstructions to dropping
those costs. No unqualified extensive-source mixing theorem.\\[1mm]
V38--V56 &
Native compact conditional landscapes, coherent boundary minima,
reflected source laws and paid pressure comparisons.
Rare or low-action regions are not automatically harmless for spectra.\\[1mm]
V57--V65 &
Positive Wilson transfer and vacuum response estimates in their
selected regimes; soft microscopic spectral limits.
A collapsing one-step lattice gap is not a disproof of a positive
gap after independently calibrated physical time scaling.\\[1mm]
V66--V81 &
Exact decimation/conditional laws, local Dirichlet and bad-region
control, full Gaussian invariant sectors, and vacuum-independent
transfer/compound-operator certificates.
A local gap or a finite matrix check is not the full native gap.\\[1mm]
V82--V91 &
Exact Gauss-boundary charge decomposition, magnetic charge walks,
charged-sector domination, compact endpoint channels and paid
spectral tails. Literature and undecidability audits import
methods, not an unverified mass-gap conclusion.\\[1mm]
V92--V97 &
Complete native coupled finite-region kernels and their growing
boundary carriers; full internal oscillator spectra and
finite-volume complement approximations.\\[1mm]
V98--V100 &
Actual spatial assembly, a sharp auxiliary coarse electric edge,
exact magnetic Gram/dressing, a fixed-volume Gaussian response,
and extensive old-oscillator excitation budgets.
Volume-uniform native dressed dynamics remain open.\\
\bottomrule
\end{tabular}
\end{center}

\subsection{The exact spatial assembly and its unresolved dynamics}

Let the spatial torus have side \(2m\), \(m\ge4\), with
\(B=m^3\) disjoint vertex cubes. A link is internal when its base
coordinate in its own direction is even; otherwise it is a bridge.
There are \(12B\) of each. Of the \(24B\) plaquettes, \(6B\)
are internal, \(12B\) cross one block interface and \(6B\) cross
in both directions. The coarse graph has four parallel bridges
between each pair of positive-direction neighbouring blocks.

A rooted spanning tree in each cube gives exact Haar coordinates
\(g\in G^{7B}\), \(X\in G^{5B}\), \(Y\in G^{12B}\), \(G=\SU(2)\).
After nonroot Gauss reduction the whole physical Hilbert space is
\begin{equation}
 \mathcal H_{\rm phys}
       =L^2(G^{5B}\times G^{12B})^{G^B}.
 \label{f16v1:physical}
\end{equation}
The root group conjugates its internal chords and acts at all its
bridge endpoints. Mixed plaquette words are exactly the two-bridge,
two-internal-edge words or the four-bridge words from the original
lattice. No plaquette is discarded.

For \(s(U)=1-\frac12\Tr U\), define the normalized central convolution
kernel and bridge operator
\begin{equation}
 k_\gamma(U)=\frac{e^{-\gamma s(U)}}{z_\gamma},\quad
 z_\gamma=\int_Ge^{-\gamma s(U)}\,dH(U),\quad
 \mathcal C_\gamma=\bigotimes_{e\ {\rm bridge}} C_\gamma^{(e)}.
 \label{f16v1:kinetic}
\end{equation}
These are the Wilson kinetic kernels, not exactly heat kernels.
The full native transfer has the exact factorization
\begin{equation}
 T_{\rm full,\gamma}
 =M_\times\left[
       \left(\bigotimes_b T_{\Box,\gamma}^{(b)}\right)
                       \otimes\mathcal C_\gamma\right]M_\times,
 \quad M_\times=e^{-\gamma S_\times/2},
 \label{f16v1:full}
\end{equation}
where \(S_\times\) sums all \(18B\) mixed plaquette actions.
The regional transfer factors and the bridge factor commute before
the magnetic sandwich, but the magnetic factor cannot be omitted.

The ground-carrier compression of the cut transfer in f15 V98
has coarse electric operator
\[
 \mathcal D=\mathcal E_{\rm bridge}
       +\sum_b\sum_{v,w}(L_\Box^+)_{vw}J_{bv}\cdot J_{bw},
 \qquad \mathcal E_{\rm bridge}\le\mathcal D
                         \le4\mathcal E_{\rm bridge}.
\]
On the coarse physical space its first nonzero eigenvalue is \(19/2\).
This is an auxiliary operator. Its proved native growing-carrier
error \(O_R(B/\sqrt\gamma)\) is larger than the \(1/\gamma\)
scale of that auxiliary heat edge. The magnetic image contributes
an additional order-one effect. No full-transfer gap follows.

\subsection{The exact normalized carrier and the leading response}

Here are sufficient definitions to continue without reprinting the
finite-region proofs. In one cube let \(\mathsf B\) be the
root-deleted edge incidence matrix, with row \(z_t-z_s\);
let \(\mathsf J\) inject its five chord coordinates; and let
\(\mathsf Z\mathsf B=0,\ \mathsf Z\mathsf J=I\) be the cycle matrix.
Use the rooted axial tree of f15 V97. Put
\[
 \mathsf F=-(\mathsf B^{\mathsf T}\mathsf B)^{-1}
                     \mathsf B^{\mathsf T}\mathsf J,\quad
 \mathsf M=\mathsf Z\mathsf Z^{\mathsf T},\quad
 U=\mathsf J+\mathsf B\mathsf F=\mathsf Z^{\mathsf T}\mathsf M^{-1}.
\]
Let \(\mathsf H\) be the Gram of the six linearized cube faces and
\(T=\mathsf M^{1/2}\mathsf H\mathsf M^{1/2}\).
The spectrum of \(T\) is \(4,4,4,6,6\).
The covariance of one colour component of the squared internal
Gaussian ground vector is
\begin{equation}
 C=\mathsf M^{1/2}[2\sqrt{T+T^2/4}]^{-1}\mathsf M^{1/2},
 \qquad C_B=\bigoplus_b C.
 \label{f16v1:covariance}
\end{equation}
For \(X_i=\operatorname{Exp}(x_i/\sqrt\gamma)\) on the complete
chord charts, let \(\Phi_\gamma(X)\) be the positive normalized
product Gaussian vector with its Haar square-root Jacobian divided
out, and with its chart restriction normalized exactly. Its
squared \(x\)-density is the corresponding truncated Gaussian.

Set \(h_{bu}=\operatorname{Exp}((\mathsf F x_b)_u/\sqrt\gamma)\)
at nonroot ports and \(h_{bo}=I\). The moving bridge coordinates are
\[
                    (\Theta_XY)_e=h_s^{-1}Y_eh_t.
\]
At fixed \(X\) this is Haar preserving. Define
\begin{align}
 \widehat S_\gamma f(X,Y)
       &=\Phi_\gamma(X)f(\Theta_XY),\notag\\
 F_{\gamma,t}(Z)
       &=\int|\Phi_\gamma(X)|^2
              M_\times(X,\Theta_X^{-1}Z)^t\,dH^{5B}(X).
 \label{f16v1:gram}
\end{align}
The identities
\[
 \widehat S_\gamma^*M_\times^t\widehat S_\gamma=M_{F_{\gamma,t}},
 \qquad
 \mathcal I_\gamma=M_\times\widehat S_\gamma
                                    M_{F_{\gamma,2}^{-1/2}},
 \qquad \mathcal I_\gamma^*\mathcal I_\gamma=I
\]
are exact on the coarse physical carrier. At each finite regulator
\(F_{\gamma,2}>0\) is bounded below. It is \(F_{\gamma,2}\), not
\(F_{\gamma,1}^2\), that normalizes the magnetic image.

Let \(C_{\rm int}\) and \(D\) be the signed mixed-plaquette incidence
matrices on internal and bridge edges. In the moving frame the
internal linear embedding is \(\bigoplus_bU\), so put
\[
 A=C_{\rm int}\bigoplus_bU,\qquad
 K=AC_BA^{\mathsf T},\quad
 L=C_B^{1/2}A^{\mathsf T}AC_B^{1/2},\quad \rho=1/\sqrt2 .
\]
F15 V99 proves
\[
 0\le K,\quad\|K\|\le\rho<1,\quad
 \Tr K=B\left(\frac3{2\sqrt2}+\frac2{\sqrt{15}}\right)>B.
\]
With \(Z_e=\operatorname{Exp}(y_e/\sqrt\gamma)\),
\begin{equation}
 F_{\gamma,t}(Z)\longrightarrow
 F_t(y)=\det(I+(t/2)K)^{-3/2}
     e^{-(t/4)y^{\mathsf T}
       [D^{\mathsf T}(I+(t/2)K)^{-1}D\otimes I_3]y},
 \label{f16v1:limit}
\end{equation}
uniformly on bounded \(y\)-windows at fixed \(B\).
The proved absolute error is \(O_t(B(1+R^3)/\sqrt\gamma)
+O_t(Be^{-c\gamma})\). It is not a relative bound uniform in volume.

The normalized conditional limit of \(\mathcal I_\gamma\) is the
positive square root of a Gaussian with covariance and mean
\begin{equation}
 C_{\rm p}=(C_B^{-1}+A^{\mathsf T}A)^{-1},\qquad
 m_y=-(C_{\rm p}A^{\mathsf T}D\otimes I_3)y .
 \label{f16v1:posterior}
\end{equation}
The inverse has a block-distance convergent Neumann expansion:
for block sets \(E,F\) at distance \(d\),
\[
 \|P_EC_{\rm p}P_F\|\le\|C\|\rho^d/(1-\rho).
\]
The leading log determinant also has a uniformly summable
per-block closed-walk expansion. These are Gaussian response
statements, not a native non-Gaussian polymer theorem.

F15 V100 counts all internal oscillator levels in this dressed
vector. It proves exponentially small capture by every
\(o(B)\) total-degree cutoff at zero bridge field and an
exponentially accurate cutoff of order
\(B+\|Dy\|^2+\log\delta^{-1}\). Its compact projection handoff
takes \(\gamma\to\infty\) at each fixed \(B\); it does not control
the norm of the full excited-state transfer complement.

\section{New V1 result: exact kinetic compression by conditional affinity}
\label{f16v1:affinity}

For the derivation it is convenient to work first before root Gauss
restriction, on \(L^2(G^{12B})\) and
\(L^2(G^{5B}\times G^{12B})\) in \((X,Z)\) coordinates.
All maps below intertwine the root action, so their exact
operator identities restrict to \eqref{f16v1:physical}.
Define normalized conditional vectors
\begin{equation}
 a_{\gamma,Z}(X)
   =\frac{\Phi_\gamma(X)M_\times(X,\Theta_X^{-1}Z)}
                    {\sqrt{F_{\gamma,2}(Z)}},
 \qquad \int|a_{\gamma,Z}(X)|^2\,dH^{5B}(X)=1,
 \label{f16v1:conditional}
\end{equation}
and their affinity
\[
 \mathcal A_\gamma(Z,Z')
       =\int a_{\gamma,Z}(X)a_{\gamma,Z'}(X)\,dH^{5B}(X).
\]

\begin{theorem}[Exact native bridge kernel on the dressed carrier]
\label{f16v1:exact-kernel}
The compression of the bridge kinetic factor alone is
\begin{equation}
 \bigl(\mathcal I_\gamma^*\mathcal C_\gamma
                      \mathcal I_\gamma\bigr)(Z,Z')
   =\left[\prod_{e\ {\rm bridge}}
             k_\gamma(Z_eZ_e'^{-1})\right]\mathcal A_\gamma(Z,Z').
 \label{eq:f16v1:exact-kernel}
\end{equation}
The affinity is symmetric, positive semidefinite as a Gram kernel,
continuous, and satisfies \(0<\mathcal A_\gamma\le1\) and
\(\mathcal A_\gamma(Z,Z)=1\).
The compressed operator is a positive contraction.
The root gauge invariance is simultaneous in \(Z,Z'\), not
a claim that the affinity is separately invariant in each argument.
\end{theorem}

\begin{proof}
The bridge operator leaves \(X\) fixed. Under
\(Y_e=h_s Z_eh_t^{-1}\) and \(Y_e'=h_s Z_e'h_t^{-1}\),
\[
                  Y_eY_e'^{-1}
                       =h_s Z_eZ_e'^{-1}h_s^{-1}.
\]
Centrality of \(k_\gamma\) and Haar invariance remove both frame
factors. In these coordinates \(\mathcal I_\gamma f=a_{\gamma,Z}f(Z)\).
Integrating the internal coordinate in its compression gives
\eqref{eq:f16v1:exact-kernel}. The Gram assertion and the upper
bound are respectively the squared-norm identity and Cauchy--Schwarz.
Positivity follows also directly from compression of the positive
Wilson convolution. Its norm is at most one.

At finite \(\gamma,B\), the magnetic multiplier is strictly positive
and bounded, \(F_{\gamma,2}\) has a positive lower bound, and
\(\Phi_\gamma\) has norm one. Dominated convergence proves
continuity and strict positivity of the affinity. Equivariance
of the original carrier and kernel proves the root-action
statement and the restriction to invariant functions.
\end{proof}

This identity keeps the whole compact bridge configuration space.
It does not replace \eqref{f16v1:full} by its bridge factor.
In particular the cube kinetic dynamics and their internal
excited-state returns have not been integrated by this theorem.

\subsection{The leading affinity and its exact null directions}

In the next formulas the colour tensor \(I_3\) is suppressed:
when \(R\) acts on a \(36B\)-component bridge vector, it means
\(R\otimes I_3\).
Set
\begin{align}
 R&=D^{\mathsf T}A C_{\rm p}A^{\mathsf T}D
    =D^{\mathsf T}K(I+K)^{-1}D
    =D^{\mathsf T}D-H_{\rm eff},\notag\\
 H_{\rm eff}&=D^{\mathsf T}(I+K)^{-1}D .
 \label{f16v1:R}
\end{align}

\begin{theorem}[Conditional overlap and native local-kernel limit]
\label{f16v1:gaussian-affinity}
For the normalized Gaussian amplitudes \(\psi_y\) with squared
density \eqref{f16v1:posterior},
\begin{equation}
 \langle\psi_y,\psi_z\rangle
             =e^{-(y-z)^{\mathsf T}R(y-z)/8}.
 \label{eq:f16v1:gaussian-affinity}
\end{equation}
At fixed \(B\), \(\mathcal A_\gamma\) in rescaled bridge charts
converges to this expression uniformly on bounded pairs of
bridge windows. After including both square-root Haar chart
densities, the kernel \eqref{eq:f16v1:exact-kernel} converges there to
\begin{equation}
 (2\pi)^{-q/2}
   \exp\left[-\frac{|y-z|^2}{2}
                    -\frac{(y-z)^{\mathsf T}R(y-z)}8\right],
       \qquad q=36B.
 \label{f16v1:native-leading-kernel}
\end{equation}
Thus compression to any fixed bounded coordinate window
converges in Hilbert--Schmidt norm. No global norm or
volume-uniform spectral convergence is asserted.
\end{theorem}

\begin{proof}
Completing the common square of two Gaussian amplitudes with
covariance \(C_{\rm p}\) gives the overlap
\[
       \exp[-(m_y-m_z)^{\mathsf T}C_{\rm p}^{-1}(m_y-m_z)/8].
\]
Substitute \eqref{f16v1:posterior}; Woodbury inversion gives
\eqref{f16v1:R}. F15 V99's normalized-density convergence,
or f15 V100's \(L^2\) amplitude argument, gives the uniform
convergence of the affinities at fixed \(B\).

For completeness, in one bridge chart with Jacobian \(J_\gamma\),
ordinary Taylor expansion and Laplace normalization give
\[
 [J_\gamma(y)J_\gamma(z)]^{1/2}
 k_\gamma(\operatorname{Exp}(y/\sqrt\gamma)
                      \operatorname{Exp}(-z/\sqrt\gamma))
       \longrightarrow(2\pi)^{-3/2}e^{-|y-z|^2/2}
\]
uniformly on compact pairs. Here
\(s(\operatorname{Exp}v)=1-\cos|v|\);
\(\gamma s\) tends to \(|y-z|^2/2\).
The normalization follows by the same rescaling in
\(z_\gamma\); the inequality \(1-\cos r\ge2r^2/\pi^2\)
for \(0\le r\le\pi\) dominates the full chart integral.
The Haar Jacobian cancels to the displayed Euclidean prefactor.
Multiply over the fixed \(12B\) bridges and use
\eqref{eq:f16v1:gaussian-affinity}. Uniform convergence on
a bounded window implies square-integral kernel convergence.
This is a kinematic chart assertion before root averaging;
it does not identify a physical spectral subspace with that window.
\end{proof}

One has
\begin{equation}
 0\le R\le\frac{\rho}{1+\rho}D^{\mathsf T}D.
 \label{f16v1:Rbound}
\end{equation}
The covariance locality from \eqref{f16v1:posterior} and the
finite-range incidence factors give exponentially decaying
block entries of \(R\), uniformly in volume.

There is a stronger exact statement than simply
\(\ker D\subseteq\ker R\). Let \(\mathcal E\) consist of bridge
vectors constant on each set of four parallel bridges.
Then
\begin{equation}
 A^{\mathsf T}D y=0,\qquad R y=0,\qquad
 H_{\rm eff}y=D^{\mathsf T}Dy,\qquad \psi_y=\psi_0
                    \quad (y\in\mathcal E).
 \label{f16v1:common-bridge}
\end{equation}
Indeed a decorated digon takes the difference of two such
parallel values, so its curvature vanishes. A quadrilateral
has no internal edge, so the corresponding row of \(A\) is
zero. This proves the first identity; the others follow.
More generally \(\psi_{y+v}=\psi_y\) for \(v\in\mathcal E\).
The scalar magnetic normalizer still depends on the coarse
quadrilateral curvature; it has not disappeared.

In particular the transverse common-bridge cosine used in
f15 V99 has exactly
\[
       \frac{\langle y,H_{\rm eff}y\rangle}{\|y\|^2}
                           =2\sin^2(\pi/m),
\]
not merely the upper bound established there.
The leading fast conditional vector cannot generate a restoring
term in these coarse common-bridge directions.
This concerns a static linearized Hessian, not a family of
gapless native physical Hamiltonian states.

\section{Finite kinetic time is not fixed by the static normalization}
\label{f16v1:time}

The Gaussian limit allows a complete check of this issue.
Let \(G\) be any positive definite \(q\)-by-\(q\) constant matrix,
and let \(H_t^G\) act on the bridge variable, leaving the internal
variable fixed, with kernel
\[
 h_t^G(y-z)
 =(4\pi t)^{-q/2}\det(G)^{-1/2}
        e^{-(y-z)^{\mathsf T}G^{-1}(y-z)/(4t)}.
\]
The isometry \(\mathcal J f(u,y)=\psi_y(u)f(y)\) compresses it
to \(\mathcal B_t=\mathcal J^*H_t^G\mathcal J\).
The actual one-step local limit \eqref{f16v1:native-leading-kernel}
is the case \(G=I,t=1/2\).

\begin{theorem}[Exact finite-time diffusion change and its composition defect]
\label{f16v1:heat}
For every \(t>0\),
\begin{equation}
 \mathcal B_t=c_t H_t^{G_t},\qquad
 G_t=(G^{-1}+tR/2)^{-1},\qquad
 c_t=\det(I+(t/2)G^{1/2}RG^{1/2})^{-1/2}.
 \label{f16v1:heat-factor}
\end{equation}
Its operator norm on \(L^2(\mathbb R^q)\) is \(c_t\).
If \(R\ne0\), the norm-normalized family
\(c_t^{-1}\mathcal B_t\) is not a semigroup.
\end{theorem}

\begin{proof}
Multiply \(h_t^G(y-z)\) by the affinity
\eqref{eq:f16v1:gaussian-affinity}.
The exponent becomes
\(-(y-z)^{\mathsf T}(G^{-1}+tR/2)(y-z)/(4t)\).
Comparing Gaussian prefactors gives \(c_t\).
The Fourier multiplier is
\(c_t e^{-t p^{\mathsf T}G_t p}\), whose essential supremum is
\(c_t\).

For \(R\ne0\), positivity and inversion show
\(G_t-G_{2t}\) is nonzero positive semidefinite.
Thus for some \(p\),
\[
 e^{-2t p^{\mathsf T}G_t p}
                   \ne e^{-2t p^{\mathsf T}G_{2t}p}.
\]
These are the multipliers of
\((c_t^{-1}\mathcal B_t)^2\) and
\(c_{2t}^{-1}\mathcal B_{2t}\). Continuity in \(p\) makes this
a genuine operator inequality, proving the failure of the
semigroup law.
\end{proof}

The nonzero case actually occurs in the assembly.
For example take a bridge perturbation supported only on
the positive first-direction bridge from fine vertex \((1,0,0)\).
In its incident block, \(A^{\mathsf T}D\) applied to this vector
is nonzero. To see this without a numerical rank test, its
internal-edge source is, up to overall sign, the sum of the
second- and third-direction edges starting at port \((1,0,0)\).
The first/second-direction cube face at third coordinate zero
contains just the former edge, with nonzero signed coefficient.
Thus the source has nonzero circulation. Since
\(\ker U^{\mathsf T}\) is the cut space, its image under
\(U^{\mathsf T}\) cannot vanish. Positivity of \(C_{\rm p}\)
then implies \(R\ne0\).
The exact incidence checker supplies this literal relative-port
example on sides eight and twelve.

The failure is also transparent before taking limits:
\(\mathcal J^*H_s^G H_t^G\mathcal J\) differs from
\(\mathcal J^*H_s^G\mathcal J\mathcal J^*H_t^G\mathcal J\)
by the omitted intermediate internal complement.
A static isometry alone cannot remove that return.

\subsection{The infinitesimal check and its limited physical meaning}

There is a useful independent derivative check.
For bridge indices \(i,j\), real positivity and normalization give
\[
 \langle\psi_y,\partial_i\psi_y\rangle=0,\qquad
 \langle\partial_i\psi_y,\partial_j\psi_y\rangle=R_{ij}/4.
\]
Differentiate the Gaussian mean in
\eqref{f16v1:posterior} to verify the second identity.
For Schwartz \(f\), Fubini's theorem and this zero cross term yield
\begin{equation}
 \int \bigl|\nabla_y(\psi_y f)\bigr|_G^2\,du\,dy
  =\int|\nabla f|_G^2\,dy
                     +\frac14\Tr(GR)\int|f|^2\,dy .
 \label{f16v1:metric}
\end{equation}
Thus the compressed differential quadratic form has a constant
scalar addition, not a bridge-dependent mass potential.
It agrees with the derivative of \eqref{f16v1:heat-factor}:
\[
       \left.\frac d{dt}\log c_t\right|_{t=0}
                           =-\frac14\Tr(GR),\qquad G_0=G.
\]
Exponentiating that compressed differential form is nevertheless
not equal to the finite-time compression in
\eqref{f16v1:heat-factor}. The latter has \(G_t\), not \(G\).
For \(G=I\), the scalar addition is a vacuum normalization of
order at most \(B\), using \eqref{f16v1:Rbound} and the bounded
plaquette incidence. None of these chart statements fixes the
physical lattice time unit or the full interacting transfer gap.

\section{What must be controlled next}
\label{f16v1:next}

The direction is now more precise. The exact native affinity
in \eqref{eq:f16v1:exact-kernel} retains every bridge configuration
and its correct magnetic normalization. Its leading response
has a computable diffusion correction on relative-port directions,
while the fast conditional state is blind to all common-bridge
directions at this order. Repeated ground projection, even after
normalization, discards a definite intermediate return.

The next required calculation is the joint dressed compression
with the cube transfer factors in \eqref{f16v1:full}, retaining
their excited internal bands and the noncommuting moving frame
at both time ends. It must bound the discarded return relative
to the same full native top eigenvalue, uniformly in spatial
volume, or retain it through a controlled memory/Schur term.
A new scalar normalization or another fixed-volume Gaussian
determinant does not supply that estimate.

Beyond it remain nonlinear compact collective-mode coercivity,
bridge-field tail control, an independent physical-time calibration,
and nontrivial interacting continuum reconstruction. These are
requirements of the unchanged objective, not assumptions being
declared solved by the restart.

\section{Verification and preservation}
\label{f16v1:verification}

The diagnostic
\path{scripts/verify_ym_f16_v1_dressed_affinity.py}
checks independent noncommuting matrix completions,
the affinity and derivative metric, finite-time heat precisions,
the composition defect, a direct finite conditional-kernel
compression, and the complete mixed-plaquette incidence on
two actual spatial tori. It checks the common-bridge null
response and an explicit nonzero relative-port response.
It also reruns the inherited f15 V100 diagnostic chain.
All 10,918 new exact fixtures passed, together with the inherited
f15 V100 chain. These finite fixtures supplement proofs, not replace them.
The checker has 5,600 bytes and SHA--256
\begin{center}\scriptsize\ttfamily
CCBCBC50A687E88B74465051A0C7E9B283891911ADA928D910FA5C896B3EFC7B.
\end{center}

F15 V100 is preserved as an archive, with V99's body recoverable
byte for byte. Its source has 2,310,095 bytes and SHA--256
\begin{center}\scriptsize\ttfamily
D9349B9F87D70400E7A140A3BDC38D1BC60F8E3A12F830D8C37693CD6D3C30C5.
\end{center}
This V1 is the sole active main YM source.
Only \texttt{.tex} files are stored in \texttt{latex\_notes};
build products and visual checks are outside it.

\begin{firewall}
The new exact compact result is the conditional-affinity formula
for the dressed bridge kinetic factor. The subsequent Gaussian
overlap, finite-time diffusion change, null directions and
composition defect have proved fixed-volume local limits.
They are not a volume-uniform approximation of the full native
transfer or a positive physical continuum mass-gap theorem.
The full Yang--Mills goal remains unachieved.
\end{firewall}

\section{V2: joint cube dynamics in the dressed transfer}
\label{f16v2:context}

V1 computed the dressed bridge factor while holding the internal
variables fixed. We now include the cube motion at both time ends.
There are two separate corrections: nonroot gauge matching changes
the leading bridge covariance, and compression of the \emph{full}
magnetic sandwich uses the magnetic multiplier twice at each end.
Neither correction is contained in the static conditional affinity.

We first give an exact full-compact kernel. We then integrate its
entire leading internal Gaussian propagator, not just its ground
projection. The resulting two-end Schur matrices have
volume-independent locality and quadratic-form comparisons.
The native limit remains fixed-volume and local in the bridge
coordinates. This does not yet control the discarded native
transfer subspace or the interacting infrared modes.

\subsection{An exact two-ended compact kernel}

Write \(f_v(X)\) for V1's moving port frame, with \(f_o=I\) at
every block root, and put
\[
 \bar X_e=f_s(X)^{-1}X_ef_t(X),\qquad
 \bar X'_e=f_s(X')^{-1}X'_ef_t(X').
\]
This includes tree edges, whose original representatives are \(I\).
As in f15 V99, this is a pointwise gauge rewriting, not a change
of internal integration variables. Denote the exact mixed
multiplier in moving coordinates by
\[
 m_\gamma(X,Z)=M_\times(X,\Theta_X^{-1}Z),
\]
and the internal magnetic multiplier by
\(b_\gamma(X)=\exp[-\gamma S_{\rm int}(X)/2]\).
There are \(7B\) nonroot gauge variables.

\begin{proposition}[Full compact cut kernel with both moving frames]
\label{f16v2:exact-cut}
Before the remaining root Gauss restriction, the cut transfer has
kernel
\begin{align}
 \mathcal K_\gamma(X,Z;X',Z')
 &=b_\gamma(X)b_\gamma(X')\,\mathcal L_\gamma(X,Z;X',Z'),
 \notag\\
 \mathcal L_\gamma
 &=\int_{G^{7B}}
   \prod_{e\ {\rm internal}}
       k_\gamma(r_s^{-1}\bar X_e r_t\bar X_e'^{-1})
   \prod_{e\ {\rm bridge}}
       k_\gamma(r_s^{-1}Z_e r_t Z_e'^{-1})\,dH^{7B}(r).
 \label{f16v2:cut}
\end{align}
Here \(r_o=I\) at block roots. The exact kernel of the full
transfer is \(m_\gamma(X,Z)\mathcal K_\gamma
m_\gamma(X',Z')\). All formulas restrict to the root-invariant
physical space.
\end{proposition}

\begin{proof}
In the rooted tree variables at the two times write \(g'=gh\).
Nonroot Gauss reduction integrates \(h\) and leaves the kinetic
arguments \(h_s^{-1}X_eh_tX_e'^{-1}\) on internal edges and
\(h_s^{-1}Y_eh_tY_e'^{-1}\) on bridges.
This is the same exact Fourier-free Haar reduction used in
f15 V97; no boundary representation has been discarded.

Set \(h_v=f_v(X)r_vf_v(X')^{-1}\).
For fixed endpoints this is a product of Haar translations.
The internal argument becomes, up to conjugation by \(f_s(X')\),
\(r_s^{-1}\bar X_e r_t\bar X_e'^{-1}\).
On a bridge, substitute \(Y_e=f_s(X)Z_ef_t(X)^{-1}\) and its
primed counterpart. The same conjugation gives
\(r_s^{-1}Z_er_tZ_e'^{-1}\). Centrality of \(k_\gamma\) removes
the conjugations. The internal and mixed magnetic traces are
unchanged by the pointwise gauge rewriting.
This proves the kernel and the full sandwich.
\end{proof}

Let \(\mathcal I_\gamma\) be V1's exact normalized dressed
isometry. Distinguish
\[
 \mathcal B_{1,\gamma}=\mathcal I_\gamma^*
                  T_{\rm cut,\gamma}\mathcal I_\gamma,\qquad
 \mathcal B_{2,\gamma}=\mathcal I_\gamma^*
                  T_{\rm full,\gamma}\mathcal I_\gamma.
\]
Their exact kernels, for \(\ell=1,2\), are
\begin{equation}
 \begin{aligned}
 \mathcal B_{\ell,\gamma}(Z,Z')
 &=\frac{1}{\sqrt{F_{\gamma,2}(Z)F_{\gamma,2}(Z')}}\\
 &\quad\times\int \Phi_\gamma(X)\Phi_\gamma(X')\,
      m_\gamma(X,Z)^\ell m_\gamma(X',Z')^\ell
      \mathcal K_\gamma(X,Z;X',Z')\,dH^{5B}(X)dH^{5B}(X').
 \end{aligned}
 \label{f16v2:exact-compression}
\end{equation}
In particular, \(\ell=2\) is necessary for the full transfer:
\[
 \mathcal I_\gamma^*T_{\rm full,\gamma}\mathcal I_\gamma
 =M_{F_{\gamma,2}^{-1/2}}\widehat S_\gamma^*
      M_\times^2T_{\rm cut,\gamma}M_\times^2\widehat S_\gamma
                                      M_{F_{\gamma,2}^{-1/2}}.
\]
Using \(\ell=1\) instead would compress a different operator.

\subsection{The kinetic covariance produced by the moving nonroot vertices}

Use block-diagonal copies \(\mathsf B_B,\Sigma_B,\mathsf M_B,
\mathsf H_B,U_B\) of the single-cube matrices in V1 and f15 V97,
where \(\Sigma=(\mathsf B^{\mathsf T}\mathsf B)^{-1}\).
Let \(E\) be the \(12B\)-by-\(7B\) bridge endpoint incidence:
its row is \(z_t-z_s\), with each root coordinate set to zero.
Define
\begin{equation}
 G=I+E\Sigma_BE^{\mathsf T},\qquad
 L_h=\mathsf B_B^{\mathsf T}\mathsf B_B+E^{\mathsf T}E.
 \label{f16v2:G}
\end{equation}
The letter \(G\) in these matrix formulas denotes a covariance;
the group remains \(\SU(2)\).

\begin{theorem}[Exact leading kinetic square and a uniform covariance bound]
\label{f16v2:kinetic-square}
For one colour component, \(\delta x=x-x'\), \(\delta y=y-y'\),
and \(z\in\mathbb R^{7B}\),
\begin{align}
 &\|U_B\delta x+\mathsf B_Bz\|^2+\|\delta y+Ez\|^2
 \notag\\
 &\quad=\delta x^{\mathsf T}\mathsf M_B^{-1}\delta x
       +\delta y^{\mathsf T}G^{-1}\delta y
       +(z+L_h^{-1}E^{\mathsf T}\delta y)^{\mathsf T}
                    L_h(z+L_h^{-1}E^{\mathsf T}\delta y).
 \label{f16v2:completed-kinetic}
\end{align}
Moreover
\begin{equation}
  \det L_h=(\det\mathsf M)^B\det G,\qquad
                         I\le G\le30I.
 \label{f16v2:Gbounds}
\end{equation}
The matrix \(G\) couples only bridges incident to a common block.
\end{theorem}

\begin{proof}
The balanced internal embedding obeys
\(U_B^{\mathsf T}\mathsf B_B=0\) and
\(U_B^{\mathsf T}U_B=\mathsf M_B^{-1}\).
There is therefore no internal--port cross term.
Complete the \(z\)-square. Woodbury inversion gives
\[
 I-E L_h^{-1}E^{\mathsf T}
                    =(I+E\Sigma_BE^{\mathsf T})^{-1}=G^{-1}.
\]
The determinant lemma and the cube equality
\(\det(\mathsf B^{\mathsf T}\mathsf B)=\det\mathsf M\)
give the determinant statement.

Each row of \(E\) has absolute sum at most two. Each nonroot
port is incident to three bridges, so each column has absolute
sum three. Thus \(\|E\|^2\le6\).
The diagonal of the rooted inverse \(\Sigma\) consists of the
effective resistances from the root. For the cube these are
\(7/12\) at the three adjacent vertices, \(3/4\) at the three
face-opposite vertices and \(5/6\) at the opposite vertex.
They are the exact resistance values verified and derived
in f15 V98. Consequently
\[
 \|\Sigma_B\|=\|\Sigma\|\le\Tr\Sigma
                    =3(7/12)+3(3/4)+5/6=29/6.
\]
This proves \(G\le I+6(29/6)I=30I\).
Block diagonality of \(\Sigma_B\) gives the support assertion.
\end{proof}

The bound \(30\) is a convenient pre-root-quotient bound,
not a replacement for f15 V98's sharper physical electric
form comparison. The new point is that the bridge covariance
is \(G\), not \(I\). Fixing all temporal port variables to zero
would miss it.

For a positive \(d\)-by-\(d\) path covariance \(Q\), abbreviate
the three-colour Gaussian density by
\[
 \varphi_Q(v)
  =(2\pi)^{-3d/2}\det(Q)^{-3/2}
             e^{-v^{\mathsf T}(Q^{-1}\otimes I_3)v/2}.
\]
Write \(b_0(x)=e^{-x^{\mathsf T}(\mathsf H_B\otimes I_3)x/4}\).
The complete leading cut kernel is
\begin{equation}
 \mathcal K_0(x,y;x',y')
       =b_0(x)b_0(x')\,
           \varphi_{\mathsf M_B}(x-x')\varphi_G(y-y').
 \label{f16v2:leading-cut}
\end{equation}
The full leading kernel has in addition
\(e^{-|Ax+Dy|^2/4}\) at each end.

\begin{proposition}[Fixed-volume native convergence with all internal ends integrated]
\label{f16v2:native-limit}
At each fixed \(B\), after including all square-root Haar
chart densities, the kernel \eqref{f16v2:cut} converges uniformly
on bounded \(x,x',y,y'\) windows to \eqref{f16v2:leading-cut}.
The kernels \eqref{f16v2:exact-compression} converge uniformly
on bounded \(y,y'\) windows to the result of integrating the
complete Gaussian kernel \eqref{f16v2:leading-cut} with its
displayed endpoint weights and normalizers.
There is no internal excitation cutoff in this assertion.
It also gives Hilbert--Schmidt convergence after restricting
both bridge ends to a fixed bounded coordinate window.
\end{proposition}

\begin{proof}
Put \(r_v=\operatorname{Exp}(z_v/\sqrt\gamma)\).
The first-order kinetic words in \eqref{f16v2:cut} are exactly
\(U_B(x-x')+\mathsf B_Bz\) and \(y-y'+Ez\).
Their actions have leading exponent one half of the sum of
the two squared norms. Gaussian integration and
\eqref{f16v2:completed-kinetic} give the covariance in
\eqref{f16v2:leading-cut}; \eqref{f16v2:Gbounds} fixes its
normalizing determinant. The internal plaquette multiplier
converges to \(b_0\).

Here is the compact-integral justification.
When the endpoint charts tend to identity, the internal tree
kinetic words force every nonroot \(r_v\) to identity: each
tree has a fixed root. On the compact group the sum of their
squared geodesic distances controls the distances of all
\(r_v\), up to an endpoint error bounded by \(C/\gamma\)
on a fixed endpoint window. This follows successively along
the seven tree edges, by the bi-invariant triangle inequality
and Cauchy--Schwarz along each path. Since
\(1-\cos\theta\ge2\theta^2/\pi^2\) for \(0\le\theta\le\pi\),
the rescaled integrand is dominated by
\(C\exp(-c\sum_v|z_v|^2)\), with constants at fixed \(B\)
and window. The complement of any fixed small group
neighbourhood is exponentially small. Taylor expansion there,
followed by dominated convergence, proves the local kernel limit.

For the integration over all internal \(X,X'\), a separate
uniform bound avoids a hidden chart restriction.
Keep only the \(7B\) tree factors of the kinetic product and
bound the other \(17B\) factors by \(\|k_\gamma\|_\infty\).
Successive Haar integration of leaves gives one for the tree
product. The full integral is therefore at most
\(\|k_\gamma\|_\infty^{17B}\).
The normalization in \eqref{f16v1:kinetic} gives
\(\|k_\gamma\|_\infty\le C\gamma^{3/2}\) for \(\gamma\ge1\).
Each chart Jacobian is at most \(C\gamma^{-3/2}\).
The two half-densities on the \(17B\) reduced coordinates
cancel this power. Thus the cut chart kernel is bounded by
a constant depending on \(B\), uniformly throughout the
internal charts; magnetic multipliers are at most one.

The transformed \(\Phi_\gamma\) is exactly a normalized
truncation of the product Gaussian ground amplitude.
It therefore gives an integrable dominating product
\(C_B\Omega(x)\Omega(x')\).
On a fixed bridge window \(F_{\gamma,2}\) is bounded below
uniformly for sufficiently large \(\gamma\), by V1's
fixed-volume limit. Internal Gaussian tails are now uniform
in \(\gamma\). Apply the local limit on bounded internal
windows and then send those windows to infinity.
This proves uniform convergence of the compressed kernels.
The final Hilbert--Schmidt assertion follows on any bounded
bridge window. All constants may deteriorate with \(B\).
\end{proof}

\subsection{The two-end Schur formula retaining the full internal propagator}

The product cube Gaussian ground amplitude \(\Omega(x)\)
has squared covariance \(C_B\) from \eqref{f16v1:covariance}.
Write \(N=\mathsf M_B^{-1}\) and
\begin{equation}
 Q_0=\tfrac12(C_B^{-1}+\mathsf H_B),\qquad
 V_{\ell,+}=Q_0+\tfrac\ell2 A^{\mathsf T}A,\qquad
 V_{\ell,-}=V_{\ell,+}+2N,\qquad \ell=0,1,2.
 \label{f16v2:V}
\end{equation}
All are strictly positive. Define
\begin{align}
 P_{\ell,\pm}&=D^{\mathsf T}A V_{\ell,\pm}^{-1}A^{\mathsf T}D,
 \notag\\
 B_{\ell,\pm}
       &=\ell D^{\mathsf T}D-H_{\rm eff}
                              -\tfrac{\ell^2}{2}P_{\ell,\pm}.
 \label{f16v2:B}
\end{align}
Colour tensors are suppressed: each path matrix acts with \(I_3\)
on the actual bridge vectors.

Let \(r(q)=(1+q/2+\sqrt{q+q^2/4})^{-1}\), and put
\[
                   \lambda_\Box^0=r(4)^{9/2}r(6)^3.
\]
This is the top value of the entire single-cube Gaussian
transfer, not its physical mass. It is the
\href{https://dlmf.nist.gov/18.18.E28}{Mehler value}
already derived in f15 V97, with nine coordinates of parameter
four and six coordinates of parameter six.
Set
\begin{equation}
 \kappa_\ell
 =\frac{(\lambda_\Box^0)^B}{F_2(0)}
   \left[
    \frac{\det V_{0,+}\det V_{0,-}}
         {\det V_{\ell,+}\det V_{\ell,-}}
   \right]^{3/2},\qquad \ell=1,2.
 \label{f16v2:kappa}
\end{equation}

\begin{theorem}[Complete Gaussian dressed cut and full compressions]
\label{f16v2:schur}
The fixed-volume limits in Proposition~\ref{f16v2:native-limit} are
\begin{align}
 \mathcal B_{\ell,0}(y,y')
 &=\kappa_\ell\,\varphi_G(y-y') \notag\\
 &\quad\cdot
 \exp\left[-\tfrac18(y+y')^{\mathsf T}B_{\ell,+}(y+y')
           -\tfrac18(y-y')^{\mathsf T}B_{\ell,-}(y-y')\right],
       \qquad \ell=1,2.
 \label{f16v2:schur-kernel}
\end{align}
In particular \(\ell=2\) includes the entire internal Gaussian
propagator of the full magnetic transfer between the dressed ends.
No intermediate internal ground projection appears.
\end{theorem}

\begin{proof}
At one end the integrand of \eqref{f16v2:exact-compression}
tends to
\[
 \frac{\Omega(x)
           e^{-\ell|Ax+Dy|^2/4}b_0(x)}{\sqrt{F_2(y)}} .
\]
The two-end internal precision and source are
\[
 W_\ell=
 \begin{pmatrix}V_{\ell,+}+N&-N\\-N&V_{\ell,+}+N\end{pmatrix},
 \qquad
 \eta_\ell=-\frac\ell2
          \begin{pmatrix}A^{\mathsf T}Dy\\A^{\mathsf T}Dy'\end{pmatrix}.
\]
Changing to \((x+x')/\sqrt2,(x-x')/\sqrt2\) gives the two
precisions \(V_{\ell,+},V_{\ell,-}\). Gaussian integration
therefore contributes determinant
\((\det V_{\ell,+}\det V_{\ell,-})^{-3/2}\)
and exponent
\[
 \frac{\ell^2}{16}\left[
   (y+y')^{\mathsf T}P_{\ell,+}(y+y')
     +(y-y')^{\mathsf T}P_{\ell,-}(y-y')\right].
\]
The remaining endpoint exponent is
\[
 -\frac\ell4(y^{\mathsf T}D^{\mathsf T}Dy
                   +y'^{\mathsf T}D^{\mathsf T}Dy')
 +\frac14(y^{\mathsf T}H_{\rm eff}y
                   +y'^{\mathsf T}H_{\rm eff}y').
\]
The second term comes from \(F_2(y)^{-1/2}F_2(y')^{-1/2}\).
Combining these expressions gives exactly \eqref{f16v2:B}.

To fix the scalar without discarding any determinant, set the
endpoint magnetic power to zero in the internal integral.
It is \(\langle\Omega,K_{\Box,0}^{\otimes B}\Omega\rangle
=(\lambda_\Box^0)^B\). Dividing the general determinant by
this one and retaining \(F_2(0)^{-1}\) gives
\eqref{f16v2:kappa}. This proves the formula.
\end{proof}

This is a calculation of one-step compression, not a claim that
its spectrum equals the spectrum of \(T_{\rm full,\gamma}\).
V1's intermediate-complement warning remains in force.

\subsection{Uniform locality of the retained two-time fast response}

The inverses in \eqref{f16v2:V} are uniformly local even for the
full magnetic power \(\ell=2\).
For either sign put
\[
 Q_+=Q_0,\quad Q_-=Q_0+2N,\qquad
 Z_\pm=Q_\pm^{-1/2}A^{\mathsf T}A Q_\pm^{-1/2}.
\]
Since \(Q_\pm\ge C_B^{-1}/2\),
\[
                 0\le Z_\pm,\qquad \|Z_\pm\|\le2\rho .
\]
These matrices have only nearest-block off-diagonal entries.
Let \(c=1+2\rho\), \(\theta=1-c^{-1}<1\), and
\(T_\pm=I-(I+Z_\pm)/c\). Then \(0\le T_\pm\le\theta I\) and
\begin{equation}
 V_{2,\pm}^{-1}
   =c^{-1}Q_\pm^{-1/2}
                     \sum_{j=0}^{\infty}T_\pm^jQ_\pm^{-1/2}.
 \label{f16v2:local-inverse}
\end{equation}
For block sets \(E_0,F_0\) separated by distance \(d\), this gives
\[
 \|P_{E_0}V_{2,\pm}^{-1}P_{F_0}\|
                    \le\|Q_\pm^{-1}\|\theta^d.
\]
Indeed powers of degree less than \(d\) cannot connect the sets,
and \(c^{-1}/(1-\theta)=1\).
No volume or set-cardinality factor is needed.
Consequently \(P_{2,\pm}\) and \(B_{2,\pm}\) have exponentially
decaying block entries, with the fixed range shift of \(A,D\).
This controls the complete leading two-time fast response,
not the non-Gaussian native remainder.

\subsection{The full effective quadratic transfer and its surviving soft modes}

Write \(B_\pm=B_{2,\pm}\).
The lower bound here is relative to the actual coarse curl,
not relative to the identity.

\begin{theorem}[Volume-independent curvature and kinetic comparisons]
\label{f16v2:comparison}
One has
\begin{equation}
 \frac{D^{\mathsf T}D}{(1+\rho)(1+2\rho)}
       \le B_+\le B_-
       \le\frac{1+2\rho}{1+\rho}D^{\mathsf T}D .
 \label{f16v2:Bbounds}
\end{equation}
Define
\[
       C_{\rm eff}
       =\left[G^{-1}+\tfrac14(B_--B_+)\right]^{-1}.
\]
Then
\begin{equation}
 \frac{I}{1+16\rho/(1+2\rho)}
       \le C_{\rm eff}\le30I ,
 \label{f16v2:Ceffbounds}
\end{equation}
and, with
\(\widetilde\kappa=\kappa_2(\det C_{\rm eff}/\det G)^{3/2}\),
\begin{equation}
 \mathcal B_{2,0}(y,y')
   =\widetilde\kappa\,
      e^{-y^{\mathsf T}B_+y/4}
       \varphi_{C_{\rm eff}}(y-y')
      e^{-y'^{\mathsf T}B_+y'/4}.
 \label{f16v2:effective}
\end{equation}
\end{theorem}

\begin{proof}
Since \(V_{2,+}\ge C_B^{-1}/2+A^{\mathsf T}A\),
\[
 A V_{2,+}^{-1}A^{\mathsf T}
              \le2K(I+2K)^{-1}.
\]
The commuting functions of \(K\) obey the exact identity
\[
 2I-(I+K)^{-1}-4K(I+2K)^{-1}
                    =(I+K)^{-1}(I+2K)^{-1}.
\]
Substitute in \(B_+=2D^{\mathsf T}D-H_{\rm eff}-2P_{2,+}\)
and use \(0\le K\le\rho I\) for the lower bound.
The upper bound follows from \(P_{2,-}\ge0\) and
\(H_{\rm eff}\ge(1+\rho)^{-1}D^{\mathsf T}D\).
Inversion of \(V_{2,-}\ge V_{2,+}\) gives \(B_-\ge B_+\).

Furthermore
\[
 0\le B_--B_+
   =2D^{\mathsf T}A(V_{2,+}^{-1}-V_{2,-}^{-1})A^{\mathsf T}D
   \le\frac{4\rho}{1+2\rho}D^{\mathsf T}D.
\]
Every bridge belongs to four plaquettes; each mixed row has at
most four bridge entries. Thus \(\|D\|^2\le16\).
Together with \(I\le G\le30I\), this proves
\eqref{f16v2:Ceffbounds}. Finally split the exponent in
\eqref{f16v2:schur-kernel} into the two endpoint \(B_+\)
terms and the difference term \(B_--B_+\).
The Gaussian prefactor changes by the displayed determinant
ratio, proving \eqref{f16v2:effective}.
\end{proof}

These comparisons prove \(\ker B_+=\ker D\).
For the common-bridge subspace \(\mathcal E\) of V1,
\(A^{\mathsf T}Dy=0\) and \(H_{\rm eff}y=D^{\mathsf T}Dy\).
Thus
\[
 B_+y=B_-y=D^{\mathsf T}Dy,\qquad
            (B_--B_+)y=0\quad(y\in\mathcal E).
\]
The full leading cube propagation has not generated a constant
restoring term for those coarse collective modes.

For clarity, even the positive oscillator frequencies of
\eqref{f16v2:effective} have a soft finite-volume edge.
Let \(\nu_{\min}^+\) be the least positive eigenvalue of
\(C_{\rm eff}^{1/2}B_+C_{\rm eff}^{1/2}\).
Then
\begin{equation}
                     \nu_{\min}^+\le60\sin^2(\pi/m).
 \label{f16v2:soft}
\end{equation}
To prove it, use V1's common-bridge transverse cosine \(v\),
orthogonal in the Euclidean metric to \(\ker D\).
Its numerator is \(\langle v,D^{\mathsf T}Dv\rangle\).
Choose the representative \(v+w\), \(w\in\ker D\), orthogonal
to that kernel in the \(C_{\rm eff}^{-1}\) metric.
Its numerator is unchanged, and its denominator is at least
\(\|v+w\|^2/30\ge\|v\|^2/30\).
The generalized minimum principle and V1's exact quotient
\(2\sin^2(\pi/m)\) prove \eqref{f16v2:soft}.

The zero directions include residual linearized gauge directions
and torons. The full Euclidean Gaussian operator has free
diffusion in these directions, not a normalizable isolated vacuum.
If those directions are removed to discuss just its positive
oscillator factors, a parameter \(\nu>0\) has Mehler ratio
\((1+\nu/2+\sqrt{\nu+\nu^2/4})^{-1}\).
Equation~\eqref{f16v2:soft} concerns this specified Gaussian
factorization. It is neither a native physical gaplessness
theorem nor a choice of physical time scaling.

\subsection{The remaining full-operator return and verification}

The exact compact kernel and the full Gaussian Schur response
now include both time ends and all internal propagation within
one step. The missing spectral implication is still substantive:
a one-step compression does not bound the full complement.
With \(P_\gamma=\mathcal I_\gamma\mathcal I_\gamma^*\),
the actual two-step return is
\[
 \mathcal I_\gamma^*T_{\rm full,\gamma}
       (I-P_\gamma)T_{\rm full,\gamma}\mathcal I_\gamma .
\]
It is nonnegative, but no sufficiently small relative bound
uniform in spatial volume has been proved. The native bridge
tails and the nonlinear common-bridge dynamics must also be
kept under the same normalization. The fixed-volume domination
in Proposition~\ref{f16v2:native-limit} is not such a bound.

The next upstream target is to retain this return as a
two-time or Schur/memory operator and test a volume-uniform
relative comparison, rather than identify the compressed
Gaussian spectrum with the full transfer spectrum.
Independent physical-time calibration and interacting
four-dimensional continuum reconstruction remain open.

The diagnostic
\begin{center}\small
\path{scripts/verify_ym_f16_v2_joint_dressed_transfer.py}
\end{center}
checks noncommuting two-ended compact words, the actual
two-cube kinetic completion and determinant, full-torus port
incidences, the independent two-end Gaussian precision and
plus/minus Schur formulas, the matrix identity behind
\eqref{f16v2:Bbounds}, and the endpoint magnetic-square
distinction. Finite checks supplement the proofs; they are not
a continuum spectral certificate. All 4,582 new exact fixtures passed,
together with the inherited V1 chain. The diagnostic uses exact
elimination for the larger determinants, rather than the inherited
small-matrix cofactor routine. It has 8,154 bytes and SHA--256
\begin{center}\scriptsize\ttfamily
95087B3BF6B4AA9517C734E8A35B690B9074EF3D9355AC7125B2A22340C59DCF.
\end{center}

The predecessor V1 has 25,601 bytes and SHA--256
\begin{center}\scriptsize\ttfamily
A04BB24BC54769C31EF36F591282F6ABCF828DA25920C369DCB9DA9F1EF27097.
\end{center}
Its body is preserved byte for byte; removing this append and
restoring the title recovers that source. All 167 previously
protected files remain unchanged. The last companion
checkpoint is f15 V100 and the next is this lineage's V5.
Only \texttt{.tex} files belong in \texttt{latex\_notes};
build and visual-check products are kept outside that folder.

\begin{firewall}
V2 gives an exact full-compact joint kernel, a fixed-volume
native dressed-compression limit with all internal ends
integrated, and explicit volume-independent comparisons for
its complete leading two-time Gaussian response.
No small volume-uniform full-operator return, positive
physical continuum mass gap, or interacting continuum
Yang--Mills construction is proved. The goal remains unachieved.
\end{firewall}

\section{V3: extensive dressed-line leakage and a retained fast history}
\label{f16v3:context}

V2 computed one full step between the normalized dressed ends.
We now ask what happens before inserting another dressed projection.
The answer is not a small correction to the static conditional line:
even the complete leading Gaussian transfer changes its fast
covariance by an extensive amount. We prove a quantitative
pre-root-restriction obstruction to a small relative return.
We then integrate the fast variables along an arbitrary bridge
history without intermediate projection. Its covariance and
quadratic memory have volume- and time-length-independent decay.

This is a change in the next analytic target, not a mass-gap claim.
F15 V99--V100 concerned the magnetic image and its old oscillator
occupations; V1 concerned bridge-only affinity; V2 concerned the
joint one-step compression. The new calculation propagates the
dressed image by the full Gaussian transfer and retains arbitrarily
many fast time slices. The archive search also found spatial
Gaussian blocking memory in f11 V32--V33 and native observable
covariance estimates in f15 V61; neither is the conditional
fast-history calculation below.

\subsection{The propagated conditional vector is not the static dressed vector}

All statements in this subsection concern V2's leading Gaussian
model before the remaining root Gauss restriction.
Write \(\mathsf T_0\) for its full transfer on the internal and
bridge Euclidean coordinates, and
\[
 (\mathcal J_0 f)(x,y)=\psi_y(x)f(y),\qquad
 \psi_y(x)=\frac{\Omega(x)e^{-|Ax+Dy|^2/4}}{\sqrt{F_2(y)}} .
\]
The vector \(\psi_y\) has norm one in the internal coordinates.
Each displayed path matrix is tensored with \(I_3\) when acting
on the three colour components. Put
\begin{align}
 N&=\mathsf M_B^{-1},&
 S&=\mathsf H_B+A^{\mathsf T}A,&
 \Lambda&=A^{\mathsf T}D,\notag\\
 E_0&=N+\tfrac12(C_B^{-1}+\mathsf H_B),&
 E&=E_0+A^{\mathsf T}A,&
 P&=\tfrac12(C_B^{-1}+A^{\mathsf T}A).
 \label{f16v3:definitions}
\end{align}
Here \(E\) is V2's two-end diagonal block \(V_{2,+}+N\).
The precision in the amplitude \(\psi_z\), as opposed to its
squared probability density, is \(P\):
\[
 \psi_z(x)=c(z)\exp[-x^{\mathsf T}Px/2-x^{\mathsf T}\Lambda z/2],
                      \qquad c(z)>0.
\]
The sources and the positive scalar depend on the bridge field.

For fixed initial bridge coordinate \(y\) and final coordinate \(z\),
define the fast output feature
\[
 q_{z,y}(x)=\int \mathsf T_0(x,z;x',y)\psi_y(x')\,dx'.
\]
This definition has no intermediate fast projection.

\begin{proposition}[Full-step covariance mismatch]
\label{f16v3:propagation}
Set
\begin{equation}
 \Delta_f=N(E_0^{-1}-E^{-1})N\ge0.
 \label{f16v3:delta}
\end{equation}
There is a positive scalar \(c(z,y)\) such that
\begin{equation}
 q_{z,y}(x)=c(z,y)
 \exp\left[
   -\tfrac12x^{\mathsf T}(P+\Delta_f)x
   -x^{\mathsf T}\bigl(\tfrac12\Lambda z+NE^{-1}\Lambda y\bigr)
 \right].
 \label{f16v3:feature}
\end{equation}
Consequently the squared normalized output feature has covariance
\([2(P+\Delta_f)]^{-1}\otimes I_3\), independent of \(y,z\).
It is not the covariance \((2P)^{-1}\otimes I_3\) of the static
dressed line unless \(A=0\).
\end{proposition}

\begin{proof}
The full kernel is the product of its internal and bridge heat
densities and the two endpoint factors
\[
 \exp[-(x^{\mathsf T}\mathsf H_Bx+|Ax+Dz|^2)/4].
\]
After multiplication by \(\psi_y(x')\), the input precision is \(E\)
and its linear source is \(Nx-\Lambda y\).
Gaussian integration of \(x'\) leaves output precision
\[
 Q=N+\tfrac12S-NE^{-1}N
\]
and linear source \(-\Lambda z/2-NE^{-1}\Lambda y\).
The cube ground-state equation gives the exact identity
\begin{equation}
 N+\tfrac12(\mathsf H_B-C_B^{-1})=NE_0^{-1}N.
 \label{f16v3:riccati}
\end{equation}
For completeness, conjugate by \(\mathsf M_B^{1/2}\).
On a cube mode with parameter \(t=4\) or \(6\), the two sides
become
\[
 1+t/2-\sqrt{t+t^2/4}
       =\bigl(1+t/2+\sqrt{t+t^2/4}\bigr)^{-1};
\]
their product identity follows by squaring the radical.
This proves \eqref{f16v3:riccati}, hence \(Q=P+\Delta_f\).
Inversion reverses \(E\ge E_0\), so \(\Delta_f\ge0\).
If \(\Delta_f=0\), invertibility of \(N,E,E_0\) implies
\(A^{\mathsf T}A=0\). The covariance assertion follows by squaring
the amplitude and normalizing its positive Gaussian.
\end{proof}

\subsection{A volume-extensive obstruction to a small static return}

The covariance mismatch gives a uniform conditional overlap
bound, not just a nonzero example.

\begin{theorem}[Dressed-line retention after a full step]
\label{f16v3:retention}
For the actual parity-cube geometry with \(B=m^3\), define
\[
 \eta_B(z,y)=
       \frac{|\langle\psi_z,q_{z,y}\rangle|^2}
            {\|q_{z,y}\|_2^2}.
\]
Then, for every \(y,z\),
\begin{equation}
 0<\eta_B(z,y)\le
 \eta_B^{\rm cov}:=
 \frac{\det(I+Z)^{3/2}}{\det(I+Z/2)^3}
 \le e^{-c_*B}\le e^{-B/100000},
 \qquad
 c_*=\frac{50}{147}\frac{(25/1728)^2}{5}>10^{-5}.
 \label{f16v3:retention-bound}
\end{equation}
Here \(Z\) is the relative covariance-precision change defined
in the proof. The determinants retain all \(5B\) path modes,
with the displayed power accounting for all three colours.
\end{theorem}

\begin{proof}
Use bars to denote congruence by \(N^{-1/2}=\mathsf M_B^{1/2}\);
in particular
\[
 \overline E_0=N^{-1/2}E_0N^{-1/2},\quad
 W=N^{-1/2}A^{\mathsf T}AN^{-1/2},\quad
 \overline P=N^{-1/2}PN^{-1/2}.
\]
Then
\[
 \overline\Delta_f
   =\overline E_0^{-1}-(\overline E_0+W)^{-1},\qquad
 Z=\overline P^{-1/2}\overline\Delta_f\,\overline P^{-1/2}.
\]
The eigenvalues of \(Z\) are the generalized eigenvalues of
\((\Delta_f,P)\); no commuting-matrix assumption is being made.

The inherited cube spectrum and covariance give
\[
 5I\le\overline E_0\le8I,\qquad 2I\le\overline P\le6I.
\]
Indeed the exact bounds for \(\overline E_0\) are
\(3+2\sqrt2\) and \(4+\sqrt{15}\).
For the second assertion, the spectrum of the whitened
\(C_B^{-1}\) lies between \(4\sqrt2\) and \(2\sqrt{15}\), and
\(0\le W\le4I\).
This last inequality follows from
\(A=C_{\rm int}U_B\), \(\|C_{\rm int}\|\le2\), and
\(U_B^{\mathsf T}U_B=N\).
Moreover the exact diagonal Gram blocks
\((A^{\mathsf T}A)_{bb}=2N_{bb}\) from f15 V99 give
\[
                           \Tr W=10B.
\]
These are bounds for the actual geometry, not assumptions
on independently chosen toy matrices.

Woodbury inversion, or continuous extension if \(W\) is singular,
gives
\[
 \overline\Delta_f
 =\overline E_0^{-1}W^{1/2}
   (I+W^{1/2}\overline E_0^{-1}W^{1/2})^{-1}
                       W^{1/2}\overline E_0^{-1}
 \ge\frac59\,\overline E_0^{-1}W\overline E_0^{-1}.
\]
Taking traces and using \(\overline E_0\le8I\) yields
\[
 \Tr\overline\Delta_f\ge\frac{25}{288}B,\qquad
 \Tr Z\ge\frac{25}{1728}B,\qquad 0\le Z\le\frac1{10}I.
\]
For the last estimate use
\(\overline\Delta_f\le\overline E_0^{-1}\le I/5\) and
\(\overline P\ge2I\).
Cauchy--Schwarz on the \(5B\) eigenvalues gives
\[
                   \Tr Z^2\ge\frac{(25/1728)^2}{5}\,B.
\]

The squared overlap of two normalized Gaussian amplitudes with
precisions \(P\) and \(P+\Delta_f\) is at most the value when
their means coincide. Completing the square shows that unequal
means multiply this value by \(e^{-a}\) for a nonnegative
quadratic form \(a\). The coincident-mean determinant is precisely
\(\eta_B^{\rm cov}\) in \eqref{f16v3:retention-bound}; congruence changes
do not change it. For an eigenvalue \(0\le z\le1/10\),
\[
 \frac32\log\frac{(1+z/2)^2}{1+z}
 =\frac32\log\left(1+\frac{z^2}{4(1+z)}\right)
 \ge\frac{3z^2}{8(1+z/2)^2}
 \ge\frac{50}{147}z^2 .
\]
We used \(\log(1+u)\ge u/(1+u)\), which follows by integrating
its nonnegative derivative difference from zero.
Summation proves the stated exponential bound.
\end{proof}

Let
\[
 \mathsf B_0=\mathcal J_0^*\mathsf T_0\mathcal J_0,\quad
 \mathsf A_2=\mathcal J_0^*\mathsf T_0^2\mathcal J_0,\quad
 \mathsf R_0=\mathsf A_2-\mathsf B_0^2
 =\mathcal J_0^*\mathsf T_0(I-\mathcal J_0\mathcal J_0^*)
                              \mathsf T_0\mathcal J_0\ge0.
\]
Their diagonal kernels satisfy
\begin{align}
 \mathsf A_2(y,y)&=\int\|q_{z,y}\|_2^2\,dz,\notag\\
 \mathsf B_0^2(y,y)&=\int|\langle\psi_z,q_{z,y}\rangle|^2\,dz
                 \le e^{-c_*B}\mathsf A_2(y,y).
 \label{f16v3:diagonal}
\end{align}
All these diagonal integrals are finite. In fact the internal
cut transfer has norm \((\lambda_\Box^0)^B\); its magnetic
multipliers have norm at most one. Hence
\[
 \|q_{z,y}\|_2
       \le(\lambda_\Box^0)^B\varphi_G(z-y),
\]
whose square is integrable in \(z\), uniformly in \(y\).

For any bounded bridge window \(U\) of positive measure,
the resulting operators localized at their input are
Hilbert--Schmidt and
\begin{equation}
 \|(I-\mathcal J_0\mathcal J_0^*)\mathsf T_0
                      \mathcal J_0\,1_U\|_{\rm HS}^2
 \ge(1-e^{-c_*B})
              \|\mathsf T_0\mathcal J_0\,1_U\|_{\rm HS}^2.
 \label{f16v3:HS-loss}
\end{equation}
This follows by integrating \eqref{f16v3:diagonal}, or by
orthogonal decomposition of each output feature.
In particular a quadratic-form inequality
\(\mathsf R_0\le\varepsilon\mathsf A_2\) on the entire
pre-root-restriction coarse space requires
\(\varepsilon\ge1-e^{-c_*B}\). To see this, sandwich by \(1_U\)
and take the finite trace; the denominator is strictly positive.
Thus no fixed \(\varepsilon<1\) works for all volumes in this model.

This conclusion is stronger than a nonzero projection defect,
but its scope is important. It does not assert pointwise
positivity of every off-diagonal return kernel, nor does it
turn a diagonal comparison into a reverse operator inequality.
It does not apply automatically after physical Gauss
restriction, to a low-energy subspace, or to the native compact
transfer. In particular, V2's fixed-window one-step convergence
alone does not transfer this full two-step statement to that
native operator.

\subsection{Exact elimination along an arbitrary bridge history}

We now avoid intermediate projection altogether. Let \(n\ge1\),
fix the bridge history \(\mathbf y=(y_0,\ldots,y_n)\), and
integrate the internal coordinates \(x_0,\ldots,x_n\) in the
kernel of \(\mathcal J_0^*\mathsf T_0^n\mathcal J_0\).
Set \(E_{\rm i}=S+2N\), and define the \((n+1)\)-by-\((n+1)\)
block precision
\begin{equation}
 (W_n)_{00}=(W_n)_{nn}=E,\qquad
 (W_n)_{ii}=E_{\rm i}\ (0<i<n),\qquad
 (W_n)_{i,i+1}=(W_n)_{i+1,i}=-N,
 \label{f16v3:history-precision}
\end{equation}
with all other time blocks zero.
Its time blocks have \(5B\) path coordinates.
Let \(\mathcal L\mathbf y=(\Lambda y_0,\ldots,\Lambda y_n)\) and
\[
 a_n=\frac1{F_2(0)}
       \left[\det C_B\,(\det\mathsf M_B)^n\det W_n\right]^{-3/2}.
\]

\begin{theorem}[Complete fast-history kernel]
\label{f16v3:history}
After integration of every internal time slice, the exact
Gaussian history weight is
\begin{align}
 \mathcal W_n(\mathbf y)
 &=a_n\prod_{i=0}^{n-1}\varphi_G(y_i-y_{i+1})
       \exp\bigl(\mathcal E_n(\mathbf y)\bigr),\notag\\
 \mathcal E_n(\mathbf y)
 &=-\frac12\sum_{i=0}^n y_i^{\mathsf T}D^{\mathsf T}Dy_i
   +\frac14\bigl(y_0^{\mathsf T}H_{\rm eff}y_0
                         +y_n^{\mathsf T}H_{\rm eff}y_n\bigr)
   +\frac12(\mathcal L\mathbf y)^{\mathsf T}
                           W_n^{-1}(\mathcal L\mathbf y).
 \label{f16v3:history-weight}
\end{align}
The kernel of \(\mathcal J_0^*\mathsf T_0^n\mathcal J_0\) is
\(\int\mathcal W_n(\mathbf y)\,dy_1\cdots dy_{n-1}\).
No internal state or intervening bridge configuration has
been discarded. The determinant in \(a_n\) is part of this
identity, not a freely chosen spectral normalization.
\end{theorem}

\begin{proof}
Each interior time slice receives two magnetic endpoint
multipliers, whereas each external slice receives one
transfer multiplier and one multiplier from its dressed
end vector. Thus the mixed magnetic power is two at
\emph{every} slice. Each linear source is consequently
\(-\Lambda y_i\).
The internal magnetic power is two at interior slices;
at the two ends there is one such factor and one
\(\Omega\) amplitude. The kinetic differences contribute
\(\sum_{i=0}^{n-1}(x_i-x_{i+1})^{\mathsf T}N(x_i-x_{i+1})\).
These observations give exactly \eqref{f16v3:history-precision}.
Its quadratic form can also be written as
\begin{align*}
 \mathbf x^{\mathsf T}W_n\mathbf x
 ={}&x_0^{\mathsf T}(E-N)x_0+x_n^{\mathsf T}(E-N)x_n\\
 &+\sum_{i=1}^{n-1}x_i^{\mathsf T}Sx_i
   +\sum_{i=0}^{n-1}(x_i-x_{i+1})^{\mathsf T}N(x_i-x_{i+1}),
\end{align*}
which proves strict positivity.
Gaussian integration contributes the last term in
\(\mathcal E_n\) and \(\det(W_n)^{-3/2}\).
The scalar \((2\pi)^{15B(n+1)/2}\) from integration cancels
the \(n\) internal heat prefactors and the two ground-amplitude
prefactors. Their remaining factors are
\((\det\mathsf M_B)^{-3n/2}(\det C_B)^{-3/2}\).
Finally the two factors \(F_2^{-1/2}\) contribute \(F_2(0)^{-1}\)
and the displayed external \(H_{\rm eff}\) term.
This proves the history formula. Positivity permits the
iterated integrations; their interpretation as the stated
bounded-operator kernel follows from composition of the
positive Gaussian transfer kernels.
\end{proof}

For \(n=1\), the plus/minus transform has precisions
\(E-N=V_{2,+}\) and \(E+N=V_{2,-}\), so this is exactly V2,
including its scalar. For every \(n\) one may compute the
determinant without dropping its boundary terms by
\[
 R_0=E,\qquad R_i=(W_n)_{ii}-NR_{i-1}^{-1}N,\qquad
                   \det W_n=\prod_{i=0}^n\det R_i .
\]
All pivots are strictly positive by block Gaussian elimination.
The order of the matrix products matters.

\subsection{Uniform temporal and space--time locality of the fast memory}

Write \(\mathbb N=I_{n+1}\otimes N\) and
\(\widehat W_n=\mathbb N^{-1/2}W_n\mathbb N^{-1/2}\).
The graph for space--time locality has a vertex for each
(time slice, spatial block), with the nearest-time and
nearest-spatial-block edges.

\begin{theorem}[Length-independent conditional fast response]
\label{f16v3:locality}
For every \(B,n\),
\begin{equation}
                        4I\le\widehat W_n\le15I.
 \label{f16v3:uniform-spectrum}
\end{equation}
For sets \(\mathcal E,\mathcal F\) at graph distance \(d\),
\begin{equation}
 \|P_{\mathcal E}\widehat W_n^{-1}P_{\mathcal F}\|
                         \le\frac14(11/15)^d .
 \label{f16v3:spacetime}
\end{equation}
A sharper purely temporal bound is
\begin{equation}
 \|N^{1/2}(W_n^{-1})_{ij}N^{1/2}\|
       \le\frac{q^{|i-j|}}{p_*-2},
 \qquad p_*=3+2\sqrt2,\quad q=2/p_*<1.
 \label{f16v3:temporal}
\end{equation}
These are conditional fast-covariance bounds for arbitrary
bridge histories, not decay of the complete physical
Yang--Mills correlations.
\end{theorem}

\begin{proof}
Whitening the quadratic-form decomposition in the preceding
proof gives a time-path Laplacian, with norm at most four.
The interior on-site term is the whitened
\(\mathsf H_B+A^{\mathsf T}A\), between \(4I\) and \(10I\).
The external on-site term is the whitened \(E-N\), between
\((2+2\sqrt2)I\) and \((7+\sqrt{15})I<11I\).
This proves \eqref{f16v3:uniform-spectrum}.
All off-site entries of \(\widehat W_n\) join adjacent
space--time blocks: \(A^{\mathsf T}A\) has only nearest-block
entries, and the whitening is spatially block diagonal.
Thus \(T_n=I-\widehat W_n/15\) satisfies
\(0\le T_n\le(11/15)I\), with the same finite range.
The convergent expansion
\[
             \widehat W_n^{-1}=\frac1{15}\sum_{k\ge0}T_n^k
\]
has no connection between those sets before degree \(d\).
The geometric tail is \((11/15)^d/4\), proving
\eqref{f16v3:spacetime} without a set-cardinality factor.

For the temporal estimate, keep entire spatial slices as blocks.
The diagonal blocks of \(\widehat W_n\) are at least \(p_*I\):
the external ones are \(\overline E_0+W\), and the interior
ones are at least \(6I\). Its temporal off-diagonal blocks are
\(-I\). If \(\mathcal D_n\) is its time-block diagonal,
write \(\widehat W_n=\mathcal D_n-\mathcal A_n\), where
\(\mathcal A_n\) is the path adjacency with identity blocks.
Then
\(\|\mathcal D_n^{-1/2}\mathcal A_n\mathcal D_n^{-1/2}\|
\le2/p_*=q\).
This matrix need not be positive, but its norm-convergent
Neumann series is valid. A term of degree less than
\(|i-j|\) has zero \(ij\) block.
Multiplication by the two diagonal inverse square roots and
summation give
\(p_*^{-1}q^{|i-j|}/(1-q)=q^{|i-j|}/(p_*-2)\).
Undoing the whitening proves \eqref{f16v3:temporal}.
\end{proof}

There is also a curvature-relative truncation estimate.
Let \(\mathcal Q_n(\mathbf y)\) denote the last term in
\(\mathcal E_n\), and let \(\mathcal Q_n^{[r]}\) keep only
its ordered pairs with \(|i-j|\le r\), for an integer \(r\ge0\).
Since \(\|N^{-1/2}A^{\mathsf T}\|\le2\), the temporal bound yields
\begin{equation}
 |\mathcal Q_n(\mathbf y)-\mathcal Q_n^{[r]}(\mathbf y)|
 \le \frac{4}{p_*-2}\frac{q^{r+1}}{1-q}
                               \sum_{i=0}^n\|Dy_i\|^2 .
 \label{f16v3:memory-tail}
\end{equation}
Indeed each ordered matrix element is bounded by
\(4q^{|i-j|}\|Dy_i\|\|Dy_j\|/(p_*-2)\).
The omitted scalar matrix has row sum at most
\(2q^{r+1}/(1-q)\); its quadratic form is bounded by that
row sum times \(\sum_i\|Dy_i\|^2\). Include the factor
\(1/2\) in \(\mathcal Q_n\) to obtain the displayed constant.
No spatial volume or time-length factor multiplies this
relative curvature budget. Truncation is an estimate on the
quadratic action, not automatically an operator-norm or
positivity-preserving approximation.

For histories taking values in V1's common-bridge subspace,
\(\Lambda y_i=0\) at every time. The entire fast source
response \(\mathcal Q_n\), not merely its long-time tail,
then vanishes. The determinant \(a_n\) is independent of
the bridge history. This statement concerns the fixed-history
weight; integrating other bridge modes is a further operation.
It does not create a constant restoring term or settle the
native dynamics of those common modes.

\subsection{Direction of the next native estimate and verification}

The exact native counterpart of the return is still
\[
 \mathcal I_\gamma^*T_{\rm full,\gamma}
        (I-\mathcal I_\gamma\mathcal I_\gamma^*)
                         T_{\rm full,\gamma}\mathcal I_\gamma.
\]
V3 shows why asking for a small relative return on the entire
static dressed carrier is not a promising unrestricted
Gaussian target: \eqref{f16v3:HS-loss} goes the opposite way.
A physical low-energy restriction would require its own
estimate; it is not supplied by this obstruction.

The constructive replacement is the unprojected fast-history
integral. Its Gaussian precision remains uniformly coercive
even though the retained bridge dynamics has soft directions.
The next useful native theorem must control the non-Gaussian
conditional fast integral, its determinant or pressure
normalization, and its bridge-field tails in a norm compatible
with the physical Gauss restriction and the same full transfer
top value. The compact action is not globally convex in a
single chart, and V2's constants depend on volume. Consequently
\eqref{f16v3:memory-tail} is not yet a native uniform cluster or
spectral estimate. Physical-time calibration and an interacting
continuum construction remain additional open requirements.

The new diagnostic is
\begin{center}\small
\path{scripts/verify_ym_f16_v3_fast_history.py}.
\end{center}
Its 2,075 exact rational fixtures check the noncommuting
propagated precision and source, unequal-mean overlap
completion, full-history action and determinant recursion,
one-step agreement with V2, finite-range propagation, and
length-independent memory tails. They also check the rational
constants in the extensive leakage bound. The inherited V2
chain passed. Finite fixtures supplement the proofs and do
not certify a physical mass gap.
The checker has 7,716 bytes and SHA--256
\begin{center}\scriptsize\ttfamily
811A66019E7A8B4A78D6CF8C1F8F6B2882D318BBA5D2B57D6E4D5C274275EA02.
\end{center}

The predecessor V2 has 48,007 bytes and SHA--256
\begin{center}\scriptsize\ttfamily
D8BC72F2B7E92E94C2B8E3AF848D9580563E5F9302CA13E8843909C42AB80F6F.
\end{center}
Its body is preserved byte for byte: remove this append and
restore the title to recover it. The 168 previously protected
files are unchanged. The last companion checkpoint is f15 V100;
the next is f16 V5. Only \texttt{.tex} files belong in
\texttt{latex\_notes}; build products remain elsewhere.

\begin{firewall}
V3 proves extensive static dressed-line loss and a conditional,
unprojected fast-history formula with uniform Gaussian
space--time locality and curvature-relative temporal truncation.
The leakage obstruction is before physical Gauss restriction.
No native volume-uniform memory bound, nontrivial interacting
continuum construction, or positive physical Yang--Mills
mass gap is established. The full goal remains unachieved.
\end{firewall}

% BEGIN F16 V4 APPEND -- 2026-09-17
% Insert immediately before \end{document} in the supplied V3.
% This append uses only the theorem environments and packages of V3.

\clearpage
\section{V4: a native fast-history core and the nonlinear common-bridge return}
\label{f16v4:context}

V3 supplies an arbitrary-length Gaussian fast-history integral and rules
out a small relative return on the entire static dressed line before
root restriction. Repeating a one-step compression cannot repair that
obstruction. We therefore keep the nonroot temporal gauge variables as
\emph{additional fast variables}, rather than first integrating them into
a logarithmic kinetic kernel. This preserves the finite-range structure
of the exact nonlinear action.

The result below is a native, normalized, \emph{restricted} conditional
history theorem, uniform in spatial volume and history length. In
addition to exponential decay of its two-source temporal response, it
identifies an extra nonlinear small factor for distant common-bridge
sources. A separate pressure-density estimate pays the normalization of
the restricted integral. The passage from this core to the full compact
conditional remains an explicit, unestimated correction; it is not
absorbed into a choice of vacuum normalization.

\subsection{Dependency and non-duplication ledger}

The background audit used the four supplied sources: this lineage's V3,
f14 V100, f15 V100, and the Grand Tensor Monograph V076. It consisted of
section/result inventories and targeted reading of the relevant
interfaces, not a certification of every proof in those archives.
The following distinctions are essential.

\begin{center}\small
\begin{tabular}{@{}p{0.22\textwidth}p{0.72\textwidth}@{}}
\toprule
Inherited interface & Use here, without claiming it again as a new result\\
\midrule
f14 V64 & Nonlinear high-fibre concentration on a restricted convex
coordinate domain. This is not coverage of that domain by the full law.\\[1mm]
f15 V13 & Fixed product-ball conditionals, boundary-aware Witten
covariance, and weighted spatial response. We reuse its absolute
one-form realization.\\[1mm]
Monograph V076 & The direct-Hessian-minus-covariance identity and the
Witten representation, especially \texttt{thm:YM-Witten-router}. A return
of already integrated variables is charged once.\\[1mm]
f15 V97--V100 & Exact cube geometry, coupled compact kernels, moving
frames, and the magnetic dressed carrier.\\[1mm]
This lineage V2--V3 & The exact two-ended kernel, kinetic completion,
common-bridge cancellation, and Gaussian history precision.\\
\bottomrule
\end{tabular}
\end{center}

The new specialization is the complete nonlinear \emph{transfer-time}
conditional with the temporal port variables retained. Its two-source
common-bridge memory requires two small nonlinear couplings. Neither a
new abstract covariance identity nor a further Gaussian determinant is
being offered as the advance. No claim of general literature novelty is
made for the analytic tools.

\subsection{The complete native history density, including its scalar}

Work first before the remaining root Gauss restriction, with the exact
moving coordinates of Proposition~\ref{f16v2:exact-cut}. Fix
\(n\ge1\) and a bridge history \(\mathbf Z=(Z_0,\ldots,Z_n)\).
There is an internal chord configuration \(X_i\) at every slice and a
nonroot temporal gauge configuration \(r_t\in\SU(2)^{7B}\) on every
bond \(t\to t+1\). The root entries of \(r_t\) equal the identity.
Define the exact kinetic words
\begin{align}
 Q_{t,e}^{\rm int}
 &=r_{t,s}^{-1}\bar X_{t,e}r_{t,t(e)}\bar X_{t+1,e}^{-1},\notag\\
 Q_{t,e}^{\rm br}
 &=r_{t,s}^{-1}Z_{t,e}r_{t,t(e)}Z_{t+1,e}^{-1}.
 \label{f16v4:words}
\end{align}
Here \(s=s(e)\), \(t(e)\) denotes the spatial target vertex (not the
bond index), and \(\bar X\) has exactly V2's moving-frame meaning.
Let \(Q_{t,e}\) mean the appropriate word for any of the \(24B\)
fine spatial links.

Set
\[
 w_0=w_n=\tfrac12,\qquad w_i=1\quad(0<i<n).
\]
The factors in the native fixed-history integral are: \(m_\gamma^2\)
at every slice; \(b_\gamma\) at the two ends and
\(b_\gamma^2\) in the interior; one \(\Phi_\gamma\) at each end;
and all \(24Bn\) Wilson kinetic factors. Thus there is no internal
projection at an intermediate time.

We record the chart constants because subsequent logarithmic estimates
must refer to this same integral. In the convention of V1,
\[
 dH(\operatorname{Exp}(q/\sqrt\gamma))
  =c_{H,\gamma}j(q/\sqrt\gamma)\,dq,\qquad
 c_{H,\gamma}=\frac1{2\pi^2\gamma^{3/2}},\qquad
 j(a)=\left(\frac{\sin|a|}{|a|}\right)^2,
\]
on \(|q|<\pi\sqrt\gamma\), up to the Haar-null cut set. Write
\(X_i=\operatorname{Exp}(x_i/\sqrt\gamma)\),
\(r_t=\operatorname{Exp}(z_t/\sqrt\gamma)\), and
\(Z_i=\operatorname{Exp}(y_i/\sqrt\gamma)\), componentwise.
All path matrices below act with the colour factor \(I_3\), which is
suppressed.

The normalized full Euclidean cube amplitude and its exact chart mass are
\begin{align}
 \Omega(x)&=c_\Omega e^{-x^{\mathsf T}C_B^{-1}x/4},&
 c_\Omega&=(2\pi)^{-15B/4}\det(C_B)^{-3/4},\notag\\
 p_{\gamma,B}&=\int_{\mathcal P_x}\Omega(x)^2\,dx,&
 \Phi_\gamma(X(x))&=
       \frac{\Omega(x)}{\sqrt{p_{\gamma,B}J_{X,\gamma}(x)}}.
 \label{f16v4:amplitude}
\end{align}
Here \(\mathcal P_x\) is the complete product chord chart and
\(J_{X,\gamma}\) its Haar Jacobian. This is the carrier already chosen
in V1, not a newly substituted native vacuum.

On the complete internal and port chart product \(\mathcal P\), define
\begin{align}
 V_\gamma(x,z;\mathbf y)
 ={}&\tfrac14\bigl(x_0^{\mathsf T}C_B^{-1}x_0
                       +x_n^{\mathsf T}C_B^{-1}x_n\bigr)\notag\\
 &+\gamma\sum_{i=0}^n w_i S_{\rm int}(X_i)
   +\gamma\sum_{i=0}^n
             S_\times(X_i,\Theta_{X_i}^{-1}Z_i)\notag\\
 &+\gamma\sum_{t=0}^{n-1}\sum_{e\ {\rm fine}}s(Q_{t,e})\notag\\
 &-\sum_{i=0}^n w_i\sum_{a\ {\rm chord}}
                              \log j(x_{i,a}/\sqrt\gamma)
   -\sum_{t=0}^{n-1}\sum_{v\ {\rm nonroot}}
                              \log j(z_{t,v}/\sqrt\gamma).
 \label{f16v4:potential}
\end{align}
In particular, the nonlinear Wilson words and Haar factors are not
replaced by their Taylor polynomials in this definition. Put
\begin{align}
 Z_{\gamma,\mathrm{full}}(\mathbf y)
       &=\int_{\mathcal P}e^{-V_\gamma(x,z;\mathbf y)}\,dx\,dz,\notag\\
 J_{{\rm br},\gamma}(\mathbf y)
       &=\prod_{i=0}^n\prod_{e\ {\rm bridge}}
                         j(y_{i,e}/\sqrt\gamma)^{w_i},\notag\\
 c_{\gamma,B,n}
       &=p_{\gamma,B}^{-1}c_\Omega^2
                         (c_{H,\gamma}/z_\gamma)^{24Bn}.
 \label{f16v4:normalizers}
\end{align}

\begin{proposition}[Exact unprojected native history formula]
\label{f16v4:exact-history}
The full fixed-bridge-history weight in Euclidean bridge half-density
coordinates is
\begin{equation}
 \mathcal W_{\gamma,n}^{\mathrm{full}}(\mathbf y)
 =\frac{c_{\gamma,B,n}J_{{\rm br},\gamma}(\mathbf y)
                   Z_{\gamma,\mathrm{full}}(\mathbf y)}
        {\sqrt{F_{\gamma,2}(Z_0)F_{\gamma,2}(Z_n)}}.
 \label{f16v4:full-weight}
\end{equation}
Integrating \(y_1,\ldots,y_{n-1}\) over their complete principal charts
gives the chart kernel of
\(\mathcal I_\gamma^*T_{\rm full,\gamma}^{\,n}\mathcal I_\gamma\).
This is an identity at finite \(B,n,\gamma\), not a weak-field limit.
\end{proposition}

\begin{proof}
Compose the exact full kernel of Proposition~\ref{f16v2:exact-cut}
\(n\) times and attach the normalized amplitudes from V1 at the two
ends. Each interior magnetic multiplier occurs twice. At an external
slice its second occurrence comes from the dressed amplitude. This
produces precisely the internal and mixed action powers in
\eqref{f16v4:potential}.

At an external chord slice, the factor \(\Phi_\gamma\) cancels half
of the Haar Jacobian; interior chord slices retain the whole Jacobian.
The two external bridge half-densities and the interior bridge
integration measures have the same weights \(w_i\). Since
\(\sum_iw_i=n\), the constant Haar powers are \(5Bn\) from chords,
\(7Bn\) from temporal ports, and \(12Bn\) from bridges. Their total
is \(24Bn\), exactly the number of kinetic normalizers
\(z_\gamma^{-1}\). The two endpoint amplitudes contribute
\(p_{\gamma,B}^{-1}c_\Omega^2\). The nonconstant Haar factors are
those in \eqref{f16v4:potential} and \(J_{{\rm br},\gamma}\).
This proves the scalar as well as the action. Positivity justifies
composition and changes in integration order. The principal charts
cover the compact variables apart from null sets.
\end{proof}

\subsection{The augmented Gaussian reference and a fixed nonlinear core}

Avoid the notational collision with V3's endpoint matrix \(E\) by
writing \(E_{\rm br}\) for V2's bridge--nonroot incidence. Set
\[
 L_c=\mathsf B_B^{\mathsf T}\mathsf B_B,\qquad
 L_h=L_c+E_{\rm br}^{\mathsf T}E_{\rm br},\qquad
 \widehat L=L_c^{-1/2}L_hL_c^{-1/2}.
\]
All powers of \(L_c\) are spatially cube-block diagonal. Use the fast
coordinates
\[
       u_i=N^{1/2}x_i,\qquad v_t=L_c^{1/2}z_t,\qquad \xi=(u,v).
\]
Each \(u_{i,b}\) has dimension 15 and each \(v_{t,b}\) dimension 21.
Changing between \((x,z)\) and \(\xi\) has a constant Jacobian.
Partition functions below always use the original \(dx\,dz\) measure;
their normalized laws are pushed forward to \(\xi\). The constant
Jacobian cancels in those laws and in ratios with the same Gaussian
reference.

The quadratic reference potential, including its bridge-only part, is
\begin{align}
 V_0(x,z;\mathbf y)
 ={}&\tfrac14\bigl(x_0^{\mathsf T}C_B^{-1}x_0
                         +x_n^{\mathsf T}C_B^{-1}x_n\bigr)
       +\tfrac12\sum_{i=0}^n w_i x_i^{\mathsf T}\mathsf H_Bx_i
       +\tfrac12\sum_{i=0}^n|Ax_i+Dy_i|^2\notag\\
 &+\tfrac12\sum_{t=0}^{n-1}
       \left(\left|U_B(x_t-x_{t+1})+\mathsf B_Bz_t\right|^2
              +\left|y_t-y_{t+1}+E_{\rm br}z_t\right|^2\right).
 \label{f16v4:V0}
\end{align}
Because \(U_B^{\mathsf T}\mathsf B_B=0\), the fast Hessian in these
coordinates is the direct sum
\begin{equation}
 H_0:=D_\xi^2V_0
       =\widehat W_n\oplus(I_n\otimes\widehat L),\qquad
                         I\le H_0\le30I.
 \label{f16v4:H0}
\end{equation}
Indeed V3 gives \(4I\le\widehat W_n\le15I\); V2's bound
\(E_{\rm br}L_c^{-1}E_{\rm br}^{\mathsf T}\le29I\) gives
\(I\le\widehat L\le30I\). The \(u\)-diagonal time blocks are at
least \(p_*I\), where \(p_*=3+2\sqrt2\), and the only off-diagonal
time blocks of \(H_0\) are the adjacent \(u\)-blocks \(-I\).
There are no Gaussian \(u\)--\(v\) cross blocks.

Integrating all \(z_t\) in \eqref{f16v4:V0} by V2's kinetic completion
recovers V3's \(W_n\), bridge covariance \(G\), sources and determinant.
Equivalently, if \(Z_0=\int e^{-V_0}\,dx\,dz\) over the full Euclidean
fast space and \(c_{0,B,n}=c_\Omega^2(2\pi)^{-36Bn}\), then
\begin{equation}
       \mathcal W_n(\mathbf y)
        =\frac{c_{0,B,n}Z_0(\mathbf y)}
                       {\sqrt{F_2(y_0)F_2(y_n)}}.
 \label{f16v4:gaussian-consistency}
\end{equation}
This is a consistency check against Theorem~\ref{f16v3:history}, not a
new determinant theorem.

For \(R\ge1\) choose the \emph{separate-ball} product
\begin{equation}
 \mathcal O_R=
    \prod_{i=0}^n\prod_b\{|u_{i,b}|<R\}
       \ \times\!
    \prod_{t=0}^{n-1}\prod_b\{|v_{t,b}|<R\}.
 \label{f16v4:core}
\end{equation}
The two species must not be combined into a single ball: a combined
boundary would introduce a one-form boundary coupling absent from the
native action. The domain is fixed as the bridge history varies.
Let
\begin{equation}
 Z_{\gamma,R}(\mathbf y)=\int_{\mathcal O_R}e^{-V_\gamma}\,dx\,dz,
 \qquad
 \nu_{\gamma,R,\mathbf y}
       =\text{the normalized law on \(\mathcal O_R\) in \(\xi\)}.
 \label{f16v4:core-law}
\end{equation}
The notation \(\int_{\mathcal O_R}dx\,dz\) means integration over the
inverse linear image of \(\mathcal O_R\). Whenever this image is inside
\(\mathcal P\), this is \emph{exactly} the native augmented conditional
further conditioned on that event.

\begin{proposition}[Uniform native derivative certificate on the core]
\label{f16v4:derivatives}
There are finite constants \(C_*,C_1\ge1\), determined only by the
single-cube matrices and the finite list of native local words, such
that the following holds for every \(B,n\). Suppose
\begin{equation}
 \max_{i,e}|y_{i,e}|\le r,\qquad
 R\ge\max(1,r),\qquad
 \delta:=C_*R/\sqrt\gamma\le1/20.
 \label{f16v4:regime}
\end{equation}
The core is then inside the principal fast charts, and on its closure
\begin{align}
 |V_\gamma-V_0|&\le C_1B(n+1)R^3/\sqrt\gamma,\notag\\
 \|D_{(\xi,\mathbf y)}^2(V_\gamma-V_0)\|&\le\delta,\notag\\
       \tfrac{19}{20}I\le D_\xi^2V_\gamma&\le\tfrac{601}{20}I.
 \label{f16v4:derivative-bounds}
\end{align}
For a bridge direction \(h\) at a single time \(i\), put
\(J_{i,h}=\partial_{i,h}V_\gamma\). Then
\begin{equation}
       \|\nabla_\xi J_{i,h}\|_{L^\infty(\mathcal O_R)}
                                      \le12\|h\|.
 \label{f16v4:source-bound}
\end{equation}
Its derivative-index time support is contained in \(\{i-1,i\}\),
where unavailable indices are omitted and \(v_t\) is assigned time
\(t\). The exact Hessian has no \(v_t\)--\(v_s\) block for
\(t\ne s\); a \(u\)--\(v_t\) block can occur only at \(u_t,u_{t+1}\).
Finally, for V1's common-bridge space \(\mathcal E\),
\begin{equation}
 h\in\mathcal E\quad\Longrightarrow\quad
 \|\nabla_uJ_{i,h}\|_\infty\le\delta\|h\|,
              \qquad \|D_{uv}^2V_\gamma\|\le\delta.
 \label{f16v4:common-source}
\end{equation}
Here the second inequality does not require \(h\in\mathcal E\).
\end{proposition}

\begin{proof}
We give the local-to-global estimate so that a volume-dependent Taylor
constant is not being hidden in \(C_*\). After the fixed cube-linear
changes of variables, each action summand is a member of a finite list
of functions
\[
        a\longmapsto s\!\left(\prod_{\ell=1}^{k}
                       \operatorname{Exp}(T_\ell a)\right),
\]
where the product order, the signs and the matrices \(T_\ell\) are the
actual internal, mixed or kinetic words. Their lengths and dimensions
are bounded independently of \(B,n\). The moving frames are themselves
exponentials of fixed linear forms in the five chords of one cube.
Each summand has bounded third derivatives on a fixed neighbourhood of
zero. Its value and gradient vanish at zero, and its quadratic term is
one half of the squared sum of its signed linear arguments. Haar terms
\(-\log j\) have bounded derivatives on that neighbourhood, with value
and gradient zero at zero.

For completeness, constants can be fixed without estimating them
numerically. Choose a radius \(a_0>0\) such that all factors in this
finite word list, after the fixed linear changes, lie in their smooth
principal neighbourhood whenever every input block has norm at most
\(a_0\). Take the maxima of the third derivative norms of the action
words and the first three derivative norms of \(-\log j\) on the
closures of those local input sets. Multiply these finite maxima by
the maximal local input dimension and the maximal number of summands
incident to an input block. Choose \(C_*\) larger than the resulting
Hessian row-sum constant and \((20a_0)^{-1}\), and choose \(C_1\)
larger than the analogous value-sum constant. These constants use only
a finite local library, not a supremum over growing lattices.

Taylor's integral remainder for an action term, rescaled as
\(\gamma f(a/\sqrt\gamma)\), bounds its value remainder by
\(CR^3/\sqrt\gamma\) and its Hessian remainder by
\(CR/\sqrt\gamma\). The Haar contributions are bounded by
\(CR^2/\gamma\) and \(C/\gamma\), respectively, which fit those
bounds under \eqref{f16v4:regime}. The endpoint Gaussian terms are
already exactly quadratic. Every coordinate block meets only a fixed
number of terms: internal spatial and temporal words stay inside a
cube or an adjacent pair of times; a mixed two-bridge plaquette meets
two neighbouring cubes; a four-bridge plaquette has no fast variable;
a bridge kinetic word meets two neighbouring cubes on one time bond.
There are \(O(B(n+1))\) summands altogether. Summing the value bounds
and using the symmetric block row-sum bound for the Hessian proves
the first two estimates. Equation~\eqref{f16v4:H0} gives the third.
The chosen radius puts every individual \(x\) and \(z\) inside its
principal chart.

The Gaussian mixed source derivative has a \(u_i\) block
\(N^{-1/2}\Lambda h\), and \(v_{i-1},v_i\) blocks of the form
\(\pm L_c^{-1/2}E_{\rm br}^{\mathsf T}h\). The inherited bounds
\(\|N^{-1/2}A^{\mathsf T}\|\le2\), \(\|D\|\le4\), and
\(E_{\rm br}L_c^{-1}E_{\rm br}^{\mathsf T}\le29I\) therefore give
\[
 \|\nabla_\xi\partial_{i,h}V_0\|^2
              \le(64+2\cdot29)\|h\|^2=122\|h\|^2.
\]
The mixed Hessian remainder is at most \(\delta\), and
\(\sqrt{122}+1/20<12\), proving \eqref{f16v4:source-bound}.
The same argument and \(\Lambda h=0\) for \(h\in\mathcal E\)
prove the first part of \eqref{f16v4:common-source}; the Gaussian
\(u\)--\(v\) block vanishes, proving its second part.
All support assertions follow from the exact words, before Taylor
expansion. A temporal port variable occurs only on its own bond;
bridge differentiation meets \(u_i,v_{i-1},v_i\), and no other fast
time indices. This completes the certificate.
\end{proof}

The constants \(C_*,C_1\) are defined finite derivative constants, not
numerically optimized values. In particular, the small-field regime is
proved to exist; no practical numerical threshold for \(\gamma\) is
claimed. A bound on \(D_\xi^2V_\gamma\) outside this core has not been
used or established.

\subsection{Boundary-aware temporal response without a length loss}

We reuse f15 V13's absolute one-form covariance realization and the
monograph's Witten router. To fix the domain used here, let
\(\nu=\nu_{\gamma,R,\mathbf y}\), \(H(\xi)=D_\xi^2V_\gamma(\xi)\),
and let \(\mathfrak H\) have the closed quadratic form
\begin{equation}
 \mathfrak h[a]
  =\int|\nabla a|^2\,d\nu+\int a^{\mathsf T}H a\,d\nu
       +\sum_{\text{ball faces }\alpha}
             \int_{\partial_\alpha\mathcal O_R}
                    \mathrm{II}_\alpha(a_\alpha,a_\alpha)\,d\nu_\partial.
 \label{f16v4:witten-form}
\end{equation}
The component \(a_\alpha\) is tangential at its own ball boundary;
\(\mathrm{II}_\alpha\ge0\) is the outward second fundamental form.
This is the product absolute realization, including its corners. In
particular the derivative and boundary part, denoted \(\mathfrak D\),
is nonnegative and block diagonal in the individual coordinate-ball
labels. Multiplication by a constant projection onto any union of these
labels preserves its form domain. It need not be the same scalar
boundary realization for every component.

In this realization the inherited identity is
\begin{equation}
 \operatorname{Cov}_{\nu}(F,G)
       =\langle\nabla F,\mathfrak H^{-1}\nabla G\rangle_{L^2(\nu)}.
 \label{f16v4:covariance-identity}
\end{equation}
One obtains it by solving the mean-zero reflecting scalar equation,
using the scalar/one-form intertwining, and integrating by parts. The
strict lower bound in \eqref{f16v4:derivative-bounds} makes the one-form
inverse bounded. Smooth observables need not satisfy a boundary
condition. The product-ball construction and its nonnegative boundary
form are exactly the inherited boundary-aware mechanism; a primary
reference for the smooth-boundary Witten decomposition is
\href{https://arxiv.org/abs/1607.08714}{D. Le Peutrec,
\emph{On Witten Laplacians and Brascamp--Lieb's inequality on manifolds
with boundary}}. We now exploit the more specific temporal block
structure of the current native law.

Assign both \(u_i\) and \(v_i\), when present, the index-time \(i\).
For a set of such times \(I\), let \(\Pi_I\) project a vector field
onto those coordinate components. This is a projection in
\emph{derivative-index space}; the scalar coefficient functions of a
vector field may depend on every fast variable.

\begin{theorem}[Uniform nonlinear fast-history response]
\label{f16v4:temporal-inverse}
Under \eqref{f16v4:regime}, put \(\vartheta=9/19\). For any index-time
sets \(I,J\) at distance \(d\),
\begin{equation}
          \|\Pi_I\mathfrak H^{-1}\Pi_J\|
                                  \le2\vartheta^d.
 \label{f16v4:temporal-inverse-bound}
\end{equation}
There is also a space--time estimate. Give each \((i,b)\) the two
coordinate species when available, and use nearest spatial-cube and
nearest-time edges. For sets \(A_0,B_0\) of these labels,
\begin{equation}
 \|\Pi_{A_0}\mathfrak H^{-1}\Pi_{B_0}\|
       \le\frac{20}{19}
                 \left(\frac{601}{620}\right)^{d(A_0,B_0)}.
 \label{f16v4:spacetime-inverse-bound}
\end{equation}
Neither estimate has a cardinality, volume or duration prefactor.
Consequently \eqref{f16v4:temporal-inverse-bound} bounds a covariance
by \(2\vartheta^d\|\nabla F\|_2\|\nabla G\|_2\) whenever the two
gradients have the indicated index-time supports.
\end{theorem}

\begin{proof}
Let \(D_*\) be the constant coordinate operator equal to \(p_*I\)
on all \(u\)-components and to \(I\) on all \(v\)-components.
Write
\[
 \mathfrak H=\mathfrak A-\mathfrak C,\qquad
 \mathfrak A=\mathfrak D+\operatorname{diag}_{\rm time}H,\qquad
 \mathfrak C=-(H-\operatorname{diag}_{\rm time}H).
\]
The projection onto time-diagonal matrix blocks contracts operator
norm. Equation~\eqref{f16v4:H0} and the Hessian remainder give
\[
 \mathfrak A\ge(1-\delta)D_*,\qquad
 \|D_*^{-1/2}\mathfrak C D_*^{-1/2}\|
                   \le\frac2{p_*}+2\delta.
\]
For the second assertion, the Gaussian time-off-diagonal operator is
the path adjacency, with norm at most two, on the \(u\) species. The
remainder's off-diagonal part has norm at most \(2\delta\). The
nonnegative derivative and boundary terms preserve time blocks, so
\(\mathfrak A^{-1/2}\) preserves them too. Since \(2/p_*<7/20\),
\[
 \|\mathfrak A^{-1/2}\mathfrak C\mathfrak A^{-1/2}\|
       \le\frac{2/p_*+2\delta}{1-\delta}
       \le\frac9{19}=\vartheta<1.
\]
Bounded form perturbation and the norm-convergent Neumann series give
\[
 \mathfrak H^{-1}=\mathfrak A^{-1/2}
   \sum_{k\ge0}(\mathfrak A^{-1/2}\mathfrak C\mathfrak A^{-1/2})^k
                                                   \mathfrak A^{-1/2}.
\]
Every factor \(\mathfrak C\) changes the time label by at most one.
Terms with \(k<d\) therefore vanish between the stated projections.
Using \(\|\mathfrak A^{-1/2}\|^2\le20/19\), the tail is bounded by
\[
               \frac{20}{19}\frac{\vartheta^d}{1-\vartheta}
                           =2\vartheta^d.
\]
This proves \eqref{f16v4:temporal-inverse-bound} without assuming that
the off-diagonal operator is positive.

For the second estimate use instead
\(\mathfrak A_0=\mathfrak D+31I\) and
\(\mathfrak C_0=31I-H\). The derivative certificate implies
\(0\le\mathfrak C_0\le(601/20)I\), so
\[
 \|\mathfrak A_0^{-1/2}\mathfrak C_0\mathfrak A_0^{-1/2}\|
                                      \le601/620.
\]
The inverse square root is diagonal in all cube--time labels, and
\(\mathfrak C_0\) has range one on the displayed graph. The same
series argument has prefactor
\(31^{-1}/(1-601/620)=20/19\).
The locality assertion is about vector-component labels, not finite
propagation of the scalar diffusion in configuration space. Finally
apply \eqref{f16v4:covariance-identity} and Cauchy--Schwarz.
\end{proof}

Let
\[
 \mathcal A_{\gamma,R}(\mathbf y)=-\log Z_{\gamma,R}(\mathbf y),
 \qquad
 \mathcal W_{\gamma,n}^{R}
  =\frac{c_{\gamma,B,n}J_{{\rm br},\gamma}Z_{\gamma,R}}
            {\sqrt{F_{\gamma,2}(Z_0)F_{\gamma,2}(Z_n)}}.
\]
These definitions retain the full native endpoint normalizers.
The domain \(\mathcal O_R\) is independent of \(\mathbf y\), so
ordinary differentiation under its finite integral gives
\begin{equation}
 \partial_{i,h}\partial_{j,k}\mathcal A_{\gamma,R}
   =\mathbb E_\nu\partial_{i,h}\partial_{j,k}V_\gamma
         -\operatorname{Cov}_\nu(J_{i,h},J_{j,k}).
 \label{f16v4:pressure-hessian}
\end{equation}
There is no moving-boundary derivative in this formula. It is the
inherited exact response identity, not a second return to be subtracted
from an already integrated action.

\begin{proposition}[Two-source temporal memory of the native core]
\label{f16v4:memory}
Under \eqref{f16v4:regime}, for \(|i-j|\ge2\),
\begin{equation}
 |\partial_{i,h}\partial_{j,k}\mathcal A_{\gamma,R}|
       \le288\,\vartheta^{|i-j|-1}\|h\|\|k\|.
 \label{f16v4:generic-memory}
\end{equation}
The same statement holds with
\(\mathcal A_{\gamma,R}\) replaced by
\(-\log\mathcal W_{\gamma,n}^R\).
If \(K_{ij}\) denotes this bridge Hessian block, then for every
integer \(\ell\ge1\), and every sequence of bridge variations,
\begin{equation}
 \left|\sum_{|i-j|>\ell}\langle h_i,K_{ij}h_j\rangle\right|
       \le\frac{576}{1-\vartheta}\vartheta^\ell
                                         \sum_i\|h_i\|^2.
 \label{f16v4:generic-tail}
\end{equation}
This is a Hessian-tail estimate at the specified bridge history,
not an assertion about every higher connected response.
\end{proposition}

\begin{proof}
Every term of \(V_\gamma\) depends on at most two adjacent bridge
times. The direct mixed derivative in
\eqref{f16v4:pressure-hessian} is therefore zero when \(|i-j|\ge2\).
The source gradients lie in time sets \(\{i-1,i\}\) and
\(\{j-1,j\}\), at distance at least \(|i-j|-1\). Apply
Theorem~\ref{f16v4:temporal-inverse} and
\eqref{f16v4:source-bound}; \(2\cdot12^2=288\).
The logarithms of the bridge Jacobian and the two endpoint
normalizers are sums of single-time functions, and the remaining
scalar is history independent. None changes a distinct-time Hessian
block.

For the tail, the scalar majorant for the omitted ordered pairs has
row sum at most
\(2\cdot288\sum_{d>\ell}\vartheta^{d-1}
 =576\vartheta^\ell/(1-\vartheta)\).
The symmetric Schur bound, or \(2ab\le a^2+b^2\) term by term,
gives \eqref{f16v4:generic-tail}.
\end{proof}

\subsection{Distant common-bridge response needs two nonlinear couplings}

The generic estimate does not use \(\Lambda\mathcal E=0\).
That cancellation is not enough by itself to eliminate the temporal
ports: even at Gaussian order a bridge variation differentiates their
kinetic sources. The additional step is to retain their time-diagonal
one-form inverse before taking the \(u\)-Schur complement.

\begin{theorem}[A small common-bridge nonlocal return in the native core]
\label{f16v4:common-memory}
Assume \eqref{f16v4:regime} and set
\[
             K_{\rm c}=1+\frac{240}{19}=\frac{259}{19}.
\]
For \(h,k\in\mathcal E\) and \(|i-j|\ge2\),
\begin{equation}
 |\partial_{i,h}\partial_{j,k}\mathcal A_{\gamma,R}|
  \le2K_{\rm c}^{\,2}\delta^2
             \vartheta^{(|i-j|-2)_+}\|h\|\|k\|.
 \label{f16v4:common-memory-bound}
\end{equation}
The base bridge history need only satisfy its bound in
\eqref{f16v4:regime}; it need not lie in \(\mathcal E\).
The same estimate holds for \(-\log\mathcal W_{\gamma,n}^R\).
In particular the prefactor is \(O(R^2/\gamma)\), not merely order
one. For \(h_i\in\mathcal E\) and \(\ell\ge1\),
\begin{equation}
 \left|\sum_{|i-j|>\ell}\langle h_i,K_{ij}h_j\rangle\right|
  \le\frac{4K_{\rm c}^{\,2}}{1-\vartheta}\delta^2
                     \vartheta^{\ell-1}\sum_i\|h_i\|^2.
 \label{f16v4:common-tail}
\end{equation}
\end{theorem}

\begin{proof}
Decompose the one-form Hilbert space by species, not by scalar
conditional expectation:
\[
       \mathfrak H=
       \begin{pmatrix}\mathfrak H_{uu}&B_{uv}\\
                       B_{uv}^*&\mathfrak H_{vv}\end{pmatrix}.
\]
The derivative and boundary forms do not mix these species because
their balls are separate. Thus \(B_{uv}=D_{uv}^2V_\gamma\) is a
bounded multiplication operator with norm at most \(\delta\).
The exact native words imply that \(\mathfrak H_{vv}\) is block
diagonal in \emph{bond time}. Its inverse preserves those labels and
has norm at most \(20/19\).

Define the positive Schur operator in the form sense by
\[
 S_u=\mathfrak H_{uu}
             -B_{uv}\mathfrak H_{vv}^{-1}B_{uv}^*.
\]
For any vector fields \(g=(g_u,g_v)\), \(g'=(g'_u,g'_v)\), block
completion of the coercive form gives the exact identity
\begin{align}
 \langle g,\mathfrak H^{-1}g'\rangle
   ={}&\langle g_v,\mathfrak H_{vv}^{-1}g'_v\rangle\notag\\
     &+\langle g_u-B_{uv}\mathfrak H_{vv}^{-1}g_v,
            S_u^{-1}(g'_u-B_{uv}\mathfrak H_{vv}^{-1}g'_v)\rangle.
 \label{f16v4:species-schur}
\end{align}
All off-diagonal couplings are bounded, so this follows also by the
usual block triangular factorization on form domains. In particular
\(S_u^{-1}\) is precisely the \(uu\) corner of
\(\mathfrak H^{-1}\), not an independently normalized inverse.

Take \(g=\nabla J_{i,h}\). For common \(h\),
\eqref{f16v4:common-source} and \eqref{f16v4:source-bound} give
\[
 \|g_u-B_{uv}\mathfrak H_{vv}^{-1}g_v\|_2
       \le\left(1+\frac{20}{19}\,12\right)\delta\|h\|
       =K_{\rm c}\delta\|h\|.
\]
Its \(u\)-index support is contained in \(\{i-1,i,i+1\}\).
Indeed \(g_v\) is supported on bonds \(i-1,i\), the inverse
\(\mathfrak H_{vv}^{-1}\) preserves each bond label, and
\(B_{uv}\) only connects bond \(t\) to \(u_t,u_{t+1}\).
The same assertions hold for \(g'\) at \(j\).

For \(|i-j|\ge2\) the first term of
\eqref{f16v4:species-schur} vanishes: its two \(v\)-index supports
are disjoint. The corrected \(u\)-supports are separated by at least
\((|i-j|-2)_+\). Since \(S_u^{-1}\) is a corner of the full inverse,
Theorem~\ref{f16v4:temporal-inverse} bounds the remaining term by the
right side of \eqref{f16v4:common-memory-bound}.
The direct bridge mixed derivative is zero as before, so
\eqref{f16v4:pressure-hessian} proves the assertion. The normalizer
terms remain single-time. Summing the omitted scalar majorant gives
row sum
\[
  4K_{\rm c}^{\,2}\delta^2
       \sum_{d=\ell+1}^{\infty}\vartheta^{d-2}
    =\frac{4K_{\rm c}^{\,2}}{1-\vartheta}
                       \delta^2\vartheta^{\ell-1},
\]
which proves the tail estimate.
\end{proof}

The two small factors have different possible origins but the same
controlled size: a common source can acquire a small direct
\(u\)-component, or its order-one \(v\)-component can reach \(u\)
through the small \(uv\) coupling. To connect two nonadjacent source
times, both ends must pay one such factor. This conclusion uses the
exact nonlinear core; it is not obtained by declaring its covariance
Gaussian. It does not control the sign of a same-time restoring term,
and it is not a statement of physical gaplessness or of a mass gap.

\subsection{A paid pressure-density comparison for the restricted integral}

A covariance estimate does not fix the scalar part of a path integral.
We pay that scalar separately, without claiming a volume-independent
probability that every fast coordinate lies in the core.
Let \(\mathbb P_{0,\mathbf y}\) be the normalized Gaussian law with
potential \(V_0\), in the whitened coordinates \(\xi\), and let
\(\mu(\mathbf y)\) be its mean.

\begin{proposition}[A uniform local mean bound for the Gaussian reference]
\label{f16v4:gaussian-mean}
There is a finite geometric constant \(C_\mu\ge1\), independent of
\(B,n\), such that
\begin{equation}
       \max_\alpha|\mu_\alpha(\mathbf y)|\le C_\mu r
       \quad\hbox{when}\quad \max_{i,e}|y_{i,e}|\le r.
 \label{f16v4:mean-bound}
\end{equation}
Here \(\alpha\) runs over the separate \(u\)- and \(v\)-balls.
Every corresponding marginal Gaussian covariance is at most the
identity.
\end{proposition}

\begin{proof}
The covariance assertion follows from \(H_0\ge I\). The source
\(-D_\xi V_0(0;\mathbf y)\) is a finite-range linear function of the
bridge history and has norm at most \(C r\) on each ball, with a
fixed local constant. To obtain a block row-sum estimate for its
inverse, write
\[
       H_0^{-1}=\frac1{30}\sum_{k\ge0}(I-H_0/30)^k.
\]
The summand has norm at most \((29/30)^k\) and range at most \(k\)
on the cube--time graph. Thus the norm of an inverse block between
labels at distance \(d\) is at most \((29/30)^d\). The number of
labels at distance \(d\) from any fixed label is at most
\(C(d+1)^4\), uniformly on a three-dimensional periodic cube graph
times a finite time path, including the two species. Consequently the
sum of these block norms is at most the finite constant
\(C\sum_{d\ge0}(d+1)^4(29/30)^d\).
Multiplication by the local source bound proves
\eqref{f16v4:mean-bound}. This row-sum argument is needed: the
Euclidean norm bound \(\|H_0^{-1}\|\le1\) alone would give a
spurious volume factor for a uniformly bounded bridge field.
\end{proof}

\begin{theorem}[Core pressure per cube and time slice]
\label{f16v4:pressure}
Assume \eqref{f16v4:regime} and, in addition,
\begin{equation}
                  R\ge\max(12,2C_\mu r).
 \label{f16v4:mean-margin}
\end{equation}
Then
\begin{equation}
 \frac{1}{B(n+1)}
       \left|\log\frac{Z_{\gamma,R}(\mathbf y)}{Z_0(\mathbf y)}\right|
 \le C_1\frac{R^3}{\sqrt\gamma}
                  +4\,2^{21/2}e^{-R^2/16}.
 \label{f16v4:pressure-bound}
\end{equation}
This compares the exact native \emph{core} integral with the complete
Euclidean Gaussian fast integral. It does not compare the complete
compact native integral with that Gaussian integral.
\end{theorem}

\begin{proof}
Let \(a_R=2^{21/2}e^{-R^2/16}\). If \(g\) is a centered Gaussian
of dimension at most 21 and covariance at most \(I\), then
\[
 \mathbb P\{|g|>R/2\}
  \le e^{-R^2/16}\mathbb E e^{|g|^2/4}
  \le2^{21/2}e^{-R^2/16}=a_R.
\]
For \(R\ge12\), \(a_R<1/2\). Each event
\(\{|\xi_\alpha-\mu_\alpha|\le R/2\}\) is a symmetric convex
cylinder for the centered full Gaussian \(\xi-\mu\).
The Gaussian correlation inequality, iterated over these cylinders,
gives
\[
 \mathbb P_{0,\mathbf y}\!
     \left\{\forall\alpha:\ |\xi_\alpha-\mu_\alpha|\le R/2\right\}
                      \ge(1-a_R)^{(2n+1)B}.
\]
This use of the Gaussian correlation inequality is an external theorem,
not an independence assertion; see
\href{https://arxiv.org/abs/1408.1028}{T. Royen,
\emph{A simple proof of the Gaussian correlation conjecture extended to
multivariate gamma distributions}}. The centered event is contained in
\(\mathcal O_R\) by \eqref{f16v4:mean-bound} and
\eqref{f16v4:mean-margin}. Since \(\log(1-a)\ge-2a\) on
\([0,1/2]\),
\begin{equation}
 -4B(n+1)a_R\le
                \log\mathbb P_{0,\mathbf y}(\mathcal O_R)\le0.
 \label{f16v4:gaussian-core-cost}
\end{equation}
There is no lower bound here independent of the number of cells.

Set \(\epsilon_N=C_1B(n+1)R^3/\sqrt\gamma\).
The pointwise action remainder on the core yields
\[
 e^{-\epsilon_N}Z_0\mathbb P_{0,\mathbf y}(\mathcal O_R)
       \le Z_{\gamma,R}
       \le e^{\epsilon_N}Z_0\mathbb P_{0,\mathbf y}(\mathcal O_R).
\]
Take logarithms, use \eqref{f16v4:gaussian-core-cost}, and divide by
\(B(n+1)\). The argument controls the extensive scalar even when the
Gaussian probability of the all-good event is small.
\end{proof}

The remaining explicit scalar factors in the native weight can also be
separated without inventing a spectral normalization. Uniformly in
\(B,n\), for sufficiently large \(\gamma\),
\begin{align}
 \left|\log\frac{c_{\gamma,B,n}}{c_{0,B,n}}\right|
                &\le C Bn/\gamma+CB e^{-c\gamma},\notag\\
 |\log J_{{\rm br},\gamma}(\mathbf y)|
                &\le CBn r^2/\gamma.
 \label{f16v4:explicit-scalar-cost}
\end{align}
Here and below constants may depend on the fixed cube geometry but
not on \(B,n\). To verify the first estimate, rescale the defining
integral for \(z_\gamma\). Taylor expansion of \(1-\cos a\) and
\(j(a)\), with the bound
\(1-\cos a\ge2a^2/\pi^2\) on \([0,\pi]\), gives
\[
 \frac{z_\gamma}{c_{H,\gamma}}
                  =(2\pi)^{3/2}(1+O(\gamma^{-1})).
\]
For a direct domination argument, use Taylor expansion on
\(|q|\le\gamma^{1/8}\); the rescaled remainder is bounded by
\(C\gamma^{-1}(|q|^2+|q|^4)e^{-c|q|^2}\).
The complement is controlled by a Gaussian tail using the displayed
cosine bound. The cube Gaussian chart mass factorizes by cubes and
has \(-\log p_{\gamma,B}\le CB e^{-c\gamma}\), by fixed-dimensional
Gaussian tails for each of the five chord vectors. These facts prove
the first bound of \eqref{f16v4:explicit-scalar-cost};
\(|\log j(a)|\le C|a|^2\) on the small chart proves the second.

In particular the exact compensated comparison is
\begin{align}
 &\log\frac{\mathcal W_{\gamma,n}^{R}(\mathbf y)}
                    {\mathcal W_n(\mathbf y)}
   +\frac12\log\frac{F_{\gamma,2}(Z_0)F_{\gamma,2}(Z_n)}
                       {F_2(y_0)F_2(y_n)}\notag\\
 &\hspace{12mm}
   -\log\frac{c_{\gamma,B,n}}{c_{0,B,n}}
   -\log J_{{\rm br},\gamma}(\mathbf y)
            =\log\frac{Z_{\gamma,R}(\mathbf y)}{Z_0(\mathbf y)}.
 \label{f16v4:compensated-pressure}
\end{align}
The endpoint \(F_{\gamma,2}\) factors are kept literally. Theorem
\ref{f16v4:pressure} does not give a new uniform estimate for their
full compact-to-Gaussian ratios. They do not enter the separated-time
mixed response because their logarithms are single-time functions.

For fixed \(r\) and any fixed \(K>0\), the choice
\(R=K\sqrt{\log\gamma}\) satisfies all the stated conditions once
\(\gamma\) is sufficiently large. It gives a pressure-density error
\[
 O\!\left(\frac{(\log\gamma)^{3/2}}{\sqrt\gamma}
                             +\gamma^{-K^2/16}\right),
\]
while the common-bridge distant-time response has prefactor
\(O(\log\gamma/\gamma)\). These estimates are uniform in \(B,n\)
inside the indicated regime; no total-variation convergence of the
full growing-dimensional native law follows.

\subsection{The exact remaining correction and the next upstream target}

At every finite regulator and admissible bridge history set
\begin{equation}
 p_{\rm good}(\mathbf y)
       =\frac{Z_{\gamma,R}(\mathbf y)}
                        {Z_{\gamma,\mathrm{full}}(\mathbf y)}
       \in(0,1].
 \label{f16v4:good-probability}
\end{equation}
It is the actual probability of the specified joint \(x,z\) core under
the augmented native conditional. Positivity and the exact history
identity give
\begin{equation}
 \mathcal W_{\gamma,n}^{R}
         =p_{\rm good}\mathcal W_{\gamma,n}^{\mathrm{full}},\qquad
 -\log\mathcal W_{\gamma,n}^{\mathrm{full}}
         =-\log\mathcal W_{\gamma,n}^{R}+\log p_{\rm good}.
 \label{f16v4:full-core-split}
\end{equation}
The sign in this equation is important: restricting an integral makes
its negative logarithm larger. For separated bridge times, the full
mixed response is therefore exactly
\begin{equation}
 \partial_{i,h}\partial_{j,k}
            (-\log\mathcal W_{\gamma,n}^{\mathrm{full}})
 =\partial_{i,h}\partial_{j,k}
            (-\log\mathcal W_{\gamma,n}^{R})
          +\partial_{i,h}\partial_{j,k}\log p_{\rm good}.
 \label{f16v4:unpaid-response}
\end{equation}
Smoothness for these bridge derivatives follows at finite regulator
by differentiating the positive compact integral; the core integral
has a fixed domain. No estimate for the last term is being asserted.

This isolates a specific next interface. A useful V5 estimate would
control the \emph{local mixed derivatives} of
\(\log p_{\rm good}\), together with its pressure per cell, by a
native large-field decomposition that retains the temporal-port tails
and the actual bridge sources. It must control this correction rather
than count the same covariance return a second time. A uniform lower
bound on \(p_{\rm good}\) itself is not the appropriate demand:
even the Gaussian all-good probability can decrease exponentially in
\(B(n+1)\) while its pressure density and local response are controlled.
The next argument should therefore test local bad-region weights or
source-visible capacity, not repeat a Gaussian inverse estimate.
The archived pinned-cube compact bounds provide possible ingredients,
but do not by themselves bound this time-history correction.

There are two other boundaries on the claim. First, bounding every
bridge coordinate in an identity chart is not a gauge-invariant
coverage theorem for the physical root quotient. The full native
kernel identity is compatible with restriction to invariant
functions; the present local derivative estimate does not
\emph{automatically} survive that restriction or its boundary-field
integration. Second, \(\mathfrak H\) is an auxiliary conditional
one-form response operator. Its lower bound is not a spectral gap of
\(T_{\rm full,\gamma}\), and it has not been compared to that
transfer's top eigenvalue. Physical-time calibration and an
interacting continuum construction remain separate requirements.

\clearpage
\subsection{Verification, source preservation, and current status}

The newly supplied diagnostic is
\begin{center}\small
\path{verify_ym_f16_v4_native_history.py}.
\end{center}
Its 9,467 exact algebraic assertions passed. They check the cube
orthogonality and spectra; actual parity-cube plaquette and common-mode
incidences on fine tori of sides eight and ten; noncommuting finite
history elimination; endpoint and Haar power counts; rational temporal
contraction bounds; the two-species Schur source identity; and the universal quadratic
Wilson-word identity used in the local native expansion. The report is
\path{ym_v4_verification_report.json}.
The script has 8,442 bytes and SHA--256
\begin{center}\scriptsize\ttfamily
8C22B165F3A623A89F8E5EAB73A8B1F5291FF1575B4773344815847514EE0D06.
\end{center}
These finite fixtures check the indicated algebra and support structure.
They do not evaluate the analytic derivative suprema, independently
verify the Witten domain construction, or certify a physical mass gap.
The analytic arguments above, including their cited inherited tools,
are essential. No external proof-assistant verification was performed.

The supplied V3 source has 68,826 bytes and SHA--256
\begin{center}\scriptsize\ttfamily
05F39BD90C9796D694D109855BE56A508D2F2DF2857238F773B1F3E5B55F505E.
\end{center}
V4 appends this section before its final \verb|\end{document}| and changes
only the title version from V3 to V4. Removing this append and restoring
that title recovers the supplied V3 byte for byte. All four supplied
originals are retained unchanged. The older scripts and companion files
mentioned in the inherited notes were not supplied; their historical
checks are not represented here as newly rerun. Only the three supplied
background sources were available for the present companion audit.

\begin{firewall}
V4 proves duration- and volume-uniform nonlinear two-source memory for
a specified native fast-history \emph{core}, a smaller
\(O(R^2/\gamma)\) distant common-bridge return, and a paid core pressure
density relative to the full Gaussian fast integral. The exact
compact normalization and the remaining \(\log p_{\rm good}\)
correction are retained. No bound on that full-law correction,
gauge-invariant bridge-tail coverage, full-transfer spectral gap,
interacting continuum reconstruction, or positive physical
Yang--Mills mass gap is established by this appendix.
\end{firewall}
% END F16 V4 APPEND
% BEGIN F16 V5 APPEND -- 2026-09-17
% Insert immediately before \end{document} in the cumulative V4.
% All inherited text is preserved; only the cumulative title changes to V5.

\clearpage
\section{V5: complete compact fast pressure by rooted pins and sparse corridors}
\label{f16v5:context}

V4 leaves two different questions in its last term
\(\log p_{\rm good}\): its extensive value and its local source derivatives.
This appendix pays the first without assuming the second. The upstream
change is to undo the moving frame on the \emph{temporal integration
variables only}. The seven tree kinetic terms on a cube then become an
independent compact pin. Along with the internal magnetic pin, they remain
coercive after inter-cube and inter-time interactions are cut. This permits
an upper/lower comparison through finite boxes separated by sparse
corridors, instead of an unproved large-field cluster expansion.

The resulting estimate concerns the \emph{complete compact} fast-history
integral, at an arbitrary bridge history in a fixed rescaled bounded
window. It is uniform in the number of cubes and the history length. It
also holds for bounded one-cell quadratic fast sources. Convexity in
those auxiliary sources then gives an actual-law, spatially and temporally
averaged second-moment comparison. Neither conclusion is a local
bridge-source mixing theorem.

\subsection{The V5 checkpoint and the part that is genuinely new}

The available checkpoint inputs are the supplied f14 V100, f15 V100,
Grand Tensor Monograph V076, and this lineage through V4. Section
inventories and targeted passages were examined; this is not a claim to
have independently certified every historical proof. In particular:
\begin{center}\small
\begin{tabular}{@{}p{0.23\textwidth}p{0.71\textwidth}@{}}
\toprule
Inherited result & What is reused, and what is not inferred\\
\midrule
f15 V97 & The five-chord compact magnetic inequality and the rooted
seven-edge tree estimate are already proved. They are ingredients,
not new V5 theorems.\\[1mm]
f14 V95--V100 & Full Wilson pressure stability and thermal estimates
refer to their specified periodic measure, not to the present
fixed-bridge, dressed-end history conditional.\\[1mm]
f15 V59 & A pressure-value comparison must not be differentiated without
additional control. That limitation still applies to bridge derivatives.\\[1mm]
f15 V69--V71 & Native bad-region estimates there use different conditional
laws and an actual-vacuum exterior-energy budget. No such budget is
imported into this history law.\\[1mm]
Monograph V076 & The direct Hessian minus the returned covariance is
already available. We do not repackage that identity as a new result.\\[1mm]
This lineage V4 & The full history density, restricted-core response,
core pressure comparison and exact definition of \(p_{\rm good}\)
remain unchanged.\\
\bottomrule
\end{tabular}
\end{center}

The new conclusion is convergence of the \emph{full conditional fast
pressure} to its particular, complete Gaussian history pressure, with a
vanishing bound per cube and time slice. The proof does not assume a
uniform probability that every fast variable is small. The sparse-corridor
comparison is a proof device, not a replacement of the native transfer.
The other companion programs named in older checkpoints were not supplied
in this conversation; their historical checks have not been rerun.

For methodological context, partition-function stability via exact tree
coordinates is not new. See P.~A.~Faria da Veiga and M.~O'Carroll,
\href{https://arxiv.org/abs/1903.09829v2}{\emph{On Thermodynamic and
Ultraviolet Stability of Yang--Mills}}, and their
\href{https://arxiv.org/abs/2205.07376v2}{\emph{On Yang--Mills Stability
Bounds and Plaquette Field Generating Function}}. The former treats
\(\mathrm U(N)\) and \(\SU(N)\) with free boundaries; the latter's stated
setting is \(\mathrm U(N)\). We use neither as a black-box theorem for the
present \(\SU(2)\) conditional. The inequalities needed below are proved
for its actual words and normalization.

\subsection{Undoing temporal frames, without changing the compact integral}

Retain \(B=m^3\), \(m\ge4\), \(n\ge1\), and write
\[
 \mathcal N=B(n+1),\qquad
 \max_{i,e}|y_{i,e}|\le r,
 \qquad Z_i=\operatorname{Exp}(y_i/\sqrt\gamma),
\]
where \(r\ge0\) is fixed independently of \(B,n,\gamma\). Constants denoted
\(C_r\) may depend on this window and the fixed cube geometry, but not on
\(B,n\). All principal logarithms below are taken away from their
Haar-null cut sets.

Let \(f_{i,v}=f_v(X_i)\) be the moving port frame from V2--V4, and set
\begin{equation}
 h_{t,v}=f_{t,v}r_{t,v}f_{t+1,v}^{-1},\qquad
 h_{t,o}=I,\qquad
 Y_{i,e}=f_{i,s(e)}Z_{i,e}f_{i,t(e)}^{-1}.
 \label{f16v5:unmove}
\end{equation}
For fixed internal configurations, \(r\mapsto h\) is a product of Haar
translations. The internal kinetic word, up to conjugation, becomes
\[
 h_{t,s(e)}^{-1}X_{t,e}h_{t,t(e)}X_{t+1,e}^{-1};
\]
its bridge counterpart is
\[
 h_{t,s(e)}^{-1}Y_{t,e}h_{t,t(e)}Y_{t+1,e}^{-1}.
\]
For a tree edge, \(X_{t,e}=X_{t+1,e}=I\), so its word is exactly
\(h_{t,s(e)}^{-1}h_{t,t(e)}\). In particular it contains no internal chord
or bridge variable.

Use complete charts
\[
 X_i=\operatorname{Exp}(x_i/\sqrt\gamma),\qquad
 h_t=\operatorname{Exp}(a_t/\sqrt\gamma),\qquad \zeta=(x,a).
\]
A cell \(c=(i,b)\) contains \(x_{i,b}\in\mathbb R^{15}\) and also
\(a_{i,b}\in\mathbb R^{21}\) when \(i<n\). Thus its dimension \(d_c\)
is 36, except that the final-time cells have dimension 15. Set
\(d_{\rm f}=\sum_c d_c=15B(n+1)+21Bn\). The domain
\(\mathcal P'\) is the product of the principal three-dimensional group
charts, not a small-field restriction.

Let \(U_\gamma(\zeta;\mathbf y)\) be V4's classical exponent after
\eqref{f16v5:unmove}: it includes the two endpoint Gaussian quadratic
forms and every native Wilson term, but not the negative logarithms of
the Haar densities. Define
\begin{equation}
 J^f_\gamma(\zeta)=
  \prod_{i,b,j\ {\rm chord}}j(x_{i,b,j}/\sqrt\gamma)^{w_i}
  \prod_{t,b,v\ {\rm nonroot}}j(a_{t,b,v}/\sqrt\gamma),
 \qquad 0\le J^f_\gamma\le1.
 \label{f16v5:haar}
\end{equation}
Here the endpoint weights \(w_i=1/2\), interior weights \(w_i=1\), and
\(j\) have exactly their meanings in V4.

\begin{proposition}[Same full fast integral; different useful coordinates]
\label{f16v5:exact-unmoving}
At every finite regulator,
\begin{equation}
 Z_{\gamma,\mathrm{full}}(\mathbf y)
   =\int_{\mathcal P'}J^f_\gamma(\zeta)
                         e^{-U_\gamma(\zeta;\mathbf y)}\,d\zeta.
 \label{f16v5:full-integral}
\end{equation}
No history-dependent factor multiplies this identity. If \(U_0\) is its
complete quadratic limit, then
\begin{equation}
 Z_0(\mathbf y)=\int_{\mathbb R^{d_{\rm f}}}
                         e^{-U_0(\zeta;\mathbf y)}\,d\zeta
 \label{f16v5:Z0}
\end{equation}
is exactly V4's \(Z_0\), and hence the Gaussian integral in V3's history
formula, with its determinant intact.
\end{proposition}

\begin{proof}
Substitution of \(\bar X_{i,e}=f_{i,s}^{-1}X_{i,e}f_{i,t}\) in V4's
internal word gives its conjugate by \(f_{t+1,s}\) in the displayed
unmoving coordinates. The bridge calculation is identical. Class
functions remove these conjugations. Conditional Haar invariance gives
\(dH(r)=dH(h)\) before any chart is introduced. In complete charts this
replaces \(j(z/\sqrt\gamma)\,dz\) by
\(j(a/\sqrt\gamma)\,da\), with the same constant Haar powers on both
sides. The internal half-density weights and endpoint amplitudes have
not been changed. This proves \eqref{f16v5:full-integral}, including its
normalization. The dependence of \(h\) on the internal variables causes
no extra Jacobian: the conditional Haar changes are performed at each
fixed \(X\).

At quadratic order the coordinate change is
\begin{equation}
             a_t=z_t+\mathsf F_B(x_t-x_{t+1}).
 \label{f16v5:linear-unmoving}
\end{equation}
On the stacked fast space it is block triangular with identity diagonal,
so its determinant is one, also with all three colours. Substitute
\(z_t=a_t-\mathsf F_B(x_t-x_{t+1})\) in V4's \(V_0\).
This gives precisely \(U_0\), and the same Euclidean integral. This is a
change of variables in the full Gaussian integral, not a new choice of
its scalar prefactor.
\end{proof}

V4's event \(\mathcal O_R\) is \emph{not} a product of balls in these new
coordinates. We will not replace it by one. The new coordinates are used
to estimate the full integral; the already defined core is compared to it
only through their common \(Z_0\).

\subsection{Compact pins that survive every corridor cut}

Call the following part of the classical exponent its \emph{anchor}:
\begin{align}
 A_\gamma(\zeta)
 ={}&\frac14\left(x_0^{\mathsf T}C_B^{-1}x_0+
                            x_n^{\mathsf T}C_B^{-1}x_n\right)
      +\gamma\sum_{i,b}w_i S_\Box(X_{i,b})\notag\\
 &+\gamma\sum_{t,b}\sum_{e\in T_\Box}
                      s(h_{t,s(e)}^{-1}h_{t,t(e)}).
 \label{f16v5:anchor}
\end{align}
It is a sum of one-cell functions. Every remaining Wilson term is
nonnegative. Its fast carrier consists of cells in a box of side at most
two: the internal chord kinetic terms meet two adjacent times in one
cube; a mixed two-bridge plaquette meets neighbouring cubes; and a bridge
kinetic word meets at most two neighbouring cubes at two adjacent times.
Terms with no fast variable are retained as a separate scalar
\(b_\gamma(\mathbf y)\). There are at most \(C\mathcal N\) terms in all,
and at most a fixed number incident to any cell. The constants refer to
the literal word list, including its moving frames in \(Y\).

\begin{proposition}[A global independent pin for the augmented history]
\label{f16v5:global-pin}
On every complete chart,
\begin{equation}
 A_\gamma(\zeta)\ge a_*\|\zeta\|^2,
       \qquad a_*:=\frac1{6\pi^2},\qquad \gamma>0.
 \label{f16v5:pin}
\end{equation}
The same inequality holds for its quadratic limit \(A_0\). Deleting any
subset of the other positive terms, or setting any subset of cells to
zero in those other terms while retaining \emph{all} anchors, preserves
this pin. The zero of the anchor is the identity in every fast group.
\end{proposition}

\begin{proof}
The already proved f15 V97 cube inequality is
\(\sum_{j=1}^5s(X_j)\le3S_\Box(X)\).
For a principal angle \(\theta\in[0,\pi]\),
\(1-\cos\theta\ge2\theta^2/\pi^2\). Since \(w_i\ge1/2\),
\[
 \gamma\sum_{i,b}w_i S_\Box(X_{i,b})
                         \ge\frac1{3\pi^2}\sum_{i,b}|x_{i,b}|^2.
\]
The seven rooted tree paths have total depth 12. By the bi-invariant
quaternion chord triangle inequality and Cauchy--Schwarz along a path,
\[
 s(h_v)\le d_v\sum_{e\ {\rm on}\ o\to v}
                      s(h_{s(e)}^{-1}h_{t(e)}),\qquad
 \sum_v s(h_v)\le12\sum_{e\in T_\Box}
                      s(h_{s(e)}^{-1}h_{t(e)}).
\]
Consequently the tree terms in \eqref{f16v5:anchor} are at least
\(\sum_{t,b}|a_{t,b}|^2/(6\pi^2)\). Add the nonnegative endpoint
quadratics. This proves the pin and the zero assertion. Taking the
quadratic limit proves the Gaussian inequality on every finite
Euclidean vector. All the allowed changes leave these one-cell anchors
in place and leave the other terms nonnegative.
\end{proof}

This use of the inherited tree bound is the essential upstream step.
In the moving \(z\)-variables the tree kinetic terms appear to depend on
both internal time ends. In the full compact integral that apparent
coupling can be removed before estimating the normalization. No global
convexity of the nonlinear Wilson action is asserted.

We shall also allow an auxiliary one-cell quadratic source
\begin{equation}
 \mathcal Q_{\mathsf Q}(\zeta)=\sum_c\zeta_c^{\mathsf T}Q_c\zeta_c,
 \quad Q_c=Q_c^{\mathsf T},\quad\|Q_c\|\le1,
 \quad |t|\le t_*:=a_*/4.
 \label{f16v5:quadratic-source}
\end{equation}
It is included as \(-t\mathcal Q_{\mathsf Q}\) in the exponent. The
anchor then still exceeds \((a_*-|t|)\|\zeta\|^2\). Such a source is
kept in the one-cell anchor during every cut. It is an auxiliary moment
generating parameter; the native history law itself is always \(t=0\).
The source changes neither the integration measure nor the fast event.

\subsection{A complete finite-box comparison, with the dimension paid}

The next estimate applies to a block of \(M\ge1\) cells carrying their
full anchors, any retained subset of the native non-anchor terms, and
terms obtained by setting exterior fast cells to zero. Allow also
\eqref{f16v5:quadratic-source}. The domain of every integrated group is
its \emph{complete} principal chart. The block may include the first or
last history slice. Its number of real variables is at most \(36M\).
Fast-free terms can be split off as scalar factors; their number is
bounded by a fixed multiple of the cells or boundary cells to which they
are assigned.

Write its integral as \(\mathfrak Z_\gamma\), with Haar product as in
\eqref{f16v5:haar}, and let \(\mathfrak Z_0\) be the entire Euclidean
integral of its quadratic limit. All constants below are uniform in the
choice of retained terms, the source matrices, and \(|t|\le t_*\).

\begin{theorem}[Unrestricted finite-box Laplace comparison]
\label{f16v5:box-comparison}
There are constants \(C_r,K_r\ge1\) and \(\gamma_r\) such that, whenever
\(\gamma\ge\gamma_r\) and
\begin{equation}
                   K_r(M+\log\gamma)\le\gamma,
 \label{f16v5:box-size}
\end{equation}
one has
\begin{equation}
 \left|\log\frac{\mathfrak Z_\gamma}{\mathfrak Z_0}\right|
 \le \frac{C_r}{\sqrt\gamma}(M+\log\gamma)^{3/2}
                                      +C_r\gamma^{-4}.
 \label{f16v5:box-error}
\end{equation}
The integrated nonlinear variables are not conditioned to a small box.
The dimensional cost on the right is explicit.
\end{theorem}

\begin{proof}
Both classical exponents obey
\(U_\gamma,U_0\ge(a_*/2)|\zeta|^2\), with the quadratic source
included. We use this weaker pin for convenience. At \(|\zeta_c|\le1\) in each cell, the Gaussian exponent is
at most \(C_rM\). Indeed every word is from a fixed finite list, there
are \(O(M)\) terms, their local inputs are bounded, and their incidence
is uniformly bounded. The product of these unit balls has volume at
least \(e^{-CM}\). Thus
\begin{equation}
                  \mathfrak Z_0\ge e^{-C_rM}.
 \label{f16v5:box-lower}
\end{equation}
This product-ball lower bound is important; the volume of a unit ball
in the \emph{total} dimension would introduce an unnecessary
\(M\log M\) cost.

For a Wilson word \(f\), Taylor's integral formula for
\(\gamma f(w/\sqrt\gamma)\) gives a cubic error bounded by
\(C\gamma^{-1/2}|w|^3\). The constant can be taken uniform along the
whole real line segment of the input: differentiated exponentials have
bounds given by their fixed linear coefficients, and the remaining
matrix factors are unitary. Each factor here is an exponential of a
fixed linear form in \(x,a,\mathbf y\); no derivative of a newly chosen
logarithm of a product is used. The quadratic term is the squared
linearized word with its correct coefficient. Bounded local dimension
and bounded incidence give the global-on-the-block estimate
\begin{equation}
 |U_\gamma-U_0|
     \le\frac C{\sqrt\gamma}
               \left(\sum_c|\zeta_c|^3+Mr^3\right)
     \le\frac C{\sqrt\gamma}(|\zeta|^3+Mr^3).
 \label{f16v5:cubic-sum}
\end{equation}
Exactly quadratic endpoint and source terms cancel from this difference.
Frozen zero inputs and omitted terms do not worsen it. On
\(|\zeta|\le T\), \(T^2\le c\gamma\), the full Haar product obeys
\[
        0\le-\log J^f_\gamma(\zeta)\le C T^2/\gamma,
\]
by \(-\log j(v)\le C|v|^2\) in a fixed small chart.

Choose
\(T^2=K'_r(M+\log\gamma)\), with \(K'_r\) large enough. From the
pin, the estimate \(J^f_\gamma\le1\), and an elementary Gaussian
Chernoff bound, each full tail divided by \(\mathfrak Z_0\) is at most
\begin{align*}
 e^{C_rM}\int_{|\zeta|>T}e^{-a_*|\zeta|^2/2}\,d\zeta
 &\le\exp\{C'_rM-a_*T^2/4\}\\
 &\le e^{-M}\gamma^{-4}.
\end{align*}
This bounds the nonlinear tail on its principal domain and the Gaussian
tail on its entire Euclidean domain. Constants such as
\((4\pi/a_*)^{d/2}\), \(d\le36M\), were included in \(C'_rM\).
Enlarge \(K_r\) in \eqref{f16v5:box-size} so that the ball of radius
\(T\) lies inside every principal chart and inside the smaller chart
used for the Haar logarithm estimate.

On that total-radius ball the two integrands have logarithmic ratio at
most
\[
 \epsilon_M=\frac C{\sqrt\gamma}(T^3+Mr^3)+CT^2/\gamma
          \le\frac{C_r}{\sqrt\gamma}(M+\log\gamma)^{3/2}
\]
in absolute value. Writing \(\tau=e^{-M}\gamma^{-4}\), we obtain
\[
 e^{-\epsilon_M}(1-\tau)
       \le\frac{\mathfrak Z_\gamma}{\mathfrak Z_0}
       \le e^{\epsilon_M}+\tau.
\]
For \(\gamma\ge2\), taking logarithms costs at most
\(\epsilon_M+2\tau\). No assumption that \(\epsilon_M\) is small is
needed. This proves \eqref{f16v5:box-error}. The ball was used only to
compare integrands; the complement of the native integral was estimated,
not removed from its definition.
\end{proof}

Fast-free scalar terms satisfy the corresponding elementary bound
\(C_r\gamma^{-1/2}\) per term by the same Taylor argument. All finite-box
comparisons below retain these scalars.

\subsection{Uniform Gaussian boundary costs, not native mixing assumptions}

We need a Gaussian comparison only \emph{between two auxiliary cuts},
not between a cut spectrum and the full native spectrum. Let a corridor
be any set \(\mathcal C\) of cells. Define two modifications of the
classical exponent, always retaining every anchor and auxiliary source:
\[
 U_\gamma^{\mathrm N}=A_\gamma-t\mathcal Q_{\mathsf Q}
       +b_\gamma+\sum_{e:\,\operatorname{car}(e)\cap\mathcal C=\varnothing}
                                      V_{\gamma,e},
\]
and
\[
 U_\gamma^{\mathrm D}(\zeta)=A_\gamma(\zeta)
       -t\mathcal Q_{\mathsf Q}(\zeta)+b_\gamma
       +\sum_e V_{\gamma,e}(\zeta_{\mathcal C^c},0_{\mathcal C};\mathbf y).
\]
The symbols \(\mathrm N,\mathrm D\) denote, respectively, deletion of
terms touching the corridor and freezing its inputs \emph{in those
terms}. They are not claims about a physical boundary condition.
The corridor anchors are still integrated in the second expression.
Denote the corresponding Gaussian integrals by
\(Z_0^{\mathrm N}(t,\mathsf Q)\), \(Z_0^{\mathrm D}(t,\mathsf Q)\).

\begin{proposition}[Gaussian cost proportional to the changed cells]
\label{f16v5:gaussian-cut}
For \(|t|\le t_*\), \(\|Q_c\|\le1\), and \(\max|y_{i,e}|\le r\),
\begin{equation}
 \left|\log\frac{Z_0^{\mathrm N}(t,\mathsf Q)}{Z_0(t,\mathsf Q)}\right|
 +\left|\log\frac{Z_0^{\mathrm D}(t,\mathsf Q)}{Z_0(t,\mathsf Q)}\right|
                              \le C_r|\mathcal C|.
 \label{f16v5:gaussian-cut-bound}
\end{equation}
The same constant works for all corridor shapes, spatial volumes and
history lengths.
\end{proposition}

\begin{proof}
Every relevant interpolating Gaussian exponent retains the independent
anchor minus the source. All other terms are positive multiples of
squares of fixed finite-range linear forms plus bridge sources. The
interpolation for freezing is the convex combination of an original
square and the square with its corridor inputs zero. There are constants
\(h_-,h_+>0\), independent of all regulators and of the interpolation,
for which the full fast Hessian satisfies
\[
                          h_-I\le H\le h_+I.
\]
The lower bound follows from the anchor; the upper bound is the finite
word-incidence bound. Its fast linear source has norm at most \(C r\)
in each cell. All Hessians have range at most one in the cell graph whose
edges allow a unit change in each space--time coordinate.

Here is a local mean estimate that avoids a volume factor. With
\(h_+\) enlarged if necessary,
\[
 H^{-1}=h_+^{-1}\sum_{k\ge0}(I-H/h_+)^k,\qquad
 \|I-H/h_+\|\le q:=1-h_-/h_+<1.
\]
An inverse block between cells at graph distance \(d\) has norm at most
\(h_-^{-1}q^d\). A radius-\(d\) ball in the three-dimensional spatial
torus times the finite time path has at most \(C(d+1)^4\) cells,
uniformly in its side lengths; deleting cells or edges cannot create
shorter paths than in this ambient graph. Hence every inverse block row
has summable norm, bounded by
\(C h_-^{-1}\sum_{d\ge0}(d+1)^4q^d\). Multiplying by the local source
bound proves \(|\mu_c|\le C_r\) for every conditional Gaussian mean
cell. Every marginal covariance is at most \(h_-^{-1}I\).
It follows that the expectation of any one of the squared local words
is at most \(C_r\), including words with frozen zero inputs.

There are at most \(C|\mathcal C|\) changed terms. Differentiate the
logarithm of the Gaussian integral along the deletion interpolation:
its derivative is minus the expectation of the changed quadratic
potential. Integrating its absolute value costs at most
\(C_r|\mathcal C|\). Along the freezing interpolation, the changed
term is the difference of two nonnegative squares. Its absolute
expectation is bounded by the sum of their expectations, with the same
bound. Fast-free constants obey \(C_r\) per changed term as well.
All interpolated Gaussian integrals are finite and differentiable by the
uniform pin. This proves both estimates.
\end{proof}

These estimates use only the gapped \emph{fast Gaussian pin}. They do
not require a gap in the retained bridge quadratic form or a Witten
estimate for the full compact conditional.

\subsection{Sparse corridors and the complete native pressure}

Let \(\ell\ge8\) be an integer. On each spatial cycle and on the finite
time path, use no cut when its length is at most \(\ell\). Otherwise
partition it into consecutive intervals of length at most \(\ell\) and
at least \(\ell/2-1\), and mark the first two positions of each interval.
Let \(\mathcal C\) be the union of the resulting coordinate corridors.
Then
\begin{equation}
 |\mathcal C|\le\frac{16\mathcal N}{\ell},\qquad
     |\mathcal D_j|\le\ell^4
 \label{f16v5:corridors}
\end{equation}
for every connected interior component \(\mathcal D_j\). The first
constant is a convenient safe bound. Marking an end of the physical
time path is allowed and only strengthens the separation. Short
periodic directions remain intact and do not increase component size.
The width two ensures that a native local term cannot meet two different
interior components. Terms do not have to be pairwise interactions.

To see the count explicitly, a coordinate of length \(L>\ell\) can be
partitioned into \(k=\lceil L/\ell\rceil\) intervals of lengths
\(\lfloor L/k\rfloor\) or \(\lceil L/k\rceil\). Its marked fraction
is at most \(2k/L\le4/\ell\). Sum over the four coordinates. The
complement consists of products of the remaining coordinate intervals;
no range-one carrier crosses a width-two marked interval.

Denote the full and modified integrals, with the auxiliary source, by
\[
 Z_\gamma(t,\mathsf Q),\qquad
 Z_\gamma^{\mathrm N}(t,\mathsf Q),\qquad
 Z_\gamma^{\mathrm D}(t,\mathsf Q).
\]
Both modified integrals factor into complete compact integrals on
\(\mathcal D_j\), independent one-cell anchor integrals on
\(\mathcal C\), and their literal fast-free scalar factors.

\begin{theorem}[Full compact fast-pressure comparison, uniform in duration]
\label{f16v5:full-pressure}
Fix \(r<\infty\). There are \(C_r,K_r,\gamma_r\) such that, for all
\(B=m^3\), \(m\ge4\), \(n\ge1\), \(\gamma\ge\gamma_r\),
\(\max|y_{i,e}|\le r\), \(|t|\le t_*\), and symmetric
\(\|Q_c\|\le1\), the following holds. For every \(\ell\ge8\) with
\(K_r(\ell^4+\log\gamma)\le\gamma\),
\begin{equation}
 \frac1{\mathcal N}
    \left|\log\frac{Z_\gamma(t,\mathsf Q;\mathbf y)}
                       {Z_0(t,\mathsf Q;\mathbf y)}\right|
 \le C_r\left\{
       \frac{1+\log\gamma}{\ell}
       +\frac{\ell^2+(\log\gamma)^{3/2}}{\sqrt\gamma}
          \right\}.
 \label{f16v5:pressure-scale}
\end{equation}
In particular, for sufficiently large \(\gamma\),
\begin{equation}
 \frac1{\mathcal N}
    \left|\log\frac{Z_\gamma(t,\mathsf Q;\mathbf y)}
                       {Z_0(t,\mathsf Q;\mathbf y)}\right|
 \le C_r\varepsilon_\gamma,
 \qquad
 \varepsilon_\gamma=
       \gamma^{-1/6}(1+\log\gamma)^{2/3}.
 \label{f16v5:pressure-rate}
\end{equation}
At \(t=0\), \(Z_\gamma\) is exactly V4's
\(Z_{\gamma,\mathrm{full}}\), not its restricted core.
\end{theorem}

\begin{proof}
The upper native bound is immediate from positivity of every deleted
term:
\begin{equation}
                  Z_\gamma(t,\mathsf Q)
                            \le Z_\gamma^{\mathrm N}(t,\mathsf Q).
 \label{f16v5:native-upper}
\end{equation}
The source is retained in the anchors, so it does not affect this
ordering.

For the lower bound restrict only the corridor cells to
\(|\zeta_c|\le\eta\), \(\eta=\gamma^{-1/2}\). This small event is a
temporary lower integration bound, not the definition of the full law.
For a native Wilson term, differentiation with respect to a rescaled
fast coordinate has norm at most \(C\sqrt\gamma\) throughout the
complete principal charts. This follows by differentiating the product
of unitary exponentials; the fixed moving-frame matrices have bounded
linear coefficients. There is no need to bound any interior cell.
Changing all of a term's corridor inputs from zero to size at most
\(\eta\) therefore changes its action by at most \(C\). Only
\(C|\mathcal C|\) terms are affected. The one-cell anchor quadratics
and auxiliary source cost at most \(C\eta^2\) per corridor cell, and
its Haar density is bounded below by \(e^{-C\eta^2/\gamma}\).
Thus integration over these corridor balls gives
\begin{equation}
 Z_\gamma(t,\mathsf Q)
       \ge e^{-C(1+\log\gamma)|\mathcal C|}
                         Z_\gamma^{\mathrm D}(t,\mathsf Q).
 \label{f16v5:native-lower}
\end{equation}
Here are the normalization details. The ball in a corridor cell has
volume \(v_{d_c}\eta^{d_c}\), with \(d_c\le36\), whose negative
logarithm is at most \(C+18\log\gamma\). At frozen corridor inputs,
the remaining integral is the product of the interior Dirichlet
integrals and the fast-free factors. To turn it into
\(Z_\gamma^{\mathrm D}\), multiply by the complete independent
corridor anchor integrals. Each of these has logarithm bounded in
absolute value by a constant: the global pin gives its upper bound,
and a fixed unit ball with its positive small-chart Haar density gives
its lower bound. Its cell dimension is fixed. Dividing by these anchor
integrals only changes the constant in \eqref{f16v5:native-lower}.
All constants are uniform in \(t,\mathsf Q\). No coupling-dependent
factor proportional to \(\gamma|\mathcal C|\) is hidden here.

Apply Theorem~\ref{f16v5:box-comparison} to every factor of either
modified integral. The block sizes \(M_j\) sum to \(\mathcal N\),
with each at most \(\ell^4\); corridor anchors count as blocks of size
one. The number of nonempty blocks is at most \(\mathcal N\). Hence
\begin{align*}
 \sum_j(M_j+\log\gamma)^{3/2}
 &\le C\sum_j\{M_j^{3/2}+(\log\gamma)^{3/2}\}\\
 &\le C\mathcal N\{\ell^2+(\log\gamma)^{3/2}\}.
\end{align*}
Fast-free factors contribute at most \(C_r\mathcal N/\sqrt\gamma\)
to the comparison, also after freezing. The \(\gamma^{-4}\) errors
are absorbed by the same displayed budget. Consequently, for either
modification,
\begin{equation}
 |\log Z_\gamma^{\mathrm N,\mathrm D}
                   -\log Z_0^{\mathrm N,\mathrm D}|
 \le \frac{C_r\mathcal N}{\sqrt\gamma}
                         \{\ell^2+(\log\gamma)^{3/2}\}.
 \label{f16v5:sum-box-error}
\end{equation}
This compares \emph{complete} box integrals, not Gaussian integrals
with an unestimated nonlinear complement.

Combine \eqref{f16v5:native-upper}, \eqref{f16v5:native-lower},
\eqref{f16v5:sum-box-error} and
Proposition~\ref{f16v5:gaussian-cut}. Finally use
\eqref{f16v5:corridors} and divide by \(\mathcal N\).
This proves \eqref{f16v5:pressure-scale}. The argument also works when
there are no corridors because all four directions are shorter than
\(\ell\): there is then a single block of size at most \(\ell^4\).

Choose
\(\ell=\lceil\gamma^{1/6}(1+\log\gamma)^{1/3}\rceil\).
For large \(\gamma\), the finite-box condition holds since
\(\ell^4=O(\gamma^{2/3}(1+\log\gamma)^{4/3})=o(\gamma)\).
The first two dominant errors both have size
\(O(\gamma^{-1/6}(1+\log\gamma)^{2/3})\); the remaining logarithmic
term divided by \(\sqrt\gamma\) is smaller. This proves the rate.
The exponent is a conservative balance of the two proved costs, not a
claim of an optimal asymptotic expansion.
\end{proof}

The proof has not assumed translation invariance of the bridge history,
reflection positivity of its conditional law, convexity of the compact
potential, or an a priori small probability of bad cells. The corridors
are removed from the statement by the two-sided pressure comparison.
Their widths and the shrinking integration radius do not alter the
physical spatial or temporal regulator.

\subsection{The previously unpaid all-good pressure is now bounded}

Return to \(t=0\). Keep \(p_{\rm good}\) and \(Z_{\gamma,R}\) exactly
as defined in V4, in its moving, whitened fast variables. The following
uses the equality of the \emph{full} Gaussian integrals in
Proposition~\ref{f16v5:exact-unmoving}; it does not identify two different
small-field events.

\begin{theorem}[A vanishing full-law cost per cube and time slice]
\label{f16v5:good-pressure}
Under the hypotheses of V4's core pressure theorem and, in addition,
\(\gamma\ge\gamma_r\),
\begin{equation}
 0\le-\frac{\log p_{\rm good}(\mathbf y)}{\mathcal N}
 \le C_r\varepsilon_\gamma
       +C_1\frac{R^3}{\sqrt\gamma}
       +4\,2^{21/2}e^{-R^2/16}.
 \label{f16v5:good-pressure-bound}
\end{equation}
In particular, with \(R=4\sqrt{\log\gamma}\) and fixed \(r\), all the
required core conditions hold eventually and
\begin{equation}
 e^{-C_r\mathcal N\varepsilon_\gamma}
          \le p_{\rm good}(\mathbf y)\le1.
 \label{f16v5:good-exp}
\end{equation}
This is a full-native conditional probability, with the original
endpoint amplitudes and no changed fast normalizer.
\end{theorem}

\begin{proof}
By the exact definition,
\[
 -\log p_{\rm good}
     =\log\frac{Z_{\gamma,\mathrm{full}}}{Z_0}
                      -\log\frac{Z_{\gamma,R}}{Z_0}.
\]
Use \eqref{f16v5:pressure-rate} at \(t=0\) and
Theorem~\ref{f16v4:pressure}. Nonnegativity is exact because the core
is a restriction of the full integral. For the specified \(R\), the
additional errors are \(O((\log\gamma)^{3/2}/\sqrt\gamma)\) and
\(O(\gamma^{-1})\), both smaller than \(\varepsilon_\gamma\).
The fixed mean-margin and chart conditions in V4 hold for sufficiently
large \(\gamma\). This proves \eqref{f16v5:good-exp}.
\end{proof}

There is also an exact entropy interpretation. With
\(\nu_{\gamma,\mathrm{full},\mathbf y}\) denoting V4's augmented full
native conditional in the original moving variables,
\[
 \nu_{\gamma,R,\mathbf y}
    =\nu_{\gamma,\mathrm{full},\mathbf y}(\,\cdot\mid\mathcal O_R),
 \qquad
 D(\nu_{\gamma,R,\mathbf y}\Vert
                  \nu_{\gamma,\mathrm{full},\mathbf y})=-\log p_{\rm good}.
\]
This follows immediately from the conditional Radon--Nikodym derivative
\(1_{\mathcal O_R}/p_{\rm good}\). Thus its relative entropy per
cube and time slice vanishes at the proved rate. For each finite volume
their total-variation distance, in the convention \(\sup_A|\mu(A)-\nu(A)|\),
is exactly \(1-p_{\rm good}\). Consequently it tends to zero along
sequences with \(\mathcal N\varepsilon_\gamma\to0\). It is not
asserted to tend to zero uniformly in unrestricted growing volumes.

The scalar statement for the exact history weight retains all the
external normalizers as well:
\begin{align}
 &\log\frac{\mathcal W_{\gamma,n}^{\mathrm{full}}(\mathbf y)}
                   {\mathcal W_n(\mathbf y)}
  +\frac12\log\frac{F_{\gamma,2}(Z_0)F_{\gamma,2}(Z_n)}
                       {F_2(y_0)F_2(y_n)}\notag\\
 &\hspace{7mm}-\log\frac{c_{\gamma,B,n}}{c_{0,B,n}}
                 -\log J_{{\rm br},\gamma}(\mathbf y)
       =\log\frac{Z_{\gamma,\mathrm{full}}(\mathbf y)}{Z_0(\mathbf y)}.
 \label{f16v5:full-compensated}
\end{align}
The right side divided by \(\mathcal N\) is now bounded by
\(C_r\varepsilon_\gamma\). Equation~\eqref{f16v5:full-compensated}
is not permission to drop the exact \(F_{\gamma,2}\) endpoint factors.

\subsection{A justified consequence for actual full-law quadratic moments}

An arbitrary bridge derivative of a small pressure error is not small
by inference alone. The auxiliary sources
\eqref{f16v5:quadratic-source} provide a different, justified derivative:
the log partition function is convex in their \emph{linear coefficient}
\(t\), and the Gaussian comparison has a uniform second-derivative
bound. This yields the following averaged native statement.

Let \(\widetilde\nu_{\gamma,\mathbf y}\) be the probability law of
\(\zeta=(x,a)\) obtained by normalizing
\eqref{f16v5:full-integral}, and let \(\widetilde\nu_{0,\mathbf y}\)
be its full Gaussian reference. In a cell \(c\), define raw second-moment
matrices
\[
 M_{\gamma,c}=\mathbb E_{\widetilde\nu_\gamma}
                               \zeta_c\zeta_c^{\mathsf T},\qquad
 M_{0,c}=\mathbb E_{\widetilde\nu_0}
                               \zeta_c\zeta_c^{\mathsf T}.
\]
The tilde distinguishes these unmoving coordinates from V4's moving
whitened law. No mean is being set to zero at a nonzero bridge history.

\begin{theorem}[Full-native, averaged one-cell second moments]
\label{f16v5:second-moments}
For all sufficiently large \(\gamma\), uniformly in \(B,n\) and
\(\max|y_{i,e}|\le r\),
\begin{equation}
 \frac1{\mathcal N}\sum_c\|M_{\gamma,c}-M_{0,c}\|_1
           \le C_r\sqrt{\varepsilon_\gamma}
           =C_r\gamma^{-1/12}(1+\log\gamma)^{1/3}.
 \label{f16v5:moment-rate}
\end{equation}
Here \(\|\cdot\|_1\) is the finite-matrix trace norm. In particular
\begin{equation}
        \frac1{\mathcal N}
             \mathbb E_{\widetilde\nu_\gamma}\|\zeta\|^2\le C_r.
 \label{f16v5:actual-energy-density}
\end{equation}
This controls an average of complete one-cell moment matrices, not
every distant-time covariance block.
\end{theorem}

\begin{proof}
For fixed \(\mathsf Q=(Q_c)_c\), set
\[
 f_\gamma(t)=\mathcal N^{-1}\log Z_\gamma(t,\mathsf Q),\qquad
 f_0(t)=\mathcal N^{-1}\log Z_0(t,\mathsf Q).
\]
Theorem~\ref{f16v5:full-pressure} gives
\(\sup_{|t|\le t_*}|f_\gamma(t)-f_0(t)|
 \le\delta_\gamma:=C_r\varepsilon_\gamma\), with a constant independent
of \(\mathsf Q\). These are real, differentiable convex functions of
\(t\). Native differentiability holds at every finite regulator because
all principal fast coordinates have bounded norm there. Gaussian
integrability is uniform over this interval by the anchor.

If \(C_t,\mu_t\) are the Gaussian covariance and mean on the whole
fast space, and \(Q=\bigoplus_cQ_c\), direct Gaussian integration gives
\begin{align*}
 \operatorname{Var}_{0,t}(\mathcal Q_{\mathsf Q})
 &=2\Tr\bigl[(C_t^{1/2}QC_t^{1/2})^2\bigr]
                         +4\mu_t^{\mathsf T}QC_tQ\mu_t\\
 &\le C_r\mathcal N.
\end{align*}
For clarity, expand \(\zeta=\mu_t+g\): the centered quadratic and the
linear term in the centered Gaussian \(g\) have zero covariance by
oddness, their variances are the two displayed terms, and the constant
part does not contribute. The estimate uses
\(\|C_t\|\le C\), \(d_{\rm f}\le36\mathcal N\), and
\(\sum_c|\mu_{t,c}|^2\le C_r\mathcal N\), proved by the same uniform
Gaussian block-row argument as Proposition~\ref{f16v5:gaussian-cut}.
Thus \(0\le f_0''(t)\le C_r\) on the whole interval, including for
indefinite \(Q_c\).

For \(0<h\le t_*\), convexity traps \(f_\gamma'(0)\) between the
backward and forward secants. Comparing those secants with \(f_0\),
then using its second-derivative bound, yields
\[
 |f_\gamma'(0)-f_0'(0)|\le2\delta_\gamma/h+C_rh/2.
\]
Take \(h=\sqrt{\delta_\gamma}\) for large \(\gamma\); constants can
be absorbed into \(C_r\). Since differentiation of the partition
function gives
\[
 f_\gamma'(0)-f_0'(0)
       =\mathcal N^{-1}\sum_c\Tr\{Q_c(M_{\gamma,c}-M_{0,c})\},
\]
this expression has absolute value at most
\(C_r\sqrt{\varepsilon_\gamma}\), uniformly over all allowed \(Q_c\).
Choose in each cell the spectral sign of the symmetric matrix
\(M_{\gamma,c}-M_{0,c}\), with sign zero on its kernel. This is an
allowed norm-one symmetric matrix, and its trace pairing is the trace
norm. The uniformity in \(\mathsf Q\) permits this choice after the
moments are defined. This proves \eqref{f16v5:moment-rate}. The Gaussian
one-cell mean/covariance bounds and \(Q_c=I\) give
\eqref{f16v5:actual-energy-density}.
\end{proof}

There is a useful, weaker density consequence for V4's \emph{original}
fast labels. Let
\[
 N_R=\sum_{i,b}1_{\{|u_{i,b}|\ge R\}}
              +\sum_{t,b}1_{\{|v_{t,b}|\ge R\}},
 \qquad u=N^{1/2}x,\quad v=L_c^{1/2}z.
\]
For the full native conditional and any \(R\ge1\),
\begin{equation}
                   \frac{\mathbb E N_R}{\mathcal N}\le\frac{C_r}{R^2}.
 \label{f16v5:bad-density}
\end{equation}
Indeed, from \(r_t=f_t^{-1}h_tf_{t+1}\) and the bi-invariant geodesic
triangle inequality,
\[
 |z_{t,v}|\le|a_{t,v}|+|(\mathsf F_Bx_t)_v|
                              +|(\mathsf F_Bx_{t+1})_v|.
\]
This is a global inequality for principal angles, not a linearization
of \eqref{f16v5:linear-unmoving}. Each internal slice is used at most
twice; the cube matrices have fixed norms. Therefore
\[
 \sum_{i,b}|u_{i,b}|^2+\sum_{t,b}|v_{t,b}|^2
                                      \le C\|\zeta\|^2.
\]
Apply \(1_{\{|w|\ge R\}}\le|w|^2/R^2\) and
\eqref{f16v5:actual-energy-density}. With \(R=4\sqrt{\log\gamma}\),
the expected bad-label density is \(O_r((\log\gamma)^{-1})\).
There is no inferred exponential probability for every specified bad
set, nor a percolation or conditional-capacity conclusion.

\subsection{The remaining interface is now source-marked, not scalar}

The exact identity from V4 still reads, at nonadjacent bridge times,
\begin{align*}
 \partial_{i,h}\partial_{j,k}
       (-\log\mathcal W_{\gamma,n}^{\mathrm{full}})
  ={}&\partial_{i,h}\partial_{j,k}
       (-\log\mathcal W_{\gamma,n}^{R})\\
    &+\partial_{i,h}\partial_{j,k}\log p_{\rm good}.
\end{align*}
The value of \(-\log p_{\rm good}\) per cell is no longer unpaid.
Its \emph{local bridge mixed derivatives} are still unpaid. A small
pressure density does not permit their differentiation, and the
convex-source argument in Theorem~\ref{f16v5:second-moments} concerns
one-cell quadratic fast observables, not these nonlinear bridge writers.
Likewise, a small expected density of bad labels does not by itself
bound a source-visible connected bad region.

There is a specific constructive input for V6. All compact interior
boxes already have the full normalization estimate above; it is not
necessary to derive it again or introduce a new scalar adjustment.
A marked comparison must now preserve cancellations between the full
and core integrals near a chosen bridge-source carrier while controlling
how the sparse corridors or bad regions communicate with that carrier.
Applying the unmarked global bound to a single source would lose its
\(\mathcal N\) factor. The needed improvement is locality of the
\emph{normalizer's source response}, not another complete Gaussian
inverse or another global count.

The retained bridge window is still a coordinate window before the
remaining root Gauss restriction, not gauge-invariant coverage of the
physical bridge law. The full identity is the correct native one, but
the bound is conditional on these bridge histories. No comparison of a
full-transfer spectral gap with its top eigenvalue has been established.
Physical-time calibration and the interacting continuum construction
remain distinct requirements.

\subsection{Verification and exact source preservation}

The new diagnostic is
\begin{center}\small
\path{verify_ym_f16_v5_full_fast_pressure.py}.
\end{center}
It tests exact compact word identities, cube and tree pins, and the
determinant-one change of history coordinates. It also tests corridor
separation and counts, positive quadratic cuts, Gaussian source
derivatives, and the pressure and moment rates. All 89,339 assertions passed; most enumerate corridor support and separation
on finite grids. The separate available V4 rerun passed its 9,467
assertions. The V5 report is
\path{ym_v5_verification_report.json}. These are finite algebraic and
combinatorial diagnostics, not proof-assistant verification of the
analytic theorems. In particular they do not numerically optimize the
local derivative constants, replace the finite-box tail proof, establish
bridge-source mixing, or certify a mass gap.
The V5 checker has 14,365 bytes and SHA--256
\begin{center}\scriptsize\ttfamily
56942128B285E46AD12743BAEE41413BEA6E2B893BEEBEF4B7C74F2F9A786918.
\end{center}

The supplied V4 source has 111,673 bytes and SHA--256
\begin{center}\scriptsize\ttfamily
004C1B37DA71B3C6BF4BA25DFCDAFAF7B45B6CF744FEA8BA5A0C3A92BB5986D5.
\end{center}
This cumulative V5 changes only that source's title version and inserts
the present appendix immediately before its final
\verb|\end{document}|. Removing the appendix and restoring the title
recovers the supplied V4 byte for byte. The four original background
uploads and both V4 sources are retained unchanged. The available V4
diagnostic was extracted from the supplied package and rerun; older
scripts named in the archives were not supplied and have not been
represented as rerun. Build, diagnostic and hash results are provided
with the V5 package. This V5 checkpoint uses only the available companion
sources; the next scheduled checkpoint is V10.

\begin{firewall}
V5 proves a vanishing pressure-density error for the complete compact
fast-history conditional, with its Gaussian determinant and exact
native normalization retained. It therefore bounds the previously
unpaid value of \(-\log p_{\rm good}\) per cube and time slice.
Uniform auxiliary quadratic-source pressure also yields convergence of
averaged one-cell raw second-moment matrices under the full native law,
and a vanishing expected bad-label density at growing core radius.
No local bridge-source derivative bound for \(\log p_{\rm good}\),
full-law temporal mixing theorem, gauge-invariant bridge-tail coverage,
full-transfer spectral gap, interacting continuum construction, or
positive physical Yang--Mills mass gap is proved here.
\end{firewall}
% END F16 V5 APPEND
% BEGIN F16 V6 APPEND -- 2026-09-17
% Insert immediately before the final \end{document} in cumulative V5.
% No inherited proof is rewritten. The cumulative title alone changes to V6.

\clearpage
\section{V6: compensated likelihoods and source-marked full-law response}
\label{f16v6:context}

V5 pays the full fast pressure but correctly forbids differentiating its
error to obtain a local bridge susceptibility. We go upstream to a stronger
comparison of the \emph{normalized full laws}. A small positive power of the
native-to-Gaussian likelihood is integrable with a vanishing logarithmic
cost per cell. The negative Gaussian action in this likelihood tilt is
paid by an explicit redistribution of V5's compact anchors. It is not
silently treated as a positive Wilson interaction.

A Gaussian marginal-entropy inequality then distributes that cost over
bounded-overlap families of local source carriers. Together with global
bounds on the actual Wilson first derivatives, this yields a new
\emph{source-marked} conclusion: the mean absolute error of separated-time
bridge Hessian entries tends to zero, uniformly in spatial volume and
history length, for every fixed-overlap bank of localized source pairs.
For common-bridge sources the Gaussian entry vanishes exactly. Spatial
translation symmetry at the zero bridge history additionally gives an
individual-entry estimate uniform in spatial volume, with its explicit
history-length cost retained.

These conclusions do not assert exponentially decaying native memory.
The small error below is independent of separation, so it cannot be summed
over arbitrarily many lags. Nor does a bank average become a pointwise
bound at an arbitrary inhomogeneous bridge history.

\subsection{Audit, inherited tools, and the precise change of output}

The supplied V5 appendix was read in full. The available f14/f15 archives
and Grand Tensor Monograph V076 were inventoried and searched, with
relevant passages read rather than their entire proof collections
recertified. In particular, f15 V2 already converts controlled likelihood
moments to entropy; f15 V23 already has entropy budgets for other
conditional cores and an obstruction to inferring a functional inequality
from small entropy; f15 V37 separates baseline closeness from a new
extensive physical tilt; and f15 V59 prohibits differentiation of an
unmarked pressure error. None is claimed anew. The monograph's
first-source/Hessian identity is also inherited and charged only once.

The comparison here has a different pair of laws: the complete compact
\emph{fixed-bridge fast-history} law of V5 and its complete Euclidean
Gaussian reference, in the same unmoving fast coordinates. The new
analytic input is the compensated likelihood estimate for this pair.
The marginal-entropy step is a standard Gaussian entropy tool; a direct
proof of the exact version used below is included. Its application to the
actual nonlinear bridge sources, with their tails and centering retained,
is the marked advance. No general literature-novelty claim is made for
entropy subadditivity or the source-covariance identity.

\subsection{The same full native law and a compensated likelihood tilt}

Use V5's unmoving coordinates \(\zeta=(x,a)\), principal domain
\(\mathcal P'\), cell dimensions at most 36, and exact factors
\(J^f_\gamma,U_\gamma\). Abbreviate
\[
 \mathcal N=B(n+1),\qquad
 \epsilon_\gamma=\gamma^{-1/6}(1+\log\gamma)^{2/3},\qquad
 \omega_\gamma=\sqrt{\epsilon_\gamma},\qquad
 \max_{i,e}|y_{i,e}|\le r<\infty.
\]
All sufficiently-large-\(\gamma\) thresholds may depend on the fixed
window \(r\), but not on \(B=m^3\), \(m\ge4\), or \(n\ge1\).
Let the following be probability measures on the \emph{same} Euclidean
fast space:
\begin{equation}
 \begin{split}
  d\mu_\gamma(\zeta)
   &=Z_\gamma^{-1}\,1_{\mathcal P'}J^f_\gamma(\zeta)
                              e^{-U_\gamma(\zeta;\mathbf y)}\,d\zeta,\\
  d\mu_0(\zeta)&=Z_0^{-1}e^{-U_0(\zeta;\mathbf y)}\,d\zeta,
  \qquad Z_\gamma=Z_{\gamma,\mathrm{full}}(\mathbf y).
 \end{split}
 \label{f16v6:laws}
\end{equation}
The native law is extended by zero outside its complete charts. This does
not restrict it further. V5 proves that \(Z_\gamma\) and \(Z_0\) are
exactly V4's full integrals; in particular the Gaussian determinant is
unchanged. The external \(F_{\gamma,2}\) and bridge half-density factors
cancel in a normalized conditional on the fixed bridge history, not in
the full bridge history weight.

The Gaussian precision \(H=D_\zeta^2U_0\), covariance \(C=H^{-1}\), and
mean \(m\) obey
\begin{equation}
 h_- I\le H\le h_+I,\qquad
 \sup_c|m_c|\le C_r,
 \label{f16v6:gaussian-data}
\end{equation}
with fixed positive geometric constants \(h_-,h_+\). These are V5's
anchor, finite-incidence, and Gaussian inverse-row estimates in the
unmoving coordinates, not a new coercivity claim for the full native
law. Write \(\kappa=h_+/h_-\).

For \(\theta\ge0\), define the \emph{unnormalized} likelihood integral
\begin{equation}
 \mathcal Z_\gamma(\theta)
   =\int_{\mathcal P'}(J^f_\gamma)^{1+\theta}
       \exp\{-[(1+\theta)U_\gamma-\theta U_0]\}\,d\zeta.
 \label{f16v6:likelihood-integral}
\end{equation}
The sign of \(-\theta U_0\) is essential. V5's positive-term deletion
argument does not apply to this expression before the following
compensation.

\begin{proposition}[Positive local compensation for the likelihood integral]
\label{f16v6:compensation}
There is \(\theta_*>0\), depending only on the fixed local geometry, such
that uniformly for \(0\le\theta\le\theta_*\),
\begin{equation}
 \frac1{\mathcal N}
       |\log\mathcal Z_\gamma(\theta)-\log Z_0|
                           \le C_r\epsilon_\gamma.
 \label{f16v6:tilted-pressure}
\end{equation}
The complete native integration domain and the power
\((J^f_\gamma)^{1+\theta}\) are retained. The quadratic reference of this
integral is exactly \(U_0\), independent of \(\theta\).
\end{proposition}

\begin{proof}
Write V5's decomposition, with its fast-free scalar separated, as
\[
 U_\gamma=A_\gamma+\sum_{e\in\mathcal L}V_{\gamma,e}+b_\gamma,
 \qquad
 U_0=A_0+\sum_{e\in\mathcal L}V_{0,e}+b_0.
\]
Here \(A_\gamma\) is the sum of all one-cell anchors, every
\(V_{\gamma,e}\ge0\), and \(V_{0,e}\) is a nonnegative local square.
The index \(e\) in this proof labels a remaining interaction, not
necessarily one spatial link. Its nonempty fast carrier is \(S_e\).
There are fixed constants \(L,\Delta,a^+\) such that
\begin{equation}
 \begin{gathered}
  A_\gamma\ge a_*\|\zeta\|^2,\qquad
  a_*\|\zeta\|^2\le A_0\le a^+\|\zeta\|^2,\qquad
  a_*=(6\pi^2)^{-1},\\
  0\le V_{0,e}(\zeta;\mathbf y)
     \le Q_e(\zeta)+b_e(r),\qquad
  Q_e=L\sum_{c\in S_e}|\zeta_c|^2,\qquad b_e(r)\le Lr^2,\\
  \#\{e:c\in S_e\}\le\Delta.
 \end{gathered}
 \label{f16v6:compensation-data}
\end{equation}
Enlarge \(L\) to include the fixed number of bridge coordinates in a
local word. All these bounds are uniform over interior and endpoint
cell types. They follow respectively from V5's global pin, its
one-cell quadratic anchors, and the finite list of linearized words.

Set
\begin{align}
 A^\theta_\gamma
   &=(1+\theta)A_\gamma-\theta A_0-\theta\sum_e Q_e,\notag\\
 V^\theta_{\gamma,e}
   &=(1+\theta)V_{\gamma,e}
                     +\theta(Q_e+b_e-V_{0,e}),\notag\\
 b^\theta_\gamma
   &=(1+\theta)b_\gamma-\theta b_0-\theta\sum_e b_e.
 \label{f16v6:compensated-terms}
\end{align}
Then the following is an exact identity, including its scalar:
\begin{equation}
 (1+\theta)U_\gamma-\theta U_0
       =A^\theta_\gamma+\sum_eV^\theta_{\gamma,e}+b^\theta_\gamma.
 \label{f16v6:exact-compensation}
\end{equation}
Each new interaction is nonnegative by
\eqref{f16v6:compensation-data}. Its fast support is unchanged.
The compensated anchor is still a sum of one-cell functions. For example,
choose
\[
 \theta_*=
  \min\left(\frac12,\frac{a_*}{4(a^++L\Delta+a_*)}\right)>0.
\]
Since \(\sum_eQ_e\le L\Delta\|\zeta\|^2\), this choice gives
\begin{equation}
 A^\theta_\gamma\ge\tfrac34a_*\|\zeta\|^2,
 \qquad
 A^\theta_0=A_0-\theta\sum_e Q_e
                         \ge\tfrac34a_*\|\zeta\|^2.
 \label{f16v6:surviving-pin}
\end{equation}
The quadratic limits of the other compensated terms are
\[
 V^\theta_{0,e}=V_{0,e}+\theta(Q_e+b_e),\qquad
 b^\theta_0=b_0-\theta\sum_e b_e.
\]
They sum with \(A^\theta_0\) to \(U_0\). In particular the Gaussian
reference has not been altered to force agreement of partition functions.
The scalar \(b^\theta_\gamma\) need not be positive; it is split off,
retained exactly, and has a comparison error at most
\(C_r\mathcal N/\sqrt\gamma\).

We verify why V5's finite-box and corridor proof now applies, rather
than invoking positivity for a signed interaction. On any complete box
of \(M\) cells, after scalar factors are separated, the surviving
anchors give a fixed Gaussian lower pin in both exponents. Their
quadratic integrals have the product-unit-ball lower bound
\(e^{-C_rM}\). The termwise comparison remainders are
\[
 A^\theta_\gamma-A^\theta_0
      =(1+\theta)(A_\gamma-A_0),\qquad
 V^\theta_{\gamma,e}-V^\theta_{0,e}
      =(1+\theta)(V_{\gamma,e}-V_{0,e}).
\]
Thus V5's global cubic remainder and complete-box tail estimate have
uniform constants for \(\theta\in[0,\theta_*]\). The logarithm of the
small-chart Haar correction is merely multiplied by \(1+\theta\),
and globally \(0\le(J^f_\gamma)^{1+\theta}\le1\).
The resulting complete-box logarithmic error is still
\[
 C_r\gamma^{-1/2}(M+\log\gamma)^{3/2}+C_r\gamma^{-4}
\]
when the same constant-enlarged box-size condition holds. This includes
the entire compact complement, not just a small-field comparison.

For a corridor, deletion is now legal because all
\(V^\theta_{\gamma,e}\ge0\); all compensated anchors remain.
The lower integration bound freezes corridor inputs in these terms and
integrates those cells over radius \(\gamma^{-1/2}\), exactly as in V5.
The added and subtracted quadratics have derivatives bounded by
\(C\sqrt\gamma\) throughout the principal charts, as do the rescaled
Wilson terms. The change in a touched interaction therefore costs
\(C_r\) on that small corridor ball. The pin and a fixed unit-ball lower
bound give uniform upper and lower bounds on every independent
compensated-anchor integral. The resulting logarithmic corridor debit
is at most \(C_r(1+\log\gamma)|\mathcal C|\).

For the two Gaussian cuts, the compensated anchors retain their lower
Hessian floor and every other term is a nonnegative quadratic with
uniformly bounded finite-range Hessian. The Gaussian interpolation
argument of Proposition~\ref{f16v5:gaussian-cut} therefore costs
\(C_r|\mathcal C|\), with the same local mean bound. The extra constants
\(b_e\) are part of the literal changed terms and incur only this local
cost. No native mixing inequality has entered.

Consequently the existing corridor summation gives, for every admissible
\(\ell\),
\[
 \mathcal N^{-1}|\log\mathcal Z_\gamma(\theta)-\log Z_0|
 \le C_r\left\{\frac{1+\log\gamma}{\ell}
        +\frac{\ell^2+(\log\gamma)^{3/2}}{\sqrt\gamma}\right\}.
\]
Use V5's choice
\(\ell=\lceil\gamma^{1/6}(1+\log\gamma)^{1/3}\rceil\).
This proves \eqref{f16v6:tilted-pressure}. The work specific to this
extension is the exact compensated positive representation and the
checks above; the corridor construction itself is not a new result.
\end{proof}

\subsection{A native-to-Gaussian entropy density, in the useful direction}

For probabilities \(P\ll Q\), use
\[
 D(P\Vert Q)=\int\log\frac{dP}{dQ}\,dP,
 \qquad
 D_{1+\theta}(P\Vert Q)
   =\frac1\theta\log\int\left(\frac{dP}{dQ}\right)^{1+\theta}dQ
 \quad(\theta>0).
\]
These definitions refer to normalized probabilities.

\begin{theorem}[Complete native likelihood and relative entropy]
\label{f16v6:entropy}
For some fixed \(\theta_*>0\) and all sufficiently large \(\gamma\),
\begin{equation}
 0\le D(\mu_\gamma\Vert\mu_0)
      \le D_{1+\theta_*}(\mu_\gamma\Vert\mu_0)
      \le C_r\mathcal N\epsilon_\gamma.
 \label{f16v6:entropy-density}
\end{equation}
This is the full native conditional relative to the full Gaussian
conditional. It is neither a core entropy nor an assumption of a
uniformly small total entropy.
\end{theorem}

\begin{proof}
The density ratio is zero outside \(\mathcal P'\). Within the charts,
direct substitution of \eqref{f16v6:laws} gives the exact formula
\begin{equation}
 \int\left(\frac{d\mu_\gamma}{d\mu_0}\right)^{1+\theta}d\mu_0
       =\frac{Z_0^\theta\mathcal Z_\gamma(\theta)}{Z_\gamma^{1+\theta}}.
 \label{f16v6:renyi-identity}
\end{equation}
Its logarithm equals
\[
 [\log\mathcal Z_\gamma(\theta)-\log Z_0]
       -(1+\theta)[\log Z_\gamma-\log Z_0].
\]
Apply Proposition~\ref{f16v6:compensation} at the fixed
\(\theta=\theta_*\), and V5 at \(\theta=0\). Dividing by
\(\theta_*\) gives the last bound. Jensen under \(\mu_\gamma\),
applied to \(\exp(\theta_*\log(d\mu_\gamma/d\mu_0))\), gives the
relative-entropy bound; Jensen under \(\mu_0\) gives nonnegativity
of the likelihood logarithm. All these integrals are finite by the
compensated estimate. This is an elementary likelihood implication,
already used in f15 V2, applied to the newly controlled pair of laws.
\end{proof}

The orientation matters. At any finite \(\gamma\), the Gaussian gives
positive mass outside the principal native charts, whereas \(\mu_\gamma\)
gives that set mass zero. Hence
\(D(\mu_0\Vert\mu_\gamma)=+\infty\). A two-sided full entropy theorem
has not been asserted. V5's separate identity
\(D(\nu_{\gamma,R}\Vert\nu_{\gamma,\mathrm{full}})=-\log p_{\rm good}\)
compares different laws and has not been counted again.

\subsection{Distributing entropy over bounded-overlap marginal carriers}

For a set \(S\) of cells, let \(P_S\) be the coordinate projection onto
all real coordinates in those cells, and denote marginal probabilities
by a subscript \(S\). The diameter of \(S\) need not be bounded: a union
of two distant local source supports is allowed.

\begin{proposition}[Gaussian marginal-entropy localization; standard tool]
\label{f16v6:entropy-localization}
Let \(Q\) be any nondegenerate Gaussian with covariance \(C\), and let
\(P\ll Q\) have finite relative entropy. Suppose a finite family
\((S_\alpha)_{\alpha\in\mathcal A}\) of nonempty cell sets satisfies
\begin{equation}
        \sup_c\#\{\alpha:c\in S_\alpha\}\le q.
 \label{f16v6:overlap}
\end{equation}
Then, with \(\kappa_C=\lambda_{\max}(C)/\lambda_{\min}(C)\),
\begin{equation}
       \sum_\alpha D(P_{S_\alpha}\Vert Q_{S_\alpha})
                          \le q\kappa_C D(P\Vert Q).
 \label{f16v6:marginal-entropy}
\end{equation}
No independence among the Gaussian cell coordinates or native
coordinates is assumed.
\end{proposition}

\begin{proof}
Translate and whiten the full Gaussian. In standard Gaussian
coordinates, the standardized marginal map is
\[
 B_\alpha=C_{S_\alpha S_\alpha}^{-1/2}P_{S_\alpha}C^{1/2},
 \qquad B_\alpha B_\alpha^{\mathsf T}=I.
\]
Thus \(\Pi_\alpha=B_\alpha^{\mathsf T}B_\alpha\) is an orthogonal
projection. The overlap bound gives
\begin{align*}
 \sum_\alpha\Pi_\alpha
  &=C^{1/2}\left(\sum_\alpha P_{S_\alpha}^{\mathsf T}
              C_{S_\alpha S_\alpha}^{-1}P_{S_\alpha}\right)C^{1/2}\\
  &\le\lambda_{\min}(C)^{-1}C^{1/2}(qI)C^{1/2}
   \le q\kappa_C I.
\end{align*}
Let \(f\) be the density of the whitened \(P\) relative to standard
Gaussian measure \(g\), and put \(f_t=\mathsf P_t f\), where
\(\mathsf P_t\) is the Ornstein--Uhlenbeck semigroup with generator
\(\Delta-x\cdot\nabla\). Conditional expectation onto \(B_\alpha x\)
commutes with this semigroup. For a smooth positive density, the
marginal gradient is the conditional expectation of
\(B_\alpha\nabla f_t\). Conditional Cauchy--Schwarz therefore gives
\[
 I((f_t)_\alpha)
    \le\int\frac{|B_\alpha\nabla f_t|^2}{f_t}\,dg,
 \qquad
 \sum_\alpha I((f_t)_\alpha)\le q\kappa_C I(f_t),
\]
where \(I(f)=\int|\nabla f|^2/f\,dg\) is relative Gaussian Fisher
information. Integration by parts along the same semigroup yields
\[
 \operatorname{Ent}_g(f)=\int_0^\infty I(f_t)\,dt.
\]
Apply the identity to the full density and to every marginal, and
integrate the Fisher inequality. This proves
\eqref{f16v6:marginal-entropy}. For completeness, bounded densities suffice first. Smooth them for a
positive time, integrate the dissipation identity over \([s,T]\), and
let \(T\to\infty\): the Mehler formula gives convergence to the
constant density one, and boundedness gives entropy convergence by
dominated convergence. As \(s\downarrow0\), entropy lower
semicontinuity and semigroup Jensen imply convergence to the original
entropy. For a general finite-entropy density, use normalized
truncations \(\min(f,M)/\int\min(f,M)\,dg\). Their full entropies
converge to that of \(f\), and marginal entropy lower semicontinuity
passes the proved inequality to the limit. This also justifies the
stated dissipation use without any native Fisher-information assumption
at time zero. At each fixed regulator the full native likelihood is
in fact bounded on its compact support.
\end{proof}

This is the Gaussian relative-entropy version of the Fisher-information
projection method behind entropy Brascamp--Lieb inequalities. See
E.~A.~Carlen and D.~Cordero--Erausquin,
\href{https://arxiv.org/abs/0710.0870}{\emph{Subadditivity of the entropy
and its relation to Brascamp--Lieb type inequalities}}, especially the
Fisher-information and heat-flow argument in Section 3. The proof above
specifies the correlated Gaussian marginal constant used here. It is not
an approximate tensorization of \emph{native conditional} entropies and
not a native logarithmic Sobolev inequality.

For the laws \eqref{f16v6:laws}, write
\(d_\alpha=D((\mu_\gamma)_{S_\alpha}\Vert(\mu_0)_{S_\alpha})\).
The preceding results give
\begin{equation}
        \sum_\alpha d_\alpha\le C_r q\mathcal N\epsilon_\gamma.
 \label{f16v6:bank-entropy}
\end{equation}
If \(|S_\alpha|\le s\) for fixed \(s\), every reference marginal has
bounded dimension, uniformly bounded mean, and covariance bounds from
\eqref{f16v6:gaussian-data}, independently of the diameter of
\(S_\alpha\). This last fact is what permits distant source pairs.
The entropy budget is on marginals of the full native law, not on a
collection of separately normalized cavity laws.

\subsection{Unbounded local scores: the tail cost is included}

A total-variation comparison of small marginals would not suffice for
unbounded thermal scores. The following finite-dimensional estimate
makes the additional control explicit.

\begin{proposition}[Gaussian entropy controls affine scores and native remainders]
\label{f16v6:score-entropy}
Let \(Q\) be a Gaussian in dimension at most \(d_*\), with bounded mean
and fixed positive upper and lower covariance bounds. Let \(P\ll Q\)
have \(d=D(P\Vert Q)<\infty\). Suppose \(L_1,L_2\) are affine
functions with bounded coefficients, and \(R_1,R_2\) satisfy, for
\(j=1,2\),
\begin{equation}
 |R_j(z)|\le K(1+|z|),\qquad
 |R_j(z)|\le\frac K{\sqrt\gamma}(1+|z|^2),\qquad\gamma\ge1.
 \label{f16v6:two-remainder-bounds}
\end{equation}
Then, with constants depending only on the stated fixed bounds,
\begin{align}
 \mathbb E_P(1+|z|^2)&\le C(1+d),\notag\\
 \mathbb E_PR_j^2&\le C(d+\gamma^{-1}),\notag\\
 \big|\operatorname{Cov}_P(L_1+R_1,L_2+R_2)
                 -\operatorname{Cov}_Q(L_1,L_2)\big|
       &\le C\big(\sqrt d+d+\gamma^{-1/2}\big).
 \label{f16v6:score-entropy-bound}
\end{align}
In particular no native fourth-moment assumption is needed.
\end{proposition}

\begin{proof}
We use only the entropy variational inequality
\begin{equation}
 t\mathbb E_PF\le d+\log\mathbb E_Qe^{tF},\qquad t>0,
 \label{f16v6:entropy-variational}
\end{equation}
first for bounded truncations and then by monotone or dominated limits
when applicable. It follows from nonnegativity of relative entropy
against the normalized exponential tilt of \(Q\). Gaussian exponential
integrability of \(1+|z|^2\) at a fixed positive \(t\) proves the first
bound.

Choose \(t_0>0\) sufficiently small that the Gaussian exponential
moment below is uniformly finite. Both inequalities in
\eqref{f16v6:two-remainder-bounds} are used:
\begin{align*}
 \mathbb E_Q(e^{t_0R_j^2}-1)
  &\le t_0\mathbb E_Q[R_j^2e^{t_0R_j^2}]\\
  &\le\frac C\gamma\,
       \mathbb E_Q[(1+|z|^4)e^{C t_0(1+|z|^2)}]
   \le C\gamma^{-1}.
\end{align*}
Thus \(\log\mathbb E_Qe^{t_0R_j^2}\le C\gamma^{-1}\).
Equation~\eqref{f16v6:entropy-variational} proves the second bound,
including the contribution from arbitrarily large native coordinates.
Using only the quartic bound in \eqref{f16v6:two-remainder-bounds}
would not have justified this exponential integral.

For an affine \(L\), its centered Gaussian logarithmic moment generating
function is exactly a quadratic in the source. The entropy variational
inequality with both signs therefore gives
\[
     |\mathbb E_PL-\mathbb E_QL|\le C\sqrt d.
\]
For \(F=L_1L_2\), Gaussian exponential integrability of a quadratic
implies
\[
 \log\mathbb E_Q e^{t(F-\mathbb E_QF)}\le Ct^2
                          \quad (|t|\le t_1)
\]
for a fixed \(t_1>0\). One can verify this without a determinant formula:
use \(|F|\le C(1+|z|^2)\) and
\(e^u\le1+u+\tfrac12u^2e^{|u|}\); the linear expectation vanishes.
Optimizing \(d/|t|+C|t|\) on \((0,t_1]\) gives
\[
      |\mathbb E_PF-\mathbb E_QF|\le C(\sqrt d+d).
\]
Combining this product estimate with the two affine mean estimates,
including their product of size \(Cd\), proves
\[
  |\operatorname{Cov}_P(L_1,L_2)-\operatorname{Cov}_Q(L_1,L_2)|
                                      \le C(\sqrt d+d).
\]
Finally each cross covariance involving one \(R_j\) is bounded by
\(C\sqrt{(d+\gamma^{-1})(1+d)}\), using Cauchy--Schwarz and the
first two estimates. The covariance of the two remainders is bounded
by \(C(d+\gamma^{-1})\). Since
\[
 \sqrt{(d+\gamma^{-1})(1+d)}
                 \le C(\sqrt d+d+\gamma^{-1/2}),
\]
these bounds prove the last line of
\eqref{f16v6:score-entropy-bound}. All covariances are centered under
the indicated law; no conditional mean has been discarded.
\end{proof}

\subsection{The actual native bridge writers satisfy these bounds}

A localized bridge mark is \(\alpha=(i,h)\), where \(h\) is a real
bridge-coordinate direction at time \(i\), \(\|h\|\le1\), supported
on a fixed bounded number of bridges. It may carry all three colours.
Define, in V5's unmoving coordinates held fixed,
\[
   J_{\gamma,\alpha}(\zeta)=\partial_{i,h}U_\gamma(\zeta;\mathbf y),
   \qquad
   J_{0,\alpha}(\zeta)=\partial_{i,h}U_0(\zeta;\mathbf y).
\]
Neither \(J^f_\gamma\) nor the unmoving coordinate map depends on
\(\mathbf y\). These are therefore exactly the score insertions for
the fixed-domain fast integral, with no coordinate-derivative term
missing. A bridge-only part of a writer is a deterministic constant;
retaining or subtracting it does not affect its covariance.

\begin{proposition}[Complete-chart score certificate]
\label{f16v6:native-score}
For such a mark there is a finite fast cell carrier \(S_\alpha\), whose
cardinality is bounded by the marked bridge support and the local
geometry, on which both writers depend. The Gaussian writer is affine,
with uniformly bounded coefficients. On the complete principal charts,
\begin{equation}
 \begin{split}
  |J_{\gamma,\alpha}|+|J_{0,\alpha}|
                    &\le C_r(1+|\zeta_{S_\alpha}|),\\
  |J_{\gamma,\alpha}-J_{0,\alpha}|
                    &\le C_r\gamma^{-1/2}(1+|\zeta_{S_\alpha}|^2).
 \end{split}
 \label{f16v6:native-score-bounds}
\end{equation}
These same bounds hold for the analytic exponential-word expressions
extended to all Euclidean fast inputs. The carrier occupies at most the
times \(i-1,i,i+1\), omitting missing endpoint times. The constants are
independent of spatial volume, history length and location.
\end{proposition}

\begin{proof}
Every term involving the marked bridge variables is one of a fixed
finite list of rescaled functions
\[
    \gamma s\!\left(\prod_{j=1}^k
                       \operatorname{Exp}(T_j w/\sqrt\gamma)\right),
\]
where \(w\) contains only its fast variables and nearby bridge inputs.
The matrices \(T_j\), product order, and number of incident terms are
fixed. This is also true for the unmoving bridge words: the port frames
are exponentials of fixed linear forms in the adjoining chord inputs.
There is no principal logarithm of an intermediate product to
differentiate.

Duhamel's formula for derivatives of a matrix exponential shows that
its real skew-Hermitian linear arguments have globally bounded first,
second, and third derivatives with the corresponding fixed coefficient
norms. Indeed the derivative is an integral of products of unitary
matrices and the inserted coefficient matrices; at higher orders the
simplex volumes cancel the permutation count. Applying the product
rule to the finite word gives a uniform bound on its rescaled Hessian
and a third-derivative bound \(C/\sqrt\gamma\), throughout all real
inputs. Its gradient at zero vanishes. Integrating the Hessian from
zero bounds a first bridge derivative by \(C|w|\). Integrating the
third derivative in the Taylor expansion of that first derivative
bounds its difference from the quadratic-limit first derivative by
\(C|w|^2/\sqrt\gamma\). The endpoint Gaussian anchor is independent
of bridge variables and contributes nothing here.

Summing the fixed number of terms incident to the marked support and
using the bound \(r\) on nearby bridge inputs proves
\eqref{f16v6:native-score-bounds}. The bound on the difference by
\(C_r(1+|\zeta_{S_\alpha}|)\) also follows by adding the two first
bounds, so the remainder satisfies both hypotheses of
Proposition~\ref{f16v6:score-entropy}. Finally, a spatial mixed word
uses one bridge time, and a kinetic word uses only two adjacent bridge
times. In unmoving coordinates its frames involve the internal
coordinates at those same ends. This proves the stated fast carrier
and its duration-independent size.
\end{proof}

Let
\[
 \mathcal A_\gamma=-\log Z_{\gamma,\mathrm{full}},\qquad
 \mathcal A_0=-\log Z_0.
\]
For marks \(\alpha=(i,h),\beta=(j,k)\) with \(|i-j|\ge2\), every
native term has zero direct mixed bridge derivative. Exact
fixed-domain differentiation consequently gives
\begin{equation}
 \mathcal H_\gamma(\alpha,\beta)
       :=\partial_{i,h}\partial_{j,k}\mathcal A_\gamma
       =-\operatorname{Cov}_{\mu_\gamma}
                               (J_{\gamma,\alpha},J_{\gamma,\beta}),
 \label{f16v6:exact-marked-Hessian}
\end{equation}
and the same identity with subscript zero. This is the inherited
source/Hessian identity, not a new formal claim. Differentiation is
legitimate at each finite regulator: the principal fast domain is fixed,
the Haar weight is integrable, and the exponential-word derivatives
are bounded on its compact closure. There is an open bridge
neighbourhood of every bounded base point for large \(\gamma\).

The separated-time Hessian in
\eqref{f16v6:exact-marked-Hessian} is unchanged when \(\mathcal A_\gamma\)
is replaced by \(-\log\mathcal W_{\gamma,n}^{\mathrm{full}}\), or when
\(\mathcal A_0\) is replaced by \(-\log\mathcal W_n\). All external
logarithmic factors in V4 are sums of single-time functions or
history-independent scalars. In particular this assertion does not
require estimating or dropping either \(F_{\gamma,2}\) normalizer.

\subsection{An actual source-marked bank estimate, not a pressure derivative}

\begin{theorem}[Full-native separated-time Hessians on bounded-overlap banks]
\label{f16v6:marked-bank}
Let \((\alpha_a,\beta_a)_{a\in\mathcal I}\) be localized bridge-mark
pairs at nonadjacent times. Choose nonempty carriers
\(S_a=S_{\alpha_a}\cup S_{\beta_a}\) with \(|S_a|\le s\), and
suppose each fast cell lies in at most \(q\) such carriers, counting
multiplicity. For fixed \(s,q\), uniformly in \(B,n\), the positions
and separations of the marks, and every bounded bridge history,
\begin{equation}
 \frac1{\mathcal N}\sum_{a\in\mathcal I}
       |\mathcal H_\gamma(\alpha_a,\beta_a)
                         -\mathcal H_0(\alpha_a,\beta_a)|
                       \le C_{r,s}\,q\,\omega_\gamma.
 \label{f16v6:marked-bank-bound}
\end{equation}
The same estimate holds for the Hessians of the exact native and
Gaussian full history weights. No intermediate fast projection or
native small-field restriction occurs in the left side.
\end{theorem}

\begin{proof}
Set \(d_a=D((\mu_\gamma)_{S_a}\Vert(\mu_0)_{S_a})\). All
Gaussian marginals in this family have the uniform parameters required
in Proposition~\ref{f16v6:score-entropy}, by
\eqref{f16v6:gaussian-data} and \(|S_a|\le s\). Their diameter does
not affect these bounds. Apply that proposition on the union carrier,
with affine functions \(J_{0,\alpha_a},J_{0,\beta_a}\) and the two
native remainders supplied by
Proposition~\ref{f16v6:native-score}. The exact source identity gives
\begin{equation}
 |\mathcal H_\gamma(\alpha_a,\beta_a)
                    -\mathcal H_0(\alpha_a,\beta_a)|
                      \le C_{r,s}(\sqrt{d_a}+d_a+\gamma^{-1/2}).
 \label{f16v6:individual-entropy-bound}
\end{equation}
This inequality is also useful separately: it identifies the specific
\emph{marginal entropy} needed to control an individual pair.

Equation~\eqref{f16v6:bank-entropy} bounds
\(\sum_a d_a\) by \(C_r q\mathcal N\epsilon_\gamma\).
Nonemptiness and the overlap condition imply
\(|\mathcal I|\le q\mathcal N\). Cauchy--Schwarz therefore gives
\[
 \frac1{\mathcal N}\sum_a\sqrt{d_a}
     \le\sqrt{\frac{|\mathcal I|}{\mathcal N}
                    \frac{\sum_a d_a}{\mathcal N}}
     \le C_rq\sqrt{\epsilon_\gamma}.
\]
Sum \eqref{f16v6:individual-entropy-bound} and use
\(\epsilon_\gamma\le\omega_\gamma\) and
\(\gamma^{-1/2}\le\omega_\gamma\) for large \(\gamma\).
This proves \eqref{f16v6:marked-bank-bound}. The conclusion about
history weights follows from the exact single-time-normalizer
observation after \eqref{f16v6:exact-marked-Hessian}.

At no step was \(\log Z_\gamma-\log Z_0\) differentiated.
The pressure estimate supplied a controlled likelihood moment, which
supplied marginal entropy; the actual source insertions were then
estimated directly, including their unbounded parts.
\end{proof}

The normalization by \(\mathcal N\), rather than by
\(|\mathcal I|\), is part of the theorem. A bank with positive density
has a vanishing mean absolute error per mark. A lone inhomogeneous mark
would still inherit a possible \(\mathcal N\) loss through its
uncontrolled marginal entropy; the bank inequality does not conceal
that distinction. Repeating a fixed pair many times raises \(q\)
by the same factor and creates no pointwise improvement.

\subsection{Fixed-lag common-bridge memory and the full/core correction}

Let \(h^{b,\mu,v}\) be the common-bridge direction on the four parallel
bridges from cube \(b\) to \(b+e_\mu\), with value \(v/2\) on each,
where \(|v|=1\) is a colour vector. Thus \(\|h^{b,\mu,v}\|=1\) and
\(h^{b,\mu,v}\in\mathcal E\). Fix two such direction types, a spatial
cube displacement \(a\), and a lag \(d\) with \(2\le d\le n\).
For \(0\le i\le n-d\) and every cube \(b\), write
\[
 \alpha_{i,b}=(i,h^{b,\mu,v}),\qquad
 \beta_{i,b}=(i+d,h^{b+a,\nu,w}).
\]
The two types or colours need not agree. The displacement may depend on
the spatial torus and need not remain bounded.

\begin{theorem}[Full-law common-bridge response at any fixed lag, in density]
\label{f16v6:common-bank}
Uniformly in the above choices and in every bounded bridge history,
\begin{equation}
 \frac1{B(n+1)}\sum_{i=0}^{n-d}\sum_b
       |\partial_{\alpha_{i,b}}\partial_{\beta_{i,b}}
                (-\log\mathcal W_{\gamma,n}^{\mathrm{full}})|
                                      \le C_r\omega_\gamma.
 \label{f16v6:common-native}
\end{equation}
Moreover, for V4's \emph{original} core and its actual probability
\(p_{\rm good}\), with any admissible \(R\),
\begin{equation}
 \frac1{B(n+1)}\sum_{i=0}^{n-d}\sum_b
       |\partial_{\alpha_{i,b}}\partial_{\beta_{i,b}}
                                        \log p_{\rm good}|
 \le C_r\omega_\gamma
          +C\frac{R^2}{\gamma}
                         (9/19)^{(d-2)_+}.
 \label{f16v6:common-correction}
\end{equation}
In particular \(R=4\sqrt{\log\gamma}\) is admissible eventually and
makes the last right side at most \(C_r\omega_\gamma\).
\end{theorem}

\begin{proof}
The source carrier of a direction anchored at \((i,b)\) is contained
in the cells at times \(i-1,i,i+1\) and spatial cube
\(\ell^\infty\)-distance at most one from \(b\). This follows from
the literal mixed plaquette and kinetic words; it is a safe enlargement
of the smaller actual carrier. It has at most \(3\cdot3^3=81\)
cells. A union of the two carriers has at most 162 cells. For the
translated fixed-lag bank above, any given cell occurs in at most 81
first carriers and at most 81 second carriers, even with periodic
wrapping and time endpoints. Thus \(s=q=162\) are safe uniform choices;
large separation does not enlarge them.

The inherited Gaussian cancellation is exact. In V4's moving fast
coordinates, \(D_x\partial_{i,h}V_0=\Lambda h=0\) for a common
bridge direction. Its remaining fast source uses only the port bonds
\(i-1,i\). The Gaussian fast Hessian is block diagonal between the
internal species and the ports, and between different port bonds.
Consequently the Gaussian source covariance at times separated by at
least two is zero. This is precisely the common-bridge null response
of V1/V3 together with V4's retained-port kinetic completion, not a
new Gaussian mass argument. V5's determinant-one linear coordinate
change leaves this covariance unchanged. Hence
\[
         \mathcal H_0(\alpha_{i,b},\beta_{i,b})=0.
\]
Theorem~\ref{f16v6:marked-bank} now gives
\eqref{f16v6:common-native}.

Finally the exact full/core identity has the sign
\[
 \partial_\alpha\partial_\beta\log p_{\rm good}
   =\partial_\alpha\partial_\beta
                (-\log\mathcal W_{\gamma,n}^{\mathrm{full}})
     -\partial_\alpha\partial_\beta
                (-\log\mathcal W_{\gamma,n}^{R}).
\]
Apply the triangle inequality, the first part of this theorem, and
V4's common-bridge bound
\eqref{f16v4:common-memory-bound}. The latter has
\(\delta^2=C_*^2R^2/\gamma\), and the bank has at most
\(B(n+1)\) pairs. This proves
\eqref{f16v6:common-correction}. The stated growing radius satisfies
all V4 core and mean-margin conditions eventually at fixed \(r\).
No event in the unmoving coordinates has been substituted for that
original core.
\end{proof}

For other unit localized bridge directions the Gaussian entry need not
vanish. V3's inverse bound gives, at lag \(d\ge2\),
\[
 |\mathcal H_0(\alpha,\beta)|
       \le\frac{64}{p_*-2}\left(\frac2{p_*}\right)^d,
 \qquad p_*=3+2\sqrt2,
\]
using \(\|N^{-1/2}A^{\mathsf T}\|\le2\) and \(\|D\|\le4\).
For the corresponding fixed-overlap translated banks,
\eqref{f16v6:marked-bank-bound} thus gives a mean absolute native
bound \(C_r[(2/p_*)^d+\omega_\gamma]\). The second term is a
separation-independent error floor, not exponential mixing.

\subsection{A pointwise spatial-volume-uniform consequence at zero history}

The full native law, Gaussian reference, and original V4 core at
\(\mathbf y=0\) are invariant under spatial translations by whole
cubes. The endpoint amplitudes are identical products over cubes, and
the rooted tree and frame in each cube use the same relative port
convention. No temporal translation invariance is asserted for these
dressed-end histories.

\begin{theorem}[Individual common-source entries at zero bridge history]
\label{f16v6:homogeneous}
Fix common-bridge direction types as above, two times \(i,j\) with
\(|i-j|\ge2\), and a spatial displacement \(a\). At \(\mathbf y=0\),
for every cube \(b\),
\begin{equation}
 |\partial_{i,h^{b,\mu,v}}\partial_{j,h^{b+a,\nu,w}}
                    (-\log\mathcal W_{\gamma,n}^{\mathrm{full}})|
 \le C\left\{\sqrt{(n+1)\epsilon_\gamma}
                  +(n+1)\epsilon_\gamma+\gamma^{-1/2}\right\}.
 \label{f16v6:homogeneous-bound}
\end{equation}
The constant is independent of \(B,n,i,j,a\). In particular, if
\((n+1)\epsilon_\gamma\le1\) and \(R=4\sqrt{\log\gamma}\),
\begin{equation}
 \begin{split}
 &|\partial_{i,h^{b,\mu,v}}\partial_{j,h^{b+a,\nu,w}}
                    (-\log\mathcal W_{\gamma,n}^{\mathrm{full}})|\\
 &\quad+
 |\partial_{i,h^{b,\mu,v}}\partial_{j,h^{b+a,\nu,w}}
                                      \log p_{\rm good}|
                  \le C\sqrt{(n+1)\epsilon_\gamma}.
 \end{split}
 \label{f16v6:homogeneous-correction}
\end{equation}
Thus at each fixed history length these individual derivatives vanish
uniformly in spatial volume. The same is true along growing lengths
with \((n+1)\epsilon_\gamma\to0\).
\end{theorem}

\begin{proof}
Use the bank of the \(B\) translated pairs at these fixed two times,
with the carriers chosen as spatial translates of one another. Their
overlap is bounded by the same geometric constant. By exact spatial
translation symmetry, every marginal entropy \(d_b\) for these
carriers has the same value. Equation~\eqref{f16v6:bank-entropy}
therefore gives
\[
       d_b=B^{-1}\sum_{b'}d_{b'}
                              \le C(n+1)\epsilon_\gamma.
\]
Apply the individual estimate
\eqref{f16v6:individual-entropy-bound} and the exact common-bridge
Gaussian cancellation. This proves
\eqref{f16v6:homogeneous-bound}, with no spatial-volume factor.
It does not divide by the number of time slices: the dressed endpoints
break that symmetry.

The full/core identity and V4's common-source estimate add at most
\(C(R^2/\gamma)(9/19)^{(|i-j|-2)_+}\) to bound the second term in
\eqref{f16v6:homogeneous-correction}. If
\((n+1)\epsilon_\gamma\le1\), the linear entropy term in
\eqref{f16v6:homogeneous-bound} is bounded by its square root. At the
specified \(R\), both \(\gamma^{-1/2}\) and \(R^2/\gamma\) are
bounded by \(C\sqrt{(n+1)\epsilon_\gamma}\) for large \(\gamma\).
This proves the last assertion. Spatial symmetry is used only at the
stated zero bridge history; no symmetry of an arbitrary bounded
history has been assumed.
\end{proof}

\subsection{What remains, and the next noncircular target}

V6 has now estimated the actual mixed derivatives of
\(\log p_{\rm good}\) in two specified senses: a duration-uniform
\emph{density} bound on fixed-lag common-source banks at arbitrary
bounded histories, and an individual-entry bound uniform in spatial
volume at zero history, with its length dependence explicit. This is
stronger than the V5 scalar pressure and one-cell moment statements.
It is not the uniform pointwise, time-summable response originally
needed for the complete spectral program.

There are two precise places where the present proof loses that
stronger conclusion. For an individual pair at an inhomogeneous base
history, \eqref{f16v6:individual-entropy-bound} requires its own marginal
entropy \(d_a\), whereas the proved input is only the bank sum
\eqref{f16v6:bank-entropy}. Second, even a uniform small bound on every
such \(d_a\) would give a separation-independent covariance error via
the present argument. It would not control the sum over all remote
lags. Applying the fixed-lag bank theorem repeatedly costs the number
of lags; it is not a hidden summability theorem.

A useful V7 step is therefore a \emph{conditional, source-local}
entropy or covariance estimate for the connected part between two
separating collars, rather than another unmarked global pressure
comparison. The Gaussian marginal localization proved here provides
local normalization and unbounded-score budgets that such an argument
can reuse. A new proof must additionally show how native communication
through the intervening fast variables decays, or bound its connected
source-visible contribution. Replacing \(d_a\) by an assumed uniform
small number, using a native functional inequality not yet proved, or
differentiating the likelihood-pressure remainder would be circular.
The inherited f15 V23 counterexample already warns that small entropy
alone does not yield the missing native functional inequality; it is
not reproduced here.

All estimates remain for fixed rescaled bridge windows before the
remaining root Gauss restriction. The source directions used here are
coordinate directions, not newly identified physical Hamiltonian
states. Gauge-invariant bridge-tail coverage, a comparison with the
same full-transfer top eigenvalue, physical-time calibration and a
nontrivial interacting continuum construction remain additional
requirements.

\subsection{Verification and exact preservation}

The new diagnostic is
\begin{center}\small
\path{verify_ym_f16_v6_marked_entropy.py}.
\end{center}
Its 906 exact assertions passed. They check the compensated likelihood
identity and scalar signs, correlated Gaussian marginal-projection
matrices and overlap budgets, noncommuting history-coordinate source
cancellations, literal finite-torus source carriers, and the two source
Taylor orders. The finite carrier checks use fine spatial tori of sides
eight and ten, with endpoint and distant-pair support counts kept.
There are no floating-point fixtures in this checker. These tests
are diagnostics of the indicated algebra and support bookkeeping,
not formal verification of the analytic entropy, interpolation, or
source-tail arguments. The report is
\path{ym_v6_verification_report.json}.
The checker has 15,240 bytes and SHA--256
\begin{center}\scriptsize\ttfamily
598A2E0AAA1F37E32BFB5114B5E83B37C76F563888407158C4894E02B62ED9C5.
\end{center}

The V4 and V5 diagnostics available in the supplied V5 package were
rerun separately and passed 9,467 and 89,339 assertions, respectively.
Their exact input scripts, result reports and hashes are supplied.
No unavailable older checker is represented as rerun. The V5 predecessor
has 153,945 bytes and SHA--256
\begin{center}\scriptsize\ttfamily
7751C0A73C4344D29305D84C8180BFE25A97CF0F33696394113742B490CF4F82.
\end{center}

The cumulative V6 changes only the V5 title version and inserts this
appendix before its final \verb|\end{document}|. Removing this exact
appendix and restoring the title recovers V5 byte for byte. The input
hash manifest and the preservation/build reports accompany the
artifacts. The original uploads and earlier generated sources remain
unchanged. This is not a scheduled five-version companion checkpoint;
V5's next scheduled checkpoint remains V10.

\begin{firewall}
V6 proves a full-native likelihood and entropy-density comparison to
the complete Gaussian fast-history law, then a source-marked
separated-time Hessian comparison for bounded-overlap banks with
vanishing mean absolute error per cell. It estimates the actual
full/core common-source correction in that density sense and,
at zero bridge history, pointwise uniformly in spatial volume with
an explicit length cost. No arbitrary-history pointwise bound,
time-summable full-law memory estimate, native functional inequality,
gauge-invariant bridge-tail coverage, full-transfer spectral gap,
interacting continuum construction, or positive physical Yang--Mills
mass gap is established by this appendix.
\end{firewall}
% END F16 V6 APPEND

% BEGIN F16 V7 APPEND -- 2026-09-17
% Insert immediately before the final \end{document} in cumulative V6.
% The predecessor body is unchanged; the cumulative title becomes V7.

\clearpage
\section{V7: Haar-compensated common sources and whole-history trace-norm memory}
\label{f16v7:context}

V6 estimates the covariance of the original bridge scores by comparison
with their affine Gaussian limits. Its cross terms cost the square root
of the local entropy, and its separated-time error has no decay in the
separation. Repeating that argument at more lags is not an improvement.
Here we first change the \emph{exact source representation}, using Haar
integration by parts on the temporal ports. For common-bridge directions,
the new Gaussian score is deterministic. The complete nonlinear returned
covariance is consequently the Gram of two remainder scores, not a cross
term between a remainder and an order-one Gaussian score.

This yields a different output from V6. We bound the \emph{entire}
nonadjacent-time common-source Hessian in trace norm per cube and time
slice, simultaneously including all lags and spatial displacements.
The rate is $O_r(\epsilon_\gamma)$ rather than
$O_r(\sqrt{\epsilon_\gamma})$. The normalization by the number of cells
is essential: a trace-norm density is not a uniform operator norm, and
it is not an absolutely summable temporal row estimate. The same bound
is obtained with exactly finite spatial supports of logarithmically
growing radius. No native functional inequality is assumed.

\subsection{Dependencies and the non-duplication boundary}

The supplied V6 appendix was read in full. Section inventories and
focused searches of f14 V100, f15 V100 and the Grand Tensor Monograph
V076 were followed by reading the relevant interfaces. This is not an
independent certification of all archived proofs. Haar integration by
parts itself is already present in f14 V97, especially
\texttt{prop:v97-finite-Stein}, and in f15 V5--V6. The monograph's
\texttt{thm:YM-Wilsonian-signed-head} already keeps the full direct term
minus one covariance return. Those facts are inherited, not claimed anew.

The particular temporal-port lift below uses V2--V4's kinetic square and
V1's $A^{\mathsf T}D\mathcal E=0$. Its native energy estimate uses V6's
\emph{already proved} complete likelihood bound. There is no new corridor
pressure calculation. Gaussian logarithmic Sobolev and finite-matrix
trace-ideal arguments are standard tools; their precise versions are
proved or derived below. The new application is the exact compensated
common-source bank, its full-native quadratic energy budget, and the
whole-history trace-norm estimate for the actual full/core correction.

\subsection{A uniformly bounded temporal-port lift of every common source}

Keep the unmoving fast variables $\zeta=(x,a)$, full laws
$\mu_\gamma,\mu_0$, and potentials $U_\gamma,U_0$ of
\eqref{f16v6:laws}. Thus $a$ is the principal logarithm of V5's temporal
\emph{group} variable $h$, not V4's moving variable $r$. Set
\[
 \mathcal N=B(n+1),\qquad
 \epsilon_\gamma=\gamma^{-1/6}(1+\log\gamma)^{2/3},\qquad
 \max_{i,e}|y_{i,e}|\le r<\infty.
\]
Unless stated otherwise, all constants are independent of $B=m^3$,
$m\ge4$, $n\ge1$, and the admitted bridge history. Thresholds and
constants with subscript $r$ may depend on its fixed window.

The real common-source space is
\[
 \mathcal U=\bigoplus_{i=0}^n\mathcal E,
 \qquad \dim\mathcal U=9B(n+1)=9\mathcal N.
\]
Use the orthonormal basis from V6: at time $i$, each of the four parallel
bridges from $b$ to $b+e_\mu$ carries $e_k/2$, for
$\mu,k\in\{1,2,3\}$. Denote a basis index by $\alpha=(i,b,\mu,k)$.
The norm on $\mathcal U$ is the inherited Euclidean norm of the bridge
variations, not a physical Hamiltonian norm.

Write $E=E_{\rm br}$ in this appendix only, and retain
\[
 L_c=\mathsf B_B^{\mathsf T}\mathsf B_B,\qquad
 L_h=L_c+E^{\mathsf T}E,\qquad
 G^{-1}=I-EL_h^{-1}E^{\mathsf T}.
\]
Colour identity factors are suppressed. For a common history variation
$\mathbf v=(v_0,\ldots,v_n)$ put
$(d_t\mathbf v)_t=v_t-v_{t+1}$ and define the real linear map
\begin{equation}
 (\mathsf P\mathbf v)_t=-L_h^{-1}E^{\mathsf T}(v_t-v_{t+1}).
 \label{f16v7:lift}
\end{equation}
Its range is the $21Bn$-dimensional temporal-port Lie algebra. The matrix
$\mathsf P$ depends on the finite geometry, not on $\zeta$, $\mathbf y$
or $\gamma$.

\begin{proposition}[Common-source cancellation and finite-range lifts]
\label{f16v7:lift-prop}
Let $\partial_{\mathbf v}$ differentiate the retained bridge history.
In the complete Gaussian model,
\begin{equation}
 \left(\partial_{\mathbf v}
              +(\mathsf P\mathbf v)\cdot\partial_a\right)U_0
                  =g_{\mathbf v}(\mathbf y),
 \label{f16v7:Gaussian-cancellation}
\end{equation}
where the right side is independent of every fast variable and equals
\begin{equation}
 g_{\mathbf v}(\mathbf y)
 =\sum_{i=0}^n\langle v_i,D^{\mathsf T}Dy_i\rangle
   +\sum_{t=0}^{n-1}
       \langle v_t-v_{t+1},G^{-1}(y_t-y_{t+1})\rangle.
 \label{f16v7:g}
\end{equation}
Moreover $\|\mathsf P\|\le 2\|L_h^{-1}E^{\mathsf T}\|\le C$.
There are matrices $\mathsf P_\ell$, $\ell=0,1,\ldots$, supported within
spatial cube distance $\ell+1$ of each source and on its two adjacent
port bonds, with
\begin{equation}
 \|\mathsf P_\ell\|\le C,\qquad
 \|\mathsf P_\ell-\mathsf P\|\le C q_h^{\ell+1},\qquad 0<q_h<1.
 \label{f16v7:lift-approximation}
\end{equation}
All constants are independent of volume and duration.
\end{proposition}

\begin{proof}
Return only for this calculation to the Gaussian moving variables
$z_t=a_t-\mathsf F_B(x_t-x_{t+1})$ from V5. At fixed $x$,
$\partial_a=\partial_z$. The port-dependent part of the quadratic
potential is
\[
 \frac12\sum_t
 \{z_t^{\mathsf T}L_hz_t+
       2z_t^{\mathsf T}E^{\mathsf T}(y_t-y_{t+1})
                         +|y_t-y_{t+1}|^2\}.
\]
There is no $xz$ cross term by the balanced kinetic square. The magnetic
bridge derivative in a common direction has no $x$ term because
$A^{\mathsf T}Dv_i=0$. Its remaining term is
$\langle v_i,D^{\mathsf T}Dy_i\rangle$. Substitution of
\eqref{f16v7:lift} cancels every $z_t$ coefficient and leaves
\eqref{f16v7:g}, with the stated $G^{-1}$. This reuses the inherited
Gaussian cancellation; the next subsection implements it by a legitimate
\emph{compact} integration-by-parts field.

For the norm and support estimates take fixed geometric bounds
\[
 \lambda_-:=\lambda_{\min}(\mathsf B^{\mathsf T}\mathsf B)>0,
 \qquad \lambda_+:=\|\mathsf B^{\mathsf T}\mathsf B\|+6,
 \qquad q_h=1-\lambda_-/\lambda_+.
\]
V2 gives $\|E\|^2\le6$, so
$\lambda_-I\le L_h\le\lambda_+I$. Also $\|d_t\|\le2$.
Set
\[
 T_h=I-L_h/\lambda_+,\qquad
 Q_\ell=\lambda_+^{-1}\sum_{j=0}^{\ell}T_h^j,
 \qquad
 (\mathsf P_\ell\mathbf v)_t=-Q_\ell E^{\mathsf T}(v_t-v_{t+1}).
\]
Then $0\le T_h\le q_hI$, $\|Q_\ell\|\le\lambda_-^{-1}$ and
\[
 L_hQ_\ell=I-T_h^{\ell+1},\qquad
 \|Q_\ell-L_h^{-1}\|\le\lambda_-^{-1}q_h^{\ell+1}.
\]
The spatial blocks of $L_h$ have nearest-cube range; its powers have
the corresponding degree range. A localized common bridge source enters
$E^{\mathsf T}$ only at its endpoint cubes. These facts prove the stated
support and both norm bounds. No decay of a native covariance has been
used. Hereafter $\mathsf P_\infty=\mathsf P$ and
$q_h^{2(\infty+1)}=0$ are convenient conventions.
\end{proof}

\subsection{The exact compact source identity and its changed contact term}

For a temporal port $p=(t,b,v)$ and colour $k$, define the
\emph{left-translation derivative}
\[
 (X^\gamma_{p,k}F)(x,h;\mathbf y)
  =\left.\frac{d}{ds}
     F(x,\ldots,\operatorname{Exp}(s e_k/\sqrt\gamma)h_p,
                         \ldots;\mathbf y)\right|_{s=0}.
\]
This terminology specifies the flow, without assigning it an incorrect
Lie-bracket sign. These derivatives do not in general commute on a
single group. Their reference divergence is zero for product Haar.
They differentiate neither the chord variables nor their endpoint
amplitudes. In particular $X^\gamma_{p,k}$ must \emph{not} be replaced
by flat differentiation of $a_{p,k}$ at finite $\gamma$.

For $\ell\in\{0,1,\ldots,\infty\}$ and a basis mark $\alpha$, set
\begin{align}
 V_\alpha^{[\ell]}
   &=\sum_{p,k}(\mathsf P_\ell)_{p,k;\alpha}X^\gamma_{p,k},&
 \mathcal D_\alpha^{[\ell]}&=\partial_\alpha+V_\alpha^{[\ell]},
 \notag\\
 \widehat J_\alpha^{[\ell]}&=\mathcal D_\alpha^{[\ell]}U_\gamma,&
 r_\alpha^{[\ell]}&=\widehat J_\alpha^{[\ell]}-g_\alpha(\mathbf y),
 \label{f16v7:lifted-scores}
\end{align}
where $g_\alpha=g_{\mathbf v_\alpha}$. All the coefficients of $V$ are
fixed as the bridge history varies. Let $\mathbf r^{[\ell]}$ be the
complete vector of these $9\mathcal N$ residual scores.

\begin{theorem}[Full Haar-compensated Hessian, with all contact terms]
\label{f16v7:Haar-identity}
At each finite regulator, for the \emph{full} native conditional,
\begin{align}
 \partial_\alpha(-\log Z_\gamma)
     &=\mathbb E_\gamma\widehat J_\alpha^{[\ell]},\notag\\
 \partial_\alpha\partial_\beta(-\log Z_\gamma)
     &=\mathcal K^{[\ell]}_{\alpha\beta}
                              -\Gamma^{[\ell]}_{\alpha\beta},
 \label{f16v7:Hessian-decomposition}
\end{align}
where
\begin{align}
 \mathcal K^{[\ell]}_{\alpha\beta}
   &=\frac12\mathbb E_\gamma
      (\mathcal D_\alpha^{[\ell]}\mathcal D_\beta^{[\ell]}
        +\mathcal D_\beta^{[\ell]}\mathcal D_\alpha^{[\ell]})U_\gamma,
       \notag\\
 \Gamma^{[\ell]}_{\alpha\beta}
   &=\operatorname{Cov}_{\mu_\gamma}
                         (r_\alpha^{[\ell]},r_\beta^{[\ell]}),
       \qquad \Gamma^{[\ell]}\succeq0.
 \label{f16v7:contact-Gram}
\end{align}
If the source times obey $|i(\alpha)-i(\beta)|\ge2$, then
$\mathcal K^{[\ell]}_{\alpha\beta}=0$ \emph{exactly}. Thus every such
full-history Hessian entry is the negative of the indicated residual
covariance, also when $-\log Z_\gamma$ is replaced by
$-\log\mathcal W_{\gamma,n}^{\mathrm{full}}$.
\end{theorem}

\begin{proof}
Perform the port integrations in compact group variables before using
principal charts. At fixed $x$, their reference is product probability
Haar, up to a scalar independent of $x,h,\mathbf y$. The remaining chord
Jacobian factors depend on $x$ only. Haar invariance gives
\begin{equation}
 \mathbb E_\gamma[V_\alpha^{[\ell]}F]
    =\mathbb E_\gamma[F V_\alpha^{[\ell]}U_\gamma],
 \qquad \mathbb E_\gamma V_\alpha^{[\ell]}U_\gamma=0.
 \label{f16v7:Haar-parts}
\end{equation}
There is no artificial boundary term at the logarithm cut: it is not a
boundary of the compact port group. At finite $\gamma,B,n$, the relevant
functions and their port and bridge derivatives are bounded in the port
groups and in the closed principal chord ranges. The chord weights are
integrable. Fubini and dominated differentiation therefore justify all
identities, also if a port flow crosses its principal-chart cut.

Ordinary bridge differentiation gives the first line of
\eqref{f16v7:Hessian-decomposition}, by the second identity in
\eqref{f16v7:Haar-parts}. Differentiating that expectation in $\beta$
gives
\[
 \mathbb E_\gamma\partial_\beta\widehat J_\alpha^{[\ell]}
       -\operatorname{Cov}_\gamma
          (\widehat J_\alpha^{[\ell]},\partial_\beta U_\gamma).
\]
In the covariance substitute
$\partial_\beta U_\gamma=
 \widehat J_\beta^{[\ell]}-V_\beta^{[\ell]}U_\gamma$.
The added covariance is
$\mathbb E_\gamma V_\beta^{[\ell]}\widehat J_\alpha^{[\ell]}$
by \eqref{f16v7:Haar-parts}. This yields
\[
 \mathbb E_\gamma
       \mathcal D_\beta^{[\ell]}\mathcal D_\alpha^{[\ell]}U_\gamma
 -\operatorname{Cov}_\gamma
       (\widehat J_\alpha^{[\ell]},\widehat J_\beta^{[\ell]}).
\]
Interchange $\alpha,\beta$ and average. Subtraction of the deterministic
$g$'s does not change the covariance. This proves the symmetric contact
formula without assuming commuting group derivatives. Equivalently,
the expectation of the commutator contribution vanishes; it is another
Haar-divergence identity, not a pointwise zero.

Each native kinetic term uses ports on \emph{one} bond $t$ and bridge
variables only at $t,t+1$. A spatial magnetic term has no temporal port
and uses only one bridge time. The lift of a source at $i$ uses only
bonds $i-1,i$, even for $\ell=\infty$. For source times at least two
apart, these two bond sets are disjoint. Expanding the two $\mathcal D$'s
shows that every double derivative of every native term vanishes:
there is no common bond, no two nonadjacent bridge times in one term,
and no source-dependent lift coefficient. This proves the contact
support assertion. The endpoint Gaussian amplitude has no such
variation at all.

Finally V4's exact external factors have logarithms that are sums of
single-time functions or constants. Their Hessians vanish at these
separated times. They have not been estimated, omitted, or renormalized.
\end{proof}

The contact matrix has changed along with the score covariance. Merely
subtracting a Gaussian-looking random variable from the old scores,
while keeping the old direct Hessian, would be wrong. Equation
\eqref{f16v7:Hessian-decomposition} is a new representation of the
\emph{same} Hessian, not an extra negative return to add to V4 or to the
monograph's already integrated response. Positivity is asserted for
$\Gamma^{[\ell]}$, not for the contact or for its time-off-diagonal part.

\subsection{A whole-vector remainder certificate on complete charts}

Let $J_\gamma=(\partial_\alpha U_\gamma)_\alpha$ be the ordinary
common-source bank, and write
\[
 S_\gamma=(X^\gamma_{p,k}U_\gamma)_{p,k},\qquad
 J_0=(\partial_\alpha U_0)_\alpha,\qquad S_0=\partial_aU_0.
\]
The group-insertion formulas defining $S_\gamma$ make sense as analytic
exponential-word functions on all real $\zeta$; use those expressions
for Gaussian estimates outside the native charts. Their restriction to
the charts is precisely the compact derivative, with no change of law.
Put $e^{\rm br}=J_\gamma-J_0$ and $e^{\rm pt}=S_\gamma-S_0$.

\begin{proposition}[Small Gaussian energy and dimension-free Lipschitz norm]
\label{f16v7:remainder-certificate}
The complete residual vector obeys, uniformly in $\ell$,
\begin{equation}
 \sup_{\zeta}\|D_\zeta\mathbf r^{[\ell]}\|\le L_*,\qquad
 \mathbb E_0\|\mathbf r^{[\ell]}\|^2
       \le C_r\mathcal N
                   (\gamma^{-1}+q_h^{2(\ell+1)}),
 \label{f16v7:reference-energy}
\end{equation}
where $L_*\ge1$ is a fixed geometric constant. For $\ell=\infty$ the
second summand on the right is zero. These are statements on the
\emph{whole score vector}; there is no overlap constant growing with
$\ell$ or the number of marks.
\end{proposition}

\begin{proof}
V6 already proves the local bridge-score Taylor bound. For clarity we
check the additional port-insertion derivatives and the whole-vector
norm. Each scalar native term is a rescaled trace of a fixed product
of exponentials of real skew-Hermitian linear arguments. To evaluate a
port insertion, add one factor
$\operatorname{Exp}(s e_k/\sqrt\gamma)$ adjacent to the appropriate
$h_p$ or its inverse and differentiate in the added scalar $s$ at zero.
This avoids differentiating a logarithm or its singular chart Jacobian.
Duhamel differentiation and unitarity give global bounds on the
second derivatives of this rescaled word, and a bound
$C/\sqrt\gamma$ on its third derivatives. At zero all first
variation terms vanish. The quadratic inserted derivative is exactly
$\partial_{a_{p,k}}U_0$; both operations insert the same linear Lie
algebra variable into the linearized kinetic word.

Consequently every component $e_j$ of the two primitive remainder banks
has a fixed-size local fast carrier $C_j$ and satisfies
\begin{equation}
 |e_j(\zeta)|\le C_r\gamma^{-1/2}(1+|\zeta_{C_j}|^2),\qquad
 \|D_{\zeta_{C_j}}e_j\|\le C.
 \label{f16v7:primitive-remainders}
\end{equation}
The first estimate integrates the third-derivative bound twice; the
second uses the global second-derivative bounds for the native inserted
score and the Gaussian score separately. The same reasoning applies to
negative occurrences of a port: the adjacent inverse and its sign are
kept in the insertion. Neither the fast Haar Jacobian nor a derivative
of the endpoint amplitude is part of $S_\gamma$, since the integration
by parts was in Haar measure and varied ports only.

There are $9\mathcal N$ bridge components and $21Bn$ port components.
Each row of their Jacobian has a fixed number of entries, and each
fast coordinate appears in a fixed number of rows. The row and column
Schur bound therefore proves
\[
 \left\|D_\zeta\binom{e^{\rm br}}{e^{\rm pt}}\right\|\le C.
\]
This is an operator norm, not a sum of component Lipschitz constants.
The Gaussian law from V6 has uniformly bounded local means and
covariances. Its fourth moments on every fixed-size $C_j$ are bounded
by $C_r$. Squaring the first bound in
\eqref{f16v7:primitive-remainders} and summing the at most
$30\mathcal N$ components gives
\begin{equation}
 \mathbb E_0(\|e^{\rm br}\|^2+\|e^{\rm pt}\|^2)
                                  \le C_r\mathcal N/\gamma.
 \label{f16v7:primitive-energy}
\end{equation}

The exact Gaussian cancellation supplies the vector identity
\begin{equation}
 \mathbf r^{[\ell]}
   =e^{\rm br}+\mathsf P_\ell^{\mathsf T}e^{\rm pt}
                   +(\mathsf P_\ell-\mathsf P)^{\mathsf T}S_0.
 \label{f16v7:residual-synthesis}
\end{equation}
The last term is zero for the full lift. For a finite lift, $S_0$ is
an affine \emph{Gaussian score}: $\mathbb E_0 S_0=0$ and
\[
 \operatorname{Cov}_0(S_0)=H_{aa}=I_n\otimes L_h,
              \qquad H=D_\zeta^2U_0.
\]
Indeed $\nabla U_0=H(\zeta-m)$, so the covariance of its full gradient
is $H$, and its $a$ corner is the displayed matrix. Therefore
\[
 \mathbb E_0\| (\mathsf P_\ell-\mathsf P)^{\mathsf T}S_0\|^2
 \le \lambda_+\|\mathsf P_\ell-\mathsf P\|_{\rm HS}^2
 \le C\mathcal N q_h^{2(\ell+1)}.
\]
Use $\|\mathsf P_\ell\|\le C$ in
\eqref{f16v7:residual-synthesis} and
\eqref{f16v7:primitive-energy} to prove its squared-norm estimate.
For its Jacobian use additionally $\|D_\zeta S_0\|\le\|H\|\le h_+$.
The same operator-norm synthesis gives $L_*$ independently of $\ell$.
In particular a growing finite support has not been paid by a growing
entropy-overlap factor.
\end{proof}

\subsection{Transferring the residual energy, not differentiating pressure}

Here is the standard Gaussian estimate needed for the transfer. It is
stated separately to distinguish a Gaussian functional inequality from
an unproved native one.

\begin{proposition}[Gaussian exponential energy of a Lipschitz vector map]
\label{f16v7:Gaussian-energy}
Let $Q$ be any finite-dimensional Gaussian of covariance $C\preceq cI$.
Let $R$ be a smooth vector-valued map with $\|DR\|\le L$, and set
$F=\|R\|^2$, $a=2cL^2$. For $0<t<a^{-1}$,
\begin{equation}
 \log\mathbb E_Q e^{tF}\le\frac{t\mathbb E_QF}{1-at}.
 \label{f16v7:Gaussian-energy-mgf}
\end{equation}
The target dimension of $R$ does not enter the constant.
\end{proposition}

\begin{proof}
We use the Gaussian logarithmic Sobolev inequality in the convention
\[
 \operatorname{Ent}_Q(e^f)
       \le\frac12\mathbb E_Q[e^f\langle\nabla f,C\nabla f\rangle].
\]
For a direct derivation of this convention, whiten $Q$. V6's Gaussian
entropy dissipation identity is
$\operatorname{Ent}_g u=\int_0^\infty I(P_su)\,ds$.
The Mehler commutation $\nabla P_su=e^{-s}P_s\nabla u$ and
conditional Cauchy--Schwarz give
$I(P_su)\le e^{-2s}I(u)$. Integrating gives the factor $1/2$;
affine transport yields the displayed covariance form. This is the
classical Gaussian inequality, not a new native estimate. A primary
reference is L.~Gross,
\href{https://doi.org/10.2307/2373688}{\emph{Logarithmic Sobolev
Inequalities}}, American Journal of Mathematics 97 (1975), 1061--1083.

Now $|\nabla F|^2\le4L^2F$. Write
$\psi(t)=\log\mathbb E_Q e^{tF}$. The displayed inequality gives
\[
 t\psi'(t)-\psi(t)\le at^2\psi'(t).
\]
Hence $((1-at)\psi(t)/t)'\le0$. Its limit at zero is
$\mathbb E_QF$, proving \eqref{f16v7:Gaussian-energy-mgf}.
All differentiations are legitimate for $t<a^{-1}$: global
Lipschitz continuity gives $|R(z)|\le|R(0)|+L|z|$, and the strictly
subcritical Gaussian quadratic exponential has finite moments. Smooth
bounded truncations followed by limits justify the logarithmic Sobolev
step at the same parameters. The constant-map case is immediate.
\end{proof}

\begin{theorem}[A complete-native residual-energy and Gram budget]
\label{f16v7:native-energy}
Define
\[
 \eta_{\gamma,\ell}=\epsilon_\gamma+\gamma^{-1}
                                      +q_h^{2(\ell+1)}.
\]
For all sufficiently large $\gamma$, every admitted bridge history,
and every $\ell\in\{0,1,\ldots,\infty\}$, one has
\begin{align}
 \mathbb E_\gamma\|\mathbf r^{[\ell]}\|^2
                        &\le C_r\mathcal N\eta_{\gamma,\ell},
                              \notag\\
 0\le\operatorname{Tr}\Gamma^{[\ell]}
                        &\le C_r\mathcal N\eta_{\gamma,\ell}.
 \label{f16v7:native-energy-bound}
\end{align}
There is also a fixed $s_*>0$, independent of all these regulators and
of $\ell$, such that
\begin{equation}
 \log\mathbb E_\gamma
               e^{s_*\|\mathbf r^{[\ell]}\|^2}
                         \le C_r\mathcal N\eta_{\gamma,\ell}.
 \label{f16v7:native-mgf}
\end{equation}
In particular the probability that the residual energy exceeds
$(C_r\mathcal N\eta_{\gamma,\ell}+u)/s_*$ is at most $e^{-u}$,
$u\ge0$, with the constant from \eqref{f16v7:native-mgf}.
\end{theorem}

\begin{proof}
V6 provides, for the \emph{same full normalized laws},
\[
 D(\mu_\gamma\Vert\mu_0)
       \le D_p(\mu_\gamma\Vert\mu_0)
       \le C_r\mathcal N\epsilon_\gamma,
                 \qquad p=1+\theta_*>1.
\]
Their Gaussian covariance obeys $C\preceq h_-^{-1}I$.
Use Proposition~\ref{f16v7:Gaussian-energy} with
$R=\mathbf r^{[\ell]}$, $L=L_*$, $c=h_-^{-1}$, and fix
$t_0=(4cL_*^2)^{-1}$. Then
\[
 \log\mathbb E_0e^{t_0\|\mathbf r^{[\ell]}\|^2}
     \le 2t_0\mathbb E_0\|\mathbf r^{[\ell]}\|^2.
\]
The entropy variational inequality
$t_0\mathbb E_\gamma F\le
D(\mu_\gamma\Vert\mu_0)+\log\mathbb E_0e^{t_0F}$,
together with \eqref{f16v7:reference-energy}, proves the first bound.
Centering can only reduce the sum of squared scores, giving the trace
bound. Conditional means are not set to zero.

For the exponential bound write $p'=p/(p-1)$ and $s_*=t_0/p'$.
H\"older's inequality for the full density $f=d\mu_\gamma/d\mu_0$
gives
\begin{align*}
 \log\mathbb E_\gamma e^{s_*F}
  &\le \frac1p\log\mathbb E_0f^p
                   +\frac1{p'}\log\mathbb E_0e^{t_0F}\\
  &=\frac1{p'}D_p(\mu_\gamma\Vert\mu_0)
                   +\frac1{p'}\log\mathbb E_0e^{t_0F}\\
  &\le C_r\mathcal N\epsilon_\gamma
                   +2s_*\mathbb E_0F.
\end{align*}
Insert the same Gaussian energy bound. Exponential Markov proves the
last assertion. All constants are uniform in the polynomial lift degree.
The only logarithmic Sobolev inequality used was Gaussian; a native
logarithmic Sobolev or Poincar\'e inequality has not been assumed.
\end{proof}

For the full lift this is $O_r(\mathcal N\epsilon_\gamma)$, since
$\gamma^{-1}\le\epsilon_\gamma$ eventually. The same rate holds with
an exactly finite-range lift: choose any integer $\ell_\gamma\ge0$ with
\[
 \ell_\gamma+1\ge
           \frac{\log(\epsilon_\gamma^{-1})}{2\log(q_h^{-1})}.
\]
Its spatial radius is $O(\log\gamma)$, and its two-bond temporal
support never grows. This finite-range version concerns the actual
compact source identity, not a truncated native action or measure.

\subsection{The whole nonlocal history Hessian in trace norm}

Let $\Pi_i$ be projection onto the $9B$ common-source coordinates at
time $i$. For any matrix $M$ on $\mathcal U$, define its
nonadjacent-time part by
\[
 \operatorname{Off}_{>1}(M)
      =\sum_{|i-j|\ge2}\Pi_iM\Pi_j.
\]
We use the matrix trace norm $\|M\|_{S_1}$, the sum of singular values.
It is not the entrywise $\ell^1$ norm. Define the \emph{actual} memory
matrix and the original full/core correction matrix by
\begin{align}
 \mathsf M_\gamma
    &=\operatorname{Off}_{>1}
         \bigl(\nabla_{\mathcal U}^2
                   [-\log\mathcal W_{\gamma,n}^{\mathrm{full}}]\bigr),
                  \notag\\
 \mathsf C_{\gamma,R}
    &=\operatorname{Off}_{>1}
                  \bigl(\nabla_{\mathcal U}^2\log p_{\rm good}\bigr).
 \label{f16v7:memory-matrices}
\end{align}
The derivatives are evaluated at the chosen, possibly inhomogeneous,
bridge history. V4's $p_{\rm good}$ and its moving-coordinate core event
are unchanged.

\begin{theorem}[All-lag common-source memory and the full/core correction]
\label{f16v7:whole-memory}
Uniformly in spatial volume and history length, for all sufficiently
large $\gamma$,
\begin{equation}
 \boxed{\quad
   \frac{\|\mathsf M_\gamma\|_{S_1}}{B(n+1)}
                                  \le C_r\epsilon_\gamma.\quad}
 \label{f16v7:whole-memory-bound}
\end{equation}
This includes all time lags and all spatial displacements in one matrix.
For every core radius $R$ satisfying V4's derivative regime,
\begin{equation}
 \boxed{\quad
  \frac{\|\mathsf C_{\gamma,R}\|_{S_1}}{B(n+1)}
                  \le C_r\epsilon_\gamma+C\frac{R^2}{\gamma}.\quad}
 \label{f16v7:whole-correction-bound}
\end{equation}
In particular $R=4\sqrt{\log\gamma}$ is admissible eventually and the
second bound is also $C_r\epsilon_\gamma$. Either the full lift or the
finite lift $\ell_\gamma$ above proves these same statements about the
same full native Hessian.
\end{theorem}

\begin{proof}
The exact contact support gives, for every $\ell$,
\begin{equation}
 \mathsf M_\gamma=-\operatorname{Off}_{>1}(\Gamma^{[\ell]}).
 \label{f16v7:off-Gram}
\end{equation}
Although the full contact and Gram depend on the lift, their
nonadjacent-time combination in \eqref{f16v7:off-Gram} does not.
No sign is assigned to this off-diagonal matrix.

Here is a duration-independent trace-norm bound for this extraction.
Complexify the source space and put
$U_s=\sum_{i=0}^n e^{\mathrm i i s}\Pi_i$. The Fourier coefficients
\[
 \Phi_k(M)=\frac1{2\pi}\int_0^{2\pi}
                     e^{-\mathrm i k s}U_sMU_s^*\,ds
                 =\sum_{i-j=k}\Pi_iM\Pi_j
\]
have $\|\Phi_k(M)\|_{S_1}\le\|M\|_{S_1}$ by unitary invariance
and the triangle inequality. Thus
\begin{equation}
 \|\operatorname{Off}_{>1}(M)\|_{S_1}
   =\|M-\Phi_{-1}(M)-\Phi_0(M)-\Phi_1(M)\|_{S_1}
                           \le4\|M\|_{S_1}.
 \label{f16v7:mask}
\end{equation}
For a positive Gram, $\|\Gamma^{[\ell]}\|_{S_1}
=\operatorname{Tr}\Gamma^{[\ell]}$. Use the full lift in
Theorem~\ref{f16v7:native-energy}, or $q_h^{2(\ell_\gamma+1)}
\le\epsilon_\gamma$, to obtain \eqref{f16v7:whole-memory-bound}.
The factor four has not been multiplied by the number of lags.

For the core, V4's whole common-source tail estimate
\eqref{f16v4:common-tail}, with its cutoff index equal to one, gives
\[
 \left\|\operatorname{Off}_{>1}
     (\nabla_{\mathcal U}^2[-\log\mathcal W_{\gamma,n}^{R}])
                                     \right\|_{\rm op}
 \le\frac{4K_{\rm c}^2}{1-\vartheta}\delta^2,
 \qquad \vartheta=9/19,\quad \delta=C_*R/\sqrt\gamma.
\]
This estimate acts on arbitrary common-source history variations, not
only on localized basis marks. Since $\dim\mathcal U=9\mathcal N$,
its trace norm divided by $\mathcal N$ is at most
$36K_{\rm c}^2\delta^2/(1-\vartheta)$. Finally the exact V4 identity
has the sign
\[
 \nabla^2\log p_{\rm good}
     =\nabla^2[-\log\mathcal W_{\gamma,n}^{\mathrm{full}}]
                  -\nabla^2[-\log\mathcal W_{\gamma,n}^{R}].
\]
Apply $\operatorname{Off}_{>1}$ and the trace-norm triangle inequality.
This proves \eqref{f16v7:whole-correction-bound}. No endpoint normalizer
or chart coverage hypothesis was changed.
\end{proof}

\subsection{What the stronger matrix budget does and does not imply}

There are two useful, rigorous readings of the new bound. First it
improves V6's fixed-lag bank rate. For any family of pairs of canonical
common-source basis marks with each mark appearing in at most $q_0$
pair endpoints, Gram positivity gives
\[
 \sum_a|\Gamma^{[\infty]}_{\alpha_a\beta_a}|
 \le\frac12\sum_a
       (\Gamma^{[\infty]}_{\alpha_a\alpha_a}
                         +\Gamma^{[\infty]}_{\beta_a\beta_a})
 \le\frac{q_0}{2}\operatorname{Tr}\Gamma^{[\infty]}.
\]
Thus each translated fixed-lag, fixed-displacement common-source bank
from V6 now has mean absolute full-native response per cell at most
$C_r\epsilon_\gamma$, and its full/core correction at lag $d\ge2$
is at most
\[
 C_r\epsilon_\gamma
      +C(R^2/\gamma)(9/19)^{d-2}.
\]
The same conclusion for fixed unit colour combinations follows by their
norm-one synthesis from the three colour basis marks. At zero history,
spatial translation invariance makes
$\mathbb E_\gamma|r_{i,b,\mu,k}^{[\infty]}|^2$ independent of $b$.
The total energy bound consequently gives the individual-entry estimate
\begin{equation}
 |\partial_\alpha\partial_\beta
           [-\log\mathcal W_{\gamma,n}^{\mathrm{full}}]|\big|_{\mathbf y=0}
       \le C(n+1)(\epsilon_\gamma+\gamma^{-1}),
       \qquad |i(\alpha)-i(\beta)|\ge2.
 \label{f16v7:zero-entry}
\end{equation}
Here Cauchy--Schwarz was applied to the two \emph{residual} variances.
The length cost remains; no temporal stationarity is assumed.

Second, the whole-history bound controls the distribution of singular
values even when all lags are present. For any $\tau>0$,
\begin{equation}
 \frac{\#\{j:\sigma_j(\mathsf M_\gamma)>\tau\}}{9\mathcal N}
        \le\frac{C_r\epsilon_\gamma}{\tau}.
 \label{f16v7:singular-count}
\end{equation}
This follows by summing the singular values above $\tau$. As the matrix
is real symmetric, its spectral projection $P_\tau$ onto those absolute
eigenvalues obeys
\[
 \operatorname{rank}P_\tau\le C_r\mathcal N\epsilon_\gamma/\tau,
 \qquad \|\mathsf M_\gamma(I-P_\tau)\|_{\rm op}\le\tau.
\]
With $\tau=\sqrt{\epsilon_\gamma}$, the exceptional fraction and the
complement operator norm both vanish. The analogous statements hold for
$\mathsf C_{\gamma,R}$ with $\epsilon_\gamma$ replaced by
$\epsilon_\gamma+R^2/\gamma$.

These spectral projections live in the \emph{finite common-source
coordinate space}. They can depend on the complete base history. They
are not physical low-energy projections and cannot be discarded in a
Yang--Mills gap proof. A vanishing trace-norm density allows a coherent
exceptional direction with a large eigenvalue: a positive rank-one matrix
$\eta\mathbf1\mathbf1^{\mathsf T}$ in dimension $M$ has both trace and
operator norm $M\eta$. In particular neither
\eqref{f16v7:whole-memory-bound} nor the residual-energy exponential
moment supplies a bound on every normalized linear combination of
scores independent of volume.

\subsection{The remaining noncircular target}

The new estimate avoids V6's loss from summing its error floor over all
lags, but it does so in \emph{normalized trace norm}, not by proving
entrywise decay or temporal row summability. The small positive Gram
in \eqref{f16v7:contact-Gram} identifies the precise remaining object.
One possible stronger target is
\[
 \operatorname{Var}_{\mu_\gamma}
       \left(\sum_\alpha v_\alpha r_\alpha^{[\ell_\gamma]}\right)
                 \le a_\gamma\sum_\alpha|v_\alpha|^2,
 \qquad a_\gamma\longrightarrow0,
\]
uniformly in $v,B,n$ and the admitted histories. A time-weighted
connected-covariance estimate for these same residuals would additionally
address summability. Neither statement follows from the trace bound.

The important preparation for a collar or capacity argument is now
concrete: each residual can be built from exact compact Haar flows with
finite spatial radius $O(\log\gamma)$ and only two port times; its entire
native squared-energy density is small; and the Gaussian affine kinetic
score has already been removed with its contact cost retained. A future
argument must control \emph{coherent covariance} of these corrected
sources through the intervening native variables. Repeating the full
likelihood calculation or assuming a native functional inequality would
not supply that control. The source lift is not a physical gauge fixing
and does not establish gauge-invariant coverage of the retained bridge
window. The same-top full-transfer spectral comparison, physical-time
calibration, and interacting continuum construction remain separate. No infrared coercivity or positivity
for the changed contact matrix has been established either.

\clearpage
\subsection{Diagnostics, preservation, and claim scope}

The new diagnostic is
\begin{center}\small
\path{verify_ym_f16_v7_haar_memory.py}.
\end{center}
It checks the lift and its polynomial remainder, noncommuting Gaussian
source cancellation, exact compact Haar-derivative identities, the
contact-support bookkeeping, the whole-vector Gram synthesis, Fourier
block extraction, and the scalar logarithmic-moment inequalities used
above. 311 exact assertions passed; the family counts
and their scope are in \path{ym_v7_verification_report.json}.
Finite algebraic fixtures do not verify the analytic likelihood theorem,
Gaussian log-Sobolev theorem, or all native derivative suprema. They are
not proof-assistant verification and do not certify a physical mass gap.
The checker has 11,803 bytes and SHA--256
\begin{center}\scriptsize\ttfamily
EAA753817261B9830D904CC07586C2D4809BE46900B065AB91F063303032324F.
\end{center}

The cumulative V6 predecessor has 196,992 bytes and SHA--256
\begin{center}\scriptsize\ttfamily
CCB21705950AFCF7A33958811670B94E1E959459974DEF7E7D640F5ABA1B3D2D.
\end{center}
Only its title version changes and this exact appendix is inserted before
its final \verb|\end{document}|. Removing the appendix and restoring the
title recovers the supplied V6 byte for byte. Original and earlier
sources are checked against the input manifest. The supplied V4, V5 and V6 diagnostics were rerun unchanged and passed
9,467, 89,339 and 906 assertions respectively.
Unavailable historical scripts are not represented as rerun. This is not
the scheduled V10 companion checkpoint; only the supplied background
sources were available for the present targeted audit.

\begin{firewall}
V7 constructs exact Haar-compensated common-source derivatives and pays
their complete-native residual-energy and positive-Gram densities. It
bounds the entire nonadjacent-time common-source Hessian and the original
full/core correction in trace norm per cube and time slice, with all lags
included at once. A logarithmically growing finite spatial lift gives the
same rate. No volume-independent operator norm of that whole memory,
uniform arbitrary-history bound for every individual source, time-row
summability, native functional inequality, gauge-invariant bridge-tail
coverage, full-transfer spectral gap, interacting continuum theory, or
positive physical Yang--Mills mass gap is established here.
\end{firewall}
% END F16 V7 APPEND
% BEGIN F16 V8 APPEND -- 2026-09-17
% Insert immediately before the final \end{document} in cumulative V7.
% The predecessor body is unchanged; the cumulative title becomes V8.

\clearpage
\section{V8: a native fast-transfer estimate and time-summable coherent memory}
\label{f16v8:context}

V7's small trace-norm density permits a coherent exceptional direction.
Repeating its likelihood-to-energy argument cannot exclude that direction.
This appendix instead controls the \emph{actual propagation between two
corrected source insertions}. It obtains a small covariance operator and
geometric temporal decay, not merely a trace budget. There is a substantive
restriction: the bridge background is constant in time, and the spatial
cost is exponential in the number of cubes. That cost is displayed and
allows a logarithmically growing spatial volume at weak coupling. The
result is uniform in the entire history length, including the two original
dressed ends.

The argument is not a proof of the physical transfer gap. Its transfer
keeps the bridges fixed at every time and propagates only the complete
compact internal and temporal-port variables. No physical Gauss projection,
bridge integration, or physical time calibration is supplied by this
conditional construction.

\subsection{Audit and the restriction that makes a new estimate possible}

The supplied V7 appendix was read in full. The f14/f15 archives and Grand
Tensor Monograph V076 were searched through their result inventories and
relevant passages, rather than their whole proof collections recertified.
In particular, f15 V96 already proves fixed-region interacting spectral
convergence; f15 V97 already gives compact rooted pins and the full cube
oscillator spectrum. Finite-region Perron separation and Mehler
diagonalization are therefore \emph{not} being offered as new general
results. The monograph already distinguishes a conditional response from
the physical Perron law and forbids charging the same covariance twice.

The new object is the full assembled \emph{fixed-bridge fast transfer},
with all mixed plaquettes and all temporal port integrations present.
The new output is its quantitative, two-insertion estimate for V7's
Haar-compensated common sources under the actual dressed-end law. The
complete-chart error, the exact native top value, the dressed-end overlap,
and the source-insertion norms are all paid before taking arbitrarily
large time separation. None of V5--V7's extensive entropy estimates is
used to assert a native functional inequality.

Fix one bridge configuration in rescaled coordinates,
\begin{equation}
 \bar y\in\mathbb R^{36B},\qquad
 \max_e|\bar y_e|\le r<\infty,\qquad
 \bar Z_e=\operatorname{Exp}(\bar y_e/\sqrt\gamma),
 \qquad \mathbf y=(\bar y,\ldots,\bar y).
 \label{f16v8:static-history}
\end{equation}
There is no spatial homogeneity assumption on $\bar y$. All bridge
\emph{derivatives} are taken with respect to the independent time-slice
coordinates first, and only then evaluated at
\eqref{f16v8:static-history}. We are not differentiating only along the
diagonal submanifold of constant histories. Put
\[
 B=m^3,\quad m\ge4,\qquad n\ge1,\qquad \mathcal N=B(n+1).
\]
Constants may depend on $r$ and the fixed cube geometry. Dependence on
$B$ is never hidden in a volume-independent constant below. In particular
an estimate of the form $e^{c_rB}$ is not described as uniform in spatial
volume.

\subsection{The exact frozen-bridge transfer and its dressed-end law}

Use V5's unmoving compact temporal variables $h$, and write
$X=\operatorname{Exp}(x/\sqrt\gamma)$,
$h=\operatorname{Exp}(a/\sqrt\gamma)$ on their complete principal charts
$\mathcal P_x$, $\mathcal P_a$. The dimensions are $15B$ and $21B$.
At an internal tree edge set $X_e=I$; at a chord take its actual chord
variable. Define the original fine spatial representative
\[
 U_e(x;\bar y)=
 \begin{cases}
 X_e,&e\text{ internal},\\
 f_{s(e)}(X)\bar Z_e f_{t(e)}(X)^{-1},&e\text{ a bridge}.
 \end{cases}
\]
The moving frames $f_v$ are exactly those of V1--V7. Set
\begin{align}
 V_{\gamma,\bar y}(x)
   &=\gamma\{S_{\rm int}(X)+
                   S_\times(X,\Theta_X^{-1}\bar Z)\},\notag\\
 W_{\gamma,\bar y}(x,x',a)
   &=\gamma\sum_{e\ {\rm fine}}
       s\bigl(h_{s(e)}^{-1}U_e(x;\bar y)
                          h_{t(e)}U_e(x';\bar y)^{-1}\bigr),\notag\\
 \mathfrak E_{\gamma,\bar y}(x,x',a)
   &=\tfrac12 V_{\gamma,\bar y}(x)
         +W_{\gamma,\bar y}(x,x',a)
                          +\tfrac12 V_{\gamma,\bar y}(x').
 \label{f16v8:one-step-action}
\end{align}
Root entries of $h$ remain the identity. The sum has all $24B$ fine
spatial links; the two magnetic ends contain all $24B$ spatial plaquettes.
Let $j(v)=(\sin|v|/|v|)^2$ and abbreviate the products of these factors by
\[
 j_x(x)=\prod_{q\ {\rm chord}}j(x_q/\sqrt\gamma),\quad
 j_a(a)=\prod_{p\ {\rm nonroot}}j(a_p/\sqrt\gamma),\quad
 j_{\rm br}(\bar y)=\prod_{e\ {\rm bridge}}j(\bar y_e/\sqrt\gamma).
\]
The exact scalar for one step in internal and bridge half-density
coordinates is
\begin{equation}
 \rho_{\gamma,\bar y}
    =\left(\frac{c_{H,\gamma}}{z_\gamma}\right)^{24B}
                                    j_{\rm br}(\bar y),\qquad
 c_{H,\gamma}=(2\pi^2\gamma^{3/2})^{-1}.
 \label{f16v8:rho}
\end{equation}
Define an augmented step density and its integrated kernel by
\begin{align}
 d\mathfrak K_{\gamma,\bar y}(x,x';a)
   &=\rho_{\gamma,\bar y}
      1_{\mathcal P_x}(x)1_{\mathcal P_x}(x')1_{\mathcal P_a}(a)
      \sqrt{j_x(x)j_x(x')}j_a(a)
                   e^{-\mathfrak E_{\gamma,\bar y}(x,x',a)}\,da,\notag\\
 K_{\gamma,\bar y}(x,x')&=\int d\mathfrak K_{\gamma,\bar y}(x,x';a).
 \label{f16v8:kernel}
\end{align}
The operator $K_{\gamma,\bar y}$ acts on $L^2(\mathbb R^{15B},dx)$
with zero extension outside $\mathcal P_x$. The same letter for kernel
and operator will cause no ambiguity. Its endpoint vector is
\begin{equation}
 a_{\gamma,\bar y}(x)=
 \frac{1_{\mathcal P_x}(x)\Omega(x)
       e^{-\gamma S_\times(X,\Theta_X^{-1}\bar Z)/2}}
      {\sqrt{p_{\gamma,B}F_{\gamma,2}(\bar Z)}}.
 \label{f16v8:endpoint}
\end{equation}
All three quantities $\Omega,p_{\gamma,B},F_{\gamma,2}$ have their
inherited definitions. In particular the vector has norm one; it is not
replaced by a newly chosen native vacuum.

\begin{proposition}[An exact conditional transfer, not a compression]
\label{f16v8:exact-transfer}
The kernel \eqref{f16v8:kernel} is symmetric, positive semidefinite as
an operator, Hilbert--Schmidt, and strictly positive on the interiors of
its two principal internal charts. It has a simple positive top
eigenvalue $\lambda_{\gamma,\bar y}$ and a normalized positive top vector
$\phi_{\gamma,\bar y}$ on those charts. Moreover
\begin{equation}
 \mathcal W_{\gamma,n}^{\rm full}(\bar y,\ldots,\bar y)
   =\langle a_{\gamma,\bar y},
               K_{\gamma,\bar y}^{\,n}a_{\gamma,\bar y}\rangle.
 \label{f16v8:exact-power}
\end{equation}
Normalizing the augmented path integral on the right gives precisely
V4--V7's complete native fast-history conditional at
\eqref{f16v8:static-history}. There is no projection at an internal time.
\end{proposition}

\begin{proof}
Start with V2's exact full moving-frame kernel and perform V5's
conditional Haar change $h=f(X)r f(X')^{-1}$. It gives the kinetic
words in \eqref{f16v8:one-step-action}. At the two ends, internal
half-density factors contribute $c_{H,\gamma}^{5B}$ in total,
bridge half-densities contribute $c_{H,\gamma}^{12B}$, and the port
Haar integral contributes $c_{H,\gamma}^{7B}$. There are $24B$
kinetic normalizers $z_\gamma^{-1}$. This gives
\eqref{f16v8:rho} and all the nonconstant Haar factors in
\eqref{f16v8:kernel}. The two equal bridge half-densities give
$j_{\rm br}(\bar y)$, not its square.

Swapping $x,x'$ and replacing every $h_v$ by $h_v^{-1}$ leaves each
kinetic trace unchanged, by inversion and cyclicity; Haar measure is
invariant under this replacement. Thus the kernel is symmetric.
Its operator positivity is inherited from the full positive Wilson
transfer before the remaining root restriction, but requires a small
justification because evaluation at a bridge value is not an $L^2$
projection. In that full kernel test against
$u(x)\eta_\varepsilon(y-\bar y)$, with $u$ compactly supported inside
the internal charts and with a nonnegative bridge approximate identity
of integral one. The test has finite $L^2$ norm for each $\varepsilon$.
Positivity gives a nonnegative quadratic form. Local continuity of the
exact kernel lets $\varepsilon\downarrow0$ and yields
$\langle u,K_{\gamma,\bar y}u\rangle\ge0$. The kernel bound proved
below extends the assertion by density to all $L^2$.

For clarity the bound already follows from the inherited global pins.
The two internal magnetic ends supply the five-chord cube pins, and the
seven tree kinetic terms supply the independent rooted $h$ pin. Thus
\begin{equation}
 \mathfrak E_{\gamma,\bar y}(x,x',a)
       \ge a_*\bigl(|x|^2+|x'|^2+|a|^2\bigr),
             \qquad a_*=(6\pi^2)^{-1},
 \label{f16v8:step-pin}
\end{equation}
on the complete charts. All other terms are nonnegative. With every
Haar factor at most one and a finite prefactor, integrating $a$ gives
an integrable Gaussian square majorant in $x,x'$. This proves the
Hilbert--Schmidt assertion. Strict positivity follows directly from
the positive integrand on the open charts.

The usual compact positive-kernel argument can be given here directly.
The positive compact self-adjoint operator has a maximizing eigenvector.
Replacing a real eigenvector by its absolute value cannot decrease its
Rayleigh quotient, and increases it strictly if it has positive and
negative parts of positive measure. A nonnegative top eigenvector is
strictly positive by its integral equation. A second linearly independent
top eigenvector could be made orthogonal to it, and would change sign,
contradicting the strict inequality. This proves simplicity. The zero
extension adds zero spectrum, not another top eigenvector.

Finally expand the power in \eqref{f16v8:exact-power} using the augmented
kernel. An interior internal slice has a full Haar factor and two
magnetic half-weights. Each external slice has one internal half-weight
and the endpoint vector. The latter supplies the other mixed magnetic
half-weight and $\Omega/\sqrt{p_{\gamma,B}F_{\gamma,2}}$. The products
are exactly V4's action powers and normalizers, including its bridge
Jacobian power $j_{\rm br}(\bar y)^n$. This proves the identity and
the equality of normalized conditional laws.
\end{proof}

The positivity argument is the standard symmetric positive-kernel
Perron argument; see also R.~Jentzsch,
\href{https://doi.org/10.1515/crll.1912.141.235}
{\emph{\"Uber Integralgleichungen mit positivem Kern}} (1912).
The operator here is an exact bridge-diagonal kernel. It is not
$\mathcal I_\gamma^*T_{\rm full,\gamma}\mathcal I_\gamma$, and its
Perron value is not identified with the full physical transfer's top
value.

\subsection{The reference spectrum and the complete-chart error}

Suppress colour identities and put
\begin{align}
 S&=\mathsf H_B+A^{\mathsf T}A,&
 \Lambda&=A^{\mathsf T}D,&
 \xi_{\bar y}&=-S^{-1}\Lambda\bar y,\notag\\
 c(\bar y)&=\bar y^{\mathsf T}
            (D^{\mathsf T}D-\Lambda^{\mathsf T}S^{-1}\Lambda)\bar y,&
 T_S&=\mathsf M_B^{1/2}S\mathsf M_B^{1/2}.
 \label{f16v8:Gaussian-data}
\end{align}
The minimum of $x^{\mathsf T}\mathsf H_Bx+|Ax+D\bar y|^2$ is
$c(\bar y)\ge0$. V2's complete Gaussian kinetic integration gives
\begin{equation}
 K_{0,\bar y}(x,x')
  =\varphi_G(0)e^{-c(\bar y)/2}
    e^{-(x-\xi_{\bar y})^{\mathsf T}S(x-\xi_{\bar y})/4}
    \varphi_{\mathsf M_B}(x-x')
    e^{-(x'-\xi_{\bar y})^{\mathsf T}S(x'-\xi_{\bar y})/4}.
 \label{f16v8:Gaussian-kernel}
\end{equation}
The three-colour densities $\varphi$ are the ones already fixed in V2.
The reference endpoint is
\[
 a_{0,\bar y}(x)
       =F_2(\bar y)^{-1/2}\Omega(x)e^{-|Ax+D\bar y|^2/4}.
\]
Let $\phi_{0,\bar y}$ be the normalized positive ground vector of
\eqref{f16v8:Gaussian-kernel}.

\begin{proposition}[Quantitative frozen-background Gaussian comparison]
\label{f16v8:comparison}
There are finite constants $c_r,C_r\ge1$, independent of $B$, such that
for $\gamma\ge C_r$ and every admitted $\bar y$,
\begin{align}
 e^{-c_r B}&\le\lambda_{0,\bar y}\le e^{c_rB},&
 \frac{\lambda_1(K_{0,\bar y})}{\lambda_{0,\bar y}}
     &\le q_0:=3-2\sqrt2<\frac15,\notag\\
 \|K_{\gamma,\bar y}-K_{0,\bar y}\|_{\rm HS}
     &\le \frac{e^{c_rB}}{\sqrt\gamma},&
 \|a_{\gamma,\bar y}-a_{0,\bar y}\|_2
     &\le \frac{e^{c_rB}}{\sqrt\gamma},\notag\\
 \langle a_{0,\bar y},\phi_{0,\bar y}\rangle
     &\ge e^{-c_rB}.&&
 \label{f16v8:comparison-bounds}
\end{align}
The Gaussian top is, more precisely,
\begin{equation}
 \lambda_{0,\bar y}=\varphi_G(0)e^{-c(\bar y)/2}
             \prod_{j=1}^{5B}r(t_j)^{3/2},\qquad
 r(t)=\bigl(1+t/2+\sqrt{t+t^2/4}\bigr)^{-1},
 \label{f16v8:Gaussian-top}
\end{equation}
where $t_j$ are the eigenvalues of $T_S$ and $4\le t_j\le10$.
The errors include the complete compact complements, not a fixed-window
restriction of the internal variables.
\end{proposition}

\begin{proof}
The inherited whitened bounds for $\mathsf H_B$ and $A^{\mathsf T}A$
give $4I\le T_S\le10I$. Translate by $\xi_{\bar y}$, whiten by
$\mathsf M_B$, and diagonalize $T_S$. The scalar kernel of a mode of
parameter $t$ is
\[
 (2\pi)^{-1/2}
       \exp\{-\tfrac12(1+t/2)(u^2+v^2)+uv\}.
\]
Writing $a=1+t/2$ and $w=\sqrt{a^2-1}$, its ground precision is $w$,
its top eigenvalue is $(a+w)^{-1/2}=r(t)^{1/2}$, and its successive
Hermite eigenvalues have ratio $r(t)$. This is the inherited complete
Mehler calculation, with the general mixed magnetic matrix $S$ now
retained. For reference the identity is
\href{https://dlmf.nist.gov/18.18.E28}{NIST DLMF, equation 18.18.28}.
It proves \eqref{f16v8:Gaussian-top} and the ratio bound without deleting
non-singlet internal excitations.
Since $I\le G\le30I$, $c(\bar y)\le\|D\bar y\|^2\le CBr^2$,
and every $r(t_j)$ lies in the fixed interval $[r(10),r(4)]$, the
displayed exponential bounds on the top follow.

Here are details of the dimensional error. Write $w=(x,x',a)$, with
$51B$ real coordinates, and let $\mathfrak E_0$ be the quadratic limit of
\eqref{f16v8:one-step-action}. Both exponents obey
\eqref{f16v8:step-pin}, the second by taking the quadratic limit of the
compact anchors. Bounded-incidence differentiation of the actual
unitary exponential words gives, throughout the complete native charts,
\begin{equation}
 |\mathfrak E_{\gamma,\bar y}-\mathfrak E_{0,\bar y}|
     \le \frac{C}{\sqrt\gamma}(|w|^3+Br^3).
 \label{f16v8:global-error}
\end{equation}
Indeed each word has uniformly bounded third derivatives in its fixed
local inputs before rescaling. The sum of the resulting local cubic
norms is bounded by $C|w|^3+CBr^3$. No logarithm of a group product is
differentiated. The Haar products obey
\[
 0\le\sqrt{j_x(x)j_x(x')}j_a(a)\le1,\qquad
 1-\sqrt{j_x(x)j_x(x')}j_a(a)\le C|w|^2/\gamma.
\]
This uses $1-\sin u/u\le u^2/6$ and $0\le\sin u/u\le1$ on
$[0,\pi]$, followed by $1-\prod b_j\le\sum(1-b_j)$.
V4's one-link normalization also gives
\[
 \rho_{0}=(2\pi)^{-36B},\qquad
 \left|\log(\rho_{\gamma,\bar y}/\rho_0)\right|
                                      \le C_r B/\gamma.
\]
The bound on $j_{\rm br}(\bar y)$ uses its small fixed bridge chart;
the internal charts have not been shrunk.

Use $|e^{-u}-e^{-v}|\le|u-v|e^{-\min(u,v)}$ and the global pins.
On the common chart domain the difference of augmented integrands is
bounded by a fixed-degree polynomial in $|w|,B$, times
$e^{C_rB}\gamma^{-1/2}e^{-a_*|w|^2}$. The Gaussian complement of a
principal chart has some group coordinate of norm at least
$\pi\sqrt\gamma$. Its integral is bounded by
$e^{C_rB-c\gamma}$, and the same reasoning bounds its contribution to
the kernel's Hilbert--Schmidt norm. For example, split off half the
Gaussian exponent to pay that coordinate threshold and integrate the
other half over all $51B$ coordinates. This incurs an exponential in
$B$, not an unrecorded uniform constant.
Integrate $a$ and then take the $L^2(dx\,dx')$ norm. Fixed-degree
Gaussian moments in dimensions proportional to $B$ are a fixed
polynomial in $B$ times a constant to the power $B$. Such polynomials
are absorbed into $e^{c_rB}$. Also
$e^{-c\gamma}\le C\gamma^{-1/2}$. This proves the kernel error on the
whole zero-extended space.

For the endpoint, first compare the unnormalized vectors
\[
 q_\gamma=1_{\mathcal P_x}\Omega
           e^{-\gamma S_\times/2},\qquad
 q_0=\Omega e^{-|Ax+D\bar y|^2/4}.
\]
Their multipliers are at most one. The same cubic word remainder,
now integrated against $\Omega^2$, and its complete-chart Gaussian
tail give $\|q_\gamma-q_0\|_2\le e^{c_rB}/\sqrt\gamma$.
The exact identities are $\|q_\gamma\|_2^2=p_{\gamma,B}
F_{\gamma,2}(\bar Z)$ and $\|q_0\|_2^2=F_2(\bar y)$.
The latter is bounded below by $e^{-c_rB}$, either from its inherited
determinant formula or by integrating on a product of fixed one-cell
balls. Normalization uses the elementary inequality
\[
 \left\|\frac u{\|u\|}-\frac v{\|v\|}\right\|
                         \le\frac{2\|u-v\|}{\|v\|}
       \quad(u,v\ne0).
\]
This proves the endpoint bound with an enlarged exponent.

Finally both $a_{0,\bar y}$ and $\phi_{0,\bar y}$ are normalized
positive Gaussian amplitudes. Their squared-density covariance
spectra lie between two fixed positive constants, and the squared
norms of their means are at most $C_rB$. For $a_0$ use
$(C_B^{-1}+A^{\mathsf T}A)^{-1}$; for $\phi_0$ use
\[
 C_S=\mathsf M_B^{1/2}
           [2\sqrt{T_S+T_S^2/4}]^{-1}\mathsf M_B^{1/2},
       \qquad \text{mean }\xi_{\bar y}.
\]
Their product, integrated over $[-1,1]^{15B}$, is at least
$e^{-c_rB}$: the two determinant normalizers are exponential in $B$,
and both quadratic exponents are bounded there by $C_rB$. This proves
the overlap lower bound without asserting that the two amplitudes
coincide.
\end{proof}

\subsection{A native spectral ratio and an overlap valid for every duration}

\begin{theorem}[Frozen-background fast propagation with a paid spatial cost]
\label{f16v8:native-propagation}
There are $k_r\ge1$ and $\gamma_r$ such that, whenever
\begin{equation}
 \gamma\ge\gamma_r,\qquad
             \frac{e^{k_rB}}{\sqrt\gamma}\le\frac1{100},
 \label{f16v8:entry-regime}
\end{equation}
the exact operator and endpoint in \eqref{f16v8:kernel}--\eqref{f16v8:endpoint}
satisfy
\begin{equation}
 S_\gamma:=K_{\gamma,\bar y}/\lambda_{\gamma,\bar y}
       =P_\gamma+Q_\gamma,\qquad
 P_\gamma=|\phi_{\gamma,\bar y}\rangle
                 \langle\phi_{\gamma,\bar y}|,\qquad
 0\le Q_\gamma,\quad \|Q_\gamma\|\le\tau:=\frac13,
 \label{f16v8:spectral-split}
\end{equation}
with $P_\gamma Q_\gamma=0$. Moreover, for another fixed $b_r$,
\begin{equation}
 \beta_\gamma:=\langle a_{\gamma,\bar y},\phi_{\gamma,\bar y}\rangle
       \ge e^{-b_rB},\qquad
 \beta_\gamma^2\le
       Z_{\gamma,n}^{\rm norm}:=
          \langle a_{\gamma,\bar y},S_\gamma^n a_{\gamma,\bar y}\rangle
       \le1\quad(n\ge1).
 \label{f16v8:boundary-overlap}
\end{equation}
No factor depending on $n$ occurs in these bounds. Normalization always
uses the \emph{exact} $\lambda_{\gamma,\bar y}$.
\end{theorem}

\begin{proof}
The preceding proposition gives
\[
 \delta_{\gamma,B}:=
 \frac{\|K_{\gamma,\bar y}-K_{0,\bar y}\|}
                         {\lambda_{0,\bar y}}
              \le e^{c_rB}/\sqrt\gamma.
\]
Enlarge $k_r$ so that \eqref{f16v8:entry-regime} implies
$\delta_{\gamma,B}\le1/16$. The variational characterization of the
first two eigenvalues gives
\[
 (1-\delta_{\gamma,B})\lambda_{0,\bar y}
           \le\lambda_{\gamma,\bar y},\qquad
 \lambda_1(K_{\gamma,\bar y})
           \le(q_0+\delta_{\gamma,B})\lambda_{0,\bar y}.
\]
Thus their ratio is at most
$(1/5+1/16)/(1-1/16)=7/25<1/3$. This controls the complete compact
fast operator, including its excited states; it is not a band-compression
calculation. Operator positivity gives the nonnegative spectral split.

We must also pay the dressed ends. Let $P_0$ be the Gaussian top
projection. Projecting the native eigenvector equation onto $I-P_0$
gives
\[
 \|(I-P_0)\phi_{\gamma,\bar y}\|
       \le\frac{\delta_{\gamma,B}}{1-q_0-\delta_{\gamma,B}}.
\]
The two positive top eigenvectors have nonnegative inner product, so
$\|\phi_{\gamma,\bar y}-\phi_{0,\bar y}\|\le3\delta_{\gamma,B}$
in this regime. Combine this with the endpoint comparison and the
Gaussian overlap lower bound. Enlarging $k_r$ again makes the sum of
the two approximation errors at most half that lower bound, yielding
$\beta_\gamma\ge e^{-b_rB}$ for some $b_r$. Such a choice is possible
because all three costs are explicit exponentials in $B$ times
$\gamma^{-1/2}$.

Finally $S_\gamma^n=P_\gamma+Q_\gamma^n$ and $Q_\gamma^n\ge0$.
The lower bound in \eqref{f16v8:boundary-overlap} follows by retaining
the Perron term; the upper bound uses $\|S_\gamma\|=1$ and
$\|a_{\gamma,\bar y}\|=1$. No replacement of
$\lambda_{\gamma,\bar y}^n$ by $\lambda_{0,\bar y}^n$ has occurred.
Such a replacement would in general accumulate an error with duration.
\end{proof}

\subsection{Unbounded corrected scores as bounded transfer insertions}

Use V7's \emph{full} Haar lift $\mathsf P$ and its residuals
$r^{[\infty]}_\alpha$, with no new subtraction. A spatial common variation
$v\in\mathcal E$ at time $i$ gives
\[
 R_i(v)=\sum_{b,\mu,k}v_{b,\mu,k}
                     r^{[\infty]}_{i,b,\mu,k}.
\]
Coefficients refer to V7's orthonormal common-bridge basis, so their
Euclidean norm is $\|v\|$. Arbitrary spatially coherent $v$ are allowed.
For the current purpose the infinite lift means the exact finite-volume
inverse $L_h^{-1}$; it is bounded independently of $B$.

The variable $R_i(v)$ uses only the temporal window
\begin{equation}
 a_i=\max(0,i-1),\qquad b_i=\min(n,i+1),\qquad p_i=b_i-a_i\in\{1,2\}.
 \label{f16v8:source-window}
\end{equation}
It depends on the internal slices in that window and on its $p_i$ port
bonds. Its spatial dependence need not be confined to one cube.
Let $\mathbb A_i(v)$ be the kernel obtained by multiplying the product
of those $p_i$ augmented step densities by $R_i(v)$, then integrating all
ports and intermediate internal slices. Let $\mathbb B_i(v)$ use
$R_i(v)^2$ instead. These are signed and nonnegative-kernel insertion
operators, respectively; a nonnegative kernel is not asserted to be a
positive semidefinite operator. Define their exact normalized versions
\[
 A_i(v)=\lambda_{\gamma,\bar y}^{-p_i}\mathbb A_i(v),\qquad
 B_i(v)=\lambda_{\gamma,\bar y}^{-p_i}\mathbb B_i(v).
\]

\begin{proposition}[Two-insertion norms on the full compact domains]
\label{f16v8:insertions}
After increasing constants in \eqref{f16v8:entry-regime} if needed,
there are fixed $c_r$ such that, for every $n,i$ and $v\in\mathcal E$,
\begin{equation}
 \|A_i(v)\|\le\frac{e^{c_rB}}{\sqrt\gamma}\|v\|,\qquad
 \|B_i(v)\|\le\frac{e^{c_rB}}{\gamma}\|v\|^2.
 \label{f16v8:insertion-norms}
\end{equation}
The constants do not depend on the source time or history length. These
are full native insertion estimates, not estimates of observables under
a Gaussian probability substituted for the native law.
\end{proposition}

\begin{proof}
The additional point beyond V7's Gaussian energy budget is that a source
is now integrated against the complete \emph{step kernel pin}. We give
the pointwise and kernel estimates explicitly.

Separate one common-source derivative into its magnetic contribution at
$i$ and its kinetic contributions on bonds $i-1,i$ when present. At
Gaussian order the common magnetic fast coefficient vanishes by
$A^{\mathsf T}Dv=0$. On each port bond the exact lift cancels the
Gaussian fast coefficient separately, as in V7's kinetic square. At a
constant history the deterministic kinetic contribution is zero;
the deterministic magnetic contribution is precisely the corresponding
part of $g_\alpha(\mathbf y)$. Thus each of the three pieces of $R_i$
is a primitive nonlinear remainder or the bounded $L_h^{-1}E^{\mathsf T}$
synthesis of primitive port remainders.

For a local primitive component $e_j$, V7's complete-chart insertion
calculation gives
\[
 |e_j|\le C_r\gamma^{-1/2}(1+|w_{C_j}|^2).
\]
Here $w$ stacks the at most three internal slices and two port slices of
\eqref{f16v8:source-window}; its dimension is at most $87B$.
Each $C_j$ has fixed size and each primitive input occurs a bounded
number of times. Consequently
\[
 \sum_j|e_j|^2\le\frac{C_r}{\gamma}(B+|w|^4).
\]
For example the sum of squared local squared norms is bounded by their
sum times their maximum, and hence by $C|w|^4$. Applying the bounded
spatial lift and combining the at most three pieces proves the global
pointwise estimate
\begin{equation}
 |R_i(v)|\le\frac{C_r}{\sqrt\gamma}
                          (\sqrt B+|w|^2)\|v\|.
 \label{f16v8:window-remainder}
\end{equation}
No sum of absolute lift coefficients or support-volume factor is needed;
only its $\ell^2$ operator norm is used. The formula comes from group
insertions, so it remains valid across a principal port cut in its
compact interpretation. The analytic exponential-word expressions supply
the same polynomial majorant on all real inputs.

For $p_i=1$ or $2$, the product of augmented step densities is at most
\[
 e^{C_rB}\exp\{-a_*|w|^2\}
\]
with respect to its port and intermediate-slice Lebesgue measures, with
the endpoint internal coordinates left as kernel variables. This follows
by applying \eqref{f16v8:step-pin} to each step; interior slice pins can
be weakened to the displayed common lower bound. Every Haar factor is
at most one. The at most two factors $\rho_{\gamma,\bar y}$ cost only
$e^{C_rB}$.
Multiplication by \eqref{f16v8:window-remainder}, or by its square,
therefore leaves a fixed-degree polynomial times this Gaussian majorant.
Integrate the complete native port charts and intermediate internal
charts by enlarging them to Euclidean space in the majorant, and then
take the kernel Hilbert--Schmidt norm in its two endpoint variables.
The resulting bounds are
$e^{c_rB}\gamma^{-1/2}\|v\|$ and
$e^{c_rB}\gamma^{-1}\|v\|^2$ for the unnormalized insertions.
As in the one-step comparison, Gaussian moments of fixed degree and
dimension at most $87B$ cost an exponential in $B$. There is no time
length in this calculation.
The native top lower bound
$\lambda_{\gamma,\bar y}\ge\tfrac{15}{16}\lambda_{0,\bar y}
\ge e^{-C_rB}$ pays division by its first or second power. This proves
\eqref{f16v8:insertion-norms}. The argument integrates the entire native
domain; it has not thrown away a source-weighted large-field complement.
\end{proof}

\subsection{A two-window covariance identity with the boundary cost retained}

The following elementary transfer argument is recorded because using only
the stationary Perron law would omit the actual dressed-end cost. It
applies as well to observables of integrated auxiliary bond variables.

\begin{proposition}[Finite-history mixing from the exact normalized transfer]
\label{f16v8:two-window}
Let $S=P+Q$ be a positive self-adjoint contraction with positive kernel,
$P=|\phi\rangle\langle\phi|$, $PQ=0$, and
$0\le Q$, $\|Q\|\le\tau<1$. Let $a\ge0$, $\|a\|=1$, and
$\beta=\langle a,\phi\rangle>0$. Use the normalized length-$n$ path
law with these ends. Suppose real observables $F,G$ are inserted in
nonoverlapping bond windows, separated by $g\ge0$ unmarked steps,
and their normalized insertion operators are $A_F,A_G$. Then
\begin{equation}
 |\operatorname{Cov}(F,G)|
       \le4\beta^{-4}\|A_F\|\|A_G\|\tau^g.
 \label{f16v8:two-window-bound}
\end{equation}
For a single observable with square-insertion operator $B_F$,
\begin{equation}
             \mathbb E F^2\le\beta^{-2}\|B_F\|.
 \label{f16v8:one-window-square}
\end{equation}
No stationarity of the finite path law is required.
\end{proposition}

\begin{proof}
Write the left window as $[a_1,b_1]$, the right as $[a_2,b_2]$,
$p=b_1-a_1$, $q=b_2-a_2$, $g=a_2-b_1\ge0$. Set
\[
 u=S^{a_1}a,\quad v=S^{n-b_2}a,\quad
 u_0=S^pu,\quad u_F=A_F^*u,\quad
 v_0=S^qv,\quad v_G=A_Gv,\quad T=S^g.
\]
Then $\|u_0\|,\|v_0\|\le1$,
$\|u_F\|\le\|A_F\|$, and $\|v_G\|\le\|A_G\|$.
The common denominator is
$Z=\langle u_0,Tv_0\rangle=\langle a,S^na\rangle\ge\beta^2$.
The exact centered covariance is
\begin{equation}
 \operatorname{Cov}(F,G)=
 \frac{\langle u_F,Tv_G\rangle\langle u_0,Tv_0\rangle
       -\langle u_F,Tv_0\rangle\langle u_0,Tv_G\rangle}{Z^2}.
 \label{f16v8:covariance-minor}
\end{equation}
Replacing $T$ by the rank-one $P$ makes its numerator zero exactly.
For $g\ge1$, $\|T-P\|=\|Q^g\|\le\tau^g$; for $g=0$ use
$T=I$ and $\|I-P\|\le1=\tau^0$. Each of the two products in the
numerator changes by at most
$2\tau^g\|A_F\|\|A_G\|$, since $\|T\|,\|P\|\le1$.
Division by $Z^2\ge\beta^4$ proves
\eqref{f16v8:two-window-bound}. This cancellation includes the means
under the actual finite path law, not means borrowed from its stationary
limit.
For the square insertion, its expectation has numerator
$\langle a,S^{a_1}B_FS^{n-b_1}a\rangle$, bounded by $\|B_F\|$,
and denominator at least $\beta^2$. This proves
\eqref{f16v8:one-window-square}. Auxiliary bond integrations are part of
the inserted kernels and satisfy the same composition identities.
\end{proof}

\subsection{A coherent common-source bound and an absolutely summable time row}

Let $\Gamma_\gamma=\Gamma^{[\infty]}$ be exactly V7's full corrected
source Gram, evaluated at \eqref{f16v8:static-history}. Regard its time
blocks as operators on $\mathcal E$; there are $9B$ coordinates per block.
No spatial average or trace normalization is made in the next result.

\begin{theorem}[Complete-native coherent covariance on static backgrounds]
\label{f16v8:coherent-covariance}
There are constants $K_r\ge1$, $\gamma_r$ such that, for every $B=m^3$,
$m\ge4$, every $n\ge1$, every $\bar y$ in
\eqref{f16v8:static-history}, and
\begin{equation}
       \gamma\ge\gamma_r,\qquad
                  e^{K_rB}/\sqrt\gamma\le1/100,
 \label{f16v8:final-regime}
\end{equation}
one has, with $\tau=1/3$ and $A_{\gamma,B}=e^{K_rB}/\gamma$,
\begin{equation}
 \|\Gamma_{\gamma,ij}\|_{\mathcal E\to\mathcal E}
      \le A_{\gamma,B}\tau^{(|i-j|-2)_+},
              \qquad 0\le i,j\le n.
 \label{f16v8:covariance-blocks}
\end{equation}
In particular, for every common-source history vector
$\mathbf v=(v_0,\ldots,v_n)$,
\begin{equation}
 \boxed{
 \operatorname{Var}_{\mu_\gamma}
       \left(\sum_{i=0}^n R_i(v_i)\right)
           \le6A_{\gamma,B}\sum_{i=0}^n\|v_i\|^2.}
 \label{f16v8:coherent-bound}
\end{equation}
For every fixed $0\le\alpha<\log3$ there is also the weighted row bound
\begin{equation}
 \sup_i\sum_{j=0}^n e^{\alpha|i-j|}
                 \|\Gamma_{\gamma,ij}\|
  \le A_{\gamma,B}
       \left(1+2e^\alpha+
                    \frac{2e^{2\alpha}}{1-e^\alpha/3}\right).
 \label{f16v8:weighted-row}
\end{equation}
All duration dependence has been removed. Spatial dependence is the
explicit $e^{K_rB}$, not an unproved uniform constant.
\end{theorem}

\begin{proof}
Take unit spatial common variations $v,w$ at times $i,j$. If
$|i-j|\le1$, Cauchy--Schwarz, the square-insertion bound, and
\eqref{f16v8:boundary-overlap} give
\[
 |\operatorname{Cov}(R_i(v),R_j(w))|
                   \le\beta_\gamma^{-2}e^{c_rB}/\gamma.
\]
If $j-i\ge2$, their windows in \eqref{f16v8:source-window} have
disjoint bonds and gap $g=j-i-2$. They may share an endpoint internal
slice when $g=0$, which is allowed by
Proposition~\ref{f16v8:two-window}. Apply that proposition and the two
first-insertion norm bounds to obtain
\[
 |\operatorname{Cov}(R_i(v),R_j(w))|
       \le4\beta_\gamma^{-4}
               e^{2c_rB}\gamma^{-1}\tau^{j-i-2}.
\]
The same estimate holds in the reverse order by covariance symmetry.
Enlarge $K_r$ to absorb the finitely many constants and the powers of the
exponential overlap lower bound. Increasing it strengthens the entry
condition, so all the previous results remain valid. The estimates hold
for arbitrary unit spatial vectors, not just for one coordinate, and
therefore prove \eqref{f16v8:covariance-blocks} in block operator norm.

The scalar time majorant has row sum at most
\[
 1+4+2\sum_{d=3}^{\infty}\tau^{d-2}
                 =5+\frac{2\tau}{1-\tau}=6.
\]
Its symmetric Schur bound, or $2ab\le a^2+b^2$ applied to all time
pairs, proves \eqref{f16v8:coherent-bound}. Positivity of the full Gram
identifies its quadratic form with the displayed variance. Multiplying
the block estimate by $e^{\alpha|i-j|}$ and summing the corresponding
geometric series proves \eqref{f16v8:weighted-row}. There is no factor
$B(n+1)$ on the right of the coherent variance estimate.
\end{proof}

\begin{theorem}[Full history memory, original core correction, and time tails]
\label{f16v8:memory}
Under \eqref{f16v8:final-regime}, let $\mathsf M_\gamma$ and
$\mathsf C_{\gamma,R}$ be V7's actual nonadjacent-time Hessian matrices,
at the static background. Then
\begin{equation}
 \boxed{\quad \|\mathsf M_\gamma\|_{\rm op}
                         \le3A_{\gamma,B}.\quad}
 \label{f16v8:memory-op}
\end{equation}
More precisely, if $\mathsf M_\gamma^{>L}$ keeps only blocks with
$|i-j|>L$, for an integer $L\ge1$, then
\begin{equation}
 \|\mathsf M_\gamma^{>L}\|_{\rm op}
                         \le3A_{\gamma,B}(1/3)^{L-1}.
 \label{f16v8:memory-tail}
\end{equation}
For every radius $R$ obeying V4's derivative regime,
\begin{align}
 \|\mathsf C_{\gamma,R}\|_{\rm op}
        &\le3A_{\gamma,B}+C R^2/\gamma,\notag\\
 \|\mathsf C_{\gamma,R}^{>L}\|_{\rm op}
        &\le3A_{\gamma,B}(1/3)^{L-1}
                         +C(R^2/\gamma)(9/19)^{L-1}.
 \label{f16v8:correction-tail}
\end{align}
The constants are independent of $n$. The event, scalar factors, and
endpoint normalizers defining $p_{\rm good}$ are exactly V4's originals.
\end{theorem}

\begin{proof}
V7's full Haar-compensated identity, whose contact term has temporal range
one, gives
$\mathsf M_\gamma=-\operatorname{Off}_{>1}\Gamma_\gamma$ exactly.
At separation $d\ge2$ its block norm is at most
$A_{\gamma,B}\tau^{d-2}$. The scalar row sum of blocks with $d>L$ is
at most
\[
 2A_{\gamma,B}\sum_{d=L+1}^{\infty}\tau^{d-2}
          =\frac{2A_{\gamma,B}}{1-\tau}\tau^{L-1}
          =3A_{\gamma,B}\tau^{L-1}.
\]
This proves the first two claims by the same block Schur bound. Unlike
V7's trace-ideal extraction, this step uses actual separation decay.

V4's common-source core estimate acts on arbitrary spatial common
variations and gives a core block bound
$C(R^2/\gamma)(9/19)^{d-2}$ for $d\ge2$. Its time-tail row sum is
$C(R^2/\gamma)(9/19)^{L-1}$. Finally
\[
 \mathsf C_{\gamma,R}
   =\mathsf M_\gamma-
      \operatorname{Off}_{>1}\nabla_{\mathcal U}^2
                    [-\log\mathcal W_{\gamma,n}^{R}].
\]
Use the operator-norm triangle inequality, with the corresponding tail
mask when needed. The equality is the original full/core identity, not
a new covariance debit. External logarithmic factors depend on one time
only and have zero nonadjacent-time mixed derivatives; they have not
been dropped from either path integral.
\end{proof}

\subsection{What is gained, and the spatial limitation that remains}

For each fixed spatial $B$, the complete-native coherent covariance and
full memory bounds are $O_{r,B}(\gamma^{-1})$, with geometric time decay
and no history-length restriction. This improves the earlier fixed-lag
or zero-history bounds \emph{on the static-background subclass}; it does
not supersede V7's all-background, all-volume trace-density statement.
The present $\bar y$ may vary arbitrarily from one spatial bridge to
another inside its fixed rescaled window.

There is also a genuine growing-spatial-size consequence. With the same
$K_r$ as in the theorem, the regime
\begin{equation}
          B=m^3\le\frac{\log\gamma}{4K_r}
 \label{f16v8:mesoscopic}
\end{equation}
is admissible for sufficiently large $\gamma$, since the left side of
\eqref{f16v8:final-regime} is then at most $\gamma^{-1/4}$. It gives
$A_{\gamma,B}\le\gamma^{-3/4}$. Taking
$R=4\sqrt{\log\gamma}$ in V4's original core, all its fixed-$r$
conditions hold eventually and
\begin{equation}
 \|\Gamma_\gamma\|_{\rm op}
       +\|\mathsf M_\gamma\|_{\rm op}
       +\|\mathsf C_{\gamma,R}\|_{\rm op}
                   \le C_r\gamma^{-3/4},
 \label{f16v8:mesoscopic-conclusion}
\end{equation}
uniformly in \emph{all} $n\ge1$. The spatial side is therefore allowed
to grow like a constant times $(\log\gamma)^{1/3}$. The constants come
from conservative Gaussian integrals and derivative bounds; no useful
numerical entry coupling is claimed.

The source operator in \eqref{f16v8:coherent-bound} is the corrected
one from V7. The proof estimates every spatially coherent combination
within the specified volume, not a complement obtained by discarding
exceptional eigenvectors. All excited fast states and all compact port
tails remain in the transfer and its inserted kernels. In particular
there is no interchange of an uncontrolled $n\to\infty$ error with the
weak-field limit: the spectral split and boundary denominator hold at
each admitted finite $\gamma,B$ for every $n$.

Two restrictions are material. First, this proof uses one time-independent
kernel $K_{\gamma,\bar y}$ and its one Perron projection. An arbitrary
bounded history instead produces different consecutive kernels. Their
powers cannot be replaced by the spectral powers used here. This
appendix therefore gives no arbitrary-history temporal mixing theorem.
Second, the complete-chart and endpoint comparisons have paid costs
$e^{c_rB}$. They do not allow $B\to\infty$ at fixed $\gamma$, nor may
the small-volume covariance inequalities be tensorized across interacting
spatial regions. The fixed bridges are also not a gauge-invariant coverage
event for the integrated physical bridge law.

The next useful upstream target is a \emph{boundary-conditioned spatial
version} of the two-insertion estimate, preserving the exact normalization
while bounding how a spatial collar changes the source and endpoint
messages. The global internal pin supplies the present finite-volume
input, but arbitrary exterior fast fields can change the conditional
landscape, so the displayed torus estimate is not such a boundary-uniform
theorem. Separately, varying bridge histories require a normalized
inhomogeneous transfer estimate; a gap for each static operator alone
would not prove that estimate. These are concrete propagation problems,
not another request for scalar pressure or an assumption of native
functional inequalities. The full-transfer top comparison, physical time
calibration, infrared coercivity of the retained bridge dynamics, and an
interacting continuum construction remain outside the proved claims.

\clearpage
\subsection{Diagnostics, preservation, and scope}

The accompanying diagnostic is
\begin{center}\small
\path{verify_ym_f16_v8_static_transfer.py}.
\end{center}
It checks exact compact word reversal, noncommuting Gaussian square
completion, the scalar Mehler precision and ratio identities, finite
history insertion and centering formulas with non-Perron endpoints,
time-window and tail constants, and the logarithmic spatial entry regime.
All 609 exact assertions passed. Its finite positive-transfer examples
are tests of the transfer algebra, not numerical Wilson simulations. The
report separates the assertion families. The script has SHA--256
\begin{center}\scriptsize\ttfamily
92AA1F6C83ADE826AA35BA145CC6FDEACAB849C36626D2A0A7FB56B267009598.
\end{center} These tests do not evaluate the
suprema of native word derivatives, certify the infinite-dimensional
operator estimates with a proof assistant, or establish a physical mass
gap. The analytic proofs, with the restrictions above, are essential.

The supplied V4--V7 diagnostics were rerun unchanged and passed 9,467,
89,339, 906 and 311 assertions, respectively. Their reports are included.
Unavailable archive diagnostics are not represented as rerun. The V7
predecessor has 231,802 bytes and SHA--256
\begin{center}\scriptsize\ttfamily
36044181A0F2AB65DF5ABC97B95B0652891AAA923D08508DBF903277A3115BC3.
\end{center}
It is preserved byte for byte except for the cumulative
title version and insertion of this appendix immediately before its final
\verb|\end{document}|. Removing the exact appendix and restoring that title
recovers the supplied V7. The package includes source hashes, preservation
checks, and build/visual-review reports. This is not the scheduled V10
companion checkpoint.

\begin{firewall}
V8 proves geometric temporal decay and a small \emph{operator norm} for
V7's complete-native corrected-source covariance and nonadjacent-time
memory on time-independent bounded bridge backgrounds. The original
full/core correction has the same summable time-tail control. These
statements hold for arbitrary duration but with an explicit exponential
spatial cost, admitting $B=O_r(\log\gamma)$ in the stated weak-field
regime. No spatial-volume-uniform result at fixed coupling, arbitrary
inhomogeneous-history mixing, boundary-uniform spatial gluing,
gauge-invariant bridge coverage, full physical transfer gap, interacting
continuum Yang--Mills theory, or positive physical mass gap is established
by this appendix.
\end{firewall}
% END F16 V8 APPEND
% BEGIN F16 V9 APPEND -- 2026-09-17
% Insert immediately before the final \end{document} in cumulative V8.
% The predecessor body is unchanged; the cumulative title becomes V9.

\clearpage
\section{V9: exact history centering and inhomogeneous fast-transfer products}
\label{f16v9:context}

V8 controls coherent common-source memory for arbitrarily long histories,
but only when the bridge background is constant in time. A separate
spectral gap for each frozen operator does not control their product.
Here we remove that restriction without assuming that consecutive bridge
values are close. The upstream step is to center the \emph{whole Gaussian
history}, including its dressed ends, before comparing any native step.
Explicit linear multipliers then telescope between slices. Every resulting
Gaussian step is the same centered operator, while every resulting native
step is uniformly close to it. A direct nonautonomous product estimate
pays the exact product normalizer, rather than replacing a long native
product by a Gaussian power.

The spatial cost remains exponential in $B$. This appendix does not claim
a spatial-volume-uniform estimate at fixed coupling. What changes is the
background class: all bridge histories in the specified rescaled bounded
window are now admitted, with no bound on their duration or total temporal
variation. The covariance and time-tail estimates concern the same exact
Haar-corrected sources and the same original full/core probability as before.

\subsection{Dependency audit and the distinction from a global cone argument}

The supplied V8 appendix was read in full. The available archive inventories
were searched and the relevant interfaces inspected; this is not independent
certification of every earlier proof. We reuse V3's history precision and
Gaussian inverse locality, V5's exact unmoving coordinates and global pins,
V7's Haar source identity, and V8's complete-chart kernel and insertion
method. These are not new V9 results.

Two archived cautions matter here. F15 V57 explicitly restricts its repeated
transfer argument to a homogeneous baseline. F15 V91,
\texttt{f15v91:projective}, proves that endpoint multiplication does not
improve the global Hilbert-projective diameter of the full positive cone.
The present proof uses neither such a diameter estimate nor a claim that
multiplication alone improves it. It verifies an adapted Hilbert-space
slope ball around a specified Gaussian vector, with its exact perturbation
size and endpoint entry cost. The monograph's
\texttt{sec:YM-Perron-vacuum-metric} distinction between an actual slab law
and a stationary Perron law is retained.

For background on adapted cone methods, see H.~H.~Rugh,
\href{https://annals.math.princeton.edu/2010/171-3/p07}
{\emph{Cones and gauges in complex spaces: Spectral gaps and complex
Perron--Frobenius theory}}, Annals of Mathematics 171 (2010), 1707--1752.
The particular real-Hilbert-space product lemma needed here is proved
below, including non-self-adjoint factors and non-Perron endpoints. No
external theorem about the present native history is imported.

\subsection{The full history mean and its telescoping linear terms}

Throughout, fix
\[
 B=m^3,\quad m\ge4,\quad n\ge1,\qquad
 \mathbf y=(y_0,\ldots,y_n),\qquad \max_{i,e}|y_{i,e}|\le r<\infty.
\]
Constants with subscript $r$ may depend on this fixed window and the cube
geometry, never on $n$ or on the individual history. All path matrices act
with $I_3$ in colour. Let
\[
 N=\mathsf M_B^{-1},\quad S=\mathsf H_B+A^{\mathsf T}A,\quad
 \Lambda=A^{\mathsf T}D,\quad
 P_f=\tfrac12(C_B^{-1}+A^{\mathsf T}A).
\]
The symbol $P_f$ is an amplitude precision, not V7's port lift. Let $W_n$
be exactly V3's positive internal history precision and put
\begin{equation}
 \boldsymbol\mu=(\mu_0,\ldots,\mu_n)
       =-W_n^{-1}(\Lambda y_0,\ldots,\Lambda y_n).
 \label{f16v9:mean}
\end{equation}
Thus $\boldsymbol\mu$ is the internal mean of the complete Gaussian
conditional with its original dressed ends. It is not the list of minima
of the separate frozen potentials. In V5's unmoving temporal coordinates
$a_t$, its port mean is
\begin{equation}
 \nu_t=\mathsf F_B(\mu_t-\mu_{t+1})
               -L_h^{-1}E_{\rm br}^{\mathsf T}(y_t-y_{t+1}).
 \label{f16v9:port-mean}
\end{equation}
Use translated complete charts $u_i=x_i-\mu_i$, $b_t=a_t-\nu_t$.
These are translations in integration coordinates, not changes to a compact
configuration or to the definition of a source derivative.

Define the left and right linear terms of bond $t$ by
\begin{align}
 l_t&=N(\mu_t-\mu_{t+1})+\tfrac12(S\mu_t+\Lambda y_t),\notag\\
 r_t^{\rm lin}&=N(\mu_{t+1}-\mu_t)
                      +\tfrac12(S\mu_{t+1}+\Lambda y_{t+1}),
 \label{f16v9:bond-linear}
\end{align}
and define endpoint terms
\[
 d_L^{\rm lin}=P_f\mu_0+\tfrac12\Lambda y_0,\qquad
 d_R^{\rm lin}=P_f\mu_n+\tfrac12\Lambda y_n.
\]
The superscript distinguishes $r_t^{\rm lin}$ from a corrected score.

\begin{proposition}[Uniform mean bounds and exact telescoping]
\label{f16v9:centering}
There is $C_r<\infty$ independent of $B,n$ such that every cube block
of $\mu_i,\nu_t,l_t,r_t^{\rm lin},d_L^{\rm lin},d_R^{\rm lin}$ has norm
at most $C_r$. Moreover
\begin{equation}
 d_L^{\rm lin}+l_0=0,\qquad
 r_{i-1}^{\rm lin}+l_i=0\ (0<i<n),\qquad
 r_{n-1}^{\rm lin}+d_R^{\rm lin}=0.
 \label{f16v9:telescoping}
\end{equation}
No smallness of $y_t-y_{t+1}$ is assumed.
\end{proposition}

\begin{proof}
V3 gives a spatially block-local whitening of $W_n$ with spectrum in
$[4,15]$. Its inverse has block decay at most $C(11/15)^d$ on the
cube--time graph. A radius-$d$ ball has at most $C(d+1)^4$ labels,
uniformly on the spatial torus times the finite time interval. The inverse
therefore has a bounded sum of block norms in every row. The source
$\Lambda y_i$ has bounded norm per cube by the finite incidence and the
bound $r$. Undoing the fixed block whitening proves the assertion for
$\mu$. Likewise $L_h$ has a fixed positive lower spectral bound and
finite spatial range. Its inverse has a uniformly summable spatial block
row, by the Neumann expansion already used in V7. Equation
\eqref{f16v9:port-mean} gives the bound for $\nu$. The remaining matrices
in \eqref{f16v9:bond-linear} and the endpoint terms have fixed finite
range and bounded block rows, proving their bounds.

For an interior slice, the sum of the two displayed bond terms is
\[
 (S+2N)\mu_i-N\mu_{i-1}-N\mu_{i+1}+\Lambda y_i=0
\]
by \eqref{f16v9:mean}. At the first endpoint the coefficient of $\mu_0$
is $P_f+N+S/2=E$, exactly V3's endpoint diagonal; the source is
$\Lambda y_0$. The same argument applies at the last endpoint. This
proves all three identities, also for $n=1$ when there are no interior
slices. The port mean follows from V2's kinetic square: in moving
coordinates the minimizing port is
$-L_h^{-1}E_{\rm br}^{\mathsf T}(y_t-y_{t+1})$; V5's determinant-one
linear change adds $\mathsf F_B(\mu_t-\mu_{t+1})$.
\end{proof}

\subsection{Balanced kernels with one common Gaussian reference}

For two possibly different bridges $y,y'$, let $K_\gamma(y,y';x,x')$
be V8's augmented-step integral with these values at the two ends. More
explicitly, use its actual words with $U_e(x;y)$ and $U_e(x';y')$,
magnetic halves $V_{\gamma,y}(x)/2,V_{\gamma,y'}(x')/2$, and scalar
\begin{equation}
 \rho_\gamma(y,y')=
   (c_{H,\gamma}/z_\gamma)^{24B}
             [j_{\rm br}(y)j_{\rm br}(y')]^{1/2}.
 \label{f16v9:step-scalar}
\end{equation}
All internal and port charts are complete. This kernel is generally not
self-adjoint; its transpose is $K_\gamma(y',y)$. Its kernel is
nonnegative. No positive-semidefinite operator claim is needed for unequal
bridges. Exact composition gives
\begin{equation}
 \mathcal W_{\gamma,n}^{\rm full}(\mathbf y)
   =\langle a_{\gamma,y_0},
       K_\gamma(y_0,y_1)\cdots K_\gamma(y_{n-1},y_n)
                                      a_{\gamma,y_n}\rangle.
 \label{f16v9:inhomogeneous-path}
\end{equation}
The normalized endpoints are V8's literal vectors at their respective
bridge values. This follows by the same Haar powers as V8; it is not a
new replacement law.

The integrated Gaussian kernel is
\begin{equation}
 K_0(y,y';x,x')=
 \varphi_G(y-y')\,
 e^{-V_{0,y}(x)/4}\varphi_{\mathsf M_B}(x-x')e^{-V_{0,y'}(x')/4},
 \qquad
 V_{0,y}(x)=x^{\mathsf T}Sx+2x^{\mathsf T}\Lambda y+|Dy|^2.
 \label{f16v9:Gaussian-step}
\end{equation}
Set $K_*=K_0(0,0)$, $a_*=a_{0,0}$, and let
$\lambda_*,\phi_*$ be its exact top value and normalized positive vector.
In this appendix $a_*$ denotes this vector; the pin constant is denoted
$c_{\rm pin}=(6\pi^2)^{-1}$. V8's Mehler calculation gives
\begin{equation}
 S_*:=K_*/\lambda_*=P_*+Q_*,\quad
 P_*=|\phi_*\rangle\langle\phi_*|,\quad
 0\le Q_*,\quad P_*Q_*=0,\quad \|Q_*\|\le3-2\sqrt2<1/5.
 \label{f16v9:reference-split}
\end{equation}
No second oscillator diagonalization is being claimed.

Introduce positive scalars
\begin{equation}
 c_t=\frac{K_0(y_t,y_{t+1};\mu_t,\mu_{t+1})}{K_*(0,0)},\qquad
 d_L=\frac{a_{0,y_0}(\mu_0)}{a_*(0)},\qquad
 d_R=\frac{a_{0,y_n}(\mu_n)}{a_*(0)}.
 \label{f16v9:balancing-scalars}
\end{equation}
Define the balanced native kernels and endpoints on the same Euclidean
Hilbert space by
\begin{align}
 \widetilde K_t(u,v)
   &=c_t^{-1}e^{\langle l_t,u\rangle+\langle r_t^{\rm lin},v\rangle}
       K_\gamma(y_t,y_{t+1};u+\mu_t,v+\mu_{t+1}),\notag\\
 \widetilde a_L(u)
   &=d_L^{-1}e^{\langle d_L^{\rm lin},u\rangle}
                                      a_{\gamma,y_0}(u+\mu_0),\notag\\
 \widetilde a_R(u)
   &=d_R^{-1}e^{\langle d_R^{\rm lin},u\rangle}
                                      a_{\gamma,y_n}(u+\mu_n).
 \label{f16v9:balanced-native}
\end{align}
Their supports are the translated complete charts. They are not normalized
probability densities or unit vectors by definition.

\begin{theorem}[Exact balanced path and a uniform one-step comparison]
\label{f16v9:balanced-comparison}
The Gaussian versions of \eqref{f16v9:balanced-native} are exactly $K_*$
at every bond and $a_*$ at both ends. At finite regulator one has the
\emph{exact} identity
\begin{equation}
 \mathcal W_{\gamma,n}^{\rm full}(\mathbf y)
  =d_Ld_R\left(\prod_{t=0}^{n-1}c_t\right)\lambda_*^n
     \langle\widetilde a_L,T_0\cdots T_{n-1}\widetilde a_R\rangle,
 \qquad T_t=\widetilde K_t/\lambda_*.
 \label{f16v9:balanced-path}
\end{equation}
There are $C_r,c_r,\gamma_r$ independent of $B,n$ such that for
$\gamma\ge\gamma_r$,
\begin{align}
 \sup_t\|T_t-S_*\|_{\rm HS}
    +\|\widetilde a_L-a_*\|_2+\|\widetilde a_R-a_*\|_2
       &\le e^{C_rB}/\sqrt\gamma,\notag\\
 \langle a_*,\phi_*\rangle&\ge e^{-c_rB},\notag\\
 e^{-C_rB}\le c_t,d_L,d_R&\le e^{C_rB}.
 \label{f16v9:balanced-bounds}
\end{align}
The matrix element in \eqref{f16v9:balanced-path} is retained exactly; it
is not replaced by $1$ or by $\langle a_*,S_*^na_*\rangle$.
\end{theorem}

\begin{proof}
Expanding the quadratic kernel \eqref{f16v9:Gaussian-step} after the two
translations gives precisely the linear coefficients
\eqref{f16v9:bond-linear}. The remaining constant is fixed by evaluation
at $u=v=0$, which is the definition of $c_t$. Likewise the amplitude
precision of $a_{0,y}$ is $P_f$ and its fast linear coefficient is
$\Lambda y/2$. Its translated linear coefficient is $d_{L,R}^{\rm lin}$
and its constant is fixed by $d_{L,R}$. This proves the exact Gaussian
claims. Substitution into the native path integral cancels every linear
multiplier by \eqref{f16v9:telescoping}. The translations have Jacobian
one. This proves \eqref{f16v9:balanced-path} without a limit.

The uniform local mean bounds imply that the squared Euclidean norms of
all means and of all balancing coefficient vectors at one bond are at
most $C_rB$. The Gaussian exponent in $c_t$ is
\[
 -\tfrac14V_{0,y_t}(\mu_t)-\tfrac14V_{0,y_{t+1}}(\mu_{t+1})
 -\tfrac12|\mu_t-\mu_{t+1}|_N^2
 -\tfrac12|y_t-y_{t+1}|_{G^{-1}}^2.
\]
It lies in $[-C_rB,0]$. The Gaussian endpoint determinant, mean and
precision bounds from V8 similarly bound $\log d_L,\log d_R$ in
absolute value by $C_rB$. The overlap assertion is exactly V8's
zero-background overlap estimate.

We spell out why linear balancing is legitimate for the \emph{complete}
native integrals. Before integrating ports, translate them by $\nu_t$ as
in \eqref{f16v9:port-mean}. The port linear term in the Gaussian exponent
then vanishes, because that mean solves its bond stationarity equation.
The two remaining internal linear terms are $l_t,r_t^{\rm lin}$.
Thus the balanced Gaussian augmented exponent is exactly the zero-bridge
augmented exponent, plus the constant already removed in $c_t$.
Both original augmented actions have the global pin
\[
 c_{\rm pin}(|x|^2+|x'|^2+|a|^2),
\]
by the internal magnetic ends and independent rooted tree, as in V8.
After translation and the balancing linear multipliers, both integrands
are bounded by
\begin{equation}
 e^{C_rB}\exp[-c(|u|^2+|v|^2+|b|^2)]
 \label{f16v9:balanced-pin}
\end{equation}
for a fixed $c>0$. This follows twice from
$|z+m|^2\ge|z|^2/2-|m|^2$ and then from Young's inequality for the
balancing linear terms. All Haar factors are at most one and all scalar
powers cost at most $e^{C_rB}$.

The native action minus its Gaussian quadratic action obeys V8's global
word remainder; in the translated variables it is bounded by
$C_r\gamma^{-1/2}(B^{3/2}+|u|^3+|v|^3+|b|^3)$.
The nonconstant Haar-factor error is at most
$C_r\gamma^{-1}(B+|u|^2+|v|^2+|b|^2)$, and the scalar logarithmic
error is at most $C_rB/\gamma$. Using
$|e^{-a}-e^{-b}|\le|a-b|e^{-\min(a,b)}$ and
\eqref{f16v9:balanced-pin} gives a polynomial times a Gaussian majorant
for the difference. A translated principal-chart complement has a group
coordinate of norm at least $\pi\sqrt\gamma-C_r$, hence at least
$\pi\sqrt\gamma/2$ once $\gamma\ge\gamma_r$.
Its Gaussian tail costs $e^{C_rB-c\gamma}$, not an omitted remainder.
Integrate all ports and then take the kernel $L^2(du\,dv)$ norm.
Fixed-degree Gaussian moments in dimension $51B$ cost an exponential
in $B$. The resulting kernel error is $e^{C_rB}/\sqrt\gamma$;
$\lambda_*^{-1}\le e^{C_rB}$ pays the displayed normalization.

For the endpoints use the unnormalized native and Gaussian functions
\[
 q_\gamma=1_{\mathcal P_x}\Omega e^{-\gamma S_\times/2},\qquad
 q_0=\Omega e^{-|Ax+Dy|^2/4}.
\]
Their normalizing norms are at least
$e^{-C_rB}$: integrate their squares over a product of fixed small
one-cube balls, using bounded local inputs and $\gamma\ge\gamma_r$.
After the translation and the linear exponential multiplication,
the factor $\Omega$ still supplies a Gaussian majorant of the form
\eqref{f16v9:balanced-pin}. The same polynomial remainder proof bounds
the weighted unnormalized difference by $e^{C_rB}/\sqrt\gamma$.
The unweighted difference bounds the difference of the normalizing norms.
Divide by the just-proved lower bounds and by $d_{L,R}\ge e^{-C_rB}$.
This proves the two endpoint estimates. In particular we have not applied
an unbounded multiplication operator to an unweighted $L^2$ estimate.
\end{proof}

The centering depends on the complete finite bridge history and its two
ends. It is a comparison coordinate system, not a causal coarse state or
a new physical vacuum. All bridge derivatives defining the observables
are taken in the original coordinates \emph{before} this transformation.
Differentiating the balanced formula while omitting derivatives of its
means or multipliers would be incorrect; no such step is used below.

\subsection{A product lemma with an exact nonautonomous normalizer}

Here is the analytic step replacing the use of one static spectral power.
It is stated on a real Hilbert space, so it also applies to the
infinite-dimensional internal space. Write $P=|\phi\rangle\langle\phi|$.

\begin{theorem}[Uniform rank-one decomposition of near-reference products]
\label{f16v9:product}
Suppose $S=P+Q$, $\|\phi\|=1$, $Q=Q^*\ge0$, $PQ=0$, and
$\|Q\|\le1/5$. Let $(T_j)$ be any finite sequence of real bounded
operators satisfying $\|T_j-S\|\le\delta\le1/100$.
Neither self-adjointness nor commutation of the $T_j$ is assumed.
For an ordered interval product $T_I$ of length $k$, set
$p_I=\langle\phi,T_I\phi\rangle$. Then $p_I>0$ and
\begin{equation}
 \frac{T_I}{p_I}=|\phi+x_I\rangle\langle\phi+y_I|+E_I,
 \quad x_I,y_I\perp\phi,\quad
 \|x_I\|,\|y_I\|\le2\delta,\quad
 PE_I=E_IP=0,\quad \|E_I\|\le3^{-k}.
 \label{f16v9:product-decomposition}
\end{equation}
Consequently $\|T_I\|\le2p_I$. For adjoining intervals $I,J$ in that
product order,
\begin{equation}
 \frac{p_{IJ}}{p_Ip_J}=1+\langle y_I,x_J\rangle
            \in[1-4\delta^2,1+4\delta^2].
 \label{f16v9:quasi-multiplication}
\end{equation}
An empty interval has $T_I=I$, $p_I=1$, $x_I=y_I=0$, $E_I=I-P$;
the same assertions with $k=0$ hold. If $1\le k\le2$, then $p_I\ge1/2$.
\end{theorem}

\begin{proof}
Decompose a single factor relative to $\mathbb R\phi\oplus\phi^\perp$:
\[
 T_j=\begin{pmatrix}a_j&c_j^*\\ b_j&D_j\end{pmatrix},\qquad
 |a_j-1|,\|b_j\|,\|c_j\|\le\delta,\quad \|D_j\|\le1/5+\delta.
\]
Its projective slope map on the unit ball of $\phi^\perp$ is
\begin{equation}
 F_j(z)=\frac{b_j+D_jz}{a_j+\langle c_j,z\rangle}.
 \label{f16v9:slope}
\end{equation}
The denominator is at least $1-2\delta\ge49/50$. On this ball
\[
 \|F_j(z)\|\le\frac{1/5+2\delta}{1-2\delta}\le\frac{11}{49},
 \qquad
 \|DF_j(z)\|\le
 \frac{1/5+\delta+\delta(11/49)}{1-2\delta}<\frac13.
\]
Thus the unit ball is invariant and these maps contract its norm distance.
The smaller ball of radius $2\delta$ is also invariant, since
\[
 \frac{\delta+(1/5+\delta)2\delta}{1-\delta-2\delta^2}
                                                        \le2\delta.
\]
Exactly the same estimates hold for every adjoint factor.

Iterate from slope zero in the order in which the product acts on $\phi$.
All scalar coefficients are positive, proving $p_I>0$ and
$x_I=P^\perp T_I\phi/p_I$ with $\|x_I\|\le2\delta$.
The adjoint product gives $y_I=P^\perp T_I^*\phi/p_I$ with the same
bound. Differentiating the composed slope map at zero gives
\[
 DF_I(0)=P^\perp T_IP^\perp/p_I-x_Iy_I^*.
\]
The chain rule bounds its norm by $3^{-k}$. Extending this operator by
zero on $\phi$ gives $E_I$ and proves the exact block identity
\eqref{f16v9:product-decomposition}. For $k\ge1$ its norm is at most
$1+4\delta^2+1/3<2$; the empty case is immediate. Multiplying two such
block matrices in their top-left corner gives
\eqref{f16v9:quasi-multiplication} exactly. Finally at each slope step
the scalar factor is at least $1-\delta-2\delta^2\ge1-2\delta$;
for $k\le2$ their product is at least $(49/50)^2>1/2$.
\end{proof}

The scalar $p_I$ can grow or decay with $k$. It is \emph{not} replaced by
one. The theorem bounds the normalized remainder, rather than the generally
false estimate $\|T_I-S^k\|=O(\delta)$ independent of $k$. Nor is the
slope ball the full positive-function cone from f15 V91.

\begin{proposition}[Original-type endpoints and a denominator for all lengths]
\label{f16v9:endpoints}
In the preceding theorem suppose $\|a\|=1$,
$\langle a,\phi\rangle\ge\beta$ with $0<\beta\le1$, and
$\|a_L-a\|,\|a_R-a\|\le\eta$. Assume
\begin{equation}
 \delta\le10^{-6}\beta^4,\qquad \eta\le10^{-6}\beta^2.
 \label{f16v9:endpoint-entry}
\end{equation}
For every product of length $n\ge1$,
\begin{equation}
 \frac{\beta^2}{16}\,p_{[0,n)}
     \le\langle a_L,T_0\cdots T_{n-1}a_R\rangle
     \le8p_{[0,n)}.
 \label{f16v9:denominator}
\end{equation}
The factors need not have a common eigenvector. No stationarity of the
path ends is required.
\end{proposition}

\begin{proof}
Put $N_\beta=\lceil\log(128/\beta^2)/\log3\rceil$.
For $n\ge N_\beta$, the rank-one term in
\eqref{f16v9:product-decomposition}, paired with the endpoints, is at
least $(\beta-\eta-4\delta)^2\ge\beta^2/4$. Its remainder has absolute
value at most $4\,3^{-n}\le\beta^2/32$, since $\|a_L\|,\|a_R\|\le2$.
This proves the stated lower bound for these lengths.

For $1\le n<N_\beta$, use $N_\beta\le10/\beta^2$, which follows
from $\log x\le x-1$ for $x\ge1$. The telescoping product estimate gives
\[
 \|T_0\cdots T_{n-1}-S^n\|
   \le n\delta(1+\delta)^{n-1}\le2n\delta.
\]
Here $n\delta\le10^{-5}\beta^2$. The actual endpoint matrix element
therefore differs from $\langle a,S^na\rangle$ by at most
$6\eta+2n\delta<\beta^2/2$. The latter is at least $\beta^2$ because
$S^n=P+Q^n$ and $Q^n\ge0$. Also
$p_{[0,n)}\le(1+\delta)^n\le2$. This proves a lower bound
$\beta^2p_{[0,n)}/4$, stronger than required. Finally
$\|T_{[0,n)}\|\le2p_{[0,n)}$ and the endpoint norms give the upper
bound for all lengths. This proof treats short histories separately rather
than imposing an unproved Perron denominator on their endpoints.
\end{proof}

\subsection{Two source windows in the nonautonomous product}

The following consequence keeps the exact finite-path means. It applies
to positive integral kernels with auxiliary bond variables integrated into
insertions. It does not require their operators to be self-adjoint.

\begin{proposition}[Centered two-window estimate for the actual product law]
\label{f16v9:two-window}
Assume the preceding two results and that the kernels and endpoints define
a positive path law. Let $F,G$ be real observables on windows of one or
two bonds, with disjoint bonds and $g\ge0$ unmarked bonds between them.
Let $A_F,A_G$ be the kernels replacing their respective window products,
with the observables inserted; let $B_F$ insert $F^2$. There is an absolute
constant $C_0$ (for example $2^{28}$) such that
\begin{equation}
 |\operatorname{Cov}(F,G)|
       \le C_0\beta^{-4}\|A_F\|\|A_G\|3^{-g},\qquad
 \mathbb E F^2\le C_0\beta^{-2}\|B_F\|.
 \label{f16v9:nonautonomous-mixing}
\end{equation}
Both expectations use the \emph{actual} normalizer in
\eqref{f16v9:denominator}.
\end{proposition}

\begin{proof}
Split the unmarked full product into left segment $L$, first window $U$,
gap $G_0$, second window $V$, and right segment $R$, in that order.
Write $p_L,p_U,p_{G_0},p_V,p_R$ for their exact scalar coefficients.
At most four applications of \eqref{f16v9:quasi-multiplication}, and
$p_U,p_V\ge1/2$, give
\[
 Z':=\frac{\langle a_L,T_{[0,n)}a_R\rangle}{p_Lp_{G_0}p_R}
                                                      \ge\beta^2/128.
\]
The convention for empty products handles endpoints and $g=0$.
Set
\[
 \begin{gathered}
 u=T_L^*a_L/p_L,\quad v=T_Ra_R/p_R,\quad H=T_{G_0}/p_{G_0},\\
 u_0=T_U^*u,\quad u_F=A_F^*u,\quad
 v_0=T_Vv,\quad v_G=A_Gv.
 \end{gathered}
\]
Then $\|u\|,\|v\|\le4$, $\|u_0\|,\|v_0\|\le5$,
$\|u_F\|\le4\|A_F\|$, $\|v_G\|\le4\|A_G\|$, and
$Z'=\langle u_0,Hv_0\rangle$. Exact composition gives
\begin{equation}
 \operatorname{Cov}(F,G)=
 \frac{\langle u_F,Hv_G\rangle\langle u_0,Hv_0\rangle
       -\langle u_F,Hv_0\rangle\langle u_0,Hv_G\rangle}{(Z')^2}.
 \label{f16v9:minor}
\end{equation}
Replace $H$ by the rank-one part of
\eqref{f16v9:product-decomposition}. The numerator vanishes exactly.
The norms of both $H$ and that rank-one part are at most two, and their
difference has norm at most $3^{-g}$. Expanding the difference of the two
products bounds the numerator by
$3200\|A_F\|\|A_G\|3^{-g}$.
The displayed denominator lower bound makes $2^{28}$ a safe constant.
For a single window, the corresponding three-segment scalar comparison
gives $Z\ge\beta^2p_Lp_R/64$, while its square-insertion numerator has
absolute value at most $16p_Lp_R\|B_F\|$. This proves the second bound.
Auxiliary port integrations belong to the window kernels, so the same
identities apply to their observables. The centering in
\eqref{f16v9:minor} includes the actual two finite-history means.
\end{proof}

\subsection{Full-native remainder insertions after history centering}

Use exactly V7's full lift $\mathsf P$ and corrected common-source
residuals $r^{[\infty]}_\alpha$, evaluated at the original bridge history.
For $v\in\mathcal E$ at time $i$, set
\[
 R_i(v)=\sum_{b,\mu,k}v_{b,\mu,k}r^{[\infty]}_{i,b,\mu,k}.
\]
Its bond window is
$[\max(0,i-1),\min(n,i+1)]$, of length one or two. First form this
observable using original bridge derivatives and compact Haar insertions;
then express its arguments in the centered integration coordinates.
All centering means are held fixed during this evaluation.

Let $A_i(v),B_i(v)$ be its first and square insertion operators using the
balanced augmented kernels divided by $\lambda_*$ per bond. Thus they
replace the corresponding products of the $T_t$, with \emph{all} their
scalars $c_t$ and endpoint linear multipliers included.

\begin{proposition}[Uniform complete-chart insertion bounds at arbitrary histories]
\label{f16v9:insertions}
For sufficiently large $\gamma$ there is $C_r$ independent of $n,B,i$
and of the admitted history such that
\begin{equation}
 \|A_i(v)\|\le e^{C_rB}\gamma^{-1/2}\|v\|,\qquad
 \|B_i(v)\|\le e^{C_rB}\gamma^{-1}\|v\|^2.
 \label{f16v9:insertion-bounds}
\end{equation}
These are estimates for complete native kernels, not Gaussian expectations
used instead of native expectations.
\end{proposition}

\begin{proof}
V7's exact Gaussian cancellation holds for every bridge history. On each
adjacent bond, the lift cancels the fast affine kinetic score, leaving its
explicit deterministic contribution to $g_\alpha(\mathbf y)$; at a
common magnetic mark the fast affine coefficient is zero. Subtracting the
same $g_\alpha$ therefore leaves only primitive nonlinear remainders and
the bounded $L_h^{-1}E_{\rm br}^{\mathsf T}$ synthesis of port remainders,
also when $y_t-y_{t+1}\ne0$. The inserted exponential-word Taylor bound
from V7 is global on the complete charts. With $w$ the at most $87B$
original real fast coordinates in the window, it gives
\[
 |R_i(v)|\le C_r\gamma^{-1/2}(\sqrt B+|w|^2)\|v\|.
\]
Only the spatial lift's $\ell^2$ operator norm is used, not a sum of
absolute coefficients. After translation the window's mean vector has
squared norm at most $C_rB$. Thus the same expression is bounded by
$C_r\gamma^{-1/2}(B+|w_{\rm centered}|^2)\|v\|$.

The product of at most two balanced augmented kernels has the Gaussian
majorant \eqref{f16v9:balanced-pin}, with its constant enlarged. Multiplying
by this polynomial or its square, integrating the full translated port
and intermediate internal charts, and taking the Hilbert--Schmidt norm in
the two external variables proves \eqref{f16v9:insertion-bounds}.
The fixed-degree moments cost $e^{C_rB}$ as in V8. Division by one or two
powers of $\lambda_*\ge e^{-C_rB}$ is included. No whole-history
denominator or duration factor enters this local integration. The native
port derivatives are still Haar derivatives; translating their displayed
coordinate arguments has not replaced them by flat derivatives.
\end{proof}

\subsection{Arbitrary-history coherent memory and the unchanged core correction}

\begin{theorem}[Inhomogeneous common-source covariance, uniformly in duration]
\label{f16v9:coherent}
Fix $r<\infty$. There are finite $K_r\ge1,\gamma_r$ such that, whenever
\begin{equation}
 \gamma\ge\gamma_r,\qquad e^{K_rB}/\sqrt\gamma\le1/100,
 \qquad \max_{i,e}|y_{i,e}|\le r,
 \label{f16v9:entry}
\end{equation}
V7's exact full-native residual covariance, at \emph{this arbitrary
history}, satisfies
\begin{equation}
 \boxed{\quad
 \|\Gamma_{\gamma,ij}\|_{\mathcal E\to\mathcal E}
     \le A_{\gamma,B}(1/3)^{(|i-j|-2)_+},\qquad
 A_{\gamma,B}=e^{K_rB}/\gamma.\quad}
 \label{f16v9:covariance}
\end{equation}
Consequently, for every $n\ge1$ and every common-source sequence $v_i$,
\begin{equation}
 \boxed{\quad
 \operatorname{Var}_{\mu_{\gamma,\mathbf y}}
       \left(\sum_{i=0}^n R_i(v_i)\right)
      \le6A_{\gamma,B}\sum_{i=0}^n\|v_i\|^2.\quad}
 \label{f16v9:coherent-bound}
\end{equation}
For $0\le\alpha<\log3$, the weighted time-row bound is
\begin{equation}
 \sup_i\sum_j e^{\alpha|i-j|}\|\Gamma_{\gamma,ij}\|
 \le A_{\gamma,B}
       \left(1+2e^\alpha+\frac{2e^{2\alpha}}{1-e^\alpha/3}\right).
 \label{f16v9:weighted}
\end{equation}
No constancy, slow variation, or total-variation budget on $\mathbf y$
is assumed. The exponential spatial cost is part of the assertion.
\end{theorem}

\begin{proof}
Apply Theorem~\ref{f16v9:balanced-comparison} with the common reference
$S_*$ and $a_*$. Put $\beta=e^{-c_rB}$ using its overlap lower bound.
By enlarging $K_r$, the entry condition makes the uniform step error
at most $10^{-6}\beta^4$ and the endpoint errors at most
$10^{-6}\beta^2$. This is possible because every error is
$e^{C_rB}/\sqrt\gamma$, whereas $\beta$ is a fixed exponential in
$B$. The constants $10^{-6}$ are absorbed by enlarging a geometric
exponent, not by introducing a duration-dependent threshold.
The product and denominator results therefore apply for \emph{all} $n$.

The exact balancing identity changes no normalized path probability when
observables are translated along with their arguments. For unit spatial
$v,w$, square insertion and Cauchy--Schwarz give, when $|i-j|\le1$,
\[
 |\operatorname{Cov}(R_i(v),R_j(w))|
                         \le C\beta^{-2}e^{C_rB}/\gamma.
\]
For $j-i\ge2$ the two windows have gap $g=j-i-2$ and disjoint bonds.
Proposition~\ref{f16v9:two-window} and the two first insertion estimates
give
\[
 |\operatorname{Cov}(R_i(v),R_j(w))|
                  \le C\beta^{-4}e^{2C_rB}\gamma^{-1}3^{-(j-i-2)}.
\]
The proof allows windows sharing an internal endpoint when $g=0$.
Enlarge $K_r$ once more to absorb these finitely many spatial exponents.
The entry condition only becomes stronger. Taking the supremum over unit
$v,w$ proves the block operator bound. The scalar row sum of its majorant
is $5+2(1/3)/(1-1/3)=6$, proving the coherent variance estimate by the
block Schur bound. Summing with the indicated weights proves
\eqref{f16v9:weighted}. These last sums are the same elementary sums as
V8; the new input is that the full covariance estimate now holds without
one repeated native kernel.
\end{proof}

\begin{theorem}[Full history Hessian and original full/core tails]
\label{f16v9:memory}
Under \eqref{f16v9:entry}, let $\mathsf M_\gamma$ and
$\mathsf C_{\gamma,R}$ be V7's nonadjacent-time common-source Hessian
and its original full/core correction. For every integer $L\ge1$,
\begin{align}
 \|\mathsf M_\gamma\|_{\rm op}&\le3A_{\gamma,B},\notag\\
 \|\mathsf M_\gamma^{>L}\|_{\rm op}
                   &\le3A_{\gamma,B}(1/3)^{L-1},\notag\\
 \|\mathsf C_{\gamma,R}^{>L}\|_{\rm op}
   &\le3A_{\gamma,B}(1/3)^{L-1}
                   +C(R^2/\gamma)(9/19)^{L-1}.
 \label{f16v9:memory-tails}
\end{align}
Here $R$ satisfies V4's original derivative regime, and $>L$ keeps blocks
with $|i-j|>L$. In particular $L=1$ bounds the entire correction. These
bounds hold at every admitted history and every duration.
\end{theorem}

\begin{proof}
V7's exact Haar identity has contact range one for arbitrary bridge
histories. Therefore $\mathsf M_\gamma=-\operatorname{Off}_{>1}
\Gamma_\gamma$ before and after the present change of integration
coordinates. For $d>L\ge1$, summing the block bound gives
$2A_{\gamma,B}\sum_{d=L+1}^\infty3^{-(d-2)}
=3A_{\gamma,B}3^{-(L-1)}$. The block Schur bound proves the first two
claims. V4's common-source core bound already holds at arbitrary bounded
histories, with time tail $C(R^2/\gamma)(9/19)^{L-1}$. Subtract it using
the exact identity
\[
 \nabla^2\log p_{\rm good}
   =\nabla^2[-\log\mathcal W_{\gamma,n}^{\rm full}]
                         -\nabla^2[-\log\mathcal W_{\gamma,n}^{R}].
\]
This proves the last claim. The core event is not translated into a new
product event and is not replaced. All endpoint normalizers remain the
original ones; their separated-time Hessians vanish because they are
single-time factors, not because their values were discarded.
\end{proof}

\subsection{Scope gained and the next upstream spatial problem}

The entire static-background restriction of V8 has been removed for these
conditional common-source estimates. For each fixed $B$, the coherent
covariance and nonlocal memory are $O_{r,B}(\gamma^{-1})$ uniformly over
all bounded histories and all $n$. With
\begin{equation}
 B\le\frac{\log\gamma}{4K_r},\qquad R=4\sqrt{\log\gamma},
 \label{f16v9:mesoscopic}
\end{equation}
the entry and core conditions hold eventually, and
\begin{equation}
 \|\Gamma_\gamma\|_{\rm op}+\|\mathsf M_\gamma\|_{\rm op}
              +\|\mathsf C_{\gamma,R}\|_{\rm op}
                                    \le C_r\gamma^{-3/4}.
 \label{f16v9:mesoscopic-bound}
\end{equation}
This includes arbitrary temporal oscillation within the fixed rescaled
window. There is no $n\gamma^{-1/2}$ error in the conclusion. Individual
scalar normalizations can drift through a long product, but the exact
$p_I$ and the exact finite-path denominator were used at every step.

The remaining large spatial cost has three identifiable sources: the
complete-chart one-step Hilbert--Schmidt comparison, the weighted insertion
norms, and the overlap of the global dressed endpoint with the Gaussian
reference vector. They are all exponential in the number of spatial cubes.
A product estimate cannot remove a spatial cost already present in its
inputs. In particular the present inequality must not be tensorized across
interacting regions or extrapolated to $B\to\infty$ at fixed $\gamma$.

The next useful estimate is spatially local: a boundary-conditioned,
source-marked version of the complete-chart comparison and its endpoint
message bound, on a source region with a retained spatial collar. V9's
product lemma can then be reused without a static-background assumption;
another temporal product proof would not address the remaining cost.
Arbitrary exterior fast configurations can alter a conditional landscape,
so a torus pin by itself is not a boundary-uniform hypothesis. The original
bridge window is still not gauge-invariant coverage of the integrated
physical bridge law. Infrared contact coercivity, comparison with the same
full physical transfer top, physical-time calibration, and interacting
continuum reconstruction remain separate obligations.

\subsection{Diagnostics and preservation}

The accompanying diagnostic is
\begin{center}\small
\path{verify_ym_f16_v9_inhomogeneous_history.py}.
\end{center}
It tests noncommuting Gaussian history centering, both endpoint equations,
exact telescoping of linear multipliers and scalars, varying nonsymmetric
matrix products and their normalized Schur remainders, quasi-multiplicative
scalar identities, centered insertion formulas with unequal endpoints,
and the geometric time-tail constants. All 990 assertions passed.
Assertion families and counts are reported in \path{ym_v9_verification_report.json}. These are exact finite
algebraic checks, not Wilson simulations or proof-assistant verification
of the complete-chart estimates or the infinite-dimensional arguments.
The diagnostic source has SHA--256
\begin{center}\scriptsize\ttfamily
EF4E12653D522EB3F96547274A7AA981F64C4B4BD9398635DB39583BA5FC1657.
\end{center}

The cumulative V8 predecessor has 273,042 bytes and SHA--256
\begin{center}\scriptsize\ttfamily
D12A126849ACA4014F321B2D523BA2669C9C33503B839E86C9E13F010C54DC75.
\end{center}
Only its title version changes and this appendix is inserted before the
final \verb|\end{document}|. The package reports the byte-for-byte recovery
check, input hashes, build results, and reruns of the available V4--V8
scripts. Unavailable archive scripts are not represented as rerun. The
next scheduled companion checkpoint remains V10.

\begin{firewall}
V9 proves exact history centering, normalized nonautonomous fast-product
control, and small coherent common-source covariance with summable temporal
tails for every bounded rescaled bridge history. The original full/core
correction obeys the corresponding operator-norm bounds. The estimates are
uniform in history length but retain an exponential spatial cost, admitting
only the stated logarithmically growing spatial regime. No spatial-volume-
uniform bound at fixed coupling, arbitrary-boundary spatial gluing,
gauge-invariant bridge coverage, infrared contact coercivity, full physical
transfer gap, interacting continuum theory, or positive physical
Yang--Mills mass gap is established by this appendix.
\end{firewall}
% END F16 V9 APPEND
% BEGIN F16 V10 APPEND -- 2026-09-17
% Insert immediately before the final \end{document} in cumulative V9.
% The predecessor body is unchanged; only the cumulative title becomes V10.

\clearpage
\section{V10: finite-order local memory coefficients and a controlled native remainder}
\label{f16v10:context}

The supplied Navier--Stokes import memorandum starts from V3. The active
frontier is now V9: arbitrary bounded bridge histories have a native,
time-summable corrected-source covariance, but its comparison and endpoint
costs are exponential in the spatial volume. Starting again from the V3
Gaussian kernel would repeat work. We instead expand the \emph{already
identified exact corrected-source covariance}, retain its connected
coefficients, and improve only the remainder.

The principal new distinction is between two norms. Every fixed-order
coefficient has a volume- and duration-independent exponentially weighted
space--time row bound. Its \emph{native remainder} has a duration-independent
temporal block bound, still with a displayed exponential spatial cost.
Their combination gives a better quantitative native estimate in the
admitted growing-volume regime. It does not silently turn an asymptotic
coefficient bound into arbitrary-volume control of the full compact law.

\subsection{The V10 checkpoint and the external import actually used}

The primary source examined is OpenAI,
\href{https://cdn.openai.com/pdf/32d9f210-8b73-45\varepsilon0-91bc-82a30aef8a9a/navier-stokes.pdf}
{\emph{Finite Time Blowup for Navier--Stokes}}, released September 8, 2026.
Its Theorem 1.1 constructs smooth-forced finite-time blowup with bounded
kinetic energy. Proposition 9.6 increases the residual order by $1/10$;
Lemma 9.7 fixes the principal inverses and records finite derivative
dependencies. Lemma 8.7 separates two preserved moments from three defect
corrections. Section 5.4 and Lemma 5.4 use cutoff summation rather than
assuming convergence of the raw correction series. The accompanying
\href{https://github.com/openai/NavierStokesAndEuler}{Lean repository}
was inspected at its documentation level, not built or independently
verified here.

The useful import is a \emph{proof architecture}, not an NS estimate on a
YM measure. In particular, the native Wilson operator is fixed. Its genuine
response coefficients cannot be removed by selecting a different forcing,
adding arbitrary counterterms, or changing its probability law. Here a
``correction'' means adding the next coefficient to a specified
approximation of that fixed response. The exact residual, including its
normalizer, is then estimated anew.

The supplied V9 appendix and import memorandum were read in full. The
available f14/f15 archives and Grand Tensor Monograph V076 were inventoried
and their relevant passages checked. F14 V38 already treats connected
noise derivatives and regression for a different, uniformly convex shell
extension; it does not prove the full compact history theorem below.
The monograph's fixed-order Wilson moment bounds explicitly distinguish
finite constants at each order from an all-order summability estimate.
V6--V9 already supply the normalized laws, Haar sources, and history
balancing used here. None is presented again as a new construction.
Only the supplied companion sources were available at this V10 checkpoint;
other programs named in older ledgers were not supplied or recertified.

Several proposed imports need limits. Gauss constraints and the
$9B(n+1)$ common-source coordinates are not five fixed global moments and
are not projected out. Positive covariance does not make its nonadjacent-time
part positive. Auxiliary labels may organize overlapping terms but cannot
delete their native cross interactions. We establish arbitrary \emph{fixed}
order, not an order-uniform convergence theorem or a cutoff-summed replacement
for the native measure. The distinction between order gain and
order-dependent constants will remain explicit.

\subsection{One fixed law, one Gaussian inverse, and measured polynomial coefficients}

Use V6's full unmoving law, V7's full Haar lift, and V9's exact history
centering. To avoid collision with the earlier entropy rate, put
\[
 \varepsilon=\gamma^{-1/2},\qquad \mathcal N=B(n+1),\qquad
 \max_{i,a}|y_{i,a}|\le r<\infty.
\]
This $\varepsilon$ is the expansion parameter, distinct from V6's
entropy-rate sequence $\epsilon_\gamma$. The fast coordinates are $\zeta=(x,a)$, and
\begin{equation}
 d\mu_\varepsilon=Z_\varepsilon^{-1}1_{\mathcal P_\varepsilon}J_\varepsilon^f \exp\!\left(-U_\varepsilon\right)\,d\zeta,
 \qquad d\mu_0=Z_0^{-1}\exp\!\left(-U_0\right)\,d\zeta.
 \label{f16v10:law}
\end{equation}
These are the complete compact conditional and its complete Gaussian
reference, not a restricted core. The Gaussian precision $H$, covariance
$C=H^{-1}$, and mean $m$ satisfy V6's bounds
\[
 h_-I\le H\le h_+I,\qquad \sup_c|m_c|\le C_r.
\]
The covariance is independent of the bridge values; the mean generally is
not. Both depend on the finite geometry and its dressed time ends. No
inverse or Gaussian carrier is changed as the expansion order increases.

By a coefficient in this subsection we mean a finite Taylor coefficient at
$\varepsilon=0$ of the actual exponential-word expression. We do not yet assert a
convergent series. With $V_\varepsilon=U_\varepsilon-\log J_\varepsilon^f$ on its positive chart interior,
write formally
\begin{equation}
 V_\varepsilon\sim U_0+\sum_{q\ge1}\varepsilon^q V_q,\qquad
 r_{\varepsilon,\alpha}\sim\sum_{k\ge1}\varepsilon^k r_{\alpha,k}.
 \label{f16v10:jets}
\end{equation}
Here $r_{\varepsilon,\alpha}=\mathcal D^\varepsilon_\alpha U_\varepsilon-g_\alpha$ is precisely V7's
Haar-compensated residual. The measured coefficients $V_q$ include Haar;
the Haar insertion $\mathcal D^\varepsilon_\alpha U_\varepsilon$ does not acquire a second
Haar logarithmic derivative. The reference divergence in that identity
is Haar divergence.

\begin{proposition}[Actual local action and Haar-source coefficients]
\label{f16v10:local-jets}
For each fixed $q$, $V_q$ is a sum of fixed-range polynomials of total
fast/bridge degree at most $q+2$, with a bounded number of terms incident
to each fast cell. At fixed bounded bridge history their coefficients have
bounds depending on $q,r$, not on $B,n$.
For each $k$, $r_{\alpha,k}$ is the sum of a local bridge coefficient and
the same spatially exponentially local $\mathsf P^{\mathsf T}$ synthesis
of local port coefficients as in V7. These polynomials have degree at most
$k+1$ and occupy only the two port bonds adjacent to the source time,
and their incident internal slices.

In particular, if $U_q$ denotes the classical coefficient, then
\begin{align}
 V_1&=U_1,\notag\\
 V_2&=U_2+\frac13\sum_{i,b,j\ {\rm chord}}w_i|x_{i,b,j}|^2
                  +\frac13\sum_{t,b,v\ {\rm port}}|a_{t,b,v}|^2.
 \label{f16v10:haar-order-two}
\end{align}
For $\eta_\alpha=\mathsf P\mathsf e_\alpha$, with $\mathsf e_\alpha$ the canonical common-source basis vector, the first corrected score is
\begin{equation}
 r_{\alpha,1}
   =(\partial_\alpha+\eta_\alpha\cdot\partial_a)U_1
       +\frac12\sum_p[\eta_{\alpha,p},a_p]\cdot\partial_{a_p}U_0.
 \label{f16v10:first-score}
\end{equation}
The bracket is the Lie bracket in the fixed $\mathfrak{su}(2)$ convention.
In particular the second term in \eqref{f16v10:first-score} is not omitted.
\end{proposition}

\begin{proof}
Each classical summand is a fixed multiple of
$\varepsilon^{-2}s(\prod_j\operatorname{Exp}(\varepsilon T_jw))$, with fixed local linear maps
and the literal product order. Its constant and linear unscaled terms
vanish. The coefficient of $\varepsilon^q$ therefore has degree $q+2$. Endpoint
Gaussian terms are already quadratic. Every differentiated word has a
fixed carrier, with the same uniformly bounded incidence at every order.
For Haar, $-\log j(z)=|z|^2/3+O(|z|^4)$, and the powers $w_i$ are the
ones in V4--V9. This proves \eqref{f16v10:haar-order-two} and the asserted
local structure at higher fixed orders.

For a port, the flow is $h_p\mapsto\operatorname{Exp}(s\varepsilon\eta_p)h_p$.
At $h_p=\operatorname{Exp}(\varepsilon a_p)$, its Taylor jet in logarithmic
coordinates follows from
\[
 \varepsilon^{-1}\log\{\operatorname{Exp}(s\varepsilon\eta_p)\operatorname{Exp}(\varepsilon a_p)\}
   =a_p+s\eta_p+\frac{s\varepsilon}{2}[\eta_p,a_p]+O(\varepsilon^2s)+O(s^2).
\]
This is used only for coefficients at zero, not as a global coordinate
formula for the compact flow. Applying it to $U_\varepsilon$ proves the displayed
first coefficient; the order-zero part is exactly $g_\alpha$ by V7.
Higher coefficients can instead be obtained by inserting
$\operatorname{Exp}(s\varepsilon\eta_p)$ directly in each word and differentiating,
which avoids all logarithm singularities. A coefficient of order $k$ of
that inserted score has degree at most $k+1$. The fixed matrix $\mathsf P$
then synthesizes these primitive port coefficients. Its spatial decay and
two-bond temporal support were proved in V7 and do not deteriorate with $k$.
\end{proof}

There is a simple concrete recipe for $U_1$. Identify $\mathfrak{su}(2)$
with imaginary quaternions as in the inherited action convention. For one
ordered word $\prod_{j=1}^d\operatorname{Exp}(\varepsilon v_j)$,
\begin{equation}
 \varepsilon^{-2}s\left(\prod_j\operatorname{Exp}(\varepsilon v_j)\right)
 =\frac12\left|\sum_jv_j\right|^2
        +\varepsilon\sum_{i<j<k}\det(v_i,v_j,v_k)+O(\varepsilon^2).
 \label{f16v10:word-cubic}
\end{equation}
Indeed the scalar part of $v_iv_jv_k$ is
$-(v_i\mathbin\times v_j)\cdot v_k$; all other cubic scalar terms
vanish. Each negative occurrence is its signed $v_j$, in its actual
position. Summing with the original magnetic and kinetic weights gives
$U_1$. This recipe and \eqref{f16v10:first-score} make the leading memory
coefficient computable from the existing incidence and word data.

\subsection{Connected coefficients: the normalizer is part of every order}

Write $\kappa_0(F_1,\ldots,F_s)$ for the joint cumulant under the
\emph{same} Gaussian law $\mu_0$. It can be defined by the finite moment
partition formula; no exponential integrability of a cubic polynomial
at a nonzero source is assumed. For $j\ge2$ define
\begin{equation}
\begin{split}
 G_j(\alpha,\beta)
 =\sum_{\substack{k,l\ge1\\k+l\le j}}
   \sum_{m'=0}^{j-k-l}\frac{(-1)^{m'}}{m'!}
   \sum_{\substack{q_1,\ldots,q_{m'}\ge1\\
                   q_1+\cdots+q_{m'}=j-k-l}}
 \kappa_0(r_{\alpha,k},r_{\beta,l},V_{q_1},\ldots,V_{q_{m'}}).
\end{split}
 \label{f16v10:connected-coefficients}
\end{equation}
For $m'=0$ the inner sum has one term when $k+l=j$ and is otherwise
zero. The extensive $V_q$ in this formula is expanded as a sum of its
actual local terms when estimating it, not bounded by an extensive norm.

\begin{proposition}[Coefficient ownership and the first two corrections]
\label{f16v10:coefficient-identity}
At each fixed finite $B,n$, the asymptotic coefficients of the exact
normalized covariance $\Gamma_\varepsilon(\alpha,\beta)=
\operatorname{Cov}_{\mu_\varepsilon}(r_{\varepsilon,\alpha},r_{\varepsilon,\beta})$ are
\eqref{f16v10:connected-coefficients}. In particular
\begin{align}
 G_2(\alpha,\beta)&=\operatorname{Cov}_0(r_{\alpha,1},r_{\beta,1}),
 \notag\\
 G_3(\alpha,\beta)&=\operatorname{Cov}_0(r_{\alpha,1},r_{\beta,2})
    +\operatorname{Cov}_0(r_{\alpha,2},r_{\beta,1})
    -\kappa_0(r_{\alpha,1},r_{\beta,1},V_1).
 \label{f16v10:first-coefficients}
\end{align}
The $-\kappa_0$ term is a native measure correction, not an optional
observable correction. Adding any fast-independent scalar to any $V_q$
changes none of these covariance coefficients.
\end{proposition}

\begin{proof}
As an identity of formal series, differentiate twice in the two marking
variables $s,t$ the expression
\[
 \log\mathbb E_0\exp\left\{
 -\sum_{q\ge1}\varepsilon^qV_q+s\sum_{k\ge1}\varepsilon^kr_{\alpha,k}
                          +t\sum_{l\ge1}\varepsilon^lr_{\beta,l}\right\}.
\]
At each fixed coefficient only finitely many Gaussian polynomial moments
occur. The moment partition definition of the logarithm is exactly the
cumulant formula. Selecting the two marks and $m'$ ordered action
insertions gives the factor $(-1)^{m'}/m'!$ and the displayed degree
constraint. This is a formal identity, not a convergent integral of a
cubic exponential. Any cumulant of at least two entries with one constant
entry vanishes; thus fast-independent action scalars cancel, as they must
in a normalized covariance. Source constants cancel for the same reason.

To identify these with actual asymptotic coefficients at fixed $B,n$, use
the complete native pin, a smooth cutoff inside its principal charts, and
finite-order Taylor expansion under the resulting Gaussian majorant.
The detailed short-window version is proved below. The identical
finite-dimensional argument applies here at fixed $B,n$ with constants
allowed to depend on both. The omitted compact and Gaussian tails are
$O(\exp\!\left(-c/\varepsilon^2\right))$ times a finite constant. The partition function has positive
limit $Z_0$, so finite series division is valid. Uniqueness of finite
asymptotic coefficients proves the assertion. No duration-uniform error
is inferred in this identification step.
\end{proof}

For an explicit finite-matrix form of the leading answer, let
$\xi=\zeta-m$ and write the quadratic polynomial
\[
 r_{\alpha,1}=\xi^{\mathsf T}A_\alpha\xi
                         +b_\alpha^{\mathsf T}\xi+c_\alpha,
 \qquad A_\alpha=A_\alpha^{\mathsf T}.
\]
Wick contraction, with all means retained in $b_\alpha,c_\alpha$, gives
\begin{equation}
 G_2(\alpha,\beta)
       =2\Tr(A_\alpha C A_\beta C)+b_\alpha^{\mathsf T}C b_\beta.
 \label{f16v10:leading-matrix}
\end{equation}
The odd centered Gaussian cross terms are zero. No commutation of
$A_\alpha,A_\beta,C$ is required. The full matrix $G_2$ is positive
semidefinite. Set $M_j=-\operatorname{Off}_{>1}G_j$. These are the
coefficients of the actual nonadjacent-time common-source memory by V7's
exact contact support, not new transfers or counterterms.

\subsection{A fixed locality exponent for every finite coefficient}

Give fast cells their ambient space--time distance: nearest time and nearest
spatial cube steps, with a fixed enlargement of the step range when a
literal word meets a diagonal pair. Source coordinates inherit their
cube/time anchor; orientation and colour have fixed multiplicities. The
resulting graph has uniformly polynomial ball growth. Let $d(\alpha,\beta)$
denote its source distance. For a scalar source matrix define
\[
 \|G\|_{\mathrm{loc},b}
   =\max\left\{\sup_\alpha\sum_\beta \exp\!\left(b d(\alpha,\beta)\right)
                                      |G_{\alpha\beta}|,
               \sup_\beta\sum_\alpha \exp\!\left(b d(\alpha,\beta)\right)
                                      |G_{\alpha\beta}|\right\}.
\]
For time blocks acting on the full spatial common-source space set
\begin{equation}
 \|G\|_{\mathrm{time},b}
       =\sup_i\sum_j \exp\!\left(b|i-j|\right)\|G_{ij}\|_{\mathcal E\to\mathcal E}.
 \label{f16v10:time-norm}
\end{equation}
For symmetric block matrices this bounds their operator norm on the whole
history source space. We will use both norms, not identify them.

\begin{theorem}[Uniform linked-coefficient locality]
\label{f16v10:coefficient-locality}
There exists $a>0$, fixed independently of coefficient order, spatial volume
and duration, such that for every fixed $j\ge2$,
\begin{equation}
 \|G_j\|_{\mathrm{loc},4a}\le C_{j,r},\qquad
 \|G_j\|_{\mathrm{time},2a}+\|M_j\|_{\mathrm{time},2a}\le C_{j,r}.
 \label{f16v10:coefficient-norm}
\end{equation}
The constants may grow with $j$. We may fix this $a$ also to satisfy
$2a<\log(19/9)$ and $2a<\log3$.
\end{theorem}

\begin{proof}
The finite-range positive precision admits a norm-convergent Neumann
expansion. As in V5--V6 it gives covariance blocks
$\|C_{cd}\|\le C \exp\!\left(-\nu d(c,d)\right)$ for some fixed $\nu>0$.
The same argument for $L_h$ gives exponentially decaying spatial lift
entries; its temporal support is only the two adjacent bonds. By polynomial
ball growth, covariance and lift entries have bounded weighted absolute
row and column sums at some fixed exponent $\nu_1>0$. Take
$8a<\nu_1$, decreasing $a$ to satisfy the two stated temporal conditions.
Finite colour and coordinate multiplicities only change constants.

Expand each polynomial in \eqref{f16v10:connected-coefficients} in centered
Gaussian coordinates. Local means are bounded by $C_r$, so every
coefficient of these polynomials has a fixed-order bound. Wick's formula
represents their moments by pairings of these coordinates. The cumulant
partition formula cancels exactly every pairing whose graph on the
\emph{polynomial factors} is disconnected: its coefficient is the sum
of the partition coefficients over partitions coarser than its connected
components, which is zero unless there is one component. Self-pairings
within a factor do not connect factors. This argument proves the connected
pairing rule also when some polynomials have nonzero means.

Expand each $V_q$ into local terms. A local source term sits within a fixed
radius of its anchor; a synthesized port-source term has in addition its
lift coefficient, which is treated as an exponentially summable edge
from the source anchor to the port anchor. Every surviving pairing graph,
with these possible additional edges, connects the two marked anchors and
all its interaction anchors. Choose a spanning tree. Each tree edge carries
either one covariance with exponentially decaying block entries or one
lift coefficient. Other contractions have bounded absolute values and
can be absorbed into a constant depending on the finite total degree.
The finite stencil displacements change weights only by fixed-order
constants.

The triangle inequality on the tree path between the two marks gives
\[
 \exp\!\left(4a d(\alpha,\beta)\right)
       \le C_j\prod_{uv\ {\rm in\ the\ tree}}\exp\!\left(4a d(u,v)\right).
\]
Fix $\alpha$ and sum successively over leaf anchors, including $\beta$.
Every such sum is bounded by the same finite weighted edge row bound.
All unmarked interaction sites are summed in this way. At fixed $j$ there
are finitely many degree decompositions, monomials, contractions and trees,
because each action order is at least one and each marked order at least
one. Their number and all polynomial coefficients enter $C_{j,r}$, not
an exponent that worsens with $j$. This proves the first assertion; the
column assertion follows identically or by symmetry.

For fixed times $i,j'$, the block row and column sums are at most
$C_{j,r}\exp\!\left(-4a|i-j'|\right)$. Their Schur bound gives the same block operator
bound. Summing it against $\exp\!\left(2a|i-j'|\right)$ proves the second assertion with
a fixed geometric factor. Masking out contact times cannot increase any
of the absolute row bounds. This proves the assertion for $M_j$ as well.
The spatial dependence of $m(\mathbf y)$ does not create an uncontracted
edge: it is a bounded polynomial coefficient at the fixed base history,
not a variable differentiated in this cumulant estimate.
\end{proof}

This is linked Gaussian coefficient control for this \emph{native history
and its exact Haar-corrected writers}. It is not a native cluster expansion:
no assertion about convergence of the sum over $j$, or about a full native
large-field remainder, has entered its proof.

\subsection{Complete native short-window jets, with their spatial cost paid}

We next prove a finite-order error, rather than merely list formal
coefficients. Use V9's balanced step operators $T_t$, balanced endpoints,
and common Gaussian reference $S_*,a_*$. Their centering means and
balancing multipliers are fixed functions of the base bridge history,
independent of $\varepsilon$. Sources are formed in the original coordinates first,
as in V9. An insertion is the integral of that source observable against
one, two or three balanced augmented kernels, divided by $\lambda_*$ per
bond. All chart and normalization factors belong to these operators.

\begin{proposition}[Finite jets of complete native kernels and insertions]
\label{f16v10:short-jets}
For every fixed integer $J\ge2$ there are $C_{J,r},\varepsilon_{J,r}>0$ such that
the following estimates hold uniformly in $n$, the time location, and the
bounded bridge history when $0<\varepsilon\le \varepsilon_{J,r}$ and
$\varepsilon\,\exp\!\left(C_{J,r}B\right)$ is sufficiently small. Each balanced step and endpoint
has a Taylor polynomial through degree $J$, with zeroth terms $S_*,a_*$,
coefficient norms at most $\exp\!\left(C_{J,r}B\right)$, and remainders at most
\begin{equation}
                    \exp\!\left(C_{J,r}B\right)\varepsilon^{J+1}.
 \label{f16v10:short-error}
\end{equation}
The kernel estimates hold in Hilbert--Schmidt norm and endpoint estimates
in $L^2$. A first corrected-source insertion has such a jet starting at
order one, with an additional factor $\|v\|$. An insertion of the product
of two corrected sources on overlapping windows has a jet starting at
order two, with factor $\|v\|\|w\|$. Its union window has at most three
bonds. The estimates integrate the \emph{complete native domains}.
\end{proposition}

\begin{proof}
Here are the finite-differentiability and tail details. If $f(w)$ is an
unscaled Wilson word action, $f(0)=Df(0)=0$, and
\[
 \varepsilon^{-2}f(\varepsilon w)=\int_0^1(1-s)D^2f(s\varepsilon w)[w,w] ds.
\]
Unitary exponential differentiation bounds every fixed real derivative of
$f$ by a constant depending only on that order and the local word. Hence
its $k$th $\varepsilon$ derivative has absolute value at most $C_k|w|^{k+2}$,
also at $\varepsilon=0$ in this integral representation. Differentiating inserted
words gives the corresponding polynomial bounds for bridge and Haar-port
scores. These are direct insertions, not differentiated principal
logarithms. At most three bonds have dimension $O(B)$; bounded incidence
bounds their summed derivative polynomials by a fixed-order polynomial in
$|w|$ and $B$. The linear lift uses its bounded operator norm. Thus
coherent source directions cost $\|v\|$, or $\|v\|\|w\|$, rather than
an unrecorded number of source coordinates.

Choose a smooth product cutoff equal to one when each original fast group
coordinate has $|\varepsilon z|\le\pi/4$, and zero when any has
$|\varepsilon z|\ge\pi/2$. This cutoff is used only to prove the asymptotic estimate.
The original compact integral is not redefined. On its support all powers
of $j$ and their positive square roots are smooth functions of $ez$;
all their fixed derivatives have polynomial bounds. Differentiating a
cutoff contributes another polynomial in the fast coordinates. Its
nonconstant derivative support lies at $|z|\ge c/\varepsilon$.

V9's complete balanced pin bounds both the native and Gaussian integrands
by $\exp(C_rB-c|w_{\rm centered}|^2)$. It also applies at every real
parameter between zero and $\varepsilon$ wherever the cutoff is nonzero. The bounded
mean per group makes a cutoff or chart complement satisfy
$|w_{\rm centered}|\ge c/\varepsilon-C_r$. The integrals on those complements,
including any fixed polynomial source factor, are therefore bounded by
\[
                     \exp\!\left(C_{J,r}B-c'/\varepsilon^2\right).
\]
This bounds the difference between the cutoff and the \emph{full} native
integrals as well as the relevant Gaussian tails. Cutoff derivatives have
the same flat bound. Derivatives of the regularized whole-space integrand
are a fixed-order polynomial times a Gaussian majorant. Integrating its
Taylor remainder through degree $J$ thus costs
$\exp\!\left(C_{J,r}B\right)\varepsilon^{J+1}$. Fixed-degree Gaussian moments in dimensions
proportional to $B$ cost a fixed polynomial in $B$ times a constant to the
power $B$; both are included in the displayed exponential. Port and
intermediate-slice integrations followed by the endpoint kernel $L^2$
norm obey the same estimate by the Gaussian majorant. The flat term is
absorbed by increasing the fixed-order constant and decreasing $\varepsilon_{J,r}$.

The one-link scalar normalizer is treated by this same finite-dimensional
Laplace argument. Its normalized leading value is nonzero, so reciprocal
powers have finite jets; raising them to $24B$ costs at most an exponential
in $B$ at each fixed order. Endpoint norms are positive integrals with
lower bounds $\exp\!\left(-C_rB\right)$ from a product of one-cube balls, as in V9.
Finite series division for their reciprocal square roots is legitimate
under the stated entry condition. For balanced endpoints use the
\emph{weighted} Gaussian majorant before taking the norm, not an
unbounded exponential multiplier applied to an unweighted error.
The factors $F_{\gamma,2}$, the original endpoint chart mass, and the
bridge half densities are all included this way. The Gaussian balancing
scalars themselves are independent of $\varepsilon$ and are retained exactly.

V7's Gaussian cancellation makes a residual insertion zero at order zero.
A product of two such insertions has no term below order two. This remains
true on overlapping windows and at either time end. The same derivative
and tail estimates with these polynomial insertions prove all assertions.
No analyticity of the native integrals in complex $\varepsilon$ has been asserted.
\end{proof}

\subsection{Analytic dependence on bounded kernel arrays, not on a complex coupling}

Finite jets alone would give a factor proportional to duration if one
expanded an unnormalized long product term by term. The next lemma avoids
that loss by applying V9's exact product normalization before differentiating.
The analytic variable here is an array of \emph{bounded operators}. It is
not a complex-valued compact Wilson coupling.

\begin{proposition}[A uniform analytic two-window functional]
\label{f16v10:analytic-array}
Complexify the Hilbert space of V9. Keep $S_*=P_*+Q_*$ fixed, and let the
fixed real endpoint $a_*$ have $\|a_*\|=1$,
$\langle a_*,\phi_*\rangle\ge\beta>0$. There is an absolute $c>0$
such that on the complex array ball
\begin{equation}
 \sup_t\|T_t-S_*\|<\rho_B,
 \quad\|\ell_L-a_*^{\mathsf T}\|<\rho_B,
 \quad\|a_R-a_*\|<\rho_B,
 \qquad \rho_B=c\beta^4,
 \label{f16v10:array-ball}
\end{equation}
the exact denominator $\ell_LT_0\cdots T_{n-1}a_R$ never vanishes.
Here $\ell_L$ is a bounded complex-linear functional, not an
antiholomorphically varied adjoint vector.

The normalized centered two-window insertion functional is holomorphic
on this ball. For disjoint windows of at most two bonds separated by $g$
unmarked bonds it obeys
\begin{equation}
 |\mathcal F(T,\ell_L,a_R;A_F,A_G)|
       \le C\beta^{-4}\|A_F\|\|A_G\|3^{-g}.
 \label{f16v10:analytic-mixing}
\end{equation}
Normalized one-window expectations of length at most three satisfy the
corresponding bound without $3^{-g}$. Cauchy bounds for derivatives and
Taylor remainders on a smaller array ball cost powers of $\rho_B^{-1}$,
but no number of time slices.
\end{proposition}

\begin{proof}
We give the complex-domain check, which is not a consequence of a real
spectral gap alone. Split every factor relative to
$\mathbb C\phi_*\oplus\phi_*^\perp$. Its scalar block is within
$\delta$ of one, its row and column blocks have norm at most $\delta$,
and its lower block has norm at most $1/5+\delta$, where
$\delta=\sup\|T_t-S_*\|$. In the slope formula use the complex-linear
row functional instead of a conjugated variable. For $\|z\|\le1$,
\[
 |a+cz|\ge1-2\delta,\quad
 \|F(z)\|\le\frac{1/5+2\delta}{1-2\delta},\quad
 \|DF(z)\|<1/3\qquad(\delta\le1/100).
\]
The ball of radius $2\delta$ is invariant by the same scalar inequalities
as V9. Iterate columns and, separately, row functionals. No scalar factor
vanishes. For every interval,
\[
 T_I/p_I=(\phi_*+x_I)(\phi_*^{\mathsf T}+y_I)+E_I,
 \quad |p_I|>0,\quad \|x_I\|,\|y_I\|\le2\delta,
 \quad\|E_I\|\le3^{-|I|},
\]
where $y_I$ is a row functional annihilating $\phi_*$. The vanishing
corner properties are unchanged. The scalar multiplication identity is
$p_{IJ}/(p_Ip_J)=1+y_Ix_J$, whose modulus is between
$1-4\delta^2$ and $1+4\delta^2$. These are exact identities over the
complex field.

For long products, the rank-one endpoint pairing has modulus at least
$(\beta-\eta-4\delta)^2$, whereas the remainder has modulus at most
$4\,3^{-n}$. For short products compare with the positive real number
$\langle a_*,S_*^na_*\rangle\ge\beta^2$ by the telescoping norm
estimate. Exactly the short/long split of V9 gives
\[
 |\ell_LT_{[0,n)}a_R|\ge\frac{\beta^2}{16}|p_{[0,n)}|
\]
when $\delta\le10^{-6}\beta^4$ and
$\eta\le10^{-6}\beta^2$. Choose $c$ smaller than $10^{-6}$.
Positivity is used only for this fixed real reference in the short-product
comparison. The perturbed factors are not claimed to be positive operators
or kernels. The displayed inequality proves nonvanishing for every length.

The centered insertion is the same two-by-two numerator minor as V9,
divided by the square of this actual denominator. A rank-one middle
factor makes that numerator zero algebraically over $\mathbb C$ as well.
The error in the normalized middle product is at most $3^{-g}$.
The scalar quasi-multiplication inequalities, with absolute values, pay the
five segments and the two short windows exactly as in V9. This proves
\eqref{f16v10:analytic-mixing}. A single length-three window has
$|p_I|\ge(1-2\delta)^3>1/2$, so the one-window estimate follows by
the same three-segment calculation.

Products and insertions are continuous multilinear functions of their
bounded operator entries. Division by the nonzero denominator therefore
makes the functional holomorphic. Equip the finite array with its maximum
operator norm, not its sum norm. A complex disc in any unit array direction
of radius comparable to $\rho_B$ remains inside the ball. Cauchy's integral
formula bounds derivatives there by the uniform supremum above times the
corresponding inverse radius powers. The usual polarization, if mixed
derivatives are used, introduces only order-dependent constants. Neither
the radius nor the supremum contains the array length. In particular the
centered $3^{-g}$ factor remains in these derivative bounds.
\end{proof}

This step is deliberately different from assuming a complex strip for the
full compact partition function as a function of $\varepsilon$. The latter has not
been proved here. Real finite jets feed an analytic \emph{normalized
operator-product functional} on an independently established ball.

\subsection{A finite-order native remainder, with no duration loss}

\begin{theorem}[All fixed orders of the native coherent covariance]
\label{f16v10:native-expansion}
Fix $r<\infty$ and an integer $J\ge2$. There are finite
$K_{J,r},\gamma_{J,r}$, independent of $B,n$, such that if
\begin{equation}
 \gamma\ge\gamma_{J,r},\qquad
              \exp\!\left(K_{J,r}B\right)/\sqrt\gamma\le1/100,
 \qquad\max_{i,a}|y_{i,a}|\le r,
 \label{f16v10:entry}
\end{equation}
then for all time blocks
\begin{equation}
 \left\|\Gamma_{\gamma,ij}
       -\sum_{q=2}^{J}\gamma^{-q/2}(G_q)_{ij}\right\|
 \le \exp\!\left(K_{J,r}B\right)\gamma^{-(J+1)/2}
                         3^{-(|i-j|-2)_+}.
 \label{f16v10:block-remainder}
\end{equation}
After enlarging $K_{J,r}$, the same estimate summed in the fixed weighted
temporal norm is
\begin{align}
 \left\|\Gamma_\gamma-\sum_{q=2}^J\gamma^{-q/2}G_q
                             \right\|_{\mathrm{time},a}
       &\le \exp\!\left(K_{J,r}B\right)\gamma^{-(J+1)/2},\notag\\
 \left\|\mathsf M_\gamma-\sum_{q=2}^J\gamma^{-q/2}M_q
                             \right\|_{\mathrm{time},a}
       &\le \exp\!\left(K_{J,r}B\right)\gamma^{-(J+1)/2}.
 \label{f16v10:remainder-norm}
\end{align}
All orders use the same $C,\mathsf P$ and the same locality exponent $a$.
The constants and entry thresholds are allowed to depend on $J$.
\end{theorem}

\begin{proof}
Use V9's balanced representation. Its reference endpoint overlap is at
least $\beta=\exp\!\left(-c_rB\right)$. Proposition~\ref{f16v10:short-jets} supplies
polynomials $T_t^{[J]}(z)$, endpoint polynomials, and insertion polynomials
in an auxiliary \emph{complex} scalar $z$, using their real finite jets.
Their coefficients are bounded by $\exp\!\left(C_{J,r}B\right)$, uniformly over the
entire history array. Choose $R_B=\exp\!\left(-c_{J,r}B\right)$ with $c_{J,r}$ large
enough. On $|z|\le R_B$ all unmarked polynomials and endpoints lie in a
ball strictly smaller than \eqref{f16v10:array-ball}; this follows by
summing finitely many coefficient bounds. First insertion polynomials
have norm at most $\exp\!\left(C_{J,r}B\right)|z|\|v\|$, and overlapping-product
insertion polynomials at most $\exp\!\left(C_{J,r}B\right)|z|^2\|v\|\|w\|$.
The arbitrary-order constants only decrease this auxiliary radius by
another spatial exponential; no duration factor is present.

For separated source times use the analytic centered functional of the
preceding proposition. It has an analytic zero of order at least two at
$z=0$, and on this disc its modulus divided by $|z|^2$ is at most
$\exp\!\left(C_{J,r}B\right)3^{-(|i-j|-2)}\|v\|\|w\|$. Cauchy expansion of that
quotient shows that its remainder after covariance degree $J$ at
$0<\varepsilon\le R_B/2$ is at most
\[
 \exp\!\left(C'_{J,r}B\right)\varepsilon^{J+1}3^{-(|i-j|-2)}\|v\|\|w\|.
\]
This assertion concerns the polynomial kernel arrays evaluated at $\varepsilon$.
To replace them by the \emph{actual} arrays, use the short-jet errors
\eqref{f16v10:short-error} and the derivative bounds on the smaller
array ball. The latter have no duration factor and retain the same
separation factor. Insertion errors are multiplied by the other first
insertion, whose norm is $O(\varepsilon\,\exp\!\left(CB\right))$, and unmarked-array errors
by the two insertion norms. They are no larger than the displayed
$\varepsilon^{J+1}$ budget after increasing its spatial exponent. The real line
segments joining the actual arrays to the polynomial arrays remain inside
the larger ball by \eqref{f16v10:entry}. Thus no derivative of an
uncontrolled long-product error has been taken.

For $|i-j|\le1$, write the covariance as the normalized one-window
expectation of the product of the two sources, minus the product of their
two normalized expectations. The union of their windows has at most three
bonds. The same holomorphic and finite-jet argument gives the displayed
bound with separation factor one. The short-jet assertion covers this
product insertion, including its true lower-order zeros. It also covers
windows cut off by either dressed endpoint.

At each fixed $B,n$, this expansion computes the covariance of the exact
native residuals: balanced coordinates and scalar multipliers are only
changes in the evaluation of that same normalized path law. By
Proposition~\ref{f16v10:coefficient-identity}, uniqueness of fixed-finite
asymptotic coefficients identifies the coefficients with $G_q$, including
the action and normalizer terms. This identification does not transfer a
volume-dependent fixed-finite error; the uniform remainder was just proved
through the analytic array argument. All estimates hold for arbitrary
unit spatial $v,w$, so their supremum gives the block operator estimate.

Since $a<\log3$, the scalar sum
$\sum_j \exp\!\left(a|i-j|\right)3^{-(|i-j|-2)_+}$ is bounded independently of $i,n$.
Sum \eqref{f16v10:block-remainder} and absorb this fixed constant into
$K_{J,r}$. V7's exact identity
$\mathsf M_\gamma=-\operatorname{Off}_{>1}\Gamma_\gamma$ holds for
every coupling. Applying the same mask to each coefficient and to the
block bounds proves the second line of \eqref{f16v10:remainder-norm}.
All coupling-entry, overlap and jet costs are finitely many exponentials
in $B$ at fixed $J$, and are paid by enlarging $K_{J,r}$. None depends on
$n$ or on temporal variation of the admitted bridge history.
\end{proof}

The result gives a precise finite-order correction cycle. Define
$\mathsf M^{[J]}=\sum_{q=2}^J\gamma^{-q/2}M_q$ and
$\mathsf E^{[J]}=\mathsf M_\gamma-\mathsf M^{[J]}$. Then
\begin{equation}
 \mathsf E^{[J+1]}
       =\mathsf E^{[J]}-\gamma^{-(J+1)/2}M_{J+1},\qquad
 \|\mathsf E^{[J+1]}\|_{\mathrm{time},a}
       \le \exp\!\left(K_{J+1,r}B\right)\gamma^{-(J+2)/2}.
 \label{f16v10:cycle}
\end{equation}
The gain is one power of $\varepsilon=\gamma^{-1/2}$ at each fixed cycle. This is
an \emph{absolute finite-order hierarchy}, not the claim
$\|\mathsf E^{[J+1]}\|\le C\varepsilon\|\mathsf E^{[J]}\|$, which need not
hold if a lower-order residual happens to vanish. Every $M_q$ is retained
as a real response of the native theory; none is set to zero by definition.

\subsection{The first explicit correction and a local exported operator}

\begin{theorem}[Leading local native memory with a smaller coherent remainder]
\label{f16v10:first-cycle}
For fixed $r$ there are $K_r,\gamma_r$ such that in the entry regime
\eqref{f16v10:entry} for $J=2$,
\begin{equation}
 \mathsf M_\gamma=\gamma^{-1}M_2+\mathsf E_\gamma,
 \qquad \|M_2\|_{\mathrm{loc},4a}\le C_r,
 \qquad
 \|\mathsf E_\gamma\|_{\mathrm{time},a}
                           \le \exp\!\left(K_rB\right)\gamma^{-3/2}.
 \label{f16v10:first-native}
\end{equation}
The analogous expansion holds for $\Gamma_\gamma$ with $G_2$ in place
of $M_2$. In particular, throughout this regime,
\begin{equation}
 \|\Gamma_\gamma\|_{\mathrm{time},a}
              +\|\mathsf M_\gamma\|_{\mathrm{time},a}
                           \le C_r\gamma^{-1}.
 \label{f16v10:improved-coherent}
\end{equation}
The constants in this conclusion are independent of $B,n$ \emph{within
that coupling-dependent spatial regime}, not for unrestricted volumes.
If $M_2^{[L]}$ keeps only source entries at space--time distance at most
$L$, then
\begin{equation}
 \|\mathsf M_\gamma-\gamma^{-1}M_2^{[L]}\|_{\rm op}
       \le C_r\gamma^{-1}\exp\!\left(-4aL\right)
                              +\exp\!\left(K_rB\right)\gamma^{-3/2}.
 \label{f16v10:local-export}
\end{equation}
\end{theorem}

\begin{proof}
Apply the two preceding theorems with $J=2$. The leading coefficient is
explicitly \eqref{f16v10:leading-matrix} with the Haar term from
\eqref{f16v10:first-score}. Since
$\exp\!\left(K_rB\right)\gamma^{-1/2}\le1/100$, the remainder in
\eqref{f16v10:first-native} is at most $\gamma^{-1}/100$.
The uniform leading temporal row bound gives
\eqref{f16v10:improved-coherent}. Truncating its spatially and temporally
weighted entry rows has absolute row and column error at most
$C_r\exp\!\left(-4aL\right)$, by \eqref{f16v10:coefficient-norm}. Schur's bound and
the native temporal remainder prove \eqref{f16v10:local-export}.
This truncation changes the exported approximation, not the exact action,
measure, source definition or core event. Finite range here refers to pairs
of source indices. Coefficient values can still depend on the whole base
history through $m(\mathbf y)$; replacing exterior data requires its own
response estimate.
\end{proof}

For example, if $B\le(\log\gamma)/(4K_r)$, then the native remainder in
\eqref{f16v10:first-native} is at most $\gamma^{-5/4}$. The leading
memory and coherent covariance are $O_r(\gamma^{-1})$, for \emph{all}
history lengths and bounded histories. This improves the prior
$O_r(\gamma^{-3/4})$ mesoscopic bound without declaring its exponential
spatial error to be a uniform one.

The original full/core correction is also still controlled. With the same
fixed small $a$, V4's bound and V7's exact full/core identity give
\begin{equation}
 \|\mathsf C_{\gamma,R}\|_{\mathrm{time},a}
                   \le C_r\gamma^{-1}+C R^2/\gamma
 \label{f16v10:core-correction}
\end{equation}
for an admissible original core radius. With $R=4\sqrt{\log\gamma}$ this
is $O_r(\log\gamma/\gamma)$ in the stated spatial regime, uniformly in
$n$. The weighting is permissible because $a<\log(19/9)$. No expansion
or new event for $p_{\rm good}$ has been substituted. This last inequality
uses the existing core estimate rather than claiming a new all-order
expansion of a moving or growing cutoff boundary.

\subsection{Compatibility, the spatial obstruction, and the next useful step}

The finite-order hierarchy has three different safeguards. First, all
coefficients are those of the same fixed normalized compact conditional;
connected cumulants retain the actual source-dependent normalization.
For instance $-\log(Z_\varepsilon/Z_0)$ starts with
$\varepsilon\mathbb E_0V_1+\varepsilon^2(\mathbb E_0V_2-\operatorname{Var}_0(V_1)/2)$
at fixed finite volume. That scalar is not arbitrarily made zero.
Second, the changed contact matrix from V7 remains part of the exact
Hessian. Our expansion of nonadjacent times uses its \emph{exact} range,
not a missing contact correction. Third, all common-source coordinates
remain present; covariance under simultaneous colour rotations and spatial
translations, where the base history has those symmetries, is preserved by
Taylor coefficients of the exact expressions. No new physical Gauss
restriction is inferred from coordinate covariance.

The positive-cone suggestion has a precise limit. $G_2\succeq0$ is an
actual Gaussian covariance, but $M_2=-\operatorname{Off}_{>1}G_2$ has
zero time-diagonal and hence zero trace. A nonzero real symmetric
zero-trace matrix cannot be positive semidefinite. Therefore a nonzero
$M_2$ cannot be a positive combination of positive native sandwiches.
Higher $G_j$ may themselves be signed. A truncated series is not asserted
to define a positive transfer or probability law. Positivity stays with
the exact native law and its complete covariance.

What remains is not resolved by increasing $J$ formally. Although every
coefficient has a uniform local row bound, the proved full-native remainder
is only
\[
             \exp\!\left(K_{J,r}B\right)\gamma^{-(J+1)/2}
\]
in a temporal block norm, and its entry condition also depends on $J$.
For any fixed $J$ this does not allow $B\to\infty$ at fixed $\gamma$.
No common all-order entry domain, growth bound on $C_{j,r}$, convergent
series, or unique identification from a cutoff-summed asymptotic series
has been proved. A flat remainder invisible to every fixed Taylor
coefficient can still matter outside the admitted spatial regime. Unlike
a construction in which the forcing can be selected, the native Wilson
object cannot be replaced by a different realization having the same jets.

The new concrete input for the spatial problem is the local operator
$M_2$ and, at any fixed higher order, its linked successors. The next target
is a \emph{boundary-conditioned marked remainder estimate} with a retained
spatial collar, rather than another global Gaussian inverse or temporal
product lemma. It must pay the actual boundary and endpoint messages while
replacing the spatial exponential in the remainder by a source-local
bound. These coefficients allow that remainder to start beyond a chosen
finite order without forgetting what was subtracted. They do not provide
arbitrary-boundary compact convexity or a source-visible capacity bound.
The proposed $o(a_t)$ spectral handoff additionally requires an identified
physical-time scale and a comparison in the norm of the same physical
transfer; our source-coordinate covariance norm is not that norm.
Gauge-invariant bridge coverage, infrared contact coercivity and
interacting continuum reconstruction remain separate obligations.

\subsection{Diagnostics, preservation, and scope}

The accompanying diagnostic is
\begin{center}\small
\path{verify_ym_f16_v10_finite_order_memory.py}.
\end{center}
Its report separates exact quaternion/BCH and Haar coefficients, correlated
Gaussian polynomial moments, connected normalization coefficients, complex
near-reference products, finite path insertion jets, and the order/locality
bookkeeping. All 608 exact assertions passed; the assertion families are supplied in
the package report. The checker SHA--256 is
\begin{center}\scriptsize\ttfamily
F6F55A4534AFBA1A64E094B1D96D52CB6D5473762E3633C29E8998182BAA97B0.
\end{center} These checks are not formal verification of the analytic
kernel-jet, linked-coefficient, or Cauchy estimates. The native derivative
suprema and stage-dependent constants are proved finite but are not
numerically optimized. The NS Lean development has not been run here.

The V9 predecessor has 309,137 bytes and SHA--256
\begin{center}\scriptsize\ttfamily
26FF80027553957199D6DADEBBD7E5C6BE28E0F6B97FE5E2D070ECB57362BC50.
\end{center}
The cumulative V10 changes only V9's title version and inserts this exact
appendix before its final \verb|\end{document}|. The package records the
byte-for-byte predecessor recovery, unchanged input hashes, compilation,
render review and reruns of the supplied V4--V9 diagnostics. Those
unchanged scripts passed 9,467, 89,339, 906, 311, 609 and 990 assertions,
respectively. Unavailable
historical or companion scripts are not described as rerun. The next
scheduled five-version companion checkpoint is V15.

\begin{firewall}
V10 proves fixed-order, spatially local linked coefficients for the actual
Haar-corrected native memory, and a finite-order full-native remainder
uniform in duration with an explicit exponential spatial cost. The first
coefficient yields an $O_r(\gamma^{-1})$ coherent bound and a local exported
memory operator in the admitted logarithmic spatial regime. The original
full/core correction remains controlled there. The native measure and its
normalizers are unchanged. No unrestricted-volume remainder, arbitrary
spatial-boundary gluing theorem, convergent all-order realization,
gauge-invariant bridge coverage, physical transfer gap, or interacting
continuum Yang--Mills mass gap is established by this appendix.
\end{firewall}
% END F16 V10 APPEND

% BEGIN F16 V11 APPEND -- 2026-09-17
% Insert immediately before the final \end{document} in cumulative V10.
% No inherited theorem is edited. The cumulative title becomes V11.

\clearpage
\section{V11: complete conditional cylinders and spatial screening of the marked remainder}
\label{f16v11:context}

V10 localizes the coefficients, not the full native remainder, and explicitly
leaves exterior-data replacement unpaid. We now integrate a \emph{proper
spatial cylinder}, with its actual exterior fast history held fixed. All
crossing Wilson terms are retained. The interior still ranges over its
complete compact coordinates. Its comparison cost depends on the number of
\emph{interior cubes}, not on the size of the ambient torus. A boundary-marked
Gaussian comparison then identifies its leading memory with the already
computed full-history coefficient, up to an exponentially small collar error.
This gives a native conditional screening statement, not just a coefficient
with a finite matrix carrier.

Two limitations are declared at the outset. The exterior values that enter
the cylinder must lie in a fixed \emph{rescaled fast-field window}, at every
time. We do not prove that this event is typical under the actual exterior
posterior. Also a full covariance is not the average of conditional
covariances: the covariance of the conditional means must be retained. The
last subsection writes that additional return with its actual measure.

\subsection{Audit and the new interface}

The V10 appendix was read in full, and the supplied f14/f15 and monograph
inventories were searched with targeted passages read. In particular, f15
V13--V14 already prove boundary response and posterior-weighted transfer for
another source and another conditioned chart. The monograph's
\texttt{thm:YM-Wilson-blanket} and
\texttt{thm:YM-boundary-mediation} already record the Markov-blanket and
conditional-mean mechanisms. These identities are inherited tools, not the
new theorem. No earlier all-good probability or native functional inequality
is transferred to the present history law.

The new objects are the complete compact spatial-cylinder conditional of
V5's augmented history, its \emph{interior-port} Haar lift, and its
source-marked expansion with retained crossing interactions and dressed ends.
V9's normalized product lemma and V10's short-jet method are reused after
checking their hypotheses for this different conditional operator. Only the
first nonzero memory coefficient is needed. Repeating the all-order temporal
construction would not resolve the spatial issue.

The import memorandum's separation of order gain from locality cost is kept:
one power of the native expansion parameter is gained after extracting the
leading coefficient, while the cost is localized to a declared cylinder.
There is no new appeal to an NS theorem or to convergence of a formal series.
The inverse-decay argument used below is the elementary finite-range
positive-matrix argument already used in these notes; for classical context,
see S.~Demko, W.~F.~Moss and P.~W.~Smith,
\href{https://doi.org/10.1090/S0025-5718-1984-0758197-9}
{\emph{Decay rates for inverses of band matrices}}, Math. Comp. 43 (1984),
491--499. The boundary identities needed here are proved explicitly.

\subsection{The exact cylinder: condition the exterior, do not delete it}

Fix the ambient cube torus $\mathcal B=(\mathbb Z/m\mathbb Z)^3$, $B=m^3$,
$m\ge4$, and a nonempty proper set $D\subset\mathcal B$ of $M=|D|$ cubes.
The cylinder consists of all $x_{i,b}$, $0\le i\le n$, and all unmoving
$a_{t,b}$, $0\le t<n$, with $b\in D$. These are V5's five chord and seven
nonroot temporal variables per cube. Write $I$ for these free real
coordinate indices, and $O=I^c$ for the exterior indices. V5's unmoving
coordinates are important: $h_t=\operatorname{Exp}(a_t/\sqrt\gamma)$ has
independent rooted tree pins even after the exterior is fixed.

Let $\mathbf q=\zeta_O$ be the exterior fast data, held fixed in these
unmoving coordinates. Choose once and for all a geometric range $r_0\ge4$
large enough to contain the cube carriers of every primitive action word
and of every unlifted bridge mark. Let $\partial^+D$ be the exterior cubes
at spatial graph distance at most $r_0$ from $D$. It is a safe enlargement
of the actual collar. Assume
\begin{equation}
 \max_{i,e}|y_{i,e}|\le r,\qquad
 \sup_{c:\,b(c)\in\partial^+D}|q_c|\le Q,
 \qquad r,Q<\infty .
 \label{f16v11:windows}
\end{equation}
Here a cell vector has at most 36 real coordinates, with the final-time
cell having 15. The supremum includes every time and both species when
present. Exterior cubes beyond this collar are arbitrary. Constants
$C_{r,Q},K_{r,Q}$ will not depend on $B,M,n,D$, or on the individual histories.
Their dependence on the two fixed windows is not suppressed.

Split the exact classical exponent of V5 as
\begin{equation}
 U_\gamma(\zeta;\mathbf y)
  =U_{\gamma,D}(\zeta_I;\mathbf q,\mathbf y)
       +U_{\gamma,O}(\mathbf q;\mathbf y).
 \label{f16v11:action-split}
\end{equation}
Assign to the first term all anchors on free cells and all other primitive
terms whose fast carrier meets $I$; assign the remaining terms to the
second. Pure bridge terms can be assigned by their geometric carrier.
For definiteness assign every such term with carrier meeting $D$ to the
first term. Each original term is assigned once. Let $J^f_{\gamma,D}$ be
the product of V5's free chord Haar powers $w_i$ and free port Haar powers
one. Define
\begin{equation}
 Z_{\gamma,D}(\mathbf q;\mathbf y)
 =\int_{\mathcal P'_I}J^f_{\gamma,D}
                       e^{-U_{\gamma,D}}\,d\zeta_I,
 \qquad
 d\mu_{\gamma,D}^{\mathbf q}
 =\frac{1_{\mathcal P'_I}J^f_{\gamma,D}e^{-U_{\gamma,D}}}
        {Z_{\gamma,D}(\mathbf q;\mathbf y)}\,d\zeta_I.
 \label{f16v11:conditional-law}
\end{equation}
The domain is the product of the \emph{complete} principal charts for the
free groups. In particular this is not V4's core. Fast-independent factors
assigned in \eqref{f16v11:action-split} are part of $Z_{\gamma,D}$; they
cancel only in its normalized conditional probability.

Let $\mathcal A_D$ be the common-bridge basis marks whose spatial anchors
have distance greater than $r_0$ from $D^c$, at every available time. Their
entire unlifted action-source carrier is interior. The part
$U_{\gamma,O}$ has zero derivatives in these directions. We consider the
associated source space $\mathcal U_D$, with dimension at most $9M(n+1)$.
Empty source sets give vacuous estimates, not an exceptional case.

\begin{proposition}[An actual conditional with a surviving independent pin]
\label{f16v11:conditional-pin}
The measure \eqref{f16v11:conditional-law} is the regular conditional of
V6's full native fast-history law given $\zeta_O=\mathbf q$, with its
explicit positive-density version at every interior chart point. It depends
on the exterior only through the actual word collar, hence through
$\partial^+D$. All crossing interactions are retained, and
\begin{equation}
 U_{\gamma,D}\ge c_{\rm pin}\|\zeta_I\|^2,
 \qquad c_{\rm pin}=(6\pi^2)^{-1},
 \label{f16v11:pin}
\end{equation}
for its indicated nonnegative-term assignment, on the complete free charts.
The same pin holds for its quadratic limit. The number of primitive terms
per free cell and the real dimension per slice are bounded independently
of the ambient torus.
\end{proposition}
\begin{proof}
Product Haar measure factorizes between $I$ and $O$. Its chord half-density
powers at the two dressed ends do so as well. Substitution of
\eqref{f16v11:action-split} in the full integral gives
\eqref{f16v11:conditional-law} by ordinary normalization. Local word
support proves the collar assertion, including the moving frames, each
of which uses only the five chords of its own cube.

Every free cube retains all its internal magnetic plaquettes with weights
$w_i\ge1/2$, its two endpoint Gaussian quadratics when applicable, and
its seven tree kinetic terms on each free port bond. In unmoving variables
these last words are exactly $h_s^{-1}h_t$ within that cube, with root
$h_o=I$. V5's cube inequality and rooted depth-sum estimate therefore give
$|x|^2/(3\pi^2)$ and $|a|^2/(6\pi^2)$ on that cube. Exterior fixing changes
neither. All other assigned Wilson terms, including crossing terms with
fixed exterior entries, are nonnegative. Add these estimates. The quadratic
limit preserves the inequality. Each word touches a bounded number of
cubes, and each cube meets a bounded number of words, so all counts involve
$M$, not $B$.
\end{proof}

The pin is \emph{not} a claim that the full conditional potential is globally
convex. Its usefulness is complete-chart integration after the exterior has
been fixed. Also the kinetic normalizers of the original full transfer and
its external $F_{\gamma,2}$ factors have not changed: they cancel from the
fast conditional. They remain in the ambient path weight.

\subsection{Conditional Gaussian data and boundary-reaching inverse identities}

Use the same full Gaussian $N,S,\Lambda,H,C,m$ as V9--V10. Restrict the
spatial internal and port matrices by cube labels; subscripts $D,O$ in
spatial formulas denote these restrictions. In particular
\[
 S_D=S_{DD}=\mathsf H_D+A_D^{\mathsf T}A_D,\qquad
 L_D=(L_h)_{DD},\qquad E_D=E_{\rm br}|_{\text{port columns in }D}.
\]
The columns of $A_D$ are restricted, but its rows include \emph{all}
mixed plaquettes that meet those columns. Crossing plaquettes therefore
contribute to $S_D$. This is not the matrix obtained by deleting them.

The conditional Gaussian in the original unmoving coordinates has precision
$H_{II}$, covariance $C^D=(H_{II})^{-1}$, and mean
\begin{equation}
 m^D=m_I-C^DH_{IO}(\mathbf q-m_O).
 \label{f16v11:Gaussian-conditional}
\end{equation}
It is important that $C^D\ne C_{II}$ in general. Write, at exterior ports,
\[
 q^z_t=q^a_t-\mathsf F_O(q^x_t-q^x_{t+1}).
\]
This is the \emph{Gaussian linear} moving coordinate of the fixed exterior;
it is not asserted to be the exact finite-coupling logarithm of the moving
group. Its cell norms are bounded by $C Q$.

Inside $D$, change Gaussian variables by
$a_{t,D}=z_{t,D}+\mathsf F_D(x_{t,D}-x_{t+1,D})$. Conditional on the exterior,
the fast Gaussian precision separates as
\begin{equation}
 W_n^D\ \oplus\ (I_n\otimes L_D),
 \label{f16v11:separated-precision}
\end{equation}
where $W_n^D$ is the spatial principal restriction of V3's $W_n$. Its
internal and port linear sources, respectively, are
\begin{equation}
 b^x_i=\Lambda_Dy_i+(A^{\mathsf T}A)_{DO}q^x_i,
 \qquad
 b^z_t=E_D^{\mathsf T}(y_t-y_{t+1})+(L_h)_{DO}q^z_t.
 \label{f16v11:boundary-sources}
\end{equation}
The free internal Gaussian mean is $-(W_n^D)^{-1}\mathbf b^x$. All its
endpoint entries are retained; in particular $W_n^D$ is not made periodic
in time.

\begin{proposition}[Uniform conditional means, lifts, and screening identities]
\label{f16v11:Gaussian-screening}
There exist $c,C,\nu>0$, independent of $D,B,n$, such that all the indicated
conditional precisions and covariances have the inherited positive spectral
bounds, and
\begin{equation}
 \sup_{c\in I}|m^D_c|\le C_{r,Q},\qquad
 |m^D_c-m_c|\le C_{r,Q}e^{-\nu s_D(c)}.
 \label{f16v11:mean-screening}
\end{equation}
Here $s_D(c)$ is the spatial distance of its cube to $D^c$. The covariance
and lift differences are given by boundary-reaching identities
\begin{equation}
 \begin{aligned}
 C_{II}-C^D&=-C^D H_{IO}C_{OI},\\
 (L_h^{-1})_{DD}-L_D^{-1}
       &=-L_D^{-1}(L_h)_{DO}(L_h^{-1})_{OD}.
 \end{aligned}
 \label{f16v11:resolvent-boundary}
\end{equation}
Every factor inverse in these formulas has uniformly summable exponentially
weighted block rows and columns at a common smaller exponent. In particular
\begin{equation}
 \sum_{d\in I}e^{\nu d(c,d)/4}
       \|(C_{II}-C^D)_{cd}\|
       \le C e^{-\nu s_D(c)/4}.
 \label{f16v11:covariance-screening}
\end{equation}
The analogous spatial estimate holds for the port inverse, including its
zero extension outside $D$ when its row starts inside $D$.
\end{proposition}
\begin{proof}
Conditioning a positive Gaussian quadratic gives
\eqref{f16v11:Gaussian-conditional} by completing its $I$ square. The
block-diagonal, determinant-one change between unmoving and moving variables
is local to each spatial cube and uses only adjacent times. It therefore
commutes with the division into entire spatial cylinders. V4's separated
Gaussian square, with the exterior fixed, gives
\eqref{f16v11:separated-precision}--\eqref{f16v11:boundary-sources}.
Principal restriction preserves the bounds $h_-I\le H\le h_+I$ and the
spatially whitened bounds $4I\le\widehat W_n\le15I$. Similarly
$\lambda_-I\le L_D\le\lambda_+I$, using V7's constants. Also
$4I\le\mathsf M_D^{1/2}S_D\mathsf M_D^{1/2}\le10I$.

For any of these finite-range positive matrices $K$, with fixed spectral
bounds $k_-I\le K\le k_+I$, write
\[
 K^{-1}=k_+^{-1}\sum_{j\ge0}(I-K/k_+)^j.
\]
The summands have norm at most $(1-k_-/k_+)^j$ and range at most $j$ in
an appropriate fixed enlargement of the ambient cube--time graph. A ball
has at most $C(j+1)^4$ labels even with periodic spatial directions and
arbitrary time length. This proves exponential inverse block decay and
weighted row/column summability, also for every principal restriction.
This is a bound in the ambient distance; disconnected restricted regions
do not invalidate it. Local sources in
\eqref{f16v11:boundary-sources} are bounded by $C_{r,Q}$ per cell. The
row bound proves the first estimate in \eqref{f16v11:mean-screening}.

The $II$ block of $HC=I$ reads
$H_{II}C_{II}+H_{IO}C_{OI}=I$; multiply by $C^D$. This proves the first
identity in \eqref{f16v11:resolvent-boundary}; the second is identical.
The only nonzero entries of $H_{IO}$ join the fixed-width boundary.
Every product term in the first identity therefore travels from $c$ to
that boundary, and then to $d$. The sum of its two inverse-edge lengths,
up to a fixed range shift, is at least $s_D(c)$ and at least $d(c,d)$.
Reserve a quarter of the inverse exponent for each of these two distances
and sum the remaining factors by their weighted row bounds. This proves
\eqref{f16v11:covariance-screening} with constants enlarged for the range
shift. For the mean use \eqref{f16v11:Gaussian-conditional} and the same
boundary path, with $|q_o-m_o|\le C_{r,Q}$ on the actual collar. No value
of $q$ away from that collar is used. The port proof is spatially the same;
its exterior columns start beyond the boundary and have the same tail
bound. Choose a single smaller $\nu$ for the finitely many matrices.
\end{proof}

\subsection{A conditional Haar lift: do not differentiate a frozen port}

V7's unrestricted lift generally varies ports outside $D$. It cannot be
inserted into a conditional integration by parts after those ports have
been fixed. Use instead
\begin{equation}
 (\mathsf P_D\mathbf v)_t
     =-L_D^{-1}E_D^{\mathsf T}(v_t-v_{t+1}),
 \qquad \mathbf v\in\mathcal U_D.
 \label{f16v11:Dirichlet-lift}
\end{equation}
It acts only on the free \emph{unmoving} temporal groups $h_{t,b,v}$.
Its $\ell^2$ norm is uniformly bounded, its time support is the two adjacent
bonds, and its spatial coefficients have a uniform exponential row bound.
These follow from the preceding proposition and $\|d_t\|\le2$.
For a canonical source $\alpha$ let
\[
 \mathcal D^D_\alpha
  =\partial_\alpha+
     \sum_{p\in D,k}(\mathsf P_D)_{p,k;\alpha}X^\gamma_{p,k},
 \qquad
 g^D_\alpha=\partial_\alpha[-\log Z_{0,D}(\mathbf q;\mathbf y)],
 \qquad
 r^D_{\gamma,\alpha}=\mathcal D^D_\alpha U_{\gamma,D}-g^D_\alpha.
\]
Here $X^\gamma$ is V7's actual compact left-translation derivative. The
coefficients of $\mathsf P_D$ are independent of $\mathbf q,\mathbf y$
and all fast variables. All derivatives keep the exterior \emph{unmoving}
coordinates fixed.

\begin{theorem}[Exact conditional source representation]
\label{f16v11:Haar-conditional}
The Gaussian part of the lifted source is deterministic:
\begin{equation}
 (\partial_\alpha+(\mathsf P_D\mathsf e_\alpha)\cdot\partial_{a_D})
                          U_{0,D}=g^D_\alpha.
 \label{f16v11:Gaussian-cancellation}
\end{equation}
With $\Gamma^D_\gamma(\mathbf q)=
\operatorname{Cov}_{\mu_{\gamma,D}^{\mathbf q}}(\mathbf r^D_\gamma)$,
the exact conditional Hessian is
\begin{equation}
 \nabla^2_{\mathcal U_D}[-\log Z_{\gamma,D}]
       =\mathcal K^D_\gamma-\Gamma^D_\gamma,
 \qquad
 (\mathcal K^D_\gamma)_{ij}=0\quad (|i-j|\ge2).
 \label{f16v11:conditional-contact}
\end{equation}
The contact is the symmetrized expectation of
$\mathcal D^D_\alpha\mathcal D^D_\beta U_{\gamma,D}$, not the old direct
Hessian. In particular the conditional nonadjacent-time memory is exactly
$\mathsf M^D_\gamma(\mathbf q)=-\operatorname{Off}_{>1}\Gamma^D_\gamma(\mathbf q)$.
\end{theorem}
\begin{proof}
In Gaussian moving variables, common-source magnetic differentiation has
no free $x$ coefficient, because $\Lambda\mathcal E=0$ before restriction.
All action terms involving a free $x$ were retained, so exterior fixing
does not change that cancellation. The free port coefficient in a direction
$\mathbf v$ is $\sum_t(v_t-v_{t+1})^{\mathsf T}E_D z_{t,D}$.
The port Gaussian gradient is
$L_D z_{t,D}+b^z_t$. Substitution of
\eqref{f16v11:Dirichlet-lift} cancels every free port coefficient. The
remaining quantity is independent of the free variables. Its expectation
is the derivative of $-\log Z_{0,D}$, because the mean Gaussian score
$\partial_{a_D}U_{0,D}$ is zero. This proves
\eqref{f16v11:Gaussian-cancellation}, including its scalar value.

At finite coupling, condition the exterior first and perform only the free
port integrations in Haar variables. The field in
\eqref{f16v11:Dirichlet-lift} has zero Haar divergence. It does not vary
free chord weights or endpoint amplitudes. Thus
$\mathbb E_D[V_\alpha F]=\mathbb E_D[F V_\alpha U_{\gamma,D}]$ and
$\mathbb E_D V_\alpha U_{\gamma,D}=0$. Differentiating a normalized
expectation and using this identity is precisely V7's argument, with the
sum of ports now restricted to $D$. It gives
\eqref{f16v11:conditional-contact} with the symmetrized contact; it does
not assume that two group derivatives commute. Each kinetic word uses one
port bond and its two adjacent bridge times, even with fixed exterior
entries. Each magnetic word uses one bridge time. The two lift bond sets
are disjoint at nonadjacent source times, so the corresponding double
contact vanishes exactly. At every finite regulator the compact derivatives
and integrable chord weights justify differentiation, with no cut-boundary
term from the port logarithms.
\end{proof}

For an interior source the Gaussian scalar in
\eqref{f16v11:Gaussian-cancellation} can also be written explicitly. With
$E_O$ the exterior port columns of $E_{\rm br}$, it is the coefficient of
$\mathbf v$ in
\begin{align}
 g^D_{\mathbf v}
 ={}&\sum_i v_i^{\mathsf T}D^{\mathsf T}Dy_i\notag\\
 &+\sum_t(v_t-v_{t+1})^{\mathsf T}
 \left[(I-E_D L_D^{-1}E_D^{\mathsf T})(y_t-y_{t+1})
       +(E_O-E_D L_D^{-1}(L_h)_{DO})q^z_t\right].
 \label{f16v11:conditional-g}
\end{align}
All terms differentiated by these interior marks belong to $U_{\gamma,D}$.
There is no missing derivative of $U_{\gamma,O}$. The formula also shows
that the Gaussian exterior source message uses only the adjacent port
bonds. By the port inverse boundary identity,
\begin{equation}
 |g^D_\alpha-g_\alpha|\le C_{r,Q}e^{-\nu s_D(\alpha)/C},
 \label{f16v11:g-screen}
\end{equation}
where a fixed source-carrier range shift is absorbed into $C$. This is a
Gaussian screening estimate, not yet an estimate on the nonlinear mean.

\subsection{The leading coefficient with the exterior retained}

Put $\varepsilon=\gamma^{-1/2}$ as in V10. Taylor coefficients are taken
at fixed rescaled $\mathbf q,\mathbf y$. The exact conditional lifted score
has $r^D_{\gamma,\alpha}=\varepsilon r^D_{\alpha,1}+O(\varepsilon^2)$
at each fixed finite dimension. Its first coefficient is given by V10's
ordered-word and Haar-insertion formula, using $U_{1,D}$, the free port
variables, and $\mathsf P_D$. All frozen exterior coordinates remain in
the polynomial coefficients. For example the half-bracket term is
\[
 \frac12\sum_{p\in D}[\eta^D_{\alpha,p},a_p]
                                    \cdot\partial_{a_p}U_{0,D},
 \qquad \eta^D_\alpha=\mathsf P_D\mathsf e_\alpha.
\]
It is not discarded at a boundary port. Define
\[
 G^D_2(\mathbf q)_{\alpha\beta}
   =\operatorname{Cov}_{\mu^\mathbf q_{0,D}}
                     (r^D_{\alpha,1},r^D_{\beta,1}),\qquad
 M^D_2(\mathbf q)=-\operatorname{Off}_{>1}G^D_2(\mathbf q).
\]
If $\xi_D=\zeta_I-m^D$ and
$r^D_{\alpha,1}=\xi_D^{\mathsf T}A^D_\alpha\xi_D+
 b^{D\mathsf T}_\alpha\xi_D+c^D_\alpha$, the exact evaluation is
\begin{equation}
 G^D_2(\mathbf q)_{\alpha\beta}
 =2\Tr(A^D_\alpha C^D A^D_\beta C^D)
                  +b^{D\mathsf T}_\alpha C^D b^D_\beta.
 \label{f16v11:leading-conditional}
\end{equation}
The matrices can be noncommuting. This is the old Wick identity on the
\emph{new conditional} covariance and mean, not on $C_{II}$ and $m_I$.

Let $\mathcal A\subset\mathcal A_D$ have all its spatial anchors at least
$L+r_0$ from $D^c$, and let $\chi_{\mathcal A}$ project onto these sources
at every time. Use the source-distance and time-block norms from V10.

\begin{theorem}[Boundary-marked screening of the leading memory coefficient]
\label{f16v11:coefficient-screening}
There is $b>0$, fixed independently of $D,B,n,r,Q$, which can be chosen
smaller than V10's $a$, such that
\begin{equation}
 \|G^D_2(\mathbf q)\|_{\mathrm{loc},4b}
       +\|M^D_2(\mathbf q)\|_{\mathrm{loc},4b}\le C_{r,Q}.
 \label{f16v11:conditional-coefficient-locality}
\end{equation}
Moreover, with $G_2,M_2$ the \emph{full-history} coefficients of V10 at
this same bridge history,
\begin{align}
 &\left\|\chi_{\mathcal A}
   (G^D_2(\mathbf q)-G_2)\chi_{\mathcal A}\right\|_{\mathrm{time},b}
       \le C_{r,Q}e^{-bL},\notag\\
 &\left\|\chi_{\mathcal A}
   (M^D_2(\mathbf q)-M_2)\chi_{\mathcal A}\right\|_{\mathrm{time},b}
       \le C_{r,Q}e^{-bL}.
 \label{f16v11:coefficient-screening-bound}
\end{align}
The same estimate, with twice the constant, compares any two exterior
histories satisfying \eqref{f16v11:windows}. No bound on their difference
or temporal variation is required.
\end{theorem}
\begin{proof}
Use centered Gaussian coordinates in both polynomials. Their primitive
bridge and port pieces have bounded degree two, fixed local support and
bounded coefficients; the conditional means are bounded by the preceding
proposition. The only extended synthesis is by $\mathsf P_D^{\mathsf T}$,
whose entries have uniformly summable spatial exponential rows and two-bond
time support. Wick contraction gives one or two covariance edges between
the two nonconstant factors. It must connect them: internal contractions
and constant pieces cancel in a covariance. The same weighted-tree
summation as V10, now with $C^D$, proves
\eqref{f16v11:conditional-coefficient-locality} at a sufficiently small
fixed $b$. This argument has only finitely many quadratic contractions.

We describe the additional \emph{boundary mark} needed for comparison,
rather than inferring boundary independence from matrix-index locality.
Expand the full and conditional quadratic source polynomials as their
primitive local terms followed by their port-lift synthesis. An interior
primitive term with all inputs away from the boundary is the same ordered
word coefficient in both expressions. The possible differences are:

\begin{enumerate}
\item a synthesized port term outside $D$, or a primitive term with an
      exterior input replaced by its fixed value;
\item replacement of a full lift entry by its restricted inverse entry;
\item replacement of a mean $m_c$ by $m^D_c$;
\item replacement of a covariance $C_{cd}$ by $C^D_{cd}$.
\end{enumerate}

This list is exhaustive because the exponent and the port vector field
are obtained by literal restriction, not by changing word weights. In the
first case the source-to-port synthesis reaches the boundary. In the second
case use the port version of \eqref{f16v11:resolvent-boundary}. In the third
use \eqref{f16v11:Gaussian-conditional}, and in the fourth use the first
identity in \eqref{f16v11:resolvent-boundary}. These formulas represent the
difference by a path to a boundary index, not just by a norm estimate.
Frozen values and full/conditional means at that index are bounded by
$C_{r,Q}$. Entries $H_{IO}$ have fixed range and bounded row sums.

Telescope products in the Wick formula so that a difference term contains
one of the four marked differences, with other factors chosen from one of
the two sets of data. Its connected graph contains both source anchors and
a boundary visit. The total edge length is therefore at least the distance
between the sources and at least the distance from the first source to the
boundary, up to fixed word-range shifts. Reserve one part of the common
inverse/lift exponent for each of these two distances. Sum the remaining
edge factors using their exponential row bounds. With $b$ smaller again,
this gives the stronger row estimate
\begin{equation}
 \sum_{\beta\in\mathcal A_D}e^{2b d(\alpha,\beta)}
   |(G^D_2(\mathbf q)-G_2)_{\alpha\beta}|
       \le C_{r,Q}e^{-2b(s_D(\alpha)-r_0)_+}.
 \label{f16v11:boundary-row}
\end{equation}
There is an identical column bound when the second source is interior.
Polynomial growth of the ambient graph makes all boundary-time sums finite
with constants independent of $n$; no count of the entire boundary is used
instead of its decaying weights. Products of means are handled by telescoping
as above, so this also covers nonzero and inhomogeneous base histories.

For both sources in $\mathcal A$, the row and column estimates bound the
$(i,j)$ block operator norm by $C_{r,Q}e^{-2bL}e^{-2b|i-j|}$. Sum with
weight $e^{b|i-j|}$ to obtain \eqref{f16v11:coefficient-screening-bound}.
The time-off-diagonal mask preserves the absolute row bounds and gives the
memory assertion. Comparison through the same full coefficient proves the
last assertion. No exterior density or probability has been substituted.
\end{proof}

This screening statement compares coefficients of different \emph{conditional}
laws. It does not identify their native covariances with the full native
covariance. The role of the full coefficient is a common, already defined
local approximation, not a newly assigned vacuum.

\subsection{A complete native cylinder remainder with a local spatial cost}

The following is the marked native estimate. Its proof verifies the new
boundary hypotheses before reusing the temporal analytic argument.

\begin{theorem}[Conditional leading covariance and a uniform temporal remainder]
\label{f16v11:native-cylinder}
For fixed $r,Q$ there are $K_{r,Q},\gamma_{r,Q}<\infty$ such that, if
\begin{equation}
 \gamma\ge\gamma_{r,Q},\qquad
             e^{K_{r,Q}M}/\sqrt\gamma\le1/100,
 \label{f16v11:local-entry}
\end{equation}
then every exterior history obeying \eqref{f16v11:windows} satisfies
\begin{equation}
 \left\|\Gamma^D_{\gamma,ij}(\mathbf q)
              -\gamma^{-1}(G^D_2(\mathbf q))_{ij}\right\|
 \le e^{K_{r,Q}M}\gamma^{-3/2}
                   3^{-(|i-j|-2)_+}.
 \label{f16v11:conditional-block-remainder}
\end{equation}
After enlarging $K_{r,Q}$ the same bound summed in the fixed temporal norm
is
\begin{align}
 \|\Gamma^D_\gamma(\mathbf q)-\gamma^{-1}G^D_2(\mathbf q)
                                     \|_{\mathrm{time},b}
      &\le e^{K_{r,Q}M}\gamma^{-3/2},\notag\\
 \|\mathsf M^D_\gamma(\mathbf q)-\gamma^{-1}M^D_2(\mathbf q)
                                     \|_{\mathrm{time},b}
      &\le e^{K_{r,Q}M}\gamma^{-3/2}.
 \label{f16v11:conditional-native-remainder}
\end{align}
In particular these two exact conditional matrices have temporal norm at
most $C_{r,Q}\gamma^{-1}$ in this regime. The ambient $B$ is unrestricted,
and all estimates hold for every $n\ge1$. The cost in $M$ is not removed.
\end{theorem}
\begin{proof}
We check the ingredients of V9--V10 for this cylinder.

\emph{Gaussian reference and endpoints.} The free Gaussian after port
integration has the internal precision $W_n^D$ in
\eqref{f16v11:separated-precision}, with bounded sources
\eqref{f16v11:boundary-sources}. Its diagonal interior entries are
$S_D+2N_D$, its adjacent time entries are $-N_D$, and its two endpoint
entries are the restrictions of V3's $E$. Center its \emph{whole} internal
history at $\mu^D=-(W_n^D)^{-1}\mathbf b^x$, and its moving ports at
$-L_D^{-1}b^z_t$. In unmoving port coordinates the centering adds
$\mathsf F_D(\mu^D_t-\mu^D_{t+1})$. All centers have a bounded norm per
cube by Proposition~\ref{f16v11:Gaussian-screening}.

Completing the Gaussian squares yields linear terms at the two ends of
each bond. The internal stationarity equations cancel consecutive terms;
the two endpoint equations cancel the remaining terms against the dressed
endpoint amplitudes. This is V9's telescoping with $W_n^D$ and the actual
boundary sources, not with the full mean restricted to $D$. After this
balancing, every Gaussian step is one centered operator whose internal
magnetic matrix is $S_D$ and kinetic precision is $N_D$. The integrated port
constant contains the full $\det L_D$; it is retained as part of that step.
Its normalized transfer has the Mehler split $P_D^*+Q_D^*$ with
$\|Q_D^*\|\le3-2\sqrt2<1/5$, by
$4I\le\mathsf M_D^{1/2}S_D\mathsf M_D^{1/2}\le10I$.
Here $P_D^*$ denotes the rank-one reference projection, not a port lift.

The reference endpoint is the normalized centered amplitude with precision
$\tfrac12((C_B^{-1})_{DD}+A_D^{\mathsf T}A_D)$. Its overlap with the
centered reference ground vector is at least $e^{-cM}$, since both have
fixed positive covariance bounds in dimension $15M$. The native left and
right endpoint amplitudes contain the actual crossing mixed terms at times
zero and $n$ and the original product-cube Gaussian carrier restricted to
$D$. Their norms, and any scalar factors used in balancing, have upper and
lower bounds $e^{\pm C_{r,Q}M}$: integrate on product unit-cube balls for
the lower bounds and use the pin for upper bounds. They can be normalized
for application of the product lemma, but their exact norms multiply the
conditional partition function and are not set to one. Frozen exterior
weights cancel from a normalized conditional, not from its exterior
posterior.

\emph{Complete compact comparison.} A one-, two-, or three-bond conditional
window contains $O(M)$ free real coordinates and $O(M)$ fixed exterior
inputs from the collar. By \eqref{f16v11:pin} its free integrations have a
Gaussian majorant after the centering and linear balancing:
\[
             e^{C_{r,Q}M}e^{-c|w|^2}.
\]
The use of Young's inequality to absorb the linear terms is valid because
the squared norms of their coefficient vectors are $O_{r,Q}(M)$.
Every conditional Wilson term is a member of the original finite ordered
word list, with some arguments fixed at bounded $q$ and $y$. Its fixed
real derivatives are uniformly bounded before rescaling. The integral
Taylor formula for $\varepsilon^{-2}f(\varepsilon w)$ therefore gives
polynomial derivative majorants of each fixed order, exactly as in V10.
There are no variable principal logarithms of exterior products to
differentiate: the exterior unmoving groups are the given
$\operatorname{Exp}(\varepsilon q)$ and their fixed-frame products.

Use V10's smooth proof cutoff on the \emph{free} principal charts only.
On its complement at least one centered free group coordinate has norm
$c/\varepsilon-C_{r,Q}$; the pin pays this tail by
$e^{C_{r,Q}M-c'/\varepsilon^2}$. On the cutoff domain the Haar square roots,
port Haar densities, and their fixed derivatives have polynomial bounds.
Taylor expansion through degree two, followed by complete free integration,
gives balanced step and endpoint remainders at most
$e^{C_{r,Q}M}\varepsilon^3$ in Hilbert--Schmidt or $L^2$ norm. The
zeroth arrays are the common Gaussian reference. Fixed-degree Gaussian
moments cost a polynomial in $M$ times a constant to the power $M$.
No integration over the ambient exterior occurs in this estimate.

\emph{Marked insertions.} The source observables are formed in the original conditional coordinates
before balancing. No derivatives of the history-dependent centers or
multipliers are omitted by differentiating a transformed formula.
The local port lift has a dimension-free operator
norm. The inserted-word remainder certificate of V7, on a window with fixed
exterior entries, thus gives for arbitrary spatial common-source vectors
$v$ a first remainder insertion with norm
$e^{C_{r,Q}M}\varepsilon\|v\|$. A product of two such insertions starts at
order two and costs $e^{C_{r,Q}M}\varepsilon^2\|v\|\|w\|$.
The first and product insertions have jets through degree two with
$e^{C_{r,Q}M}\varepsilon^3$ remainders. These assertions follow by
multiplying the same Gaussian majorant by the bounded-degree insertion
polynomials. All free port tails are integrated. Constants from exterior
fields are bounded by the fixed $Q$, not a supremum over unrestricted
compact boundary data. Only free Haar derivatives are used.

\emph{Exact long-product normalization.} The preceding arrays, on a small
complex polynomial-jet disc of radius $e^{-c_{r,Q}M}$, lie in V10's
analytic array ball around the common Gaussian operator and endpoint.
That proposition uses only the reference spectral ratio, the endpoint
overlap, and maximum array norms; all have just been verified with $M$
in place of $B$. Its exact scalar normalizers are kept. Its centered
two-window numerator retains the factor $3^{-(|i-j|-2)}$ for separated
windows. Cauchy estimates on that disc and the real short-jet errors
therefore give \eqref{f16v11:conditional-block-remainder}, with no factor
of the number of bonds. For overlapping windows the product insertion
occupies at most three bonds and gives the same estimate with factor one.
This is an application of the proved analytic product functional, not an
assumption of analyticity of the native measure at complex coupling.

At fixed finite $D,n$, the leading covariance coefficient is necessarily
\eqref{f16v11:leading-conditional}, by ordinary finite-dimensional
normalized expansion and uniqueness of coefficients. In particular a new
boundary normalization is not being silently substituted for the conditional
one. The entry condition absorbs the finitely many overlap, jet and
Cauchy-radius costs, all exponential in $M$. Summation is legitimate for
$b<\log3$ and gives \eqref{f16v11:conditional-native-remainder}.
The exact contact support gives the memory bound. Finally
\eqref{f16v11:conditional-coefficient-locality} and
$e^{K_{r,Q}M}\gamma^{-1/2}\le1/100$ give the stated $O(\gamma^{-1})$
conditional norms.
\end{proof}

\subsection{A native collar-screening theorem and a computable local approximation}

\begin{theorem}[The first source-local marked native remainder]
\label{f16v11:collar-native}
Under \eqref{f16v11:windows} and \eqref{f16v11:local-entry}, for any source
set $\mathcal A$ of collar depth $L$ as above,
\begin{equation}
 \boxed{
 \left\|\chi_{\mathcal A}
    [\mathsf M^D_\gamma(\mathbf q)-\gamma^{-1}M_2(\mathbf y)]
                                      \chi_{\mathcal A}\right\|_{\mathrm{time},b}
 \le C_{r,Q}\gamma^{-1}e^{-bL}
                   +e^{K_{r,Q}M}\gamma^{-3/2}.}
 \label{f16v11:collar-bound}
\end{equation}
The analogous estimate holds with $\Gamma^D_\gamma$ and $G_2$.
Thus, for any two admitted exterior fast histories $\mathbf q,\mathbf q'$,
\begin{equation}
 \left\|\chi_{\mathcal A}
    [\mathsf M^D_\gamma(\mathbf q)-\mathsf M^D_\gamma(\mathbf q')]
                                      \chi_{\mathcal A}\right\|_{\mathrm{time},b}
 \le 2C_{r,Q}\gamma^{-1}e^{-bL}
                   +2e^{K_{r,Q}M}\gamma^{-3/2}.
 \label{f16v11:two-boundary-screening}
\end{equation}
All durations and ambient volumes are allowed. No native interior cutoff,
stationarity assumption, or probability bound for the exterior is used.
\end{theorem}
\begin{proof}
Insert the conditional leading coefficient $\gamma^{-1}M^D_2(\mathbf q)$.
The error to the native conditional is
\eqref{f16v11:conditional-native-remainder}; its comparison to the compressed
full coefficient is \eqref{f16v11:coefficient-screening-bound}. Compression
by the same spatial projection at every time cannot increase the time-block
norm. This proves \eqref{f16v11:collar-bound}. Compare both boundary histories
through the same $M_2(\mathbf y)$ to get
\eqref{f16v11:two-boundary-screening}. The proof for the complete residual
Gram is identical. The two errors have different meanings: the first is
boundary-reaching coefficient error, and the second is the complete native
remainder on the declared cylinder.
\end{proof}

For an explicit local recipe one can choose $\mathbf q=0$, compute
$C^D,m^D,\mathsf P_D$ and the literal cubic word tensors on $D$, then use
\eqref{f16v11:leading-conditional}. This gives $M^D_2(0)$ without knowing
fast fields beyond the collar. The same proof gives
\begin{equation}
 \left\|\chi_{\mathcal A}
 [\mathsf M^D_\gamma(\mathbf q)-\gamma^{-1}M^D_2(0)]
                              \chi_{\mathcal A}\right\|_{\mathrm{time},b}
 \le C_{r,Q}\gamma^{-1}e^{-bL}+e^{K_{r,Q}M}\gamma^{-3/2}.
 \label{f16v11:local-computation}
\end{equation}
The zero boundary is an \emph{approximation used for this coefficient};
no zero boundary is imposed on the native law on the left. The relevant
bridge history is still retained wherever it enters the cylinder words.

For example place a fixed source region inside a cube collar of spatial
radius $L+r_0$. There is a geometric $c_0$ such that $M\le c_0(L+r_0+1)^3$.
Choose $L=L_\gamma\to\infty$ slowly enough that
\begin{equation}
       K_{r,Q}c_0(L_\gamma+r_0+1)^3\le\tfrac14\log\gamma.
 \label{f16v11:collar-scale}
\end{equation}
For sufficiently large $\gamma$, the local entry condition holds and
\eqref{f16v11:collar-bound} is at most
\begin{equation}
 C_{r,Q}\gamma^{-1}e^{-bL_\gamma}+\gamma^{-5/4}.
 \label{f16v11:screening-scale}
\end{equation}
One may take $L_\gamma$ of order $(\log\gamma)^{1/3}$ with a sufficiently
small fixed coefficient. The ambient torus may be arbitrarily larger than
this collar. This is not an assertion for a collar at unrestricted $M$
with fixed coupling, and the first term is not a fixed algebraic improvement
in $\gamma$. On a torus too small to contain a proper such collar, the
previous full-torus results remain the applicable statements.

\subsection{Conditional source means and the return needed for actual gluing}

The covariance estimate alone is not enough to glue cylinders. There is
a useful additional mean estimate with the same boundary convention.
Let $a_{1,\alpha}(\mathbf y)=\mathbb E_0r_{\alpha,1}$ be the full Gaussian
first residual coefficient of V10, and set
\[
 j^D_{\gamma,\alpha}(\mathbf q)
       =\partial_\alpha[-\log Z_{\gamma,D}(\mathbf q;\mathbf y)],
 \qquad j^{\rm ref}_\alpha=g_\alpha+\gamma^{-1/2}a_{1,\alpha}.
\]
The reference value is independent of the exterior fast history.

\begin{proposition}[Native conditional mean message on the admitted collar]
\label{f16v11:mean-message}
After enlargement of the same constants and, if necessary, decrease of $b$,
each unit canonical source $\alpha$ of depth at least $L$ obeys
\begin{equation}
 |j^D_{\gamma,\alpha}(\mathbf q)-j^{\rm ref}_\alpha|
       \le C_{r,Q}e^{-bL}+e^{K_{r,Q}M}\gamma^{-1}.
 \label{f16v11:mean-message-bound}
\end{equation}
This bound is independent of duration and ambient volume, on the same
admitted boundary histories. It is not a bound on their probability.
\end{proposition}
\begin{proof}
The exact Haar identity gives
$j^D_{\gamma,\alpha}=g^D_\alpha+\mathbb E_Dr^D_{\gamma,\alpha}$.
The one-window version of the conditional short-jet and analytic-array
argument in Theorem~\ref{f16v11:native-cylinder}, now expanded only through
order one, gives
\[
 \left|\mathbb E_Dr^D_{\gamma,\alpha}
      -\gamma^{-1/2}\mathbb E_{0,D}^{\mathbf q}r^D_{\alpha,1}\right|
                           \le e^{K_{r,Q}M}\gamma^{-1}.
\]
There is no measure-correction term at this order because the order-zero
residual is identically zero. At higher orders such terms would be required.
The short insertion norm is uniform in its time location, including either
dressed end. The one-window normalized functional has the same
length-independent analytic domain used above.

The Gaussian scalar difference is \eqref{f16v11:g-screen}. To compare the
two first residual means, use the same quadratic polynomial and inverse
boundary identities as in the proof of
Theorem~\ref{f16v11:coefficient-screening}. There is now one source anchor
and a marked boundary visit; a quadratic mean consists only of its constant
term and one self-contraction. The mean, covariance and lift differences
all carry a boundary path, so
$|\mathbb E_{0,D}^{\mathbf q}r^D_{\alpha,1}-\mathbb E_0r_{\alpha,1}|
\le C_{r,Q}e^{-bL}$. This also covers a primitive term synthesized from a
boundary port. Combining the three bounds proves
\eqref{f16v11:mean-message-bound}. No derivative of a pressure-error bound
has been taken.
\end{proof}

For clarity about the remaining task, the exact exterior law at fixed
bridge history is
\begin{equation}
 d\bar\mu_\gamma(\mathbf q)
 =\frac{1_{\mathcal P'_O}J^f_{\gamma,O}(\mathbf q)
         e^{-U_{\gamma,O}(\mathbf q;\mathbf y)}
         Z_{\gamma,D}(\mathbf q;\mathbf y)}{Z_\gamma(\mathbf y)}\,d\mathbf q.
 \label{f16v11:posterior}
\end{equation}
In this formula \emph{all} exterior coordinates are integrated; its
projection to the word collar gives the same conditional messages. For two
interior marks at nonadjacent times, normalized differentiation or the
inherited total-covariance identity gives
\begin{equation}
 \boxed{
 \partial_\alpha\partial_\beta[-\log Z_\gamma]
 =-\mathbb E_{\bar\mu_\gamma}
          [\Gamma^D_{\gamma,\alpha\beta}(\mathbf q)]
   -\operatorname{Cov}_{\bar\mu_\gamma}
          (j^D_{\gamma,\alpha},j^D_{\gamma,\beta}).}
 \label{f16v11:actual-gluing}
\end{equation}
Indeed $U_{\gamma,O}$ has no derivative in either interior source direction;
therefore $\partial_\alpha[-\log Z_\gamma]=
\mathbb E_{\bar\mu_\gamma}j^D_{\gamma,\alpha}$. Differentiating again gives
the conditional Hessian minus the displayed covariance of conditional
means. The conditional Hessian at these times is exactly
$-\Gamma^D_{\gamma,\alpha\beta}$. At finite regulator all derivatives and
integrals are legitimate. The same identity holds for the full history
weight at nonadjacent times because its external logarithmic factors are
single-time. It is one decomposition of the original return, not an
additional debit to subtract from a previously integrated Hessian.

Let $\mathsf G_Q$ be the event in \eqref{f16v11:windows} for the entire
exterior collar history. The estimates in this appendix hold on
$\mathsf G_Q$. Neither $\bar\mu_\gamma(\mathsf G_Q^c)$ nor a source-weighted
moment on that event has been bounded here. Moreover even on the good
event, a small supremum for each scalar mean message does not itself give
a small operator norm or a summable temporal row for the covariance of
\emph{all} messages. A coherent boundary direction could remain.

For a quantitative ledger set
$d_\alpha(\mathbf q)=j^D_{\gamma,\alpha}(\mathbf q)-j^{\rm ref}_\alpha$ and
$\delta_L=C_{r,Q}e^{-bL}+e^{K_{r,Q}M}/\gamma$. The exact posterior obeys
\begin{equation}
 |\operatorname{Cov}_{\bar\mu_\gamma}(j^D_{\gamma,\alpha},j^D_{\gamma,\beta})|
 \le
 \sqrt{(\delta_L^2+\mathbb E[1_{\mathsf G_Q^c}|d_\alpha|^2])
       (\delta_L^2+\mathbb E[1_{\mathsf G_Q^c}|d_\beta|^2])}.
 \label{f16v11:bad-message-ledger}
\end{equation}
This is Cauchy--Schwarz after subtracting the deterministic reference and
splitting the second moments. It exposes, but does not estimate, the
exceptional source moments. The first term in
\eqref{f16v11:actual-gluing} has its own exceptional conditional-covariance
integral. Existing expected bad-label densities cannot be substituted for
these quantities, especially on an arbitrarily long cylinder.

The next noncircular target is consequently the \emph{actual exterior
message law} in \eqref{f16v11:posterior}: either a source-weighted collar
escape estimate plus a coherent boundary-message contraction, or a
multiscale decomposition that controls their combined connected return.
The compact conditional interior is no longer an unspecified input: its
marked remainder and leading coefficient now have the explicit local-size
cost in \eqref{f16v11:collar-bound}. Reproving a global temporal product
bound, or replacing the posterior by a bounded deterministic boundary,
would miss the remaining term. No physical Gauss restriction, bridge-window
coverage, infrared contact coercivity or same-top physical transfer
comparison is obtained from this conditioning argument.

\clearpage
\subsection{Diagnostics, preservation, and scope}

The accompanying diagnostic is
\begin{center}\small
\path{verify_ym_f16_v11_conditional_collar.py}.
\end{center}
It checks noncommuting Gaussian conditioning and covariance differences,
principal versus marginal inverse distinctions, restricted port-lift
cancellation with nonzero exterior data, endpoint weights and centering,
actual parity-cube crossing carriers, leading Gaussian polynomial
covariances, finite-range boundary-reaching inverse paths, and the exact
posterior covariance decomposition. All 463 assertions passed; their family counts and the script hash are in
\path{ym_v11_verification_report.json}. The unchanged V4--V10 scripts were
also rerun and passed 9,467, 89,339, 906, 311, 609, 990 and 608 assertions,
respectively. These finite checks do not certify analytic derivative
suprema, complete-chart Taylor remainders, or infinite-dimensional operator
estimates with a proof assistant. The theorems above rely on their analytic
proofs and on the explicitly identified inherited inputs.

The supplied V10 source is preserved byte for byte except for the cumulative
title version and insertion of this exact appendix before its final
\verb|\end{document}|. Removing this append and restoring the title recovers
V10 exactly. The package records original-source hashes, predecessor
recovery, compilation and render checks. Only diagnostics actually available
in the supplied packages are rerun and reported; no inaccessible historical
checker or NS Lean development is described as verified. The next scheduled
five-version companion checkpoint remains V15.

\begin{firewall}
V11 proves complete-native conditional cylinder estimates on a declared
bounded exterior fast-history window. The exponential comparison cost is
in the cylinder size $M$, not the ambient volume $B$, and no duration cost
is introduced. Boundary-marked coefficient screening yields a native
conditional memory approximation with error
$C_{r,Q}\gamma^{-1}e^{-bL}+e^{K_{r,Q}M}\gamma^{-3/2}$, and a corresponding
conditional mean-message estimate. All crossing interactions, free compact
tails and dressed endpoints are retained. No typicality of the admitted
collar, source-weighted control of its complement, coherent contraction of
the actual exterior-message covariance, unrestricted full-native spatial
gluing, physical transfer gap or interacting continuum Yang--Mills mass gap
is established by this appendix.
\end{firewall}
% END F16 V11 APPEND

% BEGIN F16 V12 APPEND -- 2026-09-17
% Insert immediately before the final \end{document} in cumulative V11.
% The predecessor body is unchanged; only its title version becomes V12.

\clearpage
\section{V12: posterior return without an all-good collar}
\label{f16v12:context}

V11 controls an interior cylinder on a bounded exterior-history window.
Requiring that window at every time is not a useful route to an estimate
uniform in duration: its complement can have large probability even when
local fluctuations are well controlled. We instead estimate the \emph{actual
posterior mean vector}, integrating all exterior fields. The two inputs are
V6's full-native relative entropy and the finite-range Markov specification
shared by the native law and its Gaussian reference. Conditional Gaussian
concentration is applied to a whole source vector, before summing over
cylinders. Applying it separately to every source would lose the duration.

The output has a deliberately different norm from V11's conditional theorem.
The posterior return, including its exceptional-field part, has a small
\emph{trace per space--time cell}; no uniform small operator norm is inferred.
We also match the conditional and full Haar lifts. This pays the exceptional
conditional-covariance integrals in the same density sense, without a
bounded-collar event. The final gluing identity keeps both terms under their
one actual exterior posterior. Nothing here promotes the conditional
$\gamma^{-1}$ expansion to an unrestricted-volume expansion.

\subsection{Dependencies and the upstream choice}

The supplied V11 appendix was read in full. The available f14/f15 archives
and Grand Tensor Monograph V076 were inventoried and relevant passages
examined; their entire proof collections have not been independently
recertified. F14 V67 already estimates entropy response for a different
restricted high conditional; f15 V1 already turns exponential-source control
into mean estimates, and f15 V13--V14 and V68 retain actual posterior weights
in other conditional comparisons. The monograph's Markov-blanket and
conditional-mean identities are inherited. None is presented as a new
abstract principle.

The model-specific input reused here is V6's entropy estimate for the
\emph{complete augmented fast-history law}, not entropy of a restricted core.
We also reuse V6's Gaussian marginal-entropy localization, V7's Gaussian
vector-energy inequality and exact Haar sources, and V11's literal spatial
conditioning and restricted port inverse. The new specialization combines
these inputs on entire-time spatial cylinders, with no exterior-field
cutoff, and keeps the source-vector and overlap costs explicit.
The import memorandum's separation of power gain from locality cost remains
binding: this is a change in how a remainder is measured, not another formal
order or an imported Navier--Stokes theorem.

For attribution, the Gaussian concentration used below follows from the
classical logarithmic Sobolev inequality of L.~Gross,
\href{https://doi.org/10.2307/2373688}{\emph{Logarithmic Sobolev Inequalities}},
Amer. J. Math. 97 (1975), 1061--1083. V6's marginal-entropy argument is an
instance of the entropy/Fisher-information method discussed by E.~A.~Carlen
and D.~Cordero--Erausquin,
\href{https://arxiv.org/abs/0710.0870}{\emph{Subadditivity of the entropy and
its relation to Brascamp--Lieb type inequalities}}. The precise Gaussian
consequences needed here are derived below or explicitly reused. No native
logarithmic Sobolev inequality is assumed.

\subsection{One full law, whole source vectors, and an exact local entropy cost}

Fix the same bounded bridge history as in V6--V11 and write
\[
 \mathcal N=B(n+1),\qquad
 \epsilon_\gamma=\gamma^{-1/6}(1+\log\gamma)^{2/3},\qquad
 \max_{i,e}|y_{i,e}|\le r<\infty.
\]
Here $B=m^3$, $m\ge4$, and $n\ge1$. On the same complete unmoving Euclidean
coordinate space $\zeta=(x,a)$ let
\begin{equation}
 P=\mu_{\gamma,\mathbf y},\qquad Q=\mu_{0,\mathbf y}=N(m,H^{-1}),
 \qquad h_-I\le H\le h_+I,\qquad
 D(P\Vert Q)\le C_r\mathcal N\epsilon_\gamma.
 \label{f16v12:inherited}
\end{equation}
The native density is extended by zero outside its complete principal charts,
as in V6. The threshold for \eqref{f16v12:inherited} depends on $r$, not on
$B,n$. No hypothesis $e^{KB}/\sqrt\gamma\ll1$ is used in this appendix.
The symbol $Q$ denotes the Gaussian probability, not a bound on exterior
coordinates. All exterior coordinates are integrated under $P$.

For a proper nonempty spatial cylinder $D$, let $I_D$ be all its free
space--time coordinate cells as in V11. Choose a fixed-width exterior
collar $C_D$ containing every exterior cell appearing in an action term
that meets $I_D$, and every nonzero block of $H_{I_D O_D}$. Write
$T_D=I_D\cup C_D$ and $O_D=I_D^c$. These are sets of \emph{cells}, with all
real components and all available times included. The width can be fixed
from the original finite word list, independently of $D,B,n,\gamma$.

Let $\mathcal A_D$ be a chosen bank of distinct canonical common-bridge
marks whose ordinary score derivative carriers are inside $I_D$. Define
\begin{equation}
 F_D(\zeta_{I_D})=(J_{\gamma,\alpha})_{\alpha\in\mathcal A_D},\qquad
 J_{\gamma,\alpha}=\partial_\alpha U_\gamma,\qquad
 p_D=|\mathcal A_D|.
 \label{f16v12:source-vector}
\end{equation}
These are the \emph{ordinary} source scores, not the compensated residuals.
Their Gaussian comparison below is the analytic exponential-word extension
of this same $F_D$, not a substitution of $J_0$ for $J_\gamma$.
Choose $L\ge1$ such that the fast derivative support of this bank has
spatial distance at least $L$, after a fixed word-range allowance, from the
interior cells coupled by $H$ to $C_D$. Source anchors of depth $L+2r_0$
with V11's sufficiently large $r_0$ meet this condition.

The exact native and reference messages, and a deterministic centering, are
\begin{equation}
 j_D(q)=\mathbb E_P[F_D\mid\zeta_{O_D}=q],\qquad
 f_D(q)=\mathbb E_Q[F_D\mid\zeta_{O_D}=q],\qquad
 a_D=\mathbb E_Q F_D.
 \label{f16v12:messages}
\end{equation}
In particular $j_D$ is V11's vector of derivatives of $-\log Z_{\gamma,D}$.
No interior source differentiates $U_{\gamma,O}$. The vectors $a_D$ merely
center estimates; their values are not chosen as a new vacuum or normalizer.

\begin{proposition}[Whole-bank Lipschitz bound and the correct conditional entropy]
\label{f16v12:local-input}
There is a fixed geometric $L_0<\infty$ such that
\begin{equation}
             \sup_\zeta\|D_\zeta F_D\|_{\rm op}\le L_0
 \label{f16v12:source-Lip}
\end{equation}
for every such bank, every $B,n,D$ and every $\gamma\ge1$. Both messages in
\eqref{f16v12:messages} depend on the exterior only through $q_{C_D}$.
Put
\begin{equation}
 d_T=D(P_{T_D}\Vert Q_{T_D}),\quad
 d_C=D(P_{C_D}\Vert Q_{C_D}),\quad k_D=d_T-d_C.
 \label{f16v12:local-entropies}
\end{equation}
Then
\begin{equation}
 k_D=\mathbb E_{P_{C_D}}
 D(P_{I_D\mid C_D}\Vert Q_{I_D\mid C_D})\ge0.
 \label{f16v12:chain}
\end{equation}
These are finite quantities at every regulator where
\eqref{f16v12:inherited} holds. The conditional divergence in this identity
is averaged against the \emph{native} collar marginal.
\end{proposition}
\begin{proof}
For a primitive action term $\gamma s(\prod_j\operatorname{Exp}
(T_jw/\sqrt\gamma))$, every second derivative in its real inputs is
bounded by a fixed constant on all real inputs. Duhamel differentiation
inserts fixed coefficient matrices into products of unitary matrices; the
two factors $\gamma^{-1/2}$ cancel the prefactor $\gamma$. Word lengths and
linear coefficients are fixed. Taking one bridge derivative gives a row of
the score Jacobian supported on a fixed number of fast coordinates. Each
fast coordinate appears in a fixed number of such rows, including the
four-bridge normalization of a canonical common mark. The block row and
column Schur bounds prove \eqref{f16v12:source-Lip}. Taking a subbank cannot
increase the operator norm. Endpoint Gaussian terms are independent of the
bridge marks; the fast Haar factors have no bridge dependence. This is the
whole-vector version of V6's local score certificate. It is not the sum of
$p_D$ scalar Lipschitz constants.

Native product-coordinate measure and finite-range action imply that the
conditional density on $I_D$ given $O_D$ depends only on $C_D$: exterior-only
terms cancel on normalization. The complete principal-chart indicator and
all Haar powers factor by cells, so they preserve this statement. The
Gaussian has the same property because $H_{I_D,O_D\setminus C_D}=0$.
Conditioning first on $C_D$ therefore gives the same interior kernels for
both laws. The log-density chain rule on $T_D=I_D\sqcup C_D$ gives
\eqref{f16v12:chain}; data processing from \eqref{f16v12:inherited} makes
$d_T,d_C$ finite. This argument does not invoke a sum of native conditional
entropies bounded by a full entropy without first identifying these patches.
\end{proof}

\subsection{Gaussian mean comparison without an output-dimension loss}

We record a standard Gaussian consequence in the needed vector form.

\begin{proposition}[A vector mean bound from one conditional entropy]
\label{f16v12:vector-mean}
Let $Q'$ be a finite-dimensional Gaussian of covariance at most $cI$, and
let $F$ take values in $\mathbb R^p$ with $\|DF\|\le L_F$. For any
$P'\ll Q'$ of finite relative entropy,
\begin{equation}
       \|\mathbb E_{P'}F-\mathbb E_{Q'}F\|^2
                     \le2cL_F^2D(P'\Vert Q').
 \label{f16v12:mean-transport}
\end{equation}
The dimension $p$ does not multiply the right side. For the messages above,
\begin{equation}
             \mathbb E_{P_{C_D}}\|j_D-f_D\|^2
                           \le2h_-^{-1}L_0^2 k_D.
 \label{f16v12:conditional-mean-error}
\end{equation}
\end{proposition}
\begin{proof}
For a unit vector $v$, the scalar $v^{\mathsf T}F$ is $L_F$-Lipschitz.
The Gaussian logarithmic Sobolev inequality in the convention used in V7
and its elementary Herbst integration give
\[
 \log\mathbb E_{Q'}\exp\{t(v^{\mathsf T}F-\mathbb E_{Q'}v^{\mathsf T}F)\}
                       \le ct^2L_F^2/2,\qquad t\in\mathbb R.
\]
For clarity, if $\psi(t)$ is this logarithmic moment generating function,
log-Sobolev gives $t\psi'(t)-\psi(t)\le ct^2L_F^2/2$; integrate
$(\psi(t)/t)'$ from zero, and use both signs. Gaussian integrability of
Lipschitz functions justifies this calculation by truncation.
The entropy variational inequality bounds a signed mean difference by
$D(P'\Vert Q')/t+ctL_F^2/2$. Optimize $t>0$. Choose $v$ in the direction
of the vector mean difference, rather than sum the scalar inequalities
for a basis. This proves \eqref{f16v12:mean-transport}.

The Gaussian interior conditional has covariance $(H_{I_DI_D})^{-1}
\preceq h_-^{-1}I$, independent of $q$. Its mean may be arbitrarily large;
the preceding inequality is translation independent. Apply it at each
native exterior value $q$, to the same analytic score bank $F_D$, and then
integrate the conditional divergence using \eqref{f16v12:chain}. The native
conditional is supported on its full compact chart; absolute continuity
with respect to the everywhere-positive Gaussian is sufficient. No exterior
amplitude bound is required.
\end{proof}

\begin{proposition}[The reference message is exponentially insensitive to its collar]
\label{f16v12:reference-message}
There are $C,\nu>0$, depending only on the fixed geometry and Gaussian
spectral bounds, such that
\begin{equation}
 \operatorname{Lip}(f_D)\le C e^{-\nu L},\qquad
 \mathbb E_{Q_{C_D}}\|f_D-a_D\|^2\le C p_D e^{-2\nu L}.
 \label{f16v12:Gaussian-message}
\end{equation}
Moreover under the \emph{native} collar marginal,
\begin{equation}
 \mathbb E_{P_{C_D}}\|f_D-a_D\|^2
                          \le C e^{-2\nu L}(p_D+d_C).
 \label{f16v12:native-reference-message}
\end{equation}
All times and arbitrary exterior amplitudes are included. The constants
contain neither the area of the boundary nor the duration.
\end{proposition}
\begin{proof}
Let $I=I_D$, $C_\partial=C_D$, and let $S_D^{\rm src}\subset I$ be the fast
derivative support of $F_D$. Gaussian conditioning realizes the free
coordinates as
\[
       m_I-H_{II}^{-1}H_{IC_\partial}(q-m_{C_\partial})+g,
             \qquad g\sim N(0,H_{II}^{-1}).
\]
Its covariance is independent of $q$. Differentiate this expectation, using
the global derivative bound of $F_D$. Only rows of $H_{II}^{-1}H_{IC_\partial}$
in $S_D^{\rm src}$ can contribute.

The nonzero rows of $H_{IC_\partial}$ lie in a fixed-width interior boundary
set $T$. A Neumann expansion of $H_{II}^{-1}$ has norm ratio
$\theta=1-h_-/h_+<1$ and finite spatial range per power. Before a number of
powers proportional to $L$, no term connects $S_D^{\rm src}$ to $T$.
Summing the remaining tail proves
\[
 \|P_{S_D^{\rm src}}H_{II}^{-1}H_{IC_\partial}\|
                                      \le C e^{-\nu L}.
\]
Here $\|H_{IC_\partial}\|\le h_+$; a set-to-set operator-norm estimate, not
a sum over all boundary cells, is being used. Powers can travel in time,
but cannot cross a spatial distance $L$ in fewer spatially ranged steps.
This proves the first bound in \eqref{f16v12:Gaussian-message} for every
$q$. Differentiation is justified by the Lipschitz majorant under the
Gaussian, without a bounded-mean hypothesis.

The full Gaussian marginal $Q_{C_D}$ has covariance at most $h_-^{-1}I$.
Gaussian Poincare applied componentwise, or its vector trace form, gives
\[
 \mathbb E_{Q_{C_D}}\|f_D-\mathbb E_{Q_{C_D}}f_D\|^2
       \le h_-^{-1}\mathbb E\|Df_D\|_{\rm HS}^2
       \le h_-^{-1}p_D\operatorname{Lip}(f_D)^2.
\]
The tower property identifies the mean as $a_D$.

It remains to change the \emph{collar marginal}, rather than quietly
replace it by the reference marginal. Set $R(q)=f_D(q)-a_D$ and
$\lambda_L=Ce^{-\nu L}$, chosen as an upper bound on its Lipschitz norm.
V7's Gaussian vector-energy inequality gives, for
$t_L=(4h_-^{-1}\lambda_L^2)^{-1}$,
\[
  \log\mathbb E_{Q_{C_D}}e^{t_L\|R\|^2}
                           \le2t_L\mathbb E_{Q_{C_D}}\|R\|^2.
\]
The entropy variational inequality therefore yields
\[
 \mathbb E_{P_{C_D}}\|R\|^2
 \le4h_-^{-1}\lambda_L^2d_C+2\mathbb E_{Q_{C_D}}\|R\|^2.
\]
This proves \eqref{f16v12:native-reference-message}. A constant message is
handled directly. The exponentially large allowed $t_L$ is legitimate:
it is paired with an exponentially \emph{small Lipschitz constant}, not
with an asserted Gaussian bound for the native measure.
\end{proof}

\subsection{An actual-posterior theorem, with no all-good event}

\begin{theorem}[Unrestricted posterior mean screening on a cylinder family]
\label{f16v12:posterior-screening}
For each cylinder and source bank just specified,
\begin{equation}
 \mathbb E_P\|j_D-a_D\|^2
       \le C\{k_D+e^{-2\nu L}(p_D+d_C)\}.
 \label{f16v12:single-message}
\end{equation}
The expectation uses the actual exterior posterior of the full native law.
Let $(D_s,\mathcal A_s)$ be any finite family whose enlarged cell sets
$T_s=I_{D_s}\cup C_{D_s}$ have overlap at most $q_0$, and whose derivative
supports all have the above depth $L$. Source occurrences are counted with
multiplicity. Put $p_{\rm tot}=\sum_s|\mathcal A_s|$ and stack
$\mathbf j=(j_{D_s})_s$, $\mathbf a=(a_{D_s})_s$, viewing every message as a
function on the \emph{same full probability space}. Then
\begin{equation}
 \boxed{
  \mathbb E_P\|\mathbf j-\mathbf a\|^2
   \le C_r q_0\mathcal N\epsilon_\gamma
                          +C p_{\rm tot}e^{-2\nu L}.}
 \label{f16v12:family-message}
\end{equation}
Consequently its actual posterior-return Gram obeys
\begin{equation}
 \boxed{
  \operatorname{Tr}\operatorname{Cov}_P(\mathbf j)
   \le C_r q_0\mathcal N\epsilon_\gamma
                          +C p_{\rm tot}e^{-2\nu L}.}
 \label{f16v12:return-trace}
\end{equation}
There is no fixed bound on $|q_c|$, no restriction on $|D_s|$ relative to
$\gamma$, and no duration bound. The result is not a supremum over exterior
histories, nor a covariance-operator bound independent of $\mathcal N$.
\end{theorem}
\begin{proof}
Split $j_D-a_D=(j_D-f_D)+(f_D-a_D)$ and use
$\|u+v\|^2\le2\|u\|^2+2\|v\|^2$, the conditional mean comparison and
\eqref{f16v12:native-reference-message}. This proves
\eqref{f16v12:single-message}. The precise first cost is $k_D=d_T-d_C$,
not $d_T$ counted a second time as a boundary divergence.

For the family, V6's Gaussian marginal-entropy localization gives
\[
 \sum_sD(P_{T_s}\Vert Q_{T_s})
      \le q_0\kappa D(P\Vert Q),\qquad \kappa=h_+/h_-.
\]
Its constant depends on overlap and the full Gaussian condition number,
not on marginal dimension. The sets $T_s$ can therefore include every time
of a long cylinder. Both $k_{D_s}$ and $d_{C_s}$ are at most $d_{T_s}$.
Sum \eqref{f16v12:single-message} and insert
\eqref{f16v12:inherited}. This proves \eqref{f16v12:family-message}.
Centering by the actual mean minimizes squared error, so
$\operatorname{Tr}\operatorname{Cov}_P(\mathbf j)
\le\mathbb E_P\|\mathbf j-\mathbf a\|^2$. Cross-covariances among
messages are retained in this Gram; independence is not assumed.
\end{proof}

The cancellation of the duration cost is substantive. A scalar application
of \eqref{f16v12:mean-transport} at every time would repeatedly charge the
same cylinder entropy, potentially multiplying it by $n+1$. Instead a
\emph{single} vector comparison pays each conditional entropy once, and the
bounded-overlap marginal theorem then distributes the already known full
entropy. This is not approximate tensorization of native conditional
entropies, and it is not a native functional inequality.

\begin{proposition}[The old exceptional mean term is included]
\label{f16v12:exceptional-means}
Use V11's deterministic reference
$j^{\rm ref}_\alpha=g_\alpha+\gamma^{-1/2}\mathbb E_0r_{\alpha,1}$ on every
chosen source bank. For any exterior events $E_s$, with no probability
assumption,
\begin{equation}
 \sum_s\mathbb E_P\!
       \left[1_{E_s}\|j_{D_s}-j^{\rm ref}_{\mathcal A_s}\|^2\right]
 \le C_rq_0\mathcal N\epsilon_\gamma
             +C_rp_{\rm tot}(e^{-2\nu L}+\gamma^{-1}).
 \label{f16v12:exceptional-mean}
\end{equation}
This applies in particular to failure of V11's entire-history bounded-collar
event. It bounds source-weighted escape in density, not escape probability.
\end{proposition}
\begin{proof}
For each ordinary score, V6's global Taylor certificate and bounded local
Gaussian moments give
$|\mathbb E_Q J_{\gamma,\alpha}-\mathbb E_QJ_{0,\alpha}|
\le C_r\gamma^{-1/2}$. V7's Gaussian Haar cancellation and Gaussian
integration by parts give $\mathbb E_QJ_{0,\alpha}=g_\alpha$.
The first residual coefficient has uniformly bounded Gaussian mean by its
quadratic local word representation and the summable port-lift row.
Thus $|a_\alpha-j^{\rm ref}_\alpha|\le C_r\gamma^{-1/2}$, uniformly
in the source location and duration. Sum its squared difference over the
bank. Dropping $1_{E_s}\le1$ in the squared-error upper bound and using
\eqref{f16v12:family-message} proves the claim. This argument never bounds
the probability of an all-time event by a union bound.
\end{proof}

\subsection{Geometry with bounded overlap at growing collar depth}

The overlap need not grow with the collar radius. Here is one explicit
construction when the torus is large enough to contain local cylinders.
Let $\rho$ be a fixed integer larger than all spatial word and derivative
ranges, and put $\ell=\lceil L+4\rho\rceil$. Suppose $m\ge16\ell$.
Partition each coordinate cycle into $k=\lfloor m/\ell\rfloor$ consecutive
intervals, distributing their lengths as equally as possible. Every length
is between $\ell$ and $2\ell$. Their products are disjoint core boxes.
For each core box enlarge by $L+2\rho$ in each spatial coordinate to obtain
$D_s$, and by a further $\rho$ to obtain a collar patch containing $T_s$.
Assign each canonical spatial common mark to its core box, at every time.

Every source derivative carrier then has the required depth, and
$p_{\rm tot}=9\mathcal N$. An enlarged coordinate interval has length at
most $2\ell+2L+6\rho\le4\ell$. Its starting core positions are separated
by at least $\ell$. Thus a point can belong to at most eight enlarged
intervals in one coordinate, a safe count also on the cycle. Hence
\begin{equation}
                         q_0\le8^3=512.
 \label{f16v12:cover-overlap}
\end{equation}
There is no extra temporal overlap factor: each spatial patch already
contains all available times. These patches overlap; their messages are
not conditionally independent by virtue of this cover. The next gluing
theorem imposes separation when independence is used.

With this cover, \eqref{f16v12:return-trace} yields
\begin{equation}
 \frac{\operatorname{Tr}\operatorname{Cov}_P(\mathbf j)}{\mathcal N}
               \le C_r(\epsilon_\gamma+e^{-2\nu L}).
 \label{f16v12:cover-rate}
\end{equation}
Taking $L\ge(2\nu)^{-1}\log(\epsilon_\gamma^{-1})$ makes the bound
$C_r\epsilon_\gamma$. This radius is $O(\log\gamma)$ and is \emph{not}
subject to V11's $e^{KM}/\sqrt\gamma$ entry condition, because no conditional
kernel expansion was used here. It must still fit geometrically in the
ambient torus. This does not extend V11's coefficient approximation to
cylinders of that size. On smaller tori the general family theorem remains
valid for whatever proper cylinders and depths actually fit; no nonexistent
collar is assumed.

\subsection{The conditional residuals also have an unrestricted posterior budget}

The exceptional conditional-covariance term in V11 can be paid in density
without applying its bounded-collar expansion. The point is to compare
\emph{compatible Haar fields}, rather than compare two differently centered
scores as unrelated observables.

For a common source $\alpha$ assigned to $D_s$, let $\eta_\alpha$ be its
full port coefficient column from V7 and let $\eta^D_\alpha$ be V11's
restricted coefficient, extended by zero to exterior ports. Write
$\Delta_\alpha=\eta_\alpha-\eta^D_\alpha$. At either adjacent bond put
$L=L_h$ only within the following block identities; here it denotes the
port matrix, not the collar depth. Since $E_O^{\mathsf T}v=0$ for an
interior source,
\begin{equation}
 \Delta_{\alpha,I}=-L_D^{-1}L_{DO}\eta_{\alpha,O},\qquad
 \Delta_{\alpha,O}=\eta_{\alpha,O},\qquad
                    (L\Delta_\alpha)_I=0.
 \label{f16v12:harmonic-lift}
\end{equation}
Let $e^{\rm pt}=S_\gamma-S_0$ be V7's primitive port-error vector, with
$S_\gamma=(X^\gamma_{p,k}U_\gamma)_{p,k}$ and $S_0=\partial_aU_0$.
Every evaluation uses the full configuration, with exterior values inserted
in a conditional expression when needed.

\begin{proposition}[Exact lift matching and an exponentially small mismatch]
\label{f16v12:lift-matching}
For every exterior chart value, not only bounded values,
\begin{equation}
           r^D_{\gamma,\alpha}
                =r^{[\infty]}_{\gamma,\alpha}
                                  -\Delta_\alpha^{\mathsf T}e^{\rm pt}.
 \label{f16v12:exact-residual-match}
\end{equation}
For a family with patch overlap at most $q_0$ and source depth at least $L$,
let $\Delta$ have all these columns, counting repeated marks as separate
columns. There are fixed $C,\nu'>0$ such that
\begin{equation}
               \|\Delta\|_{\rm op}\le C\sqrt{q_0}e^{-\nu' L}.
 \label{f16v12:Delta-op}
\end{equation}
The constants are uniform in domain shapes, duration and ambient volume.
After decreasing the earlier $\nu$, all exponential estimates may use a
single fixed exponent.
\end{proposition}
\begin{proof}
Restrict the equation $L\eta_\alpha=-E^{\mathsf T}d_t v_\alpha$ to free
port rows and compare it with the Dirichlet equation for $\eta^D_\alpha$.
This proves \eqref{f16v12:harmonic-lift}, including its sign. The full and
conditional Gaussian cancellation identities give, as an identity on the
whole fast space with the exterior inserted,
\[
                 g^D_\alpha-g_\alpha
                          =-\Delta_\alpha^{\mathsf T}S_0.
\]
The ordinary source is the same in both actions for these interior marks.
A free port differentiates no exterior-only term. Subtract the two exact
compact lifted scores and then their Gaussian constants to obtain
\eqref{f16v12:exact-residual-match}. There is no nonlinear comparison of
probability laws in this identity. In particular no flat derivative is used
in place of a compact Haar derivative; the only $S_0$ is the fixed Gaussian
coefficient.

We give the row/column estimate needed for a whole family. The full port
inverse and every principal inverse have uniformly exponentially decaying
entries in ambient spatial distance, with summable rows, as in V11. For an
exterior row $p$, $\Delta_{p\alpha}=\eta_{p\alpha}$; its path from the
source reaches $D_s^c$. For an interior row, the first identity in
\eqref{f16v12:harmonic-lift} expands it as an inverse path to a boundary
edge followed by a full inverse path to the source. Such a path has length
at least the source-to-row distance and at least the source depth, up to a
fixed stencil allowance. Splitting the available exponential weight between
these distances and summing the remaining convolution gives the uniform
entry bound
\[
  |\Delta_{p\alpha}|
       \le C e^{-\nu' L}e^{-\nu' d_{
m sp}(p,\alpha)}
              1_{\{p\text{ on a bond adjacent to }i(\alpha)\}}.
\]
Colour and port multiplicities are fixed. A source occurs at most $q_0$
times because its core lies in its patch. The absolute column sums are
therefore at most $Ce^{-\nu'L}$ and the absolute row sums at most
$Cq_0e^{-\nu'L}$. Schur's bound gives \eqref{f16v12:Delta-op}. This
argument does not multiply by the number of cylinders or by the duration.
\end{proof}

\begin{theorem}[Exceptional conditional Grams under their actual posteriors]
\label{f16v12:conditional-Gram-budget}
For every such cylinder family and every collection of exterior events
$E_s$,
\begin{align}
 \sum_s\mathbb E_P\|\mathbf r^D_{\gamma,\mathcal A_s}
                  -\mathbf r^{[\infty]}_{\gamma,\mathcal A_s}\|^2
        &\le C_rq_0\mathcal N\epsilon_\gamma e^{-2\nu L},\notag\\
 \sum_s\mathbb E_{\bar P_s}
        [1_{E_s}\operatorname{Tr}\Gamma^{D_s}_\gamma]
        &\le C_rq_0\mathcal N\epsilon_\gamma.
 \label{f16v12:conditional-Gram-estimates}
\end{align}
Here $\bar P_s$ is the actual exterior posterior for $D_s$, and each
conditional Gram is restricted to $\mathcal A_s$. Its complement outside
V11's bounded collar is included. No event probability is being estimated.
\end{theorem}
\begin{proof}
The primitive port vector has global Jacobian operator norm bounded by a
geometric constant and Gaussian squared energy at most
$C_r\mathcal N/\gamma$, by V7's primitive inserted-word certificate. Apply
V7's Gaussian vector-energy inequality and the entropy variational inequality
to this vector itself, using \eqref{f16v12:inherited}. At a fixed positive
parameter they give
\[
                   \mathbb E_P\|e^{\rm pt}\|^2
                                   \le C_r\mathcal N\epsilon_\gamma.
\]
This is the same proved energy-transfer argument, not a native functional
inequality. Combine it with \eqref{f16v12:exact-residual-match} and
\eqref{f16v12:Delta-op} to obtain the first bound. V7 also gives
$\mathbb E_P\|\mathbf r^{[\infty]}_\gamma\|^2\le
 C_r\mathcal N\epsilon_\gamma$. Restricting and repeating each source
at most $q_0$ times, followed by $\|u-v\|^2\le2\|u\|^2+2\|v\|^2$,
bounds the sum of uncentered conditional-residual energies by
$C_rq_0\mathcal N\epsilon_\gamma$. Conditional centering reduces that
energy. The tower property and $1_{E_s}\le1$ prove the second bound.
All conditional expectations and exterior integrals refer to the full native
law; no deterministic bounded boundary has replaced them.
\end{proof}

This estimate and \eqref{f16v12:exceptional-mean} pay both types of exceptional
term in V11's ledger, but only as sums of squared source energies or traces.
Their small parameter is $\epsilon_\gamma$, much larger than
$\gamma^{-1}$. They cannot be used to discard those terms in a leading
$\gamma^{-1}G_2$ coefficient theorem.

\subsection{Exact spatial gluing and its unconditioned trace-norm error}

For an exact gluing statement, now assume the interiors $D_s$ are separated
so that no native primitive term meets two of them. In particular a gap wider
than the fixed word range suffices. Let $I=\bigcup_s I_{D_s}$ and let $O=I^c$.
Every chosen source carrier and free Haar lift lies in its own cylinder.
Conditional on $O$, the interiors are independent, and the conditional
partition function factors into the literal $Z_{\gamma,D_s}$ times
exterior-only factors. The common exterior posterior is
\begin{equation}
 d\bar P(q)=Z_\gamma^{-1}
       1_{\mathcal P'_O}J^f_{\gamma,O}(q)e^{-U_{\gamma,O}(q;\mathbf y)}
               \prod_s Z_{\gamma,D_s}(q;\mathbf y)\,dq.
 \label{f16v12:joint-posterior}
\end{equation}
Each message $j_{D_s}$ is its earlier function of the word collar. Its
marginal integrals under this joint posterior agree with the corresponding
marginal integrals under the full native law.

Let $\chi$ select the union of the source banks, without repeats. Denote the
full nonadjacent-time memory by $\mathsf M_\gamma$ as in V7, and place each
conditional memory on its own source block. Put
\[
 \mathsf R_{\rm post}=\operatorname{Cov}_{\bar P}(\mathbf j),\qquad
 \mathsf L_D=\bigoplus_s\mathbb E_{\bar P}\mathsf M^{D_s}_\gamma(q).
\]
These are \emph{actual full-exterior averages}; they are not averages over
V11's good event, nor averages of its asymptotic coefficients.

\begin{theorem}[A full-native spatial gluing error in normalized trace norm]
\label{f16v12:gluing}
For the separated family,
\begin{equation}
       \chi\mathsf M_\gamma\chi-\mathsf L_D
                          =-\operatorname{Off}_{>1}\mathsf R_{\rm post}
 \label{f16v12:gluing-identity}
\end{equation}
is exact. If its collar patches have overlap at most $q_0$ and source depth
at least $L$, then
\begin{equation}
 \boxed{
  \frac{\|\chi\mathsf M_\gamma\chi-\mathsf L_D\|_{S_1}}{\mathcal N}
       \le C_rq_0\epsilon_\gamma
                     +C\frac{p_{\rm tot}}{\mathcal N}e^{-2\nu L}.}
 \label{f16v12:gluing-bound}
\end{equation}
Every time lag is included simultaneously. The trace norm is the sum of
singular values, not an entrywise absolute row norm. There is no claim that
the right side bounds the unnormalized operator norm.
\end{theorem}
\begin{proof}
At a fixed regulator differentiate the full normalized integral on the
separated interior source directions. Exterior-only action terms have no
such derivatives. The first derivative is the posterior mean of the
conditional first derivatives. Differentiating it once more gives
\[
 \nabla^2[-\log Z_\gamma]
    =\mathbb E_{\bar P}\bigoplus_s\nabla^2[-\log Z_{\gamma,D_s}]
                         -\operatorname{Cov}_{\bar P}(\mathbf j)
\]
on these sources. This is V11's exact total-covariance differentiation,
with the actual joint posterior \eqref{f16v12:joint-posterior}. Compact
integrability justifies it even at exterior values outside every bounded
rescaled window. Restrict to nonadjacent time blocks and use V11's exact
conditional contact range to get \eqref{f16v12:gluing-identity}. The full
history's external logarithmic factors are single-time functions and change
none of these blocks; they have not been dropped from the original weight.

The posterior Gram is positive. V7's Fourier block-extraction argument gives
$\|\operatorname{Off}_{>1}R\|_{S_1}\le4\|R\|_{S_1}$, independently
of duration. For positive $R$, its trace norm is its trace. Apply
\eqref{f16v12:return-trace} and divide by $\mathcal N$. This proves
\eqref{f16v12:gluing-bound} without summing a separation-independent error
once for every lag. The factorization of interior integrals did not
factorize the posterior, and no cross-message covariance was discarded.
\end{proof}

The bounded-overlap covering constructed above is not itself a separated
family. Its message-energy theorem holds without separation, but
\eqref{f16v12:gluing-identity} is asserted only for a separated collection.
One may take a separated subfamily, or finitely colour the bounded-degree
patch-intersection graph and apply the identity within a colour. This does
not authorize ignoring cross-colour response or tensorizing the native law.

\subsection{What is closed, and what is still needed}

The exterior-law terms in V11 are no longer wholly unestimated. For the
current full fast-history law, the actual posterior-mean Gram, the
source-weighted exceptional mean energy, and the exceptional conditional
Gram trace now have duration- and volume-uniform \emph{density} bounds.
Neither an all-good collar probability nor bounded conditional means on
arbitrary deterministic exterior data was needed. A collar of depth $L$
separates reference Gaussian communication, costing $e^{-2\nu L}$, from
native conditional-law mismatch, costing the actual entropy difference
$k_D$. Vector concentration and bounded overlap prevent either cost from
being multiplied by the number of time slices.

This is not a new complete-native $\gamma^{-1}$ expansion at arbitrary
volume. The floor $\epsilon_\gamma$ is much larger than $\gamma^{-1}$;
V11's marked coefficient remainder is still justified only on its stated
bounded exterior window and at its local-size entry condition. In
particular \eqref{f16v12:conditional-Gram-estimates} does not prove that the
exceptional conditional-covariance average is $o(\gamma^{-1})$.
No derivative of a scalar entropy or pressure error is being taken.

A trace controls total squared message energy but can allow that energy to
concentrate in a coherent source direction. It is therefore not enough to
bound the operator norm of the actual exterior return independent of
$B(n+1)$, and no exponentially weighted temporal row bound for that return
is asserted. V7's exceptional-direction warning remains in force. The
reference Gaussian inequalities in this appendix are not native
Poincare or logarithmic Sobolev inequalities.

The next useful target can be stated without the now-unnecessary all-time
collar gate. For the actual joint posterior and these exact messages, prove
an estimate on \emph{every} coefficient vector,
\[
 \operatorname{Var}_{\bar P}
       \left(\sum_{s,\alpha}v_{s,\alpha}j^D_{\gamma,\alpha}\right)
            \le a_\gamma\sum_{s,\alpha}|v_{s,\alpha}|^2,
                    \qquad a_\gamma\longrightarrow0,
\]
or a source-weighted connected return with summable space--time rows.
The present proof identifies its loss: $k_D$ is charged once to a whole
vector, yielding its squared-norm expectation, but it does not distribute
that cost according to each coherent direction's variance. A
source-resolved conditional contraction or an actual posterior functional
inequality would address that loss; another entropy-density estimate or
another formal coefficient would not. If the leading native coefficient is
also to survive spatial gluing, the exceptional budget must additionally be
improved below the $\gamma^{-1}$ scale.

All statements remain conditional on the prescribed bounded \emph{bridge}
history, before the physical root Gauss restriction. Removing a cutoff on
exterior \emph{fast} fields does not prove gauge-invariant coverage of the
integrated bridge law. Infrared contact coercivity, comparison with the same
full physical transfer top, physical-time calibration and an interacting
continuum construction remain separate obligations.

\subsection{Diagnostics and source preservation}

The new diagnostic is
\begin{center}\small
\path{verify_ym_f16_v12_posterior_screening.py}.
\end{center}
It checks Gaussian Markov conditioning and the conditional-entropy chain,
whole-vector mean comparison on translated correlated Gaussians,
boundary-reaching inverse factors, bounded-overlap spatial covers,
restricted/full port-lift matching, and exact posterior gluing on positive
finite laws. All 1,201 exact assertions passed, including 934 spatial-cover and duration-count
fixtures. Assertion families and counts are recorded in
\path{ym_v12_verification_report.json}. These are finite algebraic and
combinatorial diagnostics, not proof-assistant verification of the Gaussian
functional inequalities, native entropy input, or derivative suprema.
The analytic estimates above are the proofs; the finite checks do not
certify a physical mass gap.
The new checker has SHA--256
\begin{center}\scriptsize\ttfamily
A38B12309C4CA8CAB457AB53F820A6301EC24AA2F413D5BD4FE6B42CFA776F78.
\end{center}
The available V4--V11 diagnostics were rerun unchanged and passed 9,467,
89,339, 906, 311, 609, 990, 608 and 463 assertions, respectively.

The cumulative V12 changes only the V11 title version and appends this exact
section before its final \verb|\end{document}|. Removing the appendix and
restoring the title recovers V11 byte for byte. The package records the
input hashes, exact recovery, compilation, render checks, and actual reruns
of the available V4--V11 diagnostics. Unavailable archive diagnostics and
any external Lean development are not described as rerun. No predecessor
mathematical text or bibliography string is silently corrected. The next
scheduled five-version companion checkpoint remains V15.

\begin{firewall}
V12 proves an unrestricted-exterior posterior-mean screening estimate for
whole source banks, pays the old exceptional mean and conditional-Gram
terms in density, and gives an exact separated-cylinder spatial gluing
identity with a vanishing normalized trace-norm error. All fast exterior
fields and their actual posterior weights are retained, with no all-good
history event. The error floor is $\epsilon_\gamma$, not an
$o(\gamma^{-1})$ coefficient remainder. No unrestricted-volume small
operator norm or time-row summability of the posterior return, native
functional inequality, physical bridge coverage, same-top physical
transfer gap, or interacting continuum Yang--Mills mass gap is established
by this appendix.
\end{firewall}
% END F16 V12 APPEND

% BEGIN F16 V13 APPEND -- 2026-09-17
% Insert immediately before the final \end{document} in cumulative V12.
% The predecessor is not silently corrected; only its title becomes V13.

\clearpage
\section{V13: conditional resampling, a local-test obstruction, and a directional reduction}
\label{f16v13:context}

V12 bounds the trace density of the actual posterior return, but not its
largest eigenvalue. The new literature memorandum describes the earlier V3
frontier. We do not restart there. We instead introduce a checkable local
route to the \emph{directional} inequality left open in V12. The auxiliary
operator is the exact one-cell conditional resampling generator of the same
full compact history law. Its square-root-density realization is genuinely
local and frustration free, unlike an unverified claim about the physical
Wilson logarithmic transfer.

This appendix proves three new interfaces for the present law: a finite
local integral certificate for that auxiliary gap, a resampling-energy
estimate for the actual Haar residuals, and an operator-norm reduction for
the actual posterior return. The Gaussian auxiliary gap is determined on
\emph{all} polynomial chaoses. An exact cube witness shows that the simplest pair-row test fails even
for the Gaussian reference and, eventually, for the native law. A uniform
native gap by another criterion and uniform local fourth moments are
\emph{not} certified here. Consequently
this is a rigorous conditional route beyond V12, not a new unconditional
unrestricted-volume mass-gap or posterior-contraction theorem.

\subsection{What the literature changes, and what it does not}

The supplied V12 appendix and wider-literature memorandum were read in full.
Primary-source statements and selected technical sections were checked;
the accompanying literature audit records the depth of each review. The
following choices are based on applicability, not on the prestige of an
analogy.

\paragraph{Conditional dependence is the closest interface.}
Caputo--Menz--Tetali,
\href{https://arxiv.org/abs/1405.0608}{\emph{Approximate tensorization of
entropy at high temperature}}, derive functional inequalities from explicit
conditions on dependence. Anari--Liu--Oveis Gharan,
\href{https://arxiv.org/abs/2001.00303}{\emph{Spectral Independence in
High-Dimensional Expanders and Applications to the Hardcore Model}}, require
correlation control for the distribution and its conditionals. Neither
hypothesis is supplied by V12's small relative entropy density. Their discrete
statements are not quoted as compact-group theorems. Below we prove the
continuous projection inequality actually needed, with all conditioning and
clock conventions explicit.

\paragraph{Compatible local constraints.}
Gosset--Mozgunov's
\href{https://arxiv.org/abs/1512.00088}{\emph{Local gap threshold for
frustration-free spin systems}}, Lemm--Mozgunov's
\href{https://arxiv.org/abs/1801.08915}{\emph{Spectral gaps of
frustration-free spin systems with boundary}}, and Anshu--Arad--Vidick's
\href{https://arxiv.org/abs/1602.01210}{\emph{Simple proof of the detectability
lemma and spectral gap amplification}} concern compatible local constraints.
We construct such constraints for the \emph{auxiliary resampling law}; we do
not identify their gap with that of the physical transfer. The elementary
pair-angle criterion below works on the infinite-dimensional local group
spaces and is proved directly. It is not a new general Knabe theorem.

\paragraph{A more direct multiscale alternative than geometric analogy.}
Bauerschmidt--Bodineau's
\href{https://arxiv.org/abs/1907.12308}{\emph{Log-Sobolev inequality for the
continuum sine-Gordon model}} uses a multiscale Bakry--\'Emery criterion
controlled through the Polchinski equation, including non-log-concave
measures. This is a concrete alternative should a worst-boundary local-angle
test fail. It would require estimates on the actual renormalized potential
of the present law, not just its finite Taylor coefficients. It supplies
neither a Yang--Mills inequality nor permission to replace the compact law by
a convex extension. The more recent continuum spectral-independence work
\href{https://arxiv.org/abs/2601.18748}{\emph{Sampling Sphere Packings with
Continuum Glauber Dynamics}} concerns repulsive point processes, not this
fixed-coordinate gauge history.

\paragraph{Keep spectral elimination, but do not rename a covariance inverse.}
Bach--Ballesteros--Geisler's
\href{https://arxiv.org/abs/2508.07643}{\emph{The Spectral Renormalization Flow
Based on the Smooth Feshbach--Schur Map}} establishes a smooth spectral flow.
The monograph already retains vacuum-corrected Feshbach metrics. A physical
application still needs the complement inverse on its actual domain and the
same physical vacuum. Our resampling operator is not that physical
Feshbach operator.

\paragraph{Other proposals remain separate proof obligations.}
Kenig--Merle's method combines concentration compactness and rigidity. Its use
here needs a nonvanishing profile and a critical compactness argument; no such Yang--Mills profile
follows from the current local estimates. Perelman's Ricci-flow entropy is
not the relative entropy in V12. Tao's Collatz first-passage argument and
Bourgain--Gamburd expansion need their particular transport or convolution
and nonconcentration hypotheses. Nash--Moser smoothing needs a tame inverse,
not merely constants at each formal order. The audit records these primary
sources and the precise missing structures. None is used as an unproved
shortcut. We retain one active next test: local conditional resampling and
its actual source energy, rather than opening all these branches at once.

The archive audit found earlier conditional heat-bath and assumed-Dobrushin
interfaces, including f14's initial transition-family theorem and the
monograph's \texttt{thm:SMFS-Dobrushin-gap}. Those general constructions are
credited rather than claimed as absent. The new specialization below uses
V12's \emph{full augmented fast-history law}, its exact residual synthesis,
and its separated-cylinder posterior. All proofs keep the native measure
fixed.

\subsection{An exact local frustration-free operator for the full history law}

Keep V12's bounded bridge history, $\max_{i,e}|y_{i,e}|\le r$, and set
$\mathcal N=B(n+1)$. The fast cell $c=(i,b)$ contains V5's unmoving
coordinates $\zeta_c=(x_{i,b},a_{i,b})$ when $i<n$, and only $x_{n,b}$ at the
last slice. Dimensions are at most 36. Work on the complete principal-chart
product, with its null boundary sets ignored. Normalize each one-cell factor
of $J^f_\gamma\,d\zeta$ to a probability $m_{\gamma,c}$, and put
$m_\gamma=\bigotimes_c m_{\gamma,c}$. These normalizations multiply the
partition function by a scalar only. In particular the endpoint chord powers
$w_i=1/2$ and the interior powers one are retained. Then the same native law is
\[
 dP=Z_m^{-1}e^{-U_\gamma}\,dm_\gamma.
\]
The Gaussian endpoint amplitudes belong to $U_\gamma$, not to a newly chosen
vacuum. At fixed finite $B,n,\gamma$, $U_\gamma$ is bounded on this product
and $P$ is equivalent to $m_\gamma$ with bounded positive density ratio.

Define the bounded operators on $L^2(P)$
\begin{equation}
 E_c f=\mathbb E_P[f\mid\zeta_{c^c}],\qquad Q_c=I-E_c,
 \qquad \mathcal L_P=\sum_c Q_c,
 \qquad
 \mathcal E_P(f)=\sum_c\mathbb E_P\operatorname{Var}_P(f\mid\zeta_{c^c}).
 \label{f16v13:generator}
\end{equation}
Each cell is updated at rate one. There is no division by $\mathcal N$.
Let $\lambda_P$ be the infimum of $\mathcal E_P(f)/\operatorname{Var}_P f$
over nonconstant $f$.

\begin{proposition}[Local projections with the exact common state]
\label{f16v13:parent}
Each $Q_c$ is an orthogonal projection. The kernel of $\mathcal L_P$ consists
exactly of constants. Under the unitary
\[
 U_P:L^2(P)\longrightarrow L^2(m_\gamma),\qquad
 U_Pf=\Psi f,\qquad \Psi=(dP/dm_\gamma)^{1/2},
\]
the transformed projections have finite spatial and temporal interaction
range and annihilate the same state $\Psi$. If no primitive action term meets
both $c$ and $d$, then $E_cE_d=E_dE_c=E_{\{c,d\}}$, where the last operator
conditions on all variables except those two cells. Moreover, at every fixed
finite regulator,
\begin{equation}
 \lambda_P\ge e^{-\operatorname{osc}U_\gamma}>0.
 \label{f16v13:crude-gap}
\end{equation}
The bound in \eqref{f16v13:crude-gap} is generally exponentially poor in
$\gamma\mathcal N$ and is not the desired uniform result.
\end{proposition}
\begin{proof}
Conditional expectation is an orthogonal projection, and
$\|Q_cf\|_2^2=\mathbb E\operatorname{Var}(f\mid\zeta_{c^c})$.
For fixed exterior $z$, write the sum of terms meeting $c$ as $V_c(t;z)$.
The conditional amplitude is
\[
 \psi_c(t;z)=\frac{e^{-V_c(t;z)/2}}
 {\left(\int e^{-V_c(s;z)}dm_{\gamma,c}(s)\right)^{1/2}}.
\]
Direct multiplication gives
\begin{equation}
 (U_PE_cU_P^{-1}F)(t,z)
     =\psi_c(t;z)\int\psi_c(s;z)F(s,z)\,dm_{\gamma,c}(s).
 \label{f16v13:local-parent-kernel}
\end{equation}
Only cells in the fixed primitive word neighbourhood of $c$ enter $\psi_c$.
Thus this rank-one-in-$c$ operator acts trivially on every other exterior
coordinate. All transformed $Q_c$ kill $\Psi$ because all $Q_c$ kill one.
If no word meets $c,d$, their conditional pair density at fixed remaining
variables is a product. Integrating the two coordinates in either order
proves the commutation and joint-conditional identity, even when their
exterior neighbourhoods overlap.

For the kernel and gap assertions, let $p_-,p_+$ be the essential lower and
upper bounds of $dP/dm_\gamma$. Product variance tensorization gives
\[
 \operatorname{Var}_P f\le p_+\operatorname{Var}_{m_\gamma}f
 \le p_+\sum_c\mathbb E_{m_\gamma}\operatorname{Var}_{m_{\gamma,c}}f.
\]
Conditional variance averaged over the exterior is the infimum of
$\int|f-g(\zeta_{c^c})|^2dP$ over all exterior functions $g$. It is therefore
at least $p_-$ times the corresponding infimum under $m_\gamma$. Combining
these inequalities proves $\mathcal E_P(f)\ge(p_-/p_+)\operatorname{Var}_P f$.
Here $p_-/p_+\ge e^{-\operatorname{osc}U_\gamma}$. It also proves that the
kernel consists of constants. Approximation in $L^2$ extends the argument
from bounded functions. The complete principal angular domains have bounded
radius, and the number of positive action terms is $O(\mathcal N)$; hence
$\operatorname{osc}U_\gamma\le C_r\gamma\mathcal N$ for $\gamma\ge1$.
\end{proof}

This is a genuine local frustration-free \emph{auxiliary} operator. Its clock
is cell resampling, not lattice Euclidean time. Its constant state in $L^2(P)$
is not identified with the physical transfer vacuum. The map $U_P$ does not
assert a unitary equivalence to $-\log T_{\rm full,\gamma}$. This distinction
is essential even if a good gap for $\mathcal L_P$ is eventually proved.

\subsection{The reference gap is an all-chaos statement}

Let $Q=N(m,H^{-1})$ be V12's Gaussian law on the same Euclidean coordinate
partition, and let $\mathcal L_Q$ be \eqref{f16v13:generator} with Gaussian
conditionals. Write $D_H=\bigoplus_c H_{cc}$.

\begin{theorem}[Exact Gaussian resampling edge on the full function space]
\label{f16v13:Gaussian-gap}
The rate-one-per-cell Gaussian generator satisfies
\begin{equation}
 \lambda_Q=\lambda_{\min}(D_H^{-1/2}HD_H^{-1/2})
                         \ge h_-/h_+>0.
 \label{f16v13:Gaussian-edge}
\end{equation}
The equality is on all of $L^2(Q)$, not just on linear observables. It is
uniform in spatial volume, duration, and the Gaussian mean. It does not
compare $\mathcal L_Q$ with $\mathcal L_P$.
\end{theorem}
\begin{proof}
Whiten the Gaussian to $g=H^{1/2}(\zeta-m)$ and identify its first chaos with
the Euclidean coefficient space. Let $I_c$ inject the coordinates of cell
$c$. The innovation subspace orthogonal to all exterior linear functions
has orthogonal projection
\[
 R_c=H^{1/2}I_cH_{cc}^{-1}I_c^{\mathsf T}H^{1/2}.
\]
Indeed the columns of $H^{1/2}I_c$ are orthogonal to the exterior coordinate
columns of $H^{-1/2}$, and their Gram is $H_{cc}$. Therefore
\[
 \sum_cR_c=H^{1/2}D_H^{-1}H^{1/2}
\]
has the same eigenvalues as $D_H^{-1/2}HD_H^{-1/2}$. On the first chaos the
generator is exactly this sum, so its smallest eigenvalue is an upper bound
on the full gap.

On the $k$th homogeneous Gaussian chaos, conditional expectation onto
exterior coordinates is the restriction of $(I-R_c)^{\otimes k}$ to the
symmetric tensor space. One can see this by taking conditional expectations
of the Wick exponential
$\exp(\langle t,g\rangle-|t|^2/2)$ and comparing its homogeneous
coefficients. Projections acting on different tensor legs commute, and their
joint eigenvalues are zero or one. Hence
\[
 I-(I-R_c)^{\otimes k}\succeq\frac1k\sum_{j=1}^k(R_c)_j.
\]
Sum over $c$. The right side is at least $\lambda_Q I$ on each $k\ge1$,
where $\lambda_Q$ denotes the first-chaos value just computed. Hermite
polynomials are dense in Gaussian $L^2$, so the same lower bound holds on the
orthogonal complement of constants. Finally $H\succeq h_-I$ and
$D_H\preceq h_+I$ prove the stated comparison.
\end{proof}

Exact Gaussian sampler rates are classical; see Roberts--Sahu,
\href{https://doi.org/10.1111/1467-9868.00070}{\emph{Updating Schemes,
Correlation Structure, Blocking and Parameterization for the Gibbs Sampler}}
(1997). The direct proof above fixes the present continuous-time convention
and checks every chaos. V3's conditional covariance lower precision and
V8's fast transfer ratio concern different operators and are not substituted
for this calculation. Small $D(P\Vert Q)/\mathcal N$ still does not transfer
\eqref{f16v13:Gaussian-edge} to $P$.

\subsection{A local native conditional-angle criterion}

Fix two cells $c,d$. Conditional on all other cells $z$, use the full native
pair law $P_{cd}^z$. Define its maximal correlation
\begin{equation}
 \theta_{cd}(z)=\sup
 \left\{|\mathbb E[f(\zeta_c)g(\zeta_d)]|:
       \mathbb Ef=\mathbb Eg=0,\ \mathbb Ef^2=\mathbb Eg^2=1\right\}.
 \label{f16v13:theta}
\end{equation}
The expectations in this definition are under this \emph{conditional pair
law}. Define $\bar\theta_{cd}=\operatorname*{ess\,sup}_z\theta_{cd}(z)$,
with $\bar\theta_{cc}=0$. Only the finite word neighbourhood of the pair
enters. Noninteracting pairs have $\bar\theta_{cd}=0$ exactly. The essential
supremum allows every exterior fast value, not merely V11's bounded window.

\begin{theorem}[Weighted local-angle certificate for the native auxiliary gap]
\label{f16v13:angle-gap}
For positive weights $0<w_c\le1$, put
\begin{equation}
 \delta_w=\min_c\left(w_c-\sum_{d\ne c}w_d\bar\theta_{cd}\right).
 \label{f16v13:weighted-margin}
\end{equation}
If $\delta_w>0$, then
\begin{equation}
 \sum_c w_cQ_c\succeq\delta_w(I-\mathbb E_P),\qquad
                       \lambda_P\ge\delta_w.
 \label{f16v13:angle-conclusion}
\end{equation}
This assertion holds on the full compact-coordinate $L^2(P)$; no
finite-dimensional local-spin assumption or chart cutoff is required.
The criterion is sufficient, not necessary. A failed margin does not imply
a small actual gap.
\end{theorem}
\begin{proof}
First condition on the complement of one pair. On its mean-zero $L^2$
space let $\mathcal U,\mathcal V$ be the subspaces of centered functions of
$c$ and $d$, respectively. Their cross inner-product norm is
$\theta=\theta_{cd}(z)$. The synthesis map $(u,v)\mapsto u+v$ from
$\mathcal U\oplus\mathcal V$ has squared norm at most $1+\theta$, because
\[
 \|u+v\|^2\le\|u\|^2+\|v\|^2+2\theta\|u\|\|v\|
                  \le(1+\theta)(\|u\|^2+\|v\|^2).
\]
Its product with its adjoint is $E_c+E_d$ on this centered space (the
subspaces occur in the opposite cell order). Thus $Q_c+Q_d\succeq
(1-\theta)I$ there. Constants form its kernel. Spectral calculus yields
$(Q_c+Q_d)^2\succeq(1-\theta)(Q_c+Q_d)$, or
\begin{equation}
 Q_cQ_d+Q_dQ_c\succeq-\theta(Q_c+Q_d).
 \label{f16v13:anticommutator}
\end{equation}
This also holds on constants. Integrate the fiber inequality and replace
$\theta$ by $\bar\theta_{cd}$. Noninteracting pairs commute and have a
nonnegative anticommutator, consistently with $\bar\theta=0$.

For $\mathcal L_w=\sum_cw_cQ_c$, expand the square and use the last
inequality for each unordered pair:
\[
 \mathcal L_w^2\succeq
 \sum_c w_c\left(w_c-\sum_{d\ne c}w_d\bar\theta_{cd}\right)Q_c
                  \succeq\delta_w\mathcal L_w.
\]
The spectral measure of this bounded positive operator is supported on
$\{0\}\cup[\delta_w,\infty)$. Its kernel is constants by
Proposition~\ref{f16v13:parent} and positivity of every weight. Since
$\mathcal L_P\succeq\mathcal L_w$, this proves both claims. No covariance
bound on native test functions was assumed in this proof.
\end{proof}

There are finitely many pair and endpoint cell types at the fixed word range.
Thus a \emph{uniformly successful} family of local pair certificates and
bounded weights would give a uniform margin independent of $B,n$. Taking all
weights one gives $1-\sup_c\sum_d\bar\theta_{cd}$. Alternatively, lower
bounds for \eqref{f16v13:weighted-margin} are linear inequalities in the
weights; a proposed rational solution can be checked directly. This is
ordinary weighted projection bookkeeping, not an import of sieve weights.
The next proposition will show why this particular one-cell criterion
cannot provide the required weak-coupling margin.

\subsection{Finite local integral certificates with their unresolved continuum paid}

A small matrix sampled at selected exterior points is not by itself an
upper bound for $\bar\theta$. We now give a finite upper certificate that
also pays discretization and the complete exterior parameter range.

Use angular coordinates $u=\zeta_c/\sqrt\gamma$, $v=\zeta_d/\sqrt\gamma$,
and angular coordinates $z$ on their finite exterior word neighbourhood.
Each domain is a product of principal three-balls of radius $\pi$, with
null boundary sets ignored. The two-cell reference factors are normalized
products of the appropriate powers of $j$, independent of $z$. The exact
conditional potential consists of all terms meeting $c$ or $d$, with the
other inputs frozen; its density is
\[
 p_z(u,v)=\frac{e^{-V_z(u,v)}}{\int e^{-V_z}\,dm_cdm_d}.
\]
Here and in this subsection $m_c,m_d$ are the angular versions of
$m_{\gamma,c},m_{\gamma,d}$. The local exponential-word formulas extend
continuously to the chosen closed angular product. Fixed derivative and
incidence bounds give finite constants
\begin{equation}
 \operatorname{osc}_{u,v}V_z\le O_\gamma,\qquad
 \operatorname{Lip}_{u,v,z}(V)\le A_\gamma,
 \qquad O_\gamma+A_\gamma\le C_r(1+\gamma).
 \label{f16v13:finite-card-constants}
\end{equation}
The constants are uniform in ambient $B,n$ and all principal exterior
angles. Endpoint Gaussian quadratics are included in these bounds. They
are not claims of useful smallness. For a certificate uniform over bridge
values, include the finitely many bridge inputs in their prescribed compact
window among the parameter net coordinates, with its corresponding metric.

Partition each of the pair domains into finitely many Borel bins of diameter
at most $h$, omitting reference-null bins. At each parameter value define
\begin{equation}
 A(z)_{ab}=
  \frac{P_z(U_a\times V_b)}{\sqrt{P_z(U_a)P_z(V_b)}}.
 \label{f16v13:probability-matrix}
\end{equation}
Use an $h$-net $\mathcal Z_h$ of the entire closed local parameter domain.
It is permitted to be very large. Let $\epsilon_{\rm int}$ be a certified
operator-norm enclosure error for the matrices used in place of the exact
matrices \eqref{f16v13:probability-matrix}.

\begin{theorem}[An upper, not merely lower, local correlation certificate]
\label{f16v13:finite-card}
With the conventions above, the full pair correlation obeys
\begin{equation}
 \boxed{
 \bar\theta_{cd}\le\min\left\{1,
  \max_{z\in\mathcal Z_h}\sigma_2(\widehat A(z))
  +2e^{4O_\gamma}(e^{4A_\gamma h}-1)+\epsilon_{\rm int}\right\}.}
 \label{f16v13:finite-upper}
\end{equation}
The error is independent of the ambient volume and duration. It is not
claimed small on a practical mesh. At an exact fixed parameter,
$\sigma_1(A(z))=1$ and $\sigma_2(A(z))\le\theta_{cd}(z)$; without the
explicit remaining error the finite matrix gives the wrong inequality for
a gap certificate. Singular values are zero-padded when necessary.
\end{theorem}
\begin{proof}
Let $p_c,p_d$ be the marginal densities at fixed $z$. The normalized
conditional-expectation kernel on the fixed reference spaces is
\[
 K_z(u,v)=p_z(u,v)/\sqrt{p_c(u)p_d(v)}.
\]
It is compact since its positive kernel is bounded. Its top singular value
is one, with normalized singular functions $\sqrt{p_c},\sqrt{p_d}$;
conditional expectation is a contraction and maps these functions to each
other. Its second singular value is exactly \eqref{f16v13:theta} by restricting
to their orthogonal complements. Positivity of the pair density excludes
nonconstant perfectly correlated functions, though that strictness is not
needed for the bound below.

Equivalently, on the marginal probability spaces the kernel is
$k_z=p_z/(p_cp_d)$. Orthogonal projection onto functions constant on the
bins gives precisely \eqref{f16v13:probability-matrix}, in the normalized
indicator bases. Both constant functions remain in the projection, so the
top singular value remains one and the second cannot exceed the full one.
This projection includes its original marginals, not arbitrarily uniform
bin weights.

The potential oscillation gives
$e^{-O_\gamma}\le p_z,p_c,p_d\le e^{O_\gamma}$, so
$k_z\le e^{3O_\gamma}$. At fixed $z$, the logarithm of a marginal density
has Lipschitz constant at most $A_\gamma$ in its own coordinate: this follows
by integrating the pointwise ratio bound for $e^{-V}$. Therefore the
oscillation of $\log k_z$ on any product of two bins is at most
$4A_\gamma h$. The doubly projected kernel is the binwise conditional
average of $k_z$ under $P_c\otimes P_d$. Its difference from $k_z$ has
absolute value at most
$e^{3O_\gamma}(e^{4A_\gamma h}-1)$. Its Hilbert--Schmidt norm, in these
probability spaces, has the same upper bound. Singular-value stability
therefore bounds the full second singular value by that of $A(z)$ plus
this error.

For $z,z'$ at distance at most $h$, normalized joint density ratios have
logarithmic absolute value at most $2A_\gamma h$; integrating gives the
same bound for each marginal ratio. Hence
$|\log(K_z/K_{z'})|\le4A_\gamma h$. Both kernels on the fixed reference
spaces have supremum at most $e^{2O_\gamma}$. Their Hilbert--Schmidt
distance is at most $e^{2O_\gamma}(e^{4A_\gamma h}-1)$. This step uses a
fixed reference space, not a comparison of operators on two different
unidentified marginal $L^2$ spaces. Apply singular-value stability once
more to pass to a nearest net point. Add the two errors and the certified
matrix enclosure error, and bound their coefficient by $2e^{4O_\gamma}$.
This proves \eqref{f16v13:finite-upper}. All domains, including large angular
inputs, are covered by the bins and net. No chart-complement probability
has been dropped.
\end{proof}

The certificate uses a fixed number of local variables but can require a
mesh exponentially fine in $\gamma$. It is an existence and validation
interface, not a claim that this computation has been performed for the
native Wilson pairs. The diagnostic supplied here tests the matrix identities
on finite probability models; it supplies neither native integral enclosures
nor a successful value of \eqref{f16v13:weighted-margin}. A failure of this
worst-boundary pair test would motivate larger blocks or a multiscale
criterion, not a declaration that the physical or auxiliary gap vanishes.

\subsection{The actual cube rejects the naive pair row, already at Gaussian order}

Before spending effort on native integral enclosures, one should test whether
the proposed sufficient criterion can even recognize the reference model.
For the unmoving cells chosen above, the answer is negative. The obstruction
below is for the \emph{actual} parity-cube matrices, not an arbitrary Gaussian
countermodel.

\begin{proposition}[A certified Gaussian and native failure of the one-cell row test]
\label{f16v13:pair-obstruction}
For $m\ge4$ and $n\ge4$, consider the three interior cells at times $1,2,3$
in one fixed spatial cube. In the Gaussian reference, each of the two
consecutive pair conditionals has maximal correlation greater than $4/5$.
Consequently no positive weights can make
\eqref{f16v13:weighted-margin} positive for this Gaussian cell partition.

There is a finite $\gamma_0$, determined by the fixed local geometry, such
that at zero bridge history the same impossibility holds for the
\emph{native} worst-exterior pair coefficients for every
$\gamma\ge\gamma_0$, $m\ge4$, $n\ge4$. This is failure of this sufficient
row certificate, not a vanishing native spectral gap.
\end{proposition}
\begin{proof}
We provide a rational witness. For one colour, let $\mathsf B,\mathsf F,N$
and $\mathsf H$ be the inherited rooted cube matrices, with the five chords
ordered $(a,b,c,d,e)$ as in f15 V97 and nonroot vertices in lexicographic
order. Put
\[
 L=\mathsf B^{\mathsf T}\mathsf B+3I_7,
 \qquad K=N+\mathsf F^{\mathsf T}L\mathsf F.
\]
The port precision has diagonal spatial cube block $L$: each nonroot port
meets three bridges, and none of these bridges has both endpoints in the
same cube for $m\ge4$. The mixed magnetic diagonal Gram is $2N$, by f15
V99. Substituting $z_t=a_t-\mathsf F(x_t-x_{t+1})$ into the exact quadratic
kinetic form gives the diagonal interior cell block and next-time block
\begin{equation}
 A=\begin{pmatrix}
       \mathsf H+2N+2K&-\mathsf F^{\mathsf T}L\\
       -L\mathsf F&L
    \end{pmatrix},\qquad
 C=\begin{pmatrix}-K&0\\ L\mathsf F&0\end{pmatrix}.
 \label{f16v13:actual-pair-blocks}
\end{equation}
Thus conditioning all other cells gives pair precision
$\left(\begin{smallmatrix}A&C\\C^{\mathsf T}&A\end{smallmatrix}\right)$.
All its matrices are rational. Conditioning affects its mean, not this
precision. Endpoint amplitude precisions do not enter these interior blocks.

For two blocks with positive precision diagonal $A$, whitening each block
and then taking a singular-value decomposition of
$A^{-1/2}CA^{-1/2}$ decomposes the precision into scalar pairs
$\left(\begin{smallmatrix}1&s\\s&1\end{smallmatrix}\right)$.
Such a pair has linear correlation $-s$. In particular the maximal
correlation is at least
\[
 \frac{|u^{\mathsf T}Cv|}{\sqrt{(u^{\mathsf T}Au)(v^{\mathsf T}Av)}}
\]
for any two coefficient vectors $u,v$.

With coordinates ordered as five chords followed by seven ports, choose
\[
 \begin{split}
 u&=(0,0,0,0,0;\ 1,-1,2,0,1,-1,3)^{\mathsf T},\\
 v&=(1,-1,-3,-2,-4;\ 1,-1,2,0,1,-1,2)^{\mathsf T}.
 \end{split}
\]
Direct rational multiplication of the displayed inherited cube matrices gives
\begin{equation}
 u^{\mathsf T}Au=86,\quad v^{\mathsf T}Av=825/8,\quad
 u^{\mathsf T}Cv=303/4,\quad
 \frac{|u^{\mathsf T}Cv|^2}{(u^{\mathsf T}Au)(v^{\mathsf T}Av)}
    =\frac{30603}{47300}>\frac{16}{25}.
 \label{f16v13:rational-witness}
\end{equation}
The last strict inequality has difference $331/47300>0$.
The checker reconstructs the incidence matrices rather than importing
floating-point blocks. One colour suffices; the other colours are independent
copies in this Gaussian calculation. This proves both pair lower bounds.

If all row margins were positive, the three corresponding weights would
satisfy $w_1>(4/5)w_2$, $w_3>(4/5)w_2$, and
$w_2>(4/5)(w_1+w_3)$. Dropping all the other nonnegative row terms is
legitimate. Substitution gives $w_2>(32/25)w_2$, a contradiction. The exact
Gaussian gap in Theorem~\ref{f16v13:Gaussian-gap} remains positive: taking
absolute pair interactions has lost information needed to see it.

For the native assertion, fix the exterior fast coordinates and bridges to
zero in this pair conditional. All terms meeting the pair and all original
Haar factors are retained. Its fixed-dimensional complete-chart integral
has the independent magnetic and rooted port pins from V5. Its normalized
polynomial moments through degree two therefore converge to those of the
Gaussian pair in \eqref{f16v13:actual-pair-blocks}: Taylor expansion on fixed
rescaled balls gives local convergence, the complete Gaussian pin controls
their complements, and the limiting partition function is strictly positive.
This is an ordinary fixed-dimensional limit, not a growing-volume norm
comparison. There are only the two fixed local pair types here; their words
are independent of the ambient $m,n$ once these interior cells exist.

Choose Gaussian linear observables whose centered correlation is greater
than $4/5$, as guaranteed by the strict inequality above. Their variances
are positive. Moment convergence makes the corresponding native correlation
greater than $4/5$ for sufficiently large $\gamma$, with a threshold depending
only on these local words. At each such finite $\gamma$ the conditional
moments vary continuously with the local exterior principal angles near
zero. Thus the inequality holds on a nonempty open exterior neighbourhood,
which has positive native marginal probability. The essential supremum in
$\bar\theta$ is therefore greater than $4/5$, not merely its value on a
null conditioning event. Apply the same three-weight contradiction. No
claim of uniform convergence over all exterior fields was used.
\end{proof}

This check changes the next direction. The all-boundary one-cell test cannot
close the weak-coupling program on this partition; a finer mesh will not fix
it. The exact local integral certificates remain valid for a different
partition or update family. Useful alternatives are larger conditional
blocks, a coordinate/update organization that respects the balanced Gaussian
split, or multiscale functional inequalities retaining the signed Gaussian
metric rather than absolute pair rows. None of those native alternatives
is assumed successful here. This is precisely why a Gaussian or small-volume
preflight should precede a claimed literature-based global closure.


\subsection{Local resampling energy of the actual corrected sources}

Return to the unscaled notation $\zeta$ for V5's rescaled fast coordinates.
Define actual full-law moments
\begin{equation}
 m_2(P)=\sup_c\mathbb E_P(1+|\zeta_c|^2),\qquad
 m_4(P)=\sup_c\mathbb E_P(1+|\zeta_c|^4).
 \label{f16v13:moment-inputs}
\end{equation}
They are finite at each regulator. Uniformity in $B,n,\gamma$ is not asserted.
V5's averaged second moments and V12's entropy density do not supply the
supremum fourth-moment bound in \eqref{f16v13:moment-inputs}.

Let $e=(e^{\rm br},e^{\rm pt})$ be V7's two primitive remainder banks, where
$e^{\rm br}=J_\gamma-J_0$ for the common sources and
$e^{\rm pt}=S_\gamma-S_0$ for the Haar port insertions. Each primitive
component has a fixed-size local fast carrier $C_j$, bounded row and column
incidence, and the inherited \emph{complete-chart} estimate
\begin{equation}
 |e_j(\zeta)|\le C_r\gamma^{-1/2}(1+|\zeta_{C_j}|^2).
 \label{f16v13:primitive}
\end{equation}
For the full exact lift, V7 gives
$\mathbf r=e^{\rm br}+\mathsf P^{\mathsf T}e^{\rm pt}$ with
$\|\mathsf P\|\le C$. All these are actual native functions.

\begin{theorem}[Directional innovation energy, without a native gap assumption]
\label{f16v13:innovation}
For every primitive coefficient vector $b$ and every common-source vector
$v$, uniformly in volume and duration,
\begin{align}
 \mathcal E_P(b^{\mathsf T}e)&\le
                  C_r\frac{m_4(P)}{\gamma}\|b\|^2,\notag\\
 \mathcal E_P(v^{\mathsf T}\mathbf r)&\le
                  C_r\frac{m_4(P)}{\gamma}\|v\|^2.
 \label{f16v13:innovation-energy}
\end{align}
For the affine Gaussian port score $S_0=\partial_aU_0$, evaluated under the
native law, one also has
\begin{equation}
 \mathcal E_P(b^{\mathsf T}S_0)\le C m_2(P)\|b\|^2.
 \label{f16v13:linear-energy}
\end{equation}
No native covariance mixing estimate enters these inequalities.
\end{theorem}
\begin{proof}
At stationarity let $\zeta^{(c)}$ be obtained by replacing cell $c$ by an
independent draw from its \emph{actual} conditional given all other cells.
Both $\zeta$ and $\zeta^{(c)}$ have law $P$, and
\[
 \mathcal E_P(f)=\frac12\sum_c
                 \mathbb E_P|f(\zeta)-f(\zeta^{(c)})|^2.
\]
Let $J(c)$ be the primitive scores whose carrier contains $c$.
Its cardinality, and the number of cells in any one carrier, are bounded
by a fixed number $k_0$. Other score components are unchanged. Thus
\[
 |b^{\mathsf T}(e(\zeta)-e(\zeta^{(c)}))|^2
 \le k_0\sum_{j\in J(c)}|b_j|^2
                   |e_j(\zeta)-e_j(\zeta^{(c)})|^2.
\]
Stationarity bounds the last squared expectation by $4\mathbb E_P e_j^2$.
Using \eqref{f16v13:primitive} and the fixed number of coordinates in its
carrier gives $\mathbb E_Pe_j^2\le C_r m_4(P)/\gamma$.
Sum over $c$; every $j$ occurs at most $k_0$ times. This proves the first
inequality without summing all primitive variances for each direction.
The bounded synthesis $b=(v,\mathsf P v)$ proves the second.

For the affine score, write $b^{\mathsf T}S_0=l^{\mathsf T}\zeta+$constant.
The linear map $b\mapsto l$ is a corner of the Gaussian precision $H$, so
$\|l\|\le h_+\|b\|$. Replacing cell $c$ changes this expression by
$l_c^{\mathsf T}(\zeta_c-\zeta_c^{(c)})$. Its squared expectation is at most
$4|l_c|^2\mathbb E_P|\zeta_c|^2$. Sum the resampling identity and use the
bound on $l$. This proves \eqref{f16v13:linear-energy}.
\end{proof}

This is a local \emph{Dirichlet-energy} estimate on every coefficient vector,
not the covariance inequality V12 sought. It differs from V7's total residual
energy and from Gaussian energy transfer. Dividing it by a proved native
auxiliary gap would make it a directional covariance estimate; dividing by
a Gaussian gap without a comparison would not.

\subsection{A coherent posterior reduction on the same actual law}

Use a separated cylinder family and source selection $\chi$ exactly as in
V12's gluing theorem. There are no repeated source marks. All source carriers
are interior, source depths are at least $L$, and enlarged cylinder patches
have overlap at most $q_0$. Put $O$ for the common exterior of all these
interiors and $E_O=\mathbb E_P[\,\cdot\mid\zeta_O]$. The posterior is V12's
actual posterior with its conditional partition functions included.

Stack the restricted lifts and extend them by zero. Let $\Delta$ be their
difference from the corresponding columns of the full lift. The inherited
identities, with all exterior fields admitted, are
\begin{equation}
 \|\Delta\|\le C\sqrt{q_0}e^{-\nu L},\qquad
 \mathbf r^D=\chi\mathbf r-\Delta^{\mathsf T}e^{\rm pt},\qquad
 \mathbf g^D-\chi\mathbf g=-\Delta^{\mathsf T}S_0.
 \label{f16v13:matched}
\end{equation}
The last vector is $O$-measurable: the lift mismatch is harmonic on each
free cylinder. Equivalently it is the difference of its two Gaussian
conditional source constants. The actual conditional mean vector is
\begin{equation}
                  \mathbf j=\mathbf g^D+E_O\mathbf r^D.
 \label{f16v13:message-identity}
\end{equation}
Separation ensures that this one common conditional expectation gives each
cylinder's message. It must not be replaced by unrelated projections for
an overlapping cover.

\begin{theorem}[Gap-and-moment reduction of the coherent posterior return]
\label{f16v13:posterior-reduction}
For the exact finite-regulator native law and the separated family above,
\begin{align}
 \|\Gamma_\gamma\|_{\rm op}
   &\le\frac{C_r}{\lambda_P}\frac{m_4(P)}{\gamma},\notag\\
 \|\mathsf R_{\rm post}\|_{\rm op}
   &\le\frac{C_r q_0}{\lambda_P}
          \left(\frac{m_4(P)}{\gamma}+m_2(P)e^{-2\nu L}\right),\notag\\
 \|\chi\mathsf M_\gamma\chi-\mathsf L_D\|_{\rm op}
   &\le\frac{C_r q_0}{\lambda_P}
          \left(\frac{m_4(P)}{\gamma}+m_2(P)e^{-2\nu L}\right).
 \label{f16v13:directional-reduction}
\end{align}
Here $\Gamma_\gamma$ is V7's full residual covariance,
$\mathsf R_{\rm post}=\operatorname{Cov}_{\bar P}(\mathbf j)$ and
$\mathsf L_D=\bigoplus_s\mathbb E_{\bar P}\mathsf M^{D_s}_\gamma$ are V12's
literal matrices. The constant absorbs the factor four in the last bound.
There is no normalization by $\mathcal N$ or by the number of sources.

In particular, a \emph{separately verified} uniform auxiliary gap
$\lambda_P\ge\delta_0>0$ and local moment bound $m_4(P)\le M_4$ would give
an $O_r(\gamma^{-1}+e^{-2\nu L})$ coherent posterior and gluing bound.
This final implication does not assert that its two inputs have been proved.
\end{theorem}
\begin{proof}
The finite-regulator gap is positive by \eqref{f16v13:crude-gap}, so its
Poincar\'e inequality applies to every present square-integrable function.
Apply it to $v^{\mathsf T}\mathbf r$ and use
Theorem~\ref{f16v13:innovation}. Taking the supremum over unit $v$ proves
the first assertion. This step makes clear which \emph{native} gap is used.

For the posterior, \eqref{f16v13:matched} expresses any
$v^{\mathsf T}\mathbf r^D$ as $b^{\mathsf T}e$, with
$\|b\|^2\le Cq_0\|v\|^2$, using the bounded full lift and the mismatch
bound. Therefore
\[
 \operatorname{Var}_P(v^{\mathsf T}\mathbf r^D)
       \le C_r q_0 m_4(P)(\lambda_P\gamma)^{-1}\|v\|^2.
\]
Conditional expectation on the \emph{one common} exterior is a contraction
on centered $L^2(P)$. Consequently the same bound applies to
$\operatorname{Var}_P(E_O(v^{\mathsf T}\mathbf r^D))$. The random part of
$v^{\mathsf T}\mathbf g^D$ is $-(\Delta v)^{\mathsf T}S_0$; the deterministic
full $g$ changes no variance. The last inequality in
Theorem~\ref{f16v13:innovation} gives
\[
 \operatorname{Var}_P(v^{\mathsf T}\mathbf g^D)
       \le Cq_0m_2(P)\lambda_P^{-1}e^{-2\nu L}\|v\|^2.
\]
Use \eqref{f16v13:message-identity} and
$\operatorname{Var}(F+G)\le2\operatorname{Var}F+2\operatorname{Var}G$.
The law of any $O$-measurable function under $P$ is exactly its law under
$\bar P$, proving the second assertion. Native large fields and their
posterior weights have not been conditioned away.

Finally V12 proves the exact identity
$\chi\mathsf M_\gamma\chi-\mathsf L_D=
-\operatorname{Off}_{>1}\mathsf R_{\rm post}$. The Fourier phase extraction
used in V7 is contractive in operator norm as well as trace norm: with
$U_t=\sum_i e^{\mathrm iit}\Pi_i$, each diagonal-offset extraction is a
Fourier coefficient of $U_t A U_t^*$ and has norm at most $\|A\|$.
Subtracting offsets $-1,0,1$ gives
$\|\operatorname{Off}_{>1}A\|\le4\|A\|$. This proves the last bound for
all lags at once. It does not give an exponentially weighted temporal row.
Since $m_2(P)\le C m_4(P)$, the final conditional implication follows.
\end{proof}

A successful alternative local-gap certificate, combined with native local
moment estimates, would supply the two inputs of the last theorem without
assuming the desired posterior covariance itself. The one-cell absolute-row
criterion has instead been ruled out by Proposition~\ref{f16v13:pair-obstruction}. This is the point of the reduction. Merely
substituting the Gaussian edge, or using an empirical conditional-correlation
matrix without its upper error, would be circular.

\subsection{Concrete next tests and the boundary of the advance}

The Gaussian all-chaos gap, the native local projection construction, the
local-angle \emph{criterion}, its finite enclosure formula, the innovation
energy, and the posterior reduction are proved here. A native uniform gap by an applicable criterion and uniform local fourth
moments remain unverified; the one-cell pair-row margin is ruled out at
large coupling, rather than left as a hoped-for hypothesis. Inserting
only the crude finite-regulator gap and the compact bound
$m_4(P)\le C(1+\gamma^2)$ into \eqref{f16v13:directional-reduction} does not
improve V12. We therefore do not report a new unconditional
unrestricted-volume contraction or a preserved leading coefficient.

The immediate test has in fact rejected the simplest proposed route:
Proposition~\ref{f16v13:pair-obstruction} certifies two adjacent-time
correlations greater than $4/5$ using the actual cube matrices, and proves
that positive weights cannot rescue this cell partition. We therefore do
not propose spending the next iteration merely tightening that row bound.
A larger-block or balanced-update Gaussian test must first reproduce the
known all-chaos reference gap. Only then is a complete native local-angle
or multiscale comparison worth attempting. Failure of the pair test is not
failure of a native gap; it is a diagnosis of this sufficient criterion.

A separate local-moment estimate must use the \emph{same} native law. The
supremum over cells in \eqref{f16v13:moment-inputs} cannot be replaced by an
average or by Gaussian moments. Even successful gap-and-moment inputs would
not supply temporal row decay, an $o(\gamma^{-1})$ exceptional coefficient
budget, gauge-invariant coverage of the retained bridge law, or a comparison
to the same physical transfer top. Those obligations are unchanged.

The wider literature thus changes the active direction: rather than infer a
coherent bound from another small trace density, construct and test local
conditional constraints whose common measure is exact, and combine their
native gap with source-local innovation energy. The established Gaussian
edge is a reference check, not the missing native theorem. Physical
Yang--Mills reconstruction and mass generation remain separate.

\clearpage
\subsection{Verification and preservation}

The accompanying diagnostic is
\begin{center}\small
\path{verify_ym_f16_v13_conditional_angles.py}.
\end{center}
It checks finite exact conditional projections and their common state,
conditional pair correlations, anticommutator signs, weighted-gap algebra,
Gaussian block innovations and tensor-chaos inequalities, the exact actual-cube
pair obstruction, normalized
probability-matrix compression, and separated-posterior centering.
All 222 exact assertions passed, including 12 actual-cube
obstruction checks. The available unchanged V4--V12 diagnostics also passed
9,467, 89,339, 906, 311, 609, 990, 608, 463 and 1,201 assertions, respectively.
The reports distinguish these rational and symbolic fixtures from a computation
of native Wilson integral enclosures. No such native enclosure, uniform
local-moment certificate, or external formalization was run. Analytic proofs,
including the complete-domain discretization error, are not certified by
finite examples.

The cumulative V13 preserves V12 byte for byte except for the title version
and insertion of this exact appendix before its final \verb|\end{document}|.
The package records hashes, predecessor recovery, compilation, rendered-page
review, the new diagnostic, and actual reruns of the available V4--V12
scripts. Unavailable historical proofs and companion programs are not
represented as recertified. The next scheduled companion checkpoint remains
V15; this targeted literature audit is recorded separately.

\begin{firewall}
V13 supplies an exact local auxiliary parent for the native history law, an
all-chaos Gaussian reference gap, a finite local conditional-angle upper
certificate and its weighted gap implication, and a directional resampling
energy and posterior-operator reduction. It rules out the naive one-cell pair-row test at weak coupling but does
not certify a native uniform gap by a different route or the local
fourth-moment input. Accordingly no new
unconditional unrestricted-volume posterior contraction, time-row decay,
physical transfer-gap comparison, or interacting continuum Yang--Mills mass
gap follows from this appendix. Its auxiliary clock is not physical time.
\end{firewall}
% END F16 V13 APPEND

% BEGIN F16 V14 APPEND -- 2026-09-17
% Insert immediately before the final \end{document} in cumulative V13.
% No predecessor mathematics is edited; only its title version becomes V14.

\clearpage
\section{V14: balanced conditional updates, compact cancellation, and a native repair event}
\label{f16v14:context}

V13 rejects the absolute pair-row criterion for its unmoving cells, even
at Gaussian order. Its recommendation was to test a different update
organization before attempting a native comparison. We do that test here.
Separating whole chord cubes from individual \emph{moving} temporal ports
removes the artificial Gaussian time coupling. For this partition a rational
weighted pair certificate succeeds, with margin $2/39$ on the entire Gaussian
function space. It is not a new proof of V13's Gaussian gap for its old
partition: the updates, and hence the generator, are different.

The native test then distinguishes a second obstruction. Actual compact
exterior ports can cancel the fields pinning a neighbouring pair. The resulting
pair conditional has a correlation at least $7/8$ for $\gamma\ge4096$, even
though its balanced Gaussian pair correlation is $1/6$. Thus changing
coordinates alone does not validate a worst-exterior perturbation argument.
This is not just another negative test: an explicit neighbouring update repairs
this cancellation family with a uniform, same-clock support-energy floor.
We retain that repair rather than freezing the very boundary variable that
provides it. An integrated defect identity specifies what remains to extend
the local repair to the complete generator.

\subsection{Dependency audit and the choice not to repeat a failed criterion}

The supplied V13 appendix and literature audit were read in full. Two semantic
archive searches returned no results; the mounted source inventories and
relevant exact passages were then examined directly. The audit is targeted,
not independent recertification of the archives. F15 V99's balanced edge
embedding and diagonal mixed Gram are inherited. F15 V86 already exhibits
zero spatial staples and the resulting near-unit one-link transfer ratio;
\emph{the general zero-staple mechanism is not new}. The calculation here is
for two temporal integration ports of the current fixed-bridge history,
including its periodic fixed roots, the revised resampling partition, and an
explicit repair port. F14's early repair-event argument is also credited:
our simpler repair bound uses an exact one-port conditional and no diffusion
comparison or changed clock.

Roberts--Sahu's
\href{https://doi.org/10.1111/1467-9868.00070}{\emph{Updating Schemes,
Correlation Structure, Blocking and Parameterization for the Gibbs Sampler}}
is the relevant classical context for changing a Gaussian update
parameterization. Its primary bibliographic record and abstract were checked;
no theorem from its full paper is imported. Qin--Wang's
\href{https://arxiv.org/abs/2208.11299}{\emph{Spectral Telescope: Convergence
Rate Bounds for Random-Scan Gibbs Samplers Based on a Hierarchical Structure}}
provides a general-state-space conditional-gap framework, not just a discrete
analogy. Its primary abstract was screened here, not its full proof. We do
not invoke it without its conditional inputs. All inequalities used below
are proved directly or attributed to the exact inherited statement.

The useful part of the supplied wider-literature memorandum is its insistence
on compatible local constraints. The native measure remains fixed. We neither
assert that different Gibbs generators have the same gap nor turn their
resampling clock into physical Euclidean time. The fixed-order corrections
of V10 are not recomputed.

\subsection{A different local generator on exactly the same native law}

Use the full law of V12--V13 at a prescribed bridge history with
$\max_{i,e}|y_{i,e}|\le r$, and let $\mathcal N=B(n+1)$, $B=m^3$, $m\ge4$,
$n\ge1$. Keep the chord groups $X_i$ and change the unmoving temporal groups
$h_t$ back to the moving groups $r_t$ by the \emph{exact} relation
\begin{equation}
 h_{t,v}=f_v(X_t)r_{t,v}f_v(X_{t+1})^{-1},\qquad r_{t,o}=I.
 \label{f16v14:coordinates}
\end{equation}
The $f_v$ are the same cube frames as in V2--V12. Write $P^{\rm bal}$ for
the pushforward of $P$ to $(X,r)$. There are two update species:
\[
 \mathcal I_x=\{(i,b):0\le i\le n\},\qquad
 \mathcal I_r=\{(t,b,v):0\le t<n,\ v\ne o\}.
\]
An $x$ update resamples all five chords of one cube at one time, with all
other $(X,r)$ variables fixed. An $r$ update resamples one three-dimensional
group port, with all other $(X,r)$ variables fixed. Each update has rate one.
For $a\in\mathcal I_x\sqcup\mathcal I_r$ set
\begin{equation}
 \begin{gathered}
 \widehat E_a f=\mathbb E_{P^{\rm bal}}[f\mid\text{all coordinates except }a],
 \qquad \widehat Q_a=I-\widehat E_a,\\
 \widehat{\mathcal L}=\sum_a\widehat Q_a,\qquad
 \widehat{\mathcal E}(f)=\sum_a\|\widehat Q_a f\|_{L^2(P^{\rm bal})}^2.
 \end{gathered}
 \label{f16v14:generator}
\end{equation}
The notation $r_t$ for group ports is distinct from the scalar bridge bound
$r$ and from V7's residual scores.

\begin{proposition}[Exact balanced updates and their local realization]
\label{f16v14:same-law}
The transformation \eqref{f16v14:coordinates} preserves conditional product
Haar measure on the ports and leaves all chord endpoint factors unchanged.
The law $P^{\rm bal}$ is therefore the same normalized native probability in
new coordinates, not a comparison law. The projections in
\eqref{f16v14:generator} have a finite-range square-root-density realization,
a common null vector, and constants as their joint kernel. Their finite-regulator
gap $\widehat\lambda_P$ is positive.

Pulled back to the unmoving coordinates, an $x_{i,b}$ update changes only
$x_{i,b}$ and the ports $h_{i-1,b},h_{i,b}$ that exist. A port update changes
only its corresponding $h_{t,b,v}$. No global update is hidden in this
parameterization.
\end{proposition}
\begin{proof}
At fixed $X$, each map $r_{t,v}\mapsto h_{t,v}$ is a left and right Haar
translation. Its inverse is explicit. Fubini therefore gives an exact
measure change, including the constant Haar factors. The chord weights
$j(x/\sqrt\gamma)^{w_i}$ and the two original Gaussian endpoint amplitudes
are unaltered. Principal cuts have Haar measure zero and do not restrict the
group translations.

Use the normalized product of the inherited one-cube chord reference factors
and the individual port Haar probabilities as reference measure. The native
density relative to this reference is proportional to the exponential of the
original classical action, expressed in $(X,r)$. It is bounded above and
below by positive finite constants at a finite regulator. The action words
still have bounded spatial and temporal range, since every frame uses one
cube only. V13's conditional square-root calculation
\eqref{f16v13:local-parent-kernel} now applies to the new coordinate blocks:
each conjugated projection depends only on its finite word neighbourhood and
kills the same square-root density. Product variance comparison proves a
positive, possibly very poor finite-regulator gap and the constant kernel,
exactly as in Proposition~\ref{f16v13:parent}.

Finally inspect \eqref{f16v14:coordinates}. Changing $X_{i,b}$ can change the
frame at that slice only, and hence just the two indicated bonds. In V13's
unmoving cell partition this touches at most two cells, at times $i-1,i$,
with unavailable endpoints omitted. Each unmoving cell is touched by only a
fixed number of these updates. This also proves finite support after pulling
back. The coordinate change is unitary between the two probability $L^2$
spaces, but the old and new lists of conditional subspaces are different;
no equality of their generators or gaps follows.
\end{proof}

\subsection{The Gaussian pair calculation now has a positive margin}

In principal moving logarithms $z$, the Gaussian precision is the separated
one from V4,
\[
                 W_n\ \oplus\ (I_n\otimes L_h).
\]
Use whole $x_{i,b}$ blocks and individual three-colour $z_{t,b,v}$ blocks,
precisely the linear limit of the new native update partition. Mean shifts
change no Gaussian conditional correlations. Put $N=\mathsf M^{-1}$ for one
cube, and let $\mathsf H$ be its six-face Gram. The diagonal chord blocks are
\begin{equation}
 K_{\rm int}=\mathsf H+4N,\qquad
 K_{\rm end}=\tfrac12\mathsf C^{-1}+\tfrac12\mathsf H+3N.
 \label{f16v14:diagonals}
\end{equation}
These include the actual mixed diagonal $2N$ and the dressed endpoints.
The temporal off-diagonal chord block is $-N$; spatial off-diagonal chord
blocks at the same time come from one crossed face only.

\begin{proposition}[A rational face bound with the endpoint cost retained]
\label{f16v14:face}
Let $U=\mathsf J+\mathsf B\mathsf F$ be the inherited balanced edge embedding,
and let $V_F$ restrict its rows to the four edges of any cube face. Then
\begin{equation}
 V_F^{\mathsf T}V_F\preceq\frac9{80}K_{\rm int},\qquad
 K_{\rm end}\succeq\frac{39}{40}K_{\rm int},\qquad K_{\rm int}\succeq8N.
 \label{f16v14:face-bounds}
\end{equation}
Consequently the Gaussian pair maximal correlations for the new blocks are
bounded by $3/26$ for spatial chord neighbours and by $5/39$ for temporal
chord neighbours. Chord--port correlations are zero. Neighbouring individual
ports on the same fine spatial bond graph have pair correlation exactly $1/6$;
other port pairs have zero pair correlation.
\end{proposition}
\begin{proof}
Here is an explicit exact certificate for the face constant. In the inherited
chord order $(a,b,c,d,e)$, the rational matrices are
\[
 N=\frac1{24}\begin{pmatrix}
10&-4&1&-5&5\\-4&10&-1&5&-5\\1&-1&10&-2&-4\\
-5&5&-2&10&-4\\5&-5&-4&-4&10
\end{pmatrix},\qquad
 \mathsf H=\begin{pmatrix}
2&-1&0&-1&1\\-1&2&0&1&-1\\0&0&2&0&-1\\
-1&1&0&2&-1\\1&-1&-1&-1&2
\end{pmatrix}.
\]
For the face with first coordinate zero, in the positive edge order,
\[
 V_F=\frac1{24}\begin{pmatrix}
 -4&-2&5&-1&1\\-1&1&-4&-4&-2\\
 -2&-4&-5&1&-1\\1&-1&10&-2&-4
 \end{pmatrix}.
\]
Direct rational multiplication gives
\begin{equation}
 \det\!\left(tI-K_{\rm int}^{-1}V_F^{\mathsf T}V_F\right)
       =t(t-1/32)^2(t-1/60)(t-9/80).
 \label{f16v14:face-polynomial}
\end{equation}
The product is similar to a positive symmetric matrix. Thus its largest
eigenvalue is $9/80$. Cube automorphisms preserve the cycle-space inner
product and the face-action form and permute the six faces, up to oriented
edge signs. Transporting the chord coordinates gives a simultaneous
congruence, so the same generalized eigenvalues hold on every face. The
supplied checker also reconstructs and verifies all six faces separately.

Whitening by $N$ makes $\mathsf H$ have eigenvalues $t=4,6$, with their
inherited multiplicities. The whitened endpoint block is the function
$\sqrt{t+t^2/4}+t/2+3$. At $t=4$, its value is $\sqrt8+5\ge39/5$,
and at $t=6$ it is $\sqrt{15}+6\ge39/4$. These follow by squaring
$\sqrt8\ge14/5$ and $\sqrt{15}\ge15/4$. They prove the second bound;
the third follows from $t\ge4$.

For any two Gaussian blocks of positive precision diagonals $K_1,K_2$ and
cross block $C_{12}$, their conditional maximal correlation is
$\|K_1^{-1/2}C_{12}K_2^{-1/2}\|$. To verify that this is an upper as well
as a lower bound, whiten and take a singular-value decomposition of the cross
block. The law becomes a product of correlated scalar Gaussian pairs. The
conditional-expectation singular values on Hermite polynomials are products
of powers of the scalar correlations. On the centered space the largest is
the largest scalar correlation; Hermite density completes the argument.
This is also the two-block instance of V13's all-chaos analysis.

Across one spatial interface the cross block is $-V_+^{\mathsf T}OV_-$,
where $O$ matches the four edge orientations and is orthogonal. The face
bound and the endpoint comparison give a correlation at most
$(40/39)(9/80)=3/26$. For the temporal block $-N$, both diagonal blocks
are at least $(39/5)N$, giving $5/39$. There are no Gaussian cross-species
blocks. Finally, for individual nonroot ports the diagonal block of $L_h$
is $6I_3$, and each fine neighbouring nonroot port contributes $-I_3$.
A two-port conditional therefore has scalar correlation $1/6$ in each
colour. Fixed root neighbours contribute to the diagonal, not additional
random coordinates. The asserted zeros follow from the exact precision
support, not from marginal independence of all chord coordinates.
\end{proof}

\begin{theorem}[A successful balanced Gaussian local certificate]
\label{f16v14:Gaussian-margin}
Give every chord-cube update weight one. For a port at relative vertex
$v\in\{0,1\}^3\setminus\{0\}$, give its update weight
\begin{equation}
 w(v)=\begin{cases}7/10,&|v|_1=1,\\9/10,&|v|_1=2,\\1,&|v|_1=3.\end{cases}
 \label{f16v14:weights}
\end{equation}
The Gaussian weighted pair-row margin for this partition is at least $2/39$.
In particular its weighted and unweighted rate-one-per-block generators have
gap at least $2/39$ on all of their Gaussian $L^2$ spaces, uniformly in
$B,n$ and the prescribed bridge history.
\end{theorem}
\begin{proof}
A chord block has at most six spatial and two temporal chord neighbours.
Its row deficit is at least
\[
        1-6(3/26)-2(5/39)=2/39.
\]
An endpoint has fewer temporal neighbours, so the same safe bound applies.
There are no cross-species contributions at Gaussian order.

A fine nonroot port of Hamming weight one has four nonroot neighbours of
weight two and two fixed roots. A port of weight two has four neighbours of
weight one and two of weight three. A port of weight three has six neighbours
of weight two. Internal and bridge edges both count; every port has fine
degree six. With pair coefficient $1/6$, the three weighted deficits are
\[
 \frac7{10}-\frac46\frac9{10}=\frac1{10},\quad
 \frac9{10}-\frac16\left(4\frac7{10}+2\right)=\frac1{10},\quad
 1-\frac66\frac9{10}=\frac1{10}.
\]
Apply the projection anticommutator proof of
Theorem~\ref{f16v13:angle-gap} to these new Gaussian blocks. It works on the
complete $L^2$ space. Its constant kernel follows from Gaussian equivalence
to product Lebesgue measure, or the inherited all-chaos proof. The unweighted
form dominates the weighted one since all weights are at most one. This
proves the claim without a claim about the \emph{native} pair angles.
\end{proof}

This calculation separates a parameterization artifact from a native
obstruction. It repairs V13's failed Gaussian preflight; it does not justify
uniform convergence of full conditional-expectation operators over arbitrary
compact exterior data. The following exact test shows that such convergence
is false here.

\subsection{A native two-port cancellation including the fixed roots}

Identify $\SU(2)$ with the unit quaternions, write $\Re U$ for the scalar
part, and retain $s(U)=1-\Re U$. At one temporal bond fix all chord groups to
$I$ at its two ends and take the bridge history to be zero. The native
port-dependent action at that bond is exactly
\begin{equation}
      \gamma\sum_{\{u,v\}\text{ fine spatial edge}}s(r_u^{-1}r_v),
                  \qquad r_o=I\text{ at every cube root}.
 \label{f16v14:spin-action}
\end{equation}
Magnetic terms and chord endpoint amplitudes have no port dependence. This
is a particular actual conditional of the present law, not a new spin-model
replacement for its general background.

For neighbouring nonroot ports $c,d$, let $\mathcal H_c^{\setminus d}$ be
the sum of its five exterior neighbour quaternions, excluding $d$, and
define $\mathcal H_d^{\setminus c}$ similarly. Adjacent vertices in the cubic
graph have disjoint exterior-neighbour sets. A nonroot vertex of parity
weight one has two root neighbours; a vertex of weight two or three has none.
These facts hold without aliasing for $m\ge4$.

\begin{proposition}[Explicit compact field cancellation]
\label{f16v14:cancellation}
For every neighbouring nonroot pair there are exterior port values, respecting
all fixed roots, such that
$\mathcal H_c^{\setminus d}=\mathcal H_d^{\setminus c}=0$.
At those values the pair conditional is exactly
\begin{equation}
 d\pi_\gamma(R,S)=z_\gamma^{-1}e^{-\gamma s(R^{-1}S)}\,dH(R)dH(S).
 \label{f16v14:free-pair}
\end{equation}
Both marginals are probability Haar. Its maximal correlation is at least
\begin{equation}
 \rho_\gamma=\int\Re V\,k_\gamma(V)\,dH(V)
                      \ge1-512/\gamma\qquad(\gamma\ge1).
 \label{f16v14:pair-lower}
\end{equation}
The last constant is conservative and is not a numerical integral enclosure.
\end{proposition}
\begin{proof}
Let
\[
 I=(1,0,0,0),\qquad
 T_+=(-1/2,1/2,1/2,1/2),\qquad
 T_-=(-1/2,-1/2,-1/2,-1/2).
\]
All are unit quaternions and $T_++T_-=-I$. If there are no root neighbours,
assign the five neighbours the values $I,-I,I,T_+,T_-$. Their sum is zero.
If there are two fixed root neighbours, leave them $I$ and assign the three
free neighbours $-I,T_+,T_-$. The sum is again zero. The two neighbour
sets are disjoint, so these assignments are compatible. All other free
ports can be set to $I$.

For a fixed neighbour $A$, $s(R^{-1}A)=1-\Re(R^{-1}A)$ is affine in $A$.
Each five-neighbour contribution is therefore the constant five. The only
remaining variable action is $\gamma s(R^{-1}S)$, proving
\eqref{f16v14:free-pair}, including its normalizer. Equivalently take $R$
Haar and $S=RV$ with independent $V$ of law $k_\gamma dH$.
Conjugation symmetry gives $\mathbb EV=\rho_\gamma I$. Thus the functions
$2\Re R$ and $2\Re S$ have zero mean, variance one, and covariance
$\rho_\gamma$. The variance follows from rotational invariance of probability
Haar on the unit three-sphere: each of its four real coordinates has second
moment $1/4$.

Here are explicit bounds for $1-\rho_\gamma=\mathbb E_{k_\gamma}s(V)$.
The normalized angular Haar density is $(2/\pi)\sin^2\theta\,d\theta$ on
$[0,\pi]$. For $\gamma\ge1$, use $\theta\le\gamma^{-1/2}$,
$\sin\theta\ge2\theta/\pi$ and $s\le\theta^2/2$ to obtain
\[
 z_\gamma\ge\frac{8e^{-1/2}}{3\pi^3}\gamma^{-3/2}.
\]
Use $s\ge2\theta^2/\pi^2$, $s\le\theta^2/2$ and
$\sin\theta\le\theta$ for the numerator. Its upper bound is
\[
 \frac1\pi\int_0^\infty\theta^4e^{-2\gamma\theta^2/\pi^2}\,d\theta
       =\frac{3\pi^{9/2}}{32\sqrt2}\gamma^{-5/2}.
\]
Their ratio is $9e^{1/2}\pi^{15/2}/(256\sqrt2\,\gamma)$.
Using $e^{1/2}<2$, $\sqrt\pi<2$, $\sqrt2>1$, and $\pi<22/7$ bounds
this coefficient by $36(22/7)^7/256<512$. This proves
\eqref{f16v14:pair-lower}. No Gaussian approximation of the pair law was used.
\end{proof}

\begin{theorem}[The balanced native worst-exterior row still fails]
\label{f16v14:native-obstruction}
For zero bridge history, $m\ge4$, $n\ge1$ and $\gamma\ge4096$, the native
worst-exterior correlations for the two port pairs on the chain
\[
              (1,0,0),\quad(1,1,0),\quad(1,1,1)
\]
at any one temporal bond are at least $7/8$. No positive update weights can
make the worst-exterior pair-row margin positive for the balanced partition.
This is true despite Theorem~\ref{f16v14:Gaussian-margin}.
\end{theorem}
\begin{proof}
Apply the preceding construction separately to the two pairs. An essential
supremum for each pair is allowed to use its own exterior values; simultaneous
maximization is not required in a row bound. At finite coupling the conditional
expectations of the bounded continuous test functions $2\Re R,2\Re S$, and
of their products, depend continuously on the local compact exterior values
near these configurations and on chord inputs near the identity. Their
variances are nonzero there. Each such exterior neighbourhood has positive
native marginal probability. Values $-I$ may be represented by a principal
chart cut, but nearby group neighbourhoods minus the cut still have positive
Haar measure. Hence the essential supremum is at least the value in
\eqref{f16v14:pair-lower}, not just a value at a null conditioning point.

At $\gamma\ge4096$ this is at least $7/8$. Positive row margins would require
$w_1>(7/8)w_2$, $w_3>(7/8)w_2$, and
$w_2>(7/8)(w_1+w_3)$ after dropping all other nonnegative terms. They imply
$w_2>(49/32)w_2$, a contradiction. The calculation involves one existing
port bond only, so it also applies to a history of length one. It does not
bound the gap of the full resampling generator from above.
\end{proof}

The general slow zero-staple channel was already known in f15 V86. The new
conclusion is that it survives the present balanced \emph{two-port}
conditional even with the periodic rooted pins, and defeats this revised
worst-exterior certificate on the shortest history. A zero-field Laplace test
cannot upper-bound those native worst-exterior correlations.

\subsection{The actual cancellation witness has a uniformly available repair}

We now let one of the exterior variables move. This is an estimate for the
native generator, not for its Gaussian approximation. For general fixed
$X,\mathbf y$ and other ports, the conditional of one moving port $r_p$ has
the exact form
\begin{equation}
 d\pi_p(R\mid\text{rest})
       =\frac{e^{\gamma\Re(R^{-1}\mathcal H_p)}}{Z_p(\mathcal H_p)}dH(R),
                         \qquad |\mathcal H_p|\le6.
 \label{f16v14:one-port}
\end{equation}
Each of its six neighbouring terms supplies one unit-quaternion summand of
$\mathcal H_p$. For an oriented fine edge $e=(p,q)$ that summand is
$\bar U_e(i)r_q\bar U_e(i+1)^{-1}$; for $e=(q,p)$ it is
$\bar U_e(i)^{-1}r_q\bar U_e(i+1)$. Here $\bar U$ is the actual moving
fine representative, whose bridges are $Z$ and whose internal edges are
$\bar X$. These formulas follow by cyclicity and inversion of the real trace.
They retain all noncommuting factors. Neither magnetic weights nor chord
amplitudes depend on $r_p$.

Let
\[
 C_-=\{R:|R+I|\le1/24\},\qquad
 C_+=\{R:|R-I|\le1/24\},\qquad
 \mathcal R_p=\{\Re\mathcal H_p\ge1,\ r_p\in C_-\}.
\]
The field condition is measurable in the complement of $p$. These are
quaternion chord distances, not rescaled logarithmic distances.

\begin{theorem}[A same-law, rate-one repair floor on a compact witness family]
\label{f16v14:repair}
For every fixed bridge history and every $\gamma\ge1$,
\begin{equation}
 \pi_p(C_-\mid\text{rest})\le e^{-3\gamma/2}
                         \quad\text{when }\Re\mathcal H_p\ge1.
 \label{f16v14:escape}
\end{equation}
For any $f\in L^2(P^{\rm bal})$ supported on $\mathcal R_p$,
\begin{equation}
 \widehat{\mathcal E}(f)\ge
  \|\widehat Q_p f\|_2^2
   \ge(1-e^{-3\gamma/2})\|f\|_2^2\ge\tfrac12\|f\|_2^2.
 \label{f16v14:repair-floor}
\end{equation}
For the weights \eqref{f16v14:weights}, the weighted form has the analogous
lower bound $7\|f\|_2^2/20$.

This repair event contains a neighbourhood of an explicit two-port
zero-staple cancellation family, uniformly over all values of its two free
pair ports. No probability lower bound on that neighbourhood is required.
\end{theorem}
\begin{proof}
When $\Re\mathcal H_p\ge1$, the one-port exponent without $\gamma$ is at
most $-1+6/24=-3/4$ on $C_-$ and at least $1-6/24=3/4$ on $C_+$.
These two caps have the same positive Haar volume, by $R\mapsto-R$.
Lower-bound $Z_p$ by its integral on $C_+$ and upper-bound the numerator on
$C_-$. Their ratio proves \eqref{f16v14:escape}, without estimating the
normalizer by a different conditional law.

At a fixed exterior where $f$ can be nonzero, conditional Cauchy--Schwarz
and its support give
\[
  |\mathbb E_p f|^2\le\pi_p(C_-)\,\mathbb E_p|f|^2,
  \qquad \operatorname{Var}_p f\ge(1-e^{-3\gamma/2})\mathbb E_p|f|^2.
\]
Integrate in the actual exterior marginal. One term of the generator proves
\eqref{f16v14:repair-floor}. Since $e^{-3/2}<1/2$, the stated constants
follow, with $w_p\ge7/10$ for the weighted form.

To verify the final statement on the literal lattice, choose the pair
$c=(1,1,0)$, $d=(1,1,1)$. Assign each of their five exterior neighbours,
in increasing coordinate order, the values $I,-I,I,T_+,T_-$. On a torus of
side $2m$ this means, at the two relevant rows,
\[
\begin{array}{c|ccccc}
 & (0,1,0)&(1,0,0)&(1,1,2m-1)&(1,2,0)&(2,1,0)\\
\text{value}&I&-I&I&T_+&T_-\\ \hline
 &(0,1,1)&(1,0,1)&(1,1,2)&(1,2,1)&(2,1,1)\\
\text{value}&I&-I&I&T_+&T_-
\end{array}
\]
Leave all remaining ports at $I$ and all chord groups at $I$, with zero
bridges. The two pair exterior fields are zero as before. Take the repair
port $p=(1,0,0)$. Its five neighbours other than $c$ have sum $3I$:
$(0,0,0),(1,0,1),(1,0,2m-1),(1,2m-1,0),(2,0,0)$ carry
$I,-I,I,I,I$. Thus
\[
                 r_p=-I,\qquad\mathcal H_p=3I+r_c,
                 \qquad\Re\mathcal H_p\ge2
\]
for \emph{every} $r_c,r_d\in\SU(2)$. The roots in this list are already
fixed $I$, consistently with the prescription. Compactness of the two free
pair groups and continuity of the finite word list give a uniform exterior
and chord neighbourhood with $\Re\mathcal H_p>1$; restrict $r_p$ to $C_-$.
This proves that the entire pair-fibre witness has a repair neighbourhood.
The neighbourhood may be chosen in angular group coordinates. It need not
lie in a bounded thermal window, and no such restriction was used in the
one-port estimate.
\end{proof}

The support qualification in \eqref{f16v14:repair-floor} is essential. It is
not the false inequality $\|1_{\mathcal R_p}f\|_2^2\le C\widehat{\mathcal E}(f)$
for arbitrary $f$, which constants would contradict. A global argument must
pay localization commutators or use a compatible block comparison. What is
proved is that the explicit slow \emph{frozen pair} does not itself form a
low-energy supported state for the full resampling dynamics: the boundary
port it froze is a fast escape direction.

\subsection{Low exterior fields have a paid density cost, not a gap consequence}

For a neighbouring free pair $c,d$ at one port bond, define the two exterior
fields by omitting their mutual edge from \eqref{f16v14:one-port}. They are
sums of five transported unit quaternions even at general bounded bridge
history. Let
\[
 \mathcal B_{cd}=\{|\mathcal H_c^{\setminus d}|\le1\}
                     \cup\{|\mathcal H_d^{\setminus c}|\le1\}.
\]
Let $N_{\rm low}$ count these events over all neighbouring nonroot port
pairs and all temporal bonds, with each unoriented pair counted once.

\begin{proposition}[Actual-law low-field pair density]
\label{f16v14:density}
For sufficiently large $\gamma$ depending on the fixed bridge window,
\begin{equation}
                \frac{\mathbb E_P N_{\rm low}}{B(n+1)}
                                      \le C_r/\gamma.
 \label{f16v14:density-bound}
\end{equation}
This is an expectation under the complete native law, with every exterior
field integrated. It is not a bound on the probability of all-time good
fields, on a specified pair uniformly in location, or on source-weighted
Dirichlet defects.
\end{proposition}
\begin{proof}
If a sum of five unit quaternions has norm at most one, its distance from
$5I$ is at least four. By telescoping products in its five summands,
\[
 4\le\sum_{q\ne d}|I-A_{cq}|
      \le C\gamma^{-1/2}
           \left(r+\sum_{j\in S_{cd}}(|x_j|+|z_j|)\right),
\]
where $S_{cd}$ is a fixed-size neighbouring carrier in the two internal
slices and the one port bond. The internal moving edges are products of
exponentials of fixed cube-linear chord combinations, so
$|I-\bar X_e|\le C\gamma^{-1/2}\sum|x_j|$ globally on their principal
chord inputs. Also $|I-r_p|\le|z_p|/\sqrt\gamma$, and bridge distances are
at most $r/\sqrt\gamma$. Thus, after increasing the fixed-$r$ threshold,
$\mathcal B_{cd}$ implies
$\sum_{j\in S_{cd}}(|x_j|^2+|z_j|^2)\ge c\gamma$.

Every coordinate belongs to a bounded number of these carriers. The exact
relation $r_t=f_t^{-1}h_tf_{t+1}$ and the bi-invariant geodesic triangle
inequality give the inherited global bound
$\sum(|x|^2+|z|^2)\le C\sum(|x|^2+|a|^2)$, as in V5's bad-label proof.
V5's \emph{actual} second-moment density then bounds its expectation by
$C_r\mathcal N$. Sum the deterministic local implication and divide by
$\gamma\mathcal N$. This proves \eqref{f16v14:density-bound}. We reuse the
moment theorem; no new pointwise moment or independence assertion is made.
\end{proof}

Even this vanishing density does not suppress a worst-exterior correlation:
a small open event of positive probability already contributes to an
essential supremum. Nor may its probability be multiplied by a uniform bound
on arbitrary normalized test functions. The repair theorem controls a
supported function by its \emph{own} resampling energy instead.

\subsection{The directional source reduction survives the changed updates}

The old covariance matrices, posterior law, and Haar-compensated sources
remain those of V7 and V12. They are simply composed with
\eqref{f16v14:coordinates} when evaluated under $P^{\rm bal}$. Define the
native moments $m_2(P),m_4(P)$ in the original unmoving coordinates exactly
as in \eqref{f16v13:moment-inputs}. Uniform bounds on these moments are not
assumed established.

\begin{proposition}[Balanced-update innovation energy and posterior handoff]
\label{f16v14:energy}
For the primitive remainder bank $e=(e^{\rm br},e^{\rm pt})$ and the full
corrected source vector $\mathbf r$ of V7,
\begin{equation}
 \begin{aligned}
 \widehat{\mathcal E}(b^{\mathsf T}e)
       &\le C_r\frac{m_4(P)}\gamma\|b\|^2,\qquad
 \widehat{\mathcal E}(v^{\mathsf T}\mathbf r)
       \le C_r\frac{m_4(P)}\gamma\|v\|^2,\\
 \widehat{\mathcal E}(b^{\mathsf T}S_0)
       &\le C m_2(P)\|b\|^2.
 \end{aligned}
 \label{f16v14:energy-bound}
\end{equation}
Consequently V13's three directional bounds
\eqref{f16v13:directional-reduction} hold with $\widehat\lambda_P$ in place
of $\lambda_P$, and an enlarged geometric constant. They concern exactly
the same $\Gamma_\gamma,\mathsf R_{\rm post},\mathsf L_D$ and full memory.
No comparison between $\widehat\lambda_P$ and $\lambda_P$ is asserted.
\end{proposition}
\begin{proof}
Use stationary exact resampling in the balanced variables. Both the original
configuration and its updated configuration, pulled back to $(x,a)$, have
law $P$. By Proposition~\ref{f16v14:same-law}, each update changes only a fixed
number of unmoving cells, and each cell is changed by a fixed number of
update types. Each primitive remainder depends on a fixed finite carrier.
Thus every update changes a bounded number of primitive components, and
each component can change under only a bounded number of updates.

For the changed components use Cauchy--Schwarz in that fixed set and
stationarity to bound
$\mathbb E|e_j(\zeta)-e_j(\zeta')|^2\le4\mathbb E_P e_j^2$.
The complete-chart bound \eqref{f16v13:primitive} makes the latter at most
$C_rm_4(P)/\gamma$. The resampling identity
$\widehat{\mathcal E}(f)=\frac12\sum_a\mathbb E|f-f^{(a)}|^2$ now gives
the first estimate without a volume factor. The fixed bounded synthesis
$b=(v,\mathsf P v)$ gives the second. For an affine port score write
$b^{\mathsf T}S_0=l^{\mathsf T}\zeta+$constant with $\|l\|\le h_+\|b\|$.
Only finitely many cell terms change at one update. Apply the same argument
to those terms using $m_2(P)$; the bounded update incidence proves the third
estimate. Principal-logarithm jumps are allowed: stationarity, not a
small-displacement approximation, bounds them.

Apply the native Poincare inequality of the new finite-regulator generator
to these energies. V13's posterior proof then uses only the unchanged exact
lift-mismatch identities, one common exterior conditional expectation, and
its $L^2(P)$ contraction. Repeating those steps, not replacing that posterior,
gives the asserted handoff. Its use still needs a uniform native gap and
moments to obtain a useful uniform conclusion.
\end{proof}

\subsection{An integrated defect criterion that does not freeze the repair variables}

Let $w_a$ be \eqref{f16v14:weights} and set
$\widehat{\mathcal L}_w=\sum_a w_a\widehat Q_a$,
$\widehat{\mathcal E}_w(f)=\langle f,\widehat{\mathcal L}_w f\rangle$.
Define a deterministic symmetric table $\Theta_{ab}$ by the Gaussian
majorants $3/26,5/39,1/6$ above, and zero elsewhere. It has weighted row
margin $\delta_0=2/39$. The \emph{native} interaction graph can have additional
pairs, including cross-species pairs; retain them, with $\Theta_{ab}=0$.
For every pair in that finite-range graph let $\theta_{ab}(z)$ be its actual
conditional maximal correlation at exterior $z$, and define
\[
 d_{ab}(z)=(\theta_{ab}(z)-\Theta_{ab})_+.
\]
These nonnegative functions are measurable in the pair exterior. At a finite
regulator the normalized conditional kernels are bounded operators; they
provide a measurable version of the correlations. Put
\begin{equation}
 \mathfrak D(f)=\sum_{a<b}w_aw_b\,
 \mathbb E_P\left[d_{ab}(\omega_{(a,b)^c})
       \bigl(|\widehat Q_af|^2+|\widehat Q_bf|^2\bigr)\right].
 \label{f16v14:defect}
\end{equation}
Here $\omega=(X,r)$ denotes the balanced coordinates. The notation $P$ identifies
the same full law, not its Gaussian reference. Noninteracting pairs have no
defect and can be omitted.

\begin{theorem}[Same-law integrated defect identity and sufficient repair target]
\label{f16v14:integrated}
At each finite regulator, on the complete native $L^2$ space,
\begin{equation}
 \|\widehat{\mathcal L}_w f\|_2^2
       \ge\delta_0\widehat{\mathcal E}_w(f)-\mathfrak D(f).
 \label{f16v14:defect-identity}
\end{equation}
If, independently, constants $0\le u<\delta_0$ and $v\ge0$ are proved to obey
\begin{equation}
 \mathfrak D(f)\le u\widehat{\mathcal E}_w(f)
                         +v\|\widehat{\mathcal L}_w f\|_2^2
                              \quad\text{for every }f,
 \label{f16v14:repair-target}
\end{equation}
then
\begin{equation}
             \widehat\lambda_P\ge\frac{\delta_0-u}{1+v}.
 \label{f16v14:repair-gap}
\end{equation}
Neither \eqref{f16v14:escape} nor \eqref{f16v14:density-bound} alone proves
\eqref{f16v14:repair-target}; it is the still-unverified global input.
\end{theorem}
\begin{proof}
Condition on the exterior of a pair. V13's two-projection argument bounds
its anticommutator below by
$-\theta_{ab}(z)(\widehat Q_a+\widehat Q_b)$ on that fibre. Multiply by
$w_aw_b$ and integrate before taking any supremum. An exterior-measurable
coefficient commutes with both pair updates, so its quadratic form is its
expectation times the corresponding squared innovation, exactly as in
\eqref{f16v14:defect}. Use $\theta_{ab}\le\Theta_{ab}+d_{ab}$, expand
$\widehat{\mathcal L}_w^2$, and sum its deterministic part. Its coefficient
on $w_a\widehat Q_a$ is at least the verified margin $\delta_0$. This proves
\eqref{f16v14:defect-identity} on all $L^2$; all operators and coefficients
are bounded at finite regulator.

Combining with \eqref{f16v14:repair-target} gives
$(1+v)\widehat{\mathcal L}_w^2\succeq
(\delta_0-u)\widehat{\mathcal L}_w$.
The joint kernel is constants. Spectral calculus yields a gap at least
$(\delta_0-u)/(1+v)$, and the unweighted generator dominates the weighted
one. This proves \eqref{f16v14:repair-gap} without assuming the desired
posterior covariance. The deficit has not been multiplied by the probability
of a bad event independently of its test function.
\end{proof}

This criterion keeps the innovation where a conditional pair is bad and
allows the energy of the \emph{other} updates to pay for it. The explicit
repair port shows why that distinction is necessary. A complete proof would
need to classify or cover the relevant defect configurations, pay the
commutators introduced by such a localization, and control the resulting
constants in \eqref{f16v14:repair-target}. We have not shown that low staple
magnitude exhausts the defects: chord--port interactions and other native
conditional landscapes remain in \eqref{f16v14:defect}.

\subsection{Status and the next noncircular calculation}

There are three distinct conclusions. The old Gaussian row obstruction is
removed by exact local balanced updates, with an explicit successful margin.
The native worst-boundary row is nevertheless disproved by a compact
cancellation compatible with the rooted geometry. That cancellation family
has a native, rate-one repair event with a uniform support-energy floor.
These facts identify which variable a useful local block must retain:
freezing the cancellation's repair port produces a slow pair that the full
update leaves with the stated uniformly positive probability.

The next test should therefore include the boundary-repair variables in a
local block or in an integrated defect estimate, rather than improve the
same impossible worst-pair bound. The defect form and the adapted directional
energy estimate supply exact objects for that calculation. This is not a
claim that a finite block theorem has already been applied successfully.
Uniform native fourth moments, a global repair/auxiliary-gap bound, temporal
weighted-row decay of the posterior, and a leading-coefficient estimate at
unrestricted volume are still unproved. The existing physical bridge-window,
same-top transfer, infrared coercivity, and continuum obligations are
unchanged. No supported repair floor is identified with a physical mass.

\subsection{Diagnostics and exact predecessor preservation}

The accompanying script is
\begin{center}\small\path{verify_ym_f16_v14_balanced_repair.py}.\end{center}
It reconstructs the inherited cube matrices; checks the six rational face
characteristic polynomials, endpoint comparisons and weighted rows; checks
literal fine-torus roots and compatible quaternion cancellation assignments;
and verifies the repair-port sums, cap constants, and the integrated defect
algebra on finite positive laws. The report distinguishes exact diagnostics
from exploratory floating-point calculations. The latter are not proof
certificates and are not used for a theorem constant. Native continuum
integrals are bounded analytically above, not by an unreported quadrature.

The new diagnostic passed 622 exact assertions. Of these, 446 check
cancellation and compatibility with fixed roots on two finite tori. The available unchanged
V4--V13 diagnostics passed 9,467, 89,339, 906, 311, 609, 990, 608, 463,
1,201 and 222 assertions, respectively. The assertion families and script
hashes are recorded in the accompanying reports. These finite fixtures do
not formally verify the analytic proofs or establish a global native gap.

The cumulative source changes only V13's title version and inserts this
appendix before the final document-ending command. Removing the exact append
and restoring that title recovers the predecessor byte for byte. The package
records source hashes, preservation checks, compilation and visual review.
No inaccessible archive checker, external formalization or complete inherited
proof collection is described as recertified. The scheduled companion
checkpoint remains V15.

\begin{firewall}
V14 proves a successful balanced Gaussian pair certificate, an exact native
compact-boundary obstruction to that worst-exterior test, and a uniform
native support-energy repair floor on an explicitly identified witness
family. It also proves a low-field pair-density bound, adapted directional
source energies, and an integrated defect reduction with its global repair
input left explicit. The native law and all original normalizers are
unchanged. No uniform full-native auxiliary gap, unrestricted-volume coherent
posterior contraction, native local fourth-moment bound, physical transfer
gap, or interacting continuum Yang--Mills mass gap is established here.
\end{firewall}
% END F16 V14 APPEND
% BEGIN F16 V15 APPEND -- 2026-09-17
% Insert immediately before the final document-ending command in V14.
% No predecessor proof or bibliography is silently revised.

\clearpage
\section{V15: a complete three-port repair and a same-clock conditional gap}
\label{f16v15:context}

V14 identifies an escape port for a slow frozen pair, but its repair energy
floor applies only to functions supported on a specified event. Localizing an
arbitrary innovation to that event introduces an unpaid term. We avoid that
localization: integrate the escape port together with the pair and prove a
Poincar\'e inequality for \emph{every} function on this three-group conditional.
The final inequality uses the original rate-one single-port updates. A
collective update appears only in an explicitly paid intermediate comparison.

The exterior hypothesis is a fixed open angular neighbourhood of the actual
cancellation witness, not an all-exterior assertion. Within it, all three
ports range over their complete groups, including antipodal configurations.
The neighbourhood does not shrink with coupling, spatial volume, or duration.
Its radius and entry coupling are proved to exist but are not numerically
enclosed. This is a local conditional gap, not a global native or physical gap.

\subsection{The V15 checkpoint and the non-duplication boundary}

The V14 appendix and analytic audit were read in full. The supplied f14 V100,
f15 V100, Grand Tensor Monograph V076, and this lineage through V14 were
inventoried and relevant dependencies checked. Two semantic archive searches
were empty; this was not treated as evidence of absence. In particular, f14
V2 already proves same-clock \emph{supported} repair mechanisms. F15 V71 proves
a thermal conditional \emph{diffusion} gap on its own good-boundary law. Neither
is the three-moving-port resampling result below. The monograph's physical
Perron and vacuum-corrected Feshbach interfaces remain separate. This is a
targeted audit, not independent recertification of the entire archives.

The active obstruction is V14's integrated conditional-angle defect, not V3's
Gaussian inverse or V10's formal coefficients. The wider-literature memorandum's
compatible-local-constraint proposal is applied here to the exact auxiliary
operator. The other memorandum's distinction between order gain and locality
cost is retained: this iteration enlarges the class of test functions and
retains a previously frozen variable; it does not append a formal order to the
same uncontrolled error.

For context, Qin--Wang, \href{https://arxiv.org/abs/2208.11299}{\emph{Spectral
Telescope: Convergence Rate Bounds for Random-Scan Gibbs Samplers Based on a
Hierarchical Structure}}, Section 3.2 and Theorem 2, develops hierarchical
conditional-gap comparisons on general spaces. Its definitions and theorem
were checked, including the theorem page. The two-level comparison needed
below is proved directly with our clock, not quoted with a discrete random-scan
normalization. Liu--Wong--Kong,
\href{https://doi.org/10.1093/biomet/81.1.27}{\emph{Covariance structure of the
Gibbs sampler with applications to the comparisons of estimators and
augmentation schemes}} (1994), provides classical grouping/collapsing context;
only its primary summary was checked. No literature theorem identifies these
sampling gaps with a physical transfer gap.

Only supplied companions were available at this scheduled checkpoint. Other
programs named in older ledgers were not supplied or recertified. The next
five-version checkpoint is V20.

\subsection{The exact native three-port conditional}

Keep V14's balanced probability $P^{\rm bal}$. At one temporal bond use the
three nonroot fine vertices
\[
 c=(1,1,0),\qquad d=(1,1,1),\qquad p=(1,0,0).
\]
Their induced graph has exactly the edges $cd,cp$ for $m\ge4$. Condition every
other port and all chord groups. Each native port-edge interaction is bilinear
in its two quaternions, with fixed left/right unit factors. Since these two
edges form a tree, separate fixed left/right translations of the two leaves
make both couplings ordinary quaternion inner products. Keep the central port
unchanged and denote the resulting variables by $U,V,W$, respectively. These
changes preserve probability Haar and the three individual conditional update
operations; they are independent of the three free variables.

Write $U\cdot V=\Re(U^{-1}V)$ and $I=(1,0,0,0)$. The exact conditional is
\begin{equation}
 d\pi_\gamma^{A,B,D}(U,V,W)
 =\frac{e^{\gamma[U\cdot V+U\cdot W+A\cdot U+B\cdot W+D\cdot V]}}
 {\mathcal Z_\gamma(A,B,D)}\,dH(U)dH(V)dH(W),
 \label{f16v15:triplet}
\end{equation}
for exterior-dependent real quaternion fields with $|A|\le4$, $|B|,|D|\le5$.
The letter $D$ here is a four-vector, not V11's spatial cylinder. Each root
remains fixed. Exterior-only constants cancel on conditional normalization,
not in the exterior posterior. Magnetic and dressed-end factors have no
free-port dependence and are not altered. All kinetic terms meeting these
ports are included.

\begin{proposition}[Retaining the actual repair variable]
\label{f16v15:normal-form}
At V14's displayed witness, with only $c,d,p$ now left free,
\begin{equation}
                         A=I,\qquad B=3I,\qquad D=0.
 \label{f16v15:witness-fields}
\end{equation}
Freezing $W=-I$ recovers the unpinned two-port law proportional to
$e^{\gamma U\cdot V}dH(U)dH(V)$. The full three-port law instead has the unique
energy minimum $(I,I,I)$. The aligned fields depend continuously on the finite
exterior word data near the witness.
\end{proposition}
\begin{proof}
V14 writes every oriented edge contribution as the dot product of one port
with a left/right translate of its neighbour. Absorb the two such translates
on the leaves. A tree has no residual loop constraint. The central vertex has
four fixed neighbours, each leaf five, giving the stated field bounds.

At the witness all local chord frames and bridges are identities. V14's five
neighbours of $c$, excluding $d$, sum to zero and include $p=-I$. Removing $p$
leaves field $I$. The five exterior neighbours of $d$ sum to zero; the five
neighbours of $p$ other than $c$ sum to $3I$. This proves
\eqref{f16v15:witness-fields}. Setting $W=-I$ cancels $U\cdot W+I\cdot U$.
With all three variables free, the deficit from the maximum $6\gamma$ is
\[
 \gamma\{s(U^{-1}V)+s(U^{-1}W)+s(U)+3s(W)\}.
\]
It is nonnegative and vanishes exactly at $(I,I,I)$. Continuity follows from
the finite quaternion products. Group neighbourhoods at $-I$ remain legitimate
even though one principal logarithm has a Haar-null cut there.
\end{proof}

The following results hold on a fixed parameter set
\begin{equation}
 \mathcal T_\eta=\{(A,B,D):|A-I|\le\eta,\ |B-3I|\le\eta,\ |D|\le\eta\}.
 \label{f16v15:parameter-window}
\end{equation}
We will choose $0<\eta_0\le1/100$ independently of $\gamma,B,n$. The preimage of
its strict interior is a genuine open exterior neighbourhood of the witness.
No condition is imposed on the three free groups. Not every low-field or
large-correlation configuration is claimed to belong to this neighbourhood.

\subsection{The marginal denominators require their own compact estimate}

For a positive unnormalized pair density $f$ on $S^3\times S^3$, set
\begin{equation}
 m_x(x)=\int f(x,y)dH(y),\quad m_y(y)=\int f(x,y)dH(x),\quad
 \mathcal C(x,y)=\frac{f(x,y)}{\sqrt{m_x(x)m_y(y)}}.
 \label{f16v15:normalized-kernel}
\end{equation}
The total partition function cancels \emph{exactly}. Between the fixed Haar
$L^2$ spaces this is the normalized conditional-expectation kernel; its top
singular value is one with singular vectors the square roots of the normalized
marginals. Its second singular value is the pair maximal correlation. This
normalization is inherited from V13, not a new covariance convention.

\begin{proposition}[Uniform complete-chart normalized-kernel limit]
\label{f16v15:Laplace}
Let $f_{\gamma,\vartheta}=a_{\gamma,\vartheta}e^{\gamma F_\vartheta}$ on
$S^3\times S^3$, with $\vartheta$ in a compact parameter space. Assume the
family $F_\vartheta$ is continuous in the $C^3$ topology, with uniformly bounded
higher derivatives whenever they are used, and that the positive amplitudes
are uniformly bounded above and below and converge uniformly in $C^2$ to
positive $a_{0,\vartheta}$. Define
\[
 M_\vartheta(x)=\max_y F_\vartheta(x,y),\quad
 N_\vartheta(y)=\max_x F_\vartheta(x,y),\quad
 \mathcal D_\vartheta(x,y)=\tfrac12(M_\vartheta(x)+N_\vartheta(y))-F_\vartheta(x,y).
\]
Suppose $\mathcal D_\vartheta$ has exactly one zero
$(x_\vartheta,y_\vartheta)$, and the negative joint Hessian of $F_\vartheta$
there is uniformly positive. Then the kernels \eqref{f16v15:normalized-kernel},
in complete normal charts centered there and scaled by $\sqrt\gamma$, with
half-Haar Jacobians and zero extension, converge uniformly in
Hilbert--Schmidt norm to the normalized joint-Gaussian kernels with that
Hessian as precision. Their maximal correlations converge uniformly to the
corresponding Gaussian correlations.
\end{proposition}
\begin{proof}
The two terms $M-F,N-F$ are nonnegative. Their simultaneous zero means that
both coordinates are conditional maximizers. Every global maximizer is such
a zero, hence is the specified unique point. The points vary continuously:
compactness and continuity identify every subsequential limit of maximizers.
Use continuous tangent frames on finitely many local parameter neighbourhoods.

Write the negative joint Hessian in these frames as
$H=\left(\begin{smallmatrix}A_0&C_0\\C_0^{\mathsf T}&B_0\end{smallmatrix}\right)$.
At the joint point each individual conditional maximum is unique: any other
partial maximum with the other coordinate fixed would also be a global
maximum. Compactness excludes distant competing partial maxima in a small
neighbourhood. The implicit function theorem then describes the local partial
maxima using the positive diagonal Hessian blocks. Their quadratic deficits give
\begin{equation}
 \mathcal D(u,v)=\frac14\left(
 |u+A_0^{-1}C_0v|_{A_0}^2+
 |v+B_0^{-1}C_0^{\mathsf T}u|_{B_0}^2\right)+O(|(u,v)|^3).
 \label{f16v15:deficit-Hessian}
\end{equation}
Both squares zero imply $H(u,v)=0$, so this quadratic form is strictly
positive. Compactness makes its lower bound uniform. Away from a small
neighbourhood, continuity of the maxima and the unique-zero condition give
a uniform positive bound. Thus, in quaternion chord distances,
\begin{equation}
 \mathcal D_\vartheta(x,y)\ge c(|x-x_\vartheta|^2+|y-y_\vartheta|^2).
 \label{f16v15:deficit-pin}
\end{equation}
This argument does not require smoothness of $M,N$ away from their joint point.

A Hessian bound about \emph{any} partial maximizer, together with the positive
amplitude lower bound and integration on a ball of radius $\gamma^{-1/2}$,
gives the global marginal lower bounds
\[
 m_x(x)\ge c\gamma^{-3/2}e^{\gamma M_\vartheta(x)},\qquad
 m_y(y)\ge c\gamma^{-3/2}e^{\gamma N_\vartheta(y)}.
\]
They remain true for degenerate or multiple partial maxima away from the joint
point. Consequently
\[
 |\mathcal C_{\gamma,\vartheta}(x,y)|
       \le C\gamma^{3/2}e^{-\gamma\mathcal D_\vartheta(x,y)}.
\]
A complete normal chart on $S^3$ centered at a unit quaternion has
$dH=(2\pi^2)^{-1}\gamma^{-3/2}j(u/\sqrt\gamma)du$ on
$|u|<\pi\sqrt\gamma$. Both square-root factors in the kernel cancel
$\gamma^{3/2}$, and $j\le1$. Chord distance satisfies
$|x-x_\vartheta|\ge2|u|/(\pi\sqrt\gamma)$ on this entire chart. Thus the
zero-extended transformed kernel has the uniform majorant
\begin{equation}
                         C e^{-c'(|u|^2+|v|^2)}.
 \label{f16v15:normalized-majorant}
\end{equation}
In particular neither small marginal densities nor compact chart tails have
been discarded.

On every fixed rescaled ball, ordinary Taylor expansion about the moving
maximizer and Laplace expansion of the two marginals give the joint Gaussian
and its \emph{own} marginal normalizers. For these marginal expansions, the
partial Hessians are uniformly positive nearby; outside a fixed neighbourhood
the compact conditional phase has a uniform loss. These observations justify
uniform Laplace convergence there. The amplitude tends to its positive value
at the joint point and cancels from the normalized limit. The Gaussian
half-Jacobian normalization is also retained. Uniform convergence on balls
and the squared majorant \eqref{f16v15:normalized-majorant} prove the uniform
Hilbert--Schmidt limit, including the full compact complement. Singular-value
stability yields the last assertion. The Gaussian second singular value is
$\|A_0^{-1/2}C_0B_0^{-1/2}\|$: a singular-value decomposition of the whitened
precision and the Hermite basis give this value on all of centered $L^2$.
\end{proof}

The normalized deficit $\mathcal D$ is essential. A pin for $\max F-F$ alone
would not exclude another simultaneous conditional maximum carrying a large
normalized correlation in a small marginal region. We verify the stronger
condition for each of the following two reductions.

\subsection{The inner conditional remains controlled for every value of the third port}

Condition $V$ in \eqref{f16v15:triplet}, and put $C=A+V$. The free pair is
\begin{equation}
 d\pi_\gamma(U,W\mid V)\ \propto\
                e^{\gamma[U\cdot W+C\cdot U+B\cdot W]}dH(U)dH(W).
 \label{f16v15:inner-pair}
\end{equation}
The term $D\cdot V$ cancels here. In particular the worst old value $V=-I$
is not omitted.

\begin{proposition}[Uniform inner pair correlation]
\label{f16v15:inner}
For $|A-I|,|B-3I|\le1/100$ there is a finite $\gamma_1$, independent of
$A,B,V$, such that for all $V\in S^3$ and $\gamma\ge\gamma_1$ the maximal
correlation of \eqref{f16v15:inner-pair} is at most $3/4$. Thus its two
rate-one single-port resampling updates have gap at least $1/4$ on their
complete conditional $L^2$ space.
\end{proposition}
\begin{proof}
Here $\Re C\ge-1/100$, $|C|\le201/100$. The unique conditional maximizer in
$W$ is $W_B(u)=(B+u)/|B+u|$. For $|u|\le1$, $|B+u|\ge199/100$.
For $B=3I$ one has
\[
 \frac{3+u_0}{|3I+u|}\ge\frac{\sqrt8}{3}\ge\frac{14}{15}.
\]
Indeed $(3+u_0)^2-8|u_\perp|^2\ge(3u_0+1)^2\ge0$. Changing $B$ by
at most $1/100$ changes the normalized vector by at most $1/100$, using
$|3I+u|\ge2$ and
$|x/|x|-y/|y||\le2|x-y|/|y|$. Hence
\begin{equation}
           \Re W_B(u)\ge14/15-1/100>23/25.
 \label{f16v15:inner-cap}
\end{equation}
The same lower real part holds on line segments between two such vectors.
Thus $|C+w|\ge91/100$ on those segments.

The derivative of vector normalization has norm $1/|x|$. Applying it first
along the segment between two $u$'s in the unit ball, and then along the
segment between their $W_B$ values, proves that the best-response composition
\[
                u\longmapsto\frac{C+W_B(u)}{|C+W_B(u)|}
\]
is a contraction of $S^3$ in chord distance with constant at most
\begin{equation}
                          q_*:=\frac{10000}{18109}<1.
 \label{f16v15:inner-contraction}
\end{equation}
The sphere is complete. The contraction theorem gives a unique fixed point,
therefore a unique simultaneous conditional maximizer. Conversely every
global maximizer supplies such a point. No conditional-maximizer set is lost
by division, since \eqref{f16v15:inner-cap} makes $C+W_B(u)$ nonzero.

At this point the negative joint tangent Hessian is
\[
 \begin{pmatrix}aI_3&-T\\-T^{\mathsf T}&bI_3\end{pmatrix},\qquad
 a=|C+W|\ge91/100,\quad b=|B+U|\ge199/100,\quad\|T\|\le1.
\]
It has a uniform positive lower bound because $ab\ge18109/10000>1$, and all
upper bounds are fixed. Its Gaussian correlation is at most
$100/\sqrt{18109}<3/4$, the latter strict inequality being
$160000<9\cdot18109$. The compact parameter family consists of the two closed
field balls and every $V\in S^3$. Proposition~\ref{f16v15:Laplace}, with
amplitude one, applies uniformly. Its strict Gaussian margin proves the native
$3/4$ bound eventually.

For completeness, the standard two-projection implication does not change
the clock: for a pair with correlation $\theta$, the two centered
single-coordinate subspaces have synthesis norm at most $\sqrt{1+\theta}$.
Thus the sum of the two conditional-expectation projections is at most
$(1+\theta)I$ on the centered space, and the sum of the two rate-one
innovations is at least $(1-\theta)I$. This is V13's pair argument on the
present complete conditional.
\end{proof}

\subsection{Collapsing the repair port retains a nontrivial generated weight}

Integrate $W$ exactly. Set
\[
 Z(t)=\int_{S^3}e^{t\Re R}dH(R),\qquad t\ge0.
\]
The unnormalized marginal of $(U,V)$ is
\begin{equation}
 f_\gamma^{A,B,D}(U,V)
   =e^{\gamma[U\cdot V+A\cdot U+D\cdot V]}Z(\gamma|B+U|).
 \label{f16v15:collapsed}
\end{equation}
This generated factor must not be replaced by one. For $|B-3I|\le1/100$,
$199/100\le|B+U|\le401/100$. Define
\begin{equation}
 \begin{split}
 F_{A,B,D}(U,V)&=U\cdot V+A\cdot U+D\cdot V+|B+U|,\\
 a_{\gamma,B}(U)&=\gamma^{3/2}e^{-\gamma|B+U|}Z(\gamma|B+U|).
 \end{split}
 \label{f16v15:collapsed-phase}
\end{equation}
Then $f_\gamma=\gamma^{-3/2}a_{\gamma,B}e^{\gamma F}$. The common scalar
$\gamma^{-3/2}$ cancels in \eqref{f16v15:normalized-kernel}, but the amplitude
$a_{\gamma,B}$ does not disappear from the finite-coupling law.

\begin{proposition}[A complete outer-pair bound near the witness]
\label{f16v15:outer}
There exist $0<\eta_0\le1/100$ and $\gamma_2<\infty$ such that, throughout
$\mathcal T_{\eta_0}$ and for $\gamma\ge\gamma_2$, the actual collapsed pair
\eqref{f16v15:collapsed} has maximal correlation at most $2/3$. At the exact
witness its limiting correlation is $2/\sqrt{11}$.
\end{proposition}
\begin{proof}
First check the amplitude, including the uniformity used below. Angular Haar
measure is $(2/\pi)\sin^2\theta\,d\theta$ on $[0,\pi]$, so
\[
 e^{-t}Z(t)=\frac2\pi\int_0^\pi
               e^{-t(1-\cos\theta)}\sin^2\theta\,d\theta.
\]
Laplace expansion at zero gives
\begin{equation}
 a_{\gamma,B}(U)\longrightarrow
             \sqrt{2/\pi}\,|B+U|^{-3/2}
 \label{f16v15:amplitude-limit}
\end{equation}
uniformly with every fixed number of derivatives in $U,B$. To justify the
needed two derivatives directly, put $t=\gamma r$, with $r$ in the compact
positive interval above, and scale $\theta=v/\sqrt\gamma$ in the small-angle
integral. Differentiation in $r$ introduces powers of
$\gamma(1-\cos\theta)$; these have uniform polynomial Gaussian majorants
using $1-\cos\theta\ge2\theta^2/\pi^2$. The complement of a fixed small-angle
interval is exponentially small, also after these derivatives. Taylor
expansion and dominated convergence give the stated uniform derivative
limits; composition with the smooth $r=|B+U|$ gives the claim. The leading
constant is $(2/\pi)\int_0^\infty v^2e^{-rv^2/2}dv
=\sqrt{2/\pi}\,r^{-3/2}$. The amplitudes therefore have fixed positive upper
and lower bounds for sufficiently large $\gamma$.

We next verify the normalized deficit, not merely the total energy minimum.
At $(A,B,D)=(I,3I,0)$, write
\[
 a=1-U_0,\quad b=1-V_0,\quad d=1-U\cdot V,\qquad
 T(U)=U_0+\sqrt{10+6U_0}.
\]
Here $F(U,V)=U\cdot V+T(U)$, and $M(U)=1+T(U)$. With
$N(V)=\max_U F(U,V)$ and $E=\mathcal D(U,V)$, the nonnegative two
conditional gaps imply $2E\ge d$. Also $N(V)\ge V_0+5$ by testing $U=I$.
Concavity of the square root yields
$T(U)\le5-(7/4)a$. Consequently
\begin{equation}
                    2E\ge2d+\tfrac74 a-b.
 \label{f16v15:outer-pin-inequality}
\end{equation}
Chord-distance Young inequalities give
$b\le\tfrac32 a+3d$ and $b\le2a+2d$. From the first, the last displayed
bound implies $2E\ge a/4-d$; combine it with $2E\ge d$ to get $E\ge a/16$.
Since also $E\ge d/2$, one obtains the global bound
\begin{equation}
                  \mathcal D(U,V)\ge\frac{a+b}{52}.
 \label{f16v15:outer-global-pin}
\end{equation}
Thus the normalized deficit has exactly the zero $(I,I)$.

In ordinary orthonormal tangent coordinates at that point, the negative
Hessian of $F$ is
\begin{equation}
       H_{\rm out}=\begin{pmatrix}11/4&-1\\-1&1\end{pmatrix}\otimes I_3.
 \label{f16v15:outer-Hessian}
\end{equation}
Indeed $U\cdot V=1-|u-v|^2/2+O(|(u,v)|^3)$ and
$T(\operatorname{Exp}u)=5-7|u|^2/8+O(|u|^4)$. The Hessian is positive
and its Gaussian correlation is $2/\sqrt{11}<5/8$.

We now make this a fixed \emph{open-field} statement. The phases in
\eqref{f16v15:collapsed-phase} are smoothly dependent on $(A,B,D)$ near the
witness, since $|B+U|$ stays positive. Uniform convergence of phases implies
uniform convergence of $M,N$ and of their normalized deficits. The strict
pin \eqref{f16v15:outer-global-pin} places every zero of a nearby deficit in
a fixed small normal-coordinate neighbourhood of $(I,I)$. Choose this
neighbourhood convex in its Euclidean coordinates and small enough that the
joint Hessian of the baseline phase is strictly negative there. It remains
strictly negative for all sufficiently small parameter perturbations.
Every zero of a normalized deficit is a critical point of the joint phase.
Strict concavity permits at most one such point there; a global maximum
always provides one. Hence nearby deficits have exactly one zero and a
uniformly positive negative joint Hessian there. By continuity, shrink a
fixed $\eta_0>0$ so that every corresponding Gaussian pair correlation is at
most $5/8$.

The compact parameter set $\mathcal T_{\eta_0}$ and the amplitudes
\eqref{f16v15:amplitude-limit} now satisfy Proposition~\ref{f16v15:Laplace}.
Uniform Hilbert--Schmidt convergence makes the native second singular value
at most $5/8+1/24=2/3$ for $\gamma\ge\gamma_2$. At the exact witness its
limit is the specified $2/\sqrt{11}$. This proves an upper correlation
bound on the \emph{whole} pair $L^2$ space, not just convergence of a few
moments. All free compact tails and both marginal denominators were paid.
\end{proof}

\subsection{The full three-port gap in the original update clock}

For a probability $\pi$ on these three variables, write
\[
 \mathcal E_U(f)=\mathbb E_\pi\operatorname{Var}_\pi(f\mid V,W),\quad
 \mathcal E_V(f)=\mathbb E_\pi\operatorname{Var}_\pi(f\mid U,W),\quad
 \mathcal E_W(f)=\mathbb E_\pi\operatorname{Var}_\pi(f\mid U,V).
\]
The block quantity $\mathcal E_{UW}(f)=
\mathbb E_\pi\operatorname{Var}_\pi(f\mid V)$ is not one of the original
single-port updates. The next proof explicitly bounds it by those updates.

\begin{theorem}[Complete three-port conditional repair, for every test function]
\label{f16v15:triplet-gap}
There are $0<\eta_0\le1/100$ and a finite $\gamma_0$, depending only on this
fixed local conditional, such that for every $\gamma\ge\gamma_0$, every
$(A,B,D)\in\mathcal T_{\eta_0}$, and every $f\in L^2(\pi_\gamma^{A,B,D})$,
\begin{equation}
 \boxed{\quad
 \operatorname{Var}_{\pi_\gamma^{A,B,D}}f
       \le12\{\mathcal E_U(f)+\mathcal E_V(f)+\mathcal E_W(f)\}.
 \quad}
 \label{f16v15:three-port-gap}
\end{equation}
In particular the three original single-port resamplings, each at rate one,
have gap at least $1/12$. No support assumption on $f$, no free-port cutoff,
and no boundary probability estimate is used. The constants are independent
of the ambient spatial volume and history length.
\end{theorem}
\begin{proof}
The exact star density gives $W\perp V\mid U$. For a centered function
$h(U,W)$ and a centered function $g(V)$, this conditional independence gives
\[
 \mathbb E[h(U,W)g(V)]
   =\mathbb E[\mathbb E(h\mid U)\,g(V)].
\]
Conditional expectation decreases the $L^2$ norm. Hence maximal correlation
between $(U,W)$ and $V$ is no greater than that between $U$ and $V$ under the
collapsed law. The reverse inequality follows by taking $h$ to depend only
on $U$, so the two correlations are equal. Proposition~\ref{f16v15:outer}
and the two-projection argument therefore give
\begin{equation}
          \operatorname{Var}_\pi f\le
                  3\{\mathcal E_{UW}(f)+\mathcal E_V(f)\}.
 \label{f16v15:outer-block-gap}
\end{equation}
At every fixed $V$, Proposition~\ref{f16v15:inner} applies to the actual
conditional $(U,W)\mid V$, uniformly over all $V\in S^3$. Integrating its
variance inequality gives
\[
                 \mathcal E_{UW}(f)
                      \le4\{\mathcal E_U(f)+\mathcal E_W(f)\}.
\]
Insert this into \eqref{f16v15:outer-block-gap}. The coefficient on
$\mathcal E_V$ is actually three, so the common constant twelve is safe.
Take $\gamma_0=\max(\gamma_1,\gamma_2)$, increasing it if necessary for the
amplitude bounds. Conditional variance disintegration and continuity of
conditional projections extend the inequalities from bounded functions to
all $L^2$. Haar alignment of the two leaves is a separate coordinate bijection
on each free port and hence conjugates its single-port projection exactly.
Thus the result uses the original three port updates, not a changed clock.
\end{proof}

For comparison only, the complete leading precision of the exact witness is
\[
 \begin{pmatrix}3&-1&-1\\-1&1&0\\-1&0&4\end{pmatrix}\otimes I_3.
\]
Its determinant per colour is seven. Its rate-one single-port Gaussian gap is
$1-\sqrt{5/12}$ by V13's all-chaos formula. The Schur complement after
integrating $W$ is exactly \eqref{f16v15:outer-Hessian}. The proof of the
native theorem did not simply substitute this Gaussian gap: the two complete
normalized-kernel limits and the uniform conditional bound over $V$ were
needed before the comparison. V14's slow pair froze the variable that supplies
this additional pin.

\subsection{An exterior gate that acts on all native innovations}

Let $\mathcal R=\{c,d,p\}$ be the free triple, or a translated/permuted copy
whose induced graph and exterior field conditions are as above. At every
fixed exterior perform the specified leaf alignment and compute $(A,B,D)$.
Let $G_{\mathcal R}$ be the event that these fields lie in
$\mathcal T_{\eta_0}$. It is measurable in the complement of the \emph{entire}
triple. In particular it imposes no condition on $r_p$.
Let
\[
 E_{\mathcal R}=\mathbb E_{P^{\rm bal}}[\,\cdot\mid\omega_{\mathcal R^c}],
 \qquad Q_{\mathcal R}=I-E_{\mathcal R}.
\]
Use $\widehat Q_a$ for the original individual balanced updates from V14.

\begin{theorem}[Posterior-integrated repair with no support restriction]
\label{f16v15:gated}
For $\gamma\ge\gamma_0$ and every $f\in L^2(P^{\rm bal})$,
\begin{equation}
 \boxed{
 \mathbb E_P[1_{G_{\mathcal R}}|Q_{\mathcal R}f|^2]
       \le12\sum_{a\in\mathcal R}
                   \mathbb E_P[1_{G_{\mathcal R}}|\widehat Q_a f|^2].}
 \label{f16v15:gated-form}
\end{equation}
The expectations use the actual exterior posterior with all original
normalizers. At zero bridge history $G_{\mathcal R}$ contains a fixed open
neighbourhood of V14's witness in its exterior variables. The inequality
holds at general fixed bridge histories wherever the same aligned-field
gate is satisfied.
\end{theorem}
\begin{proof}
Condition the full native probability on $\mathcal R^c$. Its conditional on
the three free ports is precisely \eqref{f16v15:triplet}, because every term
meeting them was retained. On $G_{\mathcal R}$ apply
Theorem~\ref{f16v15:triplet-gap}, and integrate with respect to the actual
exterior marginal. The conditional variance on the left is the integral of
$|Q_{\mathcal R}f|^2$; each conditional single-coordinate variance on the
right is the corresponding squared original innovation. The field gate
commutes with these four projections, since it is exterior to all of them.
The open-neighbourhood statement follows from
Proposition~\ref{f16v15:normal-form} and the strict parameter window. No
constant reference exterior law has replaced this marginal.
\end{proof}

This improves V14's supported-function floor in a specific way. Both sides of
\eqref{f16v15:gated-form} vanish on constants, and the estimate applies to
arbitrary $f$, not only $f$ supported in the gate. It does \emph{not} assert
$\mathbb E[1_G|f|^2]\le C\widehat{\mathcal E}(f)$, which would fail on
constants. Nor is a gate depending on one of the three free variables inserted
and then commuted through its update. The formerly restricted repair variable
is now among the variables being resampled.

\subsection{Compatible local blocks can be added without an unrecorded clock}

Take a finite collection of such triples, with deterministic rates
$\alpha_{\mathcal R}\ge0$, and let $w_a$ be V14's weights. Suppose
\begin{equation}
  \kappa=\sup_a w_a^{-1}
              \sum_{\mathcal R\ni a}\alpha_{\mathcal R}<\infty.
 \label{f16v15:overlap}
\end{equation}
A translated finite set of shapes with bounded rates has such a constant
independent of $B,n$. Overlap is allowed; no independence is assumed. Define
\[
 \mathcal B_{\rm rep}=\sum_{\mathcal R}\alpha_{\mathcal R}
                         1_{G_{\mathcal R}}Q_{\mathcal R},\qquad
 \widehat{\mathcal L}_{\rm aug}=\widehat{\mathcal L}_w+\mathcal B_{\rm rep}.
\]
Each gated summand is a positive orthogonal projection because its indicator
commutes with $Q_{\mathcal R}$. It annihilates constants and has finite
interaction support in the native square-root-density representation.

\begin{proposition}[Same-clock comparison for a family of repaired blocks]
\label{f16v15:clock-comparison}
At every admitted coupling,
\begin{equation}
 0\preceq\mathcal B_{\rm rep}\preceq12\kappa\widehat{\mathcal L}_w,
 \qquad
 \widehat{\mathcal L}_w\preceq\widehat{\mathcal L}_{\rm aug}
             \preceq(1+12\kappa)\widehat{\mathcal L}_w.
 \label{f16v15:operator-comparison}
\end{equation}
Consequently an \emph{independently proved} gap $g>0$ for the augmented
operator would imply a gap at least $g/(1+12\kappa)$ for the original weighted
generator, and hence for the original unweighted balanced generator.
No such uniform augmented gap is asserted here.
\end{proposition}
\begin{proof}
The quadratic form of one gated projection is exactly the left side of
\eqref{f16v15:gated-form}. Apply that inequality, drop each gate only from
the nonnegative upper bound, sum the deterministic rates, and use
\eqref{f16v15:overlap}. This gives
\[
 \langle f,\mathcal B_{\rm rep}f\rangle
    \le12\sum_a\left(\sum_{\mathcal R\ni a}\alpha_{\mathcal R}\right)
                         \|\widehat Q_af\|_2^2
    \le12\kappa\langle f,\widehat{\mathcal L}_wf\rangle.
\]
Adding the original form proves \eqref{f16v15:operator-comparison}. Its common
kernel is constants since the augmented form contains the original one.
Comparison on their common centered space proves the stated implication.
No products of noncommuting gated operators or independence between different
triples were used. In particular block acceleration is not counted as free.
\end{proof}

\subsection{What has advanced, and the next global obligation}

The main new input is a full conditional $L^2$ theorem in place of a supported
repair floor. The slow two-port cancellation is not removed by a cutoff or
by changing its probability; its repairing boundary variable is retained.
A nontrivial collapsed normalizer, two complete normalized-kernel estimates,
and the original-clock comparison give the explicit constant $1/12$ on a
fixed open exterior-field neighbourhood. This is more than a Gaussian
preflight or a positive floor for a selected supported test function.

The conclusion does not cover every exterior and does not prove that these
gates exhaust V14's defects. Other port cancellation geometries, chord--port
correlations and chord-cube conditional landscapes remain. No probability
estimate for the gate is claimed or needed for the conditional theorem. A
small gate probability would not by itself pay the omitted spectral defect.
The exact global target is still V14's
\[
 \mathfrak D(f)\le u\widehat{\mathcal E}_w(f)
             +v\|\widehat{\mathcal L}_wf\|_2^2,
                       \qquad u<2/39.
\]
Neither \eqref{f16v15:gated-form} nor the comparison proposition establishes
that inequality for the uncovered sectors. An independent uniform gap for
$\widehat{\mathcal L}_{\rm aug}$ would be another route, but is not assumed.

The next noncircular calculation is therefore to organize compatible
repair-retaining blocks and bound the remaining defect sectors, keeping
actual gates and overlap costs. The present three-port result can enter that
calculation without a new supported-function localization error. Re-estimating
the old frozen pair cannot do so. V14's directional source-energy handoff is
unchanged; a uniform global native gap and a uniform native fourth moment
are still separate inputs. No posterior operator contraction or temporal
weighted-row theorem is newly inferred from the local conditional gap.
The prescribed bridge window, physical root Gauss restriction, same-top full
transfer comparison, physical-time calibration and interacting continuum
construction remain distinct obligations.

\clearpage
\subsection{Diagnostics, checkpoint preservation, and claim scope}

The accompanying exact diagnostic is
\begin{center}\small\path{verify_ym_f16_v15_three_port_repair.py}.\end{center}
It reconstructs the literal three-port fields and fixed-root constraints,
checks noncommuting tree-edge alignment, the full three-port and collapsed
Gaussian Hessians, the normalized-deficit algebra, and conditional-projection
and rate-one comparison identities on positive finite probability models.
All 434 exact assertions passed; the report records their families and source
hash.
These tests are not native integral enclosures or proof-assistant verification
of the normalized-kernel convergence argument. In particular $\eta_0$ and
$\gamma_0$ are not numerically certified, and no simulation is used to justify
a theorem constant.

The available V4--V14 diagnostics were rerun unchanged and passed 9,467,
89,339, 906, 311, 609, 990, 608, 463, 1,201, 222 and 622 assertions,
respectively. The package records those reports, input hashes, exact predecessor
recovery, compilation and rendered-page review.
The V15 cumulative source changes only V14's title version and inserts this
appendix before its final document-ending command. Removing that insertion
and restoring the title recovers V14 byte for byte. No earlier mathematical
claim or bibliography string is silently corrected. The V15 companion audit
covers only the supplied sources; no inaccessible historical checker or
external formalization is described as rerun. The next checkpoint is V20.

\begin{firewall}
V15 proves a uniform $1/12$ conditional gap for the three original rate-one
single-port updates on a fixed open exterior neighbourhood of an actual native
cancellation witness, with all free compact groups and all $L^2$ test functions
retained. It gives a posterior-integrated gated form inequality and an explicit
same-clock comparison for overlapping families of those blocks. It does not
prove gate coverage, a uniform full-native auxiliary gap, native local fourth
moments, unrestricted-volume coherent posterior contraction, a physical transfer
gap, or an interacting continuum Yang--Mills mass gap. The clock remains an
auxiliary resampling clock, not physical Euclidean time.
\end{firewall}
% END F16 V15 APPEND
% BEGIN F16 V16 APPEND -- 2026-09-17
% Insert immediately before the final document-ending command in V15.
% The predecessor body and bibliography strings are not edited.

\clearpage
\section{V16: local superstability and uniform native fast-field moments}
\label{f16v16:context}

V15 proves a complete three-port repair on one exterior neighbourhood, not
coverage of the integrated defect. Adding more neighbourhoods without a
coverage mechanism risks repeating that distinction. We take up the other
explicit input in V13--V15: uniform local fourth moments of the \emph{actual}
full fast-history law. The independent compact anchors already available in
V5 are stronger for this purpose than the global entropy comparison used in
V6--V12. They survive arbitrary finite sets of free space--time cells, not
only entire-time spatial cylinders.

We prove uniform local exponential moments, including a cube of inhomogeneous
quadratic sources. The proof first pays a conditional normalizer by a product
of fixed-radius cell balls. This gives a local energy inequality with an
\emph{actual random boundary-energy term}. Iterating that inequality through
successive balls, and using polynomial graph growth, removes the boundary
term after averaging in the same law. No small interaction coefficient,
translation invariance, conditional gap, or Gaussian replacement is used.

Consequences include the previously missing native fourth-moment bound,
exponential suppression of the specific low-field configurations of V14,
and $O(\gamma^{-1})$ residual and posterior-return \emph{trace densities} at
unrestricted volume. The global auxiliary gap remains unproved. A rare
configuration can still carry a large conditional correlation or a normalized
spectral test function; the probability estimates below do not erase that
obstruction.

\subsection{Dependencies, literature, and the non-duplication boundary}

The supplied V15 appendix and analytic audit were read in full. The current
lineage and the mounted f14/f15 archives and Grand Tensor Monograph V076 were
searched, with the relevant normalization, pin, and moment interfaces read.
The targeted archive search returned no semantic matches; direct text search
and identified source ranges were used instead. This is not independent
recertification of the entire archive. F15 V50 explicitly distinguishes its
total-action bound from a joint local bound on an arbitrary background. F15
V90 bounds local physical energy through a specified physical transfer-source
input. F14 V67 obtains joint tails for a different restricted conditional.
None of those conclusions is transplanted to the current fixed-bridge law.

For classical context, D.~Ruelle,
\href{https://doi.org/10.1007/BF01609400}{\emph{Probability estimates for
continuous spin systems}}, Commun. Math. Phys. 50 (1976), 189--194,
extends superstability probability estimates to continuous spins. The primary
publisher abstract and bibliographic record were checked here; its complete
proof is not invoked or described as reviewed. The present finite-graph
argument is proved below from its explicit hypotheses. ``Local
superstability'' here refers to those hypotheses and their moment consequence,
not a claim that every standard continuum-particle hypothesis has been verified.
The high-temperature or superquadratic-spin functional-inequality literature
is not used to supply a missing native gap.

This is the upstream alternative to another formal correction order. V5's
one-cell anchors, V7's exact Haar residual synthesis, and V12--V14's same-law
source and posterior identities are inherited. The new analytic input is the
local-energy argument under quadratic tilts and its uniform moment conclusion.
No NS assertion is needed. The memoranda's distinction between locality,
normalization, and order gain is retained; the next companion checkpoint is
still V20.

\subsection{Complete cells, local interactions, and a uniform source domain}

Return to the \emph{unmoving} complete coordinates of V5. Let
\[
 \mathcal V=\{(i,b):0\le i\le n,\ b\in(\mathbb Z/m\mathbb Z)^3\},
 \quad B=m^3,\quad m\ge4,\quad n\ge1,\quad
 \mathcal N=|\mathcal V|=B(n+1).
\]
A cell contains $\zeta_c=(x_{i,b},a_{i,b})$ for $i<n$ and just $x_{n,b}$
at the last time. Its dimension is 36 or 15. Set
\[
 e_c=|\zeta_c|^2,\qquad e_S=\sum_{c\in S}e_c,\qquad
 a_0=\frac1{6\pi^2},\qquad t_0=\frac{a_0}{4}=\frac1{24\pi^2}.
\]
The symbol $e_c$ in this appendix denotes cell energy, not a primitive
source-error component. Retain the bridge assumption
$\max_{i,e}|y_{i,e}|\le r<\infty$ and, for definiteness, take
$\gamma\ge\gamma_r:=1+r^2$. This keeps these bridge coordinates strictly
inside their principal charts. The moment argument has no further
volume- or duration-dependent entry threshold.

Let $d\nu_{\gamma,c}$ be the \emph{unnormalized} product of the original
complete principal-chart Lebesgue factors and their Haar powers in cell $c$.
These include chord powers $w_i=1/2$ at the dressed time ends, chord powers
one elsewhere, and port powers one. Thus
\[
 d\nu_\gamma(\zeta)=\prod_{c\in\mathcal V}d\nu_{\gamma,c}(\zeta_c)
                 =1_{\mathcal P'}J^f_\gamma(\zeta)\,d\zeta.
\]
No growing Haar-volume scalar is absorbed into a supposedly uniform constant.
The classical action has V5's exact decomposition
\begin{equation}
 U_\gamma=b_\gamma(\mathbf y)+\sum_c A_{\gamma,c}(\zeta_c)
                             +\sum_{\ell\in\mathcal L}V_{\gamma,\ell}(
                                  \zeta_{S_\ell};\mathbf y),
 \quad A_{\gamma,c}\ge a_0 e_c,\quad V_{\gamma,\ell}\ge0.
 \label{f16v16:decomposition}
\end{equation}
Here $S_\ell$ is a nonempty, uniformly bounded fast carrier. Word lengths,
carrier ranges, and the number of interactions incident to a cell are bounded
by fixed geometric constants. Fast-free terms remain in $b_\gamma$; it cancels
only when normalizing the fixed-bridge fast probability, not from the full
bridge-history weight.

For a real array $s=(s_c)$ with $|s_c|\le t_0$, define the exact auxiliary tilt
\begin{equation}
 dP_s=\frac{e^{\sum_c s_c e_c}}{\mathbb E_P e^{\sum_c s_c e_c}}\,dP,
 \qquad
 Z(s)=\int e^{-U_\gamma+\sum_c s_c e_c}\,d\nu_\gamma,
 \qquad P=P_0.
 \label{f16v16:tilts}
\end{equation}
All these are finite positive integrals at a fixed regulator. The sources
are moment-generating parameters, not new physical couplings or a replacement
for $P$. They leave the complete domains, endpoint amplitudes, and original
Haar factors unchanged. The tilted one-cell anchors remain at least
$3a_0 e_c/4$.

Join distinct cells if they occur together in a primitive carrier $S_\ell$.
Call this graph $\mathcal G$, adding no edges beyond a fixed geometric range.
Let $B_k(c)$ denote its ball of radius $k$, and let
$\partial S=\{d\notin S:\ d\sim c\text{ for some }c\in S\}$.
There is a fixed $C_{\rm gr}$ such that
\begin{equation}
 |B_k(c)|\le C_{\rm gr}(k+1)^4
       \quad(k\ge0),\qquad
 \partial B_k(c)\subset B_{k+1}(c)\setminus B_k(c).
 \label{f16v16:growth}
\end{equation}
Indeed each graph step changes spatial cube and time coordinates by a bounded
amount. A containing coordinate box gives this polynomial bound, including
wrapping and the two time endpoints. Saturation of a finite ball only improves
it. This argument does not assume a periodic time direction.

\begin{proposition}[Uniform local upper bounds and complete reference masses]
\label{f16v16:local-data}
There are $K_r<\infty$, $v_*>0$, and $G_*\ge1$, independent of $\gamma,B,n,s$,
such that
\begin{equation}
 A_{\gamma,c}\le K_r e_c,\qquad
 V_{\gamma,\ell}\le K_r\left(1+\sum_{c\in S_\ell}e_c\right),
 \label{f16v16:upper-quadratic}
\end{equation}
and, with $\mathbb B_c=\{|\zeta_c|\le1\}$,
\begin{equation}
 \nu_{\gamma,c}(\mathbb B_c)\ge v_*,\qquad
 \int e^{-a_0 e_c/2}\,d\nu_{\gamma,c}\le G_*.
 \label{f16v16:cell-mass}
\end{equation}
For every nonempty $S\subset\mathcal V$, the tilted conditional exponent is
exactly
\[
 U_{S,s}=\sum_{c\in S}(A_{\gamma,c}-s_ce_c)
           +\sum_{\ell:S_\ell\cap S\ne\varnothing}V_{\gamma,\ell}.
\]
It obeys $U_{S,s}\ge3a_0 e_S/4$, and on $\prod_{c\in S}\mathbb B_c$,
\begin{equation}
 U_{S,s}(\zeta_S;q)\le K'_r|S|+K'_r e_{\partial S}(q).
 \label{f16v16:trial-bound}
\end{equation}
Every crossing interaction is retained in this conditional exponent.
\end{proposition}
\begin{proof}
For unit quaternions $g_j$, $s(g)=|I-g|^2/2$ and telescoping give
\[
 \gamma s\left(\prod_{j=1}^k\operatorname{Exp}(w_j/\sqrt\gamma)\right)
 \le\frac12\left(\sum_{j=1}^k|w_j|\right)^2
 \le\frac{k}{2}\sum_{j=1}^k|w_j|^2.
\]
Every actual word in unmoving coordinates is a product of exponentials of
fixed local linear forms in fast inputs and the bounded bridge inputs.
Frames do not require taking a new logarithm of a product in this statement.
Bounded word length and local matrix norms prove
\eqref{f16v16:upper-quadratic}, with $r$ included in its constant. Endpoint
Gaussian quadratics are cube-local with a bounded coefficient matrix. The
anchor lower bound is V5's inherited cube and rooted-tree estimate, not a
new convexity assertion. Both bounds apply cell by cell.

On the unit cell ball every group input has angle at most one because
$\gamma\ge1$. Each factor $j(v)=(\sin|v|/|v|)^2$ is at least
$(\sin1)^2$, and a power $w\in[1/2,1]$ does not decrease this lower bound.
There are at most twelve group factors. The minimum of the positive
15- and 36-dimensional unit-ball volumes, times $(\sin1)^{24}$, is an
admissible $v_*$. Conversely $J^f_{\gamma,c}\le1$ on its entire chart, so
an ordinary Euclidean Gaussian integral gives an admissible $G_*$, for
example $\max\{1,(2\pi/a_0)^{18}\}$. No bound on the unweighted total
volume of a growing principal chart is used.

Conditioning $P_s$ on $S^c$ cancels exactly the exterior-only action terms,
exterior source factors, and exterior reference factors. The remaining
exponent is $U_{S,s}$. Every interaction meeting $S$ is included once.
Each such term has its exterior fast inputs in $\partial S$. On the
product of unit cell balls, bounded incidence bounds all free energies
and all constant bridge costs by $C_r|S|$, and bounds the exterior
contribution by $C_r e_{\partial S}$. The extra source term costs at most
$t_0|S|$ there. This proves \eqref{f16v16:trial-bound}. Positivity of all
non-anchor terms and $|s_c|\le t_0$ prove the lower bound for every exterior
value, not only for a bounded exterior window.
\end{proof}

\subsection{A local conditional normalizer with its boundary energy paid}

\begin{proposition}[Conditional exponential energy and an averaged boundary inequality]
\label{f16v16:conditional-energy}
There are constants $a_r,b_r\ge1$, independent of $B,n,\gamma,s$ and the
subset $S$, such that, at every admitted exterior value $q$,
\begin{equation}
 \log\mathbb E_{P_s}\left[e^{t_0 e_S}\mid\zeta_{S^c}=q\right]
 \le t_0 a_r|S|+t_0 b_r e_{\partial S}(q).
 \label{f16v16:conditional-mgf}
\end{equation}
Consequently, putting $u_c(s)=\mathbb E_{P_s}e_c$,
\begin{equation}
 \sum_{c\in S}u_c(s)\le a_r|S|+b_r\sum_{d\in\partial S}u_d(s).
 \label{f16v16:energy-boundary}
\end{equation}
Both expectations use the exact tilted native law. The conditional bound
itself is \emph{not} uniform in arbitrary exterior energy.
\end{proposition}
\begin{proof}
Write $Z_{S,s}(q)=\int e^{-U_{S,s}}\prod_{c\in S}d\nu_{\gamma,c}$.
Restrict its denominator only for a lower bound to the product of unit cell
balls, retaining every crossing term. Proposition~\ref{f16v16:local-data} gives
\[
 Z_{S,s}(q)\ge v_*^{|S|}
                  e^{-K'_r|S|-K'_r e_{\partial S}(q)}.
\]
The numerator with $e^{t_0 e_S}$ has exponent at most $-a_0 e_S/2$,
since $3a_0/4-t_0=a_0/2$. Integrating over the \emph{full} native charts,
and then enlarging only this positive upper bound to Euclidean space, gives
an upper bound $G_*^{|S|}$. Divide these two bounds to get
\eqref{f16v16:conditional-mgf}, for instance with
\[
 a_r\ge t_0^{-1}\bigl(\log G_* -\log v_*+K'_r\bigr),
 \qquad b_r\ge t_0^{-1}K'_r,
\]
where each constant is also enlarged to be at least one. Jensen's inequality
gives $t_0\mathbb E[e_S\mid q]\le\log\mathbb E[e^{t_0e_S}\mid q]$.
Integrate under the actual exterior marginal of $P_s$. This is
\eqref{f16v16:energy-boundary}. At a finite regulator every quantity is finite
by compactness, so neither truncation nor an unproved integrability estimate
has been used to start the argument.
\end{proof}

The denominator is the reason for the boundary term. Dropping it would give
a false all-exterior thermal estimate, contradicted by the unpinned pair
conditionals in V14. The next step averages and then removes this term; it
does not claim that the term vanishes pointwise.

\subsection{Polynomial growth closes the local moment inequality}

\begin{proposition}[Boundary iteration on a finite graph of polynomial growth]
\label{f16v16:ball-lemma}
Let a finite graph satisfy \eqref{f16v16:growth}. Suppose finite nonnegative
numbers $u_c$ satisfy, for every nonempty set $S$,
\[
 \sum_{c\in S}u_c\le a|S|+b\sum_{d\in\partial S}u_d,
 \qquad a,b>0.
\]
Set $\rho=b/(1+b)\in(0,1)$. Then every vertex obeys
\begin{equation}
 u_c\le M(a,b):=
 \frac{a C_{\rm gr}}{1+b}\sum_{j=0}^\infty\rho^j(j+1)^4
 =a C_{\rm gr}\frac{1+11\rho+11\rho^2+\rho^3}{(1-\rho)^4}.
 \label{f16v16:ball-bound}
\end{equation}
The interaction coefficient $b$ need not be smaller than one. No
translation symmetry or spectral-gap assumption is required.
\end{proposition}
\begin{proof}
Fix the vertex and write $F_k=\sum_{d\in B_k(c)}u_d$ and
$v_k=|B_k(c)|$. The boundary inclusion in \eqref{f16v16:growth} implies
\[
 F_k\le a v_k+b(F_{k+1}-F_k),\qquad
 F_k\le\rho F_{k+1}+\frac{a}{1+b}v_k.
\]
Iterating this scalar inequality gives, for every integer $R\ge1$,
\begin{equation}
 u_c=F_0\le\rho^R F_R+
                 \frac{a}{1+b}\sum_{j=0}^{R-1}\rho^j v_j.
 \label{f16v16:iteration}
\end{equation}
At this \emph{fixed finite graph}, $F_R\le\sum_d u_d<\infty$, so the first
term tends to zero as $R\to\infty$. This does not assume a uniform bound on
that total and does not interchange a volume limit with the iteration.
Apply the polynomial ball bound to the remaining sum. The identity
\[
 \sum_{j=0}^\infty\rho^j(j+1)^4
       =\frac{1+11\rho+11\rho^2+\rho^3}{(1-\rho)^5}
\]
follows by applying $(\rho\,d/d\rho+1)^4$ to $(1-\rho)^{-1}$.
Since $(1+b)^{-1}=1-\rho$, it proves \eqref{f16v16:ball-bound}.
A ball that has already saturated its connected component still satisfies
the same inequality with an empty exterior boundary. Thus small tori,
disconnected interaction graphs, and finite time endpoints cause no exception.
\end{proof}

The order of operations matters. A one-cell estimate alone would invite the
unsuccessful requirement $b\deg(\mathcal G)<1$. Instead the conditional
normalizer is estimated on every ball. The geometric decrease in
\eqref{f16v16:iteration} then beats polynomial ball growth. This proves a
moment bound even when a worst-boundary mixing row cannot be small.

\begin{theorem}[Uniform native local moments and joint quadratic-source pressure]
\label{f16v16:moments}
For every fixed bridge bound $r$, set $M_r=M(a_r,b_r)$ from
\eqref{f16v16:ball-bound}. For all $m\ge4$, $n\ge1$,
$\gamma\ge1+r^2$, all admitted bridge histories, and all arrays
$s\in[-t_0,t_0]^{\mathcal V}$,
\begin{equation}
                   \sup_c\mathbb E_{P_s}|\zeta_c|^2\le M_r.
 \label{f16v16:tilted-mean}
\end{equation}
For any two such arrays $s,t$,
\begin{equation}
 \left|\log\frac{Z(t)}{Z(s)}\right|
                           \le M_r\sum_c|t_c-s_c|.
 \label{f16v16:pressure-Lip}
\end{equation}
In particular, for arbitrary $0\le t_c\le t_0$ on any set of cells,
\begin{equation}
 \boxed{\quad
 \log\mathbb E_P e^{\sum_c t_c|\zeta_c|^2}
                              \le M_r\sum_c t_c.\quad}
 \label{f16v16:joint-mgf}
\end{equation}
Thus, for every fixed integer $k\ge1$,
\begin{equation}
 \sup_c\mathbb E_P|\zeta_c|^{2k}
       \le k!\,t_0^{-k}e^{t_0M_r},\qquad
 m_4(P)\le1+2t_0^{-2}e^{t_0M_r}<\infty.
 \label{f16v16:all-moments}
\end{equation}
These constants are independent of ambient volume, duration, and coupling in
the stated range. They concern the same $m_4(P)$ left open in V13--V15, not an
averaged moment or a moment of a reference law.
\end{theorem}
\begin{proof}
Apply the ball lemma to \eqref{f16v16:energy-boundary}, at each fixed $s$.
All constants in that inequality were uniform on the entire source cube,
which proves \eqref{f16v16:tilted-mean}. This step used only finite moments
from the compact regulator to start; the conclusion now makes them uniform.

At fixed regulator, compactness permits exact differentiation of the source
partition function along the straight segment $s+u(t-s)$, $0\le u\le1$:
\[
 \frac d{du}\log Z(s+u(t-s))
           =\sum_c(t_c-s_c)\mathbb E_{P_{s+u(t-s)}}e_c.
\]
The segment stays inside the same source cube. Bound the absolute value by
\eqref{f16v16:tilted-mean} and integrate. This proves
\eqref{f16v16:pressure-Lip} and, with $s=0$, \eqref{f16v16:joint-mgf}.
This is integration of controlled \emph{actual means}, not differentiation
of an unmarked asymptotic pressure error. Finally
$u^k\le k!t_0^{-k}e^{t_0u}$ for $u\ge0$ and
\eqref{f16v16:joint-mgf} on one cell prove
\eqref{f16v16:all-moments}. The definition of $m_4(P)$ in V13 is
$\sup_c\mathbb E_P(1+|\zeta_c|^4)$, so $k=2$ is the required input.
\end{proof}

This theorem does not give a native logarithmic Sobolev inequality, nor
covariance decay. It does not control moments uniformly under arbitrary
\emph{fixed} exterior configurations. The actual marginal, with its conditional
partition functions included, is essential. Its constants are conservative
expressions in fixed local bounds, not numerically optimized entry constants.

\subsection{Joint large-field tails, including a genuinely native cancellation family}

\begin{proposition}[Joint energy tails and exponentially rare low-field pairs]
\label{f16v16:tails}
For every cell set $S$, threshold $R\ge0$, and $u\ge0$,
\begin{align}
 P\{e_S\ge M_r|S|+u\}&\le e^{-t_0u},\notag\\
 P\{e_c\ge R^2\text{ for every }c\in S\}
             &\le\exp[-t_0(R^2-M_r)|S|].
 \label{f16v16:cell-tails}
\end{align}
Let $\mathcal B_{cd}$ be precisely V14's low-five-neighbour-field event on
one moving-port bond. There are fixed $c_1>0,C_r<\infty$ and a fixed-$r$
coupling threshold such that, for \emph{every} finite set $\mathcal A$ of
distinct neighbouring nonroot port pairs, at arbitrary times,
\begin{equation}
 P\left(\bigcap_{a\in\mathcal A}\mathcal B_a\right)
                     \le\exp[-(c_1\gamma-C_r)|\mathcal A|].
 \label{f16v16:low-field-joint}
\end{equation}
In particular $\mathbb E_P N_{\rm low}/\mathcal N
\le C\exp(C_r-c_1\gamma)$, and the same bound without the counting factor
holds for each specified pair. No independence of the events is assumed.
\end{proposition}
\begin{proof}
Apply exponential Markov with $t_c=t_0$ on $S$ in
\eqref{f16v16:joint-mgf}; both assertions in \eqref{f16v16:cell-tails}
follow immediately.

We retain the actual low-field definition, not all possible conditional-angle
defects. Each five-neighbour field is a sum of five transported unit
quaternions. If its norm is at most one, its distance from $5I$ is at least
four. The telescoping and geodesic comparison in V14 therefore gives, for
each event $a$, a fixed-size \emph{unmoving} cell set $S_a$ and $c_*>0$
such that
\begin{equation}
 \mathcal B_a\ \Longrightarrow\ e_{S_a}\ge c_*\gamma.
 \label{f16v16:low-field-energy}
\end{equation}
For completeness, $|I-r_p|$ is bounded by the sum of the chord distances of
$f_t^{-1},h_p,f_{t+1}$, and these are at most a fixed multiple of
$\gamma^{-1/2}(|a_p|+|x_{t,b}|+|x_{t+1,b}|)$. Moving internal edges have
the same bound in their own cube chords; bridge distances are at most
$r/\sqrt\gamma$. There are finitely many factors in the five transported
words. Cauchy--Schwarz on that fixed list, after absorbing $r/\sqrt\gamma$
for sufficiently large fixed-$r$ coupling, proves
\eqref{f16v16:low-field-energy}. This is the local version of V14's
deterministic implication, before summing the field count.

There are fixed integers $s_*,k_*$ with $|S_a|\le s_*$ and every cell in at
most $k_*$ such sets, counting all pair types and times. Put
$T=\bigcup_{a\in\mathcal A}S_a$. On the intersection event,
\[
 k_*e_T\ge\sum_{a\in\mathcal A}e_{S_a}\ge c_*\gamma|\mathcal A|,
 \qquad |T|\le s_*|\mathcal A|.
\]
Use \eqref{f16v16:joint-mgf} on $T$ and exponential Markov. The result is
\eqref{f16v16:low-field-joint} with $c_1=t_0c_*/k_*$ and
$C_r=t_0M_rs_*$. There are at most $C\mathcal N$ pair labels. Summing the
single-label probability proves the density statement. All coordinates,
including all time slices in any selected set, retain their actual native
law; a product exterior distribution was not inserted.
\end{proof}

A further precise geometric consequence is available. Give the low-field
labels any fixed finite-range adjacency of maximum degree $d_*\ge2$, and
put $p_\gamma=\exp(C_r-c_1\gamma)$. The probability that the bad connected
component of a specified label has size at least $k$ is at most
\begin{equation}
                    d_*^{2(k-1)}p_\gamma^k.
 \label{f16v16:cluster-tail}
\end{equation}
Indeed a connected set of size at least $k$ contains a connected subset of
exactly $k$ containing that label. Such subsets number at most
$d_*^{2(k-1)}$: choose a deterministic spanning tree and its depth-first walk
of length $2(k-1)$, which records the visited set injectively. Sum
\eqref{f16v16:low-field-joint}. This is a bad-\emph{field} cluster bound,
not an identification of all spectral defects with such clusters.

\subsection{Source-weighted tails and improved unrestricted-volume trace estimates}

Let $\mathbf e^{\rm prim}=(e^{\rm br},e^{\rm pt})$ be V7's primitive bridge
and compact Haar-port error banks, and let $\mathbf r$ be its full corrected
common-source residual vector. The exact synthesis is still
\[
 \mathbf r=e^{\rm br}+\mathsf P^{\mathsf T}e^{\rm pt}.
\]
The port lift has a volume-independent operator norm and exponentially
summable absolute columns. Each primitive component $e_j^{\rm prim}$ obeys
the inherited complete-chart bound
\[
 |e_j^{\rm prim}|\le C_r\gamma^{-1/2}
                       (1+e_{T_j}),\qquad |T_j|\le C,
\]
with bounded carrier incidence. These are statements about the same compact
Haar derivatives, not flat substitutes at finite coupling.

\begin{theorem}[Actual source moments, exceptional weights, and all-lag trace density]
\label{f16v16:source-energy}
For every fixed $p\ge1$, every primitive component and every canonical common
mark $\alpha$,
\begin{equation}
 \|e_j^{\rm prim}\|_{L^p(P)}+
 \|r_\alpha\|_{L^p(P)}\le C_{p,r}\gamma^{-1/2}.
 \label{f16v16:source-Lp}
\end{equation}
For any event $E$ in the full native probability space,
\begin{equation}
 \mathbb E_P[1_E|r_\alpha|^2]
       +\mathbb E_P[1_E|e_j^{\rm prim}|^2]
                          \le C_r\gamma^{-1}P(E)^{1/2}.
 \label{f16v16:marked-event}
\end{equation}
In particular \eqref{f16v16:low-field-joint} pays these fixed source weights
on a joint low-field event with an additional factor
$\exp[-(c_1\gamma-C_r)|\mathcal A|/2]$.
The full positive Gram $\Gamma_\gamma=\operatorname{Cov}_P(\mathbf r)$ and
V7's actual nonadjacent-time Hessian satisfy
\begin{equation}
 \boxed{\quad
 \operatorname{Tr}\Gamma_\gamma\le C_r\mathcal N/\gamma,
 \qquad \frac{\|\mathsf M_\gamma\|_{S_1}}{\mathcal N}\le C_r/\gamma.
 \quad}
 \label{f16v16:trace}
\end{equation}
Also every canonical entry has $|\Gamma_{\gamma,\alpha\beta}|\le C_r/\gamma$,
independently of both mark positions. Thus the same entry bound holds for the
full-history common-source mixed Hessian at nonadjacent times. It has no
decay in the separation. For V4's original core radius in its derivative
regime,
\begin{equation}
 \frac{\|\mathsf C_{\gamma,R}\|_{S_1}}{\mathcal N}
                         \le C_r/\gamma+C R^2/\gamma.
 \label{f16v16:core-trace}
\end{equation}
All these statements permit unrestricted $B,n$ in their fixed-$r$ coupling
regime. There is no $e^{KB}/\sqrt\gamma$ entry condition.
\end{theorem}
\begin{proof}
The uniform cell moments and the bounded size of $T_j$ give
$\|1+e_{T_j}\|_p\le C_{p,r}$. Apply the primitive estimate. For a single
full residual, Minkowski's inequality and the uniform absolute column sum of
$\mathsf P$ give the same $L^p$ bound. The latter summability follows from
the finite-range positive port matrix by the inverse decay already used in
V10--V12, and is unchanged at either dressed time endpoint. Cauchy--Schwarz
with $p=4$ proves \eqref{f16v16:marked-event}; this is a bound for the stated
observables, not arbitrary normalized $L^2$ functions.

There are $O(\mathcal N)$ primitive components and $9\mathcal N$ canonical
common sources. Sum their second moments, or use the bounded whole-vector
synthesis, to obtain $\mathbb E\|\mathbf r\|^2\le C_r\mathcal N/\gamma$.
Centering reduces this quantity and proves the trace bound. Individual
covariance entries are bounded by their two standard deviations. V7's
\emph{exact} contact identity is
$\mathsf M_\gamma=-\operatorname{Off}_{>1}\Gamma_\gamma$. Its Fourier
extraction bound is $\|\operatorname{Off}_{>1}A\|_{S_1}\le4\|A\|_{S_1}$,
independent of duration. Positivity of the full Gram gives
\eqref{f16v16:trace}. V4's original core operator-norm estimate, multiplied
by its $9\mathcal N$ source dimension, and the unchanged full/core identity
give \eqref{f16v16:core-trace}, exactly as in V7 but with the new energy rate.
No new event replaces the core, and no endpoint normalizer is deleted.
\end{proof}

These are sharper consequences than the $\epsilon_\gamma$ energy floor used
in V7 and V12, and than the time-length-dependent pointwise bounds at zero
history. They still do not identify the leading coefficient $G_2$ at arbitrary
volume. Small second moments of the residuals yield an $O(\gamma^{-1})$
scale; a uniform remainder $o(\gamma^{-1})$ and connected covariance decay
require additional arguments. The source-weighted exceptional estimate in
\eqref{f16v16:marked-event} likewise does not apply to every spectral
innovation without an appropriate norm or capacity estimate.

\subsection{The posterior input improves, while its coherent gap remains explicit}

Use V12's separated-cylinder family, without repeated sources, common
exterior projection $E_O$, selection $\chi$, depth $L$, and enlarged-patch
overlap $q_0$. All exterior values are integrated. The established exact
identities are
\begin{equation}
 \mathbf r^D=\chi\mathbf r-\Delta^{\mathsf T}e^{\rm pt},\qquad
 \mathbf g^D-\chi\mathbf g=-\Delta^{\mathsf T}S_0,\qquad
 \mathbf j=\mathbf g^D+E_O\mathbf r^D,\qquad
 \|\Delta\|\le C\sqrt{q_0}e^{-\nu L}.
 \label{f16v16:posterior-identities}
\end{equation}
Here $S_0=\partial_aU_0$ is the affine Gaussian port score \emph{evaluated
under the native law}. The conditional partition functions remain in the
actual exterior posterior. These are V12's harmonic lift identities, not a
new or extra covariance subtraction.

\begin{proposition}[Sharp-scale posterior trace and the remaining gap-only reduction]
\label{f16v16:posterior}
The actual posterior return and separated-cylinder gluing error obey
\begin{align}
 \frac{\operatorname{Tr}\mathsf R_{\rm post}}{\mathcal N}
       &\le C_rq_0(\gamma^{-1}+e^{-2\nu L}),\notag\\
 \frac{\|\chi\mathsf M_\gamma\chi-\mathsf L_D\|_{S_1}}{\mathcal N}
       &\le C_rq_0(\gamma^{-1}+e^{-2\nu L}).
 \label{f16v16:posterior-trace}
\end{align}
For any exterior events $E_s$, the exceptional conditional Gram integrals also
satisfy
\begin{equation}
 \sum_s\mathbb E_{\bar P_s}
          [1_{E_s}\operatorname{Tr}\Gamma_\gamma^{D_s}]
                                \le C_rq_0\mathcal N/\gamma.
 \label{f16v16:exceptional-Gram}
\end{equation}
Let $\widehat\lambda_P>0$ be the \emph{actual} balanced auxiliary gap of V14.
The following unnormalized operator bounds now have no unproved moment input:
\begin{align}
 \|\Gamma_\gamma\|_{\rm op}
                       &\le \frac{C_r}{\widehat\lambda_P\gamma},\notag\\
 \|\mathsf R_{\rm post}\|_{\rm op}
    +\|\chi\mathsf M_\gamma\chi-\mathsf L_D\|_{\rm op}
                       &\le\frac{C_rq_0}{\widehat\lambda_P}
                                        (\gamma^{-1}+e^{-2\nu L}).
 \label{f16v16:gap-only}
\end{align}
A uniform positive lower bound on $\widehat\lambda_P$ remains unproved;
\eqref{f16v16:gap-only} does not assert one.
\end{proposition}
\begin{proof}
The primitive and full residual energies just proved are
$O_r(\mathcal N/\gamma)$. Apply the bounded synthesis and lift-mismatch
estimate in \eqref{f16v16:posterior-identities}; this gives
$\mathbb E\|\mathbf r^D\|^2\le C_rq_0\mathcal N/\gamma$.
Because $S_0$ is a fixed finite-range affine map of the unmoving coordinates,
with bounded coefficients and fixed-$r$ constants per cell, the local
second moments also give $\mathbb E_P\|S_0\|^2\le C_r\mathcal N$.
Therefore
\[
 \mathbb E_P\|\mathbf j-\chi\mathbf g\|^2
 \le2\mathbb E\|\Delta^{\mathsf T}S_0\|^2
                +2\mathbb E\|E_O\mathbf r^D\|^2
 \le C_rq_0\mathcal N(\gamma^{-1}+e^{-2\nu L}).
\]
Use conditional $L^2$ contraction and centering to get the first trace bound.
V12's same-posterior identity
$\chi\mathsf M_\gamma\chi-\mathsf L_D=
-\operatorname{Off}_{>1}\mathsf R_{\rm post}$ and the factor-four trace
extraction give the second. For \eqref{f16v16:exceptional-Gram}, conditional
centering and the tower property give an upper bound by
$\sum_s\mathbb E\|\mathbf r^{D_s}\|^2$, which is bounded by the same
combined lift synthesis. The nonnegative gates can be dropped only from that
upper bound. This improves the rate of V12's exceptional trace budget, but
not to $o(\gamma^{-1})$.

Finally V14's directional innovation-energy inequalities hold for these
same balanced updates and for the same native moment definitions. Insert
\eqref{f16v16:all-moments}, in place of a previously unverified $m_2,m_4$,
to obtain
\[
 \widehat{\mathcal E}_P(v^{\mathsf T}\mathbf r)
      \le C_r\gamma^{-1}\|v\|^2,\qquad
 \widehat{\mathcal E}_P(b^{\mathsf T}S_0)\le C_r\|b\|^2.
\]
Their actual finite-regulator Poincar\'e inequality and V13--V14's exact
posterior reduction give \eqref{f16v16:gap-only}. The current theorem proves
the moment input, not the gap input. Replacing $\widehat\lambda_P$ by the
Gaussian margin $2/39$ would still be unjustified. Its crude finite-regulator
bound alone does not yield an unrestricted-volume contraction.
\end{proof}

\subsection{What changes in the program, and what is not inferred}

The separate full-native local moment obligation in V13--V15 is now paid,
uniformly over bounded inhomogeneous bridge histories, all durations, and all
spatial tori. It is paid without using V6's entropy-density theorem, V9's
product comparison, V15's repair neighbourhood, or an assumed native
functional inequality. Their respective conclusions remain useful but were
not hidden inputs to this moment proof. The argument uses the exact pinned
fast-coordinate specification; it is not a moment theorem for an unrestricted
physical gauge connection or for the integrated bridge law.

The low-field result is also stronger than a mean density: it controls any
prescribed joint family, and its connected components in a bounded-degree
label graph have an exponential size tail at sufficiently large coupling.
It is a concrete input for an eventual source-marked cluster or repair
argument. It does \emph{not} establish that V14's entire angle defect is
supported on those low-field events. Chord--port and other conditional
landscapes remain in that defect. Even a complete bad-event probability
estimate would not permit multiplying it by the squared innovation of an
arbitrary normalized test function independently of that function.

There is no contradiction with V14's compact cancellation. At its selected
fixed exterior, a port marginal can be Haar, with principal rescaled energy
of order $\gamma$. The conditional bound
\eqref{f16v16:conditional-mgf} then has a large exterior-energy debit. The
uniform conclusion is under the \emph{averaged full native law}. Every such
exterior neighbourhood can retain positive probability and still defeat a
worst-exterior correlation criterion. Neither V13's nor V14's impossibility
result is reversed.

The remaining active global task is therefore narrower: control the native
auxiliary gap, or directly the coherent source covariance, using compatible
repairs with a source-resolved treatment of the uncovered sectors. Uniform
local moments no longer need to be assumed in that reduction. Another
trace-density estimate would not exclude a coherent exceptional direction.
The all-lag matrices above have no new time-row decay at unrestricted volume,
and their $O(\gamma^{-1})$ bounds do not identify a leading coefficient with
an $o(\gamma^{-1})$ error. The existing infrared contact, gauge-invariant
bridge coverage, same-top physical transfer, calibrated physical time, and
interacting continuum obligations remain distinct.

\subsection{Diagnostics, preservation, and claim scope}

The accompanying diagnostic is
\begin{center}\small\path{verify_ym_f16_v16_local_moments.py}.\end{center}
It checks the ball-iteration algebra and its exact generating function,
finite-graph growth including saturated balls, conditional source and
normalizer identities on positive finite laws, and the conversion of local
energy witnesses into joint-event bounds with overlap. It also checks the
source/lift and posterior quadratic-form bookkeeping used by the corollaries.
All 1,077 exact assertions passed. The report separates those tests from
analytic hypotheses and inequalities.
The proof of the full-native moment theorem is the argument above, not an
inference from those finite examples.

The cumulative V15 is preserved byte for byte except for its title version
and insertion of this exact appendix before the final document-ending
command. The package supplies the exact recovery checker, the protected input
hashes, build and render reports, the new diagnostic, and actual reruns of the
available V4--V15 diagnostics, which passed 9,467, 89,339, 906, 311, 609,
990, 608, 463, 1,201, 222, 622, and 434 assertions, respectively.
No inaccessible archive program, external
formalization, or numerical enclosure of every native derivative bound is
represented as verified. The next companion checkpoint remains V20.

\begin{firewall}
V16 proves local exponential moment bounds for the complete native fast law
at a fixed bounded bridge history. These bounds are uniform in volume and
duration. They supply the missing fourth moment and joint low-field tails,
and improve residual and posterior trace densities to order $1/\gamma$.
The operator reduction still requires a uniform native auxiliary gap.

No such gap, exhaustive defect coverage, or spatially uniform time-row
mixing is established. Nor is a leading coefficient identified with an
error smaller than $1/\gamma$ at unrestricted volume. Physical bridge
coverage, the same-top physical transfer comparison, and an interacting
continuum Yang--Mills mass gap remain unproved.
\end{firewall}
% END F16 V16 APPEND

% BEGIN F16 V17 APPEND -- 2026-09-17
% Insert immediately before the final document-ending command in V16.
% No predecessor proof, claim, or bibliography string is silently revised.

\clearpage
\section{V17: local Stein recursion and coherent estimates on growing source banks}
\label{f16v17:context}

V16 supplies uniform moments of every fixed degree for the complete native
fast law. This changes what can be done without a native resampling gap.
Instead of further enlarging a repair neighbourhood, we apply a polynomial
Gaussian divergence identity directly under the native law. Local cutoffs
make integration by parts legitimate on the principal charts; V16 pays their
errors one coordinate group at a time. An all-chart, all-volume weak
expansion follows. Its coefficients are exactly V10's, not coefficients of a
new probability measure.

The new conclusion has a different norm from V10's. At every fixed order its
entrywise covariance remainder has a constant independent of both spatial
volume and history length. It has no separation decay. Nevertheless, combined
with the inherited locality of the coefficients, it gives operator-norm
control on every source bank of any prescribed polynomial size in coupling,
inside an otherwise unrestricted ambient history. No native auxiliary gap is
assumed. A completely delocalized coherent direction remains outside this
conclusion; we quantify that limitation instead of identifying the entrywise
norm with a full operator norm.

\subsection{Dependency audit and the change of method}

The supplied V16 appendix and analytic audit were read in full. Two targeted
semantic searches in the f14/f15/monograph sources returned no results. The
mounted inventories and exact text were then examined, not treated as empty.
F14 V98, \texttt{prop:v98-polynomial-Stein}, already gives a polynomial Haar
Stein error for its adapted curvature fields. F15 V64,
\texttt{f15v64:polynomial-Stein}, already gives location-uniform linear and
cubic Stein tests for its stationary, randomly framed curvature field. These
are explicitly credited. They do not give the present finite-order expansion
of the fully pinned augmented fast history: their laws, root-frame costs,
and harmonic-mode issues differ. This is a targeted non-duplication audit,
not recertification of the archives.

The relevant classical mechanism is approximation through a Stein identity,
not transferring a spectral gap between nearby measures. Barbour,
\href{https://doi.org/10.1007/BF01197887}{\emph{Stein's method for diffusion
approximations}} (1990), provides general context; only its primary summary
and bibliographic record were checked. Braverman--Dai--Fang,
\href{https://arxiv.org/abs/2012.02824}{\emph{High-order steady-state diffusion
approximations}}, derive recursive approximations through Stein/Poisson
equations. Their abstract and introductory stationary-identity setup were
reviewed, not their complete model-specific proofs. No theorem about their
Markov chains is transplanted here. The finite polynomial right inverse and
the native error estimates used below are proved directly.

The two import memoranda distinguish order gain from locality cost, and
compatible local constraints from an assumed physical gap. We keep those
distinctions. The present proof reuses V10's actual measured action and Haar
source coefficients, V11's inverse bounds, and V16's native moments. It does
not use V6's likelihood comparison, V9's product-entry condition, or V15's
repair gate. It does not reopen the failed worst-exterior pair-row test.
The next scheduled companion checkpoint remains V20.

\subsection{The exact measure and a polynomial norm with no hidden dimension}

Put $\varepsilon=\gamma^{-1/2}$ and fix an arbitrary admitted bridge history
$\max_{i,e}|y_{i,e}|\le r$. Let
\[
 P_\varepsilon=P_{\gamma,\mathbf y},\qquad
 Q=N(m,C),\qquad C=H^{-1},\qquad \xi=\zeta-m.
\]
These are the same complete unmoving native and Gaussian laws as V10 and V16.
The Gaussian mean and precision are independent of $\varepsilon$ at fixed
rescaled $\mathbf y$ and fixed geometry. There are fixed spectral bounds
$h_-I\preceq H\preceq h_+I$, bounded means per cell, and a fixed interaction
range for $H$. Its inverse has exponentially summable coordinate rows. In
particular, with scalar real indices $i,j$ and fixed coordinate multiplicity,
\begin{equation}
 \sup_i\sum_j|H_{ij}|\le K_H,\qquad
 \sup_i\sum_j|C_{ij}|\le K_C,
 \qquad \sup_i|m_i|\le M_r^{\rm G}.
 \label{f16v17:Gaussian-input}
\end{equation}
The constants do not depend on $B,n$. For clarity, the inverse row bound is
not inferred merely from $\|C\|_{2\to2}$. The finite-range Neumann expansion
of $H^{-1}$ gives exponential entry decay; polynomial growth of the cell
graph and its at most 36 coordinates per cell make those rows summable.
This is the inherited inverse-locality argument on the complete history,
including its actual dressed endpoints.

For a scalar polynomial $f(\xi)$ of degree at most $d$, collect its ordinary
monomials and set
\begin{equation}
 \|f\|_{\rm pol}=\sum_{\mathbf k}|f_{\mathbf k}|,
 \qquad f(\xi)=\sum_{|\mathbf k|\le d}f_{\mathbf k}\xi^{\mathbf k}.
 \label{f16v17:pol-norm}
\end{equation}
For a polynomial vector $g=(g_j)$ put
$\|g\|_{{\rm pol},1}=\sum_j\|g_j\|_{\rm pol}$.
The finite number of variables can be arbitrarily large. Degree, not that
number, is kept fixed. Repeated indices are included. The norm is
submultiplicative, and
\[
 \sum_j\|\partial_j f\|_{\rm pol}\le d\|f\|_{\rm pol}.
\]
A sum of many unrelated observables pays its actual coefficient sum; no
claim that this norm is dimension free for every such sum is made.

\begin{proposition}[Native polynomial moments and coordinate-local cut tails]
\label{f16v17:poly-moments}
For every fixed $d$ and $1\le p<\infty$,
\begin{equation}
 \|f\|_{L^p(P_\varepsilon)}+\|f\|_{L^p(Q)}
       \le C_{d,p,r}\|f\|_{\rm pol}.
 \label{f16v17:poly-Lp}
\end{equation}
If $g$ is one original three-coordinate group, and
$A_g=\{|\zeta_g|\ge(2\varepsilon)^{-1}\}$, then, for some fixed $c>0$,
\begin{equation}
 \mathbb E_{P_\varepsilon}[1_{A_g}|f|]
          \le C_{d,r}e^{-c/\varepsilon^2}\|f\|_{\rm pol}.
 \label{f16v17:cut-tail}
\end{equation}
Both estimates are uniform in the locations and number of variables appearing
in $f$ when its displayed coefficient norm is used. They hold for
$\gamma\ge1+r^2$.
\end{proposition}
\begin{proof}
V16 bounds every fixed moment of each $|\zeta_c|$, uniformly in its cell.
Subtracting the bounded Gaussian mean preserves these bounds for each
$\xi_i$. H\"older applied to a monomial of degree at most $d$, allowing
repeated coordinates, bounds its $L^p$ norm by a fixed coordinate moment of
order at most $dp$. Minkowski then gives the native part of
\eqref{f16v17:poly-Lp}. The Gaussian part follows from bounded coordinate
variances and the same argument. No independence of the native coordinates
is used.

The cell containing $g$ has energy at least $1/(4\varepsilon^2)$ on $A_g$.
V16's one-cell exponential moment gives
$P_\varepsilon(A_g)\le\exp(t_0M_r-t_0/(4\varepsilon^2))$.
Cauchy--Schwarz and the $L^2$ polynomial bound prove
\eqref{f16v17:cut-tail}, for example with $c=t_0/8$.
This estimates one group tail weighted by the actual polynomial. It is not
a probability estimate for a simultaneous all-good event over the history.
\end{proof}

\subsection{A finite Gaussian divergence solver, without a native inverse}

Define the Gaussian divergence on polynomial vector fields by
\[
 \delta_Q g=\sum_j\bigl((H\xi)_j g_j-\partial_jg_j\bigr).
\]
A useful right inverse can be constructed by finite algebra. For a monomial
$M=\xi_{i_1}\cdots\xi_{i_d}$ choose its first index by one fixed ordering,
put $h=\xi_{i_2}\cdots\xi_{i_d}$, and set
\begin{equation}
 \mathcal B M=\sum_{\ell=2}^d C_{i_1i_\ell}
                \prod_{k\ne1,\ell}\xi_{i_k},\qquad
 \mathcal S M=(C_{j i_1}h)_j+\mathcal S(\mathcal B M),
 \qquad \mathcal S1=0.
 \label{f16v17:solver-recursion}
\end{equation}
Empty sums are zero; extend linearly to collected polynomials. Recursion
reduces the degree by two and terminates. The letter $\mathcal B$ in this
subsection is this degree-lowering map, not a spatial cube set.

\begin{proposition}[A coefficient-summable Gaussian right inverse]
\label{f16v17:solver}
For every polynomial $f$ of degree at most $d$,
\begin{equation}
 \delta_Q\mathcal S f=f-\mathbb E_Qf,
 \qquad \deg\mathcal S f\le d-1,
 \qquad \|\mathcal S f\|_{{\rm pol},1}\le K_d\|f\|_{\rm pol}.
 \label{f16v17:solver-bound}
\end{equation}
For constant $f$ the vector is zero. The constants $K_d$ depend on $d,K_C$,
not on $B,n$. Thus this construction requires no inverse of a native
resampling generator or native differential operator.
\end{proposition}
\begin{proof}
Since $HC=I$, direct polynomial differentiation gives
\[
 \delta_Q(C_{\cdot i_1}h)=M-\mathcal B M.
\]
Ordinary Gaussian integration by parts gives
$\mathbb E_Q M=\mathbb E_Q\mathcal B M$. Applying the degree induction to
the second term of \eqref{f16v17:solver-recursion} therefore proves the exact
divergence identity. Equivalently this is Wick's finite recursion, not a
convergence assertion for an infinite series.

For one monomial the coefficient norm of the first vector is at most $K_C$.
The norm of $\mathcal B M$ is at most $(d-1)K_C$, and its degree is $d-2$.
Thus constants satisfying $k_0=0$, $k_1=K_C$ and
$k_d=K_C+(d-1)K_C k_{d-2}$ bound the construction on degree-$d$ monomials.
Taking $K_d=\max_{0\le l\le d}k_l$ and summing absolute coefficients proves
the bound. Coordinate ordering can change the right inverse, but not its
divergence. Later normalized coefficients will be independent of that choice.
\end{proof}

Only one covariance row is summed for each algebraic step. Replacing that
sum by a sum of unrelated coordinate error bounds would lose the point of
\eqref{f16v17:solver-bound}. No semigroup continuity on an infinite-dimensional
function space is needed for this finite polynomial construction.

\subsection{A native Stein identity with the principal-chart boundary paid}

On the positive chart interior retain V10's measured potential
\[
 V_\varepsilon=U_\varepsilon-\log J^f_\varepsilon,
 \qquad V_\varepsilon\sim U_0+\sum_{q\ge1}\varepsilon^qV_q.
\]
The $V_q$ here are \emph{exactly} \eqref{f16v10:jets}, including the
endpoint Haar powers. Each $\partial_jV_q$ has a fixed local carrier and
\begin{equation}
 \sup_j\|\partial_jV_q(m+\xi)\|_{\rm pol}\le L_{q,r},
 \qquad \deg(\partial_jV_q)\le q+1.
 \label{f16v17:local-drift-coefficients}
\end{equation}
Indeed differentiating an extensive local action retains only its finitely
many terms incident to $j$. Translation by the bounded $m$ preserves the
fixed-degree coefficient bound. No norm of the entire extensive $V_q$ is
used below.

\begin{proposition}[Full-native polynomial Stein expansion]
\label{f16v17:Stein}
For any fixed integer $N\ge0$ and any polynomial vector $g$ of degree at
most $d$, there are constants $C_{d,N,r},\gamma_{d,N,r}$ independent of
volume, duration, and the vector's support, such that
\begin{equation}
 \left|\mathbb E_{P_\varepsilon}\delta_Qg+
   \sum_{q=1}^N\varepsilon^q
           \mathbb E_{P_\varepsilon}(\nabla V_q\cdot g)\right|
 \le C_{d,N,r}\varepsilon^{N+1}\|g\|_{{\rm pol},1}
 \label{f16v17:Stein-estimate}
\end{equation}
for $\gamma\ge\gamma_{d,N,r}$. The sum is empty when $N=0$.
Every expectation uses the complete native law. No all-coordinate cutoff is
imposed on that law, and no derivative of an unmarked pressure error is taken.
\end{proposition}
\begin{proof}
We give the cutoff argument because flat integration by parts through the
logarithmic cut without its Haar term would be incorrect. Let $g(j)$ denote
the original group containing scalar coordinate $j$. Choose one smooth radial
function on $\mathbb R^3$, equal to one on radius $1/2$ and zero outside
radius one, and write
\[
 \chi_j(\zeta)=\chi(\varepsilon\zeta_{g(j)}),
 \quad 0\le\chi_j\le1,\qquad |\partial_j\chi_j|\le C\varepsilon.
\]
It cuts off only this coordinate group. For fixed values of all other groups,
$\chi_jg_j$ has support strictly inside its principal ball of radius
$\pi/\varepsilon$. The native Lebesgue density on that ball is smooth and
positive. Fubini and compact finite-regulator integrability therefore give
\[
 \mathbb E[\chi_j\partial_jg_j+(\partial_j\chi_j)g_j]
                    =\mathbb E[\chi_jg_j\partial_j V_\varepsilon].
\]
Other groups need not be in their small charts. Consequently the following
identity, before estimating anything, is exact:
\begin{equation}
 \begin{split}
 \mathbb E\delta_Q g={}&
 -\sum_j\mathbb E[\chi_j g_j
            (\partial_jV_\varepsilon-(H\xi)_j)]
 +\sum_j\mathbb E[(\partial_j\chi_j)g_j]\\
 &+\sum_j\mathbb E[(1-\chi_j)
                 ((H\xi)_jg_j-\partial_jg_j)].
 \end{split}
 \label{f16v17:exact-cut-identity}
\end{equation}
Here and for the remainder of this proof expectations are under
$P_\varepsilon$. This is a chart integration identity with its measured
Jacobian. It is not an assertion that a compact Haar derivative equals a
flat derivative. The exact Haar-defined source observables remain unchanged.

For the classical drift, every local word is
$\varepsilon^{-2}f(\varepsilon w)$, with $f(0)=Df(0)=0$ and uniformly
bounded real derivatives of every fixed order, by unitary exponential
insertion. Its coordinate derivative has a Taylor remainder of order
$\varepsilon^{N+1}$ bounded by a fixed local polynomial of degree at most
$N+2$. This bound holds for \emph{all} real inputs appearing along the
Taylor segment. Fixed incidence and the bounded bridge inputs therefore give,
for every fixed $p$,
\[
 \left\|\partial_jU_\varepsilon-\partial_jU_0
                 -\sum_{q=1}^N\varepsilon^q\partial_jU_q\right\|_{L^p(P_\varepsilon)}
                         \le C_{N,p,r}\varepsilon^{N+1}.
\]
The proof uses Proposition~\ref{f16v17:poly-moments}, not moments at an
intermediate changed coupling.

The Haar logarithm need not have bounded derivatives at its cut. On the
support of $\chi_j$, however, its only differentiated group obeys
$|\varepsilon\zeta_{g(j)}|\le1$. The smooth function
$-w\log j(\varepsilon\zeta_g)$, with the original $w\in\{1/2,1\}$, has
Taylor remainder bounded there by
$C_N\varepsilon^{N+1}(1+|\zeta_g|^{N+2})$ after one coordinate derivative.
It includes the order-two term in \eqref{f16v10:haar-order-two} and no
additional copy of it. Other Haar factors have zero $j$ derivative. Thus
\[
 \left\|\chi_j\left[
 \partial_jV_\varepsilon-(H\xi)_j
       -\sum_{q=1}^N\varepsilon^q\partial_jV_q\right]\right\|_{L^p(P_\varepsilon)}
                  \le C_{N,p,r}\varepsilon^{N+1}.
\]
For $N=0$ the classical difference begins at order one and the Haar
contribution at order two, so the same estimate applies.

Pair this drift error with $g_j$ by Cauchy--Schwarz. Removing $\chi_j$ from
the finite polynomial terms costs only an expectation on $A_{g(j)}$.
Every other term of \eqref{f16v17:exact-cut-identity} is likewise supported
on that event, including $\partial_j\chi_j$. The row bound for $H$, the
polynomial product rule, and \eqref{f16v17:cut-tail} bound the sum of these
errors by
\[
 C_{d,N,r}\bigl(\varepsilon^{N+1}+e^{-c/\varepsilon^2}\bigr)
                                    \sum_j\|g_j\|_{\rm pol}.
\]
At fixed order the flat term is absorbed into $\varepsilon^{N+1}$ by
increasing the constant and, if needed, the fixed-order coupling threshold.
There is no union bound over all coordinate groups and no multiplication by
$B(n+1)$. This proves \eqref{f16v17:Stein-estimate}.
\end{proof}

\subsection{A normalized finite-order expansion without a spatial entry condition}

For a polynomial $f$ define finite algebraic operators
\begin{equation}
 \mathcal T_q f=-\sum_j(\partial_jV_q)(\mathcal Sf)_j,
 \quad q\ge1;
 \qquad b_0(f)=\mathbb E_Q f,
 \quad b_k(f)=\sum_{q=1}^k b_{k-q}(\mathcal T_q f).
 \label{f16v17:coefficient-recursion}
\end{equation}
All expressions are evaluated in $\xi$, with $\zeta=m+\xi$ in $V_q$.
They use one fixed Gaussian covariance and fixed measured coefficients.
The maps obey
\begin{equation}
 \deg\mathcal T_q f\le d+q,\qquad
 \|\mathcal T_q f\|_{\rm pol}\le L_{q,r}K_d\|f\|_{\rm pol}.
 \label{f16v17:T-bound}
\end{equation}
This follows by multiplication and the \emph{sum} in
\eqref{f16v17:solver-bound}, not a sum of $\mathcal N$ local constants.

\begin{theorem}[All fixed orders of native polynomial expectations]
\label{f16v17:expectations}
For any fixed degree $d$ and order $N\ge0$, there are finite
$C_{d,N,r},\gamma_{d,N,r}$ such that for every finite spatial torus,
every duration, every admitted bridge history, and every degree-$d$
polynomial independent of $\varepsilon$ in the centered coordinates,
\begin{equation}
 \boxed{
 \left|\mathbb E_{P_\varepsilon}f-
                  \sum_{k=0}^N\varepsilon^k b_k(f)\right|
       \le C_{d,N,r}\varepsilon^{N+1}\|f\|_{\rm pol}.}
 \label{f16v17:expectation-expansion}
\end{equation}
The condition is only $\gamma\ge\gamma_{d,N,r}$; it contains no spatial
volume or duration. Also $|b_k(f)|\le C_{d,k,r}\|f\|_{\rm pol}$.
These are coefficients of the \emph{normalized native law}; action-only
constants and all normalization corrections are treated correctly.
\end{theorem}
\begin{proof}
Apply Proposition~\ref{f16v17:Stein} to $g=\mathcal S f$ and use its exact
divergence identity. For every fixed $N$ this gives
\begin{equation}
 \mathbb E_{P_\varepsilon}f
  =\mathbb E_Q f+\sum_{q=1}^N\varepsilon^q
                  \mathbb E_{P_\varepsilon}\mathcal T_qf
       +O_{d,N,r}(\varepsilon^{N+1}\|f\|_{\rm pol}).
 \label{f16v17:native-recursion}
\end{equation}
For $N=0$ this is already the theorem. Induct on $N$, applying the
order-$N-q$ statement to $\mathcal T_q f$. Its degree increases by at most
$q$, and its coefficient norm is bounded by \eqref{f16v17:T-bound}.
There are only $N$ terms. Collecting each total power gives exactly
\eqref{f16v17:coefficient-recursion}, and all errors are of order
$\varepsilon^{N+1}$. The coefficient bound follows by the same induction
from the Gaussian moment bound. Only finitely many degrees and thresholds
occur at each fixed $(d,N)$, so their maximum is independent of $B,n$.

Here is an algebraic check of coefficient ownership which does not assume
that an exponential of a cubic polynomial is integrable. At each fixed
finite regulator form the \emph{formal} normalized functional
\[
 \mathcal M_z(f)=
 \frac{\mathbb E_Q[f\exp(-\sum_{q\ge1}z^qV_q)]}
      {\mathbb E_Q[\exp(-\sum_{q\ge1}z^qV_q)]}.
\]
Every coefficient involves only finitely many Gaussian polynomial moments.
Gaussian integration by parts applied coefficient by coefficient gives
\[
 \mathcal M_z(\delta_Q g)
       =-\sum_{q\ge1}z^q\mathcal M_z(\nabla V_q\cdot g).
\]
With $g=\mathcal S f$, its coefficients satisfy the same triangular
recursion as $b_k$. The leading value is $\mathbb E_Q f$, so the coefficients
are uniquely those of \eqref{f16v17:coefficient-recursion}, independently of
the ordering used to choose $\mathcal S$. The already proved error estimate
identifies them with the actual native asymptotic coefficients. Thus the
formal calculation identifies ownership; it does not supply the error bound.
Fast-independent additions to $V_q$ disappear from both its gradient and this
normalized functional, as required.
\end{proof}

For example these normalized coefficients obey
\begin{equation}
 b_1(f)=-\operatorname{Cov}_Q(f,V_1),\qquad
 b_2(f)=-\operatorname{Cov}_Q(f,V_2)
                        +\tfrac12\kappa_Q(f,V_1,V_1).
 \label{f16v17:first-expectation-coefficients}
\end{equation}
The recursion uses local \emph{derivatives} of the extensive action, whereas
the equivalent cumulants keep its scalar normalization. Constants $f$ have
no corrections. A native sampler gap was not used to solve the Stein
equation: the right inverse is the finite Gaussian algebra above.

\subsection{Apply the expansion to the same exact Haar residuals}

Keep the literal residuals $r_{\varepsilon,\alpha}$ of V7 and their
coefficients $r_{\alpha,k}$ of V10. In particular the order-one coefficient
still includes
\[
 \frac12\sum_p[\eta_{\alpha,p},a_p]\cdot\partial_{a_p}U_0,
 \qquad \eta_\alpha=\mathsf P\mathsf e_\alpha.
\]
The new flat-chart integration identity does not remove this noncommutative
Haar-insertion term or differentiate a different source. All retained bridge
derivatives are formed in the original coordinates first.

\begin{proposition}[Uniform native source jets]
\label{f16v17:source-jets}
For every fixed $K\ge1$ and $1\le p<\infty$,
\begin{equation}
 \left\|r_{\varepsilon,\alpha}-
                 \sum_{k=1}^K\varepsilon^k r_{\alpha,k}\right\|_{L^p(P_\varepsilon)}
        \le C_{K,p,r}\varepsilon^{K+1},
 \qquad \|r_{\alpha,k}(m+\xi)\|_{\rm pol}\le C_{k,r}.
 \label{f16v17:source-jets-bound}
\end{equation}
These bounds are uniform in the source position, ambient volume, and duration.
The same statements hold for each primitive port or bridge remainder.
\end{proposition}
\begin{proof}
Insert the scaled Lie-algebra generator directly into each ordered compact
word, as in V7 and V10. Its fixed-order real Taylor remainder is bounded by
$C_{K,r}\varepsilon^{K+1}(1+|\zeta_T|^{K+2})$ on its fixed local carrier.
Uniform native moments from V16 bound its $L^p$ norm at the actual coupling.
The primitive polynomial coefficients have degree at most $k+1$ and bounded
coefficient norms after the bounded mean translation. The full residual is
the local bridge remainder plus the same $\mathsf P^{\mathsf T}$ synthesis
of port remainders. Its absolute column sum is uniformly finite. Minkowski
and the triangle inequality for polynomial coefficients therefore preserve
both estimates. There is no new native kernel comparison or exterior window.
\end{proof}

\begin{theorem}[Unrestricted-volume entrywise expansion of the native memory]
\label{f16v17:entry-expansion}
For every fixed $J\ge2$ there are $C_{J,r},\gamma_{J,r}$ independent of
$B,n$ such that for all admitted histories and $\gamma\ge\gamma_{J,r}$,
\begin{equation}
 \boxed{
 \sup_{\alpha,\beta}
 \left|\Gamma_{\gamma,\alpha\beta}
        -\sum_{q=2}^{J}\gamma^{-q/2}(G_q)_{\alpha\beta}\right|
       \le C_{J,r}\gamma^{-(J+1)/2}.}
 \label{f16v17:entry-remainder}
\end{equation}
The $G_q$ are exactly V10's connected coefficients
\eqref{f16v10:connected-coefficients}. The identical estimate holds for
$\mathsf M_\gamma$ and $M_q=-\operatorname{Off}_{>1}G_q$.
The norm in \eqref{f16v17:entry-remainder} is $\ell^1\to\ell^\infty$,
not an unrestricted operator norm or a time-weighted row norm.
\end{theorem}
\begin{proof}
Truncate both residuals through degree $J-1$. Their $L^2$ errors are
$O(\varepsilon^J)$ by Proposition~\ref{f16v17:source-jets}; their full and
truncated $L^2$ norms are $O(\varepsilon)$. Cauchy--Schwarz controls the
resulting covariance error by $C_{J,r}\varepsilon^{J+1}$.
For each product of source coefficients $r_{\alpha,k}r_{\beta,l}$ with
$k+l\le J$, apply Theorem~\ref{f16v17:expectations} through order
$J-k-l$. Apply it also to each factor when expanding the product of its
two means. Their polynomial coefficient norms are bounded uniformly, even
when their two source positions are arbitrarily distant. The omitted terms
with $k+l>J$ are bounded by fixed moments and cost
$O(\varepsilon^{J+1})$. There are finitely many terms at this fixed order.

The coefficient-ownership calculation in
Theorem~\ref{f16v17:expectations}, followed by subtracting the two native
mean expansions, is precisely the connected-cumulant formula of V10. Thus
its Gaussian term and all measure corrections are retained. In particular,
\[
 G_2=\operatorname{Cov}_Q(r_{\alpha,1},r_{\beta,1}),
\]
and $G_3$ contains the existing term
$-\kappa_Q(r_{\alpha,1},r_{\beta,1},V_1)$; it is not an expansion of the
observables alone. This proves \eqref{f16v17:entry-remainder} uniformly.
Finally V7's exact compact Haar identity, not the new chart identity, gives
$\mathsf M_\gamma=-\operatorname{Off}_{>1}\Gamma_\gamma$.
Applying that entry mask to the estimate proves the memory assertion.
No changed contact term or endpoint normalizer has been dropped.
\end{proof}

The new $J=2$ conclusion is an actual $o(\gamma^{-1})$ error for every
entry at unrestricted ambient volume:
\[
 \Gamma_{\gamma,\alpha\beta}
     =\gamma^{-1}(G_2)_{\alpha\beta}+O_r(\gamma^{-3/2}).
\]
This does not supersede V10's stronger temporal norm in its restricted
spatial regime. V10 pays $e^{K_{J,r}B}$ to retain geometric temporal decay of
the remainder; the present argument removes that spatial cost but does not
retain remainder decay. These are different theorems in different norms.

\subsection{A coherent conclusion on polynomially growing source banks}

V10 already proves $\|G_q\|_{\mathrm{loc},4a}\le C_{q,r}$ at one fixed
$a>0$, for each fixed $q$. We reuse that result, including its full-history
Gaussian mean and covariance, rather than repeat the linked-pairing proof.
It implies $\|G_q\|_{\rm op}+\|M_q\|_{\rm op}\le C_{q,r}$.
Let $\chi_{\mathcal A}$ select any finite set of distinct common-source
coordinates, with no restriction on their relative positions, and put
$k=|\mathcal A|$.

\begin{theorem}[Mixed-norm remainder and source-size control]
\label{f16v17:coherent}
Under the fixed-order hypotheses of Theorem~\ref{f16v17:entry-expansion},
write $R_{\gamma,J}=\Gamma_\gamma-\sum_{q=2}^J\gamma^{-q/2}G_q$.
For any real vectors $v,w$,
\begin{equation}
 |v^{\mathsf T}R_{\gamma,J}w|
       \le C_{J,r}\gamma^{-(J+1)/2}\|v\|_1\|w\|_1,
 \qquad
 \|\chi_{\mathcal A}R_{\gamma,J}\chi_{\mathcal A}\|_{\rm op}
       \le C_{J,r}k\gamma^{-(J+1)/2}.
 \label{f16v17:mixed-norm}
\end{equation}
The same estimates hold for the memory remainder. In particular
\begin{equation}
 \operatorname{Var}_{P_\varepsilon}(v^{\mathsf T}\mathbf r)
   \le C_{J,r}\left(\gamma^{-1}\|v\|_2^2
                +\gamma^{-(J+1)/2}\|v\|_1^2\right).
 \label{f16v17:variance}
\end{equation}
For every \emph{fixed} $p\ge0$ there are finite
$C_{p,r},\gamma_{p,r}$ such that all source banks with $k\le\gamma^p$
satisfy, at every ambient $B,n$ and every admitted history,
\begin{equation}
 \boxed{
 \begin{aligned}
 \|\chi_{\mathcal A}(\Gamma_\gamma-\gamma^{-1}G_2)
                         \chi_{\mathcal A}\|_{\rm op}
       &\le C_{p,r}\gamma^{-3/2},\\
 \|\chi_{\mathcal A}(\mathsf M_\gamma-\gamma^{-1}M_2)
                         \chi_{\mathcal A}\|_{\rm op}
       &\le C_{p,r}\gamma^{-3/2}.
 \end{aligned}}
 \label{f16v17:polynomial-banks}
\end{equation}
Consequently the restricted native covariance and memory have operator norms
at most $C_{p,r}/\gamma$. No average over source vectors or division by
$B(n+1)$ occurs.
\end{theorem}
\begin{proof}
Sum the entrywise remainder against $|v_\alpha w_\beta|$ to get the first
bound. On a bank of size $k$, $\|v\|_1\le\sqrt{k}\|v\|_2$ proves the
operator bound. The sum of the fixed-order coefficient operators has norm
at most $C_{J,r}/\gamma$ for $\gamma\ge1$, by the inherited absolute row
bounds. Its quadratic form plus the remainder gives \eqref{f16v17:variance}.
The memory mask preserves the entrywise estimate and the absolute coefficient
row bounds; positivity of the memory matrix is not asserted.

Given fixed $p$, choose a fixed integer $J\ge\max\{2,2p+2\}$. Then
$k\gamma^{-(J+1)/2}\le\gamma^{-3/2}$. The remaining finite sum of
coefficients with $q\ge3$ has operator norm at most
$C_{J,r}\gamma^{-3/2}$. Apply \eqref{f16v17:mixed-norm} and enlarge the
constant and threshold for this chosen $J$. This proves
\eqref{f16v17:polynomial-banks}. The argument is uniform over all choices
of the bank and all bounded bridge histories. The exponent $p$ is fixed
before taking the large-coupling limit; no bound on growth of the constants
as $p\to\infty$ is claimed.
\end{proof}

For $J=2$, the raw covariance remainder on a $k$-coordinate bank is
$C_rk\gamma^{-3/2}$, and the covariance itself is $O_r(\gamma^{-1})$
when $k\le\sqrt\gamma$. Higher fixed orders remove that particular source-size
restriction: \eqref{f16v17:polynomial-banks} identifies the same leading
operator on any prescribed polynomially growing bank. The rest of the
ambient torus and the full duration are unrestricted. This is not a theorem
about a growing cylinder on which all exterior fields have been fixed.

\begin{proposition}[What a remaining coherent obstruction must look like]
\label{f16v17:delocalization}
Define $k_{\rm eff}(v)=\|v\|_1^2/\|v\|_2^2$ for $v\ne0$.
For every fixed $J\ge2$ and unit vector $v$,
\begin{equation}
 \left|\gamma v^{\mathsf T}\Gamma_\gamma v-v^{\mathsf T}G_2v\right|
 \le C_{J,r}\gamma^{-1/2}
             +C_{J,r}\gamma^{-(J-1)/2}k_{\rm eff}(v).
 \label{f16v17:effective-support}
\end{equation}
Suppose along a sequence $\gamma\to\infty$ of admitted finite histories,
unit vectors $v_\gamma$ have the left side bounded below by a fixed
$\delta>0$. Then for every fixed $p\ge0$,
\begin{equation}
                  k_{\rm eff}(v_\gamma)/\gamma^p\longrightarrow\infty.
 \label{f16v17:superpolynomial}
\end{equation}
The same implication holds for $\mathsf M_\gamma,M_2$.
\end{proposition}
\begin{proof}
Multiply \eqref{f16v17:mixed-norm} by $\gamma$ and use the uniform operator
bounds for the higher coefficients. Their sum starts at
$\gamma^{-1/2}$, proving \eqref{f16v17:effective-support}.
For any fixed $p$, select a fixed $J$ with $(J-1)/2>p$.
Eventually the first term is at most $\delta/2$; hence
$k_{\rm eff}(v_\gamma)\ge c_{J,r,\delta}\gamma^{(J-1)/2}$.
Divide by $\gamma^p$. Each selected order has its own finite threshold,
which the sequence eventually crosses. No order is chosen as a function of
the instantaneous volume. The masked memory proof is identical.
\end{proof}

This is a necessary delocalization condition for an order-one failure of the
scaled leading-coefficient comparison, not a construction of such a failure
or a proof that it is impossible. There is no upper bound on ambient source
dimension. Very large volumes still admit vectors with this effective support.
The finite-source conclusion therefore does not give a global spectral gap.

\subsection{Separation, the posterior, and the next noncircular target}

The coefficients are exponentially local, whereas the new error is uniform
but unweighted. Combining the two gives, for every fixed $J\ge2$,
\begin{equation}
 |\Gamma_{\gamma,\alpha\beta}|
 \le C_{J,r}\gamma^{-1}e^{-4a d(\alpha,\beta)}
                          +C_{J,r}\gamma^{-(J+1)/2}.
 \label{f16v17:clustering-floor}
\end{equation}
The same estimate applies to the nonadjacent-time memory entries. Thus
separated response has an error floor smaller than any chosen algebraic
order, after choosing that fixed order. It is \emph{not} an absolutely
summable remainder at a fixed coupling. Summing the last term over all
source positions would again introduce the number of positions.

In particular this appendix does not replace V16's exact posterior
reduction by an unproved global one. V12's exterior messages are nonlinear
conditional expectations over entire cylinders; their covariance still
contains the actual exterior return. We do not identify that return with a
Gaussian marginal or infer a global sampler gap from the local Stein
recursion. V15's repair theorem remains a possible tool for a genuinely
source-resolved control of that return.

What has changed is the source-size frontier. A leading native coefficient
with an error $o(\gamma^{-1})$ is now identified at unrestricted ambient
volume in the entrywise norm, and in operator norm on every prescribed
polynomially growing source bank. Another repair neighbourhood or another
extensive trace budget was not needed. The unresolved upgrade is to an
\emph{absolutely summable connected remainder}, or an independent global
native inequality controlling the remaining delocalized directions. The
polynomial divergence recursion supplies coefficients with a full native
error, but its coefficient-$\ell^1$ estimate does not manufacture that
connected decay. It is important not to rename that loss as a closed
thermodynamic estimate.

All statements still concern a prescribed bounded rescaled bridge history,
before physical root Gauss restriction and integration of the bridges.
No new bridge-coverage theorem, physical-time calibration, infrared contact
coercivity, same-top physical transfer estimate, or interacting continuum
construction is obtained here. The truncations are approximations to an
observable covariance, not positive replacement transfer operators. No
convergence or unique infinite resummation of the formal series is asserted.

\clearpage
\subsection{Diagnostics, preservation, and claim scope}

The new diagnostic is
\begin{center}\small
\path{verify_ym_f16_v17_local_Stein.py}.
\end{center}
It tests the finite Gaussian divergence recursion with correlated rational
precisions, normalized perturbative moment and source-covariance coefficients,
measured Haar jets at both endpoint weights, the cutoff integration identity,
and the mixed source-norm and order bookkeeping. All 504 exact assertions
passed; the report records their families and the distinction between
fixtures and analytic proofs.
The fixtures are not native lattice simulations, enclosures of analytic
constants, or proof-assistant verification of the full-volume estimates.

The cumulative V16 is preserved byte for byte except for its title version
and insertion of this exact appendix before the final document-ending
command. The accompanying package records hashes, predecessor recovery,
actual diagnostic reruns, compilation and rendered-page inspection. No
previous bibliography string or inherited theorem is silently corrected.
The available V4--V16 diagnostics were rerun unchanged and passed 9,467,
89,339, 906, 311, 609, 990, 608, 463, 1,201, 222, 622, 434, and 1,077
assertions, respectively. Unavailable diagnostics are not represented as
rerun. The next scheduled five-version companion checkpoint is V20.

\begin{firewall}
V17 proves fixed-order polynomial expectation and corrected-source covariance
expansions under the full native fixed-bridge law, with entrywise errors
uniform in ambient volume and duration. It identifies the leading native
coefficient in operator norm on every prescribed polynomially growing source
bank and quantifies the delocalization required of a remaining coherent
failure. The proof uses V16's native moments and an explicit Gaussian
polynomial divergence solver, not a native sampler gap. No unrestricted
full-matrix operator bound, summable connected remainder, global auxiliary
gap, physical transfer gap or interacting continuum Yang--Mills mass gap
is established by this appendix.
\end{firewall}
% END F16 V17 APPEND

% BEGIN F16 V18 APPEND -- 2026-09-17
% Insert immediately before the final document-ending command in V17.
% The predecessor body is unchanged; only its cumulative title becomes V18.

\clearpage
\section{V18: exterior screening and full-history trace asymptotics}
\label{f16v18:context}

V17 controls polynomial tests but leaves its actual exterior messages unestimated
at arbitrary order. Repeating its entrywise expansion would not supply a
summable remainder. We instead allow an \emph{arbitrary square-integrable
exterior probe} in the native Stein identity. No derivative of that probe is
taken. Exponential localization of the Gaussian divergence solver pays the
coordinates on which integration by parts is not allowed. This gives a native
$L^2$ estimate for normalized conditional means, not just for finitely many
polynomial exterior tests.

The resulting posterior return is smaller than any prescribed algebraic order
in coupling \emph{in trace density}, with a logarithmic collar. A second step
uses separated, finite space--time blocks, rather than entire-time spatial
cylinders. The new message estimate pays their off-block covariance; V17 pays
their finite source banks; positivity pays the discarded corridors. Together
these prove a full-matrix normalized trace-norm expansion with the same
coefficients as V10, at unrestricted ambient volume and duration. This is not
a dimension-independent operator-norm expansion or exponential clustering at
fixed coupling. The distinction remains explicit throughout.

\subsection{Dependency audit and the norm selected by the new argument}

The supplied V17 appendix and its audit were read in full. The relevant V16
moment proof and the earlier exact source and posterior interfaces were checked
against the cumulative source. Two scoped semantic archive searches returned
no results. The mounted inventories were then searched, and identified
passages read: f14 V93's finite-volume Stein estimate, f15 V14's conditional-copy
posterior transfer, and the monograph's Wilson Markov-blanket and boundary
mediation statements. Those general mechanisms are credited, not presented
as missing from the archive. This is a targeted dependency review, not
independent recertification of the complete historical proofs.

V16's full-native local moments, V17's polynomial divergence solver and
finite-order expectation coefficients, V10's linked-coefficient locality,
and V7's exact Haar source synthesis are the analytic inputs. We do not assume
a native gap, a typical deterministic exterior window, or the desired
connected-correlation estimate. The new statements below are deductions from
these identified inputs. The next scheduled companion checkpoint remains V20.

The wider-literature memorandum correctly keeps summable interactions and the
subsequent global functional inequality as distinct requirements. The other
memorandum separates order gain from locality cost. The following proof does
not replace either requirement by an analogy. For context, M.~Lohmann,
\href{https://arxiv.org/abs/1411.1107}{\emph{Single Scale Cluster Expansions with
Applications to Many Boson and Unbounded Spin Systems}}, studies exponential
clustering with a massive Gaussian part and controlled remaining interactions.
Its abstract and introduction were reviewed, not its complete cluster theorem.
A.~Abdesselam, A.~Procacci and B.~Scoppola,
\href{https://arxiv.org/abs/0901.4756}{\emph{Clustering Bounds on N-Point
Correlations for Unbounded Spin Systems}}, gives related cluster-expansion
context; only the primary abstract was screened. No theorem from either paper
is imported into the compact Wilson law. In particular a weighted norm on
\emph{formal coefficients} is not a verified cluster norm for the full native
remainder. The conclusion here is a different norm, proved directly.

\subsection{Anchored polynomial norms and the inherited Gaussian solver}

Write $\varepsilon=\gamma^{-1/2}$, $\mathcal N=B(n+1)$ and retain
\[
 P=P_{\gamma,\mathbf y},\qquad Q=N(m,C),\qquad C=H^{-1},\qquad
 \xi=\zeta-m,\qquad \max_{i,e}|y_{i,e}|\le r.
\]
The native probability is the complete unmoving fast law of V5 and V16. Its
principal-chart domains, Haar powers, and two dressed endpoints are unchanged.
At fixed rescaled bridge data the Gaussian mean and precision do not depend
on $\varepsilon$. Expectations without a subscript in this appendix are under
$P$. Scalar-coordinate labels are denoted by $i,j$; $c(i)$ is their cell.

Use the ambient cell distance
\[
 d_*((t,b),(u,c))=|t-u|+d_{\mathbb T_m^3}(b,c),
\]
where spatial distance is nearest-neighbour torus distance. A fixed integer
$\rho\ge1$ bounds the diameter of each primitive carrier, the range of $H$,
and the carrier radius of each ordinary source or primitive error. Increasing
$\rho$ only changes fixed constants. Gaussian inverse locality and the port
inverse locality give one $\vartheta>0$ for which
\begin{equation}
 \sup_i\sum_j |C_{ij}|e^{\vartheta d_*(c(i),c(j))}<\infty.
 \label{f16v18:weighted-C}
\end{equation}
We choose $\vartheta$ smaller if necessary than the exponent in V10's
coefficient-locality theorem and the port-lift exponent. It is fixed
independently of all expansion orders, volumes and durations.

Fix an anchor cell $a$. For a monomial $\xi^{\mathbf k}$ put
\[
 R_a(\mathbf k)=\max_{i:k_i>0}d_*(a,c(i)),\qquad R_a(0)=0.
\]
For a polynomial $f=\sum f_{\mathbf k}\xi^{\mathbf k}$ and a polynomial vector
$g=(g_j)$ define
\begin{align}
 \|f\|_{a,\vartheta}
  &=\sum_{\mathbf k}|f_{\mathbf k}|e^{\vartheta R_a(\mathbf k)},\notag\\
 \|g\|_{a,\vartheta,1}
  &=\sum_{j,\mathbf k}|g_{j,\mathbf k}|
       e^{\vartheta\max\{R_a(\mathbf k),d_*(a,c(j))\}}.
 \label{f16v18:anchored-norms}
\end{align}
These are weighted coefficient norms, not native random norms. They charge
actual coefficient sums. No claim that arbitrary extensive polynomials have
bounded norms in \eqref{f16v18:anchored-norms} is made.

\begin{proposition}[Anchored propagation and the cost of unused coordinates]
\label{f16v18:anchored-solver}
For the \emph{same} polynomial map $\mathcal S$ as in
\eqref{f16v17:solver-recursion}, every polynomial of degree at most $d$ obeys
\begin{equation}
 \delta_Q\mathcal Sf=f-\mathbb E_Qf,\qquad
 \|\mathcal Sf\|_{a,\vartheta,1}\le K_{d,\vartheta}\|f\|_{a,\vartheta}.
 \label{f16v18:solver-weighted}
\end{equation}
The same operators $\mathcal T_q=-\nabla V_q\cdot\mathcal S$ as V17 obey
\begin{equation}
 \deg\mathcal T_qf\le d+q,\qquad
 \|\mathcal T_qf\|_{a,\vartheta}\le C_{d,q,r}\|f\|_{a,\vartheta}.
 \label{f16v18:T-weighted}
\end{equation}
Let $I$ be any union of complete coordinate cells and put
$L=d_*(a,I^c)$, with $L=\infty$ if $I^c$ is empty. If $g$ has bounded degree,
then, for every fixed finite $p$,
\begin{equation}
 \left\|\sum_{j:c(j)\notin I}
       [(H\xi)_jg_j-\partial_jg_j]\right\|_{L^p(P)}
       \le C_{d,p,r}e^{-\vartheta L}\|g\|_{a,\vartheta,1}.
 \label{f16v18:outside-divergence}
\end{equation}
All constants are independent of $I,B,n$.
\end{proposition}
\begin{proof}
Only the weighted estimates are new here; the exact divergence identity is
V17's finite algebra. For a degree-$d$ monomial with first selected index $i_1$
and remaining product $h$, the first vector of that recursion is
$C_{\cdot i_1}h$. Its radius after adding component $j$ is at most
$R_a(\mathbf k)+d_*(c(i_1),c(j))$. Equation \eqref{f16v18:weighted-C}
therefore bounds its weighted coefficient sum by a constant times the
original monomial weight. Each contraction removes two indices and cannot
increase the remaining radius. Its coefficient is bounded by $\sup|C_{ij}|$.
The terminating degree induction of V17 gives
\eqref{f16v18:solver-weighted}. Constants are sent to zero.

The polynomial $\partial_jV_q$ contains only the fixed number of original
interaction terms meeting $j$. Its variables lie within radius $\rho$ of
$c(j)$, its degree is at most $q+1$, and its coefficient norm after the bounded
mean translation is at most $C_{q,r}$. Multiplication with $g_j$ increases
the radius in the vector norm by at most $\rho$. Summing $j$ proves
\eqref{f16v18:T-weighted}. The measured coefficients $V_q$ include the original
Haar powers; no new local potential is substituted.

For outside components, the unweighted coefficient sum is at most
$e^{-\vartheta L}\|g\|_{a,\vartheta,1}$. Multiplication by $(H\xi)_j$ costs the
uniform absolute row bound for $H$; differentiation costs the fixed degree.
V17's polynomial $L^p$ estimate under the actual native law, followed by
Minkowski, gives \eqref{f16v18:outside-divergence}. The same proof bounds
$\sum_{j\notin I}(\partial_jV_q)g_j$ by
$C_{d,q,p,r}e^{-\vartheta L}\|g\|_{a,\vartheta,1}$. This last bound will pay
outside drift terms as well.
\end{proof}

The radius norm is chosen for this exterior test. It is not a norm forcing a
connected tree between two arbitrary source marks. No such conclusion is
built into its definition.

\subsection{Stein's identity tested against arbitrary native exterior functions}

Let $\mathcal F_O=\sigma(\zeta_{I^c})$. The test $h\in L^2(P,\mathcal F_O)$ may
depend on \emph{every} exterior coordinate and can be nonsmooth. Its norm is
under the actual exterior marginal of $P$, including all its conditional
partition functions. It need not be a polynomial, a cylinder function, or
bounded independently of the regulator.

\begin{theorem}[An exterior-tested native identity]
\label{f16v18:tested-Stein}
For every fixed degree $d$ and order $N\ge0$, every polynomial vector $g$ of
degree at most $d$, and every $h\in L^2(P,\mathcal F_O)$,
\begin{equation}
 \left|\mathbb E[h\,\delta_Qg]
       +\sum_{q=1}^{N}\varepsilon^q
                    \mathbb E[h\,\nabla V_q\cdot g]\right|
 \le C_{d,N,r}\bigl(\varepsilon^{N+1}+e^{-\vartheta L}\bigr)
                  \|h\|_2\|g\|_{a,\vartheta,1}.
 \label{f16v18:tested-identity}
\end{equation}
The coupling threshold depends only on $d,N,r$, not on the size, shape or
number of exterior coordinates. The empty sum convention applies for $N=0$.
\end{theorem}
\begin{proof}
Use V17's group-local cutoff
$\chi_j=\chi(\varepsilon\zeta_{g(j)})$, where $\chi$ is one on radius $1/2$
and zero beyond radius one. For $j\in I$ the whole group $g(j)$ is in $I$.
The factor $h$ is constant in that group, so ordinary coordinate integration
by parts, at fixed remaining coordinates, differentiates no part of $h$.
Initially take bounded measurable $h$. With the exact measured potential
$V_\varepsilon=U_\varepsilon-\log J_\varepsilon^f$, the identity is
\begin{align}
 \mathbb E[h\delta_Qg]
 ={}&-\sum_{j\in I}\mathbb E[
       h\chi_jg_j(\partial_jV_\varepsilon-(H\xi)_j)]\notag\\
 &+\sum_{j\in I}\mathbb E[h(\partial_j\chi_j)g_j]\notag\\
 &+\sum_{j\in I}\mathbb E[h(1-\chi_j)
                ((H\xi)_jg_j-\partial_jg_j)]\notag\\
 &+\sum_{j\notin I}\mathbb E[h((H\xi)_jg_j-\partial_jg_j)].
 \label{f16v18:exact-tested-cut}
\end{align}
Here $j\in I$ abbreviates $c(j)\in I$. The final line is an algebraic
remainder, not integration by parts on the forbidden exterior coordinates.
This distinction is the reason no regularity of $h$ is required.

The classical word-drift remainder through order $N$ has $L^4(P)$ norm at
most $C_{N,r}\varepsilon^{N+1}$ per coordinate, by the all-real Taylor
bounds and V16's local moments, exactly as in V17. On $\chi_j$'s support the
only differentiated Haar factor is a smooth function of its own group inside
radius one. Its original weight, including the endpoint weight $1/2$, is
retained. Its measured drift remainder has the same bound. Thus
\[
 \|\chi_j g_j[\partial_jV_\varepsilon-(H\xi)_j
                 -\sum_{q=1}^N\varepsilon^q\partial_jV_q]\|_2
       \le C_{d,N,r}\varepsilon^{N+1}\|g_j\|_{\rm pol}.
\]
Pair with $h$ by Cauchy--Schwarz. No moment bound under a new conditional or
intermediate coupling is being used.

For every fixed-degree polynomial $p$, V16's one-cell exponential moment and
H\"older imply
\[
 \|1_{\{|\zeta_{g(j)}|\ge(2\varepsilon)^{-1}\}}p\|_2
       \le C_{d,r}e^{-c/\varepsilon^2}\|p\|_{\rm pol}.
\]
This is the $L^2$ version of V17's local cut-tail estimate: use the fourth
moment of $p$ and the fourth root of the event probability. It pays the
second and third lines of \eqref{f16v18:exact-tested-cut}, and removal of
$\chi_j$ from the finite drift polynomials. Summation costs
$\sum_j\|g_j\|_{\rm pol}$, not the total number of coordinates.

Finally \eqref{f16v18:outside-divergence} and the outside drift bound from
its proof pay the last line and all $j\notin I$ polynomial drift terms by
$C_{d,N,r}e^{-\vartheta L}\|h\|_2\|g\|_{a,\vartheta,1}$. Absorb the
coordinate-local flat term into the fixed-order $\varepsilon^{N+1}$ budget.
This proves \eqref{f16v18:tested-identity} for bounded measurable $h$.
All other factors in its terms are in $L^2(P)$ by the estimates just given.
Truncating an arbitrary exterior $L^2$ function and passing in $L^2$ proves
the stated theorem. No all-group cutoff changes $P$ at any point.
\end{proof}

\subsection{Native conditional means, including their true normalizers}

Let $b_k(f)$ be precisely V17's coefficient recursion
\eqref{f16v17:coefficient-recursion}. In particular $b_0(f)=\mathbb E_Qf$;
no coefficient is fitted to the exterior. These numbers can depend on the
entire prescribed bridge history through $m$ and $C$, but that history is a
deterministic parameter, not the random fast exterior.

\begin{theorem}[All fixed orders of exterior conditional screening]
\label{f16v18:conditional-polynomials}
For every fixed polynomial degree $d$ and order $N\ge0$,
\begin{equation}
 \left\|\mathbb E[f\mid\mathcal F_O]
                  -\sum_{k=0}^N\varepsilon^k b_k(f)\right\|_2
 \le C_{d,N,r}\bigl(e^{-\vartheta L}+\varepsilon^{N+1}\bigr)
                         \|f\|_{a,\vartheta}.
 \label{f16v18:conditional-expansion}
\end{equation}
Consequently
\begin{equation}
 \|\mathbb E[f\mid\mathcal F_O]-\mathbb E f\|_2
 \le C_{d,N,r}\bigl(e^{-\vartheta L}+\varepsilon^{N+1}\bigr)
                         \|f\|_{a,\vartheta}.
 \label{f16v18:conditional-centered}
\end{equation}
These statements are uniform in ambient volume and duration, for sufficiently
large coupling depending only on the displayed fixed degree and order.
They are averaged in the native exterior law, not suprema over deterministic
exterior values.
\end{theorem}
\begin{proof}
Apply Theorem~\ref{f16v18:tested-Stein} to $g=\mathcal Sf$. Its divergence
identity gives, for every exterior $h$,
\begin{equation}
 \mathbb E[h f]=(\mathbb E h)\mathbb E_Q f
      +\sum_{q=1}^N\varepsilon^q\mathbb E[h\mathcal T_q f]
      +O\!\left((e^{-\vartheta L}+\varepsilon^{N+1})
                       \|h\|_2\|f\|_{a,\vartheta}\right).
 \label{f16v18:tested-recursion}
\end{equation}
For $N=0$ this is the assertion in dual form. Induct on $N$, applying the
order-$N-q$ formula to each $\mathcal T_q f$, with the \emph{same} exterior
$h$, the same set $I$, and the same anchor. Its weighted coefficient norm is
bounded by \eqref{f16v18:T-weighted}. Multiplying by $\varepsilon^q$ makes
every algebraic remainder order $N+1$; its spatial error is bounded by a
constant times $e^{-\vartheta L}$ because $\varepsilon\le1$. There are only
finitely many terms at each fixed order. Their coefficients are exactly
$b_k$, not a new conditional recursion. Thus
\[
 |\mathbb E[h(f-\sum_{k=0}^N\varepsilon^k b_k(f))]|
 \le C_{d,N,r}(e^{-\vartheta L}+\varepsilon^{N+1})
                         \|h\|_2\|f\|_{a,\vartheta}.
\]
Hilbert-space duality on $L^2(P,\mathcal F_O)$ proves
\eqref{f16v18:conditional-expansion}. Apply the orthogonal projection
$I-\mathbb E_P$ on that exterior space to obtain
\eqref{f16v18:conditional-centered}; all subtracted coefficients are constants.
This proof permits coefficients involving distant Gaussian variables inside
$\mathcal T_q f$: their anchored weights, not a false finite support claim,
pay the exterior tail at each step.
\end{proof}

This extends the class of \emph{tests}, not merely the polynomial order in
V17. The dual test can be the sign of a conditional message, or its centered
normalized value, and may inspect an arbitrarily large exterior. No regularity,
polynomial approximation, or reference probability for that test is required.
The factor $\varepsilon^{N+1}$ remains when $L$ grows at a fixed coupling;
this is not a proof of strong mixing as $L\to\infty$ with $\gamma$ fixed.

\begin{proposition}[Application to ordinary scores and Haar residuals]
\label{f16v18:source-screening}
Let $J_{\gamma,\alpha}=\partial_\alpha U_\gamma$ be an ordinary common-source
score, anchored at $a(\alpha)$, and let $r_{\gamma,\alpha}$ be the exact
Haar-corrected residual of V7. For every fixed $N\ge1$,
\begin{align}
 \|\mathbb E[J_{\gamma,\alpha}\mid\mathcal F_O]
                          -\mathbb E J_{\gamma,\alpha}\|_2
   &\le C_{N,r}(e^{-\vartheta L}+\varepsilon^{N+1}),\notag\\
 \|\mathbb E[r_{\gamma,\alpha}\mid\mathcal F_O]
                          -\mathbb E r_{\gamma,\alpha}\|_2
   &\le C_{N,r}(\varepsilon e^{-\vartheta L}+\varepsilon^{N+1}).
 \label{f16v18:source-centered}
\end{align}
The second bound also holds for every component of the primitive bank
$e^{\rm prim}=(e^{\rm br},e^{\rm pt})$.

If $J_{\gamma,\alpha}\sim\sum_{k\ge0}\varepsilon^kJ_{\alpha,k}$, define
$j_{\alpha,q}=\sum_{k=0}^q b_{q-k}(J_{\alpha,k})$. Then also
\begin{equation}
 \left\|\mathbb E[J_{\gamma,\alpha}\mid\mathcal F_O]
                         -\sum_{q=0}^N\varepsilon^qj_{\alpha,q}\right\|_2
       \le C_{N,r}(e^{-\vartheta L}+\varepsilon^{N+1}).
 \label{f16v18:source-expansion}
\end{equation}
The first coefficients are $j_{\alpha,0}=g_\alpha$ and
$j_{\alpha,1}=\mathbb E_Q r_{\alpha,1}$, as in V11's deterministic reference.
\end{proposition}
\begin{proof}
The ordinary score coefficients have fixed local carriers and bounded
fixed-degree coefficient norms. Their complete-chart Taylor remainder in
$L^2(P)$ is $O(\varepsilon^{N+1})$, by the same inserted-word bounds and
V16 moments used in V17. The primitive errors have the same properties but
start at order one. Their weighted coefficient norms are therefore uniformly
bounded at each fixed order. For the full Haar residual, the spatially
exponentially summable port-lift column preserves these weighted norms after
decreasing $\vartheta$ once, independently of order. Its exact native $L^2$
jet remainder remains that of Proposition~\ref{f16v17:source-jets}.

Apply Theorem~\ref{f16v18:conditional-polynomials} to each coefficient,
through the remaining order $N-k$, and use conditional $L^2$ contraction on
the native jet remainder. The spatial error is $O(e^{-\vartheta L})$ for
the ordinary score and $O(\varepsilon e^{-\vartheta L})$ for a bank starting
at order one. Centering is a contraction and proves
\eqref{f16v18:source-centered}. Collecting coefficients gives
\eqref{f16v18:source-expansion}.

Coefficient ownership is unchanged: V17's $b_k$ are normalized native
coefficients. In particular the direct formula for the first score correction
is $\mathbb E_QJ_{\alpha,1}-\operatorname{Cov}_Q(J_{\alpha,0},V_1)$.
The exact Haar identity $\mathbb E J_{\gamma,\alpha}
=g_\alpha+\mathbb E r_{\gamma,\alpha}$ identifies it with
$\mathbb E_Qr_{\alpha,1}$ by uniqueness of fixed-order coefficients.
The half-bracket term in $r_{\alpha,1}$ and the measure correction are not
removed. No derivative of a conditional approximation is taken.
\end{proof}

\subsection{The actual spatial-cylinder posterior has arbitrary-order trace accuracy}

Return first to the literal separated-cylinder setting of V12. Source marks
are interior to their cylinders, with ordinary score carriers strictly inside,
and their anchor distance to the exterior is at least $L$. There are no
repeated selected marks. The common exterior posterior retains
\[
 d\bar P(q)=Z_\gamma^{-1}1_{\mathcal P'_O}J^f_{\gamma,O}(q)
      e^{-U_{\gamma,O}(q;\mathbf y)}
                       \prod_sZ_{\gamma,D_s}(q;\mathbf y)\,dq.
\]
No conditional partition function in this expression is replaced. Let
$\mathbf j$ be the same vector of derivatives of $-\log Z_{\gamma,D_s}$,
$k$ its number of coordinates, $\mathsf R_{\rm post}=\operatorname{Cov}_{\bar P}
(\mathbf j)$, and $\mathsf L_D$ and $\chi$ be V12's objects.

\begin{theorem}[All-order posterior return and unchanged cylinder gluing]
\label{f16v18:posterior}
For every fixed $N\ge1$, at sufficiently large coupling depending only on
$N,r$, and for all ambient volumes and durations,
\begin{equation}
 \operatorname{Tr}\mathsf R_{\rm post}
     \le C_{N,r}k\bigl(e^{-2\vartheta L}+\gamma^{-(N+1)}\bigr),\qquad
 \|\chi\mathsf M_\gamma\chi-\mathsf L_D\|_{S_1}
     \le C_{N,r}k\bigl(e^{-2\vartheta L}+\gamma^{-(N+1)}\bigr).
 \label{f16v18:posterior-trace}
\end{equation}
Every time lag is included at once. In particular, for $k\le C\mathcal N$,
$L\ge(N+1)(2\vartheta)^{-1}\log\gamma$ gives both normalized trace bounds
$C_{N,r}\gamma^{-(N+1)}$. Such a collar must fit in the chosen cylinders;
no bound on their volume relative to coupling is otherwise imposed.
\end{theorem}
\begin{proof}
For an interior mark, exact normalized differentiation gives
$j^D_{\gamma,\alpha}=\mathbb E[J_{\gamma,\alpha}\mid\zeta_{D^c}]$.
No exterior-only action term is differentiated. Separation makes this the
same function as the conditional mean under the one common exterior $O$;
the other interiors have no interaction with this interior after $O$ is fixed.
Apply the first bound of \eqref{f16v18:source-centered} to each scalar
message. The law of that message under $P$ is precisely its law under the
actual $\bar P$. Sum the squared estimates and use
$(u+v)^2\le2u^2+2v^2$. The trace is the sum of those variances, proving
the first inequality. No covariance between different messages is claimed
zero.

The inherited exact identity is
$\chi\mathsf M_\gamma\chi-\mathsf L_D
=-\operatorname{Off}_{>1}\mathsf R_{\rm post}$.
The matrix on the right before extraction is positive, so its trace norm is
its trace. V7's duration-independent bound for that extraction is four.
This proves the second inequality. The single-time endpoint normalizers
remain in the original history weight; their nonadjacent mixed derivatives
vanish for the original reason, not because they were removed.
\end{proof}

These are actual exterior-message estimates, unlike a polynomial test of a
Gaussian exterior. They remove the $1/\gamma$ \emph{floor} from V16's posterior
trace bound at a chosen logarithmic collar. They do not identify V11's
conditional covariance for each deterministic exterior, or prove a uniform
native conditional gap. For a bank of size $k\le\gamma^p$, trace domination
also gives an unnormalized operator bound $C_{p,s,r}\gamma^{-s}$ for any
fixed $p,s>0$, by choosing $N+1\ge p+s$ and
$L\ge(p+s)(2\vartheta)^{-1}\log\gamma$. The constants and threshold depend
on the chosen orders. The entire unrestricted source bank can still be much
larger than $\gamma^p$.

\subsection{Finite space--time blocks and their exact connected remainder}

For the full-matrix conclusion, work first with the primitive errors, which
have finite carriers. A full lifted residual is spatially nonlocal and must
not be declared independent of another block just because its anchor lies
there. Write
\[
 e=e^{\rm prim},\qquad K_\gamma=\operatorname{Cov}_P(e),\qquad
 r=Te,\qquad T=(I,\mathsf P^{\mathsf T}),\qquad \|T\|\le C.
\]
There are at most $30\mathcal N$ primitive scalar components: $9\mathcal N$
common bridge errors and $21Bn$ port errors. The matrix $T$ is the original
Gaussian lift synthesis, independent of coupling at fixed geometry.

\begin{proposition}[Primitive coefficient inputs and a block return identity]
\label{f16v18:primitive-blocks}
The normalized coefficients $K_q$ of $K_\gamma$ satisfy, at every fixed
order $q\ge2$,
\begin{equation}
 \|K_q\|_{\mathrm{loc},\vartheta}\le C_{q,r},\qquad
 \sup_{u,v}\left|K_{\gamma,uv}-\sum_{q=2}^K\varepsilon^q(K_q)_{uv}\right|
                  \le C_{K,r}\varepsilon^{K+1}.
 \label{f16v18:primitive-inputs}
\end{equation}
Here the local norm is the maximum weighted absolute row and column sum on
primitive anchors. Moreover $G_q=TK_qT^{\mathsf T}$ exactly.

Let $(I_s)$ be separated free cell sets such that no native interaction meets
two sets, and select primitive banks $\mathcal A_s$ whose carriers lie inside
$I_s$ and whose anchor depth is at least $L$. Write $\Pi_s$ for their disjoint
coordinate projections, $\Pi=\sum_s\Pi_s$, and
$K_{\rm blk}=\sum_s\Pi_sK_\gamma\Pi_s$. Then
\begin{equation}
 \|\Pi K_\gamma\Pi-K_{\rm blk}\|_{S_1}
 \le C_{N,r}|\mathcal A|
          \bigl(\varepsilon^2e^{-2\vartheta L}
                        +\varepsilon^{2N+2}\bigr),\qquad N\ge1,
 \label{f16v18:block-return}
\end{equation}
where $\mathcal A=\bigcup_s\mathcal A_s$. The blocks may be finite in time;
they need not be V11's cylinders.
\end{proposition}
\begin{proof}
V17's native jet statement explicitly includes both primitive banks. Applying
its expectation expansion to their two products and means gives the entry
estimate in \eqref{f16v18:primitive-inputs}, by the same finite coefficient
calculation. Its coefficients are the connected Gaussian cumulants of the
primitive jets with the measured $V_q$. The linked-pairing proof of V10
applies to these coefficients without the port-lift edges: every primitive
carrier is local, covariance rows decay exponentially, and every surviving
Gaussian contraction connects the two marks. Summing a spanning tree bounds
the weighted row at one fixed exponent and each fixed order. This is a
specialization of that inherited argument, not a new native cluster
expansion. Since $r=Te$ exactly and $T$ is coupling independent, uniqueness
of the finite-dimensional coefficients gives $G_q=TK_qT^{\mathsf T}$.

Condition on $O=(\bigcup_s I_s)^c$. The product coordinate reference factors
and the local native action imply exact conditional independence of the
$I_s$. This includes the fractional Haar powers and the original dressed
endpoints, which are cell local. Put
$m_s=\mathbb E[e_{\mathcal A_s}\mid\mathcal F_O]$ and
$R=\operatorname{Cov}_P((m_s)_s)$. Total covariance gives
\[
 \Pi K_\gamma\Pi=\bigoplus_s\mathbb E\operatorname{Cov}
                      (e_{\mathcal A_s}\mid\mathcal F_O)+R.
\]
The $s$ diagonal block of this equality is the \emph{unconditional} covariance
$\Pi_sK_\gamma\Pi_s$, not a chosen boundary covariance. Hence exactly
\begin{equation}
 \Pi K_\gamma\Pi-K_{\rm blk}=R-\sum_s\Pi_sR\Pi_s.
 \label{f16v18:pinching-return}
\end{equation}
Each message $m_s$ agrees with conditioning on $I_s^c$ alone, by the same
finite-range specification. The primitive estimate in
\eqref{f16v18:source-centered} bounds the trace of $R$ by the right side
of \eqref{f16v18:block-return}, up to a fixed factor. Both $R$ and its block
pinching are positive and have the same trace. Their trace-norm difference
is at most twice that trace, proving the claim. No product approximation of
the actual exterior posterior has been used.
\end{proof}

The estimate is genuinely connected: it bounds the covariance \emph{between}
conditionally independent interiors through their actual posterior means.
It remains a trace estimate after summing a bank, not a bound on each
pair's remainder by its separation alone.

\begin{proposition}[A geometric partition with sparse discarded corridors]
\label{f16v18:partition}
For integers $L\ge\rho+1$ and $R\ge16(L+3\rho)$ there are separated free
cell sets $I_s$ and primitive banks $\mathcal A_s$ as above such that
\begin{equation}
 |I_s|\le(4R)^4,\qquad |\mathcal A_s|\le C R^4,\qquad
 |\mathcal A^c|\le C\mathcal N\frac{L+\rho}{R}.
 \label{f16v18:partition-bounds}
\end{equation}
The last complement is within the full primitive bank. Distinct selected
banks have anchor separation at least $L$. Constants are independent of
spatial periods and of the finite time length.
\end{proposition}
\begin{proof}
In each of the three coordinate cycles, and in the finite time path, leave a
direction uncut if its length is at most $4R$. Otherwise partition it into
$k=\lfloor\text{length}/R\rfloor$ consecutive intervals, distributing the
lengths as equally as possible. Each length is between $R$ and $2R$. In a
cycle the wrap-around cut is included. In the time direction only interior
cuts are used; the actual two history endpoints are not artificial cuts.
The products of these intervals, or complete uncut directions, are tiles of
at most $(4R)^4$ cells.

Remove a strip of width $\rho$ on each side of every genuine tile cut to
obtain $I_s$. Its gap from another free set exceeds the interaction range,
so no primitive action term meets two free sets. Select primitive anchors
in the same tile with distance at least $L+3\rho$ from every genuine cut.
Their complete carriers are inside $I_s$, and their distance to $I_s^c$ is
at least $L$. Where a direction is uncut and fills its entire cycle or time
path, it creates no exterior face and no depth condition in that direction.
This also covers the case of one tile filling the entire system; its exterior
is empty.

In a cut direction the number of removed coordinate positions is at most
$C(L+\rho)$ times the number of cuts, which is at most its length divided by
$R$. Taking the union of these strips in four dimensions bounds the discarded
cell count by $C\mathcal N(L+\rho)/R$. There are only a bounded number of
primitive types per anchor cell, including at time endpoints. Increasing the
strip allowance by the fixed carrier radius proves the last assertion of
\eqref{f16v18:partition-bounds}. Source counts per tile give the first two.
A path between selected anchors of different tiles must cross a genuine cut,
so their ambient distance is at least $L$. No translation symmetry of the
probability law is used.
\end{proof}

\subsection{All fixed orders in normalized trace norm on the entire history}

\begin{theorem}[Full-matrix native trace asymptotics at unrestricted volume]
\label{f16v18:full-trace}
For every fixed integer $J\ge2$ and fixed bridge window $r$, there are finite
$C_{J,r},\gamma_{J,r}$ independent of $B,n$ such that, for all admitted bridge
histories and $\gamma\ge\gamma_{J,r}$,
\begin{equation}
 \boxed{
 \frac1{\mathcal N}
 \left\|\Gamma_\gamma-\sum_{q=2}^J\gamma^{-q/2}G_q\right\|_{S_1}
 +\frac1{\mathcal N}
 \left\|\mathsf M_\gamma-\sum_{q=2}^J\gamma^{-q/2}M_q\right\|_{S_1}
       \le C_{J,r}\gamma^{-(J+1)/2}.}
 \label{f16v18:trace-expansion}
\end{equation}
Both matrices are the entire original common-source matrices; no bank is
omitted in the conclusion. There is no volume-dependent coupling-entry
condition. The norm is the sum of singular values divided by
$\mathcal N=B(n+1)$, not an unnormalized operator norm.
\end{theorem}
\begin{proof}
We prove the primitive assertion and then synthesize. All constants in this
proof can depend on fixed $J,r$. Decrease $\vartheta$ so that both anchored
screening and primitive coefficient locality use that exponent. Write
$K^{[J]}=\sum_{q=2}^J\varepsilon^qK_q$. Take
\begin{equation}
 L=\left\lceil\frac{J+1}{\vartheta}\log(\varepsilon^{-1})\right\rceil
                      +\rho+1,\qquad
 R=\lceil\varepsilon^{-2J}\rceil,\qquad K=9J.
 \label{f16v18:trace-scales}
\end{equation}
For sufficiently small $\varepsilon$, independently of $B,n$,
$R\ge16(L+3\rho)$. Apply Proposition~\ref{f16v18:partition}, and use its
coordinate projections $\Pi,\Pi_s$. There are four errors to pay.

\emph{Discarded native coordinates.} Positivity of $K_\gamma$ and V16's
component bound give $\Tr K_\gamma\le C\mathcal N\varepsilon^2$ and
$\Tr[(I-\Pi)K_\gamma]\le C|\mathcal A^c|\varepsilon^2$. The Schatten
Cauchy--Schwarz inequality, applied to the two factors of the positive square
root, gives
\begin{align}
 \|K_\gamma-\Pi K_\gamma\Pi\|_{S_1}
 &\le \Tr[(I-\Pi)K_\gamma]
      +2\sqrt{\Tr(\Pi K_\gamma)\Tr[(I-\Pi)K_\gamma]}\notag\\
 &\le C\mathcal N\varepsilon^2\sqrt{(L+\rho)/R}.
 \label{f16v18:corridor-cost}
\end{align}
If there are no discarded coordinates the error is zero. The last bound
uses $(L+\rho)/R\le1$. This is why a mere probability count is not multiplied
by an arbitrary spectral test function; positivity pays the actual covariance
cross block.

\emph{Exact off-block native return.} Apply
\eqref{f16v18:block-return} with screening order $N=J$. It gives
\begin{equation}
 \|\Pi K_\gamma\Pi-\sum_s\Pi_sK_\gamma\Pi_s\|_{S_1}
   \le C\mathcal N\bigl(\varepsilon^2e^{-2\vartheta L}
                                     +\varepsilon^{2J+2}\bigr).
 \label{f16v18:offblock-cost}
\end{equation}
The conditional covariance is not substituted for the unconditional block;
\eqref{f16v18:pinching-return} paid their distinction.

\emph{Finite-bank native expansion.} Put $k_s=|\mathcal A_s|\le CR^4$.
The order-$K$ entry error in \eqref{f16v18:primitive-inputs} has block operator
norm at most $C k_s\varepsilon^{K+1}$ and trace norm at most
$C k_s^2\varepsilon^{K+1}$. For the finite sum of coefficients with
$J<q\le K$, their uniform absolute row bounds give a trace error at most
$C k_s\varepsilon^{J+1}$. Since $\sum_s k_s\le C\mathcal N$,
\begin{equation}
 \left\|\sum_s\Pi_s(K_\gamma-K^{[J]})\Pi_s\right\|_{S_1}
     \le C\mathcal N
           \bigl(R^4\varepsilon^{K+1}+\varepsilon^{J+1}\bigr).
 \label{f16v18:bank-cost}
\end{equation}
This invokes V17 at a \emph{fixed} order $9J$, in the unrestricted ambient
law. It does not expand a conditional kernel on a fixed large exterior and
does not require a bound exponential in $R^4$.

\emph{Completion of the coefficient matrices.} Removing the discarded rows
and columns of $K_q$ changes a matrix of rank at most $2|\mathcal A^c|$;
its operator norm is at most twice the original norm. This costs at most
$C_q|\mathcal A^c|$ in trace norm. Between distinct selected banks all
anchors have distance at least $L$. The weighted row and column bounds for
$K_q$ therefore give operator norm at most $C_qe^{-\vartheta L}$ for the
remaining off-block part, and trace norm at most
$C_q\mathcal N e^{-\vartheta L}$. Sum the finitely many coefficients to obtain
\begin{equation}
 \left\|K^{[J]}-\sum_s\Pi_sK^{[J]}\Pi_s\right\|_{S_1}
      \le C\mathcal N\varepsilon^2
                     \bigl((L+\rho)/R+e^{-\vartheta L}\bigr).
 \label{f16v18:coefficient-completion}
\end{equation}
There is no positivity assumption on the higher $K_q$ in this argument.

Combining these four bounds gives the explicit finite-scale estimate
\begin{align}
 \frac{\|K_\gamma-K^{[J]}\|_{S_1}}{\mathcal N}
 \le C\bigl[&\varepsilon^2\sqrt{(L+\rho)/R}
       +\varepsilon^2 e^{-\vartheta L}
       +\varepsilon^{2J+2}\notag\\
       &+R^4\varepsilon^{K+1}+\varepsilon^{J+1}\bigr].
 \label{f16v18:four-costs}
\end{align}
For the choices \eqref{f16v18:trace-scales},
$e^{-\vartheta L}\le\varepsilon^{J+1}$,
$R^4\varepsilon^{9J+1}\le16\varepsilon^{J+1}$, and
\[
 \varepsilon^2\sqrt{(L+\rho)/R}
     \le C\varepsilon^{J+2}\sqrt{1+\log(\varepsilon^{-1})}
     \le C\varepsilon^{J+1}.
\]
This proves the primitive trace estimate. The geometric construction already
covers short directions and small tori; no nonexistent large block is assumed.
All thresholds are fixed-order thresholds, not functions of the ambient size.

Finally the exact synthesis and coefficient identity give
\[
 \Gamma_\gamma-\sum_{q=2}^J\varepsilon^qG_q
       =T(K_\gamma-K^{[J]})T^{\mathsf T}.
\]
Its trace norm is at most $\|T\|^2$ times the primitive trace norm. The
original Haar-contact identity gives
$\mathsf M_\gamma=-\operatorname{Off}_{>1}\Gamma_\gamma$ and
$M_q=-\operatorname{Off}_{>1}G_q$. That extraction has trace-norm bound four
for every duration. These two facts prove \eqref{f16v18:trace-expansion}.
Neither the measure nor the physical transfer has been changed by the proof
partition.
\end{proof}

\subsection{What the trace expansion establishes, and the remaining edge}

At its first order Theorem~\ref{f16v18:full-trace} gives
\begin{equation}
 \boxed{
 \frac{\|\gamma\Gamma_\gamma-G_2\|_{S_1}
              +\|\gamma\mathsf M_\gamma-M_2\|_{S_1}}{\mathcal N}
                           \le C_r\gamma^{-1/2}.}
 \label{f16v18:leading-trace}
\end{equation}
Thus the leading native coefficient is identified with an
$o(\gamma^{-1})$ normalized trace error on the \emph{entire} history, not
only entrywise or on a chosen small source bank. This replaces V16's
unidentified $O(\gamma^{-1})$ trace budget by a genuine fixed-order
asymptotic expansion in that norm. V17's stronger unnormalized bank estimates
remain separate and are used in the proof.

For every fixed $\delta>0$, the number of singular values of
$\gamma\Gamma_\gamma-G_2$ exceeding $\delta$ is at most
$C_r\mathcal N\gamma^{-1/2}/\delta$, and likewise for memory. This follows
by summing those singular values; it is a statement about the error matrix,
not an assertion that two noncommuting matrices share eigenvectors. A
vanishing proportion of exceptional directions can still include a single
large eigenvalue. Neither \eqref{f16v18:leading-trace} nor the analogous
higher-order formula bounds the unrestricted operator norm without its
volume factor.

The intermediate posterior result also has a more limited meaning than an
all-boundary conditional theorem. Large fast exteriors have been integrated
under the native marginal, not prohibited, and every message normalizer is
kept. But the bound contains an arbitrarily high \emph{fixed-order} coupling
floor at fixed $L$. No order-dependent constants have been controlled as
$N\to\infty$ with coupling fixed. Thus no assertion of strong mixing at all
distances, convergence of a series, or a unique infinite resummation follows.

The next unresolved question is an edge rather than an average: control the
remaining exceptional coherent directions, or prove a summable connected
native remainder at fixed sufficiently large coupling. The present block
argument cannot answer it by repeating the trace estimate. In
\eqref{f16v18:corridor-cost} it uses total covariance traces to pay discarded
rows, and in \eqref{f16v18:offblock-cost} it sums scalar message energies.
Those are the two locations where a directional, instead of trace, bound
would be needed. A source-resolved capacity/repair estimate or an independently
proved native functional inequality remains a distinct possible route.

V4's original full/core probability has not been redefined or given an
all-order boundary expansion. This appendix concerns the full native law
and its exact conditional messages. The prescribed bridge window also
remains: no integration over unrestricted physical bridges, root Gauss
restriction, physical-time calibration, infrared contact coercivity,
same-top full-transfer estimate, or interacting continuum construction is
supplied by a trace asymptotic. The auxiliary sampler gap is not used and
is not proved.

\subsection{Diagnostics, preservation, and claim scope}

The new diagnostic is
\begin{center}\small
\path{verify_ym_f16_v18_exterior_trace.py}.
\end{center}
It checks the weighted polynomial solver and its covariance path bounds.
Further tests cover the Stein identity on permitted coordinates, conditional
messages and posterior pinching on finite positive Markov laws, sparse
partitions including short directions, the four trace costs, and positive
matrix corridor and synthesis identities. All 819 exact assertions passed. Its report distinguishes exact algebraic
fixtures from analytic proofs. In particular a finite probability example is
not a native compact-integral enclosure, and no computation certifies all the
analytic constants or the inherited moment theorem.

The cumulative source changes only V17's title version and inserts this exact
appendix before the final document-ending command. Removing that insertion
and restoring the title recovers the predecessor byte for byte. The package
records the input hashes, preservation check, actual diagnostic runs,
compilation, and rendered-page inspection. The fourteen available V4--V17 diagnostics were rerun unchanged and passed.
Their individual reports and hashes are included. No inaccessible archive
checker or external formalization is described as verified. The new arguments explicitly depend on the identified V16--V17
analytic inputs, not on their finite diagnostic counts.

\begin{firewall}
V18 proves native polynomial identities with exterior tests. It gives
averaged screening estimates for conditional sources, with fixed-order
errors uniform in ambient volume and duration. These give arbitrary-order posterior trace accuracy
with a logarithmic collar. Separated finite space--time blocks then yield
fixed-order asymptotics of the full covariance and memory in trace norm per
cell, including an $o(1/\gamma)$ leading-coefficient error. No unrestricted
operator-norm remainder, fixed-coupling strong mixing or summable connected
remainder, uniform native sampler gap, physical transfer gap, or interacting
continuum Yang--Mills mass gap is established. All original native laws,
normalizers, and source definitions remain unchanged.
\end{firewall}
% END F16 V18 APPEND

% BEGIN F16 V19 APPEND -- 2026-09-17
% Insert immediately before the final document-ending command in V18.
% No inherited proof, normalization, or bibliography string is revised.

\clearpage
\section{V19: a positive shell return and one-sided coherent asymptotics}
\label{f16v19:context}

V18 identifies the full covariance in trace norm per cell. Its two losses at
an unrestricted spectral edge are explicit: discarded corridors are paid by
traces, and exterior messages are summed through their scalar squared norms.
We do not improve those same trace estimates. Instead, every primitive source
gets its own local ball, with no discarded sources. Subtracting the Gram of
its exact exterior conditional means leaves an \emph{exactly finite-range
matrix}. The subtracted Gram is positive; the finite-range matrix need not be.
This distinction gives a one-sided, unnormalized operator estimate at every
fixed order, and an explicit positive object carrying the remaining upper edge.

A second step resolves that object's change across scales through actual
annular innovations. Each scale has a positive finite-range Gram, although
cross-scale terms do not in general vanish. A dyadic synthesis estimate gives
a source-specific sufficient condition for the missing upper edge. The known
finite-order screening bounds pay transport to any prescribed polynomial
radius, but not an infinite sequence of radii at fixed coupling. We keep that
remaining hypothesis separate from the results proved here.

\subsection{Dependencies, non-duplication, and the literature distinction}

The supplied V18 appendix and audit were read in full, including the distinction
between primitive finite carriers and the full spatial Haar lift. The new
argument uses V16's native moments, V17's all-volume entry expansion, V18's
primitive exterior screening, and the exact finite-range specification of the
same law. It does not use V18's final corridor trace theorem to prove an
operator bound, nor assume a native functional inequality.

Two targeted semantic searches in the supplied f14/f15 and monograph sources
returned no matches; their mounted text inventories were then searched. The
monograph's shell--Doob covariance proposition already states the
martingale covariance formula for \emph{one common filtration}. F15 V56's
exact reference covariance decomposition already separates a retained-predictor
Gram from conditional variances for its reference law. These mechanisms are
credited, not claimed as new. Here each primitive has a \emph{different}
centered ball filtration. A common-filtration orthogonality formula cannot
be transferred to that collection without proof. An exact counterexample
below explains why. This is a targeted audit, not recertification of the
complete archive. The next scheduled companion checkpoint remains V20.

For context, Brydges--Guadagni--Mitter,
\href{https://arxiv.org/abs/math-ph/0303013}{\emph{Finite range Decomposition
of Gaussian Processes}}, construct positive finite-range decompositions for
specified Gaussian covariances. Bauerschmidt,
\href{https://arxiv.org/abs/1206.2212}{\emph{A simple method for finite range
decomposition of quadratic forms and Gaussian fields}}, gives a related
Green-form construction. Their primary abstracts and scope were checked;
neither complete theorem is imported. Our conditional means are native
non-Gaussian functions and their scales are not independent Gaussian fields.
The first page and introduction of Lohmann's
\href{https://arxiv.org/abs/1411.1107}{\emph{Single Scale Cluster Expansions}}
were also checked; no cluster-expansion hypotheses for this Wilson measure
are silently declared verified. The supplied memoranda's distinction between
power gain and locality cost, and between summable dependence and a physical
gap, remains binding.

\subsection{One native law and one finite shell for each primitive}

Keep the complete unmoving law $P=P_{\gamma,\mathbf y}$, the fixed bridge
window $\max_{i,e}|y_{i,e}|\le r$, and put
\[
 \varepsilon=\gamma^{-1/2},\qquad \mathcal N=B(n+1).
\]
Use V18's cell metric $d_*$ and a fixed integer $\rho\ge1$ larger than every
native interaction diameter and every primitive-source carrier radius. Let
$\mathcal J$ be the primitive scalar index set and $a(u)$ the anchor of $u$.
There are at most $30\mathcal N$ such indices and uniformly bounded index
multiplicity at each cell. Write
\begin{equation}
 F_u=e^{\rm prim}_u-\mathbb E_P e^{\rm prim}_u,\qquad
 K_\gamma=(\mathbb E_P F_uF_v)_{u,v\in\mathcal J},\qquad
 \Gamma_\gamma=T K_\gamma T^{\mathsf T},\quad T=(I,\mathsf P^{\mathsf T}).
 \label{f16v19:primitive}
\end{equation}
The exact synthesis $T$ has uniformly bounded operator norm and absolute row
sums. It is the original Haar-source synthesis, not a new random projection.
The original residuals and all their normalizers are unchanged.

For any integer $L\ge\rho$, define the cell ball and its actual word collar by
\[
 I_{u,L}=\{c:d_*(a(u),c)\le L\},\qquad
 C_{u,L}=\{c\notin I_{u,L}:c\text{ occurs in an interaction meeting }I_{u,L}\}.
\]
Then $C_{u,L}\subset\{L<d_*(a(u),c)\le L+\rho\}$. Balls are intersected
with the actual finite history; there is no artificial time continuation.
Let
\begin{equation}
 m_{u,L}=\mathbb E_P[F_u\mid\zeta_{I_{u,L}^c}],\qquad
 R_{\gamma,L}=(\mathbb E_P m_{u,L}m_{v,L})_{u,v},\qquad
 A_{\gamma,L}=K_\gamma-R_{\gamma,L}.
 \label{f16v19:shell-definitions}
\end{equation}
If the ball is the whole system, its conditional mean is zero. All matrices
are evaluated under the one full native probability. In particular the
conditional partition function is included in each $m_{u,L}$, and its actual
marginal is used in the Gram. Different balls can overlap arbitrarily.

\begin{theorem}[Exact finite-range subtraction with a positive return]
\label{f16v19:exact-local}
For every finite regulator and every $L\ge\rho$, $m_{u,L}$ is measurable
with respect to $\zeta_{C_{u,L}}$, and
\begin{equation}
 R_{\gamma,L}\succeq0,\qquad
 (A_{\gamma,L})_{uv}=0\quad
       \text{if }d_*(a(u),a(v))>2L+\rho.
 \label{f16v19:finite-range}
\end{equation}
Thus $K_\gamma=A_{\gamma,L}+R_{\gamma,L}$ exactly, with no discarded row or
column. This does not assert $A_{\gamma,L}\succeq0$.
For every fixed integer $N\ge1$, V18's analytic input gives
\begin{equation}
 s_\gamma(L):=\sup_u\|m_{u,L}\|_2
 \le C_{N,r}\bigl(\varepsilon e^{-\vartheta L}+\varepsilon^{N+1}\bigr),
 \label{f16v19:screening-input}
\end{equation}
for sufficiently large coupling independent of $B,n,L$. Consequently every
entry of $R_{\gamma,L}$ has absolute value at most $s_\gamma(L)^2$, and
$\Tr R_{\gamma,L}\le30\mathcal N s_\gamma(L)^2$.
\end{theorem}
\begin{proof}
The primitive carrier is inside $I_{u,L}$. Product coordinate reference
measure, including the fractional endpoint Haar powers, and the finite-range
native action imply that the conditional integral for $F_u$ given $I_{u,L}^c$
depends only on $C_{u,L}$. Every crossing term and its normalizer is retained.
Subtracting the deterministic full mean changes no measurability. This is the
same Markov-blanket identity used in V11--V18, now for overlapping balls.

Suppose the two anchors have distance greater than $2L+\rho$. The entire
carrier of $F_v$ is outside $I_{u,L}$, and $C_{u,L}$ is outside $I_{v,L}$.
Successive conditional expectations therefore give
\[
 \mathbb E_P F_uF_v=\mathbb E_P m_{u,L}F_v
                   =\mathbb E_P m_{u,L}m_{v,L}.
\]
No independence of overlapping balls or of their exterior marginals is used.
This proves the exact range statement. The return is a Gram, hence positive.
The primitive case of \eqref{f16v18:source-centered}, applied to this ball
and centered by the actual mean, proves \eqref{f16v19:screening-input}.
Cauchy--Schwarz and summation of the diagonal give the last assertions.
The proof is unchanged when a ball saturates the finite system.
\end{proof}

The matrix $A_{\gamma,L}$ is not defined by deleting interactions from $P$.
Its near entries still involve exact native integrals. Nor is it the Gram of
$F_u-m_{u,L}$ alone. If $d_u=F_u-m_{u,L}$, its correct expression is
\[
 (A_{\gamma,L})_{uv}=\mathbb E_P d_ud_v
              +\mathbb E_P d_um_{v,L}+\mathbb E_P m_{u,L}d_v.
\]
The cross terms are necessary because the conditional projections differ
with $u$. A finite-range matrix is not thereby a replacement transfer operator.

\subsection{The finite-range part has an unnormalized operator expansion}

Retain exactly the primitive coefficients $K_q$ from V18 and put
$K^{[J]}=\sum_{q=2}^J\varepsilon^qK_q$. Their exponentially weighted
absolute row and column bounds and V17's entry error are inherited as
\eqref{f16v18:primitive-inputs}. Decrease $\vartheta$ once so that these
coefficient bounds and \eqref{f16v19:screening-input} use the same exponent.

\begin{proposition}[A finite-scale operator estimate without corridors]
\label{f16v19:local-error}
For fixed integers $k\ge J\ge2$ and $N\ge1$, and $L\ge\rho$,
\begin{align}
 \|A_{\gamma,L}-K^{[J]}\|_{\rm op}
 \le C_{k,N,r}\bigl[&(1+L)^4
    (\varepsilon^{k+1}+\varepsilon^2e^{-2\vartheta L}
                                  +\varepsilon^{2N+2})\notag\\
    &+\varepsilon^{J+1}
       +\varepsilon^2 e^{-\vartheta(2L+\rho)}\bigr].
 \label{f16v19:local-error-bound}
\end{align}
The coefficient-sum term $\varepsilon^{J+1}$ may be omitted when $k=J$.
The coupling threshold and constants are independent of the ambient volume
and history length. There is no factor proportional to their source count.
\end{proposition}
\begin{proof}
Let $\mathsf B_L$ be the entry mask retaining pairs of anchor distance at
most $2L+\rho$. The exact range statement gives
\[
 A_{\gamma,L}-K^{[J]}
 =\mathsf B_L(K_\gamma-K^{[k]})
   +\mathsf B_L(K^{[k]}-K^{[J]})
   -\mathsf B_L R_{\gamma,L}-(I-\mathsf B_L)K^{[J]}.
\]
Here $I-\mathsf B_L$ denotes the complementary entry mask, not subtraction
of a source-space projection. We estimate every term by absolute rows and
columns, not by assuming that an arbitrary entry mask contracts operator norm.

Polynomial growth of the four-dimensional cell metric and bounded primitive
multiplicity imply at most $C(1+L)^4$ retained entries per row or column.
The entrywise remainder of order $k$ therefore costs
$C(1+L)^4\varepsilon^{k+1}$. The entry bound for $R_{\gamma,L}$ costs
$C(1+L)^4s_\gamma(L)^2$. The finitely many higher coefficients, when present,
have uniformly bounded absolute rows even after masking and cost
$C\varepsilon^{J+1}$. Finally the weighted coefficient rows give a tail cost
$C\varepsilon^2e^{-\vartheta(2L+\rho)}$. Schur's row-and-column bound and
\eqref{f16v19:screening-input} prove the assertion. Small systems have fewer
entries, so the same estimate applies without requiring a large ball to fit.
\end{proof}

\begin{theorem}[A positive native remainder and one-sided coherent asymptotics]
\label{f16v19:positive-remainder}
For every pair of fixed integers $J\ge2$ and $M\ge1$, there are a logarithmic radius
$L_\gamma=O_{J,M,r}(\log\gamma)$ and finite constants independent of $B,n$
such that, at sufficiently large coupling,
\begin{equation}
 \boxed{\quad
 \Gamma_\gamma=G^{[J]}+\mathsf R^{\rm sh}_{\gamma,L_\gamma}
                                      +\mathsf E_{\gamma,J,M},\qquad
 \mathsf R^{\rm sh}_{\gamma,L}:=T R_{\gamma,L}T^{\mathsf T}\succeq0,
 \quad}
 \label{f16v19:positive-split}
\end{equation}
where $G^{[J]}=\sum_{q=2}^J\varepsilon^qG_q$ uses the original coefficients,
and
\begin{equation}
 \|\mathsf E_{\gamma,J,M}\|_{\rm op}\le C_{J,M,r}\varepsilon^{J+1},\qquad
 \max_\alpha(\mathsf R^{\rm sh}_{\gamma,L_\gamma})_{\alpha\alpha}
       +\frac{\Tr\mathsf R^{\rm sh}_{\gamma,L_\gamma}}{\mathcal N}
       \le C_{J,M,r}\varepsilon^M.
 \label{f16v19:positive-budgets}
\end{equation}
In particular, on the \emph{entire} original source space,
\begin{equation}
 \boxed{\quad\Gamma_\gamma\succeq
           G^{[J]}-C_{J,r}\gamma^{-(J+1)/2}I.\quad}
 \label{f16v19:one-sided}
\end{equation}
Equivalently the negative spectral part of $\Gamma_\gamma-G^{[J]}$ has
operator norm at most $C_{J,r}\gamma^{-(J+1)/2}$. This is not an upper bound
on its positive part.
\end{theorem}
\begin{proof}
Choose $k=J+1$, $N=\max\{J+2,\lceil M/2\rceil\}$, and for
$0<\varepsilon<1/2$ take
\begin{equation}
 L_\gamma=\left\lceil\frac{N+2}{\vartheta}
                     \log(\varepsilon^{-1})\right\rceil+2\rho+1.
 \label{f16v19:log-radius}
\end{equation}
Then $e^{-\vartheta L_\gamma}\le\varepsilon^{N+2}$ and
$s_\gamma(L_\gamma)^2\le C_{N,r}\varepsilon^{2N+2}$.
In \eqref{f16v19:local-error-bound}, the first entry remainder is bounded by
$C(1+\log\varepsilon^{-1})^4\varepsilon^{J+2}$, which is at most
$C\varepsilon^{J+1}$. The other logarithm-weighted terms and the coefficient
tail are smaller after increasing fixed constants. Thus
$\|A_{\gamma,L_\gamma}-K^{[J]}\|_{\rm op}\le C\varepsilon^{J+1}$.

Apply the same deterministic synthesis $T$ to the exact primitive split.
V18's identity $G_q=TK_qT^{\mathsf T}$ identifies every coefficient.
Set $\mathsf E=T(A_{\gamma,L_\gamma}-K^{[J]})T^{\mathsf T}$; boundedness
of $T$ proves its norm estimate. Positivity and the trace ideal inequality
bound the returned trace by $C\mathcal N s_\gamma(L_\gamma)^2$.
For a canonical source, its returned variable is one row of $T$ applied to
$m_{u,L_\gamma}$. The uniform absolute row sum and Minkowski bound its
$L^2$ norm by $Cs_\gamma(L_\gamma)$, proving the diagonal estimate.
Since $2N+2\ge M$, these are \eqref{f16v19:positive-budgets}.

The error is symmetric. Its lower quadratic-form bound is
$-C\varepsilon^{J+1}I$, and the return is positive. This proves the
one-sided bound, taking one fixed $M$ to remove it from the notation there.
Spectral calculus for a finite symmetric matrix gives the equivalent
negative-part bound. There is no division by $\mathcal N$ in that operator
statement. No common eigenbasis for $\Gamma_\gamma,G^{[J]}$ is assumed.
\end{proof}

At the leading order this gives the new unrestricted-direction statement
\[
 \gamma\Gamma_\gamma\succeq G_2-C_r\gamma^{-1/2}I.
\]
It rules out a coherent \emph{downward} failure of the leading comparison.
A large upward failure can still live in the positive shell return. In fact
\begin{equation}
 \left|\|\Gamma_\gamma-G^{[J]}\|_{\rm op}
                 -\|\mathsf R^{\rm sh}_{\gamma,L_\gamma}\|_{\rm op}\right|
 \le C_{J,M,r}\varepsilon^{J+1}.
 \label{f16v19:edge-equivalence}
\end{equation}
This is the triangle inequality, not an assumption that the return is small.
The same decomposition for the original memory is
\begin{equation}
 \mathsf M_\gamma-M^{[J]}
   =-\operatorname{Off}_{>1}\mathsf R^{\rm sh}_{\gamma,L_\gamma}
                      +\mathsf E^{\rm mem},\qquad
 \|\mathsf E^{\rm mem}\|_{\rm op}\le4C_{J,M,r}\varepsilon^{J+1}.
 \label{f16v19:memory-split}
\end{equation}
The factor four is the inherited Fourier block-extraction estimate. The
extracted return need not be positive: the Loewner conclusion is for the
complete covariance, not for its time-off-diagonal mask. No contact term or
endpoint normalizer is changed.

\subsection{Two tempting positivity inferences are false}

\begin{proposition}[Different local filtrations do not supply a monotone covariance split]
\label{f16v19:counterexamples}
In a positive Gaussian Markov law, $K-R_L$ need not be positive even though
$R_L$ is a conditional-message Gram. Also $R_L-R_{L'}$ need not be positive
for $L<L'$, although every individual message variance decreases. These
statements concern general local conditional constructions, not counterexamples
to the native Yang--Mills bounds above.
\end{proposition}
\begin{proof}
For three centered Gaussian coordinates take precision
$H_{ii}=1$, $H_{ij}=3/4$ for $i\ne j$, and $C=H^{-1}$. For singleton
interiors, the vector of exterior means is $(I-H)\zeta$. Hence
\[
 C-(I-H)C(I-H)=2I-H.
\]
The precision is positive, with eigenvalues $5/2,1/4,1/4$. The displayed
local difference has eigenvalue $-1/2$ on $(1,1,1)$.

For failure of matrix monotonicity, instead use the path precision
\[
 H=\begin{pmatrix}1&-1/3&0\\-1/3&1&-1/3\\0&-1/3&1\end{pmatrix},\qquad
 C=\frac17\begin{pmatrix}8&3&1\\3&9&3\\1&3&8\end{pmatrix}.
\]
Singleton interiors give
\[
 R_0=\begin{pmatrix}1/7&2/21&1/7\\2/21&2/7&2/21\\1/7&2/21&1/7\end{pmatrix}.
\]
Radius-one interiors on this path give exterior means
$(\zeta_3/8,0,\zeta_1/8)$ and therefore
\[
 R_1=\begin{pmatrix}1/56&0&1/448\\0&0&0\\1/448&0&1/56\end{pmatrix}.
\]
All diagonal differences are nonnegative, but for $v=(1,0,-1)^{\mathsf T}$,
$v^{\mathsf T}(R_0-R_1)v=-1/32$. These are exact rational Gaussian
computations. The radius-zero convention is used only in these examples;
the native construction above continues to use $L\ge\rho$.
\end{proof}

Thus increasing $L$ is not an unproved Loewner-monotone renormalization flow.
Nor can one subtract successive $R_L$ as positive covariances and then sum
those differences. The appropriate positive objects are the Grams of actual
\emph{scalar} reverse-martingale increments, with cross-scale terms retained.

\subsection{Annular innovations: positivity at each scale, not between scales}

Introduce the actual synthesis operator into the one common probability space,
\[
 \mathcal M_L:\ell^2(\mathcal J)\longrightarrow L^2(P),\qquad
 \mathcal M_Lv=\sum_u v_u m_{u,L},\qquad R_{\gamma,L}=\mathcal M_L^*\mathcal M_L.
\]
Let $L_0\ge\rho$ be an integer and $L_j=2^jL_0$. Define
\[
 d_{u,j}=m_{u,L_j}-m_{u,L_{j+1}},\qquad
 b_j=\sup_u\|d_{u,j}\|_2,\qquad
 \mathcal D_j v=\sum_uv_ud_{u,j}.
\]
At every finite regulator these sequences are eventually zero, because all
balls eventually contain the whole system.

\begin{proposition}[Exact annular orthogonality and finite-range Grams]
\label{f16v19:annuli}
For every anchor and scale,
\begin{equation}
 \mathbb E_P[d_{u,j}\mid\zeta_{I_{u,L_{j+1}}^c}]=0,\qquad
 \|d_{u,j}\|_2^2=\|m_{u,L_j}\|_2^2-\|m_{u,L_{j+1}}\|_2^2.
 \label{f16v19:scalar-drop}
\end{equation}
The variable $d_{u,j}$ is measurable in the union of its two actual word
collars, contained in the ball of radius $L_{j+1}+\rho$. Its Gram is positive
and has range at most $2L_{j+1}+\rho$. In particular,
\begin{equation}
 \|\mathcal D_j\|_{\ell^2\to L^2(P)}
                  \le C(1+L_{j+1})^2 b_j.
 \label{f16v19:annular-norm}
\end{equation}
The exponent two is half the four-dimensional volume-growth exponent.
\end{proposition}
\begin{proof}
For a \emph{fixed} $u$, the sigma algebras of exteriors decrease with $L$.
The tower property gives
$\mathbb E[m_{u,L_j}\mid I_{u,L_{j+1}}^c]=m_{u,L_{j+1}}$.
Orthogonality of one projection difference gives both identities in
\eqref{f16v19:scalar-drop}. Markov-blanket measurability gives the collar
support. If the two anchors have distance greater than $2L_{j+1}+\rho$,
the entire support of $d_{v,j}$ is outside $I_{u,L_{j+1}}$. Conditioning
therefore makes their inner product zero. No independence assertion is needed.
The Gram is positive by definition, each entry has absolute value at most
$b_j^2$, and there are at most $C(1+L_{j+1})^4$ nonzero entries per row
and column. Schur's estimate bounds its operator norm by that count times
$b_j^2$. Taking a square root proves \eqref{f16v19:annular-norm}.
This calculation does not assert orthogonality for different source anchors
at different scales.
\end{proof}

\begin{theorem}[A source-specific dyadic criterion for the remaining coherent return]
\label{f16v19:dyadic}
At every finite regulator,
\begin{equation}
 \boxed{\quad
 \|R_{\gamma,L_0}\|_{\rm op}^{1/2}
 \le C\sum_{j\ge0}(1+L_{j+1})^2 b_j
 \le C\sum_{j\ge0}(1+L_{j+1})^2s_\gamma(L_j).
 \quad}
 \label{f16v19:dyadic-bound}
\end{equation}
For the full-source return the same right side is multiplied by $\|T\|$.
The sums are finite at each finite regulator. A uniform bound on the infinite
majorant would make the conclusion uniform in volume and duration.

In particular, the \emph{additional} source-specific estimate
\begin{equation}
 s_\gamma(L)\le A_\gamma(1+L)^{-2-\delta}
       \quad(L\ge L_0),\qquad \delta>0,
 \label{f16v19:screening-target}
\end{equation}
would imply
\begin{equation}
 \|\mathsf R^{\rm sh}_{\gamma,L_0}\|_{\rm op}
       \le C_\delta A_\gamma^2(1+L_0)^{-2\delta}.
 \label{f16v19:power-consequence}
\end{equation}
No assertion that \eqref{f16v19:screening-target} has been proved for the
native law is part of this theorem.
\end{theorem}
\begin{proof}
At a fixed finite system, the centered message at every sufficiently large
radius is zero. Hence $\mathcal M_{L_0}=\sum_j\mathcal D_j$ exactly as a
finite sum of operators into $L^2(P)$. Apply the triangle inequality in that
operator space and \eqref{f16v19:annular-norm}. The scalar variance-drop
identity gives $b_j\le s_\gamma(L_j)$, proving the second inequality.
No unproved infinite-volume tail triviality is used to remove the final term:
the telescoping is completed on each finite regulator before taking a uniform
upper bound. A use directly in infinite volume would have to retain its tail
projection unless separately shown zero.

For \eqref{f16v19:screening-target}, the summands are at most
$C A_\gamma(1+L_0)^{-\delta}2^{-j\delta}$. Sum the geometric series,
square, and use the bounded source synthesis. Cross-scale correlations are
paid by this triangle inequality, rather than set to zero.
\end{proof}

The target here is weaker in kind than a Poincar\'e inequality for every
native $L^2$ function: it involves only the two specified primitive error
banks and their own exterior conditional means. It is also an averaged
native statement, not a worst-exterior bound. If it held with
$A_\gamma\le C_r\varepsilon$ at all radii beyond a logarithmic $L_0$, the
positive split would give
\begin{equation}
 \|\Gamma_\gamma-\gamma^{-1}G_2\|_{\rm op}
 \le C_r\gamma^{-3/2}
       +C_{r,\delta}\gamma^{-1}(1+\log\gamma)^{-2\delta}
                  =o(\gamma^{-1}).
 \label{f16v19:conditional-edge}
\end{equation}
The memory operator would obey the same conclusion up to the factor four.
This implication does not supply exponential temporal row decay or a sampler
gap. A critical-power bound with
$L^{-2}(\log(2+L))^{-1-\delta}$ would also make the dyadic majorant summable,
with its corresponding slower tail. The ordinary volume-growth exponent
four has not been replaced by a claim of two-dimensional geometry.

\subsection{What is already proved across polynomially many radii}

Although \eqref{f16v19:screening-target} is unproved, the established
finite-order screening does give a directional statement across a long but
finite range of scales. This uses the annular operator bound, not a trace.

\begin{theorem}[Coherent transport from logarithmic to polynomial shells]
\label{f16v19:transport}
For fixed $p>0$ and $q>0$ there are a logarithmic integer
$L_\gamma=O_{p,q,r}(\log\gamma)$, a radius
$\gamma^p\le\mathcal R_\gamma<2\gamma^p$, and fixed constants such that
\begin{equation}
 \boxed{\quad
 \|\mathcal M_{L_\gamma}-\mathcal M_{\mathcal R_\gamma}\|_{\ell^2\to L^2(P)}
                      \le C_{p,q,r}\varepsilon^q.
 \quad}
 \label{f16v19:transport-bound}
\end{equation}
The larger radius can be chosen as a dyadic multiple of $L_\gamma$.
The estimate holds at every ambient volume and duration, for sufficiently
large coupling depending on the fixed exponents, not on the system size.
It remains valid after composition with $T^{\mathsf T}$.
\end{theorem}
\begin{proof}
Choose $N\ge\max\{1,\lceil4p+q\rceil\}$ in
\eqref{f16v19:screening-input}, and
\[
 L_\gamma=\left\lceil\frac{q+3}{\vartheta}
                        \log(\varepsilon^{-1})\right\rceil+2\rho+1.
\]
For sufficiently large coupling $L_\gamma<\gamma^p$. Choose the first dyadic
multiple at least $\gamma^p$, and telescope the annular innovations only up
to that radius. Their total operator norm is bounded by
\begin{align*}
 C_{N,r}\sum_{L_j<\mathcal R_\gamma}(1+2L_j)^2
                 (\varepsilon e^{-\vartheta L_j}+\varepsilon^{N+1})
 &\le C_{N,r}\bigl[
       \varepsilon(1+L_\gamma)^2e^{-\vartheta L_\gamma}
            +\mathcal R_\gamma^2\varepsilon^{N+1}\bigr]\\
 &\le C_{p,q,r}\varepsilon^q.
\end{align*}
The first sum is bounded by the corresponding integer exponential tail;
its logarithmic factor is absorbed by the extra power in the choice of
$L_\gamma$. The second is a finite geometric sum of squared dyadic radii.
Since $\mathcal R_\gamma^2\le4\varepsilon^{-4p}$, its power is at least
$N+1-4p\ge q$. Saturation of balls only makes later terms vanish. No factor
counting source coordinates or time slices occurs. Composing with the
bounded deterministic synthesis proves the last assertion.
\end{proof}

The theorem is an amplitude-level estimate. One cannot simply square it to
conclude $\|R_{\gamma,L_\gamma}-R_{\gamma,\mathcal R_\gamma}\|_{\rm op}
=O(\varepsilon^{2q})$: the cross term with a potentially large remote message
must be retained. The following relative enclosure pays that term.

\begin{theorem}[A relative positive-return enclosure at a remote shell]
\label{f16v19:remote-enclosure}
Fix $J\ge2$, $p>0$, and $h>0$. For suitable radii as in
Theorem~\ref{f16v19:transport}, write
$\tau=\varepsilon^h$ and
$\mathsf R^{\rm far}=T R_{\gamma,\mathcal R_\gamma}T^{\mathsf T}\succeq0$.
For all admitted geometries and sufficiently large fixed-order coupling,
\begin{equation}
 \boxed{
 \begin{aligned}
 G^{[J]}+(1-\tau)\mathsf R^{\rm far}-C\varepsilon^{J+1}I
   &\preceq\Gamma_\gamma\\
   &\preceq G^{[J]}+(1+\tau)\mathsf R^{\rm far}
                                      +C\varepsilon^{J+1}I .
 \end{aligned}}
 \label{f16v19:remote-order}
\end{equation}
The same underlying native functions define $\mathsf R^{\rm far}$; no
probability law is replaced by a reference shell law. The matrix is not
asserted small in operator norm.
\end{theorem}
\begin{proof}
Take $q\ge(J+1+h)/2$ in the transport theorem and enlarge its logarithmic
radius and fixed expansion order to also satisfy the local error estimate
of Theorem~\ref{f16v19:positive-remainder}. This is allowed: finitely many
fixed requirements are imposed, and the same estimates remain valid at a
larger logarithmic radius. Let
$\mathcal U=(\mathcal M_{L_\gamma}-\mathcal M_{\mathcal R_\gamma})T^{\mathsf T}$.
Its norm is at most $C\varepsilon^q$. For each source vector $v$, put
$X=\mathcal M_{\mathcal R_\gamma}T^{\mathsf T}v$ and $Y=\mathcal Uv$.
For $0<\tau<1$, the two elementary Hilbert-space inequalities are
\[
 (1-\tau)\|X\|^2-\tau^{-1}\|Y\|^2
 \le\|X+Y\|^2
 \le(1+\tau)\|X\|^2+(1+\tau^{-1})\|Y\|^2.
\]
Thus the two shell Grams are relatively enclosed, with additive error at most
$C\varepsilon^{2q-h}\|v\|^2$. Insert the positive split and its local
operator error. Since $2q-h\ge J+1$, these combine into
\eqref{f16v19:remote-order}. This explicitly pays the cross-scale terms
without bounding the norm of $X$ or silently changing the clock or measure.
\end{proof}

For fixed $J,p,h$, every unresolved coherent excess is therefore represented,
up to the displayed small relative and additive errors, by messages on a
shell at distance of order $\gamma^p$. If that shell already lies beyond the
finite system, the return is zero; this is consistent with V17's finite-bank
conclusion, not a new thermodynamic assumption. At unrestricted volume the
remote shell may still carry a nonzero, highly correlated vector of messages.

\subsection{What has advanced and the precise next target}

The new unconditional result is a \emph{one-sided coherent} asymptotic at
unrestricted volume, and an exact positive-return decomposition with a small
local operator error. It is not merely another trace budget. There are no
sparse discarded corridors and no omitted primitive sources. The returned
Gram has arbitrary fixed-order small diagonal and trace density, but its
largest eigenvalue is retained, not bounded by those densities.

The annular argument also gives the proved transport estimate through every
prescribed polynomial range of radii. To reach all scales at a fixed coupling,
one would need a summable bound on
\[
 \sum_{j\ge0}(1+2^{j+1}L_0)^2
 \sup_u\left\|m_{u,2^jL_0}-m_{u,2^{j+1}L_0}\right\|_2.
\]
This is a concrete source-resolved target for the actual compact law.
An averaged scalar decay faster than radius to the power minus two is a
sufficient condition, not a proved consequence of V18. Substituting its
fixed-order floor into this infinite sum would diverge. Nor can one choose
an order increasing with the scale while ignoring its constants and entry
thresholds. The counterexample to matrix monotonicity excludes another
possible shortcut.

The local repair results remain possible tools for this shell problem. The
new formulation does not require a native gap for every test function as an
input, but it does not prove that such a gap exists. A successful coherent
covariance bound would still be a statement about the fixed-bridge fast law,
not a physical Hamiltonian. Bridge-window coverage, root Gauss restriction,
infrared contact coercivity, the same-top physical transfer comparison,
physical-time calibration and an interacting continuum construction retain
their separate obligations. The positive shell return is not identified with
V12's one-common-exterior cylinder return; they are different exact matrices
under the same probability.

\subsection{Diagnostics, preservation, and scope}

The new diagnostic is
\begin{center}\small
\path{verify_ym_f16_v19_positive_shell.py}.
\end{center}
It checks exact shell measurability in an inhomogeneous finite Markov law,
finite-range covariance subtraction, the annular tower and variance-drop
identities, positive source synthesis, both rational Gaussian counterexamples,
finite-geometry support, and the relative-order and scale bookkeeping. Its
report distinguishes these finite fixtures from analytic proofs and from
native compact integral enclosures. The majority of assertions verify the
same collar value across individual configurations of the finite fixture;
the count is not a measure of proof completeness. The new diagnostic passed
3,948 exact assertions, including 3,530 collar-value comparisons. All fifteen
available V4--V18 diagnostics were rerun unchanged and passed; their actual
outputs and hashes are retained in the package.

The cumulative V19 changes only the V18 title version and appends this exact
section before the final document-ending command. The package includes a
byte-for-byte predecessor recovery check, input hashes, actual diagnostic
reruns, build reports and rendered-page review. No inherited proof is silently
corrected or independently recertified. The new analytic deductions explicitly
use the identified V16--V18 inputs, not their diagnostic assertion counts.
The next scheduled companion checkpoint remains V20.

\begin{firewall}
V19 proves an exact finite-range primitive covariance difference, a positive
native shell-return decomposition, and an unrestricted-volume lower Loewner
asymptotic with a genuine operator-norm error. It proves coherent transport
to polynomially remote shells and a relative positive-return enclosure.
The dyadic source-screening criterion for the remaining upper edge is proved
as an implication, not verified for the native law at all radii. No complete
two-sided operator expansion, summable native remainder, uniform sampler gap,
physical transfer gap or interacting continuum mass gap is established.
All original measures, normalizers, source definitions and physical caveats
remain unchanged.
\end{firewall}
% END F16 V19 APPEND


% BEGIN F16 V20 APPEND -- 2026-09-17
% Insert immediately before the final document-ending command in V19.
% No inherited mathematical statement or bibliography string is changed.

\clearpage
\section{V20: uniform boundary preparation and buffered contraction}
\label{f16v20:context}

V19 leaves a positive Gram of remote conditional messages at the unknown upper
edge. Increasing a fixed perturbative order again would not control its
infinite sequence of radii. We go upstream to a different conditional input:
\emph{prepare a deep interior under every compact exterior}, then compare its
entire marginal law, not just polynomial observables. The buffer is integrated,
not frozen. Its exterior may be at angular distance of order one from the
identity, including the cancellation configurations of V14.

The result is a volume-independent total-variation estimate for a core of
$k$ cells, with its \emph{linear} cost in $k$ explicit. It implies a conditional
maximal-correlation bound for every square-integrable core observable and every
incoming exterior law. This is stronger in test class and conditioning scope
than V18's averaged polynomial screening, but has only a first-order coupling
error. It does not give the missing all-radius bound for V19's whole source
bank. The separate roles of output size, buffer depth, and update clock are
kept explicit.

\subsection{The V20 checkpoint and the new analytic input}

The supplied V19 appendix and audit were read in full. At this scheduled
checkpoint the available f14 V100, f15 V100, Grand Tensor Monograph V076,
and current lineage were inventoried and targeted passages checked. Two scoped
semantic searches were empty; exact mounted-text searches and identified
source ranges were then used. This is not recertification of the entire
archive. In particular, f15 V22 already compares full compact core kernels in
total variation, with a confidence-event error, and constructs a separated
coupling. F14 V65 compares normalized nested chart laws. The monograph's
strong-electric Dobrushin theorem concerns a different regime and transfer.
None of these hypotheses is silently transferred to the present law.

The new preparation estimate keeps the terminal boundary term in V16's ball
recursion, rather than sending the radius to infinity after averaging the
whole native law. A second change replaces V17's polynomial right inverse
by a Gaussian divergence solver for \emph{bounded tests}. A common cutoff on
the observed core permits the proof to control the complete divergence,
without assuming a bounded Hessian of a Stein solution for arbitrary bounded
tests. Coordinate-local cutoffs are still used on the other differentiated
groups. All Haar powers and all cutoff terms are retained.

For attribution, the Gaussian Ornstein--Uhlenbeck/Stein representation is
classical; the Mehler formula and Lemma 3.3 in Nourdin--Peccati--R\'eveillac,
\href{https://arxiv.org/abs/0804.1889}{\emph{Multivariate normal approximation
using Stein's method and Malliavin calculus}}, were checked. That paper's
Wasserstein theorem is not quoted as our total-variation theorem. The
bounded-test gradient estimate actually used is proved below. Polyanskiy--Wu,
\href{https://arxiv.org/abs/1508.06025}{\emph{Strong data-processing inequalities
for channels and Bayesian networks}}, Section 2.1, Theorem 1 and (7)--(9),
records the domination of chi-square contraction by the Dobrushin coefficient
and its relation to maximal correlation on general spaces. We give a direct
Hilbert-space proof of this consequence with our precise kernels.

The two supplied import memoranda separate power gain from locality cost and
require a genuine dependence estimate before global coercivity. This appendix
provides such an estimate for a \emph{buffered marginal channel}, not for the
full collection of overlapping updates. Only supplied companions were
available at the checkpoint; other programs named in older ledgers were not
recertified. The next scheduled checkpoint is V25.

\subsection{The conditional law and finite-distance boundary preparation}

Use V16's unmoving cells, $e_c=|\zeta_c|^2$, graph $\mathcal G$, and constants
$a_0=(6\pi^2)^{-1}$, $t_0=a_0/4$. Fix
\[
 P=P_{\gamma,\mathbf y},\qquad \max_{i,e}|y_{i,e}|\le r,
 \qquad \gamma\ge1+r^2.
\]
At most twelve principal three-dimensional groups occur in a cell, so on the
complete compact coordinates, including every fixed exterior value,
\begin{equation}
                  0\le e_c\le E_*\gamma,\qquad E_*=12\pi^2.
 \label{f16v20:compact-cap}
\end{equation}
The graph joins cells sharing a non-anchor primitive interaction. It has the
polynomial ball bound of V16. Write $d_{\mathcal G}$ for its distance. Different
components do not interact. This metric is used in this appendix; every graph
step has bounded range in V18's ambient space--time metric.

Let $\Lambda$ be a nonempty set of free cells, and freeze all cells outside it
at $q$. Denote the explicit complete conditional law by
\[
 P_\Lambda^q(d\zeta_\Lambda)
  =Z_\Lambda(q)^{-1}e^{-U_\Lambda(\zeta_\Lambda;q)}
                                  \prod_{c\in\Lambda}d\nu_{\gamma,c}.
\]
Here $U_\Lambda$ contains every anchor on $\Lambda$ and every primitive
interaction meeting $\Lambda$, once each, with the original endpoints.
The factors $\nu_{\gamma,c}$ are V16's unnormalized complete chart factors.
There is no imposed fast-field window on $q$. We use this positive-density
version of the conditional integral at all principal exterior values; at
Haar-null chart boundaries this is an explicitly specified extension, not a
claim of uniqueness of a regular conditional version.

Define $d_c=d_{\mathcal G}(c,\Lambda^c)$. Set $d_c=\infty$ if the entire
component of $c$ is free, and $e^{-\kappa\infty}=0$. Tilts $s$ on the free cells
have $|s_c|\le t_0$, exactly as in V16; their law is denoted $P_{\Lambda,s}^q$.

\begin{theorem}[Conditional preparation with an exponentially attenuated boundary term]
\label{f16v20:preparation}
There are $B_r<\infty$ and $\kappa_r>0$, independent of $\Lambda,q,B,n,\gamma$
and the source array, such that
\begin{equation}
 \mathbb E_{P_{\Lambda,s}^q}e_c
             \le M_r+B_r\gamma e^{-\kappa_r d_c},\qquad c\in\Lambda.
 \label{f16v20:prepared-mean}
\end{equation}
Here $M_r$ may be V16's constant enlarged once. For every array
$0\le t_c\le t_0$ on the free cells,
\begin{equation}
 \log\mathbb E_{P_\Lambda^q}\exp\!\left(\sum_c t_c e_c\right)
       \le\sum_c t_c\bigl(M_r+B_r\gamma e^{-\kappa_r d_c}\bigr).
 \label{f16v20:prepared-mgf}
\end{equation}
In particular, at depths $d_c\ge\kappa_r^{-1}\log\gamma$, every fixed local
moment has a uniform bound, now for \emph{all} compact exterior values.
There also exist a fixed integer $D_r$ and $c_0>0$ such that every original
group $g$ in a cell of depth at least $D_r$ satisfies
\begin{equation}
 P_\Lambda^q\{|\zeta_g|\ge\sqrt\gamma/2\}
                         \le e^{t_0M_r-c_0\gamma}.
 \label{f16v20:prepared-cut-tail}
\end{equation}
The latter constant-depth bound concerns an angular cut event, not a uniform
thermal moment at constant depth.
\end{theorem}
\begin{proof}
Keep the proof of V16's conditional normalizer estimate, including all crossing
terms and its product unit-ball lower bound. It applies to any $S\subseteq
\Lambda$ conditional on all its other variables. Average only in
$P_{\Lambda,s}^q$. With
\[
 u_c=\mathbb E_{P_{\Lambda,s}^q}e_c\ (c\in\Lambda),\qquad
 u_c=e_c(q)\ (c\notin\Lambda),
\]
its conclusion is still
\begin{equation}
             \sum_{c\in S}u_c\le a_r|S|+b_r\sum_{d\in\partial S}u_d.
 \label{f16v20:conditional-set}
\end{equation}
This is not a new estimate on arbitrary fixed boundary values without a debit:
the exterior entries on the right remain their actual deterministic energies.
All $u_c$ also satisfy \eqref{f16v20:compact-cap}.

For $c\in\Lambda$, write $F_k=\sum_{d\in B_k(c)}u_d$ and
$\rho_r=b_r/(1+b_r)$. For $0\le k<d_c$, the ball $B_k(c)$ is entirely free,
so V16's same recurrence gives
\[
 F_k\le\rho_r F_{k+1}+\frac{a_r}{1+b_r}|B_k(c)|.
\]
If $d_c$ is finite, stop at $d_c$, not at infinity. The terminal term is
bounded by
\[
 \rho_r^{d_c}F_{d_c}
     \le E_*C_{\rm gr}\gamma(d_c+1)^4\rho_r^{d_c}.
\]
The sum of the preceding terms is at most V16's $M_r$. Choose
$\kappa_r=-\tfrac12\log\rho_r$ and
\[
 B_r\ge E_*C_{\rm gr}\sup_{j\ge0}(j+1)^4\rho_r^{j/2}<\infty.
\]
This proves \eqref{f16v20:prepared-mean}. If $d_c=\infty$, the recurrence
can be completed on its finite component and the terminal term tends to zero,
as in V16. No volume limit is exchanged with this iteration.

The proof is uniform over all signed free sources. Integrate the exact
finite-regulator derivatives of $\log Z_\Lambda(q;s)$ along the segment
from zero to $t$. This proves \eqref{f16v20:prepared-mgf}; no derivative of a
pressure approximation is taken. The local moment assertion follows by the
same exponential-moment inequality as V16. Finally choose $D_r$ such that
$B_r e^{-\kappa_rD_r}\le1/8$. A cut group has $e_c\ge\gamma/4$.
Exponential Markov at $t_0$ in that cell gives
\eqref{f16v20:prepared-cut-tail} with $c_0=t_0/8$.
\end{proof}

The preparation is uniform over the \emph{size} of the free set as well as
the exterior. A boundary cell may still have mean of order $\gamma$.
This is consistent with V14's unpinned pair, which had no free buffer between
the observed variables and its cancelling exterior. No global convexity of
the native compact potential has been asserted.

\subsection{A bounded-test Gaussian solver with a local core}

Let $\mathcal A\subseteq\Lambda$ be a nonempty set of $k$ complete cells, and
let $A$ denote all its real coordinates, with $d_A\le36k$. Put
\[
 L=d_{\mathcal G}(\mathcal A,\Lambda^c),\qquad
 Q=N(m,C),\quad C=H^{-1},\qquad
 Q_A=N(m_A,C_A),\quad C_A=C_{AA}.
\]
The Gaussian $Q$ is the original full-history reference. In particular
$C_A$ is a \emph{marginal} covariance, not $(H_{AA})^{-1}$ or the covariance
of a newly chosen conditional reference. It is independent of $q$.
All scalar indices below include their fixed colour multiplicities.

The inherited bounds on $H$ imply, at some fixed $\nu>0$,
\begin{equation}
 h_+^{-1}I\preceq C_A\preceq h_-^{-1}I,\qquad
 \sup_i\sum_j|C_{ij}|e^{\nu d_{\mathcal G}(c(i),c(j))}\le K_C,
 \qquad \sup_i|m_i|\le M_r^{\rm G}.
 \label{f16v20:Gaussian-local}
\end{equation}
For completeness, $H$ has range one in $\mathcal G$. The Neumann series
$H^{-1}=h_+^{-1}\sum_{l\ge0}(I-H/h_+)^l$ has geometric operator decay and
no entry before the corresponding graph distance. Polynomial ball growth
makes its exponentially weighted rows summable at a smaller exponent.
Across different components the entries vanish. Thus using the native graph
rather than the ambient metric introduces no unproved inverse bound.

For a smooth bounded $\varphi$ on $\mathbb R^{d_A}$ with bounded derivatives,
define
\begin{equation}
 u_\varphi(x)=\int_0^\infty
 \left[\mathbb E\varphi\bigl(m_A+e^{-t}(x-m_A)
                   +\sqrt{1-e^{-2t}}\,C_A^{1/2}G\bigr)-Q_A\varphi\right]dt,
 \quad G\sim N(0,I).
 \label{f16v20:OU-solver}
\end{equation}

\begin{proposition}[Bounded-test divergence without a Hessian test norm]
\label{f16v20:bounded-solver}
Set $b(x)=\nabla u_\varphi(x)$ and, on the entire coordinate space, set
$g_j(\zeta)=C_{jA}b(\zeta_A)$. Then
\begin{equation}
 \delta_Qg=\varphi(\zeta_A)-Q_A\varphi,\qquad
 \|b\|_\infty\le\sqrt{\pi/2}\,\|C_A^{-1/2}\|\,\|\varphi\|_\infty.
 \label{f16v20:bounded-solver-identity}
\end{equation}
Here $\delta_Q$ is the same full Gaussian divergence as in V17.
The components $g_j$ depend only on the core, even for $j\notin A$.
If $a_j=\|C_{jA}\|_2$, then
\begin{equation}
 \|g_j\|_\infty\le C a_j\|\varphi\|_\infty,\qquad
 \sum_j a_j\le Ck,\qquad \|g_A\|_\infty\le C\|\varphi\|_\infty.
 \label{f16v20:solver-size}
\end{equation}
None of these bounds uses a derivative norm of $\varphi$.
\end{proposition}
\begin{proof}
The generator of the Mehler semigroup in \eqref{f16v20:OU-solver} is
$\mathcal L_A=\operatorname{Tr}(C_A\nabla^2)-(x-m_A)\cdot\nabla$.
Integration of its time derivative gives
$-\mathcal L_Au_\varphi=\varphi-Q_A\varphi$. The integral converges at each
$x$; for large $t$, Gaussian translation and covariance comparison, or a
bounded derivative of the smooth test, gives an integrable exponential bound.
For any unit vector $h$, Gaussian integration by parts inside the Mehler
integral gives
\[
 |\partial_h P_t^A\varphi(x)|
 \le\sqrt{2/\pi}\,\frac{e^{-t}}{\sqrt{1-e^{-2t}}}
             |C_A^{-1/2}h|\,\|\varphi\|_\infty.
\]
The integral of the time factor is $\pi/2$, by the substitution $s=e^{-t}$.
This proves the gradient bound, uniformly in $x$ and in $d_A$. Initially the
smooth-test assumption justifies differentiations; its derivative constants
do not occur in this estimate.

Since $HC=I$, the drift term in $\delta_Qg$ is
$(\zeta_A-m_A)\cdot b$. Only indices $j\in A$ differentiate $b$, and their
sum is $\operatorname{Tr}(C_A\nabla^2u_\varphi)$. This proves the divergence
identity. The absolute column bound in \eqref{f16v20:Gaussian-local} gives
$\sum_j a_j\le\sum_{i\in A,j}|C_{ji}|\le K_Cd_A$. Finally
$\|g_A\|\le\|C_A\|\,\|b\|$ gives the last estimate.
\end{proof}

A generic multivariate bounded-test Stein solution need not have the uniform
Hessian bound needed for a Wasserstein-to-TV substitution. We do not use such
a bound. The next proof keeps all core divergence terms together; it never
estimates their individual second derivatives on a cutoff complement.

\subsection{The conditional drift error and the boundary reaching its core}

Write $V_\gamma=U_\gamma-\log J^f_\gamma$ on the open full chart, and
$\xi=\zeta-m$. Only derivatives in free coordinates are used; the exterior
is substituted after forming these derivatives. Use the same smooth group
cutoff $\chi_j=\chi(\zeta_{g(j)}/\sqrt\gamma)$ as V17: it equals one within
angular radius $1/2$ and vanishes beyond radius one. For $j\in\Lambda$ define
\[
                    D_j=\partial_jV_\gamma-(H\xi)_j.
\]
All expectations in this subsection are under $P_\Lambda^q$.

\begin{proposition}[An actual conditional drift estimate]
\label{f16v20:drift}
For every free coordinate $j$, any multiplier $0\le\eta_j\le1$ supported
where $|\zeta_{g(j)}|\le\sqrt\gamma$ obeys
\begin{equation}
 \mathbb E[\eta_j|D_j|]
                \le C_r\bigl(\gamma^{-1/2}
                    +\gamma^{1/2}e^{-\kappa_r d_{c(j)}}\bigr),
 \label{f16v20:drift-bound}
\end{equation}
after enlarging constants and, if necessary, decreasing $\kappa_r$.
There is $\mu_r>0$ such that the solver weights satisfy
\begin{align}
 \sum_{j\in\Lambda}a_j e^{-\kappa_r d_{c(j)}}
                              &\le C_r k e^{-\mu_r L},\notag\\
 \sum_{j\notin\Lambda}a_j
   +\sum_{j\in\Lambda:\ d_{c(j)}<D_r}a_j
                              &\le C_r k e^{-\mu_r L}.
 \label{f16v20:boundary-sums}
\end{align}
Membership of a scalar coordinate in $\Lambda$ abbreviates membership of its
cell. The constants do not count all exterior coordinates.
\end{proposition}
\begin{proof}
V17's global, first-order word-gradient remainder gives
\[
 |\partial_jU_\gamma-(H\xi)_j|
             \le C_r\gamma^{-1/2}(1+e_{N_j}),
\]
where $N_j$ has fixed size and lies within one graph step of $c(j)$.
This follows from differentiating the literal unitary exponential words;
quadratic endpoint terms cancel exactly. It is valid even when some inputs
are frozen exterior values. For a free neighbour apply
\eqref{f16v20:prepared-mean}; its distance to the boundary is at least
$d_{c(j)}-1$. An exterior neighbour occurs only at that boundary and has the
compact cap \eqref{f16v20:compact-cap}. This bounds the expected classical
remainder as claimed. On the multiplier's support the only differentiated
Haar logarithm is smooth inside its own angular unit ball, with its original
power $1$ or $1/2$. Its derivative is bounded by
$C\gamma^{-1}|\zeta_{g(j)}|\le C\gamma^{-1/2}$, proving
\eqref{f16v20:drift-bound}. No cut-singular logarithmic derivative is bounded
on the complete chart.

For $i\in A$ and $j\in\Lambda$, the path distance obeys
$d_{\mathcal G}(c(i),c(j))+d_{c(j)}\ge L$.
Choose $0<\mu_r\le\min(\nu,\kappa_r)/2$.
Multiply by the weighted covariance row in \eqref{f16v20:Gaussian-local}
and sum, first over $j$ and then over $i\in A$. This proves the first bound
using $a_j\le\sum_{i\in A}|C_{ji}|$. For an exterior $j$ the distance from
every $i\in A$ is at least $L$. For a free $j$ with $d_{c(j)}<D_r$ it is at
least $L-D_r$. The same row tail proves the second bound; $D_r$ is fixed.
If a core component has no boundary the corresponding terms vanish.
\end{proof}

\subsection{A full core law, uniform over all compact exteriors}

Use total variation in the convention
$d_{\rm TV}(\mu,\nu)=\sup_E|\mu(E)-\nu(E)|$.
Let $K_{\Lambda,A}(q,\cdot)$ be the marginal of $P_\Lambda^q$ on $\zeta_A$,
viewed on $\mathbb R^{d_A}$ by zero extension outside its original charts.

\begin{theorem}[Uniform native core total variation after an integrated buffer]
\label{f16v20:core-TV}
There are fixed $C_r<\infty$, $\mu_r>0$, and $L_r<\infty$ such that, for
$L\ge L_r$, $\gamma\ge1+r^2$, every finite ambient history and every
$\mathcal A\subseteq\Lambda$ as above,
\begin{equation}
 \boxed{
 \sup_q d_{\rm TV}\bigl(K_{\Lambda,A}(q,\cdot),Q_A\bigr)
 \le\min\left\{1,\ C_r k\left(\gamma^{-1/2}
                           +\gamma^{1/2}e^{-\mu_r L}\right)\right\}.}
 \label{f16v20:TV}
\end{equation}
The supremum includes every compact principal exterior value, not just values
with a bounded rescaled energy or a high-probability event. The buffer
$\Lambda\setminus\mathcal A$ is fully integrated. Both the native marginal
normalizer and the Gaussian marginal covariance are retained.
\end{theorem}
\begin{proof}
Take initially a smooth bounded test $\varphi$ with bounded derivatives and
use the solver of Proposition~\ref{f16v20:bounded-solver}. Scale so
$\|\varphi\|_\infty\le1$. Introduce \emph{one common core cutoff}
\[
 \kappa_A(\zeta_A)=\prod_{g\text{ in the core}}\chi(\zeta_g/\sqrt\gamma).
\]
It has $0\le\kappa_A\le1$, support strictly inside all core charts and
$\|\nabla\kappa_A\|\le C\gamma^{-1/2}\sqrt{k}$. Every core group has depth
at least $L$. For $L_r\ge D_r$, the union bound used \emph{only on this
observed core} and \eqref{f16v20:prepared-cut-tail} give
\begin{equation}
              \mathbb E(1-\kappa_A)\le C_r k e^{-c_0\gamma}.
 \label{f16v20:core-cut-mass}
\end{equation}
This is not an all-history event.

As $\delta_Qg=\varphi-Q_A\varphi$ is bounded by two, splitting its expectation
with $\kappa_A$ costs at most twice \eqref{f16v20:core-cut-mass}.
For $j\in A$, integrate by parts using $\kappa_Ag_j$.
For $j\in\Lambda\setminus A$, use $\kappa_A\chi_jg_j$; both $g_j$ and
$\kappa_A$ are independent of this coordinate. No exterior coordinate is
differentiated. These operations give the exact identity
\begin{align}
 \mathbb E[\kappa_A\delta_Qg]
 ={}&-\sum_{j\in A}\mathbb E[\kappa_A g_jD_j]
       -\sum_{j\in\Lambda\setminus A}
                          \mathbb E[\kappa_A\chi_jg_jD_j]\notag\\
 &+\mathbb E[\nabla\kappa_A\cdot g_A]
       +\sum_{j\in\Lambda\setminus A}
                          \mathbb E[\kappa_A(\partial_j\chi_j)g_j]\notag\\
 &+\sum_{j\in\Lambda\setminus A}
                         \mathbb E[\kappa_A(1-\chi_j)(H\xi)_jg_j]
       +\sum_{j\notin\Lambda}\mathbb E[\kappa_A(H\xi)_jg_j].
 \label{f16v20:bounded-cut-identity}
\end{align}
The last sum is an algebraic exterior term, not integration by parts at fixed
coordinates. There is no missing $\partial_jg_j$ there because $g_j$ depends
only on $\zeta_A$. The measured Haar derivatives are included in $D_j$.
Fubini on one interior group justifies every integration; the cutoff removes
its own logarithmic cut. Initially smooth tests suffice for this calculation.

The first line is at most
$C_r k(\gamma^{-1/2}+\gamma^{1/2}e^{-\mu_r L})$ by the drift proposition
and \eqref{f16v20:solver-size}. The first term on the second line is at most
$C\gamma^{-1/2}\sqrt{k}$ by the vector bound on $g_A$.
The remaining cutoff-derivative terms are at most
$C\gamma^{-1/2}\sum_j a_j\le Ck\gamma^{-1/2}$.

On the full compact coordinates, including the fixed exterior,
$|(H\xi)_j|\le C_r\sqrt\gamma$: use the bounded absolute row of $H$, the
bounded mean and \eqref{f16v20:compact-cap}.
In the third line, split free indices by $d_{c(j)}\ge D_r$ and its complement.
For the former, $1-\chi_j$ requires the group cut event; its probability is
at most $C_r e^{-c_0\gamma}$. Their total is at most
$C_r k\sqrt\gamma e^{-c_0\gamma}\le C_r k\gamma^{-1/2}$.
For the latter, and for the last exterior sum, use
\eqref{f16v20:boundary-sums} and the compact bound instead. They cost at most
$C_r k\sqrt\gamma e^{-\mu_r L}$. The discarded expectation
$\mathbb E[(1-\kappa_A)\delta_Qg]$ has the same required bound.
We have proved
\[
 |P_\Lambda^q\varphi(\zeta_A)-Q_A\varphi|
          \le C_rk(\gamma^{-1/2}+\gamma^{1/2}e^{-\mu_r L})
\]
with a constant depending on no derivative norm of the test.

Both core measures are finite regular Borel measures on a Euclidean space
and have Lebesgue densities (one extended by zero). Bounded indicators can
therefore be approximated in their joint $L^1$ norm by smooth bounded tests.
Apply the preceding uniform test bound and pass to that limit. This proves
the total-variation assertion, absorbing its factor-of-two convention in
$C_r$. Although the approximating tests can depend on $q$, the bound does not,
so taking the supremum is legitimate. At no point was a small conditional
partition function replaced by a Gaussian denominator.
\end{proof}

The common core cutoff is important. Cutting each core component separately
and estimating the individual Hessian terms would reintroduce a test-regularity
loss. Here the omitted part multiplies the \emph{whole bounded divergence}.
The proof does not assert finite relative Fisher information of a marginal
with a fractional Haar power at its compact cut.

For example,
\begin{equation}
 L\ge L_r+\mu_r^{-1}\log\gamma
 \quad\Longrightarrow\quad
 \sup_q d_{\rm TV}(K_{\Lambda,A}(q),Q_A)\le C_rk/\sqrt\gamma.
 \label{f16v20:log-buffer}
\end{equation}
A core of any fixed number of cells is thus prepared uniformly, independent
of the size or energy of the frozen exterior. A growing core pays its stated
$k$ cost; this is not a bound for an arbitrary macroscopic output bank.

\subsection{An incoming-law-independent conditional contraction}

For a Markov kernel $K:Z\longrightarrow X$ on standard Borel spaces define
its diameter
\[
                 \delta(K)=\sup_{z,z'}d_{\rm TV}(K(z),K(z')).
\]
A law $\pi$ on $Z$ gives the joint probability
$\mathbb J_\pi(dz,dx)=\pi(dz)K(z,dx)$ and the output marginal $\mu=\pi K$.

\begin{proposition}[Dobrushin diameter bounds all-observable correlation]
\label{f16v20:channel-lemma}
For every incoming law $\pi$ and every $F\in L^2(\mu)$,
\begin{equation}
 \operatorname{Var}_\pi(KF)\le\delta(K)\operatorname{Var}_\mu F,
 \qquad \operatorname{corr}_{\max}(Z,X)^2\le\delta(K).
 \label{f16v20:channel-variance}
\end{equation}
The adjoint conditional expectation has the same squared norm on centered
spaces. In particular, if $\pi'\ll\pi$ has finite chi-square divergence, then
\begin{equation}
      \chi^2(\pi'K\Vert\pi K)\le\delta(K)\chi^2(\pi'\Vert\pi).
 \label{f16v20:chi-square}
\end{equation}
This is a standard data-processing consequence; it requires no lower bound
on either incoming probability and no bound on the test function's sup norm.
\end{proposition}
\begin{proof}
Let $\mathsf K:L^2(\mu)\to L^2(\pi)$ be conditional expectation and let
$\mathsf K^*$ be its adjoint, the reverse conditional kernel.
The operator $B=\mathsf K^*\mathsf K$ on $L^2(\mu)$ is a positive
self-adjoint Markov contraction with invariant law $\mu$. Its transition
rows are mixtures of the rows of $K$, hence their total-variation diameter
is at most $\delta(K)$. This remains true for regular conditional versions
on a full-measure set; extend null rows by any fixed mixture.
For bounded $F$, the oscillation contraction gives
\[
 \|B^l(F-\mu F)\|_2\le\operatorname{osc}(B^l F)
             \le\delta(K)^l\operatorname{osc}(F).
\]
If the centered positive operator $B$ had spectral mass above $\delta(K)$,
a bounded centered $F$, by density, would have a nonzero projection onto an
interval $[a,1]$ with $a>\delta(K)$. The spectral theorem would then give
$\|B^lF\|_2\ge a^l\|1_{[a,1]}(B)F\|_2$, contradicting the last estimate.
Thus $B\preceq\delta(K)I$ on centered space. Since
$\|\mathsf KF\|_2^2=\langle F,BF\rangle$, this proves the first inequality
and the maximal-correlation statement, including unbounded $L^2$ tests.
The adjoint norm is equal. Finally the centered density of $\pi'K$ relative
to $\pi K$ is $\mathsf K^*(d\pi'/d\pi-1)$. Its squared norm is the output
chi-square divergence. Apply the same adjoint norm bound.
\end{proof}

\begin{theorem}[A complete native buffered channel, for every exterior law]
\label{f16v20:buffered-contraction}
For the exact conditional kernel $K_{\Lambda,A}$, set
\begin{equation}
 \Delta_{\gamma,k,L}
   =\min\{1,\ C_r k(\gamma^{-1/2}+\gamma^{1/2}e^{-\mu_rL})\},
 \label{f16v20:Delta}
\end{equation}
with $C_r$ enlarged from Theorem~\ref{f16v20:core-TV}. Then
$\delta(K_{\Lambda,A})\le\Delta_{\gamma,k,L}$, and for \emph{every} law
$\pi$ on all compact exterior parameters,
\begin{equation}
 \boxed{
 \operatorname{Var}_\pi\!\bigl(\mathbb E[F(\zeta_A)\mid q]\bigr)
       \le\Delta_{\gamma,k,L}
                         \operatorname{Var}_{\pi K_{\Lambda,A}}F.}
 \label{f16v20:native-channel}
\end{equation}
This holds for all output $L^2$ functions, including arbitrary rare-event
weights with finite variance. It also holds for every vector of such functions
as a Loewner inequality between its two covariance matrices, with the same
factor. In particular, in the actual native joint law of core and exterior,
\begin{equation}
 \operatorname{corr}_{\max}(\zeta_A,\zeta_{\Lambda^c})
                           \le\sqrt{\Delta_{\gamma,k,L}}.
 \label{f16v20:native-rho}
\end{equation}
The conclusion survives arbitrary reweighting of the \emph{exterior law},
or further fixing outside $\Lambda$. It does not authorize extra conditioning
inside the core or its integrated buffer.
\end{theorem}
\begin{proof}
The triangle inequality through the \emph{same} Gaussian marginal $Q_A$
bounds the distance between any two native rows by twice the right side of
\eqref{f16v20:TV}. Enlarge the constant, cap at one, and apply
Proposition~\ref{f16v20:channel-lemma}. Test that scalar inequality on every
linear combination of a finite vector to obtain the covariance order.
For the actual law, $\pi$ is precisely its exterior marginal, including
$Z_\Lambda(q)$ in its density. No replacement exterior probability is used.
For a different incoming $\pi$, the conditional kernel remains exactly the
same; the lemma was uniform in $\pi$. Further conditioning outside $\Lambda$
only restricts incoming parameters and changes their law. Fixing a buffer
variable instead changes $K_{\Lambda,A}$ itself, which the theorem has not
allowed.
\end{proof}

On the collapsed two-block probability $\mathbb J_\pi$ of $(\zeta_A,q)$,
the two rate-one conditional resamplings consequently have gap at least
$1-\sqrt{\Delta_{\gamma,k,L}}$. Indeed, the two centered coordinate subspaces
have angle norm at most $\sqrt\Delta$; V13's two-projection argument gives
this bound on their full joint $L^2$ space. For a logarithmic buffer and fixed
$k$, sufficiently large coupling makes $\Delta\le1/4$ and gives gap at least
$1/2$ for \emph{this collapsed two-block generator}.

This is not a gap for V14's original local sampler. The buffer has been
integrated out; updating the core conditional on $q$ resamples its buffer
implicitly, and the other collapsed update can resample the entire far
exterior. No comparison returning these two operations to the original
single-port/chord clock is asserted. This limitation is mathematical, not a
choice of terminology.

\subsection{Separated cores: a simultaneous same-law coupling}

\begin{proposition}[Uniform buffered-core coupling with no confidence exception]
\label{f16v20:coupling}
Let free cell sets $\Lambda_s$ be separated so that no primitive interaction
meets two of them. Let $\mathcal A_s\subset\Lambda_s$ have $k_s$ cells and
depth $L_s\ge L_r$. Set
\[
 \eta_s=\min\{1,C_r k_s(\gamma^{-1/2}
                              +\gamma^{1/2}e^{-\mu_rL_s})\}.
\]
On a probability extension of the \emph{unchanged} native law, reference core
vectors $\widehat\zeta_{A_s}$ can be added such that their joint law is
$\bigotimes_s Q_{A_s}$, independently of the common exterior, and for every
subfamily $S$,
\begin{equation}
 \mathbb P\{\zeta_{A_s}\ne\widehat\zeta_{A_s}\text{ for all }s\in S\}
                      \le\prod_{s\in S}\eta_s.
 \label{f16v20:joint-coupling}
\end{equation}
\end{proposition}
\begin{proof}
Condition on the actual common exterior of the separated free sets. Their
interior laws factor exactly by the finite-range specification. For each
core couple its conditional marginal to $Q_{A_s}$ by the common-part
(maximal) coupling, independently between blocks at this fixed exterior.
Its mismatch probability is the corresponding total variation, at most
$\eta_s$ for every exterior by Theorem~\ref{f16v20:core-TV}. The reference
marginals do not depend on the exterior, so their conditional product law
remains the stated independent product after exterior averaging. Fill every
buffer using its original conditional law given its core and the exterior;
this retains the original native full-field law. Conditional multiplication
of the mismatch probabilities, followed by exterior averaging, proves the
claim. Measurability follows from the explicit densities and the measurable
common-part construction on Euclidean spaces. The same construction was used
in f15 V22; the new input here is a uniform bound without a confidence-bad
boundary term. The reference cores are auxiliary random variables, not a
replacement for the integrated native probability.
\end{proof}

\subsection{Why a local channel certificate is not yet the whole positive return}

Two finite examples make the remaining operations precise. They are examples
of probability channels, not counterexamples to the native theorems.

\begin{proposition}[Input reweighting, output conditioning, and coherent aggregation differ]
\label{f16v20:limits}
A channel can have arbitrarily small Dobrushin diameter while conditioning on
an event in its \emph{output} produces diameter arbitrarily close to one.
Also many individually contracting output channels, conditionally independent
given one input, can have a conditional-message Gram whose largest eigenvalue
grows with their number.
\end{proposition}
\begin{proof}
For binary input $Z=\pm1$ and output $X\in\{*,+,-\}$ take the positive rows
\[
 K(+)=\bigl(1-\eta,\eta(1-\sigma),\eta\sigma\bigr),\qquad
 K(-)=\bigl(1-\eta,\eta\sigma,\eta(1-\sigma)\bigr),
 \quad0<\eta<1,\quad0<\sigma<1/2.
\]
Their diameter is $\eta(1-2\sigma)$. Conditional on the output event
$X\ne *$, the diameter is $1-2\sigma$. For $\eta=1/100$ and
$\sigma=1/1000$ these are respectively $499/50000$ and $499/500$.
Input-independent probability of the conditioning event does not prevent
this change of the channel. Thus arbitrary exterior reweighting in
Theorem~\ref{f16v20:buffered-contraction} is not arbitrary output or buffer
conditioning.

For the second assertion let $Z=\pm1$ be uniform and, conditional on $Z$,
let $X_1,\ldots,X_N\in\{\pm1\}$ be independent with
$\mathbb E[X_i\mid Z]=aZ$, $0<a<1$. Each individual channel has diameter $a$
and squared maximal correlation $a^2$. Nevertheless
\[
 \operatorname{Cov}\bigl((\mathbb E[X_i\mid Z])_{i=1}^N\bigr)
                       =a^2\mathbf1\mathbf1^{\mathsf T}
\]
has operator norm $Na^2$, while every diagonal entry is $a^2$.
Conditional independence of the outputs did not make their conditional means
independent under the actual incoming law.
\end{proof}

The new native result provides a true all-observable channel contraction after
a logarithmic buffer, even with an adversarial compact exterior. This is a
usable input for a future block or information-percolation argument: it does
not first assume the required native $L^2$ gap. But the output-size factor in
\eqref{f16v20:Delta} cannot be suppressed. A shell surrounding a large ball
has growing dimension. Applying the theorem to that whole shell eventually
makes the bound trivial at fixed coupling. Applying it separately to small
pieces does not permit a free tensorization through their correlated incoming
law, as the example shows.

In particular the estimate does not verify V19's dyadic all-radius condition.
For a fixed local primitive it gives a relative $L^2$ screening inequality
under any incoming exterior law, but the power gain for that primitive under
the unmodified native law is weaker than V18's fixed-order estimates. Its
advance is the \emph{class of tests and the boundary uniformity}, not a new
arbitrary-order expansion. No previous coefficient or shell Gram is redefined.

The next noncircular target is a \emph{local composition theorem} for these
buffered channels with the actual shared boundary messages and output counts
retained, or a comparison of a compatible buffered block dynamics to the
original update clock. Conditioning inside the buffer, multiplying unrelated
contraction factors, or treating a sparse mismatch coupling as independence
of the native error fields would each omit a required step. The original
positive shell-return norm, and the physical same-top transfer problem,
therefore remain distinct tasks.

\subsection{Diagnostics, checkpoint preservation, and claim scope}

The accompanying diagnostic is
\begin{center}\small
\path{verify_ym_f16_v20_buffered_channels.py}.
\end{center}
It checks the finite stopped boundary recurrence, finite-graph distance and
boundary paths, the core-supported Gaussian divergence, the exact combined
cutoff identity with fixed exterior parameters, channel reversals and
chi-square quadratic forms, conditional couplings, and both limiting examples.
The new checker passed 401 exact assertions. The report records their families
and distinguishes these fixtures from analytic proofs. They do not enclose every native integral or verify the
analytic estimates with a proof assistant.

The cumulative V20 changes only V19's title version and appends this exact
section before its final document-ending command. The package records input
hashes, byte-for-byte predecessor recovery, new diagnostic output, the actual
reruns of the available V4--V19 diagnostics, and compilation/render checks.
No inaccessible companion program or external formalization is represented as
run. The current proof depends on V5/V16's exact independent anchors and local
normalizer estimates, and the identified Gaussian and word-derivative bounds;
it does not use V18--V19's final asymptotic theorems to supply a native gap.
The scheduled V20 review is limited to the supplied sources; the next
checkpoint is V25.

\begin{firewall}
V20 proves deep conditional moment bounds for every compact exterior, a
complete-core total-variation comparison with the original Gaussian marginal,
and a buffered all-observable maximal-correlation/chi-square contraction
uniform in the incoming exterior law. It also supplies a separated-core
coupling without a confidence exception. The output-size cost and integrated
buffer are explicit. No all-radius contraction of an unrestricted source bank,
uniform gap for the original local native sampler, same-top physical transfer
gap, bridge-window removal, or interacting continuum mass gap is inferred.
All native conditional normalizers, free compact domains, Haar factors and
original source definitions are retained.
\end{firewall}
% END F16 V20 APPEND


% BEGIN F16 V21 APPEND -- 2026-09-17
% Insert immediately before the final document-ending command in V20.
% No predecessor proof, bibliography, source definition, or normalizer is edited.

\clearpage
\section{V21: buffered temporal composition and polynomial-width operator asymptotics}
\label{f16v21:context}

V20 proves a buffered channel estimate but warns that unrelated channels cannot
be multiplied through a shared incoming law. We take a composition for which
the required conditional factorization is exact: \emph{complete time slices}
of the augmented native history. They are Markov separators. The continuation
message at each step is retained, including its endpoint and normalization;
no stationary endpoint or Gaussian transition replaces it.

There is also an upstream improvement before composition. V20 bounded its
Stein drift by summing absolute scalar errors. Its conditional preparation
actually supplies second moments of that drift sufficiently far inside the
buffer. Keeping the drift vector intact reduces the core-size cost from
$k/\sqrt\gamma$ to $\sqrt{k/\gamma}$. This permits spatial cross-sections
$B\le c_r\gamma$, rather than the exponential-in-$B$ entry restriction of
V8--V10. Time remains unrestricted. At this maximal width the proved decay
length is $O_r(\log\gamma)$. In every strictly sublinear power-width regime,
a fixed positive temporal decay exponent is available.

The resulting covariance expansion is \emph{two-sided in a weighted temporal
operator-row norm} on the entire source history. It is not a new unrestricted
spatial-volume theorem: the cross-section restriction is essential in the
proof and is stated with every main conclusion. No gap for the original local
sampler or the physical transfer is inferred from the conditional channels.

\subsection{Dependency check and the deliberate restriction of the composition}

The supplied V20 appendix and analytic audit were read in full. Its prepared
conditional moments, bounded-test Gaussian solver, exact combined cutoff
identity, and general-state-space channel lemma are the principal inputs.
V16 supplies full-native local moments; V17 supplies entrywise jets and
polynomial-bank estimates; V10 supplies locality of the normalized coefficients;
V7 supplies the unchanged Haar sources and contact identity. The present
proof does not use V19's one-sided order bound to assume the missing upper
bound. It does not use an unproved native functional inequality.

Two scoped semantic archive searches were empty. Mounted exact-text searches
were then used, and the identified f14 V52 Dobrushin passage and the
monograph's strong-electric slab statement were checked. Those are different
coupling regimes and physical operators. Their temporal mixing principles
are credited, not claimed absent; their hypotheses do not supply the current
weak-coupling conditional input. The current V8--V10 temporal results are also
not forgotten: their exponential spatial cost is precisely what the present
buffered comparison avoids. This is a targeted dependency audit, not an
independent certification of the historical analytic collection. The next
scheduled companion checkpoint remains V25.

For literature context, Polyanskiy--Wu,
\href{https://arxiv.org/abs/1508.06025}{\emph{Strong data-processing inequalities
for channels and Bayesian networks}}, Section 2.1 and Section 4, treats
single-channel contraction and its propagation through an actual network.
The definition of serial composition, the single-channel theorem, and the
network theorem and proof were examined, including their displayed PDF pages.
We only use the elementary serial-chain implication, proved below on the
present standard Borel spaces. No finite-alphabet network or percolation
result is imported without its hypotheses. The supplied memoranda's
separation of power gain from locality cost is retained: the new locality
input is a normalized conditional channel, not another formal coefficient.

\subsection{A square-root core-size estimate from the same bounded-test identity}

Use exactly V20's notation $P_\Lambda^q$, $\mathcal A$, $A$, $k=|\mathcal A|$,
$L=d_{\mathcal G}(\mathcal A,\Lambda^c)$, $Q_A=N(m_A,C_{AA})$, and its
explicit conditional-density version on complete principal charts. In
particular all fast exterior values $q$ are admitted. Let
$\varepsilon=\gamma^{-1/2}$. V20's preparation constants are denoted by
$\kappa_r,B_r,M_r,D_r$, and its Gaussian inverse row exponent by $\nu>0$.
After enlarging $D_r$ once, define
\begin{equation}
 d_\gamma=D_r+2+\left\lceil\kappa_r^{-1}\log\gamma\right\rceil,
 \qquad
 b_\gamma=d_\gamma+L_r+2+
                   \left\lceil2\nu^{-1}\log\gamma\right\rceil.
 \label{f16v21:buffer-length}
\end{equation}
All logarithms below are natural. We may increase a fixed-$r$ lower coupling
threshold so that $\gamma\ge e$, $b_\gamma\ge1$, and
\begin{equation}
                         b_\gamma\le C_{b,r}\log\gamma.
 \label{f16v21:buffer-growth}
\end{equation}
These quantities do not depend on $B,n,\Lambda,q,k$.

\begin{proposition}[Prepared second moments of the measured drift]
\label{f16v21:drift-L2}
Let $D_j=\partial_jV_\gamma-(H(\zeta-m))_j$, with the exact measured potential
and the cutoffs of V20. For a multiplier $0\le\eta_j\le1$ supported where the
original group containing $j$ has angular radius at most one,
\begin{equation}
 \|\eta_jD_j\|_{L^2(P_\Lambda^q)}\le
 \begin{cases}
 C_r\varepsilon,&d_{c(j)}\ge d_\gamma,\\
 C_r\sqrt\gamma,&j\in\Lambda.
 \end{cases}
 \label{f16v21:drift-two-regions}
\end{equation}
Consequently, for any choice of such multipliers and $L\ge d_\gamma$, the
vector indexed by the real core coordinates,
$W_i=\sum_{j\in\Lambda}C_{ij}\eta_jD_j$, satisfies
\begin{equation}
 \mathbb E_{P_\Lambda^q}|W_A|
 \le C_r\sqrt{k}\left(\varepsilon+
                   \sqrt\gamma e^{-\nu(L-d_\gamma)}\right).
 \label{f16v21:vector-drift}
\end{equation}
The multipliers need not be independent, deterministic, or local to $j$.
\end{proposition}
\begin{proof}
Every cell in the fixed drift carrier of a coordinate with depth at least
$d_\gamma$ is free and has depth at least
$\kappa_r^{-1}\log\gamma+D_r$. V20's conditional exponential moment
\eqref{f16v20:prepared-mgf}, on each such cell, bounds its fourth moment
uniformly in every compact $q$. The inherited complete-chart classical bound
is $|\partial_jU_\gamma-\partial_jU_0|\le
C_r\varepsilon(1+e_{N_j})$. Its $L^2$ norm is therefore $C_r\varepsilon$.
On the stated support the only differentiated Haar logarithm is smooth and
its derivative has magnitude at most $C\varepsilon$; the endpoint power
$1/2$ is retained. This proves the first case of
\eqref{f16v21:drift-two-regions}.

For the second case, a scaled Lie or chart derivative of each original
classical word has magnitude at most $C\sqrt\gamma$. The bounded local
incidence, bounded bridge window, exactly quadratic endpoint terms, and
compact cap $|\zeta_c|^2\le E_*\gamma$ preserve this bound. The Gaussian
linear drift has the same bound by its absolute row sum. The cut Haar term
again costs at most $C\varepsilon$. This proves the second case without
claiming thermal moments at arbitrary near-boundary conditionals.

For each $i\in A$, Minkowski in $L^2$, followed by the inverse row bound,
gives
\[
 \|W_i\|_2\le C_r\varepsilon\sum_j|C_{ij}|
       +C_r\sqrt\gamma\sum_{j:\ d_{c(j)}<d_\gamma}|C_{ij}|
 \le C_r\left(\varepsilon+
                 \sqrt\gamma e^{-\nu(L-d_\gamma)}\right).
\]
The last step uses the distance from every core cell to every shallow cell.
Summing the \emph{squared} coordinate bounds and using
$\mathbb E|W_A|\le(\sum_{i\in A}\mathbb E W_i^2)^{1/2}$ proves
\eqref{f16v21:vector-drift}, since $|A|\le36k$. No orthogonality or independence
of drift coordinates has been used.
\end{proof}

\begin{theorem}[Square-root output cost for a complete buffered marginal]
\label{f16v21:sqrt-core}
There are constants $C_{0,r}\ge1$ and $\gamma_{0,r}<\infty$ such that, at all
$\gamma\ge\gamma_{0,r}$ and for every core at depth $L\ge b_\gamma$,
\begin{equation}
 \boxed{\quad
 \sup_q d_{\rm TV}(K_{\Lambda,A}(q),Q_A)
                  \le\min\{1,C_{0,r}\sqrt{k/\gamma}\}.
 \quad}
 \label{f16v21:sqrt-TV}
\end{equation}
The Gaussian is the original full-history \emph{marginal}; the native law,
conditional normalizer, complete free domains, and both endpoint factors are
unchanged. No restriction on ambient volume or duration occurs in this theorem.
\end{theorem}
\begin{proof}
Keep V20's bounded-test solver $b=\nabla u_\varphi$ and
$g_j=C_{jA}b(\zeta_A)$, with $\|b\|_\infty\le C\|\varphi\|_\infty$.
Take $\|\varphi\|_\infty\le1$. Use precisely the common core cutoff
$\kappa_A$ and the exact identity \eqref{f16v20:bounded-cut-identity}.
The first line of that identity can be written as
\[
       -\mathbb E\, b(\zeta_A)\cdot W_A,
\]
with multipliers $\eta_j=\kappa_A$ on core indices and
$\eta_j=\kappa_A\chi_j$ on free noncore indices. They satisfy the support
conditions just proved. This whole line therefore costs
$C_r\sqrt{k}(\varepsilon+\sqrt\gamma e^{-\nu(L-d_\gamma)})$,
not the sum of $k$ unrelated error bounds.

For the common-core derivative use the unchanged bounds
$|\nabla\kappa_A|\le C\varepsilon\sqrt{k}$ and $\|g_A\|_\infty\le C$.
For the remaining cutoff derivatives, write their sum as
\[
 \mathbb E\,\kappa_A b(\zeta_A)\cdot
       \left(\sum_{j\in\Lambda\setminus A}C_{ij}\partial_j\chi_j\right)_{i\in A}.
\]
Each vector component is at most $C\varepsilon\sum_j|C_{ij}|$, so this
term also costs $C\varepsilon\sqrt{k}$.

The noncore cutoff-complement and fixed-exterior terms have the same vector
form. For a free coordinate of depth at least the fixed $D_r$,
\[
 \|(1-\chi_j)(H(\zeta-m))_j\|_2
               \le C_r\sqrt\gamma e^{-c_0\gamma/2},
\]
by the compact drift bound and V20's prepared angular-tail estimate.
For shallower free coordinates and for fixed exterior coordinates use the
compact drift bound and the exponentially small covariance row tail. For
each core component this costs at most
$C_r\sqrt\gamma(e^{-c_0\gamma/2}+e^{-\nu(L-D_r)})$.
Taking the vector second moment bounds these two lines by $\sqrt{k}$ times
that expression. All estimates remain valid after multiplication by
$\kappa_A$.

Only the complement of the \emph{whole bounded divergence} is estimated by
its core union bound, exactly as in V20. Its cost is $C_r k e^{-c_0\gamma}$.
Combining the terms gives, for $L\ge d_\gamma$ and a core deep enough for the
fixed-depth tail bound,
\begin{equation}
 |P_\Lambda^q\varphi-Q_A\varphi|
 \le C_r\sqrt{k}\left(\varepsilon+
                   \sqrt\gamma e^{-\nu(L-d_\gamma)}\right)
                           +C_r k e^{-c\gamma}.
 \label{f16v21:sqrt-pre-bound}
\end{equation}
The exponential noncore term has been absorbed into the first term by
increasing a fixed constant. At $L\ge b_\gamma$, the boundary term is at most
$\gamma^{-3/2}$, by \eqref{f16v21:buffer-length}. If $k\le\gamma$, then
$k e^{-c\gamma}\le C\sqrt{k/\gamma}$; if $k>\gamma$, the proposed capped
bound is already the trivial bound one. Approximation of bounded indicators
by smooth bounded tests in the two measures' joint $L^1$ norm gives total
variation as in V20. No derivative norm of $\varphi$, Fisher-information
assumption at a Haar cut, or changed reference probability is needed.
\end{proof}

The square-root estimate still has an output cost. We now pay it for a
complete slice, rather than claim that infinitely many independently tested
cores form a contracting channel for free.

\subsection{Complete time slices are exact Markov separators}

Let $Z_i=(\zeta_{i,b})_b$ denote all fast cells at time $i$. For $i<n$ its
real dimension is $36B$, and for $i=n$ it is $15B$. These are entire spatial
slices, not independent spatial cubes. The actual word list in V5 has no
interaction meeting times with separation greater than one. The full
probability can therefore be written exactly as
\begin{equation}
 dP=Z^{-1}\prod_{i=0}^n\psi_i(z_i)\,d\nu_i(z_i)
                         \prod_{i=0}^{n-1}W_i(z_i,z_{i+1}),
 \label{f16v21:chain-density}
\end{equation}
where $\nu_i$ is the product of the original cell factors at time $i$,
$\psi_i>0$ contains the within-slice action, and $W_i>0$ contains all
interactions joining the two indicated slices. Source-free scalar factors
cancel only in this normalized fast probability. The original bridge-history
weight is not redefined. No spatial interaction is deleted.

\begin{proposition}[Exact nonstationary transitions with their continuation messages]
\label{f16v21:Markov}
Set $h_n=1$ and recursively define
\begin{equation}
 h_i(x)=\int W_i(x,y)\psi_{i+1}(y)h_{i+1}(y)\,d\nu_{i+1}(y),\qquad
 K_i(x,dy)=\frac{W_i(x,y)\psi_{i+1}(y)h_{i+1}(y)}{h_i(x)}\,d\nu_{i+1}(y).
 \label{f16v21:continuation}
\end{equation}
Then $P$ is the inhomogeneous Markov-chain law with initial measure
$Z^{-1}\psi_0h_0\nu_0$ and transitions $K_i$. Every $h_i$ is finite and
strictly positive at the finite regulator. The same statement, with its
actual terminal factor, holds on an interval whose right exterior is fixed.
There is an analogous backward construction.
\end{proposition}
\begin{proof}
Assign each local action term to its single time or to its unique adjacent
time pair to obtain \eqref{f16v21:chain-density}. Every original complete
coordinate factor is finite; the classical exponent is bounded at a fixed
compact regulator. Thus the recursive integrals are finite and positive.
Their displayed quotients integrate to one. Multiplication of the initial
measure and all transitions cancels each $h_i$ against its literal neighbor,
leaving exactly \eqref{f16v21:chain-density}. In particular the right dressed
endpoint belongs to $\psi_n$, not an imposed invariant measure. A fixed right
boundary contributes its actual edge factor to the terminal integrand and
the same calculation applies. Reading the product from the other end gives
the backward statement, with the original left endpoint. This argument does
not assume equality of different $K_i$, self-adjointness of an individual
transition on a fixed space, or a lower bound on $h_i$ uniform in duration.
\end{proof}

Choose once a fixed width constant
\begin{equation}
 c_r=\min\{1/36,(8C_{0,r})^{-2}\}>0.
 \qquad B\le c_r\gamma.
 \label{f16v21:width}
\end{equation}
Every theorem below that uses contraction assumes \eqref{f16v21:width} and
$\gamma\ge\gamma_{0,r}$. This permits polynomially growing spatial tori but
not arbitrary $B$ at fixed coupling. Define
\begin{equation}
 \delta_{\gamma,B}=2C_{0,r}\sqrt{B/\gamma}\le\tfrac14,
 \qquad q_0=\tfrac14.
 \label{f16v21:delta}
\end{equation}
The constant $q_0$ here is a channel bound, not a cylinder-overlap count.

\begin{theorem}[Serial contraction under the actual endpoint law]
\label{f16v21:temporal-channel}
For $0\le i<j\le n$, let $K_{i,j}$ be the actual conditional transition from
$Z_i$ to $Z_j$ in \eqref{f16v21:chain-density}. Its Dobrushin diameter obeys
\begin{equation}
 \delta(K_{i,j})\le
             \delta_{\gamma,B}^{\lfloor(j-i)/b_\gamma\rfloor}.
 \label{f16v21:diameter-products}
\end{equation}
For all square-integrable $F$ measurable at times at most $i$, and $G$
measurable at times at least $j$, under the actual full native law,
\begin{equation}
 \boxed{
 |\operatorname{Cov}_P(F,G)|\le
 \delta_{\gamma,B}^{\lfloor(j-i)/b_\gamma\rfloor/2}
                 \sqrt{\operatorname{Var}_PF\operatorname{Var}_PG}.}
 \label{f16v21:past-future}
\end{equation}
The statement is uniform in all bounded, arbitrarily time-varying bridge
histories and every duration. No Gaussian transition or stationary endpoint
law is substituted.
\end{theorem}
\begin{proof}
Given the complete past through time $t$, the conditional future is the
native conditional on the free set consisting of all cells at times $t+1$
through $n$. Its marginal at time $t+b_\gamma$ has $B$ cells. Each graph step
changes time by at most one, so this core has depth at least $b_\gamma$;
the physical end at $n$ is free, not an artificial frozen boundary.
Theorem~\ref{f16v21:sqrt-core} compares every such row with the same original
Gaussian slice marginal. Two rows therefore have distance at most
$\delta_{\gamma,B}$. By the exact Markov property this is precisely the
$b_\gamma$-step transition, including the future continuation message in
\eqref{f16v21:continuation}.

For any two probability measures $\alpha,\beta$ and any Markov kernel $K$,
\[
 d_{\rm TV}(\alpha K,\beta K)
                     \le\delta(K)d_{\rm TV}(\alpha,\beta).
\]
To verify it on general spaces, write the Jordan decomposition of
$\alpha-\beta$ as $a(\alpha_+-\alpha_-)$, where
$a=d_{\rm TV}(\alpha,\beta)$ and both positive parts are probabilities.
Mixtures of kernel rows have pairwise distance at most $\delta(K)$, by
integrating the row comparison. Multiplication by $a$ proves the assertion.
Consequently diameters multiply under serial composition. Partition
$[i,j]$ into intervals of length $b_\gamma$ and one shorter remainder;
any remaining kernel has diameter at most one. Exact chain disintegration
proves \eqref{f16v21:diameter-products}.

V20's channel lemma bounds the centered conditional-expectation norm by the
square root of this diameter, for the \emph{actual} law of $Z_i$. The Markov
separator property replaces $F$ and $G$ in their covariance by their
conditional expectations on $Z_i$ and $Z_j$, respectively. These projections
decrease their variances. This proves \eqref{f16v21:past-future}. Equivalently
one can multiply the centered $L^2$ norms between the actual, possibly
unequal, successive marginal spaces. No repeated comparison of unnormalized
kernels, accumulating scalar approximation, or independence of spatial
coordinates occurs.
\end{proof}

\subsection{Two boundary messages can also be kept, rather than conditioned away}

Let $l<a\le b<r$ be time indices. Freeze the two exterior parts up to $l$
and from $r$ onward, and integrate all cells at times $l+1,\ldots,r-1$.
A missing left or right boundary means the original natural endpoint is
retained. Let $\mathcal K_{[a,b]}^{l,r}$ be the conditional marginal of the
whole central time window $[a,b]$. Set $d_-=a-l$, $d_+=r-b$, and use
$\delta^{\lfloor\infty/b_\gamma\rfloor}=0$ for a missing boundary.

\begin{theorem}[Two-sided temporal screening without an algebraic distance floor]
\label{f16v21:two-sided}
When every existing boundary is at distance at least $b_\gamma$ from the
central window,
\begin{equation}
 \delta(\mathcal K_{[a,b]}^{l,r})\le
 \min\{1,\delta_{\gamma,B}^{\lfloor d_-/b_\gamma\rfloor}
                  +\delta_{\gamma,B}^{\lfloor d_+/b_\gamma\rfloor}\}.
 \label{f16v21:two-boundary-TV}
\end{equation}
This is the diameter over the joint pair of boundary values. Thus for every
incoming law of those boundary values and every central-window $L^2$
observable $F$, the variance of its exterior conditional mean is at most
the right side of \eqref{f16v21:two-boundary-TV} times its variance in that
incoming-law joint distribution. The entire central window may have
arbitrary length; no factor counts its cells in this conclusion.
\end{theorem}
\begin{proof}
Vary the left boundary while keeping the right one fixed. Use the exact
forward continuation kernels with that fixed right factor. Each full
$b_\gamma$ step ending no later than $a$ has its target at least
$b_\gamma$ from both existing frozen boundaries: the left boundary is its
current starting slice, and the fixed right boundary is at least $r-a\ge
b_\gamma$ away. The free future interval with its actual right factor is
therefore a conditional to which Theorem~\ref{f16v21:sqrt-core} applies.
Serial contraction bounds the change in the law of $Z_a$ by
$\delta_{\gamma,B}^{\lfloor d_-/b_\gamma\rfloor}$. Given $Z_a$ and the fixed
right boundary, the continuation through $b$ is a common Markov kernel,
independent of the changed left boundary. Adding the rest of the central
window cannot enlarge total variation. This proves the left-boundary bound
for the whole window. Reverse time, with the left boundary now fixed, for
the right-boundary bound. Change two boundary pairs through one intermediate
pair and use the triangle inequality.

The pair of exterior parts enters only through the two boundary slices by
\eqref{f16v21:chain-density}. Apply V20's general incoming-law channel lemma
to the resulting joint-boundary kernel. This proves the variance assertion.
The fixed opposite boundary is retained in each transition before its
contraction is estimated; it is not treated as a harmless later conditioning
of a previously estimated output. The missing-boundary cases use only the
remaining comparison, and both missing boundaries give a constant kernel.
\end{proof}

This is a true fixed-coupling, all-distance temporal statement in the width
regime \eqref{f16v21:width}, not a fixed-order power floor. It does not assert
an unrestricted four-dimensional mixing theorem.

\begin{proposition}[V19's ball messages now decay at every sufficiently large radius]
\label{f16v21:radial-screening}
Let $s_\gamma(L)$ denote precisely V19's centered primitive-message norm and
let $D_m=3\lfloor m/2\rfloor$ be the spatial torus diameter in its metric
$d_*$. Let $\rho$ be a fixed temporal/carrier allowance. If
$h=L-D_m-\rho\ge b_\gamma+2$, then
\begin{equation}
 s_\gamma(L)\le C_r\varepsilon\,
 q_0^{\lfloor(h-1)/b_\gamma\rfloor/2}.
 \label{f16v21:radial-bound}
\end{equation}
The bound is zero when the relevant ball fills the entire history.
In particular V19's dyadic majorant is finite for every admitted coupling,
uniformly in duration at a fixed permitted spatial width. No infinite-volume
tail-triviality premise is used.
\end{proposition}
\begin{proof}
The ball of radius $L$ around a source contains every spatial cell at times
within $L-D_m$ of that source time, intersected with $[0,n]$. Its actual
primitive carrier lies in a central time window of fixed radius $\rho$.
The exterior sigma algebra of the ball is a sub-sigma algebra of the exterior
of this full temporal slab. Conditional expectation onto the smaller algebra
cannot increase the centered $L^2$ norm. Apply
Theorem~\ref{f16v21:two-sided} to the slab. Each existing exterior face is at
least $h-1$ from the central window. V16 bounds the primitive standard
deviation by $C_r\varepsilon$, and the square root of the sum of two
boundary terms costs at most $\sqrt2$ times the indicated power. Missing
physical time faces add no term. If both are missing the conditional mean of
a centered source is zero. The fixed finite sum of dyadic scales before
$D_m+\rho+b_\gamma+2$ has finite terms; afterwards exponential decay beats
the squared-radius factor in V19's criterion. This proves the last assertion
without exchanging infinite-volume and martingale limits.
\end{proof}

\subsection{Whole-spatial-source temporal rows, with the original observables}

Group the original common sources by time, so that each time block has $9B$
coordinates. Their exact Haar residual vector at time $i$, denoted $r_i$, is
measurable at times in $[i-\sigma,i+\sigma]\cap[0,n]$ for a fixed integer
$\sigma$. This follows from the ordinary score's two incident kinetic bonds
and the same temporally local port lift; its spatial inverse is not declared
finite-range. Every component and coefficient remains the one from V7/V10.

For a symmetric time-block matrix define
\begin{equation}
 \|A\|_{{\rm tr},a}=
             \sup_i\sum_{j=0}^n e^{a|i-j|}\|A_{ij}\|_{\rm op},\qquad
                  a_\gamma=\frac{\log2}{4b_\gamma}.
 \label{f16v21:time-row}
\end{equation}
The blocks act on the \emph{entire} spatial common-source space.
Symmetry makes the analogous column bound equal, so this norm controls the
full operator norm. Its subscript means temporal row, not matrix trace.

\begin{proposition}[Coherent covariance decay between arbitrary source time blocks]
\label{f16v21:coherent-decay}
In \eqref{f16v21:width}, for sufficiently large fixed-$r$ coupling,
\begin{equation}
 \|\Gamma_{\gamma,ij}\|_{\rm op}\le
 C_r\varepsilon^2
 \delta_{\gamma,B}^{\lfloor(|i-j|-2\sigma)_+/b_\gamma\rfloor/2}
 \le C_r\varepsilon^2
             2^{-\lfloor(|i-j|-2\sigma)_+/b_\gamma\rfloor}.
 \label{f16v21:coherent-decay-bound}
\end{equation}
The constants are independent of $B,n,i,j$ in this regime.
\end{proposition}
\begin{proof}
Each single-time source bank has $9B\le\gamma/4$ coordinates by
\eqref{f16v21:width}. V17's fixed polynomial-bank theorem, for example with
its fixed exponent one, gives
$\operatorname{Var}_P(v^{\mathsf T}r_i)\le C_r\varepsilon^2|v|^2$ for every
spatial source vector $v$. This is an actual native covariance bound; it is
not inferred from summing $9B$ separate variances. The inherited coefficient
operator bound and V17's bank error supply it also at either dressed time
endpoint. For two separated time carriers apply
\eqref{f16v21:past-future} with their actual intervening time gap. For
intersecting carriers ordinary covariance Cauchy--Schwarz suffices. Take
the supremum over unit spatial vectors and use $\delta_{\gamma,B}\le1/4$.
The full spatially nonlocal lift causes no difficulty because it changes no
temporal carrier outside the stated fixed allowance.
\end{proof}

By itself this proposition bounds the temporal row norm by
$C_r b_\gamma/\gamma$. The next theorem removes that logarithmic amplitude
loss by using the actual local coefficients on the short temporal band.

\begin{theorem}[Two-sided full-history operator-row asymptotics at polynomial width]
\label{f16v21:row-asymptotics}
For every fixed integer $J\ge2$, there are $C_{J,r},\gamma_{J,r}<\infty$
such that, for all $B\le c_r\gamma$, all durations $n\ge1$, all prescribed
bridge histories in the fixed window, and $\gamma\ge\gamma_{J,r}$,
\begin{equation}
 \boxed{
 \left\|\Gamma_\gamma-\sum_{q=2}^J\gamma^{-q/2}G_q\right\|_{{\rm tr},a_\gamma}
 +\left\|\mathsf M_\gamma-\sum_{q=2}^J\gamma^{-q/2}M_q\right\|_{{\rm tr},a_\gamma}
                \le C_{J,r}\gamma^{-(J+1)/2}.}
 \label{f16v21:row-expansion}
\end{equation}
In particular both complete source matrices have operator norms at most
$C_r/\gamma$, and their leading-coefficient errors are at most
$C_r\gamma^{-3/2}$ in the displayed weighted norm. No normalization by
$B(n+1)$ and no restriction on the total number of time-source coordinates
occurs. The exponent $a_\gamma$ is positive but is of order
$1/\log\gamma$, not a uniform positive exponent at the full linear width.
\end{theorem}
\begin{proof}
Write $\varepsilon=\gamma^{-1/2}$ and $G^{[J]}=\sum_{q=2}^J\varepsilon^qG_q$.
The inherited coefficient locality gives constants $a_*>0$ and $C_{q,r}$ such
that
\[
 \|(G_q)_{ij}\|_{\rm op}\le C_{q,r}e^{-a_*|i-j|},
 \qquad \|G_q\|_{{\rm tr},a_*/2}\le C_{q,r}.
\]
Indeed for each time block the scalar row and column sums are bounded by the
weighted full source-row estimate times $e^{-a_*|i-j|}$; Schur's inequality
bounds that block, and the remaining time geometric sum is finite. This is
not the invalid exchange of a source supremum with an unbounded time sum.
Increase the fixed coupling threshold so that $a_\gamma\le a_*/2$.

Take the fixed entry expansion order $K=2J+8$ in V17 and split time lags at
\begin{equation}
 D=2\sigma+
 \left\lceil\frac{4(J+2)}{\log2}\,b_\gamma\log(\varepsilon^{-1})\right\rceil.
 \label{f16v21:lag-cut}
\end{equation}
The entry error on one $9B$ by $9B$ block is at most
$C_{K,r}B\varepsilon^{K+1}\le C_{K,r}\varepsilon^{K-1}$ in operator norm,
by scalar rows and columns. Summing this error on $|i-j|\le D$ with its
weight gives at most
\[
 C_{J,r}(1+D)e^{a_\gamma D}\varepsilon^{K-1}
 \le C_{J,r}(1+\log\varepsilon^{-1})^2
                                \varepsilon^{K-J-3}
 \le C_{J,r}\varepsilon^{J+1}.
\]
Here $b_\gamma=O_r(\log\varepsilon^{-1})$ and
$e^{a_\gamma D}\le C_r\varepsilon^{-(J+2)}$ were used. The finitely many
coefficients with $J<q\le K$ cost $C_{J,r}\varepsilon^{J+1}$ in the entire
weighted row norm, rather than once for every lag.

For $|i-j|>D$, \eqref{f16v21:coherent-decay-bound} and
$2^{-\lfloor x\rfloor}\le2e^{-(\log2)x}$ bound the weighted native row tail by
\[
 C_r\varepsilon^2 b_\gamma
 \exp\!\left[-\frac{3\log2}{4b_\gamma}(D-2\sigma)\right]
 \le C_{J,r}b_\gamma\varepsilon^{3J+8}
 \le C_{J,r}\varepsilon^{J+1}.
\]
The desired coefficient tail is at most
$C_{J,r}\varepsilon^2e^{-a_*D/2}\le C_{J,r}\varepsilon^{J+1}$,
after a fixed threshold increase. Combine short and long lags. These bounds
are valid if the actual duration is shorter than $D$: the tail is then empty.

Finally the same exact Haar-contact identity gives
$\mathsf M_\gamma=-\operatorname{Off}_{>1}\Gamma_\gamma$ and
$M_q=-\operatorname{Off}_{>1}G_q$. This mask only removes entire time blocks,
so it cannot increase the temporal absolute row norm. It proves the second
term of \eqref{f16v21:row-expansion}. Higher coefficients retain their actual
native measure corrections; the core-size argument did not change them.
\end{proof}

\subsection{A fixed temporal exponent below the maximal width}

\begin{theorem}[Uniform positive temporal weight on every sublinear power cross-section]
\label{f16v21:sublinear}
Fix $0\le p<1$ and $A_0<\infty$. There is $a_{p,r}>0$, independent of
$J,B,n,\gamma$, such that, for each fixed $J\ge2$, throughout
\begin{equation}
                   B\le A_0\gamma^p,
                   \qquad\gamma\ge\gamma_{J,p,A_0,r},
 \label{f16v21:sublinear-width}
\end{equation}
the conclusion \eqref{f16v21:row-expansion} holds with $a_{p,r}$ in place of
$a_\gamma$ and constant $C_{J,p,A_0,r}$. Thus the leading-coefficient error is
$O_{p,A_0,r}(\gamma^{-3/2})$ in a \emph{fixed} exponentially weighted temporal
row norm on the entire history. The case $p=0$ means a fixed bounded spatial
cross-section; it is not an assertion for $B=1$ only.
\end{theorem}
\begin{proof}
At a fixed-$p,A_0,r$ coupling threshold,
\[
 B\le c_r\gamma,\qquad
 \delta_{\gamma,B}\le2C_{0,r}\sqrt{A_0}\,
                \gamma^{-(1-p)/2}\le\gamma^{-(1-p)/4}.
\]
For a temporal gap $g\ge2b_\gamma$, the floor bound
$\lfloor g/b_\gamma\rfloor\ge g/(2b_\gamma)$ and
\eqref{f16v21:buffer-growth} give
\[
 \delta_{\gamma,B}^{\lfloor g/b_\gamma\rfloor/2}
                  \le e^{-\lambda_{p,r}g},
             \qquad\lambda_{p,r}=\frac{1-p}{16C_{b,r}}>0.
\]
Choose once
\[
 0<a_{p,r}\le
 \min\{\lambda_{p,r}/4,\ a_*/4,\ (4C_{b,r})^{-1}\}.
\]
This choice does not depend on $J$. Repeat the short/long lag proof with
$K=2J+8$ and
\[
 D=2b_\gamma+2\sigma+
       \left\lceil\frac{J+3}{a_{p,r}}\log(\varepsilon^{-1})\right\rceil.
\]
Then $D=O_{J,p,r}(\log\varepsilon^{-1})$ and
$e^{a_{p,r}D}\le C_r\varepsilon^{-(J+4)}$ because
$b_\gamma\le2C_{b,r}\log\varepsilon^{-1}$.
The near error is at most
$C(1+D)\varepsilon^{K-J-5}=O(\varepsilon^{J+1})$.
The higher coefficients retain their uniform weighted row bounds. On the
far band the new genuine native exponential has rate at least $4a_{p,r}$,
so its weighted tail is at most
$C\varepsilon^2e^{-3a_{p,r}(D-2\sigma)}=O(\varepsilon^{J+1})$;
the coefficient tail has the same bound. There is now no factor $b_\gamma$
from that geometric sum. The time mask is unchanged. This proves the
statement with no order-dependent loss of temporal exponent.
\end{proof}

\subsection{What has closed for the shell return, and what remains spatial}

\begin{proposition}[The positive shell edge is small in the admitted width regime]
\label{f16v21:return-closure}
Fix $J\ge2$ and use a logarithmic radius satisfying V19's local operator-error
conditions at that order. In \eqref{f16v21:width}, with a sufficiently large
fixed-order coupling threshold,
\begin{equation}
       \|\mathsf R^{\rm sh}_{\gamma,L_\gamma}\|_{\rm op}
                          \le C_{J,r}\gamma^{-(J+1)/2}.
 \label{f16v21:return-small}
\end{equation}
For any desired fixed power the radius and fixed orders can therefore be
chosen so that this \emph{actual} positive return is smaller than that power.
The choice of radius is not asserted independent of the requested power.
\end{proposition}
\begin{proof}
V19's exact positive split is
$\Gamma_\gamma-G^{[J]}=\mathsf R^{\rm sh}_{\gamma,L_\gamma}
+\mathsf E_{\gamma,J,M}$, with
$\|\mathsf E_{\gamma,J,M}\|_{\rm op}\le C\gamma^{-(J+1)/2}$ at its specified
logarithmic radius. The new two-sided row estimate bounds
$\|\Gamma_\gamma-G^{[J]}\|_{\rm op}$ by the same power. The triangle
inequality proves \eqref{f16v21:return-small}, taking one fixed $M$.
This argument bounds the previous positive matrix rather than redefining it.
It is consistent with the direct all-radius screening in
Proposition~\ref{f16v21:radial-screening}; neither proof uses an assumed native
sampler gap.
\end{proof}

The spatial improvement is substantial but explicit. V9--V10's native
operator estimates require an entry bound of the form
$e^{KB}/\sqrt\gamma\le1/100$. Here the width can be as large as $c_r\gamma$
for a deteriorating temporal exponent, or $A_0\gamma^p$, $p<1$, for a fixed
positive temporal exponent. The duration and total source count are
unrestricted in either case. At fixed $\gamma$, however, arbitrary spatial
volume is still not allowed. The constants are proved finite from the
inherited local estimates, not optimized numerical thresholds.

This is not the false tensorization warned against in V20. All outputs on a
time slice are included in one vector and their size is paid once. The
intermediate sigma algebras are actual complete-slice separators, and every
future endpoint message stays inside each conditional channel. Arbitrary
conditioning inside a step's buffer is not used. Two-sided screening
re-estimates the channels with the opposite boundary retained, rather than
conditioning a previously contracted output after the fact.

The next spatial question is sharper: replace the cost of a \emph{complete}
separator of $B$ cells by a local spatial composition with its shared
boundary data retained. The square-root bound alone cannot do that once
$B/\gamma$ is unbounded. Repeating the same temporal composition would not
remove this remaining cross-section cost. V19's unrestricted-volume positive
return, V12's separate spatial-cylinder posterior, and a gap for the original
single-port/chord sampler are therefore not declared closed by the current
restricted-width result.

Nothing here integrates the retained physical bridges outside their prescribed
window, proves root Gauss coverage, controls the infrared contact term, or
identifies a conditional fast-channel contraction with the same-top physical
Wilson transfer. A decay exponent in lattice time, especially one that can
deteriorate as $1/\log\gamma$, is not a calibrated physical mass. The
interacting continuum construction and its nontriviality remain distinct
obligations. The native measure and every original source definition are
unchanged throughout.

\subsection{Diagnostics, predecessor preservation, and proof status}

The accompanying diagnostic is
\begin{center}\small\path{verify_ym_f16_v21_temporal_composition.py}.\end{center}
It checks the vector regrouping behind the square-root size bound, exact
nonstationary path disintegrations with unequal endpoint factors, conditioned
continuation kernels, serial diameter contraction, two-boundary screening,
source-window support, and the lag/order arithmetic. The new diagnostic passed 351 exact assertions. The report separates
exact rational and symbolic fixtures from analytic estimates. No fixture is
presented as native Wilson quadrature, a numerical enclosure of the entry
constants, or a proof-assistant certification of V20's premises.

The supplied V20 cumulative source is preserved byte for byte except for one
title-version change and insertion of this exact appendix before its final
document-ending command. The package records exact recovery, input hashes,
actual reruns of all seventeen available V4--V20 diagnostics, the new diagnostic, compilation,
and rendered-page inspection. Old sources are not silently corrected. These
new analytic deductions depend on the identified inherited pin, preparation,
Stein, source, and coefficient results; the archive as a whole is not
independently recertified. The next scheduled companion checkpoint is V25.

\begin{firewall}
V21 proves a square-root core-size total-variation bound under every compact
exterior, a valid serial composition through complete time slices, and
all-distance temporal contraction in the explicit regime $B\le c_r\gamma$.
It proves full-history two-sided covariance and memory asymptotics in a
weighted temporal operator-row norm, with a fixed positive temporal exponent
on every strictly sublinear power cross-section. The old positive shell
return is correspondingly bounded in that regime. It does not prove
unrestricted spatial-volume contraction at fixed coupling, a uniform gap for
the original local sampler, a physical transfer gap, or an interacting
continuum Yang--Mills mass gap. All endpoint normalizers and all compact fast
fields are retained; the bridge-history window and physical clock obligations
are unchanged.
\end{firewall}
% END F16 V21 APPEND


% BEGIN F16 V22 APPEND -- 2026-09-17
% Insert immediately before the final document-ending command in V21.
% No predecessor theorem, normalization, source, or bibliography is altered.

\clearpage
\section{V22: native boundary comparison beyond the Gaussian channel floor}
\label{f16v22:context}

V21 pays the discrepancy between a native core and a Gaussian marginal before
composing its temporal channels. That discrepancy includes a local nonlinear
correction which is present even when the exterior has no influence. It is
therefore not the right floor for comparing \emph{two native boundary rows}.
We retain the nonlinear core integral instead. A small-angular comparison is
used only inside a finite integrated buffer; the actual conditional law pays
its complement, uniformly over every compact exterior. The resulting channel
bound has an exponentially small large-field term, not an intrinsic
$\gamma^{-1/2}$ floor.

This is an upstream change in the reference used for boundary comparison,
not another order of Gaussian perturbation theory. It permits an
arbitrary prescribed polynomial spatial cross-section, with a fixed positive
temporal exponent, for the full native covariance expansion. A separate
all-observable temporal-channel statement allows an exponentially growing
cross-section. Neither assertion is an unrestricted-spatial-volume theorem
at fixed coupling. The local spatial composition problem remains explicit.

\subsection{Dependency review and the native comparison being changed}

The supplied V21 appendix and audit were read in full. Its remaining
complete-separator cost, V20's conditional preparation, and V4's product-ball
one-form covariance identity were checked against the cumulative source.
Two scoped semantic archive searches returned no results. Exact mounted-text
searches and the identified f15 V22 passages were then used. F15 V22 already
proves a restricted core density comparison through a pinned-complement
covariance identity, and a transfer to the full law with a confidence
exception. We credit both mechanisms. The new ingredients are their
application to the current unmoving specification with \emph{uniform}
compact-boundary preparation, a square-root core-size response bound, and an
exponentially small finite-buffer exception. No archive-wide novelty or
independent recertification claim is made.

For boundary-aware analytic context, D.~Le Peutrec,
\href{https://arxiv.org/abs/1607.08714}{\emph{On Witten Laplacians and
Brascamp--Lieb's inequality on manifolds with boundary}}, derives the
Neumann/Dirichlet form identities for smooth boundaries. Its introductory
identity and boundary hypotheses were checked. The product-ball realization
with corners used below is the one already specified in V4 and f15 V13;
we do not pretend that a smooth-boundary theorem alone states that product
result. The elementary resolvent estimate needed here is supplied below.

A spatial-composition literature check also examined the definitions and
initial theorems of Spinka's
\href{https://arxiv.org/abs/1803.10578}{\emph{Finitary codings for spatial
mixing Markov random fields}}. The distinction between weak spatial mixing,
strong spatial mixing, and a grand coupling is relevant. Its finite-valued,
translation-invariant coding theorem is not applied to this inhomogeneous
compact law. The abstract of Chazottes--Redig--V\"ollering,
\href{https://arxiv.org/abs/1012.2489}{\emph{Poincar\'e inequality for Markov
random fields via disagreement percolation}}, was screened only. Its required
subcritical disagreement coupling has not been constructed here.

The import memoranda's separate ledgers for power gain and locality cost
remain useful; neither a physical gap nor an all-scale cluster theorem is
an input. The current proof uses V5/V16's exact specification and bounds,
V20's all-boundary moment preparation, and the inherited product-convex
covariance realization. V17 and V10 are used only for the subsequent source
asymptotics. The next scheduled companion checkpoint remains V25.

\subsection{A comparison domain, not a replacement for the native law}

Use V20's interaction graph $\mathcal G$, distance $d_{\mathcal G}$, and
complete unmoving cells. Write
\[
 P=P_{\gamma,\mathbf y},\quad \varepsilon=\gamma^{-1/2},\quad
 \mathcal N=B(n+1),\quad \max_{i,e}|y_{i,e}|\le r.
\]
A cell has 36 or 15 real coordinates. The full measured potential on its
open principal charts is $V_\gamma=U_\gamma-\log J^f_\gamma$. Its Gaussian
fast precision is the same $H$, with $h_-I\preceq H\preceq h_+I$.
Every non-anchor interaction has graph diameter one. No new logarithm of a
matrix product is used in differentiating the local exponential words.

Choose a fixed angular radius $0<\eta_*<1/4$, small enough as specified
below, and put
\[
 \mathcal D_{\gamma,c}=\{\zeta_c:|\zeta_c|\le\eta_*\sqrt\gamma\},
 \qquad \mathcal D_{\gamma,W}=\prod_{c\in W}\mathcal D_{\gamma,c}.
\]
Each is a ball in a \emph{whole unmoving cell}. All its original group
coordinates have angular radius at most $\eta_*$, so their Haar factors,
including the endpoint powers $1/2$, are smooth and positive nearby.
This is an auxiliary comparison domain. It is neither V4's event
$\mathcal O_R$ nor a redefinition of $p_{\rm good}$.

For a set $W$, let $C=\partial_{\mathcal G}W$ be its actual exterior word
collar, with redundant adjacent cells harmless. For $z_C$ in the corresponding
product of small balls define the \emph{restricted native} conditional
\begin{equation}
 \pi_W^z(d\zeta_W)=
 \frac{1_{\mathcal D_{\gamma,W}}e^{-U_W(\zeta_W;z)}
                         \prod_{c\in W}d\nu_{\gamma,c}}
      {Z_W^{\rm box}(z)}.
 \label{f16v22:restricted-law}
\end{equation}
Every term meeting $W$ is retained, including crossing terms. Exterior-only
constants cancel in this conditional probability, not from the full exterior
posterior. The free domain is independent of $z$. Its measured potential
includes the complete original Haar factors on $W$.

\begin{proposition}[Uniform restricted convexity and a spatial inverse]
\label{f16v22:convexity}
There are $\eta_*>0$, $\gamma_r<\infty$, $\lambda>0$, and $M>\lambda$, with
constants independent of $W,B,n$, such that for $\gamma\ge\gamma_r$ and every
admitted $z$, throughout
$\mathcal D_{\gamma,W}$,
\begin{equation}
       \lambda I\preceq\nabla_W^2V_\gamma^z\preceq MI.
 \label{f16v22:H-bounds}
\end{equation}
Scalar mixed derivatives involving a free and a collar coordinate have
uniform absolute row and column bounds and vanish beyond graph distance one.
The same Hessian statements hold when any subset of the free cells is fixed
anywhere in its small balls.

The reflecting scalar law satisfies
$\operatorname{Var}_{\pi_W^z}f\le\lambda^{-1}\mathbb E|\nabla f|^2$.
For its absolute one-form covariance realization $\mathfrak H$, and cell
component projections $P_c,P_d$ on any remaining free set,
\begin{equation}
 \|P_c\mathfrak H^{-1}P_d\|
       \le\lambda^{-1}(1-\lambda/M)^{d_{\mathcal G}(c,d)}.
 \label{f16v22:inverse}
\end{equation}
Distances in the full graph give a valid, possibly weaker, bound after
pinning. The covariance identity is the inherited one from V4.
\end{proposition}
\begin{proof}
For each original classical term, two coordinate derivatives of
$\gamma f(w/\sqrt\gamma)$ have a uniformly bounded third-derivative remainder
relative to its quadratic limit. On the displayed product, every local fast
argument has angular magnitude at most a fixed multiple of $\eta_*$; each
bridge argument is at most $r/\sqrt\gamma$. Bounded word length, local
linear-map norms, and incidence give
\[
 \|\nabla_W^2 U_W-H_{WW}\|\le C(\eta_*+r/\sqrt\gamma).
\]
The estimate is by bounded symmetric rows, not by the number of cells.
The exactly quadratic endpoint terms are retained. The Hessian of each
$-w\log j(\zeta_g/\sqrt\gamma)$ is bounded by $C/\gamma$ on this domain,
with the original $w=1$ or $1/2$. Choose $\eta_*$ first and then a fixed-$r$
threshold so that the total difference is at most $h_-/2$. One may take
$\lambda=h_-/2$ and enlarge $M$. Mixed rows with collar coordinates follow
from the same local derivative bounds. Principal submatrices preserve the
lower and upper bounds after more cells are fixed. No convexity outside
these small angular domains is asserted.

Use exactly the product absolute form described in
\eqref{f16v4:witten-form}: its derivative and ball-boundary part
$\mathfrak D$ is nonnegative and block diagonal in cell component labels.
The outward second fundamental forms are nonnegative; the component at its
own ball boundary is tangential. Cell projections preserve the form domain.
The scalar/one-form intertwining gives
\[
 \operatorname{Cov}(f,g)=
       \langle\nabla f,(\mathfrak D+\nabla^2V)^{-1}\nabla g\rangle.
\]
The inverse lower bound gives the Poincar\'e inequality. These statements use
the inherited product-ball boundary realization, including its corners,
not boundary-free integration by parts.

Put $\mathfrak A=\mathfrak D+MI$ and $\mathfrak B=MI-\nabla^2V$.
Then $0\preceq\mathfrak B\preceq(M-\lambda)I$,
$\mathfrak A^{-1/2}$ preserves cell component labels, and
\[
 \mathfrak H^{-1}=\mathfrak A^{-1/2}
 \sum_{l\ge0}(\mathfrak A^{-1/2}\mathfrak B\mathfrak A^{-1/2})^l
 \mathfrak A^{-1/2}.
\]
The series converges in operator norm, with ratio at most
$1-\lambda/M$. Its $l$th term cannot connect component labels at distance
greater than $l$. Summing the geometric tail proves \eqref{f16v22:inverse}.
The differential coefficients may depend on all free variables; the range
statement is about \emph{derivative indices}, not finite propagation in
configuration space. Product balls of either cell dimension and their
principal-subset products satisfy the same argument.
\end{proof}

\subsection{Compare native rows without paying their common nonlinear correction}

Let $A$ denote all coordinates of a nonempty cell set $\mathcal A\subset W$,
$k=|\mathcal A|$, and assume $d_{\mathcal G}(\mathcal A,C)\ge\ell+1$,
$\ell\ge3$. Write $\pi_{W,A}^z$ for the core marginal of
\eqref{f16v22:restricted-law}. When $C$ is empty this marginal is independent
of exterior data and every boundary-comparison error below is zero.

\begin{theorem}[Restricted native boundary response with square-root output cost]
\label{f16v22:box-TV}
There are $C_r<\infty$ and $a>0$, independent of $W,A,B,n,\gamma$, such that
\begin{equation}
 \sup_{z,z'\in\mathcal D_{\gamma,C}}
 d_{\rm TV}(\pi_{W,A}^z,\pi_{W,A}^{z'})
            \le\min\{1,C_r\sqrt{\gamma k}\,e^{-a\ell}\}.
 \label{f16v22:box-boundary}
\end{equation}
The distance convention is $\sup_E|\mu(E)-\nu(E)|$. The reference in this
comparison is native, not Gaussian. No expansion of the restricted
partition function is used.
\end{theorem}
\begin{proof}
Join $z$ to $z'$ by its straight segment $z_s$. It stays in the product of
collar balls and has scalar displacements at most $2\eta_*\sqrt\gamma$.
Let $T=W\setminus\mathcal A$. At fixed core $x$, disintegrate the restricted
native density over the \emph{fixed} product domain on $T$. Write $p_s(x)$
for the normalized core density and
\[
 S_s=\partial_s V_W^{z_s},\qquad
 q_s(x)=\mathbb E_{\pi_T^{x,z_s}}S_s.
\]
All denominators are positive and all differentiated integrands smooth on
these compact sets. Normalized differentiation gives, exactly,
\begin{equation}
 \partial_s\log p_s(x)=-q_s(x)+\mathbb E_{\pi_{W,A}^{z_s}}q_s,
 \qquad
 \partial_{x_i}q_s=\mathbb E\partial_{x_i}S_s
                       -\operatorname{Cov}(S_s,\partial_{x_i}V_W^{z_s}).
 \label{f16v22:score-identity}
\end{equation}
The covariance is in the pinned-complement conditional on $T$. In particular
both the inner and outer normalizers have been retained. This is the same
mixed-normalization mechanism as f15 V22, on the present conditional.

The direct derivative in the second identity vanishes by separation.
For a collar scalar coordinate $j$, the free gradients of
$\partial_{z_j}V$ and $\partial_{x_i}V$ have fixed carriers within one graph
step of their respective cells and bounded norms. The inherited covariance
identity and \eqref{f16v22:inverse} consequently give
\[
 |\operatorname{Cov}_{\pi_T^{x,z_s}}
                    (\partial_{z_j}V,\partial_{x_i}V)|
            \le C e^{-a_0d_{\mathcal G}(c(i),c(j))}
\]
for some $a_0>0$, after absorbing the fixed two-step allowance. The case of
empty $T$ has zero covariance. Fix $0<a<a_0/2$. Polynomial graph growth and
bounded coordinate multiplicity imply
\[
 \sum_{j\in C}e^{-a_0d_{\mathcal G}(c(i),c(j))}
                  \le C e^{-a\ell}.
\]
Summing the actual collar displacement therefore proves
\begin{equation}
 \sup_x|\partial_{x_i}q_s(x)|\le C_r\sqrt\gamma e^{-a\ell},
 \qquad \sup_x|\nabla_Aq_s(x)|^2\le C_r\gamma k e^{-2a\ell}.
 \label{f16v22:marginal-score-gradient}
\end{equation}
This uses a distance-summed row bound, not the collar's total cardinality.

The core marginal inherits the full restricted Poincar\'e inequality by
lifting a function of $x_A$ to $W$. Hence
$\operatorname{Var}_{\pi_{W,A}^{z_s}}q_s\le C_r\gamma k e^{-2a\ell}$.
For every bounded measurable $f$ with $|f|\le1$, differentiation of its
expectation, not of $f$, yields
\[
 \left|\partial_s\mathbb E_{\pi_{W,A}^{z_s}}f\right|
       =|\operatorname{Cov}(f,q_s)|
       \le C_r\sqrt{\gamma k}e^{-a\ell}.
\]
Finite-dimensional compact differentiation justifies this also for
measurable bounded $f$ by dominated convergence. Integrate in $s$ and take
indicator tests. This proves \eqref{f16v22:box-boundary}. No derivative norm
of the observed function and no native \emph{full-law} Poincar\'e inequality
has been assumed.
\end{proof}

\subsection{Pay the entire comparison buffer under every actual compact boundary}

Use V20's complete $P_\Lambda^q$, with an arbitrary compact exterior $q$.
For $\mathcal A\subseteq\Lambda$ and integer $\ell\ge3$, set
\[
 W=\{c:d_{\mathcal G}(c,\mathcal A)\le\ell\},\qquad
 C=\partial_{\mathcal G}W,\qquad
 G_{W,C}=\{\zeta_c\in\mathcal D_{\gamma,c}\text{ for all }c\in W\cup C\}.
\]
Balls and collars refer to the actual finite graph. Assume
\begin{equation}
       L=d_{\mathcal G}(\mathcal A,\Lambda^c)\ge\ell+D_*+2,
 \label{f16v22:buffer-geometry}
\end{equation}
where $D_*$ is fixed below. In particular $W\cup C\subseteq\Lambda$;
no artificially frozen time endpoint is introduced.

\begin{theorem}[Complete buffered rows with an exponentially small native exception]
\label{f16v22:native-TV}
There are $D_*<\infty$, $c_*>0$, and $C_r<\infty$, independent of ambient
volume, duration, core size, and frozen exterior, such that for
$\gamma\ge\max\{1+r^2,\gamma_r\}$ and geometry \eqref{f16v22:buffer-geometry},
\begin{equation}
 \boxed{
 \sup_q d_{\rm TV}(K_{\Lambda,A}(q),\lambda_{W,A})
  \le\min\{1,C_r[\sqrt{\gamma k}e^{-a\ell}
                       +k(1+\ell)^4e^{-c_*\gamma}]\},
 \quad \lambda_{W,A}:=\pi_{W,A}^{0}.}
 \label{f16v22:native-buffer-bound}
\end{equation}
The reference is the explicitly normalized restricted \emph{native} integral
at zero collar. It can depend on $\gamma,W,\mathbf y$ but not on $q$.
The target is the complete native core marginal. All compact fast fields
are retained there. Consequently its Dobrushin diameter satisfies the same
bound with $C_r$ enlarged, as does its squared maximal correlation for
\emph{every} incoming exterior law.
\end{theorem}
\begin{proof}
Let $M_r,B_r,\kappa_r$ be the conditional preparation constants from
\eqref{f16v20:prepared-mgf}. Choose a fixed $D_*$ so that
$B_r e^{-\kappa_rD_*}\le\eta_*^2/2$, and put
$c_*=t_0\eta_*^2/2>0$. At any cell of depth at least $D_*$,
exponential Markov in that actual conditional gives
\[
 P_\Lambda^q\{e_c>\eta_*^2\gamma\}
     \le\exp(t_0M_r-c_*\gamma).
\]
This is the fixed-depth angular tail, not a claim that arbitrary boundary
cells have uniform thermal moments. Every cell of $W\cup C$ has this depth
by \eqref{f16v22:buffer-geometry}. Polynomial ball growth gives
$|W\cup C|\le Ck(1+\ell)^4$. The union bound therefore yields
\begin{equation}
 \sup_q P_\Lambda^q(G_{W,C}^{\,c})
                 \le C_r k(1+\ell)^4e^{-c_*\gamma}.
 \label{f16v22:buffer-tail}
\end{equation}
This event involves only the chosen finite comparison buffer, not all cells
of the history.

Condition the \emph{actual} law further on $G_{W,C}$; it has positive
probability. Given all variables outside $W$, its conditional inside $W$ is
exactly \eqref{f16v22:restricted-law} at the observed small collar $z_C$.
Thus its core marginal is a mixture of $\pi_{W,A}^z$ over a posterior that
retains the conditional normalizers $Z_W^{\rm box}(z)$ and the actual gate
weights. No property of that posterior except its support is needed.
Every row in this mixture is within \eqref{f16v22:box-boundary} of
$\lambda_{W,A}$. The complete law and its restriction differ in total
variation by $P_\Lambda^q(G_{W,C}^{\,c})$, and taking a marginal cannot
increase this distance. The triangle inequality proves
\eqref{f16v22:native-buffer-bound}. There is no division by a small
estimated good-event probability.

Compare two native rows through this \emph{same} reference to bound the
diameter. V20's proved general-space channel lemma then gives its
all-incoming-law, all-$L^2$ maximal-correlation consequence. The good event
is not commuted through any original update or differentiated in a retained
bridge source. It is eliminated by the exact mixture argument before the
claim about the full native kernel is made.
\end{proof}

This removes a comparison floor, not all large fields from the target law.
The nonlinear core correction is present on both restricted native boundary
rows and cancels in their comparison. It was not zero when either row was
compared with $Q_A$. The estimates \eqref{f16v22:box-boundary} and
\eqref{f16v22:buffer-tail} keep these two different errors separate.

\begin{proposition}[Any fixed polynomial output and any fixed power of contraction]
\label{f16v22:polynomial-core}
Fix $p\ge0$, $s>0$ and $A_0<\infty$. There is
$b_{\gamma,p,s}=O_{p,s,r}(\log\gamma)$ such that every core with
$k\le A_0\gamma^p$ and depth at least $b_{\gamma,p,s}$ has
\begin{equation}
          \delta(K_{\Lambda,A})\le\gamma^{-s}
 \label{f16v22:arbitrary-power-channel}
\end{equation}
for $\gamma\ge\gamma_{p,s,A_0,r}$. There is no restriction on the ambient
volume or duration in this local-core statement.
\end{proposition}
\begin{proof}
Take
\[
 \ell=3+\left\lceil a^{-1}\bigl(s+(p+1)/2+1\bigr)\log\gamma\right\rceil,
 \qquad b_{\gamma,p,s}=\ell+D_*+2.
\]
The first term of \eqref{f16v22:native-buffer-bound} is at most a fixed
constant times $\gamma^{-s-1}$. The second is a fixed polynomial times
$e^{-c_*\gamma}$, hence also at most that power at a fixed threshold.
Absorb the constants and the two-row factor by increasing the threshold.
All exponents were fixed before the coupling limit. No order-dependent
Taylor expansion was used in this proof.
\end{proof}

\subsection{Valid temporal composition on much larger separators}

We now reuse, rather than reprove, V21's exact complete-slice Markov
factorization, including every continuation message and endpoint factor.
Its serial and two-boundary proofs require a uniform diameter for a
buffered slice under every compact exterior, including the fixed opposite
endpoint. Theorem~\ref{f16v22:native-TV} supplies exactly that input with a
better boundary-comparison scale.

\begin{theorem}[All polynomial spatial widths have a fixed temporal decay rate]
\label{f16v22:polynomial-time}
Fix $p\ge0$ and $A_0<\infty$, and assume
\begin{equation}
                B\le A_0\gamma^p.
 \label{f16v22:polynomial-width}
\end{equation}
At a fixed-$p,A_0,r$ coupling threshold, set $b_\gamma=b_{\gamma,p,2}$.
There is $C_{b,p,r}$ with $b_\gamma\le C_{b,p,r}\log\gamma$, and the
actual complete-slice transitions satisfy
\begin{equation}
 \delta(K_{i,j})\le\gamma^{-2\lfloor(j-i)/b_\gamma\rfloor},
 \qquad
 |\operatorname{Cov}_P(F,G)|
 \le\gamma^{-\lfloor(j-i)/b_\gamma\rfloor}
                         \sqrt{\operatorname{Var}_PF\operatorname{Var}_PG}.
 \label{f16v22:polynomial-time-bound}
\end{equation}
Here $F$ and $G$ are arbitrary past/future $L^2$ observables separated by
$i,j$, as in V21. For $j-i\ge2b_\gamma$, their normalized correlation is
at most $e^{-\lambda_{p,r}(j-i)}$ for a fixed $\lambda_{p,r}>0$.
Two-sided central-window screening holds with the sum of the two respective
powers in the diameter bound and no central-window-size factor.
\end{theorem}
\begin{proof}
Observe one complete output slice, so that its core size is $k=B$.
Given the complete past, the future interval is the original conditional
with its natural endpoint, not a transition with that endpoint discarded.
A step of length $b_\gamma$ has graph depth at least $b_\gamma$. Apply
Proposition~\ref{f16v22:polynomial-core} with $s=2$.
V21's exact serial proof multiplies these diameters and gives the first
inequality. Its all-observable channel argument gives the square-root
power in the covariance bound. For a gap $g\ge2b_\gamma$,
\[
 \gamma^{-\lfloor g/b_\gamma\rfloor}
       \le\exp[-g\log\gamma/(2b_\gamma)]
       \le\exp[-g/(2C_{b,p,r})].
\]
A smaller fixed rate can be selected for later estimates. The two-boundary
argument of Theorem~\ref{f16v21:two-sided} applies without modification:
retain the right boundary in every forward conditional, then reverse time
and retain the left boundary. No previously contracted output is
subsequently conditioned inside its buffer. All constants are independent
of the duration, which is unrestricted.
\end{proof}

\begin{theorem}[An exponential-width channel result, distinct from source asymptotics]
\label{f16v22:exponential-time}
There are $\beta_r>0$, $c_r'>0$, and a buffer $b_\gamma^{\rm exp}=O_r(\gamma)$
such that, for
\begin{equation}
                B\le e^{\beta_r\gamma},\qquad \gamma\ge\gamma_r',
 \label{f16v22:exponential-width}
\end{equation}
the actual $b_\gamma^{\rm exp}$-step slice diameter is at most
$e^{-c_r'\gamma}$. Hence all past/future $L^2$ correlations have a fixed
positive exponential decay rate at gaps at least $2b_\gamma^{\rm exp}$.
The same two-boundary conclusion holds. This is not a covariance expansion
in this exponential-width regime.
\end{theorem}
\begin{proof}
In \eqref{f16v22:native-buffer-bound}, choose
\[
 \beta_r=c_*/8,\qquad
 \ell=3+\lceil c_*\gamma/(4a)\rceil,\qquad
 b_\gamma^{\rm exp}=\ell+D_*+2.
\]
For $k=B\le e^{c_*\gamma/8}$ the two error terms are bounded respectively
by $C_r\sqrt\gamma e^{-3c_*\gamma/16}$ and
$C_r(1+\gamma)^4e^{-7c_*\gamma/8}$. After a fixed threshold increase the
two-row diameter is at most $e^{-c_*\gamma/16}$. Take $c_r'=c_*/16$.
The exact serial argument gives a correlation factor
$\exp[-c_r'\gamma\lfloor g/b_\gamma^{\rm exp}\rfloor/2]$.
Since $b_\gamma^{\rm exp}\le C_r'\gamma$, it is at most
$\exp[-c_r'g/(4C_r')]$ when $g\ge2b_\gamma^{\rm exp}$.
The opposite-boundary proof is the same as above.

For gaps shorter than this buffer the general bound is only one. An
exponentially wide source slice is not a polynomial bank in V17. Accordingly
neither a coupling-independent short-gap prefactor nor an all-source
operator expansion at this width is inferred from this theorem.
\end{proof}

\subsection{Fixed-weight operator asymptotics at every polynomial cross-section}

Retain V21's temporal-row norm, original common-source blocks, and fixed
source time-carrier allowance $\sigma$. For a symmetric time-block matrix,
\[
 \|A\|_{{\rm tr},b}=\sup_i\sum_j e^{b|i-j|}\|A_{ij}\|_{\rm op}.
\]
This is not a trace norm. Each block has $9B$ coordinates and includes the
whole spatially nonlocal Haar lift, whose time support remains bounded.

\begin{theorem}[Two-sided native expansion for every fixed polynomial width]
\label{f16v22:row-expansion}
For every fixed $p\ge0$ there is $a_{p,r}>0$, independent of expansion order,
coupling, volume, and duration, such that for every fixed $J\ge2$ and
$A_0<\infty$, throughout \eqref{f16v22:polynomial-width} at a sufficiently
large $\gamma_{J,p,A_0,r}$,
\begin{equation}
 \boxed{
 \begin{aligned}
 &\left\|\Gamma_\gamma-\sum_{q=2}^J\gamma^{-q/2}G_q
                          \right\|_{{\rm tr},a_{p,r}}\\
 &\quad+\left\|\mathsf M_\gamma-\sum_{q=2}^J\gamma^{-q/2}M_q
                          \right\|_{{\rm tr},a_{p,r}}
              \le C_{J,p,A_0,r}\gamma^{-(J+1)/2}.
 \end{aligned}}
 \label{f16v22:row-asymptotics}
\end{equation}
The coefficients, sources, complete native law and endpoints are unchanged.
Their full operator norms are $O_{p,A_0,r}(\gamma^{-1})$. The
leading-coefficient errors are $O_{p,A_0,r}(\gamma^{-3/2})$, for every
duration. There is no division by $\mathcal N$ and no source-vector average.
\end{theorem}
\begin{proof}
V17's coherent bank theorem applies to a single time block: for large enough
fixed-$p,A_0$ coupling, $9B\le\gamma^{p+1}$. Thus every spatial vector $v$
satisfies
\[
 \operatorname{Var}(v^{\mathsf T}r_i)\le C_{p,A_0,r}\varepsilon^2|v|^2.
\]
This is the native bank theorem, not the sum of separate component variances.
Together with Theorem~\ref{f16v22:polynomial-time}, it gives
\[
 \|\Gamma_{\gamma,ij}\|\le C_{p,A_0,r}\varepsilon^2,
 \qquad
 \|\Gamma_{\gamma,ij}\|\le
 C_{p,A_0,r}\varepsilon^2 e^{-\lambda_{p,r}(|i-j|-2\sigma)}
 \quad\text{if }|i-j|-2\sigma\ge2b_\gamma.
\]
The Gaussian coefficient locality of V10 gives a fixed $a_*>0$ and
$\|G_q\|_{{\rm tr},a_*/2}\le C_{q,r}$ at every fixed order, just as in V21.
Choose once
\[
 0<a_{p,r}\le\min\{\lambda_{p,r}/4,\ a_*/4,\ (4C_{b,p,r})^{-1}\}.
\]
This choice is independent of $J$ and $A_0$; the entry threshold may depend
on both. Set $b=a_{p,r}$ and
\[
 D=2b_\gamma+2\sigma+
       \left\lceil (J+4)b^{-1}\log(\varepsilon^{-1})\right\rceil,
 \qquad K=2J+2\lceil p\rceil+12.
\]
Here $D=O_{J,p,r}(\log\varepsilon^{-1})$ and
$e^{bD}\le C_{p,r}\varepsilon^{-(J+5)}$, since
$b_\gamma\le2C_{b,p,r}\log\varepsilon^{-1}$.

On $|i-j|\le D$, V17's order-$K$ entry error has block operator norm at
most $C_{K,p,A_0,r}B\varepsilon^{K+1}\le
C_{K,p,A_0,r}\varepsilon^{K+1-2p}$ by scalar rows and columns. Summing over
the near band with its weight costs at most
\[
 C(1+D)e^{bD}\varepsilon^{K+1-2p}
      \le C(1+\log\varepsilon^{-1})\varepsilon^{K-J-2p-4}
      \le C\varepsilon^{J+1}.
\]
The coefficients of orders $J+1$ through $K$ cost
$C\varepsilon^{J+1}$ in their entire weighted row norm, not once per lag.

On the far band the genuine native decay has rate at least $4b$, so its
weighted sum is at most
$C\varepsilon^2 e^{-3b(D-2\sigma)}=O(\varepsilon^{J+1})$.
The desired coefficient tail is bounded the same way, using its rate at
least $4b$. Short actual histories simply have an empty far band.
This proves the first term of \eqref{f16v22:row-asymptotics}.
The exact identity $\mathsf M_\gamma=-\operatorname{Off}_{>1}\Gamma_\gamma$
and $M_q=-\operatorname{Off}_{>1}G_q$ removes entire time blocks and cannot
increase this row norm, proving the second. Symmetry and Schur's bound give
the full operator consequences. No higher coefficient is required to be
positive and no hard comparison gate is differentiated in a bridge source.
\end{proof}

V19's same positive shell return is therefore small, at its stipulated
logarithmic radius, throughout \emph{every fixed polynomial} width regime:
subtract its unchanged local-error decomposition from
\eqref{f16v22:row-asymptotics}. The resulting norm is
$C_{J,p,A_0,r}\gamma^{-(J+1)/2}$. This is an enlargement of V21's
$B\le c_r\gamma$ closure, not a redefinition of that return or of V12's
different one-common-exterior cylinder Gram.

\subsection{What remains spatial and what the comparison has not changed}

The new conditional comparison removes an intrinsic native--Gaussian error
from a native--native boundary question. Its two actual costs are
\[
           \sqrt{\gamma k}e^{-a\ell}
                  \quad\text{and}\quad k(1+\ell)^4e^{-c_*\gamma}.
\]
The first is a boundary-response estimate within one normalized nonlinear
comparison law. The second is the probability of leaving its finite buffer,
under the complete conditional at the actual exterior. An artificial
convex extension, a new physical action, and an assumed full-native
functional inequality are not used.

The local-core theorem is uniform in ambient volume. Its complete-slice
applications are not: at fixed $\gamma$, sufficiently large $B$ makes the
exception term ineffective. The exponential-width theorem has a buffer of
order $\gamma$ and only a relative all-observable channel conclusion at that
width. Finite-order source-bank information cannot be used on exponentially
large slices without a new norm estimate. The polynomial-width expansion
keeps the exponent $p$ fixed before its coupling limit; it is not permissible
to choose $p$ from the instantaneous unrestricted volume and ignore the
resulting constants and thresholds.

Likewise, for a fixed local core, sending $\ell$ to infinity at fixed
coupling does not make our bound tend to zero: the counted finite-buffer
exception eventually dominates. The next spatial task is to replace that
whole-buffer union by a \emph{connected propagation} estimate or a compatible
local dynamics with its actual update cost. The uniform native kernel and
its exponentially small exception are inputs to such an argument, not its
conclusion. Neither pairwise maximal couplings nor separate shell tests
supply an unproved grand coupling or tensorization through a shared exterior.

The original local resampling gap remains unproved. All statements concern
the prescribed bounded-bridge fast law. They do not remove the bridge window,
perform the physical Gauss reduction, prove an infrared contact floor, identify
a conditional transition with the physical transfer at its same top, or
calibrate physical time. The interacting continuum construction and a
physical Yang--Mills mass gap retain their separate obligations.

\subsection{Diagnostics, exact preservation, and analytic status}

The accompanying script is
\begin{center}\small\path{verify_ym_f16_v22_native_boundary.py}.\end{center}
It checks the normalized marginal mixed-derivative signs, covariance
localization algebra on exact Gaussian fixtures, conditional restriction and
posterior mixture identities on finite positive nonlinear laws, graph-buffer
depths and growth, and the polynomial/exponential-width and time-row power
bookkeeping. The new diagnostic passed 258 exact rational, integer, and symbolic assertions.
Its output distinguishes finite algebra from the analytic estimates; no finite
fixture is a native compact-integral enclosure. The available V4--V21
diagnostics were rerun without source changes and all eighteen passed; their
actual outputs and hashes are included in the package.

The cumulative source changes only the single V21 title version and inserts
this exact appendix before its final document-ending command. Removing that
insertion and restoring the title recovers V21 byte for byte. The package
records the initial input hashes, recovery, compilation, and rendered-page
inspection. These checks do not certify the inherited analytic archive. The
new deductions explicitly rely on the identified pinned specification,
V20 preparation, and product-ball covariance realization. In particular the
small-angular radius, comparison constants, and coupling thresholds are
proved to exist by the estimates above, not numerically enclosed or formally
verified. The next scheduled companion checkpoint remains V25.

\begin{firewall}
V22 proves a uniform complete-native buffered boundary comparison with an
exponentially small finite-buffer exception and no native--Gaussian power
floor. It gives arbitrarily strong fixed-power channel contraction for every
polynomial-size core. Exact temporal composition then proves fixed-weight
source covariance and memory asymptotics for every fixed polynomial spatial
width, and a separate all-observable temporal channel bound for a specified
exponential width. It does not establish unrestricted spatial volume at fixed
coupling, a gap in the original local sampler clock, a physical transfer gap,
or an interacting continuum mass gap. Every native normalizer and original
source remains unchanged; the small-angular restriction is paid and removed
from the target law, not silently imposed on it.
\end{firewall}
% END F16 V22 APPEND


% BEGIN F16 V23 APPEND -- 2026-09-17
% Insert immediately before the final document-ending command in V22.
% No predecessor theorem, probability law, source, or bibliography is edited.

\clearpage
\section{V23: common regeneration and exponentially wide coherent histories}
\label{f16v23:context}

V22 supplies a native comparison with an exponentially small finite-buffer
exception. Two consequences have not yet been extracted. First, a comparison
of densities, rather than only their total variation, supplies \emph{one common
subprobability measure for every compact boundary}. This constructs a genuine
simultaneous coupling, including uncountably many incoming boundary values.
It does not by itself control the dependence of the filled buffer. Second,
the same exponential error, combined with the exact two-shell identity of V19,
gives a native spatial covariance envelope with a \emph{nonperturbative} floor.
Only that floor, not the finite-order local error, pays the size of a coherent
source bank. Exponentially large banks can therefore be treated without an
expansion order growing with the volume.

These inputs close the source-asymptotic question that V22 explicitly left
open at exponential spatial width. There are fixed $\beta_r,a_r>0$ for which
all fixed orders of the complete covariance and memory hold in an exponentially
weighted temporal operator-row norm when $B\le e^{\beta_r\gamma}$, for every
duration. The width and temporal exponents are fixed independently of expansion
order; error constants and entry thresholds may depend on that order. This is
still not arbitrary spatial volume at fixed coupling. The all-boundary
regeneration is a new input for that remaining spatial problem, not a proof
that its overlapping buffers have already been composed.

\subsection{Source review and the distinction between a pair coupling and a common part}

The complete V22 appendix and analytic audit were read. The inputs used below
are its restricted mixed-normalizer identity, the componentwise core-score
gradient estimate, and its actual full-conditional buffer tail. V16 supplies
native source moments, V17 fixed-order entry coefficients, V19 the exact local
shell identity, and V21--V22 the actual temporal Markov composition. No global
native sampler gap is assumed. All conclusions are deductions from these
identified analytic inputs, not independent certification of their full
historical proofs.

Two targeted semantic searches of the f14/f15/Grand Tensor sources returned
no results. Mounted text searches and identified passages were then examined.
The monograph already gives a finite-card positive shell minorization and
warns that its constants may collapse with the carrier. F15 V22 already uses
pairwise common-part couplings on separated cores. Neither general mechanism
is claimed new. Here the common component is \emph{quantitatively near one},
uniform over the entire compact boundary family, and is constructed before any
boundary is sampled. The current calculation does not equate the monograph's
physical shell channels with the present fixed-bridge fast kernels.

For classical attribution, Spinka,
\href{https://arxiv.org/abs/1803.10578}{\emph{Finitary codings for spatial mixing
Markov random fields}}, Section 2.2, distinguishes a grand coupling from
separate optimal pair couplings. Its common-overlap formula and the associated
warning, as well as the scope of Theorems 1.2 and 4.1, were examined. Its
finite-valued, translation-invariant finitary-coding theorem is \emph{not}
imported for this inhomogeneous continuous law. The common randomization used
here is given directly by finite-dimensional measurable kernels. The
memoranda's separate ledgers for power gain and locality cost remain binding;
no Navier--Stokes theorem is an analytic premise. The next companion checkpoint
remains V25.

\subsection{A common native density floor, not an inference from total variation}

Keep V22's notation: $\mathcal A\subset W\subset\Lambda$ is a core of $k$
complete unmoving cells, $W$ is its $\ell$-neighbourhood in $\mathcal G$,
$C=\partial_{\mathcal G}W$, and
\[
 d_{\mathcal G}(\mathcal A,\Lambda^c)\ge\ell+D_*+2,\qquad \ell\ge3.
\]
The actual full conditional is $P_\Lambda^q$, its core kernel is
$K_{\Lambda,A}(q,\cdot)$, and its prescribed bridge history satisfies
$\max_{i,e}|y_{i,e}|\le r$. All compact exterior values $q$ are admitted in
V20's explicit principal-chart version. The restricted native row
$\pi_{W,A}^z$ and reference $\lambda_{W,A}=\pi_{W,A}^{0}$ use the same
whole-cell comparison balls of radius $\eta_*\sqrt\gamma$. They are not
Gaussian laws, and this event is not V4's original $\mathcal O_R$.

\begin{proposition}[Restricted density ratios with their output cost]
\label{f16v23:density-ratio}
At V22's fixed-$r$ entry threshold, there is $C_r<\infty$ such that, for
all small collar values $z,z'$ and all core points in their common comparison
domain,
\begin{equation}
 e^{-D_{\gamma,k,\ell}}\le
 \frac{d\pi_{W,A}^z}{d\pi_{W,A}^{z'}}(x)
 \le e^{D_{\gamma,k,\ell}},
 \qquad D_{\gamma,k,\ell}=C_r\gamma k e^{-a\ell}.
 \label{f16v23:ratio}
\end{equation}
Here $a>0$ may be decreased once from V22 and is independent of size and
coupling. If the collar is empty, take $D_{\gamma,k,\ell}=0$.
\end{proposition}
\begin{proof}
Use the \emph{same} normalized density $p_s$ along the collar segment and
$q_s(x)=\mathbb E_{\pi_{W\setminus A}^{x,z_s}}\partial_sV_W^{z_s}$ as in
\eqref{f16v22:score-identity}. That identity and
\eqref{f16v22:marginal-score-gradient} give
\[
 \left|\partial_{x_i}\log\frac{p_1(x)}{p_0(x)}\right|
 =\left|\int_0^1-\partial_{x_i}q_s(x)\,ds\right|
 \le C_r\sqrt\gamma e^{-a\ell}.
\]
There are at most $36k$ real core coordinates. Join any two core points by
their segment inside the product of whole-cell balls. Each scalar displacement
is at most $2\eta_*\sqrt\gamma$, so the oscillation of the log ratio is at
most $C_r\gamma k e^{-a\ell}$. Both densities are normalized and strictly
positive on the common compact comparison domain. Hence the minimum of the
log ratio is at most zero and its maximum at least zero: otherwise its
exponential could not have integral one against $p_0$. Its absolute value is
therefore at most its oscillation. This proves \eqref{f16v23:ratio}.
Both marginal normalizers, the core Haar factors, and the pinned-complement
covariance used in the gradient estimate remain those of V22. This proof
uses no functional inequality for the complete native law.
\end{proof}

Define
\begin{equation}
 u_{\gamma,k,\ell}=C_r k(1+\ell)^4e^{-c_*\gamma},\qquad
 \alpha_{\gamma,k,\ell}=(1-u_{\gamma,k,\ell})_+
                                      e^{-D_{\gamma,k,\ell}}.
 \label{f16v23:alpha}
\end{equation}
The constants can be enlarged to include the actual V22 buffer-tail bound.
Always $0\le\alpha\le1$ and
\begin{equation}
 1-\alpha\le\min\{1,u_{\gamma,k,\ell}+D_{\gamma,k,\ell}\}.
 \label{f16v23:failure}
\end{equation}

\begin{theorem}[One common submeasure for every compact exterior]
\label{f16v23:minorization}
For every Borel core set $E$ and every compact exterior $q$,
\begin{equation}
 \boxed{\qquad
 K_{\Lambda,A}(q,E)\ge\alpha_{\gamma,k,\ell}\lambda_{W,A}(E).
 \qquad}
 \label{f16v23:common-part}
\end{equation}
The reference and weight are independent of $q$. This is an inequality of
measures, not just a total-variation statement. The target remains the
complete native marginal. No good-event probability or conditional partition
function is replaced by a Gaussian denominator.
\end{theorem}
\begin{proof}
Let $G$ be V22's event that every cell in $W\cup C$ lies in its comparison
ball. The full conditional satisfies $P_\Lambda^q(G)\ge1-u$ for every $q$.
Conditional on $G$ and the actual exterior of $W$, its core law is exactly a
restricted row $\pi_{W,A}^z$. The mixing posterior retains all the original
conditional partition functions, as in V22. Every such row dominates
$e^{-D}\lambda_{W,A}$ by \eqref{f16v23:ratio}. Therefore
\[
 P_\Lambda^q(\zeta_A\in E,G)
 \ge P_\Lambda^q(G)e^{-D}\lambda_{W,A}(E)
 \ge(1-u)_+e^{-D}\lambda_{W,A}(E).
\]
Discard only the nonnegative contribution of $G^c$. This proves the measure
inequality on every Borel set. The empty-collar case has a constant restricted
row and is covered with $D=0$. The inequality is trivial but valid when $u\ge1$.
\end{proof}

The price $\gamma k$ in a log-density oscillation is larger than the
$\sqrt{\gamma k}$ in V22's TV comparison. They prove different statements.
For fixed $k\le A_0\gamma^p$ and any fixed $s>0$, choosing
\[
 \ell=3+\left\lceil a^{-1}(s+p+2)\log\gamma\right\rceil
\]
makes $1-\alpha\le\gamma^{-s}$ at a fixed threshold. There is still no
unrestricted-output assertion at fixed coupling.

\subsection{A simultaneous native regeneration map, with its buffer retained}

\begin{theorem}[All-boundary coalescence and exact full-block filling]
\label{f16v23:grand}
There is a measurable random map with independent uniform seeds such that:
for each $q$ its output on $\Lambda$ has exactly the law $P_\Lambda^q$;
and on a measurable seed event of probability $\alpha$, its core output is
\emph{the same for every compact boundary $q$ simultaneously}, with law
$\lambda_{W,A}$. The probability is uniform also when the incoming boundary
is any measurable function of previously exposed randomness, provided the
new seeds are independent of it.

Applying this map as a full $\Lambda$ conditional resampling preserves the
original native probability. It does not assert that the filled buffer is
independent of $q$ on the coalescence event, or that replacing only the core
and keeping the old buffer would preserve that probability.
\end{theorem}
\begin{proof}
For $0\le\alpha<1$ define the exact residual kernel
\[
 K^{\rm res}(q,dx)
       =\frac{K_{\Lambda,A}(q,dx)-\alpha\lambda_{W,A}(dx)}{1-\alpha}.
\]
It is a positive probability kernel by \eqref{f16v23:common-part}. It is
measurable in $q$: the full and restricted finite-dimensional integrals have
measurable densities. Let $H(q,x,dz)$ be the actual conditional law of the
remaining free cells given the core $x$ and the exterior $q$. Choose arbitrary
measurable values at null core parameters. The exact disintegration is
\[
 P_\Lambda^q(dx,dz)
 =\{\alpha\lambda_{W,A}(dx)+(1-\alpha)K^{\rm res}(q,dx)\}H(q,x,dz).
\]
Every kernel on these finite-dimensional standard Borel spaces can be sampled
by a measurable uniform-seed map. One direct construction orders its real
coordinates, uses measurable conditional distribution functions, and takes
their generalized inverses; values on null conditional sets can be fixed
arbitrarily. Apply this construction to $\lambda$, $K^{\rm res}$, and $H$.

With a fresh Bernoulli decision of success probability $\alpha$, draw the
core from $\lambda$ on success and from $K^{\rm res}(q)$ on failure, then
fill the rest using $H(q,x)$. Use the \emph{same} success decision and the
same $\lambda$ seed for all $q$. On success the equality of core outputs is
pathwise for all $q$, so no uncountable intersection of separately almost-sure
statements is needed. The final buffer seed is allowed to depend on $q$ through
its sampling map. Integrating the seeds gives exactly the displayed
conditional disintegration for every $q$. If $\alpha=1$, omit the residual
branch; if $\alpha=0$, the claimed common event is empty.

Conditional resampling of the entire $\Lambda$ from its actual law preserves
$P$ by the tower property. A deterministic sequence of such updates, or an
update schedule chosen independently of the current field, therefore also
preserves it. A field-dependent schedule would need a separate invariance
argument. Fresh decisions in a sequence have conditional failure probabilities
$1-\alpha_t$ given the past; for deterministic weights, any prescribed family
of failures has probability $\prod_t(1-\alpha_t)$. This is a statement about
algorithmic labels, not independence of native error fields. The buffer
conditional $H$ can still transmit incoming information even on success.
\end{proof}

For separated blocks the maps can be used independently after fixing their
one common exterior. Their success labels are then independent of that
exterior, unlike a posterior estimate that merely says an event is usually
good. For overlapping updates, use fresh seeds in their actual order; do not
replace the resulting joint law by independent simultaneous buffer outputs.
These are genuine common random maps, but not yet a subcritical exploration
of all their dependencies.

There is a simple reason not to infer this theorem from V22's TV bound alone.
A positive regularization of the standard grand-coupling example in Spinka's
Section 2.2 has rows on $\{1,\ldots,N\}$
\[
 k_i(j)=\frac{\theta}{N}
              +\frac{1-\theta}{N-1}1_{\{j\ne i\}},\qquad0<\theta<1.
\]
Each pair has TV distance $(1-\theta)/(N-1)$, tending to zero with $N$.
Nevertheless the mass of \emph{any} submeasure common to all rows is at most
$\sum_j\min_i k_i(j)=\theta$. The density floor above, rather than pairwise
closeness alone, proves the near-unit simultaneous coalescence. This is a
classical coupling warning, not a counterexample to the native specification.

\subsection{A nonperturbative spatial envelope for the actual primitive sources}

Return to V19's ambient cell distance
\[
 d_*((t,b),(s,c))=|t-s|+d_{\mathbb T_m^3}(b,c),\qquad
 \varepsilon=\gamma^{-1/2},\quad\mathcal N=B(n+1).
\]
Every edge of $\mathcal G$ has $d_*$ length at most a fixed $r_0$. Every
primitive error has a carrier of fixed $d_*$ radius $r_s$ and a uniformly
bounded number of cells. Let $F_u=e_u^{\rm prim}-\mathbb E_Pe_u^{\rm prim}$,
$K_\gamma=\operatorname{Cov}_P(e^{\rm prim})$, and
\[
 m_{u,R}=\mathbb E_P[F_u\mid\zeta_{I_{u,R}^c}],
 \qquad I_{u,R}=\{c:d_*(a(u),c)\le R\}.
\]
These are precisely the centered shell messages of V19, with their actual
conditional normalizers. V16 gives $\operatorname{Var}_P F_u\le C_r\varepsilon^2$.

\begin{proposition}[Source screening with an exponentially small coupling floor]
\label{f16v23:nonpert-screen}
There are fixed $R_0,b,c,C_r>0$ such that, for all sufficiently large
fixed-$r$ coupling and every $R\ge R_0$,
\begin{equation}
 \sup_u\|m_{u,R}\|_2^2
 \le C_r\varepsilon^2
       \min\{1,\gamma e^{-bR}+(1+R)^4e^{-c\gamma}\}.
 \label{f16v23:screen}
\end{equation}
A ball covering the whole system has centered message zero. There are also
fixed $b_1,\kappa>0$ such that, for every pair of primitive marks,
\begin{equation}
 |(K_\gamma)_{uv}|
 \le C_r\varepsilon^2
       \left[\gamma e^{-b_1d_*(a(u),a(v))}+e^{-\kappa\gamma}\right].
 \label{f16v23:primitive-envelope}
\end{equation}
The same bound, after one fixed decrease of $b_1$, holds for each entry of
$\Gamma_\gamma=TK_\gamma T^{\mathsf T}$ in its original common-source anchors.
Constants do not depend on volume, duration, or the marks.
\end{proposition}
\begin{proof}
Use the complete primitive carrier as the observed core; its cell count is
uniformly bounded. Its distance in $\mathcal G$ to $I_{u,R}^c$ is at least
$(R-r_s)/r_0$. Choose
\[
 \ell=\left\lfloor(R-r_s)/r_0\right\rfloor-D_*-3.
\]
Choose $R_0$ also larger than every primitive carrier allowance. For
$R\ge R_0$ this is at least three and satisfies the exact buffer geometry
of \eqref{f16v23:common-part}. Its common component has failure probability
at most $C_r[\gamma e^{-bR}+(1+R)^4e^{-c\gamma}]$, by
\eqref{f16v23:failure}. A common component of mass $\alpha$ bounds the
kernel diameter by $1-\alpha$. V20's proved all-incoming-law channel lemma,
applied with the \emph{actual} exterior marginal and the fixed observable
$F_u$, bounds the squared conditional norm by $(1-\alpha)\operatorname{Var}F_u$.
This proves \eqref{f16v23:screen}. No native fourth moment at an adversarial
boundary or inverse of a native generator is used.

For two anchors at distance $d>2R+\rho$, V19's exact two-shell identity is
\[
 \mathbb E_P F_uF_v=\mathbb E_P m_{u,R}m_{v,R}.
\]
Here $\rho$ is one fixed ambient interaction/carrier allowance. The first
message is a function of its finite exterior word collar, which lies outside
the second ball; two conditional expectations justify the identity. It does
not assert that the collars are independent. Choose, whenever $d$ is large
enough,
\[
 R=\min\left\{\left\lfloor\frac{d-\rho-1}{3}\right\rfloor,
                                      \lceil\gamma\rceil+R_0\right\}.
\]
This obeys $d>2R+\rho$. Before the radius reaches its cap, the first term of
\eqref{f16v23:screen} is bounded by $C_r\gamma e^{-b_1d}$ for a fixed
$b_1>0$. The second is at most $C_r(1+\gamma)^4e^{-c\gamma}$, which is
bounded by a fixed exponential $C_re^{-c\gamma/2}$. After saturation, both
terms are at most $C_re^{-\kappa\gamma}$ for some fixed $\kappa>0$.
Cauchy--Schwarz in the exact two-shell identity proves
\eqref{f16v23:primitive-envelope}. The finitely many shorter distances are
covered by the original variance bound, enlarging $C_r$. Empty exteriors
cause no exception.

The original deterministic synthesis $T=(I,\mathsf P^{\mathsf T})$ has
uniform exponentially weighted absolute rows in these anchor labels, as
used in V18--V19. Decrease $b_1$ below its weight. The triangle inequality
for $d_*$ bounds
\[
 \sum_{u,v}|T_{\alpha u}T_{\beta v}|
                       e^{-b_1d_*(a(u),a(v))}
                    \le C e^{-b_1d_*(a(\alpha),a(\beta))}.
\]
The constant floor is preserved by the two absolute row sums. This proves the
full-source assertion. The spatial lift is not falsely assigned finite support.
\end{proof}

In particular, at the \emph{single} linear radius
$R_\gamma=\lceil\gamma\rceil+R_0$, decrease $\kappa$ once so that
\begin{equation}
 \sup_u\|m_{u,R_\gamma}\|_2^2\le C_r\varepsilon^2e^{-\kappa\gamma}.
 \label{f16v23:linear-screen}
\end{equation}
Unlike an arbitrarily high fixed-order floor, this estimate has specified
exponential dependence on coupling. It still does not tend to zero as
$R\to\infty$ at fixed coupling by the displayed bound alone.

\subsection{Only the nonperturbative floor pays coherent source size}

The next estimate uses the same normalized coefficients $G_q$ and
$M_q=-\operatorname{Off}_{>1}G_q$ as V10. All fixed coefficient rows decay
exponentially in $d_*$, and V17's entry expansion is uniform in ambient volume.
Let $\|v\|_1$ be the counting coefficient norm on the original canonical
common-source coordinates, without volume normalization.

\begin{theorem}[An operator error plus an exponentially small coherent remainder]
\label{f16v23:mixed}
For every fixed $J\ge2$ there are $C_{J,r},\gamma_{J,r}<\infty$ such that,
for every admitted finite history with $\gamma\ge\gamma_{J,r}$ and all real
or complex source vectors $v,w$,
\begin{equation}
 \boxed{
 \left|v^*\left(\Gamma_\gamma-\sum_{q=2}^J\gamma^{-q/2}G_q\right)w\right|
 \le C_{J,r}\gamma^{-(J+1)/2}\|v\|_2\|w\|_2
       +C_r\gamma^{-1}e^{-\kappa\gamma}\|v\|_1\|w\|_1.}
 \label{f16v23:mixed-bound}
\end{equation}
The same estimate holds for $\mathsf M_\gamma$ and $M_q$, with enlarged fixed
constants. The exponent $\kappa>0$ is independent of $J$. Consequently, for
every fixed $0<\beta<\kappa$ and any source bank $\mathcal A$ with
$|\mathcal A|\le e^{\beta\gamma}$,
\begin{equation}
 \left\|\chi_{\mathcal A}
   \left(\Gamma_\gamma-\sum_{q=2}^J\gamma^{-q/2}G_q\right)
       \chi_{\mathcal A}\right\|_{\rm op}
                  \le C_{J,\beta,r}\gamma^{-(J+1)/2},
 \label{f16v23:exponential-bank}
\end{equation}
and likewise for memory. The ambient system may be arbitrarily larger and
all source positions are allowed. The fixed order does not grow with the bank.
\end{theorem}
\begin{proof}
Choose a fixed $b_2>0$ below the envelope and coefficient decay exponents, so
polynomial graph growth gives an exponential absolute row tail at exponent
$b_2$ for either exponential envelope. Put
\[
 D_J=\left\lceil (J+3)b_2^{-1}\log(\varepsilon^{-1})\right\rceil+R_0,
 \qquad K=J+1.
\]
On pairs with $d_*\le D_J$, apply V17 at order $K$. The number of entries
per row and column is at most $C(1+D_J)^4$. The entry remainder therefore
has absolute row and column sums at most
$C_{J,r}(1+D_J)^4\varepsilon^{J+2}\le C_{J,r}\varepsilon^{J+1}$.
The extra coefficient $\varepsilon^{J+1}G_{J+1}$ has that same global row
bound without a distance count.

On pairs with $d_*>D_J$, bound the native covariance by
\eqref{f16v23:primitive-envelope} in its full-source form and the finitely
many desired coefficients by their inherited exponential rows. Their
exponential parts have absolute row and column sums at most
\[
 C_{J,r}(\varepsilon^2\gamma+\varepsilon^2)e^{-b_2D_J}
                         \le C_{J,r}\varepsilon^{J+1}.
\]
The remaining entry bound is precisely
$C_r\varepsilon^2e^{-\kappa\gamma}$. The nonnegative envelopes with small
row and column sums bound a bilinear form by their Schur operator norm times
$\|v\|_2\|w\|_2$. The constant entry floor costs instead
$C_r\varepsilon^2e^{-\kappa\gamma}\|v\|_1\|w\|_1$. This proves
\eqref{f16v23:mixed-bound} without an unrestricted sum of that floor.
All estimates used absolute coefficients, so complex vectors cause no change.

For memory the exact identity masks entries of this same error. It preserves
both nonnegative entry envelopes and their absolute row bounds; hence the
identical argument applies. It is not assumed to preserve positivity of a
matrix. On a bank of size $k$, $\|v\|_1\le\sqrt{k}\|v\|_2$. For
$k\le e^{\beta\gamma}$ the extra term is at most
$C_r\gamma^{-1}e^{-(\kappa-\beta)\gamma}\|v\|_2\|w\|_2$, smaller than
any specified fixed algebraic power at a fixed-$J,\beta,r$ threshold. This
proves \eqref{f16v23:exponential-bank}.
\end{proof}

At leading order, arbitrary banks of that exponential size obey
\[
 \|\chi_{\mathcal A}(\Gamma_\gamma-\gamma^{-1}G_2)\chi_{\mathcal A}\|
 \le C_{\beta,r}\gamma^{-3/2}.
\]
The same holds for memory. Their actual covariances are
$O_{\beta,r}(\gamma^{-1})$ in operator norm. This is not derived by
multiplying an algebraic entry error by an exponentially large bank.

A sharper necessary condition on a remaining coherent obstruction follows.
For $v\ne0$ put $k_{\rm eff}(v)=\|v\|_1^2/\|v\|_2^2$. For unit vectors,
\begin{equation}
 \left|\gamma v^*\Gamma_\gamma v-v^*G_2v\right|
 \le C_r\gamma^{-1/2}+C_re^{-\kappa\gamma}k_{\rm eff}(v).
 \label{f16v23:effective-support}
\end{equation}
An order-one failure along $\gamma\to\infty$ must therefore have
$k_{\rm eff}(v_\gamma)\ge c\,e^{\kappa\gamma}$ eventually. This improves
V17's superpolynomial necessary delocalization, but does not exclude such
vectors in unrestricted volume. The same conclusion applies to memory.

\subsection{The exponential-width temporal channel now has a coherent input}

V22's exponential-width channel theorem supplies fixed constants
$\beta_{\rm ch},d_{\rm ch},C_{\rm ch}>0$ and actual buffers $b_\gamma^{\rm exp}$
with
\begin{equation}
 B\le e^{\beta_{\rm ch}\gamma},\qquad
 c_{\rm ch}\gamma\le b_\gamma^{\rm exp}\le C_{\rm ch}\gamma,
 \qquad \delta(K_{i,i+b_\gamma^{\rm exp}})\le e^{-d_{\rm ch}\gamma}.
 \label{f16v23:channel-input}
\end{equation}
The positive lower constant $c_{\rm ch}$ follows from V22's explicit linear
choice of its comparison radius. Every continuation message and both endpoints
remain in these native transitions. Its past/future correlation bound yields
a fixed $\lambda_{\rm ch}>0$ with relative correlation at most
$e^{-\lambda_{\rm ch}g}$ whenever the gap $g\ge2b_\gamma^{\rm exp}$.
These are inherited serial-chain facts, not a new composition across spatial
pieces.

Fix once
\begin{equation}
 0<\beta_r\le\min\{\beta_{\rm ch},\kappa/8\},\qquad
                         B\le e^{\beta_r\gamma}.
 \label{f16v23:width}
\end{equation}
At sufficiently large coupling a full single-time source bank, of size $9B$,
is within the exponential-bank theorem with some exponent strictly below
$\kappa$. Thus its coherent variance is at most
$C_r\varepsilon^2\|v\|_2^2$. The full Haar source has its original finite
temporal carrier allowance $\sigma$, although its spatial lift is nonlocal.
Combining this coherent input with \eqref{f16v23:channel-input} gives
\begin{equation}
 \|\Gamma_{\gamma,ij}\|_{\rm op}
 \le C_r\varepsilon^2e^{-\lambda_{\rm ch}(|i-j|-2\sigma)}
 \quad\text{when }|i-j|-2\sigma\ge2b_\gamma^{\rm exp}.
 \label{f16v23:far-time}
\end{equation}
At every lag, spatial row summation of the new native envelope also gives
\begin{equation}
 \|\Gamma_{\gamma,ij}\|_{\rm op}
 \le C_r\varepsilon^2
              \bigl[\gamma e^{-b_3|i-j|}+B e^{-\kappa\gamma}\bigr]
 \label{f16v23:middle-time}
\end{equation}
for a fixed $b_3>0$. The spatial exponential row sum is bounded independently
of torus size; the constant floor pays exactly $B$. These two estimates have
different roles. Neither replaces the coherent short-time variance by a sum
of $B$ independent-source variances.

\begin{theorem}[All fixed orders on an exponentially wide complete history]
\label{f16v23:exponential-history}
There are fixed $a_r>0$ and $\beta_r>0$, independent of expansion order, such
that for every fixed $J\ge2$ there are finite $C_{J,r},\gamma_{J,r}$ such
that, for all $\gamma\ge\gamma_{J,r}$ in \eqref{f16v23:width} and every duration,
\begin{equation}
 \boxed{
 \begin{aligned}
 &\left\|\Gamma_\gamma-\sum_{q=2}^J\gamma^{-q/2}G_q
                                          \right\|_{{\rm tr},a_r}\\
 &\quad+\left\|\mathsf M_\gamma-\sum_{q=2}^J\gamma^{-q/2}M_q
                                          \right\|_{{\rm tr},a_r}
                  \le C_{J,r}\gamma^{-(J+1)/2}.
 \end{aligned}}
 \label{f16v23:row-asymptotics}
\end{equation}
Here $\|A\|_{{\rm tr},a_r}=\sup_i\sum_j e^{a_r|i-j|}\|A_{ij}\|_{\rm op}$
is the \emph{temporal operator-row norm}, not trace norm. Each block includes
the entire spatial common-source space. There is no normalization by
$B(n+1)$ or restriction on the total time-source count. In particular the
full covariance and memory are $O_r(\gamma^{-1})$ in this norm and their
leading-coefficient errors are $O_r(\gamma^{-3/2})$.
\end{theorem}
\begin{proof}
Let $a_*>0$ be a fixed temporal coefficient exponent, so that each fixed
$G_q$ has finite row norm at exponent $a_*$. This follows from V10's full
space--time absolute rows, as already explained in V21. Put
$D_1=2b_\gamma^{\rm exp}+2\sigma$ and choose a fixed $C_1$ with
$D_1\le C_1\gamma$. Choose once
\begin{equation}
 0<a_r\le\min\{\lambda_{\rm ch}/4,b_3/4,a_*/4,\kappa/(8C_1)\}.
 \label{f16v23:weight}
\end{equation}
For the given $J$, set
\[
 D_0=2\sigma+\left\lceil(J+5)a_r^{-1}\log(\varepsilon^{-1})\right\rceil,
 \qquad K=2J+8.
\]
At a fixed-$J,r$ threshold $D_0<D_1$. Split the temporal sum into three bands.

For $|i-j|\le D_0$, apply \eqref{f16v23:mixed-bound} at order $K$ to unit
vectors supported in time blocks $i,j$. Its error block is bounded by
\[
 C_{K,r}\varepsilon^{K+1}
                      +C_r B\varepsilon^2e^{-\kappa\gamma}.
\]
There is \emph{no} factor $B$ on the algebraic term. Summing that term with
its temporal weight costs
\[
 C_{J,r}(1+D_0)e^{a_rD_0}\varepsilon^{K+1}
 \le C_{J,r}(1+\log\varepsilon^{-1})\varepsilon^{K-J-4}
 \le C_{J,r}\varepsilon^{J+1}.
\]
The coefficients of orders $J+1,\ldots,K$ cost
$C_{J,r}\varepsilon^{J+1}$ in their \emph{entire} weighted row norm.
They are not charged separately for every near lag.

On $D_0<|i-j|\le D_1$, use \eqref{f16v23:middle-time} for the native matrix
and the coefficient row tails. The exponential part has weighted sum at most
\[
 C_{J,r}\varepsilon^2\gamma e^{-(b_3-a_r)D_0}
               +C_{J,r}\varepsilon^2e^{-(a_*-a_r)D_0}
                         \le C_{J,r}\varepsilon^{J+1}.
\]
The constant floor on \emph{both} bands up to $D_1$ has total cost at most
\begin{equation}
 C_r\varepsilon^2(1+D_1)e^{a_rD_1}B e^{-\kappa\gamma}
 \le C_r\varepsilon^2(1+C_1\gamma)e^{-3\kappa\gamma/4}.
 \label{f16v23:middle-floor}
\end{equation}
Here $\beta_r\le\kappa/8$ and $a_rC_1\le\kappa/8$ were used. This is
smaller than every prescribed fixed power of $\varepsilon$; its entry
threshold may depend on $J$.

For $|i-j|>D_1$, the genuine temporal estimate \eqref{f16v23:far-time}
has weighted tail at most
\[
 C_r\varepsilon^2e^{-(\lambda_{\rm ch}-a_r)(D_1-2\sigma)}.
\]
Since $b_\gamma^{\rm exp}\ge c_{\rm ch}\gamma$, this is exponentially small
in $\gamma$. The coefficient tail is likewise exponentially small. If the
actual duration is shorter, the missing bands are simply empty. The three
bounds prove the covariance assertion. The unchanged contact identity
$\mathsf M_\gamma=-\operatorname{Off}_{>1}\Gamma_\gamma$ removes entire
time blocks and contracts this absolute row norm, giving memory. The exponent
\eqref{f16v23:weight} was fixed before $J$ was chosen. All source and
measure-correction coefficients remain those of the original native law.
\end{proof}

\subsection{One linear shell has an exponentially small coherent return}

The following identifies the old positive object more strongly than an
order-dependent logarithmic radius. It uses exactly
\[
 \mathsf R^{\rm sh}_{\gamma,R}
 =T\operatorname{Cov}_P((m_{u,R})_u)T^{\mathsf T}
\]
from V19, not V12's different one-common-exterior cylinder return.

\begin{proposition}[Exponential return bound at one order-independent radius]
\label{f16v23:linear-return}
After a possible fixed decrease of $a_r,\beta_r$, there is $\kappa_r'>0$ such
that, with $R_\gamma=\lceil\gamma\rceil+R_0$,
\begin{equation}
 \|\mathsf R^{\rm sh}_{\gamma,R_\gamma}\|_{{\rm tr},a_r}
                      \le C_r\gamma^{-1}e^{-\kappa_r'\gamma}
 \label{f16v23:linear-return-bound}
\end{equation}
throughout \eqref{f16v23:width}, uniformly in duration. At unrestricted
ambient volume, its diagonal and trace density still obey
$C_r\gamma^{-1}e^{-\kappa\gamma}$, but that alone is not an operator bound.
\end{proposition}
\begin{proof}
By \eqref{f16v23:linear-screen}, every primitive message has variance at most
$C_r\varepsilon^2e^{-\kappa\gamma}$. Its actual collar is contained in the
space--time ball of radius $R_\gamma+\rho$. A time bank of such messages has
at most $30B$ entries, so its covariance operator norm is at most
$C_rB\varepsilon^2e^{-\kappa\gamma}$ by its trace. This use of a trace is
explicit and is now paid by an exponential coupling margin.

Two such time banks have separated temporal carriers once their time gap
exceeds $2(R_\gamma+\rho)$. Apply the actual exponential-width past/future
channel after a further gap $2b_\gamma^{\rm exp}$. Covariance
Cauchy--Schwarz handles the shorter lags. Write
$G_\gamma=2(R_\gamma+\rho)+2b_\gamma^{\rm exp}\le C_G\gamma$.
The resulting weighted row sum is at most
\[
 C_rB\varepsilon^2e^{-\kappa\gamma}
                         (1+G_\gamma)e^{a_rG_\gamma},
\]
provided $a_r\le\lambda_{\rm ch}/4$. Decrease $a_r$ so that
$a_rC_G\le\kappa/8$ and retain $\beta_r\le\kappa/8$.
The last expression is at most $C_r\varepsilon^2e^{-\kappa\gamma/2}$
after a fixed threshold increase. The original synthesis $T$ has a fixed
finite temporal range and bounded spatial block norms, hence a uniformly
bounded temporal row norm at this weight. Row-norm multiplication transfers
the bound through $T$ and $T^{\mathsf T}$. This proves the proposition.
At unrestricted volume, sum the scalar diagonal bounds and use bounded
absolute synthesis rows, or its operator norm and the trace ideal inequality,
to get the stated diagonal and trace-density bounds.
\end{proof}

\subsection{What remains spatial, and why common regeneration is not yet closure}

The new coherent remainder is
\[
 C_{J,r}\gamma^{-(J+1)/2}\|v\|_2\|w\|_2
       +C_r\gamma^{-1}e^{-\kappa\gamma}\|v\|_1\|w\|_1,
\]
not an algebraic error multiplied by the number of coordinates. It permits
exponentially large source banks and, using V22's actual temporal channels,
all-duration operator-row asymptotics at one fixed exponential spatial width.
The constants $\beta_r,a_r,\kappa$ are proved positive from the inherited
analytic bounds but have not been numerically enclosed or optimized. They
cannot be identified with a physical running-coupling coefficient.

At fixed $\gamma$, an unrestricted source vector can still have arbitrarily
large $\ell^1/\ell^2$ ratio, and $B$ can exceed $e^{\beta_r\gamma}$. The
constant floor cannot be summed over that class. The new grand coupling also
does not solve this by declaration. On a successful core refresh, the exact
buffer law $H(q,x)$ can retain dependence on the incoming boundary. Subsequent
overlapping updates can query precisely that buffer. Independent success
labels do not make those filled fields independent, and changing only a core
while retaining its old buffer is not the original conditional resampling.
A proof of subcritical \emph{dependency propagation}, or a paid comparison of
compatible block dynamics with the local clock, remains necessary.

Thus the simultaneous all-boundary map supplies an input explicitly missing
from V22, while the nonperturbative source estimate closes its separate
exponential-width covariance question. Neither is presented as unrestricted
spatial mixing. The original local sampler gap, V4's unchanged full/core
boundary derivatives, and V12's distinct spatial-cylinder posterior are not
silently redefined. Bridge-window removal, physical root Gauss restriction,
infrared contact coercivity, same-top physical transfer comparison, independent
physical-time calibration, and interacting continuum reconstruction retain
their separate obligations.

\subsection{Diagnostics, source preservation, and claim scope}

The accompanying diagnostic is
\begin{center}\small\path{verify_ym_f16_v23_regeneration_coherence.py}.\end{center}
It checks normalized density-ratio algebra, common submeasures and residual
kernels, simultaneous finite randomizations with actual conditional buffer
filling, exact preservation of a positive finite Gibbs law, and a failure of
core-only replacement. Further checks cover the two-shell covariance identity,
spatial-to-coherent envelopes, exponential support and three-band exponent
bookkeeping, and the distinction between pairwise and common overlap. The
diagnostic passes 291 exact assertions; its report records the families.
These are finite rational, integer, and symbolic diagnostics, not native compact-integral enclosures or
formal verification of the analytic estimates.

The cumulative V22 is preserved byte for byte except for one literal title
version and this exact terminal insertion. The package records predecessor
recovery, protected input hashes, build and rendering results, and the actual
reruns of available earlier diagnostics. No inaccessible historical checker
or external formalization is represented as run. The new deductions depend
on the identified V16--V22 analytic inputs, especially the preparation and
restricted product-ball covariance theorem. The entire archive has not been
independently recertified. The next scheduled companion checkpoint remains V25.

\begin{firewall}
V23 proves a uniform common native submeasure and a simultaneous all-boundary
core regeneration map, with the full buffer filled from its actual conditional.
It proves a native spatial covariance envelope with an exponentially small
coupling floor, mixed coherent-error estimates and fixed-order expansions on
exponentially large source banks. Exact temporal composition then gives
all-duration covariance and memory asymptotics in a fixed weighted temporal
operator-row norm for a specified exponential spatial width, and an
exponentially small positive return at one linear radius. It does not prove
unrestricted spatial volume at fixed coupling, subcritical composition of
all overlapping buffers, the original local sampler gap, a physical transfer
gap, or an interacting continuum Yang--Mills mass gap. All native normalizers,
compact fields and original sources remain unchanged.
\end{firewall}
% END F16 V23 APPEND

% BEGIN F16 V24 APPEND -- 2026-09-17
% Insert immediately before the final document-ending command in V23.
% No predecessor assertion, source, probability, or bibliography is edited.

\clearpage
\section{V24: a paid local-clock comparison and boundary-sensitive block contraction}
\label{f16v24:context}

V23 constructs a simultaneous core refresh but correctly retains dependence
of its filled buffer on the incoming boundary. Increasing its common-core
probability does not, by itself, control that dependence. We address the
\emph{update operator} before making another covariance expansion. For complete
time intervals the required boundary-sensitive coupling can be proved: its
expected number of affected slices is bounded by the channel buffer length,
not by the resampled interval length. Overlapping exact interval resamplings
then have a uniform spectral gap in a stated per-slice clock.

There is a second, separate calculation. Returning this operator to V14's
original rate-one balanced chord/port updates enlarges its port support by one
temporal bond. We keep that enlargement and prove the form comparison using
the actual complete conditional density. The resulting original-clock gap is
positive uniformly in duration in the admitted width regimes, but has an
explicit, pessimistic coupling/width cost. It is \emph{not} a positive gap
uniform along the weak-coupling limit. The last part identifies the local
boundary-response quantity needed to run the same construction spatially,
and shows why the coarse common-core certificate alone does not verify it.

\subsection{Dependencies, non-duplication, and which gap is under discussion}

The supplied V23 appendix and audit were read in full. Two targeted semantic
searches of the current lineage and supplied f14/f15/Grand Tensor files were
empty; mounted text searches and identified passages were then checked.
V15 already compares its gated three-port updates to the original clock; that
form-comparison principle is credited. Its gates neither cover the present
intervals nor prove their gap. The monograph's random-scan and blocking
countertheorems distinguish a per-cell clock from a single global update and
warn that taking a power of an unestimated transfer creates no gap. F14 V2's
same-clock local repair likewise concerns a different conditional and support.
These statements remain binding. This is a targeted audit, not independent
recertification of the archives; the next scheduled checkpoint remains V25.

Path coupling is a classical method; see Bubley--Dyer,
\href{https://doi.org/10.1109/SFCS.1997.646111}{\emph{Path coupling: A technique
for proving rapid mixing in Markov chains}}. Its primary abstract and metadata
were checked; no inaccessible full-paper theorem is invoked. Spinka,
\href{https://arxiv.org/abs/1803.10578}{\emph{Finitary codings for spatial mixing
Markov random fields}}, Section 4 and Theorems 4.1--4.3, makes a stronger
all-boundary-subset sensitivity requirement explicit for grand couplings.
The relevant definitions and contraction proof were reviewed. Its finite-valued,
translation-invariant coding theorem is not applied to this inhomogeneous
continuous law. We prove below the pair-coupling/spectral statement actually
used, directly on finite products of standard Borel spaces.

The analytic native input is V22's buffered \emph{temporal channel} under the
actual continuation law, including an arbitrarily fixed opposite endpoint.
V21 supplies that Markov factorization. The new spectral proof does not assume
a gap of any native generator, use V23's covariance expansion to infer one,
or substitute the Gaussian gap. The full-domain comparison to the local
clock uses only the exact coordinate transformation, product reference
factors and bounded local classical action. No small-field gate is imposed.

\subsection{A boundary-input ledger for exact conditional updates}

Let $\mu$ be a probability with strictly positive density relative to a product
of $N$ probability measures on standard Borel coordinate spaces. Null parameter
values can be assigned the explicit versions of the kernels when available.
For a coordinate block $D$, put
\[
 E_Df=\mathbb E_\mu[f\mid\omega_{D^c}],\qquad Q_D=I-E_D,
 \qquad \mathcal E_D(f)=\langle f,Q_Df\rangle_\mu.
\]
The positive generator for a finite indexed family of blocks with deterministic
rates $r_D>0$ is $\mathcal L_{\rm bl}=\sum_D r_DQ_D$. Repeated blocks may carry
separate indices and rates. These are \emph{full} conditional resamplings,
not replacement of a core while retaining its previous buffer.

Choose fixed positive coordinate weights $w_j$ and define the bounded measurable
weighted Hamming metric
\[
 d_w(x,y)=\sum_j w_j1_{\{x_j\ne y_j\}}.
\]
Use fixed measurable conditional-kernel versions on the full product carrier,
and require the following comparison for every such pair of parameters,
including the intermediate configurations in a sequence of coordinate changes.
For each $j\notin D$ suppose that when two exteriors differ only at $j$, a
coupling of the two \emph{entire} conditional outputs on $D$ has expected
inside distance at most $I_{D,j}$. Set $I_{D,j}=0$ when that exterior coordinate
cannot affect this kernel. This is a bound on a specified boundary change,
not the diameter of an arbitrary collection of core marginals. Define
\begin{equation}
 g=\inf_j\left\{\sum_{D\ni j}r_D
              -\frac1{w_j}\sum_{D\not\ni j}r_DI_{D,j}\right\}.
 \label{f16v24:influence-margin}
\end{equation}

\begin{theorem}[Boundary-sensitive block contraction in its stated clock]
\label{f16v24:general-gap}
If $g>0$, then on the complete $L^2(\mu)$ space
\begin{equation}
 \mathcal L_{\rm bl}\succeq g(I-\mathbb E_\mu),\qquad
 \operatorname{Var}_\mu f\le g^{-1}\sum_Dr_D\mathcal E_D(f).
 \label{f16v24:general-gap-bound}
\end{equation}
This holds on finite continuous product spaces as well as finite alphabets.
No global measurable grand coupling or prior spectral-gap premise is needed.
\end{theorem}
\begin{proof}
Let $R=\sum_Dr_D$ and $K=R^{-1}\sum_Dr_DE_D$, a self-adjoint Markov operator
preserving $\mu$. Consider configurations differing only at $j$. If $j\in D$,
the two exterior rows are identical, and the entire output can be coupled
identically: the distance drops from $w_j$ to zero. If $j\notin D$, the old
$w_j$ remains, and the expected additional distance is at most $I_{D,j}$.
Average the specified couplings over the chosen block. By
\eqref{f16v24:influence-margin}, the expected resulting distance is at most
$(1-g/R)w_j$. Here $0<g\le R$.

For a bounded measurable $f$, write
\[
 \operatorname{Lip}_w(f)=
 \max_j\sup_{x_{j^c}=y_{j^c},\ x_j\ne y_j}
                    \frac{|f(x)-f(y)|}{w_j}.
\]
It is bounded by $2\|f\|_\infty/\min_jw_j$. Changing one coordinate at a
time bounds $|f(x)-f(y)|$ by $\operatorname{Lip}_w(f)d_w(x,y)$. The pair
couplings therefore give
\[
 \operatorname{Lip}_w(Kf)\le(1-g/R)\operatorname{Lip}_w(f).
\]
Iteration and Poissonization, using $\mathcal L_{\rm bl}=R(I-K)$, give
\begin{equation}
 \operatorname{osc}(e^{-t\mathcal L_{\rm bl}}f)
       \le \left(\sum_jw_j\right)e^{-gt}\operatorname{Lip}_w(f).
 \label{f16v24:osc-decay}
\end{equation}
This argument needs only the existence of the pair couplings: no measurable
selection over all pairs is used to define the Markov operator.

For centered bounded $f$, invariance of $\mu$ bounds its semigroup $L^2$ norm
by the right side of \eqref{f16v24:osc-decay}. If the nonnegative self-adjoint
operator $\mathcal L_{\rm bl}$ had nonzero centered spectral subspace in
$[0,a]$, where $a<g$, density of bounded centered functions would give a test
$f$ with nonzero projection there. The spectral theorem would give a lower
bound $e^{-at}$ times its projected norm, contradicting
\eqref{f16v24:osc-decay} as $t\to\infty$. The finite constant $\sum_jw_j$ is
held fixed before this time limit; no volume-independent total-variation
prefactor is asserted. Constants are fixed, so this proves the operator
inequality and then the variance inequality for every $L^2$ function.
\end{proof}

The death term in \eqref{f16v24:influence-margin} records how often an already
different coordinate is actually refreshed. The birth term counts the
\emph{whole output}, including its buffer, for one changed incoming coordinate.
A high probability of refreshing only its central part does not supply that
birth term at a useful scale.

\subsection{A complete temporal interval has bounded boundary susceptibility}

Return to the unchanged native fixed-bridge probability $P=P_{\gamma,\mathbf y}$,
with $\max_{i,e}|y_{i,e}|\le r$. Its complete unmoving slice is $Z_i$,
$0\le i\le n$, as in \eqref{f16v21:chain-density}. A slice includes all $B$
cubes, and its size is \emph{not} counted as independent spins in the next
Hamming metric. Write $b\ge1$ for an integer channel buffer. We use the
following already proved input:
\begin{equation}
 \delta(K^{v}_{t,t+b})\le q\le\tfrac14
 \quad\hbox{whenever the output is at least $b$ from a fixed right boundary.}
 \label{f16v24:channel-input}
\end{equation}
The analogous backward statement holds with the left boundary fixed. Natural
free endpoints require no artificial boundary allowance. The superscript $v$
means the actual fixed right factor is retained inside every continuation
kernel; it is not a conditioning imposed after an estimate.

V22 proves \eqref{f16v24:channel-input} in two usable regimes. For every fixed
$p\ge0,A_0<\infty$, $B\le A_0\gamma^p$ permits $b=b_{\gamma,p,2}=O_{p,r}(\log\gamma)$
and $q\le\gamma^{-2}$ at its stated threshold. Separately,
$B\le\exp(\beta_{\rm ch}\gamma)$ permits $b=b_\gamma^{\rm exp}=O_r(\gamma)$ and
$q\le\exp(-d_{\rm ch}\gamma)$. Increase only those fixed thresholds so that
$q\le1/4$. No source-bank expansion is needed for this input.

\begin{proposition}[An exact interval coupling, including the near-boundary remainder]
\label{f16v24:interval-coupling}
Let $I=[a,c]\cap\mathbb Z\subseteq[0,n]$ be nonempty. If two full conditional
laws on $I$ differ only in the left boundary slice, with the right boundary
fixed or natural, there is a coupling of their \emph{entire} slice paths for
which
\begin{equation}
 \mathbb E\sum_{i\in I}1_{\{Z_i\ne Z_i'\}}\le4b.
 \label{f16v24:boundary-susceptibility}
\end{equation}
The same is true for a right-boundary change. The bound is independent of
$|I|,B,n$, and the actual boundary values, in the stated width regime.
Both endpoint changes can be paid with the sum of their separate budgets.
\end{proposition}
\begin{proof}
Use the actual forward Markov disintegration of the interval, with its fixed
right endpoint included in its continuation messages. Starting at its left
boundary, put $T=|I|$ and take
$M=\max\{0,\lfloor T/b\rfloor-1\}$ complete steps of length $b$.
The residual segment has fewer than $2b$ slices, and for $T\ge2b$ has
at least $b$ slices. If $T<2b$, couple the interval arbitrarily; its
Hamming cost is already less than $2b$.

For every complete step before the last segment, the target is at least $b$
from the fixed right endpoint. Given its two starting values, couple its
endpoint marginals by the usual common-part coupling; when the starting values
are equal use identical samples. Equation \eqref{f16v24:channel-input} bounds
the conditional probability of unequal step endpoints, given unequal starting
values, by $q$. Given both endpoints of a segment, fill its internal slices
from their \emph{exact} conditional bridge laws. These are the marginals of
the original interval disintegration. Their normalizers are not discarded.
If the starting values are equal, the entire segment may be drawn identically.
If they are unequal, its inside cost is at most $b$, even when its last slices
coalesce. Once the sampled step endpoints agree, the full remaining right
conditional can be coupled identically; the common fixed right factor has
not changed.

If $M$ complete steps were used, their expected total cost is at most
$b\sum_{j=0}^{M-1}q^j$. The residual segment costs at most $2bq^M$.
Thus the total is at most $b/(1-q)+2b\le(10/3)b<4b$. The construction retains
the native law in each component by Markov disintegration and conditional
bridge filling. These are measurable finite-dimensional probability kernels;
their common-part couplings use their actual densities. No simultaneous
coupling of all interval boundaries is required. The case $M=0$ is included.
Reverse time for a right change. For two changed endpoints, concatenate the
two comparisons through a path law with one endpoint from each, and use
conditional gluing of couplings and the metric triangle inequality. Natural
endpoints create no additional term.
\end{proof}

The terminal $2b$ term is essential. A forward step too close to the fixed
right boundary was \emph{not} proved contracting by V22. We stop before that
step and pay its whole output, instead of treating the opposite conditioning
as harmless.

\subsection{Overlapping interval resampling with unit coverage at every time}

Set $\ell=16b$. Use the indexed clipped intervals
\begin{equation}
 I_s=[s,s+\ell-1]\cap[0,n],\qquad 1-\ell\le s\le n,
 \qquad
 \mathcal L_{\rm slab}=\frac1\ell\sum_{s=1-\ell}^{n}(I-E_{I_s}).
 \label{f16v24:slab-generator}
\end{equation}
Here $E_{I_s}$ conditions on all unmoving cells outside the \emph{entire}
space--time slab $I_s\times\mathbb T_m^3$. Each indexed summand has rate
$1/\ell$. Some clipped intervals can coincide when $n+1<\ell$; they are
\emph{not} deduplicated. Every time slice belongs to exactly $\ell$ indexed
intervals, so its coverage rate is exactly one, including either dressed end.
There is no factor $1/(n+1)$ in this clock.

\begin{theorem}[A full-$L^2$ native slab gap with no duration loss]
\label{f16v24:slab-gap}
Under \eqref{f16v24:channel-input}, the exact reversible generator
\eqref{f16v24:slab-generator} satisfies
\begin{equation}
 \boxed{\quad \mathcal L_{\rm slab}\succeq\tfrac12(I-\mathbb E_P),\qquad
 \operatorname{Var}_P f\le\frac2\ell
           \sum_{s=1-\ell}^{n}\mathbb E_P\operatorname{Var}_P(f\mid Z_{I_s^c}).
 \quad}
 \label{f16v24:slab-gap-bound}
\end{equation}
It holds for every $f\in L^2(P)$, not only local sources or polynomial tests.
No small-field condition is imposed on a conditional exterior.
\end{theorem}
\begin{proof}
Treat complete slices as coordinates and give each weight one in
Theorem~\ref{f16v24:general-gap}. A changed slice $j$ is contained in exactly
$\ell$ indexed intervals: $j-\ell+1\le s\le j$. Its death rate is one.
For a slab not containing $j$, its conditional law can depend on this change
only if $j$ is an existing immediate left or right boundary. There are at
most two such indexed intervals, namely $s=j+1$ and $s=j-\ell$ when they
exist. All other exterior changes cancel from the conditional by the exact
temporal Markov property. Proposition~\ref{f16v24:interval-coupling} bounds
each affected output cost by $4b$. Hence the margin is at least
\[
                    1-\frac{8b}{\ell}=\frac12.
\]
Apply Theorem~\ref{f16v24:general-gap}. Short histories and clipped endpoints
were included in both counts. Strict positivity of the native density and
all its original endpoint factors ensures these are its own orthogonal
conditional projections. No independence between simultaneous overlapping
slabs and no product of their grand couplings is asserted.
\end{proof}

The gap belongs to a \emph{newly defined auxiliary sampler}, not to a power of
the physical transfer. Its time normalization has been specified explicitly.
An individual move may change $B\ell$ cells. A comparison to original local
updates is therefore still required, and is now paid rather than assumed.

\subsection{The exact cost of returning to the original balanced updates}

Let $\omega=(X,r)$ be V14's balanced compact variables. One elementary update
resamples all five chords of one cube at one time, or one nonroot moving
port. Denote their rate-one projections by $\widehat Q_\alpha$ and set
\[
 \widehat{\mathcal L}_P=\sum_\alpha\widehat Q_\alpha,
 \qquad \widehat{\mathcal E}_P(f)=\sum_\alpha
             \mathbb E_P\operatorname{Var}_P(f\mid\omega_{\alpha^c}).
\]
This is the original unweighted balanced generator, not the whole-slice
sampler. Its weighted version is a separate clock and is not substituted.
The exact relation to unmoving ports is
\begin{equation}
             r_t=f_t(X_t)^{-1}h_t f_{t+1}(X_{t+1}).
 \label{f16v24:frame-change}
\end{equation}
All frame operations are within their original spatial cubes.

\begin{proposition}[Bounded classical oscillation gives a conditional clock comparison]
\label{f16v24:finite-comparison}
For any finite set $D$ of elementary balanced update groups, its complete
conditional satisfies
\begin{equation}
 \mathbb E_P\operatorname{Var}_P(f\mid\omega_{D^c})
 \le \exp(C_r\gamma|D|)
             \sum_{\alpha\in D}\mathbb E_P
                       \operatorname{Var}_P(f\mid\omega_{\alpha^c}).
 \label{f16v24:conditional-clock}
\end{equation}
The constant is independent of $D,B,n$ and every fixed exterior. Its potentially
large exponential is retained; it is not a coupling-uniform local gap.
\end{proposition}
\begin{proof}
In balanced variables the reference measure is a product over chord-cube
blocks and nonroot ports. The chord factors retain exactly the original
principal-chart Lebesgue measure with Haar exponent $1$ or $1/2$; moving ports
retain Haar measure by conditional Haar invariance. Normalize these individual
finite reference measures only for this comparison and call their product on
$D$ $\nu_D$. Their normalization scalars cancel in the conditional density.
No ratio of a fractional endpoint Haar factor to full Haar is bounded here.

The remaining conditional density is proportional to $e^{-U_D}\nu_D$,
where $U_D$ is the \emph{classical} action of all terms meeting $D$.
Every Wilson term lies between $0$ and $2\gamma$, has fixed carrier and bounded
incidence in these coordinates. V5/V14's cube-local frame dependence preserves
these bounds. Each endpoint quadratic is bounded by $C\gamma$ per incident
chord cube on its complete principal domain. Consequently
\[
            \operatorname{osc}U_D\le C_r\gamma|D|.
\]
Exterior-only terms cancel, but every crossing term remains. This bound uses
neither convexity nor a typical-exterior event.

For completeness, if $\mu_D=p\nu_D$ and $m\le p\le M$, then
$M/m\le e^{\operatorname{osc}U_D}$. The variational characterization of
variance and the product variance inequality give
\[
 \operatorname{Var}_{\mu_D}f
 \le M\operatorname{Var}_{\nu_D}f
 \le M\sum_{\alpha\in D}\inf_{g\text{ measurable off }\alpha}
                             \int|f-g|^2d\nu_D
 \le\frac M m\sum_{\alpha\in D}
                  \mathbb E_{\mu_D}\operatorname{Var}_{\mu_D}(f\mid\omega_{D\setminus\alpha}).
\]
The product inequality follows, for example, by revealing independent
coordinates successively and applying conditional Jensen to each martingale
difference. The infimum formula proves the last inequality even for unbounded
$L^2$ tests, by truncation. Conditioning the other variables of $D$ restores
the original elementary single-group conditional. Integrate over the actual
exterior marginal to obtain \eqref{f16v24:conditional-clock}. All measures have
the same full support up to reference-null sets, so the bounded density
comparison is legitimate.
\end{proof}

\begin{proposition}[A slab changes its preceding balanced port as well]
\label{f16v24:slab-enlargement}
For an unmoving interval $I=[a,c]$, let $D(I)$ contain its balanced chord cubes
at times $a,\ldots,c$ and moving ports at times $a-1,\ldots,c$, omitting
nonexistent ports at $-1$ or $n$. Then
\begin{equation}
 \mathbb E_P\operatorname{Var}_P(f\mid Z_{I^c})
 \le\mathbb E_P\operatorname{Var}_P(f\mid\omega_{D(I)^c}),\qquad
 |D(I)|\le8B(|I|+1).
 \label{f16v24:enlargement}
\end{equation}
For the clipped indexed slabs, each chord-cube update belongs to exactly
$\ell$ sets $D(I_s)$ and each moving-port update to exactly $\ell+1$.
\end{proposition}
\begin{proof}
The unmoving exterior fixes all $X_i,h_i$ outside $I$. Formula
\eqref{f16v24:frame-change} then fixes every balanced variable outside the
stated $D(I)$: if $t\notin[a-1,c]$, both $X_t,X_{t+1}$ and $h_t$ are fixed.
Thus $\sigma(\omega_{D(I)^c})\subseteq\sigma(Z_{I^c})$. Conditional expectation
is an orthogonal projection, so conditioning on the latter, larger sigma
algebra decreases the averaged conditional variance. This proves the inequality;
it does \emph{not} identify the two conditional kernels.

There are $B|I|$ chord-cube groups and at most $7B(|I|+1)$ port groups.
For coverage, a chord time $i$ occurs in $\ell$ intervals as already counted.
A port time $t\in[0,n-1]$ occurs when $I_s$ contains $t$ or $t+1$; the union
of start indices is $t-\ell+1\le s\le t+1$, of size $\ell+1$. Both time
indices exist, so clipping changes neither count. In general omitting the
preceding port would be false: changing $X_a$ can change $r_{a-1}$ with
$h_{a-1}$ held fixed. The elementary updates are therefore not the same in
moving and unmoving coordinates.
\end{proof}

\begin{theorem}[A duration-uniform bound in the original rate-one clock]
\label{f16v24:original-gap}
Under the channel and width hypotheses above, with $\ell=16b$,
\begin{equation}
 \mathcal L_{\rm slab}\preceq
 \left(1+\ell^{-1}\right)e^{C_r\gamma B(\ell+1)}\widehat{\mathcal L}_P,
 \qquad
 \boxed{\quad\widehat\lambda_P\ge
               \frac14 e^{-C_r\gamma B(\ell+1)}.\quad}
 \label{f16v24:original-gap-bound}
\end{equation}
All inequalities hold on the complete native $L^2$ space, with original endpoint
factors. In particular, for every fixed $p,A_0$,
\begin{equation}
 B\le A_0\gamma^p\quad\Longrightarrow\quad
       \widehat\lambda_P\ge\tfrac14
                      e^{-C_{p,r}\gamma B\log\gamma},
 \label{f16v24:polynomial-gap}
\end{equation}
and in the specified exponential-width channel regime,
\begin{equation}
 B\le e^{\beta_{\rm ch}\gamma}\quad\Longrightarrow\quad
       \widehat\lambda_P\ge\tfrac14 e^{-C_r\gamma^2B}.
 \label{f16v24:exponential-gap}
\end{equation}
The entry thresholds are those of the corresponding channels, enlarged by
fixed requirements only. These are positive bounds uniform in duration, not
positive bounds uniform in $\gamma$ or unrestricted spatial volume.
\end{theorem}
\begin{proof}
Use \eqref{f16v24:enlargement} for each slab and then
\eqref{f16v24:conditional-clock} on its enlarged balanced block. Its group
count is at most $8B(\ell+1)$, absorbed into $C_r$. Sum the rates $1/\ell$.
The exact coverage counts bound the rate assigned to any elementary update
by $1+1/\ell$. This proves the operator-form comparison. Combine it with
\eqref{f16v24:slab-gap-bound} on centered space. The sharper lower constant
is $[2(1+1/\ell)]^{-1}$; $1/4$ is a uniform simplification. Insert the inherited
bounds $b=O_{p,r}(\log\gamma)$ or $b=O_r(\gamma)$ to obtain the two displayed
regimes. No source-covariance estimate was used in the proof of this gap.
\end{proof}

This pays a genuinely missing clock step for a concrete dynamics, but not at
a useful weak-coupling constant. Inserting \eqref{f16v24:original-gap-bound}
into V16's gap-only source estimate is generally \emph{worse} than V23's direct
source theorem. It is not reported as an improved covariance bound. Its new
content is all-$L^2$ coercivity in the original rate-one clock without a duration
factor, and an exact identification of its expensive finite-block comparison.
A sharp estimate for those finite conditional blocks could improve this cost;
the present proof has not established such an estimate.

\subsection{Why common-core success alone does not verify the spatial ledger}

For the spatial diagnostic below, the Hamming coordinates are complete
\emph{unmoving cells}, not the elementary balanced chord/port groups.
Returning a spatial cell-block dynamics to those elementary updates would
require the analogous frame-support enlargement as well.
The same general criterion can be tested on local space--time cell blocks. Its
missing quantity is $I_{D,j}$, the expected \emph{entire} resampled output
disagreement due to a \emph{single} changed exterior input. V23's common map
already supplies a valid but generally inadequate upper bound. If the block
has core $A$, fringe $D\setminus A$, and common core mass $\alpha$, couple the
two full outputs with its same success decision and seeds. On success the
core agrees, but every fringe coordinate may still disagree. On failure every
coordinate of $D$ may disagree. For unit weights this proves
\begin{equation}
 I_{D,j}\le U_D:=|D\setminus A|+(1-\alpha)|A|,
                  \qquad j\in\partial D.
 \label{f16v24:crude-fringe}
\end{equation}
It uses the actual conditional filling of the fringe. It is not an assertion
that different native error fields are independent.

\begin{proposition}[The whole-fringe estimate has the wrong spatial scaling]
\label{f16v24:fringe-obstruction}
For translated nearest-neighbour boxes $V_R=\{1,\ldots,R\}^4$ and core
$A_R=\{2,\ldots,R-1\}^4$, with $R\ge3$, normalize each box rate by $R^{-4}$.
The uniform influence table supplied solely by \eqref{f16v24:crude-fringe}
would give the certified margin
\begin{equation}
 1-\frac{8}{R}\left[R^4-(R-2)^4+(1-\alpha)(R-2)^4\right]<0,
 \label{f16v24:fringe-margin}
\end{equation}
even at $\alpha=1$. This is failure of this \emph{sufficient bound}, not proof
that the actual dynamics has a nonpositive gap or that better couplings cannot
work. The geometry is an explicit four-dimensional comparison template, not
an assertion that the native interaction graph has exactly these eight faces.
\end{proposition}
\begin{proof}
A site belongs to $R^4$ translated boxes and lies on the exterior nearest-neighbour
collar of $8R^3$ translated boxes, away from finite-size identifications.
Thus coverage is one, whereas the uniform birth estimate is $8U_D/R$.
Even at $\alpha=1$, the two opposite layers in one coordinate give
$R^4-(R-2)^4\ge2R^3$. The displayed margin is therefore at most
$1-16R^2<0$. Larger fringes only worsen this particular table. No lower bound
on the \emph{true} $I_{D,j}$ has been inferred from its upper bound.
\end{proof}

For comparison, if a translated finite-range four-dimensional box family
could be given \emph{full native} pair couplings satisfying
\begin{equation}
 I_{D,j}\le\Xi_\gamma\quad\hbox{for each changed collar input,}\qquad
 \frac{|\partial D|}{|D|}\le\frac{C_\rho}{R},
 \label{f16v24:spatial-target}
\end{equation}
with $\Xi_\gamma$ independent of the box side $R$ and the ambient size, then
$R\ge2C_\rho\Xi_\gamma$ would give $g\ge1/2$ in unit-coverage block time.
This follows directly from Theorem~\ref{f16v24:general-gap}. An actual cover on
a finite path or small torus must use its true coverage and boundary counts;
no translation-invariance of $P$ is needed, but those counts cannot be omitted.
A product-density comparison on correctly enlarged balanced supports would
then pay the elementary-clock cost; a cell update is not identified with an
elementary balanced update. This is a \emph{conditional criterion};
\eqref{f16v24:spatial-target} is not proved for the unrestricted native law.

The temporal argument verifies precisely such a boundary-sensitive quantity
with $\Xi=4b$ and only two incoming slices. Its success did not come from
multiplying unrelated core refresh probabilities. Spatially, the near-boundary
fringe can carry a change tangentially, and the common-core event alone places
no localized bound on that propagation. Spinka's boundary-subset condition
records the analogous issue for grand couplings; our spectral criterion needs
only pairwise single-input couplings, but it still needs their full-output cost.

\subsection{What changes in the program and what remains open}

There are two distinct new outputs. First, the complete conditional interval
has an actual coupling with affected-slice budget $4b$, uniformly in its
length and in both compact boundary values. This constructs a reversible
full-block sampler with gap $1/2$ and \emph{unit coverage per slice}. Second,
the exact moving/unmoving support change and complete density comparison return
it to the original balanced sampler with its explicit exponential price.
The latter bound is uniform in duration within the stated polynomial or
exponential spatial-width regime. It applies to every native $L^2$ function,
not just the chosen response bank.

Neither result gives an original-sampler gap uniform in coupling. Nor does it
improve V23's source asymptotics. The primitive/common-source covariances, V19's
positive shell return and V12's different cylinder return have not been
redefined. In particular the exponential cost in
\eqref{f16v24:original-gap-bound} cannot be suppressed by saying that each
single update has rate one. Overlap, coordinate conversion, and conditional
mixing are separate parts of the clock ledger, and all have been paid here.

The next spatial calculation is sharper than another common-core estimate:
bound the total resampled disagreement caused by one changed collar input,
with the genuine buffer law retained, or obtain an equivalent directional
operator estimate. The whole-fringe bound fails that test already on the
simplest four-dimensional box. Improving only its common success probability
cannot fix its surface-volume scaling. Alternatively, sharpening the native
finite-block comparison would improve the original-clock result in the existing
width regime, but would not itself remove the complete-separator cost.

All clocks in this appendix are auxiliary resampling clocks. No statement
identifies them with Euclidean transfer time. The fixed rescaled bridge window,
root Gauss restriction, infrared contact coercivity, same-top physical Wilson
transfer comparison, independently calibrated physical time, and nontrivial
interacting continuum construction remain separate obligations. No physical
Yang--Mills mass gap is claimed.

\subsection{Diagnostics, preservation, and proof status}

The new diagnostic is
\begin{center}\small\path{verify_ym_f16_v24_clock_sensitivity.py}.\end{center}
It checks exact finite conditional projections, the weighted influence ledger,
normalized inhomogeneous path kernels with a retained opposite endpoint,
clipped-interval coverage and boundary counts, the preceding-port enlargement,
product-density comparisons and the constants in the fringe and clock estimates.
The new checker passed 1,128 exact assertions. The finite path and spin
fixtures are not native Wilson integral quadratures.
Native compact integral suprema and entry thresholds are not enclosed by those
fixtures. The proof of the theorems is the argument above with its identified
analytic premises, not a consequence of the finite assertion count.

All twenty supplied V4--V23 diagnostics were rerun unchanged and passed.
The package records their actual outputs, source hashes, the new diagnostic,
predecessor checks, compilation and rendered-page inspection. The cumulative source changes
only the title version and inserts this exact appendix before the final
document-ending command; its reverse operation recovers V23 byte for byte.
No inherited claim, normalization or bibliography has been silently changed.
The scope of the inherited analytic archive has not been independently
recertified. The next scheduled companion checkpoint remains V25.

\begin{firewall}
V24 proves an all-observable gap for exact overlapping temporal-slab resampling
in a specified unit-coverage clock. It pays the conversion to the original
rate-one balanced chord/port sampler and obtains a positive duration-uniform
bound with an explicit coupling/width-dependent exponential cost. It proves
a boundary-input sensitivity criterion and diagnoses the insufficient
whole-fringe bound. It does not prove unrestricted spatial mixing, a
noncollapsing original-sampler gap as coupling increases, the missing spatial
sensitivity estimate, a physical transfer gap, or an interacting continuum
mass gap. Every native conditional, compact buffer and endpoint factor is
retained.
\end{firewall}
% END F16 V24 APPEND

% BEGIN F16 V25 APPEND -- 2026-09-17
% Insert immediately before the final document-ending command in V24.
% The predecessor body and all its probability/source conventions are unchanged.

\clearpage
\section{V25: interface-cost comparison and an original-clock finite-width baseline}
\label{f16v25:context}

V24 pays for converting a whole native slab update to elementary balanced
updates by the oscillation of the \emph{entire} conditional action. That costs
an exponential in the slab volume. It is not necessary to pay every interaction
in that comparison. Exchange the coordinates of two configurations in a fixed
order. Terms wholly on one side of the moving cut cancel \emph{exactly} between
the two configurations and their complementary recombinations. Only terms
crossing the cut remain. We give this comparison for continuous product spaces,
with the conversion from an independent marginal draw to the actual one-site
conditional included explicitly.

For the native slab this replaces a volume cost by the smaller of a temporal
and a spatial interface cost. It also makes a different, elementary temporal
input useful: compact bounded interactions give a very small but positive
one-step common component under every continuation message. Combining these
two finite arguments gives a positive lower bound for the original rate-one
balanced sampler, uniform in duration, for \emph{every} finite spatial width.
The lower bound decays exponentially with that width and with coupling. It is
not a thermodynamic spectral floor, and it does not replace V23's sharper
source estimates in their stated regimes.

\subsection{V25 checkpoint: source ownership and the changed comparison}

The supplied V24 appendix and analytic audit were read in full. The V5 exact
unmoving action and its local carriers, V14's balanced coordinates, and V24's
conditional-support enlargement were checked. The mounted f14 V100, f15 V100,
and Grand Tensor V076 sources were searched for cutwidth, coordinate-splicing,
and canonical-path comparisons after two scoped semantic searches returned no
results. The f15 hit for ``canonical path indicators'' concerns geometric
repair paths, not this Gibbs variance comparison. The monograph's random-scan
and blocking warnings remain binding. A failed search is not evidence of
archive-wide novelty; this is a targeted checkpoint, not recertification of
all historical analytic claims or of unavailable companion programs.

The method is classical. Berger--Kenyon--Mossel--Peres,
\href{https://arxiv.org/abs/math/0308284}{\emph{Glauber Dynamics on Trees and
Hyperbolic Graphs}}, Proposition 1.1 and its Section 2.1 proof, relates an Ising
relaxation bound to the number of interactions crossing an ordering cut. The
statement and the canonical-path proof were examined, including their rate
convention. We do not quote its finite-spin theorem for a continuous Wilson
law. The finite-product proof below treats arbitrary probability reference
factors and many-coordinate interactions directly. No lower bound on a point
mass of a continuous spin is required.

The two supplied import memoranda propose separating order gain from locality
cost, and retaining a genuine dependence estimate before a global functional
inequality. Here there is no new perturbative order and no new global spatial
coupling. The changed object is the cost of the \emph{same-law clock comparison}.
The new unrestricted-finite-width baseline below uses bounded classical words,
product references, exact temporal disintegration, and V24's proved block-gap
criterion. It does not use the V16 moment theorem, the small-angular one-form
realization, or the V17--V23 asymptotic conclusions. The sharper versions may
reuse V22's better temporal buffer as a separate input. The next scheduled
companion checkpoint is V30.

\subsection{A two-copy identity that charges only an interface}

Consider a finite product of standard Borel spaces, with probability reference
\(\nu=\bigotimes_{i=1}^{N}\nu_i\), and the exact probability
\[
 d\mu=Z^{-1}e^{-U}\,d\nu,\qquad U=\sum_X\Phi_X.
\]
The finitely many interaction functions are bounded and measurable. One-site
terms may first be incorporated into their own reference factor; this changes
neither \(\mu\) nor its conditional resampling. All remaining \(X\)'s can
therefore be taken to have at least two elements. A fixed exterior parameter,
if present, is substituted in each interaction before its free support is
counted. There is no restriction on those exterior values.

Fix an ordering, relabelled \(1,\ldots,N\). Write
\[
 S_i=\{1,\ldots,i-1\},\quad
 c_i=\sum_{X:\,X\cap S_i\ne\varnothing,\ X\cap S_i^c\ne\varnothing}
                  \operatorname{osc}\Phi_X,\quad
 s_i=\sum_{X\ni i}\operatorname{osc}\Phi_X,
\]
where the complement is within this free product. Define
\begin{equation}
 \mathfrak c_\prec=\max_i c_i,\qquad
 \mathfrak s=\max_i s_i.
 \label{f16v25:interface-costs}
\end{equation}
The second quantity is a one-coordinate interaction cost. It is not the
oscillation of all terms in \(U\).

For configurations \(x,y\), let
\[
 z^i=(y_1,\ldots,y_i,x_{i+1},\ldots,x_N),\quad z^0=x,
 \qquad w^{i-1}=(x_{S_i},y_{S_i^c}).
\]
Thus \(z^{i-1}=(y_{S_i},x_{S_i^c})\). These are coordinate exchanges, not
interpolations on the group or approximations of a conditional density.

\begin{proposition}[Exact two-copy cancellation]
\label{f16v25:splice}
Let \(p=d\mu/d\nu\). The coordinate exchange is a measurable
involution preserving \(\nu\otimes\nu\). Its density ratio is
\begin{equation}
 \frac{p(x)p(y)}{p(z^{i-1})p(w^{i-1})}
 =\exp\{-U(x)-U(y)+U(z^{i-1})+U(w^{i-1})\}
 \le e^{2c_i}.
 \label{f16v25:splice-ratio}
\end{equation}
Every term with support on only one side of the cut cancels, even if it is
large, nonconvex, or depends on a fixed exterior. The partition functions
cancel exactly in this two-copy ratio.
\end{proposition}
\begin{proof}
At each coordinate either exchange the two entries or leave them alone. Its
reference factor is \(\nu_i\otimes\nu_i\), so product Fubini proves invariance
and the map is its own inverse. A term supported in \(S_i\) has its two values
exchanged; one supported in \(S_i^c\) has them retained. Both contributions to
the exponent are zero. For a crossing term its two new values minus its two
old values are at most twice its oscillation. Sum only those bounds. This
proves the identity and inequality. The argument does not commute the factors
inside a noncommutative Wilson word: a whole scalar word potential is evaluated
at each of the four configurations before it is bounded.
\end{proof}

\begin{theorem}[Interface comparison in rate-one time]
\label{f16v25:path-comparison}
Let \(\mathcal E_i^\mu(f)=\mathbb E_\mu
\operatorname{Var}_\mu(f\mid\omega_{i^c})\). For every \(f\in L^2(\mu)\),
\begin{equation}
 \boxed{\quad
 \operatorname{Var}_\mu f
 \le N\sum_{i=1}^N e^{2c_i+2s_i}\mathcal E_i^\mu(f)
 \le N e^{2\mathfrak c_\prec+2\mathfrak s}
                         \sum_i\mathcal E_i^\mu(f).
 \quad}
 \label{f16v25:path-bound}
\end{equation}
This bounds the inverse gap of the sum of rate-one conditional projections.
No extra division by \(N\) is part of the generator. The factor \(N\) in this
bound is the cost of the chosen telescoping path, not a change of update clock.
\end{theorem}
\begin{proof}
Initially let \(f\) be bounded. The independent-copy variance identity and
Cauchy--Schwarz along the path give
\[
 \operatorname{Var}_\mu f
 =\tfrac12\mathbb E_{\mu\otimes\mu}|f(y)-f(x)|^2
 \le\tfrac N2\sum_i\mathbb E_{\mu\otimes\mu}
                       |f(z^i)-f(z^{i-1})|^2.
\]
Use Proposition~\ref{f16v25:splice} in the \(i\)th term. Writing \(z=z^{i-1}\)
and \(w=w^{i-1}\), that term is at most
\[
 e^{2c_i}\int\mu(dz)\mu(dw)
                          |f(z^{i\leftarrow w_i})-f(z)|^2.
\]
The incoming \(w_i\) in this expression has the \emph{unconditional marginal}
\(\mu_i^{\rm m}\). It is not the native heat-bath conditional at \(z_{i^c}\).
We pay that distinction before interpreting a Dirichlet form.

For two exteriors \(a,a'\), the difference of the one-coordinate interaction
potentials as a function of its free value has oscillation at most \(2s_i\).
The exact normalized conditional density ratio is consequently between
\(e^{-2s_i}\) and \(e^{2s_i}\). This follows by dividing the two exponential
weights and their two integrals; the normalized log ratio lies between the
minimum and maximum of the unnormalized log ratio minus one scalar in that
same interval. Averaging one of the conditional measures over its actual
exterior marginal gives
\begin{equation}
             \mu_i^{\rm m}(db)\le e^{2s_i}\mu_i(db\mid z_{i^c}).
 \label{f16v25:marginal-conditional}
\end{equation}
No lower bound on the mass of an individual \(b\) is used.

Substitute \eqref{f16v25:marginal-conditional}. The two free coordinates are
now independent draws from the \emph{same} conditional on \(z_{i^c}\); the
inner squared difference has integral twice its conditional variance.
This cancels the initial factor \(1/2\) and proves the first inequality of
\eqref{f16v25:path-bound}. The second uses the maxima. Truncating a real or
complex \(L^2\) function and using contraction of conditional expectation in
\(L^2\) extends the result. Bounded interaction weights give equivalent
finite reference and native \(L^2\) spaces, so all intermediate integrals are
legitimate. The proof covers continuous, atomic, or mixed reference factors.
\end{proof}

For an independent product this particular path estimate loses \(N\), although
product tensorization has constant one. The theorem is deliberately a robust
comparison, not a claim of sharpness. Its advantage over a whole-density bound
is the cancellation of all noncrossing interactions from the exponential.

\subsection{Actual balanced supports and two native ordering cuts}

Keep \(P=P_{\gamma,\mathbf y}\), \(B=m^3\), \(m\ge4\), \(n\ge1\), and the
prescribed bridge window. We use \(\gamma\ge1+r^2\) throughout for consistency
with the existing lineage. The original elementary coordinates are the five
chords of one cube at one time, regarded as \emph{one} update group, and each
of its seven nonroot moving ports. These are precisely V14's rate-one updates.
Their product reference factors retain all original Haar powers; the two
endpoint chord factors have power \(1/2\), not power one.

Assign each elementary group to its cube and its time. There are at most eight
such groups per cube--time cell. The exact balanced frame relation is local to
a cube. Consequently the remaining scalar classical word potentials have a
fixed bounded carrier diameter, a fixed incidence per elementary group, and
oscillation at most \(C\gamma\) per word. The endpoint quadratic is a sum of
one-cube functions and can be included in that cube's reference. After a
conditional exterior is fixed, any word with only one free elementary group
can likewise be included in its reference; all words with two or more free
groups are retained as interactions. This is a factorization of the same
conditional probability, not a cutoff or a reference law substituted for it.

\begin{proposition}[Two interface bounds for the exact enlarged slab]
\label{f16v25:native-cuts}
For an unmoving interval \(I=[a,c]\), let \(D(I)\) be the exact enlarged balanced
support of V24: chords at \(a,\ldots,c\), and moving ports at
\(a-1,\ldots,c\), with unavailable endpoints omitted. Put \(T=|I|\). Then
\begin{equation}
 |D(I)|\le8B(T+1),\qquad
 \inf_\prec\mathfrak c_\prec(D(I))+\mathfrak s(D(I))
 \le C_r\gamma\min\{B,m^2(T+1)\}.
 \label{f16v25:cuts}
\end{equation}
The bound is uniform in every fixed balanced exterior, all compact fast
coordinates, and the original endpoint factors. Periodic spatial seams are
included in the cut count.
\end{proposition}
\begin{proof}
The group count is V24's exact support count, not a newly shrunken support.
Let \(R_*\) bound the carrier range of the balanced words in every coordinate
direction, and let the fixed incidence bound be \(M_*\). Neither depends on
the regulator or the prescribed bridge values in their window. Assign every
word crossing an ordering cut to one incident group in its active layers;
there are at most \(M_*\) assigned words per such group.

First order by time, keeping all groups of a cube--time cell consecutive.
At any prefix, words crossing the cut must meet either the partially processed
time layer or at most \(R_*\) neighboring time layers on either side. There
are at most \(C B\) groups there, regardless of \(T\). A word with all free
groups earlier or all later cancels even if it also contains frozen variables.
There is no time-periodic seam. This gives \(\mathfrak c_\prec\le C\gamma B\).

Alternatively order by the first spatial cube coordinate. An ordering prefix
has fully processed spatial layers, at most one partially processed layer,
and unprocessed layers. A crossing word must meet a fixed-width band at the
moving cut or a fixed-width band at the periodic seam. Each band contains at
most \(C m^2(T+1)\) groups. Terms inside the partial layer are included by
counting that entire layer. If the period is smaller than the band width,
count the actual period once; the bound only improves after enlarging the
fixed constant. Missing chord or port groups at either temporal end can be
removed from this ordering without adding a crossing word. Thus this second
ordering gives \(C\gamma m^2(T+1)\).

Every elementary group meets at most a fixed number of remaining interactions,
so \(\mathfrak s\le C\gamma\). Since the minimum in
\eqref{f16v25:cuts} is at least one, absorb this last term and select the better
of the two orderings. No spatial translation invariance of \(P\) has been used.
\end{proof}

\begin{theorem}[Finite native block comparison with an interface, not volume, exponential]
\label{f16v25:native-comparison}
For every \(f\in L^2(P)\), the actual enlarged-block conditional obeys
\begin{align}
 &\mathbb E_P\operatorname{Var}_P(f\mid\omega_{D(I)^c})\notag\\
 &\qquad\le
 8B(T+1)\exp\!\left(C_r\gamma\min\{B,m^2(T+1)\}\right)
       \sum_{\alpha\in D(I)}
                 \mathbb E_P\operatorname{Var}_P(f\mid\omega_{\alpha^c}).
 \label{f16v25:native-block}
\end{align}
There is no restriction on the exterior values or on the ratio between \(B\)
and \(\gamma\). The target is still the complete native conditional.
\end{theorem}
\begin{proof}
Condition on the balanced exterior of \(D(I)\). Normalize the original product
reference factors, with the one-free-group terms included as just described.
Theorem~\ref{f16v25:path-comparison} applies to this finite product. Its
interaction bounds are uniform in the exterior by
Proposition~\ref{f16v25:native-cuts}. Fixing the other coordinates of \(D(I)\)
restores exactly the original elementary conditional on the right. Integrate
over the actual exterior marginal. Its conditional partition functions are
not replaced or estimated separately. This proves the statement.
\end{proof}

In particular the exponential cost is \(C\gamma B\), not
\(C\gamma B(T+1)\), when the temporal cut is chosen. When \(T+1\ll m\), the
spatial cut improves it further to \(C\gamma m^2(T+1)\). The remaining linear
factor counts elementary path steps. This does not prove a spatial coupling
of different boundary rows; it compares forms at one and the same boundary.

\subsection{A sharper paid return to the original update clock}

Let \(b\) be any valid V24 temporal channel buffer with step diameter at most
\(1/4\). Its complete-slab generator with \(\ell=16b\) has gap at least
\(1/2\), as already proved there. Its indexed clipped intervals, including
multiplicities, remain unchanged. Define
\[
       W_m(\ell)=\min\{B,m^2(\ell+1)\}.
\]

\begin{theorem}[Interface-priced original-clock comparison]
\label{f16v25:clock}
Under that channel hypothesis, on the entire native \(L^2(P)\) space,
\begin{equation}
 \mathcal L_{\rm slab}
 \preceq 8B(\ell+1)(1+\ell^{-1})
                 e^{C_r\gamma W_m(\ell)}\widehat{\mathcal L}_P.
 \label{f16v25:form-clock}
\end{equation}
\begin{equation}
 \boxed{\quad
 \widehat\lambda_P
 \ge \frac{e^{-C_r\gamma W_m(\ell)}}{32B(\ell+1)}.
 \quad} \label{f16v25:clock-gap}
\end{equation}
The unweighted balanced generator \(\widehat{\mathcal L}_P\) is still the sum
of the original rate-one five-chord and single-port innovations.
\end{theorem}
\begin{proof}
V24 proves the exact sigma-algebra inclusion
\(\sigma(\omega_{D(I)^c})\subseteq\sigma(Z_{I^c})\), including the preceding
port. Thus the slab conditional variance is bounded by the left side of
\eqref{f16v25:native-block}. Apply that estimate to each clipped interval and
use \(T\le\ell\). In the indexed family a chord group is covered \(\ell\)
times, and a moving port \(\ell+1\) times. Rates are \(1/\ell\), so the
maximum aggregate elementary rate is \(1+1/\ell\), not the number of slabs.
This proves \eqref{f16v25:form-clock}. The slab gap gives the lower constant
\([16B(\ell+1)(1+1/\ell)]^{-1}\); use \(1+1/\ell\le2\) for the displayed
simplification. No variance estimate for a selected source was used.
\end{proof}

For example, insert the already proved V22 buffers. In their respective
large-coupling regimes this gives
\begin{align}
 B\le A_0\gamma^p\quad&\Longrightarrow\quad
 \widehat\lambda_P\ge
 \frac{\exp[-C_{p,r}\gamma\min\{B,m^2\log\gamma\}]}{C_{p,r}B\log\gamma},
 \label{f16v25:polynomial-clock}\\
 B\le e^{\beta_{\rm ch}\gamma}\quad&\Longrightarrow\quad
 \widehat\lambda_P\ge
 \frac{\exp[-C_r\gamma\min\{B,m^2\gamma\}]}{C_rB\gamma}.
 \label{f16v25:exponential-clock}
\end{align}
In the first line \(p,A_0\) are fixed before taking the coupling limit; its
entry threshold may depend on both. Constants in the buffer length can be
absorbed outside and inside the exponential: for \(C\ge1\),
\(\min\{B,Cx\}\le C\min\{B,x\}\). These are better comparison costs than
V24's volume exponent, but they still vanish as \(\gamma\) increases. They
are not substituted into V23's source theorem as a better weak-coupling
covariance estimate.

\subsection{An elementary channel available at every finite spatial width}

The next argument is separate from all small-field preparation and expansion
results. Retain the exact native slice factorization
\[
 dP=Z^{-1}\prod_{i=0}^n\psi_i(z_i)\,d\nu_i(z_i)
                              \prod_{i=0}^{n-1}W_i(z_i,z_{i+1}).
\]
Each factor \(W_i=e^{-U_{i,i+1}}\) contains only the original classical terms
joining the two adjacent slices. There are at most \(C B\) such terms, each
with bounded oscillation \(C\gamma\). Fix
\begin{equation}
 A_{\gamma,B}=C_r\gamma B\ge
                     \sup_i\operatorname{osc}U_{i,i+1},\qquad
 \alpha_0=e^{-A_{\gamma,B}}.
 \label{f16v25:raw-cost}
\end{equation}
All within-slice weights, including the actual endpoint reference factors,
remain in \(\psi_i\nu_i\). Scalar shifts of a pair potential do not change
its normalized transition. Shift it so that \(e^{-A_{\gamma,B}}\le W_i\le1\).

\begin{proposition}[Bounded pair words contract under every continuation message]
\label{f16v25:raw-channel}
For any actual finite continuation message, including an arbitrarily fixed
opposite endpoint, every one-step slice transition has a common submeasure
of mass at least \(\alpha_0\). Its diameter is at most \(1-\alpha_0\).
Consequently a valid V24 buffer at \emph{every} finite width is
\begin{equation}
 b_0=\left\lceil(\log4)e^{A_{\gamma,B}}\right\rceil,
 \qquad \delta(K_{i,i+b_0})\le\tfrac14,
 \qquad b_0\le3e^{A_{\gamma,B}}.
 \label{f16v25:raw-buffer}
\end{equation}
The estimate is uniform in duration. It is usually far weaker than V22's
channel where that stronger theorem applies.
\end{proposition}
\begin{proof}
Write the transition, with its actual right continuation, as
\[
 K_i(x,dy)=\frac{W_i(x,y)\psi_{i+1}(y)h_{i+1}(y)\,d\nu_{i+1}(y)}{h_i(x)}.
\]
The measure \(\psi_{i+1}h_{i+1}\nu_{i+1}\) has a positive finite integral at
each finite regulator. Normalize it and call the resulting probability
\(\lambda_i\). This is a native probability retaining the continuation;
it is not a Gaussian or stationary replacement. Then
\[
 K_i(x,dy)=\frac{W_i(x,y)}{\int W_i(x,z)\,\lambda_i(dz)}\lambda_i(dy)
                                \ge\alpha_0\lambda_i(dy).
\]
The denominator is at most one. A fixed opposite boundary changes \(h\) and
hence \(\lambda_i\), but not the inequality for any row \(x\). Reverse time
for the other orientation. Exact serial composition gives diameter at most
\((1-\alpha_0)^b\le e^{-b\alpha_0}\), proving the middle assertion of
\eqref{f16v25:raw-buffer}. The ceiling and \(\log4<2\) give its final bound.
If an interval is too short to contain a full buffer, V24's clipped-interval
and terminal-segment argument already pays its entire length. No long interval
is assumed to fit into a shorter history.
\end{proof}

\begin{theorem}[Duration-uniform original-clock baseline at every finite cross-section]
\label{f16v25:all-width}
There are finite constants \(C_r,c_r>0\), independent of \(B,n,\gamma\), such
that the original balanced sampler satisfies
\begin{equation}
 \boxed{\qquad
 \widehat\lambda_{\gamma,\mathbf y,B,n}
                 \ge \frac{c_r}{B}e^{-C_r\gamma B}>0
 \qquad}
 \label{f16v25:all-width-gap}
\end{equation}
for every \(m\ge4\), \(B=m^3\), \(n\ge1\),
\(\max_{i,e}|y_{i,e}|\le r\), and \(\gamma\ge1+r^2\). There is no polynomial
or exponential upper restriction on \(B\) relative to coupling. The bound is
uniform in duration \emph{at a fixed cross-section}, not uniformly positive
as \(B\to\infty\) or \(\gamma\to\infty\).
\end{theorem}
\begin{proof}
Use Proposition~\ref{f16v25:raw-channel}, V24's interval coupling and its exact
clipped-slab gap with \(b=b_0\), \(\ell=16b_0\). None of their conditional
arguments needs the stronger channel's width hypothesis once the required
diameter is supplied independently. This is a new application of those
proved finite arguments, not deletion of a hypothesis from their earlier
width-specialized statements.

Apply Theorem~\ref{f16v25:clock} and choose its temporal cut, so that
\(W_m(\ell)\le B\). Since
\(\ell+1\le49e^{A_{\gamma,B}}\), we obtain
\[
 \widehat\lambda_P\ge
 \frac{1}{1568B}\exp[-C_r\gamma B-A_{\gamma,B}].
\]
Absorb the two fixed interface constants into one \(C_r\) and take, for
example, \(c_r=1/1568\). All finite reference factors, endpoint messages and
conditional normalizers were kept. In particular no exponentially small
one-step common mass has been renamed a coupling-independent quantity. The
large buffer contributes its exponential only through the remaining linear
path prefactor, not a second exponential in its volume.
\end{proof}

This theorem does not require V16--V23's analytic comparison inputs. It uses
the finite bounded-word specification and V24's separately proved Markov,
path-coupling, and support identities. Thus it supplies a more elementary
baseline against which later sharp claims can be checked. Its decay
\(e^{-C\gamma B}\) is still far too costly for the unrestricted spatial
spectral edge. The phrase ``every finite width'' must not be replaced by
``uniform in width.''

\subsection{The exact congestion left behind by the interface supremum}

There is also a useful same-law object underneath the exponential majorant.
For the general finite product above, fix \(i\), use the cut \(S_i\), and set
\[
 x=(w_{S_i},z_{S_i^c}),\qquad y=(z_{S_i},w_{S_i^c}),\qquad
 Q_i(z,w)=\frac{p(x)p(y)}{p(z)p(w)}.
\]
Let \(\mu_i^{\rm m}\) be the actual coordinate marginal. With positive
conditional-density versions, define
\begin{equation}
 a_i(z,b)=
 \frac{d\mu_i^{\rm m}}{d\mu_i(\cdot\mid z_{i^c})}(b)
       \mathbb E_\mu[Q_i(z,W)\mid W_i=b].
 \label{f16v25:congestion}
\end{equation}
It is an average over the complementary configuration under its
\emph{actual conditional law}, with no Gaussian normalization. Values on
reference-null pairs may be fixed arbitrarily. Write
\[
 d\mathbb B_i(z,b)=\mu(dz)\mu_i(db\mid z_{i^c}).
\]
This is the exact edge measure of a rate-one native conditional update.

\begin{proposition}[Normalized edge congestion, and the weight it must multiply]
\label{f16v25:edge-identity}
The nonnegative density \(a_i\) obeys
\begin{equation}
 \int a_i\,d\mathbb B_i=1,\qquad 0\le a_i\le e^{2c_i+2s_i},
 \label{f16v25:edge-normalization}
\end{equation}
and, for every bounded \(f\),
\begin{align}
 &\mathbb E_{\mu\otimes\mu}|f(z^i)-f(z^{i-1})|^2\notag\\
 &\qquad=
 \int a_i(z,b)|f(z^{i\leftarrow b})-f(z)|^2\,d\mathbb B_i(z,b).
 \label{f16v25:edge-exact}
\end{align}
Consequently the first line of the path proof can retain these exact weighted
innovations instead of replacing \(a_i\) by its supremum. Its mean one does
\emph{not} authorize that replacement by one inside a squared innovation.
\end{proposition}
\begin{proof}
Apply the reference-preserving exchange without estimating its ratio. The
resulting integral is against \(\mu(dz)\mu(dw)Q_i(z,w)\). Disintegrate the
second copy at \(w_i=b\), then convert its marginal to the exact conditional
\(\mu_i(db\mid z_{i^c})\). This is exactly
\eqref{f16v25:congestion} and proves \eqref{f16v25:edge-exact}. Apply the same
change of variables with integrand one to obtain
\(\int a_i\,d\mathbb B_i=\int d\mu\,d\mu=1\).
Proposition~\ref{f16v25:splice} and
\eqref{f16v25:marginal-conditional} give its upper bound. All these statements
refer to the same normalized law. The last warning follows because a
nonnegative factor with mean one can correlate with the displayed squared
increment; the identity includes that correlation.
\end{proof}

This identifies a sharper potential use of native geometry: control the
\emph{weighted innovations} in \eqref{f16v25:edge-exact}, or choose paths
whose interaction cuts have a better source-dependent cost. An unweighted
mean of \(a_i\), a rare bad-event probability, or an estimate for a different
core conditional does not supply that comparison. The present proof uses the
uniform supremum and retains its price. Normalization of \(a_i\) is global
under \(\mathbb B_i\), not row normalization at each \(z\); it is not
a replacement Markov transition.

For a simple check on the scope of boundedness alone, take two binary spins
with law \(\mu_J(x,y)\propto e^{Jxy}\), \(J>0\), and both updates at rate one.
The centered function \(x+y\) has generator eigenvalue \(1-\tanh J\); the
other centered eigenvalues are \(1+\tanh J\) and \(2\). Hence its gap is
exactly \(1-\tanh J\to0\). This finite-spin example is not a native
Yang--Mills counterexample. It shows why a uniform positive low-temperature
gap cannot follow from compact support, positivity and bounded finite-range
interactions alone. Native-specific information is needed for that stronger
conclusion.

\subsection{Checkpoint conclusions and the next noncircular task}

The new cost ledger distinguishes three pieces: a same-law interface ratio,
a marginal-to-conditional conversion, and a linear path length. For native
slabs the first costs the smaller of a spatial and temporal section, not the
whole slab volume. The previous port enlargement and exact coverage remain
paid. This gives sharper original-clock estimates in the already controlled
width regimes, and a baseline positive duration-uniform gap at \emph{every}
finite width without using the perturbative hierarchy.

The spatial single-input susceptibility asked for in V24 has not been proved.
The new path comparison does not couple two different boundary laws and
cannot be relabelled as that result. Nor do the elementary one-step common
components solve spatial composition: their mass decays with the full slice
width. Repeating the crude whole-fringe estimate or improving only a core
success probability would still leave the same obstruction.

The next useful alternatives are now explicit: pay the actual weighted
interface congestion using native structure, or establish a spatial
single-input full-output coupling whose affected volume stays controlled.
Either would be new information beyond boundedness. An all-boundary
conditional gap uniform in coupling is not assumed. The expensive baseline
is not inserted into a gap-only covariance bound to claim improvement over
V23. The stronger covariance and memory conclusions already proved there
remain limited to their stated width regimes, with all their original
coefficients and normalizers unchanged.

At this checkpoint the main native finite-product and update-clock inputs
are isolated from the more demanding V16--V23 analytic chain. This does not
certify that chain independently or invalidate it. The supplied memoranda
remain proposals rather than evidence for their suggested physical arrows.
No assertion about an external Navier--Stokes theorem is a premise here.
Only the supplied f14/f15/Grand Tensor and current-lineage companions were
reviewed; no unavailable program is described as checked.

All resampling clocks in this section are auxiliary. No physical Gauss
projection, unrestricted retained-bridge integration, infrared contact floor,
same-top physical-transfer comparison, physical-time calibration or
nontrivial interacting continuum limit follows from
\eqref{f16v25:all-width-gap}. In particular its width dependence has not been
exchanged for a uniform continuum mass.

\clearpage
\subsection{Diagnostics, preservation, and precise verification status}

The diagnostic is
\begin{center}\small\path{verify_ym_f16_v25_interface_clock.py}.\end{center}
The new diagnostic passed 1,292 exact assertions. It checks two-copy
cancellation for finite many-body potentials, nonuniform product references,
marginal-to-conditional conversion, and normalized edge congestion. An exact
continuous polynomial-density example is included. The native geometry checks
use a stated dominating support envelope of the literal fine-lattice words;
they include preceding ports and periodic seams, not numerical Wilson
integrals. Counts and fixture scope are recorded in the report. These finite
checks do not enclose analytic constants or constitute proof-assistant
verification. All twenty-one available V4--V24 diagnostics were rerun with
their sources unchanged and passed. Their actual outputs and hashes are
included in the package.

The cumulative V24 is preserved byte for byte except for its one title-version
change and insertion of this exact appendix before the final document-ending
command. A reconstruction script checks the reverse operation. The package
contains the input hashes, checkpoint and analytic audit, executable finite
diagnostics, their actual outputs, and compilation/render checks. No inherited
proof or bibliography is silently corrected. The next companion checkpoint
is V30.

\begin{firewall}
V25 proves a continuous-product interface-cost variance comparison and applies
it to the original balanced native supports, including their preceding ports
and periodic seams. It sharpens the paid slab-to-local clock conversion and
obtains a positive duration-uniform original-sampler bound at every finite
spatial width, with explicit exponential decay in coupling times width. It
also identifies an exactly normalized, innovation-weighted interface
congestion. It does not prove a gap uniform in width or coupling, the missing
spatial boundary-susceptibility theorem, unrestricted coherent covariance
control, a physical transfer gap, or an interacting continuum mass gap.
\end{firewall}
% END F16 V25 APPEND

% BEGIN F16 V26 APPEND -- 2026-09-17
% Insert immediately before the final document-ending command in V25.
% No inherited proof, probability, source, or normalization is edited.

\clearpage
\section{V26: single-input spatial coupling on a controlled collar}
\label{f16v26:context}

V25 leaves two distinct targets: control its actual innovation-weighted
interface density, or construct a spatial coupling that pays the whole output
for one changed collar input. We first test the density route. Its exact
factorization shows that its second moment can grow exponentially even for a
product of uniformly gapped interacting pairs. This does not disprove a
native weighted-form estimate, but it rules out interpreting an unweighted
congestion moment as a proxy for the gap. We then pursue the second target.

The new estimate is a \emph{single-input, full-output coupling} of the same
native conditionals. It is uniform in ambient volume and duration, but it
requires a declared bound $Q$ on the rescaled fields of the actual collar.
Within that window, the expected number of different output cells is at most
a constant times the changed input displacement, plus a counted exponentially
small exception. All free cells, including their large fields, are retained.
A finite colour sweep is used only to smooth test functions inside an
auxiliary convex conditional; it is not substituted for the original sampler.
A positive influence ledger on these admitted rows is not asserted to hold on
the excluded boundary rows.

\subsection{Dependencies and the direction selected by the audit}

The complete V25 appendix, analytic audit, and checkpoint were read. Two scoped
semantic searches of the current source, f14 V100, f15 V100, and Grand Tensor
V076 returned no matches. Exact mounted-text searches followed. The monograph's
pressure-flow congestion concerns graph currents, not V25's two-copy density.
F14 V99's action-density susceptibility concerns a particular gauge-invariant
observable, not a coupling of entire conditional outputs. These distinctions,
and the existing V24 single-input criterion, are retained. This is a targeted
non-duplication review, not an independent certification of the archives.

The analytic inputs to the positive result are precisely V16/V20's
conditional set-energy inequality and V22's restricted native Hessian and
product-ball covariance realization. In particular, the absolute one-form
boundary realization with corners is inherited, not newly certified. V17's
all-order expansion and V23's exponential-width covariance theorem are not
used to prove this coupling. V25's better finite-width gap remains unchanged.
The next scheduled companion checkpoint is V30.

For context, Dorlas--Savoie,
\href{https://arxiv.org/abs/2305.19371}{\emph{On the Wasserstein distance and the
Dobrushin uniqueness theorem}}, formulates block dependence using transport
costs; its Section 5 makes the all-boundary hypothesis explicit. That theorem
is not invoked for the present boundary-restricted, inhomogeneous law. Its
transport definitions, block-criterion statement, and duality statement were
examined, not its complete proof. The coloured-sweep derivative argument and
the finite-product coupling conversion actually needed here are proved below.
The two supplied import memoranda remain methodological proposals. A local
power estimate is not silently promoted to global dependence or physical time.

\subsection{The splice density can be large in a uniformly gapped product}

Use the finite-product probability $\mu$ and coordinate order of V25. For
coordinate $i$ let $S=\{1,\ldots,i-1\}$, $T=S^c$, and let
$\mu_S,\mu_T$ denote its actual marginals. At fixed finite regulator the
positive bounded density hypotheses of V25 give all Radon--Nikodym derivatives
below. For its exact edge probability
\[
 d\mathbb B_i(z,b)=\mu(dz)\mu_i(db\mid z_{i^c}),
\]
keep \emph{exactly} V25's density $a_i$.

\begin{proposition}[A marginal factorization of the actual splice congestion]
\label{f16v26:factorization}
Up to $\mathbb B_i$-null sets,
\begin{equation}
 a_i(z,b)=
 \frac{d(\mu_S\otimes\mu_T)}{d\mu}(z)
 \frac{d\mu_i(\cdot\mid z_S)}{d\mu_i(\cdot\mid z_{i^c})}(b).
 \label{f16v26:congestion-factor}
\end{equation}
The first factor destroys the correlations across the ordering cut. The
second restores the correct conditional distribution of the incoming value.
Their product is globally normalized on the update-edge space, as in V25;
neither factor is dropped from a squared innovation.
\end{proposition}
\begin{proof}
Before the $i$th replacement, $z=(y_S,x_T)$ and the incoming value is $b=y_i$,
where $x,y$ are independent with law $\mu$. Their joint law is exactly
\[
 \mu_T(dz_T)\mu_S(dz_S)\mu_i(db\mid z_S).
\]
Its density relative to $\mathbb B_i$ is the right side of
\eqref{f16v26:congestion-factor}. V25's unestimated coordinate-exchange
identity identifies the same joint law as $a_i\,d\mathbb B_i$.
Uniqueness of the density proves the assertion. This is an identity between
normalized marginals; it does not replace any conditional by a product law.
\end{proof}

\begin{proposition}[Exponential congestion moments do not diagnose a small gap]
\label{f16v26:congestion-test}
Fix $0<\theta<1$. Take $M$ independent pairs $(X_j,Y_j)\in\{-1,1\}^2$ with
probability $(1+\theta X_jY_j)/4$, together with an independent uniform spin
$V$. Order the coordinates as $X_1,\ldots,X_M,V,Y_1,\ldots,Y_M$.
For the $V$ update, V25's density is
\begin{equation}
 a_V=\prod_{j=1}^M(1+\theta X_jY_j)^{-1},\qquad
 \mathbb E_{\mathbb B_V}a_V=1,\qquad
 \mathbb E_{\mathbb B_V}a_V^2=(1-\theta^2)^{-M}.
 \label{f16v26:product-moments}
\end{equation}
Nevertheless the sum of the original rate-one single-spin updates has gap
$1-\theta$, independently of $M$. A termwise inequality replacing $a_V$ by
a uniformly bounded constant against its own $V$-innovation is impossible.
This is a finite-spin counterfixture for that proposed inference, not a
counterexample to native Yang--Mills mixing.
\end{proposition}
\begin{proof}
At the cut just before $V$, the two coordinate marginals are products of
uniform spins. The incoming $V$ law is uniform under both conditionings.
Equation \eqref{f16v26:congestion-factor} gives the product displayed above.
For each pair, $XY=\pm1$ with probabilities $(1\pm\theta)/2$. Its inverse
factor has mean one and second moment $(1-\theta^2)^{-1}$. Independence proves
\eqref{f16v26:product-moments}.

One pair has generator eigenvalues $0,1-\theta,1+\theta,2$: its linear
modes are $X+Y,X-Y$, and the remaining centered mode is $XY-\theta$.
The independent $V$ generator has nonzero eigenvalue one. Tensor products
and sums of these commuting generators therefore give gap $1-\theta$.
For the last assertion set $H=\{X_jY_j=-1\text{ for all }j\}$ and
$f=V1_H$. Its $V$ squared increment is supported on $H$, where
$a_V=(1-\theta)^{-M}$. The ratio of its weighted and unweighted $V$-edge
integrals is exactly that number. The other spin innovations are not zero;
this example does not rule out a comparison using their energies too.
\end{proof}

Thus a bound on a raw congestion moment is not selected as the next sufficient
criterion. The remaining argument follows the spatial single-input target
itself, with its boundary window stated rather than hidden in an average.

\subsection{Preparation for a bounded rescaled collar, including boundary cells}

Let $D$ be a nonempty finite set of complete \emph{unmoving} cells, let
$C=\partial_{\mathcal G}D$ be its actual interaction collar, and put $N_D=|D|$.
Let $P_D^q$ denote V20's complete native conditional, with all terms meeting
$D$ and all original endpoint factors retained. Only $q_C$ affects this law.
Assume
\begin{equation}
 \max_{c\in C}|q_c|\le Q,\qquad Q\ge0,\qquad
 \max_{i,e}|y_{i,e}|\le r.
 \label{f16v26:collar-window}
\end{equation}
No value farther away than $C$ is restricted. Write $e_c=|\zeta_c|^2$ and
retain $t_0=1/(24\pi^2)$. The parameter $Q$ is a bound on \emph{rescaled}
whole-cell Euclidean coordinates, not an angular bound independent of coupling.

\begin{proposition}[Uniform full conditional moments at a controlled collar]
\label{f16v26:prepared-Q}
For the signed source cube $|s_c|\le t_0$ on $D$, the same native conditional
satisfies
\begin{equation}
 \sup_{c\in D}\mathbb E_{P_{D,s}^q}e_c\le A_r(1+Q^2),\qquad
 \log\mathbb E_{P_D^q}\exp\!\left(\sum_{c\in D}t_ce_c\right)
       \le A_r(1+Q^2)\sum_{c\in D}t_c\quad(0\le t_c\le t_0).
 \label{f16v26:conditional-moments}
\end{equation}
Here and below constants do not depend on $D,B,n,Q,\gamma$, unless explicitly
indicated, and $\gamma\ge1+r^2$. In particular, for V22's fixed comparison
radius $\eta_*$, there are $c_*>0,A_r<\infty$ such that
\begin{equation}
 P_D^q\{\exists c\in D:\ e_c>\eta_*^2\gamma\}
 \le N_D\exp\{A_r(1+Q^2)-c_*\gamma\}.
 \label{f16v26:conditional-tail}
\end{equation}
Unlike arbitrary-compact-boundary preparation, this estimate needs no depth
allowance: it includes the first free layer next to the bounded collar.
\end{proposition}
\begin{proof}
The inherited local normalizer inequality, averaged under $P_{D,s}^q$, gives
for every $T\subseteq D$
\[
 \sum_{c\in T}u_c\le a_r|T|+b_r\sum_{d\in\partial T}\widetilde u_d,
 \qquad u_c=\mathbb E_{P_{D,s}^q}e_c.
\]
On $D$ put $\widetilde u=u$, and outside $D$ set $\widetilde u=Q^2$.
This dominates each frozen collar energy that can occur in the inequality;
values assigned beyond the collar are an auxiliary scalar extension, not a
change of the probability or its boundary. For an arbitrary ambient cell set
$S$, apply the inequality to $T=S\cap D$. Boundary terms lying in $S\setminus
D$ cost at most $b_rQ^2|S\setminus D|$; the others lie in $\partial S$.
Consequently
\[
 \sum_{c\in S}\widetilde u_c
 \le [a_r+(1+b_r)Q^2]|S|+b_r\sum_{d\in\partial S}\widetilde u_d.
\]
Iterate on ambient graph balls as in V16, with
$\rho_r=b_r/(1+b_r)<1$. Polynomial ball growth bounds the resulting series
$\sum_{k\ge0}\rho_r^k(k+1)^4$. At the fixed finite regulator the terminal
term tends to zero, since $\widetilde u\le\max(E_*\gamma,Q^2)$ and the ambient
vertex set is finite. This proves the first bound, uniformly throughout the
signed source cube. Integrating the \emph{actual} derivatives of the conditional
$\log Z_D(q;s)$ along $s=ut$, $0\le u\le1$, proves the second. Exponential
Markov in each free cell and a union bound prove
\eqref{f16v26:conditional-tail}, after absorbing $t_0$ in $A_r$ and taking
$c_*=t_0\eta_*^2$. No conditional partition function is approximated.
\end{proof}

\subsection{A finite sweep turns Hamming oscillations into local derivatives}

For the comparison only, use exactly V22's restricted native law $\pi_D^q$
on the fixed product $\mathcal D_{\gamma,D}$ of whole-cell balls of radius
$\eta_*\sqrt\gamma$. At its original coupling threshold and for collars in
these balls, retain its constants $\lambda,M>0$, bounded graph degree $\Delta$,
bounded mixed Hessian blocks, and absolute one-form covariance identity.
In particular all one-cell conditionals have Poincar\'e constant at most
$\lambda^{-1}$, and the inverse one-form blocks are at most
$\lambda^{-1}e^{-a_0d_{\mathcal G}(c,d)}$ for a fixed $a_0>0$.
These are the explicitly inherited analytic inputs, not full-law convexity.

For a real bounded function $f$ define its cell oscillation by
\[
 o_c(f)=\sup_{x_{c^c}=x'_{c^c}}|f(x)-f(x')|.
\]
Choose a proper colouring of the free interaction graph with $s\le\Delta+1$
colours $C_1,\ldots,C_s$. Each primitive carrier is a clique in that graph,
so, given their exterior, cells in one colour are conditionally independent.
Let $E_{C_l}^q$ be their exact conditional expectation and define
$f_0=f$, $f_l=E_{C_l}^qf_{l-1}$, $g_q=f_s$.

\begin{proposition}[A local derivative estimate with no derivative norm of the test]
\label{f16v26:sweep}
There are fixed $h,C_r<\infty$ such that, whenever $f$ is smooth and independent
of $q$, and $o_c(f)\le w_c$ with $w_c\ge0$,
\begin{align}
 \pi_D^qg_q&=\pi_D^qf,\notag\\
 \|\nabla_cg_q\|_\infty
       &\le C_r\sum_{d\in D:\ d_{\mathcal G}(c,d)\le h}w_d,\notag\\
 \|\nabla_{q_j}g_q\|_\infty
       &\le C_r\sum_{d\in D:\ d_{\mathcal G}(j,d)\le h}w_d,
                    \qquad j\in C.
 \label{f16v26:sweep-gradients}
\end{align}
The constants depend on the local analytic bounds, not the number of colours
actually occupied, the size of $D$, the radius $\eta_*\sqrt\gamma$, or any
initial derivative of $f$. This sweep is used for smoothing and invariance
only; neither its spectral gap nor its contraction is assumed.
\end{proposition}
\begin{proof}
For a colour $C_l$ and $j\notin C_l$, changing cell $j$ changes only the
conditional factors in $C_l\cap N(j)$. Product coupling of those factors,
with no better bound than probability one of a mismatch, gives
\begin{equation}
 o_j(E_{C_l}F)\le o_j(F)+\sum_{i\in C_l\cap N(j)}o_i(F),
       \qquad o_j(E_{C_l}F)=0\quad(j\in C_l).
 \label{f16v26:osc-propagation}
\end{equation}
A telescoping change of the output coordinates proves the same inequality
without selecting couplings. Thus after at most $s$ layers every oscillation
is bounded by a fixed local sum of the original $w$'s: the propagation range
increases by at most one per layer, with row and column sums at most
$(1+\Delta)^s$.

Differentiation of a \emph{normalized} one-cell conditional gives
\[
 \nabla_jE_iF=E_i\nabla_jF-\operatorname{Cov}_i(F,\nabla_jV),\qquad j\ne i.
\]
For a unit vector $v$ in cell $j$, conditional covariance Cauchy--Schwarz and
the one-cell Poincar\'e inequality give
\[
 |\operatorname{Cov}_i(F,\partial_{j,v}V)|
 \le \tfrac12 o_i(F)\lambda^{-1/2}
                  \|\nabla_i\partial_{j,v}V\|_\infty
 \le C_r o_i(F).
\]
This is zero for nonneighbours. For a whole colour, differentiate its product
conditional; the same terms are summed only over $C_l\cap N(j)$, with the
other conditional factors integrated. Hence
\begin{equation}
 \|\nabla_jE_{C_l}F\|_\infty
 \le\|\nabla_jF\|_\infty+C_r\sum_{i\in C_l\cap N(j)}o_i(F)
                       \quad(j\notin C_l).
 \label{f16v26:gradient-propagation}
\end{equation}
When $j\in C_l$ the output is independent of its old value, so its derivative
is zero. Each free cell has such a layer. All dependence on the initial
$\nabla_jf$ is therefore erased before the finite sweep ends. Subsequent terms
in \eqref{f16v26:gradient-propagation} and
\eqref{f16v26:osc-propagation} prove the free derivative bound with, for example,
$h=2s+2$. A collar derivative is never erased, but starts at zero because $f$
is independent of $q$. The same recursion, with only the finitely many free
neighbours of $j$ contributing, proves the collar bound. Each colour
expectation preserves $\pi_D^q$, proving invariance. All derivatives keep the
original measured Haar terms and the conditional normalizers. No uniform
Hessian bound for an arbitrary nonsmooth test has been used.
\end{proof}

\subsection{One changed collar input has a localized full-output response}

For weights $w=(w_c)_{c\in D}\ge0$, put
\[
 d_w(x,x')=\sum_{c\in D}w_c1_{\{x_c\ne x'_c\}},\qquad
 \mathsf W_w(\mu,\nu)=\inf_{\Pi\in\mathrm{Cpl}(\mu,\nu)}\int d_w\,d\Pi.
\]
This is a cost on whole-cell outputs, not Euclidean displacement cost. In
particular $\mathsf W_{\boldsymbol1}$ is the minimum expected number of unequal
cells. Separate optimal couplings of the scalar marginals would not bound it.

\begin{theorem}[Weighted boundary response in the restricted native law]
\label{f16v26:restricted-response}
For two collar values $q,q'$ in the same small-angular product, differing only
at cell $j\in C$, every bounded measurable $f$ with $o_c(f)\le w_c$ satisfies
\begin{equation}
 |\pi_D^{q'}f-\pi_D^qf|
       \le C_r|q'_j-q_j|\sum_{c\in D}w_c e^{-a d_{\mathcal G}(c,j)}
 \label{f16v26:weighted-response}
\end{equation}
for one fixed $a>0$. Its constants are independent of the block and the
collar size. No global gap of a complete native sampler is an input.
\end{theorem}
\begin{proof}
First take a smooth test. Along the collar segment $q(t)$ let $v=q'_j-q_j$,
$S_t=\partial_tV_D^{q(t)}$, and use the \emph{parameter-dependent} invariant
sweep $g_t$ above. Keeping both its derivative and the density derivative gives
\begin{equation}
 \frac{d}{dt}\pi_D^{q(t)}f
     =\pi_D^{q(t)}(\partial_tg_t)
                        -\operatorname{Cov}_{\pi_D^{q(t)}}(g_t,S_t).
 \label{f16v26:swept-response}
\end{equation}
The first term must not be dropped: the smoothing kernel depends on the
boundary. Its absolute value is at most
$C_r|v|\sum_{d_{\mathcal G}(c,j)\le h}w_c$ by
\eqref{f16v26:sweep-gradients}.

The free gradient of $S_t$ has support within one graph step of $j$ and each
block has norm at most $C_r|v|$. Insert the inherited one-form covariance
identity, its componentwise exponential inverse bound, and
\eqref{f16v26:sweep-gradients}. For a fixed $0<a<a_0$ this bounds the second
term by
\[
 C_r|v|\sum_{x\in D}\sum_{c:\ d_{\mathcal G}(c,x)\le h}
                             w_c e^{-a_0d_{\mathcal G}(x,j)}
 \le C_r|v|\sum_{c\in D}w_c e^{-a d_{\mathcal G}(c,j)}.
\]
The interchange of sums uses only a fixed-radius ball count. The finite-range
first term obeys the same bound after increasing $C_r$. Integrate in $t$.
Every intermediate collar remains in its declared convex product domain.

To extend to bounded measurable tests, extend $f$ outside each closed ball by
coordinatewise nearest-point projection and convolve with a product of smooth
approximate identities. The extension and convolution do not increase any
cell oscillation. At almost every interior point the approximation converges
to $f$; both endpoint probabilities have Lebesgue densities. Bounded dominated
convergence therefore applies in their joint $L^1$ space. The estimate has no
derivative norm of the test and survives the limit. The finite regulator is
held fixed during this approximation; its constants in
\eqref{f16v26:weighted-response} are already size independent.
\end{proof}

\begin{proposition}[From all weighted tests to one spatially localized coupling]
\label{f16v26:coupling-conversion}
There is a coupling $(X,X')$ of $\pi_D^q,\pi_D^{q'}$ such that simultaneously
for every $c\in D$,
\begin{equation}
 \mathbb P\{X_c\ne X'_c\}
       \le C_r|q'_j-q_j|e^{-a d_{\mathcal G}(c,j)}.
 \label{f16v26:restricted-coupling}
\end{equation}
Consequently
\begin{equation}
 \mathsf W_{\boldsymbol1}(\pi_D^q,\pi_D^{q'})\le C_r|q'_j-q_j|.
 \label{f16v26:restricted-total}
\end{equation}
This is one coupling for a \emph{fixed pair} of complete boundary rows, not a
grand coupling over all boundary values.
\end{proposition}
\begin{proof}
Here is a finite-product version of the transport step, to avoid changing the
measurable spaces to the nonseparable discrete topology. Partition each compact
cell ball by nested finite grids of mesh tending to zero, choosing grid faces
of zero mass for both one-cell marginals. Weighted Hamming cost on the finite
product of labels has finite Kantorovich dual equal to the supremum over
functions of coordinate oscillation at most $w_c$. Lifting such a function to
the continuous product preserves these oscillations, so
\eqref{f16v26:weighted-response} bounds the optimal finite cost by
$\sum_cw_c b_c$, where $b_c=C_r|q'_j-q_j|e^{-a d(c,j)}$.
Lift an optimal label coupling using the actual conditional distributions in
each product atom. On the compact product, extract a weak limit of these
couplings. For any fixed grid, the cost is continuous except on marginal-null
faces, so its integral in the limit is at most $\sum_cw_cb_c$. The grid costs
increase to the full weighted Hamming cost; monotone convergence proves
$\inf_\Pi\sum_cw_c\Pi\{X_c\ne X'_c\}\le\sum_cw_cb_c$.

For completeness this implies one coupling satisfying all component bounds.
The set of vectors dominating the component disagreement probabilities of some
coupling is convex and upward closed. It is closed: compactness of the
couplings and lower semicontinuity of each open-set indicator
$1_{\{X_c\ne X'_c\}}$ preserve each limiting upper constraint. If the vector
$b$ were outside this set, finite-dimensional separation would give a
nonnegative weight $w$ for which every coupling has cost strictly greater than
$\sum_cw_cb_c$, contradicting the preceding transport bound. This proves
\eqref{f16v26:restricted-coupling}. Polynomial graph growth makes
$\sup_j\sum_c e^{-ad(c,j)}$ finite, uniformly in $D$ and the finite ambient
history. Summing proves \eqref{f16v26:restricted-total}. All conditionals in
this argument retain their complete restricted joint law; scalar marginal
couplings were not independently combined.
\end{proof}

\subsection{The full native block: the restriction is paid in the coupling}

We now return to $P_D^q$ on its entire original compact domain. Assume
\eqref{f16v26:collar-window} for $q$ and $q'$, differing only at $j$, and require
\begin{equation}
 Q\le\eta_*\sqrt\gamma,
       \qquad \gamma\ge\max\{1+r^2,\gamma_r^{\rm box}\}.
 \label{f16v26:entry}
\end{equation}
This is a genuine restriction on the deterministic collar. It is not inherited
from V20's all-compact-boundary comparison and must not be omitted. Set
\[
 \varepsilon_D(Q)=
    \min\{1,N_D\exp[A_r(1+Q^2)-c_*\gamma]\}.
\]

\begin{theorem}[Single-input full-output susceptibility on a controlled collar]
\label{f16v26:full-coupling}
There is a coupling $(Y,Y')$ of the \emph{complete} native laws $P_D^q,P_D^{q'}$
for which all cells satisfy
\begin{equation}
 \mathbb P\{Y_c\ne Y'_c\}
 \le C_r|q'_j-q_j|e^{-ad_{\mathcal G}(c,j)}+2\varepsilon_D(Q).
 \label{f16v26:full-component}
\end{equation}
In particular the full expected resampled disagreement has the bound
\begin{equation}
 \boxed{
 I_{D,j}(q,q'):=\mathsf W_{\boldsymbol1}(P_D^q,P_D^{q'})
 \le \min\{N_D,\ C_r|q'_j-q_j|
                    +2N_D^2 e^{A_r(1+Q^2)-c_*\gamma}\}.}
 \label{f16v26:full-output}
\end{equation}
The constants are independent of ambient volume, duration, block shape, and
its number of cells. That number occurs only in the stated exceptional term.
No fast variable inside $D$, including its fringe, has been fixed or omitted.
\end{theorem}
\begin{proof}
Let $G_D$ be the event that every free cell belongs to its comparison ball.
It has positive probability. The exact identity
\[
       P_D^q(\,\cdot\mid G_D)=\pi_D^q
\]
holds with the same original action, crossing words, and product references.
Equation \eqref{f16v26:conditional-tail} gives
$d_{\rm TV}(P_D^q,\pi_D^q)=P_D^q(G_D^c)\le\varepsilon_D(Q)$, and the same
holds at $q'$. Couple each full row to its restricted row by a whole-vector
maximal coupling. Between the two restricted rows use the one coupling in
Proposition~\ref{f16v26:coupling-conversion}. Disintegration glues these three
couplings while preserving every full and restricted marginal. For each cell,
a mismatch of the full outputs requires either a restricted-pair mismatch or
one of the two whole-vector coupling failures. This gives
\eqref{f16v26:full-component}. Summation, the exponential graph-row bound, and
the trivial cost $N_D$ prove \eqref{f16v26:full-output}. On either exceptional
branch the full field remains drawn from its actual conditional distribution.
No independence of failures, changed exterior law, or division by an estimated
good-event probability is used.
\end{proof}

For fixed $Q$ and $N_D\le e^{\beta\gamma}$ with any fixed $0<2\beta<c_*$,
the exceptional total cost is exponentially small, at a threshold depending
on $Q,\beta,r$. More generally if $Q_\gamma=\gamma^\alpha$ with
$0<\alpha<1/2$ and $N_D$ is any prescribed polynomial in $\gamma$, then
\begin{equation}
 I_{D,j}(q,q')\le C_r|q'_j-q_j|+e^{-c_*\gamma/2}
 \label{f16v26:moderate-output}
\end{equation}
eventually, uniformly throughout that collar window. Constants multiplying
the polynomial and its exponent enter the threshold. For a single changed
input $|q'_j-q_j|\le2Q_\gamma$, this is a side-length-independent leading
susceptibility $2C_rQ_\gamma$, not the whole-fringe cost of V24.

\subsection{The admitted-row influence ledger is positive; the excluded rows remain}

For clarity we give the actual finite-geometry count used to compare the new
budget to V24. Choose a fixed integer $\rho\ge1$ bounding every native cell
interaction in the ambient coordinatewise torus/path distance. On a spatial
cycle of length $m>R$ use every translated cyclic interval of $R$ cells; if
$m\le R$ use the whole cycle once. On the time path use V24's indexed clipped
intervals of length $R$. Form product blocks $D$. Give each product index rate
$R^{-(k+1)}$, where $k$ is the number of cut spatial directions. In this
isotropic torus $k$ is zero or three; the notation also makes the counting
transparent. Repeated clipped blocks keep their indices and rates.

\begin{proposition}[A finite cover and a positive moderate-collar table]
\label{f16v26:moderate-ledger}
For $R\ge2\rho$, every cell has coverage rate one, $|D|\le R^4$, and the total
rate of blocks not containing a given cell $j$ whose conditional can depend on
$j$ is at most $C_\rho/R$. Fix $0<\alpha<1/2$. There is a fixed $K_r$ such that
with $Q_\gamma=\gamma^\alpha$ and $R_\gamma=\lceil K_r(1+Q_\gamma)\rceil$,
the single-input table supplied by \eqref{f16v26:full-output}, \emph{only on
pairs of collars within this window}, satisfies
\begin{equation}
 1-\sum_{D:\,j\notin D}r_DI_{D,j}(q,q')\ge\tfrac12
 \label{f16v26:good-ledger}
\end{equation}
at a fixed-$\alpha,r$ large-coupling threshold. Equation
\eqref{f16v26:good-ledger} is not a global native block-generator inequality:
its table is not proved on arbitrary compact exterior rows.
\end{proposition}
\begin{proof}
A site belongs to $R$ indexed intervals in each cut direction, including at
both time endpoints, and to one interval in an uncut cycle. Product rates
therefore give unit coverage. If $j$ is outside a product block but affects it,
some cut coordinate puts it outside its interval at distance at most $\rho$.
There are at most $2\rho$ such interval indices in that coordinate. In each
other cut coordinate there are at most $R+2\rho$ indices whose interval comes
within $\rho$ of $j$; an uncut cycle contributes one. Summing over at most four
coordinates and dividing by the coverage count gives
\[
 \sum_{D:\,j\notin D,\ j\in\partial D}r_D
       \le\frac{8\rho}{R}(1+2\rho/R)^3\le C_\rho/R.
\]
This counts periodic seams and clipped time intervals. It is an upper bound
for the literal word collar, not an assertion that every nearby box interacts.
The size estimate follows since each interval has length at most $R$.

Take $K_r\ge8C_\rho C_r+2\rho+1$, enlarging it if necessary for fixed
constants. Then the leading birth cost from
$|q'_j-q_j|\le2Q_\gamma$ is at most $1/4$. The exceptional cost per block is
at most $2R_\gamma^8\exp[A_r(1+\gamma^{2\alpha})-c_*\gamma]$, which tends to
zero exponentially since $2\alpha<1$. Its aggregate is at most $1/4$
eventually. The entry condition \eqref{f16v26:entry} also holds eventually.
This proves the restricted-row ledger and no more. The all-row hypothesis of
Theorem~\ref{f16v24:general-gap} is still missing, so it is not applied.
\end{proof}

Under the \emph{actual full native law} V16 does give, for each block's collar,
\begin{equation}
 P\{\max_{c\in\partial D}|\zeta_c|>Q_\gamma\}
       \le |\partial D|\exp[t_0M_r-t_0\gamma^{2\alpha}].
 \label{f16v26:bad-collar-prob}
\end{equation}
For any family with disjoint collars, its simultaneous bad-collar probability
is at most the product of the corresponding right sides: expand choices of
one witness per collar and apply V16's joint exponential moment to those
distinct cells. This needs neither independence nor a product posterior.
The statement is not made after arbitrary exterior conditioning.

Neither \eqref{f16v26:bad-collar-prob} nor the positive admitted-row table
bounds the omitted squared innovations for all spectral test functions.
Arbitrary normalized functions may concentrate on the excluded collars, and
field-dependent gating would require an invariance and commutator calculation.
The new coupling closes the tangential whole-fringe loss \emph{inside the
stated boundary class}; it does not erase that class from the hypotheses.

\subsection{Remaining target and exact scope of this continuation}

There are two genuine advances. The congestion test identifies a failure of
one tempting norm in an exactly gapped model. More importantly, the new spatial
coupling controls every output cell at once after one collar cell changes,
first with exponential localization in the restricted native comparison and
then under the full compact conditional with its exception paid. It does not
merely sum unrelated scalar total variations. Ambient volume and history
length are unrestricted throughout; the deterministic collar window and the
finite-block exception are explicit.

The next noncircular target is the remaining \emph{rough-collar contribution}
in a same-law, all-test-function estimate, or a boundary repair/exploration
that extends the single-input coupling to those rows. Increasing an expansion
order does not supply it. Replacing the true conditional rows by their
small-field restrictions, inserting their rare-event probability independently
of a squared innovation, or applying V24's all-row theorem to only the good
rows would all skip a required step. The coupled outputs do not define a new
physical transfer or a new field law.

V25's all-finite-width baseline and V23's direct covariance/memory results
retain their original scopes. No new unrestricted covariance operator bound,
uniform original-clock sampler gap, or improvement of physical spectral ratios
is asserted. The prescribed bridge window, physical Gauss reduction, infrared
contact coercivity, same-top transfer comparison, physical-time calibration,
and interacting continuum construction remain separate obligations.

\subsection{Diagnostics, preservation, and verification boundary}

The accompanying diagnostic is
\begin{center}\small\path{verify_ym_f16_v26_spatial_coupling.py}.\end{center}
It tests the exact marginal formula for congestion; the gapped product-pair
obstruction; normalized conditional derivative and sweep bookkeeping; finite
transport duality and simultaneous disagreement constraints; the full/restricted
coupling mixture; and the clipped finite-cover counts. Analytic inequalities
in the preceding proofs are not certified by the finite fixtures. In
particular no native compact-integral enclosure, numerical value of the
one-form constants, or proof-assistant verification is claimed. The precise
executed test families and available predecessor reruns are recorded in the
machine-readable reports, rather than used as a measure of proof completeness.

The cumulative source changes only V25's title version and inserts this exact
appendix before its final document-ending command. Removing that insertion and
restoring the title recovers V25 byte for byte. The package records preservation,
actual diagnostic outputs, compilation, and rendered-page inspection. These
new deductions are conditional on their identified inherited analytic inputs;
the entire historical lineage has not been independently recertified. The
next scheduled companion checkpoint remains V30.

\begin{firewall}
V26 proves a full-output single-input spatial coupling on a declared
rescaled-collar window, with a localized restricted-native part and an explicitly
paid full-native exception. It proves a positive influence ledger only for
those admitted boundary pairs, not a global native gap. It also shows that
unweighted splice congestion can grow exponentially in a uniformly gapped
product model. No arbitrary-compact-boundary spatial coupling, control of the
excluded innovation-weighted defect, unrestricted coherent covariance theorem,
physical transfer gap, or interacting continuum mass gap is established.
All original laws, normalizers, source definitions and clocks are retained.
\end{firewall}
% END F16 V26 APPEND

% BEGIN F16 V27 APPEND -- 2026-09-17
% Insert immediately before the final document-ending command in V26.
% No inherited statement, normalizer, source, or update clock is changed.

\clearpage
\section{V27: buffered rough-collar spectral localization}
\label{f16v27:context}

V26 controls a whole conditional output after one collar input changes, but
only on a declared collar window. The probability of its excluded collars is
small under the native law; an arbitrary squared spectral test can nevertheless
concentrate there. We do not multiply that probability by an unrelated
innovation. Instead, we use a buffered conditional expectation to separate an
arbitrary test into an exterior-measurable part and a part paid by an actual
resampling energy. Uniform conditional preparation controls the first part.
A local, interface-priced comparison pays the second in the original clock.

This proves an all-test-function rough-collar estimate and an operator bound
on its intersection with every low-energy spectral subspace. A family of
separated repair regions also gives a simultaneous rough-count estimate with
no factor counting the regions in its energy term. The target remains the
complete native law. These are localization and killed-form statements, not
a global gap theorem. The original-clock comparison is expensive, and an
extensive term still prevents the final rough-innovation bound from closing
V26's unrestricted spatial argument.

\subsection{Dependencies and the distinction from an unweighted tail}

The complete V26 appendix and audit were read. The new proof uses V20's
all-compact-boundary preparation, V24's exact moving/unmoving support inclusion,
and V25's finite-product interface comparison. It does not use V26's restricted
one-form response theorem to assume that rough collars are already repaired.
It needs neither a global native gap nor the arbitrary-order source expansions.
Every conclusion is a deduction from the expressly identified inherited inputs,
not an independent recertification of their full analytic lineage.

Two semantic archive searches were empty. Exact mounted-text searches and
identified source ranges were then checked. F14 V1 already discusses
original-clock killed repair on specified visible events; its negative-curvature
geometry is not transplanted here. The Grand Tensor monograph already separates
continuous killing from endpoint sampling. Those principles are credited. The
new event family consists of arbitrary rough cells and their collars in the
present fixed-bridge history, with preparation uniform over every exterior of
the repair region. The energy price and the number of repair regions are
controlled separately. The next scheduled checkpoint remains V30.

For context, Goel--Montenegro--Tetali,
\href{https://arxiv.org/abs/math/0505690}{\emph{Mixing Time Bounds via the Spectral
Profile}}, Definition 1.4 and (1.4), distinguish a restricted Dirichlet eigenvalue
from a global gap. The definitions and adjacent argument were examined; their
finite-state mixing-time theorem is not invoked on this continuous law.
Spinka's \href{https://arxiv.org/abs/1803.10578}{\emph{Finitary codings for
spatial mixing Markov random fields}}, Theorem 1.2, requires common overlap
along unbounded spatial scales. Its finite-valued coding result is not supplied
by V23's fixed-coupling finite-buffer bound. The present step estimates spectral
weights instead of asserting that this missing coupling composition holds.

\subsection{One conditional projection and a genuinely weighted rare event}

On an arbitrary probability space let $E$ be conditional expectation onto a
sigma algebra $\mathcal F$, $Q=I-E$, and let $M_H$ multiply by the indicator
of an event $H$. Suppose
\begin{equation}
                 P(H\mid\mathcal F)\le p\quad\hbox{a.s.},\qquad 0\le p\le1.
 \label{f16v27:conditional-smallness}
\end{equation}
The event need not be independent of $\mathcal F$.

\begin{proposition}[Conditional smallness with its energy debit]
\label{f16v27:projection}
For every real or complex $f\in L^2(P)$,
\begin{equation}
 \|M_HE\|\le\sqrt p,\qquad
 \|M_Hf\|_2\le\|Qf\|_2+\sqrt p\,\|Ef\|_2
                  \le\|Qf\|_2+\sqrt p\,\|f\|_2.
 \label{f16v27:projection-bound}
\end{equation}
On functions supported on $H$ there is the sharper form estimate
\begin{equation}
       \|Qf\|_2^2\ge(1-p)\|f\|_2^2,\qquad f=M_Hf.
 \label{f16v27:killed-projection}
\end{equation}
These statements neither replace a probability by an energy nor assume an
inequality for every mean-zero function.
\end{proposition}
\begin{proof}
For $g=Eg$, conditional expectation gives
$\|M_Hg\|_2^2=\mathbb E[P(H\mid\mathcal F)|g|^2]\le p\|g\|_2^2$.
Decompose $f=Ef+Qf$, use this estimate on the first component and the
contraction $\|M_HQf\|_2\le\|Qf\|_2$ on the second. This proves
\eqref{f16v27:projection-bound}. Taking adjoints gives $\|EM_H\|\le\sqrt p$.
For $f=M_Hf$, orthogonality of $E,Q$ consequently gives
$\|Qf\|_2^2=\|f\|_2^2-\|Ef\|_2^2\ge(1-p)\|f\|_2^2$.
No differentiability, finite alphabet, or lower bound on a point mass is used.
\end{proof}

For example, if $P(H)=p>0$ and $f=p^{-1/2}1_H$, then
$\|f\|_2=\|1_Hf\|_2=1$. Its unweighted event probability alone says nothing
small about that weighted mass. When \eqref{f16v27:conditional-smallness}
holds, \eqref{f16v27:killed-projection} instead forces an actual conditional
innovation of size at least $1-p$. Constants are also consistent with
\eqref{f16v27:projection-bound}: their energy is zero and their event mass is
paid by the probability term.

\subsection{Prepare each rough-cell event inside a local repair region}

Retain exactly $P=P_{\gamma,\mathbf y}$ on the complete unmoving charts, with
$B=m^3$, $m\ge4$, $n\ge1$, $\max_{i,e}|y_{i,e}|\le r$, and
$\gamma\ge1+r^2$. All norms below refer to $L^2(P)$. The original balanced
generator, transported unitarily to this space, is
\[
 \widehat{\mathcal L}=\sum_\alpha\widehat Q_\alpha,\qquad
 \widehat{\mathcal E}(f)=\langle f,\widehat{\mathcal L}f\rangle.
\]
Every original five-chord cube or single nonroot port still has rate one.
For a set $S$ of unmoving cells let $E_S$ condition on its complete exterior,
$Q_S=I-E_S$, and $\mathcal E_S(f)=\|Q_Sf\|_2^2$.

Let $e_c=|\zeta_c|^2$ and define the same type of rough-cell event as V26 by
\[
       H_c(Q)=\{e_c>Q^2\},\qquad Q\ge0.
\]
Fix a positive integer $r_0$ bounding the coordinatewise space--time displacement
of every edge of the native interaction graph $\mathcal G$. Denote V20's
preparation constants by $\kappa_r,B_r,M_r$. Choose
\begin{equation}
 b_\gamma=2+\left\lceil\kappa_r^{-1}\log\gamma\right\rceil,\qquad
 w_\gamma=r_0(b_\gamma+2),\qquad h_\gamma=2w_\gamma+5.
 \label{f16v27:repair-size}
\end{equation}
Let $S_c$ be the product of the coordinatewise spatial-cycle balls and the
clipped time interval of radius $w_\gamma$ around $c$. A short spatial cycle
is taken in full, once. Natural time endpoints are retained, not frozen.
Then $c$ is at graph distance at least $b_\gamma$ from $S_c^c$, unless its
component is entirely free, in which case the distance is infinite.

\begin{proposition}[A roughness bound uniform in the actual repair exterior]
\label{f16v27:uniform-tail}
With $\overline M_r=M_r+B_r$ and $t_0=1/(24\pi^2)$, put
\begin{equation}
       p_\gamma(Q)=\min\{1,\exp[t_0(\overline M_r-Q^2)]\}.
 \label{f16v27:p}
\end{equation}
The exact conditional-density versions satisfy, for every compact exterior $q$,
\begin{equation}
 P_{S_c}^{q}(H_c(Q))\le p_\gamma(Q),\qquad
 \mathbb E_{P_{S_c}^{q}}e_c^k
       \le M_{k,r}:=k!t_0^{-k}e^{t_0\overline M_r}\quad(k\ge1).
 \label{f16v27:uniform-preparation}
\end{equation}
There is no small-field restriction on $q$ or on the free repair output.
The radius $h_\gamma=O_r(1+\log\gamma)$ is independent of $Q,B,n$.
\end{proposition}
\begin{proof}
V20's exact conditional moment-generating estimate, with a source only at
$c$, bounds $\mathbb E_{P_{S_c}^{q}}e^{t_0 e_c}$ by
\[
       \exp\{t_0(M_r+B_r\gamma e^{-\kappa_r b_\gamma})\}
                         \le e^{t_0\overline M_r}.
\]
A graph step changes each ambient coordinate by at most $r_0$, so the repair
box has the claimed depth, including periodic identifications and clipped
natural ends. Exponential Markov and $x^k\le k!t_0^{-k}e^{t_0x}$ give the
two conclusions. The result concerns the complete conditional and includes
its exact normalizer, Haar powers, and every crossing term. Its uniformity in
$q$ is inherited from preparation, not from an averaged bad-event estimate.
\end{proof}

For $Q_\gamma=\gamma^\alpha$, $0<\alpha<1/2$ fixed, one has
$p_\gamma(Q_\gamma)\le e^{-(t_0/2)\gamma^{2\alpha}}$ eventually.
This is V26's growing rescaled threshold, but the exterior of $S_c$ is now
arbitrary: the observed rough cell sits inside an integrated repair buffer.
No all-row susceptibility estimate for the original block $D$ has thereby
been inferred.

\subsection{The repair has a local original-clock cost}

Let $\widetilde S_c$ denote the elementary balanced support needed for the
unmoving update $S_c$. It contains each chord cube in $S_c$ and the moving
ports at times $t$ for which that cube's cell at $t$ or $t+1$ is in $S_c$.
Thus its preceding ports are included exactly as in V24. All these statements
are about support inclusion, not equality of balanced and unmoving kernels.

\begin{proposition}[Local interface price and overlap, independent of ambient size]
\label{f16v27:clock}
There are constants independent of $B,n,\gamma,c$ such that, with
$h=h_\gamma$,
\begin{equation}
 \begin{split}
 |\widetilde S_c|&\le C h^4,\\
 \mathcal E_{S_c}(f)&\le K_0(\gamma)
            \sum_{\alpha\in\widetilde S_c}\|\widehat Q_\alpha f\|_2^2,
       \qquad K_0(\gamma)=C_r h^4e^{C_r\gamma h^3},\\
 \sup_\alpha\#\{c:\alpha\in\widetilde S_c\}&\le\varkappa_\gamma:=C h^4.
 \end{split}
 \label{f16v27:local-clock}
\end{equation}
Define $K_1(\gamma)=\varkappa_\gamma K_0(\gamma)$, enlarging fixed constants
once. For any nonnegative deterministic cell weights $a_c$,
\begin{equation}
 \sum_c a_c\mathcal E_{S_c}(f)
 \le K_0(\gamma)\sum_\alpha
          \left(\sum_{c:\alpha\in\widetilde S_c}a_c\right)
                          \|\widehat Q_\alpha f\|_2^2
 \le K_1(\gamma)\|a\|_\infty\widehat{\mathcal E}(f).
 \label{f16v27:weighted-clock}
\end{equation}
In particular $K_1(\gamma)\le C_rh^8e^{C_r\gamma h^3}$. This large
coefficient is retained; the original update rates are not accelerated.
\end{proposition}
\begin{proof}
The exact frame relation gives
$\sigma(\omega_{\widetilde S_c^c})\subseteq\sigma(\zeta_{S_c^c})$.
Projection monotonicity first bounds $\mathcal E_{S_c}$ by the conditional
variance on $\widetilde S_c$. Apply V25's continuous-product interface theorem
to that balanced conditional, for every value of its exterior. Its original
product reference retains the fractional endpoint Haar powers and the one-cube
endpoint factors.

Order free elementary groups by time. The spatial cross-section of this local
box has at most $C h^3$ cells. Only a fixed number of active time layers, including
a partially processed layer, can meet an ordering cut. Bounded native word
incidence and oscillation give cut cost $C_r\gamma h^3$ and incident cost
$C_r\gamma$. The path prefactor is $|\widetilde S_c|\le C h^4$.
A spatial cycle fully contained in a short box has no more than $h$ distinct
coordinates in that direction; counting a complete layer includes its periodic
seam. Missing time groups and preceding-port enlargement respect the same
bounds. V25 therefore gives the second inequality of
\eqref{f16v27:local-clock} after integrating the actual balanced exterior.

For a fixed chord group, centers whose $S_c$ contain it are in a coordinatewise
box of size at most $C h^4$. For a fixed moving port, either of its two adjacent
chord times may cause membership; the union has the same bound with a fixed
factor. Periodic saturation reduces the number of distinct centers. This proves
the overlap bound. Sum the nonnegative form estimates with weights $a_c$ to
obtain \eqref{f16v27:weighted-clock}. No factor counting all repair centers
appears in the energy coefficient.
\end{proof}

\subsection{An all-test-function rough-collar inequality and a low-energy projector}

Set $p=p_\gamma(Q)$. For a finite cell set $C$, including any actual collar,
write $H_C(Q)=\bigcup_{c\in C}H_c(Q)$. No separation or diameter restriction is
imposed on $C$. Let
\[
              \Pi_{\le E}=1_{[0,E]}(\widehat{\mathcal L}),\qquad E\ge0.
\]
This is a spectral projection of the original auxiliary sampler, and includes
its constant mode. It is not a physical energy projection.

\begin{theorem}[Rough-collar spectral weights with a paid energy term]
\label{f16v27:weighted-rough}
For every $f\in L^2(P)$ and every nonnegative deterministic weight array $a$,
\begin{equation}
 \boxed{
 \sum_c a_c\|1_{H_c(Q)}f\|_2^2
 \le\left[
       \sqrt{K_1(\gamma)\|a\|_\infty\widehat{\mathcal E}(f)}
       +\sqrt{p\sum_ca_c}\,\|f\|_2
       \right]^2.}
 \label{f16v27:weighted-rough-bound}
\end{equation}
Consequently, at unrestricted spatial volume and duration,
\begin{equation}
 \boxed{
 \|1_{H_C(Q)}\Pi_{\le E}\|_{2\to2}
 \le\min\{1,\sqrt{K_1(\gamma)E}+\sqrt{|C|p_\gamma(Q)}\}.}
 \label{f16v27:spectral-collar}
\end{equation}
The energy coefficient does not depend on $|C|$ or its diameter. Its explicit
coupling dependence is part of the assertion. No native gap is assumed in
defining or estimating $\Pi_{\le E}$.
\end{theorem}
\begin{proof}
For each $c$, Proposition~\ref{f16v27:uniform-tail} gives
$P(H_c\mid\zeta_{S_c^c})\le p$. Decompose
$1_{H_c}f=1_{H_c}Q_{S_c}f+1_{H_c}E_{S_c}f$.
Minkowski in the finite direct sum with weights $a_c$, followed by
Proposition~\ref{f16v27:projection}, gives
\[
 \left(\sum_ca_c\|1_{H_c}f\|_2^2\right)^{1/2}
 \le\left(\sum_ca_c\mathcal E_{S_c}(f)\right)^{1/2}
       +\left(p\sum_ca_c\|E_{S_c}f\|_2^2\right)^{1/2}.
\]
Use \eqref{f16v27:weighted-clock} and contraction of each conditional
expectation. This proves the first assertion for arbitrary complex tests.
For $a=1_C$, $1_{H_C}\le\sum_{c\in C}1_{H_c}$, so the same bound controls
$\|1_{H_C}f\|_2$. On the range of $\Pi_{\le E}$, spectral calculus gives
$\widehat{\mathcal E}(f)\le E\|f\|_2^2$. Taking the operator norm proves
\eqref{f16v27:spectral-collar}. All inequalities precede any spectral
restriction; rare-event weights are not excluded from the test class.
\end{proof}

For fixed $0<\alpha<1/2$ and $|C|\le A_0\gamma^s$, with $s,A_0$ fixed,
\begin{equation}
 \|1_{H_C(\gamma^\alpha)}\Pi_{\le E}\|
       \le \sqrt{K_1(\gamma)E}+C_{A_0,s,r}e^{-c_r\gamma^{2\alpha}}
 \label{f16v27:polynomial-collar}
\end{equation}
eventually. In particular an arbitrary normalized spectral vector with
$K_1(\gamma)E\to0$ cannot put order-one mass on that rough collar. This is
an operator statement on the complete low-energy band, not a statement about
a chosen finite list of moments. The small-energy scale here can be much
smaller than an algebraic power of $\gamma^{-1}$.

\begin{proposition}[A localized killed floor, not a global gap]
\label{f16v27:killed}
For $|C|p<1$ and every $f$ supported on $H_C(Q)$,
\begin{equation}
 \widehat{\mathcal E}(f)\ge
      \frac{(1-\sqrt{|C|p})^2}{K_1(\gamma)}\|f\|_2^2.
 \label{f16v27:collar-killed}
\end{equation}
This is a lower bound for the compressed operator on $L^2(H_C,P)$.
For the polynomial collars above, it is eventually at least
$1/(4K_1(\gamma))$.
\end{proposition}
\begin{proof}
Set $f=1_{H_C}f$ in \eqref{f16v27:weighted-rough-bound} with $a=1_C$,
take square roots, and rearrange. At a finite regulator the original
generator is bounded and positive, so its compression to the indicated closed
subspace has exactly the displayed restricted quadratic form. The final
assertion uses $|C|p\le1/4$. This compression describes continuous killing
on leaving the event, not observations of the unmodified process at separated
times with intervening returns allowed.
\end{proof}

The explicit original-clock lower bound supplied by this comparison has scale
$C_r^{-1}h_\gamma^{-8}e^{-C_r\gamma h_\gamma^3}$ in the present estimate.
It is positive independently of ambient size at fixed coupling, but does not
stay positive as coupling grows. An unrestricted growing union may have
$|C|p\ge1$; the theorem does not then assert a killed floor.

\subsection{Separated repair regions: a joint rough-count estimate without a count debit}

Choose a finite set $\mathcal J$ of centers such that no native interaction
meets two of their $S_c$ and their enlarged balanced supports $\widetilde S_c$
are disjoint. A sufficient condition is a fixed multiple of $h_\gamma$
separation in an actual finite-geometry packing. Set $U=\bigcup_{c\in\mathcal J}S_c$
and $M=|\mathcal J|$. This is one common exterior. Different centers need not
have the same conditional distribution.

For an increasing set $\mathcal A\subseteq\{0,1\}^{\mathcal J}$ define
\[
 H_{\mathcal A}=\{(1_{H_c(Q)})_{c\in\mathcal J}\in\mathcal A\},\qquad
 b_{\mathcal A}=\operatorname{Ber}(p)^{\otimes M}(\mathcal A).
\]
Here increasing means coordinatewise increasing. The Bernoulli product is a
comparison of indicators, not a substitute for the native field law.

\begin{theorem}[Simultaneous rough-event localization in the original clock]
\label{f16v27:joint}
For all $f\in L^2(P)$,
\begin{equation}
 \boxed{
 \|1_{H_{\mathcal A}}f\|_2
 \le\sqrt{K_0(\gamma)\widehat{\mathcal E}(f)}
                         +\sqrt{b_{\mathcal A}}\,\|f\|_2.}
 \label{f16v27:joint-weight}
\end{equation}
The coefficient $K_0(\gamma)$ is independent of $M$. Moreover,
\begin{equation}
 \|1_{H_{\mathcal A}}\Pi_{\le E}\|
       \le\min\{1,\sqrt{K_0(\gamma)E}+\sqrt{b_{\mathcal A}}\},\qquad
 \widehat{\mathcal E}(f)\ge
       \frac{1-b_{\mathcal A}}{K_0(\gamma)}\|f\|_2^2
       \quad(f=1_{H_{\mathcal A}}f).
 \label{f16v27:joint-spectral}
\end{equation}
In particular all $M$ selected cells being rough costs $b_{\mathcal A}=p^M$.
For $0<p<\theta<1$, at least $\theta M$ of them being rough obeys
\begin{equation}
 b_{\mathcal A}\le e^{-M D_{\rm B}(\theta\Vert p)},\qquad
 D_{\rm B}(\theta\Vert p)=
   \theta\log\frac\theta p+(1-\theta)\log\frac{1-\theta}{1-p}.
 \label{f16v27:chernoff}
\end{equation}
No independence of the actual exterior posterior is asserted.
\end{theorem}
\begin{proof}
At fixed $\zeta_{U^c}$, the exact product reference and the absence of a word
meeting two free components factor the conditional into the original
$P_{S_c}^{q_c}$. Its normalizer is the product of their normalizers only
\emph{at this fixed exterior}. The indicators are consequently conditionally
independent Bernoulli variables with success parameters at most $p$ by
Proposition~\ref{f16v27:uniform-tail}. Coupling each with a Bernoulli($p$)
using independent uniforms proves
$P(H_{\mathcal A}\mid\zeta_{U^c})\le b_{\mathcal A}$.

Product conditional variance tensorization at the same exterior gives
\[
 \|Q_Uf\|_2^2\le\sum_{c\in\mathcal J}\mathcal E_{S_c}(f)
       \le K_0(\gamma)\widehat{\mathcal E}(f).
\]
The last inequality uses disjoint enlarged balanced supports; there is no
factor $M$ and no triangle inequality summing $M$ norms. Apply
Proposition~\ref{f16v27:projection} with $E=E_U$. Its supported-function
version gives the stronger second inequality of
\eqref{f16v27:joint-spectral}. Spectral calculus gives the first. For the
count bound, conditional multiplication yields
$\mathbb E e^{t\sum1_{H_c}}\le(1-p+pe^t)^M$ under each exterior-tilted
comparison. Exponential Markov and minimizing over $t>0$ gives
\eqref{f16v27:chernoff}. Only this auxiliary product bound is used; averaging
over a correlated exterior does not make the actual fields independent.
\end{proof}

\subsection{What a low-energy change of density can and cannot do}

The last theorem has an exact change-of-law interpretation. Let $\|f\|_2=1$,
put $g=E_Uf$, and let $e=\|Q_Uf\|_2^2$. If $e<1$, define
\[
         d\nu_f=|f|^2dP,\qquad d\widetilde\nu_f=\frac{|g|^2}{1-e}\,dP.
\]
The second density depends only on the one common exterior.

\begin{proposition}[Exterior reweighting approximation for spectral tests]
\label{f16v27:tilt}
The measures just defined obey
\begin{equation}
        d_{\rm TV}(\nu_f,\widetilde\nu_f)\le\sqrt e
                 \le\sqrt{K_0(\gamma)\widehat{\mathcal E}(f)}.
 \label{f16v27:tilt-TV}
\end{equation}
Conditionals inside $U$ under $\widetilde\nu_f$ are exactly the original
native conditionals. Thus their rough indicators retain the comparison in
Theorem~\ref{f16v27:joint}. The corresponding claim for $\nu_f$ pays the
explicit error above; it does not hold by invariance alone.
\end{proposition}
\begin{proof}
Let $v=g/\sqrt{1-e}$. Orthogonality gives
$\|v\|_2=1$ and $\langle f,v\rangle=\sqrt{1-e}$, real and nonnegative.
Cauchy--Schwarz applied to
$||f|^2-|v|^2|\le|f-v||f+v|$ yields
\[
 \tfrac12\int\bigl||f|^2-|v|^2\bigr|\,dP
 \le\tfrac12\|f-v\|_2\|f+v\|_2=\sqrt e.
\]
Exterior measurability of $g$ leaves the conditional laws inside $U$ unchanged
on their nonzero exterior weights. The energy bound was proved above.
\end{proof}

A local moment consequence is also available. On the complete charts
$e_c\le E_*\gamma$, $E_*=12\pi^2$. For any normalized $f$ with
$\widehat{\mathcal E}(f)\le E$, decomposition with its single repair projection
gives
\begin{equation}
 \left(\mathbb E_P[e_c^k|f|^2]\right)^{1/2}
 \le \sqrt{M_{k,r}}+(E_*\gamma)^{k/2}\sqrt{K_0(\gamma)E},\qquad k\ge1.
 \label{f16v27:spectral-moments}
\end{equation}
Indeed the exterior-measurable part uses the uniform conditional moment in
\eqref{f16v27:uniform-preparation}, whereas the orthogonal part uses the
compact cap and its paid energy. These are moments under the test-weighted
law, not V16's unweighted moments substituted for them.

\subsection{The rough innovation form is bounded, but not yet absorbed}

To display the remaining extensive loss, let $D$ run over any deterministic
finite family of unmoving update blocks with rates $r_D\ge0$. Put
$C_D=\partial_{\mathcal G}D$ and
\[
 \mathfrak D_Q(f)=\sum_Dr_D\|1_{H_{C_D}(Q)}Q_Df\|_2^2,
       \qquad a_c=\sum_{D:\,c\in C_D}r_D.
\]
This is a specified rough-collar innovation form. It is not asserted equal
to every error in a future path-coupling or noncommuting projection argument.

\begin{proposition}[A same-law all-test-function bound on the rough innovations]
\label{f16v27:rough-form}
For every $f\in L^2(P)$,
\begin{equation}
 \boxed{
 \mathfrak D_Q(f)\le
 \left[
 \sqrt{K_1(\gamma)\|a\|_\infty\widehat{\mathcal E}(f)}
     +\sqrt{p\sum_ca_c\,\operatorname{Var}_P f}
 \right]^2.}
 \label{f16v27:rough-form-bound}
\end{equation}
For V26's unit-coverage product boxes of side $R$, its actual collar count
therefore gives
\begin{equation}
 \mathfrak D_Q(f)\le\frac{C_\rho}{R}
     \left[\sqrt{K_1(\gamma)\widehat{\mathcal E}(f)}
          +\sqrt{p\mathcal N\operatorname{Var}_P f}\right]^2,
          \qquad\mathcal N=B(n+1).
 \label{f16v27:extensive-loss}
\end{equation}
No bad-event probability is multiplied independently of an innovation.
Nevertheless this bound does not supply an unrestricted-volume absorption.
\end{proposition}
\begin{proof}
The event $H_{C_D}$ is exterior to $D$. Its multiplication operator commutes
with $E_D$ and $Q_D$. Writing $f_0=f-Pf$, contraction gives
\[
 \|1_{H_{C_D}}Q_Df\|_2
   =\|Q_D(1_{H_{C_D}}f_0)\|_2\le\|1_{H_{C_D}}f_0\|_2.
\]
Use $1_{H_{C_D}}\le\sum_{c\in C_D}1_{H_c}$, sum the deterministic rates,
and apply \eqref{f16v27:weighted-rough-bound} to $f_0$. The energy is unchanged
and $\|f_0\|_2^2=\operatorname{Var}_Pf$. This proves
\eqref{f16v27:rough-form-bound}, including vanishing on constants. V26's
finite-cover proof bounds $a_c\le C_\rho/R$, hence
$\sum_ca_c\le C_\rho\mathcal N/R$, with periodic seams and clipped ends
already counted. Substitution proves \eqref{f16v27:extensive-loss}.
\end{proof}

The remaining two costs are now explicit. First, $K_1(\gamma)$ is a pessimistic
interface comparison, not a coupling-uniform constant. Second, at fixed
coupling the term $p\mathcal N$ cannot be made small uniformly in the ambient
history. The separated-region theorem avoids a count factor for one specified
joint event; it does not turn every large union into a rare event. These facts
prevent application of V24's all-row contraction theorem to V26's good rows
alone. The proof supplies a real weighted estimate where there was previously
only a probability estimate, but not the missing global absorption.

\subsection{Scope, next target, and verification}

V27 excludes concentration on declared rough collars from sufficiently low
original-sampler spectral bands, with an explicit local energy price. It also
controls arbitrary increasing rough events on separated repair regions without
charging their number to that price. The repair exteriors are completely
unrestricted, and no field-dependent update schedule is introduced. All
conditional probabilities, product references, compact outputs, and original
source definitions remain unchanged. The supplied import notes' distinction
between a probability/power gain and a locality or functional-inequality cost
is implemented in the two different terms of
\eqref{f16v27:rough-form-bound}.

The next useful step must reduce or absorb those terms in the actual global
operator, or extend spatial boundary propagation through the remaining rough
rows. Another unweighted rough-collar tail would not do that. The killed
operator $1_H\widehat{\mathcal L}1_H$ excludes paths as soon as they leave $H$;
$1_He^{-t\widehat{\mathcal L}}1_H$ permits returns and is a different object.
No improvement of the global gap, V23's covariance norms, or the physical
transfer follows simply from the killed floor. In particular a low-energy
projection here is not a Yang--Mills physical Hilbert-space projection.

The accompanying diagnostic is
\begin{center}\small\path{verify_ym_f16_v27_rough_spectral.py}.\end{center}
It checks finite conditional projections and weighted inequalities, original
clock and overlap counts, conditional product domination with correlated
exteriors, spectral-band and density-change formulas, and a distinction between
killed and returning dynamics. Its counts record exact finite fixtures, not
native integral enclosures or formal certification of analytic constants.
The actual outputs, predecessor diagnostic reruns, hashes, reconstruction,
and build/render records are packaged separately.

The cumulative V26 source is preserved except for its one title version and
this terminal insertion. Removing the insertion and reversing that title change
recovers V26 byte for byte. The analytic deductions rely on the identified
preparation and bounded-word comparison inputs; the full archive has not been
independently recertified. The next scheduled companion checkpoint is V30.

\begin{firewall}
V27 proves an all-test-function rough-collar estimate, localized original-clock
killed bounds, and low-energy spectral-projection estimates. It proves a
simultaneous rough-count bound on separated repair regions and an exact
exterior-reweighting approximation with its energy error. The original-clock
coefficient is explicit and can grow rapidly with coupling; the global rough
innovation bound retains an extensive remainder. No all-boundary spatial
susceptibility, unrestricted global gap, new unrestricted covariance theorem,
physical Gauss coverage, bridge-window removal, infrared contact coercivity,
same-top physical transfer comparison, physical-time calibration, or interacting
continuum mass gap is inferred.
\end{firewall}
% END F16 V27 APPEND
% BEGIN F16 V28 APPEND -- 2026-09-18
% Insert immediately before the final document-ending command in V27.
% No predecessor proof, native probability, source, or elementary clock is edited.

\clearpage
\section{V28: nonextensive rough-innovation absorption by local repair dynamics}
\label{f16v28:context}

V27's final rough-innovation estimate contains $p\mathcal N\operatorname{Var}f$.
That loss occurs after replacing each block innovation by the whole centered
function. We retain the innovation instead. Before applying conditional
smallness, subtract the expectation outside the \emph{union} of the updated
block and its repair region. That component is annihilated by the block
innovation. The probability term then multiplies a local excess conditional
energy, not the full variance. No commutation of overlapping conditional
expectations is needed.

The resulting estimate has no ambient-volume remainder. With two explicitly
normalized auxiliary update families it gives a vanishing relative form bound
for the \emph{same} rough-collar innovation form defined in V27. In particular,
rough-collar updates may be suppressed in that augmented dynamics at a small
relative form cost, while the repair and enlarged-block updates remain active.
The probability law is unchanged, and all output fields remain integrated.
This is not a proof of a gap for either augmented dynamics. Returning their
energies to the original elementary clock still has a large local cost, which
is stated separately.

\subsection{Dependencies and the upstream cancellation}

The supplied V27 appendix and audit were read completely. The new main
inequality uses its conditional preparation theorem
\ref{f16v27:uniform-tail}, and the exact finite-cover geometry of V26. It does
not use V26's restricted-collar coupling to assume that rough rows contract.
V25's interface comparison is used only when returning the new auxiliary
forms to the original balanced clock. V17--V23's covariance expansions and a
native sampler gap are not inputs to the absorption argument.

Two scoped semantic searches returned no matches. Mounted-text searches and
identified passages were then checked. The Grand Tensor monograph already
uses nested conditional-expectation Pythagoras in its score-cutoff loss formula
(NL.12)--(NL.13), and its nested-source identity (NS.6). That general principle
is credited, not claimed new. Its use here subtracts a particular common
exterior inside an \emph{event-weighted block innovation}. V27's estimate
\eqref{f16v27:extensive-loss} remains valid; the present estimate is a different,
stronger interface, not a silent correction of the predecessor.

Reversible Dirichlet-form comparison and monotone-system censoring are distinct
methods. Peres--Winkler,
\href{https://arxiv.org/abs/1112.0603}{\emph{Can extra updates delay mixing?}},
addresses monotone spin systems and specified initial laws. Its primary abstract
and scope were checked; no theorem from its full proof is used here. We prove
reversibility and the form inequalities for the present exterior-gated updates
directly. No order on compact group fields, stochastic domination of trajectories,
or mixing-time censoring inequality is assumed. The review is targeted, not
independent certification of the analytic archive. The next checkpoint remains
V30. The supplied import memoranda's distinction between probability gain,
locality cost, and a subsequent global spectral theorem is retained.

\subsection{A union projection removes the irrelevant exterior component}

On the unchanged finite native probability space, for a set $A$ of unmoving
cells write
\[
 E_Af=\mathbb E_P[f\mid\zeta_{A^c}],\qquad Q_A=I-E_A,
 \qquad \mathcal E_A(f)=\|Q_Af\|_2^2.
\]
All conditional expectations are orthogonal projections on the same $L^2(P)$.
The next proposition holds on any probability space with these sigma-algebra
inclusions; it does not require product independence or a positive-density
comparison.

\begin{proposition}[Nested localization of an actual innovation]
\label{f16v28:nested}
Let $D,S$ be cell sets, let $U=D\cup S$, and let $H$ be measurable outside
$D$. Suppose $P(H\mid\zeta_{S^c})\le p$ almost surely, with $0\le p\le1$.
Then for every real or complex $f\in L^2(P)$,
\begin{equation}
 \boxed{
 \|1_HQ_Df\|_2
 \le \sqrt{\mathcal E_S(f)}+
       \sqrt{p\{\mathcal E_U(f)-\mathcal E_S(f)\}}.}
 \label{f16v28:nested-bound}
\end{equation}
The quantity in braces is nonnegative and equals
$\|(E_S-E_U)f\|_2^2$. In particular, for every $t>0$ the following is a
quadratic-form inequality:
\begin{equation}
 1_HQ_D\preceq (1+t)Q_S+(1+t^{-1})p(E_S-E_U).
 \label{f16v28:nested-Young}
\end{equation}
No identity $E_SE_D=E_DE_S$ is required.
\end{proposition}
\begin{proof}
Since $U$ contains $S$ and $D$, $E_SE_U=E_U=E_DE_U$ and their reversed
products have the same values. Multiplication $M=1_H$ commutes with $E_D$,
because $H$ is exterior to $D$. Decompose
\[
 f=Q_Sf+(E_S-E_U)f+E_Uf.
\]
The last component disappears from $MQ_Df$. Using the commutation only with
$M$ gives the exact two-term identity
\begin{equation}
 MQ_Df=Q_DM Q_Sf+Q_DM(E_S-E_U)f.
 \label{f16v28:two-term}
\end{equation}
The first term has norm at most $\|Q_Sf\|_2$. The second intermediate vector
$(E_S-E_U)f$ is exterior-measurable for $S$. Conditional smallness therefore
bounds its image under $M$ by $\sqrt p\|(E_S-E_U)f\|_2$. The contraction
$Q_D$ does not increase that norm. This proves \eqref{f16v28:nested-bound}.
Nested orthogonal projections give
\[
 \|(E_S-E_U)f\|_2^2
 =\|E_Sf\|_2^2-\|E_Uf\|_2^2
 =\mathcal E_U(f)-\mathcal E_S(f).
\]
Finally $MQ_D$ is itself a positive orthogonal projection, since its two
factors commute. Its quadratic form is $\|MQ_Df\|_2^2$. Apply
$(a+b)^2\le(1+t)a^2+(1+t^{-1})b^2$ to obtain
\eqref{f16v28:nested-Young}. These arguments are bounded-operator identities
on the whole $L^2$ space, including constants and rare-event-weighted tests.
\end{proof}

The cancellation of $E_Uf$ is essential. Its norm can be large, but it has no
innovation in $D$. Charging that norm independently for every collar witness
would recreate V27's extensive term. The subtraction changes neither the
observable $Q_Df$ nor the conditional probability used to estimate $H$.

\subsection{The same rough form now has a local-energy remainder}

Use V27's actual repair regions $S_c$, rough events $H_c(Q)$, and probability
\[
 p=p_\gamma(Q)=\min\{1,e^{t_0(\overline M_r-Q^2)}\},\qquad
 t_0=(24\pi^2)^{-1}.
\]
The preparation theorem gives $P(H_c(Q)\mid\zeta_{S_c^c})\le p$ with all
compact exterior fields retained. Let $D$ run over a deterministic indexed
family, with rates $r_D\ge0$ and actual collars $C_D=\partial_{\mathcal G}D$.
Keep exactly V27's form
\[
 \mathfrak D_Q(f)=\sum_Dr_D\|1_{H_{C_D}(Q)}Q_Df\|_2^2,
 \qquad a_c=\sum_{D:c\in C_D}r_D.
\]
Repeated blocks keep their indices. Define two positive forms
\begin{align}
 \mathcal A(f)&=\sum_c a_c\mathcal E_{S_c}(f),\notag\\
 \mathcal B(f)&=\sum_Dr_D\sum_{c\in C_D}
       \{\mathcal E_{D\cup S_c}(f)-\mathcal E_{S_c}(f)\}.
 \label{f16v28:AB}
\end{align}
Each difference in the second line is a positive nested-projection form.

\begin{theorem}[Nonextensive localization of rough innovations]
\label{f16v28:local-form}
At every admitted finite regulator and for all $f\in L^2(P)$,
\begin{equation}
 \boxed{\quad
 \mathfrak D_Q(f)^{1/2}\le\mathcal A(f)^{1/2}
                                      +p^{1/2}\mathcal B(f)^{1/2}.
 \quad}
 \label{f16v28:local-form-bound}
\end{equation}
There is no $\operatorname{Var}_Pf$ or total-cell count on the right. If
$D^+$ is any declared set containing $D\cup S_c$ for every $c\in C_D$, then
\begin{equation}
 \mathcal B(f)\le\sum_Dr_D|C_D|\mathcal E_{D^+}(f).
 \label{f16v28:B-upper}
\end{equation}
The target rough form, its gates, and its update rates have not been changed.
\end{theorem}
\begin{proof}
The event $H_c$ is exterior to $D$ for $c\in C_D$. Apply
Proposition~\ref{f16v28:nested} to each pair $(D,S_c)$. The indicator of a
union is at most the sum of its member indicators, so
\[
 \mathfrak D_Q(f)\le\sum_Dr_D\sum_{c\in C_D}
                                      \|1_{H_c}Q_Df\|_2^2.
\]
Minkowski in this finite weighted direct sum gives
\eqref{f16v28:local-form-bound}; the two sums of squared terms are exactly
\eqref{f16v28:AB}. If $D\cup S_c\subseteq D^+$, monotonicity of averaged
conditional variance gives $\mathcal E_{D\cup S_c}\le\mathcal E_{D^+}$.
Dropping the nonnegative subtracted energy proves \eqref{f16v28:B-upper}.
There is no independence assertion and no interchange of a growing-volume
limit with a conditional expectation.
\end{proof}

In particular, the probability multiplies an energy of a \emph{local union},
not the variance of the entire history. Converting that energy to the original
clock is still an additional estimate. The next construction keeps its local
update rate explicit before paying that comparison.

\subsection{Two repair families with bounded coverage}

Retain V27's box parameters $w=w_\gamma$, $h=h_\gamma=2w+5$, and choose an
integer $\rho\ge1$ dominating both the native interaction range and the frame
support allowance in ambient coordinatewise distance. Put $v=w+\rho+2$.
Choose an integer $R\ge2(v+\rho+1)$.

Use V26's indexed product cover $\mathcal B_R$. In a spatial cycle with
$m>R$ take every cyclic interval of length $R$; when $m\le R$ take the whole
cycle once. In time take all nonempty clipped intervals
$[s,s+R-1]\cap[0,n]$, $1-R\le s\le n$, without deduplication. Product rates
are $r_D=R^{-(k+1)}$, where $k$ is the number of cut spatial directions.
Let $D^+$ expand each coordinate interval of $D$ by $v$, saturating a short
cycle and clipping at natural time endpoints. Set
\[
 \kappa_{\rm rep}=h^4,\qquad \kappa_+=(1+2v/R)^4,
\]
\begin{equation}
 \begin{split}
 \mathcal L_R&=\sum_{D\in\mathcal B_R}r_DQ_D,\\
 \mathcal L_{\rm rep}&=\kappa_{\rm rep}^{-1}\sum_cQ_{S_c},\qquad
 \mathcal L_+=\kappa_+^{-1}\sum_{D\in\mathcal B_R}r_DQ_{D^+}.
 \end{split}
 \label{f16v28:generators}
\end{equation}
Their forms have the matching symbols $\mathcal E_R,\mathcal E_{\rm rep},
\mathcal E_+$. These are exact complete conditional updates in the same
unmoving coordinates. They are not elementary balanced updates.

\begin{proposition}[Finite geometry and actual rate bounds]
\label{f16v28:geometry}
Every cell has coverage exactly one under $\mathcal L_R$ and coverage at most
one under each of $\mathcal L_{\rm rep}$ and $\mathcal L_+$. Moreover,
\begin{equation}
 \begin{split}
 D\cup S_c&\subseteq D^+\quad(c\in C_D),\qquad |D^+|\le(R+2v)^4,\\
 a_*:=\sup_c a_c&\le \frac{8\rho}{R}(1+2\rho/R)^3,\\
 b_*:=\max_D|C_D|&\le8\rho(R+2\rho)^3.
 \end{split}
 \label{f16v28:geometry-bounds}
\end{equation}
All bounds include periodic seams, short cycles, clipped temporal ends and
repeated indices.
\end{proposition}
\begin{proof}
V26's product count gives unit coverage for $D$, and the stated upper bound
on $a_*$. A collar cell lies at coordinatewise distance at most $\rho$ from
$D$. Its repair box has coordinatewise radius $w$, so the declared enlargement
contains $S_c$, including the temporal and periodic saturation cases.

For a fixed cell, at most $(2w+1)^4\le h^4$ repair centers have a box containing
it. In a cut spatial direction, at most $R+2v$ starting indices have an enlarged
interval containing that cell. If the enlargement saturates the cycle, the
number is $m\le R+2v$. In an uncut direction there is only one interval. The
same $R+2v$ bound holds for indexed clipped time intervals by restricting the
possible start indices. Product rates therefore give enlarged coverage at most
$(1+2v/R)^{k+1}\le\kappa_+$. The displayed normalizations prove the rate claims.

The actual collar is contained in the coordinatewise $\rho$-enlargement minus
$D$. Assign a point of that difference to one coordinate in which it lies
outside the corresponding interval. There are at most $2\rho$ positions in
that coordinate and at most $R+2\rho$ in each other coordinate, for at most
four choices. Saturated directions and missing physical time faces only remove
positions. This proves the bound on $b_*$. No invariance of the probability
under spatial translations was used in these counts.
\end{proof}

\begin{theorem}[Rough-update absorption in explicitly normalized local forms]
\label{f16v28:absorption}
Define
\begin{equation}
 u_{\gamma,R}=\kappa_{\rm rep}a_*,\qquad
 z_{\gamma,R,Q}=p_\gamma(Q)b_*\kappa_+,\qquad
 \theta_{\gamma,R,Q}=u_{\gamma,R}+z_{\gamma,R,Q}.
 \label{f16v28:theta}
\end{equation}
Then, on the entire native $L^2(P)$ space,
\begin{equation}
 \boxed{
 \begin{aligned}
 \mathfrak D_Q(f)^{1/2}
 &\le\sqrt{u_{\gamma,R}\mathcal E_{\rm rep}(f)}
                  +\sqrt{z_{\gamma,R,Q}\mathcal E_+(f)},\\
 \mathfrak D_Q(f)
 &\le\theta_{\gamma,R,Q}
                  \{\mathcal E_{\rm rep}(f)+\mathcal E_+(f)\}.
 \end{aligned}}
 \label{f16v28:absorption-bound}
\end{equation}
In particular,
\begin{equation}
 \theta_{\gamma,R,Q}\le
 C_\rho\left(\frac{h_\gamma^4}{R}
           +p_\gamma(Q)R^3(1+2v/R)^4\right).
 \label{f16v28:theta-scale}
\end{equation}
There is no $B$, $n$, $\mathcal N$, or number of update indices in this
coefficient. All compact field values, including rough collars, remain covered.
\end{theorem}
\begin{proof}
By definition of $a_*$,
$\mathcal A(f)\le a_*\sum_c\mathcal E_{S_c}(f)
=u_{\gamma,R}\mathcal E_{\rm rep}(f)$.
Equation \eqref{f16v28:B-upper} gives
$p\mathcal B(f)\le p b_*\sum_Dr_D\mathcal E_{D^+}(f)
=z_{\gamma,R,Q}\mathcal E_+(f)$.
Insert these two bounds in Theorem~\ref{f16v28:local-form}. Cauchy--Schwarz
in the two-dimensional vector of square roots proves the second line.
Use \eqref{f16v28:geometry-bounds} and $R\ge2\rho$ for
\eqref{f16v28:theta-scale}. The price for conditioning an event on its repair
exterior has been paid as an actual energy before this comparison; an
unweighted event probability has never been multiplied by a separate test norm.
\end{proof}

\subsection{A polylogarithmic choice compatible with the controlled-row coupling}

\begin{proposition}[One small coefficient with two finite local update scales]
\label{f16v28:scales}
At sufficiently large fixed-$r$ coupling, take
\begin{equation}
 R_\gamma=\lceil h_\gamma^5\rceil,\qquad
 Q_\gamma^2=\overline M_r+\frac{20}{t_0}\log h_\gamma.
 \label{f16v28:scale-choice}
\end{equation}
Then $p_\gamma(Q_\gamma)=h_\gamma^{-20}$, all the required geometry holds, and
\begin{equation}
 \theta_\gamma\le C_r(h_\gamma^{-1}+h_\gamma^{-5})
             =O_r((\log\gamma)^{-1})\longrightarrow0.
 \label{f16v28:theta-small}
\end{equation}
The repair side is $O_r(\log\gamma)$; the enlarged blocks have side
$O_r((\log\gamma)^5)$ and at most $O_r((\log\gamma)^{20})$ cells.
In the separate V26 controlled-collar theorem, these same $R_\gamma,Q_\gamma$
give a positive admitted-row margin eventually. That statement is still only
about its admitted pairs of deterministic collars.
\end{proposition}
\begin{proof}
The probability identity follows directly from its definition. Here $h_\gamma$
is bounded above and below by positive fixed-$r$ multiples of $\log\gamma$
eventually; $v=O(h_\gamma)$. Thus $R_\gamma\ge2(v+\rho+1)$ eventually and
$\kappa_+$ is bounded. In \eqref{f16v28:theta-scale} the two terms are at most
$C h_\gamma^{-1}$ and $C h_\gamma^{15-20}$, respectively.
All update sizes follow from the actual side lengths, not a volume limit.

For the last statement $Q_\gamma=O_r(\sqrt{\log h_\gamma})$, so
$Q_\gamma\le\eta_*\sqrt\gamma$ eventually. V26's full-output estimate on
admitted rows has budget
\[
 2C_rQ_\gamma+2R_\gamma^8
                  e^{A_r(1+Q_\gamma^2)-c_*\gamma}.
\]
Multiply it by the actual affected-block rate $C_\rho/R_\gamma$.
The first part tends to zero. The second is a fixed power of $h_\gamma$
times $e^{-c_*\gamma}$ and also tends to zero. Thus the admitted-row margin
is eventually at least $1/2$. This calculation does not extend the coupling
to excluded collars or identify its error with an arbitrary spectral defect.
\end{proof}

Alternatively, retain V26's $Q_\gamma=\gamma^\alpha$ and
$R_\gamma=\lceil K_r(1+\gamma^\alpha)\rceil$, $0<\alpha<1/2$ fixed.
Equation \eqref{f16v28:theta-scale} then gives
$\theta_\gamma\le C_r\gamma^{-\alpha}(1+\log\gamma)^4+
C_r\gamma^{3\alpha}e^{-c_r\gamma^{2\alpha}}$ eventually. The polylogarithmic
choice above uses smaller blocks at the expense of a slower vanishing form
error. Neither choice changes the native measure.

\subsection{Exterior-gated updates are reversible, with a relative form bound}

Let $M_D=1_{H_{C_D}(Q)}$ and define
\begin{equation}
 \begin{split}
 \mathcal L_{\rm bad}&=\sum_Dr_DM_DQ_D,\qquad
 \mathcal L_{\rm good}=\sum_Dr_D(1-M_D)Q_D,\\
 \mathcal L_{\rm aug}&=\mathcal L_R+\mathcal L_{\rm rep}+\mathcal L_+,\\
 \mathcal L_{\rm keep}&=\mathcal L_{\rm good}+\mathcal L_{\rm rep}+\mathcal L_+
                     =\mathcal L_{\rm aug}-\mathcal L_{\rm bad}.
 \end{split}
 \label{f16v28:gated-generators}
\end{equation}
The label ``keep'' means that repair and enlarged-block moves are always
retained. It does not mean that only good configurations occur in the process.

\begin{theorem}[Small relative deletion of rough block updates]
\label{f16v28:relative}
These are bounded nonnegative self-adjoint generators at each finite regulator,
with invariant probability $P$. The gate changes an update rate, not its
conditional output law. Each cell has total coverage rate at most three under
$\mathcal L_{\rm aug}$, and no larger under $\mathcal L_{\rm keep}$.
When $\theta=\theta_{\gamma,R,Q}<1$,
\begin{equation}
 \boxed{
 \begin{aligned}
 0\preceq\mathcal L_{\rm bad}
      &\preceq\theta(\mathcal L_{\rm rep}+\mathcal L_+),\\
 (1-\theta)\mathcal L_{\rm aug}
      &\preceq\mathcal L_{\rm keep}\preceq\mathcal L_{\rm aug}.
 \end{aligned}}
 \label{f16v28:relative-order}
\end{equation}
In fact
$\mathcal L_{\rm keep}\succeq
\mathcal L_R+(1-\theta)(\mathcal L_{\rm rep}+\mathcal L_+)$.
Both full augmented generators have constants as their common null space.
These comparisons do not give a new size-uniform lower bound on their gaps.
\end{theorem}
\begin{proof}
The indicator $M_D$ is measurable outside $D$, so it commutes with $E_D,Q_D$.
Thus $M_DQ_D$ and $(1-M_D)Q_D$ are positive orthogonal projections, and
\[
 \langle f,M_DQ_Df\rangle=\|M_DQ_Df\|_2^2.
\]
Their jump forms are obtained by multiplying the ordinary reversible
conditional-update edge measure by the exterior indicator. That indicator
has the same value before and after the update. Detailed balance, preservation
of constants and the Markov property therefore hold without a field-independent
rate assumption. Equivalently their semigroups make the declared conditional
refresh at rate $r_D$ only while the exterior gate allows it. The remaining
two families are ordinary deterministic-rate conditional refreshes.

The first operator bound is Theorem~\ref{f16v28:absorption}, since the form
of $\mathcal L_{\rm bad}$ is exactly $\mathfrak D_Q$. Subtract it from
$\mathcal L_{\rm aug}$ to get the stronger lower bound and then the second
line of \eqref{f16v28:relative-order}. Coverage was proved above. Both
generators contain every repair $Q_{S_c}$ at positive rate. A function in
their common kernel is measurable outside each $S_c$, hence outside each
individual cell. The strictly positive finite-product density of the native
law makes such a function constant, up to null sets. This can also be seen
by successive single-cell conditional expectations under the equivalent
product reference. None of these statements estimates a thermodynamic gap.
\end{proof}

\begin{proposition}[Gap and shifted-resolvent comparison]
\label{f16v28:resolvent}
Let $\lambda_{\rm aug},\lambda_{\rm keep}$ be the variational gaps on the
common centered space. Then
\begin{equation}
 (1-\theta)\lambda_{\rm aug}\le\lambda_{\rm keep}\le\lambda_{\rm aug}.
 \label{f16v28:gap-relative}
\end{equation}
For every $s>0$, put $C_s=\mathcal L_{\rm aug}+sI$. On the full space,
\begin{equation}
 \begin{split}
 0\preceq C_s^{1/2}\bigl[(\mathcal L_{\rm keep}+sI)^{-1}
            -(\mathcal L_{\rm aug}+sI)^{-1}\bigr]C_s^{1/2}
 \preceq \frac{\theta}{1-\theta}I.
 \end{split}
 \label{f16v28:resolvent-bound}
\end{equation}
No operator ordering of their finite-time semigroups is inferred.
\end{proposition}
\begin{proof}
Take infima of the form comparison over centered unit vectors to get
\eqref{f16v28:gap-relative}. Also
$(1-\theta)C_s\preceq\mathcal L_{\rm keep}+sI\preceq C_s$.
Conjugate by $C_s^{-1/2}$ and invert these strictly positive operators. The
result is between $I$ and $(1-\theta)^{-1}I$; subtract $I$ to obtain
\eqref{f16v28:resolvent-bound}. No commutation of different local projections is
used. A small relative error alone does not supply a positive lower gap for
either side of the comparison.
\end{proof}

\subsection{The original clock and the remaining global test}

\begin{proposition}[The auxiliary rates have a paid original-clock comparison]
\label{f16v28:clock}
With $L=R+2v+3$ there is an ambient-size-independent bound
\begin{equation}
 \mathcal L_{\rm keep}\preceq\mathcal L_{\rm aug}
       \preceq\mathcal C_r(\gamma,R,h)\widehat{\mathcal L},\qquad
 \mathcal C_r(\gamma,R,h)=C_r L^4e^{C_r\gamma L^3}.
 \label{f16v28:clock-bound}
\end{equation}
Consequently a separately proved gap $g>0$ for either augmented generator
would imply an original-clock gap at least $g/\mathcal C_r$.
No such new uniform augmented gap is proved here.
\end{proposition}
\begin{proof}
For each unmoving update set, include its chord cubes and every moving port
whose time or following time at that cube lies in the set. This is the exact
preceding-port enlargement from V24. Its balanced exterior sigma algebra is
contained in the unmoving exterior sigma algebra, so its averaged conditional
variance dominates that of the unmoving update. Each resulting support has
at most $CL^4$ elementary groups and a time-major ordering cut of at most
$CL^3$ groups. The bounded classical word oscillations and V25's interface
comparison therefore give the local factor $C_rL^4e^{C_r\gamma L^3}$.
All product references retain the endpoint Haar exponents and endpoint factors.

Each of the three normalized unmoving families has coverage at most one.
An elementary chord belongs only through its own cell. An elementary moving
port can belong through either of two adjacent time cells, so the corresponding
aggregate balanced rate is at most two per family. Sum the three nonnegative
comparisons and absorb the factor six in $C_r$. This proves
\eqref{f16v28:clock-bound}. If one of the augmented forms has a lower bound
$g\operatorname{Var}f$, its upper comparison gives the stated implication.
The comparison factor is generally large; it has not been declared free.
\end{proof}

For the polylogarithmic scales, this bound is at worst
$C_rh_\gamma^{20}\exp(C_r\gamma h_\gamma^{15})$ after fixed constants are
enlarged. It need not improve V25's finite-width bound. The new conclusion is
\emph{relative local-form absorption without an extensive remainder}, not a
better direct source covariance obtained from this crude clock comparison.

Two finite checks identify what the new conclusion does not allow. For an
independent uniform sign $X$ and a Bernoulli $Y$ with $P(Y=1)=p$, take
$D=\{X\}$, $H=\{Y=1\}$, and $f=X1_H$. Then
\[
 \|1_HQ_Df\|_2^2=\mathcal E_D(f)=p.
\]
Thus a bound $\mathfrak D_Q\le o(1)\mathcal E_R$ cannot follow from event
rarity alone. The repair energy in \eqref{f16v28:absorption-bound} is necessary.

Also, even when an independent $X$ has identical conditional laws at all
values of $Y$, an $X$ update gated by $1_{\{Y=0\}}$ has different rates at
$Y=0$ and $Y=1$. Comparing configurations with the same $X=+1$ and those two
values of $Y$, an allowed refresh creates an $X$ mismatch with probability
$1/2$ while the other configuration does not refresh. The ungated boundary
influence was zero, but this gate-switching contribution is not. Hence V26's
admitted-row table is not by itself a path-coupling proof for
$\mathcal L_{\rm keep}$. Detailed balance proved above does not erase this
term. The ungated repair and enlarged-block families also need to be included
in any global spectral argument.

The next concrete target is therefore an independently proved coercive or
boundary-response estimate for the \emph{displayed retained generator}, with
its gate switching and its two ungated families included. Equivalently one may
seek a sharper all-boundary argument for the original blocks. The rough form
has now been absorbed relative to explicit local auxiliary energies at all
volumes, but it has not been identified with every defect in such a future
argument. No infinite-order expansion, unweighted rough probability, or
unverified commutator cancellation is used to bridge that gap.

\subsection{Diagnostics, preservation, and claim boundary}

The accompanying exact diagnostic is
\begin{center}\small\path{verify_ym_f16_v28_nested_absorption.py}.\end{center}
It checks noncommuting conditional projections, their nested union and energy
difference, event-weighted Young forms, positive rough/good generators,
finite-cover and enlargement rates, relative form and resolvent inequalities,
and the two limiting examples above. A continuous positive polynomial-density
fixture checks the normalized identities without assigning point masses to a
continuous coordinate. These are finite algebraic checks, not native compact
integral enclosures, formal verification, or numerical certification of the
inherited preparation constants.

The package records actual diagnostic outputs and any executed predecessor
reruns. The cumulative source changes only V27's title version and adds this
exact terminal section. Removing it and restoring that title recovers V27
byte for byte. No previous theorem or bibliography entry is silently edited.
The new estimates rely on the specified V27 preparation and finite conditional
identities; the optional good-row calculation retains V26's analytic hypotheses.
The full inherited archive has not been independently recertified.

\begin{firewall}
V28 removes the extensive full-variance remainder from the specified V27
rough-innovation estimate. It proves small relative form absorption into two
explicitly normalized local repair families, a reversible exterior-gated
retained dynamics, and gap/resolvent comparisons with the augmented dynamics.
Its update radii can be chosen polylogarithmically in coupling, with an
ambient-size-independent vanishing relative error. It pays a separate,
pessimistic conversion to the original elementary clock. It does not prove a
new uniform augmented or original sampler gap, all-boundary spatial
susceptibility, unrestricted coherent covariance control, or a physical mass
gap. The native measure, compact outputs, endpoint factors, original sources,
and original elementary rates are unchanged. Bridge-window removal, physical
Gauss restriction, infrared contact coercivity, same-top transfer comparison,
physical-time calibration and interacting continuum construction remain separate.
\end{firewall}
% END F16 V28 APPEND

% BEGIN F16 V29 APPEND -- 2026-09-18
% Insert immediately before the final document-ending command in V28.
% No inherited theorem, probability, source, or elementary rate is altered.

\clearpage
\section{V29: fixed-depth angular repairs and a sharper original-clock bound}
\label{f16v29:context}

V28 controls the removal of rough updates relative to two helper families, but
relative closeness supplies no spectral edge. We return to the cost of mixing
one complete native conditional. For this purpose its exceptional event can be
chosen at a \emph{fixed angular radius}, rather than at the shrinking angular
radius used to control every local boundary response. Fixed-depth buffers then
suffice. Their original-coordinate comparison costs $e^{C_r\gamma}$, not an
exponential in the volume of a growing slab.

There are two separate steps. A restricted convex native conditional has a
uniform whole-cell heat-bath variance bound. A weighted repair estimate transfers
that bound to the complete conditional, including its large fields, when the
\emph{actual} collar has a smaller fixed angular radius. We do not remove this
collar condition by probability alone. For complete temporal slabs, V28's nested
innovation identity and an independent two-overlapping-slab comparison absorb
the rough-collar contribution back into the same slab form. This yields
\[
 \widehat\lambda_P\ge c_r e^{-C_r\gamma},\qquad
 B\le e^{\beta_r\gamma},\quad n\ge1,
\]
with $\beta_r>0$ fixed and no $B$ or duration in the bound. Its cross-section
restriction is essential; its lower bound still vanishes with coupling. A
consequence controls V28's displayed retained generator in that same regime.
No unrestricted spatial gap or physical mass is inferred.

\subsection{Dependencies and the change of exceptional scale}

The supplied V28 appendix and audit were read completely. The decisive new
inputs are proved below; their inherited premises are the V16/V20 conditional
set-energy inequality on complete charts, V22's restricted measured Hessian
bounds on products of whole-cell balls, the bounded-word product reference, and
V22/V24's all-boundary temporal channels and exact slab clock. V28's nested
identity is reused, not reasserted as a new result. The higher-order source
expansions are not needed in this proof of a sampler gap.

Two semantic searches returned no results. Targeted mounted-text searches and
identified passages were then examined. F15 V70 already uses a Neumann Bochner
calculation on a convex spherical cap and a separate full-law argument. The
monograph's Bochner--Schur construction concerns a different reduced-fibre
response operator. Neither result is a theorem about the current all-test
whole-cell conditional update, and neither is substituted for it. V25's
all-finite-width baseline remains valid outside the new, stronger regime.
This is a dependency review, not independent certification of the archives.
The next scheduled companion checkpoint remains V30.

For contemporary attribution, Ascolani--Lavenant--Zanella,
\href{https://arxiv.org/abs/2410.00858v2}{\emph{Entropy contraction of the Gibbs
sampler under log-concavity}}, Corollary 3.7, proves an approximate variance
tensorization with a condition-number cost. The paper's assumptions, rate
convention, and that corollary were checked. Its Euclidean theorem is not
applied to a hard product boundary without argument. We give the weaker
product-ball estimate needed here directly, using the explicitly inherited
reflecting domain and its nonnegative curvature terms. No entropy contraction
for the complete native law is claimed. The supplied import memoranda remain
method proposals: their separation of locality cost from power gain, and of an
auxiliary spectral edge from the physical transfer, is preserved.

\subsection{Fixed-depth repair at two declared angular radii}

Retain $P=P_{\gamma,\mathbf y}$, $B=m^3$, $m\ge4$, $n\ge1$,
$\max_{i,e}|y_{i,e}|\le r$, and $\gamma\ge1+r^2$. Let
$e_c=|\zeta_c|^2\le E_*\gamma$, $E_*=12\pi^2$, and
$t_0=(24\pi^2)^{-1}$. Write
\[
 \mathcal L_{\rm cell}=\sum_c Q_{\{c\}},\qquad
 \mathcal E_{\rm cell}(f)=\sum_c\|Q_{\{c\}}f\|_2^2.
\]
A coordinate here is one complete \emph{unmoving} cell, of dimension 36 or 15.
This is not yet V14's elementary balanced chord/port clock.

Fix V22's small angular comparison radius $\eta_0:=\eta_*>0$ and its constants
$0<\lambda\le M$. We will choose a smaller fixed radius $\eta_1>0$. For any
finite free cell set $\Omega$, let $\mu=P_\Omega^q$ denote its complete native
conditional, and set
\[
 G_\Omega=\{e_c\le\eta_0^2\gamma\text{ for every }c\in\Omega\}.
\]
This event is new notation for this comparison only. It does not redefine V4's
$\mathcal O_R$, $p_{\rm good}$, or the rough threshold in V28.

\begin{proposition}[Fixed repair boxes with a controlled actual boundary]
\label{f16v29:angular-preparation}
There exist fixed $\eta_1\in(0,\eta_0)$, a fixed integer $v_0$, and positive
constants $A_r,c_0,C_r$, all independent of coupling and sizes, with the
following property. Suppose
\begin{equation}
       |q_d|\le\eta_1\sqrt\gamma\quad(d\in\partial_{\mathcal G}\Omega).
 \label{f16v29:angular-collar}
\end{equation}
For each $c\in\Omega$ let $S_c^\Omega$ be the intersection of $\Omega$ with a
coordinatewise box of radius $v_0$ about $c$, including short cycles in full
and clipping only at the natural ends. Then
\begin{equation}
 \mu(e_c>\eta_0^2\gamma\mid\zeta_{\Omega\setminus S_c^\Omega})
       \le p_0(\gamma):=\min\{1,A_r e^{-c_0\gamma}\}
 \label{f16v29:angular-tail}
\end{equation}
for every compact value of the other free variables, in the explicit
conditional-density versions. The original frozen $q$ remains fixed. Moreover
\begin{equation}
 \sum_{c\in\Omega}\mathcal E_{S_c^\Omega}^{\mu}(f)
       \le K_0(\gamma)\mathcal E_{\rm cell}^{\mu}(f),
       \qquad K_0(\gamma)\le C_r e^{C_r\gamma}.
 \label{f16v29:constant-repair-energy}
\end{equation}
The free set can have holes or disconnected components. No fast coordinate
inside it is restricted in these statements.
\end{proposition}
\begin{proof}
We spell out the boundary preparation, because a cell near the actual boundary
of $\Omega$ has no full free buffer in all directions. Denote the constants in
\eqref{f16v20:conditional-set} by $a_r,b_r$, and put
$\rho=b_r/(1+b_r)<1$. Fix $c$, condition on the exterior of $S_c^\Omega$,
and allow arbitrary signed sources of size at most $t_0$ on this free set.
Within it let $u_d$ be the actual mean energy. At frozen cells in $\Omega$
use their actual deterministic energy, bounded by $E_*\gamma$; outside
$\Omega$ extend the scalar array by $Q^2=\eta_1^2\gamma$. Only collar values
of this extension enter any native interaction, and these dominate the
actual prescribed $q$.

For an ambient set $A$ lying inside the coordinate box, apply the native set
inequality to $A\cap\Omega$. As in V26's scalar extension argument, adding the
entries outside $\Omega$ gives
\[
 \sum_{d\in A}\widetilde u_d
 \le [a_r+(1+b_r)Q^2]|A|+
                           b_r\sum_{d\in\partial A}\widetilde u_d.
\]
For graph balls of radii $k<K$ contained in that coordinate box, iteration gives
\begin{equation}
 u_c\le C_r^{(0)}+C_r^{(1)}Q^2
       +E_*C_{\rm gr}\gamma(K+1)^4\rho^K.
 \label{f16v29:stopped-mixed-boundary}
\end{equation}
Here $C_r^{(0)},C_r^{(1)}$ come from the convergent series
$\sum_{k\ge0}\rho^k(k+1)^4$ and do not depend on $K$ or the regulator.
The terminal compact cap is legitimate even when the ball crosses the
artificial repair boundary. If the finite component is already filled, one
may continue the same recurrence or take its zero terminal limit.

First choose $\eta_1$ so small that
$C_r^{(1)}\eta_1^2\le\eta_0^2/8$. Then choose a \emph{fixed} $K$ so large
that $E_*C_{\rm gr}(K+1)^4\rho^K\le\eta_0^2/8$. The native interaction has
fixed coordinatewise range, so a fixed $v_0$ contains all these graph balls.
Thus the right side of \eqref{f16v29:stopped-mixed-boundary} is at most
$C_r^{(0)}+\eta_0^2\gamma/4$. It holds throughout the signed source cube.
Integrating the actual conditional $\log Z$ derivative for a source only at
$c$ yields
\[
 \mathbb E e^{t_0e_c}\le
       \exp[t_0C_r^{(0)}+t_0\eta_0^2\gamma/4].
\]
Exponential Markov proves \eqref{f16v29:angular-tail}, for example with
$c_0=t_0\eta_0^2/2$ and a fixed prefactor. This argument keeps both the
moderate \emph{actual} exterior and the arbitrary artificial repair exterior.

For \eqref{f16v29:constant-repair-energy}, work in unmoving whole cells. Each
$S_c^\Omega$ has a fixed maximum number of cells and every cell occurs in only
a fixed number of these sets. Relative to the original product cell reference,
with one-cell classical terms included in their own factors, all remaining
words meeting a repair have total oscillation at most $C_r\gamma$. The bounded-density variance comparison of V24, or V25's
interface version, therefore bounds one repair energy by $C_re^{C_r\gamma}$
times the sum of its complete single-cell energies. This uses original cell
factors, including fractional endpoint Haar powers; no ratio to pure Haar at
its cut is taken. Sum the bounded overlap. Conditional on the remaining cells
these single-cell laws are exactly those of $\mu$. The radius and overlap are
fixed before coupling varies, so all fixed geometric factors can be included
in $C_r$. Arbitrary holes only reduce each repair's group count.
\end{proof}

The radius in this proposition is constant because the event being excluded
has energy of order $\gamma$. It does not supply uniform thermal moments at
constant depth, and it does not replace V27's logarithmic repair for a much
smaller rough threshold.

\subsection{A restricted heat-bath estimate and an all-test transfer to the full law}

\begin{proposition}[A nonsharp variance tensorization on the restricted product]
\label{f16v29:box-heatbath}
Let $\pi$ have smooth measured potential $V$ on a finite product of the
whole-cell balls, with $\lambda I\preceq\nabla^2V\preceq MI$. For its
rate-one whole-cell conditional updates,
\begin{equation}
 \operatorname{Var}_\pi f\le K_{\rm box}
                 \sum_c\mathbb E_\pi\operatorname{Var}_\pi(f\mid\zeta_{c^c}),
       \qquad K_{\rm box}=1+M/\lambda,
 \label{f16v29:box-tensorization}
\end{equation}
for every $f\in L^2(\pi)$. The constant has no number-of-cells factor.
\end{proposition}
\begin{proof}
Use the reflecting scalar realization on this product, with positive generator
$\mathcal A=-\Delta+\nabla V\cdot\nabla=\sum_c\mathcal A_c$. Its gap is at
least $\lambda$ by the inherited product-ball Brascamp--Lieb identity. For a
smooth function $u$ satisfying the Neumann condition on each ball face, the
full and blockwise integrated Bochner identities have the same nonnegative
boundary-curvature sum. The full squared Hessian contains every block-diagonal
squared Hessian as part of a sum of squares. Since $(\nabla^2V)_{cc}\preceq MI$
and $\nabla^2V\succeq0$, subtracting these identities gives
\begin{equation}
       \sum_c\|\mathcal A_cu\|_2^2
       \le\|\mathcal Au\|_2^2+M\langle u,\mathcal Au\rangle.
 \label{f16v29:partial-bochner}
\end{equation}
In particular no row-wise absolute diagonal-dominance condition is imposed.
The boundary curvature on a cell face acts on that cell's tangential
components only; cross-cell tangents have zero curvature. This is why the
boundary terms agree rather than introduce a number-of-cells loss.

The domain passage uses the same product reflecting realization as V4/V22.
At each finite regulator $V$ and its derivatives are bounded on the compact
product. Finite tensor sums of smooth Neumann eigenfunctions of the individual
balls form a graph core for the product Neumann Laplacian; the bounded smooth
first-order drift is relatively bounded with arbitrarily small relative bound.
Thus the same core is a graph core for $\mathcal A$ on weighted $L^2$.
Inequality \eqref{f16v29:partial-bochner} makes the partial generators Cauchy
under graph-core approximation. It therefore extends to $u$ in the full
scalar generator domain. Conditional integration by parts also extends
$E_c\mathcal A_cu=0$. This argument neither imposes a boundary condition on
$f$ nor substitutes the smooth-boundary theorem for an unmentioned corner
realization.

For centered $f\in L^2(\pi)$ solve $u=\mathcal A^{-1}f$. Then
$\langle u,\mathcal Au\rangle\le\lambda^{-1}\|f\|_2^2$ and
\[
 \|f\|_2^2
   =\sum_c\langle (I-E_c)f,\mathcal A_cu\rangle
   \le\left(\sum_c\|(I-E_c)f\|_2^2\right)^{1/2}
        \left(\sum_c\|\mathcal A_cu\|_2^2\right)^{1/2}.
\]
Insert \eqref{f16v29:partial-bochner} and divide by $\|f\|_2$, unless it is
zero. This proves \eqref{f16v29:box-tensorization}; complex tests follow by
real and imaginary parts. The constant is deliberately not optimized.
\end{proof}

\begin{proposition}[Variance transfer with the exceptional test mass paid]
\label{f16v29:variance-gluing}
Let $\mu$ be a finite-product probability, $G=\bigcap_cG_c$ a positive-probability
product event, and $\pi=\mu(\cdot\mid G)$. Suppose $\pi$ has whole-cell variance
tensorization constant $K_{\rm box}$ and, for every $g\in L^2(\mu)$,
\begin{equation}
 \|1_{G^c}g\|_2\le
       \sqrt{K\mathcal E_{\rm cell}^\mu(g)}+\sqrt\epsilon\|g\|_2,
             \qquad \epsilon\le\tfrac14.
 \label{f16v29:gluing-input}
\end{equation}
Then
\begin{equation}
 \operatorname{Var}_\mu f\le
       (2K_{\rm box}+4K)\mathcal E_{\rm cell}^\mu(f).
 \label{f16v29:gluing-output}
\end{equation}
No density lower bound of $\pi$ on $G^c$, and no invariance of $G$ under the
full conditional updates, is assumed.
\end{proposition}
\begin{proof}
For every coordinate $c$, conditional variance decomposition gives the exact
one-sided comparison
\begin{equation}
 \mu(G)\mathbb E_\pi\operatorname{Var}_\pi(f\mid\zeta_{c^c})
          \le\mathbb E_\mu\operatorname{Var}_\mu(f\mid\zeta_{c^c}).
 \label{f16v29:restriction-energy}
\end{equation}
Indeed condition on $z=\zeta_{c^c}$, restrict to $G_{c^c}$, and retain the
factor $\mu(G_c\mid z)$ multiplying the variance under the normalized
restricted conditional. It is one term of the full conditional variance.
Averaging proves \eqref{f16v29:restriction-energy}; no conditional normalizer
has been suppressed.

Put $a=\pi f$, $g=f-a$, and $T=\|g\|_{L^2(\mu)}$. Tensorization and
\eqref{f16v29:restriction-energy} bound
$\|1_Gg\|_2^2\le K_{\rm box}\mathcal E_{\rm cell}^\mu(f)$. The other part
is bounded by \eqref{f16v29:gluing-input}; all energies are unchanged by $a$.
Thus, with $\mathcal E=\mathcal E_{\rm cell}^\mu(f)$,
\[
 T^2\le K_{\rm box}\mathcal E+
                  (\sqrt{K\mathcal E}+\sqrt\epsilon T)^2
       \le (K_{\rm box}+2K)\mathcal E+2\epsilon T^2.
\]
Rearrange and use $\operatorname{Var}_\mu f\le T^2$. For an arbitrary $L^2$
test its restriction belongs to $L^2(\pi)$ because $\mu(G)>0$. The same
calculation, or truncation, justifies every conditional variance. Constants
and tests supported entirely on $G^c$ are both included.
\end{proof}

\begin{theorem}[Complete native conditional gap without an exponential size cost]
\label{f16v29:conditional-gap}
There are $A_r,c_0,C_r>0$ and a fixed-$r$ entry threshold such that, under
\eqref{f16v29:angular-collar}, for every finite $\Omega$ satisfying
\begin{equation}
                 |\Omega|A_re^{-c_0\gamma}\le\tfrac14,
 \label{f16v29:conditional-size}
\end{equation}
the \emph{complete} native conditional $\mu=P_\Omega^q$ obeys
\begin{equation}
 \boxed{\quad
 \operatorname{Var}_\mu f\le K_{\rm nat}(\gamma)
                     \mathcal E_{\rm cell}^{\mu}(f),
          \qquad K_{\rm nat}(\gamma)\le C_re^{C_r\gamma}.
 \quad}
 \label{f16v29:conditional-native-gap}
\end{equation}
The constants have no dependence on the shape or size of $\Omega$, ambient
volume, history length, or the admitted collar values. The collar condition
and \eqref{f16v29:conditional-size} are genuine hypotheses.
\end{theorem}
\begin{proof}
Use $G_\Omega$ defined above. The restricted law is exactly V22's native
product-ball conditional, and its Hessian has the same $\lambda,M$ when its
actual collar obeys \eqref{f16v29:angular-collar}. Hence
Proposition~\ref{f16v29:box-heatbath} applies with a fixed $K_{\rm box}$.
For every $c\in\Omega$, \eqref{f16v29:angular-tail} and V27's conditional
projection estimate give
\[
 \|1_{\{e_c>\eta_0^2\gamma\}}g\|_2
       \le\|Q_{S_c^\Omega}g\|_2+\sqrt{p_0}\|E_{S_c^\Omega}g\|_2.
\]
Direct-sum Minkowski and the union bound therefore give
\[
 \|1_{G_\Omega^c}g\|_2
       \le\sqrt{K_0(\gamma)\mathcal E_{\rm cell}^{\mu}(g)}
                                 +\sqrt{|\Omega|p_0}\|g\|_2.
\]
Apply Proposition~\ref{f16v29:variance-gluing} and
\eqref{f16v29:conditional-size}. This proves the theorem with
$K_{\rm nat}=2K_{\rm box}+4K_0$, bounded by the displayed exponential.
The probability term is used on the \emph{same test} before it is absorbed.
The interior is not replaced by its restriction. No all-field convexity or
pre-existing native gap is a premise.
\end{proof}

For instance every prescribed polynomial $|\Omega|$ is allowed eventually,
and an exponential $|\Omega|\le e^{\beta_0\gamma}$ is allowed for fixed
$\beta_0<c_0$. The cost $e^{C_r\gamma}$ remains. This theorem is not the
unproved all-compact-boundary conditional gap: surface configurations outside
\eqref{f16v29:angular-collar} have not been included.

\subsection{Fixed-depth repair of the actual slab collar}

We now work under the full native law $P$, with no conditioning or restriction
on its fast fields. Put
\[
 H_c^{(1)}=\{e_c>\eta_1^2\gamma\}.
\]
This is a fixed angular threshold, not V28's $Q_\gamma$. V20's complete
conditional preparation gives fixed-radius boxes $T_c$ and constants
$A_r',c_1>0$ such that
\begin{equation}
 P(H_c^{(1)}\mid\zeta_{T_c^c})\le
       p_1(\gamma):=\min\{1,A_r'e^{-c_1\gamma}\}.
 \label{f16v29:collar-smallness}
\end{equation}
To see this explicitly, choose a fixed graph depth $d_1$ with
$B_re^{-\kappa_rd_1}\le\eta_1^2/4$ in
\eqref{f16v20:prepared-mgf}. Exponential Markov at $t_0$ gives
\eqref{f16v29:collar-smallness}, for example with
$c_1=t_0\eta_1^2/2$. A fixed coordinate radius contains that depth, including
saturated cycles and natural ends. All compact $T_c$ exteriors are admitted.
As in \eqref{f16v29:constant-repair-energy}, bounded size and overlap yield
\begin{equation}
       \sum_c\mathcal E_{T_c}(f)
          \le K_{\rm rep}(\gamma)\mathcal E_{\rm cell}(f),
          \qquad K_{\rm rep}(\gamma)\le C_re^{C_r\gamma}.
 \label{f16v29:angular-repair-clock}
\end{equation}
Fix an integer $v$ larger than this coordinate radius plus two. It is fixed
independently of $\gamma,B,n$.

Use V22's exponential-width temporal channel with buffer $b=b_\gamma^{\rm exp}$,
$\ell=16b$, and $b\ge2v+2$. This is arranged by increasing the fixed entry
threshold. In the admissible channel regime $b\le C_r\gamma$. Let
\[
 I_s=[s,s+\ell-1]\cap[0,n],\quad
 J_s=[s-v,s+\ell-1+v]\cap[0,n],\quad 1-\ell\le s\le n.
\]
The corresponding updates contain all $B$ spatial cubes. Use the indexed
slab generator of V24, denoted here by $\mathcal L_{\rm sl}$, with rates
$1/\ell$. Let $C_s$ be its actual cell collar. Since every native interaction
meets at most two adjacent times,
\begin{equation}
 |C_s|\le2B,\qquad
       \sup_c\frac1\ell\#\{s:c\in C_s\}\le\frac2\ell,
       \qquad I_s\cup T_c\subseteq J_s\quad(c\in C_s).
 \label{f16v29:slab-collar-counts}
\end{equation}
Spatial components are never declared independent in these statements.

\begin{proposition}[A fixed enlargement is dominated by the original slab family]
\label{f16v29:enlargement-return}
In the above temporal channel regime, as quadratic forms on the full native
space,
\begin{equation}
       \frac1\ell\sum_s Q_{J_s}\preceq4\mathcal L_{\rm sl}.
 \label{f16v29:enlargement-domination}
\end{equation}
All compact exteriors and both endpoint factors are retained.
\end{proposition}
\begin{proof}
For each $s$ set $A_s=I_{s-v}$ and $B_s=I_{s+v}$, assigning the empty set
and zero innovation to an index outside the nonempty indexed family. Their
union is $J_s$ because $\ell>2v$. If one exclusive part is empty, one block
already equals $J_s$ and the needed inequality is immediate.

Otherwise condition on the actual exterior of $J_s$. Inside it the conditional
is the same inhomogeneous native time chain with the two fixed outside factors,
when present. The exclusive left and right parts are separated by an overlap
of length at least $\ell-2v$. Choose a $b$-step target inside this overlap,
at least $b$ from the right frozen endpoint. V22's actual forward channel,
with that right factor retained, has diameter at most $1/4$. Its
all-incoming-law $L^2$ consequence and Markov conditioning bound the maximal
correlation of the two exclusive parts by $1/2$. Appending the remaining
right path is a common conditional kernel and cannot increase this bound.
No estimate is applied to an output too close to the right frozen boundary.

On this fixed-exterior fibre, $E_{A_s}$ and $E_{B_s}$ are the projections onto
functions of the right and left exclusive parts. On centered space their
ranges have angle at most $1/2$. For vectors $x,y$ in those two centered
ranges, $|\langle x,y\rangle|\le\tfrac12\|x\|\|y\|$. Thus the
synthesis $(x,y)\mapsto x+y$ has squared norm at most $3/2$, and its
product with its adjoint gives $E_{A_s}+E_{B_s}\preceq3I/2$ on centered
space. Consequently $Q_{A_s}+Q_{B_s}\succeq\tfrac12 I$ there; constants
have zero on both sides.
Equivalently, before averaging, $Q_{J_s}\preceq2(Q_{A_s}+Q_{B_s})$ as a
conditional form. Integrate the actual exterior. Finally each shifted index
$s-v$ or $s+v$ occurs at most once in its own shifted sum of nonempty intervals.
Summing rates therefore gives \eqref{f16v29:enlargement-domination}. Repeated
clipped intervals retain their distinct indices, so very short histories are
covered by the same calculation.
\end{proof}

The proposition is independent of the gap comparison we are about to prove.
Its input is the already established conditional channel, not a presumed
original-clock lower bound. It is precisely the comparison that would be
missing for analogous overlapping spatial blocks.

\subsection{Rough-slab absorption closes without an elementary volume exponential}

Define exterior gates $M_s=1_{H_{C_s}^{(1)}}$, and write the actual slab energy
as $\mathcal E_{\rm sl}=\mathcal E_{\rm sl}^{\rm good}
+\mathcal E_{\rm sl}^{\rm bad}$, with
\[
 \mathcal E_{\rm sl}^{\rm bad}(f)
          =\frac1\ell\sum_s\|M_sQ_{I_s}f\|_2^2.
\]
These are exterior gates for \emph{these slabs}; none of V28's existing gates
is redefined.

\begin{proposition}[Local repair of all rough temporal boundary rows]
\label{f16v29:slab-absorption}
For every native $L^2$ test,
\begin{equation}
 \bigl(\mathcal E_{\rm sl}^{\rm bad}(f)\bigr)^{1/2}
 \le
 \sqrt{\frac{2K_{\rm rep}(\gamma)}{\ell}\mathcal E_{\rm cell}(f)}
       +\sqrt{8Bp_1(\gamma)\mathcal E_{\rm sl}(f)}.
 \label{f16v29:slab-rough}
\end{equation}
If also $B\ell A_re^{-c_0\gamma}\le1/4$, then
\begin{equation}
           \mathcal E_{\rm sl}^{\rm good}(f)
                    \le K_{\rm nat}(\gamma)\mathcal E_{\rm cell}(f).
 \label{f16v29:slab-good}
\end{equation}
The first inequality includes arbitrary rough collars. The second is integrated
only on its declared good collars, using their actual posterior weights.
\end{proposition}
\begin{proof}
Apply V28's nested inequality to $D=I_s$, $S=T_c$ and $H=H_c^{(1)}$ for each
$c\in C_s$, with conditional probability \eqref{f16v29:collar-smallness}.
The first sum of local energies is at most
$2K_{\rm rep}\mathcal E_{\rm cell}/\ell$ by
\eqref{f16v29:slab-collar-counts}. For the second, keep the local union until
using $I_s\cup T_c\subseteq J_s$ and drop only its nonnegative subtracted
energy. Its upper bound is
\[
 p_1\frac1\ell\sum_s |C_s|\mathcal E_{J_s}(f)
       \le 2Bp_1\frac1\ell\sum_s\mathcal E_{J_s}(f)
       \le8Bp_1\mathcal E_{\rm sl}(f).
\]
Direct-sum Minkowski proves \eqref{f16v29:slab-rough}. No full variance or
number of time slices enters this estimate.

On a good collar, the conditional inside $I_s$ satisfies
\eqref{f16v29:angular-collar} and has at most $B\ell$ cells. Apply
Theorem~\ref{f16v29:conditional-gap} at that exterior. Multiply its inequality
by the exterior indicator and integrate. Conditioning further on other cells
in the slab restores the original full single-cell conditional on the right.
Dropping its nonnegative gate gives an upper bound by those actual cell
energies. Each cell is in exactly $\ell$ indexed slabs, so summing their
rates gives \eqref{f16v29:slab-good}. The restriction to the good boundary was
not removed from a conditional theorem without payment: the preceding rough
form is exactly the complementary contribution.
\end{proof}

\begin{theorem}[An original-clock gap with no cross-section in its exponential cost]
\label{f16v29:main}
There exist fixed $\beta_r,c_r,C_r>0$ and $\gamma_r<\infty$ such that for every
\begin{equation}
 \gamma\ge\gamma_r,\qquad B=m^3\le e^{\beta_r\gamma},\quad m\ge4,\quad n\ge1,
                  \quad\max_{i,e}|y_{i,e}|\le r,
 \label{f16v29:width}
\end{equation}
the full native whole-cell sampler and the original unweighted rate-one
balanced sampler satisfy
\begin{equation}
 \boxed{\qquad
   \lambda_{\rm cell}\ge c_re^{-C_r\gamma},\qquad
   \widehat\lambda_P\ge c_re^{-C_r\gamma}.
 \qquad}
 \label{f16v29:main-gap}
\end{equation}
The constants do not depend on $B$ or $n$ within \eqref{f16v29:width}.
All native $L^2$ observables, complete compact coordinates, and original
endpoint factors are included. The rate of every elementary balanced update
remains one.
\end{theorem}
\begin{proof}
Let $\beta_{\rm ch}>0$ be the width exponent of V22's actual exponential-width
channel. Choose once
\[
        0<\beta_r<\min\{\beta_{\rm ch},c_0/4,c_1/4\}.
\]
Its buffer obeys $b\le C_r\gamma$. Increase the fixed threshold so that
$b\ge2v+2$, $B\ell A_re^{-c_0\gamma}\le1/4$, and
$16Bp_1(\gamma)\le1/2$ throughout \eqref{f16v29:width}. These follow from
strictly positive exponential margins; the factors $B$ and $B\ell$ have not
been ignored.

Square \eqref{f16v29:slab-rough} with $(x+y)^2\le2x^2+2y^2$, add
\eqref{f16v29:slab-good}, and rearrange:
\begin{equation}
 (1-16Bp_1)\mathcal E_{\rm sl}(f)
 \le\left(K_{\rm nat}+\frac{4K_{\rm rep}}\ell\right)
                                    \mathcal E_{\rm cell}(f).
 \label{f16v29:closed-slab-comparison}
\end{equation}
Both energy coefficients are at most $C_re^{C_r\gamma}$. V24's independently
proved slab gap is at least $1/2$ in its unit-coverage clock, hence
\[
 \operatorname{Var}_P f\le2\mathcal E_{\rm sl}(f)
                  \le C_re^{C_r\gamma}\mathcal E_{\rm cell}(f).
\]
This proves the first inequality of \eqref{f16v29:main-gap} without
conditioning away rough boundary rows or using a presumed native gap.

To return to the original balanced coordinates, one free unmoving cell
$(i,b)$ changes its five-chord group and, when present, the seven moving
ports at times $i-1$ and $i$ in that cube. Thus at most fifteen elementary
balanced groups are needed. The exact frame identity gives the same
exterior-sigma-algebra inclusion as V24. On this fixed support the original
product references and bounded classical words give a conditional variance
comparison of cost $C_re^{C_r\gamma}$. Each elementary chord occurs through
one cell and each moving port through at most two. Summing therefore gives
\begin{equation}
                    \mathcal L_{\rm cell}
                       \preceq C_re^{C_r\gamma}\widehat{\mathcal L}.
 \label{f16v29:cell-balanced}
\end{equation}
No fractional Haar factor is replaced by full Haar in this comparison.
Combining with the first inequality, and enlarging $C_r$, proves the second.
There is no hidden factor $1/[B(n+1)]$ in either generator. In particular,
\eqref{f16v29:main-gap} is not the old whole-slab density estimate with its
size deleted from the exponent; that estimate has been replaced by the
conditional gap, local rough repair, and overlap comparison proved above.
\end{proof}

The last theorem improves the \emph{scaling of the proved clock cost} inside
an explicit exponential-width regime. It does not assert that its unevaluated
constants are numerically optimal. At fixed coupling it allows only a bounded
set of spatial cross-sections. As $\gamma\to\infty$ its lower bound still
tends to zero; it is not a coupling-uniform gap or a calibrated physical mass.
V25's weaker all-finite-width baseline remains available for widths beyond
\eqref{f16v29:width}.

\subsection{A consequence for V28's actual retained generator}

\begin{proposition}[A paid spectral edge for the displayed retained dynamics]
\label{f16v29:retained}
Keep \emph{exactly} V28's $R_\gamma=\lceil h_\gamma^5\rceil$, gates, normalized
helper families, and generators $\mathcal L_{\rm aug},\mathcal L_{\rm keep}$.
For sufficiently large coupling in \eqref{f16v29:width},
\begin{equation}
 \boxed{\quad
    \lambda_{\rm keep}\ge R_\gamma^{-4}c_re^{-C_r\gamma},\qquad
    \lambda_{\rm aug}\ge\lambda_{\rm keep}.
 \quad}
 \label{f16v29:retained-gap}
\end{equation}
In particular both lower bounds can be written $c_r'e^{-C_r'\gamma}$ after
changing fixed constants. This is a statement in the width regime, not an
unrestricted spatial gap for either generator.
\end{proposition}
\begin{proof}
V28 proves, once its $\theta_\gamma<1$, the stronger form inequality
\[
 \mathcal L_{\rm keep}\succeq
           \mathcal L_R+(1-\theta_\gamma)
                         (\mathcal L_{\rm rep}+\mathcal L_+)
                 \succeq\mathcal L_R.
\]
For a cell $c\in D$, nested projections give $Q_c\preceq Q_D$. Thus
$\sum_{c\in D}Q_c\preceq |D|Q_D\preceq R_\gamma^4Q_D$. Sum the original
indexed unit-coverage rates to get
$\mathcal L_{\rm cell}\preceq R_\gamma^4\mathcal L_R$.
Theorem~\ref{f16v29:main} and variational gap comparison prove
\eqref{f16v29:retained-gap}. The polynomial logarithmic factor can be absorbed
into $e^{C_r'\gamma}$ at a fixed threshold.

This proof uses the full retained generator and V28's already established
rough-form absorption. It does not analyze only its good output rows, omit
gate switching from a path coupling, or switch off either helper family.
Detailed balance and the coverage clocks remain V28's. The new lower edge
comes from an independent ungated temporal-slab argument followed by form
comparison; it is not deduced circularly from relative closeness alone.
\end{proof}

\subsection{What has closed and what remains spatial}

The three comparisons have different jobs: fixed-depth angular repair bounds
the full conditional mixing cost by $e^{C_r\gamma}$; the overlapping-slab
argument pays the enlarged union energy; and nested rough-innovation
localization absorbs it into the same slab form. The final conversion to
balanced updates has fixed support. No duration factor or exponential cost
proportional to $B$ occurs in the resulting bound.

The argument still counts $B$ collar cells and $B\ell$ cells inside one slab.
Their exponentially small angular exceptions are paid only while
$Bp_1$ and $B\ell p_0$ are small. Arbitrary spatial volume at fixed coupling
would destroy this certificate. No common scale is taken to infinity with
those factors silently held fixed. Spatially, the analogue of
\eqref{f16v29:enlargement-domination} with useful, arbitrary-boundary constants
is not supplied here. Nor is a gap for a full conditional at every compact
surface field inferred from the controlled angular-collar theorem.

The remaining productive directions are therefore an unrestricted spatial
union-energy comparison or a direct spatial dependence estimate, and,
separately, a sharper fixed-support repair cost if a nonvanishing
weak-coupling sampler floor is desired. Another temporal power iteration
cannot remove the $B$-dependent admissibility conditions. Substitution of the
present $e^{-C_r\gamma}$ bound into a gap-only source inequality is not
advertised as an improvement over V23's direct $1/\gamma$ covariance bound.
All source definitions, $p_{\rm good}$, and prior conditional return matrices
are unchanged.

Every clock here is auxiliary. The fixed bridge window is retained. Root
Gauss coverage, unrestricted bridge integration, infrared contact coercivity,
comparison to the physical transfer at its own top, independent physical-time
calibration, and nontrivial interacting continuum reconstruction remain
separate obligations. No Yang--Mills physical mass gap follows from an
uncalibrated original-sampler lower bound.

\subsection{Diagnostics, preservation, and verification scope}

The accompanying diagnostic is
\begin{center}\small\path{verify_ym_f16_v29_angular_clock.py}.\end{center}
It checks finite conditional projections and restriction energies, the
variance-gluing algebra, Gaussian block Bochner identities, stopped-boundary
bookkeeping, exact interval unions and multiplicities, retained-endpoint
conditional angles, nonextensive slab absorption, the fifteen-group frame
support, and the original-rate conversions. The diagnostics are rational,
integer, or symbolic fixtures, not native compact-integral enclosures or
proof-assistant verification. They do not certify the inherited product-ball
domain or evaluate the constants $\eta_1,v_0,c_0,c_1,\beta_r,C_r$.

The package records the actual tests and any executed predecessor reruns,
source hashes, exact reconstruction, final compilation, and rendered-page
inspection. V28 is preserved except for the one title version and this exact
terminal insertion. The next companion checkpoint remains V30. The new
analytic deductions rely on the explicitly identified set-energy and
reflecting-domain inputs; the historical archive has not been independently
recertified.

\begin{firewall}
V29 proves full-conditional mixing on a declared angular collar and an
original-clock lower bound $c_re^{-C_r\gamma}$ for
$B\le e^{\beta_r\gamma}$, uniformly in duration. It also controls V28's
retained generator in that regime. These deductions retain every native
normalizer, compact output, endpoint factor, source, and elementary rate.
They do not give unrestricted spatial volume, a coupling-uniform sampler gap,
or a physical transfer or interacting continuum mass gap.
\end{firewall}
% END F16 V29 APPEND

% BEGIN F16 V30 APPEND -- 2026-09-18
% Insert immediately before the final document-ending command in V29.
% The predecessor, its native laws, sources, and elementary clock are preserved.

\clearpage
\section{V30: fixed-range block coercivity and spatial screening}
\label{f16v30:context}

V29's exponential elementary-clock cost enters when a fixed angular repair is
resolved into individual cell updates. It is not necessary to pay that cost
when the repair is itself an allowed, exactly normalized block update. We keep
such blocks at one fixed radius, chosen before coupling varies. A buffered-cover
variance inequality then has a coupling-independent constant. Combining it with
V29's independent temporal overlap argument gives a genuine positive gap for a
\emph{fixed-range} native block generator. No block radius or rate grows with
coupling, cross-section, or duration.

The spatial-width condition remains: the result holds for
$B\le e^{\beta_r\gamma}$, with a fixed $\beta_r>0$, and every duration. Within
that regime the new local gap also gives a finite-layer contraction and
exponential \emph{space--time} screening at every radius, without an algebraic
or nonperturbative distance floor. This yields the original coefficient
expansions in a weighted absolute space--time row norm. It is not an
unrestricted-spatial-volume theorem and it does not turn the block clock into
the original elementary or physical clock.

\subsection{V30 checkpoint and the precise dependency change}

The complete V29 appendix and audit were read. Its restricted product-ball
variance theorem, two angular repair scales, and retained-boundary slab
comparison are the analytic premises used here. V29's final original-clock gap
is \emph{not} an input to the new block gap. The V25 checkpoint and the supplied
f14 V100, f15 V100, and Grand Tensor V076 were included in targeted section and
text inventories. Empty semantic searches were not interpreted as absence of
prior work. Neither the whole archive nor unavailable companion programs are
represented as independently certified; the separate checkpoint records the
scope. The next scheduled checkpoint is V35.

V13 already identifies conditional projections as a local frustration-free
auxiliary parent and cites detectability. We reuse that construction. The
projection-product estimate below is the classical argument of
Anshu--Arad--Vidick,
\href{https://arxiv.org/abs/1602.01210}{\emph{A simple proof of the detectability
lemma and spectral gap amplification}}, Lemma 2 and Corollary 3. Their
statement and proof were examined. The proof uses only bounded projections,
so we give it on the present Hilbert space rather than assume a finite local
spin dimension. No new general detectability or Knabe theorem is claimed.
The new input that makes this principle quantitative here is a native gap for
blocks of fixed radius in an explicitly normalized clock.

The import memoranda's distinction between locality cost and perturbative
order remains important. The block-gap proof uses no higher-order expansion.
V16's moments and V17's entry expansion are used only afterward, to convert
spatial screening into the same source coefficients. The memoranda's final
physical-transfer arrow is not an analytic premise.

\subsection{A buffered cover of one conditional, without a microscopic comparison}

Keep $P=P_{\gamma,\mathbf y}$, $B=m^3$, $m\ge4$, $n\ge1$, and
$\max_{i,e}|y_{i,e}|\le r$. For a free set $\Omega$ let $\mu=P_\Omega^q$.
Retain V29's fixed radii $\eta_1<\eta_0$, the event
$G_\Omega=\bigcap_{c\in\Omega}\{e_c\le\eta_0^2\gamma\}$, its repairs
$S_c^\Omega$, and
\[
 p_0(\gamma)=\min\{1,A_re^{-c_0\gamma}\},\qquad
 K_{\rm box}=1+M/\lambda,\qquad K_{\rm buf}=2K_{\rm box}+4.
\]
For a free subset $A\subseteq\Omega$, conditional expectation and its
innovation under $\mu$ are denoted $E_A^\mu,Q_A^\mu$. Their form is
$\mathcal E_A^\mu(f)=\|Q_A^\mu f\|_\mu^2$. The complete actual collar is
required to obey
\begin{equation}
 |q_d|\le\eta_1\sqrt\gamma\quad(d\in\partial_{\mathcal G}\Omega),
 \qquad |\Omega|p_0(\gamma)\le\tfrac14.
 \label{f16v30:conditional-regime}
\end{equation}
There is no restriction on any free coordinate in $\mu$.

\begin{proposition}[Regrouping the restricted coordinates]
\label{f16v30:grouping}
Let $\pi=\mu(\cdot\mid G_\Omega)$ and let
$\Omega=V_1\sqcup\cdots\sqcup V_t$ be any partition into unions of complete
cells. If $V_j\subseteq A_j\subseteq\Omega$, then
\begin{equation}
 \operatorname{Var}_\pi f\le K_{\rm box}
                    \sum_{j=1}^t\mathcal E_{A_j}^\pi(f).
 \label{f16v30:restricted-cover}
\end{equation}
The constant is independent of the dimensions of the groups and their number.
\end{proposition}
\begin{proof}
Apply the proof of Proposition~\ref{f16v29:box-heatbath} with the reflecting
partial generators grouped as $\mathcal A_{V_j}=\sum_{c\in V_j}\mathcal A_c$.
The full Hessian-square still contains the sum of all within-group
Hessian-squares. Each principal measured-Hessian block is bounded above by
$MI$, while the full measured Hessian is nonnegative. Every cell face appears
in exactly one group, so the same nonnegative boundary-curvature sum cancels
in the full and grouped Bochner identities. Thus
\[
 \sum_j\|\mathcal A_{V_j}u\|_\pi^2
 \le\|\mathcal Au\|_\pi^2+M\langle u,\mathcal Au\rangle_\pi.
\]
The graph-core argument is exactly V29's inherited product reflecting
realization. With $u=\mathcal A^{-1}(f-\pi f)$,
$E_{V_j}^\pi\mathcal A_{V_j}u=0$ and Cauchy--Schwarz give the variance bound
with $V_j$ in place of $A_j$. Conditional-projection nesting gives
$\mathcal E_{V_j}^\pi\le\mathcal E_{A_j}^\pi$. Empty groups can be omitted.
No connectedness of a group, absolute diagonal dominance, or smooth-boundary
replacement for the product domain is required.
\end{proof}

\begin{theorem}[Complete native tensorization over a buffered cover]
\label{f16v30:buffered-cover}
Assume \eqref{f16v30:conditional-regime} and the inherited fixed-$r$ threshold.
Let $A_1,\ldots,A_t\subseteq\Omega$ admit a partition
$\Omega=V_1\sqcup\cdots\sqcup V_t$ such that
\begin{equation}
                 S_c^\Omega\subseteq A_j\quad(c\in V_j).
 \label{f16v30:buffer-cover-condition}
\end{equation}
Then, on the complete conditional space,
\begin{equation}
 \boxed{\quad
 \operatorname{Var}_\mu f\le K_{\rm buf}
                         \sum_{j=1}^t\mathcal E_{A_j}^\mu(f).
 \quad}
 \label{f16v30:buffered-cover-bound}
\end{equation}
No factor depending on coupling, the number of covering blocks, or the size
of $\Omega$ occurs in $K_{\rm buf}$. The two hypotheses in
\eqref{f16v30:conditional-regime} have not been removed.
\end{theorem}
\begin{proof}
Put $H_j=\bigcup_{c\in V_j}\{e_c>\eta_0^2\gamma\}$.
V29's \emph{uniform conditional} angular repair bound and
$S_c^\Omega\subseteq A_j$ imply, by the tower property,
\[
            \mu(H_j\mid\zeta_{\Omega\setminus A_j})
                          \le |V_j|p_0.
\]
This does not require independence of the repairs or their events. Decompose
a test $g$ into $Q_{A_j}^\mu g+E_{A_j}^\mu g$. The conditional projection
estimate from V27, the union bound, and direct-sum Minkowski give
\begin{equation}
 \|1_{G_\Omega^c}g\|_\mu
 \le\left(\sum_j\mathcal E_{A_j}^\mu(g)\right)^{1/2}
                       +\sqrt{|\Omega|p_0}\,\|g\|_\mu.
 \label{f16v30:block-exception}
\end{equation}
Crucially, each repair is already contained in an allowed update. We have not
estimated that update by the sum of elementary cell energies.

The exact restriction-energy comparison of V29 holds for any block $A_j$:
\[
 \mu(G_\Omega)\mathcal E_{A_j}^\pi(f)
                                    \le\mathcal E_{A_j}^\mu(f).
\]
To check its normalizer, condition on the exterior of $A_j$, restrict that
exterior to $G_{\Omega\setminus A_j}$, and retain
$\mu(G_{A_j}\mid\zeta_{\Omega\setminus A_j})$ multiplying the restricted
conditional variance. It is one nonnegative part of the full conditional
variance. The comparison uses the same product event, also when the $A_j$
overlap.

Set $g=f-\pi f$, $T=\|g\|_\mu$, and
$\mathcal F=\sum_j\mathcal E_{A_j}^\mu(f)$. Proposition~\ref{f16v30:grouping}
and this restriction comparison give $\|1_{G_\Omega}g\|_\mu^2\le
K_{\rm box}\mathcal F$. Equation~\eqref{f16v30:block-exception} then yields
\[
 T^2\le K_{\rm box}\mathcal F+
       (\sqrt{\mathcal F}+\sqrt{|\Omega|p_0}T)^2
 \le(K_{\rm box}+2)\mathcal F+2|\Omega|p_0T^2.
\]
Absorb the last term and use $\operatorname{Var}_\mu f\le T^2$.
Restriction of an $L^2(\mu)$ function lies in $L^2(\pi)$ because
$\mu(G_\Omega)>0$. Thus the proof includes constants, complex functions,
and tests supported on exceptional configurations.
\end{proof}

\begin{proposition}[A spatial union comparison, with the rough exterior retained]
\label{f16v30:spatial-union}
Suppose $U=A\cup B$ has a two-block cover satisfying
\eqref{f16v30:buffer-cover-condition}. On each admitted angular exterior of $U$,
provided $|U|p_0\le1/4$, its complete conditional obeys
\begin{equation}
                 Q_U\preceq K_{\rm buf}(Q_A+Q_B).
 \label{f16v30:good-union}
\end{equation}
For example, split an ordinary box along one coordinate into two overlapping
boxes with overlap of at least $2v_0+2$ cells. Assign each center to a side
containing its $v_0$-repair. This includes spatial splits, not just temporal
slabs.

Under the actual, unrestricted native law $P$, let $C=\partial_{\mathcal G}U$,
use V29's fixed-radius $T_c$ and $p_1=\min(1,A'_re^{-c_1\gamma})$, and let
$U^+\supseteq U\cup T_c$ for every $c\in C$. For any $t>0$,
\begin{equation}
 \begin{split}
 Q_U\preceq {}&K_{\rm buf}(Q_A+Q_B)
             +(1+t)\sum_{c\in C}Q_{T_c}\\
             &+(1+t^{-1})p_1|C|Q_{U^+}.
 \end{split}
 \label{f16v30:rough-union}
\end{equation}
This is an all-test-function form inequality, not an all-boundary version of
\eqref{f16v30:good-union}. The condition $|U|p_0\le1/4$ remains.
\end{proposition}
\begin{proof}
The first claim is Theorem~\ref{f16v30:buffered-cover} on the conditional fibre;
conditional variance on $U$ gives its operator interpretation. For the box
example choose an assignment plane in the middle of the overlap. Its distance
from the opposite side's end is at least $v_0$; clipping at the actual boundary
of $U$ removes rather than adds repair cells.

Let $M=1_{\cup_{c\in C}\{e_c>\eta_1^2\gamma\}}$. It is exterior to $U,A,B$.
Integrating \eqref{f16v30:good-union} on $M=0$ gives
$(1-M)Q_U\preceq K_{\rm buf}(Q_A+Q_B)$. For $MQ_U$, use the union bound and
V28's nested Young inequality on $(U,T_c)$, with the actual conditional
smallness $P(e_c>\eta_1^2\gamma\mid T_c^c)\le p_1$. It gives
\[
 MQ_U\preceq(1+t)\sum_{c\in C}Q_{T_c}
 +(1+t^{-1})p_1\sum_{c\in C}(E_{T_c}-E_{U\cup T_c}).
\]
Each excess projection is nonnegative and bounded above by $Q_{U^+}$.
This proves \eqref{f16v30:rough-union} without dropping the rough exterior
or assuming that overlapping conditional expectations commute.
\end{proof}

\subsection{One fixed-radius native generator and its fixed coverage clock}

Choose a fixed integer $v_*$ at least the radii of both V29 repair families
$S_c^\Omega$ and $T_c$. Enlarge it once for the fixed coordinate-range allowances.
For every ambient cell $c$, let $U_c$ be its coordinatewise box of radius $v_*$,
saturating short spatial cycles and clipping at natural time endpoints. Put
\begin{equation}
 \kappa_*=(2v_*+1)^4,\qquad
 \mathcal L_* =\kappa_*^{-1}\sum_c Q_{U_c},\qquad
 \mathcal E_*(f)=\kappa_*^{-1}\sum_c\mathcal E_{U_c}(f).
 \label{f16v30:fixed-generator}
\end{equation}
Each indexed update has rate $1/\kappa_*$. The index is the center, even if two
clipped boxes coincide. Each cell has coverage at most one and at least
$1/\kappa_*$. Both the radius and normalization are independent of coupling,
volume, and duration. A move resamples the \emph{whole} native $U_c$ conditional,
including its buffer and every compact field.

\begin{theorem}[A noncollapsing fixed-range auxiliary gap in the width regime]
\label{f16v30:fixed-gap}
There exist fixed $\beta_r>0$ and $\gamma_r<\infty$ such that, for
\begin{equation}
 \gamma\ge\gamma_r,\qquad B=m^3\le e^{\beta_r\gamma},\quad
 m\ge4,\quad n\ge1,\quad\max_{i,e}|y_{i,e}|\le r,
 \label{f16v30:width}
\end{equation}
the complete native generator \eqref{f16v30:fixed-generator} satisfies
\begin{equation}
 \boxed{\quad \mathcal L_*\succeq g_*(I-\mathbb E_P),\qquad
 g_*:=\frac{1}{4\kappa_*(K_{\rm buf}+4)}>0.\quad}
 \label{f16v30:fixed-gap-bound}
\end{equation}
The lower bound does not deteriorate with coupling, spatial width within
\eqref{f16v30:width}, or duration. It is not a gap claim at unrestricted width
or in the elementary balanced clock.
\end{theorem}
\begin{proof}
Use V29's actual temporal buffer $b=O_r(\gamma)$, $\ell=16b$, its unchanged
clipped slabs $I_s$, and its independently proved gap
$\operatorname{Var}_P f\le2\mathcal E_{\rm sl}(f)$. Retain V29's fixed
slab enlargement and its form bound
$\ell^{-1}\sum_s Q_{J_s}\preceq4\mathcal L_{\rm sl}$.
No elementary or fixed-block gap was used to establish either input.

On a good angular collar of $I_s$, apply
Theorem~\ref{f16v30:buffered-cover} with blocks $A_c=I_s\cap U_c$ and
assigned sets $V_c=\{c\}$, $c\in I_s$. Each contains $S_c^{I_s}$.
At this fixed exterior it gives conditional variance at most
$K_{\rm buf}\sum_{c\in I_s}\mathcal E_{I_s\cap U_c}^{\mu}(f)$,
provided $B\ell p_0\le1/4$. Integrating the good gate gives the corresponding
\emph{global} subset innovations; then $Q_{I_s\cap U_c}\preceq Q_{U_c}$.
Each center occurs in exactly $\ell$ indexed slabs, so
\begin{equation}
           \mathcal E_{\rm sl}^{\rm good}(f)
                        \le K_{\rm buf}\kappa_*\mathcal E_*(f).
 \label{f16v30:slab-good}
\end{equation}
The gate was dropped only from a nonnegative upper bound after conditioning.

Repeat the exact nested calculation in Proposition~\ref{f16v29:slab-absorption},
but keep $Q_{T_c}$ as a block energy instead of comparing it to cell energies.
Since $T_c\subseteq U_c$,
$\sum_c\mathcal E_{T_c}\le\kappa_*\mathcal E_*$. The unchanged collar counts
and the unchanged enlargement return give
\begin{equation}
 \sqrt{\mathcal E_{\rm sl}^{\rm bad}(f)}
 \le\sqrt{\frac{2\kappa_*}{\ell}\mathcal E_*(f)}
                          +\sqrt{8Bp_1\mathcal E_{\rm sl}(f)}.
 \label{f16v30:slab-bad}
\end{equation}
Choose $\beta_r<\min(\beta_{\rm ch},c_0/4,c_1/4)$ once. At a fixed threshold,
$B\ell p_0\le1/4$ and $16Bp_1\le1/2$ hold throughout
\eqref{f16v30:width}. Squaring \eqref{f16v30:slab-bad}, adding
\eqref{f16v30:slab-good}, and absorbing yields
\[
 (1-16Bp_1)\mathcal E_{\rm sl}(f)
 \le\kappa_*(K_{\rm buf}+4/\ell)\mathcal E_*(f).
\]
Thus $\operatorname{Var}_P f\le4\kappa_*(K_{\rm buf}+4)\mathcal E_*(f)$.
All constants used to choose $U_c$ preceded the coupling limit. No hidden
factor $1/[B(n+1)]$, exponentially large repair rate, or comparison of a full
slab with elementary updates occurs in this proof.
\end{proof}

\begin{proposition}[Locality and the elementary-clock boundary]
\label{f16v30:local-parent}
The generator $\mathcal L_*$ preserves $P$, is reversible, and has constants
as its kernel. In the exact square-root-density realization on the original
product reference it is a sum of bounded fixed-support projections with a
common null vector. Each summand fails to commute with at most $\mathfrak d_*$ others,
for a fixed finite $\mathfrak d_*$ independent of coupling and regulator.

In the original balanced clock one still has only the paid comparison
\begin{equation}
       \mathcal L_*\preceq C_re^{C_r\gamma}\widehat{\mathcal L}.
 \label{f16v30:elementary-clock}
\end{equation}
Consequently the new constant $g_*$ does not assert a coupling-independent
lower bound for $\widehat{\mathcal L}$.
\end{proposition}
\begin{proof}
Conditional projections preserve the exact measure. Every cell belongs to its
own $U_c$, so a function in the joint kernel is exterior-measurable for every
single cell. Equivalence with the finite product reference makes it constant.
For clarity, let $p_U(x\mid z)$ be the actual conditional density relative to
that product reference on $U$. Conjugating $E_U$ by multiplication by
$\sqrt{dP/d\nu}$ gives the fibre kernel
\[
 F(x,z)\longmapsto
 \sqrt{p_U(x\mid z)}\int\sqrt{p_U(y\mid z)}F(y,z)\,d\nu_U(y).
\]
The exterior marginal cancels. This kernel depends only on $U$ and its actual
interaction collar. The rank-one fibre projector fixes the global square-root state; its
complement annihilates that state and is a local positive projection. This is V13's auxiliary-parent
construction, now with fixed-radius blocks and a quantitative gap.

A fixed coordinate enlargement of $U_c$ contains its kernel support. The graph
of overlapping enlarged supports has uniformly bounded degree and admits a
fixed number of colors with disjoint supports in each color. Disjoint supports
commute, giving the asserted bound $\mathfrak d_*$ (take $\mathfrak d_*\ge1$).

For \eqref{f16v30:elementary-clock}, include all chord groups of $U_c$ and all
moving ports whose cell at either adjacent chord time lies in $U_c$. This
includes the preceding ports. The support has a fixed number of elementary
groups. Its complete classical oscillation is $C_r\gamma$ relative to the
original product factors, including fractional endpoint Haar powers. The fixed
conditional comparison and bounded overlap give \eqref{f16v30:elementary-clock}.
This factor, unlike the radius and block rates, can diverge with coupling.
\end{proof}

\subsection{A finite-layer contraction and all-radius spatial screening}

Color the enlarged operator supports just described using $s_*<\infty$ colors.
Put
\[
 \mathcal S=\left(\prod_{c\in C_{s_*}}E_{U_c}\right)\cdots
             \left(\prod_{c\in C_1}E_{U_c}\right),\qquad
 \Pi_P f=Pf.
\]
A color can be updated in parallel: its disjoint blocks have no interaction
between them after their one common exterior is fixed. All conditional laws
and the subsequent layers remain native. One sweep has a fixed number of
layers, not a single globally rate-one elementary update.

\begin{theorem}[Native sweep contraction and relative space--time clustering]
\label{f16v30:clustering}
Throughout \eqref{f16v30:width}, set
\begin{equation}
           q_*^2=\frac{\mathfrak d_*^2}{\mathfrak d_*^2+\kappa_*g_*}<1.
 \label{f16v30:sweep-rate}
\end{equation}
Then $\|\mathcal S-\Pi_P\|\le q_*$ and
$\|\mathcal S^k-\Pi_P\|\le q_*^k$. There are fixed $a_r,C_r>0$ such that
for arbitrary $F,G\in L^2(P)$ measurable in cell sets $A,B$,
\begin{equation}
 \boxed{\quad
 |\operatorname{Cov}_P(F,G)|
 \le C_re^{-a_r d_*(A,B)}
               \sqrt{\operatorname{Var}_PF\operatorname{Var}_PG}.
 \quad}
 \label{f16v30:relative-clustering}
\end{equation}
Here $d_*((i,b),(j,c))=|i-j|+d_{\mathbb T_m^3}(b,c)$ as in V19.
There is no cardinality prefactor, no distance-independent coupling floor,
and no boundary conditioning beyond the stated native law.
\end{theorem}
\begin{proof}
We give the standard projection-product argument to specify its Hilbert-space
scope. Order the blocks by the color sweep, write $Q_i=Q_{U_i}$, and put
$f_i=(I-Q_i)f_{i-1}$. Orthogonality gives
\[
 \sum_i\|Q_if_{i-1}\|^2=\|f\|^2-\|\mathcal Sf\|^2.
\]
At the end of update $i$, $Q_if_i=0$. A later commuting projection cannot
increase $\|Q_if_j\|$; a noncommuting update can increase it by at most
$\|Q_jf_{j-1}\|$. Cauchy--Schwarz over at most $\mathfrak d_*$ such updates and then
summing over $i$ therefore give
\[
 \sum_i\|Q_i\mathcal Sf\|^2
       \le \mathfrak d_*^2\{\|f\|^2-\|\mathcal Sf\|^2\}.
\]
For centered $f$, invariance keeps $\mathcal Sf$ centered. The gap of
$\sum_iQ_i=\kappa_*\mathcal L_*$ bounds the left side below by
$\kappa_*g_*\|\mathcal Sf\|^2$. Rearranging proves
\eqref{f16v30:sweep-rate}; constants are fixed by every factor, so powers follow.
This is the detectability estimate, not an additional gap assumption.

For clustering work in the exact product-reference representation, and write
$\Psi=\sqrt{dP/d\nu}$. Multiplication by a local bounded observable has its
ordinary cell support and commutes with the conjugation. Let $\upsilon_*$ denote a
fixed upper bound on the propagation distance of one full color sweep in
$d_*$. Such a bound exists since every layer has disjoint fixed-diameter
operator supports. In the expression
\[
              \langle\Psi,F^*\mathcal S^k G\Psi\rangle
\]
(remove the conjugation from the notation only), delete each projector outside
the expanding support of $F$ by commuting it to the left vacuum. After $k$
sweeps the remaining projectors lie within distance $k\upsilon_*$ of $A$. If
$k\upsilon_*<d_*(A,B)$, they commute with $G$ and can then be moved to the right
vacuum. The result is exactly $\langle\Psi,F^*G\Psi\rangle$.

Apply this identity to centered $F,G$. The sweep norm bounds their covariance
by $q_*^k\|F\Psi\|\|G\Psi\|$. Take, for example,
$k=\lfloor d_*(A,B)/(2\max(1,\upsilon_*))\rfloor$, with $k=0$ when the supports
meet. This gives \eqref{f16v30:relative-clustering} with
$a_r=-\log(q_*)/[2\max(1,\upsilon_*)]$ and an enlarged fixed prefactor. All
projector deletions use bounded local operators fixing \emph{the same} state.
No independence of disjoint native supports is asserted. Truncate local
$L^2$ functions and pass in $L^2(P)$ to remove boundedness. This retains their
supports and proves the all-$L^2$ assertion.
\end{proof}

The sweep and the resulting clustering do not refer to the physical Wilson
Hamiltonian. A sweep entails a number of local operations proportional to
the system size, organized in a fixed number of parallel layers. Nor does
\eqref{f16v30:relative-clustering} assert a bound after arbitrary new pinning:
that operation changes the native probability and its generator.

\subsection{Original source coefficients in a weighted space--time row norm}

Keep V7's primitive errors, deterministic synthesis $T=(I,\mathsf P^{\mathsf T})$,
and covariance $\Gamma_\gamma=TK_\gamma T^{\mathsf T}$. The primitive carriers
have fixed radius. V16 bounds their variances by $C_r\gamma^{-1}$. Applying
\eqref{f16v30:relative-clustering} and summing the exponentially localized
synthesis rows gives, at a fixed exponent $a_0>0$,
\begin{equation}
 |(K_\gamma)_{uv}|\le C_r\gamma^{-1}e^{-a_0d_*(u,v)},\qquad
 |(\Gamma_\gamma)_{\alpha\beta}|
                    \le C_r\gamma^{-1}e^{-a_0d_*(\alpha,\beta)}.
 \label{f16v30:native-spatial-envelope}
\end{equation}
Anchors are used in this notation; their fixed carrier allowances are included
in $C_r$. Decrease $a_0$ below the original synthesis exponent. The common
source is not falsely assigned a finite spatial carrier.

For a scalar matrix in canonical source coordinates define
\[
 \|A\|_{\mathrm{st},a}
       =\sup_\alpha\sum_\beta
          e^{a d_*(\alpha,\beta)}|A_{\alpha\beta}|.
\]
This is an absolute \emph{space--time} row norm, not a trace norm or the earlier
temporal norm of spatial operator blocks.

\begin{theorem}[Full space--time row asymptotics in the same width regime]
\label{f16v30:source-expansion}
There is a fixed $a>0$, independent of coupling and expansion order, such that
for every fixed $J\ge2$, at a sufficiently large $\gamma_{J,r}$ in
\eqref{f16v30:width},
\begin{equation}
 \boxed{
 \begin{aligned}
 &\left\|\Gamma_\gamma-\sum_{j=2}^J\gamma^{-j/2}G_j\right\|_{\mathrm{st},a}
 \\
 &\quad+
 \left\|\mathsf M_\gamma-\sum_{j=2}^J\gamma^{-j/2}M_j\right\|_{\mathrm{st},a}
                   \le C_{J,r}\gamma^{-(J+1)/2}.
 \end{aligned}}
 \label{f16v30:source-row}
\end{equation}
The coefficients are precisely V10's normalized coefficients. In particular
both native norms are $O_r(\gamma^{-1})$, and their leading-coefficient errors
are $O_r(\gamma^{-3/2})$. All rows and spatial displacements are included,
for every duration, with no source-count normalization.
\end{theorem}
\begin{proof}
Choose $a>0$ so small that \eqref{f16v30:native-spatial-envelope}, polynomial
four-dimensional ball growth, and the coefficient locality of V10 give
\[
 \|\Gamma_\gamma\|_{\mathrm{st},4a}\le C_r\varepsilon^2,
 \qquad \|G_j\|_{\mathrm{st},4a}\le C_{j,r},\qquad
 \varepsilon=\gamma^{-1/2}.
\]
For the chosen $J$ put
$D=\lceil(J+2)(3a)^{-1}\log(\varepsilon^{-1})\rceil$ and $K=2J+6$.
On pairs with distance at most $D$, use V17's order-$K$ entry expansion.
At most $C(1+D)^4$ source labels occur in each row there, so its weighted
remainder is at most
\[
 C_{K,r}(1+D)^4 e^{aD}\varepsilon^{K+1}
 \le C_{J,r}(1+\log\varepsilon^{-1})^4
                       \varepsilon^{K+1-(J+2)/3}
 =O_{J,r}(\varepsilon^{J+1}).
\]
The extra coefficients of orders $J+1$ through $K$ are bounded using their
whole row norms. On the complement use genuine native and coefficient
exponential rows: the weighted tail is at most
$C_{J,r}\varepsilon^2e^{-3aD}=O_{J,r}(\varepsilon^{J+1})$.
There is no flat remainder to sum and no full spatial slice count.
The exact identity $\mathsf M_\gamma=-\operatorname{Off}_{>1}\Gamma_\gamma$
and the same identity for $M_j$ remove entries and cannot increase this
absolute norm. They are not claimed to preserve positivity.
\end{proof}

\begin{proposition}[All-radius screening of the same shell return]
\label{f16v30:shell}
For V19's unchanged centered primitive messages
$m_{u,R}=\mathbb E_P[F_u\mid\zeta_{I_{u,R}^c}]$, there are fixed $a',C_r>0$
such that, throughout \eqref{f16v30:width} and all $R\ge R_0$,
\begin{equation}
 \sup_u\|m_{u,R}\|_2\le C_r\gamma^{-1/2}e^{-a'R},\qquad
 \|\mathsf R^{\rm sh}_{\gamma,R}\|_{\rm op}
                                  \le C_r\gamma^{-1}e^{-a'R}.
 \label{f16v30:shell-bound}
\end{equation}
A ball covering the whole system has zero centered message. No spatial-torus
diameter must first be crossed and no coupling-dependent floor remains.
\end{proposition}
\begin{proof}
Dualize \eqref{f16v30:relative-clustering} over centered exterior $L^2$ tests.
Their supports are separated from the primitive carrier by $R$ minus a fixed
allowance. This proves the first bound, using V16's variance estimate.
The message is measurable on its finite interaction collar, contained in the
ball of radius $R+\rho$. For two anchors farther apart than $2R+2\rho$, apply
\eqref{f16v30:relative-clustering} again to their messages; for closer anchors
use their small variances. Polynomial ball growth then bounds the absolute
row sum of their Gram by
$C_r\gamma^{-1}(1+R)^4e^{-2a'R}$ after initially choosing a common smaller
exponent. Absorb the polynomial by a final decrease of $a'$. The original
bounded synthesis $T$ transfers its operator bound to
$\mathsf R^{\rm sh}_{\gamma,R}=T\operatorname{Cov}(m_R)T^{\mathsf T}$.
This is V19's object, not V12's different common-exterior posterior.
\end{proof}

\subsection{Checkpoint conclusions and the unresolved extrapolation}

The coupling loss and the width restriction are different issues. The first
is removed for one fixed-range auxiliary block generator by retaining its
repairs as actual moves. A microscopic comparison still costs
$C_re^{C_r\gamma}$. The second survives: the gap proof used both
$B\ell p_0\le1/4$ and $16Bp_1\le1/2$, as well as V22's complete-slice
channel. None holds for arbitrary spatial width at fixed coupling merely by
choosing a larger constant in the final estimate.

The native spatial union estimate \eqref{f16v30:good-union} advances the local
overlap question on a declared outer collar. Its all-test extension
\eqref{f16v30:rough-union} retains a repair form and the larger-union return.
A uniform spatial iteration must control those terms, or establish an
all-boundary local finite-size criterion. The new periodic/cylindrical gap
cannot be inserted as the gap of such an arbitrarily pinned patch: the
conditional parent changes at its boundary. This is the next locality task,
not a new perturbative coefficient or another temporal blocking step.

A simple clock check prevents a different overstatement. For independent
pairs with law $(1+\theta XY)/4$, the original two single-spin rate-one updates
have gap $1-\theta$, while refreshing the entire pair at rate $1/2$ has gap
$1/2$ and unit-or-smaller coverage. The latter does not keep the former away
from zero as $\theta\uparrow1$. This generic fixture does not establish an
obstruction for the native model; it shows why
\eqref{f16v30:elementary-clock} remains a separate obligation even with fixed
block size. Likewise the colored sweep has a declared parallel-layer clock,
not physical Euclidean time.

The supplied memoranda suggested compatible local constraints and a downstream
physical transfer comparison. We now have a quantitative fixed-range native
parent in the stated regime, and a corresponding spatial source norm; neither
is the physical Wilson logarithmic transfer. The retained bridge window,
physical Gauss reduction, infrared contact coercivity, same-top physical
normalization, independent physical-time calibration, and nontrivial interacting
continuum reconstruction remain separate. The V30 checkpoint records the
current dependency map and the limits of the archive review.

\subsection{Diagnostics and predecessor preservation}

The new diagnostic is
\begin{center}\small\path{verify_ym_f16_v30_fixed_block_screening.py}.\end{center}
It tests buffered-cover conditioning, exact restriction energies, noncommuting
conditional projections, fixed-radius coverage, temporal absorption constants,
projection-product loss, exact local square-root kernels, and the separation
between a block and an elementary clock. Finite rational and symbolic fixtures
are not native compact-integral enclosures or proof-assistant verification.
The new diagnostic passed 206 exact assertions. All 26 available V4--V29
diagnostics were rerun unchanged and passed. Their actual reports and hashes
are included separately from the analytic arguments.

Only one literal title version and this terminal insertion differ from V29.
Deleting the insertion and reversing the title change recovers V29 byte for
byte. The compilation and render records refer to the final files. No earlier
proof, source, probability, or reference is silently amended. The present
proofs explicitly depend on the inherited angular preparation and reflecting
product-domain statements; the entire lineage has not been independently
recertified. The next companion checkpoint is V35.

\begin{firewall}
V30 proves a coupling-independent gap for an explicitly normalized fixed-range
native block sampler in a specified exponential spatial-width regime, uniformly
in duration. It proves a controlled-collar buffered-cover and spatial-union
inequality, with an explicit all-test rough-exterior remainder. A standard
projection-product argument then gives a finite-layer contraction, relative
space--time clustering at all radii, and the original source expansions in a
weighted absolute space--time row norm in that regime. It does not prove a
coupling-independent elementary gap, unrestricted spatial width at fixed
coupling, an arbitrary-pinning finite-size criterion, a physical transfer gap,
or an interacting continuum mass gap. All block operations use the same
complete native conditionals and retain their actual rates and normalizers.
\end{firewall}
% END F16 V30 APPEND
% BEGIN F16 V31 APPEND -- 2026-09-18
% Insert immediately before the final document-ending command in V30.
% No predecessor proof, native measure, source, or elementary rate is changed.

\clearpage
\section{V31: the complete retained-bridge action and an in-window conditional gap}
\label{f16v31:context}

V30 proves spatial screening under the complete fixed-bridge fast law, in its
specified width regime. It does not prove that arbitrarily pinned fast patches
have the same gap. Rather than use that unavailable implication, we extract a
different interface that the new screening does support: the \emph{complete}
retained-bridge action, not only the common-source memory submatrix. Its Gaussian
part is retained exactly. The nonlinear correction has a small, exponentially
summable Hessian and an exact finite-volume interaction representation. At a
sufficiently small fixed bridge window this yields a conditional-resampling gap
for the actual interior bridge-path weight, uniformly under further bridge
pinning \emph{within that window}.

Two qualifications are essential. The source space below has all
$36B(n+1)$ real bridge coordinates, rather than only the common-source bank.
Its leading Gaussian action is generally not small. Also, a bridge-window path
probability is not the unrestricted physical path law, and pinning a retained
bridge is not pinning a fast coordinate. Neither the spatial-width restriction
nor the physical-transfer obligations are removed.

\subsection{Dependencies, normalization, and non-duplication}

The full V30 appendix, analytic audit, and checkpoint were read. Its
Theorem~\ref{f16v30:clustering} is used as an inherited analytic premise, not
independently recertified here. V16 supplies native local exponential moments;
V6 supplies the complete-chart ordinary score certificate and the fixed-domain
normalization. The Gaussian precision and source matrices are those of V4--V5.
The new leading-order comparison below does not require an infinite series,
V17's higher-order remainder, or a new native gap assumption.

The supplied f14 V40 already constructs interaction germs by an ordered,
two-radius telescope. Its activation changes which fast variables are active,
and its conclusion concerns a different restricted radial remainder and its
gauge-orbit Hessian. That general telescoping principle is credited. Here all
fast variables stay integrated; only retained bridge inputs outside a declared
carrier are set to zero when defining a coefficient function. The representation
is for the exact complete fast integral on the whole specified bridge window,
not an assertion of new gauge-invariant germs or an infinite-volume potential.
The source and archive searches were targeted, not exhaustive novelty checks.

For classical context, the paper by Barbieri, G\'omez, Marcus, Meyerovitch
and Taati, \href{https://arxiv.org/abs/2001.03880}{\emph{Gibbsian representations
of continuous specifications}}, distinguishes
norm-summable interaction representations from stronger symmetry requirements.
Its abstract and model scope were checked. Its symbolic-configuration
theorem is not invoked for the present continuous variables. All required
finite telescopes and estimates are proved below. The Dobrushin conclusion uses
the already proved finite-product criterion of V24, with its rates explicit.
The next scheduled companion checkpoint remains V35.

Let $\mathcal V=\{(i,e):0\le i\le n,\ e\text{ a bridge}\}$; each $y_u$ has
three real components. There are $12B$ bridges per slice. Assign a bridge to its
source cube; the time plus torus distance of these anchors is denoted $d(u,v)$.
Distinct bridge types can have the same anchor. This bounded multiplicity only
changes fixed constants, and every anchor ball has at most $C(1+L)^4$ labels.
The same convention, with its bounded component multiplicities, is used for
fast scalar indices. All original word carriers have bounded diameter.

For a fixed $r>0$, write
\[
 \mathcal Y_r=\prod_{u\in\mathcal V}\{y_u\in\mathbb R^3:|y_u|\le r\},
 \qquad \varepsilon=\gamma^{-1/2}.
\]
Unless otherwise stated, every result is in V30's regime
\begin{equation}
 \gamma\ge\gamma_r,\qquad B=m^3\le e^{\beta_r\gamma},\quad m\ge4,
 \quad n\ge1,\quad \mathbf y\in\mathcal Y_r.
 \label{f16v31:regime}
\end{equation}
Thresholds may be increased by fixed-$r$ requirements. No estimate below permits
$B\to\infty$ at one fixed coupling outside this regime.

Use the unchanged classical unmoving potential $U_\gamma$ and fast Haar factor
$J^f_\gamma$. Define
\begin{equation}
 \begin{split}
 \mathcal F_\gamma(\mathbf y)
   &=-\log\frac{Z_{\gamma,\mathrm{full}}(\mathbf y)}
                       {Z_{\gamma,\mathrm{full}}(0)},\\
 \mathcal F_0(\mathbf y)&=-\log\frac{Z_0(\mathbf y)}{Z_0(0)},\\
 \mathcal R_\gamma&=\mathcal F_\gamma-\mathcal F_0.
 \end{split}
 \label{f16v31:actions}
\end{equation}
The normalization subtracts one exact value; it does not replace a native
normalizer by a Gaussian one. In the original fast Euclidean coordinates,
write the already defined quadratic form as
\begin{equation}
 \begin{gathered}
 U_0(\zeta;\mathbf y)=\tfrac12\zeta^{\mathsf T}H\zeta
       +\zeta^{\mathsf T}L\mathbf y+\tfrac12\mathbf y^{\mathsf T}B_0\mathbf y,\\
 C=H^{-1},\qquad m=-CL\mathbf y,\qquad K=B_0-L^{\mathsf T}CL.
 \end{gathered}
 \label{f16v31:Gaussian-blocks}
\end{equation}
These definitions are by differentiation of V4's exact quadratic reference,
transported by V5's triangular change of variables. They introduce no fitted
matrix. Gaussian integration gives $\mathcal F_0=\tfrac12\mathbf y^{\mathsf T}K
\mathbf y$; its determinant cancels only in this same-Gaussian value ratio.
$H,L,B_0$ have fixed finite range and bounded absolute rows and columns.
$H$ has a uniform positive lower bound; $C$ has exponentially summable rows.
No positive lower bound on $K$ is assumed.

\subsection{An exact covariance identity with a locally paid chart defect}

Put $V_\gamma=U_\gamma-\log J^f_\gamma$. Its derivative near a principal cut
must not be used in an uncut integration by parts. Choose a smooth radial
function $\chi$ on $\mathbb R^3$ equal to one for $|x|\le1/4$, zero for
$|x|\ge1/2$, and between zero and one elsewhere. For a fast scalar index $j$,
let $g(j)$ be its original three-coordinate group and set
\[
 \chi_j(\zeta)=\chi(\zeta_{g(j)}/\sqrt\gamma),\qquad
 \xi=\zeta-m.
\]
The same cutoff is used for the three components of one group. Define
\begin{equation}
 d_{\gamma,j}
   =\chi_j\partial_jV_\gamma-\partial_j\chi_j-(H\xi)_j,
 \qquad
 \Theta_\gamma=\operatorname{diag}_j\mathbb E_P(1-\chi_j),
 \quad \Sigma_\gamma=\operatorname{Cov}_P(\zeta).
 \label{f16v31:drift-defect}
\end{equation}
The product $\chi_j\partial_jV_\gamma$ is set to zero off its interior support.
Here $P=P_{\gamma,\mathbf y}$ is the same complete native fast law. These
$d_{\gamma,j}$ are auxiliary integration-by-parts defects, not a redefinition
of V7's Haar-corrected bridge scores.

\begin{proposition}[Complete-law Gaussian identities with the chart term retained]
\label{f16v31:Stein-identity}
At every finite regulator in the stated regime,
\begin{equation}
 \begin{split}
 H(\mathbb E_P\zeta-m)&=-\mathbb E_P d_\gamma,\\
 H\Sigma_\gamma&=I-\Theta_\gamma
                     -\operatorname{Cov}_P(d_\gamma,\zeta).
 \end{split}
 \label{f16v31:Stein-exact}
\end{equation}
Each $d_{\gamma,j}$ has a fixed local fast carrier and
\begin{equation}
 \|d_{\gamma,j}\|_{L^2(P)}\le C_r\varepsilon,
 \qquad 0\le(\Theta_\gamma)_{jj}\le C_re^{-c_r\gamma}.
 \label{f16v31:drift-bounds}
\end{equation}
The bounds are uniform in source history, position, volume, and duration in
\eqref{f16v31:regime}.
\end{proposition}
\begin{proof}
Integrate the derivative of $\chi_j f e^{-V_\gamma}$ in its own group.
The cutoff is compactly supported inside that group's principal ball, so its
boundary term is zero. All other groups remain on their full charts. For
$f=1$ or $f=\zeta_k-\mathbb E_P\zeta_k$ this gives
\[
 \mathbb E_P[(\chi_j\partial_jV_\gamma-\partial_j\chi_j)f]
                       =\mathbb E_P[\chi_j\partial_j f].
\]
Substitute \eqref{f16v31:drift-defect}. Centering in the second choice gives
$H\Sigma_\gamma+\operatorname{Cov}_P(d_\gamma,\zeta)
=\operatorname{diag}(\mathbb E_P\chi_j)$. This proves both identities with
their signs. No flat derivative crosses a logarithm cut.

For the estimates, split
\[
 d_{\gamma,j}=\chi_j(\partial_jV_\gamma-\partial_jU_0)
              +(\chi_j-1)\partial_jU_0-\partial_j\chi_j.
\]
The classical part of the first term is bounded by
$C_r\varepsilon(1+|\zeta_T|^2)$ on a fixed local carrier $T$.
Indeed every word is a fixed ordered product of exponentials of fixed linear
forms, with its original $\gamma$ prefactor. The real third-derivative bound
from V6 applies to a fast as well as a bridge derivative. The endpoint
quadratic is exact in $U_0$ and cancels. On the cutoff support the measured
Haar derivative is bounded by $C\gamma^{-1}|\zeta_{g(j)}|$; the endpoint
power $1/2$ only decreases this bound. There is no intermediate principal
logarithm to differentiate in a word. V16's native fourth moments therefore
bound the first term in $L^2$ by $C_r\varepsilon$.

The other terms are supported where $|\zeta_{g(j)}|\ge\sqrt\gamma/4$.
V16's one-cell exponential moment gives exponentially small probability.
The linear polynomial $\partial_jU_0=(H\xi)_j$ has a bounded local $L^4$ norm,
and $|\partial_j\chi_j|\le C\varepsilon$. Cauchy--Schwarz pays their $L^2$
tails without assuming that a rare event is independent of its polynomial.
The same probability bound controls $\Theta_\gamma$. All differentiated
quantities have fixed carriers; the Gaussian mean is a deterministic bounded
function of the prescribed history. Haar factors and their powers have been
retained in the integration identity.
\end{proof}

For a possibly rectangular scalar matrix define a weighted absolute-row norm
using its two anchor families. All row and column bounds asserted below hold
with the same fixed exponential weight, after a decrease if needed. Supremum
over $\mathbf y$ may be placed \emph{inside} each absolute entry: the local
estimates give a common pointwise exponential envelope first.

\begin{theorem}[A summable native mean and covariance comparison]
\label{f16v31:fast-moments}
There is $a>0$ such that
\begin{equation}
 \sup_{\mathbf y\in\mathcal Y_r}\|\mathbb E_P\zeta-m\|_\infty\le C_r\varepsilon,
 \quad
 \sup_j\sum_k e^{a d(j,k)}
       \sup_{\mathbf y\in\mathcal Y_r}|(\Sigma_\gamma-C)_{jk}|
                                      \le C_r\varepsilon.
 \label{f16v31:fast-covariance}
\end{equation}
The same bound holds for weighted columns. This is a consequence of V30's
complete-law screening, not a new assumption that the native law is Gaussian.
\end{theorem}
\begin{proof}
V16 bounds the variance of each $\zeta_k$ uniformly. V30's relative clustering
applies to the fixed local carrier of $d_{\gamma,j}$ and to $\zeta_k$.
Using \eqref{f16v31:drift-bounds}, it gives
\[
 |\operatorname{Cov}_P(d_{\gamma,j},\zeta_k)|
                       \le C_r\varepsilon e^{-a_0d(j,k)}.
\]
Carrier allowances are absorbed into $C_r$. Polynomial graph growth makes
this rectangular covariance summable in both directions at a smaller weight.
From \eqref{f16v31:Stein-exact},
\begin{equation}
 \Sigma_\gamma-C=-C\Theta_\gamma
                    -C\operatorname{Cov}_P(d_\gamma,\zeta).
 \label{f16v31:covariance-resolvent}
\end{equation}
The original $C$ has summable weighted rows and columns. Matrix multiplication
and the triangle inequality for anchor distance prove the covariance bound.
The first line of \eqref{f16v31:Stein-exact} and the absolute rows of $C$
prove the mean bound. Products of the exponential envelopes justify moving
the supremum inside the sums. The inverse used here is only $H^{-1}$; no
inverse of an unproved native generator has been introduced.
\end{proof}

\subsection{The complete bridge Hessian, with its direct term}

For bridge scalar marks $\alpha,\beta$ put
\[
 J_{\gamma,\alpha}=\partial_\alpha U_\gamma,
 \qquad B_{\gamma,\alpha\beta}=\partial_\alpha\partial_\beta U_\gamma.
\]
These are ordinary coordinate derivatives at fixed unmoving fast variables.
$J^f_\gamma$ is independent of the bridge history. Exact differentiation of
the positive compact integral therefore gives
\begin{equation}
 \partial_\alpha\partial_\beta\mathcal F_\gamma
 =\mathbb E_P B_{\gamma,\alpha\beta}
        -\operatorname{Cov}_P(J_{\gamma,\alpha},J_{\gamma,\beta}).
 \label{f16v31:Hessian-exact}
\end{equation}
This is the existing normalized Hessian identity, now used on \emph{all}
bridge directions and at all separations. The direct term cannot be dropped
on equal-time or neighboring source carriers.

For a $3$-by-$3$ bridge-block matrix field $M(\mathbf y)$ define
\begin{equation}
 \mathfrak N_{a,r}(M)=\sup_u\sum_v e^{a d(u,v)}
                    \sup_{\mathbf y\in\mathcal Y_r}\|M_{uv}(\mathbf y)\|_{\rm op}.
 \label{f16v31:Hessian-norm}
\end{equation}

\begin{theorem}[The full native bridge action is a small Hessian perturbation]
\label{f16v31:action-Hessian}
For a fixed $a>0$ and $C_r<\infty$ in \eqref{f16v31:regime},
\begin{equation}
 \boxed{\qquad
 \mathfrak N_{a,r}(D_y^2\mathcal F_\gamma-K)
              =\mathfrak N_{a,r}(D_y^2\mathcal R_\gamma)
              \le C_r\gamma^{-1/2}.
 \qquad}
 \label{f16v31:action-smallness}
\end{equation}
Moreover $\mathcal F_\gamma(0)=\mathcal R_\gamma(0)=0$ and
$D_y\mathcal F_\gamma(0)=D_y\mathcal R_\gamma(0)=0$. In particular,
\begin{equation}
 |\mathcal R_\gamma(\mathbf y)|
       \le C_r\gamma^{-1/2}\sum_u|y_u|^2,
 \qquad
 \sup_u|\partial_{y_u}\mathcal R_\gamma(\mathbf y)|
       \le C_r r\gamma^{-1/2}.
 \label{f16v31:action-size}
\end{equation}
No derivative of an approximate pressure is used. The matrix $K$ includes
all Gaussian direct and returned terms and is not replaced by the common-mode
memory coefficient $G_2$.
\end{theorem}
\begin{proof}
The unchanged V6 certificate gives
$J_\gamma=L^{\mathsf T}\zeta+B_0\mathbf y+\delta J_\gamma$, where each
$\delta J_{\gamma,\alpha}$ has fixed carrier and $L^2(P)$ norm at most
$C_r\varepsilon$. Its variance obeys the same bound. All rows and columns of
$L$ have bounded finite support. Expand the covariance in
\eqref{f16v31:Hessian-exact} using this identity. The linear--linear part
differs from $L^{\mathsf T}CL$ by
$L^{\mathsf T}(\Sigma_\gamma-C)L$, bounded by
Theorem~\ref{f16v31:fast-moments}. Each linear--remainder covariance has a
pointwise bound $C_r\varepsilon e^{-a_0d(\alpha,\beta)}$ by V30's clustering.
The remainder--remainder term has the stronger $C_r\varepsilon^2$ prefactor.
Their weighted rows and columns are summable uniformly over $\mathcal Y_r$.

If no classical word contains both marks, $B_{\gamma,\alpha\beta}$ and the
corresponding entry of $B_0$ vanish. Otherwise the same global real
third-derivative certificate yields
$|B_{\gamma,\alpha\beta}-(B_0)_{\alpha\beta}|
\le C_r\varepsilon(1+|\zeta_T|)$ on a fixed carrier. Native moments bound
its expectation; only a fixed number of such entries occur per row.
Consequently the entire Hessian is
$B_0-L^{\mathsf T}CL+O(\varepsilon)$ in the stated norm. Conversion from
scalar colors to $3$-by-$3$ operator blocks costs only a fixed factor.

Global simultaneous conjugation of every group preserves the exact fast
integral, its complete charts, Haar powers, Wilson words and color-isotropic
endpoint amplitudes. It sends every $y_u$ through the same adjoint rotation.
At $\mathbf y=0$, each of the separate three-component gradient vectors must
therefore be fixed by every rotation in $\mathrm{SO}(3)$, and hence is zero.
The Gaussian action has the same property. This argument asserts only global
conjugation, not gauge invariance of the bounded bridge window under independent
root transformations. Integrating the actual Hessian of $\mathcal R_\gamma$
along $t\mathbf y$, $0\le t\le1$, gives \eqref{f16v31:action-size}; symmetry
and the absolute-row bound control its quadratic form. The interpolation stays
inside the convex product $\mathcal Y_r$.
\end{proof}

\subsection{An exact, exponentially summable retained-input interaction}

Order the bridge sites once and for all. Let $P_u$ be the predecessors of $u$,
$B_l(u)=\{v:d(u,v)\le l\}$, and $A_{u,l}=P_u\cap B_l(u)$ for integers
$l\ge0$. For $S\subseteq\mathcal V$, let $\mathbf y^S$ keep the inputs in $S$
and set all other bridge inputs to zero. Define
\begin{equation}
 \begin{split}
 D_{u,l}(\mathbf y)&=
 \mathcal R_\gamma(\mathbf y^{A_{u,l}\cup\{u\}})
                -\mathcal R_\gamma(\mathbf y^{A_{u,l}}),\\
 \Phi_{u,0}&=D_{u,0},\qquad
 \Phi_{u,l}=D_{u,l}-D_{u,l-1}\quad(l\ge1).
 \end{split}
 \label{f16v31:interaction-definition}
\end{equation}
These functions use complete native fast integrals at the displayed histories.
Fast variables outside $B_l(u)$ have \emph{not} been frozen or removed.

\begin{theorem}[A normalized action, not only an entrywise response expansion]
\label{f16v31:interaction}
The finite-radius sum terminates and gives the exact identity
\begin{equation}
 \boxed{\quad
 \mathcal F_\gamma(\mathbf y)
       =\tfrac12\mathbf y^{\mathsf T}K\mathbf y
                               +\sum_{u,l\ge0}\Phi_{u,l}(\mathbf y).
 \quad}
 \label{f16v31:interaction-exact}
\end{equation}
Each term depends only on bridge inputs in $B_l(u)$, vanishes at zero, and for
some fixed $b>0$ obeys
\begin{equation}
 \boxed{\quad
 \sup_v\sum_{u,l:\,v\in B_l(u)}
       |B_l(u)|e^{b\operatorname{diam}B_l(u)}
                  \|\Phi_{u,l}\|_{\infty,\mathcal Y_r}
                         \le C_r\gamma^{-1/2}.
 \quad}
 \label{f16v31:interaction-norm}
\end{equation}
A tail with $l>L$ has the stronger bound $C_r\gamma^{-1/2}e^{-bL}$ after a
possible fixed decrease of $b$. Constants are uniform in the regime, but the
coefficient functions can depend on the ambient regulator and the chosen
ordering. No infinite-volume consistency or independent-root gauge symmetry
of each term is claimed.
\end{theorem}
\begin{proof}
Once $l$ reaches the diameter of the finite anchor graph, every predecessor
has entered. Telescoping over $l$, and then over the site ordering, yields
$\sum_{u,l}\Phi_{u,l}=\mathcal R_\gamma(\mathbf y)-\mathcal R_\gamma(0)$.
The value at zero is exactly zero. Dependence on the retained inputs follows
from their definition; no spatial factorization of a native integral is used.

For $l\ge1$, add the sites in $A_{u,l}\setminus A_{u,l-1}$ one at a time.
The change in the $u$ increment when a site $v$ is added is a mixed finite
difference. Its exact double-integral formula is bounded by
\[
 |y_u|\,|y_v|\sup_{\mathbf z\in\mathcal Y_r}
              \|\partial_{y_u}\partial_{y_v}\mathcal R_\gamma(\mathbf z)\|
                 \le C_r r^2\varepsilon e^{-a d(u,v)}.
\]
Intermediate backgrounds differ with $v$; the supremum is already inside the
Hessian envelope in \eqref{f16v31:Hessian-norm}. The shell contains at most
$C(1+l)^4$ sites. Decreasing the rate gives
$\|\Phi_{u,l}\|_\infty\le C_r\varepsilon e^{-a_1l}$.
At $l=0$ there are only boundedly many co-anchored bridge types. The gradient
at zero vanishes, so the Hessian bound on this finite carrier gives
$\|\Phi_{u,0}\|_\infty\le C_r\varepsilon$ as well.

For fixed $v$ and $l$, there are at most $C(1+l)^4$ possible anchors $u$ with
$v\in B_l(u)$; also $|B_l(u)|\le C(1+l)^4$ and its diameter is at most $2l$.
Choosing $2b<a_1$ proves the norm bound by a polynomial times an exponential
series. Leave a strict exponential margin to obtain the tail assertion.
These counts remain valid on saturated tori, short histories, and at both
endpoints. They do not require a translation-invariant choice of ordering.
\end{proof}

The Gaussian quadratic term itself has exponentially summable two-site
coefficients, since $K=B_0-L^{\mathsf T}CL$. Thus the theorem supplies an
interaction representation of the \emph{whole} normalized fast action, with a
small correction to a generally order-one Gaussian interaction. The individual
$\Phi$ are smooth on the admitted window, but \eqref{f16v31:interaction-norm}
is a supremum-norm bound, not an unproved derivative norm for each term.
At fixed coupling the width restriction still prevents an unrestricted spatial
limit. A family of finite-volume local carrier functions is not automatically
one consistent infinite-volume Gibbs interaction.

\subsection{The endpoint factors and the exact interior bridge-window law}

V4's complete history weight has not changed. Its exact logarithmic ratio is
\begin{equation}
 \begin{split}
 -\log\frac{\mathcal W_{\gamma,n}^{\mathrm{full}}(\mathbf y)}
                  {\mathcal W_{\gamma,n}^{\mathrm{full}}(0)}
   ={}&\mathcal F_\gamma(\mathbf y)-\log J_{{\rm br},\gamma}(\mathbf y)\\
     &+\tfrac12\log\frac{F_{\gamma,2}(Z_0)}{F_{\gamma,2}(I)}
      +\tfrac12\log\frac{F_{\gamma,2}(Z_n)}{F_{\gamma,2}(I)}.
 \end{split}
 \label{f16v31:complete-weight}
\end{equation}
The last two functions are retained exactly. This appendix does not prove an
exponentially summable spatial representation for those endpoint normalizers.
No assertion about the full endpoint-varying Hessian is obtained by dropping
them. $J_{{\rm br},\gamma}(0)=1$ in the inherited convention.

Now fix both bridge endpoints, in the stated window, and let $n\ge2$.
Normalize the original positive path integrand over
$|y_{i,e}|\le r$, $1\le i\le n-1$. The resulting probability is exactly
\begin{equation}
 d\nu_{\gamma,r}^{\mathrm{win}}(\mathbf y_{\rm int})
  =\frac1{\mathcal Z_{\gamma,r}^{\mathrm{win}}}
    e^{-\mathcal F_\gamma(\mathbf y)}
       \prod_{u\,\mathrm{interior}}
       1_{\{|y_u|\le r\}}j(y_u/\sqrt\gamma)\,dy_u.
 \label{f16v31:bridge-window-law}
\end{equation}
The two dressed-end factors, their bridge half-Haar factors and the scalar
$c_{\gamma,B,n}$ cancel here because \emph{both endpoints are fixed}. The
remaining bridge Haar factors stay in their original product reference.
This is a window-conditioned, endpoint-fixed version of the original path
weight, not the unrestricted physical path probability. A further conditioning
of any subset of interior bridge sites also remains an exact kernel of this
same window law.

\begin{proposition}[Finite-range correction with normalized conditional error]
\label{f16v31:truncation}
Replace only $\mathcal R_\gamma$ in \eqref{f16v31:bridge-window-law} by
$\sum_{u,l\le L}\Phi_{u,l}$ and keep $K$, the endpoints and the Haar reference
unchanged. For every one-site conditional, at any admissible exterior,
\begin{equation}
 d_{\rm TV}(\nu_{\gamma,r}^{\rm win}(dy_v\mid\mathbf y_{v^c}),
                    \nu_{\gamma,r}^{[L]}(dy_v\mid\mathbf y_{v^c}))
                         \le C_r\gamma^{-1/2}e^{-bL}.
 \label{f16v31:conditional-truncation}
\end{equation}
There is no corresponding ambient-size-independent assertion for total variation
of the two \emph{joint} history laws.
\end{proposition}
\begin{proof}
Only omitted terms whose declared carrier contains $v$ can vary with $y_v$.
Their total oscillation is at most twice their summed supremum norms, bounded
by Theorem~\ref{f16v31:interaction}. For two normalized densities with the
same reference, a potential-difference oscillation $\delta$ implies TV at most
$\tanh(\delta/4)\le\delta/4$, as proved immediately below. This includes both
one-site normalizers and proves the bound. Summing a small per-site error over
an arbitrarily large joint carrier would be a different estimate.
\end{proof}

\subsection{A checkable influence margin, and a small fixed-window consequence}

\begin{proposition}[Oscillation, normalization, and a one-bridge influence bound]
\label{f16v31:TV-influence}
Two positive densities proportional to $e^{-V}$ and $e^{-W}$ with a common
reference have TV distance at most $\tanh(\operatorname{osc}(V-W)/4)$.
For the exact window law, set
\[
 \mathcal A_{uv}=\sup_{\mathbf y\in\mathcal Y_r}
                  \|\partial_{y_u}\partial_{y_v}\mathcal F_\gamma(\mathbf y)\|,
 \qquad u\ne v.
\]
Its influence when only exterior bridge $v$ changes obeys
\begin{equation}
 c_{uv}\le\tanh(r^2\mathcal A_{uv})\le r^2\mathcal A_{uv}.
 \label{f16v31:influence}
\end{equation}
The bound holds for all in-window exterior values and after arbitrary further
pinning of retained bridge sites within that window.
\end{proposition}
\begin{proof}
Let $R=d\mu/d\nu$, with $\nu$ the second normalized density. Its extrema
$m\le R\le M$ satisfy $m\le1\le M$ and $M/m\le e^\delta$, where
$\delta=\operatorname{osc}(V-W)$. The chord bound on $(R-1)_+$ gives
$d_{\rm TV}(\mu,\nu)\le(M-1)(1-m)/(M-m)$. Maximizing this expression subject
to the ratio bound gives $(e^{\delta/2}-1)/(e^{\delta/2}+1)$; the degenerate
case is zero. This proves the claim without ignoring either normalizer.

Changing $y_v$ between two values in its radius-$r$ ball changes the one-site
potential at $u$ by a function whose oscillation in $y_u$ is at most
$4r^2\mathcal A_{uv}$. Integrate the actual mixed derivative along the two
straight segments; both stay in $\mathcal Y_r$. All terms independent of
$y_u$, including the fixed endpoint factors, cancel from that conditional.
The native one-site Haar reference is the same in the two rows. Apply the
first part. Additional bridge pinning only removes free coordinates; it does
not change the uniform derivative estimate. No fast coordinate is pinned by
this argument.
\end{proof}

Write
\[
 L_0=\sup_u\sum_{v\ne u}\|K_{uv}\|<\infty,\qquad
 q_0(r)=\sup_u\sum_{v\ne u}\tanh(r^2\|K_{uv}\|).
\]
The constants refer to the actual complete Gaussian Schur matrix. Suprema may
be taken over all finite histories; its local Gaussian bounds make $L_0$ finite.
By the symmetric Hessian envelope, row and column influences are both bounded
by $q_0(r)+C_r r^2\gamma^{-1/2}$. Small correction alone does not say that
$q_0(r)<1$ at an arbitrary $r$.

\begin{theorem}[An exact retained-path conditional gap inside a certified window]
\label{f16v31:bridge-gap}
Whenever
\begin{equation}
 q_0(r)+C_r r^2\gamma^{-1/2}\le q<1,
 \label{f16v31:margin}
\end{equation}
the sum of rate-one single-bridge conditional innovations for
$\nu_{\gamma,r}^{\mathrm{win}}$, and for any of its further in-window bridge
conditionals, has gap at least $1-q$.

In particular, fix the V30 constants for the ambient bound $r=1$, and put
\begin{equation}
 r_* =\min\left\{1,\,[4(1+L_0)]^{-1/2}\right\}>0.
 \label{f16v31:small-window}
\end{equation}
At a fixed sufficiently large coupling in that same width regime, every
$0<r\le r_*$, every pair of fixed endpoints in $\mathcal Y_r$, and every
further retained-bridge pinning in the window satisfy
\begin{equation}
 \boxed{\quad
 \operatorname{Var}_{\nu} f
       \le2\sum_{u\,\mathrm{free}}
              \mathbb E_{\nu}\operatorname{Var}_{\nu}(f\mid y_{u^c}).
 \quad}
 \label{f16v31:window-gap}
\end{equation}
The constant has no number-of-bridges or duration factor. This is a new auxiliary
clock on retained bridge paths, not the fast block clock, the balanced elementary
clock, or physical Euclidean time.
\end{theorem}
\begin{proof}
The matrix in \eqref{f16v31:influence} has a symmetric nonnegative dominating
table, since the mixed Hessian blocks are transposes. The function $\tanh$ is
one-Lipschitz on the nonnegative axis. The Hessian remainder bound therefore
gives the stated row and column margin. Apply V24's finite-product
boundary-sensitive contraction theorem with one bridge site per update, unit
rates and unit weights. The death rate is one and the birth column sum is at
most $q$. Every intermediate exterior in a sequence of one-site changes is
still in the product window. The same conclusion holds on any conditional
subproduct because its dominating table is a principal submatrix. Empty free
sets give the vacuous variance identity.

For \eqref{f16v31:small-window}, use $q_0(r)\le r^2L_0\le1/4$. The constants
for $r=1$ dominate all smaller windows. Increase only a fixed threshold so that
$C_1r_*^2\gamma^{-1/2}\le1/4$. Then $q\le1/2$ uniformly as claimed. No small
Gaussian eigenvalue is inverted and no diagonal direct term is discarded.
The positivity of the native Haar reference on the small bridge balls supplies
explicit conditional-density versions at every in-window parameter. The
physical probability outside this retained window is not estimated here.
\end{proof}

\subsection{What remains beyond the retained window}

The exact interaction keeps every fast field integrated and the complete
Gaussian bridge action. The all-pinning theorem concerns retained bridges
inside a small fixed window, not the arbitrary compact fast patches left
open by V30. Nothing proves that the unrestricted retained path typically
stays in this window. A small nonlinear correction alone also does not make
the leading Gaussian influence sum less than one on a larger window.

The next distinct tasks are to prove or bypass that influence margin on a
larger admissible domain, obtain gauge-invariant coverage, and control the
endpoint functions retained in \eqref{f16v31:complete-weight}. Compatibility
of the carrier functions across ambient regulators is not established.
Neither an unrestricted spatial limit nor a physical transfer or continuum
mass statement follows from this finite-window auxiliary clock.

\subsection{Verification and preservation}

The script
\begin{center}\small\path{verify_ym_f16_v31_effective_bridge_action.py}\end{center}
passes 1,040 exact assertions: compact integration by parts with a nonzero
cutoff defect, normalized Schur/Hessian identities, many-body retained-input
telescopes, conditional density ratios, all-pinning finite probability fixtures,
and weighted-norm bookkeeping. These are algebraic tests, not native Wilson
integral enclosures or formal verification. Execution reports state the scope.

Only V30's literal title version and this terminal insertion are changed;
reversing them recovers V30 byte for byte. No inherited proof, endpoint factor,
source, or reference is silently edited. The analytic deductions use V16's
moments and V30's screening together with the identified Gaussian and ordered-word
certificates. The inherited lineage is not independently recertified. The next
checkpoint remains V35.

\begin{firewall}
V31 proves the complete bridge-Hessian comparison, an exact summable interaction
correction, and normalized one-site truncation bounds. Its retained-path gap
requires a checked influence margin; a sufficiently small fixed window satisfies
it, including further bridge pinning within that window. Width and window
restrictions remain. Endpoint normalizers cancel only in the endpoint-fixed
conditional. No arbitrary-fast-pinning theorem, physical-time identification,
or Yang--Mills mass gap is inferred.
\end{firewall}
% END F16 V31 APPEND

% BEGIN F16 V32 APPEND -- 2026-09-18
% Insert immediately before the final document-ending command in V31.
% No predecessor assertion, native probability, endpoint, source, or clock is edited.

\clearpage
\section{V32: boundary coercivity beyond the small bridge window}
\label{f16v32:context}

V31 bounds the complete retained-action Hessian, but its conditional gap uses
an absolute influence sum that is small only on a sufficiently small window.
We do not try to improve that sum term by term. The leading Schur action is
nonnegative, although it need not be strictly positive. We retain this signed
matrix information and use the boundary of the already specified retained
window to control its soft directions. The negative curvature of the native
remainder is then a perturbation of a \emph{boundary} coercivity estimate,
not of a nonexistent positive lower eigenvalue of the Gaussian Schur matrix.

For every prescribed finite bridge radius, this proves a positive rate-one
retained heat-bath gap, uniform in duration and in the inherited spatial-width
regime. It survives every further retained-bridge pinning in that window.
The radius is fixed before coupling is taken large; the constants deteriorate
with it. All fast coordinates remain completely integrated. No estimate of
unrestricted bridge coverage, arbitrary fast pinning, or physical mass follows.

\subsection{Dependencies, ownership, and the boundary being used}

The complete V31 appendix and audit were read. We use its full Hessian
estimate and endpoint-fixed law. V29's partial Bochner/heat-bath conversion
is reused, but the needed coercivity is proved here from the boundary rather
than assumed from a positive Hessian. Two unsuccessful semantic searches were
followed by mounted-text checks and identified f14 V64 and f15 V70--V71
passages. Their different marginal and observable-oscillation results are not
substituted for the present theorem. This is a targeted review, not an
exhaustive novelty certificate. The next checkpoint remains V35.

For context, Esposito--Nitsch--Trombetti,
\href{https://arxiv.org/abs/1110.2960}{\emph{Best constants in Poincar\'e
inequalities for convex domains}}, records the sharp Neumann diameter bound;
we reviewed its statement and prove a weaker elementary ball estimate below.
Ascolani--Lavenant--Zanella,
\href{https://arxiv.org/abs/2410.00858v2}{\emph{Entropy contraction of the Gibbs
sampler under log-concavity}}, Corollary 3.7 and Section 6, distinguishes its
strongly convex and merely convex settings. Those statements and assumptions
were checked, not the full proof. No whole-space entropy theorem is imported
across our hard boundary, and no entropy tensorization is claimed.

Retain the V31 law with both bridge endpoints fixed. Let $I$ be any subset of
its interior bridge sites, and pin the other retained sites at values $q$ in
the same radius-$r$ balls. Write $y=(x_I,q)$ and
\begin{equation}
 d\nu_{\gamma,r}^{I,q}(x)
 =\frac{e^{-V_{\gamma}^{I,q}(x)}}{Z_{\gamma,r}^{I,q}}
                  \prod_{u\in I}1_{\{|x_u|\le r\}}\,dx_u,\qquad
 V_{\gamma}^{I,q}(x)=\mathcal F_\gamma(x,q)
                   -\sum_{u\in I}\log j(x_u/\sqrt\gamma).
 \label{f16v32:law}
\end{equation}
Each free site has three real coordinates. Constants depending only on $q$
cancel from this conditional, including the two dressed endpoint functions
because both endpoints remain fixed. The interior Haar density in
\eqref{f16v32:law} is not deleted: it is included in the measured potential.
This is exactly V31's probability, not a convex extension or a replacement
Gaussian law. An empty $I$ has only the vacuous variance identity.

\subsection{The Schur action has no negative directions; it need not have a mass}

Let $K,H,L,B_0$ remain the matrices in \eqref{f16v31:Gaussian-blocks}. Put
\[
 L_K=\sup_u\sum_v\|K_{uv}\|_{\rm op}<\infty,
 \qquad
 \delta_{\gamma,r}=C_r\gamma^{-1/2},
\]
where $C_r$ is a valid constant in \eqref{f16v31:action-smallness}, enlarged
for its unweighted block norm if necessary. $L_K$ is a uniform bound over
the finite Gaussian histories, including diagonal blocks. It is not the
small-window off-diagonal bound $L_0$.

\begin{proposition}[Signed curvature and one-site oscillation budgets]
\label{f16v32:budgets}
The original Schur matrix satisfies $K\succeq0$, without an asserted positive
lower edge. There is a fixed finite Haar constant $C_H$ such that, whenever
\[
 \gamma\ge\max\{4r^2,C_H,\gamma_r\},\qquad
 B\le e^{\beta_r\gamma},\qquad \delta_{\gamma,r}\le1,
\]
every conditional \eqref{f16v32:law} satisfies
\begin{equation}
 \begin{gathered}
 D^2V_{\gamma}^{I,q}\succeq-\delta_{\gamma,r}I,
 \qquad (D^2V_{\gamma}^{I,q})_{uu}\preceq M_*I_3,
 \qquad M_*=L_K+2,\\
 \sup_{x_{u^c}}\operatorname{osc}_{|x_u|\le r}
                  V_{\gamma}^{I,q}(x)
       \le\Omega_r,\qquad \Omega_r=(2L_K+3)r^2.
 \label{f16v32:budgets-bound}
 \end{gathered}
\end{equation}
All bounds include arbitrary values of the retained pins within the radius-$r$
window. They contain no number of free sites.
\end{proposition}
\begin{proof}
The original quadratic $U_0$ is a sum of nonnegative quadratic limits of
Wilson actions and the nonnegative endpoint quadratics. Completing its fast
square gives the exact identity
\[
 U_0(\zeta;y)=\tfrac12(\zeta+H^{-1}Ly)^{\mathsf T}
                  H(\zeta+H^{-1}Ly)+\tfrac12y^{\mathsf T}Ky.
\]
Since $H\succ0$, minimize in $\zeta$ to obtain $K\succeq0$. No determinant
or scalar normalization supplies this conclusion; it follows from the
nonnegative quadratic action. Principal restrictions preserve it.

For $t=|x|/\sqrt\gamma<\pi$, the inherited Haar factor is
$j(x/\sqrt\gamma)=(\sin t/t)^2$. Set $\phi(t)=2\log(t/\sin t)$. Its radial
Hessian eigenvalues before scaling are
\[
 \phi''(t)=2(\csc^2t-t^{-2})>0,\qquad
 \frac{\phi'(t)}t=2(t^{-2}-t^{-1}\cot t)>0.
\]
Both extend continuously at zero, with value $2/3$. The first sign follows
from $\sin t<t$; the second follows from
$\sin t-t\cos t>0$ on $(0,\pi)$. Their maximum on $[0,1/2]$ defines a
finite $C_H$. Hence the measured one-site Haar Hessian is between zero
and $C_H/\gamma$, and its gradient vanishes at zero.

V31 gives $\|D^2\mathcal R_\gamma\|_{\rm op}\le\delta_{\gamma,r}$ from its
symmetric block-row bound. Add the nonnegative Haar Hessian and restrict to
$I$ for the first two estimates. Its proof also gives
$|\partial_{y_u}\mathcal F_\gamma(y)|\le r(L_K+\delta_{\gamma,r})$
throughout the full product window: integrate the Hessian along $ty$ from its
zero gradient at zero. A change of one bridge has length at most $2r$,
so the action oscillation at that site is at most $2r^2(L_K+1)$. The Haar
potential ranges from zero to at most $C_Hr^2/(2\gamma)\le r^2/2$.
The stated $\Omega_r$ is a conservative sum. Fixed retained pins do not
invalidate the estimate along the original full-window segment. No fast
variable is pinned in this step.
\end{proof}

\subsection{A one-ball normal-trace estimate, with a local density comparison}

The next calculation supplies the lower edge missing from a semidefinite $K$.
It concerns vector fields, not just a scalar Poincar\'e inequality.

\begin{proposition}[Boundary exclusion of constant vector fields]
\label{f16v32:ball}
Let $D=B_r\subset\mathbb R^d$, $1\le d\le3$, and let $v\in H^1(D;\mathbb R^d)$
have normal trace $v\cdot n=0$ on $\partial D$. Then
\begin{equation}
 \int_D|v|^2\,dx\le19r^2\int_D|Dv|_{\rm HS}^2\,dx.
 \label{f16v32:ball-vector}
\end{equation}
For a positive continuous weight $w$ on $\overline D$ with
$\sup w/\inf w\le e^\Omega$, the same normal-trace class obeys
\begin{equation}
 \int_D|Dv|^2w\,dx
 \ge \frac{e^{-\Omega}}{19r^2}\int_D|v|^2w\,dx.
 \label{f16v32:weighted-ball}
\end{equation}
The ratio is local to this one ball. It is not a density ratio on a growing
product of balls.
\end{proposition}
\begin{proof}
For a scalar $a\in H^1(D)$, the independent uniform-copy variance formula,
the line-segment fundamental theorem, and $|x-y|\le2r$ give
\[
 \int_D|a-\overline a|^2\,dx
          \le2^{d+1}r^2\int_D|\nabla a|^2\,dx.
\]
Here is the dimension constant explicitly. In the segment parameter $t\le1/2$
change $x$ to $(1-t)x+ty$ at fixed $y$; its Jacobian cost is at most $2^d$.
For $t\ge1/2$ change $y$ at fixed $x$. Convexity keeps both image domains in
$D$. The squared diameter and the factor $1/2$ in the variance formula give
$2^{d+1}r^2$. Approximation gives the $H^1$ statement.

Apply this to each component of $v$. Normal trace zero and the divergence
theorem give $\int_Dv=-\int_Dx\operatorname{div}v$. Therefore
\[
 |D|\,|\overline v|^2\le r^2\int_D|\operatorname{div}v|^2
                         \le dr^2\int_D|Dv|^2.
\]
Adding mean and centered parts yields the constant $2^{d+1}+d\le19$.
The trace/divergence identity holds by the $H^1$ trace theorem and density.
Finally bound the weighted derivative energy below by $\inf w$ times the
unweighted energy and the weighted norm above by $\sup w$ times its
unweighted value. Normalization of $w$ cancels. This proves
\eqref{f16v32:weighted-ball}.
\end{proof}

\subsection{A reflecting estimate for nearly convex product-window laws}

Consider $\mu=Z^{-1}e^{-V}dx$ on $\mathcal D=\prod_{u=1}^N B_r^{d_u}$,
$1\le d_u\le3$, with $V$ smooth on a neighborhood of its closure. Suppose
\begin{equation}
 D^2V\succeq-\delta I,\qquad (D^2V)_{uu}\preceq MI,
 \qquad \sup_{x_{u^c}}\operatorname{osc}_{x_u}V\le\Omega,
 \quad M\ge0.
 \label{f16v32:abstract-hypotheses}
\end{equation}
Set $\kappa=e^{-\Omega}/(19r^2)$ and assume $0\le\delta<\kappa$.
The scalar reflecting generator is $\mathcal A=-\Delta+\nabla V\cdot\nabla$.
For a one-form $v=(v_u)$ its absolute form domain consists of $H^1$ vector
fields with $v_u\cdot n_u=0$ on each face belonging to ball $u$. The form is
\begin{equation}
 \begin{split}
 \mathfrak a_1(v,v)={}&\int_{\mathcal D}\bigl(|Dv|^2+
                                      v^{\mathsf T}D^2Vv\bigr)d\mu\\
 &+\sum_u\int_{\partial_u\mathcal D}
            \mathrm{II}_u(v_{u,\mathrm{tan}},v_{u,\mathrm{tan}})\,d\sigma_\mu .
 \end{split}
 \label{f16v32:oneform}
\end{equation}
Here $d\sigma_\mu=Z^{-1}e^{-V}d\sigma$ and the ball second fundamental form is
$\mathrm{II}_u=r^{-1}I$ on its own tangential directions (zero in dimension
one). Cross-cell tangent directions have zero face curvature. All boundary
terms are retained.

\begin{theorem}[Boundary coercivity and whole-$L^2$ variance tensorization]
\label{f16v32:abstract}
Under \eqref{f16v32:abstract-hypotheses}, the positive one-form realization
$\mathcal A_1$ and scalar reflecting law satisfy, with the variance inequality for $f\in H^1(\mu)$,
\begin{equation}
 \mathcal A_1\succeq(\kappa-\delta)I,\qquad
 \operatorname{Var}_\mu f\le\frac1{\kappa-\delta}
                                  \int|\nabla f|^2d\mu .
 \label{f16v32:diffusion}
\end{equation}
For the exact rate-one block conditional expectations $E_u$, with one ball
per update, every $f\in L^2(\mu)$ satisfies
\begin{equation}
 \boxed{\quad
 \operatorname{Var}_\mu f
 \le\frac{\kappa+M}{\kappa-\delta}
                   \sum_u\|(I-E_u)f\|_2^2.
 \quad}
 \label{f16v32:heatbath}
\end{equation}
No strictly positive Hessian lower bound, diagonal dominance, or division of
the update rates by $N$ is required.
\end{theorem}
\begin{proof}
Fix all variables other than $x_u$. The normal trace of $v_u$ on this ball
is zero, and its actual conditional density has ratio at most $e^\Omega$.
Equation \eqref{f16v32:weighted-ball}, followed by disintegration, gives
\[
 \int|D_uv_u|^2d\mu\ge\kappa\int|v_u|^2d\mu.
\]
Sum in $u$. The full $|Dv|^2$ contains these terms, the face terms in
\eqref{f16v32:oneform} are nonnegative, and the curvature term is bounded
below by $-\delta\|v\|^2$. This proves the one-form floor with no global
density comparison.

We specify the domain passage. At each fixed $N$, the positive smooth weight
is bounded above and below on the compact product; these particular bounds
may depend on $N$ and are used only to identify Sobolev domains. The scalar
Neumann realization is the closed form $\int|\nabla u|^2d\mu$. Finite tensor
sums of smooth Neumann functions on the individual balls form a graph core
for the unweighted product Laplacian. The smooth bounded drift is relatively
bounded with arbitrarily small relative bound, by interpolation between its
gradient and Laplacian norms, so the same graph-domain approximation applies
to $\mathcal A$. On these functions, componentwise integration by parts gives
\begin{equation}
 \|\mathcal Au\|_2^2=\mathfrak a_1(du,du).
 \label{f16v32:Bochner}
\end{equation}
There is no codimension-two corner integral: the product identity is a sum
over its ball faces. The face terms follow by tangential differentiation of
$\partial_nu=0$. In particular $du$ has exactly the normal trace required
above. The closed forms extend the identity and its inequality by graph
approximation. Equivalently, the associated absolute Witten realization
satisfies $d\mathcal A=\mathcal A_1d$ on its natural domain; the integration
identity against smooth normal-zero vector fields proves this before closure.

Thus $\|\mathcal Au\|^2\ge(\kappa-\delta)\langle u,\mathcal Au\rangle$.
The compact reflecting scalar resolvent has a complete discrete spectrum.
Its kernel is constants because the product is connected and its Dirichlet
form vanishes only on constants. Applying the last inequality to nonconstant
eigenfunctions gives its gap at least $\kappa-\delta$. This proves
\eqref{f16v32:diffusion}. The dimension-dependent compact-domain facts used
in this passage introduce no dimension-dependent constant into the inequality.

For the heat bath, write $\mathcal A=\sum_i\mathcal A_i$. Full and
blockwise Bochner identities have the same nonnegative boundary sum. The
sum of block-diagonal squared Hessians is no greater than the full squared
Hessian. The difference of their potential contributions is at most
$(M+\delta)\int|\nabla u|^2d\mu$. Consequently
\begin{equation}
 \sum_i\|\mathcal A_i u\|_2^2
 \le\|\mathcal Au\|_2^2+(M+\delta)\langle u,\mathcal Au\rangle.
 \label{f16v32:partial}
\end{equation}
The same graph approximation extends this estimate and the exact conditional
identity $E_i\mathcal A_i u=0$. For centered $f\in L^2(\mu)$ solve
$u=\mathcal A^{-1}f$. Then
\[
 \|f\|_2^2=\sum_i\langle(I-E_i)f,\mathcal A_i u\rangle,
 \qquad \langle u,\mathcal Au\rangle
                      \le(\kappa-\delta)^{-1}\|f\|_2^2.
\]
Cauchy--Schwarz and \eqref{f16v32:partial} give the factor
$1+(M+\delta)/(\kappa-\delta)=(\kappa+M)/(\kappa-\delta)$.
The calculation is in the complete $L^2$ space; the test $f$ is not required
to satisfy any reflecting boundary condition. Real and imaginary parts cover
complex tests. This is the V29 conversion with a newly established boundary
floor, not a reversal of a one-site Poincar\'e inequality.
\end{proof}

\subsection{Every fixed retained radius, including all in-window pinning}

Define the explicit, conservative boundary constant
\begin{equation}
 \kappa_r=\frac{\exp[-(2L_K+3)r^2]}{19r^2},\qquad
 g_r=\frac{\kappa_r}{2(\kappa_r+L_K+2)}>0.
 \label{f16v32:constants}
\end{equation}
They do not contain a lower eigenvalue of $K$.

\begin{theorem}[Retained heat-bath gap at every prescribed finite window]
\label{f16v32:main}
For every fixed $0<r<\infty$ there are $\beta_r>0$ and a finite
$\Gamma_r$ such that, for
\begin{equation}
 \gamma\ge\Gamma_r,\quad B=m^3\le e^{\beta_r\gamma},\quad m\ge4,\quad n\ge2,
 \label{f16v32:regime}
\end{equation}
every endpoint-fixed law \eqref{f16v32:law}, and every further retained pinning
within its radius-$r$ window, obeys
\begin{equation}
 \boxed{\quad
 \operatorname{Var}_{\nu}f\le g_r^{-1}
             \sum_{u\in I}\mathbb E_\nu\operatorname{Var}_\nu(f\mid x_{u^c}),
 \qquad f\in L^2(\nu).
 \quad}
 \label{f16v32:main-gap}
\end{equation}
The reflecting bridge diffusion has gap at least $\kappa_r/2$. Every
single-bridge conditional update has rate one. There is no source-count,
spatial-volume, or duration multiplier. In contrast to V31, no condition
$q_0(r)<1$ or $r\le r_*$ is imposed.
\end{theorem}
\begin{proof}
Use the inherited V31 constants for the actual chosen radius $r$, not its
constants for radius one at a larger domain. Enlarge the threshold so that
\[
 \gamma\ge\max\{4r^2,C_H,\gamma_r\},\qquad
 \delta_{\gamma,r}\le\min\{1,\kappa_r/2\}.
\]
For example, the further requirement
$\gamma\ge[2C_r/\min(1,\kappa_r)]^2$ suffices. No numerical size for this
threshold is asserted. Proposition~\ref{f16v32:budgets} and
Theorem~\ref{f16v32:abstract}, with $M=M_*$, prove the conclusion.
All subsequent pinning takes a principal Hessian restriction and keeps the
same one-site oscillation bound. Its domain is still a product of the same
balls, including when the set of free sites has holes. Normalization constants
depending on pinned values disappear from derivatives but are never removed
from the probability. Empty free sets are vacuous. The proof does not change
or condition any of the integrated fast coordinates.
\end{proof}

This bypasses the absolute influence test, not the retained cutoff itself.
The radius $r$ is rescaled: its angular size is $r/\sqrt\gamma$, not a
fixed positive angular fraction of the compact group. The radius is chosen
first; $g_r$ can be very small, and $\Gamma_r$ can
be large. The result does not give uniform constants as $r\to\infty$, a
quantitative growing-radius law $r=r(\gamma)$, or a larger spatial-width
regime than that inherited from V31. It also does not prove a gap for the
endpoint-varying path law by dropping the dressed endpoint functions.

\subsection{Local response on those conditionals without a Dobrushin margin}

The same argument retains spatial information. We record a derivative-class
statement rather than equating it to a total-variation or Hamming coupling.
Fix $r$. The V31 Hessian envelope, the Gaussian rows, and the diagonal Haar
term give constants $a_0,M_{a_0}<\infty$, uniform at sufficiently large
coupling, such that the off-diagonal measured Hessian blocks obey
\[
 \sup_u\sum_{v\ne u}e^{a_0d(u,v)}
                   \sup_{x,q}\|(D^2V)_{uv}(x)\|\le M_{a_0}.
\]
These are rows of the complete retained action, including its generally
nonzero Gaussian interaction. Its range need not be finite.

\begin{proposition}[Localized one-form inverse and pinned response]
\label{f16v32:response}
There are $a_r,C_r'>0$, independent of sizes and of the in-window pinning,
such that all laws in Theorem~\ref{f16v32:main} satisfy
\begin{equation}
 \|P_u\mathcal A_1^{-1}P_v\|\le
                   C_r'e^{-a_rd(u,v)}.
 \label{f16v32:inverse}
\end{equation}
For $H^1$ observables the exact covariance identity therefore gives
\begin{equation}
 |\operatorname{Cov}_\nu(F,G)|
 \le C_r'\sum_{u,v\in I}e^{-a_rd(u,v)}
                  \|\nabla_uF\|_{L^2(\nu)}\|\nabla_vG\|_{L^2(\nu)}.
 \label{f16v32:gradient-covariance}
\end{equation}
If a retained pin $q_j$ changes to $q_j'$ within its ball, while all other
pins are held fixed, a free-coordinate Lipschitz observable $F$ independent
of that pin satisfies
\begin{equation}
 |\nu^{I,q'}F-\nu^{I,q}F|
 \le C_r'|q_j'-q_j|\sum_{u\in I}e^{-a_rd(u,j)}
                                \operatorname{Lip}_u(F).
 \label{f16v32:pinned-response}
\end{equation}
Here $\operatorname{Lip}_u$ is the Euclidean Lipschitz seminorm in one
three-component bridge. This is not an all-bounded-test influence coefficient.
\end{proposition}
\begin{proof}
The diagonal gradient form, the normal-trace constraints, and the face curvature
in \eqref{f16v32:oneform} are block diagonal in the one-form component label.
They commute with multiplication of component $u$ by $e^{t d(u,v)}$.
Only the Hessian multiplication has off-diagonal component blocks. Conjugating
it by these weights changes its operator by a bounded matrix whose row and
column norms are at most
\[
 M_{a_0}\sup_{s\ge0}(e^{ts}-1)e^{-a_0s}
                    \le 2tM_{a_0}/a_0\quad(0<t\le a_0/2).
\]
Triangle inequality for anchor distance and the block Schur bound justify both
row and column estimates. At each finite regulator the conjugating weights
are bounded and invertible; their size is not used in the perturbation bound.
Choose a fixed $t>0$ with this bound at most $\kappa_r/4$. Since
$\mathcal A_1\succeq\kappa_r/2$, the inverse of the conjugated operator has
norm at most $4/\kappa_r$ by its convergent resolvent perturbation. Unweighting
one block yields \eqref{f16v32:inverse}. This uses a proved one-form floor,
not an assumed inverse for an arbitrarily pinned fast patch.

The absolute reflecting identity gives
\[
 \operatorname{Cov}_\nu(F,G)
       =\langle dF,\mathcal A_1^{-1}dG\rangle.
\]
For smooth tests it follows from the scalar Neumann Poisson equation and
$d\mathcal A=\mathcal A_1d$ proved by the form integration above. Form
approximation extends it to $H^1$ tests. Apply the block inverse bound for
\eqref{f16v32:gradient-covariance}.

For the last claim, interpolate the actual pin linearly and put
$S_t=\partial_tV^{I,q(t)}$. Differentiation of the normalized conditional
integral, on its unchanged compact product, gives
\[
       \frac d{dt}\nu^{I,q(t)}F=-\operatorname{Cov}_{\nu^{I,q(t)}}(F,S_t).
\]
The derivative of the normalizer is exactly the centering on the right.
The gradient of $S_t$ in component $v$ is the free/pinned mixed Hessian
block applied to $q_j'-q_j$. It has the inherited exponential envelope.
Convolve it with \eqref{f16v32:inverse}; polynomial anchor-ball growth bounds
the convolution after a fixed decrease of its exponent. Integrate $t$.
Approximation of Lipschitz tests inside the convex product preserves the
coordinate Lipschitz bounds and gives the stated class. All intermediate
pins remain in the same radius-$r$ balls.
\end{proof}

In particular the retained-coordinate covariance has an exponentially summable
block-row bound at each fixed $r$, uniformly in all its admitted pinning.
These coordinates have order-one scale: no $1/\gamma$ covariance bound is
claimed for them. The small $1/\gamma$ covariance in earlier sections belongs
to the fast Haar-corrected residual sources, a different observable bank.

\subsection{A local comparison with the normalized truncated Gaussian action}

Let $\nu_{0,r}^{I,q}$ denote the normalized probability on the \emph{same}
retained product with potential $\tfrac12 y^{\mathsf T}Ky$ and Lebesgue
reference. It exists even if $K$ is singular, because the domain is compact.
It is not the untruncated Gaussian measure obtained by inverting $K$.

\begin{proposition}[Normalized local bridge-law comparison]
\label{f16v32:Gaussian-local}
For each fixed $r$, in the regime above, every free-coordinate Lipschitz $F$
satisfies
\begin{equation}
 |\nu_{\gamma,r}^{I,q}F-\nu_{0,r}^{I,q}F|
       \le C_r'\gamma^{-1/2}\sum_{u\in I}\operatorname{Lip}_u(F).
 \label{f16v32:Gaussian-local-bound}
\end{equation}
Consequently the Wasserstein distance between the marginals on a selected
set $A\subseteq I$, with cost $\sum_{u\in A}|x_u-x_u'|$, is at most
$C_r'|A|\gamma^{-1/2}$. The bound is uniform in the surrounding native
volume and duration. It is not total variation of the complete history laws.
\end{proposition}
\begin{proof}
Write $h_\gamma(x)=-\log j(x/\sqrt\gamma)$ and use the normalized
interpolation
\[
 V_t=\tfrac12 y^{\mathsf T}Ky+tW,\qquad
 W=\mathcal R_\gamma(y)+\sum_{u\in I}h_\gamma(x_u),\qquad 0\le t\le1.
\]
Its two endpoints are exactly the two displayed probabilities. Intermediate
$t$ is a comparison device, not a replacement native Wilson model. All
Hessian, conditional-oscillation and one-form estimates above hold uniformly
along this segment: $K\succeq0$, the Haar Hessian is nonnegative, and the
negative remainder is at most $\delta_{\gamma,r}$. Moreover
\[
 \sup|\nabla_u W|\le C_r r\gamma^{-1/2}+C_Hr/\gamma.
\]
Differentiation keeps the actual partition function at each $t$ and yields
$d\nu_tF/dt=-\operatorname{Cov}_{\nu_t}(F,W)$. Apply
\eqref{f16v32:gradient-covariance} and sum its exponentially decaying row
against the last uniform component bound. This proves
\eqref{f16v32:Gaussian-local-bound}. Compact-metric Kantorovich duality,
or finite-grid duality followed by compactness as in V26, gives the stated
Wasserstein consequence. The number of selected coordinates is kept, rather
than being mistaken for a size-free bound on a joint law.
\end{proof}

\subsection{Why the radius cannot be removed by the new gap theorem}

For a precise non-native test, consider
\[
 d\mu_r(x,y)=Z_r^{-1}e^{-(x-y)^2/2}
             1_{[-r,r]^2}\,dx\,dy,\qquad r\ge2.
\]
Its quadratic Hessian is positive semidefinite and has an exact constant-mode
null direction. The product-window theorem gives a positive heat-bath gap at
each fixed $r$. Nevertheless that gap is at most $C/r^2$. To see this, use
$f=x+y$. Each one-coordinate conditional is a unit-variance Gaussian truncated
to an interval. The one-dimensional reflecting Bochner identity with curvature
one gives its variance at most one, so $\mathcal E_{\rm HB}(f)\le2$.

Put $s=(x+y)/2$, $t=x-y$. The Jacobian is one and
$|s|\le r-|t|/2$, $|t|\le2r$. The normalizer is at most
$2r\sqrt{2\pi}$. The subregion $|t|\le1$,
$r/2\le|s|\le3r/4$ is admissible for $r\ge2$; it gives
\[
 \operatorname{Var}_{\mu_r}(x+y)=4\mathbb E_{\mu_r}s^2
       \ge \frac{e^{-1/2}}{2\sqrt{2\pi}}r^2.
\]
Thus the Rayleigh quotient is at most $4e^{1/2}\sqrt{2\pi}/r^2$.
This example does not claim a null mode or slow dynamics in the native model.
It proves that fixed-window gaps, even with a convex quadratic action, do not
justify a uniform radius-removal theorem. In this example the unbounded
quadratic density is not even normalizable.

The hard boundary is essential: it excludes a constant vector field from
\eqref{f16v32:ball-vector}. It supplies no positive mass in $K$. The local
comparison pays $e^{\Omega_r}$, not $e^{N\Omega_r}$, but deteriorates with $r$.

V32 bypasses V31's small-window influence condition at every fixed rescaled
radius. Cutoff removal or gauge-invariant coverage is now the next retained
variable problem, with the Gaussian soft structure and the endpoint functions
still present. A quantitative growing $r=r(\gamma)$ requires growth bounds
for the inherited $C_r$ and thresholds. The arbitrary-fast-pinning problem
remains separate; none of these auxiliary clocks is physical Euclidean time.

\subsection{Validation and predecessor preservation}

The diagnostic
\begin{center}\small\path{verify_ym_f16_v32_boundary_coercivity.py}\end{center}
passes 534 exact rational and symbolic assertions. It includes normal-trace
polynomial fields, nonzero weighted Bochner face terms, semidefinite Schur
completions, normalized interpolation, and the radius-removal example.
All 28 supplied V4--V31 diagnostics were rerun unchanged and passed. These
are finite fixtures, not native Wilson integral enclosures or verification
of the Sobolev-domain passage. Actual outputs and hashes are packaged.

Only the title version and this terminal appendix differ from V31; reversing
them recovers it byte for byte. No inherited proof, endpoint, Haar power,
source or probability is silently amended. The package records the dependency
audit and final presentation checks. The historical analytic lineage has not
been independently recertified. The next checkpoint remains V35.

\begin{firewall}
V32 proves a rate-one retained heat-bath gap at every fixed finite rescaled
radius, with all further in-window retained pinning, plus derivative-class
spatial response and a local truncated-Gaussian comparison. The hard boundary,
radius-dependent constants, fixed endpoints, and inherited width restriction
remain. No positive mass in $K$, arbitrary fast-pinning theorem, quantitative
cutoff removal, physical transfer gap, or interacting continuum mass gap is
inferred. All native normalizers and fully integrated fast fields are retained.
\end{firewall}
% END F16 V32 APPEND

% BEGIN F16 V33 APPEND -- 2026-09-18
% Insert immediately before the final document-ending command in V32.
% All predecessor statements, measures, endpoints, and clocks are preserved.

\clearpage
\section{V33: exact period dynamics and the retained-cutoff obstruction}
\label{f16v33:context}

V32 obtains a gap at every fixed retained radius by using the hard boundary.
Before trying to remove that boundary, one must identify which directions it
confines in the actual Gaussian reference. We return to the complete quadratic
history, not to another absolute influence estimate. The static magnetic null
space is already known from f15 V99. The new calculation follows its spatially
constant period directions through the \emph{actual} kinetic covariance and all
history times. They form nine exactly decoupled Gaussian bridges with increment
variance two. Their normalization and diffusion cannot be replaced by a mass.

This calculation has two consequences of different types. First, the existing
native compact-window comparison gives an upper bound on the probability of a
fixed all-time retained window under the \emph{unrestricted} endpoint-fixed native
path law. This is a one-sided consequence of nested conditioning; it needs no
unproved native tail bound. Second, the complete Gaussian history controls the
combined temporal-difference and magnetic-curvature observables in a norm that
has no radius or history-length cost. These statements select a more precise
cutoff target: distinguish raw connection amplitudes from their derivatives and
retain the nonlinear period variables rather than confining them by definition.

\subsection{Sources, ownership, and separated analytic dependencies}

The V32 appendix and analytic audit were read completely. Its statement for
every fixed radius remains unchanged. The f15 V99 proposition on
$\ker D$ already identifies coarse gradients and three constant directional
periods, of dimension $B+2$ per colour. That static classification is used and
credited below, not claimed as a new theorem. The f14 V54 discussion separately
treats nonlinear constant modes at fixed four-torus volume; neither its compact
matrix integral nor its normalization is the marginal of the present history.
V2's kinetic completion and V4's quadratic action supply the matrices.

Empty semantic searches were followed by identified passage reads and mounted
text searches. This is targeted review, not certification of the archives.
The next checkpoint remains V35.
For external context, Fodor--Holland--Kuti--Nogradi--Wong,
\href{https://arxiv.org/abs/1208.1051}{\emph{The Yang--Mills gradient flow in finite
volume}}, Section 2, separates constant fields from perturbatively integrated
nonzero modes. Its introduction and mode decomposition were examined, not imported as a lattice
theorem. The external PDF screenshot was unavailable. The import memoranda's
power/locality and physical-time distinctions remain binding.

The quadratic results below use only the identified finite matrices and exact
Gaussian integration. The later \emph{native} probability bound additionally
uses V31's action-size estimate on each fixed bridge window. That estimate
inherits V16 and V30; those analytic dependencies are not recertified here.
No untruncated native Gaussian limit is assumed.

\subsection{The full quadratic history controls differences and curvature}

Suppress colour tensors $I_3$ when writing spatial matrices. Keep the original
$N=\mathsf M_B^{-1}$, $A,D$, $W_n$, $G$, and $\Lambda=A^{\mathsf T}D$.
The matrix called $K$ throughout this section is V31's \emph{complete retained
history} Schur matrix, not V1's differently named static auxiliary matrix.
Let $\mathsf T_n\mathbf y=(y_t-y_{t+1})_{t=0}^{n-1}$ and
$\mathsf C_n\mathbf y=(Dy_i)_{i=0}^{n}$. Define
\[
 \mathsf M_n=
 \mathsf T_n^{\mathsf T}(I_n\otimes G^{-1})\mathsf T_n
                   +\tfrac13\mathsf C_n^{\mathsf T}\mathsf C_n.
\]
All coordinates in these maps are the original bridge coordinates. Both external
bridge slices are included until they are explicitly fixed.

\begin{theorem}[Two-sided kinetic--curvature comparison for the complete Schur action]
\label{f16v33:quadratic}
At every finite $B=m^3$, $m\ge4$, and $n\ge1$, the exact Gaussian action obeys
\begin{equation}
 \boxed{\quad
 \mathsf M_n\preceq K\preceq
 \mathsf T_n^{\mathsf T}(I_n\otimes G^{-1})\mathsf T_n
                          +\mathsf C_n^{\mathsf T}\mathsf C_n.
 \quad}
 \label{f16v33:quadratic-comparison}
\end{equation}
In particular its kernel consists exactly of time-independent histories with
$Dy_i=0$. Its dimension is $3(B+2)$ by the inherited static classification.
The comparison has no coupling parameter and no hard retained boundary.
\end{theorem}
\begin{proof}
Eliminate the temporal ports in \eqref{f16v4:V0} using the unchanged completion
\eqref{f16v2:completed-kinetic}. This leaves the bridge kinetic term
$\frac12\sum_t(y_t-y_{t+1})^{\mathsf T}G^{-1}(y_t-y_{t+1})$ and the internal
functional
\[
 \frac12\sum_i\left\{x_i^{\mathsf T}H_i x_i+|Ax_i+Dy_i|^2\right\}
       +\frac12\sum_t(x_t-x_{t+1})^{\mathsf T}N(x_t-x_{t+1}),
\]
where $H_i=w_i\mathsf H_B+\frac12 1_{\{0,n\}}(i)C_B^{-1}$.
The original endpoint factors give $w_0=w_n=1/2$; interior weights are one.
The inherited cube bounds give
\[
 H_i\succeq2N,\qquad A N^{-1}A^{\mathsf T}\preceq4I.
\]
Indeed $N^{-1/2}\mathsf H_BN^{-1/2}\succeq4I$ and
$A=C_{\rm int}U_B$, $U_B^{\mathsf T}U_B=N$, $\|C_{\rm int}\|\le2$.

The determinant after integrating the $x$ variables is independent of the
bridge history. It cancels in $Z_0(\mathbf y)/Z_0(0)$, so
$\mathcal F_0(\mathbf y)=\frac12\mathbf y^{\mathsf T}K\mathbf y$ is the minimum
of this complete reduced quadratic functional. Drop only the nonnegative
internal kinetic term for a lower bound. At a single time its remaining minimum
is
\[
 \tfrac12(Dy_i)^{\mathsf T}(I+AH_i^{-1}A^{\mathsf T})^{-1}Dy_i
                            \ge\tfrac16|Dy_i|^2,
\]
since $AH_i^{-1}A^{\mathsf T}\preceq2I$. Choosing every $x_i=0$ instead gives
the upper bound. This proves \eqref{f16v33:quadratic-comparison}; no matrix
commutativity is assumed. Since $G\succ0$, its two lower terms vanish exactly
when $y_i$ is time independent and belongs to $\ker D$. The upper bound proves
the converse. The dimension and the gradient/period description are the already
proved f15 V99 result, transported constantly through time.
\end{proof}

A useful exact expression, consistent with this variational proof, is
\begin{equation}
 \mathcal F_0(\mathbf y)=
 \tfrac12\sum_t|y_t-y_{t+1}|_{G^{-1}}^2+
 \tfrac12\sum_i|Dy_i|^2-
 \tfrac12(\Lambda y_i)_i^{\mathsf T}W_n^{-1}(\Lambda y_i)_i.
 \label{f16v33:action-formula}
\end{equation}
It follows from the same integration, or from V3 and V4's Gaussian consistency
identity. There is no dressed-end subtraction in $\mathcal F_0$: the latter is
the fast partition-value ratio, before dividing the physical history weight by
its two dressed-end normalizers. Fixing endpoints makes those factors constant.

\subsection{The actual port metric on constant directional periods}

There are four bridge sites on each positive coarse edge $(b,\mu)$.
Let $h_{\mu,a}$ take the value $(2\sqrt B)^{-1}e_a$ on \emph{every} bridge of
direction $\mu$, and zero on the other directions, where $1\le\mu,a\le3$ and
$e_a$ is a colour unit vector. These nine vectors are orthonormal in the original
$36B$-dimensional bridge Euclidean space. Write $\mathcal H$ for their span.
They are the spatially constant periods; they do not include the $B-1$ gradient
modes in each colour.

\begin{proposition}[Exact harmonic kinetic coefficient]
\label{f16v33:metric}
For every $h\in\mathcal H$,
\begin{equation}
             Dh=0,\qquad Gh=2h,\qquad G^{-1}h=\tfrac12h.
 \label{f16v33:harmonic-metric}
\end{equation}
Consequently $\mathcal H^{n+1}$ is a reducing subspace of $K$, and, if
$y_i=\sum_{\mu,a}a_{i,\mu,a}h_{\mu,a}$,
\begin{equation}
             \mathcal F_0(\mathbf y)=\frac14
                    \sum_{t=0}^{n-1}|a_{t+1}-a_t|^2.
 \label{f16v33:period-action}
\end{equation}
The factor two in \eqref{f16v33:harmonic-metric} uses the actual nonroot ports.
Freezing those ports would incorrectly replace it by one.
\end{proposition}
\begin{proof}
The first identity follows from the inherited common-bridge curl description.
For the kinetic calculation use $G=I+E\Sigma_BE^{\mathsf T}$ from V2, where
$\Sigma_B=(\mathsf B_B^{\mathsf T}\mathsf B_B)^{-1}$ is block diagonal.
It suffices to work with one unnormalized period $h$, equal to one on its
chosen directional bridges and zero on others. In a cube write
$u=(u_1,u_2,u_3)\in\{0,1\}^3$, with root $u=0$. Set $z_{b,u}=-u_\mu$.

At a nonroot vertex, $E^{\mathsf T}h$ is $+1$ if $u_\mu=0$ and $-1$ if
$u_\mu=1$: the exterior edge is respectively incoming or outgoing. The internal
cube Laplacian applied to $z$ has exactly these values; the other two directions
contribute zero. The root value of $z$ is zero, so
$\mathsf B_B^{\mathsf T}\mathsf B_B z=E^{\mathsf T}h$ on its actual rooted
space. Across a bridge, $Ez$ is one in direction $\mu$ and zero otherwise.
Thus $E\Sigma_BE^{\mathsf T}h=h$. Periodic seams satisfy the same identity.
Restore normalization and colour to obtain \eqref{f16v33:harmonic-metric}.

In \eqref{f16v33:action-formula} the two magnetic terms annihilate the period
space at each time. The kinetic matrix preserves that space and its Euclidean
orthogonal complement because it is symmetric and has the displayed eigenvalue.
Therefore all cross terms vanish, proving both reduction and
\eqref{f16v33:period-action}. This is not merely evaluation of the quadratic
form on one selected path.
\end{proof}

\subsection{The untruncated endpoint-fixed Gaussian and its nine exact bridges}

For $n\ge2$ fix $y_0=y_n=0$ and let $K^{\rm int}$ be the principal retained
precision on the $36B(n-1)$ free scalar coordinates. Denote by $L_n$ the
$(n-1)$-dimensional Dirichlet time Laplacian, with diagonal two and adjacent
entries minus one. Since $I\preceq G\preceq30I$, Theorem~\ref{f16v33:quadratic}
gives
\begin{equation}
        K^{\rm int}\succeq\tfrac1{30}(L_n\otimes I_{36B})\succ0.
 \label{f16v33:dirichlet-lower}
\end{equation}
Thus the normalized Gaussian probability $Q_{B,n}^{00}$ on the entire interior
Euclidean bridge space exists. This does not assert the existence of a Gaussian
with free, unanchored endpoints, whose static kernel was identified above.

\begin{theorem}[Exact normalized period factor and temporal soft scale]
\label{f16v33:bridges}
Under $Q_{B,n}^{00}$ the nine period amplitudes
$a_{i,\mu,a}=\langle h_{\mu,a},y_i\rangle$ are independent of the complementary
bridge history and have joint density
\begin{equation}
 \frac{\prod_{t=0}^{n-1}p_1^{(9)}(a_{t+1}-a_t)}{p_n^{(9)}(0)}
                    \prod_{i=1}^{n-1}da_i,
 \qquad p_t^{(d)}(z)=(4\pi t)^{-d/2}e^{-|z|^2/(4t)}.
 \label{f16v33:bridge-density}
\end{equation}
Their means are zero and their covariance is
\begin{equation}
 \mathbb E[a_{i,\mu,a}a_{j,\nu,b}]
 =2\left(\min(i,j)-\frac{ij}{n}\right)\delta_{\mu\nu}\delta_{ab}.
 \label{f16v33:bridge-covariance}
\end{equation}
With nonzero fixed endpoints, their mean is the linear endpoint interpolation
and this covariance is unchanged. In particular $K^{\rm int}$ has the exact
period eigenvalues $2\sin^2(k\pi/(2n))$, $1\le k<n$, each with at least ninefold
multiplicity, and
\begin{equation}
 \frac{2}{15}\sin^2\frac\pi{2n}
 \le\lambda_{\min}(K^{\rm int})
 \le2\sin^2\frac\pi{2n}.
 \label{f16v33:soft-scale}
\end{equation}
No positive lower edge uniform in duration can hold for this untruncated
retained Gaussian precision, even at fixed spatial width.
\end{theorem}
\begin{proof}
Apply the orthogonal period/complement change of coordinates at every time.
Proposition~\ref{f16v33:metric} eliminates all cross terms, including terms
involving fixed endpoints. The density consequently factors after normalization.
The period action is \eqref{f16v33:period-action}; convolution of the displayed
heat densities gives its exact denominator $p_n^{(9)}(a_n-a_0)$. At zero endpoints
this is \eqref{f16v33:bridge-density}. It is not an adjustable determinant.

The inverse of $L_n$ has entries $\min(i,j)-ij/n$. Direct second differences,
with values zero at times zero and $n$, verify $L_nL_n^{-1}=I$.
The period precision is $L_n/2$, proving the covariance. Completing the square
about the linear interpolation proves the endpoint assertion. The sine vectors
diagonalize $L_n$, with eigenvalues $4\sin^2(k\pi/(2n))$.
Equation~\eqref{f16v33:dirichlet-lower} proves the lower bound in
\eqref{f16v33:soft-scale}; the reducing period space proves the upper bound.
\end{proof}

For completeness, the sum of rate-one three-component bridge heat baths of
$Q_{B,n}^{00}$ also has gap at most $30\sin^2(\pi/(2n))$. Indeed choose a unit
first temporal period eigenvector $v$ and test $f=v^{\mathsf T}y$.
Its variance is $[2\sin^2(\pi/(2n))]^{-1}$. Every diagonal bridge block of
$K^{\rm int}$ is at least $I_3/15$ by \eqref{f16v33:dirichlet-lower}, so the
sum of its exact Gaussian conditional variances is at most $15\|v\|^2=15$.
The Rayleigh quotient gives the stated upper bound. This sampler and the
precision are not physical Hamiltonians. The result does not contradict V32:
its hard window is absent here.

\subsection{A native upper bound on fixed all-time bridge-window coverage}

Let $\boldsymbol\nu_{\gamma,B,n}^{00}$ be the normalized \emph{original}
endpoint-fixed retained path measure with $y_0=y_n=0$, integrated over every
interior principal bridge chart. It is the positive unprojected history
integrand of V4, with all fast fields integrated; it is not a newly Gauss-projected
physical probability. In its chart it is proportional to
\[
 e^{-\mathcal F_\gamma(\mathbf y)}
           \prod_{u\ {
m interior}}j(y_u/\sqrt\gamma)\,dy_u.
\]
The scalar and dressed-end factors cancel because the endpoints are fixed.
For a finite rescaled radius $r$ put
\[
 G_r=\{|y_u|\le r\text{ for every interior bridge }u\},\qquad
 M_{B,n}=12B(n-1).
\]
When $r<\pi\sqrt\gamma$, this is an event in the original full compact path.
The event concerns the \emph{entire} retained history, not one local marginal.

\begin{proposition}[Nested-window comparison with the full native normalizer retained]
\label{f16v33:native-comparison}
Fix $B,n$ and $0<r\le R<\infty$. At the V31 entry threshold for the actual
radius $R$, and $\gamma\ge\max(4R^2,C_H)$, define
\[
 d_{\gamma,R}=M_{B,n}R^2
                  \left(C_R\gamma^{-1/2}+\frac{C_H}{2\gamma}\right).
\]
Then
\begin{equation}
 \boxed{\quad
 \boldsymbol\nu_{\gamma,B,n}^{00}(G_r)
 \le\min\left\{1,\ e^{2d_{\gamma,R}}
               \frac{Q_{B,n}^{00}(G_r)}{Q_{B,n}^{00}(G_R)}\right\}.
 \quad}
 \label{f16v33:nested-probability}
\end{equation}
In particular,
\begin{equation}
 \limsup_{\gamma\to\infty}
       \boldsymbol\nu_{\gamma,B,n}^{00}(G_r)\le Q_{B,n}^{00}(G_r).
 \label{f16v33:native-limsup}
\end{equation}
The limit holds with $B,n,r$ fixed. It is an upper bound, not convergence of
the unrestricted native path to an untruncated Gaussian.
\end{proposition}
\begin{proof}
On $G_R$, V31 gives
$|\mathcal F_\gamma-\mathcal F_0|\le C_R\gamma^{-1/2}\sum_u|y_u|^2$.
V32's elementary Haar estimate gives
$0\le-\log j(y_u/\sqrt\gamma)\le C_H|y_u|^2/(2\gamma)$.
Thus the difference of the two measured potentials is bounded in absolute
value by $d_{\gamma,R}$. The \emph{normalized} likelihood ratio between
$\boldsymbol\nu_\gamma(\cdot\mid G_R)$ and $Q_{B,n}^{00}(\cdot\mid G_R)$
is consequently between $e^{-2d_{\gamma,R}}$ and $e^{2d_{\gamma,R}}$.
Both restricted normalizers are included.

Since $G_r\subseteq G_R$,
$\boldsymbol\nu_\gamma(G_r)=\boldsymbol\nu_\gamma(G_R)
\boldsymbol\nu_\gamma(G_r\mid G_R)\le\boldsymbol\nu_\gamma(G_r\mid G_R)$.
This proves \eqref{f16v33:nested-probability}, without estimating the full
normalizer or any native tail outside $G_R$. For fixed $R,B,n$, the inherited
width condition is eventually satisfied and $d_{\gamma,R}\to0$. Take this
limsup first. Only afterward send $R\to\infty$; the proper Gaussian of
\eqref{f16v33:dirichlet-lower} has $Q_{B,n}^{00}(G_R)\to1$. This proves
\eqref{f16v33:native-limsup} in the stated order of limits.
\end{proof}

A fully explicit Gaussian denominator bound is available:
\begin{equation}
 Q_{B,n}^{00}(G_R)\ge1-6M_{B,n}e^{-R^2/(45n)}.
 \label{f16v33:outer-window}
\end{equation}
Indeed \eqref{f16v33:dirichlet-lower} bounds every scalar variance by
$30i(n-i)/n\le15n/2$. If a three-vector exceeds $R$, one coordinate exceeds
$R/\sqrt3$. The scalar Gaussian exponential bound and a union bound prove
\eqref{f16v33:outer-window}, without independence of the bridge coordinates.
For example $R^2\ge45n\log(12M_{B,n})$ makes this lower bound at least $1/2$.
The native constant $C_R$ and its threshold still depend on this chosen radius.

\begin{theorem}[An actual native obstruction to uniform all-time window coverage]
\label{f16v33:cutoff}
For fixed $B=m^3$, $m\ge4$, and $r>0$, every even $n\ge2$ satisfies
\begin{equation}
 \boxed{\quad
 \limsup_{\gamma\to\infty}
 \boldsymbol\nu_{\gamma,B,n}^{00}(G_r)
 \le\left[\min\left\{1,4r\sqrt{\frac{B}{\pi n}}\right\}\right]^9.
 \quad}
 \label{f16v33:midpoint-cutoff}
\end{equation}
There is also a temporal skeleton bound. Put
$\tau=\max\{1,\lceil16Br^2\rceil\}$ and take $n=k\tau$, $k\ge2$. Then
\begin{equation}
 \limsup_{\gamma\to\infty}
 \boldsymbol\nu_{\gamma,B,n}^{00}(G_r)
       \le\min\{1,k^{9/2}2^{-9(k-1)}\}.
 \label{f16v33:skeleton-cutoff}
\end{equation}
Consequently
\begin{equation}
 \inf_{n\ge2}\limsup_{\gamma\to\infty}
              \boldsymbol\nu_{\gamma,B,n}^{00}(G_r)=0.
 \label{f16v33:no-uniform-coverage}
\end{equation}
In particular there exist $n_j\to\infty$ and $\gamma_j\to\infty$, at fixed $B$,
for which the native probability of this fixed all-time window tends to zero.
No quantitative joint growth rate for $n_j,\gamma_j$ is asserted.
\end{theorem}
\begin{proof}
On $G_r$, for each period and each interior time,
\[
 |a_{i,\mu,a}|=\left|\frac1{2\sqrt B}
             \sum_{e:\,\mathrm{dir}(e)=\mu}y_{i,e}^a\right|
                                          \le2\sqrt B\,r=:A_r^{\rm per}.
\]
At $i=n/2$, Theorem~\ref{f16v33:bridges} gives nine independent scalar
Gaussians of variance $n/2$. Each has density at most $(\pi n)^{-1/2}$.
Their simultaneous interval probability is at most the right side of
\eqref{f16v33:midpoint-cutoff}. The full event is a subset of that event.
Apply \eqref{f16v33:native-limsup}.

For the sharper skeleton calculation, one scalar period sampled at times
$\tau,2\tau,\ldots,(k-1)\tau$ has transition density $p_\tau^{(1)}$ and
exact bridge denominator $p_{k\tau}^{(1)}(0)$. Uniformly in its incoming value,
\[
 \int_{-A_r^{\rm per}}^{A_r^{\rm per}}p_\tau^{(1)}(x-y)\,dy
    \le\frac{A_r^{\rm per}}{\sqrt{\pi\tau}}\le\frac12.
\]
Bound the final transition density by $(4\pi\tau)^{-1/2}$, integrate the
preceding $k-1$ transitions with this row bound, and divide by
$p_{k\tau}^{(1)}(0)$. The resulting probability is at most
$\sqrt k\,2^{-(k-1)}$. The nine scalar period bridges are independent, so
raise this bound to the ninth power. Again $G_r$ implies this skeleton event,
and \eqref{f16v33:native-limsup} transfers its upper bound. Finally choose
$k_j\to\infty$ and, for each resulting finite $n_j$, choose a sufficiently
large $\gamma_j$ making the limsup bound hold within $1/j$. Increase these
couplings successively. This is a diagonal choice, not a uniform-radius
estimate on $C_R$ or an exchange of the two limits.
\end{proof}

This native conclusion uses only the nested-window upper comparison; it does
not prove that a \emph{single-time} native marginal is non-tight, because that
event is not contained in every $G_R$. It does not contradict V32's positive
gap for the law already conditioned on $G_r$. An all-time coordinate window
and local gauge-invariant coverage are different targets. Linear period
coordinates are not nonlinear gauge-invariant Wilson holonomies, and no
physical mass conclusion follows from \eqref{f16v33:no-uniform-coverage}.

\subsection{A positive radius-free Gaussian interface for the next native estimate}

The same quadratic calculation supplies a norm that does not pay the period
amplitude. Let $\Delta y_t=y_t-y_{t+1}$, keep zero endpoints, and form
\[
 \mathcal B\mathbf y=
 \left((G^{-1/2}\Delta y_t)_{t=0}^{n-1},
                         (3^{-1/2}Dy_i)_{i=1}^{n-1}\right).
\]
Its components are assembled together, including cross covariances.

\begin{theorem}[Uniform difference--curvature covariance, without a retained radius]
\label{f16v33:derivatives}
Under the proper untruncated Gaussian $Q_{B,n}^{00}$,
\begin{equation}
                  \operatorname{Cov}(\mathcal B\mathbf y)\preceq I.
 \label{f16v33:derivative-Gram}
\end{equation}
In particular, for arbitrary coefficient banks $a_t,b_i$,
\begin{equation}
 \boxed{\quad
 \operatorname{Var}\left(\sum_t a_t^{\mathsf T}\Delta y_t+
                                \sum_i b_i^{\mathsf T}Dy_i\right)
 \le\sum_t a_t^{\mathsf T}Ga_t+3\sum_i|b_i|^2
 \le30\sum_t|a_t|^2+3\sum_i|b_i|^2.
 \quad}
 \label{f16v33:derivative-bound}
\end{equation}
The constants are independent of $B,n$ and there is no hard window. The same
covariance bounds hold after arbitrary further \emph{Gaussian} retained
pinning, and for nonzero fixed endpoints, with the appropriate means subtracted.
For the exact period factors,
\begin{equation}
 \operatorname{Cov}(a_{t+1}-a_t,a_{s+1}-a_s)
                  =2(\delta_{ts}-n^{-1})I_9.
 \label{f16v33:increment-covariance}
\end{equation}
Thus the period amplitudes grow through the bridge, but their increments have
uniform covariance with the endpoint constraint retained.
\end{theorem}
\begin{proof}
On the endpoint-fixed space, \eqref{f16v33:quadratic-comparison} says
$\mathcal B^{\mathsf T}\mathcal B\preceq K^{\rm int}$. Hence
$\|\mathcal B(K^{\rm int})^{-1/2}\|\le1$, and multiplication by its adjoint
proves \eqref{f16v33:derivative-Gram}. Apply it to the output coefficient vector
$((G^{1/2}a_t)_t,(\sqrt3b_i)_i)$; this proves the first inequality in
\eqref{f16v33:derivative-bound}, including all mixed terms. The second is
$G\preceq30I$. Further Gaussian pinning takes the same principal precision and
columns of $\mathcal B$; their matrix inequality and the proof persist.
Nonzero endpoint values change linear terms, not the precision.
Finally apply two temporal differences to \eqref{f16v33:bridge-covariance},
or condition $n$ independent variance-two increments on their zero sum. This
gives \eqref{f16v33:increment-covariance}, including its negative off-diagonal
entries and its exact zero mode in the sum of the increments.
\end{proof}

\paragraph{What would be sufficient, and what is not supplied.}
On any fixed product window with endpoints fixed, an independently proved
native lower bound $D^2V\succeq(1-\eta)K^{\rm int}$, $0\le\eta<1$, would
preserve \eqref{f16v33:derivative-Gram} with factor $(1-\eta)^{-1}$ by the
reflecting covariance identity and its nonnegative face terms. This is a
\emph{relative} curvature requirement. V31 supplies instead the absolute
bound $D^2V\succeq K^{\rm int}-\delta_{\gamma,r}I$. Using only that bound
requires $\delta_{\gamma,r}$ smaller than a constant times $n^{-2}$ to obtain
the relative implication. It therefore does not close a duration-uniform
native derivative theorem. No such additional relative bound is asserted here.

The next noncircular task is to formulate the native comparison in temporal
differences and curvature while retaining the compact nonlinear periods and
the exact gauge/endpoint data. The period/complement orthogonal change of
coordinates used for a Gaussian does not turn the native product domain into
a product in those variables. Its moving fibres, density, and conditional
normalizers would have to be kept. The f14 fixed-four-torus constant-mode
analysis is a useful separate input, not a ready-made history marginal.
Nothing here removes the spatial-width restriction from the native estimates,
identifies an auxiliary clock with physical time, or proves a physical gap.

\subsection{Verification and preservation}

The accompanying diagnostic is
\begin{center}\small\path{verify_ym_f16_v33_period_bridge.py}.\end{center}
It checks the literal rooted-cube and periodic bridge incidence identity,
Gaussian Schur comparisons, temporal inverse and normalization identities,
period reduction, derivative covariance and the finite-window comparison
algebra. The checker passes 226 exact assertions. All 29 available V4--V32 checkers were
rerun unchanged and passed. These are finite algebraic fixtures, not native
integral enclosures or formal verification. Actual reports are packaged.

Only V32's one title version and this terminal insertion are changed. Their
reversal recovers V32 byte for byte. The cutoff theorem depends on V31; the quadratic theorems do not depend on
V30's native gap. The archive has not been independently recertified.

\begin{firewall}
V33 identifies nine exact period Gaussian bridges in the complete retained
history, proves a two-sided kinetic--curvature comparison, and obtains
radius-free Gaussian derivative covariance bounds. A normalized nested-window
argument gives a one-sided native obstruction to uniform fixed all-time window
coverage in a stated order of limits. It does not claim an untruncated native
Gaussian limit, a native all-radius derivative inequality, local native
non-tightness, a gauge-invariant cutoff-removal theorem, a thermodynamic gap,
or a physical Yang--Mills mass gap. Every native normalizer used in the comparison
is retained, and all physical clocks and Gauss/continuum obligations remain
separate.
\end{firewall}
% END F16 V33 APPEND

% BEGIN F16 V34 APPEND -- 2026-09-18
% Insert before the final document-ending command of the cumulative V33.
% The preceding text, all original measures, sources and update clocks are preserved.

\clearpage
\section{V34: endpoint anchoring and quantitative native cutoff removal}
\label{f16v34:context}

V33 obtains only an upper comparison for the probability of a retained window:
its proof does not estimate the full native normalizer. We go upstream to that
missing estimate. With both retained endpoints fixed at the identity, V5's
compact fast pins and the original bridge kinetic words also confine the
retained coordinates. The confinement constant deteriorates as the inverse
square of the duration. This is sufficient for a complete normalized Laplace
comparison, but not for a duration-uniform mass statement.

The resulting theorem removes the retained cutoff at every fixed regulator,
and supplies an explicit, conservative simultaneous regime
$Bn\log(n+1)=o(\log\gamma)$. It upgrades V33's one-sided comparison to total
variation and polynomial-moment convergence of the full native joint law.
Consequences include a native difference--curvature covariance comparison, a
comparison for actual nonlinear Wilson-word observables, and a quantitative
period-bridge limit. The latter also shows that a single midpoint bridge cannot
remain in a fixed rescaled window along a specified slowly growing history.
Every assertion concerns the endpoint-fixed, unprojected history, not the
physical Gauss-reduced vacuum.

\subsection{Dependencies and the exact joint probability}

The supplied V33 appendix and audit were read completely. Its period metric,
Gaussian history and derivative bank are reused, not recomputed as new results.
The analytic starting points here are substantially earlier: V4's exact history
integrand, V5's complete-chart compact pin and frame relation, and V6's bounded
ordered-word derivatives. The new proof uses neither V16's conditional-moment
iteration, V30's screening, nor V31's fixed-window action comparison. In
particular no growing-radius estimate for the constants $C_r$ is assumed.
Those earlier results keep their own scopes and are not silently amended.

Scoped semantic searches returned no results; mounted exact-text searches and
identified passages were then checked. The inherited tree-pin and finite-window
Laplace arguments are credited. No exhaustive novelty or archive certification
is asserted. For context, Faria da Veiga--O'Carroll,
\href{https://arxiv.org/abs/1903.09829}{\emph{On Thermodynamic and Ultraviolet
Stability of Yang--Mills}}, uses exact tree coordinates for partition bounds.
Spokoiny, \href{https://arxiv.org/abs/2204.11038}{\emph{Dimension free
non-asymptotic bounds on the accuracy of high dimensional Laplace approximation}},
keeps dimension costs in a normalized approximation. Their primary abstracts
and scope were checked, not their complete proofs. Neither theorem is imported
for this history. The finite-dimensional domination and normalization estimates
needed below are proved directly. The next companion checkpoint remains V35.

Fix $B=m^3$, $m\ge4$, $n\ge2$, and $Z_0=Z_n=I$. Use V5's unmoving fast
coordinates $\zeta=(x,a)$ and all interior retained coordinates $y$ on their
\emph{complete} principal charts. Put $w=(\zeta,y)$ and
\begin{equation}
 d=d_{B,n}=15B(n+1)+21Bn+36B(n-1)=(72n-21)B.
 \label{f16v34:dimension}
\end{equation}
Let $\mathcal P_\gamma\subset\mathbb R^d$ be their product of
three-dimensional radius-$\pi\sqrt\gamma$ balls. Its boundary has measure zero.
With exactly V5's classical action, define
\begin{equation}
 \begin{split}
 \mathcal J_\gamma(w)&=J^f_\gamma(\zeta)
                 \prod_{i=1}^{n-1}\prod_{e\ {\rm bridge}}
                           j(y_{i,e}/\sqrt\gamma),\\
 q_\gamma(w)&=1_{\mathcal P_\gamma}(w)
                 \mathcal J_\gamma(w)e^{-U_\gamma(\zeta;y)},\\
 \mathcal Z_\gamma&=\int_{\mathbb R^d}q_\gamma(w)\,dw,
 \qquad d\mathbb P_\gamma=\mathcal Z_\gamma^{-1}q_\gamma(w)\,dw.
 \end{split}
 \label{f16v34:joint-law}
\end{equation}
Chord endpoint Haar powers are $1/2$, all other free-group powers are one.
The two endpoint Gaussian amplitudes remain in $U_\gamma$. The factors
$c_{\gamma,B,n}$ and $F_{\gamma,2}(I)^{-1}$ cancel only when this fixed-endpoint
probability is normalized. Thus its retained marginal is exactly V33's
$\boldsymbol\nu_{\gamma,B,n}^{00}$, not a new probability or a replacement of
its physical transfer normalization.

Write $U_0(w)=\tfrac12w^{\mathsf T}\mathcal H_{B,n}w$ for the original full
quadratic limit with these fixed endpoints, and set
\[
 q_0=e^{-U_0},\qquad
 \mathcal Z_0=(2\pi)^{d/2}(\det\mathcal H_{B,n})^{-1/2},\qquad
 d\mathbb Q=\mathcal Z_0^{-1}q_0\,dw.
\]
Positive definiteness and this normalizer are verified below. Its retained
marginal is V33's $Q_{B,n}^{00}$, since the fast Gaussian integration and its
determinant are unchanged.

\subsection{The compact kinetic terms transmit the endpoint pin}

Let $\operatorname{dist}(U,V)\in[0,\pi]$ be the bi-invariant geodesic distance
on $\SU(2)$ in the convention $s(U)=1-\Re U$. Thus
$\operatorname{dist}(I,\operatorname{Exp}(z/\sqrt\gamma))=|z|/\sqrt\gamma$
on a principal chart and
$\operatorname{dist}(I,U)^2\le(\pi^2/2)s(U)$.

\begin{proposition}[A complete joint pin with its duration cost]
\label{f16v34:compact-pin}
There is a fixed $c>0$, independent of $B,n,\gamma$, such that on every
complete principal chart with $Z_0=Z_n=I$,
\begin{equation}
 \boxed{\qquad U_\gamma(\zeta;y)\ge\frac{c}{n^2}
                    (\|\zeta\|^2+\|y\|^2),\qquad \gamma>0.\qquad}
 \label{f16v34:pin}
\end{equation}
The same inequality holds for $U_0$. Moreover
$\mathcal H_{B,n}\preceq H_*I$ for a fixed finite $H_*$.
The identity is the unique zero of the classical endpoint-fixed action in
these coordinates. No retained window is imposed.
\end{proposition}
\begin{proof}
V5's anchor is independent of $y$ and every other classical term is
nonnegative. Its literal proof therefore gives, for all compact retained values,
\[
 U_\gamma\ge a_*\|\zeta\|^2,\qquad a_*=(6\pi^2)^{-1}.
\]
We are not extending V5's radius-dependent pressure estimates. Only its
explicit, retained-input-independent anchor is used here.

Let $z_{t,v}$ be the principal coordinate of the original moving port $r_{t,v}$.
V5's exact relation $r_t=f_t^{-1}h_tf_{t+1}$ and its global angle estimate give
\[
 \sum_{t,v}|z_{t,v}|^2\le C_f\|\zeta\|^2,
 \qquad C_f=3+6\|\mathsf F_\Box\|^2.
\]
Indeed $|z_{t,v}|\le|a_{t,v}|+|(\mathsf F_Bx_t)_v|
+|(\mathsf F_Bx_{t+1})_v|$; each internal slice occurs at most twice.
This is a group-distance bound, not the linearized frame change.

For an actual bridge kinetic word
$Q_{t,e}=r_{t,s}^{-1}Z_{t,e}r_{t,t(e)}Z_{t+1,e}^{-1}$, the exact rearrangement
\[
 Z_{t,e}Z_{t+1,e}^{-1}
 =r_{t,s}Q_{t,e}Z_{t+1,e}r_{t,t(e)}^{-1}Z_{t+1,e}^{-1}
\]
and the bi-invariant triangle inequality imply
\[
 \gamma\operatorname{dist}(Z_{t,e},Z_{t+1,e})^2
 \le3|z_{t,s}|^2+3|z_{t,t(e)}|^2+\tfrac32\pi^2\gamma s(Q_{t,e}).
\]
Each nonroot fine vertex has at most three incident bridges. Sum over all
bridges and time bonds; roots have zero port value. Since the bridge kinetic
terms are among the nonnegative terms of $U_\gamma$,
\begin{equation}
 \gamma\sum_{t,e}\operatorname{dist}(Z_{t,e},Z_{t+1,e})^2
 \le C_{\nabla} U_\gamma,
 \qquad C_{\nabla}=9C_f/a_*+3\pi^2/2.
 \label{f16v34:metric-increments}
\end{equation}
No two noncommutative factors have been interchanged.

For each bridge apply the Dirichlet time-path inequality to the scalar sequence
$u_i=\operatorname{dist}(I,Z_{i,e})$, with $u_0=u_n=0$.
The reverse triangle inequality gives
$|u_{i+1}-u_i|\le\operatorname{dist}(Z_{i+1,e},Z_{i,e})$.
The least Dirichlet eigenvalue is $4\sin^2(\pi/(2n))\ge4/n^2$, whence
\[
 \|y\|^2=\gamma\sum_{i,e}u_i^2
 \le\frac{n^2}{4}\gamma\sum_{t,e}
                  \operatorname{dist}(Z_{t,e},Z_{t+1,e})^2.
\]
Combine with the fast anchor. One permissible constant is
$c=(a_*^{-1}+C_{\nabla}/4)^{-1}$. All root, seam and endpoint incidences are
included in the counts. Taking the quadratic limit on a fixed finite vector
proves the same lower bound for $U_0$, hence positive definiteness of
$\mathcal H_{B,n}$. The quadratic form is a sum of fixed-length local squares
with bounded incidence and the original bounded endpoint quadratics. Its
absolute block rows are uniformly bounded, giving $H_*<\infty$.
\end{proof}

The $n^{-2}$ loss is not an avoidable notation issue: V33's exact period
bridges exhibit that temporal soft scale in the quadratic limit. The estimate
pins the connection through its actual endpoints; it is not a curvature-relative
mass or a pin on a history with its endpoints left free.

\subsection{A global Taylor and Haar estimate on the original charts}

\begin{proposition}[Unrestricted chart remainder with bounded word incidence]
\label{f16v34:Taylor}
There is a fixed finite $C_W$, independent of $B,n,\gamma$, such that, on
$\mathcal P_\gamma$,
\begin{equation}
 |U_\gamma(w)-U_0(w)|\le C_W\gamma^{-1/2}\|w\|^3,
 \qquad
 0\le1-\mathcal J_\gamma(w)\le\frac{\|w\|^2}{3\gamma}.
 \label{f16v34:Taylor-bounds}
\end{equation}
The latter bound retains all half-Haar endpoint powers.
\end{proposition}
\begin{proof}
Every classical Wilson term is $\gamma s$ of a fixed ordered product of
exponentials of fixed local linear forms in $w/\sqrt\gamma$. The frames have
this same form in the specified unmoving charts. Duhamel differentiation of
unitary factors gives a globally bounded third derivative with coefficient
$C\gamma^{-1/2}$. Its value and first derivative at zero vanish, and its
quadratic term is the corresponding term of $U_0$. Taylor's integral formula
therefore bounds one remainder by $C\gamma^{-1/2}\|w_T\|^3$ on its fixed
carrier $T$. Bounded carrier size and incidence yield
$\sum_T\|w_T\|^3\le C\sum_g|w_g|^3\le C\|w\|^3$.
The endpoint quadratic is exact and has no remainder. This is the ordered-word
certificate already used in V6; there is no principal logarithm of an
intermediate product to differentiate.

For $0\le t\le\pi$, $0\le\sin(t)/t\le1$ and
$1-\sin(t)/t\le t^2/6$. A Haar factor of power $1/2$ is $\sin(t)/t$;
a full factor is its square and has defect at most $t^2/3$.
For numbers in $[0,1]$, $1-\prod a_j\le\sum(1-a_j)$.
Apply this to the complete joint product, with $t=|w_g|/\sqrt\gamma$.
No logarithm of the vanishing Haar density is estimated at the cut.
\end{proof}

Put $a=c/n^2$. For any real $s\ge0$ define the explicit Gaussian radial integral
\begin{equation}
 M_s(d,a)=\int_{\mathbb R^d}|w|^s e^{-a|w|^2}\,dw
 =\pi^{d/2}a^{-(d+s)/2}
                \frac{\Gamma((d+s)/2)}{\Gamma(d/2)}.
 \label{f16v34:radial-moment}
\end{equation}
For an integer $p\ge0$ let
\[
 A_p=C_W\{M_3+M_{p+3}\}
       +\tfrac13\{M_2+M_{p+2}\}
       +\pi^{-1}\{M_1+M_{p+1}\},
\]
where these moments have arguments $(d,a)$. All quantities are finite and the
dimension is explicit.

\begin{proposition}[Unnormalized weighted comparison and the actual partition mass]
\label{f16v34:unnormalized}
For $\gamma\ge1$,
\begin{equation}
 \int(1+|w|^p)|q_\gamma-q_0|\,dw\le A_p\gamma^{-1/2}.
 \label{f16v34:unnormalized-bound}
\end{equation}
Consequently
$|\mathcal Z_\gamma-\mathcal Z_0|\le A_0\gamma^{-1/2}$ and
$\mathcal Z_\gamma/\mathcal Z_0\to1$ at fixed $B,n$.
\end{proposition}
\begin{proof}
On the complete chart, both exponents are at least $a|w|^2$. Hence
$|e^{-U_\gamma}-e^{-U_0}|\le|U_\gamma-U_0|e^{-a|w|^2}$.
The Haar defect in Proposition~\ref{f16v34:Taylor} is paid separately. This gives
\[
 |q_\gamma-q_0|\le
 [C_W\gamma^{-1/2}|w|^3+(3\gamma)^{-1}|w|^2]e^{-a|w|^2}
 \quad\text{on }\mathcal P_\gamma.
\]
Outside that product, $|w|\ge\pi\sqrt\gamma$ and $q_\gamma=0$.
Use $1_{\mathcal P_\gamma^c}\le |w|/(\pi\sqrt\gamma)$ and
$q_0\le e^{-a|w|^2}$ there. Integrating proves the stated formula. In
particular the actual full chart mass, not a restricted mass, is compared.
\end{proof}

\subsection{Complete native convergence, including a simultaneous growing regime}

Write $p_\gamma=q_\gamma/\mathcal Z_\gamma$ and $p_0=q_0/\mathcal Z_0$.
The subscript zero here labels the Gaussian probability, not zero coupling.

\begin{theorem}[Quantitative full-native Laplace comparison]
\label{f16v34:full-limit}
For each fixed $B,n$ and every integer $p\ge0$,
\begin{equation}
 \boxed{\quad
 \int_{\mathbb R^d}(1+|w|^p)|p_\gamma-p_0|\,dw
            \le C_{p,B,n}\gamma^{-1/2}
 \quad}
 \label{f16v34:weighted-limit}
\end{equation}
at sufficiently large coupling. Thus the complete native joint probability
converges in total variation, with all fixed polynomial moments, to the exact
untruncated joint Gaussian. Its retained marginal converges to $Q_{B,n}^{00}$
in the same sense. No fast or retained cutoff remains in these probabilities.

There are constants $C_p<\infty$, independent of $B,n,\gamma$, for which the
right side may be bounded by
\begin{equation}
 \gamma^{-1/2}\exp\{C_p d_{B,n}\log(n+1)\},
 \label{f16v34:dimension-rate}
\end{equation}
provided the analogous expression with $C_0$ is at most a fixed sufficiently
small constant. In particular, along any sequence with
\begin{equation}
            Bn\log(n+1)=o(\log\gamma),
 \label{f16v34:joint-regime}
\end{equation}
the weighted difference in \eqref{f16v34:weighted-limit} tends to zero for
every fixed $p$. This is a conservative slowly growing regime, not V30's much
larger fixed-window width regime or a uniform-in-duration theorem.
\end{theorem}
\begin{proof}
If $A_0/\sqrt\gamma\le\mathcal Z_0/2$, the actual normalizer satisfies
$\mathcal Z_\gamma\ge\mathcal Z_0/2$. Let
$N_p=\int(1+|w|^p)q_0(w)\,dw$. Direct subtraction of the two normalized
densities gives
\begin{equation}
 \int(1+|w|^p)|p_\gamma-p_0|\,dw
 \le\frac2{\mathcal Z_0\sqrt\gamma}
                    \left(A_p+\frac{A_0N_p}{\mathcal Z_0}\right).
 \label{f16v34:normalized-bound}
\end{equation}
This includes both partition functions. Marginalization contracts the weighted
bound for weights depending only on retained variables.

We give the dimension estimate to specify a genuine simultaneous regime.
The upper Hessian bound gives
$\mathcal Z_0\ge(2\pi/H_*)^{d/2}$. Each ratio $M_s/\mathcal Z_0$ is at most
\[
 (H_*/(2a))^{d/2}a^{-s/2}
                         \frac{\Gamma((d+s)/2)}{\Gamma(d/2)}.
\]
For fixed integer $s$, the last ratio is at most
$C_s(d+s)^{s/2}$. For even $s$ this is the finite gamma-function product;
for odd $s$, Cauchy--Schwarz between the adjacent even radial moments proves it.
Also $\mathcal H_{B,n}\succeq2aI$, so the Gaussian moment
$N_p/\mathcal Z_0$ is at most $1+C_p a^{-p/2}(d+p)^{p/2}$.
Since $a=c/n^2$ and $d\ge1$, all these factors are bounded by
$\exp\{C_p d\log(n+1)\}$ after fixed constants are enlarged. Inserting in
\eqref{f16v34:normalized-bound} proves \eqref{f16v34:dimension-rate} and the
stated normalizer threshold. Relation \eqref{f16v34:joint-regime} makes each
fixed exponent $C_p d\log(n+1)=o(\log\gamma)$, so the error is
$\gamma^{-1/2+o(1)}$. No radius-dependent inherited threshold is used.
\end{proof}

At fixed $B,n$, the compact domination also gives a uniform native tail, once
$\mathcal Z_\gamma\ge\mathcal Z_0/2$:
\begin{equation}
 \mathbb P_\gamma\{|w|>R\}
 \le \frac2{\mathcal Z_0}(2\pi/a)^{d/2}e^{-aR^2/2}.
 \label{f16v34:tail}
\end{equation}
The dimension and $a=c/n^2$ are deliberately visible. This proves tightness
as $\gamma\to\infty$ at each fixed regulator; it is not a bound uniform as
$n\to\infty$. The estimate also gives uniform exponential moments with any
fixed exponent smaller than $a$, at that fixed regulator.

\subsection{Period diffusion now follows under the full native law}

\begin{theorem}[Two-sided window comparison and a local midpoint consequence]
\label{f16v34:periods}
For each fixed $B,n,r$, with exactly the probabilities and event of V33,
\begin{equation}
 \lim_{\gamma\to\infty}
       \boldsymbol\nu_{\gamma,B,n}^{00}(G_r)=Q_{B,n}^{00}(G_r).
 \label{f16v34:window-equality}
\end{equation}
The full vector of V33's nine period amplitudes converges in total variation
and moments to its exactly normalized Gaussian bridge law. For even $n$ and
any specified bridge $e$ at time $n/2$,
\begin{equation}
 \lim_{\gamma\to\infty}\operatorname{Var}_{\boldsymbol\nu_\gamma}
                 (y_{n/2,e}^{a})\ge\frac{n}{8B},
 \label{f16v34:local-variance}
\end{equation}
for each colour component. The corresponding one-bridge probability obeys
\begin{equation}
 \lim_{\gamma\to\infty}\boldsymbol\nu_\gamma
        \{|y_{n/2,e}|\le r\}
 \le\left[\min\left\{1,4r\sqrt{\frac B{\pi n}}\right\}\right]^3.
 \label{f16v34:local-cutoff}
\end{equation}
These are rescaled coordinate statements, not claims about nonlinear
period holonomies or the physical mass.
\end{theorem}
\begin{proof}
Total variation proves \eqref{f16v34:window-equality} for the original event,
with no exceptional-boundary qualification. Orthogonal linear projection onto
the period coordinates contracts total variation, and the weighted convergence
with $p=2$ proves covariance convergence. The native means vanish by simultaneous
global conjugation at identity endpoints; alternatively the same convergence
centers them correctly.

V33's reducing Gaussian sector has midpoint amplitude covariance $(n/2)I_9$.
A bridge of direction $\mu$ has coefficient $1/(2\sqrt B)$ in each of its
three period directions. The complementary Gaussian contributes a nonnegative
covariance and is independent of those amplitudes. Thus its three-component
marginal covariance is $\sigma^2I_3$ with $\sigma^2\ge n/(8B)$. Colour
isotropy follows from the actual quadratic colour tensor $I_3$. This proves
the variance assertion. Each independent scalar Gaussian component has
interval probability at most $2r/\sqrt{2\pi\sigma^2}\le
4r\sqrt{B/(\pi n)}$. The ball event is a subset of the three interval events.
Apply total variation to obtain \eqref{f16v34:local-cutoff}.
\end{proof}

\begin{proposition}[An explicit slowly growing native period process]
\label{f16v34:process}
Fix $B$. For sufficiently large $\gamma$, choose the even integer
\[
 n_\gamma=2\left\lfloor\frac{\sqrt{\log\gamma}}2\right\rfloor.
\]
Linearly interpolate V33's native amplitudes $a_i/\sqrt{n_\gamma}$ at times
$i/n_\gamma$. Under the \emph{full} endpoint-fixed native law this continuous
nine-component process converges weakly to the centered continuous Gaussian
bridge with covariance
\begin{equation}
              2(\min(s,t)-st)I_9,\qquad 0\le s,t\le1.
 \label{f16v34:process-limit}
\end{equation}
For every fixed $r>0$, the native probability
$\boldsymbol\nu_{\gamma,B,n_\gamma}^{00}\{|y_{n_\gamma/2,e}|\le r\}$
tends to zero, as does the all-time window probability. The native midpoint
variance is at least $n_\gamma/(8B)-o(1)$.
\end{proposition}
\begin{proof}
For this choice $Bn_\gamma\log(n_\gamma+1)=o(\log\gamma)$.
Theorem~\ref{f16v34:full-limit} therefore compares the entire increasing-dimensional
native law and its Gaussian with error $\gamma^{-1/2+o(1)}$, including second
moments. Under that Gaussian, the interpolated amplitudes have exactly the law
of the polygonal interpolation of one continuous variance-two Gaussian bridge
at the uniform mesh. On a common realization these polygonal paths converge
uniformly by continuity. Pushforward contracts total variation, so the native
path laws have the same weak limit. This is a probabilistic time parametrization
of a specified limit, not calibrated physical time.

The proof of \eqref{f16v34:local-cutoff} applies to the Gaussian at each
$n_\gamma$, and its right side tends to zero. The total-variation error tends
to zero as well. The all-time window is a subset of this single-time event.
The weighted $p=2$ bound proves the stated absolute $o(1)$ covariance error.
Angular coordinates still satisfy $y/\sqrt\gamma\to0$ in this regime; their
rescaled spreading does not mean escape from a fixed angular neighbourhood.
\end{proof}

\subsection{Positive native derivative and Wilson-word covariance statements}

Retain V33's assembled map
$\mathcal B y=((G^{-1/2}\Delta y_t)_t,(Dy_i/\sqrt3)_i)$.
It obeys $\mathcal B^{\mathsf T}\mathcal B\preceq K^{\rm int}$ and has bounded
operator norm, independent of $B,n$. Define
\[
 \varepsilon_{p,B,n,\gamma}
       =\gamma^{-1/2}\exp\{C_p d_{B,n}\log(n+1)\}.
\]
Constants can be enlarged once when passing to a fixed finite-degree observable.

\begin{theorem}[Native covariance in the derivative bank and the actual word bank]
\label{f16v34:fields}
At the normalizer threshold above,
\begin{equation}
 \operatorname{Cov}_{\boldsymbol\nu_\gamma}(\mathcal B y)
              \preceq(1+\varepsilon_{2,B,n,\gamma})I.
 \label{f16v34:native-derivative}
\end{equation}
In particular the complete mixed bank of temporal differences and magnetic
curvature has the V33 bound with factor $1+\varepsilon_{2,B,n,\gamma}$.
There is no retained radius. The error tends to zero at fixed $B,n$, and in
\eqref{f16v34:joint-regime}; no error uniform in all durations is asserted.

There is also a genuinely nonlinear observable statement on the joint law.
List each original classical Wilson word once, including its original magnetic
weight $\omega_\ell\in\{1/2,1\}$ and all kinetic words. Define its quaternion
vector observable by
\begin{equation}
 \mathcal X_{\gamma,\ell}(w)
   =\sqrt{\gamma\omega_\ell}\,\boldsymbol{\operatorname{Im}}W_\ell(w/\sqrt\gamma).
 \label{f16v34:word-fields}
\end{equation}
Then, with constants of the same explicit dimension-growth type,
\begin{equation}
 \operatorname{Cov}_{\mathbb P_\gamma}(\mathcal X_\gamma)
                         \preceq(1+\varepsilon_{4,B,n,\gamma})I.
 \label{f16v34:word-covariance}
\end{equation}
Every covariance is centered under its indicated native law and includes the
cross covariances of the assembled bank. These word vectors are not claimed
to be gauge-invariant scalar observables; their frames and order are unchanged.
\end{theorem}
\begin{proof}
The weighted second-moment comparison bounds the operator-norm difference of
retained covariances: for a unit vector $v$, $(v^{\mathsf T}y)^2\le|w|^2$.
The Gaussian derivative covariance is at most $I$ by V33. The uniform norm of
$\mathcal B$ and the weighted comparison prove \eqref{f16v34:native-derivative}.
Equivalently, testing its output against $((G^{1/2}a_t)_t,(\sqrt3b_i)_i)$
bounds the native variance by
$(1+\varepsilon_{2,B,n,\gamma})(30\sum_t|a_t|^2+3\sum_i|b_i|^2)$.

Let $D_\ell w$ be the first quaternion-vector differential of word $\ell$
at the identity, and let $\mathcal D$ stack $\sqrt{\omega_\ell}D_\ell$.
The exact quadratic action has the form
\[
 U_0(w)=\tfrac12|\mathcal D w|^2+E_{\rm end}(w),\qquad E_{\rm end}\ge0,
\]
so $\mathcal D^{\mathsf T}\mathcal D\preceq\mathcal H_{B,n}$. Gaussian
integration therefore gives
$\operatorname{Cov}_{\mathbb Q}(\mathcal D w)
=\mathcal D\mathcal H_{B,n}^{-1}\mathcal D^{\mathsf T}\preceq I$.
No word covariance is counted independently of the others in this inequality.

The ordered unitary second-derivative estimate gives the complete-chart bound
$\|\mathcal X_\gamma-\mathcal D w\|\le C\gamma^{-1/2}|w|^2$ by bounded
incidence, and $\|\mathcal X_\gamma\|+\|\mathcal D w\|\le C|w|$.
The full-law moment estimates control its covariance error using moments
through degree four. The normalizer threshold and Gaussian moment bounds
above bound all these moments by the same
$\exp\{C_p d\log(n+1)\}$ type of constant. Combine with the weighted-law
comparison for $\mathcal D w$. This proves \eqref{f16v34:word-covariance}.
Neither a relative lower Hessian for the nonlinear retained action nor a
native diffusion inverse has been assumed.
\end{proof}

\begin{proposition}[A native retained-sampler upper test, not a physical spectrum]
\label{f16v34:gap-test}
Let $\lambda_{\gamma,B,n}^{\rm br}$ be the variational gap of the sum of
rate-one \emph{full} retained-bridge heat baths under
$\boldsymbol\nu_{\gamma,B,n}^{00}$. At fixed $B,n$,
\begin{equation}
       \limsup_{\gamma\to\infty}\lambda_{\gamma,B,n}^{\rm br}
                         \le30\sin^2\frac\pi{2n}.
 \label{f16v34:native-gap-upper}
\end{equation}
On the sequence of Proposition~\ref{f16v34:process}, the gap is $O(n_\gamma^{-2})$.
This is not an assertion of update-operator convergence.
\end{proposition}
\begin{proof}
Use V33's unit first temporal period eigenvector $v$ and the linear test
$f=v^{\mathsf T}y$. Its Gaussian variance is
$[2\sin^2(\pi/(2n))]^{-1}$. At a retained site $u$, use as an admissible
\emph{native} exterior predictor the exact Gaussian conditional mean
$g_u(y_{u^c})=\mathbb E_Q[f\mid y_{u^c}]$. The minimizing property of native
conditional expectation gives
\[
 \mathbb E_{\boldsymbol\nu_\gamma}\operatorname{Var}(f\mid y_{u^c})
 \le \mathbb E_{\boldsymbol\nu_\gamma}|f-g_u|^2.
\]
These are quadratic polynomials. Their Gaussian sum is
$\sum_uv_u^{\mathsf T}(K^{\rm int}_{uu})^{-1}v_u\le15$, using V33's
$K^{\rm int}_{uu}\succeq I_3/15$. The weighted convergence gives the fixed
regulator limsup after division by the convergent positive variance.
For the simultaneous assertion, the coefficient operator of
$\sum_u|f-g_u|^2$ is uniformly bounded: the diagonal inverses are at most
15 and the rows of $K^{\rm int}$ are uniformly bounded. Hence the weighted
error still tends to zero in \eqref{f16v34:joint-regime}. The denominator is
$[2\sin^2(\pi/(2n_\gamma))]^{-1}+o(1)$, proving the stated order.
This is a variational upper estimate using fixed Gaussian predictors, not
conditional-kernel or spectral convergence inferred from total variation.
\end{proof}

\subsection{What has been removed and what still requires new information}

V34 pays the complete compact tail and the actual full normalizer at fixed
endpoint data. The Gaussian and native laws now agree in a two-sided,
untruncated comparison at fixed regulator and along the explicitly slow
simultaneous regime. V33's merely one-sided cutoff argument is thereby
strengthened, not withdrawn. The new local midpoint statement was not a
consequence of that earlier nested-window upper bound alone.

The rate has a large dimension price and a duration-dependent coercivity scale.
Neither disappears from the theorem. In particular the positive native field
covariance statement is not the sought unrestricted-duration relative
curvature estimate. The next useful target is to reduce the endpoint/dimension
loss specifically in the derivative bank, or keep the nonlinear compact period
sector nonperturbatively while controlling its conditional complement. Another
fixed-window boundary gap would not address either task.

The two supplied memoranda distinguish perturbative gain from locality and
normalization costs, and an auxiliary functional inequality from its physical
transfer implication. Here the Taylor power is $\gamma^{-1/2}$, while the
explicit cost is $\exp\{C_p Bn\log(n+1)\}$. Those ledgers are not conflated.
All root gauge and physical endpoint issues remain: the law is the original
\emph{unprojected} identity-endpoint history, not a stationary physical vacuum.
A long-history bridge-sampler upper bound does not prove a gapless physical
Hamiltonian; even massive physical transfer bridges can have algorithmic slow
modes. No physical time is set by the interpolation in
\eqref{f16v34:process-limit}.

\subsection{Verification and preservation}

The accompanying diagnostic is
\begin{center}\small\path{verify_ym_f16_v34_endpoint_cutoff.py}.\end{center}
It checks exact group-word rearrangements, the literal bridge-degree and frame
energy counts, scalar metric-path bounds, Haar endpoint exponents, Gaussian
normalizers, weighted-density algebra, period covariances and variational
predictors. The new script passes 259 exact assertions. All 30 supplied V4--V33 checkers
were rerun unchanged and passed. These are finite rational and symbolic
fixtures, not native Wilson integral enclosures or formal verification. Actual
outputs and hashes are recorded separately.

Only the single title version and this terminal insertion differ from V33.
Deleting the insertion and restoring that title recovers the supplied V33 byte
for byte. No old formula, source, proof or reference is silently edited. The
compact pin uses the inherited finite cube and frame estimates; this is not an
independent recertification of the complete archive. The next checkpoint is V35.

\begin{firewall}
V34 proves full endpoint-anchored compact coercivity, a quantitative normalized
joint native Gaussian comparison without a retained cutoff, and polynomial
moment convergence at fixed regulator and in the stated slow simultaneous
regime. It obtains native derivative and Wilson-word covariance comparisons,
a quantitative period process, local rescaled-window escape and a native
retained-sampler upper test. It does not prove uniformity at arbitrary spatial
volume or duration, a stationary gauge-invariant vacuum limit, unrestricted
relative native curvature, a physical transfer gap, or an interacting continuum
mass gap. All complete integration domains, original Haar powers, endpoint
normalizers and sampler rates used here are retained.
\end{firewall}
% END F16 V34 APPEND

% BEGIN F16 V35 APPEND -- 2026-09-18
% Insert immediately before the final document-ending command of V34.
% No predecessor assertion, native measure, endpoint factor, or clock is changed.

\clearpage
\section{V35: polynomial-regime native comparison and relative field covariance}
\label{f16v35:context}

V34 proves a complete normalized comparison, but charges its coarse radial
majorant throughout the integration domain. The resulting exponential
size factor limits its simultaneous application to logarithmically growing
histories. We separate two tasks which need different estimates. The coarse
compact pin pays only the tails. On an explicitly chosen interior ball, the
actual measured Hessian is compared with the \emph{full endpoint-fixed}
Gaussian precision. A normalized square-root-density estimate then measures
the correction by its score, rather than by the absolute oscillation of an
extensive action. This gives polynomially growing native histories and a
relative covariance estimate, including the temporal soft directions.

The comparison ball is a proof device, not a new retained window or a change
to V4's core event. Its complement is paid under both complete laws. The
endpoint and size conditions remain substantive. In particular, the new
covariance comparison does not give a native pointwise relative Hessian bound,
a physical transfer comparison, or a duration-uniform mass.

\subsection{V35 checkpoint and unchanged probability}

The complete supplied V34 appendix and audit were read. The new argument uses
its compact endpoint pin, V5's independent fast anchor, the actual ordered-word
derivative bounds, and V33's finite Gaussian matrices. It does not use the
V16 moment iteration, V30 spatial-screening theorem, or V31's radius-dependent
constants. Those results remain separate, with their original hypotheses.
The companion checkpoint records targeted inventories of the current lineage,
f14 V100, f15 V100, and Grand Tensor V076; it is not independent certification
of their proofs or a claim that unavailable companions were examined. The
next scheduled checkpoint is V40.

The interior/tail split is already present in V5's finite-box Laplace proof.
The monograph also uses normalized Hellinger embeddings. Neither generic
principle is claimed new. The additional step here is a score estimate in the
endpoint Gaussian metric, together with an explicitly removed comparison ball.
For the boundary-aware variance principle, see Kolesnikov--Milman,
\href{https://arxiv.org/abs/1310.2526}{\emph{Brascamp--Lieb type inequalities on
weighted Riemannian manifolds with boundary}}, Theorem 1.2(1) at $N=\infty$.
Its boundary statement and Reilly formula were checked; the special estimate
used below is also proved directly. Spokoiny's
\href{https://arxiv.org/abs/2204.11038}{\emph{Dimension free non-asymptotic bounds
on the accuracy of high dimensional Laplace approximation}} was screened for
context only, not used as a black-box native theorem.

Fix identity retained endpoints, $B=m^3$, $m\ge4$, $n\ge2$. Keep exactly V34's
joint probability $\mathbb P_\gamma$, its Gaussian $\mathbb Q$, and their full
normalizers $\mathcal Z_\gamma,\mathcal Z_0$. Write
\[
 w=(\zeta,y)\in\mathbb R^d,\qquad d=(72n-21)B,\qquad
 U_0(w)=\tfrac12 w^{\mathsf T}\mathcal H w,\qquad \mathcal C=\mathcal H^{-1}.
\]
Thus $\mathcal H$ is the full \emph{joint} precision, not the fast conditional
precision $H$, nor the retained Schur precision $K^{\rm int}$. All chord
endpoint Haar powers are $1/2$; all other free Haar powers are one. All fast
and interior bridge groups remain integrated over their complete principal
charts. Scalar dressed-end factors cancel only in the same endpoint-fixed
probability, not from a proposed physical transfer normalization.

Shrink V34's fixed constant $c$ and enlarge $H_*$ once so that
$0<c\le1/4$, $H_*\ge1$. Put $a=c/n^2$. Its global estimates are
\begin{equation}
 U_\gamma(w),U_0(w)\ge a|w|^2,\qquad
 2aI\preceq\mathcal H\preceq H_*I,\qquad 0\le\mathcal J_\gamma\le1.
 \label{f16v35:global-data}
\end{equation}
The first bound for $U_\gamma$ is only on its complete principal product, as
in V34; the native density is zero off that product.

\subsection{Local Gaussian scales and a relative interior Hessian}

\begin{proposition}[Gaussian coordinate scales and local measured derivatives]
\label{f16v35:local-data}
There are fixed $C_2,C_H,C_G<\infty$, independent of $B,n,\gamma$, such that
\begin{equation}
 \|\mathcal C\|\le(2a)^{-1},\qquad
 \mathcal C_{\zeta\zeta}\preceq(2a_*)^{-1}I,\qquad
 \mathcal C_{yy}\preceq30(L_n^{-1}\otimes I),\qquad
 \sup_g\|\mathcal C_{gg}\|\le C_G n.
 \label{f16v35:Gaussian-scales}
\end{equation}
Here $g$ labels an original three-component group and $a_*=(6\pi^2)^{-1}$.
On $|w|\le R\le\sqrt\gamma/2$, set
\[
 V_\gamma=U_\gamma-\log\mathcal J_\gamma,\qquad W_\gamma=V_\gamma-U_0.
\]
Then
\begin{equation}
 \|D^2W_\gamma(w)\|\le C_2R/\sqrt\gamma+C_H/\gamma,
 \qquad
 |\nabla W_\gamma(w)|^2\le
 \frac{C_2}{\gamma}\sum_g|w_g|^4+
 \frac{C_H}{\gamma^2}|w|^2.
 \label{f16v35:local-derivatives}
\end{equation}
The constants account for the fixed local carriers and their actual incidence,
not for independent words or independent coordinates.
\end{proposition}
\begin{proof}
The first bound is \eqref{f16v35:global-data}. To prove the fast marginal bound,
minimize the quadratic $U_0$ in $y$ at fixed $\zeta$. V5's anchor gives
$\inf_y U_0\ge a_*|\zeta|^2$, so its fast marginal precision is at least
$2a_*I$. This is a marginal statement, not a substitution of the conditional
fast covariance. The retained marginal is exactly V33's Gaussian, whose
precision is at least $(L_n\otimes I)/30$. The diagonal of $L_n^{-1}$ is
$i(n-i)/n\le n/4$. These two facts give the group bound, including the
unmoving temporal ports and the endpoint chord groups.

For a scalar derivative in one group, only a fixed number of ordered words
contribute. V34's unitary third-derivative argument gives a Hessian remainder
at most $C\gamma^{-1/2}|w_T|$ and a gradient remainder at most
$C\gamma^{-1/2}|w_T|^2$ on each such fixed carrier $T$. The zero and linear
Taylor terms vanish; the endpoint quadratic is exact. The block-row Schur
bound, fixed carrier size, and bounded row and column incidence prove the first
classical estimate and
$|\nabla(U_\gamma-U_0)|^2\le C\gamma^{-1}\sum_g|w_g|^4$.
This retains the local sum instead of replacing it by $|w|^4$.

On angles at most $1/2$, the measured Haar Hessian is positive and bounded
above by $C_H/\gamma$ on each three-group, with gradient bounded by
$C_H|w_g|/\gamma$. The endpoint exponent $1/2$ is included. Increasing the
fixed constants proves \eqref{f16v35:local-derivatives}. No logarithm of the
Haar density is used near its zero at the original principal cut.
\end{proof}

Fix $s\ge1$, $\gamma\ge2$, and define the explicit tail budget and radius
\begin{equation}
 \begin{split}
 \Lambda_s={}&1+\frac d2\log\frac{2H_*}{a}
       +\log\left[1+\left(\frac{d+6}{a}\right)^4\right]
       +(s+1)\log\gamma,\qquad R_s^2=\frac{2\Lambda_s}{a},\qquad
 \mathcal B_s=\{|w|\le R_s\}.
 \end{split}
 \label{f16v35:radius}
\end{equation}
The finite-parameter entry conditions used below are
\begin{equation}
 R_s\le\sqrt\gamma/2,\qquad
 C_2R_s/\sqrt\gamma+C_H/\gamma\le a.
 \label{f16v35:entry}
\end{equation}
They are checkable in the inherited fixed constants. They imply
\begin{equation}
       \tfrac12\mathcal H\preceq D^2V_\gamma\preceq2\mathcal H
            \quad\hbox{on }\mathcal B_s.
 \label{f16v35:interior-curvature}
\end{equation}
Indeed a symmetric error of norm at most $a$ is at most $\mathcal H/2$ in
absolute quadratic form. Both potentials and their gradients vanish at zero,
so $V_\gamma\le2U_0$ on this ball. This is an \emph{interior} comparison;
no such native Hessian inequality is asserted on the full principal product.

\begin{proposition}[Complete normalized tails, paid before the interior comparison]
\label{f16v35:tails}
Under \eqref{f16v35:entry},
\begin{equation}
 \mathcal Z_\gamma\ge2^{-d/2-1}\mathcal Z_0,
 \qquad
 \mathbb E_{\mathbb P_\gamma}[(1+|w|^8)1_{\mathcal B_s^c}]
 +\mathbb E_{\mathbb Q}[(1+|w|^8)1_{\mathcal B_s^c}]
 \le2\gamma^{-s}.
 \label{f16v35:tails-bound}
\end{equation}
The normalizer inequality is only a lower bound. In particular it is not a
claim that $\mathcal Z_\gamma/\mathcal Z_0$ is close to one uniformly in the
new growing regime.
\end{proposition}
\begin{proof}
The Gaussian of precision $2\mathcal H$ has
$\mathbb E|w|^2\le d/(4a)$. The radius \eqref{f16v35:radius} has
$R_s^2\ge d/(2a)$, so at least half this Gaussian lies in $\mathcal B_s$.
The ball lies strictly inside every original principal group chart.
Since $V_\gamma\le2U_0$ there,
\[
 \mathcal Z_\gamma\ge\int_{\mathcal B_s}e^{-2U_0}
           \ge2^{-d/2-1}\mathcal Z_0.
\]
This uses the actual full native mass and its original Haar powers.

For either tail, use the global classical pin and $\mathcal J_\gamma\le1$;
off the native chart the native density remains zero. For $t=|w|>R_s$,
$e^{-at^2}\le e^{-aR_s^2/2}e^{-at^2/2}$. The exact radial eighth moment gives
\[
 \int(1+|w|^8)e^{-a|w|^2/2}\,dw
 \le(2\pi/a)^{d/2}\left[1+\left((d+6)/a\right)^4\right].
\]
Also $\mathcal Z_0\ge(2\pi/H_*)^{d/2}$. Dividing by the proved native mass
lower bound bounds its weighted tail by
\[
 2\exp\left\{\frac d2\log(2H_*/a)
 +\log[1+((d+6)/a)^4]-\frac{aR_s^2}{2}\right\}
 =2e^{-1}\gamma^{-(s+1)}\le\gamma^{-s}.
\]
The Gaussian needs no factor $2^{d/2+1}$ and obeys the same weaker bound.
Their sum proves the result. All eighth-degree tail weights are explicit;
the size loss has been paid only in choosing the radius, not multiplied into
the later interior error.
\end{proof}

\subsection{Normalized square roots instead of a global likelihood-ratio bound}

We record the needed covariance and normalization calculation on a smooth
Euclidean ball. It does not use a product-ball one-form locality claim.

\begin{proposition}[A convex-domain variance bound and a normalized Hellinger estimate]
\label{f16v35:Hellinger}
Let $D$ be a bounded smooth convex domain, and let $\mu\propto e^{-V}$ be smooth
and positive there, with $D^2V\succeq A\succ0$. Then
\begin{equation}
       \operatorname{Var}_\mu f
          \le\mathbb E_\mu[\nabla f^{\mathsf T}A^{-1}\nabla f].
 \label{f16v35:convex-Poincare}
\end{equation}
If $Q_D\propto e^{-w^{\mathsf T}\mathcal H w/2}$ and
$P_D\propto e^{-w^{\mathsf T}\mathcal H w/2-W(w)}$ are two normalized laws on
that same domain, put
$h^2(P_D,Q_D)=\int(\sqrt{p_D}-\sqrt{q_D})^2$. Then
\begin{equation}
 \boxed{\quad h^2(P_D,Q_D)
       \le\tfrac12\mathbb E_{P_D}
                   [\nabla W^{\mathsf T}\mathcal H^{-1}\nabla W],
       \qquad d_{\rm TV}(P_D,Q_D)\le h(P_D,Q_D).\quad}
 \label{f16v35:Hellinger-bound}
\end{equation}
This statement includes both normalizers. It does not require an estimate
that their ratio is near one, or a small supremum norm of $W$.
\end{proposition}
\begin{proof}
For centered smooth $f$, solve the weighted Neumann Poisson equation
$\mathcal Au=f$. The integrated Bochner identity has the nonnegative outward
second-fundamental-form term on $\partial D$ and implies
$\|f\|_2^2\ge\mathbb E_\mu[\nabla u^{\mathsf T}A\nabla u]$.
Integration by parts and Cauchy--Schwarz in the constant metric $A$ give
\[
 \|f\|_2^2=\mathbb E_\mu[\nabla f\cdot\nabla u]
 \le\bigl(\mathbb E_\mu[\nabla f^{\mathsf T}A^{-1}\nabla f]\bigr)^{1/2}
                         \|f\|_2.
\]
This proves \eqref{f16v35:convex-Poincare}. At a fixed finite dimension the
smooth positive weight on the compact smooth ball gives the usual Neumann
realization and Poisson inverse on constants' complement; equivalence of
Sobolev norms is used only for this domain statement. Smooth approximation
extends the inequality to $H^1$. The boundary term was retained before its
nonnegative sign was used.

Apply this inequality under $Q_D$ to $b=\sqrt{p_D/q_D}$.
One has $\mathbb E_{Q_D}b^2=1$ and
$\nabla b=-b\nabla W/2$. If $m=\mathbb E_{Q_D}b$, then $0<m\le1$ and
\[
 h^2=2(1-m)\le2(1-m^2)=2\operatorname{Var}_{Q_D}b
       \le\tfrac12\mathbb E_{P_D}[\nabla W^{\mathsf T}\mathcal C\nabla W].
\]
The density ratio in $b$ is \emph{normalized}; its scalar derivative is zero
in $w$, not omitted from its definition. Cauchy--Schwarz applied to
$(\sqrt p-\sqrt q)(\sqrt p+\sqrt q)$ gives $d_{\rm TV}\le h$ in this convention.
\end{proof}

Let $P_s=\mathbb P_\gamma(\cdot\mid\mathcal B_s)$ and
$Q_s=\mathbb Q(\cdot\mid\mathcal B_s)$. The subscript here labels the tail
budget, not an interpolation of the native couplings. Simultaneous global
conjugation acts by a common orthogonal rotation on all three-groups and
preserves both balls and both densities. Thus $\mathbb E_{P_s}w=
\mathbb E_{Q_s}w=0$, as also holds for the two complete laws.

\begin{proposition}[Directional moments and the local-score estimate]
\label{f16v35:score}
Under \eqref{f16v35:entry}, for every real coefficient vector $u$,
\begin{equation}
 \mathbb E_{P_s}e^{t u^{\mathsf T}w}
       \le e^{t^2u^{\mathsf T}\mathcal C u},\qquad
 \mathbb E_{P_s}|u^{\mathsf T}w|^4
       \le C(u^{\mathsf T}\mathcal C u)^2.
 \label{f16v35:directional-moments}
\end{equation}
The same weaker bounds hold under $Q_s$. In particular
\begin{equation}
 \mathbb E_{P_s}[\nabla W_\gamma^{\mathsf T}\mathcal C\nabla W_\gamma]
                   \le C\frac{d n^4}{\gamma}.
 \label{f16v35:score-cost}
\end{equation}
\end{proposition}
\begin{proof}
Every linear tilt of $P_s$ has the same Hessian lower bound
$\mathcal H/2$ and the same convex ball. Apply
\eqref{f16v35:convex-Poincare} to $u^{\mathsf T}w$ under that tilted law.
The second derivative of its log moment-generating function is at most
$2u^{\mathsf T}\mathcal C u$. Its value and first derivative at zero are zero,
which proves the first estimate by two integrations. Exponential Markov and
integration of its two-sided Gaussian tail give the fourth-moment estimate.
This does not assume that the coordinates are independent. The argument for
$Q_s$ has the stronger Hessian $\mathcal H$.

The local Gaussian variance in \eqref{f16v35:Gaussian-scales} therefore gives
$\mathbb E_{P_s}|w_g|^4\le C n^2$ for each fixed three-group, and
$\mathbb E_{P_s}|w|^2\le Cdn$. Use the \emph{local} sum in
\eqref{f16v35:local-derivatives} to get
$\mathbb E_{P_s}|\nabla W_\gamma|^2\le Cdn^2/\gamma$.
Finally $\|\mathcal C\|\le Cn^2$ proves \eqref{f16v35:score-cost}.
The $n^2$ inverse cost is real; it has not been replaced by a positive mass.
\end{proof}

\subsection{Polynomially growing complete native histories}

Set
\[
            e_{d,n,\gamma,s}=\sqrt{\frac{dn^4}{\gamma}}+\gamma^{-s}.
\]

\begin{theorem}[Full native total variation with a polynomial size cost]
\label{f16v35:TV}
Under the explicit entry conditions \eqref{f16v35:entry},
\begin{equation}
 \boxed{\qquad d_{\rm TV}(\mathbb P_\gamma,\mathbb Q)
               \le C e_{d,n,\gamma,s}.\qquad}
 \label{f16v35:TV-bound}
\end{equation}
The same bound holds for their complete retained marginals and every measurable
pushforward. No proof cutoff remains in these laws.
There are fixed positive constants such that, for each fixed $s$, a sufficient
large-coupling entry condition is
\begin{equation}
 n^6\{d\log(n+1)+(s+1)\log\gamma\}\le c_0\gamma,
 \label{f16v35:polynomial-entry}
\end{equation}
with $c_0$ sufficiently small in the named fixed geometric constants.
In particular, the estimate tends to zero along any sequence with
\begin{equation}
 \frac{Bn^7\log(n+1)+n^6\log\gamma}{\gamma}\longrightarrow0.
 \label{f16v35:polynomial-regime}
\end{equation}
If $B=O(\gamma^b)$ and $n=O(\gamma^\alpha)$ with $b,\alpha\ge0$,
$b+7\alpha<1$ is a sufficient strict exponent condition.
\end{theorem}
\begin{proof}
Propositions~\ref{f16v35:Hellinger} and \ref{f16v35:score} bound
$d_{\rm TV}(P_s,Q_s)$ by $C\sqrt{dn^4/\gamma}$. For any probability $\mu$,
\[
 d_{\rm TV}(\mu,\mu(\cdot\mid\mathcal B_s))=\mu(\mathcal B_s^c).
\]
Use this identity twice and \eqref{f16v35:tails-bound} to obtain
\eqref{f16v35:TV-bound}. The principal-chart complement is included in those
tails. This is not a comparison of conditioned laws relabelled as full laws.

For the entry condition note that $a=c/n^2$ and
\[
 \Lambda_s\le C_s\{d\log(n+1)+(s+1)\log\gamma\}.
\]
Here $\log(d+6)\le C d$; there is no hidden exponential in $d$. Hence
\[
 R_s\le C_s n\sqrt{d\log(n+1)+(s+1)\log\gamma}.
\]
A sufficiently small
constant in \eqref{f16v35:polynomial-entry} ensures
$C_2R_s/\sqrt\gamma\le a/2$, $C_H/\gamma\le a/2$, and
$R_s\le\sqrt\gamma/2$. Finite fixed thresholds cover the other elementary
requirements. With $d=(72n-21)B$, \eqref{f16v35:polynomial-regime} implies this
entry and $dn^4/\gamma\to0$. The strict exponent example follows directly,
including rounded admissible spatial sides and bounded time lengths.
\end{proof}

The gain is not a claim of an optimal high-dimensional Laplace rate. One still
pays $n^2$ in the inverse and a coarse $n^{-2}$ pin in preparing the comparison
ball. The improvement over V34 is that its exponential radial ratio now
occurs inside the tail logarithm, whereas the normalized interior score pays
$\sqrt d$. No higher perturbative order or convergence of a series is used.

\begin{theorem}[Relative covariance in the original Gaussian metric]
\label{f16v35:covariance}
Let $\Sigma_\gamma=\operatorname{Cov}_{\mathbb P_\gamma}(w)$. Under the same
entry conditions,
\begin{equation}
 \boxed{\quad
 \|\mathcal H^{1/2}\Sigma_\gamma\mathcal H^{1/2}-I\|_{
 \mathrm{op}}\le C e_{d,n,\gamma,s}.\quad}
 \label{f16v35:relative-covariance}
\end{equation}
In particular, with $\eta=C e_{d,n,\gamma,s}$,
\begin{equation}
 (1-\eta)\mathcal C\preceq\Sigma_\gamma\preceq(1+\eta)\mathcal C.
 \label{f16v35:Loewner}
\end{equation}
The retained marginal obeys the same comparison with
$(K^{\rm int})^{-1}$. This is a \emph{covariance} conclusion; it is not a
pointwise lower bound on the full native Hessian outside $\mathcal B_s$.
\end{theorem}
\begin{proof}
Fix $u$ with $u^{\mathsf T}\mathcal C u=1$, and set $X=u^{\mathsf T}w$.
The fourth moments of $X$ under $P_s,Q_s$ are bounded by a numerical constant
by \eqref{f16v35:directional-moments}. Factor the density difference into
square roots and use Cauchy--Schwarz:
\[
 |\mathbb E_{P_s}X^2-\mathbb E_{Q_s}X^2|
 \le h(P_s,Q_s)
       \{2\mathbb E_{P_s}X^4+2\mathbb E_{Q_s}X^4\}^{1/2}
 \le C\sqrt{dn^4/\gamma}.
\]
Moreover $|u|^2\le H_*$, since $\mathcal C\succeq H_*^{-1}I$.
Equation~\eqref{f16v35:tails-bound} pays the missing $X^2$ integrals by
$C\gamma^{-s}$; the factors $P(\mathcal B_s)$ and $Q(\mathcal B_s)$ in the
interior expectations differ from one by the same bound. All four means are
zero by simultaneous global conjugation (for the Gaussian, also by reflection).
Thus $|\operatorname{Var}_{\mathbb P_\gamma}X-1|\le C e_{d,n,\gamma,s}$.
The estimate is uniform over every such $u$, proving both matrix statements.
Taking principal covariance blocks gives the retained assertion. No factor
counting the number of tested directions is introduced.
\end{proof}

\begin{theorem}[The same untruncated native difference and Wilson-word banks]
\label{f16v35:fields}
Retain V33's assembled map $\mathcal B$ and V34's actual nonlinear bank
$\mathcal X_{\gamma,\ell}=\sqrt{\gamma\omega_\ell}\,
\boldsymbol{\operatorname{Im}}W_\ell$, with every original word weight.
If $e_{d,n,\gamma,s}\le1$, then
\begin{equation}
 \begin{split}
 \operatorname{Cov}_{\boldsymbol\nu_\gamma}(\mathcal B y)
                     &\preceq(1+C e_{d,n,\gamma,s})I,\\
 \big\|\operatorname{Cov}_{\mathbb P_\gamma}(\mathcal X_\gamma)
       -\mathcal D\mathcal C\mathcal D^{\mathsf T}\big\|_{\rm op}
                     &\le C e_{d,n,\gamma,s}.
 \end{split}
 \label{f16v35:field-bounds}
\end{equation}
Here $\mathcal D$ is the original first-differential word bank from V34;
$\mathcal D\mathcal C\mathcal D^{\mathsf T}\preceq I$. Both estimates include
all cross covariances and refer to the complete native laws, without a retained
window. No automatic gauge invariance of the quaternion-vector words is claimed.
\end{theorem}
\begin{proof}
Use \eqref{f16v35:Loewner} and the unchanged inequalities
$\mathcal B^{\mathsf T}\mathcal B\preceq K^{\rm int}$ and
$\mathcal D^{\mathsf T}\mathcal D\preceq\mathcal H$.
The first gives the first assertion, including the bound
\[
 \operatorname{Var}_{\boldsymbol\nu_\gamma}
 (\textstyle\sum_t a_t^{\mathsf T}\Delta y_t+\sum_i b_i^{\mathsf T}Dy_i)
 \le(1+C e_{d,n,\gamma,s})
                (30\sum_t|a_t|^2+3\sum_i|b_i|^2).
\]
The second compares the \emph{linear} word-bank covariance in operator norm.
For the nonlinear bank, retain the local-carrier version of V34's Taylor bound:
\[
 \|\mathcal X_\gamma-\mathcal D w\|^2
                    \le\frac C\gamma\sum_g|w_g|^4.
\]
It holds on the full principal charts by ordered unitary differentiation and
bounded incidence. On $\mathcal B_s$, the expected sum is at most $Cdn^2$.
Off that ball, its expectation is at most
$\mathbb E_{\mathbb P_\gamma}|w|^4 1_{\mathcal B_s^c}\le\gamma^{-s}$.
Thus the vector $L^2$ error is at most $C\sqrt{dn^2/\gamma}+C\gamma^{-s/2}/
\sqrt\gamma$, which is bounded by $C e_{d,n,\gamma,s}$: the second term is already
dominated by $\sqrt{dn^4/\gamma}$ for $s\ge1$, $d\ge1$.
Centering contracts $L^2$. For each unit output vector the linear bank has
variance at most $1+C e_{d,n,\gamma,s}$; Cauchy--Schwarz pays both mixed
terms and the squared error. This proves the second assertion. No observable
is replaced before the nonlinear error is paid.
\end{proof}

\subsection{Algebraically long native period bridges and the sampler distinction}

\begin{proposition}[Period limits at polynomial durations]
\label{f16v35:periods}
Fix $B$ and $0<\alpha<1/7$, and let
$n_\gamma=2\lfloor\gamma^\alpha/2\rfloor$. Under the complete native law,
linearly interpolate V33's nine period amplitudes $a_i/\sqrt{n_\gamma}$ at
$i/n_\gamma$. Their laws converge weakly in $C([0,1];\mathbb R^9)$ to the
centered Gaussian bridge with covariance $2(\min(s,t)-st)I_9$.
For each fixed $r>0$ and specified midpoint bridge,
\begin{equation}
 \boldsymbol\nu_{\gamma,B,n_\gamma}^{00}(|y_{n_\gamma/2,e}|\le r)\to0,
 \qquad
 \operatorname{Var}_{\boldsymbol\nu_\gamma}(y_{n_\gamma/2,e}^a)
               \ge(1-o(1))\frac{n_\gamma}{8B}.
 \label{f16v35:midpoint}
\end{equation}
In the general entry regime, whenever $\eta=C e_{d,n,\gamma,s}<1$, the native
rate-one full retained-bridge sampler satisfies
\begin{equation}
 \lambda_{\gamma,B,n}^{\rm br}
      \le\frac{1+\eta}{1-\eta}\,30\sin^2\frac\pi{2n}.
 \label{f16v35:sampler-test}
\end{equation}
These are neither a physical-time calibration nor spectral convergence inferred
from total variation.
\end{proposition}
\begin{proof}
For this sequence \eqref{f16v35:polynomial-regime} holds, and
$e_{d,n_\gamma,\gamma,s}=O(\gamma^{-(1-5\alpha)/2})+\gamma^{-s}$.
The complete Gaussian period factor and its normalization are exactly V33's;
they are not recomputed or replaced here. Its polygonal rescaling is distributed
as uniform-mesh sampling of a continuous variance-two Gaussian bridge.
On a common realization these interpolations converge uniformly. Total
variation contracts under the measurable projection and interpolation at each
$\gamma$, so Theorem~\ref{f16v35:TV} transfers the weak process limit.
The limiting path law is not asserted to be a total-variation limit of
finite-dimensional polygonal path laws.

The Gaussian midpoint ball bound from V34 is
$[\min(1,4r\sqrt{B/(\pi n)})]^3$. Add the new total-variation error; both
terms vanish. The reducing period covariance contributes $n/(8B)$ to every
scalar midpoint bridge. The retained part of \eqref{f16v35:Loewner} therefore
gives the variance inequality. The all-time fixed window is a subset of this
midpoint event and also has probability tending to zero.

For \eqref{f16v35:sampler-test}, reuse V34's first temporal period test $f$ and
its exact Gaussian exterior predictors $g_u$. The native conditional-variance
minimum is at most $\mathbb E_{\boldsymbol\nu_\gamma}|f-g_u|^2$.
Each difference is linear, so the \emph{relative} retained covariance bound
gives a summed energy at most $(1+\eta)15$, with no accumulated coefficient
error. The native test variance is at least
$(1-\eta)[2\sin^2(\pi/(2n))]^{-1}$. Division proves the displayed bound.
No native update kernel is identified with its Gaussian counterpart.
\end{proof}

\subsection{Checkpoint: which costs have changed and which have not}

V34's full-density error $\gamma^{-1/2}\exp\{C_p d\log(n+1)\}$ remains a
valid predecessor estimate. The present result replaces it, for total variation
and the stated Gaussian-normalized covariance tests, by a polynomial error
under a new explicit entry condition. For example $n\asymp\gamma^{1/8}$ at
fixed $B$ is now allowed, with error $O(\gamma^{-3/16})$ in these norms.
This is not a claim that all \emph{unscaled} polynomial-weighted $L^1$ errors
vanish on every such sequence. Higher raw moments carry their actual Gaussian
scales; only the tested normalized moments and field banks were estimated here.

There is also a normalization boundary. The square-root argument compares two
already normalized interior densities. Its gradient does not see an additive
constant in the action. Consequently normalized-law closeness does not by itself
supply an $o(1)$ absolute pressure correction or a physical Perron-normalizer
comparison in the new growing regime. V34's fixed-regulator partition conclusion
is unchanged. The lower mass bound in \eqref{f16v35:tails-bound} is used only to
pay complete tails before normalization; it is not promoted to a sharp pressure
asymptotic.

The new rate still pays the temporal soft eigenvalue in its score inverse and
in its global comparison radius. A sharper derivative-relative score estimate,
or a treatment of nonlinear compact periods that does not expand their entire
history at once, is the next plausible improvement. The exponent $1/7$ is a
sufficient threshold of this proof, not a physical threshold or an optimality
claim. Neither an additional fixed-window boundary gap nor another temporal
power removes these costs.

All conclusions concern the original unprojected identity-endpoint history.
The prior fixed-bridge width theorems and arbitrary-fast-pinning problem retain
their scopes. Root Gauss reduction, physical endpoint/vacuum normalization,
independent physical time, unrestricted regulators, and interacting continuum
reconstruction remain separate. The supplied import memoranda distinguish
power gain from locality/normalization cost and keep the final physical-transfer
arrow conditional; this checkpoint retains both distinctions.

\subsection{Verification and preservation}

The diagnostic is
\begin{center}\small\path{verify_ym_f16_v35_polynomial_comparison.py}.\end{center}
It tests normalized square-root identities, a nonzero Neumann boundary term,
local-incidence derivative sums, Gaussian marginal scales and noncommuting
matrix comparisons, exact tail arithmetic, the relative covariance handoff,
and the retained Gaussian-predictor upper test. These are finite algebraic
checks, not native compact-integral enclosures or a formal proof of the analytic
lemmas. The new diagnostic passes 254 exact assertions. All 31 supplied
V4--V34 checkers were rerun unchanged and passed; their actual outputs and
hashes are reported separately from inherited reports.

Only the literal title version and this terminal insertion change V34.
Removing them recovers the supplied predecessor byte for byte. The original
sources and companion inventories remain unchanged. The analytic dependencies
are explicitly the compact-pin, finite-word, and Gaussian-matrix inputs above;
no independent certification of the full archive is claimed. The V35 companion
checkpoint records review depth and the next checkpoint V40.

\begin{firewall}
V35 gives a polynomial-regime total-variation comparison of the complete
endpoint-fixed native law, a relative native covariance estimate in its actual
Gaussian metric, and complete nonlinear field-bank bounds. It extends the
native period process and sampler upper test to specified algebraic durations.
The comparison ball is removed with quantified tails; its interior convexity
is not asserted globally. No unrestricted-duration estimate, sharp growing-regime
pressure asymptotic, physical Gauss or vacuum identification, calibrated
physical transfer gap, or interacting continuum mass gap is proved.
\end{firewall}
% END F16 V35 APPEND

% BEGIN F16 V36 APPEND -- 2026-09-18
% Insert immediately before the final document-ending command of V35.
% No predecessor proof, probability, scalar normalization, or clock is edited.

\clearpage
\section{V36: absolute normalization, cubic cancellation, and even-observable accuracy}
\label{f16v36:context}

V35 deliberately compares normalized probabilities. Its square-root score cannot
see an additive action constant, so its improved growing-regime estimate does
not determine the absolute partition ratio. We address that separate scalar
problem before trying to extrapolate another sampling gap. The action is already
anchored at the identity, and its leading cubic correction is odd under global
Euclidean reflection. Its Gaussian mean vanishes, but its variance does not.
Keeping that variance and the actual Haar correction gives the first absolute
normalization coefficient, with a remainder controlled in V35's polynomial
regime. The same parity improves the relative covariance estimate for linear
coordinates without claiming a better total-variation rate for arbitrary tests.

All statements concern the original identity-endpoint, unprojected joint history.
The comparison ball will again be removed. No physical Perron normalization,
root Gauss reduction, unrestricted regulator limit, or mass gap is inferred.

\subsection{Dependencies, ownership, and the fixed scalar convention}

The V35 appendix, analytic audit, and checkpoint were read completely. We use
its explicit comparison ball, complete tail estimates, full joint Gaussian
metric, and convex-domain variance inequality. The next companion checkpoint
remains V40. The new estimates below do not use V30's screening or any assumed
native global spectral gap. Two empty semantic searches were followed by
mounted-text inventories and an identified f14 V82 read; they are not evidence
of archive-wide novelty.

Thermodynamic interpolation already occurs in f14 V82, for a different
chart law weighted by a determinant. The cubic--quartic log-integral
coefficient is also a standard cumulant calculation. For primary-source context,
Greeff, \href{https://arxiv.org/abs/2009.14321}{\emph{Free Energy Perturbation Theory
at Low Temperature}}, Sections I--II, records the cumulant expansion and its
sensitivity to constant energy shifts. Katsevich,
\href{https://arxiv.org/abs/2306.07262v3}{\emph{The Laplace approximation accuracy
in high dimensions: a refined analysis and new skew adjustment}},
Theorem 2.17 and Corollary 2.25, explains the cancellation of the odd leading
correction for even observables. Those statements and surrounding assumptions
were inspected in primary HTML. We prove the needed finite interpolation
identities and estimates directly; neither paper supplies the native constants,
complete-chart tails, or regulator regime used here. No convergent cumulant
series is assumed.

Keep exactly V34--V35's notation
\[
 w=(\zeta,y),\quad d=(72n-21)B,\quad
 U_0=\tfrac12w^{\mathsf T}\mathcal H w,\quad
 \mathcal C=\mathcal H^{-1},\quad \epsilon=\gamma^{-1/2},
\]
with $B=m^3$, $m\ge4$, $n\ge2$, and both retained endpoints the identity.
Let $\mathbb P_\gamma$ and $\mathbb Q$ be their unchanged complete joint laws,
and $\mathcal Z_\gamma,\mathcal Z_0$ their unchanged chart masses. Thus
$U_\gamma(0)=U_0(0)=0$ and $\mathcal J_\gamma(0)=1$ are exact conventions, not
post hoc adjustments. Let $\lambda_g\in\{1/2,1\}$ be each original three-group's
Haar power; $1/2$ occurs on the dressed chord endpoints only.

Fix $s\ge1$, apply V35's radius construction with tail exponent $t=s+4$, and
write $\mathbb B=\{|w|\le R_t\}$. Assume its exact entry conditions
\eqref{f16v35:entry} at this exponent, and, for convenience,
\begin{equation}
 \rho:=\frac{dn^4}{\gamma}\le1.
 \label{f16v36:rho}
\end{equation}
No different full law is defined by this ball. On it put
\[
 W=U_\gamma-\log\mathcal J_\gamma-U_0,\qquad
 d\mu_\theta=Z_{\theta,\mathbb B}^{-1}
       1_{\mathbb B}e^{-U_0-\theta W}\,dw,\quad 0\le\theta\le1.
\]
Here $\mu_0=\mathbb Q(\cdot\mid\mathbb B)$ and
$\mu_1=\mathbb P_\gamma(\cdot\mid\mathbb B)$, including their own normalizers.
Intermediate $\theta$ is a comparison parameter, not physical coupling or time.
All constants denoted $C$ below depend only on fixed word geometry; tail-entry
thresholds may depend on the fixed exponent $s$.

\subsection{The odd cubic and the measured quartic coefficient}

For an original word let $f_\ell(w)=s(W_\ell(w))$, with its actual ordered
exponentials and action weight $\omega_\ell$. Define the homogeneous polynomials
\[
 P_3(w)=\frac1{3!}\sum_\ell\omega_\ell
                    D^3f_\ell(0)[w,w,w],\qquad
 U_4(w)=\frac1{4!}\sum_\ell\omega_\ell D^4f_\ell(0)[w,w,w,w],
\]
and the measured coefficient
\begin{equation}
 P_4(w)=U_4(w)+\frac13\sum_g\lambda_g|w_g|^2.
 \label{f16v36:P4}
\end{equation}
Exactly quadratic endpoint amplitudes contribute to $U_0$, not to these jets.
These polynomials are coefficients for the \emph{joint} endpoint-fixed integral;
they are not V10's common-source matrices $G_j$.

\begin{proposition}[Local jets with their Haar and parity terms]
\label{f16v36:jets}
On $\mathbb B$, with a fixed geometry constant,
\begin{align}
 |W-\epsilon P_3-\epsilon^2P_4|
   &\le C\epsilon^3\sum_g|w_g|^5+C\epsilon^4\sum_g|w_g|^4,
       \label{f16v36:jet-value}\\
 |\nabla(W-\epsilon P_3)|^2
   &\le C\epsilon^4\sum_g(|w_g|^6+|w_g|^2),
       \label{f16v36:jet-gradient}\\
 |\nabla P_3|^2&\le C\sum_g|w_g|^4.\notag
\end{align}
The polynomial $P_3$ is odd and $P_4$ is even. In particular
$\mathbb E_{\mu_0}P_3=\mathbb E_{\mathbb Q}P_3=0$.
\end{proposition}
\begin{proof}
The ordered-unitary derivative argument of V6/V34 applies at orders four and
five: a differentiated exponential is an integral of unitary products with the
specified inserted matrices. At every fixed order its norm is bounded by the
product of those fixed matrix norms. Taylor's integral remainder on each word
therefore has the asserted local degree, with its original power of $\epsilon$.
There is no principal logarithm of an intermediate product. Summing the finite
carriers and their bounded row and column incidences gives the displayed
bounds, rather than a factor counting all words. The same argument applies to
one derivative of the fourth-order remainder.

For $|w_g|/\sqrt\gamma\le1/2$,
\[
 -\lambda_g\log j(w_g/\sqrt\gamma)
 =\frac{\lambda_g\epsilon^2}{3}|w_g|^2
                         +O(\epsilon^4|w_g|^4).
\]
The derivative of this term is $O(\epsilon^2|w_g|)$. This proves
\eqref{f16v36:P4} and its contribution to both estimates. The endpoint half
powers are not replaced by full Haar. Homogeneity gives the two parities, and
both Gaussian measures in the last claim are reflection invariant. The full
native measure need not be reflection invariant: oddness is used only where
that symmetry is actually available.
\end{proof}

\subsection{Uniform interpolation moments and centered cumulants}

Every $\mu_\theta$ has Hessian at least $\mathcal H/2$, since this holds at
both interpolation endpoints. The same is true after any real linear tilt.
Global simultaneous conjugation preserves $\mathbb B$, $U_0$, and $W$, and
hence makes the \emph{untilted} coordinate means zero. V35's variance argument
therefore gives, for every fixed $k$,
\begin{equation}
 \mathbb E_{\mu_\theta}|u^{\mathsf T}w|^k
       \le C_k(u^{\mathsf T}\mathcal C u)^{k/2},\qquad
 \mathbb E_{\mu_\theta}|w_g|^k\le C_kn^{k/2}.
 \label{f16v36:directional}
\end{equation}
The first estimate is used only for integer $k\ge1$. It follows by integrating
the sub-Gaussian tail obtained from the linear-tilt moment-generating function.
It is uniform in $\theta$, not an identification of $\mu_\theta$ with a Gaussian.

\begin{proposition}[Local-score cumulant bounds]
\label{f16v36:cumulants}
Uniformly for $0\le\theta\le1$,
\begin{align}
 \operatorname{Var}_{\mu_\theta}W&\le C\rho,&
 \mathbb E_{\mu_\theta}|W-\mu_\theta W|^4&\le C\rho^2,&
 |\kappa_{3,\mu_\theta}(W)|&\le C\rho^{3/2},
 \label{f16v36:W-moments}\\
 \operatorname{Var}_{\mu_\theta}(W-\epsilon P_3)
      &\le C\epsilon^4dn^5,&
 \operatorname{Var}_{\mu_0}P_3&\le Cdn^4,&
 |\mu_0P_4|&\le Cdn^2.
 \label{f16v36:jet-moments}
\end{align}
Here $\kappa_3(W)=\mathbb E[(W-\mathbb EW)^3]$. No independence of local
scores or summands is used.
\end{proposition}
\begin{proof}
Write $G_A=\nabla A^{\mathsf T}\mathcal C\nabla A$. The variance inequality
under $\mu_\theta$ reads $\operatorname{Var}A\le2\mathbb EG_A$.
By V35's local score estimate, \eqref{f16v36:directional}, and
$\|\mathcal C\|\le Cn^2$,
\[
 \mathbb EG_W\le C\epsilon^2dn^4=C\rho,\qquad
 \mathbb EG_W^2\le C\epsilon^4d^2n^8=C\rho^2.
\]
For the second inequality, use Cauchy--Schwarz on pairs of local fourth powers
and their eighth moments, retaining $d^2$ explicitly. It is not a decorrelation
argument.

For any smooth $A$, put $A_c=A-\mathbb EA$, $M_4=\mathbb EA_c^4$ and
$G=G_A$. Poincar\'e applied also to $A_c^2$ gives
\[
 M_4\le4(\mathbb EG)^2+8\mathbb E[A_c^2G]
       \le4\mathbb EG^2+8\sqrt{M_4\mathbb EG^2}.
\]
Solving the scalar quadratic yields $M_4\le81\mathbb EG^2$.
H\"older gives $|\kappa_3(A)|\le M_4^{3/4}$, proving
\eqref{f16v36:W-moments}. All functions are smooth on the fixed compact
comparison ball, so these applications are within the variance form domain.

The gradient estimates in Proposition~\ref{f16v36:jets}, the local sixth and
fourth moments, and $\|\mathcal C\|\le Cn^2$ give the first two assertions of
\eqref{f16v36:jet-moments}. The local quartic and quadratic moments give its
last assertion. The identical polynomial estimates hold under the full
Gaussian $\mathbb Q$ when no $\theta$ occurs.
\end{proof}

\subsection{An absolute partition ratio with the full chart mass retained}

\begin{theorem}[Polynomial-regime absolute normalization]
\label{f16v36:pressure}
Under the entry conditions above,
\begin{equation}
 \boxed{\quad
  \left|\log\frac{\mathcal Z_\gamma}{\mathcal Z_0}\right|
       \le C(\rho+\gamma^{-s}).
 \quad}
 \label{f16v36:pressure-bound}
\end{equation}
In particular the actual complete chart masses have ratio tending to one in
V35's polynomial regime. This is a new scalar estimate, not an inference from
total variation of normalized laws.
\end{theorem}
\begin{proof}
Let $\psi(\theta)=\log(Z_{\theta,\mathbb B}/Z_{0,\mathbb B})$.
Differentiation of these finite smooth integrals keeps every normalizer and gives
\begin{equation}
 \psi'=-\mu_\theta W,\qquad
 \psi''=\operatorname{Var}_{\mu_\theta}W,\qquad
 \psi'''=-\kappa_{3,\mu_\theta}(W).
 \label{f16v36:thermodynamic}
\end{equation}
The exact second-order integral identity is
\[
 \psi(1)=-\mu_0W+\int_0^1(1-\theta)
                              \operatorname{Var}_{\mu_\theta}W\,d\theta.
\]
Oddness eliminates $\mu_0P_3$. Equations \eqref{f16v36:jet-value} and
\eqref{f16v36:directional} give
\[
 |\mu_0W|\le C\{\epsilon^2dn^2+\epsilon^3dn^{5/2}
                                      +\epsilon^4dn^2\}\le C\rho.
\]
The variance term is at most $C\rho$. Notice that smallness of the mean is
proved from the anchored action and parity; it does not follow from its gradient
alone.

For the complete masses the exact relation is
\begin{equation}
 \log\frac{\mathcal Z_\gamma}{\mathcal Z_0}
  =\psi(1)+\log\mathbb Q(\mathbb B)-\log\mathbb P_\gamma(\mathbb B).
 \label{f16v36:full-mass}
\end{equation}
V35's actual complete tails at exponent $t=s+4$ bound the last two terms by
$C\gamma^{-t}$. Thus neither the native chart complement nor either full
normalizer has been omitted. This proves the theorem.
\end{proof}

Define two coefficients by expectations in the \emph{untruncated} original
joint Gaussian, not the ball-conditioned measure:
\begin{equation}
 \mathfrak p_{B,n}=\tfrac12\mathbb E_{\mathbb Q}P_3^2
                                  -\mathbb E_{\mathbb Q}P_4,
 \qquad
 \mathfrak e_{B,n}=\tfrac12\mathbb E_{\mathbb Q}P_3^2\ge0.
 \label{f16v36:coefficients}
\end{equation}
These coefficients can depend on $B,n$. They contain all original word
incidences, kinetic terms, and endpoint Haar powers.

\begin{theorem}[The first absolute coefficient with a quantitative remainder]
\label{f16v36:coefficient}
The coefficients in \eqref{f16v36:coefficients} satisfy
$|\mathfrak p_{B,n}|+\mathfrak e_{B,n}\le Cdn^4$, and
\begin{equation}
 \boxed{\quad
 \left|\log\frac{\mathcal Z_\gamma}{\mathcal Z_0}
                  -\frac{\mathfrak p_{B,n}}\gamma\right|
       \le C(\rho^{3/2}+\gamma^{-s}).
 \quad}
 \label{f16v36:pressure-expansion}
\end{equation}
At fixed $B,n$, this gives the coefficient of $\gamma^{-1}$ with an
$O_{B,n}(\gamma^{-3/2})$ remainder by choosing $s\ge3/2$.
The bound is not asserted sharp in fixed dimension.
\end{theorem}
\begin{proof}
The coefficient bound follows from Gaussian Poincar\'e for $P_3$ and the
local fourth moments for $P_4$. The same calculation on $\mathbb B$ was proved
in Proposition~\ref{f16v36:cumulants}.
Taylor's formula for \eqref{f16v36:thermodynamic} gives
\[
 \psi(1)=-\mu_0W+\tfrac12\operatorname{Var}_{\mu_0}W
       -\tfrac12\int_0^1(1-\theta)^2\kappa_{3,\mu_\theta}(W)\,d\theta.
\]
The last integral is $O(\rho^{3/2})$. The jet mean satisfies
\[
 |\mu_0W-\epsilon^2\mu_0P_4|
       \le C(\epsilon^3dn^{5/2}+\epsilon^4dn^2).
\]
For $R=W-\epsilon P_3$, Cauchy--Schwarz and
\eqref{f16v36:jet-moments} give
\[
 |\operatorname{Var}_{\mu_0}W-\epsilon^2\mu_0P_3^2|
 \le C\epsilon^3dn^{9/2}+C\epsilon^4dn^5.
\]
Every displayed error is bounded by $C\epsilon^3d^{3/2}n^6=C\rho^{3/2}$
for $d\ge1$, $n\ge2$, and $\epsilon\le1$.

It remains to replace the Gaussian ball coefficients by the full ones.
Bounded word incidence gives $|P_3|^2\le C|w|^6$ and
$|P_4|\le C(|w|^4+|w|^2)$. The actual Gaussian eighth-degree tail from V35
therefore pays their omitted integrals by $C\gamma^{-t}$. The conditioning
normalizer adds at most $Cdn^4\gamma^{-t}$. After multiplication by
$\epsilon^2$, this is $C(\epsilon^2+\rho)\gamma^{-t}\le C\gamma^{-t}$.
Equation~\eqref{f16v36:full-mass} completes the proof. Only a finite Taylor
formula on a compact ball has been used; an unbounded cubic exponential has
not been declared integrable, and no infinite series is summed.
\end{proof}

\paragraph{Explicit Gaussian contractions.}
Use the symmetric coefficient tensors defined by
\[
 P_3=T_{ijk}w_iw_jw_k/6,\qquad U_4=A_{ijkl}w_iw_jw_kw_l/24,
\]
with repeated scalar indices summed. Set
$v_k=T_{ijk}\mathcal C_{ij}$. Gaussian pairings give
\begin{equation}
 \begin{split}
 \mathfrak p_{B,n}={}&
 \tfrac1{12}T_{ijk}T_{abc}\mathcal C_{ia}\mathcal C_{jb}\mathcal C_{kc}
 +\tfrac18v^{\mathsf T}\mathcal C v
 -\tfrac18 A_{ijkl}\mathcal C_{ij}\mathcal C_{kl}\\
 &-\tfrac13\sum_g\lambda_g\Tr\mathcal C_{gg}.
 \end{split}
 \label{f16v36:wick}
\end{equation}
There are six fully cross pairings and nine one-cross pairings in the cubic
square, and three pairings in the quartic mean. In this native identity-endpoint
setting $v=0$: $\mathbb E_{\mathbb Q}\nabla P_3=v/2$ is a vector in each colour
group fixed by every simultaneous adjoint rotation, and hence is zero. The
formula is displayed before this simplification to retain its general pairing
convention. No numerical enclosure or evaluation of the entire native tensor
contraction is claimed.

As an exact normalization check, one isolated Wilson factor with full Haar has
$P_3=0$, $U_4=-|w|^4/24$, and Gaussian covariance $I_3$. Its coefficient is
$15/24-3/3=-3/8$. Replacing that \emph{reference factor} by a half-Haar factor
instead gives $15/24-3/6=1/8$. These one-group checks are not the full-history
coefficient; they show why its actual endpoint powers cannot be discarded.

\subsection{Relative entropy of the complete law, with its support direction}

\begin{theorem}[Forward native relative entropy and its first coefficient]
\label{f16v36:entropy}
In the same regime,
\begin{equation}
 \boxed{\quad
 0\le D(\mathbb P_\gamma\Vert\mathbb Q)\le C(\rho+\gamma^{-s}),\qquad
 \left|D(\mathbb P_\gamma\Vert\mathbb Q)
                    -\frac{\mathfrak e_{B,n}}\gamma\right|
       \le C(\rho^{3/2}+\gamma^{-s}).
 \quad}
 \label{f16v36:entropy-bound}
\end{equation}
For every finite coupling the reverse entropy
$D(\mathbb Q\Vert\mathbb P_\gamma)$ is infinite, because the untruncated
Gaussian has positive mass outside the native principal product.
\end{theorem}
\begin{proof}
On the common comparison ball, exact normalization gives
\[
 D(\mu_1\Vert\mu_0)=\psi'(1)-\psi(1)
       =\int_0^1\theta\operatorname{Var}_{\mu_\theta}W\,d\theta.
\]
This proves the $O(\rho)$ bound there. Taylor's formula for $\psi'$ and $\psi$
and Proposition~\ref{f16v36:cumulants} give
\[
 D(\mu_1\Vert\mu_0)=\tfrac12\operatorname{Var}_{\mu_0}W
                                      +O(\rho^{3/2}).
\]
The coefficient substitutions in the preceding proof give
$\mathfrak e_{B,n}/\gamma$ with the same error. The quartic mean and Haar mean
cancel from this leading entropy coefficient, not from the pressure coefficient.

The full entropy needs an additional tail check: $-\log\mathcal J_\gamma$
is unbounded near a principal cut. The complete ordered-word second-derivative
bound gives $U_\gamma+U_0\le C|w|^2$ on the native chart. Its tail expectation
is paid by V35. For the logarithmic Haar term, use the pointwise identity and
bound
\[
 0\le-\mathcal J_\gamma\log\mathcal J_\gamma\le e^{-1},
 \qquad(0\log0:=0).
\]
The same radial tail integral and lower bound for $\mathcal Z_\gamma$ used in
V35 then give
\begin{equation}
 \mathbb E_{\mathbb P_\gamma}[|W|1_{\mathbb B^c}]
                                      \le C\gamma^{-t}.
 \label{f16v36:entropy-tail}
\end{equation}
No Gaussian expectation of the singular logarithm is substituted here.

Let $p=\mathbb P_\gamma(\mathbb B)$, $q=\mathbb Q(\mathbb B)$. On the ball,
$\log(d\mathbb P_\gamma/d\mathbb Q)
=\log(d\mu_1/d\mu_0)+\log(p/q)$. Off it, on the native support, the same
logarithm is $-W-\log(\mathcal Z_\gamma/\mathcal Z_0)$.
Integrating this exact split, using \eqref{f16v36:entropy-tail},
$1-p,1-q\le\gamma^{-t}$ and Theorem~\ref{f16v36:pressure}, changes the ball
entropy by at most $C\gamma^{-s}$. This proves both claims for the full law.
The reverse assertion is immediate from the support mismatch and is not
contradicted by either total variation convergence or the finite forward entropy.
\end{proof}

The chain rule, with the same actual marginals and conditionals, now yields
\begin{equation}
 \begin{split}
 D(\mathbb P_\gamma\Vert\mathbb Q)
 ={}&D(\boldsymbol\nu_\gamma\Vert Q_{B,n}^{00})\\
 &+\mathbb E_{\boldsymbol\nu_\gamma}
 D(\mathbb P_\gamma(d\zeta\mid y)\Vert\mathbb Q(d\zeta\mid y)).
 \end{split}
 \label{f16v36:entropy-chain}
\end{equation}
Both nonnegative terms are therefore at most $C(\rho+\gamma^{-s})$. This is
an average under the \emph{actual retained marginal}, not a supremum over all
retained boundary histories and not a new arbitrary-fast-pinning theorem.

\subsection{Even observables have a second-order covariance error}

\begin{theorem}[A sharper relative covariance bound in the same Gaussian metric]
\label{f16v36:covariance}
For $\Sigma_\gamma=\operatorname{Cov}_{\mathbb P_\gamma}(w)$,
\begin{equation}
 \boxed{\quad
 \|\mathcal H^{1/2}\Sigma_\gamma\mathcal H^{1/2}-I\|_{\rm op}
                  \le C(\rho+\gamma^{-s}).
 \quad}
 \label{f16v36:relative-covariance}
\end{equation}
This improves V35's $O(\sqrt\rho+\gamma^{-s})$ bound for this particular even
observable class; it does not claim an improved total-variation rate for
arbitrary events. The retained covariance satisfies the corresponding
relative bound with $(K^{\rm int})^{-1}$.
\end{theorem}
\begin{proof}
Fix $u$ with $u^{\mathsf T}\mathcal C u=1$, and set $F=(u^{\mathsf T}w)^2$.
The directional moment estimates give $\operatorname{Var}_{\mu_\theta}F\le C$
uniformly in $u,\theta$. For $f(\theta)=\mu_\theta F$ the exact derivatives are
\[
 f'(\theta)=-\operatorname{Cov}_{\mu_\theta}(F,W),\qquad
 f''(\theta)=\mathbb E_{\mu_\theta}
                 [(F-\mu_\theta F)(W-\mu_\theta W)^2].
\]
The second equality retains both derivative-of-normalizer terms through
centering. Reflection of $\mu_0$ makes $\operatorname{Cov}_{\mu_0}(F,P_3)=0$.
Consequently
\[
 |f'(0)|\le C\sqrt{\operatorname{Var}_{\mu_0}(W-\epsilon P_3)}
       \le C\epsilon^2\sqrt{dn^5}\le C\rho.
\]
Cauchy--Schwarz and \eqref{f16v36:W-moments} give
$|f''(\theta)|\le C\rho$. Taylor's integral formula therefore bounds
$|\mu_1F-\mu_0F|$ by $C\rho$.

Since $|u|^2\le H_*$, V35's complete weighted tails pay all omitted second
moments and the two conditioning probabilities uniformly in $u$ by
$C\gamma^{-t}$. The complete and conditioned coordinate means vanish by global
conjugation, so second moments are exactly the needed variances. The result is
uniform in every Gaussian-normalized direction; taking the supremum proves
\eqref{f16v36:relative-covariance}. Principal covariance blocks give the retained
statement. No spectral convergence of a sampler is inferred from it.
\end{proof}

\begin{proposition}[Derivative banks, marked Gaussian predictors, and growth rates]
\label{f16v36:handoff}
Put $\eta=C(\rho+\gamma^{-s})$. The unchanged V33 assembled difference--curvature
map and the unchanged linear word bank obey
\begin{equation}
 \operatorname{Cov}_{\boldsymbol\nu_\gamma}(\mathcal B y)\preceq(1+\eta)I,
 \qquad
 \|\operatorname{Cov}_{\mathbb P_\gamma}(\mathcal D w)
                     -\mathcal D\mathcal C\mathcal D^{\mathsf T}\|
                                                    \le\eta.
 \label{f16v36:linear-fields}
\end{equation}
The retained-sampler upper test of V35 remains valid with this smaller $\eta$
when $\eta<1$. No second-order rate for the \emph{nonlinear} word bank is
asserted by these linear statements.

A sufficient simultaneous regime remains
\begin{equation}
 \frac{Bn^7\log(n+1)+n^6\log\gamma}{\gamma}\longrightarrow0.
 \label{f16v36:regime}
\end{equation}
It has not been enlarged. Within it, the complete chart partition ratio tends
to one, the forward entropy tends to zero, and the relative covariance error
is $O(Bn^5/\gamma+\gamma^{-s})$. At fixed $B$ and $n\asymp\gamma^{1/8}$,
these bounds are $O(\gamma^{-3/8})$ for $s\ge1$, and the remainder after the
explicit pressure coefficient is $O(\gamma^{-9/16})$.
\end{proposition}
\begin{proof}
The matrix inequalities $\mathcal B^{\mathsf T}\mathcal B\preceq K^{\rm int}$
and $\mathcal D^{\mathsf T}\mathcal D\preceq\mathcal H$ were proved in V33--V34.
Apply the new relative covariance theorem to their original input spaces.
V35's upper sampler test uses only such a relative covariance comparison and
Gaussian conditional predictors as admissible native predictors, so the same
proof applies. The nonlinear word itself has a separate even second-order jet;
its products have not been estimated here at the improved rate.

Use V35's explicit polynomial entry with $t=s+4$ instead of $s$.
Condition \eqref{f16v36:regime} supplies that entry and $\rho\to0$.
Since $d=(72n-21)B$, the displayed powers follow by substitution. The radius and
its entry cost are unchanged except for this fixed tail exponent. In particular
this is not a new duration-uniform mass estimate or a removal of the soft
$n^2$ inverse from the proof.
\end{proof}

\subsection{The scalar handoff and remaining physical normalization}

Let $\mathcal A_{\gamma,B,n}^{00}$ denote the exact identity-endpoint compressed
kernel written in V4's bridge half-density chart. Its original factorization is
\[
 \mathcal A_{\gamma,B,n}^{00}
    =c_{\gamma,B,n}F_{\gamma,2}(I)^{-1}\mathcal Z_\gamma.
\]
Writing the analogous Gaussian coefficient as
$c_{0,B,n}F_2(0)^{-1}\mathcal Z_0$ gives the exact scalar separation
\begin{equation}
 \begin{split}
 \log\frac{\mathcal A_{\gamma,B,n}^{00}}{\mathcal A_{0,B,n}^{00}}
 ={}&\log\frac{\mathcal Z_\gamma}{\mathcal Z_0}
   +\log\frac{c_{\gamma,B,n}}{c_{0,B,n}}\\
 &+\log\frac{F_2(0)}{F_{\gamma,2}(I)}.
 \end{split}
 \label{f16v36:physical-scalars}
\end{equation}
The first term is now quantitatively estimated in the polynomial regime; the
last two are not silently set to zero or replaced by a chosen vacuum constant.
A diagonal endpoint kernel is also not a Perron eigenvalue. The normalized
probability, its chart mass, and the physical transfer at its own spectral top
remain distinct objects.

V35's warning about additive constants remains mathematically correct: replacing
$U_\gamma$ by $U_\gamma+a_\gamma$ preserves every normalized law and covariance
but changes $\log\mathcal Z_\gamma$ by $-a_\gamma$. V36 controls the scalar
because the original action convention is fixed and its mean correction is
computed, not because that warning has been bypassed by terminology.

The next structural task is still a source-relative treatment of the temporal
inverse or the compact nonlinear periods, together with the untouched physical
scalar factors where a transfer comparison is intended. The present second-order
cancellations do not enlarge the valid space--time regime, give a stationary
root-Gauss-reduced vacuum, calibrate physical time, or construct an interacting
continuum theory. They close a specific missing normalization estimate and
improve its accompanying even-observable comparison.

\subsection{Verification and preservation}

The diagnostic is
\begin{center}\small\path{verify_ym_f16_v36_pressure_parity.py}.\end{center}
It checks the signs and weights in exact normalized interpolation, Gaussian
cubic--quartic contractions, a genuine ordered quaternion cubic, full versus
half-Haar coefficients, forward entropy and its direction, and the even-observable
cancellation. Finite algebraic fixtures are not enclosures of the native
high-dimensional Wilson integral or formal verification of the analytic proof.
Actual test outputs, predecessor reruns, hashes, and rendering records are
reported in the package, not inferred from their historical descriptions.

Only one literal title version and this terminal insertion change V35. Removing
them recovers the predecessor byte for byte. The deductions inherit V35's
interior curvature and complete tails, with their compact-pin and finite-word
premises. The entire analytic archive has not been independently recertified.
The next companion checkpoint remains V40.

\begin{firewall}
V36 estimates the absolute endpoint-fixed chart partition ratio and its first
Gaussian coefficient in the existing polynomial regime. It proves finite
forward relative entropy, its leading coefficient, and a second-order relative
covariance estimate for the linear-coordinate bank. The reverse full entropy
is infinite because of the support mismatch. All complete probabilities,
endpoint Haar powers, Gaussian soft directions, and original scalar factors
are retained. No stronger total-variation rate, unrestricted space--time limit,
physical Perron normalization, calibrated transfer gap, or interacting
Yang--Mills continuum mass gap follows.
\end{firewall}
% END F16 V36 APPEND

% BEGIN F16 V37 APPEND -- 2026-09-18
% Insert immediately before the final document-ending command of V36.
% Every predecessor assertion, endpoint factor, measure and clock is preserved.

\clearpage
\section{V37: complete scalar normalization of the dressed endpoint kernel}
\label{f16v37:context}

V36 estimates the full joint chart mass but leaves two scalar factors in
\eqref{f16v36:physical-scalars}. Before proving a new spectral implication, we
return to their definitions. The finite-chart normalization of the chosen cube
carrier cancels \emph{exactly} against its magnetic endpoint normalization.
What remains is one static native integral and the explicitly normalized
one-group Wilson factor. Estimating these gives the complete identity-endpoint
compressed-kernel ratio, not merely the ratio of normalized path probabilities.
The carrier and the transfer themselves are unchanged.

This is a scalar closure of the displayed identity, not a spectral closure of
Yang--Mills. The result is at fixed identity endpoints, before the remaining
root Gauss restriction, and in the polynomial regulator regime inherited from
V35--V36. A diagonal kernel of a power does not determine its Perron eigenvalue
or its first excited spectral ratio. That limitation is retained quantitatively
below instead of renaming the computed scalar a vacuum energy.

\subsection{Exact endpoint cancellation before any asymptotic estimate}

The complete V36 appendix and analytic audit were read. Its joint pressure
coefficient and error estimate are used as inherited inputs. We also checked
the literal definitions \eqref{f16v1:gram}, \eqref{f16v4:amplitude}, and
\eqref{f16v4:normalizers}. The general cancellation is elementary; its
quantitative application to the two outstanding factors is the purpose here.
The static integral is not one cube in isolation: all mixed plaquettes of its
spatial assembly remain coupled. Targeted searches of the active source and the
supplied f14/f15/Grand Tensor archives are recorded in the audit, not presented
as an exhaustive novelty or proof certificate. The next checkpoint remains V40.

The one-variable Bessel formulas used for the kinetic factor are classical.
The NIST Digital Library of Mathematical Functions, equations
\href{https://dlmf.nist.gov/10.32.E2}{10.32.2} and
\href{https://dlmf.nist.gov/10.40.E1}{10.40.1}, with coefficients from
\href{https://dlmf.nist.gov/10.17.E1}{10.17.1}, were checked in primary HTML.
We give a direct real-integral remainder bound as well; no lattice theorem is
imported from those special-function identities. The supplied memoranda's
separation of scalar normalization, locality cost and physical spectral
comparison remains in force.

Let $m\ge4$, $B=m^3$, and write $d_{\rm e}=15B$. In scalar coordinates
$C_B$ and the spatial matrices carry their original colour tensor $I_3$.
For a complete compact bridge configuration $Z$, define
\begin{equation}
 \begin{split}
  E_{\gamma,B}(x;Z)&=\tfrac12x^{\mathsf T}C_B^{-1}x
       +\gamma S_\times(X(x),\Theta_{X(x)}^{-1}Z),\\
  \mathcal S_{\gamma,B}(Z)&=
       \int_{\mathcal P_x}e^{-E_{\gamma,B}(x;Z)}\,dx .
 \end{split}
 \label{f16v37:static-mass}
\end{equation}
The domain $\mathcal P_x$ is the original product of principal chord charts.
There is \emph{no Haar density} in this displayed integral. This is not an
approximation: the reason is the exact square-root Jacobian already built into
$\Phi_\gamma$. Haar factors in the full joint history are not removed.

\begin{proposition}[The carrier chart mass cancels from the full dressed weight]
\label{f16v37:cancellation}
At every finite $B,n,\gamma$ and every admitted compact retained history,
\begin{equation}
 F_{\gamma,2}(Z)
     =\frac{c_\Omega^2}{p_{\gamma,B}}\mathcal S_{\gamma,B}(Z),
 \qquad
 \mathcal W^{\rm full}_{\gamma,n}(\mathbf y)
 =\left(\frac{c_{H,\gamma}}{z_\gamma}\right)^{24Bn}
 \frac{J_{{\rm br},\gamma}(\mathbf y)Z_{\gamma,\rm full}(\mathbf y)}
 {\sqrt{\mathcal S_{\gamma,B}(Z_0)\mathcal S_{\gamma,B}(Z_n)}}.
 \label{f16v37:cancelled-weight}
\end{equation}
No chart probability or exponentially small carrier mass has been estimated in
this identity. Every mixed plaquette, endpoint Gaussian, and original kinetic
normalizer is retained.
\end{proposition}
\begin{proof}
The exact change to chord charts gives
$|\Phi_\gamma(X)|^2dH^{5B}(X)=p_{\gamma,B}^{-1}\Omega(x)^2dx$.
Also $M_\times^2=e^{-\gamma S_\times}$, with the original moving bridge input.
Substitute these identities into \eqref{f16v1:gram}; the first equality follows.
Each of the two square-root endpoint factors contains
$c_\Omega/\sqrt{p_{\gamma,B}}$. Their product cancels
$p_{\gamma,B}^{-1}c_\Omega^2$ in $c_{\gamma,B,n}$. Substitution into V4's exact
weight proves the second equality. The cancellation is of identical scalar
factors, not a replacement of $F_{\gamma,2}$ by its limiting determinant.
\end{proof}

At the identity define
\[
 E_{\gamma,B}(x)=E_{\gamma,B}(x;I),\qquad
 E_{0,B}(x)=\tfrac12 x^{\mathsf T}H_{\rm e}x,\qquad
 H_{\rm e}=C_B^{-1}+A^{\mathsf T}A,\qquad C_{\rm e}=H_{\rm e}^{-1}.
\]
Thus $H_{\rm e}^{-1}$ is the inherited static posterior covariance; it is not
the joint endpoint-history covariance $\mathcal C$ of V35. Put
\[
 \mathcal S_{0,B}=\int_{\mathbb R^{d_{\rm e}}}e^{-E_{0,B}}dx,
 \qquad Q_{\rm e}=N(0,C_{\rm e}).
\]
The original Gaussian normalization is
$F_2(0)=c_\Omega^2\mathcal S_{0,B}$. In particular its determinant is exactly
that in V1, with all colours. With both retained endpoints fixed at $I$, let
$\mathcal A_{\gamma,B,n}^{00}$ be V36's unchanged compressed diagonal kernel.
Integrating the interior bridges in \eqref{f16v37:cancelled-weight} gives
\begin{equation}
 \mathcal A_{\gamma,B,n}^{00}
    =\left(\frac{c_{H,\gamma}}{z_\gamma}\right)^{24Bn}
                   \frac{\mathcal Z_\gamma}{\mathcal S_{\gamma,B}(I)},
 \qquad
 \mathcal A_{0,B,n}^{00}
    =(2\pi)^{-36Bn}\frac{\mathcal Z_0}{\mathcal S_{0,B}}.
 \label{f16v37:kernel-factors}
\end{equation}
The joint masses $\mathcal Z_\gamma,\mathcal Z_0$ are exactly V34--V36's masses,
not the fixed-bridge fast integrals $Z_{\gamma,\rm full}(\mathbf y)$.

\subsection{The static endpoint integral has its own uniform Gaussian scale}

Fix a requested tail power $s\ge1$ and put $t=s+4$. There are fixed constants
$0<a_{\rm e}\le1/4$, $H_{\rm e,+}\ge1$, and $C_{\rm e,2}$ such that
\begin{equation}
 E_{\gamma,B},E_{0,B}\ge a_{\rm e}|x|^2,\qquad
 2a_{\rm e}I\preceq H_{\rm e}\preceq H_{\rm e,+}I,
 \qquad
 \|D^2(E_{\gamma,B}-E_{0,B})\|
             \le C_{\rm e,2}|x|/\sqrt\gamma.
 \label{f16v37:static-data}
\end{equation}
The native inequality is on the whole principal product. The first term of
$E_{\gamma,B}$ supplies the pin independently of the retained input; the mixed
action is nonnegative. The fixed cube covariance and bounded mixed-word
incidence give the two Gaussian bounds. The last bound follows from the same
ordered-unitary third derivatives used in V6. These constants do not involve
$n$ or $B$.

For explicit entry bookkeeping set
\begin{equation}
 \begin{split}
 L_{{\rm e},t}={}&1+\tfrac{d_{\rm e}}2\log(2H_{\rm e,+}/a_{\rm e})
  +\log[1+((d_{\rm e}+6)/a_{\rm e})^4]+(t+1)\log\gamma,\\
 R_{{\rm e},t}^2={}&2L_{{\rm e},t}/a_{\rm e},\qquad
 \mathbb B_{\rm e}=\{|x|\le R_{{\rm e},t}\}.
 \end{split}
 \label{f16v37:static-radius}
\end{equation}
Assume $\gamma\ge2$ and
\begin{equation}
 R_{{\rm e},t}\le\sqrt\gamma/2,\qquad
 C_{\rm e,2}R_{{\rm e},t}/\sqrt\gamma\le a_{\rm e},\qquad
 \sigma:=d_{\rm e}/\gamma\le1.
 \label{f16v37:static-entry}
\end{equation}
A sufficiently small fixed constant in
$d_{\rm e}+(t+1)\log\gamma\le c_{\rm e}\gamma$ is sufficient, after a fixed
threshold. This condition is separately stated, not inferred from a
fixed-volume limit.

\begin{proposition}[Complete static tails and interpolation moments]
\label{f16v37:static-preparation}
Let $P_{\rm e,\gamma}$ be the normalized density proportional to
$1_{\mathcal P_x}e^{-E_{\gamma,B}}$. Under \eqref{f16v37:static-entry},
\begin{equation}
 \mathcal S_{\gamma,B}(I)\ge2^{-d_{\rm e}/2-1}\mathcal S_{0,B},\qquad
 \mathbb E_{P_{\rm e,\gamma}}[(1+|x|^8)1_{\mathbb B_{\rm e}^c}]
 +\mathbb E_{Q_{\rm e}}[(1+|x|^8)1_{\mathbb B_{\rm e}^c}]
 \le2\gamma^{-t}.
 \label{f16v37:static-tails}
\end{equation}
On the ball, every normalized interpolation
$\mu_\theta^{\rm e}\propto\exp[-E_{0,B}-\theta(E_{\gamma,B}-E_{0,B})]$,
$0\le\theta\le1$, has curvature at least $H_{\rm e}/2$ and local moments
$\mathbb E|x_g|^k\le C_k$ at each original chord group, for every fixed $k$.
\end{proposition}
\begin{proof}
The Hessian error is at most $a_{\rm e}I\preceq H_{\rm e}/2$. Both exponents
and gradients vanish at zero. Hence $E_{\gamma,B}\le2E_{0,B}$ on the ball.
The Gaussian of precision $2H_{\rm e}$ has mean square norm at most
$d_{\rm e}/(4a_{\rm e})$, and at least half its mass lies in that ball.
Integrating proves the lower bound on the actual static mass.

For either weighted tail use the complete pin
$e^{-a_{\rm e}|x|^2}$, splitting it into two equal exponents outside the ball.
The radial eighth moment is at most
$(2\pi/a_{\rm e})^{d_{\rm e}/2}[1+((d_{\rm e}+6)/a_{\rm e})^4]$.
Also $\mathcal S_{0,B}\ge(2\pi/H_{\rm e,+})^{d_{\rm e}/2}$.
Division by the proved native mass gives the tail bound
$2e^{-1}\gamma^{-(t+1)}\le\gamma^{-t}$ for each law. The native density is
zero off its principal product; that complement is included. No product Haar
factor is inserted into this static integral.

Each linear tilt of an interpolation has the same Hessian floor on the same
convex ball. Apply the constant-metric variance inequality proved in V35 to a
linear observable under that tilt, and integrate the second log-moment
derivative. Untilting gives zero means by simultaneous colour conjugation.
Thus $\mathbb E\exp(u^{\mathsf T}x)\le\exp(u^{\mathsf T}C_{\rm e}u)$.
The fixed bound on $C_{\rm e}$ gives all stated local moments. This argument
requires neither spatial independence nor a global convexity claim outside
the comparison ball.
\end{proof}

For the mixed action only, define its classical cubic and quartic jets
$P_3^{\rm e},P_4^{\rm e}$ by
\[
 E_{\gamma,B}-E_{0,B}
     =\gamma^{-1/2}P_3^{\rm e}+\gamma^{-1}P_4^{\rm e}
          +\text{Taylor remainder}.
\]
They are the third and fourth homogeneous differentials of the literal mixed
Wilson words at $Z=I$, with their original weights. The confining Gaussian is
exactly quadratic. In contrast to V36's joint $P_4$, there is no additional
$\frac13\sum_g\lambda_g|x_g|^2$ here: the static Haar factors already cancelled.
Put
\begin{equation}
 \mathfrak s_B
   =\tfrac12\mathbb E_{Q_{\rm e}}(P_3^{\rm e})^2
                      -\mathbb E_{Q_{\rm e}}P_4^{\rm e}.
 \label{f16v37:static-coefficient}
\end{equation}

\begin{theorem}[Static endpoint pressure with all mixed couplings retained]
\label{f16v37:static-pressure}
Under \eqref{f16v37:static-entry}, $|\mathfrak s_B|\le C d_{\rm e}$ and
\begin{equation}
 \begin{split}
 \left|\log\frac{\mathcal S_{\gamma,B}(I)}{\mathcal S_{0,B}}\right|
     &\le C(\sigma+\gamma^{-s}),\\
 \left|\log\frac{\mathcal S_{\gamma,B}(I)}{\mathcal S_{0,B}}
                     -\frac{\mathfrak s_B}{\gamma}\right|
     &\le C(\sigma^{3/2}+\gamma^{-s}).
 \end{split}
 \label{f16v37:static-pressure-bound}
\end{equation}
In the original flat chord half-density representation, the normalized static
endpoint vector also obeys
\begin{equation}
 \left\|\frac{1_{\mathcal P_x}e^{-E_{\gamma,B}/2}}
                   {\sqrt{\mathcal S_{\gamma,B}(I)}}
             -\frac{e^{-E_{0,B}/2}}{\sqrt{\mathcal S_{0,B}}}\right\|_2^2
           \le C(\sigma+\gamma^{-s}).
 \label{f16v37:endpoint-vector}
\end{equation}
This is the original dressed endpoint vector at $Z=I$, not an isolated-cube
vacuum. The coefficient can depend on $B$, although its displayed bound is
linear in $B$.
\end{theorem}
\begin{proof}
Write $\epsilon=\gamma^{-1/2}$ and $W_{\rm e}=E_{\gamma,B}-E_{0,B}$.
Fixed word carriers and local incidence give, on the ball,
\[
 |\nabla W_{\rm e}|^2\le C\epsilon^2\sum_g|x_g|^4,\quad
 |\nabla(W_{\rm e}-\epsilon P_3^{\rm e})|^2
                       \le C\epsilon^4\sum_g|x_g|^6,
\]
\[
 |W_{\rm e}-\epsilon P_3^{\rm e}-\epsilon^2P_4^{\rm e}|
                       \le C\epsilon^3\sum_g|x_g|^5.
\]
The original unitary derivative formulas justify these estimates without a
logarithm of a product. Under every $\mu_\theta^{\rm e}$, the preceding local
moments and bounded $C_{\rm e}$ imply
\[
 \operatorname{Var}W_{\rm e}\le C\epsilon^2d_{\rm e},\quad
 \mathbb E|W_{\rm e}-\mathbb EW_{\rm e}|^4
                                      \le C\epsilon^4d_{\rm e}^2,\quad
 \operatorname{Var}(W_{\rm e}-\epsilon P_3^{\rm e})
                                      \le C\epsilon^4d_{\rm e}.
\]
For clarity, the centered fourth-moment estimate is obtained as in V36 by applying
Poincar\'e to the centered observable and its square; it is not a
factorization of local correlations. It bounds the third centered moment
by $C\sigma^{3/2}$. Gaussian Poincar\'e also gives
$\mathbb E_{Q_{\rm e}}(P_3^{\rm e})^2\le Cd_{\rm e}$, since the cubic is odd.
The quartic mean is bounded by $Cd_{\rm e}$.

Let $\psi_{\rm e}(\theta)$ be the log partition ratio of these ball integrals.
The exact identities $\psi_{\rm e}'=-\mathbb EW_{\rm e}$,
$\psi_{\rm e}''=\operatorname{Var}W_{\rm e}$ and
$\psi_{\rm e}'''=-\mathbb E(W_{\rm e}-\mathbb EW_{\rm e})^3$ are V36's
finite thermodynamic identities, now under the static laws. Oddness at
$\theta=0$ removes the cubic mean, not its variance. Taylor's formula at zero
therefore gives
\[
 \psi_{\rm e}(1)=\epsilon^2\left\{
     \tfrac12\mathbb E_{\mu_0^{\rm e}}(P_3^{\rm e})^2
                         -\mathbb E_{\mu_0^{\rm e}}P_4^{\rm e}\right\}
                         +O(\sigma^{3/2}).
\]
In detail, the fifth-order mean and the cubic/higher-jet covariance each cost
at most $C\epsilon^3d_{\rm e}$; the squared higher jet costs
$C\epsilon^4d_{\rm e}$. They are bounded by $C\sigma^{3/2}$ when
$d_{\rm e}\ge1$, $\sigma\le1$. No infinite cumulant expansion is invoked.

The eighth-degree Gaussian tails replace the two ball-conditioned coefficients
by their full $Q_{\rm e}$ expectations. Their normalizer corrections, multiplied
by $\epsilon^2$, are at most $C(\epsilon^2+\sigma)\gamma^{-t}$.
Finally the exact mass identity adds
$\log Q_{\rm e}(\mathbb B_{\rm e})-
 \log P_{\rm e,\gamma}(\mathbb B_{\rm e})$.
Proposition~\ref{f16v37:static-preparation} pays these terms and proves
\eqref{f16v37:static-pressure-bound}.

For \eqref{f16v37:endpoint-vector}, apply V35's normalized square-root
Poincar\'e calculation on the ball. Its score expectation is at most
$C\epsilon^2d_{\rm e}$. The square-root distance from a probability to its
ball conditioning is squared distance
$2(1-\sqrt{P(\mathbb B_{\rm e})})\le2P(\mathbb B_{\rm e}^c)$.
The triangle inequality and the two paid tails prove the full-space bound.
The exact flat representative follows by multiplying $a_{\gamma,I}(X)$ by
$\sqrt{J_{X,\gamma}}$ and using Proposition~\ref{f16v37:cancellation}.
\end{proof}

\subsection{The Wilson kinetic normalization is an explicit scalar}

Define
\begin{equation}
 \mathfrak b_\gamma=\frac{z_\gamma}{c_{H,\gamma}(2\pi)^{3/2}}.
 \label{f16v37:kinetic-ratio}
\end{equation}
The symbol $\mathfrak b_\gamma$ is a scalar, distinct from the internal magnetic
multiplier $b_\gamma(X)$ of V2.

\begin{proposition}[Exact one-group factor and a uniform logarithmic remainder]
\label{f16v37:kinetic}
For $\gamma>0$,
\begin{equation}
 z_\gamma=\frac{2e^{-\gamma}I_1(\gamma)}{\gamma},\qquad
 \mathfrak b_\gamma=\sqrt{2\pi\gamma}\,e^{-\gamma}I_1(\gamma).
 \label{f16v37:Bessel}
\end{equation}
There are numerical $C,\gamma_*<\infty$ such that, for $\gamma\ge\gamma_*$,
\begin{equation}
 \left|\log \mathfrak b_\gamma+\frac3{8\gamma}+\frac3{16\gamma^2}\right|
                                      \le C\gamma^{-3}.
 \label{f16v37:kinetic-expansion}
\end{equation}
In particular $-24Bn\log \mathfrak b_\gamma=9Bn/\gamma+O(Bn/\gamma^2)$,
with a constant independent of $B,n$.
\end{proposition}
\begin{proof}
The normalized class angle density is $(2/\pi)\sin^2\theta\,d\theta$.
The integral representation for $I_1$ gives the first equality; inserting the
original $c_{H,\gamma}=(2\pi^2\gamma^{3/2})^{-1}$ gives the second.
One may also set $u=\gamma(1-\cos\theta)$ directly, obtaining
\begin{equation}
 \mathfrak b_\gamma=\frac2{\sqrt\pi}\int_0^{2\gamma}
       e^{-u}u^{1/2}\left(1-\frac{u}{2\gamma}\right)^{1/2}du.
 \label{f16v37:kinetic-real}
\end{equation}
On $0\le u\le\gamma$, Taylor's formula gives the square root
$1-u/(4\gamma)-u^2/(32\gamma^2)$ with remainder at most
$C u^3/\gamma^3$. The remaining interval and the omitted tails of the three
polynomial integrals are exponentially small. The gamma moments with shape
$3/2$ are $\mathbb Eu=3/2$ and $\mathbb Eu^2=15/4$. Thus
\[
 \mathfrak b_\gamma=1-\frac3{8\gamma}-\frac{15}{128\gamma^2}+O(\gamma^{-3}).
\]
It is bounded away from zero eventually. Taylor's formula for the logarithm
gives the second coefficient $-15/128-9/128=-3/16$ and a uniform remainder.
This proves the required estimate without assuming convergence of a Bessel
asymptotic series. Raising the normalized factor to its actual $24Bn$ power
multiplies its logarithm; it does not produce independent uncontrolled errors.
\end{proof}

\subsection{The complete identity-endpoint kernel ratio}

Retain V36's joint coefficient $\mathfrak p_{B,n}$ and
$\rho=dn^4/\gamma$, $d=(72n-21)B$. Define
\begin{equation}
 \mathfrak a_{B,n}=\mathfrak p_{B,n}-\mathfrak s_B+9Bn.
 \label{f16v37:complete-coefficient}
\end{equation}
The joint $\mathfrak p_{B,n}$ includes its original half-Haar endpoint terms.
The static $\mathfrak s_B$ has no Haar term for the exact reason proved above.
The final summand is the one-group kinetic normalization, not an adjustable
vacuum subtraction.

\begin{theorem}[All scalar factors in the dressed identity-endpoint comparison]
\label{f16v37:main}
Assume V36's entry conditions at the requested $s\ge1$,
\eqref{f16v37:static-entry}, $\gamma\ge\gamma_*$, and $\rho\le1$.
Then the original kernels satisfy the exact identity
\begin{equation}
 \log\frac{\mathcal A_{\gamma,B,n}^{00}}{\mathcal A_{0,B,n}^{00}}
 =\log\frac{\mathcal Z_\gamma}{\mathcal Z_0}
  -\log\frac{\mathcal S_{\gamma,B}(I)}{\mathcal S_{0,B}}
                                      -24Bn\log \mathfrak b_\gamma.
 \label{f16v37:exact-scalar}
\end{equation}
Moreover $|\mathfrak a_{B,n}|\le Cdn^4$ and
\begin{equation}
 \boxed{\quad
 \left|\log\frac{\mathcal A_{\gamma,B,n}^{00}}{\mathcal A_{0,B,n}^{00}}
                    -\frac{\mathfrak a_{B,n}}\gamma\right|
     \le C\left\{\rho^{3/2}+\sigma^{3/2}+\frac{Bn}{\gamma^2}
                        +\gamma^{-s}\right\}
     \le C(\rho^{3/2}+\gamma^{-s}).\quad}
 \label{f16v37:kernel-expansion}
\end{equation}
In particular $|\log(\mathcal A_{\gamma,B,n}^{00}/\mathcal A_{0,B,n}^{00})|
\le C(\rho+\gamma^{-s})$. The entire ratio tends to one in V35's stated
polynomial regime. No scalar factor in V36's decomposition is left unestimated
for this identity-endpoint ratio.
\end{theorem}
\begin{proof}
Divide the two exact formulas \eqref{f16v37:kernel-factors}; this proves
\eqref{f16v37:exact-scalar} before any approximation. V36 supplies the joint
mass expansion with error $C(\rho^{3/2}+\gamma^{-s})$.
Theorem~\ref{f16v37:static-pressure} supplies the static term; the kinetic
proposition supplies $9Bn/\gamma$ with error $CBn/\gamma^2$. Their signs
in \eqref{f16v37:complete-coefficient} follow from the exact quotient.

Since $d_{\rm e}=15B\le d$ and $n\ge2$, $\sigma\le\rho$. Also
$Bn/\gamma^2\le C\rho^{3/2}$ for $\gamma\ge2$, $B\ge1$, $n\ge2$.
The coefficient bound follows from V36 and $|\mathfrak s_B|\le15CB$.
As $\rho\le1$, the last claim follows. V35's regime implies
$B/\gamma\to0$ and $\log\gamma/\gamma\to0$, so it eventually satisfies the
separate static entry condition as well. The constants and finite thresholds
in these estimates are not numerically optimized or enclosed.
\end{proof}

It is also legitimate to leave the exactly known kinetic factor unexpanded:
\[
 \left|\log\frac{\mathcal A_{\gamma,B,n}^{00}}{\mathcal A_{0,B,n}^{00}}
       +24Bn\log \mathfrak b_\gamma
       -\frac{\mathfrak p_{B,n}-\mathfrak s_B}{\gamma}\right|
                 \le C(\rho^{3/2}+\sigma^{3/2}+\gamma^{-s}).
\]
This is an exact resummation of one explicit scalar only. It is not a new
resummation of the coupled field theory. Neither $\mathfrak p_{B,n}$ nor
$\mathfrak s_B$ is claimed numerically evaluated by its Gaussian definition.

\begin{proposition}[Length ratios cancel the static endpoint integral exactly]
\label{f16v37:length-ratio}
At one common $B,\gamma$ and any two integer durations $n,k\ge2$, set
$\mathcal R_{n,k}=\mathcal A_{\gamma,B,n}^{00}\mathcal A_{0,B,k}^{00}/
(\mathcal A_{0,B,n}^{00}\mathcal A_{\gamma,B,k}^{00})$.
Then
\begin{equation}
 \log\mathcal R_{n,k}
 =\log\frac{\mathcal Z_{\gamma,B,n}}{\mathcal Z_{0,B,n}}
  -\log\frac{\mathcal Z_{\gamma,B,k}}{\mathcal Z_{0,B,k}}
                                     -24B(n-k)\log \mathfrak b_\gamma.
 \label{f16v37:length-exact}
\end{equation}
If both durations satisfy V36's entry conditions, this has coefficient
$[\mathfrak p_{B,n}-\mathfrak p_{B,k}+9B(n-k)]/\gamma$ and error at most
\[
 C\{\rho_n^{3/2}+\rho_k^{3/2}+B|n-k|/\gamma^2+\gamma^{-s}\},
 \qquad \rho_j=(72j-21)Bj^4/\gamma.
\]
No static approximation or division by a possibly tiny approximate endpoint
factor is needed for this ratio. The errors for the two durations are added;
they are not asserted to cancel or to gain a factor $|n-k|$.
\end{proposition}
\begin{proof}
The same static factor occurs in both exact diagonal kernels, since the
carrier and endpoints are unchanged. Subtract \eqref{f16v37:exact-scalar} at
the two durations and apply V36 and \eqref{f16v37:kinetic-expansion}.
\end{proof}

\subsection{What this closes, and the next spectral obligation}

At fixed $B$ and $n\asymp\gamma^{1/8}$, the full log-ratio bound is
$O(\gamma^{-3/8})$ and the error after its explicit first coefficient is
$O(\gamma^{-9/16})$, with $s\ge1$. These are the same joint scales as V36,
now for the entire original dressed endpoint quantity rather than its chart
mass alone. The result does not enlarge the admissible growth regime.

There is no remaining carrier-chart-mass problem in this scalar: its
cancellation is exact. The next normalization question is instead a spectral
one. A diagonal compressed kernel and finite-duration ratios are not the
spectral top of the full transfer. An elementary example makes the distinction
precise. For $0<a<b<1$, take
$T_0=\operatorname{diag}(a,0)$, $T_1=\operatorname{diag}(a,b)$ and
$f_\epsilon=(\sqrt{1-\epsilon},\sqrt\epsilon)$. Their moments differ by
$\epsilon b^j$ at every positive integer $j$, yet their top eigenvalues are
$a$ and $b$. At any finite list of durations the relative moment differences
can be made arbitrarily small by choosing $\epsilon>0$ sufficiently small.
This counterfixture is not a native model; it excludes an algebraic inference
from a controlled finite family of scalar moments to a spectral-top estimate.

Thus a useful next transfer result must control the relevant physical
spectral measure, its vacuum overlap and its complementary channels, or provide
an appropriate same-top operator comparison. The existing source-relative
temporal-inverse problem is also unaffected. The memoranda already separate
an auxiliary functional inequality from this physical-transfer implication;
closing the literal scalar identity does not remove that separation.
Root Gauss averaging, endpoint-varying matrix elements, independent physical
time calibration, unrestricted regulators and an interacting continuum
construction remain unpaid. The covariance and process theorems of V35--V36
are not replaced by a claim about the physical spectrum.

\subsection{Verification and predecessor preservation}

The accompanying exact diagnostic is
\begin{center}\small\path{verify_ym_f16_v37_scalar_closure.py}.\end{center}
It checks the carrier/Jacobian cancellation, non-factorized static Gaussian
normalizers, normalized interpolation and Wick coefficients without a spurious
static Haar term, the single-group real integral coefficients, signs and powers
in the full scalar quotient, length-ratio cancellation, and the finite-moment
spectral counterfixture. The actual report records assertion counts and scope;
these finite tests are not native high-dimensional integral enclosures or a
formal proof of the analytic estimates.

Only the literal version in the title and this terminal insertion change V36.
Removing those edits recovers its bytes exactly. All available predecessor
checkers are copied without modification; their rerun outputs and hashes are
reported separately from inherited reports. The static proof uses fixed-cube
Gaussian coercivity and the original mixed-word derivatives, not a new spatial
mixing hypothesis. The final joint scalar theorem explicitly uses V36's pressure
estimate. The full analytic lineage has not been independently recertified.
The next scheduled companion checkpoint remains V40.

\begin{firewall}
V37 cancels the original carrier chart mass exactly, estimates the complete
static magnetic endpoint mass, and includes the Wilson kinetic normalization
with its correct coefficient. It thereby controls every scalar factor in the
identity-endpoint dressed diagonal-kernel comparison in the inherited polynomial
regime. It also gives a static endpoint-vector comparison and exact endpoint
cancellation in length ratios. It does not prove a spectral-top estimate,
root-Gauss-reduced vacuum convergence, an unrestricted duration or volume theorem,
a physical time identification, or a Yang--Mills continuum mass gap. No original
measure, carrier, endpoint factor or update clock is changed.
\end{firewall}
% END F16 V37 APPEND

% BEGIN F16 V38 APPEND -- 2026-09-18
% Insert immediately before the final document-ending command of V37.
% The predecessor body, endpoint factors, transfer and clocks are unchanged.

\clearpage
\section{V38: endpoint-weighted spectral edges and vacuum visibility}
\label{f16v38:context}

V37 closes the scalar normalization of the identity-endpoint kernel. It does
not identify that scalar with a spectral top. We now load the entire sequence
into the spectral measure of the \emph{actual} full transfer, using one exact
transfer to regularize the endpoint distribution. This produces a new native
conclusion: at fixed spatial width and the admitted polynomial durations, the
endpoint-weighted spectrum has a Gamma edge profile relative to the Gaussian
threshold. It also proves that this particular filtered endpoint vector has
vanishing overlap with the native Perron vector along that limit. Thus scalar
accuracy is not the missing vacuum-preparation estimate.

The conclusion concerns the unprojected root-coordinate transfer used by the
preceding history integrals. Its Perron top agrees with its root-invariant top,
but its endpoint spectral measure need not be supported in that physical
subspace. We give the exact projection split. Neither a native upper spectral
edge, a first-excited ratio, nor a physical time scale is obtained by deleting
that distinction.

\subsection{Dependencies and the endpoint vector that is actually in the Hilbert space}

The supplied V37 appendix and audit were read completely. Its scalar theorem
is the only native asymptotic input below. V2 supplies the full Gaussian
transfer, with its actual kinetic metric; the static null-space classification
is the inherited f15 V99 result used in V33. The Grand Tensor's transfer-moment
and source-cyclic reconstruction, and f15 V80/V89's overlap and all-duration
warnings, already contain the relevant general spectral principles. We use
those principles, not present them as new general theorems. The new objects
are this endpoint sequence's exact Hilbert-space realization, quantitative
Perron visibility, full Gaussian edge exponent, and the native edge limit.
Targeted archive searches and passage reads are recorded separately; they do
not certify the entire archive. The next companion checkpoint remains V40.

For classical input, the Mehler identity is
\href{https://dlmf.nist.gov/18.18.E28}{DLMF 18.18.28}; its normalization is checked
below. Blossier--Della Morte--von Hippel--Mendes--Sommer,
\href{https://arxiv.org/abs/0902.1265}{\emph{On the generalized eigenvalue method
for energies and matrix elements in lattice field theory}}, provides context
for the distinction between correlation data, overlap and spectral extraction.
The abstract and introductory spectral representation were inspected; no
physical overlap or uniform gap hypothesis from that paper is imported here.
The moment-convergence argument needed below is proved directly.

Fix $B=m^3$, $m\ge4$, and $\gamma>0$. Work on the unchanged pre-root-restriction
Hilbert space
\[
 \mathcal H_{\mathrm{rt},B}=L^2(\SU(2)^{5B}\times\SU(2)^{12B},dH),\qquad
 T_\gamma=T_{\mathrm{full},\gamma}.
\]
The inherited operator is a positive self-adjoint contraction with a strictly
positive bounded integral kernel at each finite regulator. It is compact.
Let $\Lambda_\gamma=\|T_\gamma\|$ and let $\Omega_\gamma>0$ be its unit Perron
vector. Positivity here includes operator positivity, not just positive kernel
entries. Let $\Pi_{\rm G}$ be the orthogonal average over the remaining root
gauge action. The transfer commutes with that action. Uniqueness and positivity
of $\Omega_\gamma$ imply
\begin{equation}
 \Pi_{\rm G}\Omega_\gamma=\Omega_\gamma,\qquad
 \|T_\gamma|_{\operatorname{Ran}\Pi_{\rm G}}\|=\Lambda_\gamma.
 \label{f16v38:same-top}
\end{equation}
This is the inherited finite-regulator Perron argument, not a new continuum
Gauss theorem.

Let $a_{\gamma,I}(X)$ be the original unit conditional endpoint vector of V1.
A delta at the retained identity is \emph{not} an $L^2$ vector. Define instead
\begin{equation}
 b_\gamma^{\rm end}(Q)=j_{\rm br,0}^{1/2}
    \int T_\gamma(Q;X,I)a_{\gamma,I}(X)\,dH^{5B}(X),\qquad
 j_{\rm br,0}=c_{H,\gamma}^{12B},\qquad
 f_\gamma=b_\gamma^{\rm end}/\|b_\gamma^{\rm end}\|.
 \label{f16v38:loaded-vector}
\end{equation}
Here $Q=(X',Z')$ is a full root-coordinate slice. The scalar $j_{\rm br,0}$
is the coarse bridge-chart Jacobian at zero. It accounts for the two endpoint
half-densities in the diagonal chart kernel, rather than being absorbed into
an unspecified source normalization.

\begin{proposition}[One common spectral measure for the complete endpoint sequence]
\label{f16v38:moments}
The vector $b_\gamma^{\rm end}$ is nonzero, positive and in $\mathcal H_{\mathrm{rt},B}$. For every
integer $j\ge0$,
\begin{equation}
 \|b_\gamma^{\rm end}\|^2=\mathcal A_{\gamma,B,2}^{00},\qquad
 m_{\gamma,j}:=\langle f_\gamma,T_\gamma^j f_\gamma\rangle
       =\frac{\mathcal A_{\gamma,B,j+2}^{00}}
                    {\mathcal A_{\gamma,B,2}^{00}}
       =\int_{[0,1]}\lambda^j\,d\mu_\gamma(\lambda),
 \label{f16v38:moment-identity}
\end{equation}
where $\mu_\gamma$ is a probability spectral measure. No intermediate dressed
projection is inserted into $T_\gamma^j$.
\end{proposition}
\begin{proof}
The kernel is bounded at fixed regulator and $\|a_{\gamma,I}\|_1\le1$ on the
probability Haar space. Thus the displayed integral is bounded and belongs to
$L^2$; strict kernel positivity makes it nonzero. Symmetry, composition and
Tonelli give
\[
 \langle b_\gamma^{\rm end},T_\gamma^j b_\gamma^{\rm end}\rangle
 =j_{\rm br,0}\int a_{\gamma,I}(X)
 T_\gamma^{j+2}(X,I;X',I)a_{\gamma,I}(X')\,dH(X)dH(X').
\]
This is exactly the coarse diagonal kernel in Haar coordinates times its
zero-bridge chart Jacobian. V4's half-density formula identifies it with
$\mathcal A_{\gamma,B,j+2}^{00}$. Taking $j=0$ and then normalizing proves
the first two identities. The spectral theorem for the same positive
contraction gives the last one. This regularization is performed before any
use of moments; the delta itself is never treated as a unit vector.
\end{proof}

For $n\ge3$ define its endpoint-weighted probability
\begin{equation}
 d\mu_{\gamma,n}(\lambda)
   =\frac{\lambda^{n-2}\,d\mu_\gamma(\lambda)}{m_{\gamma,n-2}}.
 \label{f16v38:tilt}
\end{equation}
It has no mass at zero, and its $j$th moment is
$\mathcal A_{\gamma,B,n+j}^{00}/\mathcal A_{\gamma,B,n}^{00}$.
Equivalently it is the spectral measure of the unit vector
\[
 v_{\gamma,n}=T_\gamma^{(n-2)/2}f_\gamma/
                      \|T_\gamma^{(n-2)/2}f_\gamma\|.
\]
Fractional powers here are spectral calculus of the original positive transfer,
not a newly chosen physical interpolation. Even $n$ uses only integer powers.

\subsection{A quantitative native top weight, with its actual size cost}

The full kernel factors as
$T_\gamma(Q,Q')=b_{\rm mag}(Q)L_\gamma(Q,Q')b_{\rm mag}(Q')$, where
$b_{\rm mag}=b_\gamma(X)m_\gamma(X,Z)>0$ denotes only the original magnetic
multiplier. This notation is distinct from the vector $b_\gamma^{\rm end}$ above.
The root-fixed integral $L_\gamma$ contains exactly $E=24B$ Wilson kinetic
factors. Therefore
\[
 \ell_\gamma\le L_\gamma(Q,Q')\le u_\gamma,\qquad
 \ell_\gamma=z_\gamma^{-E}e^{-2\gamma E},\quad
 u_\gamma=z_\gamma^{-E},\quad u_\gamma/\ell_\gamma=e^{48\gamma B}.
\]
These inequalities hold on the full compact configuration space; they require
neither a small-field event nor a comparison with a Gaussian.

\begin{theorem}[Perron visibility and finite-duration top brackets]
\label{f16v38:visibility}
At every finite $B,\gamma$,
\begin{equation}
 q_\gamma:=|\langle\Omega_\gamma,f_\gamma\rangle|^2
       \ge \operatorname{sech}^2(48\gamma B)
       \ge e^{-96\gamma B}.
 \label{f16v38:overlap}
\end{equation}
The same lower bound holds after normalizing $\Pi_{\rm G}f_\gamma$. For every
integer $j\ge1$,
\begin{equation}
 m_{\gamma,j}^{1/j}\le\Lambda_\gamma
 \le\min\{1,e^{96\gamma B/j}m_{\gamma,j}^{1/j}\}.
 \label{f16v38:top-bracket}
\end{equation}
Thus these moments determine the same physical top when all durations are
available at one fixed regulator. Their finite-duration upper certificate has
logarithmic width at most $96\gamma B/j$, not a size-free width.
\end{theorem}
\begin{proof}
The quotient of the positive vector $b_\gamma^{\rm end}$ by $b_{\rm mag}$ lies between
$\ell_\gamma A$ and $u_\gamma A$ for a positive scalar $A$. The Perron equation
places $\Omega_\gamma/b_{\rm mag}$ between
$\ell_\gamma A'$ and $u_\gamma A'$ for another positive scalar. Consequently
$h=b_\gamma^{\rm end}/\Omega_\gamma$ has essential range $[m,M]$ with
$M/m\le(u_\gamma/\ell_\gamma)^2$. No bound on the magnetic multiplier itself
has been multiplied into this ratio.

For a positive random variable $m\le h\le M$,
$(h-m)(M-h)\ge0$ gives
$\mathbb Eh^2\le(m+M)\mathbb Eh-mM$.
Minimizing $(\mathbb Eh)^2/\mathbb Eh^2$ over $\mathbb Eh\in[m,M]$ yields
$4mM/(m+M)^2$. Apply this under $\Omega_\gamma^2dH$ to obtain
\[
 \frac{|\langle\Omega_\gamma,b_\gamma^{\rm end}\rangle|^2}{\|b_\gamma^{\rm end}\|^2}
 \ge\frac{4R^2}{(1+R^2)^2}=\operatorname{sech}^2(\log R),\qquad
 R=e^{48\gamma B}.
\]
The final elementary bound in \eqref{f16v38:overlap} follows from
$\cosh t\le e^t$ for $t\ge0$. The gauge projection preserves the numerator
and cannot increase the denominator; it is nonzero by the positive numerator.
Finally the top atom and the upper spectral support give
$q_\gamma\Lambda_\gamma^j\le m_{\gamma,j}\le\Lambda_\gamma^j$.
Take roots and use the overlap bound. No excited spectral ratio is estimated.
\end{proof}

For later use the projection split is exact. Put
$\theta_{\gamma,n}=\|\Pi_{\rm G}v_{\gamma,n}\|^2$. Commutation gives
\begin{equation}
 \mu_{\gamma,n}=\theta_{\gamma,n}\mu^{\rm G}_{\gamma,n}
                  +(1-\theta_{\gamma,n})\mu^{\perp}_{\gamma,n}.
 \label{f16v38:Gauss-mixture}
\end{equation}
The components are the spectral measures of the two normalized projections;
a zero-weight component is omitted. The second has no atom at $\Lambda_\gamma$.
There is no estimate here making $\theta_{\gamma,n}$ close to one. In particular,
root projection cannot be inserted into V37's scalar asymptotic without new
matrix-element information.

\subsection{The full Gaussian threshold and its polynomial endpoint factor}

We next compute a \emph{fixed-$B$} Gaussian long-duration asymptotic. In this
subsection every spatial matrix includes its three-colour tensor. On
$q=(x,y)\in\mathbb R^{51B}$ put
\begin{equation}
 \mathsf A_{\rm kin}=\operatorname{diag}(\mathsf M_B,G),\qquad
 \mathsf S_B=\begin{pmatrix}
  \mathsf H_B+A^{\mathsf T}A&A^{\mathsf T}D\\
  D^{\mathsf T}A&D^{\mathsf T}D
 \end{pmatrix},\qquad
 \mathsf V_B=\mathsf A_{\rm kin}^{1/2}\mathsf S_B\mathsf A_{\rm kin}^{1/2}.
 \label{f16v38:Gaussian-transfer-data}
\end{equation}
These are the complete \emph{uncompressed} Gaussian transfer data of V2.
Let $\nu_1,\ldots,\nu_{51B}$ be the eigenvalues of $\mathsf V_B$ and set
\begin{equation}
 \Lambda_{{\rm G},B}
 =\prod_{\nu_j>0}\left(1+\frac{\nu_j}{2}
                    +\sqrt{\nu_j+\nu_j^2/4}\right)^{-1/2},\qquad
 h_B=3(B+2),\qquad a_B=h_B/2.
 \label{f16v38:Gaussian-top}
\end{equation}
The letter $a_B$ here is an edge exponent, not a carrier or V37's coefficient
$\mathfrak a_{B,n}$.

\begin{theorem}[Gaussian edge exponent for the actual dressed endpoint]
\label{f16v38:Gaussian-edge}
For every fixed $B$, the norm of the full Gaussian transfer $T_0$ is
$\Lambda_{{\rm G},B}\in(0,1)$. There are $c_B>0$ and real $d_{B,1},d_{B,2}$
such that, as integer $n\to\infty$,
\begin{equation}
 \mathcal A_{0,B,n}^{00}
 =c_B\Lambda_{{\rm G},B}^{\,n}n^{-a_B}
       \{1+d_{B,1}/n+d_{B,2}/n^2+O_B(n^{-3})\}.
 \label{f16v38:Gaussian-asymptotic}
\end{equation}
The null count $h_B$ includes the $3(B-1)$ gradient directions and nine period
directions before root Gauss reduction. In particular the exponent is not
merely $9/2$. No uniformity of the constants as $B\to\infty$ is asserted.
\end{theorem}
\begin{proof}
By the inherited positive internal stiffness,
$q^{\mathsf T}\mathsf S_Bq=x^{\mathsf T}\mathsf H_Bx+|Ax+Dy|^2$ vanishes
exactly when $x=0$ and $Dy=0$. The f15 V99 kernel classification therefore gives
nullity $h_B$. Congruence by the positive kinetic square root preserves that
nullity. Whitening the kinetic covariance and orthogonally diagonalizing
$\mathsf V_B$ are unitary coordinate changes in $L^2$ with their determinants
included. The transfer becomes a tensor product of one-dimensional kernels
\[
 t_\nu(u,v)=(2\pi)^{-1/2}
  \exp[-(u-v)^2/2-\nu(u^2+v^2)/4].
\]
For $\nu>0$, let $\xi_\nu=\operatorname{arcosh}(1+\nu/2)$ and
$\omega_\nu=\sqrt{\nu+\nu^2/4}$. Mehler's identity gives eigenvalues
$e^{-(k+1/2)\xi_\nu}$ on the oscillator basis of frequency $\omega_\nu$.
One may check the ground factor directly by integrating
$t_\nu(u,v)e^{-\omega_\nu v^2/2}$; its factor is
$(1+\nu/2+\omega_\nu)^{-1/2}$. For $\nu=0$, unitary Fourier transformation
gives multiplier $e^{-p^2/2}$ with norm one, not a normalizable ground vector.
This proves \eqref{f16v38:Gaussian-top} and the norm assertion.

The Gaussian counterpart $b_0=T_0(\psi_0\delta_0)$ of
\eqref{f16v38:loaded-vector} is a strictly positive centered nondegenerate
Gaussian function in all $51B$ coordinates: integrate the original static
Gaussian in its internal endpoint, keeping the positive kinetic square for
the retained output. Its norm square is $\mathcal A_{0,B,2}^{00}$.
After the preceding unitary changes it is still a positive Gaussian. Denote
its coordinates by $(z,s)$, with $z\in\mathbb R^{h_B}$ free and $s$ massive,
and let $\varphi_0(s)>0$ be the unit massive ground function. Then
\[
 F_B(z)=\int\varphi_0(s)\widetilde b_0(z,s)\,ds
\]
is a positive nondegenerate Gaussian. With unitary Fourier convention,
$\widehat F_B(0)>0$ and
$|\widehat F_B(p)|^2=C_B^{\rm F}e^{-p^{\mathsf T}D_Bp}$ for some $D_B\succ0$.
These constants are determined by the original endpoint, not freely chosen.

Put $N=n-2$. The massive ground contribution to
$\langle b_0,T_0^Nb_0\rangle$ is exactly
\[
 \Lambda_{{\rm G},B}^{N}C_B^{\rm F}(2\pi)^{h_B/2}
                        \det(NI+2D_B)^{-1/2}.
\]
Every other massive component has its transfer norm at most
$\Lambda_{{\rm G},B}e^{-\xi_{\min}}$, where the minimum is over the finitely
many positive $\nu_j$. Its total contribution is bounded by
$\Lambda_{{\rm G},B}^{N}e^{-N\xi_{\min}}\|b_0\|^2$, using the orthogonal
spectral decomposition, not an independent-mode assumption for $b_0$.
Expand $\det(I+(2D_B-2I)/n)^{-1/2}$ through order two with its finite third-order
remainder. The exponentially small massive contribution can be included in
$O_B(n^{-3})$. This proves \eqref{f16v38:Gaussian-asymptotic}. It also proves
$c_B=\Lambda_{{\rm G},B}^{-2}C_B^{\rm F}(2\pi)^{h_B/2}>0$.
\end{proof}

The determinant expression, rather than an unspecified $O(1/n)$ error, also
justifies the adjacent-ratio expansions
\begin{equation}
 \begin{split}
 \frac{\mathcal A_{0,B,n+1}^{00}}{\Lambda_{{\rm G},B}\mathcal A_{0,B,n}^{00}}
        &=1-a_B/n+O_B(n^{-2}),\\
 \frac{\mathcal A_{0,B,n}^{00}\mathcal A_{0,B,n+2}^{00}}
          {(\mathcal A_{0,B,n+1}^{00})^2}-1
        &=a_B/n^2+O_B(n^{-3}).
 \end{split}
 \label{f16v38:Gaussian-ratios}
\end{equation}
Their errors are not obtained by differentiating an uncontrolled asymptotic.

\subsection{A native endpoint-weighted spectral law}

Fix $B$ henceforth. Let $n=n_\gamma\to\infty$ be integers satisfying
\begin{equation}
 \frac{Bn_\gamma^7\log(n_\gamma+1)+n_\gamma^6\log\gamma}{\gamma}\longrightarrow0.
 \label{f16v38:regime}
\end{equation}
Every fixed integer multiple of $n_\gamma$, and the two adjacent durations,
then satisfy V37's entry conditions eventually. One concrete choice is
$n_\gamma=2\lfloor\gamma^\alpha/2\rfloor$, $0<\alpha<1/7$.

\begin{theorem}[Native Gamma edge profile without a claimed top identification]
\label{f16v38:native-edge}
Under $\mu_{\gamma,n_\gamma}$, define
\begin{equation}
 X_\gamma=-n_\gamma\log(\lambda/\Lambda_{{\rm G},B}),\qquad
 S_\gamma=n_\gamma(1-\lambda/\Lambda_{{\rm G},B}).
 \label{f16v38:edge-variables}
\end{equation}
Both variables converge in distribution to the Gamma law with shape $a_B$ and
rate one, whose density is
\begin{equation}
 1_{\{x>0\}}\,\frac{x^{a_B-1}e^{-x}}{\Gamma(a_B)}\,dx.
 \label{f16v38:Gamma}
\end{equation}
In addition,
\begin{equation}
 \mathbb E_{\mu_{\gamma,n_\gamma}}S_\gamma\longrightarrow a_B,
 \qquad
 \operatorname{Var}_{\mu_{\gamma,n_\gamma}}S_\gamma\longrightarrow a_B.
 \label{f16v38:edge-moments}
\end{equation}
At finite coupling these edge variables can be negative; the native top is not
replaced by $\Lambda_{{\rm G},B}$. The limit is a theorem about the spectral
measure of the specified endpoint vector of the \emph{native} transfer.
\end{theorem}
\begin{proof}
Write $A_{\gamma,n}=\mathcal A_{\gamma,B,n}^{00}$ for this proof only.
V37 gives
$|\log(A_{\gamma,n}/A_{0,n})|\le C(\rho_n+\gamma^{-s})$,
where $\rho_n=(72n-21)Bn^4/\gamma$. Choose $s=2$.
For each fixed integer $k\ge0$, put
$Y_\gamma=(\lambda/\Lambda_{{\rm G},B})^{n_\gamma}$. The exact tilted moments are
\begin{equation}
 \mathbb E Y_\gamma^k
  =\Lambda_{{\rm G},B}^{-kn_\gamma}
                     \frac{A_{\gamma,(k+1)n_\gamma}}{A_{\gamma,n_\gamma}}
  \longrightarrow(k+1)^{-a_B}.
 \label{f16v38:edge-power-moments}
\end{equation}
Use \eqref{f16v38:Gaussian-asymptotic} and the errors at \emph{both} durations.
No cancellation of those errors is assumed, and $k$ is fixed before the limit.
All original scalar factors are included through V37; a comparison of normalized
path probabilities alone would not imply this formula.

For completeness these moments determine an actual probability limit.
The bounded first moments give tightness of $Y_\gamma\ge0$. The bounded
$(k+1)$st moments give uniform integrability of its $k$th moments. Hence every
subsequential limit has moments $(k+1)^{-a_B}$. It is supported on $[0,1]$:
positive mass above $1+\delta$ would force its high moments to grow exponentially,
contrary to that formula. On $[0,1]$, polynomials determine a probability measure
by uniform approximation of continuous functions. The candidate $e^{-G}$,
$G\sim\operatorname{Gamma}(a_B,1)$, has exactly these moments. Thus
$Y_\gamma\Rightarrow e^{-G}$. The limit has no mass at zero. The continuous
mapping theorem applied to $-\log Y_\gamma$ proves the first limit, and uniform
convergence of $n(1-e^{-x/n})$ to $x$ on bounded $x$ proves the second.

For the two moment assertions let $r_n=A_{\gamma,n+1}/A_{\gamma,n}$ and
$v_n=A_{\gamma,n}A_{\gamma,n+2}/A_{\gamma,n+1}^2-1\ge0$. They are exactly
the mean spectral value and its relative variance under $\mu_{\gamma,n}$.
Equation \eqref{f16v38:Gaussian-ratios}, with the V37 errors added, gives
\[
 r_n/\Lambda_{{\rm G},B}=1-a_B/n+O_B(n^{-2})+O(\rho_{n+2}+\gamma^{-2}),
 \qquad
 v_n=a_B/n^2+O_B(n^{-3})+O(\rho_{n+2}+\gamma^{-2}).
\]
The exact formulas are
$\mathbb E S_\gamma=n(1-r_n/\Lambda_{{\rm G},B})$ and
$\operatorname{Var}S_\gamma=n^2(r_n/\Lambda_{{\rm G},B})^2v_n$.
The regime implies $n^2(\rho_{n+2}+\gamma^{-2})\to0$, proving
\eqref{f16v38:edge-moments}. This is a finite-difference use of quantified
scalar errors, not differentiation of those errors.
\end{proof}

For each fixed $x>0$, the exact spectral probability satisfies
\[
 \mu_{\gamma,n_\gamma}
 \{\lambda\ge\Lambda_{{\rm G},B}e^{-x/n_\gamma}\}
 \longrightarrow \frac1{\Gamma(a_B)}\int_0^x t^{a_B-1}e^{-t}\,dt>0.
\]
Weight above $\Lambda_{{\rm G},B}e^{x/n_\gamma}$ tends to zero. These are
weights, not assertions that no eigenvalue exists outside the stated band.

\subsection{What the edge law proves about the native top and the filtered vacuum weight}

\begin{proposition}[A one-sided top estimate and failure of endpoint Perron preparation]
\label{f16v38:consequences}
Along \eqref{f16v38:regime}, at fixed $B$,
\begin{equation}
 n_\gamma\bigl[\log(\Lambda_{{\rm G},B}/\Lambda_\gamma)\bigr]_+
      \longrightarrow0,
 \qquad
 |\langle\Omega_\gamma,v_{\gamma,n_\gamma}\rangle|^2\longrightarrow0.
 \label{f16v38:top-and-weight}
\end{equation}
The first statement also concerns the physical root-invariant top by
\eqref{f16v38:same-top}. It is not an upper estimate for that top. The second
statement concerns the actual unprojected endpoint-filtered vector; it is not
a claim that the normalized root projection has the same limit.
For even durations, let $\pi_{\gamma,n}^{\rm mid}$ be the full slice marginal
at time $n/2$ of the original identity-endpoint return law, in root coordinates.
Then
\begin{equation}
 d\pi_{\gamma,n}^{\rm mid}=|v_{\gamma,n}|^2dH,\qquad
 d_{\rm TV}(\pi_{\gamma,n_\gamma}^{\rm mid},\Omega_\gamma^2dH)
                       \longrightarrow1.
 \label{f16v38:midpoint-nonequilibrium}
\end{equation}
This is a gauge-dependent conditioned midpoint law, not a statement about its
restriction to gauge-invariant observables.
The effective scalar energy also satisfies
\begin{equation}
 -\log\frac{\mathcal A_{\gamma,B,n_\gamma+1}^{00}}
                 {\mathcal A_{\gamma,B,n_\gamma}^{00}}
   =-\log\Lambda_{{\rm G},B}+\frac{a_B}{n_\gamma}+o(n_\gamma^{-1}).
 \label{f16v38:effective-energy}
\end{equation}
The right side is measured against the Gaussian threshold, not an identified
native vacuum energy or a calibrated physical energy.
\end{proposition}
\begin{proof}
For each $x>0$ the positive limiting weight just proved ensures
$\Lambda_\gamma\ge\Lambda_{{\rm G},B}e^{-x/n_\gamma}$ eventually. Taking an
arbitrarily small fixed $x$ gives the first assertion.

A tight sequence of real probability measures converging to a non-atomic law
has maximum atomic mass tending to zero. To verify this directly, an atom of
mass at least $\epsilon$ cannot escape every compact set by tightness. Choose a
convergent subsequence of its locations inside a compact set; the portmanteau
inequality on arbitrarily small closed intervals about the limiting location
would then give an atom of at least $\epsilon$ in the limit. This is impossible
for the Gamma density. Apply this to the law of $X_\gamma$. Its atom at
$-n_\gamma\log(\Lambda_\gamma/\Lambda_{{\rm G},B})$ has weight exactly
$|\langle\Omega_\gamma,v_{\gamma,n_\gamma}\rangle|^2$, since the Perron
projection is one-dimensional. The conclusion also applies to the total weight
of any fixed number of eigenvalues. No rate for this vanishing is claimed.
For even $n$, splitting the exact return integral at its midpoint gives two
identical positive half-path amplitudes. Their squared norm is the full return
normalizer, so the normalized midpoint density is precisely $|v_{\gamma,n}|^2$.
The Hellinger affinity of this density with $\Omega_\gamma^2$ equals
$\langle\Omega_\gamma,v_{\gamma,n}\rangle$, since both vectors are positive.
For probability densities $p,q$, the pointwise bound $\min(p,q)\le\sqrt{pq}$
gives $d_{\rm TV}(p,q)\ge1-\int\sqrt{pq}$. The vanishing squared overlap
therefore proves \eqref{f16v38:midpoint-nonequilibrium}.
Finally the adjacent-ratio estimate in the preceding proof and
$n(\rho_{n+2}+\gamma^{-2})\to0$ give \eqref{f16v38:effective-energy}.
\end{proof}

There is no contradiction with \eqref{f16v38:overlap}. At each fixed $B,\gamma$,
all-duration filtering eventually selects the positive Perron vector. Here
$\gamma$ changes while $n_\gamma=o(\gamma)$, and the proved overlap lower bound
is exponentially small. Its finite-duration top bracket has logarithmic width
$96\gamma B/(n_\gamma-2)$, which does not vanish in this regime. This diagnoses
the weakness of that \emph{certificate}, not a necessary mixing-time lower bound
or an upper bound on the true overlap.

Equation \eqref{f16v38:Gauss-mixture} is especially important now. If
$\theta_{\gamma,n}$ is small, root projection followed by renormalization can
radically change the spectral law and its vacuum weight. The fixed-$B$ Gaussian
exponent above counts linearized gauge-gradient directions before a physical quotient.
It cannot be relabelled a physical particle multiplicity, a toron-only exponent,
or a root-Gauss-reduced Gamma law. No such projected law is proved here.

\subsection{Direction after the spectral calculation and verification boundary}

The immediate unresolved task is now more specific than obtaining another
scalar coefficient: control the root-projected endpoint spectral measure and
its overlap, or analyze the compact soft sector on durations beyond the present
Gaussian entrance. V37's scalar sequence does have a genuine native spectral
meaning, but its polynomial-duration limit is not Perron dominated. Another
unprojected finite-time coefficient will not by itself reverse
\eqref{f16v38:top-and-weight}. The older complementary-channel and same-top
operator obligations remain, as do physical time and continuum reconstruction.

The accompanying diagnostic is
\begin{center}\small\path{verify_ym_f16_v38_endpoint_spectrum.py}.\end{center}
It checks the endpoint moment offset and Jacobian cancellation, normalized
spectral tilting, root projection, the sharp positive-ratio overlap inequality,
Mehler normalization, Gaussian heat powers, moment identification and finite
error exponents. These are finite algebraic fixtures, not numerical enclosures
of the native transfer spectrum or a formal proof of the analytic limit.
Execution and preservation reports state the actual scope. Only the title
version and this terminal insertion change V37. No earlier claim is silently
corrected, and no archive-wide proof certification is asserted.

\begin{firewall}
V38 realizes the exact dressed endpoint sequence as native transfer moments,
quantifies its Perron visibility with a paid coupling/volume cost, and derives
a Gamma endpoint-weighted edge law in the inherited polynomial regime at fixed
spatial width. It proves a one-sided native top estimate and vanishing Perron
weight of that unprojected filtered vector, not a physical excited-state gap.
Root-projected spectral weights, a native upper edge, unrestricted regulators,
physical vacuum identification, time calibration and an interacting continuum
mass gap remain unproved. No intermediate dressed projection or changed scalar
normalization is inserted into the actual transfer powers.
\end{firewall}
% END F16 V38 APPEND

% BEGIN F16 V39 APPEND -- 2026-09-18
% Insert before the final document-ending command in the cumulative V38.
% All predecessor proofs, measures, endpoint factors and clocks remain unchanged.

\clearpage
\section{V39: root projection, orbit weights, and physical endpoint preparation}
\label{f16v39:context}

V38 leaves the weight of the root-invariant component uncontrolled. Its
unprojected endpoint edge law cannot therefore be assigned a physical exponent
by subtracting the number of gauge directions. We return to the compact root
action itself. An exact forest disintegration separates amplitude projection
from averaging a probability density. It also allows a quantitative estimate
of the actual projected norm. At fixed spatial width that norm has a polynomial,
not merely exponentially small, coupling scale.

The resulting estimates do not identify a projected Gamma law. They instead
give two-sided power bounds on the projected endpoint moments. Comparing three
admitted durations proves that, on a smaller but explicit polynomial range,
the normalized \emph{root-projected} filtered endpoint still has vanishing
Perron overlap. Unlike the corresponding V38 conclusion, its midpoint and
stationary densities are both gauge invariant. This is a failure of the
specified source to prepare the finite-regulator vacuum on those durations,
not a conclusion about the physical mass gap.

\subsection{Dependencies and the unchanged compact root action}

The V38 appendix and audit were read completely. The native inputs below are
V35's complete-law tails and local linear-tilt moments, V37's scalar comparison,
and V38's endpoint loading and fixed-width Gaussian asymptotic. We do not assume
an asymptotic for a root-twisted endpoint matrix element. The forest/Haar
mechanism already appears in f15 V10 and is credited. Its application here is
to the remaining \emph{coarse root} group and to a generally non-invariant
endpoint density, not to a Wilson measure asserted to factor along these
orbits. Targeted searches of the active source and supplied archives were
followed by identified passage reads; this is not an exhaustive novelty audit.
The next scheduled checkpoint remains V40.

For context, Sections II--III.1 of Burbano--Bauer's
\href{https://arxiv.org/html/2409.13812}{maximal-tree formulation} distinguish
based gauge fixing from the remaining global constraint. Bahr--Thiemann's
\href{https://arxiv.org/html/0709.4636}{non-Abelian projected-state construction},
equations (2.4)--(2.5), uses normalized compact group averaging. These classical
constructions are context, not native gap hypotheses. We prove the identities
and estimates used here directly, in the inherited Wilson orientation.

Let $G=\SU(2)$, $B=m^3$, $m\ge4$, and retain
\[
 \mathcal H_{\rm rt}=L^2(G^{5B}\times G^{12B},dH),\quad
 T_\gamma=T_{{\rm full},\gamma},\quad
 \Pi=\Pi_{\rm G},\quad \Lambda_\gamma=\|T_\gamma\|.
\]
Write $\mathsf U_g$ for the root action, with
\[
 X_{b,a}\mapsto g_bX_{b,a}g_b^{-1},\qquad
 Z_e\mapsto g_{s(e)}Z_eg_{t(e)}^{-1}.
\]
The moving frames are equivariant under these conjugations. The full transfer
commutes with this action and has the same Perron top on $\operatorname{Ran}\Pi$,
as already established in V38. Let $f_\gamma$ be its \emph{same} once-loaded
unit endpoint vector. For even $n\ge2$ put
\begin{equation}
 v_{\gamma,n}=
 \frac{T_\gamma^{(n-2)/2}f_\gamma}
 {\|T_\gamma^{(n-2)/2}f_\gamma\|},\qquad
 \theta_{\gamma,n}=\|\Pi v_{\gamma,n}\|^2,
 \qquad v^{\rm G}_{\gamma,n}=\frac{\Pi v_{\gamma,n}}{\sqrt{\theta_{\gamma,n}}}.
 \label{f16v39:vectors}
\end{equation}
All are positive; the first is invariant under simultaneous diagonal
conjugation, because the identity endpoint and its original carrier are.
We use even durations so these vectors are obtained by positive integer
transfer powers. No new continuum-time interpolation is required.

\subsection{An exact forest disintegration of projected amplitudes}

Choose a spanning tree $\mathcal F$ on the connected coarse multigraph, using
one actual bridge for each tree edge, and choose one root $o$. Set
\[
 r_B=B-1,\qquad D_B=3(B-1),\qquad \ell_B=D_B/2.
\]
Here $D_B$ is a dimension, not the Fourier covariance matrix used in V38.
Let $t_b$ be the ordered product from $o$ to $b$ along $\mathcal F$, with
$t_o=I$ and inverses on reverse traversals. Define
\[
 \widetilde X_{b,a}=t_bX_{b,a}t_b^{-1},\qquad
 \widetilde Z_e=t_{s(e)}Z_et_{t(e)}^{-1}.
\]
Every tree value is $I$. Let $\eta$ collect all $5B$ transformed chords and
$12B-r_B$ non-tree bridges, with product Haar measure $d\nu(\eta)$.
There is still a global simultaneous conjugation on $\eta$.

\begin{proposition}[Haar fibres and the actual orbit weight]
\label{f16v39:fibres}
This is a global bijection and an exact probability-Haar identity:
\begin{equation}
 dH(X,Z)=\prod_{b\ne o}dH(t_b)\,d\nu(\eta).
 \label{f16v39:Haar}
\end{equation}
For any positive, diagonal-conjugation-invariant unit vector $v$, set
\[
 r(\eta)=\int v(t,\eta)^2\,dt,\qquad
 a(\eta)=\int v(t,\eta)\,dt,\qquad
 k(\eta)=\frac{a(\eta)^2}{r(\eta)},\qquad
 \theta=\int a(\eta)^2\,d\nu(\eta).
\]
Then $0<k\le1$ almost everywhere and $\theta=\|\Pi v\|^2=\int kr\,d\nu$.
The unprojected density restricted to invariant observables and the normalized
projected density have respective quotient representatives
\begin{equation}
 r(\eta)\,d\nu(\eta),\qquad
 \frac{k(\eta)}{\theta}r(\eta)\,d\nu(\eta).
 \label{f16v39:reweight}
\end{equation}
In particular their total variation on invariant events is exactly
\begin{equation}
 \frac12\int r(\eta)\left|\frac{k(\eta)}{\theta}-1\right|d\nu(\eta).
 \label{f16v39:quotient-TV}
\end{equation}
The normalized projected density uniquely maximizes Hellinger affinity with
$v^2dH$ among invariant probability densities, and that maximum is
$\sqrt\theta$.
\end{proposition}
\begin{proof}
The inverse is $X_{b,a}=t_b^{-1}\widetilde X_{b,a}t_b$ and
$Z_e=t_{s(e)}^{-1}\widetilde Z_et_{t(e)}$. In particular a tree bridge is
$t_s^{-1}t_t$. Successive tree Haar translations and then the left/right
translations of all non-tree variables prove \eqref{f16v39:Haar}.
No word is reordered and no determinant is introduced.

The based root group $g_o=I$ changes $t_b$ to $t_bg_b^{-1}$ and leaves $\eta$
fixed. Thus its average sends $v$ to $a(\eta)$. Write a general root element
as $g_b=h_bq$, with $h_o=I$. Diagonal invariance of $v$ shows that full root
averaging equals this based average; $a,r,k$ are themselves invariant under
the remaining simultaneous conjugation of $\eta$. Cauchy--Schwarz on the
probability Haar fibre gives $a^2\le r$. Integration yields the formula for
$\theta$ and both normalized quotient densities. The sign set of $k/\theta-1$
is invariant, so the ordinary quotient total variation is also the supremum
over invariant events. Fibrewise the coefficient is
\[
 k(\eta)=\left(\int\sqrt{v(t,\eta)^2/r(\eta)}\,dt\right)^2.
\]
It is the squared affinity of the actual conditional orbit density with Haar;
that density is not asserted to be Haar.

If $q$ is any invariant probability density, its square root is an invariant
unit vector. Hence $\int v\sqrt q\,dH=\langle\Pi v,\sqrt q\rangle\le\sqrt\theta$.
Equality in Cauchy--Schwarz requires $\sqrt q=\Pi v/\sqrt\theta$ almost
everywhere, proving the final assertion. The transfer and its kinetic energy
have not been changed by the coordinate identity.
\end{proof}

Small $\theta$ alone does not bound \eqref{f16v39:quotient-TV}. For example, if
$v(t,\eta)=u(t)b(\eta)$ with both factors normalized, $k$ is constant and the
two quotient probabilities coincide, even if $\theta$ is small. With an
$\eta$-dependent orbit profile the reweighting is genuine. Probability averaging
and amplitude projection must therefore remain distinct.

\subsection{A polynomial lower bound from the literal kinetic factors}

\begin{proposition}[Orbit Fisher bound and root-projection visibility]
\label{f16v39:lower}
At every finite $B$, $\gamma\ge1$ and even $n\ge2$, there is a constant
$c_B>0$, independent of $n,\gamma$, such that
\begin{equation}
 \theta_{\gamma,n}\ge c_B\gamma^{-\ell_B}.
 \label{f16v39:lower-theta}
\end{equation}
More precisely, for a unit Lie algebra generator acting at a single coarse
root, with skew derivative $\mathcal D_b$,
\begin{equation}
 \|\mathcal D_b v_{\gamma,n}\|_2^2\le48\gamma B.
 \label{f16v39:Fisher}
\end{equation}
These estimates use no small-field event or Gaussian comparison.
\end{proposition}
\begin{proof}
Consider any positive output $v=T_\gamma h/\|T_\gamma h\|$, allowing the
initial endpoint to be a positive finite measure. The vector in
\eqref{f16v39:vectors} has this form. Along a one-root path $e^{t\tau}$,
$|\tau|=1$, the output magnetic factor is exactly invariant. In the kinetic
integrand each of its $24B$ words contains at most two factors $e^{\pm t\tau}$.
The second derivative of its scalar quaternion part has absolute value at
most four: the two second derivatives and two cross derivatives each have
operator norm at most one. Hence the logarithm of the complete kinetic
integrand has second derivative at least $-96\gamma B$.

Integrating the nonroot Haar variables and the positive input preserves this
lower bound on the logarithm, since
\[
 \frac{d^2}{dt^2}\log\int e^{l_t}\,d\sigma
     =\mathbb E_t l_t''+\operatorname{Var}_t(l_t')\ge-96\gamma B.
\]
Here $d\sigma$ includes the unchanged positive input and all other factors.
At a fixed regulator derivatives in these compact group directions are bounded
relative to the integrand, so differentiation is justified. Thus
$\mathcal D_b^2\log v\ge-96\gamma B$.
The action is Haar preserving. Integrating $\mathcal D_b(v^2\mathcal D_b\log v)$
gives zero and proves
\[
 \|\mathcal D_bv\|_2^2
       =-\tfrac12\int v^2\mathcal D_b^2\log v\,dH\le48\gamma B.
\]
This is integration on the compact group, not through a principal-chart cut.
Along an orbit the original equivariant frames have no new cut derivative.

For $h_o=I$, write $h_b=\operatorname{Exp}(x_b)$. Unitarity, the fundamental
theorem along each one-parameter subgroup, and a product telescoping estimate give
\[
 \|\mathsf U_hv-v\|_2\le\sqrt{48\gamma B}\sum_{b\ne o}|x_b|.
\]
On $|x_b|\le\delta=(16r_B\sqrt{B\gamma})^{-1}$ the squared distance is at most
$3/16$, so the real inner product is at least $29/32$. It is nonnegative for
all $h$, by positivity of $v$. The normalized $\SU(2)$ cap satisfies, for
$0<\delta\le1$,
\[
 H\{|x|\le\delta\}
 =\frac2\pi\int_0^\delta\sin^2t\,dt\ge\frac8{3\pi^3}\delta^3.
\]
The $r_B$ based group factors are independent Haar variables. Since
$\theta=\int\langle v,\mathsf U_hv\rangle dh$, integrating over these caps
proves \eqref{f16v39:lower-theta}; for example one may take
\[
 c_B=\frac12\left[\frac8{3\pi^3}(16r_B\sqrt B)^{-3}\right]^{r_B}.
\]
The constant is conservative, but is explicit and contains no duration.
\end{proof}

\subsection{The complementary upper bound from the actual midpoint concentration}

Fix $B$ for the remainder of the asymptotic statements. Choose an integer
$p_B>\ell_B$, for instance $p_B=\lceil\ell_B\rceil+1$, and use V35's complete
comparison-ball entry at tail exponent $s_B=p_B+2$. This is a fixed exponent
once $B$ is fixed. Its asymptotic regime is still V38's fixed-$B$ polynomial
regime. Constants below may depend on $B$; no uniform-in-$B$ estimate is claimed.

\begin{theorem}[The actual root-component scale]
\label{f16v39:upper}
For all sufficiently large couplings satisfying that explicit V35 entry, and
all admitted even durations,
\begin{equation}
 \boxed{\qquad
 c_B\gamma^{-\ell_B}\le\theta_{\gamma,n}
       \le C_B(n/\gamma)^{\ell_B},\qquad \ell_B=\tfrac32(B-1).
 \qquad}
 \label{f16v39:theta-bounds}
\end{equation}
In particular, along the fixed-$B$ regime of V38, $\theta_{\gamma,n_\gamma}\to0$.
At the fixed loading duration $n=2$ it has matching upper and lower powers:
$c_B\gamma^{-\ell_B}\le\theta_{\gamma,2}\le C_B\gamma^{-\ell_B}$.
No leading coefficient or exact $n$-dependence of $\theta_{\gamma,n}$ is asserted.
\end{theorem}
\begin{proof}
For a full slice define the positive tree weight
\[
 A(Z)=1+\frac\gamma n\sum_{e\in\mathcal F}\operatorname{dist}(I,Z_e)^2.
\]
In the forest coordinates this depends only on $t$, not on $\eta$. Its negative
power has an exact Haar upper bound. Tree edge values are independent Haar
under the product fibre reference; their change from $t$ is triangular. Using
complete principal charts only for this reference integral gives
\begin{equation}
 \begin{split}
 J_p:=\int A(t)^{-p}\,dt
 &\le (2\pi^2)^{-r_B}(n/\gamma)^{D_B/2}
           \int_{\mathbb R^{D_B}}(1+|z|^2)^{-p}dz\\
 &= (2\pi^2)^{-r_B}\pi^{D_B/2}
       \frac{\Gamma(p-D_B/2)}{\Gamma(p)}(n/\gamma)^{\ell_B},
       \qquad p>\ell_B.
 \end{split}
 \label{f16v39:negative-moment}
\end{equation}
The Haar Jacobian is bounded above by one; the full native density is not
replaced by this reference.

Weighted Cauchy--Schwarz separately on every orbit gives
\[
 \left(\int v(t,\eta)dt\right)^2
       \le J_p\int v(t,\eta)^2A(t)^pdt.
\]
After integrating $\eta$ this yields the general inequality
$\theta\le J_p\,\mathbb E_{v^2dH}A^p$.
By V38 the native midpoint density for even $n$ is exactly
$v_{\gamma,n}^2dH$. In the complete joint path its tree weight is
\[
 A=1+n^{-1}\sum_{e\in\mathcal F}|y_{n/2,e}|^2.
\]
On V35's joint comparison ball, its proved linear-tilt bound implies
$\mathbb E|y_{n/2,e}|^{2p_B}\le C_{p_B}n^{p_B}$. Summing the fixed $r_B$ sites
therefore bounds the conditional expectation of $A^{p_B}$ by $C_B$.
On the complement, compactness gives $A^{p_B}\le C_B\gamma^{p_B}$ for $n\ge2$,
whereas the complete native tail probability is at most $\gamma^{-s_B}$.
Consequently
\[
 \mathbb E_{v_{\gamma,n}^2dH} A^{p_B}\le C_B+C_B\gamma^{p_B-s_B}\le C_B'.
\]
This uses high \emph{interior} moments derived from the already proved
sub-Gaussian bound; it does not claim an unproved high-degree complete-tail
estimate. The finite cap pays the complement. Equations
\eqref{f16v39:negative-moment} and \eqref{f16v39:lower-theta} prove the theorem.
\end{proof}

The upper bound also implies that the full unprojected midpoint density is far
in total variation from \emph{every} root-invariant probability density: the
bound is at least $1-\sqrt{\theta_{\gamma,n}}$ by Proposition~\ref{f16v39:fibres}.
This observation is gauge dependent and does not alone say that invariant
observables fail to be stationary. The following spectral argument is needed
to obtain a statement for the projected source itself.

\subsection{Root-projected native moments with both polynomial factors retained}

Define the exact normalized physical source and its moments by
\begin{equation}
 f_\gamma^{\rm G}=\Pi f_\gamma/\sqrt{\theta_{\gamma,2}},\qquad
 M_{\gamma,n}^{\rm G}
    =\langle f_\gamma^{\rm G},T_\gamma^{n-2}f_\gamma^{\rm G}\rangle,
 \qquad n\ge2.
 \label{f16v39:projected-moments}
\end{equation}
There is no insertion of a new carrier between transfer powers. Commutation
and \eqref{f16v38:moment-identity} give, for even $n$,
\begin{equation}
 M_{\gamma,n}^{\rm G}
   =\frac{\mathcal A_{\gamma,B,n}^{00}}{\mathcal A_{\gamma,B,2}^{00}}
                         \frac{\theta_{\gamma,n}}{\theta_{\gamma,2}}.
 \label{f16v39:exact-projected-moment}
\end{equation}
The normalization $\theta_{\gamma,2}$ is essential.

\begin{theorem}[Two-sided projected power bounds, not a projected edge law]
\label{f16v39:moment-envelope}
Let $a_B=3(B+2)/2$ and $\Lambda_{{\rm G},B}$ retain exactly their V38 meanings.
At fixed $B$, there exist $0<c_B'\le C_B'<\infty$ and $n_0(B)$ such that,
for even $n\ge n_0(B)$ satisfying the above entry and V37's scalar error bound
with that error at most one,
\begin{equation}
 \boxed{\quad
 c_B'\Lambda_{{\rm G},B}^{n-2}n^{-a_B}
 \le M_{\gamma,n}^{\rm G}
 \le C_B'\Lambda_{{\rm G},B}^{n-2}n^{-9/2}.
 \quad}
 \label{f16v39:projected-envelope}
\end{equation}
The constants are independent of $n,\gamma$ in this admitted set.
\end{theorem}
\begin{proof}
The new bounds at $n$ and two give
$c_B''\le\theta_{\gamma,n}/\theta_{\gamma,2}\le C_B''n^{\ell_B}$.
V37's complete scalar comparison, including all kinetic and endpoint factors,
and V38's fixed-$B$ Gaussian asymptotic give
\[
 c_B'''\Lambda_{{\rm G},B}^{n-2}n^{-a_B}
 \le\frac{\mathcal A_{\gamma,B,n}^{00}}{\mathcal A_{\gamma,B,2}^{00}}
 \le C_B'''\Lambda_{{\rm G},B}^{n-2}n^{-a_B}.
\]
Errors at both durations are paid by their absolute bounds. Insert them into
\eqref{f16v39:exact-projected-moment} and use $a_B-\ell_B=9/2$.
This identity of exponents gives the upper envelope, \emph{not} an asymptotic
formula with exponent $9/2$ or a root-projected Gamma distribution.
\end{proof}

\subsection{A three-duration test excludes projected vacuum preparation}

The next step uses only positive spectral measures, the same native top, and
the envelope just proved. It does not transfer V38's unprojected edge law
through a projection of vanishing norm.

\begin{proposition}[Finite three-duration Perron-weight certificate]
\label{f16v39:three-times}
Suppose even $n_0(B)\le n\le M\le L$ all satisfy
\eqref{f16v39:projected-envelope}. Enlarge a fixed $C_B^0$ if necessary. Then
\begin{equation}
 \boxed{\quad
 |\langle\Omega_\gamma,v^{\rm G}_{\gamma,n}\rangle|^2
 \le C_B^0 n^{a_B}M^{-9/2}
 \exp\left[\frac{(M-n)(a_B\log L+C_B^0)}{L-2}\right].
 \quad}
 \label{f16v39:atom-bound}
\end{equation}
The bound may be replaced by its minimum with one. Every moment and top in
this inequality belongs to the root-invariant restriction of the original
$T_\gamma$; that top is $\Lambda_\gamma$.
\end{proposition}
\begin{proof}
The variational moment bound at $L$ yields
\[
 \Lambda_\gamma\ge (M_{\gamma,L}^{\rm G})^{1/(L-2)}
 \ge\Lambda_{{\rm G},B}
        \exp[-(a_B\log L+C_B^0)/(L-2)].
\]
If $q_\gamma^{\rm G}=|\langle\Omega_\gamma,f_\gamma^{\rm G}\rangle|^2$, then
positivity of its spectral measure gives
$M_{\gamma,M}^{\rm G}\ge q_\gamma^{\rm G}\Lambda_\gamma^{M-2}$.
The actual filtered top weight is
$q_\gamma^{\rm G}\Lambda_\gamma^{n-2}/M_{\gamma,n}^{\rm G}$.
Therefore it is at most
\[
 \frac{M_{\gamma,M}^{\rm G}}
 {M_{\gamma,n}^{\rm G}\Lambda_\gamma^{M-n}}.
\]
Apply the upper envelope at $M$, the lower envelope at $n$, and the just-proved
lower bound on the \emph{native} top. The powers of $\Lambda_{{\rm G},B}$ cancel
exactly, giving \eqref{f16v39:atom-bound}. No upper estimate for $\Lambda_\gamma$
or assumed rate of projected filtering has been used.
\end{proof}

\begin{theorem}[Projected endpoint nonequilibrium in an explicit smaller entrance]
\label{f16v39:physical-nonequilibrium}
Fix $B=m^3$, $m\ge4$, and
\begin{equation}
 0<\alpha<\frac{3}{7(B+2)},\qquad
 n_\gamma=2\lfloor\gamma^\alpha/2\rfloor.
 \label{f16v39:short-regime}
\end{equation}
Then for the normalized root-projected filtered source in
\eqref{f16v39:vectors},
\begin{equation}
 |\langle\Omega_\gamma,v^{\rm G}_{\gamma,n_\gamma}\rangle|^2\longrightarrow0,
 \qquad
 d_{\rm TV}\left(|v^{\rm G}_{\gamma,n_\gamma}|^2dH,
                                      \Omega_\gamma^2dH\right)\longrightarrow1.
 \label{f16v39:physical-limit}
\end{equation}
Both densities in the second expression are root invariant, so the total
variation is attained by invariant events. It is thus a genuine distinction
on the invariant configuration algebra for this projected return source.
\end{theorem}
\begin{proof}
Choose fixed exponents
\[
 \frac{B+2}{3}\alpha<\beta<\chi<\frac17,
 \quad M_\gamma=2\lfloor\gamma^\beta/2\rfloor,
 \quad L_\gamma=2\lfloor\gamma^\chi/2\rfloor.
\]
They exist precisely under the strict bound \eqref{f16v39:short-regime}.
At fixed $B$ all three durations satisfy the inherited V35--V37 entries
at every fixed tail exponent eventually. Moreover
\[
 \frac{M_\gamma\log L_\gamma}{L_\gamma}\to0,\qquad
 n_\gamma^{a_B}M_\gamma^{-9/2}
           =O_B(\gamma^{a_B\alpha-(9/2)\beta})\to0.
\]
Proposition~\ref{f16v39:three-times} proves the vanishing projected overlap.
The projected vector is positive, so its probability affinity with
$\Omega_\gamma^2dH$ is that overlap before squaring. The inequality
$\min(p,q)\le\sqrt{pq}$ gives the total-variation conclusion. Since both
densities are invariant, the set where one exceeds the other is invariant.

For interpretation, the same vector is the exact midpoint amplitude of the
return prepared by projecting the original once-loaded endpoint on \emph{both}
sides. Indeed $\Pi T_\gamma=T_\gamma\Pi$ and the normalizing power is
\eqref{f16v39:projected-moments}. This is not obtained by averaging the density
of the unprojected bridge: Proposition~\ref{f16v39:fibres} describes that
additional distinction explicitly. The definition changes the source by the
required Gauss projection, not the transfer or its spectral top.
\end{proof}

One explicit exponent choice is
$\alpha=1/[7(B+2)]$, $\beta=1/14$, $\chi=3/28$. It gives the proved bound
$C_B\gamma^{-3/28}(1+o(1))$ on the squared projected Perron overlap.
The duration exponent is small and not asserted optimal or uniform in $B$.
At fixed coupling and regulator, sufficiently long filtering still selects the
Perron vector; this theorem changes the coupling and duration together.

\subsection{What is resolved and the next distinct estimate}

V39 determines that the root-invariant component in V38 really vanishes, with
explicit polynomial upper and lower bounds. The projected moment sequence is
therefore not approximated by the unprojected one at a relative error inferred
from ordinary total variation. The orbit affinity $k(\eta)$ is the exact missing
weight on invariant observables. Controlling its dependence on the remaining
variables would give information beyond the present power envelope.

The shorter-duration theorem also shows that root projection alone does not
prepare the vacuum for this source. A new projected edge law, a useful longer-time
vacuum-overlap estimate, or a compact nonlinear-period analysis is still needed.
The number $9/2$ above is an envelope exponent, not a proved projected density
of states, particle multiplicity, or physical spectral gap. Likewise none of
the clocks used in earlier sampling estimates is substituted for Euclidean
transfer time. This maintains the supplied memoranda's separate normalization,
locality, and physical-transfer requirements. No excited spectral ratio,
unrestricted-regulator theorem, time calibration, or interacting continuum
construction follows from the present endpoint result.

\subsection{Verification and preservation}

The new diagnostic is
\begin{center}\small\path{verify_ym_f16_v39_root_orbits.py}.\end{center}
It passes 144 exact assertions in 48 named families, including noncommutative
forest identities, orbit weights, the kinetic derivative bound, projected
moment normalization, and three-duration arithmetic. All 35 supplied V4--V38
checkers were rerun unchanged and passed. These are finite fixtures, not native
spectral enclosures or formal proof verification; actual outputs are packaged.

Removing this insertion and reversing the one title version recovers V38 byte
for byte. The new spectral conclusions inherit V35's midpoint moments and
V37--V38's scalar/Gaussian comparisons. The orbit and kinetic arguments have
separate compact proofs. The analytic lineage has not been independently
recertified. The next checkpoint remains V40.

\begin{firewall}
V39 proves polynomial bounds for the actual root-projected norm, an exact
orbit-affinity reweighting, native projected moment envelopes, and failure of
projected endpoint Perron preparation on a stated smaller polynomial duration
range at fixed spatial width. It does not derive a projected Gamma law, remove
the remaining invariant orbit weights, estimate a physical excited-state gap,
or identify a continuum Hamiltonian or its time scale. The source projection
is explicit; the native transfer, its same physical top, all scalar factors,
and the remaining nonlinear compact variables are retained.
\end{firewall}
% END F16 V39 APPEND

% BEGIN F16 V40 APPEND -- 2026-09-18
% Insert before the final document-ending command of the supplied cumulative V39.
% All predecessor bodies, native measures, endpoint factors, and clocks are preserved.

\clearpage
\section{V40: a direct root-orbit integral and the projected spectral edge}
\label{f16v40:context}

V39 estimates the norm of the root projection between two different powers of
its duration. Dividing an ordinary density approximation by that small norm
would not improve the estimate. We instead include the compact root average in
the original positive endpoint integral before making any approximation. The
based root variables become additional, explicitly normalized integration
variables. One anchored endpoint still confines them. At Gaussian order their
integral can be evaluated exactly by a temporal affine shift in the inherited
gradient subspace. This supplies the missing relative normalization rather than
an assumed gauge-volume factor.

At fixed spatial width we obtain the leading projected norm, a normalized
projected moment asymptotic, and a native root-invariant Gamma edge law of shape
$9/2$ throughout the existing polynomial-duration entrance. This strengthens
V39's power envelope and its smaller-duration non-preparation theorem. The
transfer and its physical top are unchanged. The result is not a density of
states, an estimate of the excited-state gap, or a physical time calibration.

\subsection{V40 checkpoint, sources, and the exact additional integral}

The supplied V39 appendix and analytic audit were read completely. The current
cumulative source, f14 V100, f15 V100, and Grand Tensor V076 were inventoried and
searched; identified forest, scalar, and Gaussian-history passages were checked.
This is a targeted checkpoint, not independent certification of those archives.
The next scheduled checkpoint is V45. The forest quotient and the distinction
between probability averaging and amplitude projection remain owned by V39 and
its cited f15 V10 construction. We do not repeat them as new results.

For primary-source context, Burbano--Bauer,
\href{https://arxiv.org/html/2409.13812}{\emph{Gauge Loop-String-Hadron Formulation
on General Graphs}}, Sections II--III.1, separates based gauge fixing from the
remaining global conjugation. Bahr--Thiemann,
\href{https://arxiv.org/html/0709.4636}{\emph{Gauge-invariant coherent states II}},
Section 2, defines compact group projection; its discussion of exceptional orbit
geometry is a reason not to assume a uniform stationary-phase formula. Those
primary HTML statements were checked. Neither a coherent-state asymptotic nor a
Hamiltonian truncation is imported. The compact domination and Gaussian
integral needed here are proved below. The supplied import memoranda remain
methodological proposals, not premises supplying a physical spectral theorem.

Fix $B=m^3$, $m\ge4$, and one coarse root $o$. Put
\[
 r_B=B-1,\qquad k_B=3r_B,\qquad \ell_B=k_B/2.
\]
For $h=(h_b)_{b\ne o}$, set $h_o=I$ and
\begin{equation}
 Z(h)_e=h_{s(e)}h_{t(e)}^{-1}.
 \label{f16v40:pure-endpoint}
\end{equation}
Use the original orientation, all four parallel bridges, and every periodic
seam. Let $a_{\gamma,Z}(X)$ be V1's normalized dressed endpoint vector. Its exact
equivariance gives
\[
 F_{\gamma,2}(Z(h))=F_{\gamma,2}(I),\qquad
 a_{\gamma,Z(h)}(hXh^{-1})=a_{\gamma,I}(X).
\]
Indeed the carrier's five chord vectors in a cube undergo one common adjoint
rotation, which preserves their Gaussian quadratic and their Haar Jacobians;
the mixed multiplier and its actual integral transform covariantly.

Let $\mathcal A^{\rm G}_{\gamma,B,n}$ denote the original identity-endpoint
return with a root projection on both endpoints. More precisely it is
$j_{\rm br,0}\langle \Pi s_\gamma,T_\gamma^n\Pi s_\gamma\rangle$, where
$s_\gamma=a_{\gamma,I}\delta_I$ is interpreted through V38's one-step loading.
For $n\ge2$ this is a finite positive scalar. Equivalently it is
$\langle b_\gamma^{\rm end},T_\gamma^{n-2}\Pi b_\gamma^{\rm end}\rangle$.
This uses no inner product of an unsmoothed delta. Since $s_\gamma$ is invariant
under simultaneous conjugation, full root averaging equals based averaging.
Commutation and $\Pi^2=\Pi$ reduce the two averages to one:
\begin{equation}
 \mathcal A^{\rm G}_{\gamma,B,n}
 =j_{\rm br,0}\int_{G^{r_B}}
    \langle s_\gamma,T_\gamma^n\mathsf U_hs_\gamma\rangle\,dH^{r_B}(h).
 \label{f16v40:one-average}
\end{equation}
All identities involving endpoint distributions are justified by the bounded
kernel after one loading, then Tonelli. No residual global orientation is fixed.

Write $h_b=\operatorname{Exp}(\xi_b/\sqrt\gamma)$ on its complete principal
chart, and let $w=(\zeta,y_1,\ldots,y_{n-1})$ have exactly its V34 meaning.
Substitute $Z_0=I$, $Z_n=Z(h)$ into the \emph{original} classical action in
unmoving coordinates, leaving the endpoint group products ordered. Denote the
result by $U^+_\gamma(w,\xi)$. Define
\begin{equation}
 \begin{split}
 d_+&=(72n-21)B+3(B-1)=(72n-18)B-3,\\
 \mathcal J^+_\gamma(w,\xi)
   &=J^f_\gamma(\zeta)
     \prod_{i=1}^{n-1}\prod_{e\ {\rm bridge}}j(y_{i,e}/\sqrt\gamma)
     \prod_{b\ne o}j(\xi_b/\sqrt\gamma),\\
 \mathcal Z^+_{\gamma,B,n}
   &=\int_{\mathcal P^+_\gamma}
              \mathcal J^+_\gamma e^{-U^+_\gamma}\,dw\,d\xi .
 \end{split}
 \label{f16v40:root-mass}
\end{equation}
Here $\mathcal P^+_\gamma$ is the product of complete three-group charts.
There is no retained-bridge Haar factor at the final pure-gauge endpoint: that
endpoint is not integrated against bridge Haar. Its variables are the $r_B$
root groups, whose full Haar factors and constant density are retained instead.
Chord endpoint powers remain $1/2$.

\begin{proposition}[Exact projected scalar before the weak-field limit]
\label{f16v40:exact}
At every finite regulator with $n\ge2$,
\begin{equation}
 \boxed{\quad
 \mathcal A^{\rm G}_{\gamma,B,n}
 =c_{H,\gamma}^{r_B}
   \left(\frac{c_{H,\gamma}}{z_\gamma}\right)^{24Bn}
   \frac{\mathcal Z^+_{\gamma,B,n}}{\mathcal S_{\gamma,B}(I)}.
 \quad}
 \label{f16v40:exact-scalar}
\end{equation}
For even $n$, this equals
$\mathcal A^{00}_{\gamma,B,n}\theta_{\gamma,n}$. Consequently
\begin{equation}
 \theta_{\gamma,n}
   =c_{H,\gamma}^{r_B}\frac{\mathcal Z^+_{\gamma,B,n}}{\mathcal Z_{\gamma,B,n}}.
 \label{f16v40:theta-exact}
\end{equation}
The same scalar ratio defines the corresponding spectral projection weight at
other integer durations. This is an identity for the original transfer, not a
replacement endpoint ensemble.
\end{proposition}
\begin{proof}
In the root average, the right endpoint distribution is exactly the dressed
vector at $Z(h)$ by equivariance. V37's cancellation of $p_{\gamma,B}$ and
$c_\Omega^2$ therefore applies, and both static endpoint denominators equal
$\mathcal S_{\gamma,B}(I)$. The kernel in \eqref{f16v40:one-average} uses Haar
endpoint coordinates. Multiplication by $j_{\rm br,0}$ gives the two prescribed
zero-chart half-densities. If the nonzero final endpoint is temporarily written
in its bridge chart, the conversion back divides by its square-root Jacobian.
It cancels precisely the final product
$\prod_e j(y_{n,e}/\sqrt\gamma)^{1/2}$ in V4's chart formula. At cut values one can instead use the Haar kernel directly; the same
identity holds without a final bridge logarithm.

The initial endpoint has unit nonconstant Jacobian. All interior bridge Haar
factors and the fast endpoint half powers remain those in
\eqref{f16v40:root-mass}. The $r_B$ based root integrations contribute
$c_{H,\gamma}^{r_B}$ and the displayed full $j(\xi_b/\sqrt\gamma)$ factors.
This proves \eqref{f16v40:exact-scalar}, including the unchanged $24Bn$ kinetic
normalizers. Comparing with V37's unprojected formula and V38's loaded-vector
identity proves \eqref{f16v40:theta-exact}. No intermediate carrier projection
is introduced. Positivity permits every rearrangement of these integrals.
\end{proof}

\subsection{The Gaussian root integral is an exact temporal affine shift}

Let $\mathsf d_{\rm rt}:\mathbb R^{k_B}\to\mathbb R^{36B}$ be the based
root-gradient injection
\[
 (\mathsf d_{\rm rt}\xi)_e=\xi_{s(e)}-\xi_{t(e)},\qquad \xi_o=0.
\]
All matrices include colour. The inherited static classification gives
$D\mathsf d_{\rm rt}=0$, and connectedness makes $\mathsf d_{\rm rt}$ injective.
Define the positive matrix
\begin{equation}
 \mathsf J_{\rm rt}=\mathsf d_{\rm rt}^{\mathsf T}G^{-1}\mathsf d_{\rm rt}
 \succ0,
 \qquad
 \kappa_B=\frac{(2\pi)^{k_B/2}}
                  {(2\pi^2)^{r_B}\sqrt{\det\mathsf J_{\rm rt}}}.
 \label{f16v40:root-metric}
\end{equation}
This is the original kinetic $G$, with the nonroot ports integrated, not an
identity metric chosen on the root graph. The determinant includes all colours.
Let $U_0^+$ be the quadratic limit of $U^+_\gamma$ and
$\mathcal Z^+_{0,B,n}=\int e^{-U_0^+}\,dw\,d\xi$.

\begin{theorem}[Exact Gaussian orbit mass]
\label{f16v40:Gaussian}
At every fixed finite $B,n$,
\begin{equation}
 \boxed{\quad
 \frac{\mathcal Z^+_{0,B,n}}{\mathcal Z_{0,B,n}}
       =\frac{(2\pi n)^{k_B/2}}{\sqrt{\det\mathsf J_{\rm rt}}}.
 \quad}
 \label{f16v40:Gaussian-ratio}
\end{equation}
After fast Gaussian integration, the variables $\xi$ and
$\widetilde y_i=y_i-(i/n)\mathsf d_{\rm rt}\xi$, $1\le i<n$, are independent;
the latter have the zero-endpoint retained Gaussian and
$\xi\sim N(0,n\mathsf J_{\rm rt}^{-1})$. This is a statement about the exact
quadratic integral. It is not a factorization of nonlinear native orbit fibres.
\end{theorem}
\begin{proof}
At quadratic order the final endpoint is
$y_n=\mathsf d_{\rm rt}\xi$. The conjugation of an endpoint chord by $h_b$ has
no first-order change at the identity and does not change its original endpoint
quadratic. V33's complete action formula applies with these free final data.
The magnetic terms and their integrated fast return depend on $Dy_i$ only;
they are unchanged by the displayed shift because $D\mathsf d_{\rm rt}=0$.
The kinetic cross term vanishes exactly:
\[
 \sum_{t=0}^{n-1}
 (\widetilde y_t-\widetilde y_{t+1})^{\mathsf T}
                          G^{-1}\mathsf d_{\rm rt}\xi=0,
 \qquad \widetilde y_0=\widetilde y_n=0.
\]
Thus the retained action splits as
\begin{equation}
 \mathcal F_0(0,y_1,\ldots,y_{n-1},\mathsf d_{\rm rt}\xi)
 =\mathcal F_0(0,\widetilde y_1,\ldots,\widetilde y_{n-1},0)
                  +\frac1{2n}\xi^{\mathsf T}\mathsf J_{\rm rt}\xi.
 \label{f16v40:affine-split}
\end{equation}
The fast Gaussian determinant is independent of the retained history, including
its endpoint, and is the same determinant in the two integrals. The triangular
change $(y,\xi)\mapsto(\widetilde y,\xi)$ has determinant one. Integrating the
last Gaussian factor proves \eqref{f16v40:Gaussian-ratio} with its exact
normalization and the independence statement. No simultaneous diagonalization
of $G$ with the magnetic matrices is needed.
\end{proof}

In particular, set
\begin{equation}
 M_{0,n}^{\rm G}:=
 \left(\frac n2\right)^{\ell_B}
              \frac{\mathcal A^{00}_{0,B,n}}{\mathcal A^{00}_{0,B,2}}.
 \label{f16v40:Gaussian-moment}
\end{equation}
V38's full Gaussian endpoint determinant, not a gauge-dimension guess, now gives
\begin{equation}
 M_{0,n}^{\rm G}
 =\widetilde c_B\Lambda_{{\rm G},B}^{n-2}n^{-9/2}
  \{1+\widetilde d_{B,1}/n+\widetilde d_{B,2}/n^2+O_B(n^{-3})\},
 \qquad \widetilde c_B>0.
 \label{f16v40:Gaussian-edge}
\end{equation}
The residual simultaneous conjugation has not been divided out as an additional
three-dimensional translation. The remaining nine free Gaussian directions
transform under it, while the endpoint source is already invariant.

\subsection{Compact domination for the root-averaged integral}

The new integral must be estimated in its own right. A small absolute error
for an unprojected law cannot be divided by $\gamma^{-\ell_B}$.

\begin{proposition}[One anchored endpoint confines the based root variables]
\label{f16v40:pin}
For each fixed $B$ there are $c_B^+>0$, $H_B^+<\infty$, independent of $n,\gamma$,
such that on the complete root-extended principal product, with $u=(w,\xi)$,
\begin{equation}
 U^+_\gamma(u),\ U^+_0(u)\ge\frac{c_B^+}{n^2}|u|^2,
 \qquad 2c_B^+n^{-2}I\preceq\mathcal H^+\preceq H_B^+I,
 \label{f16v40:pin-bound}
\end{equation}
where $\mathcal H^+=D^2U_0^+$. The identity is the unique zero in these based
coordinates. For $\mathcal C^+=(\mathcal H^+)^{-1}$,
\begin{equation}
 \|\mathcal C^+\|\le C_Bn^2,
 \qquad \sup_g\|\mathcal C^+_{gg}\|\le C_Bn,
 \qquad \mathcal C^+_{\zeta\zeta}\preceq(2a_*)^{-1}I.
 \label{f16v40:Gaussian-scales}
\end{equation}
The subscript $g$ in the middle bound denotes an original three-group, including
one of the new based root coordinates, not a random gauge transformation.
\end{proposition}
\begin{proof}
V5's exact independent anchor gives $U^+_\gamma\ge a_*|\zeta|^2$ at every
retained input. V34's metric-increment estimate used neither final anchoring
nor a retained window, so it still gives
\[
 \gamma\sum_{t,e}\operatorname{dist}(Z_{t,e},Z_{t+1,e})^2
                                      \le C_\nabla U^+_\gamma.
\]
Starting only from $Z_0=I$, telescoping and Cauchy--Schwarz bound the sum of all
interior squared distances by $n^2$ times the sum of squared increments, and
bound the final squared distances by $n$ times that sum. Along a coarse spanning
tree, the ordered product of the final values \eqref{f16v40:pure-endpoint} from
$o$ to $b$ is $h_b^{-1}$. The triangle inequality consequently gives
\[
 \sum_{b\ne o}\operatorname{dist}(I,h_b)^2
 \le(B-1)^2\sum_{e\in\mathcal F}\operatorname{dist}(I,Z_{n,e})^2.
\]
Principal coordinates satisfy $|\xi_b|^2=\gamma\operatorname{dist}(I,h_b)^2$.
Combining these inequalities proves the native lower bound, for example with a
constant no larger than
$[a_*^{-1}+C_\nabla(1+(B-1)^2)]^{-1}$. This constant is independent of duration
but not spatial width. Zero action forces all fast anchors and increments to
vanish; then $Z_i=I$ and every $h_b=I$ by the based tree. No additional saddle
or global-conjugation volume has been deleted.

The fixed-vector quadratic limit gives the same lower bound for $U_0^+$.
Original finite word carriers and the extra endpoint products give the upper
Hessian bound. Minimizing over $(y,\xi)$ retains the fast anchor and proves the
last bound in \eqref{f16v40:Gaussian-scales}. For the retained blocks use
\eqref{f16v40:affine-split}: their covariance is the original zero-endpoint
covariance, with local size $O(n)$, plus
$(i/n)^2n\mathsf d_{\rm rt}\mathsf J_{\rm rt}^{-1}\mathsf d_{\rm rt}^{\mathsf T}$.
The root covariance is $n\mathsf J_{\rm rt}^{-1}$. At fixed $B$ these give the
middle bound. This is a marginal calculation, not a conditional covariance
substituted for the full one.
\end{proof}

\begin{theorem}[Normalized root-extended comparison and absolute mass]
\label{f16v40:mass}
Fix $B$ and a tail power $s\ge2$. There are constants depending on $B,s$ but
not on $n,\gamma$ such that the following statements hold under the entry
conditions specified in the proof and $\rho_+=d_+n^4/\gamma\le1$:
\begin{equation}
 \left|\log\frac{\mathcal Z^+_{\gamma,B,n}}{\mathcal Z^+_{0,B,n}}\right|
       \le C_B(\rho_++\gamma^{-s}),\qquad
 d_{\rm TV}(\mathbb P^+_\gamma,\mathbb Q^+)
       \le C_B(\sqrt{\rho_+}+\gamma^{-s}).
 \label{f16v40:mass-bounds}
\end{equation}
Here the laws normalize their complete respective integrals; all root groups
remain integrated. A sufficient asymptotic entry, at fixed $B$, is
\begin{equation}
 \frac{Bn^7\log(n+1)+n^6\log\gamma}{\gamma}\longrightarrow0.
 \label{f16v40:regime}
\end{equation}
This is not a theorem uniform over growing spatial widths.
\end{theorem}
\begin{proof}
We verify the new integral's hypotheses rather than applying V36 to a different
law by name. Shrink $c_B^+\le1/4$, enlarge $H_B^+\ge1$, and put
$a=c_B^+/n^2$, $t=s+4$. In the formula for V35's radius replace $(d,c,H_*)$ by
$(d_+,c_B^+,H_B^+)$:
\[
 R^2=\frac2a\left\{1+\frac{d_+}{2}\log\frac{2H_B^+}{a}
  +\log\left[1+\left(\frac{d_++6}{a}\right)^4\right]+(t+1)\log\gamma\right\}.
\]
The precise entry is $R\le\sqrt\gamma/2$ and
$C_{2,B}R/\sqrt\gamma+C_{H,B}/\gamma\le a$. The constants are the finite
ordered-word and Haar derivative bounds for this explicit endpoint
substitution. All new factors are $\operatorname{Exp}(\pm\xi_b/\sqrt\gamma)$;
there is no logarithm of their product to differentiate. Bounded incidence
gives on this ball, for $W^+=U^+_\gamma-\log\mathcal J^+_\gamma-U_0^+$,
\[
 \|D^2W^+\|\le C_{2,B}R/\sqrt\gamma+C_{H,B}/\gamma,\qquad
 |\nabla W^+|^2\le\frac{C_B}{\gamma}\sum_g|u_g|^4
                       +\frac{C_B}{\gamma^2}|u|^2.
\]
Thus $D^2(U_0^++\vartheta W^+)\succeq\mathcal H^+/2$ for $0\le\vartheta\le1$.
All these laws and all linear tilts are on one unchanged convex ball. Their
untilted means vanish by simultaneous colour conjugation, which preserves the
based condition $h_o=I$. The convex-boundary variance identity proved in V35
yields directional moments in metric $\mathcal C^+$ and therefore local fourth
moments $O_B(n^2)$, using \eqref{f16v40:Gaussian-scales}. It follows that
\[
 \mathbb E_{\vartheta}
 [\nabla W^+{}^{\mathsf T}\mathcal C^+\nabla W^+]
             \le C_B\rho_+,\qquad
 \operatorname{Var}_{\vartheta}W^+\le C_B\rho_+.
\]

The compact lower bound \eqref{f16v40:pin-bound}, $0\le\mathcal J^+_\gamma\le1$,
and the Gaussian of precision $2\mathcal H^+$ give
$\mathcal Z^+_\gamma\ge2^{-d_+/2-1}\mathcal Z^+_0$ by integrating over the
ball. The same radial eighth-moment calculation as V35, with the displayed
parameters, bounds each complete weighted tail by $\gamma^{-t}$. This includes
the entire root-chart complement. No exterior good event is assumed.

For the scalar, the classical cubic $P_3^+$ is odd under simultaneous Euclidean
reflection. Ordered fourth derivatives and the Haar quadratic give
\[
 |W^+-\gamma^{-1/2}P_3^+|
       \le\frac{C_B}{\gamma}\sum_g(|u_g|^4+|u_g|^2).
\]
The centered Gaussian on the ball has zero cubic mean. Hence
$|\mathbb E_0W^+|\le C_Bd_+n^2/\gamma\le C_B\rho_+$. Differentiating the
\emph{normalized} interpolating integrals gives the exact identity
\[
 \log(Z_{1,\mathbb B}^+/Z_{0,\mathbb B}^+)
 =-\mathbb E_0W^+
   +\int_0^1(1-\vartheta)\operatorname{Var}_{\vartheta}W^+\,d\vartheta.
\]
The two complete tail probabilities restore the full scalar by adding
$\log\mathbb Q^+(\mathbb B)-\log\mathbb P^+_\gamma(\mathbb B)$. This proves
the first bound, including its actual normalizers. The normalized
square-root-density argument of V35 bounds total variation on the ball by
$C_B\sqrt{\rho_+}$; adding the two paid tails proves the other bound.

Finally $R\le C_{B,s}n\sqrt{d_+\log(n+1)+\log\gamma}$. A sufficiently small
constant in
$n^6[d_+\log(n+1)+(s+5)\log\gamma]\le c_{B,s}\gamma$
ensures the stated finite entry. At fixed $B$, \eqref{f16v40:regime} implies
this and $\rho_+\to0$. This verifies the parameter range without importing
an unproved estimate for a root-twisted boundary row.
\end{proof}

\subsection{Relative projected normalization, without dividing an absolute error}

\begin{theorem}[The leading native root weight and projected moments]
\label{f16v40:projection}
At fixed $B$, at the preceding entry and the original joint entry, for even
$n\ge2$,
\begin{equation}
 \boxed{\quad
 \left|\log\frac{\theta_{\gamma,n}}
                  {\kappa_B(n/\gamma)^{\ell_B}}\right|
       \le C_B(\rho_++\gamma^{-s}).
 \quad}
 \label{f16v40:theta-asymptotic}
\end{equation}
For V39's \emph{same} normalized projected loading $f_\gamma^{\rm G}$, its
moments at all integer $n\ge2$ satisfy
\begin{equation}
 \boxed{\quad
 \left|\log\frac{M_{\gamma,n}^{\rm G}}{M_{0,n}^{\rm G}}\right|
        \le C_B(\rho_++\gamma^{-s}).
 \quad}
 \label{f16v40:moment-asymptotic}
\end{equation}
In particular the leading duration dependence of $\theta$ is now identified,
and $M_{\gamma,n}^{\rm G}$ has the $n^{-9/2}$ leading factor in
\eqref{f16v40:Gaussian-edge} when \eqref{f16v40:regime} holds and $n\to\infty$.
Neither assertion is obtained by dividing a total-variation error by $\theta$.
\end{theorem}
\begin{proof}
Insert \eqref{f16v40:Gaussian-ratio} into \eqref{f16v40:theta-exact} and use
$c_{H,\gamma}=(2\pi^2\gamma^{3/2})^{-1}$. The remaining logarithm is exactly
\[
 \log(\mathcal Z^+_\gamma/\mathcal Z^+_0)
                       -\log(\mathcal Z_\gamma/\mathcal Z_0).
\]
The new theorem pays the first term and V36 pays the second; its original
$\rho\le\rho_+$. This proves \eqref{f16v40:theta-asymptotic}.

More economically, the projected moment can be compared without either static
endpoint approximation or the unprojected joint-mass expansion. The exact
formula \eqref{f16v40:exact-scalar} at $n$ and two gives
\[
 M_{\gamma,n}^{\rm G}
  =\left(\frac{c_{H,\gamma}}{z_\gamma}\right)^{24B(n-2)}
                \frac{\mathcal Z^+_{\gamma,B,n}}{\mathcal Z^+_{\gamma,B,2}}.
\]
Both $c_{H,\gamma}^{r_B}$ and $\mathcal S_{\gamma,B}(I)$ cancel exactly. Its
Gaussian counterpart is \eqref{f16v40:Gaussian-moment}. Their logarithmic ratio
is the difference of the two new mass errors minus
$24B(n-2)\log\mathfrak b_\gamma$, with precisely V37's scalar
$\mathfrak b_\gamma$. Since $|\log\mathfrak b_\gamma|\le C/\gamma$ eventually,
this last cost is $O(Bn/\gamma)$ and is bounded by $C_B\rho_+$. All three
errors are paid. This proves \eqref{f16v40:moment-asymptotic}, including the
loading normalization, with no assumed cancellation of approximation errors.
\end{proof}

The result identifies an orbit-integrated weight, not the pointwise function
$k(\eta)$ of V39. The latter need not become constant. It also leaves the
original midpoint probability and the square of its projected amplitude distinct.
For an invariant bounded midpoint multiplication $F$ and even $n$, however,
commutation gives the exact useful identity
\[
 \frac{\langle\Pi s_\gamma,T_\gamma^{n/2}F T_\gamma^{n/2}\Pi s_\gamma\rangle}
      {\langle\Pi s_\gamma,T_\gamma^n\Pi s_\gamma\rangle}
 =\frac{\int dh\,\langle s_\gamma,T_\gamma^{n/2}F
                       T_\gamma^{n/2}\mathsf U_hs_\gamma\rangle}
        {\int dh\,\langle s_\gamma,T_\gamma^n\mathsf U_hs_\gamma\rangle}.
\]
Thus the added positive integral retains the correct invariant-observable
weight; it is not probability averaging substituted for amplitude projection.

\subsection{The native root-invariant Gamma edge and the longer non-preparation range}

Let $\mu_\gamma^{\rm G}$ be the probability spectral measure of
$f_\gamma^{\rm G}$ for the unchanged root-invariant restriction of $T_\gamma$.
Define
\[
 d\mu_{\gamma,n}^{\rm G}(\lambda)
  =\frac{\lambda^{n-2}\,d\mu_\gamma^{\rm G}(\lambda)}{M_{\gamma,n}^{\rm G}}.
\]
Its top atom is at the same native $\Lambda_\gamma$ as in V38. The reference
$\Lambda_{{\rm G},B}$ is not substituted for it.

\begin{theorem}[Projected endpoint edge law at the original polynomial durations]
\label{f16v40:edge}
Fix $B$. Let even $n=n_\gamma\to\infty$ satisfy \eqref{f16v40:regime}, using
fixed tail powers large enough for the adjacent-ratio tests. Under the
\emph{native root-invariant} probability $\mu_{\gamma,n_\gamma}^{\rm G}$,
\begin{equation}
 -n_\gamma\log(\lambda/\Lambda_{{\rm G},B})
     \ \Rightarrow\ \operatorname{Gamma}(9/2,1),\qquad
 n_\gamma(1-\lambda/\Lambda_{{\rm G},B})
     \ \Rightarrow\ \operatorname{Gamma}(9/2,1).
 \label{f16v40:Gamma}
\end{equation}
The second variable has mean and variance both tending to $9/2$. The Gamma
parameter convention is shape, then rate. For every $0<\alpha<1/7$ the
admissible example $n_\gamma=2\lfloor\gamma^\alpha/2\rfloor$ gives
\begin{equation}
 \begin{split}
 |\langle\Omega_\gamma,v^{\rm G}_{\gamma,n_\gamma}\rangle|^2&\longrightarrow0,\\
 d_{\rm TV}(|v^{\rm G}_{\gamma,n_\gamma}|^2dH,\Omega_\gamma^2dH)
                                                   &\longrightarrow1.
 \end{split}
 \label{f16v40:nonpreparation}
\end{equation}
Both densities in the second line are invariant. This extends V39's smaller
range $\alpha<3/[7(B+2)]$ at each fixed $B$, not uniformly in growing $B$.
\end{theorem}
\begin{proof}
The spectral identity and \eqref{f16v40:moment-asymptotic} imply, for every
fixed integer $j\ge0$,
\[
 \mathbb E_{\mu_{\gamma,n}^{\rm G}}
       (\lambda/\Lambda_{{\rm G},B})^{jn}
 =\Lambda_{{\rm G},B}^{-jn}
         \frac{M_{\gamma,(j+1)n}^{\rm G}}{M_{\gamma,n}^{\rm G}}
                  \longrightarrow(j+1)^{-9/2}.
\]
Every fixed multiple of the duration satisfies the entry eventually. The
native scalar errors at both times tend to zero; they are not differentiated
or assumed to cancel. As in V38's proved moment argument, the first moment
gives tightness of $Y=(\lambda/\Lambda_{{\rm G},B})^n$, and the next moment
provides uniform integrability of each preceding one. The limiting moments
force support in $[0,1]$ and determine the law there by polynomial density.
They are the moments of $e^{-G_*}$ for $G_*\sim\operatorname{Gamma}(9/2,1)$.
There is no atom at zero, so taking $-\log Y$ and then the locally uniform map
$n(1-e^{-x/n})$ proves \eqref{f16v40:Gamma}. Negative edge deficits at finite
coupling are permitted; this argument estimates spectral weights, not absence
of eigenvalues above the Gaussian threshold.

For the claimed moments, the determinant expansion
\eqref{f16v40:Gaussian-edge} gives adjacent-ratio errors $O_B(n^{-2})$ and
second-difference errors $O_B(n^{-3})$, with leading constants $9/2$. Adding
the new native logarithmic errors costs $O_B(d_+n^4/\gamma+\gamma^{-s})$.
The regime gives $n^2(d_+n^4/\gamma+\gamma^{-s})\to0$ when, for example,
$s\ge2$. The exact mean and variance formulas used in V38 therefore prove
mean and variance convergence for the linear deficit. This step is separate
from weak convergence.

A tight sequence converging to a nonatomic law has maximal atomic weight
tending to zero, by the compact-subsequence and closed-interval argument in
V38. The Perron eigenvalue of the actual restricted transfer is simple, so its
atom is exactly the squared overlap in \eqref{f16v40:nonpreparation}. At even
durations the filtered projected vector is positive. Its probability affinity
with $\Omega_\gamma^2$ is that overlap before squaring, and
$d_{\rm TV}(p,q)\ge1-\int\sqrt{pq}$ proves the second line. The maximizing
sign event is invariant because both densities are invariant. No spectral
convergence of operators or a native upper-top estimate is inferred.
\end{proof}

For completeness the corresponding projected scalar energy is
\[
 -\log\frac{M_{\gamma,n+1}^{\rm G}}{M_{\gamma,n}^{\rm G}}
   =-\log\Lambda_{{\rm G},B}+\frac{9}{2n}+o_B(n^{-1})
\]
along the same regime. This is measured against the Gaussian threshold, not an
identified physical ground energy. The function $k(\eta)$ remains relevant to
other quotient observables, but a direct positive root integration has paid
its total contribution to these moments without assuming its constancy.

\subsection{Checkpoint: the compact soft sector is now the outstanding scale}

V40 obtains a projected Gamma law by integrating the based group before the
limit. It does not subtract gauge dimensions from V38's exponent without a
calculation: \eqref{f16v40:Gaussian-ratio} supplies the exact determinant and
the native mass theorem supplies its relative error, including the root Haar
volume. The root variables are not assigned an independent Gaussian law in the
native model. The approximate Gaussian separation is used only after complete
compact domination has been proved for the new integral.

The remaining nine-dimensional entrance sector is still nonlinear and compact
at longer scales. The theorem leaves arbitrary durations, growing spatial
width, quantitative vacuum-preparation times, the first excited spectral ratio,
and independently calibrated physical time unpaid. A useful next target is
an exact treatment of that compact period sector together with its integrated
complement, or a same-top spectral comparison valid beyond this entrance.
Another unprojected or projected finite-order scalar coefficient alone would
not supply it. The supplied memoranda's distinction between power gain and
locality/normalization cost, and between auxiliary estimates and physical
transfer, remains binding. No Yang--Mills continuum mass-gap claim is made.

\subsection{Verification and preservation}

The diagnostic
\begin{center}\small\path{verify_ym_f16_v40_root_integral.py}\end{center}
checks finite noncommutative endpoint identities, all root/Haar powers,
noncommuting Gaussian affine shifts and determinants, normalization before
projection, complete-tail and parameter arithmetic, and the spectral moments.
It distinguishes literal periodic incidence checks from finite reference
fixtures. These checks are not native integral enclosures or formal proof
verification. The audit records actual executions and review depth, separately
from inherited reports.

Removing this insertion and reversing the title version recovers V39 byte for
byte. The analytic inputs are the compact anchor, ordered-word bounds, finite
comparison mechanism, and V38's Gaussian asymptotic, with the new hypotheses
checked above. The full lineage is not independently recertified. The next
checkpoint is V45.

\begin{firewall}
V40 proves the exact positive root-averaged endpoint representation, an explicit
Gaussian root determinant, and a relative compact-to-Gaussian mass estimate at
fixed spatial width. It identifies the leading native projected norm and the
root-invariant endpoint Gamma edge of shape $9/2$, with vanishing projected
Perron weight throughout the stated polynomial entrance. It does not prove a
full physical density of states, a gap or gaplessness of the calibrated
Hamiltonian, stationary vacuum convergence, unrestricted regulators, or an
interacting continuum construction. All original transfer powers, endpoint
carriers, root Haar factors, and physical normalizations are retained.
\end{firewall}
% END F16 V40 APPEND

% BEGIN F16 V41 APPEND -- 2026-09-18
% Insert before the final document-ending command of the supplied cumulative V40.
% Predecessor assertions, endpoint measures, normalizations, and clocks are preserved.

\clearpage
\section{V41: the native spectral top and central-flat vacuum localization}
\label{f16v41:context}

V40 controls a projected endpoint spectral measure, but does not supply the
matching upper bound on the native spectral top. Extending that moment
calculation to longer durations would require new estimates. We instead return
to the original one-step Wilson operator on the full fine-link space. Its
positive kernel and its sum-of-squares magnetic action permit a local operator
comparison near every flat connection, including singular flat strata. The
comparison never inverts a small magnetic eigenvalue uniformly over the flat
set. A positive winding expansion then identifies which flat backgrounds have
the largest harmonic transfer norm.

The result is a fixed-spatial-regulator theorem: the actual physical Perron
value converges to V38's Gaussian threshold, and its stationary probability
concentrates on the eight central-flat gauge orbits. This is not an estimate of
the first excited ratio. No rate sufficient to center V40's spectral variable
at the native top is asserted. The localization lengths and constants below
are allowed to depend on the spatial regulator.

\subsection{Sources, exact operator, and the change of viewpoint}

The supplied V40 appendix, analytic audit, and checkpoint were read completely.
The f15 V91 fine-link transfer formula and the f14 discussion of flat backgrounds
were checked at identified passages. The f14 flat-holonomy cancellation concerns
a different quotient determinant and must not be identified with the present
transfer threshold. Searches of the current lineage and supplied archives are
recorded in the audit; these are targeted checks, not an exhaustive novelty or
proof certificate. The next scheduled companion checkpoint remains V45.

Harmonic localization is a classical method. Klein--Rosenberger,
\href{https://arxiv.org/abs/1706.06357}{\emph{Harmonic Approximation of Difference
Operators}}, supplies related context, but its lattice Schr\"odinger hypotheses
are not substituted for a positive compact integral operator with a singular
zero set. We prove the required norm statement directly. For the Bessel
identities in the winding calculation we use the primary formulas
\href{https://dlmf.nist.gov/10.32.E1}{DLMF 10.32.1} and
\href{https://dlmf.nist.gov/10.35.E2}{DLMF 10.35.2}; their relevant integral and
series manipulations are also justified below. No external mass-gap theorem is
an input.

Put $L=2m$, $m\ge4$, $B=m^3$, $V=L^3=8B$ and $E=3V=24B$.
On $\mathcal X_L=\SU(2)^E$, with probability product Haar $dH$, retain
\begin{equation}
 \begin{split}
 S(U)&=\sum_{p\ {\rm spatial}}\left(1-\Re U_p\right),\qquad
 b_\gamma(U)=e^{-\gamma S(U)/2},\\
 \overline T_\gamma&=M_{b_\gamma}C_\gamma^{\otimes E}M_{b_\gamma},\qquad
 C_\gamma(U,U')=z_\gamma^{-1}e^{-\gamma(1-\Re UU'^{-1})}.
 \end{split}
 \label{f16v41:native}
\end{equation}
Here $\Re$ is the scalar quaternion part, so it is the original normalized
fundamental trace. Every original spatial plaquette and kinetic scalar is
included. At finite $L,\gamma$, the operator is compact, positive self-adjoint,
and positivity improving. Its unit Perron vector $\Omega_\gamma>0$ is gauge
invariant by uniqueness. Consequently its norm is the physical top. The exact
rooted-tree reduction in f15 V98/V1 identifies that physical operator with the
root-invariant restriction used in V38--V40. Hence
\[
 \Lambda_\gamma=\|\overline T_\gamma\|
   =\|T_\gamma|_{\operatorname{Ran}\Pi_{\rm G}}\|.
\]
We use the full-link space to estimate this same top, not to replace it by a
row-normalized Markov kernel or a carrier compression. The stationary vacuum
probability in these coordinates is $\Omega_\gamma^2dH$.

\subsection{A narrow-kernel localization identity with its error paid}

Let $\varepsilon=\gamma^{-1/2}$, and use the round product metric for which a
unit imaginary quaternion has norm one. Write $D_L=3E$ for the dimension of
$\mathcal X_L$. Near a point $p$, choose local coordinates whose derivative at
$p$ is an orthonormal identification with $\mathbb R^{D_L}$. In the flat
half-density representation, the normalized kinetic kernel has the following
local properties. Given $\delta>0$, after shrinking the coordinate neighborhood,
for all $x,y$ in it and sufficiently small $\varepsilon$,
\begin{equation}
 c_\gamma^{\rm flat}(x,y)
 \le(1+\delta)(2\pi\varepsilon^2)^{-D_L/2}
       \exp\left[-\frac{(1-\delta)|x-y|^2}{2\varepsilon^2}\right].
 \label{f16v41:local-kernel}
\end{equation}
The corresponding rescaled kernel at the diagonal converges to the normalized
Euclidean Gaussian. These assertions include the local measure density and the
actual $z_\gamma^E$: V37's one-group formula gives
$z_\gamma=(2\pi^2)^{-1}(2\pi\varepsilon^2)^{3/2}(1+O(\varepsilon^2))$.
Also $\sum_e(1-\Re U_eU_e'^{-1})=\frac12\sum_e|U_e-U_e'|^2$ in quaternion
coordinates. Its pullback, and the smooth product volume, give
\eqref{f16v41:local-kernel} by a local $C^1$ near-isometry. The chart can be
chosen with a convex coordinate image. There is no undetermined leading
normalization in this estimate.

\begin{proposition}[Finite partition localization for the actual transfer]
\label{f16v41:IMS}
Let $\chi_0,\ldots,\chi_N$ be real smooth functions on $\mathcal X_L$ with
$\sum_j\chi_j^2=1$. Then
\begin{equation}
 \left\|\overline T_\gamma-
       \sum_jM_{\chi_j}\overline T_\gamma M_{\chi_j}\right\|
       \le C_{L,\chi}\gamma^{-1}.
 \label{f16v41:IMS-bound}
\end{equation}
The partition and $L$ are fixed before $\gamma\to\infty$.
\end{proposition}
\begin{proof}
The difference kernel is exactly
\[
 \frac12\sum_j(\chi_j(U)-\chi_j(U'))^2
 b_\gamma(U)c_\gamma(U,U')b_\gamma(U').
\]
It has nonnegative entries; operator positivity of this difference is not
asserted. The squared differences are bounded by $C_\chi\operatorname{dist}(U,U')^2$.
Since $b_\gamma\le1$, the Schur bound is at most $C_\chi$ times the second
moment of the product Wilson displacement. For one group,
$1-\cos\theta\ge2\theta^2/\pi^2$ on $[0,\pi]$, its normalization is at least
$c\gamma^{-3/2}$, and its weighted radial numerator is at most
$C\gamma^{-5/2}$. Thus that moment is $O(\gamma^{-1})$. Summing the $E$ group
moments proves \eqref{f16v41:IMS-bound}, with its fixed-regulator cost.
\end{proof}

\subsection{Normal coordinates without a regular-flat-stratum assumption}

Let $\mathcal F_L=\{U:S(U)=0\}$. It is a compact set, not presumed smooth.
In the ambient quaternion space define the smooth constraint map
\[
 \mathcal W(U)=(U_p-I)_{p\ {\rm spatial}},\qquad
 S(U)=\tfrac12|\mathcal W(U)|^2.
\]
For $p\in\mathcal F_L$, let $H_p=(D\mathcal W_p)^*D\mathcal W_p$ on its
orthonormal tangent space. Define the continuous harmonic threshold
\begin{equation}
 \mathfrak h(p)=
 \prod_{\nu\in\operatorname{spec}(H_p),\ \nu>0}
 \left(1+\frac\nu2+\sqrt{\nu+\nu^2/4}\right)^{-1/2}.
 \label{f16v41:harmonic-threshold}
\end{equation}
A zero eigenvalue contributes one. Thus the formula is continuous also where
$\operatorname{rank}H_p$ changes. Its positive factors are V38's normalized
Mehler factors, not an exponentiated Schr\"odinger frequency substituted for the
discrete transfer.

\begin{proposition}[Exact magnetic lower bound in a pointwise normal chart]
\label{f16v41:normal-chart}
For every $p\in\mathcal F_L$ there is a smooth local coordinate system
$x=(x_+,x_0)$ centered at $p$, orthonormal to first order, such that
\begin{equation}
 S(U(x))\ge\tfrac12\sum_{i=1}^{r}\nu_i x_{+,i}^2,
 \qquad r=\operatorname{rank}H_p,
 \label{f16v41:normal-lower}
\end{equation}
where $\nu_1,\ldots,\nu_r$ are the positive eigenvalues of $H_p$.
This is an exact inequality on the neighborhood, with no negative remainder.
\end{proposition}
\begin{proof}
Write the nonzero singular decomposition of $D\mathcal W_p$ as
$D\mathcal W_p v_i=\sqrt{\nu_i}e_i$, with both systems orthonormal. Set
$x_{+,i}=\langle e_i,\mathcal W(U)\rangle/\sqrt{\nu_i}$, and take as $x_0$
the components of an initial logarithmic chart along $\ker D\mathcal W_p$.
The combined derivative at $p$ is an orthogonal isomorphism. The finite
inverse function theorem gives a local coordinate system. Orthogonal
projection of $\mathcal W(U)$ onto the span of the $e_i$ proves
\eqref{f16v41:normal-lower}. If $r=0$, use the original orthonormal chart and
$S\ge0$. No inverse of a positive eigenvalue uniform over $p$ is required;
only a chart at the individual point is constructed.
\end{proof}

\begin{theorem}[Fixed-regulator harmonic principle for the native top]
\label{f16v41:harmonic-principle}
For the original operator \eqref{f16v41:native},
\begin{equation}
 \boxed{\qquad
 \lim_{\gamma\to\infty}\Lambda_\gamma
       =\max_{p\in\mathcal F_L}\mathfrak h(p).
 \qquad}
 \label{f16v41:top-principle}
\end{equation}
Let $\mathcal F_L^*$ be the maximizing set. For every open neighborhood
$O\supset\mathcal F_L^*$,
\begin{equation}
 \int_{\mathcal X_L\setminus O}\Omega_\gamma(U)^2\,dH(U)\longrightarrow0.
 \label{f16v41:abstract-localization}
\end{equation}
No positive Hessian floor in the null directions, Morse--Bott hypothesis, or
uniform rank on $\mathcal F_L$ is assumed.
\end{theorem}
\begin{proof}
For the upper bound fix a small $\delta>0$. At each flat $p$, use the preceding
chart, shrunk so \eqref{f16v41:local-kernel} holds. Kernelwise domination and
\eqref{f16v41:normal-lower} bound the localized operator norm by
\begin{equation}
 (1+\delta)(1-\delta)^{-D_L/2}
 \prod_{i=1}^{r}
 \left(1+\frac{\nu_i}{2(1-\delta)}+
 \sqrt{\frac{\nu_i}{1-\delta}+\frac{\nu_i^2}{4(1-\delta)^2}}
 \right)^{-1/2}.
 \label{f16v41:local-norm}
\end{equation}
Indeed this is the norm of the dominating whole-space Gaussian after scaling
$x$ by $\sqrt{1-\delta}/\varepsilon$. The zero directions give normalized heat
convolution of norm one; the raw kinetic prefactor gives
$(1-\delta)^{-D_L/2}$. For complex tests use their absolute values before this
positive-kernel comparison. No Loewner ordering is inferred from pointwise
ordering alone.

A finite subcover of the compact flat set suffices. Add a smooth partition
piece supported away from it; its restricted norm is at most $e^{-\gamma s_0}$
for some $s_0>0$, because both magnetic factors are bounded by
$e^{-\gamma s_0/2}$ and $C_\gamma^{\otimes E}$ has norm one. Take a smooth
quadratic partition subordinate to this cover. Proposition~\ref{f16v41:IMS}
then bounds the global Rayleigh quotient by the largest local norm plus
$O_{L,\chi}(\gamma^{-1})$. At fixed $L$ the positive eigenvalues are bounded,
and the expression \eqref{f16v41:local-norm} converges uniformly in those
values, including zero, to \eqref{f16v41:harmonic-threshold} as $\delta\downarrow0$.
First send $\gamma\to\infty$ for the chosen finite cover, then
$\delta\downarrow0$. This proves the upper bound without localizing every
configuration into a coupling-dependent tube around a singular stratum.

For the lower bound fix a flat $p$ and split an orthonormal logarithmic tangent
chart into the range and kernel of $H_p$, of dimensions $r$ and $k$.
Choose real smooth compactly supported unit functions $f$ on $\mathbb R^r$
and $g$ on $\mathbb R^k$. When $k=0$ omit the second factor. Put
$t_\varepsilon=\varepsilon^{3/4}$ and use the half-density trial function
\[
 \phi_\varepsilon(x_+,x_0)=
 \varepsilon^{-r/2}t_\varepsilon^{-k/2}
 f(x_+/\varepsilon)g(x_0/t_\varepsilon).
\]
It has norm one for small $\varepsilon$ and support inside the chart. Taylor's
formula for the constraint map, not a uniform quadratic lower bound, gives
\[
 \varepsilon^{-1}\mathcal W(U(\varepsilon u,t_\varepsilon z))
    =D\mathcal W_p(u,0)+O(t_\varepsilon^2/\varepsilon+\varepsilon)
    =D\mathcal W_p(u,0)+o(1)
\]
uniformly on its rescaled compact support. Thus both magnetic factors tend
to the required normal Gaussian multipliers. In the kinetic integral put
$x'=x+\varepsilon v$. The soft argument changes by
$(\varepsilon/t_\varepsilon)v_0\to0$, while the hard argument changes by
$v_+$. The actual normalized kernel and volume converge to the Euclidean
Gaussian displacement. On every bounded set of $v$ the convergence is
uniform; \eqref{f16v41:local-kernel} and the Schur estimate for the tail
$|v|>R$ bound the omitted bilinear form by $C e^{-cR^2}$, uniformly for small
$\varepsilon$. Integrate first at bounded $v$, then send $R\to\infty$.
The Rayleigh quotient tends to the normal Gaussian transfer quotient of $f$,
times $\|g\|_2^2=1$. Smooth compactly supported $f$ approximate its massive
Gaussian ground function. Its supremum is $\mathfrak h(p)$. This proves the
lower bound. The soft trial scale is a variational device, not an asserted
native fluctuation scale. The trial need not be gauge invariant: equality of
the full and physical tops at each coupling is used only for the largest
eigenvalue, not for excited spectral values.

For localization, choose $O_0$ with
$\mathcal F_L^*\subset O_0\subset\overline O_0\subset O$. Compactness and
continuity give a strict threshold loss $\eta>0$ on
$\mathcal F_L\setminus O_0$. In the preceding upper construction choose charts
centered in $\mathcal F_L\cap O_0$ to lie in $O$, and use charts with loss at
least $\eta$ for the remaining flat set. The nonflat partition piece is also
in the lower-norm family. The sum of squared lower-norm pieces is at least one
on $\mathcal X_L\setminus O$. With chart errors made less than an arbitrary
$\delta'>0$ and less than $\eta/2$, the Perron Rayleigh identity gives
\[
 \Lambda_\gamma\le h_*+\delta'
       -\tfrac\eta2\int_{O^c}\Omega_\gamma^2\,dH+o_\gamma(1),
 \qquad h_*:=\max\mathfrak h.
\]
Use the already proved top limit, then let $\delta'\downarrow0$. This proves
\eqref{f16v41:abstract-localization}. All limits in this proof fix $L$ first.
\end{proof}

\subsection{All compact flat backgrounds and their exact transverse thresholds}

The periodic lattice has no twisted boundary conditions. Flatness makes
parallel transport path independent on the universal cover. Its three
fundamental holonomies commute, and conversely a commuting triple specifies a
flat gauge orbit. For $\SU(2)$ write it, up to simultaneous conjugation, as
\[
 H_\mu=\operatorname{Exp}(\vartheta_\mu e_3),\qquad
 \varphi_\mu=2\vartheta_\mu\pmod {2\pi},\qquad \mu=1,2,3.
\]
The factor two is the adjoint charged character. Choosing commuting $L$th
roots gives a constant-link representative, or one may put all twists on
three seams. Gauge changes are isometries of the fine-link tangent metric.
Define
\begin{equation}
 \begin{split}
 \lambda_L(k,\varphi)&=4\sum_{\mu=1}^3
      \sin^2\!\left(\frac{2\pi k_\mu+\varphi_\mu}{2L}\right),
       \qquad k\in\{0,\ldots,L-1\}^3,\\
 \xi(t)&=\operatorname{arcosh}(1+t/2),\qquad
 F_L(\varphi)=\sum_k\xi(\lambda_L(k,\varphi)).
 \end{split}
 \label{f16v41:dispersion}
\end{equation}

\begin{proposition}[The flat-background transfer threshold with all colors]
\label{f16v41:flat-spectrum}
For a flat connection with adjoint twists $\varphi$,
\begin{equation}
 \mathfrak h(\varphi)=e^{-\mathcal E_L(\varphi)},\qquad
 \mathcal E_L(\varphi)=F_L(0)+2F_L(\varphi).
 \label{f16v41:flat-energy}
\end{equation}
At the trivial flat connection,
\begin{equation}
 e^{-3F_L(0)}=\Lambda_{{\rm G},B},
 \label{f16v41:reference-match}
\end{equation}
with exactly V38's Gaussian threshold and its original kinetic normalization.
\end{proposition}
\begin{proof}
At a flat point, differentiating the literal plaquette gives the covariant
lattice curl $d_{1,p}$; hence $H_p=d_{1,p}^*d_{1,p}$. Gauge the background
away on the universal cover. The neutral color has periodic boundary condition
and the two complex charged color lines have twists $\varphi$ and $-\varphi$.
In a Fourier mode the covariant derivative is
$d_\mu=e^{i(2\pi k_\mu+\varphi_\mu)/L}-1$. The curl symbol obeys
\[
 d_1^*d_1=|d|^2 I_3-dd^*,
\]
with eigenvalues $\lambda_L,\lambda_L,0$ if $|d|>0$, and three zeros
otherwise. Complexification preserves the trace of a real spectral function;
the two charged lines must both be counted. One half of the positive-mode sum
of $\xi$ is therefore $F_L(0)+F_L(\varphi)+F_L(-\varphi)$.
The change $k\mapsto-k$ shows $F_L(-\varphi)=F_L(\varphi)$, proving
\eqref{f16v41:flat-energy}. Gauge zero directions contribute a factor one to
the Gaussian norm, not a fictitious oscillator ground state.

To match the inherited rooted coordinates, let $R$ be the linear quotient map
from all fine-link tangent coordinates to V2's balanced chord/bridge slice
coordinates. Its kernel is the nonroot gauge-gradient space. The exact
Gaussian integration of those directions says its kinetic covariance is
$RR^*=\operatorname{diag}(\mathsf M_B,G)$; this is the original kinetic
completion, not identity on the quotient. The fine curl factors through $R$,
so $H_I=R^*\mathsf S_BR$ with the same magnetic matrix as V38. The nonzero
spectra of $R^*\mathsf S_BR$ and
$(RR^*)^{1/2}\mathsf S_B(RR^*)^{1/2}$ coincide, by the nonzero singular-value
identity. The deleted eigenvalues are zero and have norm factor one. Thus
V38's product of Mehler factors is precisely the fine-link product just
computed. This proves \eqref{f16v41:reference-match}, including normalization.
\end{proof}

\subsection{A positive winding expansion selects the central holonomies}

\begin{theorem}[Strict comparison of every flat threshold]
\label{f16v41:winding}
There are strictly positive coefficients $a_L(z)$, $z\in\mathbb Z^3\setminus\{0\}$,
with $\sum_{z\ne0}a_L(z)<\infty$, such that
\begin{equation}
 \boxed{\quad
 F_L(\varphi)-F_L(0)
   =\sum_{z\ne0}a_L(z)\{1-\cos(z\cdot\varphi)\},\qquad
 a_L(z)=L^3\int_0^\infty
   \frac{e^{-8t}I_0(2t)}{t}
       \prod_{\mu=1}^3 I_{Lz_\mu}(2t)\,dt.
 \quad}
 \label{f16v41:winding-formula}
\end{equation}
Here $I_{-j}=I_j$. Consequently $F_L(\varphi)\ge F_L(0)$, with equality
exactly at $\varphi\in(2\pi\mathbb Z)^3$. In particular, putting
$A_L=a_L(1,0,0)>0$,
\begin{equation}
 \mathcal E_L(\varphi)-\mathcal E_L(0)
       \ge4A_L\sum_{\mu=1}^3(1-\cos\varphi_\mu).
 \label{f16v41:strict-loss}
\end{equation}
The coefficients and this contrast are at fixed $L$, not uniform ultraviolet
or thermodynamic constants.
\end{theorem}
\begin{proof}
The integral for $I_0$ gives, for $s>0$,
\[
 \int_0^\infty e^{-(s+2)t}I_0(2t)\,dt
 =\frac1\pi\int_0^\pi\frac{d\theta}{s+2-2\cos\theta}
 =\frac1{\sqrt{s(s+4)}}=\xi'(s).
\]
The middle integral is elementary after the tangent half-angle substitution.
Integrating in $s$, using positivity and $\xi(0)=0$, gives
\begin{equation}
 \xi(s)=\int_0^\infty(1-e^{-st})\frac{e^{-2t}I_0(2t)}{t}\,dt.
 \label{f16v41:subordination}
\end{equation}
The integral is finite: at zero its integrand is bounded, and the elementary
bound $e^{-2t}I_0(2t)\le C(1+t)^{-1/2}$ follows from the same angle integral
and $1-\cos\theta\ge2\theta^2/\pi^2$.

The Bessel Fourier series, whose integer coefficients are positive for $t>0$,
gives the exact twisted heat trace
\[
 H_L(t,\varphi):=\sum_k e^{-t\lambda_L(k,\varphi)}
 =L^3e^{-6t}\sum_{z\in\mathbb Z^3}
       \left(\prod_\mu I_{Lz_\mu}(2t)\right)e^{iz\cdot\varphi}.
\]
The series is absolutely convergent at every positive $t$. Pairing $z$ and
$-z$ shows $H_L(t,0)-H_L(t,\varphi)$ is a sum of nonnegative terms
$L^3e^{-6t}\prod I_{Lz_\mu}(2t)(1-\cos(z\cdot\varphi))$.
Insert this difference into \eqref{f16v41:subordination} and use Tonelli to
obtain \eqref{f16v41:winding-formula}.

Absolute summability requires a separate check. The sum of the nonzero winding
coefficients before integration equals
$H_L(t,0)-L^3e^{-6t}I_0(2t)^3$. For $t\le1$ it is $O_L(t^L)$:
expand the Bessel generating series, and note that a nonzero displacement
multiple of $L$ needs at least $L$ steps. The sum of the remaining positive
series is bounded by its exponential majorant. For $t\ge1$ the difference
is between zero and $H_L(t,0)\le L^3$. Multiplication by
$e^{-2t}I_0(2t)/t$ is integrable in both ranges. This proves
$\sum a_L<\infty$, hence uniform convergence in $\varphi$. Each coefficient
is strictly positive by its positive integral. Keeping the six windings
$z=\pm e_\mu$ proves \eqref{f16v41:strict-loss}. Their equality condition
forces all $\varphi_\mu$ to be multiples of $2\pi$, and the converse is immediate.
\end{proof}

The threshold has a useful singularity check. Near $\varphi=0$,
\begin{equation}
 \mathcal E_L(\varphi)-\mathcal E_L(0)
             =\frac{2}{L}|\varphi|+O_L(|\varphi|^2).
 \label{f16v41:cusp}
\end{equation}
Indeed the $k=0$ charged term is
$\xi(\lambda_L(0,\varphi))=|\varphi|/L+O_L(|\varphi|^3)$.
The remaining finite sum is analytic and even, so its change is
$O_L(|\varphi|^2)$. This cusp comes from quantizing charged modes whose
harmonic frequency vanishes at the central point. It is not a new independent
classical potential to add to a model that still retains those same modes.
In particular a smooth Gaussian elimination with one positive gap on the
entire compact flat family would miss this rank change.

\subsection{The physical top and its stationary central-flat limit}

Let $\mathcal C_\sigma$ be the flat gauge orbit with fundamental holonomies
$\sigma_\mu I$, where $\sigma\in\{+1,-1\}^3$, and put
$\mathcal C_*:=\bigcup_\sigma\mathcal C_\sigma$. These are eight disjoint
compact orbits, not eight isolated points of the fine-link space.

\begin{theorem}[Two-sided native top identification and fixed-volume vacuum localization]
\label{f16v41:main}
At every fixed $B=m^3$, $m\ge4$, for the original normalized Wilson transfer,
\begin{equation}
 \boxed{\qquad
 \lim_{\gamma\to\infty}\Lambda_\gamma
     =\Lambda_{{\rm G},B}
     =\exp\left[-3\sum_{k\in\{0,\ldots,2m-1\}^3}
       \operatorname{arcosh}\left(1+2\sum_{\mu=1}^3
                      \sin^2\frac{\pi k_\mu}{2m}\right)\right].
 \qquad}
 \label{f16v41:top-limit}
\end{equation}
For every open neighborhood $O\supset\mathcal C_*$,
\begin{equation}
 \int_{O^c}\Omega_\gamma^2\,dH\longrightarrow0.
 \label{f16v41:vacuum-concentration}
\end{equation}
More precisely the complete stationary probabilities have the weak limit
\begin{equation}
 \boxed{\qquad
 \Omega_\gamma^2dH\ \Rightarrow\
             \frac18\sum_{\sigma\in\{\pm1\}^3}\mathfrak m_\sigma,
 \qquad}
 \label{f16v41:vacuum-limit}
\end{equation}
where $\mathfrak m_\sigma$ is normalized gauge-orbit Haar measure on
$\mathcal C_\sigma$. Equivalently every continuous gauge-invariant observable
$F$ converges in expectation to $\frac18\sum_\sigma F(U_\sigma)$, for any
representatives $U_\sigma\in\mathcal C_\sigma$.
\end{theorem}
\begin{proof}
Every compact periodic $\SU(2)$ flat connection has the commuting-holonomy
form used above. The winding theorem shows that its harmonic norm is at most
the trivial value, with equality exactly when every fundamental holonomy is
central. Thus the maximizing set in Theorem~\ref{f16v41:harmonic-principle}
is precisely $\mathcal C_*$. Apply that theorem and
\eqref{f16v41:reference-match}; these prove \eqref{f16v41:top-limit} and
\eqref{f16v41:vacuum-concentration}.

For the weights, multiply every link crossing one chosen dual periodic plane
by $-I$. Every spatial plaquette contains zero or two such factors. Both its
action and the normalized kinetic kernel are unchanged when the operation is
applied to both endpoints. These three commuting transformations generate the
usual $\mathbb Z_2^3$ center symmetry and permute the eight $\mathcal C_\sigma$
transitively. They preserve product Haar and positivity. Uniqueness of the
positive normalized Perron vector implies its invariance under them, as well
as under the local gauge group.

Every subsequential weak limit exists by compactness of $\mathcal X_L$, is
supported on $\mathcal C_*$ by the concentration assertion, and inherits these
continuous symmetries. Each $\mathcal C_\sigma$ is a single compact gauge
orbit, whose unique invariant probability is the Haar pushforward
$\mathfrak m_\sigma$. The center symmetry forces all eight weights equal.
This identifies every subsequential limit and proves \eqref{f16v41:vacuum-limit}.
No singular gauge section or Faddeev--Popov determinant is introduced. In
particular the limiting probability is singular relative to full link Haar;
no $L^2$ limit of normalized vacuum vectors is asserted.
\end{proof}

\subsection{What the new top estimate does and does not close}

The equality in \eqref{f16v41:top-limit} supplies the missing fixed-$B$ upper
comparison of V38--V40. It is a theorem about the same native physical top,
not about a surrogate local sampler. The proof uses no endpoint moment
approximation and no assumed vacuum overlap. It also identifies the weak
stationary probability, whereas V40 studies the nonstationary return prepared
by one particular endpoint source. Those laws can have total variation tending
to one while both concentrate near the same central-flat set on different
shrinking scales. Weak concentration does not determine that scale.

The theorem is deliberately qualitative in $\gamma$. In particular it does
\emph{not} imply $n_\gamma\log(\Lambda_\gamma/\Lambda_{{\rm G},B})\to0$.
Therefore V40's Gamma variables must still be centered at their declared
Gaussian threshold unless a rate is proved separately. Neither a rate for the
first Perron correction nor an estimate of the first excited ratio follows
from the normal-mode maximum principle. A positive finite limiting top is not
a positive mass above that top.

The next useful calculation is a uniform effective analysis in neighborhoods
of the central-flat orbits, keeping their nonquadratic soft directions and
center sectors. The cusp \eqref{f16v41:cusp} identifies why a Gaussian normal
elimination valid only off those orbits is insufficient there. Obtaining a
subleading spectral scale, its low-level operator, and a comparison to the
actual same-top transfer is a different task from evaluating another endpoint
coefficient. No physical time is defined by $\gamma$ or by an auxiliary update
clock. Fixed-lattice weak coupling, spatial removal, physical time calibration,
and an interacting continuum construction remain separate limits.

\subsection{Verification and preservation}

The accompanying diagnostic is
\begin{center}\small\path{verify_ym_f16_v41_flat_top.py}.\end{center}
Its exact fixtures check the localization identity, normal-coordinate
sum-of-squares bound, Gaussian prefactors, covariant-curl symbols, quotient
nonzero-spectrum identity, winding-series coefficients, color/center counts,
and the limiting statement's scope. Floating-point exploratory checks, when
used, are reported separately and are not integral or spectral enclosures.
The new checker passes 98 exact assertions. All 37 available V4--V40 checkers
were rerun unchanged and passed. No finite checker is described as formal
verification of the analytic localization proof.

Removing this exact insertion and reversing the one title-version change
recovers the supplied V40 byte for byte. The earliest exact fine-link transfer,
normalization, and rooted Gaussian identities are inherited inputs; the
singular-stratum operator comparison and winding argument are proved in this
section. The entire prior analytic archive is not independently recertified.
The next companion checkpoint remains V45.

\begin{firewall}
V41 proves the fixed-spatial-regulator limit of the original physical Perron
value, identifies its maximizing flat backgrounds, and proves weak convergence
of its stationary probability to the equal mixture of eight central-flat gauge
orbits. It does not estimate the first excitation, the subleading Perron scale,
a uniform-in-volume gap, a stationary fluctuation rate, a physical time, or an
interacting continuum limit. Its Gaussian threshold is exactly the inherited
one. All original kinetic scalars, magnetic words, and gauge constraints used
in the top comparison are retained; no endpoint source or sampling clock is
substituted for the transfer.
\end{firewall}
% END F16 V41 APPEND

% BEGIN F16 V42 APPEND -- 2026-09-18
% Insert before the final document-ending command of the cumulative V41.
% No predecessor proof, measure, scalar normalization, or clock is edited.

\clearpage
\section{V42: forbidden-region decay and exponentially close center sectors}
\label{f16v42:context}

V41 identifies the native top and eight symmetry-related regions supporting its
stationary probability, but does not estimate an excitation. Before constructing
a subleading oscillator model, we examine an implication of those eight regions
for the \emph{same} transfer. A fixed neighborhood of all maximizing flat orbits
has a complement whose compressed transfer norm stays strictly below the top.
Exponential weighting of the actual narrow Wilson kernel then gives quantitative
forbidden-region decay. The resulting center-character quasimodes lie in the
periodic-gauge-invariant Hilbert space and produce exponentially close sector
edges. No sampling generator or new kinetic normalization is introduced.

The conclusion is at fixed spatial regulator. It is an \emph{upper} bound on
seven center-charged excitation energies, not a positive lower mass bound, an
exact tunneling asymptotic, or a result for the center-neutral excitation gap.
An exact Schur reduction is also available on the union of the eight regions;
its entries remain operators on their complete local spaces, not an assumed
eight-dimensional effective Hamiltonian.

\subsection{Dependencies, exact symmetries, and the spectral question}

The V41 appendix and audit were read completely. The inherited input is its
fixed-regulator harmonic principle, including the strict maximizing set and
the original normalized fine-link kernel. The extension of its localization
argument to a compressed forbidden set is made explicit below. We do not use
V35--V40's endpoint asymptotics or any earlier auxiliary sampling gap. Targeted
searches of the supplied archives were followed by mounted-text inspection;
they are not an exhaustive novelty or proof audit. The reversible ground-state
transform already occurs in f15 V78 and its cited predecessors. The monograph's
Schur/Feshbach calculus is likewise credited, rather than renamed a new general
construction. The next scheduled companion checkpoint remains V45.

For context, Klein--L\'eonard--Rosenberger,
\href{https://arxiv.org/abs/1309.3918}{\emph{Agmon-type estimates for a class of
jump processes}}, Section 1, relates exponentially weighted estimates for
nonlocal operators to multiwell spectral splitting. Its hypotheses include
nondegenerate point wells in a different space. They are not assumed here:
our wells are compact gauge orbits in a singular flat set. The bounded-kernel
and resolvent argument needed below is proved directly. The center/electric-flux
sector terminology is consistent with de Forcrand--von Smekal,
\href{https://arxiv.org/abs/hep-lat/0107018}{\emph{'t Hooft Loops, Electric Flux
Sectors and Confinement in SU(2) Yang--Mills Theory}}; its numerical or thermal
conclusions are not inputs to the fixed-lattice theorem below.

Fix $L=2m$, $m\ge4$, $B=m^3$, $E=24B$, and
$\mathcal X_L=\SU(2)^E$ with probability Haar measure. Keep
\[
 T_\gamma=\overline T_\gamma
   =M_{b_\gamma}C_\gamma^{\otimes E}M_{b_\gamma},\qquad
 b_\gamma=e^{-\gamma S/2},\qquad
 \Lambda_\gamma=\|T_\gamma\|,\qquad T_\gamma\Omega_\gamma
       =\Lambda_\gamma\Omega_\gamma,\quad \|\Omega_\gamma\|_2=1.
\]
Here $T_\gamma$ is V41's full fine-link operator, not its root-coordinate
matrix before a further restriction. Its unique positive Perron vector is
periodic-gauge invariant, so its top is the same physical top. Let
\[
 \lambda_*:=\Lambda_{{\rm G},B}>0,\qquad
 \mathcal C_*:=\bigcup_{\sigma\in\{\pm1\}^3}\mathcal C_\sigma
\]
retain V41's meanings. Its theorem gives $\Lambda_\gamma\to\lambda_*$ and
identifies exactly these maximizing flat orbits.

For $\eta\in\mathsf Z:=\{\pm1\}^3$, let $z_\eta$ multiply every positive
$\mu$-link crossing the chosen periodic $\mu$-seam by $\eta_\mu I$.
These are commuting Haar-preserving isometries. They commute with periodic
gauge transformations, preserve each original plaquette and the two-end
kinetic kernel, and send $\mathcal C_\sigma$ to $\mathcal C_{\eta\sigma}$.
They are not identified with periodic gauge transformations: fundamental
Wilson-loop traces distinguish their actions. Write
\[
 (\mathsf Z_\eta f)(U)=f(z_\eta^{-1}U),\qquad
 \chi_a(\eta)=\prod_{\mu=1}^3\eta_\mu^{a_\mu},\quad a\in\{0,1\}^3,
\]
and define the physical character spaces
\begin{equation}
 \mathcal H_a=\{f\in L^2(\mathcal X_L)^{\rm gauge}:
            \mathsf Z_\eta f=\chi_a(\eta)f\ \hbox{for every }\eta\}.
 \label{f16v42:sectors}
\end{equation}
They are mutually orthogonal reducing subspaces. The Perron vector belongs
to $\mathcal H_0$. All constants below may depend on this fixed $L$ and on
fixed neighborhoods; none is claimed uniform in spatial volume.

\subsection{Exponential weighting of the literal normalized Wilson kernel}

Let $d_L$ be the round product geodesic distance. The raw kinetic density,
before the two magnetic factors, is denoted $c_\gamma(U,V)$; each of its rows
integrates to one with the \emph{original} $z_\gamma^E$.

\begin{proposition}[Uniform displacement weights and a bounded conjugation]
\label{f16v42:weighted-kernel}
There are finite $K_L,A_L$ and $b_L>0$ such that, for all $\gamma\ge1$,
\begin{align}
 \sup_U\int \sqrt\gamma\,d_L(U,V)e^{\sqrt\gamma d_L(U,V)}
                      c_\gamma(U,V)\,dH(V)&\le K_L,
                      \label{f16v42:exp-moment}\\
 \sup_U\int_{d_L(U,V)\ge r}c_\gamma(U,V)\,dH(V)
                      &\le A_Le^{-b_L\gamma r^2},\qquad r\ge0.
                      \label{f16v42:jump-tail}
\end{align}
If $\phi$ is a real $1$-Lipschitz function on $\mathcal X_L$, and
$0\le\alpha\le1$, multiplication by $W_\alpha=e^{\alpha\sqrt\gamma\phi}$
is bounded and invertible at each fixed coupling, and
\begin{equation}
 \|W_\alpha T_\gamma W_\alpha^{-1}-T_\gamma\|
             \le K_L\alpha.
 \label{f16v42:conjugation}
\end{equation}
For measurable sets $A,D$ separated by distance at least $r$,
$\|1_A T_\gamma1_D\|\le A_Le^{-b_L\gamma r^2}$.
\end{proposition}
\begin{proof}
For one displacement group, write its angle as $\theta\in[0,\pi]$.
The exact class density is proportional to
$e^{-\gamma(1-\cos\theta)}\sin^2\theta\,d\theta$.
The inequalities $1-\cos\theta\ge2\theta^2/\pi^2$ and
$z_\gamma\ge c\gamma^{-3/2}$, the latter obtained by integrating
$0\le\theta\le\gamma^{-1/2}$, show that the rescaled radial density is
bounded by $C x^2e^{-2x^2/\pi^2}\,dx$, $x=\sqrt\gamma\theta$.
The product displacement distance is $(\sum_e\theta_e^2)^{1/2}$.
Its exponential moment and its Gaussian tail follow by integrating this
fixed-dimensional product bound and splitting a Gaussian exponent in half.
Haar invariance makes the estimates independent of $U$. Symmetry supplies the
same column bounds. No large dimension is hidden as a volume-independent
constant.

The absolute difference kernel in \eqref{f16v42:conjugation} is bounded by
\[
 c_\gamma(U,V)\bigl(e^{\alpha\sqrt\gamma d_L(U,V)}-1\bigr)
 \le\alpha\sqrt\gamma d_L(U,V)e^{\sqrt\gamma d_L(U,V)}c_\gamma(U,V),
\]
because $0<b_\gamma\le1$. The row and column Schur estimates prove the norm
bound. This is not a Loewner inequality for an asymmetric weighted kernel.
The separated-support estimate follows from the same Schur argument using
\eqref{f16v42:jump-tail}. Only a Lipschitz weight is required; no derivative
through a principal logarithm cut is taken.
\end{proof}

\subsection{The complement of fixed central neighborhoods has a genuine barrier}

\begin{proposition}[A forbidden-region norm gap at the same native top]
\label{f16v42:barrier}
Let $\mathcal O$ be a fixed open neighborhood of $\mathcal C_*$, put
$P=1_{\mathcal O}$ and $Q=I-P$. There exist
$0<\delta<\lambda_*/8$ and $\gamma_0<\infty$ such that, for
$\gamma\ge\gamma_0$,
\begin{equation}
 \|QT_\gamma Q\|\le\lambda_*-4\delta,
       \qquad \Lambda_\gamma\ge\lambda_*-\delta.
 \label{f16v42:barrier-gap}
\end{equation}
The first norm is on the complete fine-link space; it consequently bounds its
physical restriction too.
\end{proposition}
\begin{proof}
The compact set $\mathcal F_L\setminus\mathcal O$ contains no maximizing
flat point. V41's continuous harmonic threshold therefore has a strict maximum
below $\lambda_*$ there, if that set is nonempty. Repeat the upper-bound
partition in Theorem~\ref{f16v41:harmonic-principle} for functions supported
in $\mathcal O^c$: use finitely many normal charts centered on these remaining
flat points, a piece inside $\mathcal O$ which contributes zero, and pieces
whose relevant supports are separated from the full flat set. The last pieces
have a positive action floor. The local norm errors can be fixed smaller than
one quarter of the strict threshold difference; the partition commutator then
tends to zero by Proposition~\ref{f16v41:IMS}. It follows that
$\limsup\|QT_\gamma Q\|<\lambda_*$. If no flat point remains, compactness
gives a positive action floor directly. The already-proved top limit and a
smaller positive $\delta$ give \eqref{f16v42:barrier-gap}.

This is not obtained from the small probability of $\mathcal O^c$. It is an
operator norm bound on all functions supported there. Nor is $S$ asserted
positive on that whole complement: noncentral \emph{flat} connections can
remain there, and their lower harmonic thresholds are essential.
\end{proof}

\begin{theorem}[Exponential localization of a whole native spectral band]
\label{f16v42:Agmon}
With $\mathcal O,\delta$ as above, put
\[
 \lambda_c=\lambda_*-\delta,\qquad
 \mathsf E_\gamma=1_{[\lambda_c,1]}(T_\gamma),\qquad
 \alpha=\min\{1,\delta/(1+K_L)\}>0.
\]
There is $C_{L,\mathcal O}<\infty$, independent of coupling and spectral-band
rank, such that
\begin{equation}
 \boxed{\quad
 \big\|e^{\alpha\sqrt\gamma d_L(\cdot,\mathcal O)}\mathsf E_\gamma\big\|
       \le C_{L,\mathcal O}.
 \quad}
 \label{f16v42:band-bound}
\end{equation}
In particular,
\begin{equation}
 \|1_{\{d_L(\cdot,\mathcal O)\ge r\}}\mathsf E_\gamma\|
       \le C_{L,\mathcal O}e^{-\alpha r\sqrt\gamma},\qquad
 \int e^{2\alpha\sqrt\gamma d_L(U,\mathcal O)}\Omega_\gamma(U)^2dH(U)
       \le C_{L,\mathcal O}^2.
 \label{f16v42:vacuum-tail}
\end{equation}
All these statements concern the actual full transfer and its own top.
\end{theorem}
\begin{proof}
Write $W=e^{\alpha\sqrt\gamma d_L(\cdot,\mathcal O)}$ and
$T_\alpha=WT_\gamma W^{-1}$. On $Q L^2$ set
$A=QT_\gamma Q$ and $A_\alpha=QT_\alpha Q$.
The preceding propositions give
$\|A_\alpha\|\le\lambda_*-3\delta$ and $\|T_\alpha\|\le1+\delta$.
Let $S_\gamma=T_\gamma|_{\operatorname{Ran}\mathsf E_\gamma}$, so
$S_\gamma\succeq\lambda_c I$. Regard
$X=WQ\mathsf E_\gamma$ as an operator from that spectral space to $QL^2$.
Since $WP=P$, the exact intertwining equation is
\[
 X S_\gamma-A_\alpha X
       =QT_\alpha P\mathsf E_\gamma=:B_\alpha.
\]
It has the norm-convergent solution
\begin{equation}
 X=\int_0^\infty e^{tA_\alpha}B_\alpha e^{-tS_\gamma}\,dt,
       \qquad \|X\|\le\frac{1+\delta}{2\delta}.
 \label{f16v42:Sylvester}
\end{equation}
Indeed differentiate $e^{tA_\alpha}Xe^{-tS_\gamma}$. Its limit at infinity
is zero, since $\lambda_c-\|A_\alpha\|\ge2\delta$; at each fixed coupling
$X$ is bounded before this uniform estimate is taken. This proves the integral
identity and bound without summing separate eigenvector estimates. Adding
$P\mathsf E_\gamma$ proves \eqref{f16v42:band-bound}. The Perron vector is
in the band by \eqref{f16v42:barrier-gap}. The tail conclusions follow by
multiplication. The integration parameter in \eqref{f16v42:Sylvester} is a
resolvent proof parameter, not a new physical time.
\end{proof}

\subsection{An exact eight-region Schur map, with all returning paths retained}

Choose $r_0>0$ so small that the closed $4r_0$-neighborhoods of the eight
$\mathcal C_\sigma$ are disjoint. This is possible because they are disjoint
compact subsets of the fixed compact metric space. Let
$\mathcal N_{j,\sigma}=\{d_L(U,\mathcal C_\sigma)<jr_0\}$ and
$\mathcal N_j=\bigcup_\sigma\mathcal N_{j,\sigma}$.
Apply the preceding results to $\mathcal O=\mathcal N_1$ and keep their
$\delta,\lambda_c,\alpha$. Put $P_\sigma=1_{\mathcal N_{1,\sigma}}$,
$P=\sum_\sigma P_\sigma$ and $Q=I-P$.

\begin{theorem}[Uniform forbidden complement and exponentially small inter-region returns]
\label{f16v42:Feshbach}
For $\gamma\ge\gamma_0$ and real $z\ge\lambda_c$, the operator
$z-QT_\gamma Q$ is positive and invertible on $QL^2$. On $PL^2$ define
\begin{equation}
 \mathfrak F_\gamma(z)=PT_\gamma P+
    PT_\gamma Q(z-QT_\gamma Q)^{-1}QT_\gamma P.
 \label{f16v42:Feshbach-map}
\end{equation}
Then $z$ is an eigenvalue of $T_\gamma$ if and only if it is an eigenvalue
of $\mathfrak F_\gamma(z)$, with multiplicities preserved. More generally,
$z-T_\gamma$ is invertible if and only if
$zP-\mathfrak F_\gamma(z)$ is invertible. Eigenvectors satisfy
\begin{equation}
 \psi=u+(z-QT_\gamma Q)^{-1}QT_\gamma P u,\quad u=P\psi,
 \qquad
 \|\psi\|^2=\langle u,(P-\mathfrak F_\gamma'(z))u\rangle.
 \label{f16v42:reconstruction}
\end{equation}
For distinct regions $\sigma,\tau$, uniformly in such $z$,
\begin{equation}
 \boxed{\quad
 \|P_\sigma\mathfrak F_\gamma(z)P_\tau\|
 \le C_{L,r_0}
 e^{-\alpha\sqrt\gamma\operatorname{dist}(\mathcal N_{1,\sigma},
                                           \mathcal N_{1,\tau})}.
 \quad}
 \label{f16v42:interwell}
\end{equation}
All statements restrict to periodic-gauge-invariant functions. Center symmetry
permutes the diagonal blocks by exact unitary equivalence.
\end{theorem}
\begin{proof}
The unweighted complement inverse has norm at most $(3\delta)^{-1}$ by
\eqref{f16v42:barrier-gap}. Block Gaussian elimination of $z-T_\gamma$
gives \eqref{f16v42:Feshbach-map} and the equivalences. The complement component
of an eigenvector is forced by its second block equation. Differentiating the
inverse gives
\[
 \mathfrak F_\gamma'(z)
   =-PT_\gamma Q(z-QT_\gamma Q)^{-2}QT_\gamma P\preceq0,
\]
which proves the norm identity. In particular the returning term is positive,
and its energy derivative is not omitted from reconstruction.

For the off-diagonal estimate take the different weight
$W_\tau=e^{\alpha\sqrt\gamma d_L(\cdot,\mathcal N_{1,\tau})}$.
The same conjugation bound gives
$\|QW_\tau T_\gamma W_\tau^{-1}Q\|\le\lambda_*-3\delta$.
Its resolvent at $z\ge\lambda_c$ has norm at most $(2\delta)^{-1}$.
Conjugating \eqref{f16v42:Feshbach-map}, using that all multiplication
projections commute with the weight, gives norm at most
\[
 (1+\delta)+(1+\delta)^2/(2\delta).
\]
On $P_\tau$ the weight is one; on $P_\sigma$ its inverse is at most the
exponential in \eqref{f16v42:interwell}. This proves that estimate, including
arbitrarily many returns through the forbidden region. Gauge and center
commutation follows because the regions and the original transfer have those
symmetries. No replacement probability or conditional row normalization is used.
\end{proof}

The Schur identity itself is standard. The new input is a \emph{proved} native
complement interval and its exponential cross-region control. Each diagonal
block still acts on an infinite-dimensional space and retains every soft and
hard coordinate inside that region. This is not an eight-state spectral model;
no inverse on the complement of just eight chosen vectors is supplied.

\subsection{Physical center-character quasimodes and sector edges}

There are smooth periodic-gauge-invariant functions $g_\sigma$ with
$0\le g_\sigma\le1$, equal to one on $\mathcal N_{3,\sigma}$, supported in
$\mathcal N_{4,\sigma}$, and satisfying
$g_\sigma(z_\eta U)=g_{\eta\sigma}(U)$.
To construct them, start with a smooth bump between the closures of the two
neighborhoods at $\sigma=(1,1,1)$, average over the compact gauge group, and
translate by the center action. Isometry preserves its plateau and support.
The eight supports are disjoint. For nonzero $a\in\{0,1\}^3$ put
\begin{equation}
 F_a(U)=\sum_\sigma\chi_a(\sigma)g_\sigma(U),\qquad
 u_{\gamma,a}=\frac{F_a\Omega_\gamma}{\|F_a\Omega_\gamma\|_2},
 \quad u_{\gamma,0}=\Omega_\gamma.
 \label{f16v42:quasimodes}
\end{equation}
The denominator is nonzero at all sufficiently large couplings.

\begin{theorem}[At least eight same-top eigenvalues with an exponential splitting bound]
\label{f16v42:sector-edges}
At fixed $L$ there are $C_L,\kappa_L>0$ and a finite threshold such that the
eight vectors in \eqref{f16v42:quasimodes} are orthonormal physical vectors,
$u_{\gamma,a}\in\mathcal H_a$, and
\begin{equation}
 \|(T_\gamma-\Lambda_\gamma)u_{\gamma,a}\|_2
                   \le C_Le^{-\kappa_L\sqrt\gamma}.
 \label{f16v42:quasimode-error}
\end{equation}
Let $\Lambda_{\gamma,a}=\|T_\gamma|_{\mathcal H_a}\|$. Then for every
nontrivial character $a$,
\begin{equation}
 \boxed{\quad
 0<1-\frac{\Lambda_{\gamma,a}}{\Lambda_\gamma}
       \le C_Le^{-\kappa_L\sqrt\gamma},\qquad
 0<-\log\frac{\Lambda_{\gamma,a}}{\Lambda_\gamma}
       \le C_Le^{-\kappa_L\sqrt\gamma}.
 \quad}
 \label{f16v42:sector-splitting}
\end{equation}
In particular, if $\lambda_{0,\gamma}=\Lambda_\gamma\ge\lambda_{1,\gamma}
\ge\cdots$ are the physical eigenvalues counted with multiplicity, then
\begin{equation}
 0<-\log(\lambda_{j,\gamma}/\Lambda_\gamma)
       \le C_Le^{-\kappa_L\sqrt\gamma},\qquad 1\le j\le7.
 \label{f16v42:eight-eigenvalues}
\end{equation}
The estimate asserts at least one near-top eigenvalue in each center sector,
not that there are exactly eight eigenvalues in this interval.
\end{theorem}
\begin{proof}
Center invariance of $\Omega_\gamma$ and covariance of the $g_\sigma$ give
$\mathsf Z_\eta(F_a\Omega_\gamma)=\chi_a(\eta)F_a\Omega_\gamma$.
Characters with different labels are orthogonal, including their orthogonality
to the Perron vector for $a\ne0$. Since $F_a^2=1$ on $\mathcal N_3$ and
$|F_a|\le1$ everywhere, \eqref{f16v42:vacuum-tail} gives
\[
 1-C_Le^{-4\alpha r_0\sqrt\gamma}
              \le\|F_a\Omega_\gamma\|_2^2\le1.
\]
In particular no small approximate vector norm is divided out.

The eigenvector equation gives the exact commutator identity
\[
 (\Lambda_\gamma-T_\gamma)(F_a\Omega_\gamma)
                         =[M_{F_a},T_\gamma]\Omega_\gamma.
\]
Split the input into $1_{\mathcal N_2}\Omega_\gamma$ and its complement.
The second part has norm at most $C_Le^{-\alpha r_0\sqrt\gamma}$, and the
commutator has norm at most two. If the first input lies in
$\mathcal N_{2,\sigma}$, the difference $F_a(U)-F_a(V)$ is zero when the
output lies in $\mathcal N_{3,\sigma}$. Otherwise $d_L(U,V)\ge r_0$.
The separated-jump Schur estimate therefore bounds this restricted commutator
by $2A_Le^{-b_L\gamma r_0^2}$. Combining the two parts and the norm bound proves
\eqref{f16v42:quasimode-error}, for a smaller fixed $\kappa_L>0$.

The sector restriction is compact positive self-adjoint. Its Rayleigh quotient
on $u_{\gamma,a}$ is at least
$\Lambda_\gamma-C_Le^{-\kappa_L\sqrt\gamma}$ and at most
$\Lambda_\gamma$. For sufficiently large coupling this is positive, and hence
the sector norm is an attained eigenvalue. It is strictly smaller than
$\Lambda_\gamma$ for $a\ne0$, by simplicity and center invariance of the
Perron vector. Since $\Lambda_\gamma\ge\lambda_*/2$ eventually, division and
$-\log(1-x)\le2x$ for $0\le x\le1/2$ prove
\eqref{f16v42:sector-splitting}. The eight mutually orthogonal eigenvectors,
one from each sector, give \eqref{f16v42:eight-eigenvalues} by the min--max
principle. They need not be the only near-top vectors.
\end{proof}

The proof uses the existing vacuum as a variational factor; it does not provide
a numerical representation of that vector. Its trial multipliers are genuinely
gauge invariant. This is why the result concerns physical sector eigenvalues,
unlike a non-invariant trial justified only by equality of the two Perron tops.
The positive constants and entry threshold are not optimized or numerically
enclosed. No matching tunneling lower bound or leading tunneling action is
claimed.

\subsection{Stationary center memory and the neutral-sector distinction}

Use the \emph{existing} ground-state transform, now on the literal fine-link
space, with $d\mu_\gamma=\Omega_\gamma^2dH$:
\begin{equation}
 (\mathcal P_\gamma f)(U)=\frac{1}{\Lambda_\gamma\Omega_\gamma(U)}
          \int T_\gamma(U,V)\Omega_\gamma(V)f(V)\,dH(V).
 \label{f16v42:Doob}
\end{equation}
This is a reversible positive contraction with stationary law $\mu_\gamma$;
multiplication by $\Omega_\gamma$ is unitary to $T_\gamma/\Lambda_\gamma$.
Its transition is not obtained by dividing rows of $T_\gamma$ by their sums.
It is one original transfer step at its \emph{actual} Perron top.

\begin{proposition}[Exponentially long stationary center correlations]
\label{f16v42:memory}
For each nontrivial $a$, $\mu_\gamma F_a=0$. For every integer $N\ge0$,
\begin{equation}
 0\le 1-
 \frac{\mathbb E_{\mu_\gamma}[F_a(U_0)F_a(U_N)]}
      {\mathbb E_{\mu_\gamma}F_a^2}
 \le\min\{1,C_LN e^{-\kappa_L\sqrt\gamma}\},
 \label{f16v42:correlation}
\end{equation}
where $(U_j)$ is the stationary chain of \eqref{f16v42:Doob}. In particular,
the normalized correlation tends to one for
$N_\gamma=o(e^{\kappa_L\sqrt\gamma})$.
\end{proposition}
\begin{proof}
The normalized correlation is
$\langle u_{\gamma,a},(T_\gamma/\Lambda_\gamma)^N u_{\gamma,a}\rangle$.
The operator is a positive contraction. Its spectral theorem and
$0\le1-x^N\le N(1-x)$, $0\le x\le1$, bound the deficit by $N$ times the
one-step Rayleigh deficit. Apply \eqref{f16v42:quasimode-error} and
$\Lambda_\gamma\ge\lambda_*/2$. This is a stationary correlation estimate,
not a claim that a finite-time endpoint source has prepared that stationary
law or that an entire path stays in a well without a separate event estimate.
\end{proof}

\begin{proposition}[The center-charged band is invisible to center-neutral insertions]
\label{f16v42:selection}
If a bounded gauge-invariant multiplication $F$ is invariant under every
$z_\eta$, then $F\Omega_\gamma\in\mathcal H_0$ and
\[
 \langle\psi_a,F\Omega_\gamma\rangle=0
      \quad(a\ne0,\ \psi_a\in\mathcal H_a).
\]
In particular this holds for every bounded function of a fixed collection of
contractible Wilson-loop traces. Conversely a multiplication with character
$\chi_a$ creates a vector in $\mathcal H_a$ from the vacuum.
\end{proposition}
\begin{proof}
Apply $\mathsf Z_\eta$ inside the inner product. The Perron vector is invariant,
and a nontrivial character has a group element with value $-1$. This forces the
claimed inner product to vanish. The center seam is a closed mod-two cocycle;
its pairing with a contractible loop is zero. Equivalently every elementary
plaquette crosses each seam an even number of times, and a contractible loop
bounds a lattice surface. Centrality allows the signs to cancel without
reordering noncommutative link factors. Its trace is therefore center invariant.
The last assertion is the same character multiplication identity.
\end{proof}

There are thus seven nontrivial center-charge channels with small transfer
energies, but no estimate here for the first excited value in $\mathcal H_0$.
This distinction cannot be erased by calling all eight sectors gauge equivalent.
The center-neutral spectral problem may contain additional near-top states;
no isolated eight-state cluster or complement gap after removing just these
quasimodes is proved. The full operator-valued map
\eqref{f16v42:Feshbach-map} deliberately retains those states.

\subsection{What changes at the spectral frontier}

V42 supplies an actual same-top excited-sector comparison at fixed spatial
regulator, together with exponential forbidden-region localization and an exact
regional elimination whose cross-region entries are exponentially small. It
does not determine the common shift $\Lambda_\gamma-\lambda_*$. In particular
V40's entrance law still cannot be recentered at the native top on its moving
$1/n_\gamma$ scale without a separate rate for that common shift.

If an independently specified physical time step is $a_t(\gamma)>0$, then the
center-sector energies obey only the converted upper bound
\[
 0<E_a(\gamma):=-a_t(\gamma)^{-1}
       \log(\Lambda_{\gamma,a}/\Lambda_\gamma)
       \le\frac{C_Le^{-\kappa_L\sqrt\gamma}}{a_t(\gamma)}.
\]
Whether this tends to zero cannot be decided without that time scale; a fixed
number of spatial lattice sites is not a fixed physical-volume continuum
limit either. No continuum gaplessness or positive continuum mass is inferred.
The charged characters are often described as finite-volume electric-flux
sectors. They must not be substituted for the center-neutral local excitation
channel.

The next useful problem is now inside the diagonal blocks of
\eqref{f16v42:Feshbach-map}: identify their nonquadratic soft operator and a
quantitative remainder at the same spectral top, or obtain a lower splitting
estimate with the actual returning paths. The growing-regulator and
physical-time costs remain separate from this fixed-regulator exponential
bound. This respects the supplied memoranda's separate locality, normalization,
and physical-transfer obligations; no external Navier--Stokes theorem is a
premise of this iteration.

\subsection{Verification and preservation}

The diagnostic is
\begin{center}\small\path{verify_ym_f16_v42_center_splitting.py}.\end{center}
It checks exact center/plaquette and winding-loop parities on the actual finite
geometry, character projectors, a noncommuting weighted Sylvester fixture,
Schur reconstruction with its derivative metric, and a strictly positive
finite multiwell transfer with a known Perron vector. It also tests stationary
correlation and sector-selection identities, and a counterfixture showing why
an eighth near-top eigenvalue does not isolate an eight-state cluster. These
are finite algebraic tests, not native Wilson integral or spectral enclosures,
and not formal verification of the analytic localization argument.

The package records actual diagnostic executions, source hashes, proof-review
scope, and rendering checks. Only the literal title version and this terminal
insertion differ from V41. Removing them recovers its exact bytes. The new
native conclusions depend on V41's harmonic principle and strict maximizing
set; those inputs are stated rather than independently recertified for the
whole preceding archive. The next scheduled checkpoint remains V45.

\begin{firewall}
V42 proves exponential localization outside fixed central-flat neighborhoods,
a same-native-top Schur reduction with exponentially small inter-region blocks,
and at least one exponentially close eigenvalue in each of eight physical
center-character sectors at fixed lattice size. It gives upper bounds on the
seven charged excitation ratios and stationary center correlations, not a lower
mass bound, a center-neutral gap, an exact tunneling rate, an isolated
eight-dimensional low-energy model, or a continuum result. All magnetic words,
kinetic scalars, gauge constraints, returning paths, and the actual Perron
normalization used in these deductions are retained.
\end{firewall}
% END F16 V42 APPEND

% BEGIN F16 V43 APPEND -- 2026-09-18
% Insert before the final document-ending command of cumulative V42.
% Predecessor bodies, native kernels, probabilities, and clocks are unchanged.

\clearpage
\section{V43: a single-region spectral pencil and the neutral return metric}
\label{f16v43:context}

V42 proves exponentially close charged sector tops, but does not separate a
neutral excitation from an additional local soft state. Its regional Schur map
retains exactly the space needed for that question. We now use the center
symmetry on this \emph{operator-valued} map, without selecting one vector per
region. All near-top levels, counted with multiplicity, can be compared with
one common single-region nonlinear eigenvalue problem. Its error is exponential
and has no level-count factor. A second calculation identifies the correct
same-top neutral gap on that region: its norm is the returning-state metric,
not the flat norm of a bare compression.

These are deductions about the original finite-lattice spectrum, conditional
on V42's stated forbidden-complement bounds. They do not evaluate the remaining
single-region soft spectrum. In particular no algebraic neutral-gap scale,
finite-dimensional eight-state model, or continuum mass is asserted.

\subsection{Inputs, ownership, and the retained physical spaces}

The complete supplied V42 appendix and analytic audit were read. Its complement
interval and exponential inter-region estimate are the native inputs. The
Grand Tensor's signed Feshbach inertia and two-scale gap compiler, and f15
V78's warning about an unpaid vertical norm, were checked at identified
passages. Those general mechanisms are credited, not claimed new. The present
step folds the actual eight regions into their physical character spaces,
counts every level in a fixed spectral interval, and pays the returning norm
in a neutral-gap and source-resolvent comparison. Two empty semantic searches
were followed by mounted exact-text inventories; they do not establish absence
of related results. The next companion checkpoint remains V45.

For primary context, Dusson--Sigal--Stamm,
\href{https://arxiv.org/html/2105.02058}{\emph{The Feshbach--Schur map and
perturbation theory}}, Theorem 1.2, Corollary 1.3 and Proposition 1.4,
describe isospectral reconstruction and spectral fixed points. Their setting
uses a finite-rank retained projection for perturbation estimates. Our retained
space is infinite dimensional. We prove the compact-operator counting and
comparison statements below directly, and do not assume differentiability of
ordered eigenvalue branches at crossings. The supplied wider-literature
memorandum recommends keeping the full return; it does not supply the local
neutral gap that remains to be computed here.

Fix $L=2m$, $m\ge4$. Use V42's unchanged physical transfer on
\[
 T=T_\gamma,\qquad
 \mathcal H^{\rm phys}=L^2(\mathcal X_L,dH)^{\rm gauge}.
\]
Let $\Lambda=\Lambda_\gamma$ and $\Omega=\Omega_\gamma$ be its actual top and
unit positive vacuum. Keep V42's disjoint regions
$\mathcal O_\sigma=\mathcal N_{1,\sigma}$, their projections $P_\sigma$,
$P=\sum_\sigma P_\sigma$, and $Q=I-P$. All these projections preserve the
physical space. Write
\[
 c=\lambda_*-\delta>0,\qquad A_Q=QTQ,\qquad
 R(z)=(z-A_Q)^{-1}\quad\hbox{on }Q\mathcal H^{\rm phys}.
\]
For sufficiently large coupling the inherited bounds are
\begin{equation}
 \|A_Q\|\le c-3\delta,\qquad \Lambda\ge c,
 \qquad z\ge c\ \Longrightarrow\ \|R(z)\|\le(3\delta)^{-1}.
 \label{f16v43:input}
\end{equation}
Let $\mathsf Z=\{\pm1\}^3$, with identity $e=(1,1,1)$ and character
$\chi_a(\eta)=\prod_\mu\eta_\mu^{a_\mu}$. Its exact unitary action is
$\mathsf Z_\eta$, and $\mathcal H_a$ has V42's meaning. No center operation
is quotiented as a periodic gauge transformation.

Define the one-region physical space and its extension by zero by
\[
 \mathfrak h=L^2(\mathcal O_e,dH)^{\rm gauge},\qquad
 \iota:\mathfrak h\longrightarrow P\mathcal H^{\rm phys}.
\]
It retains every soft and hard variable of that open region. The isometries
\begin{equation}
 W_a u=8^{-1/2}\sum_{\eta\in\mathsf Z}
                    \chi_a(\eta)\mathsf Z_\eta\iota u
 \label{f16v43:Wa}
\end{equation}
have ranges $P\mathcal H_a$, are mutually orthogonal, and together exhaust
$P\mathcal H^{\rm phys}$. This follows from the disjoint translated supports
and the eight-character orthogonality identity. No local state is truncated.

\subsection{An exact character decomposition with differentiated return control}

Retain V42's map
\[
 \mathfrak F(z)=PTP+PTQ R(z)QTP.
\]
On the common space $\mathfrak h$ put
\begin{equation}
 A(z)=\iota^*\mathfrak F(z)\iota,\qquad
 B_\eta(z)=\iota^*\mathfrak F(z)\mathsf Z_\eta\iota,
 \qquad F_a(z)=W_a^*\mathfrak F(z)W_a.
 \label{f16v43:blocks}
\end{equation}
Here $B_e=A$. The letter $A(z)$ is a regional operator, not an endpoint scalar.

\begin{proposition}[Exact sector pencils and their exponentially small differences]
\label{f16v43:Fourier}
For every real $z\ge c$,
\begin{equation}
 \boxed{\quad F_a(z)=A(z)+\sum_{\eta\ne e}\chi_a(\eta)B_\eta(z).
 \quad}
 \label{f16v43:Fourier-formula}
\end{equation}
The operators $F_a(z),A(z)$ are positive, compact and self-adjoint, and
\begin{equation}
 0\preceq-F_a'(z)\preceq K I,\qquad
 0\preceq-A'(z)\preceq K I,\qquad K=(3\delta)^{-2}.
 \label{f16v43:monotonicity}
\end{equation}
There are fixed $C,\kappa>0$, depending on $L$ and the regions but not on
coupling or level count, such that, with
$\epsilon_\gamma=C e^{-\kappa\sqrt\gamma}$,
\begin{equation}
 \sup_{z\ge c}\bigl(\|F_a(z)-A(z)\|+
                         \|F_a'(z)-A'(z)\|\bigr)\le\epsilon_\gamma.
 \label{f16v43:exponential}
\end{equation}
No positivity of an individual $B_\eta$ in operator order is required.
\end{proposition}
\begin{proof}
Insert \eqref{f16v43:Wa} twice. Center commutation turns each term into
$\iota^*\mathfrak F\mathsf Z_{\eta^{-1}\tau}\iota$; each relative label occurs
eight times. This proves \eqref{f16v43:Fourier-formula}, including its
normalization. Since the group has exponent two, $\mathsf Z_\eta^*=\mathsf Z_\eta$.
Commutation makes $B_\eta$ self-adjoint, although it need not be positive.
Positivity of $F_a,A$ follows by compression of $\mathfrak F\succeq0$.
Compactness follows from compactness of $T$ and boundedness of $R(z)$.

Differentiate the actual inverse:
\[
 \mathfrak F'(z)=-PTQ R(z)^2QTP.
\]
The bound $\|T\|\le1$ proves \eqref{f16v43:monotonicity}. V42 already bounds
each off-diagonal block of $\mathfrak F$. Its same exponential conjugation
also bounds the derivative. In detail, with a weight of distance from the
source region, the conjugated complement inverse has norm at most
$(2\delta)^{-1}$ and the conjugated transfer has norm at most $1+\delta$.
Consequently the conjugate of $\mathfrak F'$ has norm at most
$(1+\delta)^2/(2\delta)^2$. Restriction to a distinct target region restores
the same exponential separation factor as for $\mathfrak F$. The minimum
separation of the fixed regions is positive. Sum the seven bounds and enlarge
$C$ to obtain \eqref{f16v43:exponential}. This calculation retains the inverse
square and all its complementary returns; it does not differentiate an
unspecified big-O remainder.
\end{proof}

At a sector eigenvalue $z\ge c$, the exact reconstruction from $u\in\mathfrak h$
is
\begin{equation}
 J_a(z)u=W_a u+R(z)Q T W_a u,
 \qquad J_a(z)^*J_a(z)=I-F_a'(z).
 \label{f16v43:reconstruction}
\end{equation}
Its residual is $(T-z)J_a(z)u=W_a(F_a(z)-z)u$. More generally, for two distinct
real parameters $z,w\ge c$,
\begin{equation}
 J_a(z)^*J_a(w)
     =I-\frac{F_a(z)-F_a(w)}{z-w}.
 \label{f16v43:cross-metric}
\end{equation}
These identities follow by multiplication and the resolvent identity. They
keep the norm and mutual inner products of the returning full vectors. An
orthonormal list of regional vectors in the flat norm is not automatically
an orthonormal list of reconstructed physical states.

\subsection{Counting every near-top level, without an isolated-cluster premise}

For a compact positive operator $C$, let $n_+(C-z)$ count its eigenvalues
strictly larger than $z$, with multiplicity. For $z\ge c$ define
\[
 N_a(z)=\dim 1_{(z,1]}(T|_{\mathcal H_a}),\qquad
 N_{\rm reg}(z)=n_+(A(z)-zI).
\]
These numbers are finite. Their threshold convention excludes eigenvalues
exactly equal to $z$.

\begin{theorem}[Exact inertia and uniform sector counting brackets]
\label{f16v43:counting}
For $z\ge c$,
\begin{equation}
 \boxed{\quad N_a(z)=n_+(F_a(z)-zI).\quad}
 \label{f16v43:inertia}
\end{equation}
For $z-\epsilon_\gamma\ge c$,
\begin{equation}
 \boxed{\quad
 N_{\rm reg}(z+\epsilon_\gamma)\le N_a(z)
                       \le N_{\rm reg}(z-\epsilon_\gamma).
 \quad}
 \label{f16v43:count-bracket}
\end{equation}
In particular, if the two regional counts equal $k$, every physical center
sector contains exactly $k$ eigenvalues above $z$, and the entire physical
space contains exactly $8k$ there. No estimate of $k$ is assumed or supplied
by symmetry alone.
\end{theorem}
\begin{proof}
Restrict $z-T$ to $\mathcal H_a=(P\mathcal H_a)\oplus(Q\mathcal H_a)$.
Its second diagonal block is positive and invertible by \eqref{f16v43:input}.
The bounded triangular congruence implementing Schur elimination transforms
its form into the direct sum of $z-F_a(z)$ and that positive block. An
invertible change of variables preserves the maximal dimension of a negative
definite subspace: apply the change and its inverse to such subspaces. Thus
the finite negative index of $z-T$ is exactly that of $z-F_a(z)$, proving
\eqref{f16v43:inertia} without a finite-dimensional retained projection.
This is the inherited inertia principle in its present compact setting.

The norm estimate and monotonicity give
\[
 A(z+\epsilon_\gamma)-(z+\epsilon_\gamma)I
 \preceq F_a(z)-zI
 \preceq A(z-\epsilon_\gamma)-(z-\epsilon_\gamma)I.
\]
Monotonicity of the finite positive index and \eqref{f16v43:inertia} prove the
bracket. Summing the eight orthogonal reducing sector counts proves its last
claim. Strict inequalities in the spectral intervals are retained at both
ends; no generic absence of a threshold eigenvalue is presumed.
\end{proof}

Let $\nu_j(z)$ be the nonincreasingly ordered eigenvalues of $A(z)$, repeated
with multiplicity. For each $j$ with $\nu_j(c)>c$, define $r_j$ to be the unique
solution
\begin{equation}
                      \nu_j(r_j)=r_j,\qquad r_j>c.
 \label{f16v43:roots}
\end{equation}
These roots have an energy-independent interpretation as well. Define the
native one-region compression
\begin{equation}
 T^{\rm one}=(P_e+Q)T(P_e+Q)
       \quad\hbox{on }(P_e+Q)\mathcal H^{\rm phys}.
 \label{f16v43:one-region-transfer}
\end{equation}
It retains the chosen region and the entire forbidden complement, but deletes
the other seven retained regions. Its Schur map is exactly $A(z)$, so its
eigenvalues above $c$ are precisely the $r_j$. Thus the reference below is an
actual compressed native transfer, not an assumed oscillator Hamiltonian.
Its top is not substituted for the full physical top without the stated
error comparison. Powers of this compression restrict slice-time paths; no
continuous-time killing or new physical clock is implied.

These roots are in descending order. Existence follows because $A(z)$ is
bounded as $z\to\infty$, while $z$ increases without bound. Uniqueness follows
from monotonicity. In fact
\[
 t-s\le[t-\nu_j(t)]-[s-\nu_j(s)]\le(1+K)(t-s),\qquad t>s\ge c.
\]
This uses the min--max principle and \eqref{f16v43:monotonicity}; it remains
valid at crossings and repeated eigenvalues. At $z=\Lambda$, positivity of
$\Lambda-T$ implies $\Lambda P-\mathfrak F(\Lambda)\succeq0$, so every root
is at most $\Lambda\le1$.

\begin{theorem}[One regional ladder for all eight native sectors]
\label{f16v43:ladders}
If $r_j>c+\epsilon_\gamma$, the $j$th eigenvalue $\lambda_{a,j}$ of
$T|_{\mathcal H_a}$ exists above $c$ and satisfies
\begin{equation}
 \boxed{\qquad |\lambda_{a,j}-r_j|\le\epsilon_\gamma
                         \quad\hbox{for every }a.\qquad}
 \label{f16v43:ladder-bound}
\end{equation}
The indexing here starts at $j=1$ in \emph{each} sector; in the neutral sector
$\lambda_{0,1}=\Lambda$. It counts multiplicity and introduces no factor $j$.
In particular corresponding sector levels differ by at most $2\epsilon_\gamma$.
For two consecutive roots above this threshold,
\begin{equation}
 |(\lambda_{a,j}-\lambda_{a,j+1})-(r_j-r_{j+1})|
                                      \le2\epsilon_\gamma.
 \label{f16v43:level-spacing}
\end{equation}
At sufficiently large coupling $r_1>c+\epsilon_\gamma$, and
\begin{equation}
 \left|\log\frac{\Lambda}{\lambda_{a,j}}-
                       \log\frac{r_1}{r_j}\right|
                       \le\frac{2\epsilon_\gamma}{c}.
 \label{f16v43:log-ladder}
\end{equation}
\end{theorem}
\begin{proof}
Let $\nu_{a,j}(z)$ be the ordered eigenvalues of $F_a(z)$. They satisfy
$|\nu_{a,j}(z)-\nu_j(z)|\le\epsilon_\gamma$. Each function
$z-\nu_{a,j}(z)$ is continuous and increases by at least the increase in $z$.
At $r_j-\epsilon_\gamma$ it is nonpositive, and at
$r_j+\epsilon_\gamma$ it is nonnegative. Thus its unique root lies in that
interval. Exact inertia identifies that root with the $j$th physical sector
eigenvalue, including possible multiplicities. One can make the interval
strictly larger and then decrease it to include an equality at an endpoint.
This avoids assuming simple or differentiable eigenbranches.

The two triangle inequalities give \eqref{f16v43:level-spacing} and the
cross-sector bound. For the top-root existence, $F_0(c)\succeq F_0(\Lambda)$ and
$\nu_{0,1}(\Lambda)=\Lambda$. Thus
$\nu_1(c)\ge\Lambda-\epsilon_\gamma>c$ eventually, since
$\Lambda\to c+\delta$. The already proved root comparison then puts
$r_1$ a fixed distance above $c$. All arguments of logarithms in
\eqref{f16v43:log-ladder} are at least $c$. The mean value bound for $\log$
and two applications of \eqref{f16v43:ladder-bound} prove it.
\end{proof}

The theorem applies to the second \emph{neutral} eigenvalue when that level is
in the declared interval. It is not limited to the seven charged tops already
controlled in V42. For example, if $r_2>c+\epsilon_\gamma$, the additive
neutral deficit satisfies
\[
 |(\Lambda-\lambda_{0,2})-(r_1-r_2)|\le2\epsilon_\gamma.
\]
If a regional counting calculation instead gives
$N_{\rm reg}(z-\epsilon_\gamma)\le1$ with $c+\epsilon_\gamma\le z<\Lambda$,
then $\lambda_{0,2}\le z$, so $\Lambda-\lambda_{0,2}\ge\Lambda-z$.
These are proved implication rules, not an evaluation of either regional
count or the difference $r_1-r_2$ for the Wilson operator.

\subsection{The neutral excitation uses the returning-state norm}

A frozen spectral pencil still has an exact same-top meaning if its metric is
retained. Work now only in $\mathcal H_0$, the periodic-gauge-invariant,
center-neutral sector. All complementary inverses below are restricted to
$Q\mathcal H_0$. Write
\[
 D=\Lambda-QTQ,\quad B_0=QTW_0,\quad
 S=\Lambda I-F_0(\Lambda),\quad
 M=I-F_0'(\Lambda)=I+B_0^*D^{-2}B_0.
\]
The label $0$ on $F_0,W_0$ means the neutral character, not a Gaussian model.
Put
\begin{equation}
 J=W_0+D^{-1}B_0,\qquad
 u_0=W_0^*P\Omega,\qquad e_0=M^{1/2}u_0,
 \qquad \mathcal K_{\rm reg}=M^{-1/2}SM^{-1/2}.
 \label{f16v43:neutral-data}
\end{equation}
Then $J^*J=M$, $Ju_0=\Omega$, and $\|e_0\|=1$. Define
\[
 g_\gamma=\inf_{v\perp e_0,\,\|v\|=1}
                    \langle v,\mathcal K_{\rm reg}v\rangle,
 \qquad d_\gamma=\Lambda-\|QTQ|_{Q\mathcal H_0}\|\ge3\delta,
\]
with $d_\gamma=\Lambda$ when the complementary restriction is zero.
At each fixed coupling the compactness and the simple native Perron kernel
imply $g_\gamma>0$. This fixed-regulator positivity is not a useful uniform
lower bound. Let
$\Delta_\gamma^{\rm neu}=\Lambda-\lambda_{0,2}$ be the actual neutral
additive excitation gap.

\begin{theorem}[A paid same-top comparison for the complete neutral gap]
\label{f16v43:neutral-gap}
The quantities above obey
\begin{equation}
 \boxed{\qquad
 \frac{d_\gamma g_\gamma}{d_\gamma+g_\gamma}
       \le\Delta_\gamma^{\rm neu}\le g_\gamma,
 \qquad I\preceq M\preceq(1+K)I.\qquad}
 \label{f16v43:neutral-gap-bound}
\end{equation}
In particular, along any fixed-$L$ sequence on which $g_\gamma\to0$,
\begin{equation}
 \frac{\Delta_\gamma^{\rm neu}}{g_\gamma}\longrightarrow1,\qquad
 -\log(\lambda_{0,2}/\Lambda)
                      =\frac{g_\gamma}{\Lambda}(1+o(1)).
 \label{f16v43:neutral-relative}
\end{equation}
No hypothesis that $g_\gamma\to0$, or a formula for its rate, is proved here.
The theorem holds on all neutral form vectors, not just a selected source bank.
\end{theorem}
\begin{proof}
Block completion at the \emph{native} top gives, for every $u\in\mathfrak h$
and $w\in Q\mathcal H_0$,
\begin{equation}
 \langle Ju+w,(\Lambda-T)(Ju+w)\rangle
                       =\langle u,Su\rangle+\langle w,Dw\rangle.
 \label{f16v43:completion}
\end{equation}
The graph range of $J$ need not be orthogonal to $Q\mathcal H_0$.
Nevertheless its distance from the actual vacuum satisfies
\[
 \operatorname{dist}(Ju,\mathbb C\Omega)^2
   =\operatorname{dist}(M^{1/2}u,\mathbb C e_0)^2
   \le g_\gamma^{-1}\langle u,Su\rangle.
\]
Every neutral $\psi$ has a unique decomposition $\psi=Ju+w$ by taking
$u=W_0^*P\psi$. The triangle inequality for the distance to the vacuum,
followed by scalar Cauchy--Schwarz, yields
\[
 \operatorname{dist}(\psi,\mathbb C\Omega)^2
 \le\left(\sqrt{\langle u,Su\rangle/g_\gamma}
          +\sqrt{\langle w,Dw\rangle/d_\gamma}\right)^2
 \le(g_\gamma^{-1}+d_\gamma^{-1})
               \langle\psi,(\Lambda-T)\psi\rangle.
\]
This proves the lower bound. It does not assert the false full-vector bound
$\Lambda-T\succeq d_\gamma Q$.

For the upper bound take a unit $v\perp e_0$ approaching the infimum defining
$g_\gamma$, and put $\psi=JM^{-1/2}v$. It is a unit neutral vector orthogonal
to $\Omega$, and its Rayleigh energy is exactly
$\langle v,\mathcal K_{\rm reg}v\rangle$. The variational principle gives
$\Delta_\gamma^{\rm neu}\le g_\gamma$. The metric bounds follow from
\eqref{f16v43:monotonicity}. Finally $M-I$ and $F_0(\Lambda)$ are compact,
so $\mathcal K_{\rm reg}$ is a compact perturbation of $\Lambda I$.
Its kernel corresponds exactly to the one-dimensional native kernel by
\eqref{f16v43:completion}; zero is isolated. This also justifies the claimed
fixed-coupling positivity of $g_\gamma$. Since $d_\gamma\ge3\delta$ and
$\Lambda\to\lambda_*>0$, the two limits in
\eqref{f16v43:neutral-relative} follow from the bracket and the scalar logarithm.
\end{proof}

\paragraph{Why the metric is not optional.}
Consider the finite reference form, unrelated to native matrix elements,
\[
 L_\tau=\begin{pmatrix}
 0&0&0\\ 0&\tau+b^2/d&-b\\ 0&-b&d
 \end{pmatrix},\qquad \tau,d>0.
\]
Retain the first two coordinates. The Schur form is
$S=\operatorname{diag}(0,\tau)$, but its returning norm is
$M=\operatorname{diag}(1,1+b^2/d^2)$. Its first positive eigenvalue satisfies
\[
 \lambda_1(L_\tau)=\frac{\tau}{1+b^2/d^2}+O_{b,d}(\tau^2).
\]
The metric gap has exactly that leading coefficient. Using the bare gap
$\tau$ would miss an arbitrary fixed factor. This is consistent with the
f15 V78 counterfixture; the present proof pays the graph norm and uses the
full completed form rather than assuming an unpaid vertical inequality.

\subsection{Neutral observables retain their complementary source and vacuum subtraction}

Let $f\in\mathcal H_0$ satisfy $\langle\Omega,f\rangle=0$. For example
$f=(F-\mu_\gamma F)\Omega$ for a bounded gauge- and center-invariant
multiplication $F$, with $d\mu_\gamma=\Omega^2dH$. For real $z>\Lambda$ put
\[
 b_f(z)=W_0^*Pf+W_0^*TQ R(z)Qf.
\]
The symbol $b_f$ denotes an effective source, not the magnetic multiplier.

\begin{proposition}[Exact source resolvent on the single-region pencil]
\label{f16v43:source}
For every real $z>\Lambda$,
\begin{equation}
 \boxed{\quad
 \langle f,(z-T)^{-1}f\rangle
 =\langle Qf,R(z)Qf\rangle+
 \langle b_f(z),(zI-F_0(z))^{-1}b_f(z)\rangle.
 \quad}
 \label{f16v43:source-resolvent}
\end{equation}
At the native top $b_f(\Lambda)=J^*f$ is orthogonal to $u_0$. With the
Moore--Penrose inverse on the corresponding vacuum complement,
\begin{equation}
 \langle f,(\Lambda-T)^\dagger f\rangle
 =\langle Qf,D^{-1}Qf\rangle+
 \langle M^{-1/2}b_f(\Lambda),
       \mathcal K_{\rm reg}^{\dagger}M^{-1/2}b_f(\Lambda)\rangle.
 \label{f16v43:source-top}
\end{equation}
All inverses in this identity exist on their declared complements at each
fixed finite coupling. No bound uniform in coupling on
$\mathcal K_{\rm reg}^{\dagger}$ is inferred.
\end{proposition}
\begin{proof}
Solve the complementary block equation of $(z-T)\psi=f$, substitute it into
the retained equation, and expand $\langle f,\psi\rangle$. The two cross
terms assemble into $b_f(z)$ and give \eqref{f16v43:source-resolvent}; keeping
only $W_0^*Pf$ would discard a source return.
At $z=\Lambda$ the exact relation $Ju_0=\Omega$ gives
$\langle u_0,b_f(\Lambda)\rangle=\langle\Omega,f\rangle=0$.
Thus $Su=b_f(\Lambda)$ is solvable, with ambiguity only in $\mathbb C u_0$.
Substitution into the completed equations yields
$\langle Qf,D^{-1}Qf\rangle+\langle b_f,S^\dagger b_f\rangle$;
the kernel ambiguity contributes zero. Changing variable by $M^{1/2}$ in
this variational inverse gives \eqref{f16v43:source-top}. The transformed
source is orthogonal to $e_0$. This argument proves the quadratic inverse
identity; it does not use an invalid general congruence formula for
Moore--Penrose matrices. No physical time or new stationary probability is
introduced.
\end{proof}

\subsection{The next spectral estimate is now confined to one complete region}

The exponential estimate is no longer limited to seven character tops. It
matches every regional root separated from the lower cutoff to one native
eigenvalue in each sector, with multiplicity and a uniform error. An isolated
regional cluster of $k$ roots with a gap larger than that error would therefore
produce $8k$ physical eigenvalues. Whether $k=1$, or whether a second neutral
level approaches the top, is \emph{not} decided by this statement. Center
symmetry by itself gives neither answer.

The return-metric formula identifies a precise alternative target: bound or
asymptotically identify $g_\gamma$ on $\mathfrak h$, including all its
nonquadratic coordinates and its actual vacuum. The forbidden complement
already has the positive constant $d_\gamma\ge3\delta$, so if the regional
metric gap is small, its relative transfer to the native neutral gap loses
only $O(g_\gamma/d_\gamma)$. A proposed quartic or other local operator must
still be compared with this pencil and metric; it is not defined as the answer
by matching a classical Taylor polynomial.

Nothing here improves the common shift $\Lambda_\gamma-\lambda_*$. All bounds
fix $L$ and the separating regions before the coupling limit. At an independently
specified physical step $a_t$, the neutral energy is still
$-a_t^{-1}\log(\lambda_{0,2}/\Lambda)$. The supplied memoranda's exact-return
proposal and separate physical-transfer requirement remain distinct from the
unproved local soft spectrum, regulator removal, and continuum reconstruction.
No new numerical neutral-gap bound or continuum mass is claimed.

\subsection{Verification and preservation}

The diagnostic is
\begin{center}\small\path{verify_ym_f16_v43_sector_pencil.py}.\end{center}
It checks character isometries, a positive finite transfer with two retained
states per region, differentiated returns, rational inertia and root brackets,
the reconstruction metric, and source inverses. The finite reference is not
a discretization or enclosure of the native Wilson spectrum. Proofs and actual
execution records are reported separately.

Only the literal title version and this terminal insertion change V42.
Removing them recovers its bytes exactly. Available predecessor checkers are
retained unchanged; the audit states which were actually rerun. Native
conclusions use V42's proved-as-stated complement bounds as inherited premises;
the full earlier analytic lineage is not independently recertified. The next
scheduled companion checkpoint remains V45.

\begin{firewall}
V43 proves complete character-pencil, spectral-count and level comparisons,
and a same-top neutral-gap bound in the exact return metric. It retains
complementary sources. It does not evaluate the regional soft spectrum,
isolate exactly eight states, determine a neutral scaling law, remove the
spatial regulator, calibrate physical time, or construct a continuum mass gap.
The native operator and its physical spectral top remain unchanged.
\end{firewall}
% END F16 V43 APPEND

% BEGIN F16 V44 APPEND -- 2026-09-18
% Insert before the final document-ending command of cumulative V43.
% The predecessor, original transfer, endpoint laws and clocks are unchanged.

\clearpage
\section{V44: physical soft trial spaces and a confining quartic comparison}
\label{f16v44:context}

V43 leaves open whether its second regional level approaches the top, and its
relative neutral-gap comparison is conditional on $g_\gamma\to0$. We first
settle that qualitative question without assigning a quartic Hamiltonian to the
native operator. The lower trial in V41 proves only the largest eigenvalue:
an arbitrary fine-link trial is not a physical excited-state test. Here a
based-gauge tube, with its actual volume density, produces arbitrarily large
\emph{gauge-invariant} soft trial spaces in one central region. Their center
translates give such a space in every character sector. Every fixed level
consequently approaches the same native limiting top, including every fixed
neutral level. This also proves the previously unverified premise
$g_\gamma\to0$ in V43.

A separate calculation records the constant-field quartic model with the
original kinetic coefficient. Its commuting valleys do not destroy quantum
compactness: an explicit transverse-oscillator inequality proves compact
resolvent, also on the residual conjugation-invariant space. This is a theorem
about that declared comparison operator, not its identification with the native
regional pencil. The two conclusions have different logical status and are
kept separate throughout.

\subsection{Scope, source review, and the unchanged fixed-regulator problem}

The supplied V43 appendix and analytic audit were read completely. The native
inputs used below are V41's exact fine-link transfer and fixed-$L$ top limit,
together with the elementary normalized Wilson displacement estimates used in
V42. The regional consequences additionally use V43's proved-as-stated ladder
and return-metric comparison. No endpoint asymptotic or auxiliary sampling gap
is required. The f14 V54 local based-gauge chart and constant commutator formula
are precedents; their four-dimensional static integral is not substituted for
this three-dimensional slice transfer. The f15 V98 compact kinetic result is
likewise a different operator. Exact archive searches and those identified
passages were reviewed, not the entire archives. The next scheduled checkpoint
remains V45.

For literature context, L\"uscher's
\href{https://doi.org/10.1016/0550-3213(83)90436-4}{\emph{Some analytic results
concerning the mass spectrum of Yang--Mills gauge theories on a torus}}
addresses a small-volume low-energy expansion. Simon's
\href{https://doi.org/10.1016/0003-4916(83)90057-X}{\emph{Some quantum operators
with discrete spectrum but classically continuous spectrum}} is a classical
precedent for confinement by transverse quantum motion. Only their primary
abstracts and bibliographic scope were checked in this iteration. No expansion,
constant, or spectral theorem from either paper is imported. The elementary
trial and oscillator arguments needed here are written out. Neither the
finite-dimensional quartic model nor the general min--max principle is claimed
as a new general construction.

Fix $L=2m$, $m\ge4$, put $V=L^3=8B$, $E=3V$, and retain
\[
 T_\gamma=M_{b_\gamma}C_\gamma^{\otimes E}M_{b_\gamma},\qquad
 b_\gamma=e^{-\gamma S/2},\qquad
 S(U)=\tfrac12|\mathcal W(U)|^2,\quad
 \mathcal W(U)=(U_p-I)_p.
\]
The reference measure is probability Haar on $\mathcal X_L=\SU(2)^E$.
Let $\mathcal H_a$ be V42's physical center-character spaces, and enumerate
$T_\gamma|_{\mathcal H_a}$ in decreasing order, starting at one:
$\lambda_{a,1}\ge\lambda_{a,2}\ge\cdots$. In the neutral sector
$\lambda_{0,1}=\Lambda_\gamma$. Write
\[
 \varepsilon=\gamma^{-1/2},\qquad \lambda_*:=\Lambda_{{\rm G},B}>0.
\]
V41 supplies $\Lambda_\gamma\to\lambda_*$. All limits below fix $L$ first;
no bound uniform over growing lattice sizes or a physical time scale is given.

\subsection{A physical central tube with no division by a stabilizer}

Use real color-valued lattice cochains at the identity configuration. In the
orthonormal fine-link metric let $d_0$ be the vertex-to-link gradient and $d_1$
the link-to-plaquette curl. Set
\[
 \mathcal N=\ker d_0^*,\qquad
 \mathcal Z=\ker d_1\cap\mathcal N,\qquad
 \mathcal Y=\mathcal N\ominus\mathcal Z.
\]
The inherited periodic cochain calculation gives
\begin{equation}
 \dim\mathcal N=6V+3,\qquad \dim\mathcal Z=9,\qquad
 r:=\dim\mathcal Y=6(V-1),\qquad H_+=d_1^*d_1|_{\mathcal Y}\succ0.
 \label{f16v44:dimensions}
\end{equation}
The nine coordinates in $\mathcal Z$ are the constant three-direction,
three-color fields. They are not nine gauge-invariant coordinates individually.
The residual simultaneous adjoint action acts on them, and will be imposed on
functions rather than divided out as a translation.

\begin{proposition}[An exact based-gauge tube and its invariant half-density]
\label{f16v44:tube}
Let $\mathcal G_o=\{g\in\SU(2)^V:g_o=I\}$. For some fixed $r_0>0$ the map
\begin{equation}
 \Phi:\mathcal G_o\times\{s\in\mathcal N:|s|<r_0\}\longrightarrow\mathcal X_L,
 \qquad \Phi(g,s)=g\cdot\operatorname{Exp}(s)
 \label{f16v44:tube-map}
\end{equation}
is a diffeomorphism onto a tube about the trivial central-flat gauge orbit.
It can be chosen inside V42's selected trivial central region. Its exact
measure is
\begin{equation}
 dH(\Phi(g,s))=dH_{\mathcal G_o}(g)\,J_L(s)\,ds,
 \qquad J_L>0\ \hbox{smooth}.
 \label{f16v44:tube-measure}
\end{equation}
The density $J_L$ is invariant under simultaneous adjoint rotation of $s$.
Every compactly supported invariant $\phi\in L^2(\mathcal N,ds)$ in this ball
has an isometric physical extension
\begin{equation}
 (\mathcal E\phi)(\Phi(g,s))=J_L(s)^{-1/2}\phi(s),
 \qquad \mathcal E\phi=0\ \hbox{off the tube}.
 \label{f16v44:embedding}
\end{equation}
No value of $J_L$ is replaced by one in this definition.
\end{proposition}
\begin{proof}
The based group acts freely on a connected lattice: a transformation fixing
all links and one root is identity successively along a spanning tree. Its
orbit at the identity is a compact embedded submanifold; its tangent is
$\operatorname{Ran}d_0$. The action is isometric for the product round metric,
and its normal space at the identity is $\mathcal N$. The normal bundle along
this orbit is trivialized by $Dg\mathcal N$. The normal exponential of this
bundle is exactly \eqref{f16v44:tube-map}, because the product exponential at
the identity is the componentwise group exponential and the action is
isometric. The local inverse theorem gives invertibility near every zero
section point. Compactness of the orbit gives a common small radius and
injectivity: otherwise two distinct normal representatives tending to the
zero section would contradict a local inverse at their common limiting orbit
point. Shrink the radius to lie in the prescribed open region.

The pulled-back positive smooth volume has a positive density. Invariance
under the left action of the based group makes that density independent of
$g$ when its reference is probability Haar, proving
\eqref{f16v44:tube-measure}. Global conjugation fixes the identity, preserves
both cochain subspaces, and sends $(g,s)$ to $(kgk^{-1},\operatorname{Ad}_k s)$.
It preserves $ds$ and based Haar, proving the stated symmetry of $J_L$.
Any full gauge transformation is a based transformation followed by a global
one. Thus \eqref{f16v44:embedding} is invariant under the complete periodic
gauge group and is isometric by direct integration. The global stabilizer
has not been fixed by a singular orientation slice. At a central orbit the
based action is still free, which is the only freeness used.
\end{proof}

Let $\pi(U)=s$ be the tube coordinate modulo the based action. At the identity,
$D\pi$ is the orthogonal projection $\Pi_{\mathcal N}$, not the identity on
all link variables. For bounded $z$ and small $s,\varepsilon$, smoothness gives
\begin{equation}
 \pi(\operatorname{Exp}(s)\operatorname{Exp}(\varepsilon z))
 =s+\varepsilon\Pi_{\mathcal N}z+
 O_L(\varepsilon |s||z|+\varepsilon^2|z|^2).
 \label{f16v44:coordinate-step}
\end{equation}
The product and exponential are componentwise. This formula is local near the
identity representative; the associated gauge-invariant integral below makes
it sufficient on the whole orbit tube.

\subsection{The native transfer on arbitrarily many invariant soft tests}

Write $s=(y,z)\in\mathcal Y\oplus\mathcal Z$. The normalized hard Gaussian
transfer on $L^2(\mathcal Y,dy)$ is
\begin{equation}
 (\mathcal T_+h)(u)=(2\pi)^{-r/2}
 \int e^{-|u-u'|^2/2-u^{\mathsf T}H_+u/4-u'^{\mathsf T}H_+u'/4}
                         h(u')\,du'.
 \label{f16v44:hard-transfer}
\end{equation}
V41's matching of fine and rooted kinetic metrics gives
$\|\mathcal T_+\|=\lambda_*$. Its unique normalized positive Gaussian ground
function is invariant under global adjoint rotations. Smooth invariant compact
cutoffs approximate it in norm and in its bounded transfer quotient.

Fix any $\beta$ with $1/2<\beta<1$, put $t_\varepsilon=\varepsilon^\beta$, and
for smooth compactly supported invariant $h$ and $g$ define
\begin{equation}
 \phi_{\varepsilon,h,g}(y,z)
   =\varepsilon^{-r/2}t_\varepsilon^{-9/2}
                 h(y/\varepsilon)g(z/t_\varepsilon),\qquad
 \Psi_{\varepsilon,h,g}=\mathcal E\phi_{\varepsilon,h,g}.
 \label{f16v44:trials}
\end{equation}
For sufficiently small $\varepsilon$ their entire support is in the tube.
Their norms and inner products are exactly the product norms, before any limit.
There are arbitrarily large orthonormal families of such invariant soft $g$'s:
choose normalized radial functions on disjoint annuli of $\mathbb R^9$.

\begin{theorem}[A physical hard--soft trial limit for the original kernel]
\label{f16v44:trial-limit}
For every fixed pair of tests as above,
\begin{equation}
 \boxed{\quad
 \lim_{\gamma\to\infty}
 \langle\Psi_{\varepsilon,h,g},T_\gamma\Psi_{\varepsilon,k,f}\rangle
       =\langle h,\mathcal T_+k\rangle\langle g,f\rangle.
 \quad}
 \label{f16v44:trial-bilinear}
\end{equation}
The convergence is uniform over unit vectors in any fixed finite-dimensional
span of the chosen tests. This is a variational bilinear limit, not a claim of
norm-resolvent convergence on the full changing state space.
\end{theorem}
\begin{proof}
At $s=(\varepsilon u,t_\varepsilon v)$ on a fixed rescaled support, the
constraint map has the expansion
\begin{equation}
 \varepsilon^{-1}\mathcal W(\operatorname{Exp}s)
 =d_1u+O_L\bigl(\varepsilon+t_\varepsilon+t_\varepsilon^2/\varepsilon\bigr)
                     \longrightarrow d_1u.
 \label{f16v44:magnetic-limit}
\end{equation}
Here $d_1v=0$; Taylor's remainder is bounded on the original finite-dimensional
compact neighborhood. In particular the true multiplier tends to
$e^{-u^{\mathsf T}H_+u/4}$. No nonlinear magnetic word is deleted from the
native trial quotient.

For clarity, evaluate the inner kinetic integral as an expectation of the
actual product Wilson displacement:
$U'=U\operatorname{Exp}(\varepsilon Z)$, with its full principal product and
its original normalized density. This density converges on bounded sets to
the standard Gaussian on the full $3E$-dimensional fine-link tangent. Its
probability outside $|Z|>R$ is at most $A_Le^{-b_LR^2}$, uniformly for small
$\varepsilon$, by V42's elementary displacement estimate. The same bound is a
row and column Schur bound for the kernel truncated to these displacements.
Since both magnetic factors are at most one, the discarded bilinear form is
bounded by $A_Le^{-b_LR^2}$ times the exact test norms. This tail estimate is
independent of their shrinking supports.

On bounded $Z$, use \eqref{f16v44:coordinate-step}. Its hard coordinate after
rescaling tends to $u+Z_+$; its soft coordinate tends to $v$ because
$\varepsilon/t_\varepsilon\to0$. The coordinate errors in these two rescalings
also tend to zero. Both input and output stay in the common tube. When the
outer Haar integral is disintegrated by \eqref{f16v44:tube-measure}, the two
half-densities leave the ratio
$\sqrt{J_L(s)/J_L(s')}\to1$ in the displacement integral. Their constant
normalization is therefore accounted for, rather than omitted. The exact
powers $\varepsilon^r t_\varepsilon^9$ from $ds$ cancel the two trial
normalizations.

Orthogonal projection of the full standard Gaussian onto
$\mathcal Y\oplus\mathcal Z$ has identity covariance on that subspace. Its
vertical part and its soft displacement each integrate to one. The hard
part, with the two limiting magnetic multipliers, gives precisely
\eqref{f16v44:hard-transfer}. Dominated convergence on the bounded rescaled
supports and bounded displacements proves the claimed limit with cutoff $R$.
Then send $R\to\infty$ using the paid Schur tail. Both requirements on
$\beta$ were used: $t_\varepsilon^2/\varepsilon\to0$ for magnetism and
$\varepsilon/t_\varepsilon\to0$ for the soft displacement. Entrywise convergence
of a fixed finite matrix gives its operator-norm convergence and the final
assertion. No uniformity in the number of test functions is claimed.
\end{proof}

\subsection{Every fixed physical level approaches the top}

\begin{theorem}[An unbounded number of near-top levels in every center sector]
\label{f16v44:all-levels}
For each fixed $L$ and every fixed integer $j\ge1$, the original physical
sector eigenvalues satisfy
\begin{equation}
 \boxed{\qquad
 \lim_{\gamma\to\infty}\lambda_{a,j}(\gamma)=\lambda_*
              \quad\hbox{for every }a\in\{0,1\}^3.
 \qquad}
 \label{f16v44:sector-limits}
\end{equation}
They exist and are positive for all sufficiently large coupling. In particular,
for $j\ge2$ in the neutral sector,
\begin{equation}
 \Lambda_\gamma-\lambda_{0,j}\longrightarrow0,
 \qquad -\log(\lambda_{0,j}/\Lambda_\gamma)\longrightarrow0.
 \label{f16v44:neutral-limits}
\end{equation}
For every fixed $0<q<1$, the number of neutral eigenvalues strictly above
$q\Lambda_\gamma$ tends to infinity. This gives no rate or counting asymptotic.
\end{theorem}
\begin{proof}
Fix $N$ and $\eta>0$. Choose a normalized invariant hard cutoff $h$ with
$\langle h,\mathcal T_+h\rangle>\lambda_*-\eta$. Choose $N$ orthonormal
invariant soft tests $g_1,\ldots,g_N$. Their physical extensions
$\Psi_1,\ldots,\Psi_N$ are exactly orthonormal in the one-region tube. The
preceding theorem says that the matrix of $T_\gamma$ on their span tends to
$\langle h,\mathcal T_+h\rangle I_N$.

For a center character use the exact folded tests
\[
 \Psi_{a,i}=8^{-1/2}\sum_{\sigma\in\{\pm1\}^3}
                    \chi_a(\sigma)\mathsf Z_\sigma\Psi_i.
\]
Their eight translated supports lie in disjoint fixed central regions, so they
are orthonormal in $i$ and lie in $\mathcal H_a$. Each diagonal translated
matrix element is the same as in the trivial tube. Different regions have a
fixed positive distance; the actual separated-displacement Schur estimate makes
each cross-region matrix element $O_L(e^{-c_L\gamma})$. There are finitely many
of these for the chosen $N$. Thus the infimum of the Rayleigh quotient on this
$N$-dimensional physical sector space is at least $\lambda_*-2\eta$ eventually.

The compact self-adjoint min--max principle gives
$\liminf\lambda_{a,N}\ge\lambda_*-2\eta$. Its upper bound is
$\lambda_{a,N}\le\Lambda_\gamma\to\lambda_*$. Send $\eta\downarrow0$, with
$N$ fixed, to obtain \eqref{f16v44:sector-limits}. Positivity follows from the
positive lower bound eventually. For each prescribed count $N$, these $N$
levels exceed $q\Lambda_\gamma$ eventually; this is exactly the asserted
divergence of the count. No non-invariant trial is used for an excited level,
and no physical excited level is identified merely from equality of full and
physical Perron tops.
\end{proof}

The trial scale can, for example, be $t_\varepsilon=\varepsilon^{3/4}$.
It is not claimed to be a stationary fluctuation scale. The limits fix $j,L$
before coupling; they do not assert an estimate uniform in a growing index.
At each finite coupling compactness still makes every count above a positive
threshold finite.

\begin{proposition}[Consequences for V43's actual regional metric]
\label{f16v44:regional}
Every fixed ordered eigenvalue of V43's one-region compression tends to
$\lambda_*$. Consequently all its fixed roots $r_j$ eventually lie in V43's
matching interval. Its exact metric gap obeys
\begin{equation}
 \boxed{\qquad
 g_\gamma\longrightarrow0,\qquad
 \frac{\Delta_\gamma^{\rm neu}}{g_\gamma}\longrightarrow1,\qquad
 -\log(\lambda_{0,2}/\Lambda_\gamma)
       =\frac{g_\gamma}{\Lambda_\gamma}(1+o(1)).
 \qquad}
 \label{f16v44:metric-limit}
\end{equation}
No fixed-rank neutral truncation can leave a coupling-independent additive
complement gap at the native top.
\end{proposition}
\begin{proof}
The unfolded $N$ trial functions are supported in the chosen region and hence
in the space of $T^{\rm one}$. Their Rayleigh matrix is unchanged by that
compression. The same lower argument, and its upper norm bound
$\|T^{\rm one}\|\le\Lambda_\gamma$, prove the root limits. They are eventually
above the fixed V43 cutoff by a fixed positive margin.

The preceding theorem gives $\Delta_\gamma^{\rm neu}\to0$. V43's bracket,
with $d_\gamma\ge3\delta$, gives, once $\Delta_\gamma^{\rm neu}<d_\gamma$,
\[
 \Delta_\gamma^{\rm neu}\le g_\gamma
 \le\frac{d_\gamma\Delta_\gamma^{\rm neu}}
                {d_\gamma-\Delta_\gamma^{\rm neu}}.
\]
This proves all of \eqref{f16v44:metric-limit}, including the formerly unpaid
premise $g_\gamma\to0$. It retains the metric $I-F_0'(\Lambda_\gamma)$ and the
actual vacuum; no bare regional norm is substituted.

For any orthogonal neutral projection $R_\gamma$ of rank at most a fixed $N$,
min--max gives
$\|(I-R_\gamma)T_\gamma(I-R_\gamma)|_{\mathcal H_0}\|
\ge\lambda_{0,N+1}$. Its distance below $\Lambda_\gamma$ tends to zero. This is
why an isolated eight-state, or any fixed finite-state, global cluster cannot
have a fixed positive complement gap in this limit. The spatially supported
forbidden complement of V42 is different and its proved gap is not contradicted.
An eight-state cluster in a \emph{shrinking} spectral interval is not excluded;
no comparison of the neutral spacing with the charged exponential splitting
has been established here.
\end{proof}

\subsection{The quartic constant-field comparison: its coefficient and its limits}

The qualitative native theorem does not determine a scaling law. The following
separate calculation identifies a concrete nonquadratic candidate and proves
its own analytic properties, without asserting a native spectral approximation.
For constant fine links $U_{x,\mu}=\operatorname{Exp}(a_\mu)$, $a_\mu\in\mathbb R^3$,
write $v(a)=(\sin|a|/|a|)a$. The exact quaternion commutator identity gives
\begin{equation}
 S(\operatorname{Exp}a)=2V\sum_{\mu<\nu}|v(a_\mu)\times v(a_\nu)|^2
              =2V\mathcal Q(a)+O_L(|a|^6),\qquad
 \mathcal Q(a)=\sum_{\mu<\nu}|a_\mu\times a_\nu|^2.
 \label{f16v44:constant-action}
\end{equation}
The identity and its f14 V54 analogue are not new general quaternion formulas.
Here there are three spatial directions, nine coordinates and $V=L^3$ copies
of each direction pair, not the twelve-coordinate four-torus static integral.

The tangent metric on constant fields is $V\sum_\mu|da_\mu|^2$.
The normalized kinetic displacement has covariance $\gamma^{-1}I$ per fine
link at leading order. Its constant tangent component therefore has generator
$(2\gamma V)^{-1}\Delta_a$. Combining this principal term with the classical
quartic defines the following \emph{comparison}, not an exact restriction:
\begin{equation}
 H_{\rm cmp}(\gamma,V)
 =-\frac1{2\gamma V}\Delta_a+2\gamma V\mathcal Q(a).
 \label{f16v44:comparison-raw}
\end{equation}
The constant configurations are Haar null and are not preserved by the native
kinetic operator. Thus \eqref{f16v44:comparison-raw} is not obtained by conditioning
on them, freezing their complement, or equating a classical Hessian to V43's
spectral pencil.

Under the exact unitary dilation $a=(\gamma V)^{-1/3}q$ this comparison becomes
\begin{equation}
 H_{\rm cmp}(\gamma,V)\simeq(\gamma V)^{-1/3}H_{\rm mat},\qquad
 H_{\rm mat}=-\tfrac12\Delta_q+2\mathcal Q(q)
                \quad\hbox{on }L^2(\mathbb R^9,dq).
 \label{f16v44:model-scaling}
\end{equation}
The residual Gauss condition is simultaneous rotation of all three $q_\mu$.
It is imposed by restriction to invariant functions, not by deleting three
coordinates with an uncomputed Jacobian.

\begin{theorem}[Quantum confinement of the declared quartic model]
\label{f16v44:model-confinement}
The Friedrichs realization of $H_{\rm mat}$ has compact resolvent, a simple
strictly positive ground function and a strictly positive ground energy. In
quadratic-form sense,
\begin{equation}
 \boxed{\quad
 H_{\rm mat}\succeq2\sum_{\mu=1}^3|q_\mu|,
 \qquad
 H_{\rm mat}\succeq-\tfrac14\Delta_q+\sum_{\mu=1}^3|q_\mu|.
 \quad}
 \label{f16v44:oscillator-floor}
\end{equation}
Its simultaneous-rotation-invariant restriction has compact resolvent and
infinitely many eigenvalues $e_0<e_1\le e_2\le\cdots$, with $e_1-e_0>0$.
These are model eigenvalues, not evaluated native levels or numerical enclosures.
\end{theorem}
\begin{proof}
For each $i\in\{1,2,3\}$ define on smooth compact tests
\[
 A_i=-\tfrac14\sum_{j\ne i}\Delta_{q_j}
                        +\sum_{j\ne i}|q_i\times q_j|^2.
\]
Their sum is exactly $H_{\rm mat}$: each derivative occurs twice and each
unordered potential pair occurs twice. Hold $q_i$ fixed and put $r=|q_i|$.
For each $j\ne i$, an orthogonal rotation in the $q_j$ variable splits one
nonnegative free derivative and the two-dimensional oscillator
$-\tfrac14\Delta_y+r^2|y|^2$. Its lower form bound is $r$, including $r=0$.
For example expand the two nonnegative squares
$\|\tfrac12\partial_{y_k}f+r y_k f\|^2$ and integrate their cross terms.
Thus $A_i\succeq2|q_i|$. No derivative in $q_i$ is taken in this disintegration,
so the parameter-dependent rotation adds no derivative term. Sum the three
bounds to obtain the first inequality. Average it with
$H_{\rm mat}\succeq-\tfrac12\Delta_q$ to obtain the second.

The closed form therefore controls a local $H^1$ norm and
$\int |q|\,|f(q)|^2dq$. A bounded form sequence has uniformly small $L^2$ mass
outside large balls, and Rellich compactness gives convergence on each ball.
A diagonal extraction and the tail bound make the full form embedding compact;
this is equivalent to compact resolvent. Zero energy would force a constant
$L^2(\mathbb R^9)$ function and hence the zero vector, so the attained ground
energy is positive. The heat kernel is strictly positive: along any finite
Brownian bridge the continuous polynomial potential has finite integral, so
its Feynman--Kac weight is strictly positive. Positivity improvement makes the
ground eigenvalue simple with a positive eigenfunction. Rotational invariance
and uniqueness force that function to be invariant. Restriction to the closed
invariant subspace preserves compact resolvent. That subspace is infinite
dimensional (it contains all radial smooth compact tests), so its eigenvalues
form an infinite discrete sequence; simplicity of the ground gives its positive
first invariant gap. This proves properties of the specified model only.
\end{proof}

There is a coefficient check against V41's cusp, not a second potential to add.
At a commuting rank-one triple $a_\mu=c_\mu n$, $|n|=1$, the Hessian of
$2\mathcal Q$ has four positive eigenvalues $4|c|^2$. The transverse harmonic
ground contribution for $-\frac12\Delta+2\mathcal Q$ is therefore $4|c|$.
V41's flat-threshold variation has precisely this leading term when
$\varphi=2L c$:
$2|\varphi|/L=4|c|$. Keeping the full quartic coordinates and also adding that
cusp as a separate potential would count those transverse modes twice.
The model still has classical commuting valleys; the proved form bound shows
why classical flatness alone is not a test for its quantum spectrum.

\subsection{The remaining native-to-model estimate}

The new native result settles the qualitative regional alternative:
there are arbitrarily many near-top physical levels in every sector, and the
neutral return-metric gap tends to zero. It does not determine how rapidly.
The model calculation makes $(\gamma V)^{-1/3}$ a natural \emph{candidate}
scale, but this dilation is not a theorem about the native neutral gap.

For an actual scaling limit one must compare the complete regional metric
operator of V43, after its actual vacuum subtraction, with
$H_{\rm mat}-e_0$ (or a corrected operator if elimination produces leading
terms). This needs an upper trial comparison at the proposed scale and a
matching lower form estimate with compactness of low-energy states. The
nonzero-mode response, the true quotient density, the energy-dependent return,
and its metric must all be retained. In particular a lower bound on a discarded
restriction does not control the norm of a reconstructed full vector. The
forbidden outer region already has its V42 inverse; it is not the missing
within-region hard/soft comparison.

The intermediate trial width used above satisfies
$\varepsilon\ll t_\varepsilon\ll\sqrt\varepsilon$ so its nonlinear magnetic
cost simply tends to zero. At the candidate balanced scale those terms are
\emph{of the same order} as the soft kinetic deficit; the present bilinear
limit cannot be divided by that smaller scale. No big-O remainder at that
scale, tunneling prefactor, rate of the common top shift, volume-uniform bound,
or continuum time is supplied. With any independently specified $a_t(\gamma)$,
the physical neutral energy remains
$-a_t(\gamma)^{-1}\log(\lambda_{0,2}/\Lambda_\gamma)$. A vanishing uncalibrated
one-step deficit does not decide its continuum limit.

\subsection{Verification and preservation}

The new diagnostic is
\begin{center}\small\path{verify_ym_f16_v44_physical_soft_spaces.py}.\end{center}
It checks finite periodic cochains and harmonic dimensions, the exact gauge
normal/tangent decomposition, noncommutative constant plaquettes, character
isometries, hard--soft Gaussian trial limits, the oscillator allocation and
scaling, and the algebra taking V43's bracket to its now established limit.
Finite fixtures do not certify a tubular-neighborhood argument, a compact
Wilson spectral limit, or a native-to-matrix-model approximation. Actual
execution records and scope are reported separately in the audit.

Only the title version and this terminal insertion change V43. Removing those
changes recovers the predecessor byte for byte. The main native conclusion
inherits V41's fixed-regulator top limit; the invariant trial construction is
proved here. The earlier analytic lineage is not independently recertified.
The next scheduled companion checkpoint remains V45.

\begin{firewall}
V44 proves that every fixed physical level, in every center sector, approaches
the same limiting native top at fixed spatial regulator. It establishes
$g_\gamma\to0$ for V43's actual neutral return metric and rules out a fixed-rank
near-top reduction with a uniform complementary gap. Separately it proves
compact resolvent and a positive invariant excitation gap for an explicitly
defined quartic comparison operator. No native scaling law is inferred from
that model, no original transfer is replaced, and no physical time, unrestricted
regulator limit, or interacting continuum mass-gap conclusion is supplied.
\end{firewall}
% END F16 V44 APPEND

% BEGIN F16 V45 APPEND -- 2026-09-18
% Insert immediately before the final document-ending command of V44.
% No predecessor assertion, native measure, scalar, or clock is edited.

\clearpage
\section{V45: critical-scale recovery and corrected physical quasimodes}
\label{f16v45:context}

V44 proves qualitative collapse of every fixed physical level, but its trial
error is only $o(1)$. It cannot be divided by the proposed quartic scale. We
now expand the \emph{actual} transfer on that scale. A term linear in the soft
field changes the hard Gaussian operator. Its ground expectation vanishes by
color symmetry, but the term need not vanish as a vector. A fixed hard-band
inverse corrects that vector defect. This gives physical quasimodes with error
smaller than the quartic scale, not just a favorable Rayleigh quotient.

The result is spectral inclusion, with multiplicities, and a one-sided ordered
level bound. It does not prove that the included levels exhaust the native low
spectrum. In particular it does not identify the native Perron correction or
the first neutral spacing with the two lowest quartic eigenvalues. The missing
lower estimate and low-energy compactness are stated at the V45 checkpoint.

\subsection{Checkpoint, dependencies, and the balanced variables}

The supplied V44 appendix and audit were read in full. We retain its physical
tube, exact density, hard Mehler transfer, and separately proved quartic model.
The f14 V54 constant-field stationarity argument was checked; its different
four-dimensional static integral is not imported. The Grand Tensor's signed
compression and return packet were also checked. Searches of the active source,
f14 V100, f15 V100 and Grand Tensor V076 are recorded separately from proof
verification. No archive-wide novelty claim is made. The next checkpoint is V50.

The correction principle is classical. Panati--Spohn--Teufel,
\href{https://arxiv.org/html/math-ph/0201055}{\emph{Space-Adiabatic Perturbation
Theory}}, separates an isolated fast band from the effective slow dynamics.
M\'atyus--Teufel,
\href{https://arxiv.org/html/1904.00042}{\emph{Effective non-adiabatic
Hamiltonians}}, Sections III--IV, explains why a bare diagonal compression
retains an off-diagonal defect and how a reduced inverse corrects it. Their
whole-space Hamiltonian assumptions are not applied to our compact integral
operator. The needed one-step calculation is given below. L\"uscher's torus
expansion remains relevant context, but no continuum coefficient or time
normalization is substituted for the present Wilson convention.

Fix $L=2m$, $m\ge4$, and $V=L^3$. All constants in this section may depend on
$L$. Retain
\[
 T_\gamma=M_{b_\gamma}C_\gamma^{\otimes3V}M_{b_\gamma},\qquad
 b_\gamma=e^{-\gamma S/2},\qquad \lambda_* =\Lambda_{{\rm G},B}>0.
\]
The Hilbert space, Haar measure, gauge action and eight center sectors are
unchanged. Use V44's orthogonal decomposition
$\mathcal N=\mathcal Y\oplus\mathcal Z$, with $r=6(V-1)$ hard coordinates and
nine constant soft coordinates. Write $\mathsf K:\mathbb R^9\to\mathcal Z$ for
constant extension to all fine links, so
\[
 (\mathsf K a)_{x,\mu}=a_\mu,\qquad \mathsf K^*\mathsf K=V I_9.
\]
The normalized isometry $\mathsf K/\sqrt V$ identifies $\mathbb R^9$ with the
orthonormal soft coordinates. Define
\begin{equation}
 h=(\gamma V)^{-1/3},\qquad
 \varepsilon=\gamma^{-1/2}=\sqrt V\,h^{3/2},\qquad
 \tau=\sqrt V\,h,\qquad
 s=\varepsilon u+h\mathsf K q.
 \label{f16v45:scales}
\end{equation}
Thus $a=hq$ is the constant link angle and $z=\tau q$ its orthonormal coordinate.
In particular
\begin{equation}
 \varepsilon^2/\tau^2=h,\qquad \gamma Vh^4=h,
 \qquad \varepsilon=o(h),\quad h^2=o(h).
 \label{f16v45:balance}
\end{equation}
The symbol $h$ is an asymptotic parameter, not physical time.

Let $\mathcal T_+$ be exactly \eqref{f16v44:hard-transfer}, and let
$\varphi>0$ be its unit Gaussian ground function. Put $P_+=|\varphi\rangle
\langle\varphi|$. At this fixed lattice,
\begin{equation}
 \mathcal T_+\varphi=\lambda_*\varphi,\qquad
 R_+=(\lambda_* -\mathcal T_+)^{-1}(I-P_+),\qquad
 \|R_+\|\le[\lambda_*(1-e^{-\xi_{\min}})]^{-1}<\infty.
 \label{f16v45:hard-inverse}
\end{equation}
Here $\xi_{\min}>0$ is the smallest positive hard Mehler frequency. This inverse
is only the fixed Gaussian hard-band inverse; it is not an inverse on a native
neutral complement or on the nine soft coordinates.

\subsection{The two first variations that must be retained}

Use V44's exact tube coordinate $\pi(U)=s$ modulo the based gauge group and its
positive density $J_L(s)$. For small $a\in\mathbb R^9$ set
\[
 H(a)=D_y^2S(\operatorname{Exp}(\mathsf K a+y))\big|_{y=0,\ y\in\mathcal Y},
 \qquad
 R(a)Z=\left.\frac d{dt}\pi(\operatorname{Exp}(\mathsf K a)
                                      \operatorname{Exp}(tZ))\right|_{t=0}.
\]
Let $R_Y(a)=\Pi_{\mathcal Y}R(a)$ and $A(a)=R_Y(a)R_Y(a)^*$. These are smooth
finite-dimensional matrices on a fixed neighborhood. They obey
$H(0)=H_+$ and $A(0)=I_{\mathcal Y}$, and are positive for a sufficiently small
neighborhood. Define their directional derivatives
\[
 H_1(q)=DH(0)[q],\qquad A_1(q)=DA(0)[q].
\]
No quotient covariance is assigned the value one away from $a=0$.

\begin{proposition}[Exact hard stationarity and the normalized hard derivative]
\label{f16v45:jets}
For every small constant field $a$, including noncommuting triples,
\begin{equation}
 D_yS(\operatorname{Exp}(\mathsf K a+y))\big|_{y=0}=0
                 \quad\hbox{on }\mathcal Y.
 \label{f16v45:stationarity}
\end{equation}
Consequently, on fixed rescaled compact sets,
\begin{equation}
 \frac{S(\operatorname{Exp}(\varepsilon u+h\mathsf Kq))}{2\varepsilon^2}
 =\frac14u^{\mathsf T}H_+u+
 h\left\{\mathcal Q(q)+\frac14u^{\mathsf T}H_1(q)u\right\}
 +O_{L}(\varepsilon P(u,q)+h^2P(u,q)),
 \label{f16v45:magnetic-expansion}
\end{equation}
where $P$ is a fixed polynomial bound and
$\mathcal Q(q)=\sum_{\mu<\nu}|q_\mu\times q_\nu|^2$.

For the frozen, normalized hard Gaussian operator
\[
 \mathcal T(a)(u,v)=
 \frac{\exp[-(u-v)^{\mathsf T}A(a)^{-1}(u-v)/2
                 -u^{\mathsf T}H(a)u/4-v^{\mathsf T}H(a)v/4]}
 {(2\pi)^{r/2}\sqrt{\det A(a)}},
\]
its derivative $\mathcal C(q)=D\mathcal T(0)[q]$ has kernel
\begin{equation}
 \begin{split}
 \mathcal C(q)(u,v)=\mathcal T_+(u,v)\bigl\{
 &-\tfrac12\Tr A_1(q)+\tfrac12(u-v)^{\mathsf T}A_1(q)(u-v)\\
 &-\tfrac14u^{\mathsf T}H_1(q)u-\tfrac14v^{\mathsf T}H_1(q)v\bigr\}.
 \end{split}
 \label{f16v45:hard-derivative}
\end{equation}
It is a bounded self-adjoint Hilbert--Schmidt operator, linear in $q$, and
\begin{equation}
 \langle\varphi,\mathcal C(q)\varphi\rangle=0.
 \label{f16v45:singlet-cancellation}
\end{equation}
\end{proposition}
\begin{proof}
Translation invariance makes the logarithmic first derivative of $S$ at a
constant field independent of the base vertex, in each direction and color.
Every $y\in\mathcal Y$ has zero mean in each such component, so its pairing
with that derivative vanishes. This proves \eqref{f16v45:stationarity} exactly;
it is the same stationarity mechanism as f14 V54, now on the spatial slice.
Taylor expansion in $y$ therefore starts with its Hessian, not a linear hard
force. Its cubic remainder is bounded by $C_L|y|^3$ on the fixed chart. V44's
exact constant-field formula gives
$S(\operatorname{Exp}(\mathsf K a))=2V\mathcal Q(a)+O_L(|a|^6)$.
Expand $H(a)=H_++H_1(a)+O_L(|a|^2)$ and use
\eqref{f16v45:balance}. The constant sixth-order contribution is $O_L(h^3)$.
These facts give \eqref{f16v45:magnetic-expansion}.

Differentiate the complete Gaussian kernel, including its determinant. This
gives \eqref{f16v45:hard-derivative}; Gaussian domination proves the claimed
operator differentiability and boundedness. For a simultaneous color rotation
$k$, let $O_k$ denote its orthogonal action on $\mathcal Y$. Equivariance of the
tube and action gives $A_1(kq)=O_kA_1(q)O_k^*$ and the same formula for $H_1$.
The ground function $\varphi$ is invariant. Hence the scalar
$q\mapsto\langle\varphi,\mathcal C(q)\varphi\rangle$ is a rotation-invariant
linear functional on three copies of $\mathbb R^3$. It is zero, since none of
those vector representations has a fixed nonzero covector. This proves
\eqref{f16v45:singlet-cancellation}. It does not assert
$\mathcal C(q)\varphi=0$.
\end{proof}

The determinant term $-\Tr A_1/2$ is retained even when a choice of normal
coordinates makes $A_1$ vanish. Likewise the mean-zero stationarity is essential:
without it a constant-field quadratic hard force could appear on a larger scale.
Neither cancellation follows merely from smallness of the original action.

\subsection{A strong local expansion at the balanced scale}

Choose a smooth invariant cutoff $\chi(y)$, equal to one near zero, whose hard
support and a small fixed soft ball have product compactly contained in V44's
tube and in V42's central region $\mathcal O_e$. For a function $F(u,q)$ smooth
and compactly supported in $q$, with a polynomial times the hard ground Gaussian
in $u$, define
\begin{equation}
 (\mathcal U_\gamma F)(\Phi(g,y+z))=
 J_L(y+z)^{-1/2}\varepsilon^{-r/2}\tau^{-9/2}
 \chi(y)F(y/\varepsilon,z/\tau),
 \label{f16v45:dilation}
\end{equation}
with zero extension. The formula uses the orthonormal identification of
$\mathcal Z$ specified above. For small $h$ it is supported inside the original
tube. If $F(O_ku,kq)=F(u,q)$ it is a physical vector. Its exact Gram is
\begin{equation}
 \langle\mathcal U_\gamma F,\mathcal U_\gamma G\rangle
       =\int\chi(\varepsilon u)^2\overline{F(u,q)}G(u,q)\,du\,dq.
 \label{f16v45:Gram}
\end{equation}
Thus on any fixed finite family of these functions it differs from the full
product Gram by $O(e^{-c/\varepsilon^2})$. No density factor is dropped.

\begin{theorem}[The complete one-step expansion on the test core]
\label{f16v45:strong-expansion}
For every such $F$, with $\mathcal T_+$ and $\mathcal C(q)$ acting only in $u$,
\begin{equation}
 \begin{split}
 T_\gamma\mathcal U_\gamma F
 =\mathcal U_\gamma\bigl[\mathcal T_+F+h\{\tfrac12\mathcal T_+\Delta_qF
       -2\mathcal Q(q)\mathcal T_+F+\mathcal C(q)F\}\bigr]
       +o_{L,F}(h)
 \end{split}
 \label{f16v45:strong-local}
\end{equation}
in the original physical $L^2$ norm for invariant $F$, and in the based-invariant
space otherwise. Convergence is uniform on a fixed finite-dimensional span of
test functions. This is a core-wise expansion, not operator-norm convergence on
an entire changing Hilbert space.
\end{theorem}
\begin{proof}
We give the error accounting needed beyond V44's $o(1)$ limit. In the actual
kinetic integral write $U'=U\operatorname{Exp}(\varepsilon Z)$. Its normalized
scaled density is the standard full-link Gaussian density times
$1+O_L(\varepsilon^2P(Z))$ on logarithmically growing $Z$-balls. Its raw tail has
both row and column Schur bounds $C_Le^{-c_LR^2}$ outside $|Z|\le R$.
These follow directly from the original Wilson cosine and Haar factors and
V37's normalization; no heat kernel is substituted for the native kernel.

Write $s=\varepsilon u+h\mathsf Kq$. Smoothness of the tube coordinate gives,
uniformly on a fixed neighborhood,
\[
 s'=s+\varepsilon R(hq)Z+
 O_L(\varepsilon^2(1+|u|+|Z|)^2).
\]
Here the additional $\varepsilon u$ changes the differential by $O_L(\varepsilon
|u|)$, and is included in the remainder. In hard rescaled coordinates this is
\[
 u'=u+R_Y(hq)Z+O_L(\varepsilon P(u,Z)).
\]
In orthonormal soft coordinates it gives
\[
 q'=q+\sqrt h\,Z_0+O_L(h^{3/2}P(u,q,Z)),
 \qquad Z_0=(\mathsf K/\sqrt V)^*\Pi_{\mathcal Z}Z.
\]
At $a=0$ the random variables $Z_0$ and $\Pi_{\mathcal Y}Z$ are independent
standard Gaussians; all based-vertical components are integrated as well.
The change of their covariance at $a=hq$ is $O_L(h)$, which affects a first
soft derivative only at order $h^{3/2}$. The second differential of $\pi$ gives
an even smaller soft error $O_L(\varepsilon^2/\tau)=O_L(h^2)$.

After the outer Haar integral and the two physical half-densities have been
assembled, their exact ratio is $\sqrt{J_L(s)/J_L(s')}$. Since
$s'-s=O_L(\varepsilon |Z|)$ and $J_L$ is smooth positive, its difference from
one is $O_L(\varepsilon P(u,q,Z))=o(h)$ on the sets used below. This is why no
separate order-$h$ Jacobian potential occurs in this flat half-density
representation. It is not an assertion that the quotient Haar density is flat.

Taylor-expand the soft test through order two. Its first displacement has zero
Gaussian mean; its second gives $(h/2)\Delta_qF$. At this order the hard integral
is exactly $\mathcal T_+$. At both ends the constant magnetic multiplier is
$1-h\mathcal Q(q)+o(h)$, since the difference between $q'$ and $q$ changes this
term only at order $h^{3/2}$. The two factors give
$-2h\mathcal Q\mathcal T_+F$. Freeze the remaining hard displacement at $a=hq$.
Its hard marginal is the normalized Gaussian of covariance $A(hq)$, including
$\det A(hq)^{-1/2}$. Together with the two hard magnetic factors of
\eqref{f16v45:magnetic-expansion}, its first variation is exactly
$h\mathcal C(q)F$. No term proportional to $h^{1/2}$ survives. This yields the
three displayed terms, with all normalizations still included.

Here is a norm justification of the truncations. A polynomial times the hard
Gaussian has $L^2$ tails, and finitely many derivative tails, bounded by
$C_Fe^{-c_FR^2}$. Cut that input and the raw displacement to radius $R$.
Contraction of $T_\gamma$ pays the first omitted input; the raw row/column Schur
bound pays the second. A fixed physical cutoff boundary is separated from the
remaining input and contributes $O(e^{-c/\varepsilon^2})$. On the retained
pieces the output has hard coordinate of size $O_L(R)$ and soft coordinate in
a fixed compact enlargement of the test support. Taylor remainders, including
three soft derivatives, the hard cubic, smooth coordinate errors and the second
frozen hard variation, have $L^2$ norm at most
\[
 C_{L,F}(h^{3/2}+\varepsilon+h^2)(1+R)^{N_F}.
\]
The finite exponent also includes the volume of this rescaled output set.
All intermediate Gaussian integrals have uniformly positive hard quadratic
parts. Choose $R=(K_F\log(1/h))^{1/2}$ with fixed $K_F$ large enough that both
omitted tails are $o(h)$. Equation \eqref{f16v45:balance} makes the last display
$o(h)$. The native principal cut lies outside these sets and is part of the
same paid displacement or input tail. This proves a strong $L^2$ expansion,
not merely pointwise convergence. For a fixed finite family the constants and
truncation can be chosen in common, proving uniformity on its span.
\end{proof}

\subsection{Correcting the hard-band leakage rather than discarding it}

Gaussian integration in \eqref{f16v45:hard-derivative} shows that
$\mathcal C(q)\varphi$ is a hard polynomial of degree at most two times
$\varphi$, linear in $q$. Equation \eqref{f16v45:singlet-cancellation} removes
its ground component. The finite Hermite space of degrees at most two is
preserved by $\mathcal T_+$, so the following corrector has the same regularity:
\begin{equation}
 k(q)=R_+\mathcal C(q)\varphi,
 \qquad \langle\varphi,k(q)\rangle=0,
 \qquad (\lambda_* -\mathcal T_+)k(q)=\mathcal C(q)\varphi.
 \label{f16v45:corrector}
\end{equation}
It is equivariant under simultaneous color rotations. For an invariant
$f\in C_c^\infty(\mathbb R^9)$ put
\begin{equation}
 \Psi_\gamma f=\mathcal U_\gamma[(\varphi+h k(q))f(q)],
 \qquad H_{\rm mat}=-\tfrac12\Delta_q+2\mathcal Q(q).
 \label{f16v45:corrected-embedding}
\end{equation}
No native complement inverse is assumed in this definition.

\begin{theorem}[Corrected physical recovery at the quartic scale]
\label{f16v45:recovery}
For invariant compact smooth $f,g$,
\begin{align}
 \langle\Psi_\gamma f,\Psi_\gamma g\rangle
 &=\langle f,g\rangle+
 h^2\int\overline{f(q)}g(q)\|k(q)\|_2^2\,dq
                       +O_{L,f,g}(e^{-c/\varepsilon^2}),
 \label{f16v45:corrected-Gram}\\
 \big\|T_\gamma\Psi_\gamma f-
       \lambda_*\Psi_\gamma(f-hH_{\rm mat}f)\big\|_2&=o_{L,f}(h).
 \label{f16v45:intertwining}
\end{align}
Both conclusions are uniform on every fixed finite test space. In particular
\begin{equation}
 \frac{\lambda_*\langle\Psi_\gamma f,\Psi_\gamma g\rangle
          -\langle\Psi_\gamma f,T_\gamma\Psi_\gamma g\rangle}
 {\lambda_*h}
       \longrightarrow\langle f,H_{\rm mat}g\rangle.
 \label{f16v45:form-recovery}
\end{equation}
The same statements hold in each center sector after the exact character
folding of V44.
\end{theorem}
\begin{proof}
The Gram follows from \eqref{f16v45:Gram} and the pointwise orthogonality of
$k(q)$ to $\varphi$. Its order-$h$ cross term is exactly zero before a limit.
Apply Theorem~\ref{f16v45:strong-expansion} to $\varphi f$. Since
$\mathcal T_+\varphi=\lambda_*\varphi$, its order-$h$ term is
$-\lambda_*\varphi H_{\rm mat}f+\mathcal C(q)\varphi f$.
Apply the zeroth-order part of the same theorem to $k(q)f$ and multiply by $h$.
The identity in \eqref{f16v45:corrector} cancels the entire vector
$\mathcal C(q)\varphi f$, not just its expectation. The extra term from
$h\,k(q)H_{\rm mat}f$ on the right side of \eqref{f16v45:intertwining} has
coefficient $h^2$ and is $o(h)$. This proves the strong residual estimate and
then its bilinear consequence.

Every constructed vector is physical by the exact equivariant embedding.
Choose the tube closure inside a smaller neighborhood of the central orbit.
Its eight translates have a fixed positive separation. Folding by
$8^{-1/2}\sum_\sigma\chi_a(\sigma)\mathsf Z_\sigma$ is an exact isometry on
this supported space. The same residual estimate holds by commutation; cross
matrix elements between distinct translates are $O_L(e^{-c_L\gamma})$ by the
native displacement estimate. No regional trial state is substituted into an
intermediate transfer power. No conclusion about an almost invariant band for
all native low-energy states is made.
\end{proof}

\paragraph{A finite check on the role of the correction.}
Take two isotropic three-component hard oscillators with distinct Mehler ratios
$r_1,r_2\in(0,1)$, and let $p_q(u)=u^{\mathsf T}H_1(q)u$ be a rotation-covariant
cross quadratic proportional to $q\cdot(u_1\times u_2)$. With $A_1=0$,
\[
 \mathcal C(q)\varphi
   =-\frac{\lambda_*}{4}(1+r_1r_2)p_q\varphi,
 \qquad
 k(q)=-\frac{1+r_1r_2}{4(1-r_1r_2)}p_q\varphi.
\]
The expectation of the first vector is zero, but its norm is generally nonzero.
This example checks the sign and denominator in \eqref{f16v45:corrector}; it is
not a numerical evaluation of the native derivative matrices. A bare Rayleigh
calculation would miss the uncancelled order-$h$ residual.

\subsection{Native spectral inclusion and its precise one-sided ordering}

Write $e_0<e_1\le e_2\le\cdots$ for the invariant eigenvalues of V44's same
quartic model, repeated with multiplicity. They are not numerically evaluated
here. A compactly supported invariant smooth graph core is available: for a
model eigenfunction, multiplication by radial cutoffs converges in graph norm
because $\nabla f\in L^2$ and
\[
 [H_{\rm mat},\chi_R]f=-\nabla\chi_R\cdot\nabla f
                              -\tfrac12(\Delta\chi_R)f\longrightarrow0.
\]
Local elliptic regularity and mollification on the compact support finish the
approximation. Averaging over simultaneous rotations preserves it and the
operator. This proof also treats a finite eigenspace in one common graph core.

\begin{theorem}[Every fixed quartic level is realized to smaller than its scale]
\label{f16v45:inclusion}
Fix $L$, a model eigenvalue $e$ of invariant multiplicity $d_e$, and a center
character $a$. There is $\eta_\gamma\to0$ such that the original sector
$T_\gamma|_{\mathcal H_a}$ has at least $d_e$ eigenvalues, counted with
multiplicity, in
\begin{equation}
 [\,\lambda_*(1-he)-h\eta_\gamma,
                  \lambda_*(1-he)+h\eta_\gamma\,].
 \label{f16v45:spectral-window}
\end{equation}
One may choose a common sequence for any fixed finite list of model levels and
all eight sectors. The corresponding eigenvalues lie in V43's near-top interval
for sufficiently large coupling. They need not be the first $d_e$ or any
specified ordered eigenvalues of their native sector.
\end{theorem}
\begin{proof}
Approximate an orthonormal model eigenbasis at $e$ by invariant compact smooth
tests in graph norm. For each fixed approximation Theorem~\ref{f16v45:recovery}
gives its native residual divided by $h$, bounded by the model graph error plus
a term tending to zero. Choose the approximation index to increase sufficiently
slowly with coupling. This diagonal choice is also required to keep the physical
supports inside the common tube. The resulting native Gram tends to identity;
orthonormalizing on this fixed-dimensional space preserves a residual bound
$o(h)$ uniformly on its unit sphere. Thus there is $\zeta_\gamma\to0$ with
\[
 \|(T_\gamma-\lambda_*(1-he))u\|\le h\zeta_\gamma\|u\|
\]
on a $d_e$-dimensional physical sector space. The spectral theorem bounds the
norm outside a window of radius $h\sqrt{\zeta_\gamma}$ by
$\sqrt{\zeta_\gamma}\|u\|$. Its spectral projection is injective on that space
for large coupling, proving the multiplicity assertion. A zero residual needs
no enlargement; increasing $\eta_\gamma$ slightly handles both cases. The
windows are eventually positive and above the fixed cutoff $c<\lambda_*$, so
the compact operator has only discrete eigenvalues there. A common finite
maximum of errors and thresholds handles finitely many levels and characters.
No assertion of spectral exhaustion enters this argument.
\end{proof}

\begin{theorem}[One-sided ordered level estimate at the balanced scale]
\label{f16v45:minmax}
With indexing starting at one in each native center sector, for every fixed $j$,
\begin{equation}
 \boxed{\quad
 \limsup_{\gamma\to\infty}
     \frac{1-\lambda_{a,j}(\gamma)/\lambda_*}{h}
                  \le e_{j-1}.
 \quad}
 \label{f16v45:ordered-upper}
\end{equation}
The same statement holds for V43's ordered one-region eigenvalues $r_j$.
In particular
\begin{equation}
 \Lambda_\gamma\ge\lambda_*[1-h(e_0+o(1))],\qquad
 \limsup_{\gamma\to\infty}
     \frac{\log(\lambda_*/\Lambda_\gamma)}{h}\le e_0.
 \label{f16v45:top-one-side}
\end{equation}
These expressions are not assumed nonnegative. No upper estimate for
$\Lambda_\gamma-\lambda_*$ is proved.
\end{theorem}
\begin{proof}
Take a $j$-dimensional model graph-core space whose maximum Rayleigh quotient
is at most $e_{j-1}+\eta$. Equation \eqref{f16v45:form-recovery} and the paid
Gram give a physical sector space on which the native transfer quotient is at
least $\lambda_*[1-h(e_{j-1}+\eta)+o(h)]$. Min--max proves the displayed
limsup after $\eta\downarrow0$. The unfurled space is supported in the chosen
region and hence has the same quadratic form in the one-region compression.
This proves its assertion too. Take $j=1$ and the neutral character, and use
$-\log(1-hc)=hc+O(h^2)$ on the lower bound, for the final inequality. This
argument uses a lower bound on the top, not convergence of its first correction.
\end{proof}

\begin{proposition}[Explicit variational constants without a model diagonalization]
\label{f16v45:explicit-bounds}
For the invariant model eigenvalues and the native neutral ordered levels,
\begin{equation}
 e_0\le\frac{27}{4},\qquad e_1\le\frac{47}{4},\qquad
 \limsup_{\gamma\to\infty}\frac{1-\Lambda_\gamma/\lambda_*}{h}
       \le\frac{27}{4},\qquad
 \limsup_{\gamma\to\infty}\frac{1-\lambda_{0,2}/\lambda_*}{h}
       \le\frac{47}{4}.
 \label{f16v45:explicit-values}
\end{equation}
These are variational upper bounds on deficits relative to $\lambda_*$, not an
upper bound $5h$ on the actual neutral gap and not evaluated model eigenvalues.
\end{proposition}
\begin{proof}
Let $d=9$, $r^2=|q|^2$, and
\[
 \phi_0=(2/\pi)^{9/4}e^{-r^2},\qquad
 \phi_1=\frac{2\sqrt2}{3}(r^2-9/4)\phi_0.
\]
They are orthonormal, invariant, and in the model form domain. Under
$\phi_0^2dq$, the coordinates are independent with variance $1/4$, and $r^2$
has the gamma moments with shape $9/2$ and rate $2$. Direct integration gives
\[
 \big(\langle\phi_i,-\tfrac12\Delta\phi_j\rangle\big)_{i,j=0}^1
 =\begin{pmatrix}9/2&-3/\sqrt2\\-3/\sqrt2&13/2\end{pmatrix},\qquad
 \big(\langle\phi_i,2\mathcal Q\phi_j\rangle\big)_{i,j=0}^1
 =\begin{pmatrix}9/4&3/\sqrt2\\3/\sqrt2&21/4\end{pmatrix}.
\]
For the potential moments, homogeneity of $\mathcal Q$ separates the angular
factor, and the remaining radial moments are those just specified; its base
expectation is $9/4$. The summed matrix is
$\operatorname{diag}(27/4,47/4)$. Model min--max proves the first two bounds.
Theorem~\ref{f16v45:minmax} gives the native inequalities. Subtracting two
one-sided ordered bounds would not bound their difference, which explains the
qualification in the statement.
\end{proof}

\subsection{The existing outer return is negligible on these packets only}

Let $Q$ be V42's forbidden complement. The supports of the newly constructed
packets lie in a fixed compact subset of the union of retained regions. The
same is true for the slowly varying graph-core approximants above, by their
diagonal choice. A fixed distance $d_0>0$ therefore separates their support
from $Q$. The original displacement estimate gives
\begin{equation}
 \|QT_\gamma\Psi_{\gamma,a}f\|
                \le C_Le^{-c_L\gamma}\|\Psi_{\gamma,a}f\|.
 \label{f16v45:outer-leakage}
\end{equation}
For $u=W_a^*\Psi_{\gamma,a}f$ and $z\ge c$, V43's exact reconstruction then
satisfies
\begin{equation}
 \|J_a(z)u-W_a u\|+\|(I-F_a'(z))u-u\|
                   \le C_Le^{-c_L\gamma}\|u\|.
 \label{f16v45:outer-metric}
\end{equation}
Indeed the first term contains one bounded complement inverse and the leakage
in \eqref{f16v45:outer-leakage}; the second contains its actual inverse-square
return. This proves that the \emph{outer} metric does not alter the leading
coefficient on these recovery vectors. It is not operator-norm convergence of
that metric to identity on the whole regional space, and does not remove
within-region hard-band correction \eqref{f16v45:corrector}.

\subsection{V45 checkpoint: spectral inclusion is not a complete scaling theorem}

The new estimates establish the native upper-recovery half at its proposed
scale, including a corrected vector residual. They turn the comparison-model
levels into actual native spectral points with $o(h)$ errors and supply ordered
one-sided bounds. They do not prove the converse inclusion or exclude additional
near-top states with a different concentration scale. The all-level $o(1)$
limit of V44 alone cannot provide that exclusion.

A simple compact-operator example shows the logical distinction. On
$\ell^2\oplus\mathbb C$ let a positive contraction have eigenvalues
$\lambda_*e^{-h e_j}$ on the model coordinates, but eigenvalue
$\lambda_*(1+\sqrt h)$ on the final coordinate, for small $h$ and fixed
$\lambda_*<1$. Its model coordinate vectors have all the recovery and
spectral-inclusion properties above, and every fixed ordered eigenvalue tends
to $\lambda_*$. Nevertheless its top correction is of order $\sqrt h$, not
$h$, and the lowest ordered gap is not given by $e_1-e_0$. This is not a native
counterexample; it rules out inferring exhaustion or vacuum subtraction from
the proved one-sided data. Center copies can be added without changing the
point of the example.

A sufficient next native target is a lower bound on the rescaled deficit and
compactness of every sequence with bounded deficit after centering at
$\lambda_*$. In tube variables this must pay: tightness of the soft coordinate
$q$ along the commuting valleys, exclusion of hard excited mass on the same
scale, and all regions or coordinate tails not represented by the recovery
map. V44's coercive model inequality is not a substitute for those estimates
for the changing native operators. The exact regional Feshbach map and its
return metric remain the correct objects for such a proof.

Only after that completeness step could one identify
$\Lambda_\gamma=\lambda_*[1-he_0+o(h)]$ and then the neutral deficit with
$\lambda_*h(e_1-e_0)$. Neither formula is a theorem of V45. Nor is an independent
physical time step supplied. In particular the scale $h$ in
\eqref{f16v45:scales} is not defined to be physical Euclidean time. The supplied
memoranda's separation between residual power and locality cost, and their
requirement to keep the full Feshbach return, remain operative. The checkpoint
file records the current branch without recertifying the accumulated archive.

\subsection{Verification and preservation}

The new exact diagnostic is
\begin{center}\small\path{verify_ym_f16_v45_critical_recovery.py}.\end{center}
It checks the balanced powers and normalization, finite translation
stationarity, normalized Gaussian covariance derivatives, color-singlet
cancellation with nonzero hard leakage, the reduced-inverse sign, corrected
Grams, spectral projection counting, and the non-exhaustion example. Its
finite reference matrices are not discretizations or enclosures of the native
Wilson spectrum. Actual executions, source hashes and review scope are in the
audit and checkpoint; rendering checks concern presentation, not proof validity.

The supplied V44 source is unchanged except for its literal title version and
this terminal insertion. Removing those edits recovers it byte for byte. The
hard-band construction uses the exact V44 tube and the original one-step
Wilson kernel; it does not use the endpoint pressure branch. V43 enters only
the outer-return interpretation. The complete earlier analytic lineage has not
been independently recertified. The next scheduled checkpoint is V50.

\begin{firewall}
V45 proves critical-scale physical recovery, an explicit hard-band correction,
and native spectral inclusion at the quartic scale with multiplicities. It
also proves one-sided ordered-level bounds. It does not prove complete
rescaled spectral convergence, identify the native Perron correction or first
neutral coefficient, exclude additional soft states, calibrate physical time,
remove the fixed spatial regulator, or establish an interacting continuum
mass gap. All original probabilities, transfer normalizations, gauge and center
constraints, and the exact returning-state metric are retained.
\end{firewall}
% END F16 V45 APPEND

% BEGIN F16 V46 APPEND -- 2026-09-18
% Insert before the final document-ending command of the supplied V45.
% No predecessor proof, normalization, probability, or clock is changed.

\clearpage
\section{V46: discrete quartic confinement and a native profile alternative}
\label{f16v46:context}

V45 constructs native recovery vectors, but its error estimate is on a test
core. It does not control arbitrary states near the spectral top. We go
upstream to the form that a lower comparison must control. A finite transfer
step has a bounded, frequency-saturating deficit; it cannot simply be replaced
by an unbounded differential energy on every test. We prove an all-vector
confinement estimate for the exact Gaussian-step quartic sandwich, including
its commuting valleys. It gives asymptotic compactness, a lower form limit,
and complete spectral convergence for that explicitly declared comparison.

A separate native deduction uses V45's strong recovery and V42's outer
localization: every nonzero weak critical profile of a native eigenvector is
an eigenfunction of the same quartic model. Missing native states must
therefore lose compactness, or have deficits unbounded below on this scale.
We finish with one precise nonlocal domination inequality sufficient to
exclude both possibilities. That inequality is a remaining native target,
not an estimate proved by identifying the comparison with Wilson transfer.

\subsection{Sources, nonduplication, and the two different operators}

The supplied V45 appendix, analytic audit, and checkpoint were read completely.
The current source and supplied f14/f15/Grand Tensor archives were searched.
The Grand Tensor already states abstract Mosco and compact-screen principles;
its identified toron passage concerns an $\SU(3)$ family with a separate
localization premise. Neither is used as a missing $\SU(2)$ Wilson bound.
V44 already proves differential quartic confinement. The new estimate below
is for the \emph{finite nonlocal step}, before its differential limit, and
pays both position and frequency tails. General compactness and min--max
principles are credited rather than claimed new. The next checkpoint remains V50.

For primary context, Ichinose--Tamura,
\href{https://arxiv.org/abs/0905.3458}{\emph{Results on Convergence in Norm of
Exponential Product Formulas}}, discusses norm convergence of symmetric
products. Its introductory setting and Theorems 2.1--2.2 were inspected in
parsed primary text. We do not import a rate from that theorem; the lower
and compactness argument required here is proved explicitly. The normalized
oscillator calculation can also be checked against
\href{https://dlmf.nist.gov/18.18.E28}{DLMF 18.18.28}. Simon's
\href{https://doi.org/10.1016/0003-4916(83)90057-X}{discrete-spectrum work}
is the classical precedent already cited in V44. No new Navier--Stokes
statement is a premise. The supplied memoranda's separate compactness and
residual-scale requirements remain in force.

On $\mathcal K=L^2(\mathbb R^9,dq)$ write
\[
 \mathcal Q(q)=\sum_{i<j}|q_i\times q_j|^2,\qquad
 H_{\rm mat}=-\tfrac12\Delta+2\mathcal Q,\qquad q_i\in\mathbb R^3.
\]
For an independent numerical step $h>0$ define
\begin{equation}
 m_h=e^{-h\mathcal Q},\qquad G_h=e^{h\Delta/2},\qquad
 S_h=m_hG_hm_h,\qquad D_h=h^{-1}(I-S_h),\qquad
 d_h[f]=\langle f,D_hf\rangle.
 \label{f16v46:comparison}
\end{equation}
Multiplication is meant in $m_h$. The heat kernel is exactly
$(2\pi h)^{-9/2}e^{-|q-q'|^2/(2h)}$. The sandwich is positive self-adjoint,
$0\preceq S_h\preceq I$. It is not $e^{-hH_{\rm mat}}$, is not the original
Wilson transfer, and is not a conditional law on constant fine links. The
Gaussian step is explicit here; it is never substituted for a native Wilson
step without an additional comparison. Every assertion also restricts to the
closed simultaneous-rotation-invariant subspace $\mathcal K_{\rm inv}$.

\subsection{The exact deficit retains its saturated kinetic part}

Use the unitary Euclidean Fourier transform, with variable $p$.
\begin{proposition}[A positive two-term identity on all of $L^2$]
\label{f16v46:identity}
For every $f\in L^2$ and every $h>0$,
\begin{equation}
 \boxed{\quad
 d_h[f]=\int\frac{1-e^{-2h\mathcal Q(q)}}h|f(q)|^2dq
       +\int\frac{1-e^{-h|p|^2/2}}h
                         |\widehat{m_hf}(p)|^2dp.
 \quad}
 \label{f16v46:exact-energy}
\end{equation}
In particular $\|f-m_hf\|_2^2\le h d_h[f]$. If $hR^2\le1$, then
\begin{equation}
 \int_{|p|>R}|\widehat f(p)|^2dp
             \le 2h d_h[f]+8R^{-2}d_h[f].
 \label{f16v46:Fourier-tail}
\end{equation}
Neither estimate requires an $H^1$ test or a bounded quartic potential.
\end{proposition}
\begin{proof}
Add and subtract $\|m_hf\|^2$ in
$\|f\|^2-\langle m_hf,G_hm_hf\rangle$ and apply Plancherel. For
$0\le m\le1$, $(1-m)^2\le1-m^2$, proving the norm estimate. On $|p|>R$,
the second integrand multiplier is at least
$(1-e^{-hR^2/2})/h$. For $hR^2\le1$ this is at least $R^2/4$.
Apply the squared triangle inequality to
$\widehat f=\widehat{m_hf}+\widehat{f-m_hf}$. This gives
\eqref{f16v46:Fourier-tail}. All terms are finite since the multipliers are
bounded at fixed $h$.
\end{proof}

The kinetic multiplier tends to $|p|^2/2$ only at fixed frequency. It is at
most $1/h$. Thus a bound $d_h[f]\ge c\|\nabla f\|^2$ on every test, with
$c>0$, would be false even at one fixed $h$. The following confinement
estimate has the corresponding necessary saturation in position.

\subsection{A transverse-oscillator estimate before the differential limit}

Put
\[
 b(t)=\exp[-4\operatorname{arsinh}(t)]
       =\bigl(\sqrt{1+t^2}-t\bigr)^4,\qquad t\ge0.
\]
\begin{theorem}[Nonlocal confinement through all commuting valleys]
\label{f16v46:confinement}
For each $i=1,2,3$, as bounded quadratic forms,
\begin{equation}
 \boxed{\qquad S_h\preceq M_{b(h|q_i|)}.\qquad}
 \label{f16v46:fibre-bound}
\end{equation}
Consequently, for every $f\in L^2$,
\begin{align}
 d_h[f]&\ge\frac1{3h}\sum_{i=1}^3
       \int[1-b(h|q_i|)]|f(q)|^2dq,\label{f16v46:radial-exact}\\
 d_h[f]&\ge\frac16\int\sum_{i=1}^3
                  \min\{|q_i|,h^{-1}\}|f(q)|^2dq.\label{f16v46:radial-floor}
\end{align}
In particular, if $hR\le1$,
\begin{equation}
 \int_{|q|>R}|f(q)|^2dq\le\frac{6\sqrt3}{R}\,d_h[f].
 \label{f16v46:position-tail}
\end{equation}
No small-field window or restriction away from rank-one commuting triples is
used in these bounds.
\end{theorem}
\begin{proof}
Fix $i$. The heat factors in different three-vector variables commute and
are positive contractions. Hence
\[
 G_h\preceq I_i\otimes\exp\bigl[(h/2)\sum_{j\ne i}\Delta_{q_j}\bigr].
\]
Congruence by $m_h$ preserves this operator inequality. The right side now
disintegrates over $q_i=x$. At fixed $x$, put $r=|x|$ and remove only the
remaining pair potential $|q_j\times q_k|^2$, where $j,k\ne i$. The exact
fibre operator has the form $M_a B_x M_a$, with
\[
 0<a=e^{-h|q_j\times q_k|^2}\le1,\qquad
 B_x=e^{-h\sum_{j\ne i}|x\times q_j|^2}
       e^{(h/2)\sum_{j\ne i}\Delta_{q_j}}
       e^{-h\sum_{j\ne i}|x\times q_j|^2}.
\]
Thus its norm is at most $\|B_x\|$. This step is a contraction in operator
norm; it does not claim that increasing a multiplication factor orders two
sandwiches in Loewner order.

For $r>0$, rotate the two other vectors so that $x$ is axial. Each contributes
a free longitudinal heat factor of norm one and two transverse factors
\[
 e^{-hr^2y^2}e^{h\partial_y^2/2}e^{-hr^2y^2}.
\]
After $y=\sqrt h\,u$, one such factor has normalized kernel
\[
 (2\pi)^{-1/2}
 e^{-(u-v)^2/2-h^2r^2(u^2+v^2)}.
\]
Its ground eigenvalue is $e^{-\operatorname{arsinh}(hr)}$ by the Mehler
formula: its parameter is $\nu=4h^2r^2$ and
$\operatorname{arcosh}(1+\nu/2)=2\operatorname{arsinh}(hr)$.
This normalization can be checked by inserting the positive Gaussian of
frequency $\sqrt{\nu+\nu^2/4}$ and integrating. Its higher eigenvalues have
strictly smaller modulus. The four transverse factors therefore give
$\|B_x\|=b(hr)$. At $r=0$ all factors are free and the norm is one, the
same formula by continuity. The rotation is used only inside this fixed
fibre; no derivative in $q_i$ is taken or lost.

For any complex fibre test its Rayleigh quotient is bounded by that norm
times its squared norm. Integrating over $x$ proves
\eqref{f16v46:fibre-bound}. Average its three consequences for $I-S_h$.
For $0\le t\le1$, $\operatorname{arsinh}t\ge t/2$ and
$1-e^{-2t}\ge(1-e^{-2})t\ge t/2$; for $t\ge1$ monotonicity gives the same
lower bound $1/2$. Thus $1-b(t)\ge\frac12\min(t,1)$, proving
\eqref{f16v46:radial-floor}. If $|q|>R$, at least one $|q_i|>R/\sqrt3$.
When $hR\le1$, the sum of the saturated weights is at least $R/\sqrt3$ on
that event. This proves the last estimate.
\end{proof}

\begin{proposition}[Compact transfer, but not a compact resolvent at fixed step]
\label{f16v46:compact-step}
For every fixed $h>0$, $S_h$ is compact, injective, and positivity improving,
with $0<\|S_h\|<1$. The bounded operator $D_h$ does not have compact
resolvent on the infinite-dimensional space: its essential spectrum is
$\{h^{-1}\}$.
\end{proposition}
\begin{proof}
Partition $\{|q|>R\}$ into three disjoint measurable sets $E_i$ on which
$|q_i|>R/\sqrt3$. For $f=\sum_i1_{E_i}f$, positivity and the Hilbert-space
Cauchy--Schwarz inequality give
\[
 \|S_h^{1/2}f\|^2\le3\sum_i\|S_h^{1/2}1_{E_i}f\|^2
               \le3b(hR/\sqrt3)\|f\|^2.
\]
This tends to zero at fixed $h$ as $R\to\infty$. For $P_R=1_{\{|q|\le R\}}$,
it follows that $\|S_h-P_RS_hP_R\|\le2\sqrt{3b(hR/\sqrt3)}\to0$.
The bounded continuous kernel restricted to two bounded balls is
Hilbert--Schmidt. Hence $S_h$ is compact. Strict positivity of its kernel
gives positivity improvement. Its factors are injective, and
$\langle f,S_hf\rangle=\|G_h^{1/2}m_hf\|^2$ gives injectivity directly.
If its attained norm were one, \eqref{f16v46:exact-energy} would force
$\widehat{m_hf}$ to be supported at $p=0$, impossible for a nonzero $L^2$
function. Finally $D_h=h^{-1}I-h^{-1}S_h$ is a compact perturbation of
$h^{-1}I$. Its resolvent contains a nonzero scalar multiple of identity at
fixed $h$. The same conclusions hold in the invariant subspace.
\end{proof}

\subsection{All-vector compactness and the lower form limit}

\begin{theorem}[Asymptotic compactness, including the nonlocal ultraviolet tail]
\label{f16v46:asymptotic-compactness}
If $h_k\downarrow0$ and
\[
 \sup_k\bigl(\|f_k\|_2^2+d_{h_k}[f_k]\bigr)<\infty,
\]
then a subsequence converges strongly in $L^2(\mathbb R^9)$. If
$f_k\rightharpoonup f$, then
\begin{equation}
 \boxed{\quad
 \langle f,H_{\rm mat}f\rangle_{\rm form}
                 \le\liminf_k d_{h_k}[f_k].
 \quad}
 \label{f16v46:liminf}
\end{equation}
There are strongly convergent recovery sequences for every vector in the
model form domain. These are statements on all bounded-deficit sequences,
not only Gaussian or compactly supported trial functions.
\end{theorem}
\begin{proof}
For fixed large $R$, eventually $h_kR\le1$ and $h_kR^2\le1$.
Equations \eqref{f16v46:position-tail} and \eqref{f16v46:Fourier-tail} give
uniformly small position and Fourier tails in the successive limits
$k\to\infty$, then $R\to\infty$. For a direct compactness argument, the
operator
\[
 K_R=1_{\{|q|\le R\}}\mathcal F^{-1}1_{\{|p|\le R\}}\mathcal F
\]
is Hilbert--Schmidt. The norm of $f_k-K_Rf_k$ is bounded by the sum of the
two tail norms. At a fixed $R$ the sequence $K_Rf_k$ is relatively compact;
a diagonal extraction proves strong compactness of $f_k$. This argument
requires $h_k\to0$. Bounded form balls at one fixed $h$ are not compact.

For the lower bound select a subsequence attaining the finite liminf. It has
a further strongly convergent subsequence, whose limit must be the stated
weak limit $f$. Put $g_k=m_{h_k}f_k$. The first proposition gives $g_k\to f$
strongly. On a bounded position ball the potential multiplier in
\eqref{f16v46:exact-energy} converges uniformly to $2\mathcal Q$; on a bounded
Fourier ball its kinetic multiplier converges uniformly to $|p|^2/2$.
Retain both nonnegative truncated integrals, pass to the strong limit, and
then increase the two radii. This proves \eqref{f16v46:liminf}, including
membership in the model form domain.

If $f\in C_c^\infty$, $m_hf\to f$ in $H^1$, and
$(1-e^{-h|p|^2/2})/h\le|p|^2/2$. The exact energy identity therefore gives
$d_h[f]\to\frac12\|\nabla f\|^2+2\int\mathcal Q|f|^2$.
Compact smooth functions are a form core: spatial cutoff followed by local
mollification suffices for this nonnegative continuous polynomial potential.
A slow diagonal approximation supplies recovery for every form-domain
vector. Rotation averaging preserves that construction in the invariant
subspace. No derivative of an arbitrary $f_k$ was assumed.
\end{proof}

\begin{theorem}[Complete spectral convergence for the declared discrete comparison]
\label{f16v46:norm-resolvent}
On the full space, and separately on $\mathcal K_{\rm inv}$,
\begin{equation}
 \boxed{\quad
 \|(D_h+1)^{-1}-(H_{\rm mat}+1)^{-1}\|\longrightarrow0.
 \quad}
 \label{f16v46:resolvent}
\end{equation}
Let $\mu_0(h)\ge\mu_1(h)\ge\cdots>0$ be the invariant eigenvalues of $S_h$.
Then for each fixed $j$,
\begin{equation}
 \frac{1-\mu_j(h)}h\longrightarrow e_j,
 \qquad -h^{-1}\log\mu_j(h)\longrightarrow e_j.
 \label{f16v46:complete-eigenvalues}
\end{equation}
Isolated finite spectral clusters have convergent projections in norm and
exactly the limiting multiplicity eventually. Also, for $0<t_0<T<\infty$,
\begin{equation}
 \sup_{t\in[t_0,T]}
 \|S_h^{\lfloor t/h\rfloor}-e^{-tH_{\rm mat}}\|\longrightarrow0.
 \label{f16v46:step-semigroup}
\end{equation}
No rate is asserted, and these are not Wilson eigenvalue or physical-time
statements.
\end{theorem}
\begin{proof}
The lower bound and recovery just proved give convergence of minimizers of
$d_h[u]+\|u\|^2-2\Re\langle g,u\rangle$. This is the usual form-to-resolvent
argument already recorded in the Grand Tensor, here with its hypotheses
verified for the present operators. To prove the stronger norm statement,
let $g_h$ be any unit-bounded sequence. The resolvent vectors
$u_h=(D_h+1)^{-1}g_h$ have uniformly bounded norm and energy. Select
$g_h\rightharpoonup g$ and $u_h\to u$ strongly by the preceding theorem.
The minimizer argument works with these varying right sides because every
recovery test converges strongly; it identifies
$u=(H_{\rm mat}+1)^{-1}g$. V44's compact-resolvent theorem for $H_{\rm mat}$
also gives $(H_{\rm mat}+1)^{-1}g_h\to u$ strongly. Failure of
\eqref{f16v46:resolvent} would contradict this conclusion on a sequence of
almost norm-maximizing right sides. This avoids asserting that each
$(D_h+1)^{-1}$ is compact; its essential scalar part is $h/(1+h)$ and vanishes.

Apply norm-continuous spectral projection calculus to the bounded positive
resolvents, at intervals separated from the limiting spectrum. Min--max for
their decreasing eigenvalues gives
$[1+(1-\mu_j(h))/h]^{-1}\to(1+e_j)^{-1}$. This proves the first limit in
\eqref{f16v46:complete-eigenvalues}; the scalar logarithm proves the second.
It also proves the counting and projection assertions.

Norm resolvent convergence of these nonnegative operators gives convergence
of $e^{-tD_h}$ for $t>0$, uniformly on compact positive time intervals. This
can be seen by polynomial approximation of $e^{-t(x^{-1}-1)}$, extended by
zero at $x=0$, on the bounded resolvent spectral interval. The approximation
is uniform for $t\in[t_0,T]$. Finally, uniformly for $0\le x\le h^{-1}$,
\[
 |(1-hx)^{\lfloor t/h\rfloor}-e^{-tx}|\le C_{t_0,T}h.
\]
For example split $hx\le1/2$ from its complement and use
$0\le-\log(1-y)-y\le y^2$ on the first part; the other part is exponentially
small in $\lfloor t/h\rfloor$. The rounding error in $t$ is bounded using
$\sup_{x\ge0}xe^{-t_0x/2}<\infty$. Functional calculus proves
\eqref{f16v46:step-semigroup}.
\end{proof}

\subsection{A reference that also detects hard-band loss}

Retain V45's exact hard Gaussian operator and let
\[
 A_+=\mathcal T_+/\lambda_*,\qquad
 P_+=|\varphi\rangle\langle\varphi|,\qquad
 r_+=\|A_+|_{\varphi^\perp}\|<1.
\]
On $\mathcal X=L^2(\mathcal Y\times\mathbb R^9)$, or its joint invariant
subspace, define the bounded nonlocal comparison form
\begin{equation}
 \mathbb D_h=h^{-1}(I-A_+\otimes S_h),\qquad
 \mathcal I f=\varphi\otimes f.
 \label{f16v46:hard-reference}
\end{equation}
\begin{proposition}[Exact hard separation for the comparison]
\label{f16v46:hard-separation}
Writing $F=\mathcal I f+F_\perp$ gives
\begin{equation}
 \boxed{\quad
 \langle F,\mathbb D_hF\rangle
 \ge d_h[f]+\frac{1-r_+}{h}\|F_\perp\|^2.
 \quad}
 \label{f16v46:hard-floor}
\end{equation}
Every sequence with bounded norm and bounded $\mathbb D_h$ energy as
$h\downarrow0$ has a strongly convergent subsequence of the form
$F_h\to\varphi\otimes f$. Moreover
\begin{equation}
 \left\|(\mathbb D_h+1)^{-1}
       -\mathcal I(H_{\rm mat}+1)^{-1}\mathcal I^*\right\|
 \le\|(D_h+1)^{-1}-(H_{\rm mat}+1)^{-1}\|
       +\frac{h}{1-r_++h}\longrightarrow0.
 \label{f16v46:hard-resolvent}
\end{equation}
\end{proposition}
\begin{proof}
The ground projection reduces the exact tensor comparison. On its range the
form is $D_h$, and on the complement $A_+\otimes S_h\preceq r_+I$.
This proves \eqref{f16v46:hard-floor} and the resolvent bound. The hard
complement norm is $O(\sqrt h)$ for a bounded-energy sequence, and its ground
profile has bounded $d_h$ energy. Apply Theorem~\ref{f16v46:asymptotic-compactness}.
The joint symmetry commutes with both projections; its ground component is
exactly the invariant soft subspace. No native hard-band floor is asserted.
\end{proof}

\subsection{A native profile theorem without an assumed lower bound}

We now return to the original $T_\gamma$ and fix $L$ throughout. Use
$h=(\gamma V)^{-1/3}$, $\varepsilon=\gamma^{-1/2}$ and
$\tau=\sqrt Vh$, exactly as in V45. A new symbol $\mathbb D_h$ above denotes a
comparison only. It is not identified with
\[
 \mathfrak A_\gamma=h^{-1}(I-T_\gamma/\lambda_*).
\]
The latter need not be nonnegative. Neither V41 nor V45 proves a uniform
lower bound for it.

Choose fixed nested, center-translated invariant cutoffs in V44's tubes,
$\chi_\sigma$, equal to one near each central orbit and with disjoint supports.
Put $\chi=\sum_\sigma\chi_\sigma$. Use the same factor $8^{-1/2}$ as in V43
to unfold a vector of character $a$ to one regional vector. Remove the exact
tube half-density and dilate by $s=\varepsilon u+h\mathsf Kq$. More explicitly,
if $\widehat\psi=W_a^*(\chi\psi)$, define
\begin{equation}
 (\mathfrak E_{\gamma,a}\psi)(u,q)
 =\varepsilon^{r/2}\tau^{9/2}J_L(s)^{1/2}
                 \widehat\psi(\Phi(1,s)),
 \qquad s=\varepsilon u+h\mathsf Kq,
 \label{f16v46:extraction}
\end{equation}
with zero extension outside the dilated tube. The reference is
$du\,dq$ in orthonormal hard and soft coordinates; all based orbit variables
have already been integrated against their probability Haar measure. This
is an exact contraction satisfying
$\|\mathfrak E_{\gamma,a}\psi\|=\|\chi\psi\|$. Its outputs are invariant under
joint color rotations.

V42 can be applied to fixed still smaller neighborhoods inside the plateaux.
It supplies a fixed cutoff $c<\lambda_*$ such that on the native band
$E_{\gamma,a}=1_{[c,1]}(T_\gamma|_{\mathcal H_a})$,
\begin{equation}
 \|(1-\chi)E_{\gamma,a}\|\le C_Le^{-c_L\sqrt\gamma}.
 \label{f16v46:extraction-mass}
\end{equation}
Thus extraction is asymptotically isometric on this entire band, with no
rank factor. The neighborhoods and cutoff are fixed before coupling; no
shrinking-neighborhood bound is imported from V42.

\begin{theorem}[Every nonzero native critical profile is a model eigenfunction]
\label{f16v46:native-profile}
Let $\psi_\gamma\in\mathcal H_a$ be normalized native eigenvectors with
\[
 T_\gamma\psi_\gamma=\lambda_*(1-hE_\gamma)\psi_\gamma,
 \qquad E_\gamma\to E\in\mathbb R.
\]
Along every subsequence on which
$\mathfrak E_{\gamma,a}\psi_\gamma\rightharpoonup F$ in $\mathcal X$,
there is $f\in\mathcal K_{\rm inv}$, $\|f\|\le1$, such that
\begin{equation}
 \boxed{\qquad F=\varphi\otimes f,\qquad H_{\rm mat}f=Ef.\qquad}
 \label{f16v46:profile-equation}
\end{equation}
The value $f=0$ is allowed. If $f\ne0$, then $E$ is an invariant model
eigenvalue and $E\ge e_0$. In particular, any native branch with bounded
scaled deficit converging to $E<e_0$ has only zero weak critical profiles.
This theorem does not assert that any such extra native branch exists.
\end{theorem}
\begin{proof}
The eigenvalues tend to $\lambda_*$, so \eqref{f16v46:extraction-mass} applies
and the extracted norms tend to one. Weak limits nevertheless can lose norm.
For every invariant test consisting of a hard Hermite polynomial times
$\varphi$ and a compact smooth soft test, the zeroth-order part of V45's
strong expansion gives
\[
 (T_\gamma-\lambda_*)\mathcal U_{\gamma,a}G
   =\mathcal U_{\gamma,a}[(\mathcal T_+-\lambda_*)G]+O_G(h).
\]
The tests' non-plateau hard tails are exponentially small, and the soft
support lies in the plateau for small $h$. Thus pairing with $\psi_\gamma$
and passing to its weak extracted limit gives
$\langle F,(\mathcal T_+-\lambda_*)G\rangle=0$.
These tests are dense in the jointly invariant space. The hard inverse on
$\varphi^\perp$ is bounded, so $F=\varphi\otimes f$; joint invariance forces
$f$ to be invariant. This uses no estimate on the hard mass of
$\psi_\gamma$ before passage to a weak limit.

For invariant $g\in C_c^\infty(\mathbb R^9)$ use V45's \emph{corrected}
recovery vector. Its strong residual gives exactly
\[
 \langle\psi_\gamma,\Psi_{\gamma,a}H_{\rm mat}g\rangle
 -E_\gamma\langle\psi_\gamma,\Psi_{\gamma,a}g\rangle=o(1).
\]
The $hk(q)$ correction disappears in each fixed pairing; the exact extracted
half-density normalization has already been retained. Passing to the weak
limit yields $\langle f,H_{\rm mat}g\rangle=E\langle f,g\rangle$.
Averaging any compact smooth test over rotations extends this identity to
all such tests, since $f$ is invariant.

For completeness this distributional equation is an operator eigenvalue
equation, not just a formal identity. Local elliptic regularity gives
$f\in H^2_{\rm loc}$. Testing with $\zeta_R^2f$, for standard cutoffs, yields
\[
 \tfrac12\|\nabla(\zeta_Rf)\|^2
 +2\int\mathcal Q\,\zeta_R^2|f|^2
 =E\|\zeta_Rf\|^2+\tfrac12\int|\nabla\zeta_R|^2|f|^2.
\]
The right side is bounded above uniformly in $R$. Letting $R$ increase proves
membership in the Friedrichs form domain. The distributional identity
extends by form density, giving $f\in D(H_{\rm mat})$ and
$H_{\rm mat}f=Ef$. V44's compact resolvent identifies $E$ as a model eigenvalue
when $f\ne0$. Neither strong convergence of the native profiles nor preservation
of their full norm has been used or established.
\end{proof}

\subsection{A specific sufficient native comparison, still to be proved}

The all-vector nonlocal reference is now available to formulate one missing
inequality without an impossible unsaturated gradient bound. The following is
an implication theorem with an additional hypothesis, not a native estimate.
\begin{proposition}[A coarse lower comparison would complete the fixed-volume scaling]
\label{f16v46:conditional-completeness}
Suppose there are fixed constants $C>0$, $b\ge0$, independent of $\gamma$,
such that on the neutral band $\operatorname{Ran}E_{\gamma,0}$,
\begin{equation}
 \boxed{\quad
 \langle\mathfrak E_{\gamma,0}\psi,
            \mathbb D_h\mathfrak E_{\gamma,0}\psi\rangle
 \le C\left[\langle\psi,\mathfrak A_\gamma\psi\rangle
                                    +b\|\psi\|^2\right]
 \quad\hbox{for every }\psi.
 }
 \label{f16v46:missing-domination}
\end{equation}
Then every fixed ordered native level in every center sector satisfies
\begin{equation}
 \frac{1-\lambda_{a,j}(\gamma)/\lambda_*}{h}\longrightarrow e_{j-1}.
 \label{f16v46:conditional-levels}
\end{equation}
In particular this additional hypothesis would imply
\begin{equation}
 \begin{split}
 \Lambda_\gamma&=\lambda_*[1-he_0+o(h)],\\
 -h^{-1}\log(\lambda_{0,2}/\Lambda_\gamma)&\longrightarrow e_1-e_0,\\
 \frac{g_\gamma}{\lambda_*h}&\longrightarrow e_1-e_0,
 \end{split}
 \label{f16v46:conditional-gap}
\end{equation}
with V43's actual returning-state metric retained in $g_\gamma$.
\end{proposition}
\begin{proof}
The left side of \eqref{f16v46:missing-domination} is nonnegative, so it
first gives $\mathfrak A_\gamma\succeq-bI$ on the band. This excludes the
unbounded negative-deficit possibility left by V45. For any fixed number $j$
of ordered neutral eigenvectors, V45 supplies an upper bound on their scaled
deficits, and the new hypothesis supplies the lower bound and uniformly
bounded $\mathbb D_h$ energies of their extracted profiles. Proposition
\ref{f16v46:hard-separation} gives strong subsequential compactness. The exact
band localization \eqref{f16v46:extraction-mass} preserves their entire Gram,
so their limiting profiles are $j$ orthonormal invariant functions.
Theorem~\ref{f16v46:native-profile} makes these model eigenfunctions with
the limiting ordered deficits. Model min--max gives the missing liminf
$e_{j-1}$, while V45 supplies the matching limsup. This proves the neutral
part of \eqref{f16v46:conditional-levels} on every subsequence, hence in full.
V43's sector ladders differ by $O(e^{-c_L\sqrt\gamma})=o(h)$ for every fixed
level, so the same conclusion holds in all sectors. Finally subtract the
first two neutral expansions, use the scalar logarithm, and apply V43--V44's
proved $\Delta_\gamma^{\rm neu}/g_\gamma\to1$.
\end{proof}

The constants in \eqref{f16v46:missing-domination} need not approach one.
The sharp coefficient would come from V45's local equation; the coarse bound
would only prevent loss of norm and energies escaping below the model.
Its left side has the correct bounded-transfer saturation and detects both
hard excited mass and soft escape. This is weaker in purpose than a global
$o(h)$ operator approximation, but it is not a consequence of the recovery
map. The factors arising from the actual quotient kinetic covariance,
nonlinear magnetic action, half-density, and within-region return remain to
be controlled when proving it. The Gaussian reference has not been installed
as the native conditional transition.

\subsection{Current scope and verification}

The new unconditional model result supplies complete nonlocal compactness and
spectral convergence without adding a boundary or deleting commuting valleys.
The new native profile result classifies every nonzero weak critical profile,
but permits zero profiles. Neither statement proves
\eqref{f16v46:missing-domination}. Thus the native ordered asymptotics and
neutral coefficient in \eqref{f16v46:conditional-gap} remain conditional,
not accomplishments of this iteration. The next calculation should establish
that inequality, or its separate position, frequency, and hard-band bounds,
on arbitrary native near-top states rather than on another selected packet.

The appended diagnostic checks the exact deficit algebra, normalized Mehler
factor, finite-step saturation, hard-sector decomposition, and the ordering
logic on finite fixtures. The audit distinguishes these tests from the
analytic compactness proofs. Source preservation and PDF rendering are
reported separately. Only this terminal insertion and the literal title
version change the supplied V45; no historical proof has been rewritten.
The current native profile deduction explicitly inherits the V45 strong core
expansion and V42 band localization as stated. The complete lineage has not
been independently recertified, and the next checkpoint remains V50.

\begin{firewall}
V46 proves all-vector confinement and spectral completeness for the declared
finite-step quartic comparison, and a nonzero-profile rigidity statement for
actual native critical eigenvectors. It does not prove native low-energy
tightness or the stated sufficient domination inequality, identify a native
Perron correction or neutral scaling coefficient unconditionally, remove the
fixed spatial regulator, or calibrate physical time. The comparison step $h$
is not a replacement Wilson kernel or a supplied physical clock.
\end{firewall}
% END F16 V46 APPEND

% BEGIN F16 V47 APPEND -- 2026-09-18
% Insert before the final document-ending command of the supplied cumulative V46.
% The predecessor, original Wilson operator, probabilities and clocks are preserved.

\clearpage
\section{V47: native hard-band capture and local soft compactness}
\label{f16v47:context}

V46 leaves possible norm loss in native critical profiles. Here a local
operator-norm comparison captures the hard ground band for every native
near-top vector. On bounded critical soft regions its error is $O_{L,R}(h)$;
the exact nonlocal localization identity then gives bounded comparison energy
for arbitrary localized native eigenvectors. Frequency escape at bounded soft
position is excluded. Only escape to large critical soft positions remains.
The coarser $O_L(r+\varepsilon)$ error on a physical angular strip cannot be
divided by $h$ globally. No radial tightness or native scaling coefficient is
inferred without that remaining estimate.

\subsection{Inputs and the unchanged physical extraction}

V46's appendix and audit were read in full. The native inputs are V44's
measured tube, V45's exact stationarity, and V42's band localization. V46
supplies the nonlocal comparison estimates and the finite-profile equation.
Targeted archive searches, ownership and review depth are recorded in the
audit; they do not independently certify the lineage. The next checkpoint
remains V50.

For primary context, Panati--Spohn--Teufel,
\href{https://arxiv.org/html/math-ph/0201055}{\emph{Space-Adiabatic Perturbation
Theory}}, Section 2, distinguishes an isolated fast band from the dynamics
remaining within it. Its global symbol and gap assumptions are not applied to
the Wilson kernel below. Hanche-Olsen--Holden,
\href{https://arxiv.org/html/0906.4883}{\emph{The Kolmogorov--Riesz compactness
theorem}}, describes the classical position/translation compactness criterion.
Here its Fourier-screen version is already implemented in V46; the new work
is to verify local bounds for native eigenvectors before using that criterion.
No external Navier--Stokes theorem is a premise.

Fix $L=2m$, $m\ge4$, $V=L^3$, and keep exactly
\[
 T_\gamma=M_{b_\gamma}C_\gamma^{\otimes3V}M_{b_\gamma},\qquad
 b_\gamma=e^{-S/(2\varepsilon^2)},\qquad
 \varepsilon=\gamma^{-1/2},\quad h=(\gamma V)^{-1/3},\quad
 \tau=\sqrt Vh.
\]
Thus $\varepsilon^2=Vh^3$. All constants below may depend on the fixed $L$.
Let $s=y+\mathsf K a$ be V44's tube coordinate, with
$y\in\mathcal Y$, $a\in\mathbb R^9$ and $\mathsf K^*\mathsf K=VI$.
The critical dilation is $y=\varepsilon u$, $a=hq$.
The tube measure is the \emph{same} exact measure
\[
 dH(\Phi(g,s))=dg\,J_L(s)\,ds.
\]
Only the based gauge variables are removed. Joint simultaneous color invariance
of functions of $(u,q)$ imposes the residual Gauss condition.

Write $\mathcal X=L^2(\mathcal Y\times\mathbb R^9)$, restricted to that joint
invariant subspace when required. Retain V46's operators
\[
 A_+=\mathcal T_+/\lambda_*,\quad A_+\varphi=\varphi,\quad
 r_+=\|A_+|_{\varphi^\perp}\|<1,\quad
 \mathcal I f=\varphi\otimes f,
\]
\[
 S_h=e^{-h\mathcal Q}e^{h\Delta_q/2}e^{-h\mathcal Q},\qquad
 \mathbb D_h=h^{-1}(I-A_+\otimes S_h),\qquad
 d_h[f]=h^{-1}\langle f,(I-S_h)f\rangle.
\]
Here $\lambda_*$ remains the original Gaussian threshold, not an assigned
native top. No statement that $T_\gamma/\lambda_*\preceq I$ is assumed.

Choose a fixed product tube $|y|<d_0$, $|a|<r_0$ with compact closure inside
V44's tube and the selected central region. Let $\mathcal V_\gamma$ be the
partial isometry from $\mathcal X$ to this one-region physical space:
\[
 (\mathcal V_\gamma F)(\Phi(g,s))
 =J_L(s)^{-1/2}\varepsilon^{-r_Y/2}\tau^{-9/2}F(u,q),
 \qquad r_Y=\dim\mathcal Y,
\]
with zero extension and with the input restricted to that product tube.
Its Gram is exactly the corresponding multiplication projection. Let
$\mathcal V_{\gamma,a}$ be its normalized eight-character folding in sector
$a$. In the rest of this section a sector label on $\mathcal V$ or an
extraction is distinguished from the constant angle $a$ by its subscript.
Fix V46's extraction in these tubes:
\[
 \mathfrak E_{\gamma,a}\psi=\mathcal V_{\gamma,a}^*(\chi\psi),
 \qquad 0\le\chi\le1,
\]
where $\chi$ is fixed, invariant, has disjoint center-translated supports,
and equals one on smaller tubes. Shrinking the auxiliary radii below does not
change this final extraction.

\subsection{Uniform hard tails from the actual magnetic multiplier}

\begin{proposition}[A hard pin in the central tube and a whole-band moment bound]
\label{f16v47:hard-tail}
After decreasing $d_0,r_0$, there is $c_L>0$ such that
\begin{equation}
 S(\operatorname{Exp}(y+\mathsf K a))
       \ge S(\operatorname{Exp}(\mathsf K a))+c_L|y|^2
       \ge c_L|y|^2.
 \label{f16v47:hard-pin}
\end{equation}
For any fixed $c>0$, every spectral band
$\mathsf P_\gamma\le1_{[c,1]}(T_\gamma|_{\mathcal H_a})$ satisfies
\begin{equation}
 \big\|e^{\beta_L|u|^2}\mathfrak E_{\gamma,a}\mathsf P_\gamma\big\|
       \le c^{-1}
 \label{f16v47:hard-moment}
\end{equation}
for some $\beta_L>0$, independent of $\gamma$ and of the rank of the band.
The same bound holds for smaller supported cutoffs. It concerns rescaled hard
coordinates only; no critical soft moment is asserted.
\end{proposition}
\begin{proof}
V45's exact mean-zero stationarity gives $D_yS(\mathsf K a)=0$. At the origin,
$D_y^2S=H_+\succ0$. Smoothness on a fixed small product neighborhood makes this
Hessian uniformly positive along every segment from $0$ to $y$, also for
noncommuting constant triples. Taylor's integral formula proves
\eqref{f16v47:hard-pin}. Center translation preserves it.

On the band, write $\mathsf P_\gamma=T_\gamma\mathsf P_\gamma
(T_\gamma|_{\operatorname{Ran}\mathsf P_\gamma})^{-1}$.
The inverse has norm at most $c^{-1}$. In the physical tube, multiplication by
$\chi e^{\beta_L|y|^2/\varepsilon^2}b_\gamma$ has norm at most one when
$\beta_L\le c_L/2$. The remaining two factors $C_\gamma^{\otimes3V}$ and
$b_\gamma$ are contractions. The exact extraction isometry and center folding
therefore give \eqref{f16v47:hard-moment}. This is an operator estimate on the
whole band, not a sum of individual eigenvector bounds.
\end{proof}

\subsection{The reduced kinetic kernel, with its orbit normalization}

Let $d=\dim\mathcal N=6V+3$ and $D=9V$. On the fixed tube, the flat
half-density kernel of the \emph{raw} kinetic compression is exactly
\begin{equation}
 k_\varepsilon^{\rm red}(s,t)
  =J_L(s)^{1/2}J_L(t)^{1/2}
    \int_{\mathcal G_o}
       c_\gamma(\operatorname{Exp}s,g\cdot\operatorname{Exp}t)\,dg.
 \label{f16v47:quotient-kernel}
\end{equation}
This follows by disintegrating the original input Haar integral; $c_\gamma$
contains all original kinetic normalizers. The based integral is not a new
conditional transition. Formula \eqref{f16v47:quotient-kernel} also acts on
non-invariant quotient tests before restricting to joint color invariants.

\begin{proposition}[A normalized local kernel comparison with an integrable error]
\label{f16v47:kinetic}
In a sufficiently small fixed tube there are constants $C_L,c_L>0$ and an
integer $M_L$ such that, with
$g_\varepsilon(s-t)=(2\pi\varepsilon^2)^{-d/2}
 e^{-|s-t|^2/(2\varepsilon^2)}$,
\begin{equation}
 \begin{split}
 |k_\varepsilon^{\rm red}(s,t)-g_\varepsilon(s-t)|
 &\le C_L(|s|+|t|+\varepsilon)\varepsilon^{-d}
       (1+|s-t|/\varepsilon)^{M_L}e^{-c_L|s-t|^2/\varepsilon^2}\\
 &\quad+C_L\varepsilon^{-D}e^{-c_L/\varepsilon^2}.
 \end{split}
 \label{f16v47:kernel-error}
\end{equation}
The constants are uniform as a smaller soft angular strip shrinks inside this
tube. Both row and column integrals of the Gaussian-polynomial factor are
bounded independently of $\varepsilon$.
\end{proposition}
\begin{proof}
Choose local based-group coordinates $v\in\mathbb R^{D-d}$ whose orbit
differential at the identity is an isometry into $\operatorname{Ran}d_0$.
The remaining differential is the orthogonal inclusion of $\mathcal N$.
Let $c_X$ be the full Haar density in an orthonormal fine-link chart and let
$c_G$ be the based Haar density in these orbit-orthonormal coordinates.
The exact tube Jacobian gives $J_L(0)c_G=c_X$.
The normalized raw kernel has leading constant
$c_X^{-1}(2\pi\varepsilon^2)^{-D/2}$. Therefore integration of its vertical
Gaussian gives exactly
\[
 J_L(0)c_Gc_X^{-1}(2\pi\varepsilon^2)^{-D/2}
       \int e^{-|v|^2/(2\varepsilon^2)}dv
       =(2\pi\varepsilon^2)^{-d/2}.
\]
All orbit metric determinants are included in $c_G$ and $J_L(0)$; none is
replaced by one separately.

Put $t=s+\xi$ and let $A(s,\xi,v)$ be the sum of the actual displacement
cosines. Smoothness, its zero and zero first derivative at $(\xi,v)=0$, and
orthogonality at $s=0$ give
\[
 A(s,\xi,v)=\tfrac12(|\xi|^2+|v|^2)+\mathcal R,
 \quad
 |\mathcal R|\le C_L\{|s|(|\xi|^2+|v|^2)+(|\xi|+|v|)^3\}.
\]
After one fixed shrink, $A\ge c_L(|\xi|^2+|v|^2)$ there. The exact smooth
volume factors differ from their value at zero by
$O_L(|s|+|t|+|v|)$, and the original normalization has relative error
$O_L(\varepsilon^2)$. The mean-value estimate for two exponentials, using the
positive quadratic lower bounds for both of them, bounds their difference by
$|\mathcal R|/\varepsilon^2$ times a Gaussian majorant. Integrate in $v$;
each cubic vertical moment contributes one factor $\varepsilon$ and a fixed
polynomial in $|\xi|/\varepsilon$. This proves the first term in
\eqref{f16v47:kernel-error}.

Outside the selected based-group coordinate neighborhood, freeness of the
based action and compactness give a fixed positive displacement from the
identity orbit representative. This stays positive for all small $s,t$.
The original cosine bound then gives the last term. A fixed finite number of
charts bounds its integral. The row and column assertions follow by the
change $\xi=\varepsilon z$; all integrations are over a fixed bounded tube.
This proves an absolute kernel error and its Schur bounds, not a Loewner
ordering from pointwise kernel comparisons.
\end{proof}

\subsection{An all-vector strip comparison, including shrinking critical strips}

Let $\Pi_{\gamma,r}$ be multiplication in profile space by the strip
$|\varepsilon u|<d_0$, $|hq|<r$, with $0<r\le r_0$. Let
$K_{\gamma,a}^{\rm red}=\mathcal V_{\gamma,a}^*T_\gamma\mathcal V_{\gamma,a}$.

\begin{theorem}[Native strip norm comparison before dividing by the scale]
\label{f16v47:strip}
There are $C_L<\infty$ and a fixed coupling threshold such that, uniformly for
$0<r\le r_0$ and all center sectors,
\begin{equation}
 \boxed{\quad
 \big\|\Pi_{\gamma,r}
       (K_{\gamma,a}^{\rm red}-\mathcal T_+\otimes S_h)
          \Pi_{\gamma,r}\big\|
       \le C_L(r+\varepsilon).
 \quad}
 \label{f16v47:strip-bound}
\end{equation}
The norm is on the complete supported $L^2$ space, not a chosen test core.
In particular, on $|q|\le2R$, with fixed $R$ and $2hR\le r_0$, the error is
at most $C_L(2hR+\varepsilon)$.
\end{theorem}
\begin{proof}
Write $S_c(a)=S(\operatorname{Exp}(\mathsf K a))$. Exact stationarity and
bounded third derivatives give
\[
 S(\operatorname{Exp}(y+\mathsf K a))
 =S_c(a)+\tfrac12 y^{\mathsf T}H_+y+\mathcal R_m,
 \qquad |\mathcal R_m|\le C_L(|a||y|^2+|y|^3).
\]
Both the actual exponent and the split exponent are bounded below by
$S_c(a)+c_L|y|^2$ on the chosen tube. Their half-multiplier difference hence
has supremum at most $C_L(r+\varepsilon)$ on the strip: set $y=\varepsilon u$
and use the boundedness of each polynomial times $e^{-c_L|u|^2}$.
The constant-field identity gives, with $Q(a)=\mathcal Q(a)$,
\[
 (1-C_Lr^2)Q(a)
 \le\sum_{i<j}|v(a_i)\times v(a_j)|^2\le Q(a),
 \qquad v(a)=\frac{\sin|a|}{|a|}a.
\]
Since $\gamma Vh^4=h$, the supremum difference between
$e^{-S_c(hq)/(2\varepsilon^2)}$ and $e^{-h\mathcal Q(q)}$ is at most $C_Lr^2$.
This bound is uniform even when $h\mathcal Q(q)$ is large: use
$xe^{-c x}\le C$ instead of bounding the exponent difference alone.

Now sandwich \eqref{f16v47:kernel-error} by the actual magnetic factors.
On the strip, $|s|+|t|\le |y|+|y'|+C_Lr$. The hard pin gives
$|y|b_\gamma(s)\le C_L\varepsilon$ and the same estimate at $t$.
The row and column Gaussian-polynomial integrals therefore bound the weighted
kinetic error by $C_L(r+\varepsilon)$. The remote based-group term is
exponentially small. Comparing next the two magnetic factors with their split
versions uses only their supremum bounds and contraction of the flat Gaussian
convolution. Thus the complete difference has the claimed norm.

Finally dilate by $y=\varepsilon u$, $a=hq$. The hard convolution and its two
quadratic multipliers give exactly $\mathcal T_+$. The orthonormal soft step
has covariance $\varepsilon^2/\tau^2=h$; its two multipliers give exactly
$S_h$. These are comparisons after assembling the original measure and
kernel, not replacements of the original transfer. Center folding adds only
cross terms between fixed separated tubes, of norm $O_L(e^{-c_L\gamma})$.
They are absorbed in the displayed error. No derivative of an arbitrary
supported test is used.
\end{proof}

The estimate is deliberately coarse. The first hard variation is charged in
operator norm, so it costs $O(r)$, even though its ground expectation vanishes.
On a critical strip $r=2hR$ this is $O(hR)$, which suffices for local energy
bounds. It does not yield V46's global $O(h)$ domination as $R\to\infty$.

\subsection{Capture of the hard ground band for every near-top vector}

Fix $K\ge0$ and put
\[
 \mathsf B_{\gamma,a}(K)
       =1_{[\lambda_*(1-Kh),\,1]}(T_\gamma|_{\mathcal H_a}),\qquad
 F_\gamma(\psi)=\mathfrak E_{\gamma,a}\psi,\quad
 f_\gamma(\psi)=\mathcal I^*F_\gamma(\psi).
\]
A zero band gives zero operator norms. The native top is not assumed to be
below $\lambda_*$; negative scaled deficits are included in this band.

\begin{theorem}[Whole-band hard capture and unscaled soft-step regularity]
\label{f16v47:capture}
At fixed $L,K$, uniformly over unit vectors in $\operatorname{Ran}\mathsf B_{\gamma,a}(K)$,
\begin{equation}
 \boxed{\quad
 \|F_\gamma-\mathcal I f_\gamma\|\longrightarrow0,\qquad
 \|f_\gamma\|\longrightarrow1,\qquad
 h\,d_h[f_\gamma]\longrightarrow0.
 \quad}
 \label{f16v47:capture-result}
\end{equation}
Equivalently the first assertion is convergence to zero of
$\|(I-\mathcal I\mathcal I^*)\mathfrak E_{\gamma,a}\mathsf B_{\gamma,a}(K)\|$.
There is no spectral-rank factor. No $O(h)$ bound for the squared hard loss,
or boundedness of $d_h[f_\gamma]$, is asserted globally.
\end{theorem}
\begin{proof}
Choose, for each fixed small $r>0$, an invariant center-translated cutoff
$\chi_r$ supported in $|y|<d_0$, $|a|<r$, equal to one near the central orbits,
and with support inside the plateau of $\chi$. V42 applied to still smaller
fixed neighborhoods gives, eventually,
\[
 \|(1-\chi_r)\mathsf B_{\gamma,a}(K)\|
       \le C_{L,r}e^{-c_{L,r}\sqrt\gamma}=:t_{\gamma,r}.
\]
This uses $\lambda_*(1-Kh)\to\lambda_*$, so the shrinking spectral band lies
above the fixed neighborhood's cutoff. It does not use a shrinking spatial
neighborhood estimate with uniform constants.
For a unit band vector let $F_r=\mathcal V_{\gamma,a}^*(\chi_r\psi)$.
The extraction norm differs from one by at most $O(t_{\gamma,r})$, and
$\|F_r-F_\gamma\|\le t_{\gamma,r}$. Changing its transfer quotient to the
uncut one costs at most $2t_{\gamma,r}$ because $\|T_\gamma\|\le1$.
By \eqref{f16v47:strip-bound} and the lower spectral endpoint,
\[
 0\le\langle F_r,(I-A_+\otimes S_h)F_r\rangle
 \le Kh+C_L(r+\varepsilon+t_{\gamma,r}).
\]
The bounded reference form changes by $O(t_{\gamma,r})$ when $F_r$ is replaced
by $F_\gamma$. Its exact hard separation is
\[
 \langle F_\gamma,(I-A_+\otimes S_h)F_\gamma\rangle
 \ge(1-r_+)\|F_\gamma-\mathcal I f_\gamma\|^2
       +h d_h[f_\gamma].
\]
First let $\gamma\to\infty$ at each fixed $r$, and then $r\downarrow0$.
This proves both vanishing terms, uniformly on the band. V46's extraction-mass
bound, followed by Pythagoras, proves $\|f_\gamma\|\to1$.
\end{proof}

For normalized band sequences the extracted joint probability consequently
factorizes in the following precise sense:
\[
 \big\||F_\gamma(u,q)|^2-|\varphi(u)|^2|f_\gamma(q)|^2\big\|_{L^1(du\,dq)}
       \longrightarrow0.
\]
This is convergence of a joint density after the specified extraction. It is
not a claim that the compact native measure factors exactly, nor a control of
the remaining soft density. The exponential hard moment bound additionally
prevents hard-position mass from escaping behind this norm limit.

\subsection{A localized native energy bound at order \texorpdfstring{$h$}{h}}

Let $\vartheta\in C_c^\infty(\mathbb R^9)$ be real, radial, equal to one on
the unit ball and zero outside the ball of radius two, with $0\le\vartheta\le1$.
For $R\ge1$ put $\vartheta_R(q)=\vartheta(q/R)$. On each tube define the
physical multiplier
\[
 \chi_{\gamma,R}(U)=\chi(U)\vartheta\bigl(a(U)/(hR)\bigr),
\]
and extend by zero and center translation. It is gauge and center invariant.
Its Lipschitz constant in the original compact metric is at most
$C_L[1+(hR)^{-1}]$. The corresponding extracted vector is exactly
$\vartheta_R F_\gamma$.

\begin{theorem}[Local comparison energies for arbitrary native eigenvectors]
\label{f16v47:local-energy}
Let $\|\psi_\gamma\|=1$ and
\[
 T_\gamma\psi_\gamma=\lambda_*(1-hE_\gamma)\psi_\gamma,
 \qquad E_\gamma\le K.
\]
There is no lower bound assumed for $E_\gamma$. For every fixed $R\ge1$, at
sufficiently large coupling,
\begin{equation}
 \boxed{\quad
 \langle\vartheta_RF_\gamma,\mathbb D_h\vartheta_RF_\gamma\rangle
 \le [E_\gamma+C_L(R+1)]\|\vartheta_RF_\gamma\|^2
                       +C_L(R^{-2}+h^2).
 \quad}
 \label{f16v47:local-energy-sharp}
\end{equation}
In particular the left side is bounded above by a finite $C_{L,K,R}$.
It follows that
\begin{align}
 \|(I-\mathcal I\mathcal I^*)\vartheta_RF_\gamma\|^2
      &\le \frac{h C_{L,K,R}}{1-r_+},\\
 d_h[\vartheta_Rf_\gamma]&\le C_{L,K,R},
 \label{f16v47:local-consequences}\\
 \int_{|p|>P}|\widehat{\vartheta_R f_\gamma}(p)|^2dp
      &\le C_{L,K,R}(2h+8P^{-2}),\qquad hP^2\le1.
 \label{f16v47:local-frequency}
\end{align}
Thus every such sequence has a subsequence with strongly convergent soft
profiles on every bounded $q$-region. High-frequency loss at bounded soft
position is excluded. None of the constants is asserted bounded as $R\to\infty$.
\end{theorem}
\begin{proof}
For a real Lipschitz multiplier $b$, symmetry of the actual kernel gives
\[
 \langle b\psi,(\lambda_*-T_\gamma)b\psi\rangle
 =\lambda_*hE_\gamma\|b\psi\|^2+
 \frac12\iint(b(U)-b(V))^2\overline{\psi(U)}
             t_\gamma(U,V)\psi(V)\,dH(U)dH(V).
\]
The last integral is real but need not be nonnegative for a complex test.
Its absolute value is at most $C_L\varepsilon^2\operatorname{Lip}(b)^2$
by the actual displacement second moment and a row/column Schur estimate.
For $b=\chi_{\gamma,R}$ this is at most
$C_L(h/R^2+h^3)$ because $\varepsilon^2=Vh^3$.
This is a localization estimate for the bounded native kernel, not a diffusion
Dirichlet identity.

The localized vector is supported in $|a|\le2hR$, so
Theorem~\ref{f16v47:strip} bounds its quotient difference from the tensor
comparison by $C_L(2hR+\varepsilon)\|b\psi\|^2$.
The exact extraction preserves $\|b\psi\|$ and its quadratic form in the
compressed tube. Divide the resulting inequality by $\lambda_*h$, and use
$\varepsilon/h=\sqrt V\sqrt h\le C_L$ eventually. This proves
\eqref{f16v47:local-energy-sharp}, including its dependence on $E_\gamma$.
Its uniform upper bound uses only $E_\gamma\le K$ and $\|b\psi\|\le1$.

V46's exact hard separation and Fourier estimate give the next displays.
The position support of $\vartheta_R f_\gamma$ is bounded; the Fourier bound
makes it precompact by V46's Hilbert--Schmidt phase-space screen. Alternatively
apply its nonlocal compactness theorem. A diagonal extraction over integer
$R$ gives strong local convergence, with consistent limits on overlapping
balls. No derivative of $\psi_\gamma$ was required.
\end{proof}

If $E_\gamma\to-\infty$, nonnegativity of the left side of
\eqref{f16v47:local-energy-sharp} gives the additional quantitative fact
\begin{equation}
 \|\vartheta_RF_\gamma\|^2
 \le\frac{C_L(R^{-2}+h^2)}{|E_\gamma|-C_L(R+1)}
 \longrightarrow0
 \label{f16v47:negative-escape}
\end{equation}
for each fixed $R$, once the denominator is positive. Such a branch, if present,
must therefore escape in soft position; hard or purely local frequency escape
cannot account for it.

\subsection{Only soft position remains in the compactness alternative}

\begin{proposition}[Native profile decomposition in local norm]
\label{f16v47:profiles}
For the eigenvector sequences of Theorem~\ref{f16v47:local-energy}, a subsequence
has
\[
 F_\gamma-\mathcal I f_\gamma\to0\quad\hbox{in }L^2,
 \qquad f_\gamma\to f\quad\hbox{locally in }L^2,
 \qquad \|f_\gamma\|\to1,\quad \|f\|\le1.
\]
If $E_\gamma\to E\in\mathbb R$, V46's corrected profile equation gives
$H_{\rm mat}f=Ef$. If $E_\gamma\to-\infty$, then $f=0$.
The convergence is global and norm preserving precisely when this sequence
is tight in soft position, namely
\begin{equation}
 \lim_{R\to\infty}\limsup_{\gamma\to\infty}
             \int_{|q|>R}|f_\gamma(q)|^2dq=0.
 \label{f16v47:position-tightness}
\end{equation}
No separate Fourier-tightness hypothesis is needed.
\end{proposition}
\begin{proof}
The whole-band theorem applies since $E_\gamma\le K$; it supplies hard capture
and the total soft norm. The local-energy theorem supplies local compactness.
For finite $E$, its limit is also the weak extracted limit, so V46's theorem
applies with its original domain justification. The negative case is
\eqref{f16v47:negative-escape}. Local strong convergence and
\eqref{f16v47:position-tightness} give global strong convergence by a cutoff
and triangle inequality. Conversely, a strongly convergent $L^2$ sequence has
uniformly small tails, and its norm here tends to one. This proves the claim.
\end{proof}

\begin{proposition}[A position-only completeness criterion and an escape witness]
\label{f16v47:completion}
Suppose \eqref{f16v47:position-tightness} holds uniformly, for each fixed $K$,
over the normalized neutral native eigenvectors with $E_\gamma\le K$.
Then the ordered scaling and same-top conclusions in
Proposition~\ref{f16v46:conditional-completeness} follow. This position-tightness
premise is \emph{not proved here}.
In particular, failure of those ordered asymptotics yields, after subsequence
selection, normalized native eigenvectors with $E_\gamma\le K$, radii
$R_j\to\infty$ and $\eta>0$ such that
\begin{equation}
 \int_{|q|>R_j}|f_{\gamma_j}(q)|^2dq\ge\eta.
 \label{f16v47:escape-witness}
\end{equation}
Their hard-ground loss tends to zero, and their soft profiles are compact on
every bounded region. The witness is a statement about actual eigenvectors,
not evidence that such an escaping branch exists.
\end{proposition}
\begin{proof}
V45 bounds above the scaled deficits of every fixed number of ordered neutral
eigenvectors. Position tightness and Proposition~\ref{f16v47:profiles} make
those profiles strongly precompact. A sequence of deficits tending to
$-\infty$ would force every local profile to vanish and contradict norm
preservation. Thus the deficits are also bounded below along those sequences.
Take a common convergent subsequence of any finite orthonormal eigenbasis.
Extraction and hard capture preserve its entire Gram, so the limiting soft
profiles are orthonormal model eigenfunctions. V46's profile equation and
model min--max give the missing ordered liminf; V45 supplies the limsup.
V43 transfers the neutral conclusion to every center sector with its $o(h)$
ladder error, and its exact return metric supplies the neutral $g_\gamma$
conclusion. This is the same spectral implication as V46, now with a single
position-tightness premise instead of its stronger unproved global form
comparison. The contrapositive and the failure of the displayed tail condition
give \eqref{f16v47:escape-witness} by a diagonal choice of radii.
\end{proof}

\paragraph{Why local compactness is not enough.}
Even in the declared comparison a normalized invariant annular profile can be
$f_h(q)=R_h^{-9/2}f(q/R_h)$, where $f$ is smooth radial, supported in
$1<|q|<2$, and $R_h=h^{-1/8}$. Then $f_h\rightharpoonup0$ and all mass escapes
in position, although $\langle f_h,(I-S_h)f_h\rangle\to0$. Indeed the positive
energy identity and $1-e^{-x}\le x$ bound this last quantity by
\[
 C_f\{hR_h^4+hR_h^{-2}+h^3R_h^6\}\longrightarrow0.
\]
This is not a native eigenvector or a counterexample to the remaining estimate.
It shows why the proved unscaled step regularity cannot itself be divided by
$h$ or used as position tightness.

\subsection{Remaining task and verification}

The new native results remove hard-band norm loss and frequency loss at bounded
critical positions from the completeness problem. They do not bound the radial
tail in \eqref{f16v47:position-tightness}. The strip estimate pays a first hard
variation by its absolute norm, giving $C_LR$ in the localized energy; its
constant cannot be sent to infinity. A useful next bound must exploit the
singlet cancellation or a controlled hard-band return on arbitrary states,
while retaining transverse soft confinement along the commuting valleys.
Another fixed finite recovery space would not supply this radial estimate.
This keeps the supplied memoranda's no-vanishing requirement distinct from
residual improvement and from physical transfer calibration.

The diagnostic, audit and reconstruction report accompany the source. Only
this appendix and the literal title version differ from V46; removing them
recovers its exact bytes. Finite checks do not certify the general analytic
proofs or the inherited lineage. The next checkpoint remains V50.

\begin{firewall}
V47 proves a native strip operator-norm comparison, whole-band hard-ground
capture, and local compactness of critical native eigenprofiles. It reduces
the remaining completeness problem to escape in soft position. It does not
prove that radial tightness, a global lower bound at scale $h$, the native
Perron or neutral-gap coefficient, a growing-volume theorem, a physical-time
calibration, or an interacting continuum mass gap. The original transfer,
Gauss constraint, center sectors and returning-state norm remain unchanged.
\end{firewall}
% END F16 V47 APPEND

% BEGIN F16 V48 APPEND -- 2026-09-18
% Insert before the final document-ending command of cumulative V47.
% This insertion leaves every predecessor body and the native operator unchanged.

\clearpage
\section{V48: quadratic hard-return costs and fixed-volume spectral completeness}
\label{f16v48:context}

V47 reduces possible loss of critical spectral mass to escape in soft position.
Its strip estimate pays the linear hard variation in absolute operator norm,
which gives an error proportional to the angular radius. That error competes
with, rather than improves on, transverse confinement. We return to the exact
quotient metric and keep the first magnetic variation as an operator. The
quotient metric has zero first jet in the already chosen exponential normal
slice. The remaining linear operator has zero hard-ground compression. Its
coupling to the separated hard complement costs quadratically in the radius.
On a soft annulus the finite-step transverse confinement is linear in that
radius, so it dominates this quadratic return cost.

A dyadic localization pays the original kinetic cutoff errors without a factor
counting annuli. It supplies the position tightness missing in V47, and hence
the fixed-spatial-regulator ordered spectral limit. The proof does not replace
the Wilson transfer by its comparison, identify the asymptotic parameter with
physical time, or pass to growing spatial volume. Its new analytic inputs are
the second-order quotient-kernel estimate and the complete strip first jet;
the concluding spectral identification also uses V45's corrected core equation
and V46--V47's proved-as-stated profile results.

\subsection{Dependencies and unchanged notation}

The supplied V47 appendix and audit were read in full. V44's exact measured
based-gauge tube, V45's stationarity and hard derivative, V46's finite-step
valley bound, and V42's fixed-neighborhood forbidden norm are retained.
Targeted searches of the current cumulative source and the supplied f14/f15
and Grand Tensor archives were recorded; they are not an exhaustive novelty
or proof audit. The two import memoranda remain methodological proposals, not
premises supplying any missing estimate. The next scheduled checkpoint is V50.

Normal-slice geometry, quadratic elimination of an off-diagonal perturbation,
and compactness plus min--max are classical mechanisms. For context,
M\'atyus--Teufel, \href{https://arxiv.org/html/1904.00042}{\emph{Effective
non-adiabatic Hamiltonians}}, Section III, distinguishes diagonal compression
from the effect of the omitted coupling. Hanche-Olsen--Holden,
\href{https://arxiv.org/html/0906.4883}{\emph{The Kolmogorov--Riesz compactness
theorem}}, Corollary 7, gives the position/Fourier criterion already used in
V46--V47. No molecular gap assumption or quantitative compactness rate is
imported. The metric and kernel calculations needed here are given directly
for the native based action.

Fix $L=2m$, $m\ge4$, and $V=L^3$. Keep
\[
 T=T_\gamma=M_{b_\gamma}C_\gamma^{\otimes3V}M_{b_\gamma},\qquad
 b_\gamma=e^{-S/(2\varepsilon^2)},\qquad
 \varepsilon=\gamma^{-1/2},\quad h=(\gamma V)^{-1/3},\quad
 \varepsilon^2=Vh^3.
\]
All constants may depend on this fixed $L$. Write $\lambda_*=\Lambda_{{\rm G},B}$,
not an assigned native top. In the unchanged tube put
\[
 s=y+\mathsf K a,\quad y\in\mathcal Y,\quad a\in\mathbb R^9,\quad
 \mathsf K^*\mathsf K=VI,\quad y=\varepsilon u,\quad a=hq,
 \qquad dH=dg\,J_L(s)\,ds.
\]
Only the based action is removed. Residual simultaneous conjugation is imposed
on functions, not eliminated by a singular orientation slice. Let
\[
 \mathfrak P=|\varphi\rangle\langle\varphi|\otimes I,\quad
 \mathcal T_+\varphi=\lambda_*\varphi,\quad
 \kappa_+=\lambda_*\bigl(1-r_+\bigr)>0,
\]
where $r_+=\|\mathcal T_+|_{\varphi^\perp}\|/\lambda_*$. Retain
\[
 S_h=e^{-h\mathcal Q}e^{h\Delta_q/2}e^{-h\mathcal Q},\qquad
 \mathcal Q(q)=\sum_{i<j}|q_i\times q_j|^2.
\]
This is V46's explicit comparison. It is never declared to be a native
conditional transition. Write $K_\gamma^{\rm red}$ for the one-tube compression
in the exact flat half-density and critical coordinates, extended by zero in
the hard variable outside its fixed physical tube. Sector-folded compressions
have the same estimates below, with exponentially small cross-tube terms.

\subsection{The quotient metric has no linear error}

The following calculation strengthens the first-order bound in V47 without
changing its coordinates. Let $X_s$ be the differential of
$s\mapsto\operatorname{Exp}(s)$ on $\mathcal N$, and let $Y_s$ be the differential
of the based action at that point, using a fixed basis of its Lie algebra.
Inner products in the target are the original product round metric. Set
\begin{equation}
 V_s=Y_s^*Y_s,\qquad C_s=X_s^*Y_s,\qquad
 \mathsf g_s=X_s^*X_s-C_sV_s^{-1}C_s^*.
 \label{f16v48:quotient-metric}
\end{equation}
The letter $V_s$ here is a vertical Gram matrix, not the number of vertices.

\begin{proposition}[Second-order metric and the normalized quotient kernel]
\label{f16v48:metric}
On one sufficiently small fixed tube,
\begin{equation}
 \mathsf g_0=I,\qquad D\mathsf g_0=0,\qquad
 \|\mathsf g_s-I\|+\|\mathsf g_s^{-1}-I\|\le C_L|s|^2.
 \label{f16v48:metric-jet}
\end{equation}
For the actual quotient kernel of V47 there are fixed $C_L,c_L>0$ and $M_L$
such that
\begin{align}
 |k_\varepsilon^{\rm red}(s,t)-g_\varepsilon(s-t)|
 &\le C_L(|s|^2+|t|^2+\varepsilon)\varepsilon^{-d}
 (1+|s-t|/\varepsilon)^{M_L}e^{-c_L|s-t|^2/\varepsilon^2}
 \nonumber\\
 &\quad+C_L\varepsilon^{-D}e^{-c_L/\varepsilon^2},
 \label{f16v48:kinetic-error}
\end{align}
where $d=6V+3$, $D=9V$ and
$g_\varepsilon(\xi)=(2\pi\varepsilon^2)^{-d/2}e^{-|\xi|^2/(2\varepsilon^2)}$.
Both Gaussian-polynomial Schur integrals are uniformly bounded. In V45's
specified normal slice the complete horizontal covariance is
$\mathsf g_s^{-1}$; consequently its hard block satisfies $DA(0)=0$.
This last identity refines, rather than assumes away, the covariance derivative
that V45 retained.
\end{proposition}
\begin{proof}
Based freeness gives $V_0\succ0$ and a uniform inverse on a small fixed tube.
At the identity the slice is orthogonal to the orbit, so $C_0=0$. The product
exponential is a Riemannian normal chart at the identity. More explicitly, on
one round $\SU(2)$ factor its pulled-back metric is one in the radial direction
and $(\sin|s|/|s|)^2$ in the two transverse directions. Thus
$X_s^*X_s=I+O_L(|s|^2)$, including after restriction to $\mathcal N$.
Smoothness gives $C_s=O_L(|s|)$. Equation \eqref{f16v48:quotient-metric}
proves \eqref{f16v48:metric-jet}. The covariance assertion follows because the
differential of the quotient is an isometry on the horizontal tangent and
annihilates the vertical tangent. Equivalently its product with its adjoint is
the inverse Schur metric. No freeness of the residual global conjugation is
used.

Here are the measure details needed to improve a kernel estimate as well.
Take the same fixed local based-group coordinates $v$ as in V47, with Haar
density $c_G$ at the identity, and let $c_X$ be the full probability-Haar density
in an orthonormal tangent frame. The exact Jacobian identity at $(1,s)$ is
\begin{equation}
 J_L(s)c_G=c_X\sqrt{\det V_s}\sqrt{\det\mathsf g_s}.
 \label{f16v48:coarea}
\end{equation}
It is the determinant of the Gram matrix of $[Y_s\ X_s]$, followed by its
Schur determinant. This retains the varying orbit volume; $J_L(s)$ is not
assumed flat or constant.

Put $t=s+\xi$. Taylor expansion of the actual Wilson displacement cost in
$(\xi,v)$ gives its exact quadratic part
\[
 A_s(\xi,v)=\tfrac12\|X_s\xi+Y_sv\|^2+
ho_s(\xi,v),\qquad
 |\rho_s(\xi,v)|\le C_L(|\xi|+|v|)^3.
\]
A sign change of $v$ would give the alternative convention. On the selected
small chart both this quadratic form and the exact cost are bounded below by
$c_L(|\xi|^2+|v|^2)$. All constants are uniform in $s$. Complete the vertical
square by $v\mapsto v+V_s^{-1}C_s^*\xi$. Its integral has determinant
$(\det V_s)^{-1/2}$ and horizontal quadratic form $\xi^*\mathsf g_s\xi$.
Use \eqref{f16v48:coarea} and the two exact half-densities. The resulting
principal kernel is
\[
 \sqrt{J_L(t)/J_L(s)}\sqrt{\det\mathsf g_s}\,
 (2\pi\varepsilon^2)^{-d/2}
 e^{-\xi^*\mathsf g_s\xi/(2\varepsilon^2)}.
\]
Thus neither a vertical determinant nor an endpoint density remains as an
uncomputed order-$|s|$ scalar. The density ratio differs from one by
$O_L(|\xi|)$. Smooth based Haar factors differ from $c_G$ by $O_L(|v|)$.
The true kinetic normalizer has relative error $O_L(\varepsilon^2)$.
The cubic cost remainder divided by $\varepsilon^2$, integrated against a
positive Gaussian majorant, gives $\varepsilon$ times a fixed polynomial in
$|\xi|/\varepsilon$. The same conclusion follows directly from the mean-value
identity for two exponentials; no small supremum of their extensive exponents
is required.

Finally \eqref{f16v48:metric-jet} bounds the difference of the principal
Gaussian from $g_\varepsilon$ by $C_L|s|^2$ times such a polynomial majorant.
For $s,t$ in the fixed tube and $v$ outside the chosen based chart, compactness
and freeness give the same positive displacement floor as V47; its entire
contribution is the last term of \eqref{f16v48:kinetic-error}. Integrating the
majorants in either endpoint proves their Schur bounds. This proves an
absolute kernel estimate, not a Loewner inference from entrywise signs.
\end{proof}

\subsection{Keep the linear hard operator until its return is bounded}

Write $H(a)=H_++\sum_{\ell=1}^9a_\ell H_\ell+O_L(|a|^2)$ for the actual hard
Hessian in V45. Define hard multiplication operators
\[
 B_\ell(u)=\tfrac14u^*H_\ell u.
\]
The products $B_\ell\mathcal T_+$ and $\mathcal T_+B_\ell$ are bounded
Hilbert--Schmidt operators. Let $M_\ell$ mean multiplication by $hq_\ell$ in
the soft variable. On a supported angular strip define the bounded first jet
\begin{equation}
 \mathcal L_\gamma=-\sum_{\ell=1}^9
 \left[B_\ell\mathcal T_+\otimes(M_\ell S_h)
       +\mathcal T_+B_\ell\otimes(S_hM_\ell)\right].
 \label{f16v48:linear-operator}
\end{equation}
Products with the strip projection on both sides are understood; this avoids
asserting boundedness of $M_\ell$ on the whole soft space.

\begin{theorem}[Complete strip first jet with a quadratic remainder]
\label{f16v48:strip-jet}
Let $P_r$ multiply by $|hq|\le r$ on profile space, with no hard-coordinate
cutoff. There is a fixed $r_*>0$ such that, for $0<r\le4r_*$ and small
$\varepsilon$,
\begin{equation}
 \left\|P_r(K_\gamma^{\rm red}-\mathcal T_+\otimes S_h
                    -\mathcal L_\gamma)P_r\right\|
                 \le C_L(r^2+\varepsilon).
 \label{f16v48:quadratic-remainder}
\end{equation}
Moreover
\begin{equation}
 \|P_r\mathcal L_\gamma P_r\|\le C_Lr,\qquad
 \mathfrak P\mathcal L_\gamma\mathfrak P=0
 \label{f16v48:linear-cancellation}
\end{equation}
on every such strip. These are all-vector operator estimates, not core-wise
expansions or assertions that the off-diagonal first jet vanishes.
\end{theorem}
\begin{proof}
Exact hard stationarity and the uniform hard pin give, for $y=\varepsilon u$,
\[
 b_\gamma(y+\mathsf Ka)
 =m_{\gamma,c}(a)e^{-u^*H_+u/4}
       \left(1-\sum_\ell a_\ell B_\ell(u)\right)
       +\mathcal E_\gamma(a,u),
\]
where $m_{\gamma,c}(a)=e^{-S(\operatorname{Exp}(\mathsf Ka))/(2\varepsilon^2)}$
and
\[
 |\mathcal E_\gamma(a,u)|
 \le C_L(r^2+\varepsilon)m_{\gamma,c}(a)
                    (1+|u|)^M e^{-c_L|u|^2},\qquad |a|\le r.
\]
To obtain this bound on the whole hard tube, first compare the exact multiplier
with the one using $H(a)$, charging the cubic hard remainder by
$C_L\varepsilon|u|^3e^{-c_L|u|^2}$. Expand the latter Gaussian twice in $a$;
uniformly positive intermediate Hessians bound its second derivatives by a
polynomial times a Gaussian. This yields the $r^2$ term without requiring
$|u|$ bounded. The omitted hard exterior has exponentially small Gaussian
operator tails at the fixed physical hard radius.

The exact constant commutator formula gives the supremum multiplier bound
\[
 \sup_{|a|\le r}
 |m_{\gamma,c}(a)-e^{-h\mathcal Q(a/h)}|\le C_Lr^2.
\]
Indeed its action lies between $(1-C_Lr^2)\mathcal Q(a)$ and
$\mathcal Q(a)$, with the common factor $\gamma V$, and
$xe^{-cx}\le C$ pays arbitrarily large quartic exponents. This replacement
in a linear term costs at most $C_Lr^3$.

Sandwich \eqref{f16v48:kinetic-error} by the native magnetic factors. The hard
pin turns each $|y|^2$ factor into $O_L(\varepsilon^2)$ in its row/column bound;
the soft contribution is $O_L(r^2)$. The resulting norm error is
$C_L(r^2+\varepsilon)$. Expanding the two magnetic factors around the split
Gaussian now gives exactly \eqref{f16v48:linear-operator}; their product of
linear corrections is $O_L(r^2)$ in Schur norm. The hard quadratic polynomial
factors are controlled by their Gaussian kernels, while the soft Gaussian
step and both quartic multipliers are contractions. This proves
\eqref{f16v48:quadratic-remainder}, including zero extension in the hard
coordinate. No factor is proportional to the soft strip volume. Center-folded
cross terms are $O_L(e^{-c_L/\varepsilon^2})$ and are absorbed as well.

The norm bound in \eqref{f16v48:linear-cancellation} follows from
$\|P_rM_\ell\|\le r$. For the ground compression, the scalar
\[
 a\longmapsto\int\tfrac14u^*\Bigl(\sum a_\ell H_\ell\Bigr)u
                                      |\varphi(u)|^2du
\]
is a simultaneous-color-invariant linear functional on three vectors. It is
zero. Hence $\langle\varphi,B_\ell\varphi\rangle=0$ for each $\ell$ and
$P_+B_\ell\mathcal T_+P_+=P_+\mathcal T_+B_\ell P_+=0$.
This proves the exact compression cancellation for every pair of soft
endpoints. It does not require their equality or slow variation of a test.
\end{proof}

\begin{proposition}[Quadratic hard-return cost and annular coercivity]
\label{f16v48:annulus}
After a fixed decrease of $r_*$, there are $c_L,C_L>0$ such that every physical
vector supported in one tube with $r\le|a|\le4r$, $0<r\le r_*$, satisfies
\begin{equation}
 \boxed{\quad
 \langle\psi,(\lambda_*-T)\psi\rangle
            \ge(c_Lr-C_L\varepsilon)\|\psi\|^2.
 \quad}
 \label{f16v48:annular-floor}
\end{equation}
Every physical vector supported in $|a|\le r\le4r_*$ satisfies
\begin{equation}
 \langle\psi,(\lambda_*-T)\psi\rangle
                   \ge-C_L(r^2+\varepsilon)\|\psi\|^2.
 \label{f16v48:inner-floor}
\end{equation}
The supports in the hard variable lie in the same fixed small tube. Neither
estimate requires hard-ground support or a bound on derivatives of $\psi$.
\end{proposition}
\begin{proof}
In the exact profile coordinates write $F=\mathcal I f+F_\perp$. Soft strip
and annulus projections commute with $\mathfrak P$. The reference has hard
deficit at least $\kappa_+\|F_\perp\|^2$, and its ground deficit is
$\lambda_*\langle f,(I-S_h)f\rangle$. By
\eqref{f16v48:linear-cancellation}, the entire first-jet contribution is bounded
in absolute value by
\[
 C_Lr\|F_\perp\|^2+2C_Lr\|f\|\|F_\perp\|.
\]
Young's inequality charges the mixed term to
$\kappa_+\|F_\perp\|^2/4+C_Lr^2\|f\|^2$.
Choose the fixed radius small enough to absorb the first term as well.
Together with \eqref{f16v48:quadratic-remainder}, this proves
\begin{equation}
 \langle F,(\lambda_*-K_\gamma^{\rm red})F\rangle
 \ge\lambda_*\langle f,(I-S_h)f\rangle
     +\tfrac12\kappa_+\|F_\perp\|^2
     -C_L(r^2+\varepsilon)\|F\|^2
 \label{f16v48:block-payment}
\end{equation}
for a strip of radius bounded by a fixed multiple of $r$.

On the annulus V46's exact finite-step bound gives
\[
 \langle f,(I-S_h)f\rangle
 \ge\tfrac16\int\sum_i\min\{|a_i|,1\}|f(q)|^2dq
 \ge\tfrac r6\|f\|^2,
\]
after choosing $4r_*<1$. The hard complementary constant is also at least a
fixed constant times $r$ on this interval. A further fixed decrease of $r_*$
absorbs $C_Lr^2$ into half that radial lower bound, giving
\eqref{f16v48:annular-floor}. Without the annulus, discard the nonnegative soft
term in \eqref{f16v48:block-payment} to get \eqref{f16v48:inner-floor}.
Exact half-density isometry returns both statements to the native Hilbert
space. This is a paid quadratic coupling estimate; no native hard inverse or
unproved global soft gap has been used.
\end{proof}

\subsection{Dyadic localization pays every annulus once}

Choose fixed, center-translated, gauge-invariant tube cutoffs $\zeta_\sigma$,
with disjoint supports, equal to one on smaller neighborhoods of the central
orbits. Their supports have $|a|<r_*$ and lie in the plateau of V47's final
extraction cutoff. They may be chosen with a smooth exterior member
$\zeta_{\rm out}$ so that
\[
 \zeta_{\rm out}^2+\sum_\sigma\zeta_\sigma^2=1.
\]
For example use a sine/cosine transition on each of the disjoint tubes.
V42's supported forbidden-norm result, applied to their fixed plateaux, gives
\begin{equation}
 \langle\zeta_{\rm out}\psi,(\lambda_*-T)\zeta_{\rm out}\psi\rangle
                 \ge\delta_L\|\zeta_{\rm out}\psi\|^2
 \label{f16v48:outer-floor}
\end{equation}
eventually, with $\delta_L>0$ independent of coupling. No shrinking-neighborhood
version of that result is required. Define, for $0<\rho<r_*/8$,
\[
 \mathcal W_\rho[\psi]
 =\sum_\sigma\int \zeta_\sigma(U)^2|a_\sigma(U)|
       1_{\{|a_\sigma(U)|>4\rho\}}|\psi(U)|^2\,dH(U).
\]
The constant angles $a_\sigma$ are the exact translated tube coordinates.

\begin{theorem}[A global native radial inequality]
\label{f16v48:radial-form}
For all physical $\psi\in L^2$, all $0<\rho<r_*/8$, and sufficiently large
coupling, with the threshold independent of $\rho$,
\begin{equation}
 \boxed{\quad
 c_L\mathcal W_\rho[\psi]+\delta_L\|\zeta_{\rm out}\psi\|^2
 \le\langle\psi,(\lambda_*-T)\psi\rangle
       +C_L\left(\rho^2+\varepsilon+
                    \frac{\varepsilon^2}{\rho^2}\right)\|\psi\|^2.
 \quad}
 \label{f16v48:radial-form-bound}
\end{equation}
There is no factor counting the dyadic annuli. In particular
\begin{equation}
 h^{-1}(I-T/\lambda_*)\succeq-C_LI
 \label{f16v48:semibounded}
\end{equation}
eventually on the complete physical space.
\end{theorem}
\begin{proof}
Inside each tube use a smooth radial quadratic partition at scales
$\rho,2\rho,4\rho,\ldots$, restricted to the support of $\zeta_\sigma$.
The first member is supported on $|a|<2\rho$; each other member is supported
on $r_j\le|a|\le4r_j$, with $r_j$ a dyadic multiple of $\rho$ and
$r_j\le r_*$. Start from fixed radial bump functions with finite overlap and
divide by the square root of the sum of their squares. On every point where
a cutoff is needed that denominator has a fixed positive lower bound. Their
squares sum to one, their gradients obey the pointwise bound
$\sum_j|\nabla\theta_j|^2\le C_L\rho^{-2}$, and
\[
 \sum_{j\ge1}r_j\theta_j(a)^2
       \ge c|a|1_{\{|a|>4\rho\}}.
\]
Only finitely many members meet the fixed tube; their number is not used in
these estimates. Multiply them by the fixed $\zeta_\sigma$ and add
$\zeta_{\rm out}$ to get a global quadratic partition $(\chi_j)$.
The smooth vector-valued map $U\mapsto(\chi_j(U))_j$ has Lipschitz constant
at most $C_L/\rho$, independent of the number of its components. This follows
by integrating its pointwise squared-gradient bound along a minimizing path.

The exact native localization error is the kernel
\[
 T-\sum_j\chi_jT\chi_j
   =\tfrac12\sum_j(\chi_j(U)-\chi_j(U'))^2t_\gamma(U,U').
\]
Its operator norm is at most $C_L\varepsilon^2/\rho^2$ by the original
normalized displacement second moment and the row/column Schur bound.
Its Loewner sign is not asserted. Apply \eqref{f16v48:inner-floor} to the
inner members, \eqref{f16v48:annular-floor} to every annular member, and
\eqref{f16v48:outer-floor} to the exterior member. The $C_L\varepsilon$
errors sum against squares of a partition, hence to at most
$C_L\varepsilon\|\psi\|^2$, not their number times this norm.
The positive annular terms dominate $c_L\mathcal W_\rho$, and the inner loss
is at most $C_L\rho^2\|\psi\|^2$. This proves the displayed inequality.
Take $\rho=h$ and use $\varepsilon^2=Vh^3$, $\varepsilon=o(h)$, to obtain
\eqref{f16v48:semibounded}. No eigenvector equation was used in this theorem.
\end{proof}

\begin{theorem}[Native position tightness on the whole critical band]
\label{f16v48:tightness}
For fixed $K\ge0$ let
$\mathsf B_{\gamma,a}(K)=1_{[\lambda_*(1-Kh),1]}(T|_{\mathcal H_a})$.
Retain V47's unchanged extraction $F_\gamma=\mathfrak E_{\gamma,a}\psi$ and
soft profile $f_\gamma=\mathcal I^*F_\gamma$. Then, uniformly over unit vectors
in this band and every sector,
\begin{equation}
 \boxed{\quad
 \limsup_{\gamma\to\infty}
 \sup_{\psi\in\operatorname{Ran}\mathsf B_{\gamma,a}(K),\ \|\psi\|=1}
 \int_{|q|>4R}|F_\gamma(u,q)|^2\,du\,dq
       \le C_L\left(\frac K R+\frac1{R^3}\right),\qquad R\ge1.
 \quad}
 \label{f16v48:tail}
\end{equation}
The same bound holds for $|f_\gamma(q)|^2$. In fact, at one common sufficiently
large coupling threshold for fixed $L,K$,
\begin{equation}
 \sup_{\psi\in\operatorname{Ran}\mathsf B_{\gamma,a}(K),\ \|\psi\|=1}
 \int |q|\,|F_\gamma(u,q)|^2\,du\,dq\le C_L(K+1).
 \label{f16v48:first-moment}
\end{equation}
An empty band contributes zero. In particular V47's uniform
position-tightness condition is satisfied.
\end{theorem}
\begin{proof}
The band gives
$\langle\psi,(\lambda_*-T)\psi\rangle\le\lambda_*Kh$.
Set $\rho=hR$ in \eqref{f16v48:radial-form-bound}; at each fixed $R$ it is
admissible eventually. The right side is at most
\[
 h\left[\lambda_*K+C_L(hR^2+\sqrt h+R^{-2})\right],
\]
with the fixed $\sqrt V$ absorbed into $C_L$. On the part of an extraction
tube where $|q|>4R$, its $\zeta_\sigma^2$-weighted mass is bounded by
$\mathcal W_{hR}/(4hR)$. Its remaining mass is bounded by
$\|\zeta_{\rm out}\psi\|^2$, because the fixed cutoff squares partition one.
This argument also covers the portions of the unchanged extraction tube
outside the smaller cutoffs used in this proof. Both terms are controlled by
\eqref{f16v48:radial-form-bound}. Divide the first by $hR$, let
$\gamma\to\infty$ at fixed $R$, and note that the second tends to zero.
Exact extraction and center normalization give \eqref{f16v48:tail}.
Cauchy--Schwarz in the hard variable gives the same bound for $f_\gamma$.
Finally send $R\to\infty$. For \eqref{f16v48:first-moment}, use
$\rho=h$ instead. The same inequality bounds $\mathcal W_h/h$ and
$h^{-1}\|\zeta_{\rm out}\psi\|^2$ by $C_L(K+1)$. In the inner part
$|q|\le4$, while on the fixed extraction support $|q|\le C_L/h$.
Split by the same quadratic cutoff pair; these two bounds pay its entire
first moment. No support restriction is imposed on the original band vector
and no rank is counted.
\end{proof}

\subsection{The native fixed-volume spectrum now has the model coefficients}

Let $e_0<e_1\le e_2\le\cdots$ be the invariant eigenvalues of the unchanged
$H_{\rm mat}$, with multiplicity. Their existence and $e_1-e_0>0$ are V44's
model theorems; they are not numerical evaluations. Index native eigenvalues
from one in each center sector, so $\lambda_{0,1}=\Lambda_\gamma$.

\begin{theorem}[Ordered critical spectrum and the same-top neutral gap]
\label{f16v48:spectrum}
At every fixed spatial regulator and every fixed $j\ge1$,
\begin{equation}
 \boxed{\quad
 \frac{1-\lambda_{a,j}(\gamma)/\lambda_*}{h}\longrightarrow e_{j-1}
 \quad\hbox{in every center sector}.\quad}
 \label{f16v48:levels}
\end{equation}
The ordered one-region eigenvalues of V43 have the same limits. In particular
\begin{equation}
 \begin{split}
 \Lambda_\gamma&=\lambda_*[1-he_0+o_L(h)],\\
 \Lambda_\gamma-\lambda_{0,2}
        &=\lambda_*h(e_1-e_0)+o_L(h),\\
 -\log(\lambda_{0,2}/\Lambda_\gamma)
        &=h(e_1-e_0)+o_L(h),\\
 g_\gamma/(\lambda_*h)&\longrightarrow e_1-e_0.
 \end{split}
 \label{f16v48:gap}
\end{equation}
The last quantity is V43's actual reconstruction-metric gap, not a bare
regional compression. These statements fix $L$ and $j$ before the limit.
\end{theorem}
\begin{proof}
V45 bounds above the scaled deficits of any fixed number of ordered neutral
eigenvectors. Equation \eqref{f16v48:semibounded} bounds them below. On any
subsequence, choose a further subsequence on which those finitely many
deficits converge. V47 gives hard capture and strong local compactness of
their profiles; Theorem~\ref{f16v48:tightness} upgrades this to global strong
compactness. Extraction is asymptotically isometric on the native band, so
all inner products of this finite eigenbasis are preserved in the limit.
V46's corrected profile equation therefore produces orthonormal invariant
model eigenfunctions with these limiting deficits. Model min--max gives the
ordered liminf $e_{j-1}$; V45 gives the matching limsup. This proves the
neutral limits on every subsequence and hence in full. It is exactly the
previously conditional position-only completion argument of V47, with its
hypothesis now supplied by a native inequality.

Every fixed level lies in V43's matching interval by V44. Its exponential
sector and one-region differences are $o(h)$, so all the other limits follow.
Subtract the first two neutral expansions and use $-\log(1-x)=x+O(x^2)$.
V43--V44's exact $\Delta_\gamma^{\rm neu}/g_\gamma\to1$ gives the final
line of \eqref{f16v48:gap}. No identity metric is installed on the entire
regional space.
\end{proof}

\begin{proposition}[Exact low-cluster counts and critical vacuum profiles]
\label{f16v48:clusters}
For each fixed $E\notin\operatorname{spec}(H_{\rm mat}|_{\rm inv})$, the number
of eigenvalues in one sector satisfying
$(1-\lambda/\lambda_*)/h<E$ is eventually exactly the number of invariant
model eigenvalues below $E$. On isolated finite model eigenspaces, subsequential
extracted orthonormal native bases converge to orthonormal model bases. For
the actual neutral vacuum the whole extracted sequence satisfies
\begin{equation}
 \mathfrak E_{\gamma,0}\Omega_\gamma
       \longrightarrow\varphi\otimes\phi_0
       \quad\hbox{in }L^2(du\,dq),
 \label{f16v48:vacuum-profile}
\end{equation}
where $\phi_0$ is the unit positive invariant ground function of $H_{\rm mat}$.
Consequently the corresponding extracted density converges in $L^1$ to
$|\varphi|^2|\phi_0|^2$. This is a statement after the specified center
unfolding and critical coordinate scaling, not an unscaled $L^2$ limit of
native vacuum vectors.
\end{proposition}
\begin{proof}
Choose model indices immediately below and above the fixed threshold and use
\eqref{f16v48:levels}. Since ordered native eigenvalues exhaust the spectrum
above any positive threshold, these two indices give the exact count; possible
negative scaled deficits have already been bounded below. For a finite cluster,
extract a basis along subsequences as in the preceding proof. Compactness,
Gram preservation and the profile equation identify its limiting eigenspace
and multiplicity. For the ground, simplicity and positivity remove phase and
basis ambiguity: extraction, neutral unfolding and $\varphi$ are nonnegative.
Every subsequential limit is the same positive unit $\phi_0$. This proves
full-sequence convergence. The $L^1$ claim follows from
$\||F|^2-|G|^2\|_1\le(\|F\|+\|G\|)\|F-G\|$.
\end{proof}

\begin{proposition}[Two spectral scales and recentering the entrance law]
\label{f16v48:consequences}
Let $g_*=e_1-e_0>0$. For sufficiently large coupling the full physical transfer
has exactly eight eigenvalues above
\begin{equation}
 z_\gamma=\Lambda_\gamma-\tfrac12\lambda_*h g_*.
 \label{f16v48:eight-cut}
\end{equation}
There is one in each center sector. Their charged top splittings retain V42's
upper bound $C_Le^{-\kappa_L\sqrt\gamma}$, while their next sector levels are
separated from the top by $\lambda_*h g_*+o_L(h)$.

On every fixed-$B$ duration sequence admitted by V40, its projected spectral
law may now be centered at the actual Perron value:
\begin{equation}
 -n_\gamma\log(\lambda/\Lambda_\gamma)
       \ \Rightarrow\ \operatorname{Gamma}(9/2,1),\qquad
 n_\gamma(1-\lambda/\Lambda_\gamma)
       \ \Rightarrow\ \operatorname{Gamma}(9/2,1).
 \label{f16v48:recentered}
\end{equation}
The probability is V40's unchanged projected endpoint-weighted spectral
probability, not the stationary density of states.
\end{proposition}
\begin{proof}
The first two ordered sector limits and V42's same-top exponentially small
charged bounds place the first level of every sector above
\eqref{f16v48:eight-cut} and its second below. Ordered monotonicity excludes
any further level above that threshold. No tunneling prefactor or matching
charged lower bound is deduced.

V40's regime gives $n_\gamma h\to0$ at fixed $B$. The new top expansion gives
$n_\gamma\log(\Lambda_\gamma/\lambda_*)\to0$ and
$\Lambda_\gamma/\lambda_*\to1$. The logarithmic deficit differs from V40's
by the first deterministic vanishing quantity. The linear deficit is its
old value divided by $\Lambda_\gamma/\lambda_*$ plus a deterministic
$O(n_\gamma h)$ term. The probability limits follow. Unlike the deficits
centered at $\lambda_*$, these actual-top deficits are nonnegative at every
finite coupling. The probability measure and all endpoint normalizations
are unchanged.
\end{proof}

\clearpage
\subsection{Scope, next estimate, and validation}

The new conclusion is fixed-volume spectral completeness for the \emph{existing}
critical branch. The proof rests on three inputs of different kinds: the new
quotient-metric/first-jet and annular estimates, V42's supported forbidden-region
bound, and V45--V47's corrected recovery and profile results. It does not
independently recertify the entire analytic lineage. In particular the sharp
coefficients use V45's strong core equation as an inherited premise; the
new coarse bounds supply compactness rather than recomputing that equation.

There is no uniformity here as $L\to\infty$, no numerical enclosure of
$e_0,e_1$ or of the native thresholds, and no rate in the remaining $o_L(h)$.
The model and native kinetic normalizations are exactly those already fixed.
For a separately supplied physical time spacing $a_t(\gamma,L)$ the neutral
energy has the expression
\[
 E_{\rm neu}^{\rm phys}
   =\frac{(\gamma L^3)^{-1/3}}{a_t(\gamma,L)}
        [e_1-e_0+o_L(1)].
\]
No value of $a_t$ is chosen by this formula, and a fixed number of lattice
sites is not a fixed-physical-volume continuum limit. The next substantive
problem is uniform control of this comparison along a specified growing
regulator and physical-time trajectory, or a quantitative higher-order and
tunneling analysis at fixed regulator. Another fixed-$L$ recovery space is no
longer the missing completeness step. The next scheduled checkpoint is V50.

The diagnostic records exact finite tests of the quotient Schur metric,
coarea cancellation, color-singlet compression, quadratic hard-return cost,
annular inequalities, dyadic overlap, critical powers, and ordered spectral
consequences. These are not native spectral enclosures or formal verification
of the analytic kernel remainder. The audit records which earlier scripts
were actually executed. Removing this exact insertion and reversing the
literal title version recovers V47 byte for byte.

\begin{firewall}
V48 proves a quadratic-radius hard-return estimate, a native annular barrier,
and radial tightness at the critical scale. Together with the explicitly
identified inherited recovery and profile theorems this yields the ordered
fixed-spatial-regulator spectrum and its center-neutral gap coefficient.
It also separates the eight exponentially split sector tops from the next
algebraic sector levels. It does not establish a volume-uniform estimate,
calibrated physical mass, interacting continuum construction, or a matching
center tunneling asymptotic. No comparison operator or auxiliary clock
replaces the native transfer.
\end{firewall}
% END F16 V48 APPEND

% BEGIN F16 V49 APPEND -- 2026-09-18
% Insert before the final document-ending command of cumulative V48.
% All predecessor bodies, native operators, measures and clocks are preserved.

\clearpage
\section{V49: volume-uniform winding asymptotics and the shrinking forbidden gap}
\label{f16v49:context}

V48 closes its critical spectral comparison with the spatial lattice fixed.
Its constants cannot be reused on a growing lattice without further estimates.
We return to an earlier input whose volume dependence can be determined exactly:
V41's harmonic threshold over the compact flat family. The positive winding
coefficients have a uniform inverse-fourth-power law, with a summable second-order
lattice error. This determines a common periodic limiting potential, its cusp,
and the finite-size correction to the original Gaussian top. It also shows that
a natural native forbidden-complement margin really shrinks like $L^{-1}$ after
normalization at the same Perron top. That loss is not merely a proof artifact.

The uniform conclusions below concern explicit harmonic data. The native
compression conclusion takes the coupling limit separately at each lattice.
No unproved uniform remainder in V48, joint continuum limit, or new physical
clock is supplied. The useful output is a scale-resolved input for such a
comparison, not another fixed-volume quartic recovery space.

\subsection{Sources and the unchanged harmonic data}

The supplied V48 appendix and audit were read completely. The current lineage,
Grand Tensor V076 and f14/f15 V100 were searched for winding, heat-trace,
large-volume and related terms; these are targeted checks, not an exhaustive
novelty or proof audit. V41 already proves the positive winding representation
and strict central maximizers. Those facts are reused rather than reproved as
new. The compactness and spectral assertions of V48 are not premises of the
coefficient estimates below. Only the final native interpretation additionally
uses V41's fixed-lattice harmonic localization principle.

The Poisson representation of finite-volume vacuum-valley energies is classical.
The primary text of van Baal, \href{https://arxiv.org/abs/hep-ph/0008206}
{\emph{QCD in a Finite Volume}}, Section 3, discusses that representation and
its conic origin singularity. Its continuum spectral and coupling conventions
are not substituted for the Wilson normalization here. The coefficient below
is derived from the actual discrete heat kernel. The Bessel integral and Fourier
identities are those of \href{https://dlmf.nist.gov/10.32.E3}{DLMF 10.32.3}
and \href{https://dlmf.nist.gov/10.35.E2}{10.35.2}; the local-limit error needed
here is proved directly. The next scheduled checkpoint remains V50.

For even $L=2m\ge8$, retain V41's exact definitions
\begin{equation}
 \begin{split}
 \lambda_L(k,\phi)&=4\sum_{\mu=1}^3
       \sin^2\!\left(\frac{2\pi k_\mu+\phi_\mu}{2L}\right),\\
 \xi(s)&=\operatorname{arcosh}(1+s/2),\\
 F_L(\phi)&=\sum_{k\in\{0,\ldots,L-1\}^3}\xi(\lambda_L(k,\phi)),\\
 \mathcal E_L(\phi)&=F_L(0)+2F_L(\phi),\\
 \mathfrak h_L(\phi)&=e^{-\mathcal E_L(\phi)},\qquad
 \lambda_{*,L}=e^{-3F_L(0)}=\Lambda_{{\rm G},B},\quad B=(L/2)^3.
 \end{split}
 \label{f16v49:definitions}
\end{equation}
Here $\phi\in\mathbb T^3=\mathbb R^3/(2\pi\mathbb Z)^3$ is the adjoint twist.
The neutral color and both charged colors remain in $\mathcal E_L$.
V41's coefficients satisfy exactly
\begin{equation}
 \begin{split}
 F_L(\phi)-F_L(0)&=\sum_{z\in\mathbb Z^3\setminus\{0\}}
             a_L(z)[1-\cos(z\cdot\phi)],\\
 a_L(z)&=L^3\int_0^\infty
   \frac{e^{-8t}I_0(2t)}{t}\prod_{\mu=1}^3 I_{Lz_\mu}(2t)\,dt>0.
 \end{split}
 \label{f16v49:inherited-winding}
\end{equation}
Unlike earlier fixed-$L$ bounds, constants denoted $c,C,C_N$ in the harmonic
estimates below are independent of $L,z,\phi$. The analytic estimates in fact
hold for every integer $L\ge2$; the native statements keep the admitted even
geometry. Write $r(\phi)=\operatorname{dist}(\phi,(2\pi\mathbb Z)^3)$.

\subsection{A four-dimensional lattice kernel controls every winding coefficient}

Put $p_t(j)=e^{-2t}I_j(2t)$, and for $x\in\mathbb Z^4$ set
\[
 p_t^{(4)}(x)=\prod_{\nu=0}^3p_t(x_\nu),\qquad
 A(x)=\int_0^\infty p_t^{(4)}(x)\,\frac{dt}{t}\quad(x\ne0).
\]
This auxiliary heat kernel has generator
$\sum_{\nu=0}^3[f(x+e_\nu)+f(x-e_\nu)-2f(x)]$.
It is not a time evolution installed on the Yang--Mills field. In particular
$t$ here is a transform variable, not calibrated Euclidean time. Directly,
\begin{equation}
                  a_L(z)=L^3 A(0,Lz).
 \label{f16v49:walk-identification}
\end{equation}

\begin{proposition}[Uniform logarithmic heat-kernel asymptotics]
\label{f16v49:heat}
For every $x\in\mathbb Z^4\setminus\{0\}$,
\begin{equation}
 \boxed{\quad
 \left|A(x)-\frac1{\pi^2|x|^4}\right|\le\frac C{|x|^6},
 \qquad \frac c{|x|^4}\le A(x)\le\frac C{|x|^4}.
 \quad}
 \label{f16v49:heat-asymptotic}
\end{equation}
An exact positive alternative is
\begin{equation}
               A(x)=\sum_{n=1}^{\infty}\frac{P^n(0,x)}n,
 \label{f16v49:walk-series}
\end{equation}
where $P$ is the simple discrete-time walk with probability $1/8$ for each of
its eight neighboring sites. All sums and integrals here are finite for $x\ne0$.
\end{proposition}
\begin{proof}
We give the error estimate because its spatial decay, not merely pointwise
convergence, is needed when windings are summed. The Fourier symbol is
\[
 \omega(\theta)=2\sum_{\nu=0}^3(1-\cos\theta_\nu),\qquad
 p_t^{(4)}(x)=(2\pi)^{-4}\int_{[-\pi,\pi]^4}
           e^{-t\omega(\theta)}e^{i x\cdot\theta}\,d\theta.
\]
For every fixed integer $N$, every $t\ge1$, and every lattice $x$,
\begin{equation}
 \left|p_t^{(4)}(x)-(4\pi t)^{-2}e^{-|x|^2/(4t)}\right|
 \le C_Nt^{-3}(1+|x|/\sqrt t)^{-N}.
 \label{f16v49:local-limit}
\end{equation}
To verify this, choose a smooth cutoff $\eta$ supported inside the Fourier cube,
equal to one near its sole zero of $\omega$. On its sufficiently small support,
$\omega(\theta)=|\theta|^2+O(|\theta|^4)$ and
$\omega(\theta)\ge c|\theta|^2$. After $\theta=v/\sqrt t$, put
\[
 R_t(v)=\eta(v/\sqrt t)e^{-t\omega(v/\sqrt t)}-e^{-|v|^2}.
\]
For each pair of multiindices $\alpha,\beta$,
\[
 \|v^\alpha\partial_v^\beta R_t\|_{L^1(\mathbb R^4)}\le C_{\alpha,\beta}/t.
\]
Indeed expand the difference of the two exponentials by its integral mean-value
formula; $t\omega(v/\sqrt t)-|v|^2$ and its derivatives cost $t^{-1}$ times a
fixed polynomial. Both interpolated quadratic exponents have a Gaussian lower
bound. Derivatives of the cutoff are supported at $|v|\asymp\sqrt t$ and cost
only exponentially small Gaussian tails, which obey the same estimates.
Fourier integration by parts now gives the right side of
\eqref{f16v49:local-limit} for the cutoff piece, including the factor $t^{-2}$
from the change of variables. The complementary periodic multiplier is supported
where $\omega\ge c>0$. Any fixed number of its derivatives has $L^1$ norm at
most $C_N(1+t)^Ne^{-ct}$. Periodic integration by parts is legitimate because
$x$ is integral and the multiplier is smooth on the torus. Its Fourier
coefficient is bounded by
$C_N(1+t)^Ne^{-ct}(1+|x|)^{-N}$, which is also bounded by the right side of
\eqref{f16v49:local-limit}. This proves that estimate uniformly in $x,t$.

Choose $N>6$ and write $R=|x|\ge1$. Integrating the error for $t\ge1$ gives
\[
 \int_1^\infty C_Nt^{-4}(1+R/\sqrt t)^{-N}dt
 \le CR^{-6}\int_0^\infty u^{-4}(1+u^{-1/2})^{-N}du
 \le CR^{-6}.
\]
For $0<t<1$, at least $\ell=|x|_1$ jumps are required. The continuous-time walk
has Poisson jump count of mean $8t$, so
$p_t^{(4)}(x)\le e^8(8t)^\ell/\ell!$.
Its integral against $dt/t$ is at most $e^8 8^\ell/(\ell\ell!)$, bounded by
$C_N R^{-N}$ for any fixed $N$ after increasing the constant. The corresponding
Gaussian integral over $0<t<1$ has the same rapid bound. Finally
\[
 \int_0^\infty(4\pi t)^{-2}e^{-R^2/(4t)}\frac{dt}{t}
   =\frac1{16\pi^2}\frac{16}{R^4}\int_0^\infty u e^{-u}du
   =\frac1{\pi^2R^4}.
\]
This proves the first estimate. It gives both comparisons for sufficiently
large $R$. The remaining nonzero lattice points form a finite set, and $A(x)>0$
there, so adjustment of two absolute constants proves them for every $x\ne0$.
Poissonization gives
$p_t^{(4)}=e^{-8t}\sum_{n\ge0}(8t)^nP^n/n!$.
Tonelli and $\int_0^\infty e^{-8t}(8t)^n dt/t=\Gamma(n)$ give
\eqref{f16v49:walk-series}. Its finiteness also follows from the estimates just
proved. The missing $n=0$ term is exactly why $x\ne0$ was stipulated.
\end{proof}

\begin{theorem}[Uniform winding coefficients with a summable lattice remainder]
\label{f16v49:coefficients}
For every $L\ge2$ and $z\in\mathbb Z^3\setminus\{0\}$,
\begin{equation}
 \boxed{\quad
 \frac c{L|z|^4}\le a_L(z)\le\frac C{L|z|^4},\qquad
 \left|La_L(z)-\frac1{\pi^2|z|^4}\right|
       \le\frac C{L^2|z|^6}.
 \quad}
 \label{f16v49:coefficient-bounds}
\end{equation}
In particular, for $R\ge1$,
\begin{equation}
 \sum_{|z|>R}La_L(z)\le C/R,\qquad
 \sum_{|z|>R}|z|^j
 \left|La_L(z)-\frac1{\pi^2|z|^4}\right|
       \le C_jL^{-2}R^{j-3},\quad j=0,1,2.
 \label{f16v49:coefficient-tails}
\end{equation}
No winding cutoff or spatial volume is absorbed into the constants.
\end{theorem}
\begin{proof}
Insert $x=(0,Lz)$ into \eqref{f16v49:heat-asymptotic} and multiply by $L^4$.
This proves all pointwise estimates. Integer shell counts in three dimensions
are at most $Ck^2$ on $k\le|z|<k+1$. Summing the powers $k^{-4}$ and
$k^{j-6}$ against that count proves the two tail bounds, with a harmless
adjustment for $1\le R<2$. In particular absolute summation is now uniform
after multiplication by $L$; V41 had only established finiteness at fixed $L$.
\end{proof}

\subsection{A common periodic potential with a controlled cusp}

Define the explicit continuous periodic function
\begin{equation}
 \mathcal P(\phi)=\frac1{\pi^2}
     \sum_{z\in\mathbb Z^3\setminus\{0\}}
          \frac{1-\cos(z\cdot\phi)}{|z|^4},\qquad
 U_L(\phi)=L[F_L(\phi)-F_L(0)].
 \label{f16v49:potential}
\end{equation}
The sum defining $\mathcal P$ is absolutely convergent, but it must not be
termwise differentiated twice without further analysis. The following statement
differentiates only the more summable lattice error.

\begin{theorem}[Uniform potential, second-order lattice error, and linear contrast]
\label{f16v49:potential-theorem}
The difference $R_L=U_L-\mathcal P$ belongs to $C^2(\mathbb T^3)$ and obeys
\begin{equation}
 \boxed{\quad
 \|R_L\|_{C^2}\le CL^{-2},\qquad
 |R_L(\phi)|\le CL^{-2}r(\phi)^2.
 \quad}
 \label{f16v49:C2-error}
\end{equation}
There are absolute $c,C>0$ such that, on the whole twist torus,
\begin{equation}
 \boxed{\qquad c\,r(\phi)\le U_L(\phi)\le C\,r(\phi).
 \qquad}
 \label{f16v49:global-contrast}
\end{equation}
On the representative ball $|\phi|\le1$,
\begin{equation}
 \mathcal P(\phi)=|\phi|+G(\phi),\qquad
 G\in C^\infty(\{|\phi|<\pi\}),\quad G(0)=0,\quad DG(0)=0,
 \label{f16v49:cusp-decomposition}
\end{equation}
where $G$ is even. Consequently
\begin{equation}
 \boxed{\quad U_L(\phi)=|\phi|+O(|\phi|^2)\quad(|\phi|\le1),\quad}
 \label{f16v49:uniform-cusp}
\end{equation}
with an error constant independent of $L$. Neither $U_L$ nor $\mathcal P$ is
asserted differentiable at the origin; it is their difference that is $C^2$.
\end{theorem}
\begin{proof}
Let $b_L(z)=La_L(z)-(\pi^2|z|^4)^{-1}$. The series
$\sum b_L(z)(1-\cos(z\cdot\phi))$ and its derivatives through order two
converge absolutely and uniformly by \eqref{f16v49:coefficient-bounds}. They
prove the first assertion. At zero the function and its first derivative vanish.
Integrate the Hessian bound along the line to a nearest representative of
$\phi$ to prove its quadratic bound.

For the cusp, use $|z|^{-4}=\int_0^\infty t e^{-t|z|^2}dt$ and the Gaussian
Poisson formula
\[
 \sum_{z\in\mathbb Z^3}e^{-t|z|^2}e^{iz\cdot\phi}
   =(\pi/t)^{3/2}\sum_{k\in\mathbb Z^3}
                       e^{-|\phi-2\pi k|^2/(4t)}.
\]
This formula follows by computing the Fourier coefficients of the periodized
Gaussian; both sides converge absolutely at every $t>0$. In the integral over
$t\ge1$, the original $z$ series can be differentiated any finite number of
times, and gives a smooth even function of $\phi$. On $0<t\le1$, apply Poisson
and split $k=0$ from $k\ne0$. For $|\phi|<\pi$ the nonzero images have a fixed
separation on every smaller compact ball. Their derivatives are bounded by
summable Gaussian tails, even after integration at $t=0$, and so give a smooth
function there. The remaining part is
\[
 \frac1{\sqrt\pi}\int_0^1 t^{-1/2}
                         (1-e^{-|\phi|^2/(4t)})dt.
\]
The same integral over $0<t<\infty$ is exactly $|\phi|$, by substitution and
one integration by parts. Its removed $t\ge1$ part is smooth in $|\phi|^2$ near
zero. This proves \eqref{f16v49:cusp-decomposition}; parity and the value at
zero give the two stated zero jets. It also proves the uniform cusp using
\eqref{f16v49:C2-error}.

For clarity the global lower bound is not inferred just from local convergence.
The coefficient comparisons make $U_L$ comparable, with absolute constants,
to $\mathcal P$ pointwise. For small $r$, the cusp gives
$\mathcal P\ge r/2$. Away from zero, $\mathcal P>0$: keeping $z=e_\mu$ forces
all components of a zero to be multiples of $2\pi$. Compactness of the remaining
torus supplies a positive lower bound for $\mathcal P/r$. For the upper bound
at $0<r\le1$, split the series at $|z|=1/r$. The inner part is bounded by
$Cr^2\sum_{|z|\le1/r}|z|^{-2}\le Cr$; the outer part is at most $Cr$ by its
tail. On $r\ge1$, absolute summation gives the same inequality with a larger
constant. This proves \eqref{f16v49:global-contrast}. The limiting function is
classical; the new information used here is the uniform normalized lattice
remainder and the explicit transfer convention.
\end{proof}

\subsection{The original Gaussian scalar: bulk and finite-size terms}

Let
\begin{equation}
 \mathcal I_3=(2\pi)^{-3}\int_{[-\pi,\pi]^3}
       \xi\!\left(4\sum_{\mu=1}^3\sin^2(p_\mu/2)\right)dp>0,
 \qquad Z_3(2)=\sum_{z\in\mathbb Z^3\setminus\{0\}}|z|^{-4}<\infty.
 \label{f16v49:bulk-constants}
\end{equation}
The notation $Z_3(2)$ denotes this explicit convergent sum, not a native
partition function or an analytically continued quantity.

\begin{proposition}[An exact bulk identity and its first finite-size correction]
\label{f16v49:bulk}
For every $L$,
\begin{equation}
 F_L(\phi)=L^3\mathcal I_3-\sum_{z\ne0}a_L(z)\cos(z\cdot\phi).
 \label{f16v49:bulk-identity}
\end{equation}
Consequently, uniformly as $L\to\infty$,
\begin{equation}
 \boxed{\quad
 -\log\lambda_{*,L}=3L^3\mathcal I_3-
             \frac{3Z_3(2)}{\pi^2L}+O(L^{-3}).
 \quad}
 \label{f16v49:bulk-asymptotic}
\end{equation}
More generally, uniformly in $\phi$,
\begin{equation}
 \mathcal E_L(\phi)
 =3L^3\mathcal I_3-
   \frac1{\pi^2L}\sum_{z\ne0}\frac{1+2\cos(z\cdot\phi)}{|z|^4}
       +O(L^{-3}).
 \label{f16v49:full-threshold-asymptotic}
\end{equation}
The errors here have constants independent of lattice size. No scalar is
subtracted from the native transfer by these identities.
\end{proposition}
\begin{proof}
Average $F_L(\phi)$ over the twist torus with its normalized measure. The
changes $p=(2\pi k+\phi)/L$, as $k$ varies, tile a full Brillouin torus, so the
average is exactly $L^3\mathcal I_3$. The average of
\eqref{f16v49:inherited-winding} is $\sum_{z\ne0}a_L(z)$, justified by absolute
summation. This gives \eqref{f16v49:bulk-identity}. Sum the coefficient error
in \eqref{f16v49:coefficient-bounds}; $\sum|z|^{-6}<\infty$ makes the error
$O(L^{-3})$ before multiplication by $L$. Evaluate at zero and use
$-\log\lambda_{*,L}=3F_L(0)$, or combine one neutral and two charged copies,
to prove the remaining formulas. The bulk coefficient is strictly positive
because its nonnegative integrand vanishes only at an isolated momentum.
\end{proof}

\subsection{Two uniform scale comparisons, without a growing-volume native claim}

Let $h=(\gamma L^3)^{-1/3}=\gamma^{-1/3}/L$. The nonzero hard momentum with
smallest curl eigenvalue has $\nu_{\min}=4\sin^2(\pi/L)$. Thus the original
hard Mehler ratio satisfies
\begin{equation}
 r_+(L)=e^{-\xi_{\min}(L)},\qquad
 \xi_{\min}(L)=2\operatorname{arsinh}(\sin(\pi/L)).
 \label{f16v49:hard-frequency}
\end{equation}
This refers to the complete hard complementary space used for the band inverse,
not an assertion about a physical one-particle state.

\begin{proposition}[Hard separation and the commuting soft contrast in common units]
\label{f16v49:scales}
There are absolute positive constants such that, for all $L\ge8,\gamma\ge1$,
\begin{equation}
 \boxed{\quad
 c\gamma^{1/3}\le\frac{1-r_+(L)}h\le C\gamma^{1/3},\qquad
 L(1-r_+(L))\longrightarrow2\pi\quad(L\to\infty).
 \quad}
 \label{f16v49:hard-scaled}
\end{equation}
For a commuting constant triple $a_\mu=h b_\mu n$ with $|n|=1$ and
$2\gamma^{-1/3}|b|\le1$, its adjoint twist is $\phi=2Lh b$. Uniformly in $L$,
\begin{equation}
 \boxed{\quad
 \frac{-\log[\mathfrak h_L(2Lhb)/\lambda_{*,L}]}h
      =4|b|+O(\gamma^{-1/3}|b|^2).
 \quad}
 \label{f16v49:critical-flat}
\end{equation}
This last statement is an exact harmonic-threshold consequence. It is not a
native eigenvalue expansion on a joint $(\gamma,L)$ trajectory.
\end{proposition}
\begin{proof}
Write $s=\sin(\pi/L)$. Then
$1-e^{-2\operatorname{arsinh}s}=2s(\sqrt{1+s^2}-s)$.
For $L\ge8$, $2/L\le s\le\pi/L$ and the last parenthesis is bounded above
and below by positive absolute constants. This proves the two-sided estimate;
Taylor expansion at $L^{-1}=0$ gives its limit. The nonzero momentum formula
and the hard Mehler factors are exactly the inherited ones, with no identity
metric substituted on rooted variables.

For the second assertion, the constant links have fundamental holonomies
$\operatorname{Exp}(L a_\mu)$; the adjoint angle is twice this angle. The stated
condition keeps the representative within the unit twist ball. Since
$-\log(\mathfrak h_L(\phi)/\lambda_{*,L})=2U_L(\phi)/L$, the uniform cusp gives
\[
 \frac{2U_L(2\gamma^{-1/3}b)}{Lh}
 =2\gamma^{1/3}\left[2\gamma^{-1/3}|b|
                         +O(\gamma^{-2/3}|b|^2)\right].
\]
This is \eqref{f16v49:critical-flat}. It uniformizes V44's fixed-lattice cusp
check, not its nonlinear spectral comparison. The cusp is generated by charged
transverse modes and must not be added again as a separate potential to the
quartic model retaining those same modes.
\end{proof}

\subsection{The native forbidden margin cannot stay constant in volume}

We now draw a consequence for a specifically declared family of native
compressions. Choose three fundamental lattice cycles at one base vertex and
let $H_\mu(U)$ be their actual ordered holonomies. Define on the full compact
configuration space
\begin{equation}
 d_{\rm cen}(U)=2\min_{\sigma\in\{\pm1\}^3}
       \left(\sum_{\mu=1}^3
         \operatorname{dist}_{\SU(2)}(H_\mu(U),\sigma_\mu I)^2\right)^{1/2}.
 \label{f16v49:holonomy-distance}
\end{equation}
This is continuous, gauge invariant and center invariant. On a flat connection
with adjoint twist $\phi$, it equals $r(\phi)$. The factor two accounts for the
adjoint angle. The definition applies off the flat set without forcing its
holonomies to commute. For fixed $0<\rho\le1$ put
\[
 Q_{L,\rho}=1_{\{d_{\rm cen}\ge\rho\}},\qquad
 \Theta_L(\rho)=\lim_{\gamma\to\infty}
       \frac{\|Q_{L,\rho}T_{\gamma,L}Q_{L,\rho}\|_{\rm phys}}
            {\Lambda_{\gamma,L}}.
\]
The following theorem proves that this limit exists; its definition does not
assume it. These are new explicitly declared sets, not a claim of uniformity
for V42's geodesic neighborhoods as their radii change.

\begin{theorem}[Exact sequential compression limit and a shrinking native margin]
\label{f16v49:forbidden}
For every fixed admitted $L$ and $0<\rho\le1$,
\begin{equation}
 \boxed{\quad
 \Theta_L(\rho)
  =\exp\!\left[-\frac2L
       \min_{r(\phi)\ge\rho}U_L(\phi)\right],\qquad
 \frac{c\rho}{L}\le1-\Theta_L(\rho)\le\frac{C\rho}{L}.
 \quad}
 \label{f16v49:forbidden-margin}
\end{equation}
The last constants are absolute. The coupling limit is taken first at each
$L,\rho$; no common entry coupling or remainder uniform in $L$ is asserted.
The full-space supported norm equals the physical supported norm at each
finite coupling.
\end{theorem}
\begin{proof}
For a fixed $L$, apply V41's local norm upper argument to functions supported
on the closed set $\{d_{\rm cen}\ge\rho\}$. Its remaining flat set is compact.
The same finite cover and the original kinetic localization error give
\[
 \limsup_{\gamma\to\infty}\|Q_{L,\rho}T_{\gamma,L}Q_{L,\rho}\|
 \le\max_{r(\phi)\ge\rho}\mathfrak h_L(\phi).
\]
No action floor on all of that set is asserted: it contains noncentral flat
connections. For the lower bound, choose a flat connection with
$r(\phi)>\rho$. V41's local two-scale trial can be supported inside this open
set and has limiting quotient $\mathfrak h_L(\phi)$. Boundary points with
$r(\phi)=\rho\le1$ are approximated by such commuting twists by a small radial
increase of their nearest representative. Continuity of the harmonic threshold
therefore gives the matching supremum over $r(\phi)\ge\rho$. This proves the
supported norm limit.

At finite coupling the supported operator is a nonzero compact operator with
strictly positive kernel on its supported space. Its top vector is unique and
positive. The set and kernel are gauge invariant, so that vector is gauge
invariant by uniqueness. This justifies equality of the full and physical
supported tops without using a non-invariant trial as an excited-state test.
The original top tends to $\lambda_{*,L}>0$ by V41. Division, followed by
$\mathfrak h_L(\phi)/\lambda_{*,L}=e^{-2U_L(\phi)/L}$, proves the exact formula.

The global contrast gives
$\min_{r\ge\rho}U_L\ge c\rho$. Testing $\phi=(\rho,0,0)$ gives the upper
bound $C\rho$. Thus
$e^{-C\rho/L}\le\Theta_L(\rho)\le e^{-c\rho/L}$ after adjusting constants.
On $0<\rho\le1,L\ge8$, $1-e^{-x}$ is comparable to $x$ throughout the
resulting bounded interval. This proves \eqref{f16v49:forbidden-margin}.
All uses of native harmonic localization fix the spatial regulator first.
\end{proof}

The absolute limiting margin of this compression is therefore comparable to
$\lambda_{*,L}\rho/L$, not to a volume-independent number. The scalar
$\lambda_{*,L}$ itself decays exponentially in $L^3$ by
\eqref{f16v49:bulk-asymptotic}. Even after its exact cancellation, the relative
margin still shrinks as $L^{-1}$. In units of the proposed soft scale its
harmonic size is instead comparable to $\rho\gamma^{1/3}$. Thus this result
diagnoses the correct scale; it does not prohibit a scale-resolved growing-volume
argument, nor does it prove that the native one already holds uniformly.

\subsection{The next joint-regulator estimate and preservation}

V49 separates three facts that cannot be merged by a change of normalization:
the bulk Gaussian top is exponential in the number of sites; the hard-band and
fixed-holonomy forbidden margins are of order $L^{-1}$ relative to that top;
and the retained critical scale is $h=\gamma^{-1/3}/L$. The normalized separation
from the hard band in units of $h$ is consequently of order $\gamma^{1/3}$,
independent of the growth of $L$. The explicit harmonic geometry itself does
not destroy that scale separation.

The missing estimate is nevertheless still native. A joint result must pay
quotient-coordinate derivatives, based-orbit inverse constants, the collection
of interacting hard modes, all exceptional regions and the actual same-top
return with errors $o(\lambda_{*,L}h)$ where the proof requires an absolute
transfer estimate. V48's $o_L(h)$ does not have that quantifier. The present
coefficient bounds and native \emph{sequential} compression limit neither supply
a bound on the entry function $\gamma_0(L)$ nor permit exchanging the two limits.
A coupling-dependent tube or a modified retained band must be accompanied by
its own normalization and complement estimate, not inferred from flat holonomy
geometry alone.

The convergent winding potential is also not a new physical action to be added
to a theory that still retains its transverse oscillators. No time step is
chosen. For an independently supplied $a_t(\gamma,L)$, V48's fixed-$L$ neutral
formula still reads $\gamma^{-1/3}(La_t)^{-1}[e_1-e_0+o_L(1)]$.
Neither that remainder nor its interpretation has been made uniform here.
The import memoranda's separate power/locality and physical-transfer rows
remain separate. The next scheduled checkpoint is V50.

The accompanying diagnostic checks the exact heat-kernel and Fourier
normalizations, Poissonization, lattice symbols, coefficient-tail powers,
Gaussian integral coefficient, cusp transform and scale conversions on finite
symbolic or integer fixtures. Optional floating-point winding checks are
reported separately; they are not native spectral or quadrature enclosures.
Only this insertion and the literal title version differ from V48. Removing
them recovers its exact bytes. The existing source assertions, including its
fixed-$L$ spectral theorem, are preserved rather than silently rederived or
altered. The new analytic estimates do not independently recertify that theorem
or the complete prior lineage.

\begin{firewall}
V49 proves lattice-size-uniform winding-coefficient asymptotics, a common
periodic harmonic potential with a controlled cusp, and an explicit bulk and
finite-size formula for the original Gaussian threshold. It also proves the
large-coupling limit of a declared native forbidden compression at every
fixed lattice, revealing a relative margin of order $L^{-1}$. It does not prove
uniform native spectral convergence along growing lattices, supply a physical
time calibration, add an independent cusp potential, identify a continuum
vacuum, or establish an interacting continuum mass gap. All comparisons use
the original Wilson normalization and keep harmonic and native statements
distinct.
\end{firewall}
% END F16 V49 APPEND
% BEGIN F16 V50 APPEND -- 2026-09-18
% Insert before the final document-ending command of cumulative V49.
% Predecessor assertions, physical transfers, Haar measures and clocks are unchanged.

\clearpage
\section{V50: resistance-normalized gauge geometry and a volume-uniform tube}
\label{f16v50:context}

V49 separates the extensive Gaussian top, the harmonic $L^{-1}$ margins, and
the critical scale. It leaves the volume dependence of the based-gauge geometry
unpaid. We address that input rather than assert a uniform version of V48's
spectral remainder. The grounded vertex Laplacian really has inverse norm of
order $L^3$. Its point-evaluation norms in gradient energy, however, are bounded
uniformly on the three-dimensional torus. Using the latter before taking
operator norms removes the apparent inverse-volume loss from the local quotient
metric. A trace-ideal estimate also controls its determinant and the actual
relative Haar density without multiplying by the dimension.

The new native statements concern the existing gauge action and measured tube,
not a new transfer or a joint-limit spectral theorem. Their domain is a small
ball in the \emph{total, unnormalized} link $\ell^2$ norm. That domain has a
radius independent of $L$; it is not a volume-independent bound on each link
separately. An explicit harmonic-packet calculation records when this domain is
occupied. Native spectral concentration on it along growing lattices remains a
separate task.

\subsection{Checkpoint, source ownership, and the unchanged coordinates}

The V49 appendix and audit were read completely. Targeted searches of the current
lineage, Grand Tensor V076, and f14/f15 V100 were followed by exact-text checks.
The earlier resistance calculation in f15 V98 and f16 V2 is for a single cube;
its values and kinetic application are retained, not recomputed as new results.
V44 supplies the based normal-slice construction, and V48 supplies the exact
quotient Schur metric and coarea identity. Their constants were fixed-$L$
constants. The present argument proves size-independent bounds for those same
objects by an explicit torus Green function and Schatten estimates.

Effective resistance and grounded Green identities are classical. Devriendt,
\href{https://arxiv.org/html/2010.04521}{\emph{Effective resistance is more than
distance: Laplacians, Simplices and the Schur complement}}, Sections 2.1 and 2.4,
is primary context for the pseudoinverse/resistance relation. No estimate for a
nonlinear gauge quotient is imported from that paper. All bounds needed below
are proved directly. The supplied mechanism memoranda are not mathematical
premises. The V50 checkpoint distinguishes this geometric advance from the
remaining native kernel, spectral, and physical-time estimates. The next
checkpoint is V55.

Let $L=2m\ge8$, $V=L^3$, and let the oriented positive links of
$\mathbb T_L^3$ have cardinality $3V$. Every vertex and link norm is a counting
$\ell^2$ norm, without division by $V$. Write
\[
 (d\eta)_{xy}=\eta_x-\eta_y,\qquad \Delta_L=d^*d,
 \qquad \eta_o=0,\qquad d_o=d|_{\{\eta_o=0\}},\qquad A_L=d_o^*d_o.
\]
Initially these are one-color matrices. The color copies are tensor products
with $I_3$. The physical horizontal tangent is the unchanged
$\mathcal N_L=\ker(d^*\otimes I_3)$, of dimension $6V+3$. For
$s\in\mathcal N_L$, $\operatorname{Exp}s$ means the componentwise quaternion
exponential, with a unit imaginary quaternion having length one. Retain
\[
 \Phi_L(g,s)=g\cdot\operatorname{Exp}s,\qquad
 g\in\mathcal G_o=\{g\in\SU(2)^V:g_o=I\},\qquad
 dH(\Phi_L(g,s))=dg\,J_L(s)\,ds.
\]
Here both group measures are probability Haar and $ds$ is orthonormal horizontal
Lebesgue measure. The residual simultaneous conjugation is imposed on functions;
it is not removed through a singular orientation section. All constants $c,C$
in this section are independent of $L$ unless explicitly qualified. Schatten
norms $\|\cdot\|_{\mathrm{HS}}$ and $\|\cdot\|_1$ are unnormalized
Hilbert--Schmidt and trace norms.

\subsection{The large grounded inverse and the bounded resistance evaluation}

For $k\in\{0,\ldots,L-1\}^3$ put
\[
 \nu_L(k)=4\sum_{\mu=1}^3\sin^2(\pi k_\mu/L),\qquad
 G_L(x)=\frac1V\sum_{k\ne0}\frac{e^{2\pi i k\cdot x/L}}{\nu_L(k)}.
\]
The function is real, even, and has zero sum. It is the kernel of
$\Delta_L^\dagger$ with its constant mode removed.

\begin{proposition}[Grounded Green function with two different norm scales]
\label{f16v50:grounded}
With the root identified as $0$, the inverse $K_L=A_L^{-1}$ is
\begin{equation}
 K_L(x,y)=G_L(x-y)-G_L(x)-G_L(y)+G_L(0),\qquad x,y\ne0.
 \label{f16v50:grounded-formula}
\end{equation}
There are absolute positive constants such that
\begin{equation}
 \boxed{\quad
 c\le G_L(0)\le C,\qquad
 0\le K_L(x,x)\le C,\qquad
 cL^3\le\|A_L^{-1}\|\le CL^3.
 \quad}
 \label{f16v50:two-scales}
\end{equation}
More precisely,
\begin{equation}
 \Tr K_L=2VG_L(0),\qquad
 \big|\|K_L\|-VG_L(0)\big|\le CL^2.
 \label{f16v50:grounded-leading}
\end{equation}
Every based scalar or vector field satisfies
\begin{equation}
 \boxed{\qquad \|\eta\|_\infty\le C\|d\eta\|_2.\qquad}
 \label{f16v50:anchored-evaluation}
\end{equation}
In particular the grounded least eigenvalue is comparable to $L^{-3}$,
whereas the least positive mean-zero torus eigenvalue is comparable to $L^{-2}$.
\end{proposition}
\begin{proof}
The Fourier formula gives $\Delta_LG_L=\delta_0-V^{-1}$.
The proposed kernel vanishes if either argument is zero. Applying the Laplacian
in $x$ gives $\delta_y-\delta_0$, hence the grounded inverse identity on nonroot
rows. Positivity follows from $A_L\succ0$.
Choose momentum representatives with $|k_\mu|\le L/2$. Then
$16|k|^2/L^2\le\nu_L(k)\le4\pi^2|k|^2/L^2$.
The three-dimensional shell counts give
\[
 \sum_{0<|k|\le CL}|k|^{-2}\le CL,\qquad
 \sum_{k\ne0}|k|^{-4}\le C.
\]
They imply $G_L(0)\le C$. Since $\nu_L\le12$,
$G_L(0)\ge(V-1)/(12V)\ge c$. Positive semidefiniteness of the convolution
kernel gives $|G_L(x)|\le G_L(0)$; therefore
$K_L(x,x)=2[G_L(0)-G_L(x)]\le4G_L(0)\le C$.

Extend $K_L$ by a zero root row and column, and let ${\bf1}$ have $V$ entries.
With $G$ the convolution matrix and $g=(G_L(x))_x$, the exact identity is
\[
 K_L^{\rm ext}=G-g{\bf1}^*-{\bf1}g^*+G_L(0){\bf1}{\bf1}^*.
\]
Here $\|G\|\le CL^2$ and Parseval gives
$\|g\|_2^2=V^{-1}\sum_{k\ne0}\nu_L(k)^{-2}\le CL$.
The first three terms consequently have norm at most $CL^2$; the last has
norm $VG_L(0)$. This proves the sharper norm estimate. Direct summation gives
$\Tr K_L=2VG_L(0)$ and
${\bf1}^*K_L^{\rm ext}{\bf1}=V^2G_L(0)$.
The latter Rayleigh quotient, together with the trace upper bound, gives both
uniform comparisons in \eqref{f16v50:two-scales} without relying on a small
relative asymptotic error at the smallest lattice sizes.

For a based scalar field,
$|\eta_x|=|\langle A_L^{-1/2}e_x,A_L^{1/2}\eta\rangle|
\le\sqrt{K_L(x,x)}\|d\eta\|_2$. Sum over the components of a vector
field. This proves \eqref{f16v50:anchored-evaluation}. The mean-zero least
eigenvalue is exactly $4\sin^2(\pi/L)$ by Fourier diagonalization.
\end{proof}

Thus using $\|A_L^{-1}\|$ in every perturbative estimate would pay a genuine
$L^3$ norm even when the expression only evaluates gauge parameters at incident
vertices. That is unnecessary for the latter operation. The uniform evaluation
bound is a three-dimensional assertion; no dimension-independent graph statement
is being claimed.

\subsection{Energy coordinates for the same based gauge action}

Put $P_L=A_L^{-1/2}\otimes I_3$. This is a fixed linear coordinate change on
the based Lie algebra, not a change of Haar measure or a new gauge group. In
right-trivialized link tangent coordinates let
\begin{equation}
 (Y_s\eta)_{xy}=\eta_x-\operatorname{Ad}_{\operatorname{Exp}s_{xy}}\eta_y,
 \qquad Z_s=Y_sP_L,\qquad Z_0=(d_o\otimes I_3)P_L.
 \label{f16v50:vertical-frame}
\end{equation}
Then $Z_0^*Z_0=I$. Let $X_s$ be the derivative of the same exponential slice
on $\mathcal N_L$, in the same right trivialization; $X_0$ is the horizontal
inclusion. In particular $X_0^*Z_0=0$.

\begin{proposition}[Volume-uniform perturbation of the energy-normalized orbit]
\label{f16v50:vertical}
On a fixed ball $\|s\|_2<r_0$, whose radius is independent of $L$, set
$E_s=Z_s-Z_0$ and $\mathsf V_s=Z_s^*Z_s$. Then
\begin{equation}
 \|E_s\|_{\mathrm{HS}}\le C\|s\|_2,\qquad
 \tfrac12I\preceq\mathsf V_s\preceq\tfrac32I.
 \label{f16v50:vertical-bound}
\end{equation}
For $j=1,2,3$, the derivative bounds
\[
 \|D^jE_s[t_1,\ldots,t_j]\|_{\mathrm{HS}}
                  \le C_j\prod_{i=1}^j\|t_i\|_2
\]
hold with size-independent constants. The same first-order norm bound applies
to $X_s-X_0$. No factor $\|A_L^{-1/2}\|$ occurs in these estimates.
\end{proposition}
\begin{proof}
Let $R_x$ evaluate $P_Lv$ at the vertex $x$, with $R_o=0$. Then
$R_xR_y^*=K_L(x,y)I_3$ and
$\|R_x\|_{\mathrm{HS}}^2=3K_L(x,x)\le C$.
The $xy$ block of $E_s$ is
$(I-\operatorname{Ad}_{\operatorname{Exp}s_{xy}})R_y$.
Its squared Hilbert--Schmidt norm is at most $C|s_{xy}|^2$ because the
adjoint action is rotation through angle $2|s_{xy}|$. Sum these squared block
norms over the actual links. For derivatives, smoothness of the single-group
adjoint map on a fixed ball gives the bound
\[
 C_j\left(\sum_e\prod_i|t_{i,e}|^2\right)^{1/2}
                         \le C_j\prod_i\|t_i\|_2.
\]
The covariance identity for $R_x$ follows from the grounded inverse just proved,
so the count of vertices has not been absorbed in an unspecified constant.

Right-trivialized differentiation of one exponential is
$B(s)=\int_0^1\operatorname{Ad}_{\operatorname{Exp}(ts)}dt$.
Its difference from identity has Hilbert--Schmidt norm at most $C|s|$.
The full differential is the direct sum of these blocks, restricted to the
horizontal subspace. This proves the assertion for $X_s$ and its corresponding
derivatives. Finally
$\mathsf V_s-I=Z_0^*E_s+E_s^*Z_0+E_s^*E_s$.
Choose $r_0$ so that its operator norm is at most $1/2$.
\end{proof}

\begin{theorem}[A size-independent based normal tube]
\label{f16v50:uniform-tube}
For an absolute $r_0>0$, possibly decreased from the preceding proposition,
\[
 \Phi_L:\mathcal G_o\times\{s\in\mathcal N_L:\|s\|_2<r_0\}
                     \longrightarrow\SU(2)^{3V}
\]
is a diffeomorphism onto its image for every admitted $L$. It is the same
normal-slice parametrization as V44. Its eight center translates may be chosen
disjoint with this same radius. Residual global conjugation remains unfixed.
\end{theorem}
\begin{proof}
The derivative with domain gauge norm $\|d\eta\|_2$ is represented at $s$
by $[Z_s\ X_s]$. At zero this is orthogonal; the preceding proposition makes
it invertible on one common ball. Gauge equivariance gives the same conclusion
at every $g$. This proves local invertibility without a bound in the much
larger vertex norm of $A_L^{-1/2}$.

To prove global injectivity on this domain, suppose
$\operatorname{Exp}t=k\cdot\operatorname{Exp}s$ with $k_o=I$ and both
$s,t$ horizontal of norm less than $r_0$. The ordered link relation and
bi-invariance of distance give
\[
 \operatorname{dist}(k_x,k_y)\le |s_{xy}|+|t_{xy}|.
\]
Apply \eqref{f16v50:anchored-evaluation} to the four real components of
$k_x-I$, using $|k_x-k_y|\le\operatorname{dist}(k_x,k_y)$.
It gives $\sup_x|k_x-I|\le Cr_0$. Every $k_x$ therefore lies in one fixed
principal logarithm ball if $r_0$ is small. Write $k_x=\operatorname{Exp}\eta_x$.
The logarithm is uniformly Lipschitz there, so
$\eta_o=0$, $\|d\eta\|_2\le Cr_0$, and
$\|\eta\|_\infty\le C\|d\eta\|_2$.

For $0\le u\le1$ set
\[
 \ell_{xy}(u)=\log\bigl(
   e^{u\eta_x}e^{s_{xy}}e^{-u\eta_y}\bigr).
\]
The product bound and local Lipschitz continuity imply
$\|\ell(u)\|_2\le C(\|s\|_2+\|d\eta\|_2)\le Cr_0$.
All logarithms stay in a common ball. In right trivialization,
\[
 \dot\ell_e=B(\ell_e)^{-1}
       \{(d\eta)_e+(I-\operatorname{Ad}_{e^{\ell_e}})\eta_y\}.
\]
The single-group estimates $\|B(\ell)^{-1}-I\|\le C|\ell|$ and the
anchored evaluation bound yield
$\|\dot\ell-d\eta\|_2\le Cr_0\|d\eta\|_2$.
Since $s,t\in\mathcal N_L$,
\[
 0=\langle d\eta,t-s\rangle
   =\int_0^1\langle d\eta,\dot\ell(u)\rangle du
   \ge(1-Cr_0)\|d\eta\|_2^2.
\]
Thus $\eta=0$, $k=I$, and $s=t$. This proves injectivity and hence a smooth
inverse on the image.

For disjointness, two distinct central-flat orbits differ by a central sign in
some direction. At every one of the $L^2$ disjoint parallel fundamental loops
in that direction, their ordered products have distance $\pi$. The product
triangle inequality and Cauchy--Schwarz give at least $\pi^2/L$ for the sum of
squared link distances on each loop. Consequently the distance between any
representatives of those orbits is at least $\pi\sqrt L$. This is also a
lower bound after independent periodic gauge changes, since central holonomies
are unchanged by conjugation. Each normal tube has distance at most $r_0$
from its orbit. A further absolute decrease makes all eight tubes disjoint.
\end{proof}

\subsection{Trace-ideal control of the original quotient metric}

In the energy coordinates the exact V48 metric is
\begin{equation}
 \mathsf C_s=X_s^*Z_s,\qquad
 \mathsf g_s=X_s^*X_s-\mathsf C_s\mathsf V_s^{-1}\mathsf C_s^*.
 \label{f16v50:metric}
\end{equation}
Congruence by $P_L$ in the vertical block leaves this Schur complement unchanged.
Thus it is the original quotient metric, not a redefined physical metric.

\begin{theorem}[Dimension-free trace bounds and two controlled derivatives]
\label{f16v50:trace-metric}
On the common tube,
\begin{equation}
 \boxed{\quad
 0\preceq I-\mathsf g_s,\qquad
 \|I-\mathsf g_s\|_1+\|\mathsf g_s^{-1}-I\|_1
                  \le C\|s\|_2^2,\qquad
 \tfrac12I\preceq\mathsf g_s\preceq I.
 \quad}
 \label{f16v50:metric-main}
\end{equation}
For horizontal directions $t,v$,
\begin{equation}
 \|D\mathsf g_s[t]\|_1\le C\|s\|_2\|t\|_2,\qquad
 \|D^2\mathsf g_s[t,v]\|_1\le C\|t\|_2\|v\|_2.
 \label{f16v50:metric-derivatives}
\end{equation}
The constants have no factor equal to the horizontal dimension $6V+3$.
\end{theorem}
\begin{proof}
For one round group the exponential metric has radial eigenvalue one and two
transverse eigenvalues $(\sin|s_e|/|s_e|)^2$. Hence, after restriction to
$\mathcal N_L$,
\[
 0\preceq I-X_s^*X_s,\qquad
 \|I-X_s^*X_s\|_1
 \le2\sum_e[1-(\sin|s_e|/|s_e|)^2]\le C\|s\|_2^2.
\]
Compression by an isometry does not increase the trace norm of this positive
operator. The first and second directional trace bounds are at most
$C\sum_e|s_e||t_e|$ and $C\sum_e|t_e||v_e|$, by differentiating the
smooth single-group metric. Cauchy--Schwarz supplies the desired norms.

Because $X_0^*Z_0=0$, Proposition~\ref{f16v50:vertical} gives
$\|\mathsf C_s\|_{\mathrm{HS}}\le C\|s\|_2$,
$\|D\mathsf C_s[t]\|_{\mathrm{HS}}\le C\|t\|_2$, and a similarly
bounded second derivative. The inverse $\mathsf V_s^{-1}$ and its first two
operator-norm derivatives are bounded by the same proposition and the inverse
rule. The positive return has trace norm
\[
 \|\mathsf C_s\mathsf V_s^{-1}\mathsf C_s^*\|_1
       \le2\|\mathsf C_s\|_{\mathrm{HS}}^2\le C\|s\|_2^2.
\]
Leibniz's rule and $\|AB^*\|_1\le\|A\|_{\mathrm{HS}}\|B\|_{\mathrm{HS}}$
prove its differentiated bounds. Insert these in \eqref{f16v50:metric}.
Decrease $r_0$ until $\|I-\mathsf g_s\|\le1/2$; the inverse trace estimate
then follows by diagonalization, since $\mathsf g_s^{-1}-I$ is positive and
at most twice $I-\mathsf g_s$ on each eigenvalue.
\end{proof}

\subsection{The relative Haar density, without a hidden dimension factor}

The absolute value $J_L(0)$ need not be uniform in $L$. It is retained exactly.
\begin{theorem}[Uniform relative coarea density and its exact central scalar]
\label{f16v50:density}
On the same tube, the actual density satisfies
\begin{equation}
 \boxed{\begin{aligned}
 \left|\log\frac{J_L(s)}{J_L(0)}\right|&\le C\|s\|_2^2,\\
 \left|D\log J_L(s)[t]\right|&\le C\|s\|_2\|t\|_2,\\
 \left|D^2\log J_L(s)[t,v]\right|&\le C\|t\|_2\|v\|_2.
 \end{aligned}}
 \label{f16v50:density-main}
\end{equation}
In particular the density has zero first jet at the central configuration.
With the original probability-Haar convention,
\begin{equation}
 \boxed{\quad
 J_L(0)=(2\pi^2)^{-(2V+1)}(\det A_L)^{3/2},\qquad
 \det A_L=\frac1V\prod_{k\ne0}\nu_L(k).
 \quad}
 \label{f16v50:absolute-density}
\end{equation}
These are relative geometric bounds and an exact scalar identity, not a
statement that an extensive native partition mass is one.
\end{theorem}
\begin{proof}
The same Gram determinant used in V48, now divided by its value at zero, gives
\begin{equation}
 \frac{J_L(s)}{J_L(0)}
             =\sqrt{\det\mathsf V_s}\sqrt{\det\mathsf g_s}.
 \label{f16v50:relative-coarea}
\end{equation}
The horizontal logarithmic determinant and its two derivatives are controlled
by Theorem~\ref{f16v50:trace-metric}. For example
$|\log\det\mathsf g_s|\le2\|I-\mathsf g_s\|_1$.
A crude dimension times operator norm would not suffice for the vertical term,
so we record its trace cancellation.

Put $D_s=\mathsf V_s-I$. Its Hilbert--Schmidt norm is at most $C\|s\|_2$,
and $\|D_s\|\le1/2$. By the evaluation identities in the vertical proof,
\begin{equation}
 \Tr(Z_0^*E_s)
   =\sum_{e=xy}\bigl[K_L(x,y)-K_L(y,y)\bigr]
       \operatorname{tr}_{\mathbb R^3}
                  (I-\operatorname{Ad}_{e^{s_e}}).
 \label{f16v50:vertical-trace}
\end{equation}
Root rows of $K_L$ are zero in this formula. Every coefficient is bounded by
$C$ using the diagonal bound and positive definiteness. The last trace equals
$4\sin^2|s_e|$. Therefore
\[
 |\Tr D_s|\le C\|s\|_2^2,\quad
 |\Tr D_s'[t]|\le C\|s\|_2\|t\|_2,\quad
 |\Tr D_s''[t,v]|\le C\|t\|_2\|v\|_2.
\]
The $E_s^*E_s$ contribution follows from Hilbert--Schmidt Cauchy--Schwarz;
the displayed scalar trace formula pays the terms containing $Z_0$.
Indeed orthogonality of the adjoint rotations also gives the exact simplification
\[
                    \Tr D_s=8\sum_{e=xy}K_L(x,y)\sin^2|s_e|.
\]
Its two derivatives have the bounds just used.
For a self-adjoint $D$ with $\|D\|\le1/2$,
\[
 |\log\det(I+D)-\Tr D|\le C\|D\|_{\mathrm{HS}}^2.
\]
This proves the undifferentiated vertical bound. For its first derivative,
split $(I+D_s)^{-1}$ into $I+[(I+D_s)^{-1}-I]$ before taking its trace
against $D_s'[t]$. The second factor has Hilbert--Schmidt norm $O(\|s\|_2)$,
while $D_s'[t]$ has that norm $O(\|t\|_2)$. Differentiating once more gives
\[
 \Tr[(I+D_s)^{-1}D_s''[t,v]]
 -\Tr[(I+D_s)^{-1}D_s'[v](I+D_s)^{-1}D_s'[t]].
\]
The same split, the just-proved trace bounds, and Schatten Cauchy--Schwarz
bound it by $C\|t\|_2\|v\|_2$. Combine the two determinants in
\eqref{f16v50:relative-coarea}.

At zero, standard based vertex coordinates have Haar density
$(2\pi^2)^{-(V-1)}$, and all $3V$ link groups have density
$(2\pi^2)^{-3V}$. The Gram of the based orbit differential is $A_L\otimes I_3$,
whereas the horizontal columns are orthonormal and orthogonal to it. Their
volume determinant is $(\det A_L)^{3/2}$. This proves the first equality in
\eqref{f16v50:absolute-density}. The cofactor identity for a Laplacian with
kernel $\mathbb C{\bf1}$ gives its second equality: its adjugate is
$[V^{-1}\prod_{k\ne0}\nu_L(k)]{\bf1}{\bf1}^*$, as follows directly by
spectrally diagonalizing $\Delta_L+tI$ and taking $t\downarrow0$.
The grounded determinant is the root cofactor.
\end{proof}

\begin{proposition}[Normalized tangent Gaussians have a dimension-free error]
\label{f16v50:Gaussian}
For any $\varepsilon>0$ let
$\mathcal G_s=N(0,\varepsilon^2\mathsf g_s^{-1})$ on the horizontal tangent
and $\mathcal G_0=N(0,\varepsilon^2I)$. Then
\begin{equation}
 D(\mathcal G_s\Vert\mathcal G_0)\le C\|s\|_2^4,\qquad
 d_{\mathrm{TV}}(\mathcal G_s,\mathcal G_0)\le C\|s\|_2^2.
 \label{f16v50:Gaussian-bound}
\end{equation}
Both probabilities include their own determinants. They describe the normalized
quadratic tangent kernel only, not the complete native quotient transition.
\end{proposition}
\begin{proof}
Diagonalize $\mathsf g_s$ with eigenvalues $\alpha_j\in[1/2,1]$.
The exact relative entropy is
$\frac12\sum_j(\alpha_j^{-1}-1+\log\alpha_j)$.
The scalar summand is bounded by $C(1-\alpha_j)^2$ on this interval. Hence
it is at most $C\|I-\mathsf g_s\|_{\mathrm{HS}}^2\le C\|s\|_2^4$.
The total-variation assertion follows from Pinsker's inequality, or its
elementary density proof. The scale $\varepsilon$ cancels exactly.
\end{proof}

\subsection{An explicit geometric domain for growing harmonic packets}

This subsection uses the declared hard Gaussian ground density only. It is not
a bound for the native vacuum at growing $L$. Let $\varphi_L$ be V44's normalized
hard ground function on $\mathcal Y$, and let $u$ have probability density
$|\varphi_L|^2$. Write its covariance as $\Sigma_{+,L}$. The hard Hessian has
$\nu_L(k)$, $k\ne0$, each with six real multiplicities. The covariance eigenvalues
are
\[
 \sigma(k)=\frac1{2\sqrt{\nu_L(k)+\nu_L(k)^2/4}}.
\]
Thus
\begin{equation}
 cV\le\Tr\Sigma_{+,L}\le CV,\qquad
 \|\Sigma_{+,L}\|\le CL,
 \qquad
 \log\mathbb E e^{c|u|^2/L}\le CL^2
 \label{f16v50:Gaussian-moments}
\end{equation}
for a sufficiently small absolute $c>0$. The shell estimate
$\sum_{0<|k|\le CL}|k|^{-1}\le CL^2$ proves the trace bound; the least
momentum proves the norm bound. The last inequality is the exact Gaussian
moment determinant and $-\log(1-x)\le2x$ for $0\le x\le1/2$.

\begin{proposition}[Harmonic occupation and measured density on a joint domain]
\label{f16v50:packet}
Retain $\varepsilon=\gamma^{-1/2}$, $h=\gamma^{-1/3}/L$ and constant
extension $\mathsf K^*\mathsf K=VI$. For $|q|\le R$, put
$s=\varepsilon u+h\mathsf Kq$. Then exactly
\begin{equation}
 \|s\|_2^2=\gamma^{-1}|u|^2+L\gamma^{-2/3}|q|^2.
 \label{f16v50:packet-norm}
\end{equation}
If $L\gamma^{-2/3}R^2\le r_0^2/2$, uniformly in such $q$,
\begin{align}
 \mathbb P\{\|s\|_2\ge r_0\}
   &\le \min\{1,\exp(CL^2-c\gamma r_0^2/L)\},
       \label{f16v50:packet-tail}\\
 \mathbb E\!\left[1_{\{\|s\|<r_0\}}
       \left|\frac{J_L(s)}{J_L(0)}-1\right|\right]
   &\le C\left(\frac{L^3}{\gamma}
                       +L\gamma^{-2/3}R^2\right).
       \label{f16v50:packet-density}
\end{align}
In particular for each fixed $R$, the tube is occupied with probability tending
to one and the displayed relative-density error tends to zero if
$L^3/\gamma\to0$. This includes $L=O(\gamma^b)$ with fixed $b<1/3$.
\end{proposition}
\begin{proof}
Hard and constant subspaces are orthogonal, so \eqref{f16v50:packet-norm}
is exact. The assumed soft bound implies that leaving the tube requires
$|u|^2\ge\gamma r_0^2/2$. Apply the exponential moment in
\eqref{f16v50:Gaussian-moments}, adjusting the absolute constant $c$.
On the tube, Theorem~\ref{f16v50:density} and $|e^x-1|\le e^{|x|}|x|$
give $|J_L(s)/J_L(0)-1|\le C\|s\|_2^2$. Take its expectation using the
trace bound. If $L^3/\gamma\to0$, the negative tail exponent dominates
$CL^2$, and $L\gamma^{-2/3}\to0$ as well. All probabilities in this proof
are under the specified hard Gaussian, not an unknown native distribution.
\end{proof}

\subsection{V50 checkpoint: the geometric loss removed and the losses retained}

The large norm $\|A_L^{-1}\|\asymp L^3$ has not been denied or replaced.
It is avoided where the actual operation is incident-vertex evaluation in
gradient energy. This yields a uniform normal tube, a trace-class quadratic
metric error, and a uniform relative coarea density for the original based
quotient. The absolute scalar $J_L(0)$ remains explicit. No coordinate change
has set a native partition function or physical Perron value to one.

The new radius is in total link $\ell^2$. On a constant field it imposes
$V|a|^2<r_0^2$, not a fixed per-link window. The harmonic packet condition is
therefore an honest domain statement, not a volume-uniform localization theorem
for the native spectrum. It neither transports V48's eigenvectors to a joint
limit nor bounds their mass outside this tube. The hard Mehler gap still shrinks
as $L^{-1}$, exactly as V49 states. Its derivative and perturbation costs, the
sum of interacting magnetic remainders, displacement tails in growing dimension,
and the full returning spectral operator must still be controlled on the same
trajectory. A uniform principal metric is not a uniform full kernel expansion.

In particular $b<1/3$ above is a sufficient \emph{harmonic geometric} exponent.
It is not a native spectral exponent. This section does not improve or recertify
V48's fixed-$L$ spectral theorem, and it does not supply a joint bound on its
entry coupling or its $o_L(h)$ remainder. The next useful estimate is an assembled
native normalized-kernel or form bound in these energy coordinates, with each
remaining trace and hard-mode cost recorded. It should keep the original
magnetic words, the exact scalar top, and the returning-state metric rather
than infer their control from the present Gaussian calculation.

The checkpoint file records the V41--V50 dependency split and unresolved limits.
No assertion about the externally discussed Navier--Stokes result is needed.
The methodological distinction between residual power and derivative/locality
cost is used here literally: one geometric cost is now uniform, while the
spectral and physical-time rows are still open. With an independently supplied
$a_t(\gamma,L)$, the old fixed-$L$ expression remains
$\gamma^{-1/3}(La_t)^{-1}[e_1-e_0+o_L(1)]$; this section does not make that
expression a calibrated continuum mass.

\subsection{Verification and preservation}

The accompanying diagnostic tests grounded Green identities and resistance
normalization on exact periodic fixtures, including admitted $L=8$ data; it also
tests the vertical trace cancellation, metric/coarea jets and Gaussian scales.
The proof uses the analytic Fourier shell bounds and Schatten estimates, not a
finite sample to assert an all-$L$ inequality. The audit distinguishes exact
finite checks from any unproved native spectral statement.

Only this terminal insertion and the literal title version change V49. Removing
both recovers its bytes exactly. No predecessor theorem or comparison convention
is silently edited. The complete earlier analytic lineage is not independently
recertified. The next scheduled checkpoint is V55.

\begin{firewall}
V50 proves size-independent geometric estimates on a stated total-norm tube of
the existing compact based-gauge action, including the measured density and
its absolute central normalization. Its growing-lattice packet statement is
harmonic, not native. It does not prove uniform native spectral convergence,
calibrate physical time, remove all cutoff or volume restrictions, or establish
an interacting continuum Yang--Mills mass gap. No auxiliary gauge-energy norm
is substituted for the physical link metric or for a transfer clock.
\end{firewall}
% END F16 V50 APPEND

% BEGIN F16 V51 APPEND -- 2026-09-18
% Insert before the final document-ending command of cumulative V50.
% The original Wilson transfer and all predecessor source bodies are retained.

\clearpage
\section{V51: a volume-free Wilson--heat comparison and the remaining Perron cost}
\label{f16v51:context}

V50 controls the measured gauge geometry, but not the assembled native kinetic
error on a growing lattice. We first avoid that coordinate expansion altogether.
The original one-link convolution can be compared with the round-group heat
operator in a \emph{relative deficit} norm, uniformly over every representation.
That norm tensorizes without counting links. After the original magnetic
sandwich it still controls a relative form, and a discrete telescoping argument
gives an operator-norm comparison uniform in the number of unnormalized transfer
steps. No small-field condition, gauge chart, or restriction on lattice size is
needed for these estimates.

The heat operator below is an explicitly defined comparison, not a redefinition
of Wilson transfer. Its magnetic action is the original action. The new bounds
are not automatically relative to the exponentially small Perron value. We pay
that separate scalar cost and give a conservative joint-regulator regime where
the two \emph{actual-top normalized} evolutions agree at the existing critical
scale. This compares two lattice operators; it does not identify a continuum
limit or prove a uniform version of V48's quartic theorem.

\subsection{Inputs, classical formulas, and the declared comparison}

The V50 appendix, audit and checkpoint were read completely. Targeted searches
of the current source and supplied Grand Tensor and f14/f15 archives were
followed by exact-text inspection. The Grand Tensor's fixed one-link Casimir
moments and private-bank bounds are related precedents, not the all-representation
relative comparison proved here. Its \emph{All fixed one-link Casimir moments}
and \emph{Failure debits and private-product moments} passages were checked;
the supplied f15 V91 controls a different electric-tail problem. No absence or
originality claim is inferred from empty semantic-search results. The present
proof does not assume V48's spectral completeness or V50's uniform tube.

For the scalar identities we use the integer Bessel integral and recurrence
\href{https://dlmf.nist.gov/10.32.E3}{DLMF 10.32.3} and
\href{https://dlmf.nist.gov/10.29.E1}{10.29.1}. The heat representation is the
round-sphere formula, consistent with the dimension recurrence and periodized
one-dimensional Gaussian in Nowak--Sj\"ogren--Szarek,
\href{https://arxiv.org/html/1802.09385}{\emph{Sharp estimates of the spherical
heat kernel}}, equations (1)--(2). We derive its probability-Haar normalization
below. No fixed-order Bessel asymptotic is applied uniformly in an increasing
representation label. The tensor and relative-form arguments are elementary
operator comparisons, not new definitions of a physical clock. The next
scheduled checkpoint remains V55.

Let $G=\SU(2)$ with the unit-round quaternion metric and probability Haar
measure. Write $\theta(g)=\operatorname{dist}(I,g)\in[0,\pi]$. Retain exactly
\[
 C_\gamma f(U)=\int_G k_\gamma(UV^{-1})f(V)\,dH(V),\qquad
 k_\gamma(g)=z_\gamma^{-1}e^{-\gamma(1-\cos\theta(g))},\qquad
 z_\gamma=\int_Ge^{-\gamma(1-\cos\theta)}\,dH.
\]
Let $\mathcal L_G=-\Delta_G\ge0$ be the round Laplacian, and \emph{define}
\begin{equation}
 Q_\gamma=e^{-\mathcal L_G/(2\gamma)}.
 \label{f16v51:heat-comparison}
\end{equation}
Its time $1/(2\gamma)$ matches the covariance of the original small Wilson
step. It is not calibrated physical Euclidean time. On the irreducible matrix
coefficients with highest-weight integer $\ell\ge0$, the Laplace eigenvalue is
$s_\ell=\ell(\ell+2)$. The corresponding spin is $\ell/2$; this convention
avoids a factor-four ambiguity in the Casimir.

\subsection{Normalized radial errors with one vanishing angular moment}

For class functions the probability-Haar angle density is
$(2/\pi)\sin^2\theta\,d\theta$. Denote the radial probability densities of
$k_\gamma$ and $Q_\gamma$ by $w_\gamma$ and $v_\gamma$, respectively, extending
them by zero to $\theta>\pi$.

\begin{proposition}[Two normalized radial error bounds]
\label{f16v51:radial}
There are absolute constants $C_j$ such that, for $\gamma\ge1$ and every fixed
integer $j\ge0$,
\begin{equation}
 \int_0^\infty\theta^{2j}|w_\gamma(\theta)-v_\gamma(\theta)|\,d\theta
       \le C_j\gamma^{-j-1}.
 \label{f16v51:radial-moments}
\end{equation}
In particular the unweighted error is $O(\gamma^{-1})$, and its angular
second moment is $O(\gamma^{-2})$. Both densities have their \emph{own exact}
normalization, and their difference has integral zero.
\end{proposition}
\begin{proof}
Use the normalized Maxwell density on the half-line
\[
 m_\gamma(\theta)=\sqrt{2/\pi}\,\gamma^{3/2}\theta^2
                  e^{-\gamma\theta^2/2},\qquad
 \beta_\gamma=\sqrt{2/\pi}\,\gamma^{-3/2}.
\]
On $[0,\pi]$, both $1-\cos\theta$ and $\theta^2/2$ are at least
$c\theta^2$, their difference is at most $C\theta^4$, and
$|\sin^2\theta-\theta^2|\le C\theta^4$. The integral mean-value identity for
the exponential therefore gives
\[
 \left|e^{-\gamma(1-\cos\theta)}\sin^2\theta
                    -e^{-\gamma\theta^2/2}\theta^2\right|
 \le C(\theta^4+\gamma\theta^6)e^{-c\gamma\theta^2}.
\]
Its $\theta^{2j}$ integral is $O_j(\gamma^{-j-5/2})$. The omitted Maxwell
tail beyond $\pi$ has a smaller bound of this order. Taking $j=0$ yields
$z_\gamma=\beta_\gamma[1+O(\gamma^{-1})]$, with a positive lower ratio for
$\gamma\ge1$ after adjustment on a compact coupling interval. Divide by the
actual $z_\gamma$, and retain the difference of the two normalizers. It follows
that
\begin{equation}
 \int\theta^{2j}|w_\gamma-m_\gamma|\,d\theta
                    \le C_j\gamma^{-j-1}.
 \label{f16v51:Wilson-Maxwell}
\end{equation}
This argument does not normalize an approximate density and call its mass one.

For the heat density put $t=(2\gamma)^{-1}$. The character expansion of the
kernel relative to probability Haar is
\[
 q_t(\theta)=\frac{e^t}{\sin\theta}
                 \sum_{n=1}^\infty n e^{-tn^2}\sin(n\theta).
\]
Differentiate the absolutely convergent periodized Gaussian identity
\[
 \sum_{n\in\mathbb Z}e^{-tn^2}e^{in\theta}
 =\sqrt{\pi/t}\sum_{k\in\mathbb Z}e^{-(\theta+2\pi k)^2/(4t)}.
\]
Thus, including the radial Haar density,
\begin{equation}
 v_\gamma(\theta)=
 \sqrt{2/\pi}\,\gamma^{3/2}e^{1/(2\gamma)}\sin\theta
 \sum_{k\in\mathbb Z}(\theta+2\pi k)
                  e^{-\gamma(\theta+2\pi k)^2/2},\quad 0<\theta<\pi.
 \label{f16v51:radial-heat}
\end{equation}
The endpoints are interpreted by limits. This is the kernel of the positive
heat semigroup; individual image terms need not be positive. Its total mass is
one because the trivial character is unchanged.

The $k=0$ term is $e^{1/(2\gamma)}(\sin\theta/\theta)m_\gamma(\theta)$.
Its difference from $m_\gamma$ has the bound in
\eqref{f16v51:Wilson-Maxwell}, using
$|1-\sin\theta/\theta|\le C\theta^2$ and the Maxwell moments. For $k\ne0$,
$|\theta+2\pi k|\ge\pi$ on $[0,\pi]$. Summing the absolute image integrals,
with the bounded factor $\theta^{2j}\sin\theta$, gives at most
$C_j\gamma^{3/2}e^{-c\gamma}$. The Maxwell tail has the same harmless
exponential control. Combining the two comparisons proves
\eqref{f16v51:radial-moments}. No antipodal image or native chart complement
has been dropped.
\end{proof}

\begin{theorem}[Relative kinetic deficit, uniform over every representation]
\label{f16v51:one-link}
The original Wilson multiplier on label $\ell$ is
\[
 c_{\gamma,\ell}=\frac{I_{\ell+1}(\gamma)}{I_1(\gamma)},\qquad
 q_{\gamma,\ell}=e^{-s_\ell/(2\gamma)}
\]
for the heat comparison. There are absolute $C_*,\gamma_*>0$ such that, with
$\delta_\gamma=C_*/\gamma\le1/2$ for $\gamma\ge\gamma_*$,
\begin{equation}
 \boxed{\quad
 |c_{\gamma,\ell}-q_{\gamma,\ell}|
       \le\delta_\gamma(1-q_{\gamma,\ell}),\qquad \ell\ge0.
 \quad}
 \label{f16v51:scalar-deficit}
\end{equation}
Equivalently, on the whole one-link $L^2(G)$,
\begin{equation}
 (1-\delta_\gamma)(I-Q_\gamma)
 \preceq I-C_\gamma
 \preceq(1+\delta_\gamma)(I-Q_\gamma).
 \label{f16v51:one-link-form}
\end{equation}
There is no representation cutoff, and the zero on the constant representation
is exact.
\end{theorem}
\begin{proof}
The normalized character is
\[
 b_\ell(\theta)=\frac{\sin((\ell+1)\theta)}{(\ell+1)\sin\theta}
       =\frac1{\ell+1}\sum_{j=0}^{\ell}e^{i(\ell-2j)\theta}.
\]
It is real and obeys
\begin{equation}
 0\le1-b_\ell(\theta)\le
            \min\{2,s_\ell\theta^2/6\}.
 \label{f16v51:character-vanishing}
\end{equation}
Indeed average $1-\cos((\ell-2j)\theta)$ and use
$(\ell+1)^{-1}\sum_j(\ell-2j)^2=s_\ell/3$.

Integrating $b_\ell$ against the actual Wilson radial density gives
\[
 c_{\gamma,\ell}=
 \frac{I_\ell(\gamma)-I_{\ell+2}(\gamma)}
                {(\ell+1)[I_0(\gamma)-I_2(\gamma)]}
 =\frac{I_{\ell+1}(\gamma)}{I_1(\gamma)}.
\]
This uses the exact Bessel integral and recurrence, not their asymptotic
expansion. The positive Bessel series gives $c_{\gamma,\ell}>0$;
$C_\gamma$ is a probability convolution and hence a contraction, so
$c_{\gamma,\ell}\le1$. The heat multiplier follows from the round Laplacian.

Subtract the two normalized integrals. Their zero mass permits replacing
$b_\ell$ by $b_\ell-1$. Proposition~\ref{f16v51:radial} with $j=0,1$
therefore gives
\[
 |c_{\gamma,\ell}-q_{\gamma,\ell}|
 \le C\min\{\gamma^{-1},s_\ell\gamma^{-2}\}
 =\frac C\gamma\min\{1,s_\ell/\gamma\}.
\]
The elementary inequality $1-e^{-x/2}\ge c\min\{x,1\}$ proves
\eqref{f16v51:scalar-deficit}, including $\ell=0$. Increase the fixed
coupling threshold until $\delta_\gamma\le1/2$. Peter--Weyl orthogonality
then proves the operator inequality on arbitrary superpositions of all labels.
The constants are independent of the label; a fixed-label expansion would not
have justified this step.
\end{proof}

\subsection{A deficit tensorizes without accumulating one error per link}

\begin{proposition}[A product-deficit comparison]
\label{f16v51:product-lemma}
Let $a_i,b_i\in[0,1]$, $1\le i\le M$, and $0\le\delta<1$. If
$(1-\delta)b_i\le a_i\le(1+\delta)b_i$ for all $i$, then
\begin{equation}
 (1-\delta)\left[1-\prod_i(1-b_i)\right]
 \le1-\prod_i(1-a_i)
 \le(1+\delta)\left[1-\prod_i(1-b_i)\right].
 \label{f16v51:product-deficit}
\end{equation}
The constants do not depend on $M$ or on repeated factors.
\end{proposition}
\begin{proof}
The conclusion is immediate if all $b_i=0$. Otherwise define
$f(t)=1-\prod_i(1-tb_i)$ on $0\le t\le(\max b_i)^{-1}$. On this interval
$f''(t)\le0$ and $f(0)=0$. Concavity gives
$f(1-\delta)\ge(1-\delta)f(1)$, and monotonicity in the individual deficits
gives the lower bound. For the upper bound, if
$(1+\delta)\max b_i\ge1$, its right side is at least one. Otherwise
$f(1+\delta)\le(1+\delta)f(1)$ by the same concavity, and individual
monotonicity proves the upper bound. The possible saturation at one is thus
included, rather than continuing a product through negative factors.
\end{proof}

For any finite number $M$ of links, put
\[
 K_{\gamma,M}=C_\gamma^{\otimes M},\qquad
 Q_{\gamma,M}=Q_\gamma^{\otimes M}
       =\exp[-(2\gamma)^{-1}\textstyle\sum_{e=1}^M\mathcal L_e].
\]
\begin{theorem}[All-link, all-kinetic-duration comparison]
\label{f16v51:tensor}
For every $M\ge1$, every integer $N\ge1$, and $\gamma\ge\gamma_*$,
\begin{equation}
 \boxed{\quad
 (1-\delta_\gamma)(I-Q_{\gamma,M}^{\,N})
 \preceq I-K_{\gamma,M}^{\,N}
 \preceq(1+\delta_\gamma)(I-Q_{\gamma,M}^{\,N}).
 \quad}
 \label{f16v51:tensor-form}
\end{equation}
In particular
\begin{equation}
 \sup_{M,N\ge1}\|K_{\gamma,M}^{\,N}-Q_{\gamma,M}^{\,N}\|
                          \le\delta_\gamma.
 \label{f16v51:tensor-norm}
\end{equation}
The same assertion holds for different couplings on the links, with
$\delta=\max_e\delta_{\gamma_e}$. Gauge constraints and representation
multiplicities introduce no additional factor.
\end{theorem}
\begin{proof}
On a product Peter--Weyl label use Proposition~\ref{f16v51:product-lemma}
with $a_e=1-c_{\gamma,\ell_e}$ and $b_e=1-q_{\gamma,\ell_e}$. For duration
$N$ repeat each factor $N$ times. The conclusion is uniform before taking the
supremum over labels, and therefore applies to arbitrary vectors. The difference
is bounded in both signs by $\delta_\gamma(I-Q_{\gamma,M}^{\,N})$, whose norm
is at most $\delta_\gamma$. Different couplings use the same scalar argument.
The kernels are bi-invariant on each link, so they commute with the original
vertex gauge action. Restriction of the established quadratic-form inequality
to its invariant subspace is legitimate; no independence of constrained field
coordinates is asserted.
\end{proof}

\subsection{The original magnetic action and arbitrarily many transfer steps}

For the admitted fine torus let $M=3V=3L^3$ and keep the literal spatial action
$S(U)$ and multiplier $b_\gamma=e^{-\gamma S/2}$. Distinguish the two operators
\begin{equation}
 T_{\gamma,L}^{\rm W}=M_{b_\gamma}K_{\gamma,M}M_{b_\gamma}=T_{\gamma,L},
 \qquad
 T_{\gamma,L}^{\rm H}=M_{b_\gamma}Q_{\gamma,M}M_{b_\gamma}.
 \label{f16v51:two-transfers}
\end{equation}
The first is unchanged. The second changes only the temporal kinetic factor;
it is not an operator with a substituted quadratic magnetic potential. Both
are positive compact self-adjoint contractions, with strictly positive kernels,
on the same compact Haar space. Both preserve the same gauge and center sectors.

\begin{theorem}[A full magnetic form comparison, uniform in volume]
\label{f16v51:magnetic}
At every admitted $L$ and $\gamma\ge\gamma_*$,
\begin{equation}
 \boxed{\quad
 (1-\delta_\gamma)(I-T_{\gamma,L}^{\rm H})
 \preceq I-T_{\gamma,L}^{\rm W}
 \preceq(1+\delta_\gamma)(I-T_{\gamma,L}^{\rm H}).
 \quad}
 \label{f16v51:magnetic-form}
\end{equation}
In particular $\|T_{\gamma,L}^{\rm W}-T_{\gamma,L}^{\rm H}\|
\le\delta_\gamma$. These statements hold on the complete physical space and
each center sector. They survive any \emph{same} orthogonal compression or
isometric change of coordinates on both sides, with the identity on the
resulting space. No typical-field restriction is required.
\end{theorem}
\begin{proof}
For any self-adjoint contraction $B$ and positive contraction $K$,
\[
 I-BKB=(I-B^2)+B(I-K)B.
\]
Apply \eqref{f16v51:tensor-form} at $N=1$ and take the congruence by
$B=M_{b_\gamma}$. The additional term $I-B^2$ is nonnegative and identical
for both operators. This proves \eqref{f16v51:magnetic-form}. It does not
require commutation of magnetism with kinetic convolution. The difference is
bounded in both signs by $\delta_\gamma(I-T^{\rm H})\preceq\delta_\gamma I$.
Restriction, compression and isometric pullback preserve the form inequalities.
They do not introduce a row normalization or replace the inherited exact Haar
half-density. Every original plaquette remains inside $b_\gamma$.
\end{proof}

A crude telescoping of the last norm estimate would cost $N\delta_\gamma$.
The relative form gives more even though the two magnetic transfers need not
commute.
\begin{proposition}[Relative deficits control all unnormalized discrete powers]
\label{f16v51:powers-lemma}
Let $0\preceq T_1,T_2\preceq I$ and suppose
$(1-\delta)(I-T_2)\preceq I-T_1\preceq(1+\delta)(I-T_2)$, with $\delta<1$.
Then for every integer $N\ge1$,
\begin{equation}
 \boxed{\quad
 \|T_1^N-T_2^N\|\le\frac{\pi\delta}{2\sqrt{1-\delta}}.
 \quad}
 \label{f16v51:powers-general}
\end{equation}
Consequently
\begin{equation}
 \sup_{L\in2\mathbb N,\ L\ge8}\ \sup_{N\ge1}
 \|(T_{\gamma,L}^{\rm W})^N-(T_{\gamma,L}^{\rm H})^N\|
                         \le C/\gamma.
 \label{f16v51:all-powers}
\end{equation}
These are unnormalized powers; no vacuum subtraction is included in their name.
\end{proposition}
\begin{proof}
Put $A=I-T_1$, $B=I-T_2$. The form bound gives
$A-B=B^{1/2}EB^{1/2}$ with a bounded self-adjoint $E$ on the closure of
$\operatorname{Ran}B$, $\|E\|\le\delta$. This follows by extending the bounded
form defined first on that range; it does not require $B$ invertible. The
inequality $B\preceq(1-\delta)^{-1}A$ similarly gives a factorization
$B^{1/2}=A^{1/2}R$ with $\|R\|\le(1-\delta)^{-1/2}$.
For $k\ge0$, spectral calculus implies
\[
 \|T_1^kB^{1/2}\|\le\frac{(1-\delta)^{-1/2}}{\sqrt{2k+1}},
 \qquad
 \|B^{1/2}T_2^k\|\le\frac1{\sqrt{2k+1}},
\]
since $\sqrt{x}(1-x)^k\le(2k+1)^{-1/2}$ on $[0,1]$.
In the exact telescoping identity insert the bounded form factorization:
\[
 T_1^N-T_2^N=-\sum_{j=0}^{N-1}
 T_1^{N-1-j}B^{1/2}E B^{1/2}T_2^j.
\]
Its norm is at most
\[
 \frac{\delta}{\sqrt{1-\delta}}
 \sum_{j=0}^{N-1}\frac1{\sqrt{(2j+1)(2N-2j-1)}}
 \le\frac{\pi\delta}{2\sqrt{1-\delta}}.
\]
For the last bound, $[x(N-x)]^{-1/2}$ is convex on $(0,N)$; apply the
midpoint integral bound on each unit interval and use
$\int_0^N[x(N-x)]^{-1/2}dx=\pi$. This includes $N=1$.
No pairwise commutation of $A$ and $B$ is used. Apply the result to
Theorem~\ref{f16v51:magnetic} and $\delta_\gamma\le1/2$.
\end{proof}

\subsection{The Perron scalar still costs something different}

Let $\Lambda_{\gamma,L}^{\rm W}$ and $\Lambda_{\gamma,L}^{\rm H}$ be the actual
Perron values of the two explicitly declared transfers. Their positive
Perron functions are gauge invariant, so the full and physical tops coincide
for each operator. This justifies a full-space trial for a top lower bound,
not a non-invariant trial for an excited physical level.

\begin{proposition}[A coarse native top floor uniform in coupling and lattice size]
\label{f16v51:top-floor}
For every $\gamma\ge1$ and admitted $L$, with $V=L^3$,
\begin{equation}
 \boxed{\qquad
 \Lambda_{\gamma,L}^{\rm W}\ge e^{-C_0V},\qquad C_0=30+3\log48.
 \qquad}
 \label{f16v51:uniform-top-floor}
\end{equation}
This explicit constant is deliberately conservative.
\end{proposition}
\begin{proof}
Let $D_\gamma=\{g:\theta(g)\le\gamma^{-1/2}\}$. Elementary radial bounds give
\[
 H(D_\gamma)\ge\frac1{12}\gamma^{-3/2},\qquad
 z_\gamma\le4\gamma^{-3/2}.
\]
For the first use $\sin\theta\ge2\theta/\pi$ on $[0,1]$; for the second use
$1-\cos\theta\ge2\theta^2/\pi^2$, $\sin\theta\le\theta$, and integrate the
resulting Gaussian to infinity. If $U,V\in D_\gamma$, their displacement
angle is at most $2\gamma^{-1/2}$, and
$e^{-\gamma(1-\cos\theta(UV^{-1}))}\ge e^{-2}$. The normalized one-link
indicator therefore has kinetic Rayleigh quotient at least $e^{-2}/48$.

Take the product of these unit indicators on all $3V$ links. On its support,
each original plaquette has angle at most $4\gamma^{-1/2}$, so its Wilson
action is at most $8/\gamma$. There are $3V$ spatial plaquettes, hence
$\gamma S\le24V$. Each magnetic half-factor is at least $e^{-12V}$.
The original complete Rayleigh quotient is consequently at least
\[
 e^{-24V}(e^{-2}/48)^{3V}=e^{-(30+3\log48)V}.
\]
The literal product trial was used only for the full top. Strict positivity
and gauge commutation identify that top with the physical one. No endpoint
carrier or pressure approximation enters this bound.
\end{proof}

\begin{theorem}[Actual-top comparison and a conservative joint-regulator bridge]
\label{f16v51:normalized}
Define $\widehat T^{\rm W}=T^{\rm W}/\Lambda^{\rm W}$ and
$\widehat T^{\rm H}=T^{\rm H}/\Lambda^{\rm H}$ with their own actual tops.
Then, for $\gamma\ge\gamma_*$,
\begin{align}
 |\Lambda^{\rm W}-\Lambda^{\rm H}|&\le\delta_\gamma,
 \label{f16v51:top-difference}\\
 \|\widehat T^{\rm W}-\widehat T^{\rm H}\|
     &\le 2\delta_\gamma/\Lambda^{\rm W}
      \le2\delta_\gamma e^{C_0L^3},
 \label{f16v51:normalized-step}\\
 \|(\widehat T^{\rm W})^N-(\widehat T^{\rm H})^N\|
     &\le \min\{2,2N\delta_\gamma e^{C_0L^3}\}.
 \label{f16v51:normalized-powers}
\end{align}
These statements hold on the full physical space and every center sector,
without identifying the two Perron vectors.

On a growing-lattice sequence, set $h=\gamma^{-1/3}/L$ as before. If
\begin{equation}
              L e^{C_0L^3}\gamma^{-2/3}\longrightarrow0,
 \label{f16v51:joint-regime}
\end{equation}
then, for every finite $t_1>0$,
\begin{equation}
 \boxed{\quad
 \sup_{0\le t\le t_1}
 \|(\widehat T_{\gamma,L}^{\rm W})^{\lfloor t/h\rfloor}
       -(\widehat T_{\gamma,L}^{\rm H})^{\lfloor t/h\rfloor}\|
             \longrightarrow0.
 \quad}
 \label{f16v51:joint-powers}
\end{equation}
A sufficient explicit condition is
$L^3\le c\log\gamma$ with fixed $0<c<2/(3C_0)$.
Corresponding normalized ordered eigenvalues differ by $o(h)$; where either
one is bounded away from zero, their logarithms also differ by $o(h)$.
This is an equivalence between two lattice problems, not an identification
of either growing-volume spectrum with the quartic model.
\end{theorem}
\begin{proof}
The top norm is 1-Lipschitz in operator norm. Also
\[
 \widehat T^{\rm W}-\widehat T^{\rm H}
 =\frac{T^{\rm W}-T^{\rm H}}{\Lambda^{\rm W}}
  +\frac{\Lambda^{\rm H}-\Lambda^{\rm W}}{\Lambda^{\rm W}}
                           \widehat T^{\rm H}.
\]
Both normalized operators are positive contractions, giving
\eqref{f16v51:normalized-step}. Telescoping their powers and the top floor
prove \eqref{f16v51:normalized-powers}. Here we do \emph{not} claim that the
relative deficit comparison survives the two different top normalizations.
For $N\le t_1/h$, its right side is bounded by
$C t_1 L e^{C_0L^3}\gamma^{-2/3}$, proving the uniform-time assertion.
The case $N=0$ is exact. Under the explicit logarithmic-volume condition this
bound is $O(L\gamma^{C_0c-2/3})\to0$.

Min--max applied after restriction to a common reducing sector bounds the
difference of every pair of correspondingly ordered normalized eigenvalues
by \eqref{f16v51:normalized-step}, which is $o(h)$ under
\eqref{f16v51:joint-regime}. If one eigenvalue is at least $\eta>0$,
the other is at least $\eta/2$ eventually, and the mean-value bound for the
logarithm proves the last assertion. No rate for either individual spectral
limit is needed to compare them, and no such limit is created by comparison.
\end{proof}

\subsection{Two limits of the new comparison, and the next estimate}

The operator topology above is essential. At any fixed finite $\gamma$ the
one-link Wilson and heat probabilities differ: the positive Bessel series
gives $I_{\ell+1}(\gamma)\ge(\gamma/2)^{\ell+1}/(\ell+1)!$, whose
logarithmic decay cannot agree with $-\ell(\ell+2)/(2\gamma)$ for all labels.
Choose a measurable one-link event with distinct probabilities. Applied to
independent product displacement laws, its empirical frequency and Chebyshev's
inequality imply
\begin{equation}
 d_{\rm TV}\bigl(k_\gamma^{\otimes M}dH^{\otimes M},
                   q_{(2\gamma)^{-1}}^{\otimes M}dH^{\otimes M}\bigr)
                   \longrightarrow1\quad(M\to\infty).
 \label{f16v51:TV-distinction}
\end{equation}
Thus the uniform $L^2$ operator comparison is not a volume-uniform statement
about transition rows or complete path densities. A delta input is not a
unit $L^2$ vector. No native interlink independence is asserted by this check.

Likewise a comparison of $I-T$ is not automatically one of vacuum-subtracted
energies. For a finite reference unrelated to native matrix elements, take
\[
 R_1=\begin{pmatrix}3/4&1/8\\1/8&1/4\end{pmatrix},\qquad
 R_2=\begin{pmatrix}1/4&1/8\\1/8&3/4\end{pmatrix},\qquad T_i=tR_i,\quad 0<t\le1.
\]
Both matrices are strictly positive positive-definite contractions, with common
top $t(1/2+\sqrt5/8)$. Their relative $I-T_i$ form error is $O(t)$, while
\[
 \left\|\frac{T_1}{\|T_1\|}-\frac{T_2}{\|T_2\|}\right\|
                  =\frac{1/2}{1/2+\sqrt5/8}>0
\]
for every $0<t\le1$. This demonstrates why the separate top cost cannot be removed
from \eqref{f16v51:normalized-step} just by renaming the deficit.

V51 resolves one full kinetic comparison uniformly in link count, without using
the growing-dimensional quotient expansion. The original spatial magnetic
interaction stays intact in every new native inequality. It does not resolve
the interacting hard-mode or soft-state comparison on a general growing lattice.
The logarithmic-volume bridge is deliberately conservative and is not a native
polynomial-volume scaling theorem. Improving it requires a relative comparison
at the actual small spectral top, or a more refined low-band estimate with the
original returning contribution retained. V50's uniform geometry may be used
in that analysis; it is not a premise quietly standing in for it here.

The parameter $t$ in \eqref{f16v51:joint-powers} is dimensionless critical-scale
time. Neither it nor the heat comparison time defines a physical lattice
spacing. No running-coupling law, spatial continuum limit, vacuum representation
limit, or interacting continuum mass follows from these operator comparisons.
The supplied memoranda's separate power/locality and same-top transfer
requirements therefore remain substantive.

\subsection{Verification and preservation}

The accompanying diagnostic checks the character second moment, normalized
radial coefficients, product-deficit inequality, noncommuting magnetic and
power comparisons and the separate top-normalization cost on exact finite
fixtures. The numerical Bessel exploration is labelled noncertified.
The all-label and all-volume assertions are proved above, not inferred from a
finite grid. The audit records the actual predecessor runs and their hashes.
Only the literal title version and this terminal insertion change V50; exact
removal recovers its bytes. The earlier analytic lineage is not independently
recertified, and the next scheduled checkpoint remains V55.

\begin{firewall}
V51 proves an all-representation relative Wilson--heat deficit bound which
survives arbitrary finite link products, the complete original magnetic
sandwich, physical restriction, and all unnormalized transfer durations with
no link-count loss. It also pays a separate actual-Perron normalization cost,
yielding a conservative growing-lattice equivalence to the declared
heat-kinetic comparison. It does not replace the native operator, assert
uniform transition-row total variation, identify the reference heat parameter
with physical time, or establish growing-volume quartic spectral completeness
or an interacting continuum mass gap.
\end{firewall}
% END F16 V51 APPEND

% BEGIN F16 V52 APPEND -- 2026-09-18
% Insert before the final document-ending command of supplied cumulative V51.
% No predecessor result, Wilson action, Haar measure, or physical clock is edited.

\clearpage
\section{V52: a kinetic spectral window and polynomial-volume Perron comparison}
\label{f16v52:context}

V51 controls Wilson--heat differences relative to the identity, uniformly in
link count. Its passage to the actual Perron normalization pays an exponential
volume factor against the entire error. We instead split the \emph{kinetic}
spectrum before applying the unchanged magnetic factors. Inside a total-Casimir
window, a logarithmic multiplier estimate gives a multiplicative comparison.
Outside it, both kinetic operators have an exponentially small norm. Only that
exponentially small remainder is divided by the small Perron value. The window
can grow linearly with the volume, making the remaining error polynomial rather
than exponential in that volume.

The result compares the same two lattice transfers as V51, each normalized by
its own actual top. It improves their common growing-lattice comparison domain;
it does not identify either spectrum with the quartic model on that domain.
The auxiliary spectral window is used to prove an inequality on the full space.
It is not a projection inserted into the native dynamics, and its high-label
remainder is retained explicitly throughout.

\subsection{Inputs, classical ratio bounds, and unchanged operators}

The V51 appendix and analytic audit were read in full. Targeted semantic and
exact-text searches of the current source, Grand Tensor V076, and f14/f15 V100
were recorded. Their earlier Casimir moments and electric-label tails are
related precedents, not the multiplicative comparison at the small Perron top
proved below. The only quantitative native scalar input needed from the current
lineage is V51's elementary top floor. Neither V48's spectral theorem nor V50's
uniform gauge tube is used.

Bounds for consecutive modified-Bessel ratios are classical. Hornik--Gr\"un,
\href{https://doi.org/10.1016/j.jmaa.2013.05.070}{\emph{Amos-type bounds for
modified Bessel function ratios}}, and Ruiz-Antol\'in--Segura,
\href{https://arxiv.org/abs/1606.02008}{\emph{A new type of sharp bounds for ratios
of modified Bessel functions}}, describe this family and its Riccati method.
The latter paper's Theorems 2 and 4 include the bounds used here, with its order
shifted by one. We prove the needed integer-order case directly, rather than
apply a fixed-order asymptotic series to increasing labels. The recurrence is
\href{https://dlmf.nist.gov/10.29}{DLMF 10.29}. No new general Bessel inequality
is claimed. The next companion checkpoint remains V55.

Retain probability Haar on the unit-round quaternion group $G=\SU(2)$ and
V51's exact kernels
\[
 C_\gamma:\quad c_{\gamma,\ell}=I_{\ell+1}(\gamma)/I_1(\gamma),
 \qquad Q_\gamma=e^{-\mathcal L_G/(2\gamma)}:\quad
 q_{\gamma,\ell}=e^{-s_\ell/(2\gamma)},\qquad
 s_\ell=\ell(\ell+2).
\]
The highest-weight label is the integer $\ell\ge0$, not the spin $\ell/2$.
For any finite link count $M$ put
\[
 K^{\rm W}=C_\gamma^{\otimes M},\qquad
 K^{\rm H}=Q_\gamma^{\otimes M},\qquad
 X_\gamma=(2\gamma)^{-1}\sum_{e=1}^M\mathcal L_e.
\]
Thus $K^{\rm H}=e^{-X_\gamma}$. The original fine torus has $M=3V$, $V=L^3$,
$L=2m\ge8$. Its two transfers remain
\[
 T^{\rm W}=b_\gamma K^{\rm W}b_\gamma=T_{\gamma,L},\qquad
 T^{\rm H}=b_\gamma K^{\rm H}b_\gamma,\qquad
 b_\gamma=e^{-\gamma S/2}.
\]
Only the first is the original Wilson transfer. All spatial plaquettes, group
measures, gauge restrictions and center sectors are unchanged in both. The heat
operator changes the temporal kinetic factor only.

\subsection{A logarithmic estimate with a separate high-label bound}

\begin{proposition}[Classical ratio barriers and a uniform logarithmic ledger]
\label{f16v52:ratios}
For every integer $\nu\ge1$, every $x>0$, and $a=\nu+1/2$,
\begin{equation}
 \frac{x}{a+\sqrt{x^2+(a+1)^2}}
 <\frac{I_{\nu+1}(x)}{I_\nu(x)}
 <\frac{x}{a+\sqrt{x^2+a^2}}.
 \label{f16v52:Amos}
\end{equation}
For $\gamma>0$, $\ell\ge0$, write
\[
 \Phi_{\gamma,\ell}=-\log c_{\gamma,\ell},\qquad
 x_\ell=s_\ell/(2\gamma).
\]
Then
\begin{equation}
 -\frac{s_\ell(2s_\ell+3)}{48\gamma^3}
 \le \Phi_{\gamma,\ell}-x_\ell
 \le \frac{\ell(\ell+3)}{2\gamma^2}.
 \label{f16v52:log-bound}
\end{equation}
For $\gamma\ge2$, a second bound valid at every label is
\begin{equation}
 \boxed{\qquad
 \Phi_{\gamma,\ell}\ge\frac1{16}\min\{x_\ell,\gamma\}.
 \qquad}
 \label{f16v52:large-label}
\end{equation}
All inequalities are exact at $\ell=0$ when interpreted non-strictly. The
bounds do not use a representation cutoff or label-dependent asymptotic error.
\end{proposition}
\begin{proof}
The positive Bessel series and its recurrence give, for
$r(x)=I_{\nu+1}(x)/I_\nu(x)$,
\[
 r'=1-(2\nu+1)r/x-r^2=:\mathcal R(x,r).
\]
The upper candidate $u=x/(a+\sqrt{x^2+a^2})$ is the positive root of
$\mathcal R(x,u)=0$ and has $u'>0$. The series gives
$r=x/(2a+1)+O(x^3)<u$ near zero. At a first crossing from below,
$(r-u)'=-u'<0$, which is impossible. This proves the upper bound.

For the lower candidate $l=x/(a+b)$, $b=\sqrt{x^2+(a+1)^2}$, direct
calculation gives
\[
 \mathcal R(x,l)-l'
 =\frac{(a+1)[1-(a+1)/b]}{(a+b)^2}>0.
\]
The small-$x$ series give
\[
 r-l=\frac{x^3}{(2a+1)^2(2a+2)(2a+3)}+O(x^5)>0
\]
near zero. A first downward crossing would have positive derivative by the
last identity, again impossible. There are no poles on $x>0$ because the
Bessel series is positive. This proves \eqref{f16v52:Amos} on the full interval.

Taking logarithms in that bracket, with $x=\gamma$, gives
\[
 \operatorname{arsinh}(a/\gamma)
 \le -\log r(\gamma)
 \le \operatorname{arsinh}(a/\gamma)+\frac{\nu+1}{\gamma^2}.
\]
The difference of the two square roots in the denominators is
\[
 \frac{2a+1}{\sqrt{\gamma^2+(a+1)^2}+\sqrt{\gamma^2+a^2}}.
\]
Divide this by the smaller denominator and use $\log(1+t)\le t$. For every
$t\ge0$,
\[
 0\le t-\operatorname{arsinh}t\le t^3/6.
\]
This follows by integrating
$0\le1-(1+t^2)^{-1/2}\le t^2/2$. The exact product
$c_{\gamma,\ell}=\prod_{\nu=1}^{\ell}r_\nu(\gamma)$ and the finite sums
\[
 \sum_{\nu=1}^{\ell}(\nu+1/2)=s_\ell/2,\qquad
 \sum_{\nu=1}^{\ell}(\nu+1/2)^3=s_\ell(2s_\ell+3)/8
\]
give \eqref{f16v52:log-bound}, including its two different error signs.

For the tail estimate, if $1\le\ell\le\gamma$, then
$(\nu+1/2)/\gamma\le5/4$ for every term. On $[0,5/4]$,
$\operatorname{arsinh}t\ge t/2$, so $\Phi_{\gamma,\ell}\ge x_\ell/2$.
If $\ell>\gamma$, keep only the first $m=\lfloor\gamma\rfloor$ terms.
They are nonnegative, $m\ge\gamma/2$, and their sum is at least
$m(m+2)/(4\gamma)\ge\gamma/16$. Both cases imply
\eqref{f16v52:large-label}. Positivity of every summand follows already from
the upper ratio barrier, so no cancellation is used in the tail estimate.
\end{proof}

The logarithmic error is useful when the total heat energy is moderate. It is
not a small relative error uniformly at arbitrarily high representations. At
fixed $\gamma$, $c_{\gamma,\ell}/q_{\gamma,\ell}\to\infty$ as
$\ell\to\infty$, by the factorial lower bound in the positive Bessel series.
Consequently a global bound $K^{\rm W}\preceq C K^{\rm H}$ with finite $C$
is false even on one link. The additive tail below is necessary for this route.

\subsection{A total-Casimir window, with its full complement retained}

For $R\ge1$ define $P_R=1_{[0,R]}(X_\gamma)$ and put
\begin{equation}
 \eta_\gamma(R)=\frac{2R}{\gamma}+
                  \frac{R^2}{6\gamma}+
                  \frac{R}{8\gamma^2},\qquad
 t_R=e^{-R/16}.
 \label{f16v52:parameters}
\end{equation}
This window bounds the total electric spectrum. It is neither a separate
condition on each link nor a small-field event in configuration space.

\begin{theorem}[Multiplicative kinetic window and exponentially small remainder]
\label{f16v52:kinetic}
For every finite $M$, every $\gamma\ge2$ and $1\le R\le\gamma$,
\begin{align}
 e^{-\eta_\gamma(R)}K^{\rm H}P_R
 &\preceq K^{\rm W}P_R
 \preceq e^{\eta_\gamma(R)}K^{\rm H}P_R,
 \label{f16v52:low}\\
 \|K^{\rm W}(I-P_R)\|&\le t_R,
 \qquad \|K^{\rm H}(I-P_R)\|\le e^{-R}.
 \label{f16v52:high}
\end{align}
In particular, on the \emph{complete} link Hilbert space,
\begin{equation}
 \boxed{\quad
 e^{-\eta_\gamma(R)}K^{\rm H}-t_RI
 \preceq K^{\rm W}
 \preceq e^{\eta_\gamma(R)}K^{\rm H}+t_RI.
 \quad}
 \label{f16v52:kinetic-full}
\end{equation}
The constants have no factor counting links, representation dimensions, or
multiplicities. $P_R$ commutes with the original gauge and center actions.
It need not commute with the magnetic multiplier.
\end{theorem}
\begin{proof}
For a product label, write $X=\sum_e x_{\ell_e}$ and
$\Phi=\sum_e\Phi_{\gamma,\ell_e}$. Since
$\ell(\ell+3)/2\le s_\ell$, the logarithmic ledger gives
\[
 -\frac{X^2}{6\gamma}-\frac{X}{8\gamma^2}
 \le\Phi-X\le\frac{2X}{\gamma}.
\]
For the lower bound use $\sum_e x_e^2\le(\sum_e x_e)^2$; for the upper
bound sum $s_e/\gamma^2$. Thus $|\Phi-X|\le\eta_\gamma(R)$ when $X\le R$.
Exponentiation proves \eqref{f16v52:low}.

The tail estimate of the proposition gives
\[
 \Phi\ge\frac1{16}\sum_e\min(x_e,\gamma)
       \ge\frac1{16}\min(X,\gamma).
\]
The second inequality holds because either one term exceeds $\gamma$, or
all terms are below it and their sum is unchanged. On $X>R$, $R\le\gamma$,
the Wilson multiplier is therefore at most $e^{-R/16}$. The heat multiplier
is $e^{-X}\le e^{-R}$. This includes both one very large label and many
moderate labels; a bound on their number is unnecessary.

On the low subspace the desired full inequality has just been proved. On its
complement both operators are nonnegative, the Wilson norm is at most $t_R$,
and $e^{-\eta}K^{\rm H}$ has norm at most $e^{-R}\le t_R$. This proves
\eqref{f16v52:kinetic-full} on both reducing subspaces. Peter--Weyl orthogonality
extends it to arbitrary superpositions. Bi-invariance of the Laplacians proves
the symmetry assertion. No restriction on an interacting state's electric
labels is assumed.
\end{proof}

\subsection{The complete magnetic form and actual Perron normalization}

\begin{theorem}[A multiplicative transfer estimate with a paid electric tail]
\label{f16v52:magnetic}
For the same parameters and the unchanged magnetic multiplier,
\begin{equation}
 \boxed{\quad
 e^{-\eta}T^{\rm H}-t_R b_\gamma^2
 \preceq T^{\rm W}
 \preceq e^{\eta}T^{\rm H}+t_R b_\gamma^2,
 \qquad \eta=\eta_\gamma(R).
 \quad}
 \label{f16v52:magnetic-full}
\end{equation}
It follows on the full physical space that
\begin{equation}
 \|T^{\rm W}-T^{\rm H}\|
       \le(e^{\eta}-1)\Lambda^{\rm H}+t_R,
 \label{f16v52:relative-norm}
\end{equation}
where $\Lambda^{\rm H}$ is the actual heat-kinetic Perron value. Every same
physical restriction, center-sector restriction, or orthogonal compression
preserves \eqref{f16v52:magnetic-full}, with the corresponding identity if
$b_\gamma^2$ is bounded by one. No spectral projection is inserted into either
transfer or its powers.
\end{theorem}
\begin{proof}
Take the congruence by $b_\gamma$ in \eqref{f16v52:kinetic-full}. This step
requires no commutation of $b_\gamma$ with the spectral window. It retains the
full original spatial action. In both signs, $T^{\rm W}-T^{\rm H}$ is bounded
by $(e^\eta-1)T^{\rm H}+t_RI$, since $1-e^{-\eta}\le e^\eta-1$ and
$0\preceq b_\gamma^2\preceq I$. The quadratic-form characterization of the
norm of a self-adjoint operator gives \eqref{f16v52:relative-norm}. All
restrictions mentioned preserve quadratic-form inequalities. This is not an
entrywise-kernel inference or a row normalization.
\end{proof}

\begin{proposition}[Only the discarded tail pays the reciprocal Perron value]
\label{f16v52:normalize}
Suppose the parameters above satisfy $\eta\le1/2$ and
$r=t_R/\Lambda^{\rm W}\le1/2$. Normalize each transfer by its own top:
$\widehat T^{\rm W}=T^{\rm W}/\Lambda^{\rm W}$ and
$\widehat T^{\rm H}=T^{\rm H}/\Lambda^{\rm H}$. Then
\begin{align}
 \left|\log\frac{\Lambda^{\rm W}}{\Lambda^{\rm H}}\right|
          &\le\eta+2r,
 \label{f16v52:top-log}\\
 \|\widehat T^{\rm W}-\widehat T^{\rm H}\|
          &\le12(\eta+r),
 \label{f16v52:normalized-step}\\
 \| (\widehat T^{\rm W})^N-(\widehat T^{\rm H})^N\|
          &\le\min\{2,12N(\eta+r)\},\qquad N\ge1.
 \label{f16v52:normalized-powers}
\end{align}
The same actual tops are used in every center sector. No comparison of the two
Perron vectors, or of their vacuum-subtracted form domains, is assumed.
\end{proposition}
\begin{proof}
The top variational principle in \eqref{f16v52:magnetic-full}, followed by
$b_\gamma^2\preceq I$, gives
\[
 e^{-\eta}(1-r)\le\Lambda^{\rm H}/\Lambda^{\rm W}
                    \le e^{\eta}(1+r).
\]
Taking logarithms and using $-\log(1-r)\le2r$ proves the first assertion.
The exact identity
\[
 \widehat T^{\rm W}-\widehat T^{\rm H}
 =\frac{T^{\rm W}-T^{\rm H}}{\Lambda^{\rm W}}
  +\frac{\Lambda^{\rm H}-\Lambda^{\rm W}}{\Lambda^{\rm W}}
                                     \widehat T^{\rm H}
\]
and the Lipschitz bound on the top yield a norm at most
$2\|T^{\rm W}-T^{\rm H}\|/\Lambda^{\rm W}$. By the preceding theorem this is
at most
\[
 2\{(e^\eta-1)e^\eta(1+r)+r\}\le12(\eta+r),
\]
using $e^\eta\le2$, $e^\eta-1\le2\eta$, and $r\le1/2$.
Both normalized operators are positive contractions. Their exact noncommuting
telescoping sum proves the final assertion. The physical tops are the full tops
by positivity, uniqueness and gauge commutation, as in V51. This does not
identify the two corresponding eigenfunctions.
\end{proof}

\subsection{A polynomial-volume actual-top comparison}

V51's independently proved native floor is
\[
 \Lambda^{\rm W}_{\gamma,L}\ge e^{-C_0V},\qquad
 C_0=30+3\log48,\qquad V=L^3,\quad\gamma\ge1.
\]
For a fixed $p\ge1$, choose the \emph{proof parameter}
\begin{equation}
 R_{\gamma,V,p}=16[C_0V+(p+1)\log\gamma].
 \label{f16v52:cutoff-choice}
\end{equation}
The native operators themselves do not depend on this choice.

\begin{theorem}[Actual-top comparison with polynomial, not exponential, volume cost]
\label{f16v52:polynomial}
For every $p\ge1$, suppose $\gamma\ge2$, $R=R_{\gamma,V,p}\le\gamma$, and
$\eta_\gamma(R)\le1/2$. Then
\begin{equation}
 \frac{t_R}{\Lambda^{\rm W}}\le\gamma^{-(p+1)},
 \label{f16v52:tail-paid}
\end{equation}
and, with $\epsilon_{\gamma,V,p}=\eta_\gamma(R)+\gamma^{-(p+1)}$,
\begin{equation}
 \boxed{\quad
 \left|\log\frac{\Lambda^{\rm W}}{\Lambda^{\rm H}}\right|
       \le2\epsilon_{\gamma,V,p},\qquad
 \|\widehat T^{\rm W}-\widehat T^{\rm H}\|
       \le12\epsilon_{\gamma,V,p}.
 \quad}
 \label{f16v52:polynomial-top}
\end{equation}
In particular
\begin{equation}
 \epsilon_{\gamma,V,p}
 \le C_p\left[\frac{(V+\log\gamma)^2}{\gamma}
                        +\gamma^{-(p+1)}\right],
 \label{f16v52:polynomial-rate}
\end{equation}
with $C_p$ independent of $V$ and $\gamma$. The explicit entry conditions hold
eventually on every sequence with $(V+\log\gamma)^2/\gamma\to0$.
The conclusions hold in the complete physical space and every center sector.
\end{theorem}
\begin{proof}
The cutoff was selected so that
$t_R=e^{-C_0V}\gamma^{-(p+1)}$. The actual native floor therefore proves
\eqref{f16v52:tail-paid}, not an assumption on the heat top or a Gaussian
pressure. Its right side is at most $1/4$ under the stated conditions, so the
normalization proposition applies. For $R\ge1$, $\gamma\ge2$,
$\eta_\gamma(R)\le3(R+R^2)/\gamma$. Equation
\eqref{f16v52:cutoff-choice} now gives \eqref{f16v52:polynomial-rate}. The same
inequality and $R/\gamma\to0$ prove the eventual entry assertion. No unknown
volume-dependent threshold or inherited $o_L(1)$ remainder is used.
\end{proof}

The reciprocal Perron value has not been declared harmless. It multiplies
$e^{-R/16}$, and the proved scalar floor is used to choose $R$ large enough to
pay it. The multiplicative low-window error depends quadratically on that
window, giving the polynomial volume cost. There is no factor equal to the
rank of $P_R$, which can be extremely large. High representations have not
been removed from the final inequality.

\begin{theorem}[Perron-normalized evolution on polynomially growing lattices]
\label{f16v52:critical}
Let $L=L(\gamma)\ge8$ be even and let $h=\gamma^{-1/3}/L$ be the existing
critical parameter. If
\begin{equation}
 \boxed{\qquad
 L\,[L^3+\log\gamma]^2\gamma^{-2/3}\longrightarrow0,
 \qquad}
 \label{f16v52:joint-domain}
\end{equation}
then, for every finite $t_1>0$,
\begin{equation}
 \sup_{0\le t\le t_1}
 \left\| (\widehat T_{\gamma,L}^{\rm W})^{\lfloor t/h\rfloor}
              -(\widehat T_{\gamma,L}^{\rm H})^{\lfloor t/h\rfloor}\right\|
 \longrightarrow0.
 \label{f16v52:evolution}
\end{equation}
A sufficient explicit family is
\begin{equation}
 \boxed{\qquad L=O(\gamma^b),\qquad 0\le b<2/21.\qquad}
 \label{f16v52:power-domain}
\end{equation}
Correspondingly ordered normalized eigenvalues in every center sector differ
by $o(h)$, uniformly in their index. Their logarithms differ by $o(h)$ on any
part bounded away from zero. This compares two growing-lattice problems; it
does not establish an individual quartic or continuum limit for either one.
\end{theorem}
\begin{proof}
Use $p=1$ in the preceding theorem. Condition \eqref{f16v52:joint-domain}
implies its entry conditions, and
\[
 \epsilon_{\gamma,V,1}/h
 \le C L[V+\log\gamma]^2\gamma^{-2/3}+L\gamma^{-5/3}\longrightarrow0.
\]
The second term follows from the first condition, since $L\ge8$ and $V=L^3$.
Apply \eqref{f16v52:normalized-powers} at $N\le t_1/h$; $N=0$ is exact.
For \eqref{f16v52:power-domain}, bound the left side of
\eqref{f16v52:joint-domain} by
$C[L^7+L\log^2\gamma]\gamma^{-2/3}$.
Both terms tend to zero when $7b<2/3$. No claim of optimality is made for this
exponent. For example $L=O(\gamma^{1/12})$ gives the upper bound
$O(\gamma^{-1/12})$ for the finite-critical-time norm error.

Min--max on the same reducing sector bounds every pair of correspondingly
ordered eigenvalues by the normalized operator-norm difference. This estimate
is uniform in the index even though no individual growing-index asymptotic is
supplied. On $[a,1]$ the logarithm is Lipschitz with constant $1/a$; if one
compared eigenvalue is at least $a>0$, the other is at least $a/2$ eventually.
This proves the logarithmic assertion without identifying the two vacua.
\end{proof}

\begin{proposition}[A full rescaled resolvent comparison, with the vacuum retained]
\label{f16v52:resolvent}
Define the nonnegative bounded operators
$A^{\rm W}_h=h^{-1}(I-\widehat T^{\rm W})$ and
$A^{\rm H}_h=h^{-1}(I-\widehat T^{\rm H})$. For $z>0$ and the entry conditions
of Theorem~\ref{f16v52:polynomial},
\begin{equation}
 \|(A_h^{\rm W}+zI)^{-1}-(A_h^{\rm H}+zI)^{-1}\|
       \le\frac{12\epsilon_{\gamma,V,p}}{hz^2}.
 \label{f16v52:resolvents}
\end{equation}
It tends to zero on \eqref{f16v52:joint-domain}. These are resolvents of the
stated discrete deficits, not an identification with the logarithmic physical
Hamiltonians or their vacuum-orthogonal inverses.
\end{proposition}
\begin{proof}
Both resolvents have norm at most $z^{-1}$. Apply their exact resolvent identity
and $\|A_h^{\rm W}-A_h^{\rm H}\|\le12\epsilon_{\gamma,V,p}/h$. The zero modes
remain on each side, and their projectors are not set equal. The estimate does
not divide by an unproved excitation gap.
\end{proof}

\subsection{Scope, normalization boundary, and the next estimate}

V52 improves V51's sufficient logarithmic-volume bridge to a polynomial-volume
bridge for the \emph{same two actual-top-normalized transfers}. The improvement
comes from a relative estimate on a total kinetic window and an explicitly paid
tail, not from a new spatial action, a heat-kernel identification, or an assumed
bound on the actual near-top eigenvectors' representation labels. In particular
$P_R$ is never assumed invariant under the magnetic transfer.

This does not upgrade V48's $o_L(h)$ into a remainder uniform along
\eqref{f16v52:power-domain}. The heat-kinetic reference still has the complete
nonlinear spatial Wilson action. Its collective soft dynamics, native
localization on growing lattices, and any required same-top return must be
analyzed separately. A common operator comparison transfers a limit only after
one side has that limit with the required quantifiers. There is no claim here
that the normalized Perron vectors are close: obtaining that from operator norm
would additionally require control of the relevant spectral separation.

Nor does the all-space estimate imply volume-uniform total variation for the
transition rows; V51's product-row distinction remains valid. The total-Casimir
window is a proof device whose complement is bounded in norm, not an event
asserted typical under a stationary or conditioned native probability.

More generally, for an independently supplied physical time step $a_t>0$, the
same proof compares the two normalized evolutions over bounded physical times
whenever $\epsilon_{\gamma,V,p}/a_t\to0$. This is a conditional conversion,
not a selection of $a_t$. The present use of $h$ labels dimensionless critical
powers. No running coupling, physical scale, local vacuum representation, or
interacting continuum mass is fixed by it. The supplied memoranda's request to
pay the error relative to the actual Perron value is met for this kinetic
comparison; their separate locality and physical-transfer requirements remain.

The next task can therefore be posed for the heat-kinetic reference on the
proved polynomial domain, with its original spatial action and spectral return
intact, or can improve the present logarithmic multiplier estimate and window
cost further. Another absolute one-step kinetic estimate would not address the
remaining interacting spectral problem.

The new diagnostic completes 569 exact assertions in 46 recorded families.
Most are finite saturated-sum and Bessel-series checks; its scalar Bessel
fixtures include rigorous series-tail bounds. It also tests noncommuting
magnetic congruences and the normalization algebra. These are not native
spectral enclosures or an all-parameter proof certification. All 48 supplied
V4--V51 diagnostics were executed unchanged with successful individual records.
Removing this exact insertion and reversing the literal title-version change
recovers V51 byte for byte. The next companion checkpoint remains V55.

\begin{firewall}
V52 proves a multiplicative kinetic spectral-window comparison with an explicit
high-label remainder, and uses a proved native top floor to obtain
actual-Perron-normalized Wilson--heat equivalence on a polynomial growing-lattice
domain. All original magnetic interactions and states remain in the compared
operators. It does not prove a growing-volume quartic asymptotic, align vacua
without a separation estimate, set a physical clock, or construct an interacting
continuum Yang--Mills theory with a mass gap. The earlier analytic lineage is
not independently recertified.
\end{firewall}
% END F16 V52 APPEND

% BEGIN F16 V53 APPEND -- 2026-09-18
% Insert immediately before the final document-ending command of cumulative V52.
% The original Wilson operator and every predecessor source body are unchanged.

\clearpage
\section{V53: a subordinated kinetic reference and a cutoff-free logarithmic comparison}
\label{f16v53:context}

V52 compares the original Wilson transfer with ordinary heat kinetics by a
multiplicative estimate on a total electric window and an additive tail. The
window is necessary for that particular reference: its large-representation
Gaussian decay is faster than Wilson decay. We instead construct a different,
explicitly positive reference whose dispersion retains the inverse-hyperbolic
shape before its small-momentum expansion. A midpoint estimate then bounds the
\emph{complete} Wilson logarithmic multiplier between that dispersion and a
constant multiple of it, at every representation. The constant is independent
of the number of links. Operator convexity transports this bound through the
unchanged magnetic factors and gives a power envelope for the full transfer.

The new reference is a normalized positive mixture of round-group heat kernels,
not the ordinary heat operator of V51--V52 and not a redefinition of Wilson
transfer. The resulting actual-top comparison has a linear, rather than
quadratic, volume cost. It also compares the quadratic forms of the unbounded
logarithmic transfers on their common form domain. Neither conclusion proves a
uniform interacting spectral limit for the reference or supplies physical time.

\subsection{Dependencies, classical inputs, and unchanged native normalization}

We retain V52's exact Bessel conventions and V51's native top floor. Neither
the fixed-volume spectral branch nor V50's gauge tube is a premise here.
V41's positive integral representation concerns the different harmonic function
$\operatorname{arcosh}(1+s/2)$. The present construction instead bounds the
\emph{one-link Wilson logarithm} and then the complete magnetic transfer.

The consecutive Bessel barriers are the classical ones proved directly in
Proposition~\ref{f16v52:ratios}; their primary context is Ruiz-Antol\'in--Segura,
\href{https://arxiv.org/html/1606.02008}{\emph{A new type of sharp bounds for
ratios of modified Bessel functions}}. Bernstein subordination is described in
Hiroshima--Ichinose--L\H{o}rinczi,
\href{https://arxiv.org/abs/0906.0103}{\emph{Path Integral Representation for
Schr\"odinger Operators with Bernstein Functions of the Laplacian}}, Section 2.
The contractive Jensen inequality is Hansen--Pedersen,
\href{https://arxiv.org/html/math/0204049}{\emph{Jensen's Operator Inequality}},
Corollary 2.3. We give the needed integral construction and power-compression
proof below. No external Yang--Mills or Navier--Stokes theorem is imported.

Keep probability Haar on the unit-round quaternion group $G=\SU(2)$, the
round Laplacian $\mathcal L_G\ge0$, and its highest-weight labels
$s_\ell=\ell(\ell+2)$, $\ell\ge0$. The original normalized convolution has
\[
 c_{\gamma,\ell}=I_{\ell+1}(\gamma)/I_1(\gamma),\qquad
 \Phi_{\gamma,\ell}=-\log c_{\gamma,\ell}.
\]
The original fine torus has $M=3V$ links, $V=L^3$, $L\ge8$ even, and
\[
 T^{\rm W}=b_\gamma C_\gamma^{\otimes M}b_\gamma,\qquad
 b_\gamma=e^{-\gamma S/2},
\]
with the complete original spatial Wilson action $S$. All bounds below keep the
same gauge-invariant Hilbert space and center sectors.

\subsection{A positive dispersion retaining the high-label shape}

For $\gamma>0$ and $s\ge0$, define
\begin{equation}
 \begin{split}
 F_\gamma(s)&=\int_1^{\sqrt{1+s}}\operatorname{arsinh}(x/\gamma)\,dx,\\
 &=\sqrt{1+s}\operatorname{arsinh}\!\left(\frac{\sqrt{1+s}}\gamma\right)
       -\sqrt{\gamma^2+1+s}
       -\operatorname{arsinh}(1/\gamma)+\sqrt{\gamma^2+1}.
 \end{split}
 \label{f16v53:dispersion}
\end{equation}
In particular $F_\gamma(0)=0$. Write
\[
 \mathrm E_1(x)=\int_x^\infty e^{-v}\,\frac{dv}{v},\qquad x>0.
\]
This is an explicit exponential integral, not an energy or a spectral projector.

\begin{proposition}[A normalized heat mixture with an explicit L\'evy density]
\label{f16v53:subordinator}
The function $F_\gamma$ is a Bernstein function and has the representation
\begin{equation}
 \boxed{\quad
 F_\gamma(s)=\int_0^\infty(1-e^{-st})n_\gamma(t)\,dt,\qquad
 n_\gamma(t)=\frac{e^{-t}}{4\sqrt\pi\,t^{3/2}}\mathrm E_1(\gamma^2t)>0.
 \quad}
 \label{f16v53:Levy}
\end{equation}
There is a probability measure $\mu_\gamma$ on $(0,\infty)$ such that
\begin{equation}
 \int e^{-st}\,\mu_\gamma(dt)=e^{-F_\gamma(s)},\qquad
 R_\gamma:=e^{-F_\gamma(\mathcal L_G)}
          =\int e^{-t\mathcal L_G}\,\mu_\gamma(dt).
 \label{f16v53:heat-mixture}
\end{equation}
The integral is a strong operator integral. $R_\gamma$ is a positive
self-adjoint probability convolution with a smooth strictly positive kernel.
It is compact and injective. If $T$ has law $\mu_\gamma$, then
\begin{equation}
 \mathbb ET=\tfrac12\operatorname{arsinh}(1/\gamma),\qquad
 \operatorname{Var}T=\tfrac14\left[
 \operatorname{arsinh}(1/\gamma)-(\gamma^2+1)^{-1/2}\right]>0.
 \label{f16v53:time-moments}
\end{equation}
These are moments of a reference mixture variable, not of a physical clock.
\end{proposition}
\begin{proof}
Differentiation of the defining integral gives
\[
 F_\gamma'(s)=\frac{\operatorname{arsinh}(\sqrt{1+s}/\gamma)}{2\sqrt{1+s}}
      =\frac12\int_0^1[\gamma^2+(1+s)u^2]^{-1/2}\,du.
\]
Every derivative of the last integrand has alternating sign, so $F_\gamma'$ is
completely monotone. More explicitly, the gamma integral and Tonelli give
\[
 F_\gamma'(s)=\frac1{2\sqrt\pi}\int_0^\infty t^{-1/2}e^{-(1+s)t}
           \left(\int_0^1u^{-1}e^{-\gamma^2t/u^2}\,du\right)dt.
\]
The inner integral is $\frac12\mathrm E_1(\gamma^2t)$. Integrate in $s$ and use
$F_\gamma(0)=0$ to obtain \eqref{f16v53:Levy}.
Near zero $\mathrm E_1(\gamma^2t)=O_\gamma(1+|\log t|)$; at infinity it is at
most $e^{-\gamma^2t}/(\gamma^2t)$. Thus
$\int(1\wedge t)n_\gamma(t)dt<\infty$, and all positive time moments used
below are integrable.

For completeness, truncate $n_\gamma$ to $[j^{-1},j]$ and take a compound
Poisson random variable with this finite intensity measure. Its Laplace
transform is the exponential of the truncated integral in
\eqref{f16v53:Levy}. Its first moment is bounded by
$\int t n_\gamma(t)dt=F_\gamma'(0)$, so these probabilities are tight.
Every weak limit has transform $e^{-F_\gamma(s)}$, and uniqueness of Laplace
transforms makes the limit unique. Its mass is one. As $s\to\infty$,
$F_\gamma(s)\to\infty$, so dominated convergence in the Laplace transform
shows that there is no atom at zero. This constructs \eqref{f16v53:heat-mixture}
without assuming a random-time representation of the original Wilson kernel.
The same construction for a multiple of the intensity gives the usual
subordination family; only its unit member is needed here.

Every heat kernel at positive time is strictly positive and normalized.
The mixture preserves probability, positivity and bi-invariance. For fixed
$\gamma$, $F_\gamma(s)$ grows like $\sqrt{s}\log(2\sqrt{s}/\gamma)-\sqrt{s}$.
Thus $e^{-F_\gamma(s_\ell)}$ decays faster than every inverse power of $\ell$.
The Peter--Weyl series, with any fixed number of derivatives and its polynomial
multiplicities, converges absolutely. This proves smoothness and compactness;
all multipliers are positive, giving injectivity. Positivity can alternatively
be checked directly in the mixture; a compact positive-time interval with
positive $\mu_\gamma$ mass gives a strictly positive lower kernel bound.
Differentiating the Laplace exponent at zero gives
$\mathbb ET=F_\gamma'(0)$ and $\operatorname{Var}T=-F_\gamma''(0)$, which
are \eqref{f16v53:time-moments}. Their large-$\gamma$ forms are
$1/(2\gamma)+O(\gamma^{-3})$ and $1/(12\gamma^3)+O(\gamma^{-5})$.
\end{proof}

\subsection{A global logarithmic bracket, without an electric cutoff}

\begin{theorem}[All-label Wilson dispersion between two scalar multiples]
\label{f16v53:global-scalar}
For every $\gamma\ge2$ and every $\ell\ge0$, put
$\delta_\gamma=2/\gamma$ and $p_\gamma=1+\delta_\gamma\in(1,2]$. Then
\begin{equation}
 \boxed{\qquad
 F_\gamma(s_\ell)\le\Phi_{\gamma,\ell}
                  \le p_\gamma F_\gamma(s_\ell).
 \qquad}
 \label{f16v53:all-label}
\end{equation}
Consequently, for every finite $M$,
\begin{equation}
 K^{\rm B}_{\gamma,M}:=R_\gamma^{\otimes M},\qquad
 \boxed{\quad
 (K^{\rm B}_{\gamma,M})^{p_\gamma}
       \preceq C_\gamma^{\otimes M}\preceq K^{\rm B}_{\gamma,M}.
 \quad}
 \label{f16v53:all-link}
\end{equation}
There is no additive high-label tail, label restriction or link-count factor.
The superscript $\rm B$ distinguishes this Bernstein reference from ordinary heat.
\end{theorem}
\begin{proof}
Let $f(x)=\operatorname{arsinh}(x/\gamma)$ for $x>0$. For
$a=\nu+1/2$, $\nu\ge1$, V52's exact ratio barriers imply
\[
 f(a)\le-\log\frac{I_{\nu+1}(\gamma)}{I_\nu(\gamma)}
 \le f(a)+\log\frac{a+\sqrt{\gamma^2+(a+1)^2}}
                         {a+\sqrt{\gamma^2+a^2}}.
\]
Writing $b=\sqrt{\gamma^2+a^2}$ and using $\log(1+x)\le x$ bounds the last
logarithm by
\[
 \frac{2a+1}{2b(a+b)}\le\frac{4a}{3b(a+b)}
       \le\frac4{3\gamma}f(a).
\]
Here $a\ge3/2$ and $\operatorname{arsinh}t\ge t/\sqrt{1+t^2}$; the latter
follows by integrating the decreasing derivative over $[0,t]$.
Thus the discrete sum $S_\ell=\sum_{\nu=1}^\ell f(\nu+1/2)$ satisfies
\[
 S_\ell\le\Phi_{\gamma,\ell}
                 \le(1+4/(3\gamma))S_\ell.
\]

The midpoint error has the correct sign and a relative bound. On each unit
interval $[\nu,\nu+1]$ it equals
\[
 f(\nu+1/2)-\int_\nu^{\nu+1}f(x)dx
 =\int_\nu^{\nu+1}K_\nu(x)[-f''(x)]dx,
\]
where $K_\nu(x)=\frac12(x-\nu)^2$ on the first half and
$K_\nu(x)=\frac12(\nu+1-x)^2$ on the second. In particular
$0\le K_\nu\le1/8$. Moreover
\[
 -f''(x)=\frac{x}{(\gamma^2+x^2)^{3/2}}
           \le\gamma^{-2}f(x).
\]
Sum the Peano identities to obtain
\[
 F_\gamma(s_\ell)\le S_\ell
          \le(1+1/(8\gamma^2))F_\gamma(s_\ell),
\]
since $\sqrt{1+s_\ell}=\ell+1$. Finally
\[
 (1+4/(3\gamma))(1+1/(8\gamma^2))\le1+2/\gamma
 \quad(\gamma\ge2).
\]
This proves \eqref{f16v53:all-label}, including its exact zero at $\ell=0$.
Add these nonnegative exponents over a product representation and exponentiate.
Joint Peter--Weyl diagonalization proves \eqref{f16v53:all-link} on all
superpositions, including the gauge-invariant and center-sector restrictions.
Neither a fixed-label asymptotic nor a truncated sum is used.
\end{proof}

\subsection{Operator convexity pays the complete magnetic sandwich}

\begin{proposition}[The contractive power inequality]
\label{f16v53:Jensen}
If $0\preceq K\preceq I$, $0\preceq B\preceq I$ and $1\le p\le2$, then
\begin{equation}
                         (BKB)^p\preceq BK^pB.
 \label{f16v53:power-Jensen}
\end{equation}
No commutation of $B$ and $K$ is needed.
\end{proposition}
\begin{proof}
This is the stated contractive Jensen inequality, but its needed case has a
short proof. For $t>0$, $f_t(x)=x^2/(x+t)=x-t+t^2/(x+t)$ is operator convex:
inversion on positive invertible operators has the compression inequality
$(V^*AV)^{-1}\preceq V^*A^{-1}V$ for an isometry $V$, as follows by completing
a block Schur square. Its affine terms preserve that inequality, so the same
compression inequality holds for $f_t$. Use the isometry
$V=(B,(I-B^2)^{1/2})^{\mathsf T}$ and $X=\operatorname{diag}(K,0)$.
Since $f_t(0)=0$, this gives $f_t(BKB)\preceq Bf_t(K)B$.
For $1<p<2$, integrate the positive representation
\[
 x^p=\frac{\sin(\pi(p-1))}{\pi}
              \int_0^\infty t^{p-2}\frac{x^2}{x+t}\,dt.
\]
The integrals converge in norm for bounded nonnegative operators, giving
\eqref{f16v53:power-Jensen}. At $p=1$ it is equality, and at $p=2$ the difference
is $BK(I-B^2)KB\succeq0$. In particular no monotonicity of $x^p$ for $p>1$
is being assumed.
\end{proof}

Define the new comparison transfer, retaining the complete spatial action,
\begin{equation}
 T^{\rm B}_{\gamma,L}=b_\gamma K^{\rm B}_{\gamma,3V}b_\gamma.
 \label{f16v53:new-transfer}
\end{equation}
It has a smooth strictly positive kernel on the same compact space, is a
positive self-adjoint contraction, and commutes with the original gauge and
center actions. Its positive Perron vector is unique and gauge invariant.
This justifies equality of its full and physical tops, not the use of
non-invariant excited-state trial vectors.

\begin{theorem}[A cutoff-free power envelope for the original full transfer]
\label{f16v53:power-envelope}
For every admitted lattice and every $\gamma\ge2$,
\begin{equation}
 \boxed{\qquad
 (T^{\rm B}_{\gamma,L})^{p_\gamma}
       \preceq T^{\rm W}_{\gamma,L}\preceq T^{\rm B}_{\gamma,L}.
 \qquad}
 \label{f16v53:full-envelope}
\end{equation}
It holds on the full physical space and on each center sector. On every same
orthogonal compression $P$, it also yields
$(PT^{\rm B}P)^{p_\gamma}\preceq PT^{\rm W}P\preceq PT^{\rm B}P$, with
powers taken on $\operatorname{Ran}P$.
\end{theorem}
\begin{proof}
Congruence of \eqref{f16v53:all-link} by the actual $b_\gamma$ gives
\[
 b_\gamma(K^{\rm B})^{p_\gamma}b_\gamma
       \preceq T^{\rm W}\preceq T^{\rm B}.
\]
Apply \eqref{f16v53:power-Jensen} to its first member. This proves the result
without moving a magnetic factor through a convolution. Sector restriction
commutes with the two transfers and hence with their spectral powers. For an
arbitrary same compression, contractive Jensen gives
$(PT^{\rm B}P)^p\preceq P(T^{\rm B})^pP$, which proves the last assertion.
This is not the insertion of repeated projections into the original evolution.
\end{proof}

\subsection{Logarithmic transfer forms, including their domains}

\begin{theorem}[Global logarithmic form comparison and the vacuum subtraction]
\label{f16v53:log-forms}
Put $J^{\rm W}=-\log T^{\rm W}$ and $J^{\rm B}=-\log T^{\rm B}$ by spectral
calculus. These nonnegative self-adjoint operators have the same closed
quadratic-form domain, and
\begin{equation}
 \boxed{\qquad
 J^{\rm B}\preceq J^{\rm W}\preceq(1+\delta_\gamma)J^{\rm B}
 \qquad}
 \label{f16v53:log-form}
\end{equation}
in form sense. Let $\Lambda^{\rm W},\Lambda^{\rm B}$ be their actual transfer
tops and $E_0^{\rm B}=-\log\Lambda^{\rm B}$, $E_0^{\rm W}=-\log\Lambda^{\rm W}$.
Then
\begin{equation}
 0\le E_0^{\rm W}-E_0^{\rm B}\le\delta_\gamma E_0^{\rm B},
 \qquad 0\le E_0^{\rm B}\le C_0V,\quad C_0=30+3\log48.
 \label{f16v53:top-log}
\end{equation}
For the actual vacuum-subtracted operators
$H^{\rm W}=J^{\rm W}-E_0^{\rm W}I$ and
$H^{\rm B}=J^{\rm B}-E_0^{\rm B}I$,
\begin{equation}
 \boxed{\quad
 H^{\rm B}-\delta_\gamma E_0^{\rm B}I
 \preceq H^{\rm W}
 \preceq(1+\delta_\gamma)H^{\rm B}+\delta_\gamma E_0^{\rm B}I.
 \quad}
 \label{f16v53:subtracted-form}
\end{equation}
Neither the two vacuum vectors nor their orthogonal projectors are identified.
\end{theorem}
\begin{proof}
The transfers are injective: $b_\gamma$ is bounded below at each fixed compact
regulator, and every kinetic multiplier is strictly positive. For positive
invertible bounded operators, logarithm is operator monotone; this follows
from its resolvent integral and the reversed order of inverses.
Apply this fact to
$(T^{\rm B})^p+\epsilon I\preceq T^{\rm W}+\epsilon I
\preceq T^{\rm B}+\epsilon I$ and let $\epsilon\downarrow0$ in each scalar
spectral integral. Adding $\log(1+\epsilon)I$ before this passage makes the
negative logarithms nonnegative and monotone. Thus the inequalities hold for
extended quadratic forms, with $-\log((T^{\rm B})^p)=pJ^{\rm B}$ in the
limit. The two inequalities give equality of the finite form domains and
\eqref{f16v53:log-form}; they do not assert equality of operator domains.

The top variational principle in \eqref{f16v53:full-envelope} gives
$(\Lambda^{\rm B})^p\le\Lambda^{\rm W}\le\Lambda^{\rm B}$.
This proves the first part of \eqref{f16v53:top-log}. Since
$\Lambda^{\rm B}\ge\Lambda^{\rm W}\ge e^{-C_0V}$ by V51's original compact
trial bound, the second part follows. This is the only lattice-dependent
scalar input used here. Subtract the two \emph{actual} bottom energies in
\eqref{f16v53:log-form}. Their difference lies in
$[0,\delta_\gamma E_0^{\rm B}]$, proving \eqref{f16v53:subtracted-form}.
The vacuum defect is kept as an additive form term, not silently set to zero.
\end{proof}

For every corresponding ordered sector level, min--max in
\eqref{f16v53:full-envelope} also gives
\[
 (\lambda^{\rm B}_{a,j})^{1+\delta_\gamma}
       \le\lambda^{\rm W}_{a,j}\le\lambda^{\rm B}_{a,j}.
\]
Consequently, writing
$\mathcal E^{\rm B}_{a,j}=-\log(\lambda^{\rm B}_{a,j}/\Lambda^{\rm B})$ and
similarly for $\rm W$,
\begin{equation}
 |\mathcal E^{\rm W}_{a,j}-\mathcal E^{\rm B}_{a,j}|
       \le\delta_\gamma(E_0^{\rm B}+\mathcal E^{\rm B}_{a,j}).
 \label{f16v53:log-levels}
\end{equation}
Indeed the two uncentered logarithmic increments are in
$[0,\delta_\gamma(-\log\lambda^{\rm B}_{a,j})]$ and
$[0,\delta_\gamma E_0^{\rm B}]$, respectively. The former upper endpoint is
larger. This proves the bound with its index-dependent energy retained.

\subsection{An actual-top error linear in volume}

\begin{theorem}[Normalized transfer comparison without a reciprocal Perron factor]
\label{f16v53:normalized}
Let $\widehat T^{\rm W}=T^{\rm W}/\Lambda^{\rm W}$ and
$\widehat T^{\rm B}=T^{\rm B}/\Lambda^{\rm B}$. For every $\gamma\ge2$,
\begin{equation}
 \boxed{\quad
 \|\widehat T^{\rm W}-\widehat T^{\rm B}\|
 \le \min\left\{1,\delta_\gamma(2E_0^{\rm B}+e^{-1})\right\}
 \le \frac{4C_0V+2e^{-1}}\gamma.
 \quad}
 \label{f16v53:normalized-bound}
\end{equation}
For every integer $N\ge1$ the same two powers differ in norm by at most
\begin{equation}
 \min\left\{1,N\delta_\gamma(2E_0^{\rm B}+e^{-1})\right\}.
 \label{f16v53:normalized-powers}
\end{equation}
These are full-space bounds, and hold on every common reducing sector with the
same global actual tops. There is no representation window or additive tail.
\end{theorem}
\begin{proof}
Write $A=T^{\rm B}$, $W=T^{\rm W}$ and $a=\|A\|$, $w=\|W\|$.
The full envelope gives
\[
 0\preceq A-W\preceq A-A^{1+\delta}
                  \preceq\delta A(-\log A).
\]
The last bounded operator uses the continuous value zero at spectral point zero;
it follows from $1-e^{-x}\le x$. For $0\le x\le a$, put $x=au$. Then
\[
 x(-\log x)=a\{uE_0^{\rm B}-u\log u\}
          \le a(E_0^{\rm B}+e^{-1}).
\]
Also $0\le a-w\le a-a^{1+\delta}\le\delta aE_0^{\rm B}$.
Normalize using the \emph{larger} top in the exact identity
\[
 \frac Ww-\frac Aa=\frac{W-A}{a}+\left(\frac1w-\frac1a\right)W.
\]
The second term has norm exactly $(a-w)/a$. The two estimates give
\eqref{f16v53:normalized-bound}, without replacing either top by one or dividing
an absolute error by an exponentially small lower bound. Both normalized
operators lie between zero and identity, so their difference has norm at most
one. Telescoping their generally noncommuting powers proves
\eqref{f16v53:normalized-powers}; the same unit bound applies to their powers.
No false implication $A^N\succeq W^N$ for $N>1$ is used.
\end{proof}

\begin{theorem}[A larger joint domain for the new reference]
\label{f16v53:critical}
Set $h=\gamma^{-1/3}/L$, with even $L=L(\gamma)\ge8$. If
\begin{equation}
                         L^4\gamma^{-2/3}\longrightarrow0,
 \label{f16v53:domain}
\end{equation}
then for every finite $t_1>0$,
\begin{equation}
 \boxed{\quad
 \sup_{0\le t\le t_1}
 \| (\widehat T^{\rm W})^{\lfloor t/h\rfloor}
       -(\widehat T^{\rm B})^{\lfloor t/h\rfloor}\|\longrightarrow0.
 \quad}
 \label{f16v53:critical-powers}
\end{equation}
In particular $L=O(\gamma^b)$ with fixed $0\le b<1/6$ suffices. Corresponding
normalized ordered eigenvalues differ by $o(h)$ uniformly in their index;
their logarithms differ by $o(h)$ on spectral regions bounded away from zero.
This theorem compares Wilson with the newly declared subordinated reference,
\emph{not} with V51's ordinary heat reference on this enlarged domain.
\end{theorem}
\begin{proof}
Divide \eqref{f16v53:normalized-bound} by $h$. Since $V=L^3$ and $L\ge8$, the
result is at most $C L^4\gamma^{-2/3}$. Use the bound on powers at
$N\le t_1/h$, with the zeroth power identical. The polynomial condition is
$4b<2/3$; for example $b=1/8$ gives an $O(\gamma^{-1/6})$ upper bound on
bounded critical-time intervals. Min--max on the same sector gives the
uniform ordered-eigenvalue comparison; the logarithm is Lipschitz on any
fixed interval separated from zero. No individual growing-lattice spectrum or
physical time spacing is supplied by comparison of the two operators.
\end{proof}

\begin{proposition}[The logarithmic generators also have a full resolvent comparison]
\label{f16v53:log-resolvents}
Let $\mathcal H^{\rm W}_h=H^{\rm W}/h$ and
$\mathcal H^{\rm B}_h=H^{\rm B}/h$ be the \emph{logarithmic}, vacuum-subtracted
nonnegative operators. For $z>0$ put
\[
 \epsilon_z=\delta_\gamma\left(1+\frac{E_0^{\rm B}}{hz}\right).
\]
If $\epsilon_z<1$, then
\begin{equation}
 \| (\mathcal H^{\rm W}_h+z)^{-1}
      -(\mathcal H^{\rm B}_h+z)^{-1}\|
       \le\frac{\epsilon_z}{z(1-\epsilon_z)}.
 \label{f16v53:log-resolvent}
\end{equation}
It tends to zero on \eqref{f16v53:domain}, with both vacuum modes retained.
\end{proposition}
\begin{proof}
The common-domain form bound \eqref{f16v53:subtracted-form} implies
\[
 |\mathfrak q_{\mathcal H^{\rm W}_h}[u]
          -\mathfrak q_{\mathcal H^{\rm B}_h}[u]|
 \le\delta_\gamma\left[\mathfrak q_{\mathcal H^{\rm B}_h}[u]
                         +(E_0^{\rm B}/h)\|u\|^2\right].
\]
Relative to the positive form $\mathcal H^{\rm B}_h+z$, the difference is
represented by a bounded self-adjoint operator $D$ of norm at most
$\epsilon_z$. Thus the first resolvent equals
$(\mathcal H^{\rm B}_h+z)^{-1/2}(I+D)^{-1}
(\mathcal H^{\rm B}_h+z)^{-1/2}$. Subtract the reference resolvent and use
$\|(I+D)^{-1}-I\|\le\epsilon_z/(1-\epsilon_z)$.
The top bound in \eqref{f16v53:top-log} proves the last assertion. This does not
bound a vacuum-orthogonal inverse or identify the two vacuum projections.
\end{proof}

\subsection{The reference changed, not the original operator or its clock}

The larger domain above must not be attached to the ordinary heat comparison.
For all $s\ge0$, direct integration of
$0\le x/\gamma-\operatorname{arsinh}(x/\gamma)\le x^3/(6\gamma^3)$ gives
\begin{equation}
 0\le\frac{s}{2\gamma}-F_\gamma(s)
             \le\frac{s(s+2)}{24\gamma^3}.
 \label{f16v53:heat-distance}
\end{equation}
Thus the new reference agrees with the ordinary small-energy heat tangent,
but its large-$s$ dispersion is of order $\sqrt s\log s$, not $s$.
The ratio $e^{-F_\gamma(s)}/e^{-s/(2\gamma)}$ is unbounded at fixed coupling.
Equation \eqref{f16v53:heat-distance} reproduces a quadratic total-energy cost
if one forces that additional comparison on an increasing spectral window.
V52's ordinary-heat conclusions remain as written; the present theorem is not
a proof of them on $b<1/6$.

Independence of mixture variables applies only to the declared kinetic reference.
The magnetic sandwich remains interacting. No random-clock representation of
Wilson kinetics, or subordination of the full magnetic Hamiltonian, is asserted.
The exponent $p_\gamma$ is functional calculus, not a change of native duration.

The remaining $V/\gamma$ cost is uncertainty in the extensive bottom energy,
not a reciprocal Perron factor. Vacuum subtraction therefore retains the
additive term in \eqref{f16v53:subtracted-form}. For a separately supplied
physical spacing $a_t$, the power handoff requires
$(V+1)/(\gamma a_t)\to0$; no value of $a_t$ is chosen.

The unresolved growing-lattice problem is the reference's interacting spatial
and soft spectrum, with all returning channels retained. Comparison alone does
not construct its limit or its vacuum representation. The next companion
checkpoint remains V55.

\paragraph{Verification and preservation.}
The accompanying audit records the targeted source review, exact finite tests,
actual executions and document checks. The tests are not native spectral
enclosures or formal proof verification. Removing this terminal insertion and
reversing the literal title version recovers V52 byte for byte. The earlier
analytic lineage has not been independently recertified.

\begin{firewall}
The cutoff-free power and logarithmic comparisons use a new positive
subordinated reference, not ordinary heat. They do not prove an individual
growing-lattice spectral limit, align vacuum projectors without a separation
estimate, specify physical time or running coupling, or construct an
interacting continuum Yang--Mills theory with a mass gap.
\end{firewall}
% END F16 V53 APPEND

% BEGIN F16 V54 APPEND -- 2026-09-18
% Insert before the final document-ending command of cumulative V53.
% No predecessor theorem, native operator, Haar law, or physical clock is changed.

\clearpage
\section{V54: a joint-regulator critical branch for the subordinated reference}
\label{f16v54:context}

V53 removes the electric spectral cutoff by replacing the ordinary heat
comparison with one positive subordinated kinetic reference.  Its comparison
with the original Wilson transfer is already normalized at the actual Perron
values and costs only the extensive bottom energy.  What it does not provide is
an individual growing-lattice spectral limit for that reference.  We now attack
that missing interacting problem.

The main point of this iteration is that the fixed-$L$ critical analysis does
admit a joint-regulator version once three different volume effects are kept
separate.  The based gauge geometry is controlled by V50 in trace ideals; the
hard Gaussian inverse is controlled by Fourier shell sums rather than by its
smallest gap alone; and the soft radial tail is paid by the uniform harmonic
contrast of V49.  These estimates yield collective compactness and a Mosco
limit for the actual-top shifted subordinated reference.  V53 then transfers
that limit back to the original Wilson transfer.

The result is a spectral theorem on a growing number of lattice sites, but it is
\emph{not} yet a continuum mass-gap theorem.  The lattice spacing in physical
time and space has not been calibrated, the spatial physical volume has not
been fixed independently of $L$, and no Osterwalder--Schrader continuum
reconstruction is constructed here.

\subsection{The joint parameter and the unchanged nine-dimensional model}

Let $L=2m\ge8$, $V=L^3$, and retain
\[
 \varepsilon=\gamma^{-1/2},\qquad
 h=(\gamma L^3)^{-1/3}=\gamma^{-1/3}/L.
\]
Write $T^{\rm B}_{\gamma,L}$ for V53's subordinated reference
\[
 T^{\rm B}_{\gamma,L}
 =b_\gamma R_\gamma^{\otimes3V}b_\gamma,
 \qquad b_\gamma=e^{-\gamma S/2},
 \qquad R_\gamma=e^{-F_\gamma(\mathcal L_G)},
\]
with the complete original spatial Wilson action $S$.  Let
$\Lambda^{\rm B}_{\gamma,L}$ be its actual Perron value.  In a center sector
$a\in\{0,1\}^3$ define
\begin{equation}
 \mathcal H^{\rm B}_{\gamma,L,a}
 =-h^{-1}\log\bigl(T^{\rm B}_{\gamma,L}|_{\mathcal H_a}/
                         \Lambda^{\rm B}_{\gamma,L}\bigr)\ge0 .
 \label{f16v54:generator}
\end{equation}
The logarithm is the spectral logarithm of the injective compact positive
operator.  Its quadratic form is closed.

Retain without modification V44's invariant matrix-model operator
\begin{equation}
 H_{\rm mat}=-\tfrac12\Delta_q+2\mathcal Q(q),\qquad
 \mathcal Q(q)=\sum_{\mu<\nu}|q_\mu\times q_\nu|^2
 \label{f16v54:model}
\end{equation}
on $L^2(\mathbb R^9)^{\rm inv}$, with eigenvalues
$e_0<e_1\le e_2\le\cdots$ counted with multiplicity.  The positive gap
$e_1-e_0$ and compact resolvent are inherited results; no numerical value is
inserted here.

The joint error ledger is
\begin{equation}
 \boxed{\quad
 \Theta_{\gamma,L}
 =L^2\gamma^{-1/3}+L^4\gamma^{-2/3}+L^3\gamma^{-1}.
 \quad}
 \label{f16v54:Theta}
\end{equation}
Every term has a different origin below.  In particular no one of them is
renamed as a generic $o_L(1)$ remainder.

\subsection{The hard return is summed with its full Fourier weight}

Let $\nu_L(k)=4\sum_\mu\sin^2(\pi k_\mu/L)$ for
$k\in\{0,\ldots,L-1\}^3$.  The nonzero hard Gaussian transfer frequencies are
\[
 \omega_L(k)=2\operatorname{arsinh}\bigl(\sqrt{\nu_L(k)}/2\bigr).
\]
Put $d_L(k)=1-e^{-\omega_L(k)}$.

\begin{proposition}[Volume, rather than inverse-gap, hard susceptibilities]
\label{f16v54:hard-sums}
There is an absolute $C$ such that for all $L\ge8$,
\begin{equation}
 \boxed{\qquad
 \sum_{k\ne0}d_L(k)^{-1}\le CL^3,
 \qquad
 \sum_{k\ne0}d_L(k)^{-2}\le CL^3 .
 \qquad}
 \label{f16v54:hard-sums-eq}
\end{equation}
The same estimates hold with the fixed finite polarization and color
multiplicities of the hard Gaussian space.  Consequently a local hard source
whose Fourier coefficients obey Parseval estimates can be returned through the
complete reduced hard inverse at a cost proportional to its squared local
$\ell^2$ size times $L^3$, rather than that size times the operator norm
$O(L)$ and then another factor counting modes.
\end{proposition}
\begin{proof}
For $0\le x\le2$, $1-e^{-2\operatorname{arsinh}(x/2)}$ is comparable to
$x$.  Since $\sqrt{\nu_L(k)}$ is comparable to
$\min\{1,|\widetilde k|/L\}$ for the centered momentum representative
$\widetilde k$, the two sums are bounded by
\[
 C\sum_{0<|\widetilde k|\le CL}\left(1+\frac L{|\widetilde k|}\right),
 \qquad
 C\sum_{0<|\widetilde k|\le CL}
       \left(1+\frac {L^2}{|\widetilde k|^2}\right).
\]
Three-dimensional shell counts are $O(r^2)$.  The first weighted sum is
$O(L\sum_{r\le CL}r)=O(L^3)$; the second is
$O(L^2\sum_{r\le CL}1)=O(L^3)$.  Modes away from zero contribute only their
number.  This proves both estimates.  The statement for local sources follows
by diagonalizing the fixed Gaussian hard transfer and applying Cauchy--Schwarz
before summing modes.  No bound by the smallest $d_L(k)$ is used.
\end{proof}

This is the same dimensional mechanism already visible in V49's hard-gap
ledger: the smallest hard gap is $O(L^{-1})$, but its inverse is not paid at
every one of the $O(L^3)$ modes.

\subsection{A corrected local graph estimate with uniform bookkeeping}

Use V50's size-independent based normal tube
$s=y+\mathsf K a$, with
\[
 y=\varepsilon u,\qquad a=hq,\qquad \mathsf K^*\mathsf K=VI,
\]
and its exact half-density.  Let $\varphi_L$ denote the normalized hard Gaussian
ground function.  The hard Gaussian transfer itself depends on $L$, but its
ground line and complement are explicit Fourier products.

The first variation in the constant field has zero hard-ground expectation, as
in V45--V48.  Denote its ground-to-hard block by $C_L(q)$.  The local magnetic
and kinetic words have finite incidence, so Parseval and the preceding
proposition give, on every fixed $|q|\le R$,
\begin{equation}
 \|C_L(q)\varphi_L\|^2\le C_Rh^2L^3,
 \qquad
 \langle C_L(q)\varphi_L,
     (I-A_{+,L})^{-1}C_L(q)\varphi_L\rangle
       \le C_Rh^2L^3.
 \label{f16v54:return-cost}
\end{equation}
Here $A_{+,L}$ is the normalized hard Gaussian transfer.  The inverse is on the
orthogonal complement of its ground line.

\begin{theorem}[Uniform corrected graph expansion on bounded critical sets]
\label{f16v54:graph}
For every fixed $R<\infty$ and every finite invariant graph-core set
$\mathcal D\subset C_c^\infty(\mathbb R^9)^{\rm inv}$, there are physical
recovery maps $\Psi_{\gamma,L}$, using the same exact tube measure and a
hard-return corrector, such that uniformly for $f\in\mathcal D$,
\begin{equation}
 \left\|
 \frac{T^{\rm B}_{\gamma,L}\Psi_{\gamma,L}f}
      {\Lambda^{\rm B}_{\gamma,L}}
 -\Psi_{\gamma,L}\bigl[f-h(H_{\rm mat}-e_0)f\bigr]
 \right\|
 \le C_{R,\mathcal D}\,h\,\Theta_{\gamma,L}.
 \label{f16v54:recovery}
\end{equation}
The Gram error is bounded by the same right side without the leading $h$.
The statement holds after the exact center-character folding in every one of
the eight center sectors.
\end{theorem}
\begin{proof}
We record the volume ledger because this is the new point.  The complete local
Taylor expansion is performed only after V50's relative density and quotient
metric have been assembled.  Their expectation errors on a bounded critical
soft set are bounded by
\[
 C_R\left(L^3/\gamma+L\gamma^{-2/3}\right).
\]
After division by $h$ these are respectively
$L^4\gamma^{-2/3}$ and $L^2\gamma^{-1/3}$.

Odd hard terms have zero hard-ground expectation but not zero hard residual.
Solve them with the fixed Gaussian reduced inverse.  Translation invariance and
finite incidence give a local source square whose total Fourier mass is
$O(L^3)$; Proposition~\ref{f16v54:hard-sums} pays its reduced return without an
extra inverse-gap factor.  Its contribution after the critical rescaling is
bounded by $C_RL^3/\gamma$, hence by the last term in
\eqref{f16v54:Theta} before division by the already extracted one-step scale.
The soft-dependent first variation is bounded by
\eqref{f16v54:return-cost}; since $hL^3= L^2\gamma^{-1/3}$, its relative graph
cost is the first term of \eqref{f16v54:Theta}.

The even Taylor remainder and the trace-ideal quotient corrections cost at most
the second and third terms.  V53's subordinator has mean
$\frac12\operatorname{arsinh}(1/\gamma)$ and second centered moment
$O(\gamma^{-2})$; replacing its tangent covariance by the leading covariance
therefore contributes only $O(L^3/\gamma)$ to the relative graph ledger.  All
finite protected constraints are imposed by the same exact physical recovery
map; no auxiliary projector is inserted between transfer steps.

After the hard corrector is added, the surviving order-$h$ soft operator is
exactly the kinetic Laplacian plus the original constant-field quartic, namely
$H_{\rm mat}$.  Subtracting the Rayleigh infimum fixes the scalar $e_0$ and the
actual reference top simultaneously by min--max; no external vacuum
normalization is inserted.  Center-translated tubes are separated by distance
at least $c\sqrt L$ in the link metric, while the displacement scale is
$\varepsilon$.  Their cross kernels are $O(e^{-c\gamma L})$, uniformly smaller
than the displayed ledger.  This proves the estimate.
\end{proof}

The theorem is a graph estimate on a fixed model core.  Spectral completeness
still requires an all-vector compactness statement.  We prove that next rather
than inferring it from recovery.

\subsection{A joint-regulator no-escape estimate}

Let $\mathfrak E_{\gamma,L,a}$ be the exact measured extraction obtained from
V50's tubes and center folding.  Write
$F_{\gamma}=\mathfrak E_{\gamma,L,a}\psi$ and
$f_\gamma=\langle\varphi_L,F_\gamma\rangle_{\rm hard}$.

\begin{theorem}[Uniform hard capture and soft radial tightness]
\label{f16v54:tightness}
Assume
\begin{equation}
 \Theta_{\gamma,L}\longrightarrow0 .
 \label{f16v54:joint-condition}
\end{equation}
For every fixed $K$, uniformly over all unit vectors in the actual-top band
\[
 1_{[e^{-Kh},1]}
 \left(T^{\rm B}_{\gamma,L}/\Lambda^{\rm B}_{\gamma,L}\right)\mathcal H_a,
\]
one has
\begin{align}
 \|F_\gamma-\varphi_L\otimes f_\gamma\|&\longrightarrow0,
 \label{f16v54:hard-capture}\\
 \limsup_{\gamma\to\infty}
   \int_{|q|>R}|f_\gamma(q)|^2dq
 &\le \frac{C(K+1)}R ,\qquad R\ge1.
 \label{f16v54:tail}
\end{align}
Furthermore the localized profiles are uniformly Fourier tight on every fixed
soft ball.  Hence every bounded-energy sequence has a strongly convergent
subsequence after extraction; no norm can disappear at soft infinity.
\end{theorem}
\begin{proof}
The hard complement has one-step separation $d_L(k)$ in each Fourier mode.
The same local-source estimate as in Proposition~\ref{f16v54:hard-sums}, now
applied to the quadratic form of a band vector, gives hard-complement mass at
most $C(K+1)\gamma^{-1/3}$ plus $C\Theta_{\gamma,L}$; this proves
\eqref{f16v54:hard-capture}.

For the soft tail use a quadratic dyadic partition in the exact constant
coordinate $a=hq$.  On an annulus $r\le|a|\le4r$, the transverse hard return is
quadratic in $r$, while the local harmonic threshold has the uniform linear
contrast of V49.  In the present reference the latter is obtained directly by
functional calculus of the subordinated kinetic kernel; its tangent frequency
is the same and the dispersion error is $O(\gamma^{-1})$ relative to the
critical form.  Thus, after a fixed decrease of the common V50 tube,
\[
 \langle\chi\psi,
   (\Lambda^{\rm B}-T^{\rm B})\chi\psi\rangle
 \ge c\Lambda^{\rm B} r\|\chi\psi\|^2
   -C\Lambda^{\rm B}(r^2+h\Theta_{\gamma,L})\|\chi\psi\|^2.
\]
The localization error is paid by the second moment of the subordinated step.
Its mean time is $O(\gamma^{-1})$, so the squared-gradient cost of the dyadic
partition is $C\gamma^{-1}\rho^{-2}$ in physical coordinates.  Taking
$\rho=hR$ and using
$\gamma^{-1}/(h^2R^2)=L^2\gamma^{-1/3}R^{-2}$ gives exactly the first ledger
term.  Finite overlap avoids a number-of-annuli factor.  The fixed exterior of
all central tubes has the V49 harmonic loss $c/L$ relative to the top, which is
$c\gamma^{1/3}h$ and hence dominates the $O(h)$ band scale.

The band energy is at most $C\Lambda^{\rm B}hK$.  Divide the resulting radial
form by $hR$ and pass first in the joint sequence, then in $R$.  This proves
\eqref{f16v54:tail}.  On a fixed soft ball, the exact subordinate kinetic
multiplier satisfies
$1-e^{-F_\gamma(s)}\ge c\min\{s/\gamma,1\}$ for the relevant bounded critical
range.  Together with hard capture this gives the same bounded Fourier tails
used in V46--V47.  Position and Fourier tightness therefore give strong
compactness.  The constants in this argument are the uniform constants from
V49--V50 and the displayed hard shell sums; no fixed-$L$ entry function remains.
\end{proof}

\subsection{Norm-resolvent convergence and ordered spectra}

\begin{theorem}[Joint-regulator critical spectrum of the subordinated reference]
\label{f16v54:reference-spectrum}
Under \eqref{f16v54:joint-condition}, the actual-top shifted forms
\eqref{f16v54:generator}, after the measured extraction, converge in Mosco
sense to $H_{\rm mat}-e_0$ and are collectively compact.  Equivalently, for
every $z>0$,
\begin{equation}
 \left\|
 \mathfrak E_{\gamma,L,a}
   (\mathcal H^{\rm B}_{\gamma,L,a}+z)^{-1}
 \mathfrak E_{\gamma,L,a}^*
 -\mathcal I_L(H_{\rm mat}-e_0+z)^{-1}\mathcal I_L^*
 \right\|\longrightarrow0,
 \label{f16v54:norm-resolvent}
\end{equation}
where $\mathcal I_L f=\varphi_L\otimes f$.  The complementary resolvent norm
tends to zero.

For every fixed index $j\ge1$ in every center sector,
\begin{equation}
 \boxed{\qquad
 -\frac1h\log
 \frac{\lambda^{\rm B}_{a,j}(\gamma,L)}
      {\Lambda^{\rm B}_{\gamma,L}}
 \longrightarrow e_{j-1}-e_0 .
 \qquad}
 \label{f16v54:B-spectrum}
\end{equation}
Isolated finite model clusters have the correct eventual multiplicity, and the
extracted reference vacuum converges strongly to
$\varphi_L\otimes\phi_0$ in the moving representation.
\end{theorem}
\begin{proof}
The recovery half of Mosco convergence is Theorem~\ref{f16v54:graph} followed
by graph-core density.  For the lower half, take a bounded-energy sequence.
Theorem~\ref{f16v54:tightness} supplies strong subsequential compactness and
hard-ground capture.  Pairing against the corrected core vectors identifies
every weak limit with the model form, while lower semicontinuity of the
subordinated kinetic multiplier and the nonnegative magnetic action gives the
liminf.  No mass is lost because of the uniform radial estimate.

The same compactness applies to resolvent vectors with bounded right sides.
Exactly as in the compact-screen criterion, strong resolvent convergence plus
collective compactness upgrades to norm convergence.  Spectral projection
calculus then gives fixed-index eigenvalue and multiplicity convergence.  The
scalar relation
$-\log(1-hx)=hx+o(h)$ is uniform on each fixed limiting cluster, yielding
\eqref{f16v54:B-spectrum}.  Ground-state positivity removes the phase ambiguity
for the vacuum profile.  Center cross-tube terms are exponentially smaller than
$h$, so the same local limit holds in each character sector; no tunneling
prefactor is inferred.
\end{proof}

The abstract compact-screen mechanism is already part of the supplied Grand
Tensor.  What is new here is verification of its two missing inputs for this
joint-regulator reference: the explicit recovery ledger and the no-escape
estimate.  We do not relabel the abstract theorem itself as new.

\subsection{Transfer back to the original Wilson operator}

V53 proves for the actual-top normalized Wilson and subordinated transfers
\[
 \left\|\widehat T^{\rm W}_{\gamma,L}
       -\widehat T^{\rm B}_{\gamma,L}\right\|
 \le \frac{4C_0L^3+2/e}{\gamma},
 \qquad C_0=30+3\log48.
\]
Therefore
\begin{equation}
 \frac1h
 \left\|\widehat T^{\rm W}_{\gamma,L}
       -\widehat T^{\rm B}_{\gamma,L}\right\|
 \le C L^4\gamma^{-2/3}+o(1).
 \label{f16v54:Wilson-transfer}
\end{equation}
This is already one term of \eqref{f16v54:Theta}.

\begin{theorem}[Growing-lattice critical spectrum of the original Wilson transfer]
\label{f16v54:Wilson-spectrum}
Assume \eqref{f16v54:joint-condition}.  Then for every fixed center character
$a$ and fixed $j\ge1$,
\begin{equation}
 \boxed{\qquad
 -\frac1h\log
 \frac{\lambda^{\rm W}_{a,j}(\gamma,L)}
      {\Lambda^{\rm W}_{\gamma,L}}
 \longrightarrow e_{j-1}-e_0 .
 \qquad}
 \label{f16v54:Wilson-levels}
\end{equation}
In particular, in the center-neutral sector,
\begin{equation}
 \boxed{\qquad
 -\frac1h\log
 \frac{\lambda^{\rm W}_{0,2}}
      {\Lambda^{\rm W}}
 \longrightarrow e_1-e_0>0 .
 \qquad}
 \label{f16v54:neutral-gap}
\end{equation}
A sufficient explicit joint family is
\begin{equation}
 \boxed{\qquad L=O(\gamma^b),\qquad 0\le b<\frac16.\qquad}
 \label{f16v54:power-domain}
\end{equation}
\end{theorem}
\begin{proof}
Condition \eqref{f16v54:joint-condition} is exactly satisfied on
\eqref{f16v54:power-domain}: the first two terms require $b<1/6$, and the
third only $b<1/3$.  The normalized operator difference in
\eqref{f16v54:Wilson-transfer} is $o(h)$.  Min--max in each common reducing
center sector therefore transfers every fixed normalized eigenvalue from the
reference with an $o(h)$ error.  The limiting eigenvalues lie near one, so the
logarithm is uniformly Lipschitz there and the same statement holds for their
logarithmic deficits.  Insert Theorem~\ref{f16v54:reference-spectrum}.
\end{proof}

This theorem is a growing-\emph{regulator} result.  It is not yet a theorem at
fixed physical volume, because no relation between $L$, the spatial lattice
spacing, and a physical length has been imposed.  It also does not say that the
full first physical excitation is neutral: the center-charged tops may remain
much closer to the vacuum, as the fixed-$L$ analysis already shows.

\subsection{Physical-time boundary and the next joint limit}

If a physical Euclidean time spacing $a_t(\gamma,L)$ is supplied independently,
then the neutral physical energy corresponding to the theorem is
\begin{equation}
 E_{\rm neu}^{\rm phys}
 =\frac h{a_t(\gamma,L)}
      [e_1-e_0+o(1)].
 \label{f16v54:physical-energy}
\end{equation}
Nothing in this iteration proves that $h/a_t$ converges to a positive finite
constant.  Taking $a_t=h$ would be a choice, not a deduction.

The next genuinely different problem is therefore a calibrated spatial and
time continuum trajectory.  One needs a nontrivial local-vacuum reconstruction,
a relation between $\gamma,L$ and the physical lattice spacings, and control of
local observables/source Grams under that scaling.  The growing-regulator
spectral theorem above is an input to that problem; it does not solve it.

There is also a center-sector fork.  For a vacuum representation in which
nontrivial electric-flux sectors become superselected, the neutral limit is the
relevant local spectral edge.  If those sectors remain in one finite-volume
Hilbert space, their exponentially smaller splitting must be calibrated
separately.  No assumption about that continuum superselection is made here.

The present route is complementary to the concentration--compactness and exact
Feshbach programmes proposed in the supplied broader literature sweep: the
proof uses the same no-vanishing principle, but obtains it here through an
explicit uniform annular estimate rather than a minimal-counterexample
argument.  Those alternative routes remain useful if the subsequent physical
continuum passage loses the compactness available on this joint lattice
trajectory.

\paragraph{Verification and preservation.}
The accompanying audit records the source review, scale ledger, symbolic
checks, compilation and rendering.  The checker tests the shell sums, joint
power conditions, scalar spectral transfer and finite-dimensional Mosco/min--max
logic on exact fixtures; it is not a proof assistant or a native spectral
enclosure.  Removing this exact insertion and reversing the literal title
version recovers V53 byte for byte.  The earlier analytic lineage has not been
independently recertified.  The next companion checkpoint remains V55.

\begin{firewall}
V54 proves, within the stated inherited framework, a joint-regulator critical
spectral limit for the positive subordinated reference and transfers its fixed
ordered critical levels to the original Wilson transfer on
$L=O(\gamma^b)$, $b<1/6$.  The limit is dimensionless and concerns a growing
number of lattice sites.  It does not calibrate physical time or spatial lattice
spacing, construct a continuum vacuum representation, prove a uniform charged
sector tunneling law, or establish the four-dimensional continuum Yang--Mills
mass gap.
\end{firewall}
% END F16 V54 APPEND

% BEGIN F16 V55 APPEND -- 2026-09-18
% Insert before the final document-ending command of cumulative V54.
% No predecessor theorem, transfer normalization, Haar law, or physical clock is changed.

\clearpage
\section{V55: bare Wilson calibration, toron decoupling, and the local-gap fork}
\label{f16v55:context}

V54 reaches a joint-regulator spectral theorem, but its last line leaves the
conversion to physical units deliberately open.  At this point choosing a time
spacing by hand would be circular: the isotropic Wilson action already fixes a
bare space--time calibration, and asymptotic freedom fixes the scale on which
its bare coupling approaches zero.  We therefore go upstream and compare those
facts with the critical branch that has just been constructed.

The result is a correction of target, not a retraction of the preceding spectral
theorem.  The V44--V54 branch describes the compact constant-field, or toron,
sector.  Under the standard isotropic Wilson calibration its physical energy
scale is proportional to the bare coupling to the power $2/3$ divided by the
box size.  Moreover a fixed-physical-volume continuum trajectory has a number
of lattice sites growing exponentially, rather than polynomially, in the bare
Wilson parameter.  Thus V54's polynomial joint domain does not contain the
usual continuum trajectory, and its lowest center-neutral finite-torus level
cannot be identified directly with a positive infinite-volume mass.

This is consistent with the supplied Grand Tensor warning that volume-wrapping
soft vectors and a positive local-vacuum mass are logically compatible unless
the former survive the local Osterwalder--Schrader quotient.  We make that fork
explicit and prove a soft restriction estimate for contractible Wilson loops.
The next target is consequently a renormalized local-vacuum/source-Gram theorem,
or an RG/Feshbach theorem that replaces the bare soft coupling by the coupling
at a physical scale.  Another improvement of the finite-torus lowest eigenvalue
would not settle that problem.

\subsection{The isotropic transfer fixes the bare time calibration}

Let $a_s$ and $a_t$ denote spatial and Euclidean-time lattice spacings.  Write
$E_e^2$ for the usual Kogut--Susskind electric Casimir, normalized so that on
spin $j$ it has eigenvalue $j(j+1)$.  The round group Laplacian used throughout
this lineage has eigenvalue
\[
 s_\ell=\ell(\ell+2)=4j(j+1),\qquad \ell=2j,
\]
so that $\mathcal L_G=4E^2$ on one link.

The leading logarithm of the present one-step transfer has kinetic and magnetic
parts
\begin{equation}
 -\log T_{\gamma,L}
 =\frac1{2\gamma}\sum_e\mathcal L_e
   +\gamma\sum_p\left(1-\frac12\operatorname{Tr}U_p\right)
   +\hbox{higher transfer corrections}.
 \label{f16v55:bare-log}
\end{equation}
The corresponding Kogut--Susskind Hamiltonian convention is
\begin{equation}
 H_{\rm KS}
 =\frac{g_0^2}{2a_s}\sum_e E_e^2
 +\frac{4}{g_0^2a_s}\sum_p
      \left(1-\frac12\operatorname{Tr}U_p\right).
 \label{f16v55:KS}
\end{equation}
This is the standard Wilson/Kogut--Susskind normalization; see, for example,
\href{https://inspirehep.net/files/e438c1a894d6acca313ca4bf886de405}
{the lattice-gauge review around its Eq.~(4.171)}.  An anisotropic Wilson action
has separate temporal and spatial coefficients proportional respectively to
$a_s/a_t$ and $a_t/a_s$; see
\href{https://inspirehep.net/files/b6f4dee848c43dc1d13966736ea934c0}
{the anisotropic transfer review, Eqs.~(1)--(2)}.

\begin{proposition}[Bare calibration of the present isotropic transfer]
\label{f16v55:calibration}
Matching \eqref{f16v55:bare-log} to $a_tH_{\rm KS}$ gives
\begin{equation}
 \gamma=\frac{4a_s}{g_0^2a_t}
 \qquad\hbox{from the electric term},
 \qquad
 \gamma=\frac{4a_t}{g_0^2a_s}
 \qquad\hbox{from the magnetic term}.
 \label{f16v55:anisotropic-match}
\end{equation}
Hence the single-coupling transfer used in the present lineage is the isotropic
bare theory:
\begin{equation}
 \boxed{\qquad a_t=a_s=:a,\qquad \gamma=\frac4{g_0^2}.\qquad}
 \label{f16v55:isotropic-calibration}
\end{equation}
The equality of the physical lattice spacings is also enforced by the exact
hypercubic symmetry of the isotropic Euclidean Wilson action; no independent
bare anisotropy parameter is present to tune it.
\end{proposition}
\begin{proof}
Since $\mathcal L=4E^2$, the kinetic term in
\eqref{f16v55:bare-log} is $2\gamma^{-1}\sum E^2$.  Equating it to
$a_tg_0^2(2a_s)^{-1}\sum E^2$ gives the first relation in
\eqref{f16v55:anisotropic-match}.  Equating the magnetic coefficient $\gamma$
to $4a_t/(g_0^2a_s)$ gives the second.  Their product and ratio give
$a_t=a_s$ and then $\gamma=4/g_0^2$.  This identifies the convention; it does
not claim that a bare tree-level Hamiltonian coefficient already equals a
renormalized physical observable.
\end{proof}

\subsection{The standard continuum trajectory is exponentially larger than V54's domain}

For pure four-dimensional $\SU(2)$ Wilson theory, $\beta=4/g_0^2$ in the
standard notation, hence $\beta=\gamma$ here.  The asymptotically free lattice
scale obeys
\begin{equation}
 \log\frac1{a\Lambda_L}
 =\frac{3\pi^2}{11}\gamma+O(\log\gamma),
 \label{f16v55:asymptotic-scaling}
\end{equation}
where the $O(\log\gamma)$ contains the familiar two-loop power prefactor.  The
one-loop coefficient in exactly this $\SU(2)$ Wilson convention is displayed,
for example, in
\href{https://cds.cern.ch/record/134284/files/198105311.pdf}
{the classic finite-temperature scaling paper}, which writes
$a\Lambda_L\simeq\exp[-3\pi^2\beta/11]$ with $\beta=4/g_0^2$.  Equivalently,
the standard beta-function form has $b_0=11/(24\pi^2)$ and yields the same
coefficient.

\begin{theorem}[Quantifier mismatch with a fixed physical box]
\label{f16v55:trajectory}
Let the physical spatial side length be fixed at $\ell>0$ and impose the
isotropic Wilson calibration $a=\ell/L$.  Along the usual asymptotically free
continuum trajectory,
\begin{equation}
 \boxed{
 \log L=\frac{3\pi^2}{11}\gamma+O(\log\gamma),
 \qquad
 L=\exp\!\left[\frac{3\pi^2}{11}\gamma+O(\log\gamma)\right].
 }
 \label{f16v55:L-exponential}
\end{equation}
Consequently this trajectory is outside every polynomial domain
$L=O(\gamma^b)$, in particular outside the sufficient V54 domain $b<1/6$.
More directly, V54's joint ledger
\[
 \Theta_{\gamma,L}=L^2\gamma^{-1/3}
 +L^4\gamma^{-2/3}+L^3\gamma^{-1}
\]
does not tend to zero on \eqref{f16v55:L-exponential}.
\end{theorem}
\begin{proof}
Insert $a=\ell/L$ into \eqref{f16v55:asymptotic-scaling}.  The fixed scalar
$\ell\Lambda_L$ contributes only $O(1)$ to the logarithm, giving
\eqref{f16v55:L-exponential}.  Exponential growth dominates every power of
$\gamma$, and each positive-power $L$ term in $\Theta_{\gamma,L}$ diverges.
No statement from V54 is extrapolated beyond its hypotheses.
\end{proof}

Thus the new joint-regulator theorem is mathematically nontrivial but is not yet
on the standard fixed-physical-volume continuum route.  This is a stronger
obstruction than merely not knowing the numerical relation between $a_t$ and
$h$: the two limits currently have different quantifiers.

\subsection{What the quartic branch would do under bare physical calibration}

V54 identifies, on its own admitted trajectory,
\begin{equation}
 -\log\frac{\lambda_{0,2}}{\Lambda}
 =h\,[e_1-e_0+o(1)],
 \qquad
 h=\frac{\gamma^{-1/3}}L .
 \label{f16v55:dimensionless-soft}
\end{equation}
The following implication is deliberately conditional: it asks what that same
bare formula would mean if it persisted to an isotropic fixed-box trajectory.

\begin{proposition}[Bare toron scale cannot by itself give a positive fixed-box mass]
\label{f16v55:bare-collapse}
Suppose a center-neutral finite-volume branch has
\begin{equation}
 -\log(\lambda_1/\Lambda)=h\,\mu_{\gamma,L},
 \qquad \sup|\mu_{\gamma,L}|<\infty,
 \label{f16v55:bounded-coefficient}
\end{equation}
with $h=\gamma^{-1/3}/L$.  Under the isotropic calibration $a_t=a_s=\ell/L$,
its physical energy satisfies
\begin{equation}
 \boxed{
 E_{\rm soft}^{\rm phys}
 =\frac{-\log(\lambda_1/\Lambda)}{a_t}
 =\frac{\gamma^{-1/3}}{\ell}\,\mu_{\gamma,L}
 \longrightarrow0
 }
 \label{f16v55:soft-collapse}
\end{equation}
whenever $\gamma\to\infty$.  A finite nonzero physical limit from this branch
would require
\begin{equation}
 \mu_{\gamma,L}\asymp \gamma^{1/3},
 \label{f16v55:needed-renormalization}
\end{equation}
not a bounded coefficient such as the bare matrix-model value $e_1-e_0$.
\end{proposition}
\begin{proof}
Use $a_t=\ell/L$ in \eqref{f16v55:bounded-coefficient}.  The factors of $L$
cancel exactly, leaving \eqref{f16v55:soft-collapse}.  The last assertion is
then necessary for a finite nonzero limit.  This is a scale statement, not a
claim that the native theory actually obeys \eqref{f16v55:bounded-coefficient}
outside V54's proved domain.
\end{proof}

In the conventional coupling notation,
$\gamma^{-1/3}=4^{-1/3}g_0^{2/3}$.  The scale in
\eqref{f16v55:soft-collapse} is therefore the familiar bare constant-mode scale
$g_0^{2/3}/\ell$.  A physical small-volume zero-mode Hamiltonian is instead
expected to depend on a coupling renormalized at the box scale.  The present
calculation shows exactly where such a renormalization must enter: integrating
hard modes cannot leave the bounded soft coefficient unchanged on the physical
continuum trajectory.

\subsection{Contractible Wilson loops are blind to the constant mode at leading order}

The same conclusion is visible directly at the observable level.  Let $C$ be a
fixed oriented contractible lattice loop of length $m_C$.  On a spatially
constant configuration write
\[
 U_\mu=\operatorname{Exp}(h q_\mu),\qquad q_\mu\in\mathfrak{su}(2),
\]
and let
\[
 W_C(hq)=\frac12\operatorname{Tr}\prod_{r=1}^{m_C}
       \operatorname{Exp}(h\sigma_rq_{\mu_r}),\qquad \sigma_r\in\{\pm1\}.
\]
Closure of the loop implies
$\sum_r\sigma_r q_{\mu_r}=0$.

\begin{proposition}[Soft restriction of a fixed contractible loop]
\label{f16v55:loop-soft}
For every fixed loop $C$ and every fixed $R<\infty$ there is $C_{C,R}$ such that
\begin{equation}
 \boxed{
 \sup_{|q|\le R}|1-W_C(hq)|\le C_{C,R}h^4
 }
 \label{f16v55:loop-h4}
\end{equation}
for sufficiently small $h$.  Consequently multiplication by $W_C(hq)$
converges strongly to the identity on $L^2(\mathbb R^9)$, and for any fixed
matrix-model eigenfunctions $\phi_i,\phi_j$,
\begin{equation}
 \langle\phi_i,[W_C(hq)-1]\phi_j\rangle\longrightarrow0.
 \label{f16v55:loop-matrix}
\end{equation}
The conclusion concerns the constant soft restriction of the loop.  It is not
a statement that a renormalized continuum local field is the identity.
\end{proposition}
\begin{proof}
The product BCH expansion has no linear term because the oriented count of every
direction around a closed loop vanishes.  Thus its logarithm is
$X_C(h,q)=h^2A_C(q)+O_{C,R}(h^3)$ uniformly on $|q|\le R$.  For an imaginary
quaternion $X$,
$\frac12\operatorname{Tr}e^X=\cos|X|$ and
$0\le1-\cos|X|\le|X|^2/2$.  This proves
\eqref{f16v55:loop-h4}.  The multipliers are bounded by one in absolute value.
Strong convergence follows by first restricting to $|q|\le R$, then using the
$L^2$ tail of a fixed test as $R\to\infty$.  The matrix-element assertion is a
special case.  No moment estimate uniform in a growing eigenvalue index is
used.
\end{proof}

This proposition identifies the soft branch geometrically: fixed contractible
loops approach the same value throughout the constant-field valley, whereas
fundamental winding holonomies do not.  It is therefore natural for this branch
to control toron/electric-flux finite-volume structure rather than the local
infinite-volume glueball scale.

\subsection{The correct mass target is a local-vacuum spectral edge}

The supplied Grand Tensor already records the abstract distinction: a local
limiting gap and volume-wrapping soft vectors are compatible unless the latter
are approximable in the limiting Osterwalder--Schrader norm by bounded-support
observables.  The present calculation supplies a concrete reason to activate
that fork for the V44--V54 branch.

Let $\mathcal F=(F_1,\ldots,F_m)$ be a fixed finite family of centered,
center-neutral, gauge-invariant local writers, with synthesis $V_{\mathcal F}$.
At a physical delay $\tau$ define the entrance and delayed Grams
\begin{equation}
 G_{\mathcal F}^0=V_{\mathcal F}^*V_{\mathcal F},
 \qquad
 G_{\mathcal F}^\tau=V_{\mathcal F}^*e^{-2\tau H}V_{\mathcal F}.
 \label{f16v55:local-Grams}
\end{equation}

\begin{proposition}[Wrapping soft modes need not determine the local gap]
\label{f16v55:local-fork}
Consider any sequence of finite-volume Hamiltonians with normalized soft states
$\psi_j\perp\Omega_j$ and energies $E_j\to0$.  Suppose that for every fixed
local family $\mathcal F$,
\begin{equation}
 \|V_{\mathcal F,j}^*\psi_j\|\longrightarrow0.
 \label{f16v55:local-invisible}
\end{equation}
Then these soft states contribute vanishingly to every fixed finite delayed Gram
in \eqref{f16v55:local-Grams}.  They therefore do not obstruct a strict limiting
local Gram inequality
\begin{equation}
 G_{\mathcal F}^\tau\preceq q^2G_{\mathcal F}^0,
 \qquad q<1,
 \label{f16v55:local-gap-criterion}
\end{equation}
provided that inequality is proved uniformly on the remaining local spectral
measure.  Conversely a finite-volume full gap cannot be inferred from
\eqref{f16v55:local-gap-criterion} without a density statement for the local
core in the entire finite-volume orthogonal complement.
\end{proposition}
\begin{proof}
In the spectral resolution of a delayed Gram, the contribution of one normalized
state is the positive rank-one matrix
$e^{-2\tau E_j}(V_{\mathcal F,j}^*\psi_j)
(V_{\mathcal F,j}^*\psi_j)^*$.  Its norm tends to zero by
\eqref{f16v55:local-invisible}, independently of the fact that
$e^{-2\tau E_j}\to1$.  The same argument applies to any finite collection or a
soft spectral subspace whose projected local Gram tends to zero.  The last
assertion is the standard source-cycle/complement distinction: a contraction on
a proper cyclic subspace does not control a reducing source-orthogonal
complement.
\end{proof}

The proposition is abstract; V55 does \emph{not} prove
\eqref{f16v55:local-invisible} for a renormalized continuum local core.  The
soft restriction theorem above proves only the corresponding statement for the
constant-field restriction of fixed contractible loops.  A physical continuum
source requires its own lower-frame normalization and cutoff passage.

\subsection{V55 checkpoint: the program must now renormalize or localize}

The present line has established a precise negative conclusion and a precise
replacement target.

First, the V54 joint-regulator theorem cannot be called a continuum theorem by
choosing $a_t=h$: the original isotropic Wilson action instead fixes the bare
space--time lattice calibration, and the asymptotically free fixed-box trajectory
has $L$ exponential in $\gamma$, not polynomial.  Second, even a formal
continuation of the bounded-coefficient toron formula to that trajectory would
produce a physical energy tending to zero.  The vanishing is therefore not a
minor normalization issue that can be repaired after the fact.

There are two mathematically distinct ways forward.

\begin{enumerate}
\item \emph{Renormalized soft Hamiltonian.}  Integrate the hard modes through an
exact scale-dependent Feshbach/RG flow until the coefficient of the nine-mode
operator is expressed in a coupling at a physical scale.  Proposition
\ref{f16v55:bare-collapse} shows that a bounded bare coefficient is insufficient;
a physical fixed-box limit requires the combined dimensionless coefficient to
acquire the missing growth.  This is where the exact-return programme from the
broader literature sweep becomes structurally necessary rather than optional.

\item \emph{Local-vacuum mass.}  Treat the toron/electric-flux states as
volume-wrapping modes and prove a fixed-physical-delay contraction on a dense
center-neutral local Wilson/spin-network core.  The full finite-torus lowest
neutral eigenvalue is then not the object to keep uniformly positive.  One must
instead prove noncollapse of local source Grams and the physical-time delayed
inequality before taking the Osterwalder--Schrader limit.
\end{enumerate}

These routes are complementary: a scale-dependent Feshbach flow is a natural
way to construct the renormalized local effective operator whose local Gram
could then be tested.  The Kenig--Merle/no-collapse alternative from the
literature sweep remains a useful backstop if direct local coercivity fails.

The next substantive iteration should therefore not try to improve the
polynomial exponent $b<1/6$.  It should either begin the hard-mode RG needed to
replace $g_0$ by a physical-scale coupling, or construct the first cutoff-stable
local Gram whose toron overlap can be proved negligible after normalization.

\enlargethispage{8\baselineskip}
\paragraph{Verification and preservation.}
The accompanying diagnostic checks the normalization conversion
$\mathcal L=4E^2$, the anisotropic coefficient matching, the critical-to-physical
scale algebra, the asymptotic-scaling exponent, and closed-loop cancellation on
finite noncommuting quaternion fixtures.  These are finite checks, not a proof
of asymptotic freedom or of the inherited spectral theorems.  The audit records
the primary-source convention check and separates proved implications from the
conditional extrapolation in Proposition~\ref{f16v55:bare-collapse}.  Removing
this exact insertion and reversing the title version recovers V54 byte for byte.

\begin{firewall}
V55 proves that the current isotropic bare Wilson transfer has $a_t=a_s$ and
$\gamma=4/g_0^2$ in the standard Kogut--Susskind convention, and that the usual
fixed-physical-volume asymptotically free continuum trajectory grows
exponentially in $\gamma$, outside V54's polynomial joint domain.  It also proves
that a bounded-coefficient continuation of the quartic toron branch would have
vanishing physical energy, and that fixed contractible Wilson loops are
asymptotically blind to the constant soft coordinate.  It does not prove a
continuum local mass gap, a renormalized soft coupling, local-source noncollapse,
or toron decoupling for an already renormalized physical local algebra.  The
correct next target is an RG/Feshbach or local-vacuum theorem, not another
finite-torus lowest-gap refinement.
\end{firewall}
% END F16 V55 APPEND

% BEGIN F16 V56 APPEND -- 2026-09-18
% Insert before the final document-ending command of cumulative V55.
% No predecessor theorem, transfer normalization, Haar law, or physical clock is changed.

\clearpage
\section{V56: physical-loop toron invisibility and exact local-Gram excision}
\label{f16v56:context}

V55 shows that the finite-torus constant-mode branch is not the object that can
be identified directly with a positive continuum local mass.  The next useful
question is therefore not another estimate of its lowest eigenvalue, but whether
that branch is actually visible to the local observables from which a vacuum
Osterwalder--Schrader representation is reconstructed.  We begin that local
route here.

There are two distinct statements.  First, a contractible loop of fixed
physical size contains $O(L)$ lattice edges, so V55's fixed-word estimate does
not by itself control it.  Nevertheless the critical constant angle per link is
$h=\gamma^{-1/3}/L$, and the product of $O(L)$ such increments still has total
variation $O(\gamma^{-1/3})$.  Closure removes the linear term entirely.  We
prove a bound uniform in the lattice length of the loop and, more generally, a
background-Lipschitz estimate that remains valid in the hard Gaussian tube.

Second, if a reducing family of soft finite-volume states sees every local
writer only as an asymptotic scalar, then its contribution to every centered
entrance or delayed Gram vanishes.  This is an exact finite-dimensional matrix
statement and does not require an inverse entrance Gram.  It gives a rigorous
way to excise toron/electric-flux states from a local-Gram proof once their
local invisibility has been established.

The supplied Grand Tensor already proves that common fixed-physical-delay Gram
contractions on a dense centered local core are equivalent to a local-vacuum
mass gap, and that such inequalities pass through a singular limiting entrance
Gram.  V56 does not reprove or strengthen that abstract compiler.  Its purpose
is to pay one obstruction specific to the present toron branch and to state
precisely what remains on the physical continuum trajectory.

\subsection{Contractible loops whose lattice length grows with the regulator}

Let $L\ge8$ be even and let $C_L$ be an oriented spatial lattice loop.  Write
its successive directions as $(\sigma_r,\mu_r)$, $1\le r\le m_L$, with
$\sigma_r\in\{\pm1\}$ and $\mu_r\in\{1,2,3\}$.  We call the family
\emph{physically bounded} if, for one fixed $\kappa<\infty$,
\begin{equation}
 m_L\le \kappa L
 \label{f16v56:length}
\end{equation}
for all $L$.  If the loop is contractible without wrapping around the torus,
then its lift to $\mathbb Z^3$ closes and hence
\begin{equation}
 \sum_{r=1}^{m_L}\sigma_r e_{\mu_r}=0.
 \label{f16v56:closure}
\end{equation}
This class includes lattice approximations of any fixed rectifiable
axis-aligned loop in a box of fixed physical side length.

For $q=(q_1,q_2,q_3)\in\mathfrak{su}(2)^3$ put
\begin{equation}
 h=\frac{\gamma^{-1/3}}L,
 \qquad U_\mu(q)=\operatorname{Exp}(hq_\mu),
 \qquad
 W_{C_L}^{\rm tor}(q)=\frac12\operatorname{Tr}
 \prod_{r=1}^{m_L}\operatorname{Exp}(h\sigma_rq_{\mu_r}).
 \label{f16v56:soft-loop}
\end{equation}
The order in the product is the order of the oriented loop.

\begin{proposition}[A product-exponential remainder after exact closure]
\label{f16v56:product-lemma}
Let $X_1,\ldots,X_m$ be matrices in a normed matrix algebra, let
$A=\sum_r\|X_r\|$, and suppose $\sum_rX_r=0$.  Then
\begin{equation}
 \left\|\prod_{r=1}^m e^{X_r}-I\right\|
 \le e^A-1-A
 \le \frac12A^2e^A.
 \label{f16v56:product-remainder}
\end{equation}
If the product belongs to $\SU(2)$ and
$W=\frac12\operatorname{Tr}\prod_re^{X_r}$, then
\begin{equation}
 0\le1-W\le \frac18A^4e^{2A}.
 \label{f16v56:trace-remainder}
\end{equation}
\end{proposition}
\begin{proof}
Expand the ordered product into its absolutely convergent power series.  Its
constant term is $I$ and the total degree-one term is $\sum_rX_r=0$.  The sum
of the norms of all terms of total degree at least two is bounded by
$e^A-1-A$, proving the first inequality; the scalar Taylor remainder proves
the second.

If $U\in\SU(2)$ has eigenvalues $e^{\pm i\theta}$, then
$\|U-I\|_{\rm HS}^2=4(1-\cos\theta)=4(1-W)$ and
$\|U-I\|_{\rm HS}^2\le2\|U-I\|^2$.  Insert
\eqref{f16v56:product-remainder} to obtain
\eqref{f16v56:trace-remainder}.
\end{proof}

\begin{theorem}[Uniform soft blindness of physically bounded contractible loops]
\label{f16v56:loop-blindness}
For every physically bounded contractible family satisfying
\eqref{f16v56:length}--\eqref{f16v56:closure},
\begin{equation}
 \boxed{
 0\le1-W_{C_L}^{\rm tor}(q)
 \le \frac{\kappa^4}{8}\,\gamma^{-4/3}|q|^4
       \exp\!\left(2\kappa\gamma^{-1/3}|q|\right).
 }
 \label{f16v56:loop-uniform}
\end{equation}
In particular, for every fixed $R<\infty$,
\begin{equation}
 \sup_{L\ge8}\sup_{|q|\le R}
 |1-W_{C_L}^{\rm tor}(q)|
 \le C_{\kappa,R}\gamma^{-4/3}.
 \label{f16v56:loop-compact}
\end{equation}
The estimate remains valid when $L$ grows exponentially with $\gamma$.
\end{theorem}
\begin{proof}
Use $X_r=h\sigma_rq_{\mu_r}$.  Equation
\eqref{f16v56:closure} gives $\sum_rX_r=0$, while
\[
 A\le h m_L\max_\mu|q_\mu|
   \le \kappa\gamma^{-1/3}|q|.
\]
Lemma~\ref{f16v56:product-lemma} gives the result.  No relation between $L$
and $\gamma$ remains after $Lm_L/L$ is simplified except through the displayed
$\gamma^{-1/3}$ factor.
\end{proof}

Thus the fixed-word $O(h^4)$ statement of V55 has a physically scaled version:
when the loop length itself is $O(L)$, the relevant small parameter is $Lh$,
not $h$, and the trace deficit is still $O(\gamma^{-4/3})$ on bounded soft
sets.

\subsection{The soft coordinate is also weak on an arbitrary hard background}

The preceding fourth-order gain uses exact loop closure.  For the local-Gram
application we also need a statement that does not freeze the nonzero-momentum
coordinates.  A weaker first-order estimate is sufficient and is uniform on
the existing measured tube.

Let $s=(s_e)$ be any link field for which every link on $C_L$ lies in one fixed
principal exponential ball.  On a positively oriented link in direction $\mu$
replace $s_e$ by $s_e+hq_\mu$, and on a negatively oriented traversal use the
inverse of the corresponding group element.  Denote the resulting loop trace
by $W_{C_L}(s;q)$ and the trace at $q=0$ by $W_{C_L}(s;0)$.

\begin{proposition}[Background-Lipschitz suppression of the constant coordinate]
\label{f16v56:background}
After one fixed decrease of the principal link ball, there is a universal
$C$ such that every physically bounded loop satisfies
\begin{equation}
 \boxed{
 |W_{C_L}(s;q)-W_{C_L}(s;0)|
 \le C\kappa\gamma^{-1/3}|q|,
 }
 \label{f16v56:background-bound}
\end{equation}
uniformly in $L$, $\gamma$, and the admitted hard background $s$.
Contractibility is not required for this weaker estimate.
\end{proposition}
\begin{proof}
On a fixed exponential ball the group exponential is uniformly Lipschitz, so
changing one link angle by $hq_\mu$ changes that link group element by at most
$Ch|q|$ in the bi-invariant matrix distance.  Multiplication on either side by
unitaries preserves this distance.  Telescoping the ordered loop product over
its $m_L$ factors therefore changes the final holonomy by at most
$Cm_Lh|q|$.  The normalized trace is Lipschitz in the same matrix norm.  Use
$m_Lh\le\kappa\gamma^{-1/3}$.
\end{proof}

The point of Proposition~\ref{f16v56:background} is not its power, which is
weaker than \eqref{f16v56:loop-uniform}; it shows that the constant soft
coordinate becomes invisible even after the nonzero-momentum hard variables
are retained.  This is the estimate needed to compress a local writer to a
hard ground band without first setting the hard field to zero.

\subsection{Finite physical Wilson algebras become scalar in the soft coordinate}

Let $C_{1,L},\ldots,C_{r,L}$ be finitely many physically bounded contractible
loop families and let $P$ be a fixed complex polynomial in $r$ variables.  Put
\begin{equation}
 \Phi_L^{\rm tor}(q)=P(W_{C_{1,L}}^{\rm tor}(q),\ldots,
                         W_{C_{r,L}}^{\rm tor}(q)),
 \qquad \Phi_*=P(1,\ldots,1).
 \label{f16v56:writer}
\end{equation}
Because every normalized trace lies in $[-1,1]$, both the polynomial and its
first derivatives are bounded on the relevant compact cube.

\begin{theorem}[Strong scalarization on every tight soft family]
\label{f16v56:scalarization}
The multiplication operators obey
\begin{equation}
 M_{\Phi_L^{\rm tor}}\longrightarrow \Phi_*I
 \quad\hbox{strongly on }L^2(\mathbb R^9),
 \label{f16v56:strong-scalar}
\end{equation}
uniformly with respect to the lattice size.  More generally, let
$\mathcal S_{\gamma,L}\subset L^2(\mathbb R^9)$ be any family of closed
subspaces satisfying the uniform position-tightness condition
\begin{equation}
 \lim_{R\to\infty}\limsup_{\gamma\to\infty}
 \sup_{\substack{f\in\mathcal S_{\gamma,L(\gamma)}\\\|f\|=1}}
 \int_{|q|>R}|f(q)|^2dq=0.
 \label{f16v56:tight-family}
\end{equation}
Then
\begin{equation}
 \boxed{
 \left\|(M_{\Phi_L^{\rm tor}}-\Phi_*I)
       P_{\mathcal S_{\gamma,L}}\right\|\longrightarrow0.
 }
 \label{f16v56:uniform-scalarization}
\end{equation}
The same conclusion holds for any finite family of such writers.
\end{theorem}
\begin{proof}
On every fixed $|q|\le R$, Theorem~\ref{f16v56:loop-blindness} and the
mean-value theorem for $P$ give uniform convergence to $\Phi_*$.  The
multipliers are uniformly bounded on all of $\mathbb R^9$.  This proves strong
convergence by the usual compact-ball/tail decomposition.  Under
\eqref{f16v56:tight-family}, the same decomposition is uniform over the unit
ball of the declared subspaces, proving the operator-norm statement.
\end{proof}

For a fixed finite-dimensional subspace no separate tightness proof is needed:
its unit sphere is compact and every $L^2$ function has vanishing tails.  In
particular every fixed finite spectral subspace of the quartic matrix model is
covered.  The theorem is stronger in form because it also applies to a moving
family once position tightness has been proved independently.

\subsection{Hard-ground compression on the existing critical branch}

We record how the previous geometric estimates interface with the already
proved V54 branch.  This subsection does not extend V54 to the physical
exponential trajectory.

Let $\varphi_L(u)$ be V54's normalized hard Gaussian ground function in the
measured extraction and let $\mathcal I_L f=\varphi_L\otimes f$.  For a local
writer $\Phi_L$ built from finitely many physically bounded Wilson loops, define
its hard-ground compressed symbol on the measured tube by
\begin{equation}
 m_{\gamma,L}(q)=
 \int |\varphi_L(u)|^2
       \Phi_L(\varepsilon u,hq)\,du,
 \label{f16v56:compressed-symbol}
\end{equation}
with the same tube cutoff convention as V54; the Gaussian mass omitted by the
fixed tube is retained as the existing exponentially/superalgebraically small
tail.

\begin{proposition}[The compressed local writer is asymptotically scalar in $q$]
\label{f16v56:hard-compression}
Along V54's joint domain, for every fixed $R$,
\begin{equation}
 \sup_{|q|\le R}|m_{\gamma,L}(q)-m_{\gamma,L}(0)|
 \le C_{\Phi,R}\gamma^{-1/3}+o(1),
 \label{f16v56:compressed-lipschitz}
\end{equation}
where the $o(1)$ is the already paid hard-tube tail.  Consequently, on every
fixed finite model spectral subspace $P_J^{\rm mat}$,
\begin{equation}
 \left\|P_J^{\rm mat}
      [M_{m_{\gamma,L}}-m_{\gamma,L}(0)I]
       P_J^{\rm mat}\right\|\longrightarrow0.
 \label{f16v56:model-compression}
\end{equation}
V54's norm-resolvent and spectral-projection convergence then make $\Phi_L$
an asymptotic scalar on each corresponding finite toron band.  V53 transfers
the same statement to the original Wilson transfer on V54's joint domain.
\end{proposition}
\begin{proof}
Inside the measured tube, Proposition~\ref{f16v56:background} applies pointwise
in $u$ and gives the bound $C\gamma^{-1/3}|q|$ for each loop.  The polynomial
writer obeys the same estimate with a writer-dependent constant.  Integrate
against $|\varphi_L|^2$.  V54's hard localization and V50's tube estimate pay
the omitted hard tail on the stated joint domain.  This proves
\eqref{f16v56:compressed-lipschitz}.  Fixed model spectral subspaces are tight
in $q$, so the compact-ball/tail argument of
Theorem~\ref{f16v56:scalarization} proves
\eqref{f16v56:model-compression}.  V54 gives norm convergence of isolated
finite spectral projections after extraction; conjugating the compressed
multiplication estimate by those projectors proves the final statement.  V53's
actual-top normalized Wilson/reference comparison is $o(h)$ there, so it does
not alter a fixed finite cluster.
\end{proof}

The conclusion is deliberately local to the already controlled branch.  It
establishes the mechanism that V55 anticipated: within the toron critical
sector, a physically bounded contractible local writer acts asymptotically as
a scalar.  It does not assert that V54 controls the physical continuum
trajectory.

\subsection{Exact toron excision from centered local Grams}

We now isolate the linear-algebra statement needed for the local-vacuum route.
It is independent of the special form of Yang--Mills once the scalar-compression
hypothesis has been verified.

Let $H_j\ge0$ have normalized vacuum $\Omega_j$, and let $P_j$ be a reducing
orthogonal projection with $P_j\Omega_j=\Omega_j$.  Put
\[
 P_j^\circ=P_j-|\Omega_j\rangle\langle\Omega_j|,
 \qquad Q_j=I-P_j^\circ.
\]
Let $F_{1,j},\ldots,F_{m,j}$ be bounded local writers and define their centered
source synthesis
\begin{equation}
 V_jc=\sum_{\alpha=1}^m c_\alpha
 \left(F_{\alpha,j}-\langle\Omega_j,F_{\alpha,j}\Omega_j\rangle\right)
 \Omega_j.
 \label{f16v56:synthesis}
\end{equation}

\begin{theorem}[Scalar toron compression implies Gram invisibility]
\label{f16v56:Gram-invisibility}
Suppose there are scalars $c_{\alpha,j}$ such that
\begin{equation}
 \epsilon_j:=\max_{\alpha\le m}
 \|P_j(F_{\alpha,j}-c_{\alpha,j}I)P_j\|\longrightarrow0.
 \label{f16v56:scalar-compression-hyp}
\end{equation}
Then
\begin{equation}
 \boxed{
 \|P_j^\circ V_j\|\le2\sqrt m\,\epsilon_j,
 }
 \label{f16v56:source-overlap}
\end{equation}
and, for every $t\ge0$,
\begin{equation}
 \boxed{
 0\preceq
 V_j^*P_j^\circ e^{-2tH_j}P_j^\circ V_j
 \preceq4m\epsilon_j^2 I_m.
 }
 \label{f16v56:toron-Gram}
\end{equation}
Thus the contribution of the whole reducing toron band $P_j^\circ$ to every
fixed finite entrance or delayed Gram vanishes uniformly in the delay.
\end{theorem}
\begin{proof}
Because $P_j\Omega_j=\Omega_j$,
\[
 |\langle\Omega_j,F_{\alpha,j}\Omega_j\rangle-c_{\alpha,j}|
 \le\epsilon_j.
\]
Hence
\[
 \|P_j^\circ(F_{\alpha,j}
 -\langle\Omega_j,F_{\alpha,j}\Omega_j\rangle)\Omega_j\|
 \le2\epsilon_j.
\]
The synthesis bound follows by Cauchy--Schwarz in the coefficient space.  Since
$P_j^\circ$ reduces $H_j$ and $e^{-2tH_j}$ is a positive contraction,
its delayed Gram is positive and bounded above by
$(P_j^\circ V_j)^*(P_j^\circ V_j)$, proving
\eqref{f16v56:toron-Gram}.
\end{proof}

\begin{proposition}[Quotient-stable local contraction after toron excision]
\label{f16v56:Gram-excision}
Under the hypotheses of Theorem~\ref{f16v56:Gram-invisibility}, suppose that for
one delay $t_j\ge0$
\begin{equation}
 V_j^*Q_j e^{-2t_jH_j}Q_jV_j
 \preceq q_j^2 V_j^*Q_jV_j+E_j,
 \qquad 0\le q_j<1.
 \label{f16v56:complement-contraction}
\end{equation}
Then the full Grams satisfy
\begin{equation}
 \boxed{
 G_j^{t_j}\preceq q_j^2G_j^0+E_j
 +(1-q_j^2)4m\epsilon_j^2I_m,
 }
 \label{f16v56:full-Gram}
\end{equation}
where $G_j^0=V_j^*V_j$ and
$G_j^{t_j}=V_j^*e^{-2t_jH_j}V_j$.
\end{proposition}
\begin{proof}
The reducing decomposition gives
\[
 G_j^{t_j}
 =V_j^*Q_je^{-2t_jH_j}Q_jV_j
 +V_j^*P_j^\circ e^{-2t_jH_j}P_j^\circ V_j,
\]
and likewise
$G_j^0=V_j^*Q_jV_j+V_j^*P_j^\circ V_j$.
Use \eqref{f16v56:complement-contraction} and then
$P_j^\circ e^{-2t_jH_j}P_j^\circ\preceq P_j^\circ$ to obtain
\[
 G_j^{t_j}\preceq q_j^2G_j^0+E_j
 +(1-q_j^2)V_j^*P_j^\circ V_j.
\]
Apply \eqref{f16v56:source-overlap}.
\end{proof}

This proposition is the exact handoff to the supplied fixed-physical-delay
local Gram criterion.  No inverse entrance Gram is required, so the argument
survives source directions that become null in the continuum quotient.  What
must still be proved is the complement contraction
\eqref{f16v56:complement-contraction} on the physical continuum trajectory,
together with convergence/noncollapse of the local source family itself.

\subsection{What V56 changes in the programme}

V55 identified two routes: renormalize the toron Hamiltonian, or prove a local
vacuum gap that is insensitive to wrapping soft modes.  V56 supplies a rigorous
piece of the second route.

For physically bounded contractible Wilson writers, the constant soft
coordinate is invisible uniformly in the number of lattice links in the loop.
On the existing critical branch, the hard-ground compressed writer becomes an
asymptotic scalar on every fixed toron spectral band.  The abstract Gram theorem
then shows that any reducing band with this scalar-compression property can be
removed from the centered local Gram at a vanishing matrix cost.

This does \emph{not} prove the continuum mass gap.  Three obligations remain:

\begin{enumerate}
\item extend the scalar-compression/local-invisibility estimate to the relevant
physical continuum local source normalization on the exponential
$L(\gamma)$ trajectory, rather than only to the V54 toron band;
\item prove a fixed-physical-delay contraction on the complementary local
spectral measure for \emph{every finite mixed local family}, not only diagonal
loop probes;
\item pass the entrance and delayed Grams to one nontrivial
reflection-positive OS vacuum representation, with the local core dense after
the quotient.
\end{enumerate}

The first item is now a local source problem, not a demand for a uniform positive
finite-torus full gap.  The second is where a block/Feshbach, approximate
factorization, or Kenig--Merle rigidity argument can enter.  The third is the
noncollapse requirement already isolated in the supplied Grand Tensor.

Accordingly the next iteration should attack the \emph{complement local Gram}
(or construct it via an exact local Feshbach flow), rather than return to the
wrapping toron eigenvalue.  A physical local mass can coexist with the soft
finite-volume sector only if the scalar-compression mechanism survives the
actual continuum source normalization; V56 makes that statement operational.

\paragraph{Verification and preservation.}
The accompanying diagnostic checks the ordered-product remainder, the exact
$\SU(2)$ trace identity, physical-loop scaling $Lh=\gamma^{-1/3}$, polynomial
writer scalarization on finite fixtures, and the reducing-Gram decomposition
with nonzero toron energies.  These checks do not certify the inherited V54
norm-resolvent theorem or a continuum local contraction.  Removing this exact
appendix and reversing the literal title version recovers V55 byte for byte.
The next scheduled companion checkpoint is V60.

\begin{firewall}
V56 proves that contractible Wilson loops whose lattice length is $O(L)$ are
uniformly blind to the constant critical soft coordinate as $\gamma\to\infty$,
and gives a hard-background Lipschitz version that is uniform in lattice size.
It promotes this to scalar compression on tight soft families and, on the
already controlled V54 branch, to finite toron spectral bands.  It also proves
an exact toron-excision lemma for centered delayed local Grams.  It does not
prove the complementary local contraction, continuum source noncollapse,
Osterwalder--Schrader convergence, a renormalized toron Hamiltonian, or the
interacting continuum Yang--Mills mass gap.
\end{firewall}
% END F16 V56 APPEND

% BEGIN F16 V57 APPEND -- 2026-09-18
% Append-only continuation of the attached post-fifteenth-archive V56.
% The predecessor body is unchanged; only its displayed title version is advanced.

\clearpage
\section{V57: normalized local sources, nonreducing returns, and a native resolvent certificate}
\label{f16v57:context}

V56 changes the object of the programme from a wrapping finite-torus level to
mixed local vacuum Grams at a fixed physical delay.  Its two useful inputs are
geometric suppression of a constant soft coordinate and exact excision of a
\emph{reducing} soft band.  Neither input, by itself, controls a local writer
after a diverging field normalization, or a conditional coarse projection that
does not reduce the transfer.  This section addresses precisely those two
interfaces and then replaces an assumed complementary delayed-Gram estimate by
a finite native resolvent certificate.

There are three substantive conclusions.  First, the invariant excision
quantity is a \emph{relative} soft Gram, not the absolute norm of the unnormalized
writer.  An exact constant-field rectangle calculation shows that a vanishing
Wilson-loop fluctuation can retain a nonzero toron component after variance
normalization.  Second, a nonreducing elimination transports the source as well
as the operator: its correct soft remainder contains an explicit return term.
Third, completing the native resolvent square gives a source-dependent matrix
upper bound using a trial bank and its zero-, one-, and two-step vacuum
correlations.  No Gaussian comparison, unknown logarithmic-operator domain, or
change of physical clock enters that bound.  A polynomial construction gives an
explicit cutoff-dependent approximation rate for the trial bank.

\subsection{Dependencies and the scope of the continuation}

The predecessor is the attached post-fifteenth-archive V56, not the distinct
V56 embedded in either hundred-version archive.  The parallel V67 shares an
earlier part of this lineage but takes a different terminal route.  Its toron
coordinate estimates and cubic-vertex programme are not rederived here.
Likewise, f15 V65 already studies a normalized damped magnetic probe, and f15
V79--V81 already provide whole-transfer Haar and spectral-tail certificates.
The present construction differs in its target: \emph{the mixed local source
Gram at the original physical clock}, with only source-marked approximation
errors paid.  It does not claim to remove the need for native vacuum data.

We use the inherited finite-regulator positivity of the Wilson transfer and
its actual Perron normalization.  The V54 spectral theorem is used only when
explicitly discussing that already delimited branch; it is not extrapolated to
an exponential continuum trajectory.  The Schur-complement factorization and
Chebyshev residual polynomial are classical tools, not new general discoveries.
For attribution see G.~Dusson, I.~M.~Sigal and B.~Stamm,
\href{https://arxiv.org/abs/2105.02058}{\emph{The Feshbach--Schur map and perturbation theory}}
(2021), and G.~H.~Golub and R.~S.~Varga,
\href{https://www.math.kent.edu/~varga/pub.html}{\emph{Chebyshev semi-iterative methods,
successive overrelaxation iterative methods, and second order Richardson
iterative methods}, Parts I--II}, Numerische Mathematik \textbf{3} (1961),
147--168.  All identities and estimates needed below are proved explicitly.
These references supply no missing Yang--Mills hypothesis.

\subsection{The intrinsic soft fraction and the normalization boundary}

Let $V:\mathbb C^m\to\mathcal H_0$ be a finite source synthesis into a
vacuum-orthogonal Hilbert space, and let $P$ be an orthogonal projection on
$\mathcal H_0$.  It need not reduce any operator.  Set
\begin{equation}
 G=V^*V,\qquad S=V^*PV,\qquad E_G=\text{orthogonal projection onto }\operatorname{ran}G.
 \label{f16v57:relative-Grams}
\end{equation}
The positive square root of the Moore--Penrose inverse is denoted by
$G^{\dagger/2}$.  Every use of this inverse is finite dimensional and is on the
\emph{exact} range; numerical thresholding is not part of the definition.

\begin{proposition}[Normalization-invariant soft fraction]
\label{f16v57:relative-fraction}
The number
\begin{equation}
 \eta^2:=\inf\{c\ge0:S\preceq cG\}
   =\|PVG^{\dagger/2}\|^2
   =\|G^{\dagger/2}SG^{\dagger/2}\|
 \label{f16v57:eta}
\end{equation}
belongs to $[0,1]$, with $\eta=0$ when $V=0$.  Under any coefficient map
$Z:\mathbb C^r\to\mathbb C^m$, including a regulator-dependent unbounded
normalization, the inequality $S\preceq\eta^2G$ passes exactly to $VZ$ by
congruence.  If $Z$ is invertible, the optimal constant is unchanged.

In contrast, $\|PV_j\|\to0$ alone does not imply $\eta_j\to0$.  A sufficient
additional estimate is $G_j\succeq g_j E_{G_j}$ and
$\|PV_j\|^2/g_j\to0$.  In the notation of V56, its bound gives the sufficient
condition $4m\epsilon_j^2/g_j\to0$; it does not supply $g_j$.
\end{proposition}
\begin{proof}
Since $P$ is a projection, $0\preceq S\preceq G$ and
$\ker G=\ker V\subseteq\ker S$.  Restriction to $\operatorname{ran}G$ and
conjugation by its invertible square root prove the equalities in
\eqref{f16v57:eta}.  The identity $VE_G=V$ permits extension by zero to the
coefficient kernel.  Congruence proves the assertion for $Z$, and applying
$Z^{-1}$ proves equality of the optimal constants when it exists.  Finally,
on the range of $G_j$, $S_j\preceq\|PV_j\|^2 I\preceq
(\|PV_j\|^2/g_j)G_j$.  The example $V_j=\varepsilon_j v$, with $Pv=v$ and
$\|v\|=1$, has $\|PV_j\|\to0$ but $S_j=G_j=\varepsilon_j^2$ and $\eta_j=1$.
\end{proof}

For a reducing soft projection, V56's exact decomposition immediately gives
\begin{equation}
 G^{\tau}\preceq
 \bigl[q^2+(1-q^2)\eta^2\bigr]G+E
 \quad\text{whenever}\quad S\preceq\eta^2G.
 \label{f16v57:relative-excision}
\end{equation}
This is a normalization-aware use of the preceding theorem, not a new proof of
the local-gap compiler.  If $E\preceq\delta G$, the displayed contraction factor
is increased by $\delta$.  An error small only in an ambient coefficient norm
cannot silently be declared small relative to a collapsing entrance Gram.

\subsection{A physical-size rectangle that survives fluctuation normalization}

We can make the normalization issue concrete using the same constant-field
geometry as V56.  Identify an imaginary quaternion $q_\mu$ with a vector in
$\mathbb R^3$, so that $e^{tq_\mu}=\cos(t|q_\mu|)+
\widehat q_\mu\sin(t|q_\mu|)$.  Consider a contractible rectangle with $n_L$
links in direction $\mu$ and $m_L$ links in direction $\nu$, in a region where it
does not wrap around the spatial torus.  Write
\begin{equation}
 h=\gamma^{-1/3}/L,\qquad r_L=n_L/L,\qquad u_L=m_L/L,
 \qquad r_L\longrightarrow r>0,\quad u_L\longrightarrow u>0.
 \label{f16v57:rectangle-scales}
\end{equation}
The limits may be taken along any joint sequence for which the rectangle exists;
no polynomial relation between $L$ and $\gamma$ is needed.

\begin{theorem}[Exact rectangle coefficient and a nonzero normalized toron source]
\label{f16v57:rectangle}
For the constant-field restriction of the normalized Wilson trace,
\begin{align}
 1-W_{\gamma,L}^{\rm tor}(q)
 &=2\sin^2(r_L\gamma^{-1/3}|q_\mu|)
       \sin^2(u_L\gamma^{-1/3}|q_\nu|)
       \frac{|q_\mu\times q_\nu|^2}{|q_\mu|^2|q_\nu|^2},
 \label{f16v57:rectangle-exact}
\end{align}
with the continuous value when either denominator vanishes.  Therefore
\begin{equation}
 F_{\gamma,L}(q):=\gamma^{4/3}(1-W_{\gamma,L}^{\rm tor}(q))
 \longrightarrow f(q):=2r^2u^2|q_\mu\times q_\nu|^2,
 \qquad 0\le F_{\gamma,L}(q)\le C|q|^4.
 \label{f16v57:rectangle-limit}
\end{equation}
Let $\mu$ be a fixed probability measure on the soft coordinate space with
$\int |q|^8\,d\mu<\infty$, and suppose $\operatorname{Var}_\mu(f)>0$.
Then convergence holds in $L^2(\mu)$, and
\begin{align}
 \gamma^{8/3}\operatorname{Var}_\mu(W_{\gamma,L}^{\rm tor})
 &\longrightarrow\operatorname{Var}_\mu(f)>0,\label{f16v57:rectangle-variance}\\
 \frac{1-W_{\gamma,L}^{\rm tor}
      -\mathbb E_\mu(1-W_{\gamma,L}^{\rm tor})}
      {\sqrt{\operatorname{Var}_\mu(W_{\gamma,L}^{\rm tor})}}
 &\longrightarrow
 \frac{f-\mathbb E_\mu f}{\sqrt{\operatorname{Var}_\mu(f)}}
 \quad\text{in }L^2(\mu).
 \label{f16v57:rectangle-normalized}
\end{align}
In particular, scalarization of the bounded loop need not survive
fluctuation normalization.
\end{theorem}
\begin{proof}
Put $A=e^{n_Lh q_\mu}$ and $B=e^{m_Lh q_\nu}$.  Conjugation by $A$ rotates the
imaginary part of $B$ through angle $2n_Lh|q_\mu|$ about $\widehat q_\mu$.
If $a=n_Lh|q_\mu|$ and $b=m_Lh|q_\nu|$, the scalar part of $ABA^{-1}B^{-1}$ is
\[
 \cos^2 b+\sin^2 b\bigl[\cos(2a)
       +(1-\cos(2a))(\widehat q_\mu\cdot\widehat q_\nu)^2\bigr].
\]
Subtracting from one proves \eqref{f16v57:rectangle-exact}.  The bound
$|\sin x|\le|x|$ and boundedness of $r_L,u_L$ give the global polynomial
majorant in \eqref{f16v57:rectangle-limit}.  Pointwise convergence follows from
$\sin x/x\to1$.  The eighth-moment hypothesis now gives $L^2$ convergence by
dominated convergence.  Centering is a continuous operation in $L^2(\mu)$, so
both its norm and the normalized nonzero vector converge.  Multiplication by
$\gamma^{-4/3}$ proves the variance formula.
\end{proof}

For example, $\operatorname{Var}_\mu(f)>0$ whenever $\mu$ has a strictly positive
density on a nonempty open set: the polynomial $|q_\mu\times q_\nu|^2$ is not
constant on an open set.  We have not established the eighth-moment and source
transport assumptions for the native vacuum on the physical continuum
trajectory.  The theorem is exact for the stated constant-field restriction,
not a claim that $\gamma^{4/3}$ is the physical renormalization of a Wilson
operator.

There is also an exact dynamical illustration.  Let $K\ge0$ be self-adjoint on
$L^2(\mu)$, with vacuum $1$, and let $c_j\downarrow0$.  For the normalized
centered vectors $v_j$ in \eqref{f16v57:rectangle-normalized} and
$H_j=c_jK$,
\begin{equation}
 \langle v_j,e^{-2\tau H_j}v_j\rangle\longrightarrow1
 \qquad(\tau>0).
 \label{f16v57:auxiliary-soft-delay}
\end{equation}
Indeed $v_j\to v$ in norm and $e^{-2\tau c_jK}\to I$ strongly by the spectral
theorem.  This is an auxiliary soft-sector example, not an identification of
the native Wilson Hamiltonian.  It shows exactly why an absolute vanishing
Gram is insufficient to discard a normalized soft source.

\subsection{A dependency audit: cubic Gaussian chaos is not Gaussian}

The attached parallel V67 motivates connected rather than whole-density
comparisons.  Its proof of the displayed log-normal lower bound, however,
identifies the cubic vertex sum with a centered Gaussian random variable.
Wick ordering places it in the third Gaussian chaos; it does not make its law
Gaussian.  The log-normal moment formula used in that proof is consequently
not an inherited input here.  Nor does a second-moment computation alone repair
that particular step.

\begin{proposition}[An uncut cubic exponential has no nonzero real moment]
\label{f16v57:cubic-audit}
Let $X$ be a centered finite-dimensional Gaussian random vector, and let $p$ be a real
polynomial whose degree-three homogeneous part is nonzero on the linear
support of its covariance.  Assume this is the highest nonzero degree on that
support.  Then, for every real $t\ne0$,
\begin{equation}
 \mathbb E\,e^{t p(X)}=+\infty.
 \label{f16v57:cubic-mgf}
\end{equation}
In particular a nonzero third-chaos cubic cannot be assigned Gaussian
exponential moments.  On a bounded chart its exponential is integrable, but its
law is a truncated non-Gaussian law, so the Gaussian log-normal formula still
does not apply.
\end{proposition}
\begin{proof}
Restrict to the Gaussian support, translate its mean and whiten its covariance.
The leading homogeneous cubic is still nonzero.  Because it is odd, there is a
unit vector $\omega$ at which its product with $t$ is positive.  On a sufficiently
small open cone of directions around $\omega$ that product is at least some
$c>0$.  For $x=\rho\theta$ in that cone and all sufficiently large $\rho$,
$t p(x)\ge(c/2)\rho^3$.  The Gaussian density has a quadratic negative exponent,
so its product with $e^{t p(x)}$ has a radial lower bound proportional to
$\exp((c/2)\rho^3-C\rho^2-C\rho)\rho^{d-1}$ on a cone of positive solid angle.
Its integral diverges.  The one-dimensional case uses a half-line.  This proves
\eqref{f16v57:cubic-mgf}.
\end{proof}

A minimal diagnostic is $p(X)=X_1X_2X_3$ for three independent standard normal
variables.  Its variance is $1$ and its fourth moment is $27$, rather than the
Gaussian value $3$.  The proposition is not a divergence statement about the
\emph{full compact} Wilson law: its higher-order action, compact domain and
normalization remain present.  It instead prohibits exponentiating a cubic
truncation on an unbounded Gaussian space as though it were the complete law.
The native construction below does neither.  The other parallel V67 claims are
not certified or refuted by this limited dependency audit.

\subsection{An exact physical-clock resolvent of the native transfer}

At a fixed regulator let $\mathsf T_a$ be the positive physical Wilson transfer,
$\lambda_{0,a}>0$ its actual top eigenvalue, and $\Omega_a$ its normalized vacuum.
The subscript $a$ includes the associated coupling and spatial volume.  Put
\begin{equation}
 \mathcal T_a=\mathsf T_a/\lambda_{0,a},\qquad
 0\preceq\mathcal T_a\preceq I,\qquad
 H_a=-a^{-1}\log\mathcal T_a,\qquad
 K_a=a^{-1}(I-\mathcal T_a).
 \label{f16v57:native-K}
\end{equation}
The number $a$ is the physical time spacing inherited from V55, not a fitted
soft spectral scale.  The bounded operator $K_a$, with
$0\preceq K_a\preceq a^{-1}I$, is an auxiliary \emph{discrete-step} generator;
it is not identified with $H_a$.  Working with $K_a$ avoids any assumption that
a coarse projection preserves the domain of $\log\mathcal T_a$.

All following operators may be restricted to
$\mathcal H_{0,a}=\{\Omega_a\}^{\perp}$.  A finite centered local family has
synthesis $V_a$, entrance Gram $G_a=V_a^*V_a$, and exact delayed Gram
\begin{equation}
 G_a^{[n]}=V_a^*\mathcal T_a^{2n}V_a
         =V_a^*e^{-2naH_a}V_a.
 \label{f16v57:native-delay}
\end{equation}
Sources may depend on the regulator and include a declared renormalization;
each column only has to be a Hilbert vector at that regulator.

\begin{proposition}[Native resolvent domination at a physical delay]
\label{f16v57:resolvent-domination}
For $s>0$ and an integer $n\ge1$ such that $2na\ge s^{-1}$,
\begin{equation}
 \boxed{\quad
 G_a^{[n]}\preceq sV_a^*(K_a+sI)^{-1}V_a.
 \quad}
 \label{f16v57:resolvent-domination-eq}
\end{equation}
In particular, for a fixed $\tau>0$, take $n=\lceil\tau/a\rceil$ and
$s=(2\tau)^{-1}$ without changing the native clock.
\end{proposition}
\begin{proof}
For $0\le x\le a^{-1}$,
\[
 (1-ax)^{2n}\le e^{-2nax}\le e^{-x/s}\le\frac{s}{s+x}.
\]
The endpoint $ax=1$ is immediate; elsewhere the first inequality follows from
$\log(1-ax)\le-ax$.  The last follows from $e^y\ge1+y$.  Functional calculus
for the bounded self-adjoint $K_a$ and congruence by $V_a$ prove the assertion.
\end{proof}

No spectral gap, continuum limit, or perturbative expansion is assumed in this
proposition.  Positivity of $\mathcal T_a$ as a Hilbert-space operator is used;
no positivity-preserving assertion about its fractional powers is needed.

\subsection{A nonreducing coarse projection transports the source}

Suppress $a$ temporarily and choose any orthogonal projection $P$ on
$\mathcal H_0$, with $Q=I-P$.  Define the bounded blocks
\begin{equation}
 D_s=Q(K+sI)Q\big|_{Q\mathcal H_0},\quad
 B=PKQ,\quad
 C_s=P(K+sI)P-BD_s^{-1}B^*\big|_{P\mathcal H_0}.
 \label{f16v57:Schur-blocks}
\end{equation}
The factors $D_s^{-1}$ and $C_s^{-1}$ exist for $s>0$; no gap in the $P$ sector is
required.  Set
\begin{equation}
 r_s=PV-BD_s^{-1}QV,\qquad
 \Gamma_s=r_s^*C_s^{-1}r_s.
 \label{f16v57:source-return}
\end{equation}

\begin{theorem}[Exact local-source Feshbach identity]
\label{f16v57:Feshbach}
For every $s>0$,
\begin{gather}
 D_s\succeq sQ,\qquad C_s\succeq sP,\label{f16v57:Schur-floors}\\
 \boxed{\quad
 V^*(K+sI)^{-1}V
 =(QV)^*D_s^{-1}QV+\Gamma_s.
 \quad}\label{f16v57:source-Schur-identity}
\end{gather}
Thus even if $PV=0$, the soft source after elimination is generally
$-BD_s^{-1}QV$, not zero.

If $QKQ\succeq\Delta Q$ for some $\Delta>0$, and if
$s\Gamma_s\preceq\theta G$, then the physical-delay bound is
\begin{equation}
 G^{[n]}\preceq
 \left(\frac{s}{s+\Delta}+\theta\right)G
 \qquad(2na\ge s^{-1}).
 \label{f16v57:Feshbach-contraction}
\end{equation}
A sufficient source-only return estimate is
\begin{equation}
 (PV)^*PV\preceq\eta^2G,\qquad
 (BD_s^{-1}QV)^*(BD_s^{-1}QV)\preceq\zeta^2G,
 \label{f16v57:marked-return-bound}
\end{equation}
which permits $\theta=(\eta+\zeta)^2$.  No bound on the full norm
$\|BD_s^{-1}\|$ is imposed.
\end{theorem}
\begin{proof}
The floor for $D_s$ is immediate.  For $p\in P\mathcal H_0$ complete the
quadratic form in $q\in Q\mathcal H_0$:
\[
 \langle p,C_sp\rangle
 =\inf_q\langle p+q,(K+sI)(p+q)\rangle
 \ge s\|p\|^2.
\]
This also proves bounded invertibility of $C_s$.  The exact factorization
\[
 K+sI=
 \begin{pmatrix}I&BD_s^{-1}\\0&I\end{pmatrix}
 \begin{pmatrix}C_s&0\\0&D_s\end{pmatrix}
 \begin{pmatrix}I&0\\D_s^{-1}B^*&I\end{pmatrix}
\]
then gives \eqref{f16v57:source-Schur-identity} by inversion and evaluation on
$(PV,QV)$.  If the hard compression has the stated floor,
$(QV)^*D_s^{-1}QV\preceq(s+\Delta)^{-1}G$.
Proposition~\ref{f16v57:resolvent-domination} proves
\eqref{f16v57:Feshbach-contraction}.  Under
\eqref{f16v57:marked-return-bound}, for every coefficient vector $c$,
\[
 \|r_sc\|\le(\eta+\zeta)\sqrt{c^*Gc},\qquad
 s\,c^*\Gamma_sc\le\|r_sc\|^2,
\]
by \eqref{f16v57:Schur-floors}.  This proves the final assertion.
\end{proof}

A conditional coarse space can be defined with the \emph{native} vacuum law,
not the frozen hard Gaussian.  In the ground-state representation the law is
$\pi_a(dU)=|\Omega_a(U)|^2\,dm(U)$, with $m$ product Haar measure on the physical
configuration space, and the vacuum is the constant function $1$.  For a
specified gauge-invariant coarse sigma-algebra $\mathcal A_a$, let
$\mathbb E_{\pi_a}[\cdot\mid\mathcal A_a]$ be the orthogonal conditional
expectation and restrict it to the centered subspace.  Its soft source Gram is
exactly
\begin{equation}
 S_{\alpha\beta}
 =\operatorname{Cov}_{\pi_a}\left(
 \mathbb E_{\pi_a}[F_\alpha\mid\mathcal A_a],
 \mathbb E_{\pi_a}[F_\beta\mid\mathcal A_a]\right),
 \label{f16v57:conditional-soft-Gram}
\end{equation}
where the first entry of covariance is conjugated.  This follows from the
$L^2$ orthogonality of conditional expectation and gives an exact matrix law of
total covariance.  It does not imply that the projection reduces $K_a$; the
return $r_s$ is still required.  Choosing and controlling a globally defined
coarse sigma-algebra that represents the desired toron sector is an additional
native task, not an assumption hidden in this formula.

There is a useful history form of the return.  On $Q\mathcal H_0$ write
$T_Q=Q\mathcal T_aQ$.  Since $B=-a^{-1}P\mathcal T_aQ$,
\begin{equation}
 \boxed{\quad
 r_s=PV+\sum_{k=0}^{\infty}(1+as)^{-k-1}
       P\mathcal T_aQ\,T_Q^kQV.
 \quad}
 \label{f16v57:discounted-return}
\end{equation}
Indeed $D_s=a^{-1}[(1+as)Q-T_Q]$, and its inverse has the displayed
norm-convergent geometric series because $\|T_Q\|\le1<1+as$.
These are projected, source-marked returns, not positive first-hitting
probabilities: $Q$ can introduce cancellations.  Termwise absolute norms may
lose the very cancellation that the exact remainder retains.

The full hard-compression floor in Theorem~\ref{f16v57:Feshbach} is a sufficient
input, not a proved Yang--Mills bound.  Alternatively one can bound its exact
source-compressed first term directly.  In either formulation, a small bare
$PV$ without the return term is not a nonreducing excision theorem.

\subsection{An a posteriori resolvent bound without an exact elimination}

The preceding Schur identity is useful conceptually, but evaluating an exact
inverse in either block is not a prerequisite for a rigorous upper bound.  Let
$W:\mathbb C^m\to\mathcal H_0$ be any finite trial synthesis and define
\begin{align}
 R_s(W)&=V-(K+sI)W,\label{f16v57:trial-residual}\\
 \mathcal U_s(W)
 &=s\bigl(V^*W+W^*V-W^*(K+sI)W\bigr)+R_s(W)^*R_s(W).
 \label{f16v57:trial-upper}
\end{align}
All operators are bounded at the fixed regulator, so no trial-domain condition
is missing.

\begin{theorem}[Native source-residual certificate]
\label{f16v57:residual-certificate}
For every trial $W$,
\begin{align}
 \mathcal U_s(W)-sV^*(K+sI)^{-1}V
 &=R_s(W)^*\bigl[I-s(K+sI)^{-1}\bigr]R_s(W)
 \succeq0.\label{f16v57:residual-exact-error}
\end{align}
Consequently, if $2na\ge s^{-1}$ and
$\mathcal U_s(W)\preceq q^2G$ with $0\le q<1$, then
\begin{equation}
 \boxed{\qquad G^{[n]}\preceq q^2G.\qquad}
 \label{f16v57:residual-contraction}
\end{equation}
The certificate is an exact matrix inequality on mixed sources and remains
valid when $G$ is singular.
\end{theorem}
\begin{proof}
Writing $A=K+sI$ and $R=V-AW$ gives the exact completed square
\[
 V^*A^{-1}V
 =V^*W+W^*V-W^*AW+R^*A^{-1}R.
\]
Subtract $s$ times this identity from \eqref{f16v57:trial-upper}.  Because
$A\succeq sI$, the operator $I-sA^{-1}$ is positive.  This proves the first
claim, and Proposition~\ref{f16v57:resolvent-domination} proves the second.
If $Gc=0$, then $Vc=0$ and the same inequalities still hold; no division by an
entrance eigenvalue has been made.
\end{proof}

The hard trial $W=QD_s^{-1}QV$ has residual precisely $r_s$ in the $P$ space.
For this choice the bound reduces to
\begin{equation}
 \mathcal U_s(QD_s^{-1}QV)
 =s(QV)^*D_s^{-1}QV+r_s^*r_s,
 \label{f16v57:hard-trial}
\end{equation}
recovering the floor-based estimate of the return, without evaluating $C_s^{-1}$.
An approximate hard solve is equally admissible; its residual in the $Q$ sector
is simply not discarded.

There is also a second-level enclosure for the exact Schur return.  For any
$Y:\mathbb C^m\to P\mathcal H_0$, put $e=r_s-C_sY$.  Completing the same square
gives
\begin{equation}
 r_s^*Y+Y^*r_s-Y^*C_sY
 \preceq\Gamma_s\preceq
 r_s^*Y+Y^*r_s-Y^*C_sY+s^{-1}e^*e.
 \label{f16v57:Schur-residual-enclosure}
\end{equation}
The only inverse estimate used is $C_s\succeq sP$.  This is a two-sided
source-matrix enclosure, not a scalar estimate on selected diagonal probes.

\subsection{A finite-bank certificate from native short-lag correlations}

Fix a trial bank $X:\mathbb C^r\to\mathcal H_0$ and restrict $W=XA$, where
$A$ is an $r\times m$ coefficient matrix.  Define
\begin{equation}
 N=X^*KV,\qquad M_s=X^*K(K+sI)X\succeq0.
 \label{f16v57:bank-matrices}
\end{equation}
Expanding \eqref{f16v57:trial-upper} cancels the pure $s$ terms and yields
\begin{equation}
 \mathcal U_s(XA)=G-A^*N-N^*A+A^*M_sA.
 \label{f16v57:bank-upper}
\end{equation}

\begin{theorem}[Optimal finite-bank local-Gram certificate]
\label{f16v57:finite-bank}
One has $\operatorname{ran}N\subseteq\operatorname{ran}M_s$.  With
\begin{equation}
 A_*=M_s^\dagger N,\qquad
 \mathcal C_s(X)=G-N^*M_s^\dagger N,
 \label{f16v57:optimal-certificate}
\end{equation}
the exact identities and inequalities are
\begin{gather}
 \mathcal U_s(XA)=\mathcal C_s(X)
   +(A-A_*)^*M_s(A-A_*),\label{f16v57:bank-completed-square}\\
 0\preceq sV^*(K+sI)^{-1}V\preceq\mathcal C_s(X)\preceq G.
 \label{f16v57:bank-bracket}
\end{gather}
For $2na\ge s^{-1}$, any verified capture bound
\begin{equation}
 \boxed{\qquad
 N^*M_s^\dagger N\succeq\vartheta G,\qquad 0<\vartheta\le1,
 \qquad}
 \label{f16v57:capture-target}
\end{equation}
implies $G^{[n]}\preceq(1-\vartheta)G$.
The certificate decreases in Loewner order when the trial span is enlarged.
It transforms by exact congruence under any source renormalization $V\mapsto VZ$.
\end{theorem}
\begin{proof}
For $z\in\ker M_s$,
\[
 0=\|K^{1/2}(K+sI)^{1/2}Xz\|^2
 \quad\Longrightarrow\quad K^{1/2}Xz=0
 \quad\Longrightarrow\quad z^*N=0.
\]
Here $(K+sI)^{1/2}$ is boundedly invertible and commutes with $K$.
Finite-dimensional orthogonal complementation proves the range inclusion.
Thus $M_sA_*=N$, and expansion proves
\eqref{f16v57:bank-completed-square}.  The residual theorem supplies the lower
bound in \eqref{f16v57:bank-bracket}; taking $A=0$ supplies the upper bound.
The capture implication follows by the same theorem.

To prove monotonicity without taking a scalar minimum separately in each source
direction, embed the smaller bank's minimizing synthesis $XA_*$ in the larger
trial span.  Its matrix upper bound remains available there, while the completed
square gives a simultaneous matrix minimizer for the larger bank.  Finally,
$V\mapsto VZ$ sends $(G,N)$ to $(Z^*GZ,NZ)$ and leaves $M_s$ unchanged, proving
congruence of the certificate.
\end{proof}

The inverse in \eqref{f16v57:optimal-certificate} is not necessary for
certification: one may instead exhibit a coefficient matrix $A$ and verify
\eqref{f16v57:bank-upper} directly.  This is preferable when the exact kernel of
$M_s$ is known symbolically or when an interval computation would make a
pseudoinverse unstable.

The required matrices are native vacuum correlations at lags zero, one and two
\emph{relative to the chosen trial bank}.  Define
\begin{equation}
 C_{XV}^{(j)}=X^*\mathcal T_a^jV,\qquad
 C_{XX}^{(j)}=X^*\mathcal T_a^jX.
 \label{f16v57:short-lag-correlations}
\end{equation}
Then exactly
\begin{align}
 N&=a^{-1}\bigl(C_{XV}^{(0)}-C_{XV}^{(1)}\bigr),
 \label{f16v57:short-lag-N}\\
 M_s&=a^{-2}\bigl(C_{XX}^{(0)}-2C_{XX}^{(1)}+C_{XX}^{(2)}\bigr)
       +sa^{-1}\bigl(C_{XX}^{(0)}-C_{XX}^{(1)}\bigr).
 \label{f16v57:short-lag-M}
\end{align}
For local slice writers these are the actual centered mixed Euclidean
correlations in the same normalized vacuum.  There is no reset to a product
Haar or frozen Gaussian law.  If $X$ itself contains transfer iterates, its
entries require correspondingly longer original-source correlations; the
short-lag notation does not erase that cost.

A small illustrative case shows the distinction from a full-space gap claim.
On a two-dimensional centered space take
\[
 K=\begin{pmatrix}\varepsilon&0\\0&\Delta\end{pmatrix},\qquad
 0<\varepsilon\ll\Delta,\quad a\Delta\le1,\qquad V=X=\binom01.
\]
The certificate is $s/(s+\Delta)$ although the full discrete-generator gap is
$\varepsilon$.  If instead $V=X=\binom10$, it is $s/(s+\varepsilon)\to1$.
Thus the procedure keeps a genuinely visible soft source rather than removing
it by a normalization convention.  This finite matrix is only an algebraic
diagnostic, not a Wilson regulator.

\subsection{The omitted trial space is paid only on the source}

The difference between the optimized certificate and the exact resolvent has
an especially simple geometric form.  Define
\begin{equation}
 J=K^{1/2}(K+sI)^{1/2}X,\qquad
 Z_s=K^{1/2}(K+sI)^{-1/2}V,
 \label{f16v57:energy-bank}
\end{equation}
and let $\Pi_J$ be the orthogonal projection onto the finite-dimensional range
of $J$.  Notice that $0\preceq Z_s^*Z_s\preceq G$.

\begin{proposition}[Exact source-marked approximation error]
\label{f16v57:source-marked-error}
One has
\begin{equation}
 \boxed{\quad
 \mathcal C_s(X)-sV^*(K+sI)^{-1}V
 =Z_s^*(I-\Pi_J)Z_s.
 \quad}
 \label{f16v57:source-marked-error-eq}
\end{equation}
Consequently a bound
$\|(I-\Pi_J)Z_sc\|^2\le\epsilon\,c^*Gc$ for all coefficients $c$ pays the entire
omitted trial space by $\epsilon G$.  It requires approximation of the marked
source range, not a volume-uniform operator-norm approximation to the full
transfer or the full vacuum density.
\end{proposition}
\begin{proof}
The commuting functions of $K$ give $J^*J=M_s$ and $J^*Z_s=N$.  Hence
$N^*M_s^\dagger N=Z_s^*\Pi_JZ_s$, since
$J(J^*J)^\dagger J^*=\Pi_J$.  Also
\[
 G-Z_s^*Z_s
 =V^*\bigl[I-K(K+sI)^{-1}\bigr]V
 =sV^*(K+sI)^{-1}V.
\]
Subtracting proves the identity and its stated consequence.
\end{proof}

This identity is the positive counterpart of the dependency audit above.  It
does not assert that a large-volume Wilson density is globally close to a
Gaussian.  It identifies exactly the smaller approximation problem sufficient
for this local certificate.  Controlling that problem uniformly at interacting
physical scales remains necessary.

\subsection{An explicit trial bank and its cutoff-dependent depth}

The source-marked approximation can be made constructive at every fixed
regulator.  For an integer $k\ge1$, let $X_k$ have trial span
\begin{equation}
 \mathcal X_k=\operatorname{span}\{\mathcal T_a^jV c:
                 0\le j<k,\ c\in\mathbb C^m\}.
 \label{f16v57:Krylov-bank}
\end{equation}
Linear dependence is allowed.  Any polynomial of degree at most $k-1$ in
$K+sI$ applied to $V$ belongs to this span.

\begin{theorem}[A volume-free Chebyshev bound with the cutoff cost exposed]
\label{f16v57:Chebyshev}
Put
\begin{equation}
 \kappa_a=1+\frac1{as},\qquad
 \rho_a=\frac{\sqrt{\kappa_a}-1}{\sqrt{\kappa_a}+1}\in(0,1).
 \label{f16v57:Chebyshev-condition}
\end{equation}
For the optimized certificate of the bank \eqref{f16v57:Krylov-bank},
\begin{equation}
 \boxed{\quad
 0\preceq\mathcal C_s(X_k)-sV^*(K+sI)^{-1}V
 \preceq4\rho_a^{2k}G.
 \quad}
 \label{f16v57:Chebyshev-error}
\end{equation}
There is no Hilbert-space dimension or spatial-volume factor.  A sufficient
depth for an error at most $\epsilon G$, with $0<\epsilon<1$, is
\begin{equation}
 k\ge
 \left\lceil\frac{\log(4/\epsilon)}{2\log(1/\rho_a)}\right\rceil.
 \label{f16v57:Chebyshev-depth}
\end{equation}
For fixed physical $s$ and $a\downarrow0$, this sufficient depth is
$O((as)^{-1/2}\log(4/\epsilon))$; it is not a cutoff-independent finite-bank
claim.
\end{theorem}
\begin{proof}
The spectrum of $A=K+sI$ lies in $[\alpha,\beta]$, where
$\alpha=s$ and $\beta=s+a^{-1}$.  Let $\mathsf T_k$ denote the Chebyshev
polynomial, distinguished here from the transfer.  Define the residual
polynomial
\begin{equation}
 p_k(z)=
 \frac{\mathsf T_k((\beta+\alpha-2z)/(\beta-\alpha))}
      {\mathsf T_k((\beta+\alpha)/(\beta-\alpha))}.
 \label{f16v57:Chebyshev-polynomial}
\end{equation}
It has $p_k(0)=1$, so $(1-p_k(z))/z$ is a polynomial of degree at most $k-1$.
Therefore
$W_k=[(1-p_k(A))/A]V$ belongs to the trial span and has exact residual
$R_s(W_k)=p_k(A)V$.

For $z\in[\alpha,\beta]$, the numerator in
\eqref{f16v57:Chebyshev-polynomial} has absolute value at most one.  If
$d=(\beta+\alpha)/(\beta-\alpha)>1$, the elementary identity
\[
 \mathsf T_k(d)=\tfrac12\bigl[(d+\sqrt{d^2-1})^k
                           +(d-\sqrt{d^2-1})^k\bigr]
\]
shows that
\[
 \sup_{[\alpha,\beta]}|p_k|
 \le\frac{2\rho_a^k}{1+\rho_a^{2k}}
 \le2\rho_a^k,
 \qquad d-\sqrt{d^2-1}=\rho_a.
\]
The residual identity and $0\preceq I-sA^{-1}\preceq I$ now imply
\[
 0\preceq\mathcal U_s(W_k)-sV^*A^{-1}V
 \preceq4\rho_a^{2k}G.
\]
Optimization over the bank can only decrease this error, proving
\eqref{f16v57:Chebyshev-error}.  Solving $4\rho_a^{2k}\le\epsilon$ gives
\eqref{f16v57:Chebyshev-depth}.  Finally,
$\log(1/\rho_a)=2\operatorname{arctanh}(\kappa_a^{-1/2})
\sim2\sqrt{as}$ as $as\downarrow0$, proving the stated order.
\end{proof}

For comparison, the direct geometric trial
\[
 W_k^{\rm geom}=a\sum_{j=0}^{k-1}(1+as)^{-j-1}\mathcal T_a^jV
\]
has residual $(1+as)^{-k}\mathcal T_a^kV$ and thus pays an error at most
$(1+as)^{-2k}G$.  Its sufficient depth is of order $(as)^{-1}\log(1/\epsilon)$.
The Chebyshev construction improves this worst-case sufficient estimate; no
optimality assertion for all source-adapted banks is made.

The estimate does not by itself create a gap.  If the exact source resolvent
ratio approaches one, the certificates approach one as well.  If a strict
resolvent margin exists on a declared mixed source family, the finite bank
approximates it with the explicit cost above.  The trial vectors need not be
spatially local merely because the original source is local.  Neither
coefficient stability nor certified evaluation of high-lag correlations follows
from the polynomial approximation theorem alone.

\subsection{Certified input errors and exact source kernels}

The native moments in \eqref{f16v57:short-lag-N}--\eqref{f16v57:short-lag-M}
contain differences divided by $a$ and $a^2$.  An unscaled numerical error in a
short-lag correlation is therefore not automatically harmless on the physical
trajectory.  A direct coefficient witness gives a simple explicit bookkeeping
rule.

Suppose approximate matrices $\widetilde G,\widetilde N,\widetilde M_s$ obey
\begin{equation}
 \|G-\widetilde G\|\le e_G,\qquad
 \|N-\widetilde N\|\le e_N,\qquad
 \|M_s-\widetilde M_s\|\le e_M.
 \label{f16v57:input-error}
\end{equation}
For a \emph{fixed, exhibited} coefficient matrix $A$, form
$\widetilde{\mathcal U}_s(A)=\widetilde G-A^*\widetilde N
-\widetilde N^*A+A^*\widetilde M_sA$.  Then
\begin{equation}
 \widetilde{\mathcal U}_s(A)+
 \bigl[(1+q^2)e_G+2\|A\|e_N+\|A\|^2e_M\bigr]I
 \preceq q^2\widetilde G
 \label{f16v57:certified-robust-test}
\end{equation}
is sufficient for the true $\mathcal U_s(XA)\preceq q^2G$.
Indeed the two cross errors cost at most $2\|A\|e_NI$, the quadratic error at
most $\|A\|^2e_MI$, and the errors of $G$ on the two sides cost
$(1+q^2)e_GI$.  This proves the claim without differentiating a pseudoinverse.

An absolute error term proportional to $I$ will usually preclude this test on
an unresolved null direction.  One must either keep the exact algebraic source
kernel and test on a verified quotient, or prove relative error inequalities
that vanish on that kernel.  Deleting small eigenvalues after inspecting the
numerical outcome is not a proof of either condition.  The diagnostic supplied
with this version checks floating-point identities on finite fixtures; it does
not claim interval-certified native Wilson moment enclosures.

\subsection{What the physical contraction would require next}

The new operational alternative to V56's assumed complementary contraction is
now the following inequality, on the actual mixed physical source family:
\begin{equation}
 \boxed{\qquad
 N_a^*M_{s,a}^{\dagger}N_a\succeq\vartheta G_a,
 \quad s=(2\tau)^{-1},\quad0<\vartheta<1
 \ \text{independent of the regulator}.
 \qquad}
 \label{f16v57:next-native-target}
\end{equation}
An exhibited-coefficient version of
\eqref{f16v57:certified-robust-test} is equally sufficient.  The bank and the
source must use the same native vacuum, the same Perron normalization, and the
same physical time spacing.  This target is not supplied by the attached files
or by the finite diagnostics below.

A first meaningful native test would retain a mixed family of centered,
gauge-invariant writers in one fixed physical ball, with its source
renormalization declared before the test.  Mixed plaquette orientations and
noncoplanar contractible loop writers are useful because a diagonal-only test
cannot certify the matrix inequality.  One must then obtain the native
$G_a,C_{XV}^{(0)},C_{XV}^{(1)},C_{XX}^{(0)},C_{XX}^{(1)},C_{XX}^{(2)}$ with errors
small enough after the explicit $a^{-1}$ and $a^{-2}$ scalings.  A coarse
construction may use the exact source return
\eqref{f16v57:source-return}; it may not replace it with $PV_a$ alone.

For a fixed finite family, convergence of its entrance and delayed Grams passes
any such common inequality to the limiting quotient by testing each coefficient
vector.  To infer the programme's local-vacuum gap, the same $\tau$ and a common
positive capture margin are required for \emph{every finite family} in the
declared local core, together with noncollapse and density in one
reflection-positive Osterwalder--Schrader representation.  Under precisely
those additional inputs, the existing V55--V56 compiler gives the physical
lower bound
\begin{equation}
 m_{\rm local}\ge-\frac{1}{2\tau}\log(1-\vartheta).
 \label{f16v57:conditional-mass}
\end{equation}
This is a conditional consequence, not a mass estimate obtained in V57.  It
controls that local vacuum representation, not an unrelated source-orthogonal
finite-volume sector.  The infinite-volume, physical-renormalization and
Osterwalder--Schrader tasks are not bypassed.

The remaining mathematical target for V58 is consequently concrete: construct
or bound the first \emph{native mixed source/trial packet} in
\eqref{f16v57:next-native-target}, including its exact null space and source
normalization, or prove a source-marked connected estimate that yields that
packet.  Repeating the finite-torus lowest eigenvalue calculation, assuming a
cutoff-independent Krylov depth, or inserting an uncut cubic Gaussian
exponential would not advance this target.

\subsection{Verification, preservation, and result ledger}

The companion diagnostic
\begin{center}\small
\path{verify_ym_f16_v57_local_resolvent.py}
\end{center}
checks the exact rectangle identity on quaternion fixtures, relative source
fractions including a collapsing and a rank-deficient family, the nonreducing
Schur identity and its positive floors, native physical-delay domination,
residual completed squares, the optimized finite-bank identity, source
congruence, the source-marked projection formula, short-lag reconstruction,
and the Chebyshev residual estimate.  It also records the elementary
third-chaos moment counterexample.  These are finite numerical or exact
algebraic consistency checks, not samples from the native Wilson vacuum, and
not a formal proof assistant verification of the analytical arguments.

The preservation audit removes the single delimited V57 insertion, reverses
the literal displayed title change from V57 to V56, and compares the resulting
bytes with the uploaded predecessor.  This tests the complete predecessor,
not only a selection of theorem labels.  The other supplied files are read-only
comparison inputs.  No assertion is made that every historical theorem in them
has been independently audited.

\begin{record}[V57 result ledger]
\label{f16v57:ledger}
\emph{Proved here under the displayed finite-regulator hypotheses:} relative
source-normalization invariance; the exact physical-size constant-field
rectangle coefficient and its $L^2$ normalized limit; the cubic exponential
obstruction; native physical-clock resolvent domination; the nonreducing
source-return identity; residual and finite-bank upper certificates; the exact
source-marked omitted-bank error; and the Chebyshev depth bound.

\emph{Inherited, not recertified here:} finite-regulator positive Wilson
transfer/Perron normalization; V55's clock convention; V56's raw-loop estimates
and reducing-band excision; and the earlier local-core-to-gap implication.

\emph{Not established:} physical continuum source renormalization and toron
invisibility after that normalization; a regulator-uniform native capture
margin; a uniform physical hard-compression floor; certified native vacuum
moment data; nontrivial Osterwalder--Schrader convergence; or the interacting
continuum Yang--Mills mass gap.  The parallel V67 log-normal step is explicitly
excluded from the dependency set.
\end{record}

\begin{firewall}
V57 advances the local-Gram branch by identifying the normalization-invariant
soft fraction, retaining the source correction in a nonreducing elimination,
and proving a finite native moment certificate with an explicit source-marked
approximation error.  The displayed rectangle example prevents raw scalarization
from being confused with normalized invisibility.  The native certificate is
proved but its Yang--Mills input matrices and a uniform strict margin have not
been populated.  The mathematical body of V1--V56 is preserved without changes.
The positive physical continuum mass-gap objective remains the target, not a
conclusion of this iteration.
\end{firewall}
% END F16 V57 APPEND

% BEGIN F16 V58 APPEND -- 2026-09-18
% Append-only continuation of the uploaded V57; predecessor body unchanged.

\clearpage
\section{V58: a native private-link source packet and stable physical-time certificates}
\label{f16v58:context}

V57 ends with a specific demand: populate a native mixed source/trial packet,
including its exact source kernel and normalization, instead of repeating an
abstract gap compiler.  We do so here for a finite family of genuine Wilson
writers.  The construction uses the original compact four-dimensional Wilson
measure, not a Gaussian replacement, a toron restriction, or a fitted vacuum.
A set of mutually noninteracting private links makes the conditional source
covariance explicitly coercive.  An elementary spherical estimate bounds the
cost by a polynomial in the bare Wilson coefficient.  This supplies an actual
finite-regulator capture inequality, although its displayed margin vanishes
as that coefficient tends to infinity.

A separate conditioning improvement replaces the microscopic difference
operator in the certificate by a difference over a fixed physical interval.
The original delayed transfer and physical clock are unchanged.  A geometric
trial then has a finite convex-correlation representation, so errors in its
input correlations are not multiplied by negative powers of the microscopic
time spacing.  This improves the conditioning of an input certificate; it
does not manufacture new information about long times.  We finish by proving
an exact innovation identity that identifies the missing time-transport
estimate and prevents an illicit iteration of the private-link bound.

\subsection{Inputs, source ownership, and nonduplication}

The predecessor is the post-fifteenth-archive V57.  Its definitions of the
native Perron-normalized transfer $\mathcal T_a$, physical spacing $a$,
centered source synthesis, and mixed-Gram order remain in force.  In particular
$0\preceq\mathcal T_a\preceq I$, its vacuum is $\Omega_a$, and
$\mathcal H_{0,a}=\{\Omega_a\}^{\perp}$.  We inherit the equivalence between its
gauge-invariant correlations and those of the original Wilson Euclidean
cylinders.  The finite-spatial-volume vacuum limit is taken first below; no
continuum or spatial thermodynamic limit is interchanged with it.

The construction of the positive Wilson transfer is classical; see
M.~L\"uscher, \emph{Construction of a selfadjoint, strictly positive transfer
matrix for Euclidean lattice gauge theories}, Commun. Math. Phys.
\textbf{54} (1977), 283--292,
\href{https://doi.org/10.1007/BF01614090}{doi:10.1007/BF01614090}, and
K.~Osterwalder and E.~Seiler, \emph{Gauge field theories on a lattice},
Ann. Phys. \textbf{110} (1978), 440--471,
\href{https://doi.org/10.1016/0003-4916(78)90039-8}{doi:10.1016/0003-4916(78)90039-8}.
The compact one-link conditional law is the usual Wilson single-link law,
associated with, for example, M.~Creutz, \emph{Monte Carlo study of quantized
SU(2) gauge theory}, Phys. Rev. D \textbf{21} (1980), 2308--2315,
\href{https://doi.org/10.1103/PhysRevD.21.2308}{doi:10.1103/PhysRevD.21.2308}.
Here that law is derived from the action and used in an analytic inequality,
not sampled.  None of these references supplies the missing continuum margin.

Earlier V23 minorization concerns a different conditional block law and its
buffer; it is not reapplied as a theorem about physical Euclidean time.  V57's
finite-bank square and its geometric trial are inherited algebraic inputs.
The new parts are the native private-link packet and its exact kernel, the
physical-interval version with positive correlation weights and stable input
errors, and the explicit distinction between a one-source innovation and its
transported descendants.  The excluded parallel V67 cubic-exponential step
remains excluded.

\subsection{An elementary lower covariance for a compact Wilson link}

Identify $\SU(2)$ with the unit quaternions $x\in S^3\subset\mathbb R^4$,
with normalized Haar measure $\sigma$.  For $h\in\mathbb R^4$, let
\begin{equation}
 d\nu_h(x)=Z_h^{-1}e^{h\cdot x}\,d\sigma(x),\qquad \kappa=|h|.
 \label{f16v58:link-law}
\end{equation}
All constants below are deliberately conservative.  An optimized Bessel-ratio
bound is unnecessary for the assertions made here.

\begin{proposition}[Uniform lower covariance for every staple direction]
\label{f16v58:spherical-covariance}
For all $h\in\mathbb R^4$ and $v\in\mathbb R^4$,
\begin{equation}
 \operatorname{Var}_{\nu_h}(v\cdot x)
 \ge \frac{|v|^2}{32000(1+\kappa)^2}.
 \label{f16v58:spherical-lower}
\end{equation}
The bound includes $h=0$, all directions parallel or perpendicular to $h$,
and the full compact link, without a chart or exceptional event.
\end{proposition}
\begin{proof}
Rotate $h$ to the first coordinate and put $t=x_0$, $y=1-t$, and
$b=(1+\kappa)^{-1}\in(0,1]$.  The density of $y$ on $[0,2]$, up to its
normalizer, is $e^{-\kappa y}\sqrt{y(2-y)}$.  Its normalizer $D_\kappa$ obeys
\begin{equation}
 D_\kappa\le5b^{3/2}.
 \label{f16v58:radial-normalizer}
\end{equation}
Indeed, if $\kappa\le1$, then
$D_\kappa\le\pi/2\le\pi\sqrt2\,b^{3/2}<5b^{3/2}$.
If $\kappa\ge1$, extend the integral to $[0,\infty)$ after replacing
$\sqrt{2-y}$ by $\sqrt2$.  This gives
$D_\kappa\le\sqrt{\pi/2}\,\kappa^{-3/2}
\le2\sqrt\pi\,b^{3/2}<5b^{3/2}$.

On the intervals $I_1=[b/4,b/2]$ and $I_2=[b,3b/2]$, respectively, the
unnormalized masses are at least $b^{3/2}/(16\sqrt2)$ and
$b^{3/2}/(10\sqrt2)$.  To see this, use $2-y\ge1/2$ on both intervals,
$e^{-\kappa y}\ge1/2$ on $I_1$, and $e^{-\kappa y}\ge1/5$ on $I_2$;
the latter two inequalities follow from $\kappa b\le1$.
Thus their probabilities satisfy
\[
 p_1\ge\frac1{80\sqrt2},\qquad p_2\ge\frac1{50\sqrt2}.
\]
For an independent copy $y'$, the two orders of the event
$y\in I_1,y'\in I_2$ yield
\[
 \operatorname{Var}(y)=\tfrac12\mathbb E(y-y')^2
 \ge p_1p_2(b/2)^2\ge b^2/32000.
\]
This is the longitudinal covariance.  Rotational symmetry in the remaining
three coordinates makes their common covariance
$\mathbb E[y(2-y)]/3$.  Its contribution from $I_2$ alone is at least
$p_2b/6\ge b/(300\sqrt2)\ge b^2/32000$.  All off-diagonal covariances vanish
in these coordinates.  The two lower eigenvalue bounds and rotation back to
$h$ prove the result, including $\kappa=0$.
\end{proof}

Consider now the original four-dimensional Wilson weight, before gauge fixing,
\begin{equation}
 \exp\left(\sum_p\gamma_p w_p(U)\right)\prod_e dH(U_e),
 \qquad w_p(U)=\tfrac12\operatorname{Tr}U_p\in[-1,1],\quad \gamma_p\ge0.
 \label{f16v58:Wilson-weight}
\end{equation}
The additive constant in the Wilson action has canceled.  Assume the lattice
is large enough that an elementary plaquette does not use the same link twice.
Conditional on all links except one link $e$, every plaquette containing $e$
is a linear functional of its quaternion.  Inverting $U_e$ merely applies the
orthogonal quaternion-conjugation map.  Its coefficient has norm one, since
the other three factors form a unit quaternion.  Consequently the conditional
law is \eqref{f16v58:link-law}, with
\begin{equation}
 |h_e|\le\sum_{p\ni e}\gamma_p.
 \label{f16v58:staple-bound}
\end{equation}
For the isotropic action of V55, $\gamma_p=\gamma$, an interior link belongs to
six plaquettes and $|h_e|\le6\gamma$.  Set
\begin{equation}
 d_\gamma=\frac1{32000(1+6\gamma)^2}.
 \label{f16v58:d-gamma}
\end{equation}
If a collection of links has no two members in one plaquette, its joint
conditional law given the remaining links is the product of these one-link
laws.  This follows directly by separating the exponent; gauge fixing or an
independent-link approximation has not been introduced.  Gauge-invariant
observables are evaluated in the original, gauge-redundant integral.

\subsection{Private links give an actual native mixed source packet}

Let $W_1,\ldots,W_r$ be normalized real Wilson loop traces on time slice zero.
For each $i$ choose a spatial link $e_i$ such that $W_i$ uses $e_i$ exactly once,
none of the other $W_j$ uses $e_i$, and no elementary spacetime plaquette
contains two of the selected links.  These are the \emph{private-link
conditions}.  All loops are contained in the chosen local spatial region.
Define native centered sources and correlation matrices by
\begin{equation}
 Vc=\left(\sum_{i=1}^r c_i(W_i-\langle W_i\rangle_{\pi_a})\right)\Omega_a,
 \quad G=V^*V,\quad C_j=V^*\mathcal T_a^jV,\quad j\ge0.
 \label{f16v58:native-packet}
\end{equation}
Here $C_0=G$ and $\pi_a=|\Omega_a|^2dm$ is the actual slice vacuum law.
The sources may equivalently be represented as centered functions in its
$L^2$ space.

\begin{theorem}[Native covariance, exact kernel, and a visible decrement]
\label{f16v58:private-packet}
At every finite spatial regulator and every finite $\gamma\ge0$,
\begin{align}
 d_\gamma I_r&\preceq G\preceq rI_r,\label{f16v58:Gram-bounds}\\
 \frac{d_\gamma}{2}I_r&\preceq G-C_j\preceq G
       \quad(j\ge1),\label{f16v58:decrement-bounds}\\
 G-C_j&\succeq\frac{d_\gamma}{2r}G\quad(j\ge1).
 \label{f16v58:relative-decrement}
\end{align}
The constants do not depend on the spatial volume or the length of the
Euclidean cylinder.  The product-Haar centered source Gram is exactly
$H=I_r/4$.  In particular the native source kernel is zero, not an empirically
thresholded small-eigenvalue space.

More generally, for any declared coefficient map $Z:\mathbb C^m\to\mathbb C^r$,
including regulator-dependent or rank-deficient maps, the corresponding
matrices are obtained by congruence.  Their entrance kernel is exactly
$\ker Z$, and
\begin{equation}
 d_\gamma Z^*Z\preceq G_Z\preceq rZ^*Z,\qquad
 G_Z-C_{j,Z}\succeq(d_\gamma/2)Z^*Z.
 \label{f16v58:normalized-packet}
\end{equation}
\end{theorem}
\begin{proof}
First work in a finite open Euclidean cylinder, with the observed slices in
its interior and positive constant boundary wavefunctions.  Condition on all
links except the selected $e_i$ at time zero.  The preceding factorization
makes these selected links independent, with conditional concentration at
most $6\gamma$.  The scalar writer $f_c=\sum_i c_iW_i$ is the sum of functions
of separate selected links; each has the form $c_i v_i\cdot x_i$, with
$v_i\in\mathbb R^4$ of norm one, determined by the conditioned links.
Thus, also for complex coefficients,
\begin{equation}
 \operatorname{Var}(f_c\mid\hbox{remaining links})
   \ge d_\gamma\sum_i|c_i|^2.
 \label{f16v58:conditional-private-variance}
\end{equation}
The law of total variance gives the lower bound on the covariance of $f_c$.
For the upper bound use
$\operatorname{Var}(f_c)\le\mathbb E|f_c|^2
\le r\sum_i|c_i|^2$, because $|W_i|\le1$.

For every $j\ge1$, the value $f_c$ on slice $j$ is measurable with respect to
the conditioned complement.  Hence the conditional variance of
$f_c(U_0)-f_c(U_j)$ is still bounded below by
\eqref{f16v58:conditional-private-variance}.  Its unconditional second
moment has the same lower bound.  Send the two temporal boundaries to
infinity at fixed spatial regulator.  The positive compact transfer and its
simple Perron vector, inherited here, make these finite inserted correlations
converge to the native vacuum correlations: applying
$(T/\lambda_0)^N$ to the nonorthogonal positive boundary vector converges to
its Perron projection.  All writers are bounded and all displayed constants
are independent of $N$.  Thus the inequalities pass to the limit.  Stationarity
and self-adjointness now give exactly
\[
 \mathbb E_{\rm vac}|f_c(U_0)-f_c(U_j)|^2
       =2c^*(G-C_j)c.
\]
Positivity and contraction of $\mathcal T_a^j$ give $0\preceq C_j\preceq G$.
These facts prove \eqref{f16v58:Gram-bounds} and
\eqref{f16v58:decrement-bounds}; the upper Gram bound gives
\eqref{f16v58:relative-decrement}.

Under product Haar measure, integrating a private link gives
$\mathbb E_H W_i=0$ and kills $\mathbb E_H W_iW_k$ for $i\ne k$.
Conditional on the other links, $W_i$ is one coordinate of a uniform unit
quaternion, whose square has mean $1/4$.  Thus $H=I_r/4$.  The positive lower
bound on $G$ proves the zero kernel, and congruence proves every assertion for
$Z$.  There is no assumption that a source normalization remains bounded.
\end{proof}

This argument uses conditional probabilities only in the full compact Wilson
cylinder.  It does not assert that a spatial heat-bath update equals the
physical transfer.  The link resampling is a device for bounding the native
Euclidean correlation; the time variable in $C_j$ remains the original one.
The vacuum limit is not taken at a rate assumed uniform in spatial volume.

\subsection{A concrete three-orientation packet with an exact redundant column}

Use spatial coordinate directions $1,2,3$ and a vertex $0$ in a torus whose
local neighborhood has no periodic identifications.  Set
\[
 W_1=w_{12}(0),\qquad W_2=w_{13}(0),\qquad W_3=w_{23}(0).
\]
Choose the directed private links
\begin{equation}
 e_1=(\widehat1,2),\qquad
 e_2=(\widehat3,1),\qquad
 e_3=(\widehat2,3),
 \label{f16v58:three-private-links}
\end{equation}
where $(x,\mu)$ denotes the link from $x$ to $x+\widehat\mu$.
Their endpoint pairs are disjoint, and their directions are distinct.  Two
edges of different directions in one elementary plaquette must share a
vertex.  Therefore no plaquette contains two selected links.  Each is in its
own displayed writer and not in either other writer, proving all the
private-link hypotheses without an independence assumption about the vacuum.
The three plaquettes together use nine distinct spatial links; their
interaction neighborhood contains 42 spacetime plaquettes (24 spatial and
18 temporal), but only the six staples of each private link enter the sharper
bound \eqref{f16v58:d-gamma}.

Include a redundant fourth writer $W_4=W_1+W_2$, and let
\begin{equation}
 B=\begin{pmatrix}1&0&0&1\\0&1&0&1\\0&0&1&0\end{pmatrix},\qquad
 H_4=\tfrac14 B^*B.
 \label{f16v58:explicit-B}
\end{equation}
The actual four-column packet satisfies
\begin{equation}
 d_\gamma B^*B\preceq G_4\preceq3B^*B,\qquad
 G_4-C_{j,4}\succeq\tfrac12d_\gamma B^*B\quad(j\ge1),
 \label{f16v58:four-column-bounds}
\end{equation}
and its exact null space is
\begin{equation}
 \ker G_4=\ker H_4
 =\operatorname{span}\{(-1,-1,0,1)^{\mathsf T}\}.
 \label{f16v58:four-column-kernel}
\end{equation}
For an additional normalization $Z_a$, replace $B$ by $BZ_a$ everywhere.
This specifies the full mixed packet, not just its diagonal entries.  One can
work on its known three-dimensional quotient by selecting the first three
writers before any numerical computation.  Multiplying these elementary
plaquettes by a power of $a^{-1}$ does not, by itself, establish a nontrivial
continuum field: that remains a separate source-convergence question.

\subsection{A difference over a physical interval, with the clock unchanged}

Let $d\ge1$ be an integer, write $b=da$, and define
\begin{equation}
 S=\mathcal T_a^d,\qquad K_b=(I-S)/b,\qquad
 0\preceq K_b\preceq b^{-1}I.
 \label{f16v58:physical-block}
\end{equation}
$K_b$ is an auxiliary difference operator, not the physical Hamiltonian and
not the one-step $K_a$ of V57.  In particular, at fixed positive $b$ its
continuum spectral function would be $(1-e^{-bH})/b$, not $H$.

\begin{proposition}[Block resolvent domination at the original delay]
\label{f16v58:block-domination}
For every $s>0$ and every integer $n\ge1$ with $2na\ge s^{-1}$,
\begin{equation}
 V^*\mathcal T_a^{2n}V
 \preceq sV^*(K_b+sI)^{-1}V.
 \label{f16v58:block-domination-eq}
\end{equation}
There is no requirement that $d$ divide $n$.  For fixed $\delta,\tau>0$, one
may take $d=\lceil\delta/a\rceil$, $n=\lceil\tau/a\rceil$, and
$s=(2\tau)^{-1}$.  Then $b\in[\delta,\delta+a)$ and the measured delayed Gram
is still at the original physical delay $na\to\tau$.
\end{proposition}
\begin{proof}
For a spectral value $t\in(0,1]$ of $\mathcal T_a$, set
$x=(1-t^d)/b$.  Since $\log(t^d)\le-(1-t^d)$,
\[
 t^{2n}\le e^{-(2n/d)(1-t^d)}=e^{-2na x}
      \le e^{-x/s}\le s/(s+x).
\]
At $t=0$ the inequality is immediate.  Functional calculus and congruence by
$V$ prove the result.  No replacement of the native Perron scalar occurs.
\end{proof}

All V57 residual-square identities hold with $K_b$ in place of $K_a$.
For the exhibited self-bank $X=V$, its native input matrices are
\begin{align}
 N_b&=b^{-1}(G-C_d),\label{f16v58:Nb}\\
 M_{b,s}&=b^{-2}(G-2C_d+C_{2d})+sb^{-1}(G-C_d)
          =V^*K_b(K_b+sI)V.\label{f16v58:Mb}
\end{align}
This packet is bounded analytically by the private-link theorem, rather than
left as an unpopulated list of vacuum moments.

\begin{theorem}[An explicit native capture certificate]
\label{f16v58:native-capture}
For a private-link family of size $r$ and its arbitrary coefficient map $Z$,
write $V_Z=VZ$ and form $N_b,M_{b,s}$ using $V_Z$ in place of $V$.  Then
\begin{gather}
 N_b\succeq\frac{d_\gamma}{2b}Z^*Z
       \succeq\frac{d_\gamma}{2rb}G_Z,\qquad
 sN_b\preceq M_{b,s}\preceq(b^{-1}+s)N_b,\label{f16v58:native-NM-bounds}\\
 \ker M_{b,s}=\ker N_b=\ker G_Z=\ker Z.\label{f16v58:native-NM-kernel}
\end{gather}
With the fixed coefficient witness $A=b(1+bs)^{-1}I_m$,
\begin{equation}
 \boxed{\quad
 \mathcal U_{b,s}(V_ZA)
 :=G_Z-A^*N_b-N_b^*A+A^*M_{b,s}A
 \preceq\left(1-\frac{d_\gamma}{2r(1+bs)}\right)G_Z.
 \quad}
 \label{f16v58:explicit-native-capture}
\end{equation}
The optimized capture therefore satisfies
\begin{equation}
 N_b^*M_{b,s}^{\dagger}N_b
 \succeq\frac{d_\gamma}{2r(1+bs)}G_Z.
 \label{f16v58:optimized-native-capture}
\end{equation}
Consequently the original physical-delay Gram satisfies the same upper bound
as \eqref{f16v58:explicit-native-capture} whenever $2na\ge s^{-1}$.  For the four-column three-orientation packet, $r=3$
and $Z=B$; its displayed capture margin is
\begin{equation}
 \vartheta_{\gamma,b,s}
   =\frac1{192000(1+6\gamma)^2(1+bs)}>0.
 \label{f16v58:explicit-margin}
\end{equation}
\end{theorem}
\begin{proof}
The decrement bound proves the lower bound on $N_b$.  Functional calculus gives
$0\preceq K_b^2\preceq b^{-1}K_b$, proving both inequalities for $M_{b,s}$.
They and the strict lower bound on $N_b$ on the range of $Z^*$ give the stated
kernels; vectors in $\ker Z$ are exactly zero source vectors.
Put $\alpha=b/(1+bs)$.  Since $N_b=N_b^*$ and
$\alpha M_{b,s}\preceq N_b$,
\[
 \mathcal U_{b,s}(V_Z\alpha)
       =G_Z-2\alpha N_b+\alpha^2M_{b,s}
       \preceq G_Z-\alpha N_b
       \preceq G_Z-\frac{d_\gamma}{2(1+bs)}Z^*Z.
\]
Now $G_Z\preceq rZ^*Z$ proves \eqref{f16v58:explicit-native-capture}.
The residual certificate of V57, or its completed-square proof with $K_b$,
and Proposition~\ref{f16v58:block-domination} prove the delayed-Gram assertion.
The optimized capture inequality follows because the V57 completed square
makes its optimized upper certificate no larger than this exhibited witness.
No pseudoinverse or sampled eigenvalue is needed to verify the witness.
\end{proof}

For this particular bare packet the direct bound
$C_{2n}\preceq[1-d_\gamma/(2r)]G$ from
Theorem~\ref{f16v58:private-packet} is slightly stronger.  The purpose of
\eqref{f16v58:explicit-native-capture} is to fill V57's specific source/trial
interface with controlled native data and exact kernels, not to claim that its
resolvent detour optimizes this elementary estimate.  The bound is uniform in
spatial volume for the fixed family, but not in the weak-coupling trajectory:
its stated margin is of order $(1+\gamma)^{-2}$ at fixed $b,s$ and therefore
vanishes as $\gamma\to\infty$.  This is a limitation of the proved lower bound,
not a proof that the true capture ratio vanishes.

\subsection{A positive-weight correlation certificate without microscopic cancellation}

Write $D_j=V^*S^jV=C_{jd}$, set $q=(1+bs)^{-1}\in(0,1)$, and let
\begin{align}
 \mathcal R_{b,s}&=sV^*(K_b+sI)^{-1}V,\label{f16v58:Rbs}\\
 \mathcal B_L&=(1-q)\sum_{j=0}^{L-1}q^jD_j+q^LD_L,
                 \qquad L\ge1.\label{f16v58:convex-certificate}
\end{align}
All coefficients in $\mathcal B_L$ are nonnegative and sum to one.

\begin{theorem}[Stable correlation envelope and its exact residual origin]
\label{f16v58:convex-envelope}
One has
\begin{gather}
 0\preceq\mathcal B_L-\mathcal R_{b,s}\preceq q^LG,\qquad
 \mathcal B_L\preceq G,\label{f16v58:envelope-error}\\
 \mathcal B_L-\mathcal B_{L+1}
        =q^{L+1}(D_L-D_{L+1})\succeq0.\label{f16v58:envelope-monotone}
\end{gather}
For even $L=2k$, $\mathcal B_{2k}$ is exactly the V57 residual upper bound for
the exhibited trial
\begin{equation}
 W_k=b\sum_{j=0}^{k-1}q^{j+1}S^jV,
 \qquad V-(K_b+sI)W_k=q^kS^kV.
 \label{f16v58:positive-trial}
\end{equation}
Thus it is a finite native trial-bank certificate, not just a bound for a
modified transfer.  For fixed $\delta,s>0$, $0<\epsilon<1$, and
$b=\lceil\delta/a\rceil a$, the sufficient number of correlation lags
\begin{equation}
 L\ge\left\lceil\frac{\log(1/\epsilon)}{\log(1+\delta s)}\right\rceil
 \label{f16v58:stable-lag-count}
\end{equation}
is independent of $a$ and gives error at most $\epsilon G$.
\end{theorem}
\begin{proof}
The norm-convergent series
\[
 s(K_b+sI)^{-1}=(1-q)(I-qS)^{-1}
       =(1-q)\sum_{j=0}^\infty q^jS^j
\]
shows that replacing its tail after $L-1$ by $q^LS^L$ is an upper bound,
because $0\preceq S^{L+j}\preceq S^L\preceq I$.  The difference is positive
and at most $q^LI$, proving \eqref{f16v58:envelope-error}.  Direct subtraction
proves \eqref{f16v58:envelope-monotone}.  The coefficients sum to one, so
$D_j\preceq G$ gives $\mathcal B_L\preceq G$.

For the even case the geometric sum in \eqref{f16v58:positive-trial} gives its
stated residual.  Since all its factors commute, V57's exact residual-error
identity becomes
\[
 \mathcal U_{b,s}(W_k)
 =V^*\left[s(K_b+sI)^{-1}
     +q^{2k}S^{2k}\{I-s(K_b+sI)^{-1}\}\right]V.
\]
Expand the first resolvent series through $2k-1$; its tail cancels the
subtracted resolvent tail in the second term, leaving exactly
\eqref{f16v58:convex-certificate}.  Finally,
$q^L\le(1+\delta s)^{-L}$ proves the lag count.
\end{proof}

The number in \eqref{f16v58:stable-lag-count} counts physical correlation lags,
not microscopic transfer applications.  The latest required original lag is
$Ld$, and $d$ grows like $\delta/a$.  Its physical time is $Lb$, which remains
finite for fixed $L$ but can exceed the target delay when high precision is
requested.  Also $D_L\preceq\mathcal B_L$.  Therefore, when a shorter-lag
correlation already has a certified strict margin, this positive average is
not a stronger estimate than that correlation.  The advantage here is a
stable realization of the resolvent interface, not free long-time data,
spatial locality of the trial vectors, or a stronger spectral inference.

\subsection{Error budgets on the exact source quotient}

Let $H\succeq0$ be a declared algebraic reference Gram with the exact source
kernel; for the concrete packet use $H_4$ in
\eqref{f16v58:explicit-B}, or its congruence after normalization.  Suppose
Hermitian input enclosures satisfy, for $0\le j\le L$,
\begin{equation}
 -e_jH\preceq\widetilde D_j-D_j\preceq e_jH,\qquad e_j\ge0.
 \label{f16v58:relative-input-errors}
\end{equation}
These are assumptions on certified correlations, not a property asserted of
ordinary Monte Carlo error bars.  They retain consistent native centering.
Define weights $w_j=(1-q)q^j$ for $j<L$ and $w_L=q^L$.

\begin{proposition}[A null-preserving, nonamplifying robust test]
\label{f16v58:robust-test}
Let $0\le\chi<1$ and
$\widetilde{\mathcal B}_L=\sum_{j=0}^Lw_j\widetilde D_j$.  If
\begin{equation}
 \boxed{\quad
 \widetilde{\mathcal B}_L-\chi\widetilde D_0
 +\left(\sum_{j=0}^Lw_je_j+\chi e_0\right)H\preceq0,
 \quad}
 \label{f16v58:stable-robust-test}
\end{equation}
then $\mathcal B_L\preceq\chi G$ and hence
$V^*\mathcal T_a^{2n}V\preceq\chi G$ whenever $2na\ge s^{-1}$.
If $e_j\le e$, the displayed error budget is at most $(1+\chi)eH$,
with no $a^{-1}$, $a^{-2}$, or bank-size multiplier.
\end{proposition}
\begin{proof}
Positive weights give
$\mathcal B_L\preceq\widetilde{\mathcal B}_L+(\sum_jw_je_j)H$, while
$-\chi G\preceq-\chi\widetilde D_0+\chi e_0H$.
Add the inequalities and use \eqref{f16v58:stable-robust-test}.  The envelope
and block domination imply the delayed statement.  The last assertion uses
$\sum_jw_j=1$.  If $Hc=0$, the two-sided enclosures imply that their error
operators annihilate $c$; exact source vectors also vanish there.  Thus the
test respects the full coefficient kernel rather than adding an artificial
ambient identity error on null directions.
\end{proof}

This removes the microscopic cancellation penalty in V57's short-lag
reconstruction, not the need for precise correlations.  Converting errors
from $H$ to the true $G$ still costs its actual lower bound; for the displayed
private-link packet $G\succeq4d_\gamma H$.  Regulator-dependent normalization
must be applied to both the data and their enclosures.  Neither vacuum
subtraction nor a poorly conditioned normalization is made exact by the
positive weights.

\subsection{Why one native decrement cannot be iterated, and the next transport target}

For the original transfer, an exact telescoping identity gives
\begin{equation}
 \boxed{\quad
 G-V^*\mathcal T_a^{2n}V
 =\sum_{j=0}^{n-1} V^*\mathcal T_a^j
       (I-\mathcal T_a^2)\mathcal T_a^j V.
 \quad}
 \label{f16v58:innovation-sum}
\end{equation}
Each summand is positive, but it involves the \emph{transported} source
$\mathcal T_a^jV$, not the original private-link-affine writer.  To verify the
identity, subtract consecutive powers $\mathcal T_a^{2j}$ and
$\mathcal T_a^{2j+2}$ and sum.  There is no reducing-source-space assumption.

In the native ground-state representation, let $P_a$ denote the Markov
operator unitarily equivalent to $\mathcal T_a$, and $U_0,U_1$ a stationary
pair.  For a finite family $f$ of centered functions, reversibility gives
\begin{equation}
 \mathbb E\operatorname{Cov}\bigl(f(U_1)\mid U_0\bigr)
 =\langle f,f\rangle_{\pi_a}
   -\langle P_af,P_af\rangle_{\pi_a}.
 \label{f16v58:Markov-innovation}
\end{equation}
Here covariance is a matrix, with complex conjugation in its first argument.
The formula follows from conditional expectation
$\mathbb E[f(U_1)\mid U_0]=P_af(U_0)$ and stationarity; taking
$f=P_a^jf_0$ proves the probabilistic form of
\eqref{f16v58:innovation-sum}.  It does not identify private spatial-link
resampling with this Markov step.

There is a simple exact warning against reusing the same decrement $n$ times.
On a two-dimensional centered space take
\[
 S=\operatorname{diag}(1-\eta,\eta),\qquad
 v=(\sqrt{1-\alpha},\sqrt\alpha)^{\mathsf T},\qquad
 0<\eta<1/2,\quad0<\alpha<1.
\]
With $G=1$, the correlations are
$c_j=(1-\alpha)(1-\eta)^j+\alpha\eta^j$.  They satisfy
\begin{equation}
 c_2-c_1^2=\alpha(1-\alpha)(1-2\eta)^2>0.
 \label{f16v58:no-illicit-iteration}
\end{equation}
Thus even a strict exact one-step source contraction does not imply its
square at the next step.  Replacing $S$ by $S^2$ gives the same obstruction
for squared-norm delays.  Adjoining a vacuum eigenvalue $1$ gives a positive
transfer with a simple vacuum and the same centered example.  This is an
algebraic counterexample to an inference, not a Wilson model or a claim about
its phase.  It also shows why the large number of lattice steps in a fixed
physical interval cannot simply multiply the $d_\gamma$ margin proved here.

A concrete next target for V59 is a source-marked lower bound on the sum in
\eqref{f16v58:innovation-sum}, for $n=\lceil\tau/a\rceil$, after the declared
physical source normalization.  One possible sufficient form is to exhibit
positive matrices $J_{a,j}$ such that
\begin{equation}
 V_a^*\mathcal T_a^j(I-\mathcal T_a^2)\mathcal T_a^jV_a\succeq J_{a,j},
 \qquad \sum_{j<n}J_{a,j}\succeq\vartheta_*G_a,
 \quad\vartheta_*>0\ \hbox{independent of }a.
 \label{f16v58:next-transport-target}
\end{equation}
The first private-link estimate supplies information at the entrance, but
not control of the affine components, conditional variances, or spatial
spread of the evolved sources at later $j$.  Establishing such control is
where native time transport must enter.  Replacing $J_{a,j}$ by the entrance
bound for every $j$ would assume precisely what has not been proved.

\subsection{Verification, preservation, and result ledger}

The companion file
\begin{center}\small
\path{verify_ym_f16_v58_native_packet.py}
\end{center}
checks exact lattice incidences, the algebraic source kernel, positive-weight
polynomial identities, native one-link covariance integrals at representative
staple strengths, finite-dimensional block and residual inequalities,
quotient error bounds, and the innovation counterexample.  The one-link integrals
are numerical checks on the exact conditional law, not samples of the full
native vacuum.  Random finite-dimensional tests are diagnostics, not proofs
or interval-certified Wilson correlation data.  The analytical arguments
above, rather than those diagnostics, establish the native bounds.

The preservation check removes this single delimited insertion and reverses
the displayed title change from V58 to V57, then compares all remaining bytes
with the uploaded V57.  The other uploaded versions are read-only comparison
inputs.  Their complete historical theorem sets have not been recertified.

\begin{record}[V58 result ledger]
\label{f16v58:ledger}
\emph{Proved here:} a global compact single-link covariance lower bound; native
mixed private-link covariance and decrement estimates; an explicit
three-orientation source/trial packet with its exact redundant-source kernel;
a polynomial-in-coupling positive finite-regulator capture margin; block
resolvent domination at the original physical delay; a positive-weight finite
correlation certificate with cutoff-independent lag count and null-preserving
error bookkeeping; and the innovation identity and noniteration obstruction.

\emph{Inherited, not recertified:} finite-regulator Wilson transfer positivity,
Perron normalization and its Euclidean-cylinder realization; V55's physical
clock; V57's residual-square algebra; and the earlier conditional
local-core-to-gap implication.

\emph{Not established:} a capture margin uniform as $\gamma\to\infty$;
transport bounds for evolved local sources through a fixed physical slab;
nontrivial continuum convergence of these normalized writers; a common strict
margin on every finite family of a dense local core; certified many-link
vacuum correlation values; nontrivial Osterwalder--Schrader convergence; or
the interacting continuum Yang--Mills mass gap.
\end{record}

\begin{firewall}
V58 populates a native source/trial interface and resolves its exact kernel
without numerical thresholding.  It improves the input conditioning by using
a fixed physical interval while retaining the original clock.  Neither fact
turns the explicit $(1+\gamma)^{-2}$ lower margin into a continuum mass.
The remaining obstruction is a quantitatively controlled accumulation of
native source innovations after time transport, together with the stated
continuum and dense-core inputs.  The mathematical body of V1--V57 is
preserved unchanged.
\end{firewall}
% END F16 V58 APPEND

% BEGIN F16 V59 APPEND -- 2026-09-18
% Append-only continuation of the attached post-fifteenth-archive V58.
% Prior mathematical text is preserved; only the displayed title is advanced.

\clearpage
\section{V59: buffered Wilson ribbons, transported innovations, and exterior predictability}
\label{f16v59:context}

The target left by V58 is a lower bound on an accumulated native innovation,
not another contraction of the original writer at a single time.  We obtain
such a bound by freeing selected spatial links along a \emph{time ribbon}.
Conditional on the remaining Wilson links, the ribbon is an exact compact
one-dimensional chain.  Its boundary-dependent transition operators have a
uniform maximal-correlation bound, which can be iterated on their actual
transported functions.  No affine closure of the native evolved source is
assumed.

The resulting estimate separates two costs.  Conditional transmission along
an independent ribbon has factor $q_\gamma=\tanh\gamma$, regardless of the
surrounding fields.  The other cost is the fraction of the source variance
that those surrounding fields already predict.  The latter must not be
omitted.  In particular, freezing the first full slice after the source
forces this conditional variance to vanish on any norm-preserving,
short-time-continuous physical source core.  A temporal buffer is therefore
part of the source construction, not a dispensable technical convenience.
For the private writers of V58, we also prove an explicit bound on the error
from a finite buffer.

\subsection{Dependencies, the fixed clock, and the new object}

We retain the finite-spatial-regulator stationary Wilson vacuum, its
Perron-normalized positive transfer $\mathcal T_a$, and its ground-state
Markov representation $P_a$.  A finite family of centered slice functions
$f=(f_1,\ldots,f_r)$ has synthesis $Vc=\sum_i c_i f_i\Omega_a$ and matrices
\begin{equation}
 G=V^*V,\qquad C_j=V^*\mathcal T_a^jV.
 \label{f16v59:source-Grams}
\end{equation}
Covariance is conjugate-linear in the first entry.  The physical time of
$P_a^n$ remains $na$; none of the conditional kernels below is declared to be
a new physical transfer.  Sources are bounded at each fixed regulator;
regulator-dependent bounds or normalizations are allowed.

V58's exact compact single-link lower covariance, private-link geometry,
and native innovation identity are inherited.  The global positive-cone
diameter obstruction in f15 V91 is not contradicted: here the comparison is
made on separate conditional ribbons, followed by an explicit exterior
variance payment.  Earlier local resampling estimates do not supply that
payment.  The classical positive-kernel contraction mechanism is reviewed by
B.~Lemmens and R.~Nussbaum, \emph{Birkhoff's version of Hilbert's metric and its
applications in analysis}, Theorem~2.9 and pp.~9--10,
\href{https://arxiv.org/abs/1304.7921}{arXiv:1304.7921}.  A direct
maximal-correlation proof with the precise constants needed here is included
below.  This is not a novelty claim about that general mechanism.

The archive f15 V62 contains a native thermal-moment estimate on its
\emph{selected} coupling sequence.  We state separately exactly what that
estimate would add when its regulator and source conventions match.  It is
not silently transplanted to V55's physical continuum trajectory.  The
excluded parallel V67 cubic-exponential argument remains excluded.

\subsection{The exact conditional law of buffered private ribbons}

Let $S$ be a set of spatial links with no two members in one elementary
spatial plaquette.  For integers $n\ge1$ and $\ell\ge0$, free all temporal
copies of these links at times
\begin{equation}
 I_{n,\ell}=\{-\ell,-\ell+1,\ldots,n+\ell\}.
 \label{f16v59:free-interval}
\end{equation}
All temporal links, all other spatial links, and all copies of links in $S$
outside this interval remain in the exterior sigma-field
$\mathcal E_{S,n,\ell}$.  Work first in a finite open Euclidean cylinder
containing the interval and its two neighboring slices.  Use the original
isotropic Wilson weight \eqref{f16v58:Wilson-weight}, with coefficient
$\gamma\ge0$, and no gauge fixing.

\begin{proposition}[Exact independent ribbon laws]
\label{f16v59:ribbon-law}
Conditional on $\mathcal E_{S,n,\ell}$, the free variables belonging to
different spatial links in $S$ are independent.  The chain for one link,
written in unit-quaternion coordinates, has density proportional to
\begin{equation}
 \prod_{t=-\ell}^{n+\ell}e^{h_t\cdot x_t}\,d\sigma(x_t)
 \prod_{t=-\ell}^{n+\ell-1}
       e^{\gamma x_t\cdot R_t x_{t+1}},\qquad x_t\in S^3,
 \label{f16v59:chain-density}
\end{equation}
where $R_t\in\mathrm{SO}(4)$ is fixed by the two temporal links of the
corresponding temporal plaquette.  At nonendpoint times $|h_t|\le4\gamma$;
at the two endpoints $|h_t|\le5\gamma$.  The conditional concentration of
any one free link after its neighboring free links are also fixed is at
most $6\gamma$.

The same conditional specification holds in the stationary native vacuum
at each fixed spatial regulator.  Its constants do not depend on the
spatial volume, the length of the surrounding cylinder, or $|S|$.
\end{proposition}
\begin{proof}
A spatial plaquette uses at most one member of $S$, so its dependence on a
free link is linear in that link's quaternion, with coefficient norm one.
Each spatial link belongs to four spatial plaquettes.  A temporal plaquette
uses two copies of the \emph{same} spatial link at consecutive times; it
never couples different members of $S$.  For fixed temporal links $A,B$ its
trace can be written
\[
 \tfrac12\operatorname{Tr}(U_t B U_{t+1}^{-1}A^{-1})
   =\langle U_t,A U_{t+1}B^{-1}\rangle_{\mathbb R^4}.
\]
Left and right multiplication by unit quaternions give an orthogonal,
orientation-preserving map $R_t$.  At each end of the free interval one
additional temporal plaquette contains a fixed outside copy of the spatial
link and contributes a linear field of norm at most $\gamma$.  Splitting
the action into these terms proves the density, the product over links,
and all stated bounds.  Additive terms independent of the free variables
cancel in conditional normalization.

For completeness, at fixed spatial regulator this finite-set conditional
kernel is a bounded continuous function of its finitely many boundary
links when tested against a continuous function.  Its denominator is
strictly positive on a compact space.  The finite-cylinder conditional
expectation identity therefore passes to the stationary vacuum cylinder
limit, using the inherited convergence of finite-history marginals.
Cylinder tests in the complement and a monotone-class argument identify
the resulting regular conditional law.  The cylinder limit is taken for
fixed $n,\ell$ and fixed spatial regulator; these integers may subsequently
depend on that regulator.  No uniform cylinder-convergence rate or exchange
with the continuum limit is used.
\end{proof}

In particular, the three links in \eqref{f16v58:three-private-links} satisfy
this proposition.  The actual three Wilson writers at any fixed slice are
affine functions of their respective ribbon variables, with unit
quaternion coefficient fixed by the exterior.  Their native time-evolved
writers need not be affine; the next estimates apply to arbitrary functions
of an entire endpoint configuration.

\subsection{A positive-kernel bound that survives arbitrary boundary messages}

The following lemma is classical in content.  Including its proof fixes the
factor in the exponential time cost and prevents a total-variation bound
from being silently substituted for a relative $L^2$ estimate.

\begin{proposition}[Cross ratio controls maximal correlation]
\label{f16v59:crossratio}
Let a strictly positive continuous joint density on compact spaces be
\[
 Z^{-1}A(x)k(x,y)B(y)\,d\mu(x)d\nu(y),
\]
where $A,B,k$ are positive continuous and the reference probability measures
have full support.  Suppose
\begin{equation}
 \sup_{x,x',y,y'}\log
 \frac{k(x,y)k(x',y')}{k(x,y')k(x',y)}\le\Delta.
 \label{f16v59:crossratio-assumption}
\end{equation}
For conditional expectation $K:L^2(Y)\to L^2(X)$,
\begin{equation}
 \|K\|_{L^2_0(Y)\to L^2_0(X)}
       \le \tanh(\Delta/4).
 \label{f16v59:maxcorr}
\end{equation}
Thus the ribbon edge in \eqref{f16v59:chain-density}, even after arbitrary
past and future integrations, has centered $L^2$ norm at most
$q_\gamma:=\tanh\gamma$.
\end{proposition}
\begin{proof}
We first record the elementary total-variation estimate.  If the density
ratio $r=dp/dq$ lies in $[m,M]$, $0<m\le1\le M$, with
$M/m\le R$, then
\[
 \|p-q\|_{\rm TV}=\mathbb E_q(r-1)_+
 \le\frac{(M-1)(1-m)}{M-m}
 \le\frac{\sqrt R-1}{\sqrt R+1}.
\]
The first bound follows by bounding the convex function $(r-1)_+$ by its
chord between $m$ and $M$.  For the second, fix $u=M/m$, maximize over
$m\in[1/u,1]$, and obtain $m=u^{-1/2}$ and
$(\sqrt u-1)/(\sqrt u+1)$; this is increasing in $u$.  Degenerate intervals
follow by continuity.

The ratio of two conditional densities of $Y$ given $x,x'$ has logarithmic
oscillation at most $\Delta$: both $B$ and the normalization constants
cancel from its oscillation.  Consequently the Dobrushin total-variation
coefficient of $K$ is at most $q=\tanh(\Delta/4)$.  The same reasoning
applies to the reverse conditional expectation $K^*$, since $A$ also
cancels.  A Markov operator with coefficient $q$ contracts the oscillation
of real bounded functions by $q$.  It follows that $K^*K$ contracts that
oscillation by $q^2$.

As an operator on $L^2(Y)$, $K^*K$ is positive, self-adjoint and compact;
compactness follows from its bounded continuous kernel relative to the
strictly positive marginal densities.  Its centered subspace is invariant.
Every eigenfunction with a nonzero eigenvalue has a continuous bounded
representative.  For a nonconstant real eigenfunction with eigenvalue
$\lambda$, the oscillation estimate gives
$\lambda\operatorname{osc}(g)\le q^2\operatorname{osc}(g)$, so
$\lambda\le q^2$.  If the centered norm is nonzero it is attained by such
an eigenfunction, by compact self-adjoint spectral theory.  This proves
\eqref{f16v59:maxcorr}; complex functions follow by complexification of the
real operator.

Finally, for $k(x,y)=e^{\gamma x\cdot Ry}$ the logarithm in
\eqref{f16v59:crossratio-assumption} is
$\gamma(x-x')\cdot R(y-y')\le4\gamma$.  Integrating any portion of a
nearest-neighbor chain outside this pair changes only its positive
one-sided factors $A,B$, not this cross ratio.  This proves the ribbon
assertion, with all boundary messages retained.
\end{proof}

No lower bound on a marginal density uniform in the exterior is required.
The marginal densities enter only the fixed-regulator compactness argument;
the numerical contraction constant is independent of their size.

\subsection{Conditional time transport and its genuine innovation sum}

Let $X_t$ denote all the free ribbon variables at time $t$, conditional on a
fixed exterior.  Its marginal spaces can differ with $t$.  Write
$K_tg=\mathbb E[g(X_{t+1})\mid X_t,\mathcal E]$.  The past and future
messages from the rest of the buffered interval are included in $K_t$.

\begin{proposition}[Tensorization and transported conditional innovations]
\label{f16v59:chain-transport}
For any finite number of independent ribbons,
\begin{equation}
 \|K_t\|_{L^2_0(X_{t+1}\mid\mathcal E)
                 \to L^2_0(X_t\mid\mathcal E)}\le q_\gamma,
 \qquad
 \|K_0\cdots K_{n-1}\|_0\le q_\gamma^n.
 \label{f16v59:chain-product}
\end{equation}
There is no factor proportional to the number of ribbons.  For a finite
family $F$ of terminal functions and
$g_t=\mathbb E[F(X_n)\mid X_t,\mathcal E]$, put
$A_t=\operatorname{Cov}(g_t(X_t)\mid\mathcal E)$.  Then
\begin{align}
 A_t&\preceq q_\gamma^2 A_{t+1},\label{f16v59:chain-step-Gram}\\
 \sum_{t=0}^{n-1}\mathbb E\!\left[
   \operatorname{Cov}(g_{t+1}(X_{t+1})\mid X_t,\mathcal E)
                      \mid\mathcal E\right]
 &=A_n-A_0\succeq(1-q_\gamma^{2n})A_n.
 \label{f16v59:conditional-innovation-sum}
\end{align}
\end{proposition}
\begin{proof}
For each ribbon its $L^2$ marginal space is the orthogonal sum of constants
and centered functions.  Decompose the finite tensor product according to
which factors are centered.  The constant factor has norm one.  On every
nonconstant summand at least one factor is centered, so the product
conditional operator has norm at most $q_\gamma$ by
Proposition~\ref{f16v59:crossratio}.  This proves the first assertion for all
finite sums of tensor functions, hence by density for the full $L^2$ space.
Consistency of successive marginals means each $K_t$ maps centered
functions to centered functions.  Their norms therefore multiply even
though the marginal Hilbert spaces change with $t$.

Apply the first assertion to every linear combination of the columns of
$g_{t+1}$.  The Markov property gives $g_t=K_tg_{t+1}$ and proves the matrix
inequality.  The law of total covariance identifies each summand in
\eqref{f16v59:conditional-innovation-sum} with $A_{t+1}-A_t$.
Telescoping and $A_0\preceq q_\gamma^{2n}A_n$ finish the proof.
\end{proof}

These transported conditional functions are the functions for which an
iteration is proved.  They are not asserted to be the functions $P_a^jf$
in V58's unconditioned innovation identity.  The comparison between the two
sums is instead obtained by conditioning, as follows.

\subsection{A native accumulated-innovation theorem with the exterior paid}

For any endpoint family $f$ as above, define
\begin{align}
 \mathcal A_{S,n,\ell}
  &=\mathbb E_{\rm vac}\operatorname{Cov}
               (f(U_n)\mid\mathcal E_{S,n,\ell}),
       \label{f16v59:capture-Gram}\\
 \mathcal R_{S,n,\ell}
  &=\operatorname{Cov}_{\rm vac}
       \bigl(\mathbb E[f(U_n)\mid\mathcal E_{S,n,\ell}]\bigr).
       \label{f16v59:predictor-Gram}
\end{align}
Here $U_n$ is the entire spatial configuration on slice $n$.  Conditional
on the exterior it is a function of $X_n$ and fixed remaining links.  Thus
the preceding proposition applies even when $f$ is not affine, not a sum
of independent writers, or has been obtained by applying $P_a^k$ to a local
writer.  Total covariance gives
$\mathcal A_{S,n,\ell}+\mathcal R_{S,n,\ell}=G$.

\begin{theorem}[Native transport with an exterior-predictability cost]
\label{f16v59:native-transport}
At every finite spatial regulator, $n\ge1$, $\ell\ge0$, and finite
$\gamma\ge0$,
\begin{equation}
 \boxed{\begin{aligned}
 G-C_{2n}&\succeq(1-q_\gamma^{2n})\mathcal A_{S,n,\ell},\\
 C_{2n}&\preceq q_\gamma^{2n}G+
              (1-q_\gamma^{2n})\mathcal R_{S,n,\ell}.
 \end{aligned}}
 \label{f16v59:native-main}
\end{equation}
The left-hand side of the first inequality is exactly
\begin{equation}
 \sum_{j=0}^{n-1}
 V^*\mathcal T_a^j(I-\mathcal T_a^2)\mathcal T_a^jV.
 \label{f16v59:actual-native-sum}
\end{equation}
The theorem bounds this \emph{sum}; it makes no term-by-term comparison
between its summands and those in
\eqref{f16v59:conditional-innovation-sum}.
\end{theorem}
\begin{proof}
Stationarity, the ground-state Markov representation and self-adjointness
give the exact identity
\[
 \mathbb E\operatorname{Cov}(f(U_n)\mid U_0)=G-C_{2n}.
\]
Conditional covariance decreases in expectation when more information is
supplied.  In matrix form the difference between the two sides of
\[
 \mathbb E\operatorname{Cov}(f(U_n)\mid U_0)
 \succeq
 \mathbb E\operatorname{Cov}(f(U_n)\mid U_0,\mathcal E_{S,n,\ell})
\]
is an expected covariance of a conditional mean, and is positive.
Given the exterior, the only unspecified variables of $U_0$ are $X_0$.
The right-hand side is therefore the expected conditional innovation sum
in \eqref{f16v59:conditional-innovation-sum}, with terminal covariance
$\mathcal A_{S,n,\ell}$.  This proves the first inequality.  Substitute
$\mathcal A=G-\mathcal R$ for the second.  The telescoping identity in V58
is precisely \eqref{f16v59:actual-native-sum}.
\end{proof}

\begin{proposition}[Explicit private-packet capture and exact coefficient kernels]
\label{f16v59:private-capture}
For the $r$ private writers of V58 and their associated ribbons,
\begin{equation}
 d_\gamma I_r\preceq\mathcal A_{n,\ell}\preceq G\preceq rI_r,
 \qquad
 d_\gamma=\frac1{32000(1+6\gamma)^2}.
 \label{f16v59:private-A}
\end{equation}
The matrix $\mathcal A_{n,\ell}$ is diagonal before any coefficient mixing.
Consequently
\begin{equation}
 G-C_{2n}\succeq
 (1-q_\gamma^{2n})d_\gamma I_r
 \succeq\frac{(1-q_\gamma^{2n})d_\gamma}{r}G.
 \label{f16v59:private-native-lower}
\end{equation}
For any declared $Z:\mathbb C^m\to\mathbb C^r$, all statements descend
by congruence; in particular
$\ker G_Z=\ker\mathcal A_{n,\ell,Z}=\ker Z$.
For the fourth redundant writer of V58, take $Z=B$ in
\eqref{f16v58:explicit-B}, giving the same exact one-dimensional kernel.
\end{proposition}
\begin{proof}
The terminal writers depend on separate ribbons once the exterior is fixed,
so their conditional cross covariances vanish.  Condition further on all
free links except the one at time $n$ in a given ribbon.  Its writer has a
unit quaternion coefficient, and its total staple concentration is at most
$6\gamma$.  V58's spherical lower covariance gives variance at least
$d_\gamma$.  Integrating this further conditioning proves the lower bound.
Total covariance and V58's entrance upper bound prove the rest of
\eqref{f16v59:private-A}.  Apply
Theorem~\ref{f16v59:native-transport}, then congruence.  The strictly
positive lower coefficient Gram determines the kernels algebraically.
\end{proof}

The bare bound \eqref{f16v59:private-native-lower} still has a vanishing
$\gamma^{-2}$ margin.  Its advance over the entrance estimate is the
justified accumulated-transport factor and the identifiable relative
capture matrix $\mathcal A$, not an unproved uniform source normalization.

\subsection{Why microscopic padding fails, and how a physical buffer converges}

A source at the terminal edge of a freed slab sees a full frozen slice only
one step later.  The following exact upper bound measures that problem.

\begin{proposition}[The frozen-future upper bound]
\label{f16v59:future-bound}
For any of the source families covered by
Theorem~\ref{f16v59:native-transport},
\begin{equation}
 0\preceq\mathcal A_{S,n,\ell}
       \preceq G-C_{2(\ell+1)}.
 \label{f16v59:future-upper}
\end{equation}
In particular, suppose a regulated source core has the relative
short-time continuity property
\begin{equation}
 G_a-C_{2k_a,a}\preceq\epsilon_a G_a,\qquad
 \epsilon_a\longrightarrow0,
 \label{f16v59:short-continuity}
\end{equation}
whenever the specified $k_aa\to0$.  A ribbon with
$\ell_a+1=k_a$ cannot have a regulator-uniform positive capture fraction on
that core.  This includes unbuffered ribbons $\ell_a=0$.
\end{proposition}
\begin{proof}
The entire slice $U_{n+\ell+1}$ is measurable in
$\mathcal E_{S,n,\ell}$.  Therefore
\[
 \mathcal A_{S,n,\ell}
 \preceq\mathbb E\operatorname{Cov}(f(U_n)\mid U_{n+\ell+1}).
\]
Reversibility and stationarity identify the right-hand side with
$G-C_{2(\ell+1)}$, just as in the proof of
Theorem~\ref{f16v59:native-transport}.  Combining with
\eqref{f16v59:short-continuity} proves the final assertion.  The continuity
condition is stated explicitly; convergence of a few fixed positive-time
correlations alone would not imply it.
\end{proof}

Increasing temporal padding does not change the factor $q_\gamma^{2n}$
in \eqref{f16v59:native-main}.  It also improves, rather than worsens, the
capture: if $\ell'\ge\ell$, then
$\mathcal E_{S,n,\ell'}\subseteq\mathcal E_{S,n,\ell}$ and total covariance
gives $\mathcal A_{S,n,\ell'}\succeq\mathcal A_{S,n,\ell}$.  For the
private writers this monotonicity has a quantitative remainder.

\begin{theorem}[Private-ribbon padding error]
\label{f16v59:padding-error}
For V58's $r$ private writers, $\ell'\ge\ell\ge0$, and fixed $n\ge1$,
\begin{equation}
 0\preceq\mathcal A_{n,\ell'}-\mathcal A_{n,\ell}
 \preceq\varepsilon_{n,\ell} I_r,
 \qquad
 \varepsilon_{n,\ell}
  =(q_\gamma^{n+\ell+1}+q_\gamma^{\ell+1})^2
  \le4q_\gamma^{2(\ell+1)}.
 \label{f16v59:padding-bound}
\end{equation}
The monotone limit $\mathcal A_{n,\infty}=\lim_{\ell\to\infty}
\mathcal A_{n,\ell}$ exists and satisfies the same bound with $\ell'=\infty$.
After any coefficient map $Z$, $I_r$ is replaced by $Z^*Z$.
\end{theorem}
\begin{proof}
The claim is immediate if $\ell'=\ell$.  Otherwise put
$m_\ell=\mathbb E[f(U_n)\mid\mathcal E_{n,\ell}]$.  Nested conditional
expectations give
\begin{equation}
 \mathcal A_{n,\ell'}-\mathcal A_{n,\ell}
 =\mathbb E\operatorname{Cov}(m_\ell\mid\mathcal E_{n,\ell'}).
 \label{f16v59:padding-Pythagoras}
\end{equation}
For each private writer, conditional on the larger free slab's exterior,
$m_{\ell,i}$ depends on the smaller slab's outside variables only through
its two adjacent boundary spins, at times $-\ell-1$ and $n+\ell+1$.
The one-sided fields throughout the smaller slab are fixed by the same
nonribbon links.  Hold the right boundary fixed and vary the left boundary.
Forward conditional kernels, with the right boundary included in their
future messages, each have total-variation coefficient at most $q_\gamma$
by Proposition~\ref{f16v59:crossratio}.  Their product from time
$-\ell-1$ to $n$ has coefficient at most $q_\gamma^{n+\ell+1}$.
Since the writer has range $[-1,1]$, its conditional mean changes by at
most $2q_\gamma^{n+\ell+1}$.  Reversing time gives at most
$2q_\gamma^{\ell+1}$ for a change of the right boundary with the left one
fixed.  Triangle inequality bounds the total range of $m_{\ell,i}$ by
$2(q_\gamma^{n+\ell+1}+q_\gamma^{\ell+1})$.

A real random variable of range length $D$ has variance at most $D^2/4$,
as follows by centering at the interval midpoint.  Moreover, these
conditional means belong to independent ribbons when
$\mathcal E_{n,\ell'}$ is fixed.  Their conditional covariance is diagonal.
Apply this range bound to each diagonal entry in
\eqref{f16v59:padding-Pythagoras}, proving
\eqref{f16v59:padding-bound}.  The finite matrices increase and are bounded
above by $G$, so their limit exists; the bound passes to it.  Congruence
proves the assertion about $Z$.
\end{proof}

The infinite-padding notation denotes this exact monotone limit.  It does
not assume an exchange of limits with the spatial regulator or discard a
tail sigma-field.  Equivalently, reverse martingale convergence identifies
the limiting conditional mean with projection onto the completed
intersection of the exterior sigma-fields.  No further identification of
that intersection is needed for the theorem.

On the exact private source quotient,
\eqref{f16v59:private-A} and \eqref{f16v59:padding-bound} give the useful
relative version
\begin{equation}
 \mathcal A_{n,\infty,Z}-\mathcal A_{n,\ell,Z}
   \preceq\frac{\varepsilon_{n,\ell}}{d_\gamma}G_Z.
 \label{f16v59:relative-padding}
\end{equation}
Thus a separately established
$\mathcal A_{n,\infty,Z}\succeq\eta G_Z$, $0\le\eta\le1$, would imply
\begin{equation}
 C_{2n,Z}\preceq
 \left[1-(1-q_\gamma^{2n})
       \left(\eta-\frac{\varepsilon_{n,\ell}}{d_\gamma}\right)_+\right]G_Z.
 \label{f16v59:paid-infinite-capture}
\end{equation}
The value of $\eta$ is a native source input, not supplied by this rearrangement.

\subsection{The physical-time cost and the distinction between time and width}

\begin{proposition}[A conditional transmission cost compatible with V55's trajectory]
\label{f16v59:physical-time}
For $\gamma>0$,
\begin{equation}
 q_\gamma^{2n}\le\exp[-2n e^{-2\gamma}].
 \label{f16v59:q-time}
\end{equation}
Suppose $\gamma_a\to\infty$ and
\begin{equation}
 a e^{2\gamma_a}\longrightarrow0.
 \label{f16v59:clock-condition}
\end{equation}
Then for every fixed $\tau>0$, $n_a=\lceil\tau/a\rceil$ satisfies
$q_{\gamma_a}^{2n_a}\to0$, uniformly in all conditioned exterior fields
and in the finite spatial volume.  If, more specifically,
\begin{equation}
 \log(1/a)=c\gamma_a+O(\log\gamma_a),\qquad c>2,
 \label{f16v59:exponential-clock}
\end{equation}
then for every fixed physical buffer $\delta>0$ and
$\ell_a=\lceil\delta/a\rceil$,
\begin{equation}
 \frac{\varepsilon_{n_a,\ell_a}}{d_{\gamma_a}}
       \longrightarrow0.
 \label{f16v59:padding-physical-limit}
\end{equation}
V55's displayed asymptotic-scaling relation has $c=3\pi^2/11>2$.
This last comparison uses that relation as a trajectory hypothesis; it is
not a proof of a nonperturbatively calibrated continuum trajectory.
\end{proposition}
\begin{proof}
Since $1-q_\gamma=2/(e^{2\gamma}+1)\ge e^{-2\gamma}$ and
$\log q_\gamma\le-(1-q_\gamma)$, \eqref{f16v59:q-time} follows.
The first limit follows from
$n_ae^{-2\gamma_a}\ge\tau/(ae^{2\gamma_a})\to\infty$.
For the buffered limit use
\[
 \frac{\varepsilon_{n_a,\ell_a}}{d_{\gamma_a}}
 \le128000(1+6\gamma_a)^2
       \exp[-2(\ell_a+1)e^{-2\gamma_a}].
\]
Under \eqref{f16v59:exponential-clock}, the expression
$(\ell_a+1)e^{-2\gamma_a}$ grows like
$\exp[(c-2)\gamma_a+O(\log\gamma_a)]$, which dominates
$\log(1+6\gamma_a)$.  This proves the second limit.  At $\gamma=0$ the
ribbon law has no temporal coupling and $q_0=0$ directly.
\end{proof}

The weaker condition \eqref{f16v59:clock-condition} alone is not asserted to
pay the polynomial relative-padding prefactor: one needs, for example,
$(\ell_a+1)e^{-2\gamma_a}/\log(2+\gamma_a)\to\infty$.  The stronger
trajectory \eqref{f16v59:exponential-clock} provides this automatically.
This distinction matters for borderline clocks.

There is a separate spatial-width cost.  If a free set has connected
components in the graph that joins selected spatial links sharing a spatial
plaquette, a component of $w$ links is a single chain with state space
$(S^3)^w$.  Its same-time potential need not be linear, but its temporal
pair factor is
\[
 \exp\left(\gamma\sum_{e=1}^{w}x_e\cdot R_e y_e\right).
\]
Its logarithmic cross ratio is at most $4w\gamma$.  The proof of
Proposition~\ref{f16v59:crossratio} therefore gives
$q_{\gamma,w}=\tanh(w\gamma)$.  Different components again tensorize, so
only the largest component size enters this particular estimate.  A
sufficient physical mixing condition is now $ae^{2w\gamma_a}\to0$.
For V55's coefficient $3\pi^2/11<4$, this worst-boundary estimate does not
vanish for even a two-link connected component.  This is a limitation of
the displayed sufficient estimate, not evidence of slow native mixing.
Temporal padding, unlike increasing connected spatial width, leaves the
one-link cross-ratio exponent unchanged.

\subsection{Twelve independent masks cover arbitrary slice sources}

The preceding method is not restricted to three elementary plaquettes.
On an even spatial torus, color each positively oriented spatial link
$e=(x,\mu)$ by
\begin{equation}
 \operatorname{col}(e)
   =\bigl(\mu,(x_\nu\bmod2)_{\nu\in\{1,2,3\}\setminus\{\mu\}}\bigr).
 \label{f16v59:12-coloring}
\end{equation}
There are twelve colors.  In any spatial plaquette, its two directions have
different colors, and the parallel pair in one direction differs in the
transverse parity.  Thus each color class $S_\alpha$ satisfies the
independent-ribbon hypothesis.  No bound proportional to the number of
links in a color class is introduced by conditional $L^2$ tensorization.

For an arbitrary finite family of bounded slice functions, including a
finite-regulator time-regularized family, define
\begin{equation}
 \overline{\mathcal A}_{n,\ell}
    =\frac1{12}\sum_{\alpha=1}^{12}
             \mathcal A_{S_\alpha,n,\ell}.
 \label{f16v59:color-capture}
\end{equation}
Averaging the twelve valid matrix inequalities yields
\begin{equation}
 \boxed{\quad G-C_{2n}
       \succeq(1-q_\gamma^{2n})\overline{\mathcal A}_{n,\ell}.
       \quad}
 \label{f16v59:color-transport}
\end{equation}
More generally any nonnegative weights summing to one may replace $1/12$.
One cannot sum twelve full copies of the lower bound: the same native
innovation would then be charged twelve times.  Nor is
$\sum_\alpha\mathcal A_{S_\alpha,n,\ell}\succeq G$ a consequence of the
coloring.  Such a variance-factorization estimate is an additional property
of the correlated native law.  Independent conditional color updates do
not imply independence of the color classes after the exterior is averaged.
For general writers the capture matrices may also have additional kernels;
only the entrance kernel is automatically annihilated by every capture.

There is an exact source-marked way to express each input without numerical
inversion of a Gram.  Given the \emph{same native exterior}, draw two
conditionally independent copies of the free ribbons and evaluate the
terminal family as $F^{(1)},F^{(2)}$.  Entrywise,
\begin{equation}
 (\mathcal A_{S,n,\ell})_{ij}
  =\frac12\mathbb E\left[
       \overline{F_i^{(1)}-F_i^{(2)}}
                  (F_j^{(1)}-F_j^{(2)})\right].
 \label{f16v59:replica-capture}
\end{equation}
Conditional independence cancels the conditional mean and proves the
identity.  The exterior is distributed by its actual Wilson marginal, not
by independent Haar links or an auxiliary chain.  Equation
\eqref{f16v59:replica-capture} is an exact integral representation, not an
assertion that its full native values have been sampled or enclosed.

\subsection{A matched moment input and the ultraviolet continuity obstruction}

For the elementary private plaquettes, write $s_i=1-W_i$.  Suppose that on
an explicitly matched sequence of native regulators
\begin{equation}
 \mathbb E_{\rm vac}(\gamma s_i)^2\le M,\qquad i=1,\ldots,r,
 \label{f16v59:thermal-moment-input}
\end{equation}
with fixed $M$.  This is a moment assumption on the actual vacuum, not on
the conditional one-link law.

\begin{proposition}[What the thermal moment pays, and what it cannot pay]
\label{f16v59:thermal-corollary}
Under \eqref{f16v59:thermal-moment-input}, for $\gamma\ge1$,
\begin{align}
 G&\preceq\frac{rM}{\gamma^2}I_r,\label{f16v59:thermal-G}\\
 \mathcal A_{n,\ell}
 &\succeq\frac{\gamma^2}{32000rM(1+6\gamma)^2}G
 \succeq\frac1{1568000rM}G.
 \label{f16v59:thermal-capture}
\end{align}
These implications and their source congruences are valid on that same
sequence.  For the thermally normalized synthesis $\widehat V=\gamma V$,
regardless of whether the moment upper bound is available,
\begin{equation}
 \widehat V^*\widehat V-
          \widehat V^*\mathcal T_a^2\widehat V
 \succeq\frac1{3136000}I_r\qquad(\gamma\ge1).
 \label{f16v59:UV-innovation}
\end{equation}
Consequently this complete unsmeared packet cannot have both a
norm-preserving strong source limit and a strongly continuous realization
of the microscopic transfer at physical time $a\to0$ on that limit.
\end{proposition}
\begin{proof}
For any coefficient vector $c$, Cauchy--Schwarz gives
\[
 \operatorname{Var}\Big(\sum_i c_iW_i\Big)
 =\operatorname{Var}\Big(\sum_i c_i s_i\Big)
 \le\mathbb E\Big|\sum_i c_i s_i\Big|^2
 \le\|c\|^2\sum_i\mathbb E s_i^2
 \le rM\gamma^{-2}\|c\|^2.
\]
Combine this with \eqref{f16v59:private-A} and
$\gamma^2/(1+6\gamma)^2\ge1/49$ for $\gamma\ge1$.
For \eqref{f16v59:UV-innovation}, use V58's stronger entrance estimate
$G-C_2\succeq(d_\gamma/2)I_r$ and multiply by $\gamma^2$.
If, under norm-preserving identifications into a common Hilbert space,
$\widehat V_ac$ converged strongly and
$\mathcal T_a\widehat V_ac-\widehat V_ac$ converged to zero for a fixed
nonzero $c$, the difference of their squared norms would tend to zero.
Equation \eqref{f16v59:UV-innovation} contradicts this.  Norm-preserving
strong convergence and compatible short-time transfer convergence are the
properties excluded; distributional limits or limits that discard variance
are not being ruled out.
\end{proof}

The comparison archive's f15 V62 thermal-moment proposition gives
$\mathbb E(\gamma_Ms_p)^q\le e^3 16^q q!$ on its selected vacuum sequence.
Taking $q=2$ would provide $M=512e^3$ in
\eqref{f16v59:thermal-moment-input} \emph{on exactly that sequence}.
The selected-coupling results are inherited, not re-proved here, and they
do not identify their physical $a_M$ with
\eqref{f16v59:exponential-clock}.  Even a compatible choice would not turn
these shrinking elementary plaquettes into a nontrivial strongly
continuous local core: \eqref{f16v59:UV-innovation} is an independent
obstruction.  Canceling the displayed $\gamma^{-2}$ variance by source
normalization is therefore not, by itself, a continuum mass argument.

\subsection{A reversible exterior example and the next source-level task}

The exterior payment is not a technical artifact of the proof.  Here is
an exact finite-state example showing why even complete conditional
ribbon decorrelation cannot replace it.  Let $Y_t\in\{-1,1\}$ be a
stationary reversible two-state chain with
\[
 Q_r(y,y')=\tfrac12(1+r yy'),\qquad 0<r<1.
\]
Given the entire $Y$ trajectory, let $X_t\in\{-1,1\}$ be independent with
$\Pr(X_t=x\mid Y_t=y)=(1+bxy)/2$, $0<b<1$.  The pair process has transition
\begin{equation}
 P((y,x),(y',x'))=Q_r(y,y')\frac{1+bx'y'}2
 \label{f16v59:exterior-example-kernel}
\end{equation}
and stationary law $\pi(y,x)=(1+bxy)/4$.  Detailed balance follows by
multiplying the displayed factors.  If $Jg(y,x)=g(y)$, then
$P=JQ_rJ^*$, where $J^*$ is conditional averaging over $x$ given $y$.
Thus $P$ is a positive self-adjoint contraction, with eigenvalues
$1,r,0,0$ and simple invariant constants.  Its transition entries are
strictly positive.  This is an algebraic reversible model, not a Wilson
measure or a claim about its phase.

For the centered writer $f(y,x)=x$,
\begin{equation}
 G=1,\qquad C_j=b^2r^j\ (j\ge1),\qquad
 \mathbb E\operatorname{Var}(X_n\mid (Y_t)_t)=1-b^2.
 \label{f16v59:exterior-example-data}
\end{equation}
The conditioned $X$ chain has zero transmission at every step.  Nevertheless,
with $r_a=e^{-a^2}$, $b_a^2=1-a$, and
$n_a=\lceil\tau/a\rceil$, its native $C_{2n_a}$ tends to one.
The missing variance is predicted by the retained $Y$ process.  All
quantities in \eqref{f16v59:exterior-example-data} follow from
$\mathbb E[X_t\mid Y_t]=bY_t$ and
$\mathbb E[Y_0Y_j]=r^j$.  No transfer eigenvalue inferred from a conditional
update has been identified with the native one.

A useful V60 source bank should consequently include an actual physical
time regularization, rather than renormalized elementary traces alone.
For example, for a declared finite family of spatially smeared Wilson
writers $F_{a,h}$, with all weights and normalizations retained, one may
study
\begin{equation}
 V_{a,h,\sigma}
   =\mathcal T_a^{\lceil\sigma/a\rceil}
           (F_{a,h}-\langle F_{a,h}\rangle)\Omega_a,
 \qquad \sigma>0.
 \label{f16v59:regularized-target}
\end{equation}
At fixed regulator this is again a centered slice function in the
$P_a$ representation, so the twelve-color transport theorem applies
without an affine or locality assumption on its evolved function.
Its entrance Gram is the original smeared source correlation at lag
$2\lceil\sigma/a\rceil$; noncollapse is therefore an explicit additional
requirement, not guaranteed by applying the transfer.

The remaining native task is to prove a positive relative lower bound on
$\overline{\mathcal A}_{n_a,\ell_a}$ for such a noncollapsing physically
regularized finite family, with $n_aa\to\tau>0$ and
$\ell_aa\to\delta>0$, and to make that margin common across the finite
families of a dense physical core.  The exact conditional-replica integral
\eqref{f16v59:replica-capture} specifies the source-dependent object that
must be bounded.  Further manipulation of the already proved one-ribbon
cross ratio would not provide that exterior-variance estimate.  V60 is
also the next inherited fifth-version companion checkpoint.

\subsection{Verification, preservation, and result ledger}

The companion
\begin{center}\small
\path{verify_ym_f16_v59_buffered_ribbons.py}
\end{center}
checks the literal four-dimensional ribbon incidence and quaternion
conditional action, the twelve-color property, exact rational
cross-ratio/maximal-correlation fixtures, inhomogeneous finite-chain
transport and mixed innovation identities, tensor-product contraction,
finite-padding monotonicity and its boundary estimate, the reversible
exterior example, physical-clock exponents, and the inherited source
kernel.  Discrete compact-chain fixtures are diagnostics for the proved
operator statements, not quadrature certificates for the continuous
Wilson vacuum.  Neither they nor the algebraic checks supply the missing
many-link source capture or continuum construction.

Removing this delimited appendix and reversing only the displayed title
change from V59 to V58 recovers the uploaded predecessor byte for byte.
The uploaded comparison archives remain read-only.  The report records
hashes, exact and floating-point test categories, and the tests actually
run.  Their historical theorem sets have not been recertified.

\begin{record}[V59 result ledger]
\label{f16v59:ledger}
\emph{Proved here:} the exact buffered compact Wilson ribbon law;
boundary-message-uniform conditional $L^2$ contraction and its genuine
transported innovation sum; an aggregate lower bound on the original
native innovation with the exterior predictor retained; native private
capture and exact source-kernel compatibility; the frozen-future upper
bound and its short-buffer obstruction; quantitative private-ribbon
padding convergence; the stated physical-clock and spatial-width
comparisons; twelve-color source transport and its exact conditional-replica
input; the implication from a matched thermal moment and the unsmeared
strong-continuity obstruction; and the reversible exterior counterexample.

\emph{Inherited or explicitly conditional:} finite-regulator Wilson
transfer/vacuum construction, the $6\gamma$ spherical covariance lower
bound, V58's entrance decrement, V55's asymptotic trajectory when used,
and the archive's selected thermal moment only on a matching sequence.

\emph{Still unproved:} a regulator-uniform positive exterior-variance
fraction for a noncollapsing physically regularized source core;
nontrivial continuum convergence of that core; a common dense-core
contraction; the required nontrivial Osterwalder--Schrader limit and its
interacting Yang--Mills identification; and the continuum mass gap.
\end{record}

\begin{firewall}
The native time-transport estimate is no longer based on repeating an
entrance private-link inequality.  Its conditional chain iteration is
proved with changing boundary messages and is transferred to the native
sum by conditioning.  What remains is source variance hidden in the
exterior.  A microscopic temporal buffer makes that obstruction decisive
on a strongly continuous core, whereas a fixed physical buffer has a
controlled error on the displayed exponential clock.  None of these
facts makes conditional ribbon mixing equal to a native physical mass.
The mathematical body of V1--V58 is preserved unchanged.
\end{firewall}
% END F16 V59 APPEND

% BEGIN F16 V60 APPEND -- 2026-09-18
% Append-only continuation of the uploaded post-fifteenth-archive V59.
% Prior mathematical text is preserved; only the displayed title is advanced.

\clearpage
\section[V60: the microscopic-ribbon obstruction and a source-preserving exterior reduction]{V60: the microscopic-ribbon obstruction\\and a source-preserving exterior reduction}
\label{f16v60:context}

The final target in V59 was a positive, regulator-uniform fraction of
conditional ribbon variance for a noncollapsing, physically regularized
source family.  That target requires an upstream correction.  The native
comparison in V59 is valid not only over the full physical delay but also
over every shorter delay inside the very same freed interval.  Consequently,
if the conditional ribbons mix in a vanishing physical time, their captured
variance must vanish on every source family with relative short-time
continuity.  Increasing the temporal buffer does not remove this obstruction.
This does not disprove a physical mass gap.  It shows that the proposed
conditional variance is measuring a fast component which a physical source
core must lose, rather than a fixed fraction of its surviving variance.

We prove the obstruction and then use the same conditional laws for a
different purpose.  They give explicit, gauge-invariant exterior predictors
which approximate regular source families in $L^2$, with a temporal window
that can shrink in physical units.  Integrating the ribbons out gives an
exact retained marginal.  Its induced two-score response has a controlled
temporal tail, but its signed local part and the collective dynamics of the
retained variables are not discarded.  The continuation therefore changes
the target instead of seeking another formulation of the impossible
fixed-fraction estimate.

\subsection{Inputs, the fifth-version checkpoint, and the precise correction}

The native transfer, its Perron normalization, the stationary Wilson vacuum,
and the physical step $a$ remain those of V55--V59.  We use V59's exact
independent-ribbon conditional law and its boundary-message-uniform
contraction $q_\gamma=\tanh\gamma$, not a new heat-bath interpretation of
physical time.  The physical scaling condition used below is explicitly
\begin{equation}
 a e^{2\gamma_a}\longrightarrow0.
 \label{f16v60:fast-clock}
\end{equation}
When stronger rates are needed we separately assume V59's displayed
trajectory hypothesis
\begin{equation}
 \log(1/a)=c\gamma_a+O(\log\gamma_a),\qquad c>2.
 \label{f16v60:strong-clock}
\end{equation}
Neither condition is established here as a nonperturbative calibration of
four-dimensional Yang--Mills theory.

The attached comparison checkpoint rereads the f14 archive's final
susceptibility/source-control section group, the f15 archive's
small-staple persistence and source-weighted transfer section group, and
the parallel V67 cubic-vertex conclusion.  Those are different objects:
selected-law susceptibility and small-staple persistence do not give a
physical-source ribbon capture floor.  The parallel cubic-exponential
argument excluded in V57 remains excluded.  Duplicate archive copies are
compared by hashes; no unavailable or newer companion is claimed to have
been reviewed.  The cumulative V57--V59 proofs are the direct dependencies
of this appendix.  This checkpoint does not recertify the archives.

The general relation between conditional expectation and elimination of
unresolved variables is classical; see A.~J.~Chorin, \emph{Conditional
Expectations and Renormalization}, Multiscale Modeling \& Simulation
\textbf{1} (2003), 105--118,
\href{https://doi.org/10.1137/S1540345902405556}{doi:10.1137/S1540345902405556}.
The positive-kernel mechanism used in V59 is surveyed by B.~Lemmens and
R.~Nussbaum, \href{https://arxiv.org/abs/1304.7921}{arXiv:1304.7921}.
These references supply context, not an imported continuum theorem.  All
new inequalities needed for the present Wilson reduction are proved below.

\subsection{The same exterior can be tested at every interior subdelay}

Let $I=\{L,\ldots,R\}$ be a finite interval of lattice times.  As in V59,
free the copies of a spatial link set $S$ throughout $I$, where no spatial
plaquette contains two members of $S$.  All remaining links generate
$\mathcal E_{S,I}$.  Fix a terminal time $t\in I$ and a finite family of
centered slice functions $f$.  Write
\begin{equation}
 G=V^*V,\quad C_j=V^*\mathcal T_a^jV,\quad
 A_{S,I,t}=\mathbb E\operatorname{Cov}
                   (f(U_t)\mid\mathcal E_{S,I}).
 \label{f16v60:source-data}
\end{equation}
The definitions make sense for $L^2$ functions.  The bounded-function
conditional estimates in V59 extend to this class by $L^2$ approximation
and the contraction of conditional expectation.  Every regulator-dependent
normalization is part of $V$.

\begin{theorem}[Subdelay comparison with the exterior held fixed]
\label{f16v60:subdelay}
For every integer $m\ge1$ with $t-m\in I$,
\begin{equation}
 \boxed{\quad
 (1-q_\gamma^{2m})A_{S,I,t}
       \preceq G-C_{2m}.
 \quad}
 \label{f16v60:subdelay-eq}
\end{equation}
In particular, if $t-1\in I$,
\begin{equation}
 A_{S,I,t}\preceq\cosh^2\gamma\,(G-C_2).
 \label{f16v60:one-step-upper}
\end{equation}
Both estimates hold for convex averages over independent masks, with no
factor depending on the number of masks or their cardinalities.
\end{theorem}
\begin{proof}
The stationary ground-state Markov representation gives
\[
 G-C_{2m}=\mathbb E\operatorname{Cov}(f(U_t)\mid U_{t-m}).
\]
Condition further on the \emph{same} $\mathcal E_{S,I}$.  Conditional on it,
the unfixed variables in $U_{t-m}$ and $U_t$ are respectively the two
ribbon configurations at those times.  Integrations over all other times
in $I$ produce boundary messages, without changing the centered $L^2$
norm bound of any transition.  Their $m$-step product has norm at most
$q_\gamma^m$.  The law of total conditional covariance therefore gives
\[
 \mathbb E\operatorname{Cov}(f(U_t)\mid U_{t-m},\mathcal E_{S,I})
 \succeq(1-q_\gamma^{2m})A_{S,I,t}.
\]
Expected conditional covariance decreases on supplying more information.
This proves \eqref{f16v60:subdelay-eq}.  For $m=1$ use
$1-\tanh^2\gamma=\cosh^{-2}\gamma$.  Finally average valid Loewner
inequalities with nonnegative weights of total mass one.  In particular,
no twelvefold copy of a native decrement is introduced.
\end{proof}

There is a useful operator statement.  Identify the native slice space with
$L^2(\pi_a)$ through the ground-state transform.  If $J_t$ is the isometric
inclusion of a slice function into the stationary path space and $\Pi_E$
is conditional expectation onto $\mathcal E_{S,I}$, then
\begin{equation}
 B_{S,I,t}=J_t^*(I-\Pi_E)J_t,\qquad
 0\preceq B_{S,I,t}\preceq I,
 \qquad
 B_{S,I,t}\preceq\cosh^2\gamma\,(I-P_a^2).
 \label{f16v60:operator-upper}
\end{equation}
This follows by applying the theorem to each slice vector and polarizing.
$B_{S,I,t}$ need not commute with $P_a$.  No pointwise spectral minimum of
the two upper bounds in \eqref{f16v60:operator-upper} is being asserted.

For V59's original interval $I=\{-\ell,\ldots,n+\ell\}$ and terminal
time $t=n$, every $m\le n+\ell$ is available.  In particular, the upper
bound is independent of the distance from $t$ to the right-hand frozen
slice.  This is the point not used in V59's frozen-future obstruction.

\subsection{Fast conditional mixing forces physical capture to vanish}

A finite source family has \emph{relative short-time continuity} if, for
every sequence of integers $m_a\ge1$ with $a m_a\to0$, there are
$\varepsilon_a\to0$ such that
\begin{equation}
 0\preceq G_a-C_{2m_a,a}\preceq\varepsilon_a G_a.
 \label{f16v60:relative-continuity}
\end{equation}
This is an order statement on the exact source quotient; it does not use a
numerical threshold for $\ker G_a$.  A noncollapsing norm-preserving physical
source limit with compatible short-time transfers must have this property
on each finite nondegenerate source quotient.  Noncollapse alone does not
imply it.

\begin{theorem}[No fixed ribbon fraction on a continuous physical core]
\label{f16v60:no-go}
Assume \eqref{f16v60:fast-clock} and
\eqref{f16v60:relative-continuity}.  Let $I_a$ and $t_a$ satisfy
\[
 t_a-L_a\ge\lceil e^{2\gamma_a}\rceil.
\]
Then, uniformly over all independent masks $S_a$ and all choices of the
right endpoint $R_a\ge t_a$,
\begin{equation}
 0\preceq A_{S_a,I_a,t_a}\preceq\delta_aG_a,
 \qquad \delta_a\longrightarrow0.
 \label{f16v60:capture-vanishes}
\end{equation}
The same statement holds for arbitrary convex mask averages.  Thus a
positive constant $\eta_*$ cannot satisfy
$\overline A_a\succeq\eta_*G_a$ on any nonzero family obeying these
hypotheses.  In particular this rules out V59's proposed fixed-fraction
target for a genuine continuous physical core when $n_aa\to\tau>0$,
regardless of whether $\ell_aa$ tends to zero, to a positive constant, or
to infinity.
\end{theorem}
\begin{proof}
Set $m_a=\lceil e^{2\gamma_a}\rceil$.  The physical delay $a m_a$
tends to zero.  The elementary estimate used in V59,
$q_\gamma^{2m}\le\exp(-2m e^{-2\gamma})$, yields
\[
 1-q_{\gamma_a}^{2m_a}\ge1-e^{-2}.
\]
Apply Theorem~\ref{f16v60:subdelay} at $m_a$, and then
\eqref{f16v60:relative-continuity}.  One may take
$\delta_a=\varepsilon_a/(1-e^{-2})$.  Neither estimate depends on the
mask or on $R_a$.  Averaging preserves the bound.  If also
$A_a\succeq\eta_*G_a$, choose any coefficient with positive entrance
variance to obtain $\eta_*\le\delta_a$, a contradiction for small $a$.
For V59's interval, $m_a\le n_a+\ell_a$ eventually since
$a m_a\to0$ and $a n_a\to\tau>0$.
\end{proof}

The conclusion is stronger than a failure to prove the requested fraction:
that fraction is incompatible with the stated physical continuity and
clock.  It is not a claim that every sequence named ``regularized'' has
physical continuity, nor a claim that the native theory is massless.

\begin{proposition}[Two checkable sufficient continuity conditions]
\label{f16v60:continuity-tests}
Each of the following suffices for
\eqref{f16v60:relative-continuity} on a finite source quotient.

First, let
\[
 E_{a,M}=\mathbf1_{[e^{-aM},1]}(\mathcal T_a),\qquad
 F_a(M)=V_a^*(I-E_{a,M})V_a.
\]
Suppose there are bounds $F_a(M)\preceq r_a(M)G_a$ with
\begin{equation}
 \lim_{M\to\infty}\limsup_{a\to0}r_a(M)=0.
 \label{f16v60:spectral-tightness}
\end{equation}
Second, after choosing a fixed exact coefficient quotient, suppose
$G_a\to G_\infty\succ0$ and, for every $h>0$,
\begin{equation}
 C_{2\lceil h/a\rceil,a}\longrightarrow
 V_\infty^*e^{-2hH_\infty}V_\infty,
 \qquad V_\infty^*V_\infty=G_\infty,
 \label{f16v60:fixed-time-convergence}
\end{equation}
where $H_\infty\ge0$ generates a strongly continuous semigroup.
Neither sufficient condition assumes a positive gap for $H_\infty$.
\end{proposition}
\begin{proof}
Splitting the spectral integral at $e^{-aM}$ gives
\[
 G_a-C_{2m,a}\preceq(1-e^{-2amM})G_a+F_a(M).
\]
First let $a\to0$ with $am\to0$, and then $M\to\infty$.
For the second assertion, for each fixed $h>0$ and eventually
$m_a\le\lceil h/a\rceil$, positivity of $\mathcal T_a$ gives
\[
 0\preceq G_a-C_{2m_a,a}
   \preceq G_a-C_{2\lceil h/a\rceil,a}.
\]
Pass to the finite-matrix limit.  Strong continuity makes the resulting
right-hand side tend to zero as $h\downarrow0$.
$G_a\to G_\infty\succ0$ converts this norm convergence into the claimed
relative order bound.  Restore any exact redundant columns by congruence.
\end{proof}

This proposition prevents a moving ultraviolet tail from being silently
removed by a mere declaration of source convergence.  The complete thermal
elementary packet excluded in V59 need not satisfy either condition.

\subsection{Quantitative suppression after physical time regularization}

The obstruction can be quantified without constructing a continuum limit.
Let $W$ be a centered source synthesis, $G_0=W^*W$, and set
\begin{equation}
 V_k=\mathcal T_a^kW,\qquad G_k=W^*\mathcal T_a^{2k}W,
 \qquad k\ge1.
 \label{f16v60:smoothed-bank}
\end{equation}
Let $A_k$ be its capture matrix for any interval containing the terminal
slice and its predecessor, or a convex average of such captures.

\begin{theorem}[A native smoothing upper bound with normalization retained]
\label{f16v60:smoothed-upper}
For every finite regulator,
\begin{equation}
 \boxed{\quad
 A_k\preceq b_k\cosh^2\gamma\,G_0,
 \qquad b_k=\frac{k^k}{(k+1)^{k+1}}\le\frac1{e k}.
 \quad}
 \label{f16v60:smoothed-upper-eq}
\end{equation}
In particular, for $k=\lceil\sigma/a\rceil$, $\sigma>0$,
\begin{equation}
 A_k\preceq\frac{a\cosh^2\gamma}{e\sigma}G_0.
 \label{f16v60:physical-smoothing-upper}
\end{equation}
If the \emph{additional} retention bound $G_0\preceq R_aG_k$ holds,
the relative capture is at most
$R_a a\cosh^2\gamma/(e\sigma)$.  This tends to zero whenever its
numerator does.  No uniform bound on $R_a$ is inferred from time smoothing.
\end{theorem}
\begin{proof}
The one-step comparison gives
\[
 A_k\preceq\cosh^2\gamma\,
 W^*\mathcal T_a^{2k}(I-\mathcal T_a^2)W.
\]
For $u\in[0,1]$, the maximum of $u^k(1-u)$ occurs at
$u=k/(k+1)$ and equals $b_k$.  Moreover
$k b_k=(k/(k+1))^{k+1}\le e^{-1}$.
Functional calculus proves the result.  All source coefficient maps and
normalizations pass through the inequalities by congruence.
\end{proof}

For a source family spectrally supported in $E_{a,M}$ there is also the
relative bound
\begin{equation}
 A\preceq\cosh^2\gamma(1-e^{-2aM})G
       \preceq2Ma\cosh^2\gamma\,G.
 \label{f16v60:low-energy-upper}
\end{equation}
Thus every uniformly bounded physical-energy band becomes invisible to
these fast conditional updates under \eqref{f16v60:fast-clock}.  There is
no implication in the opposite direction: exterior-predictable states
need not have low energy.

For comparison, if $H_a=-a^{-1}\log\mathcal T_a$ is taken on its spectral
domain, \eqref{f16v60:operator-upper} implies the quadratic-form inequality
\begin{equation}
 \frac{B_{S,I,t}}{2a\cosh^2\gamma}\preceq H_a
 \label{f16v60:rescaled-form}
\end{equation}
after the ground-state identification.  Indeed
$I-e^{-2aH_a}\preceq2aH_a$.  This specifies a possible rescaling of the
vanishing form; it does not establish a positive lower bound for the
rescaled form or identify it with a physical Hamiltonian.

The same no-go proof applies to independent connected ribbon components of
at most $w_a$ spatial links when the available contraction is
$\tanh(w_a\gamma_a)$ and $a e^{2w_a\gamma_a}\to0$.
Under a trajectory with $2<c<4$, it rules out $w_a=1$; the worst-boundary
estimate alone does not decide $w_a=2$.  Failure of this sufficient
fast-mixing condition is neither a capture lower bound nor evidence of
slow native motion.

\subsection{A massive slow sector has the same vanishing ribbon fraction}

V59's reversible exterior model also separates the present obstruction
from a statement about the physical spectrum.  Retain its notation:
$Y_t\in\{-1,1\}$ has transition $Q_r(y,y')=(1+r yy')/2$, and, given
its complete trajectory, the $X_t$ are independent with
$\Pr(X_t=x\mid Y_t=y)=(1+bxy)/2$, $0<b<1$.  Its positive reversible
transfer is $P=JQ_rJ^*$ with spectrum $1,r,0,0$.
For $f(y,x)=x$ and every $k\ge1$, direct conditional expectation gives
\begin{equation}
 P^kf=b r^k y,\qquad
 \|P^kf\|_{L^2(\pi)}^2=b^2r^{2k},\qquad
 \mathbb E\operatorname{Var}(P^kf\mid(Y_t)_t)=0.
 \label{f16v60:massive-fixture}
\end{equation}
Choose $r_a=e^{-\mu a}$ with fixed $\mu>0$ and
$k_a=\lceil\sigma/a\rceil$.  The entrance norm tends to
$b^2e^{-2\mu\sigma}>0$ and the surviving nonconstant physical energy is
exactly $\mu$, while the ribbon capture is identically zero.  Alternatively
$r_a=e^{-a^2}$ gives physical energy $a\to0$ with the same zero capture.
The equality in \eqref{f16v60:massive-fixture} is the new regularized-source
calculation; the reversible model itself is inherited from V59.  Neither
choice is a Wilson simulation.  Together they show that zero fast-variable
capture cannot distinguish a positive physical gap from a collapsing one.

\subsection{The source survives as an exterior predictor}

The vanishing variance has a constructive meaning.  Work in one finite
freed interval and use one common exterior sigma-field $\mathcal E$.
For a finite set of terminal sources $F_i=f_i(U_{t_i})$, put
$\widehat F_i=\mathbb E[F_i\mid\mathcal E]$.  The times need not coincide.
All variables are centered and square integrable.

\begin{proposition}[Exact preservation identity for mixed source Grams]
\label{f16v60:predictor-gram}
With $\Gamma_{ij}=\mathbb E\overline F_iF_j$ and
$\widehat\Gamma_{ij}=\mathbb E\overline{\widehat F_i}\widehat F_j$,
\begin{equation}
 \Gamma-\widehat\Gamma
   =\mathbb E\operatorname{Cov}((F_i)_i\mid\mathcal E)\succeq0,
 \qquad
 |\Gamma_{ij}-\widehat\Gamma_{ij}|
 \le\sqrt{A_iA_j},
 \label{f16v60:predictor-gram-eq}
\end{equation}
where $A_i=\mathbb E|F_i-\widehat F_i|^2$.
Consequently, when the individual captures tend to zero and the entrance
norms remain bounded, every such finite mixed two-point matrix is retained
in the limit by the exterior predictors.
\end{proposition}
\begin{proof}
Conditional expectation is an orthogonal projection in the common path
$L^2$ space.  Its residual is orthogonal to every $\widehat F_i$, giving
the equality.  The residual Gram is positive, and its two-by-two principal
minors give the displayed bound.  It is important here that the same
$\mathcal E$ is used in both factors; predictors from different exteriors
do not obey this exact Gram identity.
\end{proof}

This is preservation of Euclidean $L^2$ source correlations, not yet a
claim about an Osterwalder--Schrader source class.  A predictor using the
whole exterior may depend on both past and future.  The next construction
keeps its temporal support under explicit control.

\subsection{Gauge-invariant predictors in shrinking physical windows}

Let $f(U_t)$ be a bounded gauge-invariant slice function, with
$\|f\|_\infty\le M$.  Suppose it depends on $r$ members of $S$; all
other links on which it depends remain fixed when the ribbons are
conditioned.  At a finite spatial regulator $r$ is finite, but it may
grow with that regulator.  The case $r=0$ requires no elimination.
Fix an integer $R\ge1$ and free the interval
$I_{t,R}=\{t-R,\ldots,t+R\}$.

Define $\mathsf P^0_{S,R}f$ as follows.  For each of the $r$ relevant
ribbons, use its exact conditional density inside $I_{t,R}$, including
all four spatial staples at every time and all temporal plaquettes between
times in $I_{t,R}$.  Omit only the two temporal plaquettes connecting that
ribbon to the outside copies at $t-R-1$ and $t+R+1$.  Integrate the ribbon
variables against this normalized \emph{free-end} chain, taking their
product over the relevant spatial links, and evaluate $f$ at time $t$.
This is a deterministic function of the retained exterior links in this
window.  The operation is linear in $f$ and is a predictor, not a
modification of the native probability law.

\begin{theorem}[An explicit local-in-time exterior approximation]
\label{f16v60:local-predictor}
At every fixed regulator and every exterior configuration,
\begin{equation}
 \left|\mathbb E[f(U_t)\mid\mathcal E_{S,I_{t,R}}]
                   -\mathsf P^0_{S,R}f\right|
       \le4rM q_\gamma^{R+1}.
 \label{f16v60:predictor-boundary}
\end{equation}
Therefore
\begin{equation}
 \|f(U_t)-\mathsf P^0_{S,R}f\|_2
 \le\sqrt{\mathbb E\operatorname{Var}
                   (f(U_t)\mid\mathcal E_{S,I_{t,R}})}
                  +4rM q_\gamma^{R+1}.
 \label{f16v60:predictor-total}
\end{equation}
The predictor is gauge invariant.  Its spatial support is contained in the
original support together with the plaquette neighborhoods of its $r$
freed links, and its temporal support is inside $I_{t,R}$.  In particular,
no all-time boundary message is required to evaluate it.
\end{theorem}
\begin{proof}
For one ribbon, first fix the right boundary spin and vary the left one.
The conditional chain includes all fields and the right-boundary message.
By the total-variation contraction proved in V59, propagating from
$t-R-1$ to $t$ makes the diameter of the resulting central marginals at
most $q_\gamma^{R+1}$.  The reverse-time argument gives the same bound
on changing the right boundary with the left one fixed.  Thus any two
central marginals with fixed boundary spins have total-variation distance
at most $2q_\gamma^{R+1}$.

The free-end central law is a mixture of these fixed-boundary central
laws.  To see this without an approximation, adjoin two dummy boundary
spins with Haar measure and the usual temporal pair factors.  Integrating
a dummy spin gives the constant
$\int_{S^3}e^{\gamma x\cdot R y}\,d\sigma(y)$, independent of $x$.
After normalization this is precisely the free-end chain.  Disintegration
with respect to the dummy spins gives the asserted mixture; its mixing
law need not be uniform.  Hence the same $2q_\gamma^{R+1}$ bound compares
any actual central marginal with the free-end one.

The $r$ ribbon laws are conditionally independent.  The total-variation
distance between their two product central laws is at most the sum of the
$r$ one-ribbon distances.  For a bounded complex function the expectation
difference is bounded by twice its supremum norm times total variation.
This proves \eqref{f16v60:predictor-boundary}.  The triangle inequality in
$L^2$ and the conditional-variance identity give
\eqref{f16v60:predictor-total}.

Every plaquette retained in the free-end integral is an entire Wilson
plaquette.  Under a gauge transformation of the retained links, change
variables in the integrated links by that same vertex transformation.
Haar measure, the retained plaquette traces and $f$ are invariant.  The
normalizing integral transforms in the same way.  Thus the predictor is
gauge invariant.  Inspection of the retained plaquettes gives exactly the
stated support.  The omission of the two outside temporal plaquettes
serves only to define this predictor; no assertion of equality of the
free-end and native measures is used.
\end{proof}

\begin{proposition}[A shrinking-window choice with all size costs paid]
\label{f16v60:shrinking-window}
Assume $\gamma_a\to\infty$, \eqref{f16v60:strong-clock}, bounded entrance
norms, and relative short-time continuity for a finite source family.
For each of its columns let the number of relevant freed links and its
supremum bound obey $r_aM_a\le C a^{-p}$ for fixed $C,p$.
Choose
\begin{equation}
 L_a^\sharp=\max\{1,\,2\log(2+r_aM_a)+\gamma_a\},\qquad
 R_a=\left\lceil e^{2\gamma_a}L_a^\sharp\right\rceil.
 \label{f16v60:window-choice}
\end{equation}
Then $aR_a\to0$ and
\begin{equation}
 \|f_a(U_t)-\mathsf P^0_{S,R_a}f_a\|_2\longrightarrow0.
 \label{f16v60:local-predictor-limit}
\end{equation}
A common maximum of the columnwise sizes works for the complete finite
family and preserves declared exact linear source relations.
\end{proposition}
\begin{proof}
The polynomial size hypothesis and \eqref{f16v60:strong-clock} give
$L_a^\sharp=O(\gamma_a)$, hence
$aR_a\le a e^{2\gamma_a}O(\gamma_a)+a\to0$.
The bound $q_\gamma^R\le e^{-R e^{-2\gamma}}$ yields
\[
 4r_aM_a q_{\gamma_a}^{R_a+1}
 \le \frac{4r_aM_a}{(2+r_aM_a)^2}e^{-\gamma_a}\longrightarrow0.
\]
The interval has left depth $R_a\ge\lceil e^{2\gamma_a}\rceil$.
The no-go theorem therefore makes the conditional-variance term in
\eqref{f16v60:predictor-total} tend to zero.  Linearity gives exact
preservation of any zero relation.  For a finite family, columnwise $L^2$
convergence implies operator-norm convergence of its synthesis map.
\end{proof}

The polynomial bound is a stated sufficient condition, not a conclusion
about an unspecified source renormalization.  More generally the proof
requires $a e^{2\gamma_a}L_a^\sharp\to0$ and
$r_aM_a e^{-L_a^\sharp}\to0$.  Time smoothing by $P_a^k$ preserves an
$L^\infty$ bound at each regulator, but it need not preserve a small spatial
support.  No spatial locality is claimed for a time-smoothed writer unless
its own support or spatial tail has been controlled independently.

If the original sources are placed at a fixed positive physical distance
from a reflection plane, these shrinking-window predictors eventually
remain on that side of the plane.  Their $L^2$ approximation then also
controls each reflected two-point pairing by Cauchy--Schwarz and reflection
invariance.  This is a support-sensitive transfer of source pairings,
not the construction of a continuum Osterwalder--Schrader space.

\subsection{The exact retained marginal and its induced temporal response}

The source-preserving map must be accompanied by the \emph{actual} exterior
law.  In a finite open Euclidean cylinder, split the full Wilson action into
terms without a freed variable and the ribbon terms:
\begin{equation}
 S^{\rm W}(\xi,x)=S_{\rm ext}(\xi)+\sum_{e\in S}S_e(x_e;\xi),
 \qquad Z_e(\xi)=\int e^{-S_e(x_e;\xi)}\,dH(x_e).
 \label{f16v60:exact-decimation}
\end{equation}
Boundary terms are retained here; the free-end convention used to construct
a predictor does not replace them.  Fubini gives the exact marginal
\begin{equation}
 d\mu_{\rm ext}(\xi)=Z^{-1}e^{-S_{\rm eff}(\xi)}dH(\xi),
 \qquad S_{\rm eff}=S_{\rm ext}-\sum_{e\in S}\log Z_e.
 \label{f16v60:effective-action}
\end{equation}
This factorization uses local Wilson boundary terms, or a finite freed
interval with its outside links fixed.  It is not a factorization of an
arbitrary nonlocal endpoint density: a finite-cylinder Perron endpoint
weight, for example, may not be absorbed into independent $Z_e$ factors.
Such a weight must remain explicit or be reached through the fixed-regulator
cylinder limit.  For the stationary vacuum, one may condition on the links
outside a finite freed interval first, or pass from ordinary open cylinders
to their fixed-regulator limit.  No Perron scalar is replaced and no product
exterior law is postulated.

For a reflection-positive full cylinder or stationary vacuum law, if the
set of removed variables is reflection invariant, the exact marginal is
reflection positive on the retained positive-time algebra:
for such an $F$, $\mathbb E_{\rm ext}\overline{F(\theta\xi)}F(\xi)$
is the identical nonnegative expectation in the full law.  Gauge invariance
also follows by a change of integrated variables.  A marginal process of
retained slice variables need not be Markov, and
\eqref{f16v60:effective-action} does not assert that it has a nearest-slice
Wilson action.

\begin{theorem}[Exact signed response and a uniform time-tail bound]
\label{f16v60:response-tail}
For one ribbon write $\Psi=-\log Z_e$.  For two real smooth exterior
parameters $u,v$, with all other exterior data fixed,
\begin{equation}
 \partial_u\partial_v\Psi
   =\mathbb E_e[\partial_u\partial_vS_e]
        -\operatorname{Cov}_e(\partial_uS_e,\partial_vS_e).
 \label{f16v60:signed-response}
\end{equation}
Suppose the two score functions depend on ribbon times in disjoint
intervals $[s_0,s_1]$ and $[t_0,t_1]$ with $t_0>s_1$, and their
supremum bounds are $L_u,L_v$.  Then
\begin{equation}
 |\operatorname{Cov}_e(\partial_uS_e,\partial_vS_e)|
        \le L_uL_v q_\gamma^{t_0-s_1}.
 \label{f16v60:score-covariance}
\end{equation}
The constants are uniform over all exterior fields and finite cylinder
lengths.  In particular, for time-indexed local parameters whose scores are
supported in $[s,s+1]$, with bounds $L$, and whose mixed contact derivative
vanishes for $|s-t|\ge2$, the Hessian tail retaining only entries
$|s-t|\ge R\ge2$ has Euclidean matrix norm at most
\begin{equation}
 \|\mathsf H^{\rm tail}_R\|_{2\to2}
       \le\frac{2L^2 q_\gamma^{R-1}}{1-q_\gamma}.
 \label{f16v60:Hessian-tail}
\end{equation}
For $q_\gamma=0$ this tail is zero.
\end{theorem}
\begin{proof}
Compact finite-dimensional integration permits differentiation under the
integral.  The first derivative is $\partial_u\Psi=\mathbb E_e\partial_uS_e$;
differentiating a normalized expectation gives
\eqref{f16v60:signed-response}, including its negative covariance sign.
For the covariance estimate, center both scores.  Conditional expectations
onto $X_{s_1}$ and $X_{t_0}$ have norms no larger than the corresponding
$L^2$ norms, each at most its stated supremum bound.  The Markov property
reduces the covariance to the pairing of these two endpoint functions
through the conditional transition product from $s_1$ to $t_0$.
V59's centered product bound gives \eqref{f16v60:score-covariance}.

For indices separated by $|s-t|\ge2$ there is no contact term, and the
score supports give the bound $L^2q_\gamma^{|s-t|-1}$.  Each row and each
column of the tail has absolute sum at most
$2L^2\sum_{d=R}^\infty q_\gamma^{d-1}$.
The Schur row/column bound proves \eqref{f16v60:Hessian-tail}.
\end{proof}

A derivative in a fixed normalized Lie-algebra direction of a Wilson
plaquette trace has magnitude bounded by that direction's matrix norm.
Thus, for a declared exterior variation incident to at most $N_0$ link occurrences in ribbon
plaquette terms, a valid score bound is $L\le N_0\gamma$ when those
norms are at most one.  Shared contact plaquettes must still be counted;
they are not part of the separated-time tail estimate.  If several exterior
coordinates are marked at one time, their incidence multiplicity must
likewise be retained rather than hidden in a volume-independent constant.

For one such bounded-incidence marked ribbon and fixed $\delta>0$, take
$R_a=\max\{2,\lceil\delta/a\rceil\}$.  Since
$(1-q_\gamma)^{-1}\le e^{2\gamma}$, the tail is bounded by
\begin{equation}
 2N_0^2\gamma_a^2 e^{2\gamma_a}
   \exp[-(R_a-1)e^{-2\gamma_a}],
 \label{f16v60:physical-response-tail}
\end{equation}
which tends to zero under \eqref{f16v60:strong-clock}.  This is a proved
local two-score response bound.  It is not a bound on all higher cumulants,
on the full norm of an extensive effective action, or on a continuum
spectral gap.

Induced couplings cannot be dropped even where the original conditional
ribbons were independent.  Already a single Haar link gives
\[
 Z(u,v)=\int_{S^3}e^{u\,p\cdot x+v\,r\cdot x}\,d\sigma(x),\qquad
 \left.\partial_u\partial_v[-\log Z]\right|_{u=v=0}
                 =-\tfrac14 p\cdot r.
\]
This is the exact second moment $\int xx^{\mathsf T}d\sigma=I_4/4$.
Distinct retained marks can therefore interact through one eliminated
link.  After a color has been integrated out, another original color need
not have an independent-ribbon conditional law under the new marginal.
Repeating the original twelve-color argument on that marginal without
recomputing its induced interactions would be invalid.

\subsection{The revised upstream task and the noncircular boundary}

The impossible target is now retired: under
\eqref{f16v60:fast-clock}, an order-one independent-ribbon fraction cannot
be common to a continuous physical source core.  This does not invalidate
V59's finite-regulator lower bounds for unsmeared elementary writers,
whose ultraviolet continuity was already excluded there.  It corrects
the proposed use of those bounds after physical source regularization.

The constructive replacement is an exact exterior reduction, not another
attempt to force the fast component to contain the surviving field.
For a declared source bank, the source loss is bounded by
\eqref{f16v60:predictor-total}; the retained measure is
\eqref{f16v60:effective-action}; and its differentiated local interaction
must retain both terms of \eqref{f16v60:signed-response}.
The new physical information is carried by this retained law.

The next concrete calculation is a source-marked retained slab after
\emph{one} actual color elimination: keep the native exterior links, the
plaquette contact terms, and the induced mixed-score block, and establish
a signed estimate for its finite-physical-scale source correlations.
The response-tail theorem pays only the part separated in time; it does
not pay the local negative covariance or justify a second independent
elimination.  A useful next result must control that signed local block or
another specified native retained correlation, rather than repackage the
vanishing conditional capture.  Noncollapse, physical-time compatibility,
spatial source support when required, and dense-core common constants
remain separate obligations.

\subsection{Verification, append-only preservation, and the result ledger}

The companion
\begin{center}\small
\path{verify_ym_f16_v60_exterior_reduction.py}
\end{center}
checks finite-state reversible source/subdelay inequalities, exact rational
smoothing extrema and source relations, spectral continuity diagnostics,
regularized massive and gap-collapsing exterior fixtures, finite-chain
boundary replacement and free-end predictors, covariance response and
Hessian-tail inequalities, native quaternion gauge-covariance identities,
and the stated asymptotic scale comparisons.  Exact algebra, floating-point
matrix diagnostics and finite-chain enumeration are reported separately.
They are not samples or interval enclosures of the full Wilson vacuum.
The analytical arguments above are the proofs; the tests are independent
checks on their identities, constants and limiting claims.

Removing this single delimited appendix and reversing the displayed title
change from V60 to V59 recovers the uploaded predecessor byte for byte.
The verification report also records hashes of the attached comparison
inputs and the duplicate-copy comparisons.  Those files are not modified.
No new Lean formalization or native many-link numerical certificate is
claimed.

\begin{record}[V60 result ledger]
\label{f16v60:ledger}
\emph{Proved here:} the fixed-exterior subdelay comparison and its operator
form; vanishing independent-ribbon capture on a relatively continuous
physical source family under the stated fast clock; spectral-tightness and
fixed-time-limit sufficient continuity conditions; exact quantitative
suppression after transfer smoothing, with the retention cost explicit;
low-energy and rescaled-form comparisons; the massive regularized
exterior fixture; the exact mixed predictor-Gram identity; a gauge-invariant
free-end predictor with its finite-regulator error and shrinking-window
limit under explicit size hypotheses; exact retained-marginal preservation
of gauge invariance and reflection positivity at their stated scopes;
and the signed ribbon-response identity with its boundary-uniform
separated-time Hessian tail.

\emph{Inherited or conditional:} the finite-regulator Wilson transfer and
vacuum; V59's conditional ribbon contraction; the physical trajectory
where used; relative physical source continuity or its stated sufficient
conditions; the source-size and noncollapse hypotheses; and reflection
positivity of the original law for the marginal statement.

\emph{Ruled out at the stated scope:} V59's proposed positive uniform
independent-ribbon capture fraction on a genuine continuous physical core
under \eqref{f16v60:fast-clock}; temporal buffering alone as a remedy;
and unmodified reuse of the independent-color law after marginalization.

\emph{Still unproved:} a suitable nontrivial physically regularized native
source limit; a coercive signed estimate for the retained local interaction;
physical-scale mixing of that retained source family with common dense-core
constants; the required nontrivial interacting continuum construction;
and the continuum Yang--Mills mass gap.
\end{record}

\begin{firewall}
This continuation goes upstream because the previous positive-fraction
target is incompatible with its own fast conditional clock and a continuous
physical core.  The replacement retains, rather than discards, the exterior
that carries the surviving source.  Short temporal range of an induced
response is not a spectral gap, and a local predictor is not a new physical
transfer.  The mathematical body of V1--V59 is preserved unchanged;
the correction to V59's future target is stated explicitly in this appendix.
\end{firewall}
% END F16 V60 APPEND

% BEGIN F16 V61 APPEND -- 2026-09-18
% Append-only continuation of the uploaded post-fifteenth-archive V60.
% The preceding mathematical body is unchanged.

\clearpage
\section[V61: exact retained-ribbon curvature and transverse Ward cancellation]{V61: exact retained-ribbon curvature\\and transverse Ward cancellation}
\label{f16v61:context}

V60 isolated the contact-minus-covariance term of the actual retained
Wilson action but did not control its local sign.  This continuation
computes that term without replacing the compact ribbon by a Gaussian.
There are two different conclusions.  At every exterior configuration,
the exact ribbon free energy controls a specified collection of retained
Wilson-loop deficits.  At the coherent exterior, its complete second
variation is a positive weighted difference form, with an exact rotational
null space.  The first conclusion is nonlinear and global in the exterior;
the second is a Hessian statement at the coherent configuration only.
They must not be conflated: a native maximally frustrated one-link example
below has a negative second variation.

The principal new input is a positive Haar-moment expansion, followed by a
clipping argument that retains configurations with a small staple sum.
It gives an all-exterior bound for a whole temporally coupled ribbon,
with constants independent of its length.  A rotational integration by
parts then pays the signed local Hessian exactly.  These are statements
about the retained interaction, not physical-time contraction of the
retained vacuum.  V60's obstruction to a positive microscopic-ribbon
capture fraction remains in force.

\subsection{Dependency and novelty boundary}

We retain V59--V60's literal Wilson plaquettes, independent spatial link
mask, finite freed time interval, and local boundary convention.  In
particular, a nonlocal Perron endpoint density is not inserted into a
product of ribbon integrals.  V52's proved Bessel-ratio barrier supplies
one elementary one-link estimate used below.  No cubic exponential-moment
claim from the parallel V67 branch is used.

The f15 archive, V66--V67, already treated an independent \emph{single-link}
radial integral and smooth retained comparisons.  Those results are not
reproved as new.  The present object keeps all temporal couplings of a
ribbon, allows arbitrary retained temporal transports, and obtains both
an all-exterior generated-loop lower bound and an exact all-order compact
Hessian at its coherent exterior.  The f14 gauge-Ward calculations and
f15 covariant dressing costs concern different operators; they are not
substitutes for the conditional covariance computed here.

For orientation, the sign-coordinate representation used in the proof
is the elementary embedding underlying continuous-spin cluster methods;
see U.~Wolff, \emph{Collective Monte Carlo Updating for Spin Systems},
Phys.\ Rev.\ Lett.\ \textbf{62} (1989), 361--364,
\href{https://doi.org/10.1103/PhysRevLett.62.361}{doi:10.1103/PhysRevLett.62.361}.
The modified-Bessel recurrence convention agrees with
\href{https://dlmf.nist.gov/10.29}{NIST DLMF, Section 10.29}.
All inequalities needed below are proved here or explicitly reduced to
V52; neither reference imports a Wilson-vacuum mixing result.

\subsection{The exact finite ribbon and its retained geometric data}

Identify a Wilson link with a unit quaternion $x\in S^3\subset\mathbb R^4$,
and write $w(U)=\tfrac12\operatorname{Re}\Tr U$.  The convention is
$x\cdot p=w(UP^{-1})$.  Free one spatial link at times $t=1,\ldots,T$,
keeping its copies at times $0,T+1$ and every other link fixed.  Its
conditional partition function, with exterior-independent Wilson
constants omitted, is exactly
\begin{equation}
 Z(p,R)=\int\exp\left\{\gamma\sum_{t=1}^T\sum_{\alpha=1}^{n_t}
                   p_{t\alpha}\cdot x_t
             +\gamma\sum_{t=1}^{T-1}x_t\cdot R_t x_{t+1}\right\}
                \prod_{t=1}^T d\sigma(x_t),
 \label{f16v61:ribbon}
\end{equation}
where $p_{t\alpha}\in S^3$, $R_t\in\mathrm{SO}(4)$, and $d\sigma$ is
normalized Haar measure.  There are four spatial staples per time and
one additional field from each frozen temporal endpoint:
\begin{equation}
 n_t=4+\mathbf1_{t=1}+\mathbf1_{t=T},\qquad
 d_t=\mathbf1_{t>1}+\mathbf1_{t<T},\qquad n_t+d_t=6.
 \label{f16v61:counts}
\end{equation}
The indicators add even when $T=1$, so then $n_1=6$.  The two retained
temporal links in a temporal plaquette act by left and right quaternion
multiplication, giving the displayed orthogonal $R_t$.  Each $p_{t\alpha}$
is a three-link retained path with the same endpoints as the freed link.
Thus these are literal Wilson data, not independent random staple fields.

Put
\begin{equation}
 r_t=\left|\sum_\alpha p_{t\alpha}\right|,\qquad
 F_t=n_t-r_t,\qquad
 \Psi(p,R)=-\log Z(p,R),\qquad \Psi_0=-\log Z(e_0,I),
 \label{f16v61:frustration}
\end{equation}
where $e_0=(1,0,0,0)$ and every slot in $Z(e_0,I)$ equals $e_0$.
All $F_t$ are nonnegative and gauge invariant.  No physical spacing is
changed in defining these quantities.

\subsection{Alignment domination from positive compact moments}

\begin{proposition}[Alignment bound, including arbitrary orthogonal transports]
\label{f16v61:alignment}
On a finite graph, let the edge factors be
$\exp(J_{ij}x_i\cdot R_{ij}x_j)$ with $J_{ij}\ge0$ and
$R_{ij}\in\mathrm O(4)$, and the field factors be $\exp(h_i\cdot x_i)$.
Then
\begin{equation}
 0<Z((h_i),(R_{ij}))
     \le Z((|h_i|e_0),(I)).
 \label{f16v61:alignment-eq}
\end{equation}
For aligned fields $h_i=H_i e_0$, $H_i\ge0$, and identity transports,
$Z$ is nondecreasing in every $H_i$.
\end{proposition}
\begin{proof}
The odd Haar moments vanish, while
\begin{equation}
 \int x_{a_1}\cdots x_{a_{2k}}\,d\sigma(x)
   =\frac{\displaystyle\sum_{\text{pairings }\pi}
                \prod_{\{u,v\}\in\pi}\delta_{a_u a_v}}
          {4\cdot6\cdots(2k+2)}.
 \label{f16v61:Haar-moments}
\end{equation}
For example, write a four-dimensional standard Gaussian as $RX$, with
$X$ uniform on $S^3$ and independent of $R$, compare its Gaussian pairing
moments, and use $\mathbb E R^{2k}=4\cdot6\cdots(2k+2)$.
This also proves the normalization in \eqref{f16v61:Haar-moments}.

Expand every exponential.  For each multi-index of edge and field powers,
apply \eqref{f16v61:Haar-moments} at every vertex.  Each pairing diagram
consists of closed index lines and open lines with field endpoints.
A closed line contributes the trace of a product of orthogonal matrices,
whose absolute value is at most $4$.  An open line contributes
$h_i^{\mathsf T}Oh_j$, whose absolute value is at most $|h_i||h_j|$.
All remaining coefficients are nonnegative.  Replacing every transport
by $I$ and every field by its norm times $e_0$ attains these upper bounds
in the same diagrams.  Termwise domination proves
\eqref{f16v61:alignment-eq}.  There is no conditional convergence issue:
the dominating aligned expansion sums to a finite compact integral.

For the monotonicity claim, fix the absolute value $a_i=|x_i^0|$ and the
other coordinates of each spin.  Conditional on these data, the signs
$\varepsilon_i$ of $x_i^0$ have an Ising law with couplings
$J_{ij}a_i a_j\ge0$ and fields $H_i a_i\ge0$.  Its one-spin expectations
are nonnegative.  To check this directly, expand each bond as
$\cosh K(1+\varepsilon_i\varepsilon_j\tanh K)$ and each field as
$\cosh b(1+\varepsilon_i\tanh b)$.  The sign sums contributing to the
numerator of $\mathbb E\varepsilon_i$ have nonnegative coefficients;
all surviving parity terms are positive.  Hence
$\partial_{H_i}\log Z=\mathbb E x_i^0\ge0$ after averaging the fixed
absolute-value data.
\end{proof}

Only the proof of the monotonicity clause uses an aligned field law.
The first clause is an all-exterior inequality, including frustrated
orthogonal transports.  On an open ribbon the transports may alternatively
be removed exactly: take $O_1=I$, $O_{t+1}=R_t^{\mathsf T}O_t$, and change
variables $x_t=O_t z_t$.  The fields become
$\widetilde p_{t\alpha}=O_t^{\mathsf T}p_{t\alpha}$.  This is a change
of integrated variables, not a gauge fixing of the retained measure.
A periodic ribbon has an additional holonomy and is not silently replaced
by this open-interval convention.

\subsection{An all-exterior retained-curvature lower bound}

Let $a_4(u)=I_2(u)/I_1(u)$ for $u>0$.  The mean of the one-link law
proportional to $e^{h\cdot x}$ is
$a_4(|h|)h/|h|$.  V52's lower ratio barrier implies
\begin{equation}
 \chi_4(u):=\frac{a_4(u)}u\ge\frac1{4+u},\qquad
 \chi_4(0)=\frac14.
 \label{f16v61:radial-lower}
\end{equation}
Indeed the V52 denominator for $I_2/I_1$ is
$3/2+\sqrt{u^2+(5/2)^2}\le u+4$.  The value at zero follows either
from the Haar second moment or the Bessel series.  This is a consequence
of an inherited global bound, not a new large-$u$ asymptotic assumption.

\begin{theorem}[Global ribbon frustration with no discarded small-staple set]
\label{f16v61:global-curvature}
For every $T\ge1$, $\gamma>0$, and every set of retained data in
\eqref{f16v61:ribbon},
\begin{equation}
 \boxed{\qquad
 \Psi(p,R)-\Psi_0
   \ge \lambda_\gamma\sum_{t=1}^T F_t,
 \qquad
 \lambda_\gamma=\frac{\gamma^2}{4(4+6\gamma)}.
 \qquad}
 \label{f16v61:global-curvature-eq}
\end{equation}
If $r_t\ge3$ at every time, $4\lambda_\gamma$ is valid in place of
$\lambda_\gamma$.  The constants do not depend on ribbon length,
spatial volume, or exterior fields.  In particular, zero staple sums
are included in the first inequality.
\end{theorem}
\begin{proof}
Combine the slots at time $t$ into $\gamma\sum_\alpha p_{t\alpha}$.
Proposition~\ref{f16v61:alignment} gives
\[
 Z(p,R)\le Z((\gamma r_t e_0)_t,I)
          \le Z((\gamma\max\{3,r_t\}e_0)_t,I).
\]
The second inequality is the proved monotonicity, not a deletion of a
bad event.  Interpolate each aligned amplitude from
$\widetilde r_t=\max\{3,r_t\}$ to $n_t$.  At any point of this
interpolation, conditioning on all spins except $x_t$ gives the field
\[
 h_t=\gamma\bigl(u_t e_0+\sum_{s\sim t}x_s\bigr),
 \qquad 3\le u_t\le n_t.
\]
It obeys $h_t^0\ge\gamma(3-d_t)\ge\gamma$ and
$|h_t|\le\gamma(n_t+d_t)=6\gamma$.  Therefore
\eqref{f16v61:radial-lower} yields
\[
 \mathbb E[x_t^0\mid(x_s)_{s\ne t}]
       =\chi_4(|h_t|)h_t^0\ge\frac\gamma{4+6\gamma}.
\]
The logarithmic derivative in $u_t$ is $\gamma\mathbb E x_t^0$.
Integration along the amplitude segment proves
\begin{equation}
 \log Z((\gamma n_t e_0)_t,I)
   -\log Z((\gamma\widetilde r_t e_0)_t,I)
 \ge\frac{\gamma^2}{4+6\gamma}\sum_t(n_t-\widetilde r_t).
 \label{f16v61:clipped-estimate}
\end{equation}
For $4\le n\le6$ and $0\le r\le n$,
$n-\max\{3,r\}\ge(n-r)/4$: it is equality in the needed bound when
$n=4,r=0$, and follows separately on $r\le3$ and $r\ge3$.
This proves \eqref{f16v61:global-curvature-eq}.  When all $r_t\ge3$,
there is no clipping loss, giving the stronger constant.
\end{proof}

The coefficient is of order $\gamma$, not exponentially small:
$\lambda_\gamma/\gamma\to1/24$ as $\gamma\to\infty$.
This does not assert an order-$\gamma$ physical spectral gap.  The
quantity controlled is the specified retained curvature, which has
coherent and gauge directions and is measured in a finite-regulator action.

\subsection{The generated terms are literal retained Wilson loops}

For slots at the same time, define the retained closed path
$C_{t,\alpha\beta}=P_{t\alpha}P_{t\beta}^{-1}$.  Its length is at
most six, with any backtracking cancellation understood.  Then
\begin{equation}
 p_{t\alpha}\cdot p_{t\beta}=w(C_{t,\alpha\beta}),\qquad
 F_t=\frac{2\sum_{\alpha<\beta}(1-w(C_{t,\alpha\beta}))}{n_t+r_t}
 \ge\frac1{n_t}\sum_{\alpha<\beta}(1-w(C_{t,\alpha\beta})).
 \label{f16v61:generated-loops}
\end{equation}
These identities remain valid for the temporal-endpoint slots.  All slot
paths at one site transform by the same endpoint rotations, so neither
the inner products nor $F_t$ depend on the frame used to untwist a ribbon.

For an independent spatial mask $S$, the actual action decomposes exactly
as in V60:
\[
 S_{\rm eff}(\xi)=S_{\rm ext}(\xi)+\sum_{e\in S}\Psi_e(\xi)
                    +\text{exterior-independent constant}.
\]
Each plaquette without a freed variable occurs once in $S_{\rm ext}$;
each other plaquette belongs to one ribbon.  For flat prescribed boundary
links, or when the boundary links are included among $\xi$, write $\mathbf1$
for the all-identity exterior.  The theorem gives the native bound
\begin{equation}
 \boxed{\begin{aligned}
 S_{\rm eff}(\xi)-S_{\rm eff}(\mathbf1)
 \;&\ge\gamma\sum_{p\text{ retained}}(1-w(U_p))\\
 &\quad+\lambda_\gamma\sum_{e\in S}\sum_{t=1}^T\frac1{n_t}
               \sum_{\alpha<\beta}(1-w(C_{e,t,\alpha\beta})).
 \end{aligned}}
 \label{f16v61:native-curvature}
\end{equation}
For incompatible fixed nonflat boundary data, the comparison to an
admissible all-identity exterior is not asserted; the individual ribbon
bound relative to its coherent reference still holds as stated.

Equation~\eqref{f16v61:native-curvature} does not replace $S_{\rm eff}$
by its right-hand side.  The temporal and higher induced interactions
remain in the exact action.  In particular, a pointwise lower action
bound is not a normalized-measure comparison, a correlation inequality
between the two measures, or permission to perform a second independent
color elimination.  No partition-function ratio has been discarded.

\subsection{The transverse Ward identity pays the local negative covariance}

Return to the coherent reference $p_{t\alpha}=e_0$, $R_t=I$, and let
$\langle\cdot\rangle_0$ denote its complete compact ribbon law.  Set
\begin{equation}
 m_t=\langle x_t^0\rangle_0,\qquad
 c_{ts}=\langle x_t^i x_s^i\rangle_0\quad(i\in\{1,2,3\}).
 \label{f16v61:mc}
\end{equation}
The value of $c_{ts}$ is independent of $i$, the transverse means vanish,
and different transverse components have zero covariance.

\begin{theorem}[Exact positive difference form for the full compact Hessian]
\label{f16v61:Ward-Hessian}
For $\gamma>0$ the matrix $C=(c_{ts})$ is positive semidefinite, every
$c_{ts}$ is strictly positive, and
\begin{equation}
 m_t=\gamma\sum_s n_s c_{ts}.
 \label{f16v61:Ward}
\end{equation}
Let $p_{t\alpha}(\varepsilon)\in S^3$ be smooth with
$p_{t\alpha}(0)=e_0$ and transverse velocity
$b_{t\alpha}=p'_{t\alpha}(0)\in\mathbb R^3$.  With transports fixed at
$I$, the exact second variation of $\Psi$ is
\begin{align}
 Q(b)
 &=\gamma\sum_{t,\alpha}m_t|b_{t\alpha}|^2
       -\gamma^2\sum_{t,s}c_{ts}
                  \Big(\sum_\alpha b_{t\alpha}\Big)\cdot
                  \Big(\sum_\beta b_{s\beta}\Big)
 \label{f16v61:contact-minus-cov}\\
 &=\boxed{\frac{\gamma^2}{2}\sum_{t,\alpha,s,\beta}
                  c_{ts}|b_{t\alpha}-b_{s\beta}|^2}\ \ge0.
 \label{f16v61:difference-form}
\end{align}
Its kernel consists exactly of slot velocities that are all equal.
Thus the kernel has dimension three in the unrestricted slot space.
\end{theorem}
\begin{proof}
Simultaneous orthogonal rotations of the transverse three coordinates
preserve the law, proving the stated component symmetries.  Fix one
transverse component $i$.  Rotate every spin in the $(0,i)$ plane with
generator $Lx^i=x^0$, $Lx^0=-x^i$.  Haar measure and pair interactions
are invariant, whereas the derivative of the field exponent is
$-\gamma\sum_s n_s x_s^i$.  Integration of $L(x_t^i\,d\mu_0)$ gives
\eqref{f16v61:Ward}, including its coefficient and sign.

For positivity of individual entries, condition on $a_u=|x_u^i|$ and all
other coordinates.  The remaining signs form a zero-field Ising chain
with bonds $\gamma a_u a_{u+1}$.  Consequently, for $t<s$,
\begin{equation}
 c_{ts}=\left\langle a_ta_s\prod_{u=t}^{s-1}
                  \tanh(\gamma a_u a_{u+1})\right\rangle_0>0.
 \label{f16v61:sign-path}
\end{equation}
The formula follows by successively summing a free-end zero-field chain;
its two-spin correlation is the product of the intervening bond tangents.
The integrand is positive almost everywhere for $\gamma>0$.
Also $c_{tt}=\langle (x_t^i)^2\rangle_0>0$, and $C$ is a covariance
matrix, so $C\succeq0$.

The first derivative of $\Psi$ is zero because the spin mean is parallel
to $e_0$.  The constraint $|p_{t\alpha}(\varepsilon)|=1$ gives
$e_0\cdot p''_{t\alpha}(0)=-|b_{t\alpha}|^2$.  The signed derivative
identity in V60 now gives exactly \eqref{f16v61:contact-minus-cov}.
Substituting \eqref{f16v61:Ward} and symmetrizing proves
\eqref{f16v61:difference-form}.  Strict positivity of all $c_{ts}$
identifies its kernel.  In particular, the contact term is not estimated
separately from a supposedly negligible covariance term.
\end{proof}

The constant-velocity kernel is necessary: a simultaneous rotation of
all field slots leaves the integral unchanged.  The negative covariance
pays exactly the contact energy in that direction.  Erasing it would
manufacture a spurious restoring term.

\subsection{Quantitative susceptibility and a paid neighboring-time block}

Write $\bar b_t=n_t^{-1}\sum_\alpha b_{t\alpha}$ and
$d_{t\alpha}=b_{t\alpha}-\bar b_t$.  The preceding theorem gives the
\emph{exact} splitting
\begin{equation}
 Q(b)=\gamma\sum_t m_t\sum_\alpha|d_{t\alpha}|^2
       +\frac{\gamma^2}{2}\sum_{t,s}n_tn_s c_{ts}
                           |\bar b_t-\bar b_s|^2.
 \label{f16v61:contrast-split}
\end{equation}
In particular, within-time zero-sum slot contrasts have no covariance
score: their Hessian is precisely the displayed contact contribution.

\begin{proposition}[Length-independent compact bounds]
\label{f16v61:quantitative}
Define
\begin{equation}
 \kappa_\gamma=\frac{3}{4(4+6\gamma)}.
 \label{f16v61:kappa}
\end{equation}
At the coherent exterior,
\begin{gather}
 \frac{2\gamma}{4+6\gamma}\le m_t\le1,\qquad
 c_{tt}\ge\kappa_\gamma,\qquad
 c_{t,t+1}\ge\frac{\gamma}{1+\gamma}\kappa_\gamma^2,
 \label{f16v61:entry-bounds}\\
 0\preceq C\preceq\frac1\gamma
              \operatorname{diag}\left(\frac{m_t}{n_t}\right)
                \preceq\frac1{4\gamma}I.
 \label{f16v61:susceptibility}
\end{gather}
For every slot family,
\begin{align}
 Q(b)&\le\gamma\sum_{t,\alpha}|b_{t\alpha}|^2,
 \label{f16v61:upper-Q}\\
 Q(b)&\ge\frac{2\gamma^2}{4+6\gamma}
                      \sum_{t,\alpha}|d_{t\alpha}|^2
       +\frac{\gamma^3\kappa_\gamma^2}{1+\gamma}
                      \sum_{t=1}^{T-1}n_tn_{t+1}
                                  |\bar b_t-\bar b_{t+1}|^2.
 \label{f16v61:lower-Q}
\end{align}
These are finite-ribbon estimates, without a volume- or length-dependent
constant hidden in their coefficients.
\end{proposition}
\begin{proof}
In the coherent law, conditioning on the other spins gives
$h_t=\gamma(n_t e_0+\sum_{s\sim t}x_s)$.  Hence
$h_t^0\ge\gamma(n_t-d_t)\ge2\gamma$ and $|h_t|\le6\gamma$.
Equation~\eqref{f16v61:radial-lower} proves the bound on $m_t$.

The covariance of a one-link law has transverse eigenvalue
$\chi_4(|h_t|)$ relative to its field direction, and a nonnegative
longitudinal eigenvalue.  This follows by differentiating its radial
log partition function.  The field lies in the ball centered at
$n_t\gamma e_0$ of radius $d_t\gamma$.  Its angle $\theta$ from $e_0$
satisfies $\sin\theta\le d_t/n_t\le1/2$: intersect a ray from the
origin with that ball, or minimize its distance from the center.
For every fixed transverse unit vector $v\perp e_0$, its projection onto
$h_t^\perp$ therefore has squared length at least $3/4$.
Its conditional variance is at least
$3\chi_4(|h_t|)/4\ge\kappa_\gamma$.  Conditional second moments are
at least their variances, proving the bound on $c_{tt}$.

For $u\ge0$, $\tanh u\ge u/(1+u)$, since
$(1+u)\sinh u-u\cosh u=\sinh u-u e^{-u}\ge0$.
As $a_ta_{t+1}\le1$, \eqref{f16v61:sign-path} gives
\[
 c_{t,t+1}\ge\frac\gamma{1+\gamma}
             \langle (x_t^i)^2(x_{t+1}^i)^2\rangle_0
          \ge\frac\gamma{1+\gamma}\kappa_\gamma^2.
\]
For the last step condition on all spins except $x_t$, then use the
same uniform second-moment lower bound at $t+1$.

Apply $Q\ge0$ to $b_{t\alpha}=v_t/n_t$ in one transverse direction.
Equation~\eqref{f16v61:contact-minus-cov} gives
\eqref{f16v61:susceptibility}.  Its positive covariance subtraction
also proves \eqref{f16v61:upper-Q}.  Finally use
\eqref{f16v61:contrast-split}, keep only the two symmetric
neighboring-time entries in its second term, and insert the preceding
lower bounds.  The retained neighboring terms and the contrast term
are distinct parts of that exact splitting; no energy is counted twice.
\end{proof}

The susceptibility norm in \eqref{f16v61:susceptibility} is of order
$\gamma^{-1}$ rather than a worst-boundary exponential bound.  Its scope
is the coherent reference.  Also the neighboring-time term in
\eqref{f16v61:lower-Q} is a discrete gradient form, not a uniform
positive multiple of the squared norm of all time profiles.
A Poincare estimate on a growing interval would still carry its actual
length dependence.  It cannot be reinterpreted as a physical mass term.

\subsection{Pullback to actual retained links and source coefficient maps}

Let $\xi(\varepsilon)$ be any smooth path of retained Wilson links with
$\xi(0)=\mathbf1$.  Untwist each open ribbon as above and define
\begin{equation}
 b_{e,t\alpha}
   =\left.\frac{d}{d\varepsilon}
             \bigl(O_{e,t}(\varepsilon)^{\mathsf T}
                         p_{e,t\alpha}(\varepsilon)\bigr)
                                      \right|_{\varepsilon=0}.
 \label{f16v61:pullback}
\end{equation}
These are the actual derivatives of retained path products and temporal
transports; they are not independent slot coordinates chosen after the
fact.  The vanishing first derivative at the coherent point ensures that
second derivatives of this change of parameters do not add an unrecorded
term.  If $A$ denotes the retained link velocity in unit-quaternion tangent
normalization, the exact native identity is
\begin{equation}
 D^2S_{\rm eff}(\mathbf1)[A,A]
       =\gamma\sum_{p\text{ retained}}|(dA)_p|^2
                        +\sum_{e\in S}Q_e(b_e).
 \label{f16v61:native-Hessian}
\end{equation}
Here $(dA)_p$ is the oriented linearized plaquette sum.  The first term
uses the expansion
\[
 1-w(U_p(\varepsilon))
       =\tfrac12\varepsilon^2|(dA)_p|^2+O(\varepsilon^3).
\]
Gauge velocities give zero in the full expression, as follows either
from exact gauge invariance or from the coherent-slot kernel after
untwisting.  Global flat directions are not excluded.

For any declared finite coefficient map $B$ into these velocities,
the Hessian on that family is $B^*HB$.  All exact zero relations and
regulator-dependent normalizations therefore pass by congruence.
This is a marked effective-interaction matrix, not the entrance Gram
of an unspecified physical source family.  A comparison between those
two matrices would require a separate theorem.

There is also an exact local source-response statement in covariant
coordinates.  Differentiating a transverse ribbon spin expectation gives
\begin{equation}
 \delta\langle x_t^i\rangle_0
    =\gamma\sum_{s,\alpha}c_{ts}b_{s\alpha}^i,
 \qquad
 W_{t,s\alpha}=\frac{\gamma c_{ts}}{m_t}\ge0,
 \qquad\sum_{s,\alpha}W_{t,s\alpha}=1.
 \label{f16v61:barycentric-response}
\end{equation}
Thus the linear response of its orientation is a normalized positive
average of the slot rotations.  A common rotation is reproduced, not
contracted.  These component coordinates are gauge covariant auxiliary
marks; gauge-invariant sources must still be assembled from the actual
Wilson path products.  Neither this response nor the curvature bound
claims a continuum source construction.

\subsection{Exact fixtures and the failure of global Hessian positivity}

For $T=1$, $n_1=6$ and
\begin{equation}
 Z(p)=z_4\left(\gamma\left|\sum_{\alpha=1}^6p_\alpha\right|\right),
 \qquad z_4(u)=\frac{2I_1(u)}u,\quad z_4(0)=1.
 \label{f16v61:single-link}
\end{equation}
At coherence, $m=a_4(6\gamma)$, $c_{11}=m/(6\gamma)$, and
\begin{equation}
 Q(b)=\gamma m\left(\sum_\alpha|b_\alpha|^2
                       -\frac16\left|\sum_\alpha b_\alpha\right|^2\right).
 \label{f16v61:single-link-Hessian}
\end{equation}
This is an exact compact fixture for the signs, factors, and kernel.
It is not a replacement of the many-time theorem by a one-link calculation.

Conversely choose three slots equal to $e_0$ and three equal to $-e_0$.
Rotate one of the positive slots to $\cos\varepsilon\,e_0+
\sin\varepsilon\,e_1$, keeping the other five fixed.  The sum of the
slots has squared norm $2(1-\cos\varepsilon)$.  The Haar moment expansion
gives $\log z_4(u)=u^2/8+O(u^4)$, so
\begin{equation}
 \left.\frac{d^2}{d\varepsilon^2}\Psi(p(\varepsilon))
                                                   \right|_{0}
                         =-\frac{\gamma^2}{4}<0.
 \label{f16v61:negative-Hessian}
\end{equation}
Such central staple choices are realizable by retained Wilson links:
the parallel displaced link in each incident plaquette can set its staple
sign, with the remaining short-path links equal to the identity.
This concerns the ribbon contribution; other retained plaquettes are
not asserted to have zero Hessian along that path.
The global lower bound \eqref{f16v61:global-curvature-eq} remains true.
A positive action deficit relative to its minimum plainly does not imply
global convexity on the compact exterior space.

\subsection{Where the first connected enlargement loses this argument}

The clipping proof has a useful scope test.  For a finite pair-interaction
graph with $n_i\ge m$ field slots, degree at most $D<m$, and
$n_i+\deg(i)\le B$, the same proof gives
\begin{equation}
 \Psi-\Psi_0\ge
   \frac{\gamma^2(m-D)^2}{4m(4+B\gamma)}
                         \sum_i(n_i-r_i).
 \label{f16v61:pinning-balance}
\end{equation}
To verify the constant, clip $r_i$ upward at $u=(m+D)/2$.
During the aligned interpolation the conditional longitudinal field is
at least $\gamma(u-D)$ and its norm at most $B\gamma$.
The retained interpolation distance is at least
$(1-u/m)(n_i-r_i)$.  Multiplication yields
\eqref{f16v61:pinning-balance}; alignment domination applies even to
orthogonal transports around cycles.  This is the exact place where
strictly more fixed staples than free neighbors entered the argument.

For one native ribbon, $m=4$, $D=2$, $B=6$, recovering
$\lambda_\gamma$.  Now free jointly two neighboring spatial links which
share exactly one spatial plaquette, through the same time interval.
At an interior time, each has three remaining fixed spatial staples,
two temporal neighbors, and one free neighbor across that spatial
plaquette.  The conditional interaction graph is a two-rail ladder:
$m=D=3$ in its interior.  The spatial plaquette is bilinear in its two
free quaternion links, and its fixed complementary links give an
orthogonal transport.  There are closed temporal-spatial cycles.
The general alignment upper bound still applies, but
\eqref{f16v61:pinning-balance} no longer gives a positive coefficient.
This is a failure of this sufficient estimate, not a proof that the
ladder lacks a positive retained-curvature bound.

It would be incorrect to replace that joint calculation by reusing the
independent-ribbon law after one color has already been marginalized.
The next concrete upstream cell is this native two-link ladder, with its
three staple fields per vertex and its actual cycle transports.  A useful
continuation must control its coherent contrasts and transport-holonomy
terms, or bound a specified normalized retained-source correlation, while
retaining the induced interactions.  The first-color bound above does
not supply that second-color input automatically.

\subsection{Verification, preservation, and result ledger}

The companion diagnostic
\begin{center}\small
\path{verify_ym_f16_v61_retained_curvature.py}
\end{center}
checks exact Haar moments, the clipping constant, algebraic Ward
cancellation, quaternion path and frame identities, and finite Ising
sign sums.  It also evaluates genuine compact one- and two-spin integrals
and aligned multi-time $S^3$ ribbons by deterministic quadrature, with
resolution comparisons.  These are conditional compact integrals, not
samples of the full retained Wilson vacuum.  Floating-point quadrature
is not an outward-rounded interval certificate, and numerical agreement
is not the proof of the inequalities above.

The aligned-ribbon quadrature keeps the sphere exactly at the analytic
level.  Write $x_t=(u_t,\sqrt{1-u_t^2}\,\omega_t)$ with $\omega_t\in S^2$.
The $u$ measure is $(2/\pi)\sqrt{1-u^2}\,du$.  After integrating angular
variables along the open chain, the scalar bond is
\[
 e^{\gamma uv}\frac{\sinh z}{z},\qquad
 z=\gamma\sqrt{1-u^2}\sqrt{1-v^2}.
\]
For a transverse two-point insertion, each bond between the marked times
is multiplied by $\coth z-1/z$, with its continuous value zero at $z=0$,
and the endpoint factor is
$\sqrt{1-u_t^2}\sqrt{1-u_s^2}/3$.
This provides an independent finite-dimensional integration check on the
Ward identity and all three pieces of \eqref{f16v61:contrast-split}.

Removing this single delimited appendix and reversing the displayed title
from V61 to V60 recovers the uploaded predecessor byte for byte.  The
comparison manuscripts are not modified.  No new Lean formalization,
full-vacuum simulation, native physical-source Gram certificate, or
nonperturbative calibration of the physical spacing is claimed.

\begin{record}[V61 result ledger]
\label{f16v61:ledger}
\emph{Proved here:} alignment domination for compact vector-spin integrals
with orthogonal transports; the length-independent, all-exterior ribbon
frustration bound including zero staple sums; its literal generated
retained Wilson-loop lower bound after one color elimination; the exact
transverse Ward susceptibility identity and positive slot-difference
Hessian; its exact kernel, contrast decomposition, susceptibility norm,
and paid neighboring-time contribution; the pullback to actual retained
link variations and coefficient relations; the normalized orientation
response; the native negative-Hessian fixture away from coherence; and
the explicit pinning-balance criterion and its first connected-ladder
scope limitation.

\emph{Inherited:} the compact Wilson ribbon factorization, local boundary
convention, and actual retained action of V59--V60; the global one-link
Bessel barrier of V52; and the unchanged physical transfer.  No continuum
trajectory hypothesis is needed for the new finite-regulator inequalities.

\emph{Not asserted:} global Hessian positivity, replacement of the exact
retained law by a six-loop Wilson law, normalized-vacuum comparison from
a pointwise action inequality, independence after a second elimination,
or contraction of a common rotational source direction.

\emph{Still unproved:} the required interacting physical source limit;
source-weighted control of the nonlinear retained collective dynamics;
physical-time contraction with common dense-core constants; and the
continuum Yang--Mills mass gap.
\end{record}

\begin{firewall}
This iteration pays a previously unsigned native local interaction and
keeps its exact symmetry cancellation.  Its nonlinear all-exterior bound
controls explicit retained curvature rather than an unspecified source
norm.  A positive effective-action deficit and a positive coherent Hessian
are not a physical spectral gap.  V1--V60 remain unchanged; the continuation
does not reinstate V59's excluded positive ribbon-capture target.
\end{firewall}
% END F16 V61 APPEND

% BEGIN F16 V62 APPEND -- 2026-09-18
% Append-only continuation of V61.  The preceding mathematical body is unchanged.
\clearpage
\section[V62: joint Wilson ladders and paid transport holonomy]{V62: joint Wilson ladders\\and paid transport holonomy}
\label{f16v62:context}

V61 ended at the first connected two-link ladder: three fixed staples
and three integrated neighbors at an interior vertex.  Its clipping
argument required more fixed slots than free neighbors.  The present
continuation removes that requirement rather than introducing another
hypothesis at the same interface.  A transverse Ward identity and a
\emph{raw second-moment} bound replace the worst-boundary longitudinal
estimate.  They give a nonlinear, all-exterior retained-curvature bound
on any finite pair-interaction graph of bounded incidence.  In
particular, the native two-link ladder has a positive coefficient.

A second calculation retains the transport holonomies around its squares.
Differentiating the positive Haar-pairing expansion at the coherent
exterior gives an exact sum of open-path mismatch and closed-loop
holonomy energies.  Marking and removing a short component of that
expansion gives polynomial, not empty-patch exponential, lower bounds
for its coefficient.  The resulting single lower form simultaneously
pays slot contrasts, edge mismatches and square holonomies, with an
exact local-frame kernel.  Both calculations extend to a rectangular
strip of parallel Wilson links of arbitrary finite width.

These are statements about the actual joint conditional integral and its
retained interaction.  They do not assert positivity of a Hessian away
from coherence, comparison of normalized vacuum correlations from an
action bound, or physical-time contraction.  The short-time obstruction
of V60 is not bypassed by increasing this conditional block.

\subsection{Inputs, comparison with earlier work, and the literal joint cell}

The only analytic inputs inherited here are the Wilson incidence and
local boundary convention of V59--V61, the compact alignment domination
of Proposition~\ref{f16v61:alignment}, and the global radial barrier
\eqref{f16v61:radial-lower}.  V61's fixed-transport Ward calculation
is extended to cyclic graphs and changing transports, not repeated as
an independent claim.  The independent single-link integrals in the
f15 archive do not integrate the joint cell below.  Neither the cubic
exponential-moment step of the parallel V67 branch nor a continuum
trajectory assumption is used.

The use of paired strands in compact spin integrals has a substantial
literature.  For context, see B.~Lees and L.~Taggi,
\emph{Exponential decay of transverse correlations for $O(N)$ spin systems
and related models}, Probab.\ Theory Relat.\ Fields \textbf{180}
(2021), 1099--1133,
\href{https://doi.org/10.1007/s00440-021-01053-5}{doi:10.1007/s00440-021-01053-5},
and A.~Quitmann and L.~Taggi, \emph{Macroscopic Loops in the Bose Gas,
Spin $O(N)$ and Related Models}, Commun.\ Math.\ Phys.\ \textbf{400}
(2023), 2081--2136,
\href{https://doi.org/10.1007/s00220-023-04633-9}{doi:10.1007/s00220-023-04633-9}.
The labelled insertion and Wilson-transport estimates needed here are
proved below; no external spin-model decay theorem is imported as a
Wilson-vacuum conclusion.

Choose two \emph{parallel}, positively oriented spatial $\mu$ links at
$x$ and $x+\hat\nu$, with $\mu\ne\nu$, and free their copies at times
$t=1,\ldots,T$.  Their copies outside this interval, all complementary
links, and all local boundary data remain fixed.  Parallel orientation
ensures that the two free links in their shared plaquette occur with
opposite signs.  Quaternion coordinates therefore give the exact factor
$x_i\cdot R_{ij}x_j$, with $R_{ij}\in\mathrm{SO}(4)$, rather than an
unrecorded conjugation of one spin.

The graph $\mathcal G=(\mathcal V,\mathcal E)$ has vertices
$i=(r,t)$, $r\in\{0,1\}$, temporal rail edges and one rung at each time.
Its conditional partition function is
\begin{equation}
 Z(p,R)=\int\exp\left\{\gamma\sum_{i\in\mathcal V}
                  \sum_{\alpha=1}^{n_i}p_{i\alpha}\cdot x_i
          +\gamma\sum_{\{i,j\}\in\mathcal E}x_i\cdot R_{ij}x_j\right\}
                         \prod_i d\sigma(x_i),
 \label{f16v62:joint-integral}
\end{equation}
where $R_{ji}=R_{ij}^{\mathsf T}$ and an orientation is selected once for
each undirected edge.  All field slots are actual retained three-link
paths.  For the ladder,
\begin{equation}
 n_{r,t}=3+\mathbf1_{t=1}+\mathbf1_{t=T},\qquad
 n_i+\deg(i)=6.
 \label{f16v62:ladder-incidence}
\end{equation}
Thus $n=5$ when $T=1$, $n=4$ at distinct time endpoints, and $n=3$ in
the time interior.  Constants in the Wilson action independent of the
exterior have been omitted from $Z$, not from a normalization comparison.

More generally, free the parallel $\mu$ links at $x+r\hat\nu$,
$r=0,\ldots,W-1$, through this same time interval.  The graph is the
rectangular $W\times T$ grid, with
\begin{equation}
 n_i=6-\deg(i)\ge2.
 \label{f16v62:strip-incidence}
\end{equation}
At an interior vertex its two fixed staples are in the remaining spatial
direction.  Each Wilson plaquette has zero, one or two free links; the
two-free terms are precisely grid edges.  There is no four-free plaquette
inside this \emph{one-orientation} strip.  This incidence fact is essential
to the scope of the result.

Put $r_i=|\sum_\alpha p_{i\alpha}|$, $F_i=n_i-r_i$, $\Psi=-\log Z$, and
$\Psi_0=-\log Z(e_0,I)$, with $e_0=(1,0,0,0)$.  The broader graph
statements below assume $n_i+\deg(i)\le B$, and use
\begin{equation}
 D_{\gamma,B}=4+B\gamma,\qquad D_\gamma=4+6\gamma.
 \label{f16v62:D}
\end{equation}
All statements are at a fixed regulator and use normalized Haar measure.
A nonlocal Perron endpoint weight is not absorbed into these local factors.

\subsection{The raw second moment removes the pinning-balance restriction}

\begin{proposition}[One-site second moments and graph Ward positivity]
\label{f16v62:raw-Ward}
For the one-spin law proportional to $e^{h\cdot x}d\sigma(x)$ on $S^3$,
\begin{equation}
 \mathbb E_h[xx^{\mathsf T}]\succeq\frac{1}{4+|h|}I_4.
 \label{f16v62:raw-second}
\end{equation}
On any finite graph with identity transports and aligned fields
$\gamma u_i e_0$, $0\le u_i\le n_i$, let
$m_i=\langle x_i^0\rangle_u$ and
$c_{ij}=\langle x_i^a x_j^a\rangle_u$, for a fixed transverse coordinate
$a\in\{1,2,3\}$.  Then
\begin{equation}
 c_{ij}\ge0,\qquad c_{ii}\ge D_{\gamma,B}^{-1},\qquad
 m_i=\gamma\sum_j u_jc_{ij}\ge\frac{\gamma u_i}{D_{\gamma,B}}.
 \label{f16v62:Ward-lower}
\end{equation}
No inequality between $n_i$ and $\deg(i)$ is required.
\end{proposition}
\begin{proof}
Write $u=|h|$ and, for $u>0$, $\widehat h=h/u$.  Rotational symmetry and
one-site rotation integration by parts give
\[
 \mathbb E_h[xx^{\mathsf T}]
 =\chi_4(u)I_4+(1-4\chi_4(u))\widehat h\widehat h^{\mathsf T},
 \qquad \chi_4(u)=\frac{I_2(u)}{uI_1(u)}.
\]
For completeness, a transverse rotation gives
$\mathbb E_h(x^a)^2=\mathbb E_h x^0/u=\chi_4(u)$ in the field frame.
The longitudinal second moment is at least its Haar value $1/4$:
by sign symmetry it is the Haar expectation of $(x^0)^2$ reweighted by
$\cosh(u|x^0|)$, and these are two increasing functions of $|x^0|$.
Their covariance is nonnegative by the identity expressing it as one
half the expectation of the product of their two-copy differences.
It follows that $0<\chi_4(u)\le1/4$.  The least raw second-moment
eigenvalue is consequently $\chi_4(u)$, not the possibly much smaller
longitudinal \emph{variance}.  The inherited lower barrier
$\chi_4(u)\ge1/(4+u)$ proves \eqref{f16v62:raw-second}, with its Haar
value at $u=0$.

For the graph, condition on all spins other than $x_i$.  Its conditional
field has norm at most $\gamma(u_i+\deg(i))\le B\gamma$.
Equation~\eqref{f16v62:raw-second} gives the asserted diagonal bound.
For off-diagonal positivity, condition on $|x_i^a|$ and the other three
coordinates at every vertex.  The signs of $x_i^a$ have a zero-field
ferromagnetic Ising law, including when the graph has cycles.  Its
two-sign numerator is a sum of nonnegative high-temperature parity
weights, and its denominator is positive.  Hence $c_{ij}\ge0$ after
averaging.  On a connected graph with positive edge coefficients the
entries are strictly positive: one simple path between $i$ and $j$
provides a positive numerator term for almost every set of amplitudes.
No product formula valid only for a chain is used.

Finally rotate all spins simultaneously in the $(0,a)$ plane.  The pair
interactions and Haar measure are invariant.  Applying the generator to
$x_i^a$ times the density gives
$m_i=\gamma\sum_j u_j c_{ij}$.  Every summand is nonnegative, so keeping
$j=i$ proves the last bound in \eqref{f16v62:Ward-lower}.
\end{proof}

\begin{theorem}[All-exterior graph curvature without degree dominance]
\label{f16v62:global}
For \eqref{f16v62:joint-integral} on a finite simple graph of incidence
at most $B$, with arbitrary unit field slots and arbitrary orthogonal
transports,
\begin{equation}
 \boxed{\quad
 \Psi(p,R)-\Psi_0
 \ge\frac{\gamma^2}{2D_{\gamma,B}}\sum_i(n_i^2-r_i^2)
 =\frac{\gamma^2}{D_{\gamma,B}}
                  \sum_i\sum_{\alpha<\beta}
                            (1-p_{i\alpha}\cdot p_{i\beta}).\quad}
 \label{f16v62:global-bound}
\end{equation}
For the native ladder this also implies
\begin{equation}
 \Psi-\Psi_0\ge\frac{3\gamma^2}{2D_\gamma}\sum_i F_i,
 \label{f16v62:ladder-frustration}
\end{equation}
and for any rectangular strip it implies
$\Psi-\Psi_0\ge(\gamma^2/D_\gamma)\sum_iF_i$.
All constants are independent of graph size.  Zero staple sums and
frustrated cycle transports are included.
\end{theorem}
\begin{proof}
Combine the fields at each vertex.  V61's alignment inequality gives
$Z(p,R)\le Z((\gamma r_i e_0)_i,I)$.  Interpolate the aligned amplitudes
by $u_i(s)=r_i+s(n_i-r_i)$, $0\le s\le1$.  By the preceding proposition,
\[
 \frac{d}{ds}\log Z((\gamma u_i(s)e_0)_i,I)
 =\gamma\sum_i(n_i-r_i)m_i(s)
 \ge\frac{\gamma^2}{D_{\gamma,B}}
                       \sum_i(n_i-r_i)u_i(s).
\]
Integrating gives the first inequality.  The exact identity
$n_i^2-r_i^2=2\sum_{\alpha<\beta}(1-p_{i\alpha}\cdot p_{i\beta})$
gives the equality.  Also
$n_i^2-r_i^2=(n_i+r_i)F_i\ge n_iF_i$; use $n_i\ge3$ for the ladder
and $n_i\ge2$ for the strip.
\end{proof}

For actual slots, $p_{i\alpha}\cdot p_{i\beta}
=w(P_{i\alpha}P_{i\beta}^{-1})$, so these are retained Wilson loops of
length at most six.  With the same admissible flat-boundary convention
as V61, the actual retained action therefore satisfies
\begin{equation}
 \begin{split}
 S_{\rm eff}(\xi)-S_{\rm eff}(\mathbf1)
 \ge{}&\gamma\sum_{p\text{ retained}}(1-w(U_p))\\
 &+\frac{\gamma^2}{D_\gamma}
       \sum_i\sum_{\alpha<\beta}
                    (1-w(P_{i\alpha}P_{i\beta}^{-1})).
 \end{split}
 \label{f16v62:native-global}
\end{equation}
This is an improvement of V61's sufficient bound, not a correction to
its valid weaker inequality.  The lost coefficient at $m=D=3$ was a
limitation of its proof.  The present interpolation does not discard
small-staple configurations and never requires a positive longitudinal
field for every conditioning configuration.

Equation~\eqref{f16v62:native-global} still does not compare the two
\emph{normalized} measures.  The exact retained action contains the
whole joint integral, not just the displayed six-link terms.

\subsection{The coherent Hessian with changing transports}

Fix the coherent point $p_{i\alpha}=e_0$, $R_{ij}=I$.  A tangent consists
of $b_{i\alpha}\in e_0^\perp$ and $A_{ij}\in\mathfrak{so}(4)$, with
$A_{ji}=-A_{ij}$.  Choose any smooth unit-slot and orthogonal-transport
curves with these tangents.  The first derivative of $\Psi$ vanishes at
coherence; its second derivative $Q(b,A)$ is therefore independent of
curve accelerations.

For an oriented path $P:i\leadsto j$, write
$A(P)=\sum_{e\in P}\operatorname{sgn}_P(e)A_e$, counting repeated
traversals.  For a closed walk $C$ the same sum is denoted $A(C)$.
These are first-order transport and holonomy data, not replacements
of the noncommuting finite products.

\begin{theorem}[Exact marked-strand representation of the compact Hessian]
\label{f16v62:Hessian-expansion}
There is a probability measure $\mathbb P_0$ on the coherent Haar-pairing
diagrams of \eqref{f16v62:joint-integral} such that
\begin{equation}
 \boxed{\quad
 Q(b,A)=\mathbb E_0\left[
   \sum_{P\text{ open}}
       |b_{i(P)\alpha(P)}-b_{j(P)\beta(P)}-A(P)e_0|^2
       +\frac14\sum_{C\text{ closed}}\|A(C)\|_{\rm F}^2
                         \right].\quad}
 \label{f16v62:exact-Q}
\end{equation}
An open component has two field-slot endpoints.  A closed component has
none.  Components, including multiplicities, are each counted once;
the orientation of a closed component does not affect its energy.
In particular $Q\ge0$ and the formula retains all cycle transports.
\end{theorem}
\begin{proof}
Expand every bond exponential and every individual slot exponential.
Let $k_e$ be the bond multiplicities, $\ell_{i\alpha}$ the slot
multiplicities, and
\begin{equation}
 N_i=\sum_{e\ni i}k_e+\sum_\alpha\ell_{i\alpha}.
 \label{f16v62:local-count}
\end{equation}
Odd $N_i$ have zero integral.  For even $N_i$, pair the $N_i$ local
half-edges using the Haar moment tensor of V61.  Each pairing has factor
\[
 u(N_i)=\frac1{4\cdot6\cdots(N_i+2)},\qquad u(0)=1.
\]
Give repeated bond and slot occurrences labels $1,\ldots,k_e$ and
$1,\ldots,\ell_{i\alpha}$, retaining the exponential factorials.
The pairings form closed index lines and open index lines terminated
by two slot vectors.  At coherence every closed component contributes
$4$ and every open component contributes $1$.  A diagram $D$ thus has
positive weight
\begin{equation}
 a_D=\frac{\gamma^{\sum_e k_e+\sum_{i,\alpha}\ell_{i\alpha}}}
                 {\prod_e k_e!\prod_{i,\alpha}\ell_{i\alpha}!}
                 \prod_i u(N_i)\,4^{\#\text{ closed components}}.
 \label{f16v62:diagram-weight}
\end{equation}
The weights sum to $Z_0$, so $a_D/Z_0$ defines $\mathbb P_0$.

For arbitrary external data the same diagram has relative factor
\[
 \prod_{P\text{ open}}p_{i(P)\alpha(P)}^{\mathsf T}R(P)
                           p_{j(P)\beta(P)}
 \prod_{C\text{ closed}}\frac{\Tr R(C)}4,
\]
where transports are ordered along each index line, with transposes on
reversed edges.  At coherence each factor equals one and its first
derivative is zero.  The second derivative of an open factor is
$-|b_{i\alpha}-b_{j\beta}-A(P)e_0|^2$, because it is an inner product
of two unit-vector curves agreeing at zero.  For a closed factor,
orthogonality of $R(C)$ gives
$\Tr R(C)''(0)=-\|A(C)\|_{\rm F}^2$.
Thus the negative second derivative of the product is exactly the sum
inside \eqref{f16v62:exact-Q}.

Termwise differentiation is legitimate.  On a finite graph the series
of absolute diagram weights is the finite coherent integral, and a
second derivative is bounded by a constant times the square of the
total number of occurrences.  Its expectation is finite, as follows
by differentiating the same compact exponential integral with a common
positive scale on its coefficients.  Since $Z'(0)=0$, the Hessian of
$-\log Z$ is $-Z''(0)/Z_0$, proving the result.
\end{proof}

For example, on geodesic representative curves the same Hessian is
\begin{equation}
 \begin{split}
 Q(b,A)={}&\gamma\sum_{i,\alpha}m_i|b_{i\alpha}|^2
   -\gamma\sum_{\{i,j\}\in\mathcal E}
                     \langle x_i\cdot A_{ij}^2x_j\rangle_0\\
 &-\operatorname{Var}_0\left[
       \gamma\sum_{i,\alpha}b_{i\alpha}\cdot x_i
          +\gamma\sum_{\{i,j\}\in\mathcal E}
                                      x_i\cdot A_{ij}x_j\right].
 \end{split}
 \label{f16v62:signed-Q}
\end{equation}
The variance includes all field--current and current--current cross
terms.  Formula~\eqref{f16v62:exact-Q} pays their cancellation, rather
than deleting them.  Setting $A=0$ reduces to the graph version of
V61's slot-difference Hessian; changing transports is the new part.

\subsection{Component insertion pays short holonomies without empty patches}

It is not sufficient to select a diagram having no other activity in
an entire strip: its normalized weight could be exponentially small
in the strip area.  The next insertion leaves \emph{all} unmarked
activity in place.

\begin{proposition}[Exact local component intensities]
\label{f16v62:insertion}
Under $\mathbb P_0$, put
\[
 J(S)=\mathbb E_0\prod_{i\in S}(4+N_i)^{-1}
\]
for any finite vertex set $S$.  Let $K_{i,\alpha\beta}$ count open
components having no graph edge and joining distinct slots $\alpha<\beta$
at $i$.  Let $K_{ij,\alpha\beta}$ count open components that traverse
exactly the one edge $\{i,j\}$ and have endpoints in the indicated slots.
For a specified unoriented simple graph cycle $C$ of length $L\ge3$,
let $K_C$ count closed components traversing it once.  Then
\begin{align}
 \mathbb E_0K_{i,\alpha\beta}&=\gamma^2 J(\{i\}),
       \label{f16v62:insert-zero}\\
 \mathbb E_0K_{ij,\alpha\beta}&=\gamma^3 J(\{i,j\}),
       \label{f16v62:insert-edge}\\
 \mathbb E_0K_C&=4\gamma^LJ(\mathcal V(C)).
       \label{f16v62:insert-cycle}
\end{align}
Moreover
\begin{equation}
 \mathbb E_0N_i\le\gamma(n_i+\deg(i))\le B\gamma,
 \qquad J(S)\ge D_{\gamma,B}^{-|S|}.
 \label{f16v62:count-Jensen}
\end{equation}
No independence of the $N_i$ is assumed.
\end{proposition}
\begin{proof}
Consider a diagram together with one marked component of a specified
type on the left of \eqref{f16v62:insert-zero}--\eqref{f16v62:insert-cycle}.
Remove its occurrences and the local pair joining its two half-edges
at each visited vertex; then order the remaining occurrence labels.
The result is an arbitrary background diagram.  Conversely insert the
prescribed component into any background diagram and choose its position
among the old labels for each affected bond or slot.  An added occurrence
has $k+1$ label positions.  This factor cancels exactly the factorial
ratio $k!/(k+1)!$ in \eqref{f16v62:diagram-weight}.  The same cancellation
holds for each affected slot.  Distinct endpoint slots are used for the
zero-edge insertion, so there is no unrecorded two-identical-endpoint
symmetry factor.

Every visited vertex acquires two half-edges, and its Haar factor changes
by $u(N_i+2)/u(N_i)=1/(4+N_i)$.  The inserted component contributes two
slot coefficients for the first case, two slot and one edge coefficients
for the second, and $L$ edge coefficients and one extra trace $4$ for
the closed cycle.  It is paired internally at every vertex, so the
background components are unchanged.  This is a bijection of marked
diagrams with background diagrams and insertion-label choices, proving
all three identities.  For a cycle fix a starting vertex and orientation
once to describe the insertion; each unoriented component on that fixed
geometric cycle is marked once, not once per start or per orientation.

Allow independent coefficients for each exponential before setting them
to $\gamma$.  Differentiating the positive series and the compact integral
in the same coefficient gives
\[
 \mathbb E_0 k_{ij}=\gamma\langle x_i\cdot x_j\rangle_0\le\gamma,
 \qquad
 \mathbb E_0\ell_{i\alpha}=\gamma m_i\le\gamma.
\]
These equalities prove the first part of \eqref{f16v62:count-Jensen}.
For the second, two applications of Jensen's inequality yield
\[
 \mathbb E_0\exp\left[-\sum_{i\in S}\log(4+N_i)\right]
 \ge\exp\left[-\sum_{i\in S}\mathbb E_0\log(4+N_i)\right]
 \ge\prod_{i\in S}(4+\mathbb E_0N_i)^{-1}
 \ge D_{\gamma,B}^{-|S|}.
\]
The estimates retain the complete normalized coherent partition function.
\end{proof}

\begin{theorem}[A simultaneously paid local Hessian form]
\label{f16v62:paid-form}
For a ladder or rectangular strip, orient its edges once and let
$\mathcal C_\square$ be its distinct elementary squares, each included
once.  Then the exact compact Hessian obeys
\begin{equation}
 \boxed{\begin{aligned}
 Q(b,A)\ge{}&\frac{\gamma^2}{D_\gamma}
           \sum_i\sum_{\alpha<\beta}|b_{i\alpha}-b_{i\beta}|^2\\
 &+\frac{\gamma^3}{D_\gamma^2}
       \sum_{\{i,j\}\in\mathcal E}\sum_{\alpha=1}^{n_i}
           \sum_{\beta=1}^{n_j}
                  |b_{i\alpha}-b_{j\beta}-A_{ij}e_0|^2\\
 &+\frac{\gamma^4}{D_\gamma^4}
                 \sum_{C\in\mathcal C_\square}\|A(C)\|_{\rm F}^2.
 \end{aligned}}
 \label{f16v62:paid-bound}
\end{equation}
For a more general graph the same statement holds with $D_{\gamma,B}$,
and with each selected distinct simple length-$L$ cycle contributing
$\gamma^L D_{\gamma,B}^{-L}\|A(C)\|_{\rm F}^2$.
\end{theorem}
\begin{proof}
In \eqref{f16v62:exact-Q}, keep precisely the zero-edge open components,
the one-edge open components, and the closed components traversing one
of the selected simple cycles once.  These component classes are disjoint.
A component cannot be two different elementary square components; multiple
components on one square are correctly counted with their multiplicity.
The insertion identities evaluate their expected counts, and
\eqref{f16v62:count-Jensen} bounds them.  The factor $4$ for inserting
a closed component cancels the $1/4$ in its trace energy.  All omitted
component energies are nonnegative.  Thus all three lines are paid by
one partition of the same exact form, not by adding unrelated lower
bounds for $Q$.
\end{proof}

As $\gamma\to\infty$, the first two coefficients are asymptotic to
$\gamma/6$ and $\gamma/36$, while the square coefficient tends to
$1/6^4$.  These are finite-regulator effective-action coefficients.
There is no strip-area probability, $e^{-c\gamma}$ isolation penalty,
or assumed decorrelation of the surrounding diagrams in this estimate.
There is equally no conclusion about a gap in physical time.

\subsection{Exact kernel and the distinction between open mismatch and holonomy}

\begin{theorem}[The local-frame kernel is complete]
\label{f16v62:kernel}
For $\gamma>0$ on a connected rectangular strip, the kernel of the full
Hessian consists exactly of the tangents
\begin{equation}
 b_{i\alpha}=\Omega_i e_0\quad\text{for every slot},\qquad
 A_{ij}=\Omega_i-\Omega_j,\qquad
 \Omega_i\in\mathfrak{so}(4).
 \label{f16v62:frame-kernel}
\end{equation}
The kernel is already detected by the paid local form
\eqref{f16v62:paid-bound}.  In the unrestricted slot-and-transport
parameter space its dimension is $6|\mathcal V|-3$.
\end{theorem}
\begin{proof}
The change of integrated variables $x_i=O_i z_i$ leaves $Z$ unchanged
when $p_{i\alpha}=O_i e_0$ and $R_{ij}=O_iO_j^{\mathsf T}$.  Differentiating
at $O_i=I$ shows that \eqref{f16v62:frame-kernel} is in the kernel.

Conversely $Q=0$ forces every term in \eqref{f16v62:paid-bound} to vanish.
Thus all slots at $i$ have a common tangent $b_i$,
$b_i-b_j=A_{ij}e_0$ on each edge, and $A(C)=0$ on every elementary square.
The elementary squares span the cycle space of a rectangular grid.
The matrix-valued one-cochain $A$ consequently has zero sum along every
closed path.  Choose a root and a skew matrix $\Omega_*$ satisfying
$\Omega_*e_0=b_*$; define $\Omega_i$ by integrating $A$ from the root,
with signs chosen so that $A_{ij}=\Omega_i-\Omega_j$.  Zero cycle sums
make the definition path independent.  The edge equations imply
$\Omega_i e_0=b_i$ everywhere, proving \eqref{f16v62:frame-kernel}.
For a tree the same proof omits the cycle condition.

The map from $(\Omega_i)$ to $(b,A)$ has kernel precisely the constant
skew matrices annihilating $e_0$, a three-dimensional $\mathfrak{so}(3)$.
Its image therefore has dimension $6|\mathcal V|-3$.  All vertices have
at least one slot here; an unpinned graph needs a separate kernel analysis.
\end{proof}

The holonomy line is necessary even after all open mismatches vanish.
On one square set all $b=0$, and put a nonzero transverse-plane rotation
$A_e$ on one edge, with $A_e e_0=0$; set the other $A$ to zero.  Both open
lines of \eqref{f16v62:paid-bound} vanish, whereas the square line is
strictly positive.  This is an exact tangent example, not an assumption
about a particular source coefficient map.  It identifies the response
missed by a reduction keeping only the motion of the mean spin direction.

Conversely a common rotation is reproduced exactly by the integral; it
is not contracted.  The kernel dimension just computed counts auxiliary
slot-and-transport parameters.  It is not a physical number of Wilson
polarizations, since actual exterior links impose relations among those
parameters.

\subsection{The square holonomies are literal retained Wilson data}

For the parallel-link construction, each graph edge comes from a Wilson
plaquette with two free opposite sides and two retained complementary
sides.  Its orthogonal transport has the exact quaternion form
\begin{equation}
 R_{ij}(q)=\ell_{ij}q\,r_{ij}^{-1},\qquad
 \ell_{ij},r_{ij}\in\SU(2),
 \label{f16v62:Spin4-transport}
\end{equation}
where the two retained links lie at the left and right endpoints of the
free $\mu$ links.  Indeed
$w(U_i r_{ij}U_j^{-1}\ell_{ij}^{-1})
= x_i\cdot R_{ij}x_j$.

For an elementary graph square $C$ in the $(\nu,\text{time})$ plane,
ordered multiplication gives
\begin{equation}
 R(C)(q)=L_Cq\,R_C^{-1},\qquad
 \Tr_{\mathbb R^4}R(C)=4w(L_C)w(R_C).
 \label{f16v62:literal-holonomy}
\end{equation}
Here $L_C$ and $R_C$ are the actual retained plaquette holonomies in
that plane at the two endpoints of the free $\mu$ links, with basepoints
and orientation inherited from the ordered graph path.  This follows
by composing \eqref{f16v62:Spin4-transport}; right factors accumulate in
the inverse order exactly as displayed.  The real trace identity follows
by evaluating left and right quaternion multiplication on
$1,\mathbf i,\mathbf j,\mathbf k$.

At an all-identity exterior write their tangent quaternions as
$\lambda_C,\rho_C\in\operatorname{Im}\mathbb H$.  Then
\begin{equation}
 A(C)q=\lambda_Cq-q\rho_C,\qquad
 \|A(C)\|_{\rm F}^2=4(|\lambda_C|^2+|\rho_C|^2).
 \label{f16v62:native-holonomy-norm}
\end{equation}
Left and right multiplication by a pure imaginary unit are orthogonal
in Frobenius inner product, and each has squared Frobenius norm $4$.
The last line of \eqref{f16v62:paid-bound} therefore pays a specific
retained-plaquette curvature at both endpoint planes.  It does not
replace an $\mathrm{SO}(4)$ holonomy by a single fundamental $\SU(2)$
Wilson trace.  Globally the simultaneous center pair $L_C=R_C=-I$
acts trivially on quaternions; no global separation of those two center
holonomies follows just from $\Tr R(C)$.

Let $\xi(\varepsilon)$ be a smooth curve of actual retained links through
$\mathbf1$, and obtain $b_{i\alpha}$ and $A_{ij}$ by differentiating
its actual path products and transports.  With $v$ the retained link
velocity in unit-quaternion normalization, the exact pullback is
\begin{equation}
 D^2S_{\rm eff}(\mathbf1)[v,v]
 =\gamma\sum_{p\text{ retained}}|(dv)_p|^2+Q(b(v),A(v)).
 \label{f16v62:actual-pullback}
\end{equation}
No tree untwisting is allowed to erase the remaining square transports.
The first derivative of the conditional action vanishes at coherence,
so the pullback carries no extra acceleration term.  The retained
plaquettes on the first line and the induced Hessian $Q$ are distinct
terms of the actual action.  A declared finite coefficient map pulls
\eqref{f16v62:actual-pullback} and its lower bound back by congruence,
preserving exact source relations and regulator-dependent normalization.
This marked interaction Hessian is not an entrance Gram of an unspecified
physical source family.

\subsection{Consequences for joint elimination and the next upstream cell}

The two-link calculation is a joint native integral.  Performing the
first link integral and then integrating the second with its full induced
interaction gives the same $Z$ by Fubini.  Replacing that induced
interaction by another independent-ribbon Wilson law would give a
different object.  The same boundary and retained-law conditions apply
to the arbitrary-width strip extension.

There is now no need to stop at the inequality $m=D=3$: the all-exterior
curvature and the coherent signed response are paid for the joint ladder,
including its nontrivial cycle directions.  The coefficients are uniform
in the number of rows and columns, but their controlled quantities are
local curvature defects.  A Poincare estimate converting these differences
to a norm still carries the appropriate size and boundary information.
A positive conditional Hessian does not restore the microscopic capture
fraction excluded by V60, and neither an arbitrary exterior law nor the
normalized retained vacuum was replaced by the coherent reference.

A useful next calculation must leave the pair-interaction class or
actually control a specified normalized retained-source correlation.
The next native object selected here frees all four oriented boundary
links of one spatial plaquette through time.  At each
time it contains the literal compact term
\begin{equation}
 \gamma\,w(U_{1,t}U_{2,t}U_{3,t}^{-1}U_{4,t}^{-1}),
 \label{f16v62:next-quartic}
\end{equation}
in addition to its temporal couplings and remaining staple fields.
This is a four-free-link plaquette, not an orthogonal two-spin bond.
The positive pair-strand insertion proved above cannot simply be asserted
for its quaternion multiplication tensor.  The next source-marked cell
should retain this plaquette product, derive its complete mixed response,
and either establish a signed compact comparison or expose a precise
sign obstruction.  Repeating the already paid two-link degree estimate
would not advance that step.

Separately, promotion to the continuum objective requires a noncollapsing
physical source family, its normalized correlations under the actual
retained law, physical-time compatibility, and common dense-core bounds.
No one of those obligations is supplied by the local Hessian kernel or
the pointwise action lower bound.

\subsection{Verification, preservation, and result ledger}

The companion diagnostic is
\begin{center}\small
\path{verify_ym_f16_v62_joint_ladders.py}.
\end{center}
It distinguishes exact rational coefficient and kernel fixtures, literal
Wilson incidence and quaternion transport identities, labelled-component
insertion checks, and floating-point compact-$S^3$ integration.  The
compact integration uses positive product quadrature on each sphere and
a two-rail transfer contraction, retaining the rung and both temporal
bonds.  Derivative jets keep the complete contact-minus-covariance
expression; they do not approximate the spins by unconstrained Gaussians.
Resolution comparisons are diagnostics, not outward-rounded error bounds.
No full Wilson-vacuum sample, physical-source Gram certificate, or
proof-assistant verification is claimed.

Removing this delimited appendix and changing the displayed title back
from V62 to V61 recovers the predecessor byte for byte.  All comparison
manuscripts remain unchanged.

\begin{record}[V62 result ledger]
\label{f16v62:ledger}
\emph{Proved here:} the raw second-moment graph Ward lower bound without
fixed-field degree dominance; the all-exterior joint-ladder and strip
curvature bounds, including zero staple sums; the exact coherent Hessian
with both open-path mismatches and closed transport holonomies; labelled
short-component insertion intensities with all background activity
retained; polynomial local coefficient lower bounds paid simultaneously;
the complete local-frame kernel; the literal two-endpoint Wilson
holonomy identity; and pullback of the joint interaction to actual
retained-link tangents.

\emph{Inherited:} the positive compact alignment expansion and global
one-link Bessel barrier; literal Wilson plaquettes, local boundary data,
and the actual retained marginal; the unmodified physical transfer and
V60's source-continuity obstruction.

\emph{Not asserted:} global Hessian positivity; removal of induced
interactions after a sequential elimination; a normalized-measure
comparison from the six-link action bound; physical contraction of the
frame kernel; a four-free-link plaquette analogue of the pair-strand
positivity; or a uniform norm bound from local curvature without its
geometry and boundary costs.

\emph{Still unproved:} the necessary interacting physical source limit;
source-weighted control of the nonlinear retained vacuum; physical-time
contraction with common dense-core constants; and the continuum
Yang--Mills mass gap.
\end{record}

\begin{firewall}
This continuation removes an actual sufficient-condition obstruction and
pays cycle responses that a ribbon calculation omits.  The conclusions
concern exact finite compact interactions, not an inferred continuum gap.
The covariance subtraction, cycle transports, induced retained law, and
physical clock remain explicit.  V1--V61 are preserved unchanged.
\end{firewall}
% END F16 V62 APPEND


% BEGIN F16 V63 APPEND -- 2026-09-18
% Append-only continuation of V62.  The preceding mathematical body is unchanged.
\clearpage
\section[V63: the four-link cell and signed plaquette closure]{V63: the four-link cell\\and signed plaquette closure}
\label{f16v63:context}

The next object selected in \eqref{f16v62:next-quartic} contains a genuine
four-free-link Wilson plaquette.  It is not an $O(4)$ pair interaction.
This continuation integrates one complete time layer of that object,
with its five actual complementary staples per link, by a character
convolution.  It then performs the corresponding simultaneous
alternate-layer elimination in the full tube.  No quaternion product is
replaced by a pair bond, and no independent-ribbon law is assigned to the
resulting marginal.

The calculation pays both staple contrasts and the signed cyclic closure
response.  In particular, the four cyclically oriented link spins have
\emph{strictly negative} off-diagonal transverse correlations at the
coherent cell.  Thus the nonnegative transverse-correlation argument of
V62 really does stop at this cell, even though its complete local Hessian
is positive.  The positive quantity is a squared plaquette-closure
mismatch, not a ferromagnetic sum of pair differences.  A five-writer
conditional covariance matrix and its nondegenerate scaled limit are
also computed.  These are conditional compact results; they are not
correlations in an unspecified interacting physical vacuum.

\subsection{Literal geometry and the five-factor character convolution}

Work at a fixed four-dimensional Wilson regulator, on a nondegenerate
open patch or with all links outside the chosen finite interval fixed.
Orient the four free spatial links \emph{cyclically} around one square:
\[
 X_1=U_\mu(x),\quad X_2=U_\nu(x+\hat\mu),\quad
 X_3=U_\mu(x+\hat\nu)^{-1},\quad X_4=U_\nu(x)^{-1}.
\]
Their plaquette is $Q=X_1X_2X_3X_4$.  Each link belongs to six Wilson
plaquettes.  Precisely one is this four-free-link square; the other
five have only that free link when the neighboring time layers are
held fixed.  Their traces are $p_{i\alpha}\cdot x_i$, where
$p_{i\alpha}\in S^3$ is the quaternion of the actual complementary
three-link path, in the same orientation as $X_i$.

For clarity first allow $n_i\ge1$ slots and a marked plaquette twist
$B\in\SU(2)$ at its basepoint.  The actual one-layer cell has $n_i=5$
and $B=I$.  Put
\begin{align}
 Z_\square(p,B)&=\int_{\SU(2)^4}
  \exp\left\{\gamma\,w(B^{-1}Q)
       +\gamma\sum_{i=1}^4\sum_{\alpha=1}^{n_i}
                                      p_{i\alpha}\cdot x_i\right\}
                                      \prod_{i=1}^4 dH(X_i),
 \label{f16v63:cell}\\
 S_i&=\sum_\alpha p_{i\alpha},\quad r_i=|S_i|,\quad
 h_i=\gamma r_i,\quad P_i=S_i/r_i\quad(r_i>0),\quad
 H=B^{-1}P_1P_2P_3P_4.
 \label{f16v63:aggregate}
\end{align}
Here $w(U)=\tfrac12\operatorname{Re}\Tr U$ and Haar measure is normalized.
Every statement at $r_i=0$ below is interpreted without choosing $P_i$.
Write $\Psi_\square=-\log Z_\square$.

Use a new symbol for the \emph{unshifted} one-link normalization, to avoid
confusion with the earlier $e^{-t(1-w)}$ normalization:
\begin{equation}
 \mathfrak z(t)=\int e^{t w(U)}dH(U)=\frac{2I_1(t)}t,\qquad
 \mathfrak z(0)=1,\qquad k_t(U)=e^{t w(U)}/\mathfrak z(t).
 \label{f16v63:z}
\end{equation}
The standard modified-Bessel integral representations can be found in
NIST DLMF, \href{https://dlmf.nist.gov/10.32}{\S10.32}; the particular
convolution and response estimates used here are proved below.  Character
integration in non-Abelian lattice gauge theory also has a substantial
literature, for example R.~Anishetty, S.~Cheluvaraja, H.~S.~Sharatchandra
and M.~Mathur, \emph{Dual of three-dimensional pure SU(2) lattice gauge
theory and the Ponzano--Regge model}, Phys.\ Lett.\ B \textbf{314}
(1993), 387--390,
\href{https://doi.org/10.1016/0370-2693(93)91254-K}{doi:10.1016/0370-2693(93)91254-K}.
No positivity theorem for a general non-Abelian spin-network contraction
is imported from that context.

\begin{theorem}[Exact compact four-link integration]
\label{f16v63:convolution}
For $t=(t_0,\ldots,t_4)=(\gamma,h_1,\ldots,h_4)$, define
\begin{align}
 d_k&=k+1,\qquad
 \lambda_k(t)=I_{k+1}(t)/I_1(t)\quad(t>0),\qquad
 \lambda_0(0)=1,\quad \lambda_k(0)=0\ (k>0),
 \label{f16v63:lambda}\\
 b_k(t)&=\prod_{j=0}^4\lambda_k(t_j),\qquad
 F_t(H)=\sum_{k=0}^\infty d_k b_k(t)\chi_k(H),\qquad
 W_t=F_t(I)=\sum_{k=0}^\infty d_k^2 b_k(t).
 \label{f16v63:F}
\end{align}
The character $\chi_k$ has dimension $d_k$ and
$\chi_k(e^{\theta u})=\sin((k+1)\theta)/\sin\theta$ for a unit imaginary
quaternion $u$.  Then
\begin{equation}
 \boxed{\quad Z_\square(p,B)=
       \left(\prod_{j=0}^4\mathfrak z(t_j)\right)F_t(H).\quad}
 \label{f16v63:exact-Z}
\end{equation}
Equivalently $F_t=k_{t_0}*\cdots*k_{t_4}$.  The series and its derivatives
on compact parameter sets converge.  It is a strictly positive function,
although its individual character values need not be nonnegative.
If some $h_i=0$, $F_t\equiv1$ and the answer is independent of the
undefined direction $P_i$.
\end{theorem}
\begin{proof}
For a class function, normalized Haar integration is
$(2/\pi)\int_0^\pi (\cdot)\sin^2\theta\,d\theta$.  Character orthogonality
and the cosine integral for $I_k$ give
\[
 \int e^{t w(U)}\chi_k(U)dH(U)
 =I_k(t)-I_{k+2}(t)=\frac{2(k+1)I_{k+1}(t)}t.
\]
Thus $k_t=\sum_k d_k\lambda_k(t)\chi_k$.  Schur orthogonality gives
$\chi_k*\chi_l=\mathbf1_{k=l}\chi_k/d_k$, so five convolutions have
exactly the coefficients in \eqref{f16v63:F}.

For an equally direct interpretation, under its normalized one-link
field law $X_i=P_iY_i$, where $Y_i$ has the central density $k_{h_i}$.
Move the deterministic $P_i$ through the product by conjugating the
independent $Y_i$.  Centrality preserves their laws.  The resulting
random product has deterministic center $P_1P_2P_3P_4$ and central
noise density $k_{h_1}*\cdots*k_{h_4}$.  Integrating the remaining
$e^{\gamma w(B^{-1}Q)}$ factor gives \eqref{f16v63:exact-Z}; central
inversion symmetry removes any convolution-orientation convention.
The product of positive densities is strictly positive.

Finally the Bessel power series gives
\begin{equation}
 0\le\lambda_k(t)\le\min\left\{1,
                          \frac{(t/2)^k}{(k+1)!}\right\}.
 \label{f16v63:lambda-upper}
\end{equation}
For the factorial bound, compare the coefficient of $(t^2/4)^s$ in
$I_{k+1}(t)/(t/2)^{k+1}$ and $I_1(t)/(t/2)$, using
$(s+k+1)!\ge(s+1)!(k+1)!$.  The bound by one follows, for example,
from the successive ratios proved in the next proposition.
These bounds give absolute convergence, including any fixed number of
angular derivatives.  Smoothness at zero and differentiation in the
field vectors also follow directly from the finite compact integral.
If $h_i=0$, that link is Haar and all nontrivial Fourier coefficients
vanish.  No direction is introduced there.
\end{proof}

\subsection{An all-exterior curvature comparison, including zero staples}

\begin{proposition}[Radial monotonicity and explicit one-link bounds]
\label{f16v63:radial}
For $\nu\ge1$, $R_\nu(t)=I_{\nu+1}(t)/I_\nu(t)$ is strictly increasing
on $t>0$, is less than one, and satisfies
\begin{equation}
 R_\nu(t)\ge\frac{t}{2\nu+2+t}.
 \label{f16v63:ratio-barrier}
\end{equation}
Consequently every $\lambda_k$ is nondecreasing, and
\begin{equation}
 \partial_{t_j}\log Z_\square(I;t)
 \ge R_1(t_j)\ge\frac{t_j}{4+t_j}.
 \label{f16v63:radial-slope}
\end{equation}
Also
\begin{equation}
 e^{-t}\mathfrak z(t)\ge
       \frac{4}{3\pi e}(1+t)^{-3/2}.
 \label{f16v63:cap}
\end{equation}
\end{proposition}
\begin{proof}
The integral representation
$I_\nu(t)=c_\nu t^\nu\int_{-1}^1 e^{tx}(1-x^2)^{\nu-1/2}dx$
shows that $R_\nu(t)$ is the mean of $x$ in this tilted probability
law.  Its derivative is the strictly positive variance; the mean is
less than one.  The Bessel recurrence gives
$R_\nu'=1-R_\nu^2-(2\nu+1)R_\nu/t$.
Put $d=2\nu+2$ and $r(t)=t/(d+t)$.  At $R_\nu=r$ the right-hand side
minus $r'$ equals $(d+1)t/(d+t)^2>0$.  Near zero,
$R_\nu=t/d+O(t^3)>r(t)$.  A first downward crossing is therefore
impossible, proving \eqref{f16v63:ratio-barrier}.

The product $\lambda_k=\prod_{\nu=1}^kR_\nu$ is nondecreasing.
Differentiate \eqref{f16v63:exact-Z} at $H=I$.  The factors
$\mathfrak z$ contribute $R_1$, and the derivative of the positive sum
$W_t$ is nonnegative.  This proves \eqref{f16v63:radial-slope}.
For the last assertion, set $v=1-x$ in the scalar Haar integral and
restrict to $0\le v\le(1+t)^{-1}$:
\[
 e^{-t}\mathfrak z(t)
 =\frac2\pi\int_0^2e^{-tv}\sqrt{v(2-v)}\,dv
 \ge\frac{2}{\pi e}\int_0^{(1+t)^{-1}}\sqrt v\,dv.
\]
\end{proof}

\begin{theorem}[Amplitude and cyclic-holonomy costs in one exact action]
\label{f16v63:global}
Let $\Psi_{\square,0}$ be the value with all slots equal to $e_0$ and
$B=I$, and put $D_*=4+\gamma\max_i n_i$.  When all $r_i>0$, set
$\kappa_{\rm f}(t)=4b_1(t)/W_t$.  Then
\begin{equation}
 \boxed{\quad
 \Psi_\square(p,B)-\Psi_{\square,0}
 \ge\frac{\gamma^2}{2D_*}\sum_i(n_i^2-r_i^2)
                    +\kappa_{\rm f}(t)(1-w(H)).\quad}
 \label{f16v63:global-bound}
\end{equation}
The first term equals
$(\gamma^2/D_*)\sum_{i,\alpha<\beta}(1-p_{i\alpha}\cdot p_{i\beta})$.
It remains valid if a staple sum is zero, with the second term zero.
For the native $n_i=5$ cell, an everywhere-defined closure term is
\begin{align}
 \Delta_B(S)&=\prod_i r_i-w(B^{-1}S_1S_2S_3S_4)\ge0,
 \label{f16v63:raw-closure}\\
 \Psi_\square-\Psi_{\square,0}
 &\ge\frac{\gamma^2}{2(4+5\gamma)}\sum_i(25-r_i^2)
                                  +c_\gamma\Delta_B(S),
 \label{f16v63:continuous-bound}\\
 c_\gamma&=4R_1(\gamma)
       \left(\frac{\gamma}{4+5\gamma}\right)^4
                            e^{-\gamma}\mathfrak z(\gamma)>0.
 \label{f16v63:raw-coefficient}
\end{align}
In particular this closure coefficient has a polynomial lower bound
from \eqref{f16v63:cap}, not an empty-cell exponential penalty.
\end{theorem}
\begin{proof}
Unitary characters satisfy $\chi_k(H)\le d_k$.  Hence
$0<F_t(H)\le W_t$ and, retaining $k=1$,
\[
 W_t-F_t(H)=\sum_k d_kb_k(t)(d_k-\chi_k(H))
                    \ge4b_1(t)(1-w(H)).
\]
The inequality $-\log u\ge1-u$ for $u>0$ pays the second term of
\eqref{f16v63:global-bound}.  For the first, interpolate the aligned
amplitudes from $h_i=\gamma r_i$ to $\gamma n_i$ and apply
\eqref{f16v63:radial-slope}.  On that interval the derivative with
respect to $r_i$ is at least $\gamma^2r_i/D_*$.  Integration and the
identity $n_i^2-r_i^2=2\sum_{\alpha<\beta}(1-p_{i\alpha}\cdot p_{i\beta})$
prove the assertion.  The two costs are the two successive differences
of the \emph{same} logarithm, so they are not competing bounds added
without payment.

Convolution by a probability density contracts the supremum norm, so
$W_t\le k_\gamma(I)=e^\gamma/\mathfrak z(\gamma)$.  Moreover
$b_1=R_1(\gamma)\prod_iR_1(\gamma r_i)$ and
$R_1(\gamma r_i)\ge\gamma r_i/(4+5\gamma)$.
It follows that
$\kappa_{\rm f}(t)\ge c_\gamma\prod_i r_i$.
Multiplying by $1-w(H)$ gives \eqref{f16v63:continuous-bound}.
Quaternion norm multiplicativity proves \eqref{f16v63:raw-closure};
if one sum is zero, both its product norm and its product quaternion
are zero.  Thus no polar-coordinate singularity remains in that bound.
\end{proof}

The slot contrasts are literal retained closed paths of length at most
six.  The normalized product $H$, in contrast, is the holonomy of the
\emph{aggregate} staple directions; it is not generally one retained
Wilson path.  The polynomial numerator in \eqref{f16v63:raw-closure}
is a specified sum of products of actual paths.  This distinction must
be kept in a pullback to retained-link or source coefficients.

\subsection{The complete mixed Hessian and the actual sign obstruction}

At coherence $p_{i\alpha}=e_0$, $B=I$, put
$t=(\gamma,\gamma n_1,\ldots,\gamma n_4)$ and
\begin{equation}
 \pi_k=\frac{d_k^2b_k(t)}{W_t},\quad
 a_k=\frac{k(k+2)}3,\quad
 \kappa(t)=\sum_k\pi_k a_k,\quad
 m_j(t)=\partial_{t_j}\log Z_\square(I;t).
 \label{f16v63:stiffness}
\end{equation}
Here $m_0=\langle w(Q)\rangle$ and $m_i=\langle x_i^0\rangle$ for
$i=1,\ldots,4$.  They are positive and at most one.  A tangent consists
of $b_{i\alpha}\in\mathbb R^3$ and a twist velocity
$\beta\in\operatorname{Im}\mathbb H$.  Write
$\bar b_i=n_i^{-1}\sum_\alpha b_{i\alpha}$.

\begin{theorem}[Signed closure, full kernel, and mixed covariance]
\label{f16v63:Hessian}
For arbitrary smooth unit-slot and group-twist curves through coherence,
\begin{equation}
 \boxed{\quad
 \mathcal Q_\square(b,\beta)
 =\gamma\sum_i m_i\sum_\alpha|b_{i\alpha}-\bar b_i|^2
               +\kappa(t)\left|\sum_i\bar b_i-\beta\right|^2.
 \quad}
 \label{f16v63:Q}
\end{equation}
This is the complete Hessian of $-\log Z_\square$, including its covariance
subtraction.  All its displayed coefficients are strictly positive for
$\gamma>0$.  With the twist fixed, its kernel consists exactly of
\begin{equation}
 b_{i\alpha}=b_i\quad\hbox{for every slot},\qquad \sum_i b_i=0.
 \label{f16v63:kernel}
\end{equation}
Equivalently $b_i=\omega_i-\omega_{i+1}$ with cyclic vertex velocities
$\omega_i\in\mathbb R^3$.  The unrestricted slot-space kernel has
dimension nine.  Allowing the twist to vary replaces the last constraint
by $\sum_i b_i=\beta$ and gives dimension twelve.

For one transverse coordinate, let
$Y=(\operatorname{Im}Q,\operatorname{Im}X_1,\ldots,
\operatorname{Im}X_4)$ denote the five scalar coordinate functions and
let $s=(-1,1,1,1,1)$.  Their exact covariance is
\begin{equation}
 \boxed{\quad
 \operatorname{Cov}(Y_j,Y_l)
 =\mathbf1_{j=l}\frac{m_j}{t_j}
                    -\kappa(t)\frac{s_js_l}{t_jt_l}.
 \quad}
 \label{f16v63:signed-C}
\end{equation}
In particular
\begin{equation}
 \langle x_i^a x_j^a\rangle
           =-\frac{\kappa(t)}{h_i h_j}<0\qquad(i\ne j).
 \label{f16v63:negative-cross}
\end{equation}
\end{theorem}
\begin{proof}
At a coherent slot family $r_i'=0$ and
$r_i''=-\sum_\alpha|b_{i\alpha}-\bar b_i|^2$.
This follows by differentiating $r_i^2=S_i\cdot S_i$, using
$e_0\cdot p_{i\alpha}''=-|b_{i\alpha}|^2$.
The normalized direction has first derivative $\bar b_i$; the first
velocity of $H$ is $\sum_i\bar b_i-\beta$.
The radial part of the second derivative of $-\log Z$ is therefore
exactly the first term of \eqref{f16v63:Q}.  No second derivative of an
undefined direction is used, since this calculation is at $r_i=n_i>0$.

For the angular part, the weights of $\chi_k(e^{\theta u})$ are
$k,k-2,\ldots,-k$.  Thus $\chi_k'(I)=0$ and
$-\chi_k''(I)=d_k k(k+2)/3$.  Differentiating
\eqref{f16v63:exact-Z} gives the second term in \eqref{f16v63:Q}.
The first differential vanishes, so accelerations of the original
parameter curves contribute nothing.  Positivity of $m_i$ follows from
\eqref{f16v63:radial-slope}; positivity of $\kappa$ follows already from
$\pi_1>0$.  The kernel now follows directly.  The cyclic difference
map from four vertex vectors has only the three-dimensional constant
kernel, proving the dimension nine statement.

To check the sign subtraction explicitly, take geodesic curves.  The
first exponent derivative is
\[
 J=\gamma\sum_{i,\alpha} b_{i\alpha}\cdot\operatorname{Im}X_i
                      +\gamma\beta\cdot\operatorname{Im}Q,
\]
and minus its mean second derivative is
$\gamma\sum_i m_i\sum_\alpha|b_{i\alpha}|^2+
\gamma m_0|\beta|^2$.  Consequently
\begin{equation}
 \mathcal Q_\square
 =\gamma\sum_i m_i\sum_\alpha|b_{i\alpha}|^2
                      +\gamma m_0|\beta|^2-\operatorname{Var}(J).
 \label{f16v63:contact-minus}
\end{equation}
All field--plaquette and field--field cross terms are inside this one
variance.  Set all slots at each link to the same tangent, and compare
coefficients in \eqref{f16v63:Q} and \eqref{f16v63:contact-minus}.
This gives \eqref{f16v63:signed-C}.  Simultaneous conjugation of all
four variables rotates every imaginary part; their means vanish.
It also makes the three transverse-coordinate blocks identical and
the mixed-coordinate blocks zero.  This proves the full stated covariance.
\end{proof}

The failure of the inherited sign argument is now exact, not inferred
from the appearance of a quaternion tensor.  Even assigning separate
sign conventions to the four link coordinates cannot make every
nonzero off-diagonal entry in \eqref{f16v63:negative-cross} nonnegative:
the product of the three signs around any triangle is invariant under
such changes and is negative.  This does not contradict positivity of
the \emph{matrix} covariance or of \eqref{f16v63:Q}.

A low-order independent fixture makes the tensor sign equally visible.
Multilinearity and $\int xx^{\mathsf T}d\sigma=I_4/4$ give
\begin{equation}
 \int w(X_1X_2X_3X_4)\prod_{i=1}^4(v_i\cdot x_i)
                  \prod_i dH(X_i)
       =4^{-4}w(v_1v_2v_3v_4).
 \label{f16v63:tensor-sign}
\end{equation}
Taking $(v_1,v_2,v_3,v_4)=(\mathbf i,\mathbf i,1,1)$ gives $-1/256$.
The character convolution groups the complete compact integral; it does
not claim that the coordinate tensor or every expansion monomial is
nonnegative.

\subsection{A paid linear large-coefficient stiffness bound}

\begin{proposition}[Uniform stiffness bounds and the native limiting coefficient]
\label{f16v63:stiffness-bounds}
For five strictly positive parameters, put
\[
 s_t=\left(\sum_{j=0}^4t_j^{-1}\right)^{-1},\qquad
 c_5=\frac{2e^{-11}}{9\pi\,6^{3/2}}.
\]
Then
\begin{equation}
 0<\kappa(t)\le
       \left(\sum_{j=0}^4(t_jm_j)^{-1}\right)^{-1}\le s_t,
 \qquad \kappa(t)\ge c_5s_t\quad(s_t\ge1).
 \label{f16v63:stiffness-two-sided}
\end{equation}
These constants have no time-length or surrounding-volume dependence.
They are not optimized.  For the native coherent layer,
$t=(\gamma,5\gamma,5\gamma,5\gamma,5\gamma)$,
\begin{equation}
 \kappa(t)/\gamma\longrightarrow5/9\qquad(\gamma\to\infty).
 \label{f16v63:stiffness-limit}
\end{equation}
More generally $t=(\gamma,\alpha_1\gamma,\ldots,\alpha_4\gamma)$ with
fixed positive $\alpha_i$ gives the limit
$(1+\sum_i\alpha_i^{-1})^{-1}$.
\end{proposition}
\begin{proof}
The matrix in \eqref{f16v63:signed-C} is positive semidefinite.
Writing it as a positive diagonal matrix minus a rank-one matrix gives
$\kappa\sum_j(t_jm_j)^{-1}\le1$.  Since $m_j\le1$, the upper bounds follow.

For the lower bound, the ratio barrier gives
\[
 \lambda_k(t_j)\ge
 \prod_{\nu=1}^k(1+(2\nu+2)/t_j)^{-1}
 \ge\exp[-k(k+3)/t_j],\qquad
 b_k(t)\ge\exp[-k(k+3)/s_t].
\]
Let $s=s_t\ge1$.  Take the $\lfloor\sqrt s\rfloor$ integers starting at
$\lceil\sqrt s\rceil$.  Their number is at least $\sqrt s/2$, each is
between $\sqrt s$ and $2\sqrt s$, and $b_k\ge e^{-10}$ on this set.
The numerator of \eqref{f16v63:stiffness} is therefore at least
$e^{-10}s^{5/2}/6$.
For $t_{\min}=\min_jt_j$, convolution contraction and
\eqref{f16v63:cap} give
\[
 W_t\le\frac{e^{t_{\min}}}{\mathfrak z(t_{\min})}
 \le\frac{3\pi e}{4}(1+t_{\min})^{3/2}
 \le\frac{3\pi e}{4}(6s)^{3/2}.
\]
The last step uses $t_{\min}\le5s$ and $s\ge1$.
Division yields exactly the stated $c_5$.

For the limit, keep the original compact density and write its
nonnegative energy as
\[
 E(X)=1-w(Q)+\sum_i\alpha_i(1-w(X_i)).
\]
It has the unique minimum $X_i=I$.  In a small chart with positive
scalar parts, write
\[
 X_i=(\sqrt{1-|x_i|^2},x_i),\qquad y_i=\sqrt\gamma x_i.
\]
Taylor expansion of the literal quaternion product gives
\[
 \gamma E(X)\longrightarrow
    \tfrac12\left(\sum_i\alpha_i|y_i|^2+
                                     |\textstyle\sum_i y_i|^2\right).
\]
This limiting quadratic integral is a consequence of the compact law,
not its substitution at finite $\gamma$.  Indeed outside a fixed chart
the pinning energy has a strictly positive minimum.  Its integral,
even after multiplying by any fixed power of $\gamma$, is exponentially
small relative to the $c\gamma^{-6}$ lower bound on the partition
function obtained from $|y_i|\le1$.  Inside the chart,
$1-w(X_i)\ge|x_i|^2/2$ and the Haar Jacobians are bounded.
Thus $e^{-\frac12\sum\alpha_i|y_i|^2}$ times a fixed polynomial is an
integrable majorant for every moment needed here.  Dominated convergence
therefore proves convergence of the normalized second and fourth moments.
Also $\sqrt\gamma\operatorname{Im}Q\to\sum_i y_i$ with the same
moment control, and $m_0\to1$.

For each color the limiting covariance of $(y_i)$ is
$A^{-1}$, where $A=\operatorname{diag}(\alpha_i)+\mathbf1\mathbf1^{\mathsf T}$.
If $S=\sum_i\alpha_i^{-1}$, the rank-one inverse identity gives
$\mathbf1^{\mathsf T}A^{-1}\mathbf1=S/(1+S)$.
By the twist-only case of \eqref{f16v63:contact-minus},
$\kappa/\gamma=m_0-\gamma\operatorname{Var}(\operatorname{Im}Q^a)$.
Its limit is $1/(1+S)$.  Taking $\alpha_i=5$ gives $5/9$.
\end{proof}

\subsection{Five actual conditional writers and their exact scaled covariance}

The preceding Hessian marks directions of the retained action; it is
not a source entrance Gram.  The latter can also be specified explicitly
for this cell.  At the coherent conditional boundary let
\begin{equation}
 A_0=w(Q),\qquad A_i=w(X_i)\quad(1\le i\le4),\qquad
 G_\square(t)=\operatorname{Cov}_t(A_0,\ldots,A_4).
 \label{f16v63:source-bank}
\end{equation}
Each $A_i$ with $i>0$ can be written as
$w(P_{i1}^{-1}X_i)$ using one chosen complementary retained path; at
coherence that path is $I$.  Thus it is the restriction of a gauge-invariant
Wilson writer, not an unclosed link trace in the full theory.
Independent variations of its slot coefficient generate the parameter
$t_i$ at coherence.  Consequently
\begin{equation}
 G_\square(t)=D_t^2\log Z_\square(I;t).
 \label{f16v63:source-generating}
\end{equation}
For a directly evaluable expression set
\[
 u_k(t)=R_{k+1}(t)+k/t,\qquad
 u_k'(t)=1-R_{k+1}(t)^2-(2k+3)R_{k+1}(t)/t-k/t^2.
\]
Differentiating the positive character sum gives
\begin{equation}
 (G_\square)_{jl}
 =\operatorname{Cov}_{\pi}(u_k(t_j),u_k(t_l))
                      +\mathbf1_{j=l}\mathbb E_{\pi}u_k'(t_j).
 \label{f16v63:source-series}
\end{equation}
The last term need not be nonnegative term by term.  Positivity belongs
to the complete original covariance, not to a discarded diagonal
contribution in this representation.

\begin{theorem}[Exact source rank and a noncollapsing conditional limit]
\label{f16v63:source-Gram}
For every finite positive parameter vector $t$, $G_\square(t)$ is
strictly positive definite.  For the native ray $(\gamma,5\gamma,5\gamma,
5\gamma,5\gamma)$,
\begin{equation}
 \gamma^2G_\square(t)\longrightarrow
 G_*:=\frac1{1350}
 \begin{pmatrix}
 400&25&25&25&25\\
 25&64&1&1&1\\
 25&1&64&1&1\\
 25&1&1&64&1\\
 25&1&1&1&64
 \end{pmatrix},\qquad G_*\succeq\frac{18}{467}I_5.
 \label{f16v63:source-limit}
\end{equation}
The centered writers $\gamma(A_j-\mathbb E A_j)$ therefore have a
nondegenerate five-dimensional conditional limiting Gram.  Adding
$A_5=A_1+A_2$ gives exactly the kernel
\[
 \operatorname{span}\{(0,-1,-1,0,0,1)^{\mathsf T}\}
\]
at every regulator.
\end{theorem}
\begin{proof}
The compact cell density is everywhere positive.  A zero variance
would imply a linear combination of the five writers is constant
Haar-almost everywhere.  Under independent Haar links they are centered,
mutually orthogonal, and each has squared norm $1/4$.  For example,
in a product $w(Q)w(X_i)$, integrating any other link first gives zero.
The remaining orthogonality statements follow from independent Haar
integration.  Thus all five coefficients must vanish.  Congruence gives
the assertion for the redundant sixth writer.

The dominated moment argument in the preceding proposition applies also
to $\gamma(1-A_j)$ and their products.  They converge in moments to
$\tfrac12|\sum_i y_i|^2$ and $\tfrac12|y_i|^2$, respectively, with three
independent color copies and covariance
\[
 \Sigma=(5I_4+\mathbf1\mathbf1^{\mathsf T})^{-1}
                         =\tfrac15I_4-\tfrac1{45}\mathbf1\mathbf1^{\mathsf T}.
\]
For $v_0=\mathbf1$, $v_i=e_i$, direct Gaussian fourth-moment pairing
makes the limiting covariance equal to
$\tfrac32(v_j^{\mathsf T}\Sigma v_l)^2$.
Here $v_0^{\mathsf T}\Sigma v_0=4/9$,
$v_0^{\mathsf T}\Sigma e_i=1/9$, $\Sigma_{ii}=8/45$ and
$\Sigma_{ij}=-1/45$ for $i\ne j$, proving the displayed matrix.
This evaluation is of the rigorously obtained limit, not a Gaussian
assumption about the finite cell.

On the three link-contrast directions the matrix has eigenvalue
$7/150$.  On the orthonormal span of the plaquette coordinate and
one half the sum of the four link coordinates, its block is
$1350^{-1}\left(\begin{smallmatrix}400&50\\50&67\end{smallmatrix}\right)$.
Its determinant is $1/75$ and its trace is $467/1350$.
The smaller eigenvalue is at least determinant divided by trace,
namely $18/467$; the contrast eigenvalue is larger.  This proves the
claimed uniform lower bound on the limiting matrix.
\end{proof}

This noncollapse is conditional on the declared coherent exterior.
It supplies neither a lower bound on the probability of that exterior
nor physical-time continuity of the unsmeared writers.  In particular
it is not a way around the source-regularization obstruction in V60.
A declared source coefficient map, including a regulator-dependent one,
pulls every finite covariance identity back by congruence; its own
limit and null relations must still be specified.

\subsection{Analytic truncation bounds and numerical evaluation}

There is a scalar series rather than a twelve-dimensional integral to
evaluate in \eqref{f16v63:exact-Z}.  Its tails can be bounded without
assuming that a floating-point summand is zero.  Let $c=t_{\min}/2$,
truncate at $k=K$, and for $p\in\{2,4\}$ put
\begin{equation}
 \rho_p(K)=c\frac{(K+3)^{p-1}}{(K+2)^p},\qquad
 T_p(K)=\frac{(K+2)^p c^{K+1}}{(K+2)!}\,
                          \frac1{1-\rho_p(K)}
                   \quad(\rho_p(K)<1).
 \label{f16v63:tails}
\end{equation}

\begin{proposition}[Paid character tails]
\label{f16v63:tail-proposition}
Writing $W_K=\sum_{k\le K}d_k^2b_k$ and
$N_K=\sum_{k\le K}d_k^2b_k a_k$, the tails satisfy
\begin{equation}
 0\le W_t-W_K\le T_2(K),\quad
 |F_t(H)-F_{t,K}(H)|\le T_2(K),\quad
 0\le N_t-N_K\le T_4(K)/3.
 \label{f16v63:tail-bounds}
\end{equation}
In particular
\begin{equation}
 \frac{N_K}{W_K+T_2(K)}\le\kappa(t)\le
                  \frac{N_K+T_4(K)/3}{W_K}.
 \label{f16v63:stiffness-enclosure}
\end{equation}
These are analytic truncation bounds.  Evaluating the finite sums with
ordinary floating-point Bessel routines does not by itself provide an
outward-rounded interval certificate.
\end{proposition}
\begin{proof}
By \eqref{f16v63:lambda-upper}, $b_k\le c^k/(k+1)!$.
For $f_p(k)=(k+1)^pc^k/(k+1)!$, the successive ratio
$c(k+2)^{p-1}/(k+1)^p$ decreases with $k$.  Starting at $K+1$, it is
at most $\rho_p(K)$, so its tail is bounded by the first term divided
by $1-\rho_p(K)$.  Use $|\chi_k|\le d_k$ and
$a_k\le(k+1)^2/3$.  Positivity of $W_K$ then gives the final enclosure.
\end{proof}

For orientation, direct character-sum evaluations on the native ray are
shown below.  These rounded values are diagnostics; the preceding bounds
and identities, not this table, are the proofs.
\begin{center}\small
\begin{tabular}{@{}rrrr@{}}
\toprule
$\gamma$ & $\kappa$ & $\kappa/\gamma$ & $\langle x_1^a x_2^a\rangle$\\
\midrule
$0.5$ & $0.0320078$ & $0.0640157$ & $-0.00512125$\\
$1$   & $0.2274370$ & $0.2274370$ & $-0.00909748$\\
$2$   & $0.7586043$ & $0.3793022$ & $-0.00758604$\\
$5$   & $2.4168929$ & $0.4833786$ & $-0.00386703$\\
$20$  & $10.7471345$ & $0.5373567$ & $-0.00107471$\\
$100$ & $55.1908490$ & $0.5519085$ & $-0.000220763$\\
\bottomrule
\end{tabular}
\end{center}

\subsection{One actual alternate-layer elimination, with sources retained}

Return to the four-link plaquette tube at times $1,\ldots,T$.
There are $T$ four-free-link plaquettes, $4(T-1)$ two-free-link temporal
plaquettes, and $12T+8$ one-free-link plaquettes.  Every integrated link
has total incidence six.  In particular, there is no hidden pair bond
between two different edges of the square in one layer: their common
plaquette is the retained four-link product itself.

Let $\mathcal I$ be one parity of the time layers.  Keep the opposite
parity and all exterior links.  No Wilson plaquette connects two distinct
layers in $\mathcal I$, so those four-link cells are conditionally
independent.  Each has exactly five complementary slots per link:
three spatial staples and two temporal staples, the latter including
fixed endpoint links when necessary.

\begin{theorem}[Exact temporal decimation and its coherent pullback]
\label{f16v63:decimation}
For the original finite Wilson law with local boundary data, let
$S_{\rm rem}$ contain exactly the plaquette terms having no link in an
eliminated layer.  Up to a configuration-independent additive constant,
the retained action after integrating all layers $t\in\mathcal I$ is
\begin{equation}
 \boxed{\quad S_{\rm ret}(\xi)=S_{\rm rem}(\xi)
             +\sum_{t\in\mathcal I}\Psi_\square(p_t(\xi),I).
 \quad}
 \label{f16v63:retained-action}
\end{equation}
Here every term is the exact formula \eqref{f16v63:exact-Z}, not a new
Wilson or independent-ribbon ansatz.  Under admissible all-identity
boundary comparison, its nonlinear action difference is bounded below
by the surviving Wilson defects plus the sum of
\eqref{f16v63:continuous-bound} over the eliminated layers.

For an actual retained-link tangent $v$ at the identity exterior and
identity retained links, let $b_{it\alpha}(v)$ be the derivative of the
corresponding three-link path, in the cyclic orientation.  Then
\begin{equation}
 \begin{split}
 D^2S_{\rm ret}(\mathbf1)[v,v]
  ={}&\gamma\sum_{p\text{ surviving}}|(dv)_p|^2\\
 &+\sum_{t\in\mathcal I}\left\{
   \gamma m_{\rm link}\sum_{i,\alpha}
                    |b_{it\alpha}-\bar b_{it}|^2
       +\kappa(\gamma,5\gamma,5\gamma,5\gamma,5\gamma)
                         |\textstyle\sum_i\bar b_{it}|^2\right\}.
 \end{split}
 \label{f16v63:retained-Hessian}
\end{equation}
Here $m_{\rm link}\ge5\gamma/(4+5\gamma)$.  The closure coefficient
obeys \eqref{f16v63:stiffness-two-sided} and the limit
\eqref{f16v63:stiffness-limit}.  None of these coefficients depends on
the number of layers.  No size-independent norm estimate on the full
retained tangent space is asserted.
\end{theorem}
\begin{proof}
Each eliminated link occurs in its one four-link spatial plaquette and
its five complementary one-link factors once neighboring layers have
been retained.  The chosen parity has no nearest-time neighbors inside
it.  The Wilson action consequently separates into $S_{\rm rem}$ and
one conditional cell action per selected layer.  Fubini proves
\eqref{f16v63:retained-action}, retaining its normalizing partition
function.  The global bound follows by adding the inequalities for
these distinct additive terms; no plaquette or induced term is counted
twice.  Differentiation at coherence uses
\eqref{f16v63:Q} with $\beta=0$.  Its first differential is zero, so
accelerations of the actual retained paths produce no extra term.
Each surviving Wilson plaquette has quadratic action
$\gamma|(dv)_p|^2$.  This proves \eqref{f16v63:retained-Hessian}.
\end{proof}

All path sums transform by their endpoint gauge transformations;
normalized nonzero sums do likewise.  Their cyclic product transforms
by conjugation at the square basepoint.  At zero sums the original
integral is used.  Thus the retained law is gauge invariant.  If the
removed set is reflection invariant and the original boundary law is
reflection positive, its retained marginal is reflection positive on
the retained positive-time algebra, since those expectations are the
identical expectations in the original law.  An arbitrary nonlocal
Perron endpoint density cannot be silently absorbed into the separate
cell factors.  For a stationary vacuum the assertion is instead used
conditional on outside links of a finite interval, or after an
appropriate fixed-regulator cylinder limit with boundary weights explicit.

Sources can be retained without postulating a new transfer.  For any
finite vector $F$ of square-integrable original writers, disintegration gives
its exact retained predictor $\widehat F=\mathbb E(F\mid\xi)$ and
\begin{equation}
 \operatorname{Cov}(F)=
 \mathbb E_{\rm ret}\operatorname{Cov}(F\mid\xi)
                       +\operatorname{Cov}_{\rm ret}(\widehat F).
 \label{f16v63:total-source}
\end{equation}
Source insertions supported in one cell are derivatives of that cell
integral.  At coherence the five specified writers give
\eqref{f16v63:source-generating}; elsewhere the same insertions are taken
in the unsimplified vector-field integral, including cancellations.
Exact linear source relations are preserved by conditional expectation.
For time-separated writers the analogous mixed covariance identity
must use the same retained sigma-field for both, not two incompatible
conditional laws.

\subsection{What is now paid, and the next noncircular comparison}

The former four-link interface is resolved in a definite sense: a whole
compact plaquette cell, with all its complementary Wilson slots, has an
exact all-exterior integral, a positive signed coherent response, and
an explicit conditional source Gram.  It can be eliminated on alternate
time layers of the actual tube.  Negative transverse correlations are
part of that resolution, not exceptional terms to suppress.

The next integration is \emph{not} a repetition of the five-factor
convolution.  A surviving layer appears in the aggregate field vectors
of its adjacent eliminated layers.  Their functions $F_t(H_t)$ and
radial factors couple its four links nonlinearly.  A general conditional
law of that layer is therefore not four independent one-link fields
plus just one Wilson plaquette.  Its induced interactions must remain.
For a retained parameter $u$, integration of further variables $Y$
satisfies the exact signed identity
\begin{equation}
 \partial_u\partial_v\left[-\log\int e^{-S_{\rm ret}(Y;u,v)}dH(Y)\right]
 =\mathbb E[\partial_u\partial_vS_{\rm ret}]
     -\operatorname{Cov}(\partial_uS_{\rm ret},\partial_vS_{\rm ret}).
 \label{f16v63:next-signed}
\end{equation}
Positivity of the Hessian only at the all-identity retained configuration
does not control either term of this averaged expression.  Likewise,
\eqref{f16v63:source-limit} concerns neither the retained predictor term
in \eqref{f16v63:total-source} nor the probability of its conditional
boundary configurations.

There is now a small exact starting integral for an actual two-time
comparison, rather than an unspecified effective interaction.  For two
adjacent layers with flat complementary links, retain $Y_i$ in the second
layer.  Each link has four fixed slots outside its temporal bond.
With plaquette marks $J_1,J_2$, the complete two-layer generating
partition function is
\begin{equation}
 \begin{split}
 Z_2(J_1,J_2)=\int &\exp\{(\gamma+J_2)w(Y_1Y_2Y_3Y_4)
                              +4\gamma\textstyle\sum_iw(Y_i)\}\\
 &\times Z_\square\bigl(t_0=\gamma+J_1,
                  \ h_i=\gamma|4e_0+y_i|,\\
 &\hspace{5.5em}P_i=(4e_0+y_i)/|4e_0+y_i|,\ B=I\bigr)
                                           \prod_i dH(Y_i).
 \end{split}
 \label{f16v63:two-layer-target}
\end{equation}
In this formula the notation $Z_\square$ allows its plaquette coefficient
$t_0$ to differ from its field amplitudes, as in
\eqref{f16v63:exact-Z}.  The orientation convention for a reversed
cyclic edge reverses both its temporal path and its spin, so at the
flat boundary the temporal factor is still $x_i\cdot y_i$.
Then $\partial_{J_1}\partial_{J_2}\log Z_2(0,0)$ is precisely the
conditional connected covariance of the two native plaquette writers.
This is a concrete next local calculation, but the upstream target
for the continuum program is stronger: a source-weighted estimate
under the \emph{actual retained marginal}, including exterior-predictor
variation and the induced mixed-score covariance.  Another flat-cell
coercivity calculation alone will not supply it.

Neither the one-layer integration nor the alternate-layer marginal
rescales $a$, identifies a conditional update with physical evolution,
or supplies a new physical Hamiltonian.  Uniform dense-core physical-time
contraction and a nontrivial interacting continuum source limit remain
separate unproved obligations.

\subsection{Verification, append-only preservation, and result ledger}

The independent companion diagnostic is
\begin{center}\small
\path{verify_ym_f16_v63_four_link_cell.py}.
\end{center}
It separates exact rational quaternion-moment and covariance-kernel
fixtures, literal Wilson incidence and gauge identities, scalar
character calculations, and numerical compact integration.  The latter
integrates the fourth link analytically and uses deterministic
low-discrepancy Haar points for the other three; no spin is replaced
by a noncompact Gaussian.  Limiting Gaussian
moments are checked separately against the proved compact limit.
Analytic series tails are reported separately from floating-point
roundoff.  These checks are not full Wilson-vacuum simulations,
outward-rounded interval certificates, or proof-assistant verification.

Removing this delimited appendix and reverting the single displayed
version change recovers V62 byte for byte.  The earlier appendices and
the comparison archives are not rewritten.

\begin{record}[V63 result ledger]
\label{f16v63:ledger}
\emph{Proved here:} the exact five-factor character convolution for a
four-free-link Wilson cell; amplitude and cyclic-holonomy lower bounds
including zero staple sums; the complete signed coherent mixed Hessian;
its cyclic gauge kernel and strictly negative cross-link transverse
covariances; uniform stiffness bounds and the native coefficient $5/9$;
a full-rank five-writer conditional covariance and its explicit positive
scaled limit; analytic character-tail bounds; and an exact alternate-time
elimination retaining every induced interaction and source predictor.

\emph{Inherited:} the literal Wilson convention, compact Haar law,
local boundary discipline, actual retained-marginal interpretation,
physical clock, and V60's distinction between conditional fast variables
and surviving physical sources.  The V62 pair-strand positivity theorem
is not used for the four-link tensor.

\emph{Not asserted:} all-exterior Hessian positivity; a positive-entry
quaternion expansion; another independent cell law after marginalization;
a source covariance under the full native vacuum inferred from a coherent
conditional calculation; or a physical gap from the local $O(\gamma)$
closure coefficient.

\emph{Still unproved:} source-weighted control of the interacting retained
law and its induced signed response; a nontrivial physical continuum
source limit; physical-time contraction with common dense-core constants;
and the continuum Yang--Mills mass gap.
\end{record}

\begin{firewall}
The quartic product is integrated, not linearized away.  Its exact
negative correlations and the full covariance subtraction are retained.
The conditional five-writer limit is explicitly distinguished from a
physical source limit.  A local action bound is not a normalized-vacuum
comparison.  V1--V62 remain unchanged.
\end{firewall}
% END F16 V63 APPEND

% BEGIN F16 V64 APPEND -- 2026-09-18
% Append-only continuation of V63.  The preceding mathematical body is unchanged.
\clearpage
\section[V64: retained two-time sources and a native cancelling exterior]{V64: retained two-time sources\\and a native cancelling exterior}
\label{f16v64:context}

The two-layer integral in \eqref{f16v63:two-layer-target} is a normalized
retained-source problem, not another Hessian at a prescribed retained
configuration.  We first evaluate its leading complete mixed covariance
and its retained predictor.  The induced interaction is essential: its
quadratic contribution is a Schur complement, and discarding it changes
the answer.  The compact localization method of V63 is reused; a generic
finite-dimensional Laplace theorem is not claimed as a new result.

We then move upstream from the coherent boundary.  There is a literal
Wilson exterior in which the three spatial staple paths of each link
cancel.  In its actual conditional tube law, four bounded, closed Wilson
writers have an exact noncollapsing Gram and a one-step correlation ratio
tending to one.  This gives a concrete obstruction to extending a coherent
boundary contraction uniformly to all exteriors.  It is not a counterexample
to a continuum physical gap.  Finally, an elementary pressure estimate
bounds the occurrence of the same cancellation event in the unconditioned,
translation-invariant Wilson measure, uniformly in volume.  The distinction
between this probability estimate and a relative source-weighted estimate
is kept explicit.

\subsection{Ownership, normalization, and the literal two-layer law}

Only the fixed-regulator Wilson convention, cyclic link orientation,
positive character coefficients, local boundary discipline, and compact
moment argument of V63 are inherited.  None of the parallel V67 cubic
exponential-moment argument is used.  Previous generic Gaussian posterior,
Schur, and rare-region observations are not reproved as independent
advances.  The new objects are the specified quartic two-layer covariance,
its predictor fraction, the cancelling native spatial exterior and its
conditional source spectrum, and a native probability bound for that
particular exterior defect.

Write $Q(X)=X_1X_2X_3X_4$ in the cyclic convention of V63.  At a flat
complementary boundary the full two-layer probability density is
\begin{equation}
 \begin{split}
 d\nu_\gamma(X,Y)={}&Z_{2,\gamma}^{-1}\exp\!\Bigl\{\gamma\Bigl[
 w(Q(X))+w(Q(Y))\\
 &\qquad+4\sum_{i=1}^4(w(X_i)+w(Y_i))
                  +\sum_{i=1}^4x_i\cdot y_i\Bigr]\Bigr\}
 \,dH^4(X)dH^4(Y).
 \end{split}
 \label{f16v64:two-law}
\end{equation}
The four fixed slots per link consist of three spatial slots and the
outer temporal slot.  The temporal bond between $X_i$ and $Y_i$ is counted
once.  Integrating $X$ gives precisely \eqref{f16v63:two-layer-target},
including the dependence of every aggregate $4e_0+y_i$ on $Y$.  Thus the
$Y$ marginal in this section is that induced marginal, not its value at
$Y=I$ and not a product one-link law.

Let $A(X)=w(Q(X))$.  Its restriction here is a closed plaquette writer.
One can also use the five writers $A_0(X)=A(X)$ and
$A_i(X)=w(P_{i,\mathrm{out}}^{-1}X_i)$, $1\le i\le4$, where the selected
outer three-link paths are retained; at the flat boundary they are $I$.
All centering below is with respect to \eqref{f16v64:two-law} or its exact
marginal, as specified.

\begin{proposition}[A strict normalized two-time comparison at every coefficient]
\label{f16v64:positive-cross}
For any finite square-integrable real source family $F$ whose components are independent
modulo constants under Haar measure, define
\[
 G_\gamma=\operatorname{Cov}_{\nu_\gamma}(F(X)),\qquad
 C_\gamma=\operatorname{Cov}_{\nu_\gamma}(F(X),F(Y)).
\]
Then, for $\gamma>0$,
\begin{equation}
 0\prec C_\gamma\prec G_\gamma.
 \label{f16v64:strict-Gram}
\end{equation}
The assertions apply in particular to the five actual writers above.
For a redundant coefficient map they hold on its exact source quotient.
No coefficient-uniform strict margin is asserted by this proposition.
\end{proposition}
\begin{proof}
Put $a_\gamma(X)=\exp\{\gamma[w(Q(X))+4\sum_iw(X_i)]\}$ and let
$\mathcal C_\gamma$ have kernel $\exp(\gamma\sum_i x_i\cdot y_i)$.
Each one-link convolution has strictly positive Fourier multipliers:
its normalized multipliers are $\lambda_k(\gamma)$ from
\eqref{f16v63:lambda}.  Their fourfold tensor product is a positive,
injective operator on $L^2(dH^4)$.  If $f$ is centered in the actual
marginal, its cross covariance is
\[
 Z_{2,\gamma}^{-1}\langle a_\gamma f,
                  \mathcal C_\gamma(a_\gamma f)\rangle>0
\]
unless $f=0$.  On the other hand
$\operatorname{Var}(f(X))-\operatorname{Cov}(f(X),f(Y))
=\tfrac12\mathbb E(f(X)-f(Y))^2>0$ for a nonconstant $f$, since the
joint density is strictly positive everywhere.  Apply both assertions
to every real source combination.  Independence for the five writers
follows from the same Haar orthogonality argument as in V63.  Congruence
preserves exactly the prescribed linear relations.
\end{proof}

\subsection{The induced marginal is paid: the two-time covariance and predictor}

For a scalar color coordinate put $J_4=\mathbf1\mathbf1^{\mathsf T}$,
$B=5I_4+J_4$, and
\begin{equation}
 \mathcal A=\begin{pmatrix}B&-I_4\\-I_4&B\end{pmatrix}.
 \label{f16v64:precision}
\end{equation}
The three color copies use this same matrix.  Set
$P_0=J_4/4$, $P_\perp=I_4-P_0$.  Its inverse has diagonal and off-diagonal
blocks
\begin{equation}
 S=\frac5{24}P_\perp+\frac9{80}P_0,
 \qquad R=\frac1{24}P_\perp+\frac1{80}P_0.
 \label{f16v64:inverse-blocks}
\end{equation}
In particular $S=(B-B^{-1})^{-1}$, not $B^{-1}$.

\begin{theorem}[Complete two-layer source limit under the actual retained law]
\label{f16v64:two-source}
As $\gamma\to\infty$ in \eqref{f16v64:two-law},
\begin{equation}
 \boxed{\quad
 \gamma^2\operatorname{Cov}_{\nu_\gamma}
       \begin{pmatrix}A(X)\\A(Y)\end{pmatrix}
 \longrightarrow\frac3{800}
       \begin{pmatrix}81&1\\1&81\end{pmatrix}.
 \quad}
 \label{f16v64:two-limit}
\end{equation}
Consequently the normalized adjacent plaquette correlation tends to
$1/81$, not to zero and not to the one-layer coherent variance.

More generally, take $v_0=\mathbf1$, $v_i=e_i$ for $1\le i\le4$.
For the ten writers on the two layers the limiting covariance has
$5\times5$ diagonal and off-diagonal blocks
\begin{equation}
 (G_*)_{jl}=\tfrac32(v_j^{\mathsf T}Sv_l)^2,\qquad
 (C_*)_{jl}=\tfrac32(v_j^{\mathsf T}Rv_l)^2.
 \label{f16v64:ten-limit}
\end{equation}
The complete $10\times10$ limit is strictly positive definite.

Let $\widehat A_\gamma(Y)=\mathbb E_{\nu_\gamma}[A(X)\mid Y]$.
The following assertion uses the exact induced $Y$ marginal:
\begin{equation}
 \left\|\gamma(\widehat A_\gamma-\mathbb E A(X))
       -\frac{\gamma}{81}(A(Y)-\mathbb E A(Y))\right\|_{L^2(\nu_{\gamma,Y})}
 \longrightarrow0.
 \label{f16v64:predictor-limit}
\end{equation}
In particular
\begin{align}
 \gamma^2\operatorname{Var}(\widehat A_\gamma)&\longrightarrow\frac1{21600},
       \label{f16v64:predictor-var}\\
 \gamma^2\mathbb E\operatorname{Var}(A(X)\mid Y)&\longrightarrow\frac{41}{135}.
       \label{f16v64:innovation-var}
\end{align}
The retained predictor carries the fraction $1/6561$ of the entrance
variance in this particular two-layer limit.
\end{theorem}
\begin{proof}
Subtract the maximum exponent in \eqref{f16v64:two-law}.  Its exact
nonnegative compact energy is
\[
 E_2=1-w(Q(X))+1-w(Q(Y))
       +4\sum_i[(1-w(X_i))+(1-w(Y_i))]
       +\sum_i(1-x_i\cdot y_i).
\]
Its unique minimum is $X_i=Y_i=I$.  In fixed northern charts write
$X_i=(\sqrt{1-|z_i|^2/\gamma},z_i/\sqrt\gamma)$ and similarly $Y_i$
with variable $u_i$.  Expansion of the actual quaternion products gives
\[
 \gamma E_2\longrightarrow\frac12\left[
 4\sum_i(|z_i|^2+|u_i|^2)+|\textstyle\sum_i z_i|^2+
 |\textstyle\sum_i u_i|^2+\sum_i|z_i-u_i|^2\right].
\]
This is \eqref{f16v64:precision}.  The localization and normalized moment
steps can be checked directly as follows.  Outside the fixed charts the
pinning part of $E_2$ is bounded below by a positive constant; the full
partition function has a lower bound $c\gamma^{-12}$ from a ball of
radius $\gamma^{-1/2}$.  Polynomially weighted outside integrals are
therefore negligible.  Inside, the pinning part gives an integrable
Gaussian majorant and the Haar Jacobians are bounded.  The product
inequality $|Q(X)-I|\le\sum_i|X_i-I|$ also bounds
$\gamma(1-w(Q(X)))$ by a fixed quadratic polynomial in the scaled chart
variables.  Dominated convergence applies to all fixed moments used
below, including fourth and eighth moments of those scaled deficits.
This proves the compact, normalized moment limit, not just an expansion
of an unnormalized integrand.

The contrast and common-link sectors of $B$ have eigenvalues $5$ and
$9$.  Inverting their two-layer matrices gives \eqref{f16v64:inverse-blocks}.
The five scaled deficits on the first layer converge in moments to
$\tfrac12|v_j^{\mathsf T}z|^2$, with analogous expressions for $u$.
For three independent color copies, fourth-moment pairing yields
\eqref{f16v64:ten-limit}.  In particular
$\mathbf1^{\mathsf T}S\mathbf1=9/20$ and
$\mathbf1^{\mathsf T}R\mathbf1=1/20$, which give
\eqref{f16v64:two-limit}.  A zero variance of a combination of the ten
limiting quadratic polynomials would make that polynomial constant on
$\mathbb R^{24}$.  On each layer, the off-diagonal link monomials first
force its $v_0$ coefficient to vanish; the four remaining diagonal
coefficients then vanish.  Thus the complete limiting Gram is positive
definite.

To justify the conditional assertion rather than infer it from joint
weak convergence, fix a compact set of scaled $u$ variables.  The inner
$X$ integral has field vectors $\gamma(4e_0+y_i)$, whose scalar parts
remain uniformly positive there.  The same localized dominated argument
is uniform on this compact set and gives the conditional Gaussian law
\[
 z\mid u\ \sim\ N(B^{-1}u,B^{-1})
\]
for each color.  The conditional mean of the limiting plaquette deficit
is therefore
\begin{equation}
 \mathbb E\left[\tfrac12|\textstyle\sum_i z_i|^2\mid u\right]
        =\frac23+\frac1{81}\,\frac12|\textstyle\sum_i u_i|^2.
 \label{f16v64:conditional-polynomial}
\end{equation}
Indeed $B^{-1}\mathbf1=\mathbf1/9$ and
$\mathbf1^{\mathsf T}B^{-1}\mathbf1=4/9$.

This compact-uniform conditional convergence upgrades to $L^2$ under
the actual $Y$ marginal.  To see the tail payment explicitly, for
$D_X=\gamma(1-A(X))$ Jensen gives
$\mathbb E|\mathbb E(D_X\mid Y)|^4\le\mathbb E|D_X|^4$, uniformly bounded
by the already proved moment estimates.  Those estimates also make
the probability of leaving a large scaled $Y$ ball, or its northern
chart, uniformly small.
Cauchy--Schwarz bounds the squared conditional-mean contribution on
that complement; the same argument applies to $D_Y$.  First restrict
to a fixed scaled ball, take $\gamma\to\infty$, and then let the ball
grow.  Centering \eqref{f16v64:conditional-polynomial} now proves
\eqref{f16v64:predictor-limit}.  Its variance is
$(243/800)/81^2=1/21600$.  Subtracting from the entrance limit $243/800$
by the exact conditional-variance identity gives $41/135$.
\end{proof}

The $Y$ precision $B-B^{-1}$ in this proof is the signed response after
integrating $X$.  Freezing the induced factor or keeping only its contact
term would instead use a different precision and would not give
\eqref{f16v64:two-limit}.  This is a calculation under one actual
normalized retained law.  Its flat exterior is still prescribed; no
probability of that exterior under the four-dimensional vacuum has been
assigned.  The delay here is one lattice step, not a fixed physical time.

\subsection{A literal spatial exterior with cancelling staples}

Return to the tube of four cyclic spatial links through time.  With its
neighboring free time layers not frozen, each tube link has three fixed
spatial staple paths and two temporal neighbors, in addition to the
central four-free-link plaquette.  Write the three spatial paths as
$p_{i1},p_{i2},p_{i3}$ and their sum as $s_i$.

\begin{proposition}[The cancelling spatial boundary is an actual Wilson exterior]
\label{f16v64:native-cancellation}
On a nondegenerate four-dimensional lattice patch there is a
time-independent retained-link assignment such that all temporal
transports of the tube are identity and
\begin{equation}
 (p_{i1},p_{i2},p_{i3})=(1,\omega,\omega^2),\qquad
 \omega=\cos(2\pi/3)+\mathbf i\sin(2\pi/3),\qquad s_i=0
 \label{f16v64:cancel-boundary}
\end{equation}
for each cyclic edge $i=1,\ldots,4$.  The path with value $1$ is an
actual complementary closing path for that edge.
\end{proposition}
\begin{proof}
Use a square in the $(\mu,\nu)$ plane and the remaining spatial direction
$\rho$.  The three complementary spatial plaquettes of each free edge
are its outward coplanar plaquette and the plaquettes displaced toward
$+\rho$ and $-\rho$.  Each complementary three-link path has an opposite
parallel retained link absent from all the other eleven selected paths.
The other two links of the path are connectors, which may be shared but
are never these private opposite links.  Set every connector to identity.
Set the cyclically oriented opposite link of each outward coplanar path
to $1$, those in the $+\rho$ copy of the square to $\omega$, and those in
the $-\rho$ copy to $\omega^2$.  Reverse the stored positive-link value
when the cyclic orientation requires it.  This produces exactly the
three indicated path values.  Set all retained temporal links to
identity and repeat the same spatial assignment at each time.  No two
instructions concern the same private link.  Remaining exterior links
can be identity.  The resulting exterior has nonzero Wilson curvature;
flatness or a small exterior action is not asserted.  The quaternion
identity $1+\omega+\omega^2=0$ finishes the construction.
\end{proof}

At this boundary the spatial one-link terms cancel pointwise, not merely
in expectation.  Define on $\mathcal X=\SU(2)^4$
\begin{equation}
 V(X)=1-w(Q(X)),\quad D(X,Y)=\sum_{i=1}^4(1-x_i\cdot y_i),\quad
 K_\gamma(X,Y)=e^{-\gamma[V(X)/2+D(X,Y)+V(Y)/2]}.
 \label{f16v64:cancel-kernel}
\end{equation}
This is the homogeneous conditional tube transfer obtained from the
literal Wilson action, up to a configuration-independent scalar.
Let $\rho_\gamma$ be its largest eigenvalue, let $\phi_\gamma>0$ be its
$L^2(dH^4)$-normalized eigenfunction, and put
$P_\gamma=K_\gamma/\rho_\gamma$.
The stationary slice and adjacent-pair laws are
\begin{equation}
 d\pi_\gamma=\phi_\gamma^2dH^4,\qquad
 d\Pi_\gamma(X,Y)=\rho_\gamma^{-1}
       \phi_\gamma(X)K_\gamma(X,Y)\phi_\gamma(Y)dH^4(X)dH^4(Y).
 \label{f16v64:cancel-vacuum}
\end{equation}
They are the fixed-regulator long-cylinder limits of this conditional
tube with nonzero nonnegative local endpoint data.  They are not the
full four-dimensional vacuum or a substituted physical transfer.

\begin{theorem}[Four noncollapsing native writers with slow conditional transfer]
\label{f16v64:slow-tube}
The kernel \eqref{f16v64:cancel-kernel} is positive, compact, injective,
and positivity improving.  Its Perron eigenvalue is simple and
$0\preceq P_\gamma\preceq I$.  For the four closed complementary writers
\[
 f_i(X)=w(P_{i,\mathrm{out}}^{-1}X_i)=w(X_i),\qquad 1\le i\le4,
\]
let $V_f c=\phi_\gamma\sum_i c_i f_i$.  Then
\begin{equation}
 \mathbb E_{\pi_\gamma}f_i=0,\qquad
 G_f=V_f^*V_f=\tfrac14I_4,\qquad
 V_f^*P_\gamma V_f=\alpha_\gamma G_f.
 \label{f16v64:exact-source-Gram}
\end{equation}
Set
\begin{equation}
 c_H=\frac8{3\pi^3},\qquad C_0=16-4\log c_H,\qquad
 \delta_\gamma=\frac{3\log\gamma+C_0/2}{\gamma}\quad(\gamma\ge1).
 \label{f16v64:slow-constants}
\end{equation}
Then
\begin{equation}
 0<\alpha_\gamma<1,\qquad
 \alpha_\gamma\ge1-\delta_\gamma,\qquad
 V_f^*P_\gamma^nV_f\succeq(1-\delta_\gamma)_+^nG_f
       \quad(n\ge1).
 \label{f16v64:slow-bound}
\end{equation}
There are at least four eigenvalues, counting multiplicity and excluding
the Perron eigenvalue, whose normalized values are at least
$(1-\delta_\gamma)_+$.  Thus no strictly subunit, coefficient-uniform
contraction factor holds on these actual conditional sources over all
admissible exteriors.
\end{theorem}
\begin{proof}
The temporal convolution in \eqref{f16v64:cancel-kernel} is a product of
four strictly positive-character convolutions.  Sandwiching it by the
strictly positive multiplier $e^{-\gamma V/2}$ proves positivity and
injectivity.  Its continuous kernel is strictly positive on a compact
space.  Compactness, existence of a strictly positive top eigenfunction,
and simplicity can also be seen directly from the variational principle:
replacing a sign-changing maximizing function by its absolute value
strictly increases the quadratic form; two independent top eigenvectors
would yield a sign-changing one orthogonal to a positive top vector.
The spectral theorem gives $0\preceq P_\gamma\preceq I$.

For $g=(g_1,\ldots,g_4)\in\SU(2)^4$, with $g_5=g_1$, act by
\begin{equation}
 X_i\longmapsto g_iX_i g_{i+1}^{-1}.
 \label{f16v64:cycle-action}
\end{equation}
The plaquette trace, simultaneous temporal bond, and Haar measure are
invariant.  Uniqueness makes $\phi_\gamma$ invariant.  Averaging one
endpoint group shows that each $X_i$ is Haar distributed under
$\pi_\gamma$, so its scalar part has mean zero and second moment $1/4$.
For different edges, an endpoint belonging to only one of them makes
the mixed moment zero.  The same argument applies to different-edge
moments in the adjacent-pair law.  On the diagonal the elementary Haar
identity gives
\[
 \mathbb E_{\Pi_\gamma}[f_i(X)f_i(Y)]
       =\tfrac14\mathbb E_{\Pi_\gamma}(x_i\cdot y_i).
\]
Cyclic permutation of the four links preserves the kernel (the real
trace of $Q$ is cyclic), so the last expectation is independent of $i$.
This proves \eqref{f16v64:exact-source-Gram}, with
\begin{equation}
 1-\alpha_\gamma=\tfrac14\mathbb E_{\Pi_\gamma}D(X,Y).
 \label{f16v64:alpha-defect}
\end{equation}
Strict positivity follows from injectivity; strict inequality from one
follows from the simple Perron eigenspace and centering.

Here is a nonperturbative bound on the right side.  A one-link Haar cap
of geodesic radius $r\le1$ has mass at least $c_Hr^3$, because
$\sin\theta\ge2\theta/\pi$ for $0\le\theta\le1$.  Take the product cap
$\mathcal B_\gamma$ of radius $r=\gamma^{-1/2}$.  It has mass at least
$c_H^4\gamma^{-6}$.  Quaternion norm multiplicativity and telescoping
give $V(X)\le8/\gamma$ there; for $X,Y$ in the cap,
$D(X,Y)\le8/\gamma$.  Thus the kernel is at least $e^{-16}$ on the
product cap.  Its normalized indicator is a Rayleigh test, giving
\begin{equation}
 e^{-16}c_H^4\gamma^{-6}\le\rho_\gamma\le1.
 \label{f16v64:rho-cap}
\end{equation}
The upper bound follows from $K_\gamma\le1$ and normalized Haar measure.
Put $\mathcal E(X,Y)=V(X)/2+D(X,Y)+V(Y)/2$.  Testing $K_{\gamma/2}$ on
$\phi_\gamma$ and applying Jensen under the exact pair law yields
\[
 1\ge\rho_{\gamma/2}
 \ge\rho_\gamma\mathbb E_{\Pi_\gamma}e^{\gamma\mathcal E/2}
 \ge\rho_\gamma e^{(\gamma/2)\mathbb E_{\Pi_\gamma}\mathcal E}.
\]
Therefore
\[
 \mathbb E_{\Pi_\gamma}D\le
 \frac2\gamma\log(\rho_\gamma^{-1})
 \le\frac{12\log\gamma+2C_0}{\gamma}.
\]
Combine with \eqref{f16v64:alpha-defect}.

For every vector in the four-dimensional source range, the normalized
spectral measure of $P_\gamma$ has first moment $\alpha_\gamma$.
Convexity of $t^n$ on $[0,1]$ gives the last inequality in
\eqref{f16v64:slow-bound}, including all mixed coefficients.  The source
range is orthogonal to $\phi_\gamma$.  The five-dimensional span of the
ground vector and that range has minimum Rayleigh quotient
$\alpha_\gamma$.  The min--max principle proves the stated multiplicity
claim without identifying an invariant four-dimensional subspace.
\end{proof}

In particular, when $n_\gamma\log\gamma/\gamma\to0$,
$V_f^*P_\gamma^{n_\gamma}V_f\to G_f$.  No upper bound on a physical energy
follows unless $n_\gamma a$ is also controlled in the required direction.
At a fixed physical delay $n\simeq\tau/a$, the error term $n\delta_\gamma$
need not be small.  A closing \emph{lattice} gap is not a disproof of a
positive continuum physical gap.  These slowly changing writers close
through a prescribed curved exterior; their norms under the native
unconditioned vacuum have not been replaced by $I_4/4$.

\begin{proposition}[The spectral obstruction survives small spatial-boundary perturbations]
\label{f16v64:robust-cancellation}
Keep the temporal transports identity and perturb the time-independent
spatial paths, with new aggregate vectors $s_i$.  Let $K_{\gamma,s}$
be the corresponding symmetric conditional transfer, using the same
configuration-independent scalar convention as $K_\gamma$.  Put
$\epsilon=\gamma\sum_i|s_i|$ and $d=e^\epsilon-1$.  Its fifth largest
normalized eigenvalue, counting the Perron eigenvalue as the first,
satisfies
\begin{equation}
 \frac{\lambda_4(K_{\gamma,s})}{\rho(K_{\gamma,s})}
 \ge e^{-\epsilon}(\alpha_\gamma-d)_+.
 \label{f16v64:robust-bound}
\end{equation}
In particular it tends to one if $\epsilon\to0$ and $\gamma\to\infty$.
\end{proposition}
\begin{proof}
The kernel ratio is
$\exp\{(\gamma/2)\sum_i s_i\cdot(x_i+y_i)\}$, between
$e^{-\epsilon}$ and $e^\epsilon$.  Pointwise domination gives
$|(K_{\gamma,s}-K_\gamma)f|\le dK_\gamma|f|$, hence
$\|K_{\gamma,s}-K_\gamma\|\le d\rho_\gamma$.
The min--max principle bounds its fifth eigenvalue below by
$\rho_\gamma(\alpha_\gamma-d)$ and its largest eigenvalue above by
$(1+d)\rho_\gamma=e^\epsilon\rho_\gamma$.  The perturbed operator remains
positive by the same multiplier--convolution factorization.  This proves
\eqref{f16v64:robust-bound}.  Continuity of the actual finite path maps
gives nonempty open neighborhoods of the constructed spatial-boundary
parameters satisfying a prescribed small $\epsilon$.  No probability of
an indefinitely time-independent exterior is inferred from that fact.
\end{proof}

This last statement is spectral.  It does not assert that the four old
writers keep the exact Gram $I_4/4$ at the perturbed boundary: the exact
cyclic symmetry used for that identity may be lost.  Nor does it silently
fix a time-varying boundary by the same time-homogeneous transfer.

\subsection{How often does the marked cancellation occur in the native law?}

The construction above cannot by itself decide the averaged problem.
It is possible to put a first, volume-uniform estimate on precisely its
bad exterior event without assuming a small-field theorem.  Work now
with the \emph{unconditioned} isotropic Wilson probability measure on a
periodic four-dimensional hypercubic lattice of common side at least four.
Write $N_L=4N_v$, $N_P=6N_v$, and
\begin{equation}
 d\mu_{\gamma,L}=Z_{\gamma,L}^{-1}e^{-\gamma E_L}dH^{N_L},\qquad
 E_L=\sum_p(1-w(U_p)).
 \label{f16v64:native-measure}
\end{equation}
This normalized law is different from both conditional laws used above.

\begin{theorem}[A native marked-staple probability bound uniform in volume]
\label{f16v64:native-probability}
For $\gamma\ge1$, put
\begin{equation}
 B_\gamma=8+\log\gamma-\tfrac23\log c_H.
 \label{f16v64:pressure-constant}
\end{equation}
Every plaquette satisfies
\begin{equation}
 \mathbb E_{\mu_{\gamma,L}}[1-w(U_p)]\le B_\gamma/\gamma.
 \label{f16v64:plaquette-pressure}
\end{equation}
For a marked tube edge let $s_e=p_1+p_2+p_3$ be its three complementary
\emph{spatial} staple paths, and let
$\mathcal B_e(r_*)=\{|s_e|\le r_*\}$ for $0\le r_*<3$.  Then
\begin{equation}
 \boxed{\quad
 \mu_{\gamma,L}(\mathcal B_e(r_*))
 \le \min\left\{1,\frac{3B_\gamma}{\gamma(3-r_*)}\right\}.
 \quad}
 \label{f16v64:bad-probability}
\end{equation}
For a union over $M$ marked edges the right side is multiplied by $M$,
then capped at one.  The estimate passes to any thermodynamic limit of
these periodic measures, for the same bounded local event.
\end{theorem}
\begin{proof}
Restrict every positively stored link to its Haar cap of radius
$\gamma^{-1/2}$.  For every plaquette, telescoping its four unit quaternion
factors gives $1-w(U_p)\le8/\gamma$.  Therefore
\[
 Z_{\gamma,L}\ge e^{-8N_P}c_H^{N_L}\gamma^{-3N_L/2},\qquad
 -\log Z_{\gamma,L}/N_P\le B_\gamma.
\]
The relative entropy with respect to normalized product Haar measure is
\[
 0\le D(\mu_{\gamma,L}\Vert H^{N_L})
      =-\gamma\mathbb E E_L-\log Z_{\gamma,L}.
\]
Hence $\mathbb E E_L\le-\log Z_{\gamma,L}/\gamma$.
Translation and hypercubic symmetry make the plaquette expectations
identical, proving \eqref{f16v64:plaquette-pressure}.  This is a normalized
native estimate; the extensive cap cost was divided by the actual number
of plaquettes, not declared volume independent before division.

For any value $x_e$ of the shared free link, its three adjacent plaquette
defects have the exact sum
\begin{equation}
 \sum_{\alpha=1}^3(1-p_\alpha\cdot x_e)
       =3-s_e\cdot x_e\ge3-|s_e|.
 \label{f16v64:bad-deterministic}
\end{equation}
Consequently $(3-r_*)\mathbf1_{\mathcal B_e(r_*)}$ is at most this sum.
Take expectations using \eqref{f16v64:plaquette-pressure}.  This step uses
no product distribution for the three paths and no assumption on the
conditional free link.  Union bounds prove the marked-set assertion.
For a limit measure, one can first use the continuous bounded function
$\min\{1,[\sum_\alpha(1-p_\alpha\cdot x_e)]/(3-r_*)\}$, which majorizes
the event indicator and has the same expectation bound.  Weak convergence
of local continuous observables then proves the event bound without
assuming that its threshold has zero mass.
\end{proof}

There is also a purely retained curvature witness:
\begin{equation}
 \sum_{\alpha<\beta}(1-w(P_\alpha P_\beta^{-1}))
           =\tfrac12(9-|s_e|^2).
 \label{f16v64:retained-witness}
\end{equation}
These are actual closed paths of length at most six.  Exact cancellation
requires their summed defect to be $9/2$; the constructed exterior is not
a hidden coherent configuration.  Nevertheless the theorem gives only
an $O(\log\gamma/\gamma)$ probability bound, not exponential suppression,
not an estimate for every curved exterior, and not a bound uniform under
arbitrary boundary conditions.  In a physical-time window the number
$M$ of marked edges may grow like $a^{-1}$; that factor cannot be dropped.

\subsection{The relative source-weighted cost is a separate, now explicit quantity}

Let $F$ be any finite vector of native square-integrable writers, centered
under \eqref{f16v64:native-measure}, with Gram $G=\mathbb E FF^*$ on its
exact quotient.  For any marked cancellation event $\mathcal B$ put
$G_{\mathcal B}=\mathbb E[FF^*\mathbf1_{\mathcal B}]$.
The following elementary consequence specifies what extra source input
would make the preceding probability bound sufficient.

\begin{proposition}[Paid source weighting, without a hidden noncollapse assumption]
\label{f16v64:source-weight}
Suppose the finite source family has relative fourth-moment constant
\begin{equation}
 \mathbb E|c^*F|^4\le\mathcal K\,(c^*Gc)^2
       \quad\hbox{for every coefficient vector }c.
 \label{f16v64:kurtosis}
\end{equation}
If $\mu(\mathcal B)\le p$, then
\begin{equation}
 0\preceq G_{\mathcal B}\preceq\sqrt{\mathcal Kp}\,G.
 \label{f16v64:weighted-bad}
\end{equation}
This statement is invariant under any declared source coefficient map,
including regulator-dependent ones, and retains the exact null space.
For merely bounded $F$ with $\|F\|\le M_F$, the unconditional conclusion
is only $G_{\mathcal B}\preceq M_F^2p\,I$.
\end{proposition}
\begin{proof}
For each $c$, Cauchy--Schwarz gives
$\mathbb E[|c^*F|^2\mathbf1_{\mathcal B}]
\le(\mathbb E|c^*F|^4)^{1/2}\mu(\mathcal B)^{1/2}$.
Apply \eqref{f16v64:kurtosis}; pullback by congruence proves the map
assertion.  A zero Gram direction is zero almost surely and contributes
no fourth moment or bad-event mass.  The supremum bound proves the last
assertion independently of any entrance lower bound.
\end{proof}

For example, a coefficient-uniform $\mathcal K$ and a fixed marked set
would turn \eqref{f16v64:bad-probability} into a vanishing relative cost
$O(\sqrt{\log\gamma/\gamma})$.  Such a native dense-core fourth-moment
bound is \emph{not} proved here.  It cannot be supplied by the coherent
two-layer conditional computation.  A bounded rare-event source already
shows the distinction sharply: if $p=\mu(\mathcal B)\in(0,1)$ and
$f=\mathbf1_{\mathcal B}-p$, then
\begin{equation}
 \frac{\mathbb E[f^2\mathbf1_{\mathcal B}]}{\mathbb E f^2}=1-p.
 \label{f16v64:rare-source}
\end{equation}
The event can be rare and still contain almost all of the normalized
source variance.  The cancellation event is gauge invariant, so this is
also a gauge-invariant local bounded test when the marked set is local.
It is not a claim that this test lies in a chosen physical continuum core.

\subsection{The next noncircular task and the scope of this iteration}

The actual two-layer retained covariance and its predictor have now been
computed.  A second isolated flat-cell Hessian would add little.  The
cancelling exterior proves that a uniform contraction over all spatial
boundary data is an unsuitable replacement for the averaged problem.
The native marked-event bound establishes that this explicit obstruction
is uncommon in translation-invariant Wilson measures, but its frequency
alone does not pay its source weight or a growing physical-time window.

The next calculation should use a declared physically regularized source
bank under the actual retained marginal and estimate the joint quantity
$G_{\mathcal B}$, or its conditional-predictor analogue, with normalization
and the number of marked time layers retained.  On the complementary
boundary set one must still control the induced mixed-score covariance;
noncancellation of three staples is not by itself a temporal contraction
theorem.  Establishing \eqref{f16v64:kurtosis} for that bank would close
one explicit part of the present boundary cost, but is not asserted as
a stand-in for the remaining dynamics.

Neither $\rho_\gamma$ in the cancelling conditional tube nor the factor
$1/81$ in the coherent two-layer law is the vacuum normalization of the
full four-dimensional transfer.  The original physical lattice spacing
$a$ is unchanged.  A nontrivial physical source limit, source-weighted
control of the retained collective law, and common dense-core
physical-time contraction remain separate obligations.

\subsection{Verification, preservation, and result ledger}

The companion diagnostic is
\begin{center}\small
\path{verify_ym_f16_v64_retained_sources.py}.
\end{center}
It separates exact rational block inversions and moment identities,
literal lattice-path and quaternion checks, finite positive-kernel
variational diagnostics, and importance-weighted compact two-layer
integration.  The latter retains the compact integrand and also uses
a coupled Gaussian control variate with an exactly evaluated expectation.
Both raw and variance-reduced results are floating-point diagnostics with
stated resolution and sampling sensitivity, not rigorous quadrature enclosures.
The analytic proofs above do not depend on its numerical values.  No
four-dimensional vacuum simulation, interval certificate, or
proof-assistant verification is claimed.

Removing the delimited V64 appendix and reverting the one displayed
version change recovers V63 byte for byte.  All earlier mathematical
sections and comparison archives remain unchanged.

\clearpage
\begin{record}[V64 result ledger]
\label{f16v64:ledger}
\emph{Proved here:} strict positive normalized two-layer cross Grams;
the explicit two-plaquette limit and complete ten-writer covariance;
the retained predictor limit, its $1/6561$ variance fraction, and the
paid innovation remainder; a literal cancelling spatial Wilson exterior;
four exact noncollapsing conditional source directions with slow transfer
and a quantified spectral-neighborhood extension; a volume-uniform native
marked-cancellation probability bound; and an explicit sufficient
relative-fourth-moment payment for its source weight.

\emph{Inherited:} the compact Wilson action and cyclic path convention,
positive one-link character multipliers, the compact normalized moment
method, actual retained-marginal interpretation, exact source quotients,
and the distinction between lattice and physical time.

\emph{Not asserted:} a coherent-boundary estimate uniform in the exterior;
a probability of a whole static exterior from a one-time event bound;
a conditional Gram identified with the native vacuum Gram; a uniform
relative fourth-moment bound for an unspecified physical source core;
an invariant four-dimensional slow subspace; exponential suppression of
the cancellation event; or a physical gap conclusion from either
conditional time scale.

\emph{Still unproved:} normalized source-weighted control of the full
retained law across physical time, a nontrivial interacting continuum
source limit, common dense-core contraction constants, and the continuum
Yang--Mills mass gap.
\end{record}

\begin{firewall}
The retained two-layer marginal is integrated rather than frozen.
The cancelling exterior is built from actual links rather than
independent fictitious staple vectors.  Its spectral estimate and its
native probability estimate concern different measures, explicitly
identified above.  No relative source estimate is obtained by deleting
its normalization.  V1--V63 are preserved unchanged.
\end{firewall}
% END F16 V64 APPEND

% BEGIN F16 V65 APPEND -- 2026-09-18
% Append-only continuation of V64.  The preceding mathematical body is unchanged.
\clearpage
\section[V65: cancelling-boundary exclusion with relative source weights]{V65: cancelling-boundary exclusion\\with relative source weights}
\label{f16v65:context}

V64 distinguishes a small probability from a small fraction of a source
Gram.  Its proposed fourth-moment payment was sufficient but unpopulated.
This continuation supplies a different payment for a specified class:
finite Wilson-polynomial families under the actual native measure, followed
by common Markov smoothing or conditional prediction.  A uniform
fourth-moment constant is not assumed.  Instead, a polynomial
supremum-to-native-norm bound is paired with exponential suppression of
the particular cancelling-staple geometry.  The resulting estimates hold
in every source direction, not just after taking a trace.

The joint four-edge event is important.  Its twelve distinct complementary
plaquettes impose the full action cost $12-\sum_i|s_i|$.  Disseminating
that joint event keeps this cost intact; bounding its constituent
plaquettes separately would lose the exponent needed for a growing
physical region.  We prove an explicit exclusion for the near-cancelling
neighborhood of V64's example, including a region with $O(a^{-4})$ marks
under a stated trajectory.  This is not a contraction theorem for the
complementary boundary configurations.

\subsection{Archive audit, measures, and the inherited reflection input}

The f14 archive already proves bare-Wilson site reflection positivity and
chessboard estimates, thermal exponential moments in V98, and scalar
cell-source bounds in V99--V100.  In particular, its labels
\texttt{eq:v98-site-RP}, \texttt{thm:v98-exponential}, and
\texttt{prop:v99-chessboard-source} are prior results, not discoveries
of this continuation.  Its selected susceptibility and Gaussian-limit
couplings are not identified with a physical trajectory.  We use the
same finite reflection argument, with a new joint cancellation witness,
and V64's elementary cap normalization.  No susceptibility selection or
limiting Gaussian field is needed.

For reference, the classical closed-block chessboard inequality is
Theorem 5.8 in M.~Biskup, \emph{Reflection Positivity and Phase Transitions
in Lattice Spin Models},
\href{https://arxiv.org/abs/math-ph/0610025}{arXiv:math-ph/0610025v2},
Section 5.3.  Its blocks share their boundary variables.  We use this
algebraic reflection theorem, not a spin-model decay conclusion.

Let $\mathbb T=\prod_{\alpha=1}^4\mathbb Z/L_\alpha\mathbb Z$ be a
rectangular dyadic torus, with the side lengths large enough for the blocks
specified below.  The normalized native law is
\begin{equation}
 d\mu_{\gamma,\mathbb T}=Z_{\gamma,\mathbb T}^{-1}
       e^{-\gamma E(U)}\prod_e dH(U_e),\qquad
 E(U)=\sum_p d_p(U),\qquad d_p=1-w(U_p),
 \label{f16v65:measure}
\end{equation}
with $\gamma\ge1$.  There are $N_v=\prod_\alpha L_\alpha$ sites,
$4N_v$ stored links and $6N_v$ plaquettes.  Put
\begin{equation}
 c_H=\frac8{3\pi^3},\qquad
 B_\gamma=8+\log\gamma-\frac23\log c_H.
 \label{f16v65:B}
\end{equation}
The cap argument of Theorem~\ref{f16v64:native-probability} uses only
these incidence counts, not equal side lengths, and gives
\begin{equation}
 -\log Z_{\gamma,\mathbb T}\le6N_vB_\gamma.
 \label{f16v65:cap-input}
\end{equation}
The law is not a conditional tube law or a flat-boundary surrogate.

We recall precisely the reflection input.  For an integer coordinate
plane and its opposite plane, every unit plaquette lies in one closed
half or in the boundary.  With $U_0$ the tangent links on the two planes,
$E=E_++\theta E_++E_0$.  The inherited identity is
\begin{equation}
 \mathbb E[F\,\theta\overline F]
 =Z^{-1}\int e^{-\gamma E_0}
       \left|\int e^{-\gamma E_+}F\,dH(U_+)\right|^2dH(U_0)\ge0.
 \label{f16v65:RP-input}
\end{equation}
In particular no bond-plane factorization without a character expansion
is asserted.  Reflection inverts the stored link when its orientation
is reversed; the plaquette trace is unchanged.

For a closed rectangular block with side lengths $b_\alpha$ dividing
$L_\alpha$, let $v_b=\prod b_\alpha$ and $N_b=N_v/v_b$.  Assume each
$L_\alpha/b_\alpha$ is a power of two of at least two.  If $\vartheta_z A$ denotes
the alternating reflected copy of a block event, the inherited
chessboard inequality reads
\begin{equation}
 \mu\!\left(\bigcap_{j=1}^m\vartheta_{z_j}A_j\right)
 \le\prod_{j=1}^m
  \mu\!\left(\bigcap_z\vartheta_z A_j\right)^{1/N_b},
 \qquad z_1,\ldots,z_m\ \hbox{distinct}.
 \label{f16v65:chessboard-input}
\end{equation}
Its hypotheses are exactly the block algebras, coordinate symmetries
and positive forms in \eqref{f16v65:RP-input}; a product distribution
of blocks is not a hypothesis.  Coordinate-by-coordinate application
of the closed-block theorem gives the rectangular version.  Equivalently,
one normalizes each block label by its fully disseminated expectation;
reflection Cauchy--Schwarz duplicates either half.  A maximizing label
array with the fewest disagreements must be constant, since reflecting
across a disagreement preserves maximality and reduces that count.
This is the same scalar-label argument as f14 V99, now with block
indicators.  Replacing zero indicators by a positive regularization and
taking its limit covers zero disseminated expectations.

\subsection{An unsplit action cost for one edge and for the complete cell}

Use coordinates $(\mu,\nu,\rho,0)$, where $0$ denotes time.  For a tube
edge let $p_{e1},p_{e2},p_{e3}$ be its three complementary spatial paths
in V64, all oriented like the edge, and write
\begin{equation}
 s_e=p_{e1}+p_{e2}+p_{e3},\quad
 \mathcal B_e(r)=\{|s_e|\le r\},\quad 0\le r<3.
 \label{f16v65:edge-event}
\end{equation}
For one complete spatial square $q$, let $s_1,\ldots,s_4$ be the
corresponding sums on its cyclic edges and set
\begin{equation}
 \mathcal C_q(R)=\left\{\sum_{i=1}^4|s_i|\le R\right\},
 \qquad 0\le R<12.
 \label{f16v65:cell-event}
\end{equation}
These events depend only on retained complementary paths.  Their norms
are gauge invariant.  The event $\mathcal C_q(R)$ is a joint condition;
it is not the union of four single-edge conditions.

\begin{proposition}[Literal block placement and disjoint witnesses]
\label{f16v65:geometry}
The event $\mathcal B_e(r)$ can be placed in a block of sides $(2,2,2,2)$.
Its three witnessing plaquettes are distinct and no witness lies entirely
in a face of that block.  The event $\mathcal C_q(R)$ can be placed in a
block of sides $(4,4,2,2)$.  Its twelve witnessing plaquettes have the same
properties.  In either case different reflected blocks therefore have
disjoint sets of witnessing plaquettes.  Pointwise,
\begin{align}
 \mathcal B_e(r)&\ \Longrightarrow\ \sum_{p\in W_e}d_p\ge3-r,
          \label{f16v65:edge-cost}\\
 \mathcal C_q(R)&\ \Longrightarrow\ \sum_{p\in W_q}d_p\ge12-R.
          \label{f16v65:cell-cost}
\end{align}
\end{proposition}
\begin{proof}
For the edge block take the edge from $(0,1,1,1)$ to $(1,1,1,1)$.
Use its outward $-\nu$ plaquette and the two $\pm\rho$ plaquettes.
All their vertices lie in $[0,2]^4$.  Each plaquette either varies in a
coordinate or has an interior value $1$ there, so none is contained in a
block boundary hyperplane.  Their complementary paths are exactly the
three paths defining the event.

For the square block use the cyclic square with $(\mu,\nu)$ corners
$(1,1),(2,1),(2,2),(1,2)$ and $(\rho,0)=(1,1)$.  For each edge take its
outward coplanar plaquette and its two $\pm\rho$ plaquettes.  Their vertices
lie in $[0,4]^2\times[0,2]^2$, and none of these plaquettes lies entirely
in a boundary face.  Each contains exactly one of the four central edges.
Consequently the twelve plaquettes are distinct.  Two closed blocks can
meet only in their boundary faces; the assertion about reflected copies
follows, including at the periodic identification.

For every actual value $x_e$ of a shared edge,
$\sum_{p\in W_e}d_p=3-s_e\cdot x_e\ge3-|s_e|$.
For the square, add these four identities.  The twelve plaquettes are
counted once, not four times.  This proves both costs without integrating
or assigning independent values to the paths.
\end{proof}

\begin{theorem}[Native exponential exclusion of the cancelling geometry]
\label{f16v65:exponential}
For the law \eqref{f16v65:measure} on compatible tori,
\begin{align}
 \mu(\mathcal B_e(r))&\le
  \zeta_e(\gamma,r):=\min\{1,\exp[-(3-r)\gamma+96B_\gamma]\},
       \label{f16v65:edge-prob}\\
 \mu(\mathcal C_q(R))&\le
  \zeta_q(\gamma,R):=\min\{1,\exp[-(12-R)\gamma+384B_\gamma]\}.
       \label{f16v65:cell-prob}
\end{align}
For $m$ distinct compatible reflected block copies of the same event,
the probability of their intersection is at most $\zeta^m$.
For an arbitrary union of $M$ translated or rotated marks of either fixed
type, the bound is $\min\{1,M\zeta\}$.  No independence is asserted.
The one-mark and finite-union bounds pass to local limits of these
periodic laws, including limits taking the time side to infinity first.
\end{theorem}
\begin{proof}
Let $\Delta=3-r$ or $12-R$.  On the fully disseminated event,
Proposition~\ref{f16v65:geometry} gives $E\ge\Delta N_b$, because the
witness sets are disjoint.  Therefore
\[
 \mu\!\left(\bigcap_z\vartheta_z A\right)
 \le Z^{-1}e^{-\gamma\Delta N_b}.
\]
Take the $N_b$th root and use \eqref{f16v65:cap-input}:
$Z^{-1/N_b}\le\exp(6v_bB_\gamma)$.  Here $v_b=16$ or $64$, giving
$96$ or $384$.  The one-event case and then the multi-event case of
\eqref{f16v65:chessboard-input} prove the assertions.  Translation and the same block argument with
permuted coordinate roles handle arbitrary individual marks; the union
bound retains their factor $M$.

For limits, approximate the indicator at threshold $r$ from above by a
continuous function supported at threshold $r+\epsilon$, and likewise
for $R$.  Local expectations converge, and letting $\epsilon\downarrow0$
recovers the displayed bounds.  This does not assume that the exact
threshold has zero probability.  Rectangular reflection and cap constants
are independent of aspect ratio, so temporal limits are included.
\end{proof}

The exponential rate in \eqref{f16v65:cell-prob} is the full twelve-plaquette
cost.  It is not obtained by applying a one-plaquette estimate and then
multiplying dependent probabilities.  The large prefactors are explicit
and not optimized: $e^{96B_\gamma}$ and $e^{384B_\gamma}$ grow like
$\gamma^{96}$ and $\gamma^{384}$, respectively, times fixed constants.
These estimates are asymptotic statements, not claims of useful numerical
margins at moderate Wilson coefficients.  General exponential defect
suppression is inherited machinery; its unsplit application to V64's
particular exterior is the point here.

For completeness, on $M$ compatible block positions let $N_A$ count bad
blocks.  Expanding $\prod_z(1+t\mathbf1_{A_z})$ gives the useful consequence
\begin{equation}
 \mathbb E(1+t)^{N_A}\le(1+t\zeta)^M,\qquad t\ge0.
 \label{f16v65:count-mgf}
\end{equation}
This controls joint occurrences without postulating a Bernoulli law.
It is not a statement that differently placed blocks are stochastically
independent or that arbitrary exterior conditionings preserve the bound.

\subsection{A degree-controlled native norm bound, with no fourth-moment input}

A Wilson polynomial means a finite linear combination of products of
matrix elements of closed link words, restricted to $\SU(2)$ links.
In quaternion coordinates it is a polynomial separately in each stored
link.  Inversion is linear in these coordinates.  A bound $D_e$ on the
number of occurrences of link $e$ gives a valid separate degree bound.
Constants, centering, and arbitrary coefficient combinations do not
increase these degrees.

We first record a compact one-link estimate.  Its proof is included to
make the dependence on degree and concentration explicit, rather than
invoke an unquantified finite-dimensional norm equivalence.  For $D\ge1$
let
\begin{equation}
 \mathcal I_D=\{\alpha\in\mathbb N_0^3:|\alpha|\le2D\},\quad
 n_D=\binom{2D+3}{3},\quad
 (M_D)_{\alpha\beta}=\prod_{j=1}^3
       \int_{-1}^1 z^{\alpha_j+\beta_j}\,dz.
 \label{f16v65:cube-Gram}
\end{equation}
This is a positive definite rational matrix.  Write $m_D=\lambda_{\min}(M_D)$
and define a usable positive constant
\begin{equation}
 \eta_D=\frac{m_D}
  {\pi^2e^2\,2^{2D+3}3^{2D+3/2}n_D},\qquad \eta_0=1.
 \label{f16v65:eta}
\end{equation}
One can replace $m_D$ by the explicit rational lower bound
$\det(M_D)/(\operatorname{tr}M_D)^{n_D-1}$.  Thus the constant does not
require a numerically thresholded eigenvalue.

\begin{proposition}[Polynomial norms in every one-link Wilson conditional law]
\label{f16v65:one-link}
Let $d\nu_h(x)=e^{h\cdot x}d\sigma(x)/\mathfrak z(|h|)$ on $S^3$.
For a complex spherical polynomial $f$ of degree at most $D$,
\begin{equation}
 \int|f|^2d\nu_h\ge
        \eta_D(1+|h|)^{-2D}\|f\|_\infty^2.
 \label{f16v65:one-link-bound}
\end{equation}
For $D=1$ the stronger usable choice $\eta_1=1/256002$ follows from
the one-link covariance estimate of V58.  Either choice is valid.
\end{proposition}
\begin{proof}
Rotate the field to $he_0$, with $h\ge0$.  The scalar Haar integral gives
\[
 e^{-h}\mathfrak z(h)=\frac2\pi\int_0^2e^{-hv}\sqrt{v(2-v)}\,dv
       \le4(1+h)^{-3/2}.
\]
For $h\le1$ this follows from $e^{-h}\mathfrak z(h)\le1$.
For $h\ge1$, replace $\sqrt{v(2-v)}$ by $\sqrt{2v}$ and extend the
integral to infinity, obtaining $\sqrt{2/\pi}\,h^{-3/2}$, which is
smaller than the displayed bound.

In stereographic coordinates
\[
 x(y)=\frac{(1-|y|^2,2y)}{1+|y|^2},\qquad
 d\sigma=\frac4{\pi^2}(1+|y|^2)^{-3}dy,
\]
write $f(x(y))=P(y)/(1+|y|^2)^D$, with $\deg P\le2D$.
This remains true after the initial rotation.  Put $r=(1+h)^{-1/2}$
and integrate only over $[-r/\sqrt3,r/\sqrt3]^3$.  On this cube,
$h(1-x^0)\le2$ and $1+|y|^2\le2$.  If $a$ is the coefficient vector
of $P$, the preceding normalization bound yields
\[
 \int|f|^2d\nu_h\ge
 \frac{r^{-3}}{\pi^2e^2\,2^{2D+3}}
                  \int_{[-r/\sqrt3,r/\sqrt3]^3}|P(y)|^2dy.
\]
Scaling the cube and using $M_D\succeq m_DI$ bounds its last integral
below by
\[
 3^{-3/2}r^3m_D(r/\sqrt3)^{4D}|a|^2.
\]
Indeed the scaled coefficient of degree $k\le2D$ is multiplied by
$(r/\sqrt3)^k$, whose square is at least $(r/\sqrt3)^{4D}$.
For every $y$, each monomial satisfies
$|y^\alpha|/(1+|y|^2)^D\le1$, so
$\|f\|_\infty^2\le n_D|a|^2$, including the omitted south pole by
continuity.  Combining these inequalities proves the stated constant.
The moment matrix is positive definite because a nonzero polynomial
cannot vanish on the open cube.  Its eigenvalues also give the stated
determinant/trace lower bound.

For the sharper linear choice write $f=a+b\cdot x$, let
$\bar f=\mathbb E f$, and use the inherited bound
$\operatorname{Var}(b\cdot x)\ge[32000(1+h)^2]^{-1}|b|^2$.
Since $|x-\mathbb Ex|\le2$,
$\|f\|_\infty^2\le2|\bar f|^2+8|b|^2
\le256002(1+h)^2\mathbb E|f|^2$.
The same calculation applies to complex $b$ using its Hermitian variance.
\end{proof}

The power $2D$ cannot be replaced by zero in a bound of this kind.
Already $f(x)=(1-x^0)^D$ has supremum $2^D$ and
$\mathbb E_{he_0}|f|^2$ of order $h^{-2D}$ as $h\to\infty$.
This follows by setting $v=1-x^0$ in the displayed scalar integral and
scaling $v$ by $h$.  The degree cost is real, even for a single link.

\begin{theorem}[Native leverage bound on every fixed polynomial source space]
\label{f16v65:native-leverage}
Let $\mathcal S$ be a finite set of integrated stored links and let $f$ have separate
degrees at most $D_e$, $e\in\mathcal S$.  Under the native measure,
\begin{equation}
 \boxed{\quad
 \|f\|_\infty^2\le\mathcal L_{\gamma,\mathbf D}\mathbb E|f|^2,
 \qquad
 \mathcal L_{\gamma,\mathbf D}
 =\prod_{e\in\mathcal S}\eta_{D_e}^{-1}
               (1+6\gamma)^{2\sum_{e\in\mathcal S}D_e}.
 \quad}
 \label{f16v65:leverage}
\end{equation}
The bound is independent of the surrounding volume, positions of the
links, and polynomial coefficients.  It applies in every finite-volume
Wilson boundary law with this local incidence, and in local limits of
the native periodic measures.  The event bounds above, in contrast,
are not claimed uniformly under arbitrary boundary laws.
\end{theorem}
\begin{proof}
Given all other links, a stored link has the exact law $\nu_{h_e}$ with
$|h_e|\le6\gamma$.  Each incident plaquette is linear in that link;
no product approximation or condition on the staple direction is used.
Order $\mathcal S=\{e_1,\ldots,e_m\}$, and put
$g_j=\sup_{x_{e_1},\ldots,x_{e_j}}|f|^2$, with $g_0=|f|^2$.
Condition on all links except $e_j$.  For each fixed choice of the first
$j-1$ dummy link values, the one-link bound applies to the polynomial
in $x_{e_j}$.  The bound is uniform in the actual conditional field,
which still depends on the actual, not dummy, neighboring links.
Taking the supremum after integrating gives
\[
 \mathbb E[g_{j-1}\mid U_{e_j^c}]
 \ge\sup_{u_{e_1},\ldots,u_{e_{j-1}}}
       \mathbb E[|f(u_{e_1},\ldots,u_{e_{j-1}},x_{e_j},\ldots)|^2
                            \mid U_{e_j^c}]
 \ge\eta_{D_{e_j}}(1+6\gamma)^{-2D_{e_j}}g_j.
\]
Average this inequality under the original law and iterate over $j$.
The final $g_m=\|f\|_\infty^2$ is constant.  Importantly, all expectations
remain under the original measure; there is no sequential claim that an
already marginalized link still has a Wilson conditional law.  This
proves \eqref{f16v65:leverage}.  Compact local moment convergence passes
the inequality to limit laws.
\end{proof}

For a centered finite column of such polynomials $F$, write
$G=\mathbb E FF^*$.  Applying \eqref{f16v65:leverage} to $c^*F$ gives
\begin{equation}
 \|c^*F\|_\infty^2\le\mathcal L_{\gamma,\mathbf D}\,c^*Gc
             \qquad\hbox{for every }c.
 \label{f16v65:all-directions}
\end{equation}
Thus a zero Gram direction is an exact polynomial identity, not a small
sample eigenvalue.  Arbitrary regulator-dependent coefficient maps pull
this inequality back by congruence without altering its constant.
For a fixed support and degree profile,
$\log\mathcal L_{\gamma,\mathbf D}=O(\log\gamma)$.
For a growing profile its full product of constants must be retained.

\subsection{Relative bad-source forms and exact retained predictors}

\begin{theorem}[A populated relative source-weight bound]
\label{f16v65:source-weight}
Let $F$ and $G$ be as above, and let $\mathcal A$ be a union of $M$ marks
of one of the two cancelling events.  Write $(\Delta,v_b)=(3-r,16)$ or
$(12-R,64)$, and define
\begin{equation}
 \epsilon_\gamma=
  \min\{1,\mathcal L_{\gamma,\mathbf D}
                  M e^{-\Delta\gamma+6v_bB_\gamma}\}.
 \label{f16v65:epsilon}
\end{equation}
Then under the actual native measure,
\begin{equation}
 \boxed{\qquad
 0\preceq\mathbb E[FF^*\mathbf1_{\mathcal A}]
                         \preceq\epsilon_\gamma G.
 \qquad}
 \label{f16v65:weighted-form}
\end{equation}
If $\Xi$ is any retained sigma-field containing $\mathcal A$ and
$\widehat F=\mathbb E(F\mid\Xi)$, the exact decomposition on that event is
\begin{equation}
 \mathbb E[FF^*\mathbf1_{\mathcal A}]
 =\mathbb E[\operatorname{Cov}(F\mid\Xi)\mathbf1_{\mathcal A}]
      +\mathbb E[\widehat F\widehat F^*\mathbf1_{\mathcal A}].
 \label{f16v65:weighted-disintegration}
\end{equation}
Both terms are positive and their \emph{sum} obeys
\eqref{f16v65:weighted-form}.  No relative fourth-moment bound is assumed.
\end{theorem}
\begin{proof}
For every $c$, use
$\mathbb E[|c^*F|^2\mathbf1_{\mathcal A}]
\le\|c^*F\|_\infty^2\mu(\mathcal A)$, then
\eqref{f16v65:all-directions} and Theorem~\ref{f16v65:exponential}.
The bound by $G$ is automatic and supplies the minimum in
\eqref{f16v65:epsilon}.  Conditional covariance plus the squared conditional
mean equals the conditional second moment; multiply by the
$\Xi$-measurable indicator and average.  This proves the identity and
pays both terms once.  Congruence preserves all exact source relations.
\end{proof}

There is a much smaller explicit constant for the three-orientation
private-plaquette packet of Theorem~\ref{f16v58:private-packet}.  Its
variance argument uses only the exact conditional links and therefore
also holds on the present periodic measures.  With $F_i=W_i-\mathbb EW_i$,
\begin{equation}
 G\succeq d_\gamma I_3,\qquad
 |c^*F|\le2\sqrt3|c|,\qquad
 \mathcal L_{\rm priv}(\gamma)=\frac{12}{d_\gamma}
                        =384000(1+6\gamma)^2.
 \label{f16v65:private-leverage}
\end{equation}
Hence this explicit packet satisfies \eqref{f16v65:weighted-form} with
$\mathcal L_{\rm priv}$ in place of the general product.  Adding
$W_4=W_1+W_2$ retains exactly the kernel
$\operatorname{span}\{(-1,-1,0,1)^{\mathsf T}\}$.
This use of the native entrance Gram does not reuse V58's microscopic
transfer decrement as a physical-time contraction.

\begin{proposition}[Physical-time smoothing pays no new support product]
\label{f16v65:smoothing}
Let $\mathcal K$ be a common linear, constant-preserving contraction from
$L^\infty(\mu)$ to $L^\infty(\nu)$, also preserving expectations and
contracting $L^2$.  Put $F^{\mathcal K}=\mathcal KF$,
$G_{\mathcal K}=\mathbb E_\nu F^{\mathcal K}(F^{\mathcal K})^*$.
If a marked event under $\nu$ has probability at most $p$, then
\begin{equation}
 \mathbb E_\nu[F^{\mathcal K}(F^{\mathcal K})^*\mathbf1_{\mathcal A}]
 \preceq\min\{1,\mathcal L_{\gamma,\mathbf D}p\}\,G.
 \label{f16v65:smooth-original-Gram}
\end{equation}
This includes conditional prediction, common probability-weighted time
averaging in a stationary law, and any number $k$ of steps of the native
vacuum Markov operator.  The product in \eqref{f16v65:leverage} is charged
to the original polynomial family, not to the possibly nonlocal smoothed
function.  If in addition, on the declared original quotient,
\begin{equation}
                    G_{\mathcal K}\succeq r_\gamma G,
                    \qquad r_\gamma>0,
 \label{f16v65:retention}
\end{equation}
then the left side of \eqref{f16v65:smooth-original-Gram} is at most
$\min\{1,\mathcal L_{\gamma,\mathbf D}p/r_\gamma\}\,G_{\mathcal K}$.
\end{proposition}
\begin{proof}
The common linear map commutes with coefficient combinations, so
$\|c^*F^{\mathcal K}\|_\infty\le\|c^*F\|_\infty$.
Multiply its square by $p$ and use \eqref{f16v65:all-directions}.
The $L^2$ contraction gives $G_{\mathcal K}\preceq G$, which permits
capping the original-Gram bound by one.  The retention inequality gives
the relative output bound; the output's own event form is also bounded
by $G_{\mathcal K}$.  Conditional expectation and stationary Markov
kernels have the required contractions by Jensen.  A probability-weighted
average of time shifts has the same properties.
\end{proof}

For the actual vacuum Markov operator $P_a$, the ground-state unitary
map identifies $P_a$ with the original vacuum-normalized transfer
$\mathcal T_a$.  Thus $\mathcal K=P_a^{\lceil\sigma/a\rceil}$ uses the
original physical delay $\sigma$; it is not a heat-bath update or a
replacement Hamiltonian.  The required finite-mark probabilities are
inherited from temporal limits of the periodic law, not from an arbitrary
Perron endpoint density inserted in a chessboard block.

Neither Markov smoothing nor the present norm bound proves
\eqref{f16v65:retention}.  If the original packet is Gram-normalized and
the smoothed packet has a positive definite limiting Gram, a fixed
$r_\gamma\ge r_*>0$ follows.  Without such noncollapse the correct
conclusion is \eqref{f16v65:smooth-original-Gram}, not a relative estimate
for an output whose norm might vanish.  Different operators applied to
different columns are not silently treated as one common contraction.

\subsection{A growing physical region: the exponent beats the number of marks}

\begin{theorem}[Source-weighted exclusion on a declared physical trajectory]
\label{f16v65:physical-exclusion}
Suppose $\gamma_a\to\infty$ and impose the trajectory hypothesis
\begin{equation}
 \log(1/a)=c\gamma_a+O(\log\gamma_a),\qquad c>0.
 \label{f16v65:trajectory}
\end{equation}
Let $M_a\le C a^{-p}\gamma_a^q$ for fixed $C,p,q\ge0$.  For a fixed
polynomial support and degree profile, the relative bad-source form in
\eqref{f16v65:weighted-form} tends to zero whenever
\begin{equation}
                         \Delta>pc.
 \label{f16v65:exponent-condition}
\end{equation}
More explicitly its bound is
$\exp[-(\Delta-pc)\gamma_a+O(\log\gamma_a)]$ once below one.
The same conclusion holds relative to the original Gram after any of
the common smoothing operations above, uniformly in their number of
steps.  It holds relative to the smoothed Gram if
$\log(r_{\gamma_a}^{-1})=o(\gamma_a)$.
For a growing polynomial profile the corresponding sufficient condition
is
\begin{equation}
 \limsup_{a\downarrow0}
  \frac{\log\mathcal L_{\gamma_a,\mathbf D_a}
                +\log M_a+\log(r_{\gamma_a}^{-1})}{\gamma_a}
       <\Delta,
 \label{f16v65:growing-profile}
\end{equation}
where the retention term is omitted for an original-Gram conclusion.
\end{theorem}
\begin{proof}
The logarithm of the uncapped bound is
$-\Delta\gamma_a+6v_bB_{\gamma_a}
 +\log M_a+\log\mathcal L_{\gamma_a,\mathbf D_a}$.
For a fixed profile the last term and $B_{\gamma_a}$ are
$O(\log\gamma_a)$, while $\log M_a\le pc\gamma_a+O(\log\gamma_a)$.
This proves the first assertion.  The smoothing proposition adds only
the stated retention cost when a relative output estimate is requested.
The same calculation proves \eqref{f16v65:growing-profile}.
\end{proof}

The coefficient $c=3\pi^2/11$ displayed in
\eqref{f16v55:asymptotic-scaling} gives two concrete applications:
\begin{center}\small
\begin{tabular}{@{}p{0.31\textwidth}p{0.31\textwidth}p{0.29\textwidth}@{}}
\toprule
Marked set & Event at a mark & Positive exponent margin\\
\midrule
Fixed physical tube, $M=O(a^{-1})$ &
One edge, $|s_e|\le1/4$ & $11/4-3\pi^2/11$\\[1mm]
Fixed physical four-volume, $M=O(a^{-4})$ &
One square, $\sum_i|s_i|\le1$ & $11-12\pi^2/11$\\
\bottomrule
\end{tabular}
\end{center}
The margins are approximately $0.05829$ and $0.23316$, respectively.
All large powers of $\gamma$ remain in the bound; the assertion is their
eventual domination by these strictly positive exponential margins.
This uses V55's scaling relation as a hypothesis, not a new calibration
of the nonperturbative continuum theory.

For the second row the explicit private packet has the entirely native
bound
\begin{equation}
 \mathbb E[FF^*\mathbf1_{\mathcal A_a}]
 \preceq
 \min\{1,384000(1+6\gamma_a)^2
             M_a e^{-11\gamma_a+384B_{\gamma_a}}\}\,G_a.
 \label{f16v65:explicit-physical-packet}
\end{equation}
Any boundary sequence in Proposition~\ref{f16v64:robust-cancellation}
with $\gamma\sum_i|s_i|\to0$ eventually lies in
$\mathcal C_q(1)$ at a marked slice.  Thus the estimate covers an actual
near-cancelling neighborhood of that spectral obstruction.  It does not
assign positive probability to one exact static exterior, nor treat a
whole time-dependent boundary as a homogeneous tube transfer.

\subsection{What is removed, and the remaining source-resolved dynamics}

The exceptional-set term for this explicit obstruction is now populated:
for fixed Wilson-polynomial families, its relative native Gram tends to
zero, even on the physical regions just specified.  The complete retained
predictor and conditional-innovation terms on that same event are paid
by one positive decomposition.  A coefficient-uniform fourth-moment
hypothesis is no longer needed for this conclusion.

The restriction to source complexity is essential.  V64's bounded
rare-event test $\mathbf1_{\mathcal A}-\mu(\mathcal A)$ still carries
almost all of its normalized variance on $\mathcal A$.  It does not lie
in a fixed finite-degree Wilson-polynomial space.  Approximating such
cutoffs with increasing degree must pay
\eqref{f16v65:growing-profile}; density of polynomials at each fixed
regulator does not make that cost uniform.  The union of the microscopic
polynomial spaces is not thereby identified with a nontrivial physical
continuum source core.

Nor is the complement of $\mathcal C_q(1)$ a coherent exterior.  It
includes partially cancelling staples, changing transports, correlated
boundary predictors, and induced interactions from previous integration.
Excluding this small exceptional class gives no lower bound on a native
physical-time innovation there.  In particular, conditional prediction
need not commute with time transfer: local moment and tail estimates
cannot be multiplied into a native semigroup contraction.

The next upstream calculation should therefore keep a common actual
retained sigma-field, source normalization, and physical delay, and
estimate the source-directed \emph{collective predictor return} on this
complement.  One concrete formulation is to bound, for a declared
smoothed bank $F_a$ and fixed $\tau>0$,
\[
 \operatorname{Cov}_{\rm ret}
 \bigl(\mathbb E[F_a(U_0)\mid\Xi_a],
       \mathbb E[F_a(U_{\lceil\tau/a\rceil})\mid\Xi_a]\bigr)
\]
with the accompanying conditional mixed covariance in the same
law-of-total-covariance identity.  The exceptional pieces from the
specific cancelling cells now have the bounds above.  Noncancellation
alone does not pay the remaining signed return.  Repeating a coherent
cell Hessian, an unweighted defect density, or the bare chessboard theorem
would not advance that missing step.

\subsection{Verification, preservation, and result ledger}

The companion diagnostic is
\begin{center}\small
\path{verify_ym_f16_v65_source_exclusion.py}.
\end{center}
It separates exact block-incidence and cube-moment calculations,
quaternion and gauge identities, normalized one-link compact quadrature,
finite-probability source/predictor identities, and a separately labelled
$\mathbb Z_2$ gauge analogue of reflected defect bounds.  Finite numerical
checks do not prove the all-coupling Wilson statements.  No full
four-dimensional Wilson-vacuum simulation, outward-rounded interval
certificate, or proof-assistant verification is claimed.

Removing this delimited appendix and reversing the single displayed
version change recovers V64 byte for byte.  The f14 and f15 comparison
archives and all earlier mathematical sections remain unchanged.

\begin{record}[V65 result ledger]
\label{f16v65:ledger}
\emph{Proved here:} disjoint three- and twelve-plaquette witnesses for the
specified cancelling events; their unsplit native exponential bounds
and compatible joint-event bounds; a quantitative one-link polynomial
norm estimate; a coefficient-uniform native leverage bound obtained
without a sequential marginal ansatz; relative bad-source forms for
finite Wilson-polynomial families and the explicit private packet;
exact weighted predictor/innovation payment; propagation through common
Markov smoothing; and the stated source-weighted physical-region
exclusions with every mark-count and retention cost retained.

\emph{Inherited:} bare-Wilson site reflection positivity and the classical
closed-block chessboard argument already used in f14; V64's cap
normalization and cancellation geometry; V58's compact one-link variance
and private-packet entrance bounds; the original physical transfer and
clock.  General chessboard suppression is not claimed as new.

\emph{Not asserted:} arbitrary-boundary rare-event bounds; independent
bad blocks; a uniform source bound over increasing polynomial degree;
noncollapse after physical smoothing; elimination of every adverse
boundary; or contraction of the actual retained collective predictor
from the good-event label alone.

\emph{Still unproved:} a nontrivial interacting physical source limit,
source-resolved control of the retained dynamics at a fixed physical
delay with common dense-core constants, and the continuum Yang--Mills
mass gap.
\end{record}

\begin{firewall}
The event estimate and the source norms use the same native law.
The joint witness is not split into fictitiously independent plaquettes.
Source normalization, increasing support, time smoothing, and the number
of physical-region marks all retain their stated costs.  The result
removes an explicit exceptional-source obstruction, not the remaining
collective dynamics.  V1--V64 remain unchanged.
\end{firewall}
% END F16 V65 APPEND

% BEGIN F16 V66 APPEND -- 2026-09-18
% Append-only continuation of V65.  The preceding mathematical body is unchanged.
\clearpage
\section[V66: temporal windows and the reflected source-retention test]{V66: temporal windows\\and the reflected source-retention test}
\label{f16v66:context}

V65 pays a specified exceptional geometry under the native Wilson law,
including source coefficients, but leaves a retention factor after
physical-time smoothing.  We go upstream to the source construction
rather than estimate another conditional cell.  An actual stationary
time window has a universal, polynomial retention bound and a relative
translation modulus.  These statements populate V65's exceptional-source
estimate for normalized \emph{history} sources without assuming retention.
A source-attached deletion also avoids a union over every time mark.
Neither operation changes the probability law.

There is an essential second calculation.  The ordinary history Gram is
not the reflected Gram of the same positive-time observable.  The latter
is exactly a native boundary-predictor Gram, with a different spectral
weight.  We compute both weights and their mixed-time versions.  A literal
compact Wilson ribbon then shows that window normalization can produce a
continuous, nonzero limiting history covariance while every reflected
source norm vanishes.  A smooth temporal profile has the same problem;
it is not an artifact of a sharp window.  Thus replacing the missing
physical retention by the history bound would be an invalid closure.

\subsection{Ownership and the native stationary setup}

The reflection quotient and positive transfer are established machinery,
not new reconstruction theorems.  For background see K.~Osterwalder and
R.~Schrader, \emph{Axioms for Euclidean Green's functions}, Commun.\ Math.\
Phys.\ \textbf{31} (1973), 83--112, and its corrected continuation,
\textbf{42} (1975), 281--305,
\href{https://doi.org/10.1007/BF01608978}{doi:10.1007/BF01608978};
K.~Osterwalder and E.~Seiler, \emph{Gauge field theories on a lattice},
Ann.\ Phys.\ \textbf{110} (1978), 440--471,
\href{https://doi.org/10.1016/0003-4916(78)90039-8}{doi:10.1016/0003-4916(78)90039-8}.
The finite identities needed below are proved directly.  We do not import
a continuum existence or gap conclusion from these references.

The f14 archive's V52, label \texttt{thm:v52-physical-source}, already
separates a slice-local reflected norm at fixed spatial depth from
cofinal noncollapse.  Its V53 treats a different, strong-coupling clock
obstruction.  V60 in the present lineage requires relative short-time
continuity; V65 proves native polynomial norm and event estimates.  None
of these statements identifies an ordinary, overlapping temporal-window
Gram with its positive-time reflected Gram.  That distinction, the exact
window bounds and source-attached error, and the compact window example
are the new calculations here.  The historical bodies are not recertified.

At a fixed spatial regulator let $\pi_a$ be the actual vacuum slice law
and $P_a$ its stationary ground-state Markov operator.  We use the inherited
unitary equivalence to the \emph{same} normalized physical Wilson transfer:
\begin{equation}
 P_a1=1,\qquad P_a=P_a^*,\qquad 0\preceq P_a\preceq I,
 \qquad P_a\simeq\mathcal T_a.
 \label{f16v66:native-P}
\end{equation}
The lattice time spacing remains $a$.  This is not a conditional update
operator.  Let $(U_j)_{j\in\mathbb Z}$ be the corresponding two-sided
stationary Markov history.  Use a finite family of real, centered,
gauge-invariant slice polynomials $f=(f_1,\ldots,f_d)$, and write
\[
 Vc=\sum_i c_i f_i,\qquad G=V^*V,\qquad
 C_j=V^*P_a^jV\quad(j\ge0).
\]
Complex coefficients use the Hermitian extension.  All inverses below
are taken on the declared exact source quotient; null vectors are never
removed by a numerical threshold.  Native finite-mark probability bounds
are obtained by the temporal limits admitted in V65, not by inserting an
arbitrary endpoint density into a reflection block.

For $N\ge1$ define the actual history source and two scalar polynomials
\begin{align}
 A_{t,N}&=\frac1N\sum_{j=0}^{N-1}f(U_{t+j}),
     &D_N&=\mathbb E[A_{0,N}A_{0,N}^*],
       \label{f16v66:window}\\
 b_N(z)&=\frac1N\sum_{j=0}^{N-1}z^j,
     &r_N(z)&=\frac1{N^2}
          \left[N+2\sum_{j=1}^{N-1}(N-j)z^j\right].
       \label{f16v66:weights}
\end{align}
The letter $r_N$ denotes this explicit polynomial, not V65's unknown
retention constant.  We take $b_N(1)=r_N(1)=1$.

\begin{proposition}[Exact window Grams and a universal history retention]
\label{f16v66:window-Gram}
Under \eqref{f16v66:native-P},
\begin{equation}
 \boxed{\quad D_N=V^*r_N(P_a)V,\qquad
            \frac1N G\preceq D_N\preceq G.\quad}
 \label{f16v66:retention}
\end{equation}
If $B_N=b_N(P_a)V$ is the endpoint-predictor synthesis, then
\begin{equation}
 \frac1N D_N\preceq B_N^*B_N\preceq D_N,
 \qquad \ker D_N=\ker G=\ker(B_N^*B_N).
 \label{f16v66:endpoint-retention}
\end{equation}
For two nonoverlapping windows with start times $0$ and $N+m$, $m\ge0$,
\begin{equation}
 \mathbb E[A_{0,N}A_{N+m,N}^*]
             =V^*b_N(P_a)P_a^{m+1}b_N(P_a)V.
 \label{f16v66:disjoint-return}
\end{equation}
All statements hold simultaneously in every source direction.
\end{proposition}
\begin{proof}
Stationarity, reversibility and the Markov property give
$\mathbb E[\overline{Vc(U_i)}Vc(U_j)]=\langle Vc,P_a^{|j-i|}Vc\rangle$.
Counting equal lags proves the expression for $D_N$.  Every $C_j$ is
positive semidefinite; retaining its diagonal contribution proves the
lower bound, and Jensen for the average proves the upper bound.

For $0\le z\le1$, $z^{i+j}\le z^{|i-j|}$, hence $b_N(z)^2\le r_N(z)$.
For the other inequality put $S_N=N^2r_N$ and $T_N=Nb_N$.
Coefficientwise $S_N\le NT_N^2$: at degree zero both coefficients are $N$;
at degree $1\le j<N$, $2(N-j)\le N(j+1)$; all further coefficients of
$S_N$ vanish.  Thus $b_N^2\ge r_N/N$.  Functional calculus proves
\eqref{f16v66:endpoint-retention}, including the kernels.  Finally,
\[
 \frac1{N^2}\sum_{i,j=0}^{N-1}P_a^{N+m+j-i}
       =P_a^{m+1}b_N(P_a)^2.
\]
The smallest exponent is $m+1$, not $m$; this proves
\eqref{f16v66:disjoint-return} without an overlap approximation.
\end{proof}

\begin{theorem}[A relative translation modulus without a physical gap]
\label{f16v66:modulus}
For every integer $m$,
\begin{equation}
 \mathbb E[(A_{m,N}-A_{0,N})(A_{m,N}-A_{0,N})^*]
       \preceq2\min\{1,|m|/N\}\,D_N.
 \label{f16v66:translation}
\end{equation}
Consequently $N_a=\lceil\sigma/a\rceil$ with fixed $\sigma>0$, followed
by normalization on $D_{N_a}$, has a uniform relative $L^2$ translation
modulus bounded by $2(|h|+a)/\sigma$ for lattice approximations to a
physical shift $h$ with $|h|\le\sigma$.  This is a statement about history
observables, not yet a strongly continuous physical Hilbert-space limit.
\end{theorem}
\begin{proof}
It suffices to take $1\le m\le N$.  The unnormalized difference of the
two windows is the last $m$ terms minus the first $m$ terms.  The covariance
between these two disjoint sums is positive semidefinite, since it is a
sum of $C_j\succeq0$.  The difference Gram is therefore at most
$2N^{-2}V^*S_m(P_a)V$.  For $0\le z\le1$,
\[
 \frac{S_m(z)}m=1+2\sum_{j=1}^{m-1}(1-j/m)z^j
       \le1+2\sum_{j=1}^{N-1}(1-j/N)z^j=\frac{S_N(z)}N.
\]
This proves the bound for $m\le N$.  At every other lag the window cross
Gram is again a sum of positive powers, so the difference Gram is at most
$2D_N$.  Negative shifts follow by stationarity.  The statement in physical
units follows from $|m|a\le |h|+a$ and $N_aa\ge\sigma$.
\end{proof}

\subsection{A populated exceptional-source estimate after history normalization}

Assume the native input of V65 in the form
\begin{equation}
 \|Vc\|_\infty^2\le\mathcal L_\gamma\,c^*Gc.
 \label{f16v66:leverage-input}
\end{equation}
For a fixed polynomial profile this has $\log\mathcal L_\gamma=O(\log\gamma)$;
for the private three-orientation packet one can use
$\mathcal L_\gamma=384000(1+6\gamma)^2$.

\begin{theorem}[Relative native history exclusion, with its full clock cost]
\label{f16v66:history-exclusion}
For any history event $\mathcal A$ with probability at most $p$,
\begin{equation}
 \boxed{\quad
 \mathbb E[A_{0,N}A_{0,N}^*\mathbf1_{\mathcal A}]
       \preceq\min\{1,N\mathcal L_\gamma p\}\,D_N.
 \quad}
 \label{f16v66:history-bad}
\end{equation}
In particular, for a union of $M$ of V65's complete-cell events
$\sum_i|s_i|\le1$, the bound is
\begin{equation}
 \min\{1,N\mathcal L_\gamma M
                 e^{-11\gamma+384B_\gamma}\}\,D_N.
 \label{f16v66:history-native}
\end{equation}
No unproved retention factor is used.  Under
$\log(1/a)=c\gamma_a+O(\log\gamma_a)$,
$N_a=\lceil\sigma/a\rceil$, $M_a\le Ca^{-p_0}\gamma_a^q$,
and a fixed polynomial profile, this bound tends to zero provided
\begin{equation}
                       11>(p_0+1)c.
 \label{f16v66:history-exponent}
\end{equation}
\end{theorem}
\begin{proof}
For each $c$, the supremum of its window average is at most $\|Vc\|_\infty$.
Multiply its square by $p$, apply \eqref{f16v66:leverage-input}, then use
$G\preceq ND_N$.  The event form is also bounded by $D_N$, giving the
minimum.  V65 gives $p\le\min\{1,M e^{-11\gamma+384B_\gamma}\}$.
Taking logarithms pays $\log N=c\gamma+O(\log\gamma)$ in addition to
$\log M$.  This is the extra one-power clock cost in
\eqref{f16v66:history-exponent}.
\end{proof}

For $c=3\pi^2/11$ as a \emph{trajectory hypothesis}, a fixed physical tube
has $p_0=1$ and margin $11-6\pi^2/11>0$.  A fixed physical four-volume has
$p_0=4$ and $11-15\pi^2/11<0$.  The second case is \emph{not} covered by
this universally normalized window estimate.  V65's original-Gram
four-volume bound remains valid; it cannot be silently promoted to a
window-relative bound without paying the new factor.  The constants and
large polynomial prefactors are retained, not claimed numerically useful
at moderate coefficients.

\begin{theorem}[Source-attached deletion avoids the all-time union]
\label{f16v66:deletion}
Let $\mathcal B_j$ be a stationary family of local history events, with
$\mathbb P(\mathcal B_j)\le\zeta$.  Define
\begin{equation}
 R_{t,N}=\frac1N\sum_{j=0}^{N-1}
              f(U_{t+j})\mathbf1_{\mathcal B_{t+j}},\qquad
 \widetilde A_{t,N}=A_{t,N}-R_{t,N}+\mathbb E R_{t,N}.
 \label{f16v66:delete-source}
\end{equation}
For $\epsilon=N\mathcal L_\gamma\zeta$,
\begin{align}
 \mathbb E[(A_{t,N}-\widetilde A_{t,N})
                  (A_{t,N}-\widetilde A_{t,N})^*]
                      &\preceq\epsilon D_N,
       \label{f16v66:delete-error}\\
 (1-\sqrt\epsilon)_+^2D_N\preceq
        \mathbb E[\widetilde A_{t,N}\widetilde A_{t,N}^*]
                      &\preceq(1+\sqrt\epsilon)^2D_N.
       \label{f16v66:delete-frame}
\end{align}
For any two translated windows, the difference between their original
and deleted cross Grams has normalized operator norm at most
$2\sqrt\epsilon+\epsilon$, where normalization is on $D_N$.
The same single bound controls the \emph{sum} of retained-predictor and
conditional-innovation cross Grams for any one common retained
sigma-field.
\end{theorem}
\begin{proof}
Jensen and \eqref{f16v66:leverage-input} give, without independence,
\[
 \mathbb E|c^*R_{t,N}|^2
 \le\frac1N\sum_{j=0}^{N-1}
       \mathbb E[|c^*f(U_{t+j})|^2\mathbf1_{\mathcal B_{t+j}}]
 \le\mathcal L_\gamma\zeta\,c^*Gc
 \le\epsilon c^*D_Nc.
\]
Centering $R$ only decreases this form.  The triangle and reverse triangle
inequalities give \eqref{f16v66:delete-frame}.  On the normalized source
quotient each original synthesis has norm one and each error synthesis
has norm at most $\sqrt\epsilon$.  Expanding the difference of two cross
Grams proves $2\sqrt\epsilon+\epsilon$.

Let $Q$ be conditional expectation on the common retained sigma-field.
The map $J:g\mapsto(Qg,(I-Q)g)$ into the orthogonal direct sum is an
isometry.  Apply the same synthesis estimate after $J$ to both windows.
The first cross pairing is the predictor covariance; the second is the
averaged conditional mixed covariance.  Their sum is the original cross
pairing.  This proves the bound on their sum once, rather than adding two
independent copies of the available error budget.
\end{proof}

For the cell event of V65 take
$\zeta=\min\{1,e^{-11\gamma+384B_\gamma}\}$.  A physical window then has
$\epsilon=\exp[-(11-c)\gamma+O(\log\gamma)]$ for a fixed source profile.
There is no further count $M=N$: the deletion follows the source through
time rather than require the absence of every bad mark.  This is an
observable approximation under the \emph{unchanged} native law.  It does
not condition the whole history to be good, bound the remaining predictor,
or turn the deleted source into a finite-degree polynomial.  Reflection
supports must still be checked when the local indicators themselves extend
in time.  The conclusions extend by congruence to the exact redundant
private-packet relation of V65.

\subsection{The physical reflection form is a different Gram}

Let time reflection fix slice zero.  For $k\ge0$, $A_{k,N}$ is a
positive-time source, and its reflection has support
$\{-k-N+1,\ldots,-k\}$.  Denote its reflected Gram by $H_{k,N}$.

\begin{theorem}[Complete native boundary predictor and reflected return]
\label{f16v66:reflected}
In the stationary native Markov law,
\begin{align}
 \mathbb E[A_{k,N}\mid U_0]&=P_a^kb_N(P_a)f(U_0),
            \label{f16v66:boundary-predictor}\\
 \boxed{\quad H_{k,N}
       =V^*P_a^{2k}b_N(P_a)^2V.\quad}
            \label{f16v66:OS-Gram}
\end{align}
For a further translation by $m\ge0$ in the positive half,
\begin{equation}
 H_{k,N}(m)=V^*P_a^{2k+m}b_N(P_a)^2V.
 \label{f16v66:OS-delay}
\end{equation}
The conditional mixed covariance of the positive and negative half
histories, given $U_0$, is zero.  Their entire reflected pairing is the
boundary-predictor pairing.  In particular $0\preceq H_{k,N}\preceq D_N$,
but \eqref{f16v66:retention} does not lower-bound $H_{k,N}$ at a positive
physical buffer.
\end{theorem}
\begin{proof}
The Markov property gives the first formula term by term.  Past and future
are conditionally independent given $U_0$; reversibility identifies the
reflected predictor with the same $P_a^kb_N(P_a)f$.  Taking their pairing
proves \eqref{f16v66:OS-Gram}.  Translating just the positive source adds
$P_a^m$, proving \eqref{f16v66:OS-delay}.  Alternatively the double sum
has exponent $2k+i+j$ (or $2k+m+i+j$), never $|i-j|$.
The upper bound follows from $P_a^{2k}\preceq I$ and $b_N^2\le r_N$.
The lower window bound was obtained from diagonal terms with exponent
zero.  Those terms do not occur in the reflected sum when $k>0$.
\end{proof}

This theorem uses the full native slice as a Markov separator.  A spatially
retained sigma-field need not be such a separator, so its conditional
mixed covariance is not set to zero by this calculation.  It instead
provides a direct upstream test on the same physical transfer.  It is
also precisely the common Markov smoothing $P_a^kb_N(P_a)$ allowed in
V65, now with its output Gram explicitly identified.  The history norm,
the endpoint-predictor norm and the positive-buffer reflected norm must
not be interchanged.

\begin{proposition}[A finite-physical-energy test with both directions paid]
\label{f16v66:finite-energy}
Let $N_a a\to\sigma>0$ and $k_a a\to\delta>0$.  On the exact quotient,
put $\mathsf H_a=-a^{-1}\log P_a$ and define the positive matrix measure
\begin{equation}
 \mathsf M_a(J)=D_{N_a}^{-1/2}V^*r_{N_a}(P_a)
       \mathbf1_J(\mathsf H_a)V D_{N_a}^{-1/2}.
 \label{f16v66:energy-measure}
\end{equation}
Its total mass is $I$.  A zero spectral value of $P_a$, if present, is
assigned energy $+\infty$.  If $k_aa\ge\delta_0>0$ and
$a(k_a+N_a-1)\le T_0$, then for every $E>0$,
\begin{equation}
 \boxed{\quad
 e^{-2T_0E}\mathsf M_a([0,E])
 \preceq D_{N_a}^{-1/2}H_{k_a,N_a}D_{N_a}^{-1/2}
 \preceq \mathsf M_a([0,E])+e^{-2\delta_0E}I.\quad}
 \label{f16v66:two-sided-test}
\end{equation}
Consequently a regulator-uniform positive lower bound on the normalized
reflected Gram is equivalent to a regulator-uniform positive lower bound
on $\mathsf M_a([0,E])$ for some fixed finite $E$.  This requires only a
positive \emph{fraction} of bounded physical energy, not spectral tightness
of the entire family, and is not a mass-gap conclusion.
\end{proposition}
\begin{proof}
All factors in \eqref{f16v66:energy-measure} are bounded commuting spectral
functions except for the harmless notation for their spectral parameter.
There is no logarithmic-domain assumption on $V$.  The normalized reflected
Gram is the integral against $\mathsf M_a$ of
\[
 h_a(E')=e^{-2k_aaE'}
       \frac{b_{N_a}(e^{-aE'})^2}{r_{N_a}(e^{-aE'})}.
\]
For $E'\ge0$, $0\le h_a(E')\le e^{-2\delta_0E'}$.
On $0\le E'\le E$, use $b_N(z)\ge z^{N-1}$ and $r_N(z)\le1$ to get
$h_a(E')\ge e^{-2T_0E}$.  Split the integral at $E$ to prove
\eqref{f16v66:two-sided-test}.  If the middle matrix is at least $rI$,
choose $E$ with $e^{-2\delta_0E}\le r/2$ to obtain
$\mathsf M_a([0,E])\succeq(r/2)I$.  Conversely a lower bound $mI$ on
that measure gives $me^{-2T_0E}I$ on the reflected Gram.
\end{proof}

For bounded physical $E$ the weights have the explicit limits, with
$x=\sigma E$ and continuous values at $x=0$,
\begin{equation}
 b_{N_a}(e^{-aE})\longrightarrow\frac{1-e^{-x}}x,\qquad
 r_{N_a}(e^{-aE})\longrightarrow\frac{2(x-1+e^{-x})}{x^2}.
 \label{f16v66:finite-E-weights}
\end{equation}
These are Riemann sums.  They show that the history and reflection weights
remain comparable on any prescribed bounded physical-energy interval.
They do not prove that a native source places any norm in that interval.
Conversely a positive reflected Gram may be carried by energies approaching
zero; it supplies no exclusion of $(0,m)$ and hence no positive mass.

The deletion theorem controls reflected errors in \emph{history} units,
since conditional expectation at slice zero is an $L^2$ contraction.
Converting those errors to relative reflected units again requires the
lower bound in \eqref{f16v66:two-sided-test}.  Thus source-attached
exceptional suppression does not conceal the remaining physical retention.

\subsection{A literal compact ribbon exposes the overlap-only limit}

The following example is an exact conditional Wilson law, not the native
unconditioned four-dimensional vacuum.  It diagnoses what the new general
history estimates can and cannot imply.

Free one spatial $\mu$ link at each integer time.  Its four complementary
spatial plaquettes have private opposite $\mu$ links at displacements
$\pm\nu,\pm\rho$.  Put the connectors and temporal retained links at
identity, and assign the opposite links so their oriented staple paths are
\begin{equation}
                         (I,-I,I,-I).
 \label{f16v66:single-cancel}
\end{equation}
All four instructions concern distinct retained links.  The spatial
one-link field cancels identically, and the two temporal plaquettes give
neighboring factors $e^{\gamma x_j\cdot x_{j+1}}$.  With the same local
boundary convention as V64, the stationary limit is exactly
\begin{equation}
 (\mathsf P_\gamma g)(X)=\frac1{\mathfrak z(\gamma)}
      \int e^{\gamma w(XY^{-1})}g(Y)\,dH(Y),
 \qquad d\pi_\gamma=dH.
 \label{f16v66:compact-chain}
\end{equation}
Here $\mathfrak z$ is the unshifted one-link factor in V63.  Choosing a
complementary path with value $I$, the closed writer
$f(X)=2w(P_{\rm out}^{-1}X)=2w(X)$ has mean zero and variance one.
The notation is a conditional restriction of a closed loop.  A triple
chosen from \eqref{f16v66:single-cancel} has norm one; the near-zero
three-staple event in V65 was never an exclusion of every slow conditional
ribbon.

\begin{theorem}[Continuous normalized history covariance with zero reflected norm]
\label{f16v66:compact-example}
For \eqref{f16v66:compact-chain},
\begin{equation}
 \mathsf P_\gamma f=\lambda_\gamma f,\qquad
 \lambda_\gamma=\frac{I_2(\gamma)}{I_1(\gamma)}\in(0,1),\qquad
 1-\lambda_\gamma=\frac{3}{2\gamma}+O(\gamma^{-2}).
 \label{f16v66:compact-eigen}
\end{equation}
Take $\gamma_a\to\infty$, $a\gamma_a\to0$ and
$N_a=\lceil\sigma/a\rceil$, $\sigma>0$.  Then
\begin{equation}
 D_{N_a}=r_{N_a}(\lambda_{\gamma_a})
                    \sim\frac{4\gamma_a}{3N_a}.
 \label{f16v66:compact-window-size}
\end{equation}
For window averages normalized to variance one, at shifts with $m_aa\to t$,
\begin{equation}
 \frac{\mathbb E[A_{0,N_a}A_{m_a,N_a}]}{D_{N_a}}
                    \longrightarrow(1-|t|/\sigma)_+.
 \label{f16v66:triangle-limit}
\end{equation}
Nevertheless, even with no initial buffer,
\begin{equation}
 \frac{H_{k_a,N_a}}{D_{N_a}}
  =\lambda_{\gamma_a}^{2k_a}
     \frac{b_{N_a}(\lambda_{\gamma_a})^2}
          {r_{N_a}(\lambda_{\gamma_a})}
  \sim\lambda_{\gamma_a}^{2k_a}\frac{\gamma_a}{3N_a}
       \longrightarrow0\qquad(k_a\ge0).
 \label{f16v66:compact-OS-zero}
\end{equation}
These are exact two-point asymptotics of the compact law; no Gaussian
replacement or full-law central limit theorem is used.
\end{theorem}
\begin{proof}
Central convolution acts on the fundamental matrix coefficients with
multiplier $I_2/I_1$, as also follows by integrating $Y$ in coordinates
with $X=I$ and rotating back.  Thus $f$ has the stated eigenvalue.
For completeness its large-$\gamma$ value is its one-spin mean $x^0$.
Set $v=1-x^0$ in the scalar Haar density.  The ratio giving
$\mathbb E v$ has numerator and denominator proportional to
\[
 \int_0^2 v e^{-\gamma v}\sqrt{v(2-v)}\,dv,
 \qquad \int_0^2 e^{-\gamma v}\sqrt{v(2-v)}\,dv.
\]
Scaling $u=\gamma v$ gives leading ratio $3/(2\gamma)$.
On $v\le1$, expand $\sqrt{2-v}$ with a uniform first-order remainder;
on $v\ge1$ the integral is exponentially small.  The remaining gamma
integrals bound the error by $O(\gamma^{-2})$.

Write $\lambda=\lambda_{\gamma_a}$ and
$L_\lambda=\sum_{j\in\mathbb Z}\lambda^{|j|}=(1+\lambda)/(1-\lambda)$.
Let $S_{N,m}=\sum_{i,j=0}^{N-1}\lambda^{|m+j-i|}$.  Reindexing gives
\[
 S_{N,m}=\sum_{j\in\mathbb Z}(N-|m-j|)_+\lambda^{|j|}.
\]
The triangular coefficient is one-Lipschitz, so
\begin{equation}
 \left|S_{N,m}-(N-|m|)_+L_\lambda\right|
              \le\sum_{j\in\mathbb Z}|j|\lambda^{|j|}
              =\frac{2\lambda}{(1-\lambda)^2}.
 \label{f16v66:triangle-error}
\end{equation}
Since $N(1-\lambda)\to\infty$, this error is $o(NL_\lambda)$,
uniformly in $m$.  Divide first by $N^2$, then by $S_{N,0}$, to obtain
\eqref{f16v66:compact-window-size} and \eqref{f16v66:triangle-limit}.
Finally $\lambda^N\to0$, so $b_N(\lambda)\sim[N(1-\lambda)]^{-1}$ and
\[
 \frac{b_N(\lambda)^2}{r_N(\lambda)}
          \sim\frac1{N(1-\lambda^2)}\sim\frac\gamma{3N}.
\]
Multiplication by the exact factor $\lambda^{2k}$ proves the last claim,
including all nonnegative buffer sequences.
\end{proof}

In this example the physical source energy is exactly
$-a^{-1}\log\lambda_{\gamma_a}\sim3/(2a\gamma_a)\to\infty$.
The measure \eqref{f16v66:energy-measure} is concentrated at that energy.
The missing finite-energy participation is therefore explicit.
The family has a continuous limiting \emph{history} covariance after
normalization, but its physical reflected classes disappear.  The assumed
trajectory with logarithmic $\gamma_a$ lies in this example's regime;
this does not identify its conditional ribbon with the physical YM vacuum.

\begin{proposition}[Smooth profiles do not remove the obstruction]
\label{f16v66:smooth-example}
In the same compact model let $g,h\in C_c^1(\mathbb R)$ and define
$A_{g,a}=a\sum_j g(aj)f(U_j)$.  If $a\gamma_a\to0$, then
\begin{equation}
 \frac{\mathbb E[A_{g,a}A_{h,a}]}
           {aL_{\lambda_{\gamma_a}}}\longrightarrow
                       \int_{\mathbb R}g(t)h(t)\,dt.
 \label{f16v66:smooth-contact}
\end{equation}
For a nonzero $g$ supported strictly in positive physical time, the ratio
of its reflected norm to its ordinary variance tends to zero.  Choosing
$g\in C_c^\infty$ gives a smooth limiting translation covariance; no
positive physical reflected norm follows from that smoothness.
\end{proposition}
\begin{proof}
Put $J_a(k)=a\sum_j g(aj)h(a(j+k))$.  The covariance divided by
$aL_\lambda$ equals
$L_\lambda^{-1}\sum_k\lambda^{|k|}J_a(k)$.
Compact support and bounded first derivatives give
$|J_a(k)-J_a(0)|\le C a|k|$ for $a\le1$, and Riemann sums give
$J_a(0)\to\int gh$.  The remaining error is at most
\[
 \frac{Ca}{L_\lambda}\sum_k |k|\lambda^{|k|}
       =\frac{2Ca\lambda}{1-\lambda^2}\longrightarrow0.
\]
Take $h(t)=g(-t)$.  Its support is disjoint from that of $g$, so the
reflected numerator is $o(aL_\lambda)$ while the variance is asymptotic
to $aL_\lambda\int g^2>0$.  The translation statement follows by taking
$h(t)=g(t-s)$, whose limiting autocorrelation is smooth when $g$ is smooth.
\end{proof}

Support cannot be repaired by relabelling a window as a point source.
For the centered-window limiting covariance
$K_\sigma(t)=(1-|t|/\sigma)_+$, treating all positive center times as an
Osterwalder--Schrader positive-time algebra would produce, at
$t_1=\sigma/8$, $t_2=3\sigma/8$, the matrix
\begin{equation}
 (K_\sigma(t_i+t_j))_{i,j=1}^2
          =\begin{pmatrix}3/4&1/2\\1/2&1/4\end{pmatrix},
       \qquad \det=-1/16.
 \label{f16v66:false-RP}
\end{equation}
There is no violation of the native reflection positivity: those centered
windows cross the reflection plane and are not positive-time observables.
For windows wholly on the positive side, the reflected limiting pairings
are zero.  This exact check rules out interpreting overlapping-window
continuity as a new time-local physical field without its support map.

\subsection{The next source calculation and the result ledger}

For actual native time windows, we have proved retention and translation
bounds, together with estimates for exceptional source contributions.
The source-attached deletion controls the error in complete mixed
covariances, including their predictor and innovation decomposition,
without a growing all-good event.  These are new applications of V65's
native estimates rather than another proof of its chessboard or
finite-degree norm inequality.

For a physical construction the prior question is now sharply identified:
for the chosen spatially supported, temporally smeared Wilson bank, does
any regulator-uniform fraction of its normalized source measure remain
at bounded \emph{physical} energy?  The equivalent reflected matrix is
\eqref{f16v66:OS-Gram} with $k_aa\to\delta>0$, not $D_N$ and not the
coherent conditional cell Gram.  It is a normalized mixed derivative of
the full native generating functional with the source and its reflected
source, not a derivative of an unmarked pressure bound.  No lower bound
on that mixed derivative has been populated here.

A next genuinely dynamical result must establish such a bound for an
explicit native bank, or exhibit its ultraviolet escape and replace that
bank by a specified spatially regularized one.  Only after a surviving
reflected bank has been identified should one claim a bound on
$H_{k,N}(m)$ relative to $H_{k,N}$ at $ma\to\tau>0$.  A bound relative to
$D_N$ can instead describe a vector that disappears in the physical
quotient, as the compact example proves.  Conversely finite-energy
participation by itself permits energies near zero and is not a gap.
Spatially growing source complexity, the probability law of the full
retained exterior, the original physical clock and dense-core common
constants remain explicit obligations.

The companion diagnostic is
\begin{center}\small
\path{verify_ym_f16_v66_temporal_windows.py}.
\end{center}
It separates exact rational window kernels and covariance identities,
a finite reversible Markov enumeration, literal quaternion conditional
plaquettes, and compact one-link integral diagnostics.  Large-$\gamma$
series evaluations and quadrature are ordinary floating-point checks,
not outward-rounded certificates.  The continuum conclusions above are
two-point statements in an explicitly conditional example; no full
four-dimensional Wilson-vacuum simulation or interacting field-law limit
is claimed.  Removing this delimited appendix and restoring the title
version recovers V65 byte for byte.

\clearpage
\begin{record}[V66 result ledger]
\label{f16v66:ledger}
\emph{Proved here:} exact ordinary, endpoint and reflected time-window
Gram formulas with mixed delays; universal history retention and a relative
physical translation modulus; a populated native exceptional-source bound
in history-relative units with its extra clock exponent; source-attached
deletion with a single predictor/innovation error budget; an explicit
two-sided bounded-physical-energy test for reflected noncollapse; a literal
compact Wilson ribbon with nonzero continuous window covariance but
vanishing reflected norms; and the persistence of that distinction for
smooth temporal profiles, with an explicit support-misidentification test.

\emph{Inherited:} the actual vacuum-normalized Wilson transfer and its
ground-state Markov law, finite-polynomial native leverage, the specific
cancelling-event bounds, local Wilson boundary discipline and exact source
relations.  The spectral theorem, Markov reflection quotient and general
reconstruction framework are not claimed as new results.

\emph{Not asserted:} physical noncollapse from a normalized history Gram;
a point-local reflection algebra for extended temporal windows; a
four-volume window-relative exclusion when its exponent condition fails;
a full-law limit from covariance convergence; or replacement of the
native vacuum by the conditional ribbon used to test the inference.

\emph{Still unproved:} bounded-physical-energy participation for a chosen
nontrivial native continuum source bank; the interacting reflected source
limit; source-directed physical-time contraction with common dense-core
constants; and the continuum Yang--Mills mass gap.
\end{record}

\begin{firewall}
Temporal averaging is evaluated under the native law, and its support and
normalization are not erased.  The ordinary variance and physical reflected
norm are different matrices.  Exceptional-source control does not supply
the missing low-energy participation.  The compact example tests a proposed
inference, not the existence of the sought four-dimensional theory.
V1--V65 remain unchanged.
\end{firewall}
% END F16 V66 APPEND

% BEGIN F16 V67 APPEND -- 2026-09-18
% Append-only continuation of V66; not the older parallel V67 attachment.
\clearpage
\section[V67: contact divergence and spatial constituent sources]{V67: contact divergence\\and spatial constituent sources}
\label{f16v67:context}

V66 distinguishes an ordinary history Gram from its reflected Gram.  Its
bounded-energy criterion is exact for the \emph{history-normalized} family
used there.  It must not be promoted into a necessary condition for every
renormalization of a composite field: an ordinary contact variance may
diverge while separated reflected pairings survive.  This continuation
makes that distinction quantitative in the \emph{selected native Wilson
sequence already treated in f15 V65}.  Its positive-time comparison is
inherited, not reproved as a new continuum result.  What is new is the
complete leading ordinary-history variance, the resulting normalization
comparison, and the incompatibility of that mesoscopic window with the
present proposed exponential physical clock.

We then specify a different source, rather than prescribe another
unsupplied retention constant.  A spatial adjoint covariant heat operator
is applied to the plaquette curvature \emph{before} taking its quadratic
invariant.  It is defined on one native slice, admits an exact parallel-path
expansion and a uniform spatial tail estimate, and does not diffuse through
the time-reflection plane.  A free-curvature calculation proves finite
ordinary and nonzero reflected norms for this construction at fixed spatial
radius.  That calculation is explicitly a comparison model, not an estimate
of the interacting Wilson vacuum.  A final reflected-error identity identifies
which native positive-separation estimate is still needed.

\subsection{Ownership audit and the selected native input}

The f15 archive's V65 already gives dimension-four spatial sources
(\path{f15v65:smeared} and \path{f15v65:actual-spectrum}),
nonzero time-damped norms, and a soft limiting spectral measure on a
particular selected sequence.  In particular it explicitly leaves the
undamped contact norm uncontrolled.  We use precisely its covariance
comparison below.  We neither recertify its earlier hypotheses nor replace
its selected sequence by an arbitrary weak-coupling trajectory.  The older
parallel file named V67 is unchanged and is not the predecessor of this
appendix.  The predecessor is the current Post-Fifteenth-Archive V66.

Spatial covariant smearing is established practice.  Relevant background
includes G.~M.~von Hippel et al., \emph{The Shape of Covariantly Smeared
Sources in Lattice QCD}, \href{https://arxiv.org/abs/1306.1440}{arXiv:1306.1440}.
K.~Sakai and S.~Sasaki, \emph{Glueball spectroscopy in lattice QCD using
gradient flow}, \href{https://doi.org/10.1103/PhysRevD.107.034510}{Phys.\ Rev.\ D
\textbf{107} (2023), 034510}, separate spatial from space--time diffusion.  No numerical
spectroscopy result, full gradient-flow finiteness theorem, or
boundary-uniform ground-state overlap is used here.  The adjoint heat
source and its estimates below are defined and proved directly.

Use the f15 selected vacuum on
$\mathbb Z\times(\mathbb Z/M\mathbb Z)^3$.  Let $X_{M,(j,x)}$ denote its
thermal spatial $12$-plaquette energy $\gamma_M s_{(j,x;12)}$, and use a bar
for centering in this same vacuum.  Its inherited input is
\begin{equation}
 \left|\operatorname{Cov}(X_{M,(0,x)},X_{M,(n,y)})
                    -R(n,x-y)\right|\le C\delta_M,
 \qquad R(z)=\tfrac32 A(z)^2,
 \label{f16v67:archive-input}
\end{equation}
for all time differences and spatial differences in the admitted lift,
including zero separation.  Here
\begin{align}
 \delta_M&=\sqrt{\vartheta_M}+\gamma_M^{-1/16}+M^{-2}\longrightarrow0,
       \label{f16v67:delta}\\
 A(z)&=\int_{[-\pi,\pi]^4}e^{ip\cdot z}m(p)\frac{d^4p}{(2\pi)^4},
 &m(p)&=\frac{\widehat p_1^2+\widehat p_2^2}
                         {\sum_{\mu=0}^3\widehat p_\mu^2},
 &\widehat p_\mu&=2\sin(p_\mu/2).
       \label{f16v67:projection}
\end{align}
The value of $m$ at the single origin is immaterial.  This is the diagonal
spatial-curvature projection multiplier in the archive, not a new
covariance model assigned to the native law.  Select integers $b=b_M$ with
\begin{equation}
 b\to\infty,\qquad b/M\to0,\qquad \delta_M b^8\to0.
 \label{f16v67:mesoscopic}
\end{equation}
All results in the next three subsections are conditional on these inherited
selected-sequence inputs.

\subsection{The contact susceptibility and the actual ordinary source norm}

\begin{proposition}[A finite, strictly positive lattice contact susceptibility]
\label{f16v67:susceptibility}
The nonnegative kernel $R$ in \eqref{f16v67:archive-input} is summable on
$\mathbb Z^4$, and
\begin{equation}
 \chi_{12}:=\sum_{z\in\mathbb Z^4}R(z)
      =\frac32\int_{[-\pi,\pi]^4}m(p)^2\frac{d^4p}{(2\pi)^4},
 \qquad \frac38<\chi_{12}<\frac34.
 \label{f16v67:chi}
\end{equation}
This constant includes all microscopic space--time lags, not just the
single coincident plaquette.
\end{proposition}
\begin{proof}
The multiplier is real, even and lies between zero and one, so $A$ is real.
Plancherel gives the equality and summability.  Permutation symmetry of the
four momentum coordinates gives $\int m=1/2$.  The multiplier is neither
constant nor almost everywhere in $\{0,1\}$.  Strict Jensen and
$m^2<m$ on a set of positive measure give the two strict inequalities.
No spatial or temporal mixing theorem for the native law is inferred from
this summability of its \emph{comparison} kernel.
\end{proof}

Fix linearly independent real $f_1,\ldots,f_d\in C_c^\infty(\mathbb R^3)$
and a nonnegative, nonzero $h\in C_c^\infty((\delta,T))$, where
$0<\delta<T<\infty$.  Define the actual native history sources
\begin{align}
 S_{M,i}&=\sum_{j\in\mathbb Z}\sum_{x\in\mathbb Z^3}
             h(j/b)f_i(x/b)\,\overline X_{M,(j,x)},
       \label{f16v67:native-source}\\
 D_M&=\mathbb E[S_M S_M^*],
 &H_M(n)&=\mathbb E[(\theta S_M)(\tau_n S_M)^*],\qquad n\ge0.
       \label{f16v67:two-Grams}
\end{align}
The spatial support is embedded in the admitted lift for large $M$.
The coefficient convention is exactly the temporal Riemann sum of the
archive's $b\sum_x f_i(x/b)\overline X_{M,(j,x)}$: its additional factor
$1/b$ cancels the displayed $b$.  No variance normalization has been applied.
The time reflection is about slice zero and $\tau_n$ translates into the
positive half.

For $\tau>0$ retain the archive's kernel
\begin{align}
 a_{12}(\tau,z)&=\int_{\mathbb R^3}e^{ip\cdot z}
              w_{12}(p)e^{-\tau|p|}\frac{d^3p}{(2\pi)^3},
 &w_{12}(p)&=\frac{p_1^2+p_2^2}{2|p|},
       \label{f16v67:positive-kernel}\\
 K_{ij}(\tau)&=\frac32\int f_i(x)f_j(y)a_{12}(\tau,x-y)^2\,d^3x\,d^3y,
       \label{f16v67:K}\\
 H^*_{ij}(u)&=\int h(s)h(t)K_{ij}(s+t+u)\,ds\,dt,
 \qquad J_{ij}=\int f_i(x)f_j(x)\,d^3x.
       \label{f16v67:Hstar}
\end{align}
All integrals defining $H^*$ stay strictly away from coincident time.

\begin{theorem}[Native contact growth with a surviving reflected bank]
\label{f16v67:native-contact}
Under \eqref{f16v67:archive-input}--\eqref{f16v67:mesoscopic},
\begin{equation}
 \boxed{\quad b^{-4}D_M\longrightarrow
                D_*:=\chi_{12}\|h\|_2^2 J\succ0,\qquad
       H_M(n_M)\longrightarrow H^*(u)\quad(n_M/b\to u\ge0).\quad}
 \label{f16v67:native-limits}
\end{equation}
In particular $H^*(0)\succ0$ and
\begin{equation}
 b^4D_M^{-1/2}H_M(0)D_M^{-1/2}
            \longrightarrow D_*^{-1/2}H^*(0)D_*^{-1/2}\succ0.
 \label{f16v67:ratio}
\end{equation}
Thus ordinary-variance normalization makes this entire reflected bank
vanish, although the declared dimension-four normalization has finite,
strictly positive reflected pairings on the same selected native sequence.
\end{theorem}
\begin{proof}
Put $\psi_i(t,x)=h(t)f_i(x)$ and
$c_{i,b}(z)=\psi_i(z/b)$ on $\mathbb Z^4$.  The native comparison error in
any entry of $D_M$ or $H_M(n)$ is at most
\begin{equation}
 C\delta_M\|c_{i,b}\|_1\|c_{j,b}\|_1=O(\delta_M b^8)=o(1).
 \label{f16v67:full-error}
\end{equation}
The time range in the reflected comparison is allowed by the input;
spatial differences are $O(b)=o(M)$.

For the ordinary comparison form, reindex by the lattice displacement $r$:
\[
 b^{-4}\sum_{z,w}c_{i,b}(z)c_{j,b}(w)R(z-w)
 =\sum_rR(r)\left[b^{-4}\sum_z\psi_i(z/b)\psi_j((z-r)/b)\right].
\]
For fixed $r$ the bracket tends to $\int\psi_i\psi_j$ by Riemann sums.
Its absolute value is bounded uniformly in $b\ge1$ and $r$, by
Cauchy--Schwarz and the scaled discrete $\ell^2$ bounds for the fixed
compactly supported tests.  Summability in \eqref{f16v67:chi} therefore
permits dominated convergence.  This proves the first limit, with
$\int\psi_i\psi_j=\|h\|_2^2J_{ij}$.  This argument keeps every microscopic
lag; it does not integrate a positive-time continuum kernel through its
singularity at zero.

For the reflected form, write it as the double temporal Riemann sum of
the archive's canonically spatially smeared covariance.  That comparison
converges uniformly on positive compact time intervals.  Every argument
$s+t+n_M/b$ is at least $2\delta+o(1)$ and remains bounded for bounded $u$.
The temporal Riemann sums consequently converge to
\eqref{f16v67:Hstar}.  Alternatively the uniform positive-time lattice
kernel scaling in f15 V65 gives the same eight-dimensional Riemann sum.
Equation~\eqref{f16v67:full-error} pays the native remainder in either proof.

Here is also a direct proof of the strict matrix assertion.  Put
$\ell_h(E)=\int h(t)e^{-tE}dt>0$.  For $f_c=\sum_i c_i f_i$, Fourier
inversion in \eqref{f16v67:K} gives
\begin{equation}
 c^*H^*(u)c=\frac32\int
   |\widehat f_c(p+q)|^2 w_{12}(p)w_{12}(q)
       |\ell_h(|p|+|q|)|^2 e^{-u(|p|+|q|)}
                \frac{d^3p\,d^3q}{(2\pi)^6}.
 \label{f16v67:positive-spectral}
\end{equation}
It is finite because $\ell_h(E)\le\|h\|_1e^{-\delta E}$ and
$w_{12}(p)\le |p|/2$.  If $c\ne0$, $f_c$ is not zero, so its Fourier
transform is nonzero on a set of positive measure.  For almost every such
sum $p+q$, both weights are positive for almost every $p$.  The integral
is strictly positive.  The same independence makes $J\succ0$.
Continuity of the inverse square root on positive definite matrices now
proves \eqref{f16v67:ratio}.
\end{proof}

The conclusion does not contradict Proposition~\ref{f16v66:finite-energy}.
That proposition fixes normalization on its ordinary Gram.  Theorem
\ref{f16v67:native-contact} changes neither its proof nor its exact
criterion; it excludes treating that particular normalization as mandatory
for all composite sources.  Nor do finite reflected two-point limits
construct a full Euclidean random field with finite ordinary variance.
The divergent ordinary contact term is not declared to be zero.

\subsection{The mesoscopic survivor still fails the proposed physical clock}

Let $P_M$ be the same vacuum-normalized transfer and
$\mathsf L_M=-\log P_M$ in lattice units.  The reflected source synthesis
is obtained by conditional expectation at slice zero.  For example,
its $i$th vector in the ground-state Markov representation is
\[
 u_{M,i}=\sum_{j,x}h(j/b)f_i(x/b)P_M^j\overline X_{M,(0,x)},
 \qquad H_M(n)=(\langle u_{M,i},P_M^nu_{M,j}\rangle)_{ij}.
\]
All $j$ are positive for sufficiently large $M$.

\begin{proposition}[Clock audit and actual source-weight escape]
\label{f16v67:clock}
The selected window \eqref{f16v67:mesoscopic} requires
\begin{equation}
                         b_M=o(\gamma_M^{1/128}).
 \label{f16v67:polynomial-window}
\end{equation}
For any nonzero fixed source combination, its normalized spectral measure
for $b_M\mathsf L_M$ converges to the finite measure in
\eqref{f16v67:positive-spectral}, normalized by its mass at $u=0$.
This measure has no atom at zero.

If an independently specified spacing satisfies $a_Mb_M\to0$, then for
every fixed $E,\tau>0$ that same reflected-normalized combination satisfies
\begin{align}
 \frac{\langle u_M,\mathbf1_{[0,E]}(\mathsf L_M/a_M)u_M\rangle}
      {\|u_M\|^2}&\longrightarrow0,
       \label{f16v67:physical-escape}\\
 \frac{\langle u_M,P_M^{\lceil\tau/a_M\rceil}u_M\rangle}
      {\|u_M\|^2}&\longrightarrow0.
       \label{f16v67:physical-delay}
\end{align}
In particular $a_Mb_M\to0$ for any putative trajectory
$\log(1/a_M)=c\gamma_M+O(\log\gamma_M)$ with $c>0$.
\end{proposition}
\begin{proof}
Since $\delta_M\ge\gamma_M^{-1/16}$, the condition
$\delta_Mb_M^8\to0$ implies \eqref{f16v67:polynomial-window}.
The Laplace transforms of the source measures for $b_M\mathsf L_M$
converge by Theorem~\ref{f16v67:native-contact}, including their total
mass.  For a real rescaled time $u>0$, bracket $P_M^{b_Mu}$ between
$P_M^{\lceil b_Mu\rceil}$ and $P_M^{\lfloor b_Mu\rfloor}$ in each
source quadratic form; both integer-time limits are $H^*(u)$.  Tightness follows from
\[
 \nu_M([L,\infty))\le
       \frac{\nu_M([0,\infty))-\int e^{-tE'}d\nu_M(E')}{1-e^{-tL}},
\]
first taking the limit in $M$, then small $t>0$, then large $L$.
Uniqueness of finite Laplace transforms identifies the limit.  This is the
same spectral argument already used in f15 V65, now applied to the displayed
history source.  The limiting integral is absolutely continuous before
pushforward by $(p,q)\mapsto|p|+|q|$; its zero level set has zero measure.

In the rescaled spectral variable the event in
\eqref{f16v67:physical-escape} is $[0,a_Mb_ME]$.  For any fixed
$\epsilon>0$ it is eventually contained in $[0,\epsilon]$.  Weak convergence,
then $\epsilon\downarrow0$, gives zero because there is no zero atom.
For the delayed pairing split at $\epsilon$: the upper part is at most
$\exp[-\epsilon\tau/(a_Mb_M)]$ times the total normalized mass, while the
lower part tends to zero after the same two limits.  The last assertion
follows from \eqref{f16v67:polynomial-window}.
\end{proof}

Thus the archived nonzero mesoscopic reflected vector cannot be relabelled
as a fixed-physical-scale one on the proposed exponential clock.  This is
an audit of a stated bound and an explicitly retained source, not proof
that the actual covariance error outside the Gaussian window diverges.
Both errors must be avoided: ordinary normalization may erase a valid
\emph{mesoscopic} vector, and reflected normalization does not by itself
change that vector's energy units.

\subsection{A slice-local native construction: smooth the curvature before squaring}

We now return to an arbitrary finite native spatial Wilson regulator,
without using the selected Gaussian comparison.  Spatial periodic boundary
conditions are convenient; the same definitions hold with a specified
spatial boundary.  Let $U_j(x)\in\SU(2)$ be the three spatial links of one
slice.  On real adjoint fields $v(x)\in\operatorname{Im}\mathbb H$ define
\begin{equation}
 (\mathcal L_{a,U}v)(x)=\frac1{a^2}\sum_{j=1}^3
 \left[2v(x)-\operatorname{Ad}_{U_j(x)}v(x+a\hat j)
       -\operatorname{Ad}_{U_j(x-a\hat j)^{-1}}v(x-a\hat j)\right].
 \label{f16v67:covlap}
\end{equation}
The inner product is $a^3\sum_x v(x)\cdot w(x)$.  Choose a declared scalar
normalizer $\zeta_a>0$ and a spatial smoothing radius $r>0$, and put
\begin{align}
 \mathcal F_{a,U}(x)&=\zeta_a a^{-2}\operatorname{Im}U_{12}(x),
 &B_{a,r}(U;x)&=e^{-(r^2/2)\mathcal L_{a,U}}\mathcal F_{a,U}(x),
       \label{f16v67:heat-field}\\
 \mathcal O_{a,r,f}(U)&=\frac{a^3}{2}\sum_x f(x)|B_{a,r}(U;x)|^2,
 &F_{a,r,f}&=\mathcal O_{a,r,f}-\mathbb E_{\pi_a}\mathcal O_{a,r,f}.
       \label{f16v67:heat-source}
\end{align}
Here $f\ge0$ is a nonzero spatial test, sampled on the lattice, and $\pi_a$
is the \emph{actual} Wilson vacuum slice law.  Its expectation has not been
computed by substituting a flat connection.  For a classical linearization
$U_j=e^{g_a a A_j}$ one can take $\zeta_a=g_a^{-1}$; no nonperturbative
renormalization theorem is asserted by that convention.

\begin{theorem}[Exact gauge covariance, path representation and time support]
\label{f16v67:heat-native}
At every finite spatial regulator the following assertions hold.
The operator $\mathcal L_{a,U}$ is self-adjoint and nonnegative.
The field $B_{a,r}$ is gauge covariant at its basepoint, and
$F_{a,r,f}$ is a bounded continuous gauge-invariant slice observable.
Let $s=r^2/2$ and let $p_s^a(x,y)$ be the scalar continuous-time nearest-neighbor
random-walk kernel with rate $a^{-2}$ along each oriented edge.  Then
\begin{equation}
 \boxed{\quad
 \big\|e^{-s\mathcal L_{a,U}}(x,y)\big\|_{\mathrm{op}}
                       \le p_s^a(x,y),\qquad
       |B_{a,r}(U;x)|\le M_a:=\zeta_a a^{-2}.\quad}
 \label{f16v67:diamagnetic}
\end{equation}
There is an absolutely convergent positive-weight path expansion, with
orthogonal parallel transport on each path; the quadratic invariant is
therefore an absolutely convergent sum of closed gauge-invariant
plaquette--path contractions.  At each fixed regulator polynomial path
truncations approximate it uniformly.

A positive-time history made from these sources has exactly its stated
time support.  Its reflected pairing is the pairing of its conditional
predictors on slice zero under the \emph{same} native transfer $P_a$.
No physical time step, transfer eigenvalue, or Wilson measure is changed.
For a nondegenerate lattice with $f$ positive at a site, this single source
is nonconstant, so its native entrance variance is positive; its finite-delay
reflected norm is positive for the injective Wilson transfer.
\end{theorem}
\begin{proof}
A direct summation gives
\[
 \langle v,\mathcal L_{a,U}v\rangle
      =a\sum_{x,j}|v(x)-\operatorname{Ad}_{U_j(x)}v(x+a\hat j)|^2\ge0.
\]
Under a vertex gauge transformation, both the plaquette imaginary part
and the output field transform by its adjoint action at $x$, and
$\mathcal L_{a,U^g}=\mathcal G\mathcal L_{a,U}\mathcal G^{-1}$.
The exponential inherits this identity and the squared norm is invariant.

Let $T_U$ average the six oriented covariant shifts.  Exactly,
\begin{equation}
 e^{-s\mathcal L_{a,U}}
   =e^{-\lambda}\sum_{n=0}^\infty\frac{\lambda^n}{n!}T_U^n,
 \qquad \lambda=6s/a^2.
 \label{f16v67:Poisson}
\end{equation}
Each term is an average over $n$-step paths and the corresponding adjoint
parallel transports.  Every transport is orthogonal.  Taking its operator
norm and summing gives the scalar random-walk kernel, proving
\eqref{f16v67:diamagnetic}.  The sum converges uniformly in the compact link
configuration because $\|T_U\|_{\infty\to\infty}\le1$.
For two paths from $x$ to $y,z$, their cross term is the inner product of
the two transported plaquette curvatures.  Its matrix trace is a closed
plaquette--path network.  Products, sums and uniform convergence preserve
gauge invariance.  A finite path truncation is polynomial in link matrix
entries and their adjoints.

All ingredients use spatial links at the same time; applying this map
slice by slice therefore commutes with time reflection and never extends
support across its plane.  Native reflection positivity is inherited by
this subalgebra.  The Markov proof of
Theorem~\ref{f16v66:reflected} applies to its history sources without a
polynomial restriction.  It uses the old transfer, not $e^{-s\mathcal L}$
as physical evolution.

To check nonconstancy, the observable is zero at the identity configuration.
In an arbitrarily small Abelian perturbation with a nonzero spatial curl,
its leading term is a nonzero squared scalar-heat-smoothed curl.  On a
periodic lattice a nonzero Fourier curl with a suitable phase makes that
term nonzero at any chosen site with $f>0$.  The scalar heat multiplier is
strictly positive at every finite momentum and finite $s$.  Hence the
observable is not identically zero or constant.  The positive native slice
density has full support, giving positive variance.  Injectivity of the
native normalized Wilson transfer then gives positive reflected norm for
any finite buffer.  No uniform lower bound in $a$ is implied.
\end{proof}

The flat-connection first variation can also be read without a continuum
assumption.  If $U_e(\varepsilon)=\exp(\varepsilon v_e)$ and
$(d_{\rm lat}v)_{12}$ is its oriented plaquette sum, then
\begin{equation}
 DB_{a,r}(\mathbf1)[v]=\zeta_a a^{-2}e^{-s\mathcal L_{a,\mathbf1}}
                                  (d_{\rm lat}v)_{12}.
 \label{f16v67:linearization}
\end{equation}
Indeed $\mathcal F_{a,\mathbf1}=0$, so the derivative of the heat operator
multiplies zero; only the derivative of the plaquette remains.  At identity
$\mathcal L_{a,\mathbf1}$ is the scalar lattice Laplacian on each adjoint
component.  This proves the constituent heat multiplier at linear order
without replacing the native distribution by a Gaussian.

Smoothing the final scalar $\operatorname{Im}(U_{12})^2$ is a different
construction.  At a flat connection its Fourier multiplier acts on the
\emph{sum} of the two constituent momenta.  In \eqref{f16v67:heat-source}
each constituent has its own heat factor.  Opposite large momenta can
survive the first procedure but not the second.

\subsection{Spatial tails and a truncation cost that is not hidden in the radius}

For $s>0$, $R>0$, define
\begin{equation}
 I_a(R,s)=\frac{R}{a}\operatorname{arsinh}\!\left(\frac{Ra}{2s}\right)
       -\frac{2s}{a^2}\left[\sqrt{1+\left(\frac{Ra}{2s}\right)^2}-1\right].
 \label{f16v67:rate}
\end{equation}
For an exit estimate, $R$ is the physical $\ell^\infty$ distance from the
source support to the exterior of a chosen spatial box; take the box below
the torus wrapping distance.  Scalar walks can equivalently be lifted
until their first exit.  The retained-link region also includes the one-plaquette
collar needed to evaluate $\mathcal F_{a,U}$ at sites inside the box;
this enlarges spatial support by at most two lattice edges, not in time.

\begin{proposition}[Paid spatial locality and Poisson path truncation]
\label{f16v67:locality}
Replace the covariant heat operator by the one killed when its path leaves
that box, and denote the resulting centered source by $F^{[R]}_{a,r,f}$.
With $\|f\|_{1,a}=a^3\sum_x|f(x)|$,
\begin{equation}
 \|F_{a,r,f}-F^{[R]}_{a,r,f}\|_\infty
    \le2M_a^2\|f\|_{1,a}\min\{1,6e^{-I_a(R,s)}\}.
 \label{f16v67:locality-bound}
\end{equation}
Moreover $I_a(R,s)\to R^2/(4s)$ for fixed $R,s$ as $a\downarrow0$.
If \eqref{f16v67:Poisson} is truncated at path length $m$, its centered
source error is at most
\begin{equation}
 2M_a^2\|f\|_{1,a}\,
              \mathbb P\{\operatorname{Poisson}(6s/a^2)>m\}.
 \label{f16v67:truncate}
\end{equation}
For $m+1>\lambda=6s/a^2$ a valid bound on this last probability is
\[
 \exp[-(m+1)\log((m+1)/\lambda)+(m+1)-\lambda].
\]
\end{proposition}
\begin{proof}
Couple full and killed path expansions.  They agree on paths not exiting
the box.  The field error is therefore at most $M_a$ times the scalar exit
probability; both fields have norm at most $M_a$.
The error in one half of their squared norms is at most $M_a$ times the
field error.  Sum against $a^3f$ and allow a second copy of the bound for
centering under the same native law.

For one coordinate of the lifted scalar walk,
\[
 \exp\{\theta(X_t^j-X_0^j)-2t a^{-2}(\cosh(a\theta)-1)\}
\]
is a martingale.  Its maximal inequality on $0\le t\le s$ gives
\[
 \mathbb P\{\max_{0\le t\le s}(X_t^j-X_0^j)\ge R\}
 \le\exp[-\theta R+2s a^{-2}(\cosh(a\theta)-1)].
\]
Optimize at $a\theta=\operatorname{arsinh}(Ra/(2s))$, and sum over the
six coordinate directions.  Taylor expansion at $Ra/(2s)=0$ gives the
stated limit.  Truncating the Poisson mixture omits a mass equal to its
scalar tail, so the same field and quadratic-source argument proves
\eqref{f16v67:truncate}.  Exponential Markov inequality for the Poisson
variable, optimized at $e^t=(m+1)/\lambda$, gives its displayed tail.
\end{proof}

At fixed physical $r$, the mean path length is $3r^2/a^2$, not a constant.
The uniform spatial error retains $M_a^2=\zeta_a^2a^{-4}$.  Thus a Gaussian
spatial tail alone is not an absolute cutoff-uniform error for a growing
normalizer.  Nor can V65's fixed-degree norm constant be applied unchanged:
the heat source is not a fixed-degree polynomial, and its finite truncations
have increasing path length.  Scalar heat domination is an upper estimate,
not a lower native source-norm or physical-energy estimate.

\subsection{A separate free-curvature benchmark with both norms controlled}

This subsection is not a replacement for the interacting native measure.
Take three independent color copies of a Gaussian spatial magnetic
component with covariance
\begin{equation}
 \mathbb E[B^A(t,x)B^B(s,y)]
   =\delta^{AB}\int e^{ip\cdot(x-y)}w_{12}(p)e^{-|p||t-s|}
                                           \frac{d^3p}{(2\pi)^3}.
 \label{f16v67:free-field}
\end{equation}
It is the linearized curvature component used in the archive.
A momentum cutoff $|p|\le\Lambda$ makes every expression an ordinary
Gaussian random variable.  Apply spatial heat $e^{(r^2/2)\Delta}$ to each
$B^A$ and set $\mathcal E_r=\tfrac12\sum_A:(B_r^A)^2:$, with its exact
Gaussian mean subtracted.  For $f\ge0$, $\int f=1$, and $h\ge0$ with
$\int h=1$ and support in $[\delta,T]$, consider
$S_r=\int h(t)f(x)\mathcal E_r(t,x)\,dt\,d^3x$.

\begin{theorem}[Constituent smoothing gives finite ordinary and positive reflected norms]
\label{f16v67:free-norms}
For fixed $r>0$, all the ordinary and reflected two-point matrices of a
finite test family converge as $\Lambda\to\infty$.  For the specified
single normalized test let $D_{r,\Lambda}=\mathbb E|S_r|^2$ and
$H_{r,\Lambda}=\mathbb E[(\theta S_r)S_r]$.  Put
$M_1=\int |x|f(x)\,d^3x$ and
$\kappa=\min\{r^{-1},(4M_1)^{-1}\}$.  Then for $\Lambda\ge\kappa$,
\begin{equation}
 \boxed{\quad
 H_{r,\Lambda}\ge
   \frac{\kappa^8}{1536\pi^4}e^{-2r^2\kappa^2-4T\kappa}>0,
 \qquad
 D_{r,\Lambda}\le\frac{1}{96\pi^4r^8}.
 \quad}
 \label{f16v67:free-explicit}
\end{equation}
In particular the reflected-to-ordinary ratio is bounded below by
$(\kappa r)^8e^{-2r^2\kappa^2-4T\kappa}/16$, independently of the cutoff.
The reflected spectral measures have uniform Gaussian upper tails in
physical energy.  This gives neither an interacting limit nor a mass gap.
\end{theorem}
\begin{proof}
Wick contraction, with its factor of three colors, gives the slice kernel
\begin{equation}
 C_{r,\Lambda}(t)=\frac32\int_{|p|,|q|\le\Lambda}
  |\widehat f(p+q)|^2w_{12}(p)w_{12}(q)
       e^{-r^2(|p|^2+|q|^2)}e^{-|t|(|p|+|q|)}
                                  \frac{d^3p\,d^3q}{(2\pi)^6}.
 \label{f16v67:free-kernel}
\end{equation}
The reflected expression replaces the last factor by
$|\ell_h(|p|+|q|)|^2$.  The ordinary expression uses
$\int h(t)h(s)e^{-|t-s|(|p|+|q|)}dt\,ds\le1$.
Both are bounded by $C_{r,\infty}(0)$, and Gaussian momentum decay gives
dominated convergence for a finite test family.

For the upper bound use $|\widehat f|\le1$ and
\[
 \int w_{12}(p)e^{-r^2|p|^2}\frac{d^3p}{(2\pi)^3}
                         =\frac1{12\pi^2r^4}.
\]
Multiplying its square by $3/2$ proves the upper constant.  For the lower
bound restrict $|p|,|q|\le\kappa$.  The inequality
$|\widehat f(k)-1|\le M_1|k|$ gives $|\widehat f(p+q)|\ge1/2$.
Also $\ell_h(E)\ge e^{-TE}$ and the heat product is at least
$e^{-2r^2\kappa^2}$.  Finally
\[
 \int_{|p|\le\kappa}w_{12}(p)\frac{d^3p}{(2\pi)^3}
                           =\frac{\kappa^4}{24\pi^2}.
\]
The factor $(3/2)(1/4)$ times its square gives the lower constant in
\eqref{f16v67:free-explicit}.

For $E'=|p|+|q|>E$, use
$|p|^2+|q|^2\ge (E')^2/2$ to split the heat exponent in half.  The
unnormalized reflected tail is at most
\begin{equation}
 \frac{16}{96\pi^4r^8}\,e^{-r^2E^2/4}.
 \label{f16v67:free-tail}
\end{equation}
Divide by the strictly positive lower bound to obtain uniform tightness.
There is positive spectral mass arbitrarily near zero: $\widehat f(0)=1$,
$\ell_h(0)=1$, and both curvature weights are positive almost everywhere
on every small pair of momentum balls.  Thus this benchmark has no
positive spectral threshold.
\end{proof}

A cutoff on $\widehat f(p+q)$ alone would not give the Gaussian majorant:
$p$ and $q$ can be large and nearly opposite.  Spatial smearing of the
\emph{composite} and of its \emph{constituents} are therefore not
interchangeable ultraviolet controls.  The all-connection native
construction above has the latter linearized kernel, but the native
connection inside $\mathcal L_{a,U}$ is random and correlated with the
curvature.  Theorem~\ref{f16v67:free-norms} does not average it out.

\subsection{The remaining comparison is blind to contact terms, not to dynamics}

Let a centered translation-invariant lattice source have connected kernel
$C_a(z)$, and let a declared reference kernel be $C_a^{\rm ref}(z)$.
For sampled positive-time tests $c_{i,a}(z)=\psi_i(az)$, with time support in
$[\delta,T]$, the reflected pairing uses separations whose time coordinate
is at least $2\delta/a$ up to the specified integer endpoints.
There is no factor $a^4$ in $c_{i,a}$ when the source is the same
dimension-four thermal plaquette convention used in
\eqref{f16v67:native-source}.

\begin{proposition}[A contact-blind reflected comparison with its normalization paid]
\label{f16v67:contact-blind}
Suppose, on all pair separations of the declared tests,
\[
 C_a-C_a^{\rm ref}=Q_a+E_a,
\]
where $Q_a$ vanishes whenever the temporal separation is at least the
smallest positive-negative test separation.  Then $Q_a$ contributes
\emph{exactly zero} to every reflected Gram entry, regardless of its size.
If the remaining reflected kernel matrix from the negative support to the
positive support has absolute row and column sums at most $\epsilon_a$,
then for every source coefficient vector $v$,
\begin{equation}
 |v^*(H_a-H_a^{\rm ref})v|
       \le\epsilon_a\Big\|\sum_i v_i c_{i,a}\Big\|_{\ell^2(\mathbb Z^4)}^2
       \le C_{\psi}\epsilon_a a^{-4}|v|^2.
 \label{f16v67:offplane-row}
\end{equation}
Consequently $\epsilon_a=o(a^4)$ suffices to transfer a fixed strictly
positive reference reflected matrix.  Contact terms need not obey that
bound.  An entrywise error $o(a^8)$ on the entire pair-support is a distinct,
stronger sufficient input; no row improvement is inferred from such a
scalar error without additional structure.
\end{proposition}
\begin{proof}
The support condition gives the exact zero assertion before summation.
For the remainder, the Schur row/column test bounds the operator from
negative-support coefficients to positive-support coefficients by
$\epsilon_a$.  Reflection preserves the counting $\ell^2$ norm.  Applying
the operator bound to the reflected and unreflected coefficient arrays
gives the first inequality.  Fixed compact smooth tests have
$\|\sum_i v_i\psi_i(a\cdot)\|_2^2\le C_\psi a^{-4}|v|^2$, by finite-family
Cauchy--Schwarz and Riemann-sum bounds.  The entrywise alternative instead
uses two $\ell^1$ norms, each $O(a^{-4})$, and thus costs $a^{-8}$.
No derivative of an unmarked pressure or cancellation probability appears.
\end{proof}

The decomposition is an estimate of a pairing, not permission to replace
the probability law by a contact-subtracted measure.  The statement also
has the obvious nontranslation-invariant row/column version for a specified
retained law.  It does not estimate $\epsilon_a$ for the interacting theory.
In particular the earlier global $\delta_M$ bound cannot be turned into
an exponential-clock estimate simply by relabelling it a contact error.

The next concrete native bank is now \eqref{f16v67:heat-source}, with $r$
fixed in physical spatial units and temporal test support a fixed positive
distance from the reflection plane.  Its normalization, sampling weights,
and exact path or spatial truncation must be declared before comparison.
The remaining task is to bound its \emph{native positive-separation}
reflected covariance, including the random connection in the heat operator.
Either a direct lower reflected Gram or a controlled comparison in those
pairings would establish noncollapse.  Decay relative to that surviving
Gram, a common dense-core bound, and an interacting continuum construction
remain separate tasks.  Neither contact suppression, a positive finite-cutoff
norm, nor the free benchmark supplies them.

\subsection{Verification, preservation, and the result ledger}

The companion diagnostic is
\begin{center}\small
\path{verify_ym_f16_v67_contact_sources.py}.
\end{center}
It separates exact rational source and contact-support identities,
finite Fourier projection and susceptibility checks, literal quaternion
covariance and adjoint-heat calculations, Poisson and killed-path tests,
and positive free-curvature quadrature.  The latter two are finite
matrix or floating-point diagnostics, not full native Wilson-vacuum
samples, interval certificates, or proof-assistant verification.
The native selected-sequence theorem uses its stated archived analytic
comparison; a numerical free kernel does not verify that input.

The cumulative file is placed in a branch-specific directory so that the
older attached V67 is not overwritten.  Removing this delimited appendix
and reverting the single displayed title change recovers V66 byte for
byte.  Both older parallel V67 copies and both comparison archives remain
unchanged.

\begin{record}[V67 result ledger]
\label{f16v67:ledger}
\emph{Proved here from the stated archived input:} the finite lattice
contact susceptibility; the $b_M^4$ leading ordinary Gram of the selected
native dimension-four history bank; finite strictly positive reflected
Grams for the same normalization and their $b_M^{-4}$ relative ratio;
and the reflected-normalized bank's ultraviolet escape when
$a_Mb_M\to0$, including the incompatibility of the admitted mesoscopic
window with an exponential clock.

\emph{Proved here at a native finite spatial regulator:} a slice-local
adjoint-heat curvature source, gauge covariance and invariance, unchanged
time-reflection support and native transfer, exact path expansion,
scalar heat domination, paid spatial exit and path-length truncation,
and nonconstancy without a uniform cutoff lower bound.

\emph{Proved in an explicitly separate comparison model:} finite ordinary
and positive reflected norms after constituent spatial smoothing, explicit
cutoff-independent constants, and a Gaussian upper spectral tail with no
positive mass threshold.  The contact-blind reflected row estimate is a
general exact comparison tool; its interacting remainder is not estimated.

\emph{Inherited rather than reproved:} f15 V65's selected native
covariance and positive-time kernel comparison; the physical Wilson
transfer and stationary reflection framework; V65--V66's source and
exceptional-set distinctions.  No claim in V66's fixed-normalization
spectral criterion is silently changed.

\emph{Still unproved:} a nonzero native reflected continuum limit for the
new constituent-smoothed bank on the specified physical trajectory;
control of the correlated spatial transport in its source; uniform
source-directed physical-time contraction on a dense core; and the
interacting four-dimensional continuum Yang--Mills mass gap.
\end{record}

\begin{firewall}
The ordinary contact norm is not a required universal physical
normalization.  Conversely, changing to a reflected normalization does
not change the physical clock.  The archive supplies a selected mesoscopic
native input, while the constituent-smoothed explicit continuum calculation
is only a free benchmark.  Their laws and scales are not interchanged.
V1--V66 are preserved unchanged.
\end{firewall}
% END F16 V67 APPEND

% BEGIN F16 V68 APPEND -- 2026-09-18
% Append-only continuation of the Post-Fifteenth-Archive V67.
\clearpage
\section[V68: rough-connection screening and a coupled spatial regulator]{V68: rough-connection screening\\and a coupled spatial regulator}
\label{f16v68:context}

The source in \eqref{f16v67:heat-source} smooths a curvature constituent
with the \emph{unflowed} random connection.  Scalar heat domination gives
an upper bound, not preservation of that constituent.  Before trying to
compare its reflected covariance with the free benchmark, we test this
unpaid point on actual compact Wilson configurations.  Small plaquette
defects alone do not suffice: two noncommuting, spatially constant links
have vanishing plaquette defects but a covariant-Laplacian lower edge
that diverges in physical spatial units.  We prove a volume-independent
lower bound on that edge and calculate the resulting source suppression.
No probability or typicality of these configurations in the Wilson vacuum
is inferred.

A constructive comparison then evolves the spatial connection before
applying the constituent heat.  Spatial Wilson flow is an established
method, not introduced here as a new algorithm.  What is calculated here
is its exact action on the same obstructing configurations, its literal
compact finite-regulator source map, and the costs that remain in a
spatial-locality comparison.  The flow changes observables, not the Wilson
probability measure or the physical-time transfer.  A finite nonzero
\emph{deterministic curvature} after flow is not a nonzero centered vacuum
variance, still less a reflected continuum lower Gram.

\subsection{Scope and the two different spectral variables}

Only V67's definitions of the adjoint Laplacian, the constituent source,
and the fixed-regulator Wilson time-reflection framework are inherited.
No selected Gaussian covariance estimate from f15 is used in the results
below.  This continuation does not repeat V67's free-curvature integral,
path expansion, or contact-blind comparison.  Its new comparison tests the
connection in that path expansion rather than replacing it by identity.

For external context, numerical studies of covariant-Laplacian modes
include J.~Greensite et al., \emph{Localized eigenmodes of covariant
Laplacians in the Yang--Mills vacuum},
\href{https://doi.org/10.1103/PhysRevD.71.114507}{Phys.\ Rev.\ D \textbf{71}
(2005), 114507}, and \emph{Peculiarities in the spectrum of the adjoint
scalar kinetic operator in Yang--Mills theory},
\href{https://arxiv.org/abs/hep-lat/0606008}{arXiv:hep-lat/0606008}.
Their numerical conclusions are not assumed here and do not establish a
statement about our three-dimensional slice operator.  The spatial flow
used below is the spatial-action construction discussed by K.~Sakai and
S.~Sasaki, \href{https://doi.org/10.1103/PhysRevD.107.034510}{Phys.\ Rev.\ D
\textbf{107} (2023), 034510}; its normalization is fixed explicitly below.
No spectroscopy or continuum-finiteness conclusion is imported.

Work on a finite cubic spatial torus of side at least four, with spacing
$a$.  Set $s=r^2/2>0$ and use the inner product
$\langle v,w\rangle_a=a^3\sum_x\overline{v(x)}\cdot w(x)$.
The operator $\mathcal L_{a,U}$ is exactly \eqref{f16v67:covlap}.
Its spectrum has units of inverse \emph{spatial length squared}.
It is not the spectrum of $(-\log P_a)/a$, which is the physical Hamiltonian
in the native transfer construction.  Complexification of the real adjoint
space below is only for the spatial Fourier calculation.

\subsection{Small Wilson defects with a large covariant kinetic edge}

Let $\mathbf i,\mathbf j,\mathbf k$ be the standard imaginary quaternion
basis.  At every spatial site put
\begin{equation}
 U_1=\cos\theta+\mathbf i\sin\theta,\qquad
 U_2=\cos\theta+\mathbf j\sin\theta,\qquad U_3=1,
 \qquad 0<\theta<\pi/2,
 \label{f16v68:background}
\end{equation}
and write $u=\sin^2\theta$, $c=\cos\theta$, $v=\sin\theta$.
This notation $v$ is used only in the following quaternion formula.
A time-independent four-dimensional extension sets every temporal link
to identity.  Its only nontrivial plaquette orientation is $12$.

\begin{proposition}[Literal plaquette and a volume-independent spatial lower edge]
\label{f16v68:edge}
For \eqref{f16v68:background},
\begin{align}
 U_{12}&=1-2u^2+2cv^3\mathbf i-2cv^3\mathbf j+2c^2v^2\mathbf k,
       &1-w(U_{12})&=2u^2, \label{f16v68:plaquette}\\
 \mathcal L_{a,U}&\succeq
             \frac{7u(1-u)}{4a^2}I.
       \label{f16v68:spatial-edge}
\end{align}
The operator bound is independent of the number of spatial sites.
Over continuous Bloch momenta, its leading coefficient is sharp:
\begin{equation}
 \lim_{\theta\downarrow0}\frac{a^2}{\theta^2}
       \inf_{p\in[-\pi,\pi]^3}\lambda_{\min}
          (\mathcal L_{a,U}(p))=\frac74.
 \label{f16v68:sharp-edge}
\end{equation}
The latter assertion is about the continuous symbol; a fixed finite torus
need not contain momenta realizing this infimum.
\end{proposition}
\begin{proof}
Quaternion multiplication gives \eqref{f16v68:plaquette} directly.  The
adjoint transports $R_1,R_2$ are rotations by $2\theta$ about the first
and second coordinate axes.  The spatial Fourier symbol is
\[
 a^2\mathcal L(p)=\sum_{j=1}^3
       (2I-e^{ip_j}R_j-e^{-ip_j}R_j^{\mathsf T}),\qquad R_3=I.
\]
Each summand is positive semidefinite.  For a unit complex vector $z$, put
$t_j=|z_j|^2$, let $C_jz=e_j\mathbin{\times}z$, and put $J_j=-iC_j$.
Then $J_j$ is Hermitian, $J_j^2=I-e_je_j^{\mathsf T}$, and
$m_j=z^*J_jz$ is real.  With $q_j=z^*R_jz$, Rodrigues' formula gives
\[
 |q_j|^2=[1-2u(1-t_j)]^2+4u(1-u)m_j^2.
\]
It follows that
\begin{equation}
 1-|q_j|^2=4u\{(1-u)\operatorname{Var}_z(J_j)
                                      +u\,t_j(1-t_j)\}.
 \label{f16v68:unitary-variance}
\end{equation}
Here $\operatorname{Var}_z(J_j)=z^*J_j^2z-m_j^2$ is a vector-state variance,
not a Wilson-vacuum expectation.

For completeness the needed two-axis uncertainty bound is elementary.
Write $z=b+id$ with real $b,d$.  Up to an irrelevant common sign,
$(m_1,m_2,m_3)=2b\times d$.  If $t=t_3=b_3^2+d_3^2$, then
\[
 m_1^2+m_2^2
 =4\{(b_2d_3-b_3d_2)^2+(b_3d_1-b_1d_3)^2\}
 \le4t(1-t).
\]
Consequently
\begin{equation}
 \operatorname{Var}_z(J_1)+\operatorname{Var}_z(J_2)
 =1+t-m_1^2-m_2^2
 \ge1-3t+4t^2\ge\frac7{16}.
 \label{f16v68:planar-uncertainty}
\end{equation}
For either of the first two symbol terms,
\[
 z^*(2I-e^{ip_j}R_j-e^{-ip_j}R_j^{\mathsf T})z
 \ge2(1-|q_j|)\ge1-|q_j|^2.
\]
Sum these bounds and use \eqref{f16v68:unitary-variance}--
\eqref{f16v68:planar-uncertainty}; the third symbol term is nonnegative.
This proves \eqref{f16v68:spatial-edge} for every momentum and hence for
every finite torus.

For sharpness take $z=\sqrt{3/8}\,e_3+i\sqrt{5/8}\,e_1$.
It attains equality in \eqref{f16v68:planar-uncertainty}.
Choose $p_j=-\arg q_j$ for $j=1,2$ and $p_3=0$.  Expansion of the exact
formula above gives
$2\sum_{j=1}^2(1-|q_j|)=4\theta^2\sum_{j=1}^2
\operatorname{Var}_z(J_j)+O(\theta^4)=7\theta^2/4+O(\theta^4)$.
This Rayleigh upper bound and \eqref{f16v68:spatial-edge} prove
\eqref{f16v68:sharp-edge}.
\end{proof}

The estimate is not confined to an isolated link configuration.

\begin{proposition}[Stability under actual link perturbations]
\label{f16v68:robust}
Suppose in one common gauge that
$\sup_e\|\operatorname{Ad}_{U'_e}-\operatorname{Ad}_{U_e}\|_{\mathrm{op}}
\le\epsilon$.  Then
\begin{equation}
 \mathcal L_{a,U'}\succeq a^{-2}
 \left(\sqrt{7u(1-u)/4}-\sqrt3\,\epsilon\right)_+^2 I.
 \label{f16v68:robust-edge}
\end{equation}
In particular $\epsilon\le\sqrt{7u(1-u)}/(4\sqrt3)$ retains the lower
edge $7u(1-u)/(16a^2)$.  These conditions define nonempty open neighborhoods
at each fixed regulator.  No lower probability for them is claimed.
\end{proposition}
\begin{proof}
Use the covariant gradient on positively oriented spatial edges:
\[
 \|d_Uz\|^2=\langle z,\mathcal L_{a,U}z\rangle_a.
\]
Each site is the endpoint of three such edges, so
$\|(d_{U'}-d_U)z\|\le\sqrt3\epsilon\|z\|_a/a$.
The reverse triangle inequality and \eqref{f16v68:spatial-edge} give the
assertion.  Simultaneous gauge transformation of both configurations
preserves all norms in this argument.
\end{proof}

\subsection{The original curvature source is screened on this configuration}

\begin{theorem}[Exact frozen-connection heat loss at vanishing plaquette defect]
\label{f16v68:heat-loss}
For the original V67 field
$B_s=e^{-s\mathcal L_{a,U}}\zeta_a a^{-2}\operatorname{Im}U_{12}$
on \eqref{f16v68:background}, it is spatially constant and
\begin{equation}
 |B_s|^2=4\zeta_a^2a^{-4}u^2(1-u)
 \left\{2u\,e^{-8su/a^2}+(1-u)e^{-16su/a^2}\right\}.
 \label{f16v68:exact-screening}
\end{equation}
If $u_a\to0$, $u_a/a^2\to\infty$, and
\begin{equation}
 \log_+(\zeta_a a^{-2})=o(u_a/a^2),
 \label{f16v68:normalizer-condition}
\end{equation}
then $B_s\to0$ for every fixed $s>0$, even though the Wilson plaquette
defect tends to zero.  For $\zeta_a=1$, the unheated curvature has norm
asymptotic to $2u_a/a^2$ and diverges.

One explicit choice is
\begin{equation}
 u_a=\gamma_a^{-1/2},\qquad
 \log(1/a)=c\gamma_a+O(\log\gamma_a),\quad c>0.
 \label{f16v68:thermal-test}
\end{equation}
It has $\gamma_a(1-w(U_{12}))=2$, while its spatial lower edge diverges
like at least $7/(4a^2\sqrt{\gamma_a})$.
Any polynomial $\zeta_a$ in $\gamma_a$ satisfies
\eqref{f16v68:normalizer-condition} on this test sequence.
\end{theorem}
\begin{proof}
On constant adjoint fields,
\[
 \mathcal L_{a,U}(0)=a^{-2}\operatorname{diag}(4u,4u,8u).
\]
Apply its exponential to the three imaginary components in
\eqref{f16v68:plaquette}.  Squaring and summing gives
\eqref{f16v68:exact-screening}; at $s=0$ it equals
$4\zeta_a^2a^{-4}u^2(1-u^2)$.
The exponential factors beat \eqref{f16v68:normalizer-condition}, giving
the limit.  The last assertions follow by substitution into
\eqref{f16v68:plaquette} and \eqref{f16v68:spatial-edge}.
\end{proof}

Thus a small plaquette-defect estimate, even at the displayed thermal
order, cannot \emph{by itself} justify replacing the heat operator in
V67 by its flat-connection multiplier.  The new estimate concerns an
auxiliary spatial operator, not a physical mass gap.  The backgrounds
have not been sampled from the native vacuum, and their extensive
probability is not estimated by their local action density.
The theorem does not show that V67's source collapses in that vacuum.

There is a precise native implication, should an appropriate spatial
spectral event be established.  It specifies what is still missing rather
than assuming a typical lower edge.

\begin{proposition}[A source-collapse certificate with its native probability unpaid]
\label{f16v68:native-certificate}
On a finite torus of physical volume $\mathcal V_a$, use the actual
Wilson slice law $\pi_a$.  For $f\ge0$, set
\[
 O_s=\frac12a^3\sum_x f(x)|B_s(x)|^2.
\]
If $\pi_a\{\lambda_{\min}(\mathcal L_{a,U})<m_a^2\}\le p_a$, then
\begin{equation}
 \operatorname{Var}_{\pi_a}(O_s)
 \le \frac14\|f\|_\infty^2\mathcal V_a^2\zeta_a^4a^{-8}
           \left(e^{-4sm_a^2}+p_a\right).
 \label{f16v68:native-collapse}
\end{equation}
The ordinary norm, and hence the reflected norm, of any probability-weighted
stationary temporal average of its centered source obeys the same bound.
No such $p_a$ is proved here.  On growing volume, the lowest eigenvalue
may in particular be affected by spatially remote regions.
\end{proposition}
\begin{proof}
On the specified spectral event,
$\|B_s\|_a^2\le e^{-2sm_a^2}\|\mathcal F_{a,U}\|_a^2
\le e^{-2sm_a^2}\mathcal V_a\zeta_a^2a^{-4}$.
Off that event use heat contraction without the exponential.
Bound $O_s$ by $\|f\|_\infty\|B_s\|_a^2/2$, square and integrate.
Variance is at most the second moment.  Jensen and stationarity control
a probability-weighted temporal average, and Cauchy--Schwarz with reflection
invariance bounds its reflected norm by its ordinary norm.
\end{proof}

\subsection{A coupled spatial regulator that keeps the original probability law}

Instead of freezing the rough connection, define its spatial evolution
with an explicit metric and normalization.  For a spatial link
configuration $W$ let
\begin{equation}
 E_{\rm sp}(W)=\sum_x\sum_{1\le i<j\le3}(1-w(W_{ij}(x))),
 \qquad
 \partial_\tau W_e\,W_e^{-1}
                  =-a^{-2}\nabla_e E_{\rm sp}(W),\quad W(0)=U.
 \label{f16v68:flow}
\end{equation}
The right-trivialized gradient uses the unit-quaternion metric
$\langle \xi,\eta\rangle=-\frac12\Tr(\xi\eta)$ on imaginary quaternions.
Equivalently it is the round metric on each $S^3$.
The factor $\gamma$ of the statistical Wilson action is not included in
$E_{\rm sp}$; the chosen flow time $\tau$ has units of spatial length squared.
Write $W=\Phi_{a,\tau}(U)$ and define, with $s\ge0$,
\begin{align}
 \mathcal B_{a,\tau,s}(U)
   &=e^{-s\mathcal L_{a,W}}\zeta_a a^{-2}\operatorname{Im}W_{12},
          \label{f16v68:coupled-field}\\
 O_{a,\tau,s,f}(U)
   &=\frac{a^3}{2}\sum_x f(x)|\mathcal B_{a,\tau,s}(U;x)|^2,
 & F_{a,\tau,s,f}&=O_{a,\tau,s,f}-\mathbb E_{\pi_a}O_{a,\tau,s,f}.
          \label{f16v68:coupled-source}
\end{align}
At $\tau=0$ this is precisely V67's source.  At $s=0$ it uses the flowed
curvature without an additional frozen-connection heat step.

\begin{theorem}[Finite-regulator flow and exact slice support]
\label{f16v68:flow-native}
For every finite spatial regulator, \eqref{f16v68:flow} exists for all real
$\tau$, is a smooth diffeomorphism at each finite $\tau$, and is gauge
equivariant.  For $\tau\ge0$,
\begin{equation}
 \frac{d}{d\tau}E_{\rm sp}(\Phi_{a,\tau}U)
          =-a^{-2}\sum_e|\nabla_eE_{\rm sp}|^2\le0.
 \label{f16v68:energy-law}
\end{equation}
The source \eqref{f16v68:coupled-source} is bounded, continuous, gauge
invariant and uses only the original spatial slice.  Its history sources
retain their literal time-reflection support and the same physical transfer
$P_a$.  It is nonconstant for a nondegenerate lattice and a nonzero
nonnegative sampled test.  Therefore its finite-regulator entrance variance
and its finite-delay reflected norm are positive.  No cutoff-uniform bound
or spatial-locality limit is asserted by those strict inequalities.
\end{theorem}
\begin{proof}
The negative gradient is a smooth tangent vector field on a finite compact
product of groups.  Boundedness in every coordinate and local uniqueness
extend solutions to all real times.  Uniqueness gives
$\Phi_{a,-\tau}\Phi_{a,\tau}=I$ and smooth dependence gives the
diffeomorphism.  Gauge transformations act isometrically and leave
$E_{\rm sp}$ invariant, so they commute with its flow.  Differentiating
$E_{\rm sp}$ along that vector field gives \eqref{f16v68:energy-law}.
V67's all-connection covariance and boundedness apply to the output links
$W$.  No temporal link or neighboring slice occurs in \eqref{f16v68:flow}.
Reflection positivity is therefore inherited on the resulting
positive-time subalgebra under the original law.

The uncentered heat-curvature observable as a function of $W$ is nonconstant
by V67's finite-regulator argument, also when $s=0$.
Composition with the surjective map $\Phi_{a,\tau}$ remains nonconstant.
The original positive slice density has full support, giving positive
variance; injectivity of $P_a$ gives the finite-delay reflected assertion.
\end{proof}

All expectations in \eqref{f16v68:coupled-source} remain under $\pi_a(U)$.
Equivalently one could use the pushforward measure
$(\Phi_{a,\tau})_*\pi_a$ on $W$, with its exact change-of-variables density.
It is \emph{not} a Wilson density at a changed coupling.  The product flow
is not Haar-volume preserving in general, and its Jacobian must not be
discarded.  Invertibility at fixed $a$ also does not control the inverse
map uniformly as $a\downarrow0$.

\subsection{The obstructing background has an exactly solvable spatial flow}

\begin{theorem}[Connection evolution replaces exponential screening on the test family]
\label{f16v68:solved-flow}
Starting from \eqref{f16v68:background}, the spatial flow preserves its
form, with a common angle $\theta(\tau)$ and $u(\tau)=\sin^2\theta(\tau)$.
Exactly,
\begin{align}
 \theta'&=-4a^{-2}\sin^3\theta\cos\theta,
 &u'&=-8a^{-2}u^2(1-u), \label{f16v68:angle-ode}\\
 T(u(\tau))&=T(u_0)+8\tau/a^2,
 &T(u)&=u^{-1}+\log((1-u)/u), \label{f16v68:implicit-flow}\\
 \frac1{u_0^{-1}+8\tau/a^2}
 &\le u(\tau)\le
 \frac1{u_0^{-1}+8(1-u_0)\tau/a^2}. &&\label{f16v68:flow-bracket}
\end{align}
If $u_0\to0$ and $u_0/a^2\to\infty$, then for every fixed $\tau>0$,
\begin{equation}
 \frac{u(\tau)}{a^2}\longrightarrow\frac1{8\tau},
 \qquad
 a^{-2}\operatorname{Im}W_{12}\longrightarrow\frac{\mathbf k}{4\tau}.
 \label{f16v68:flow-curvature-limit}
\end{equation}
For $\zeta_a=1$ and fixed $s\ge0$, this gives
\begin{equation}
 \boxed{\quad
 \mathcal B_{a,\tau,s}\longrightarrow
             \frac{e^{-s/\tau}}{4\tau}\mathbf k,\qquad
 O_{a,\tau,s,f}\longrightarrow
             \frac{e^{-2s/\tau}}{32\tau^2}\int f(x)\,d^3x.
 \quad}
 \label{f16v68:coupled-limit}
\end{equation}
For general $\zeta_a$ the field divided by $\zeta_a$, and the uncentered
observable divided by $\zeta_a^2$, have these limits.  The spatial torus or
test sampling is chosen so the indicated Riemann sum converges.
\end{theorem}
\begin{proof}
Use quaternion coordinates for the spatial flow.  If $S_e$ is the sum of
its four complementary spatial staple vectors, then
\begin{equation}
 \partial_\tau x_e=a^{-2}\{S_e-(x_e\cdot S_e)x_e\}.
 \label{f16v68:ambient-flow}
\end{equation}
For constant links $X=\cos\alpha+\mathbf i\sin\alpha$ and
$Y=\cos\beta+\mathbf j\sin\beta$, the first-link staple is
$YXY^{-1}+Y^{-1}XY+2X$.  Its scalar part is $4\cos\alpha$ and its only
imaginary component is
$(2+2\cos2\beta)\sin\alpha\,\mathbf i$.
Projection onto the unit tangent
$-\sin\alpha+\mathbf i\cos\alpha$ gives
$\alpha'=-4a^{-2}\sin\alpha\cos\alpha\sin^2\beta$.
The second equation is obtained by exchanging $\alpha,\beta$; the third
link has zero tangent.  Translation symmetry and uniqueness make these
the actual full-lattice flow, not a separately constrained minimization.
If $\alpha=\beta$ initially, uniqueness preserves equality and gives
\eqref{f16v68:angle-ode}.

Direct differentiation yields $T'(u)=-1/[u^2(1-u)]$, proving
\eqref{f16v68:implicit-flow}.  Since $u$ decreases,
$8(1-u_0)/a^2\le(1/u)'\le8/a^2$; integration gives
\eqref{f16v68:flow-bracket}.  Its two endpoints imply the first limit in
\eqref{f16v68:flow-curvature-limit}.
The $\mathbf i,\mathbf j$ components of the rescaled plaquette are
$O(u^{3/2}/a^2)=O(a)$, whereas its $\mathbf k$ component is
$2u(1-u)/a^2\to1/(4\tau)$.

The constant-field Laplacian of the evolved connection is
\[
 \mathcal L_{a,W}(0)=a^{-2}\operatorname{diag}(4u(\tau),4u(\tau),8u(\tau))
 \longrightarrow \operatorname{diag}((2\tau)^{-1},(2\tau)^{-1},\tau^{-1}).
\]
Applying the heat multiplier gives \eqref{f16v68:coupled-limit} and the
spatial Riemann sum gives its second statement.
\end{proof}

The comparison uses exactly the same initial links as
Theorem~\ref{f16v68:heat-loss}.  With a frozen connection the field tends
to zero; with a fixed positive spatial preflow time its rescaled curvature
has the finite limit above.  This resolves the deterministic test, not the
native source-norm problem.  In particular subtracting a vacuum mean is
not evaluated by a single deterministic configuration.  A nonzero value
in \eqref{f16v68:coupled-limit} is not proof that the centered random source
has positive limiting variance.

\subsection{The same linear benchmark, and an explicit nonlinear locality cost}

\begin{proposition}[First variation at a flat connection]
\label{f16v68:linearization}
Let $U_e(\epsilon)=\exp(\epsilon v_e)$ and let $d_{\rm lat}$ be the
oriented lattice curl.  At the identity connection,
\begin{equation}
 D\mathcal B_{a,\tau,s}(\mathbf1)[v]
   =\zeta_a a^{-2}e^{-(s+\tau)\mathcal L_{a,\mathbf1}}
                                     (d_{\rm lat}v)_{12}.
 \label{f16v68:flat-linear}
\end{equation}
Consequently its free linearized constituent multiplier is the V67 one
with $r^2$ replaced by $2(s+\tau)$.  This is an equality of first
variations, not an equality of interacting distributions.
\end{proposition}
\begin{proof}
The quadratic expansion of $E_{\rm sp}$ is
$\frac12\|d_{\rm lat}v\|_{\ell^2}^2$.  Its linearized flow is
$v'=-a^{-2}d_{\rm lat}^*d_{\rm lat}v$.
The linearized curvature is closed: $d_{\rm lat}(d_{\rm lat}v)=0$.
On the cubic torus the discrete Hodge Laplacian on these cochains is the
componentwise scalar lattice Laplacian, with multiplier
$\sum_j4a^{-2}\sin^2(p_j/2)$.
Thus the curvature evolves by $e^{-\tau\mathcal L_{a,\mathbf1}}$.
The derivative of the final heat operator multiplies zero curvature at
identity, exactly as in V67.  Only the differentiated curvature remains,
and the two scalar heat factors commute, proving the formula.
\end{proof}

Using a physical flow parameter does not by itself prove cutoff-uniform
spatial locality.  The following elementary all-configuration bound
records the cost rather than assigning the linearized Gaussian radius
to the nonlinear native observable.

\begin{proposition}[An all-configuration influence bound with its cutoff cost]
\label{f16v68:locality}
Connect two positively stored spatial links when they belong to one
spatial plaquette.  Let $N_{ef}$, $e\ne f$, count the plaquettes containing
both; $N_{ee}=0$ and every row has sum twelve.
For two initial configurations let
$d_e(\tau)=|\Phi_{a,\tau}(U)_e-\Phi_{a,\tau}(U')_e|$ in quaternion norm.
Entrywise,
\begin{equation}
 d(\tau)\le e^{\tau(8I+N)/a^2}d(0).
 \label{f16v68:influence-matrix}
\end{equation}
If the configurations agree on every link at graph distance less than
$d$ from $e$, then
\begin{equation}
 d_e(\tau)\le
 \min\left\{2,\,
    2e^{20\tau/a^2}\,
        \mathbb P\{\operatorname{Poisson}(12\tau/a^2)\ge d\}\right\}.
 \label{f16v68:influence-tail}
\end{equation}
The constants are volume independent, but this bound does not yield a
fixed-physical-distance small tail as $a\downarrow0$ at fixed $\tau>0$.
\end{proposition}
\begin{proof}
In \eqref{f16v68:ambient-flow}, each $S_e$ is a sum of four quaternion
products of three other links, each factor with norm one.  Extend the
same polynomial expression to factors of norm at most one.
Its derivative with respect to $x_e$ has norm at most eight:
$S_e$ is independent of $x_e$ and differentiation of
$-(x_e\cdot S_e)x_e$ has norm at most $2|S_e|$.
Differentiation with respect to any other link costs at most one per
shared plaquette, since quaternion multiplication is norm multiplicative
and $I-x_ex_e^{\mathsf T}$ has norm at most one.
Straight interpolation lies inside these unit balls, proving the
componentwise Lipschitz bound with matrix $(8I+N)/a^2$.
There are four incident spatial plaquettes, each with three other links,
so the row sum of $N$ is twelve.  The integral inequality for the distances
and iteration of its nonnegative matrix kernel prove
\eqref{f16v68:influence-matrix}.

For the distinguished row, $N^k$ cannot reach an initially different link
in fewer than $d$ steps.  Also $d_f(0)\le2$ and $N^k\mathbf1=12^k\mathbf1$.
The exponential series therefore gives
$2e^{8\tau/a^2}\sum_{k\ge d}(12\tau/a^2)^k/k!$, which is
\eqref{f16v68:influence-tail}.  Quaternion distance is at most two.
At a fixed physical separation the available graph distance is only of
order $a^{-1}$, while the Poisson mean is of order $a^{-2}$.
The displayed sufficient bound is then generally trivial; no assertion
about failure of the actual spatial locality follows.
\end{proof}

At every fixed regulator the source remains strictly local in \emph{time}.
Its spatial tails and its increasingly long path or flow dependence need
a separate estimate under the native law.  In particular neither V65's
fixed-degree source bound nor V67's original-connection killed-path bound
can be applied unchanged to the preflow map.  The all-configuration
factor $e^{20\tau/a^2}$ above is not paid by simply quoting an exceptional
probability of order $e^{-c\gamma_a}$ on an exponential physical trajectory.
This statement is a limitation of those two displayed bounds, not a lower
bound on the true average sensitivity.

\subsection{The noncircular continuation and result ledger}

The new obstruction is upstream of a reflected-covariance comparison:
the spatial regulator itself can remove the field in small-defect
configurations.  The explicit preflow construction passes that test and
keeps the original time support and law.  The next native calculation is
therefore the joint correlation of
$F_{a,\tau,s,f}(U_t)$ and its reflected copy under the actual stationary
Wilson history, for declared $\tau>0$, $s\ge0$, $\zeta_a$, and spatial tests.
Its connection is $\Phi_{a,\tau}(U_t)$, correlated with its flowed
curvature.  Replacing that connection by identity or replacing its
pushforward law by a Wilson density at another coupling is not authorized.

A useful advance must estimate these joint positive-separation pairings
or show an actual native failure for this bank.  Finite-regulator
nonconstancy, the deterministic test limit, and the first variation do not
supply a cutoff-uniform reflected lower bound.  Noncollapse alone would
still leave physical-time decay and common dense-core constants unpaid.
No additional coherent-cell Hessian or scalar heat upper estimate should
be counted as closing those steps.

The companion diagnostic is
\begin{center}\small
\path{verify_ym_f16_v68_connection_regulator.py}.
\end{center}
It tests exact quaternion and uncertainty identities, Fourier symbols,
perturbed covariant quadratic forms, the full spatial staple force,
gauge covariance, the implicit homogeneous flow and its brackets,
linearized curvature evolution, and finite-flow influence bounds.
Floating-point ODE and matrix checks are distinguished from exact algebra.
They do not sample the Wilson vacuum, certify typical spatial eigenvalues,
provide outward-rounded interval bounds, or constitute proof-assistant
verification.  The analytic arguments do not depend on their numerical
outcomes.  Removing this delimited appendix and reversing the single
displayed version change recovers the current lineage's V67 byte for byte.
The older parallel V67 attachments remain unchanged.

\begin{record}[V68 result ledger]
\label{f16v68:ledger}
\emph{Proved here:} a literal noncommuting compact-link family with
vanishing Wilson defects; the volume-independent $7u(1-u)/(4a^2)$
covariant-Laplacian bound, its sharp small-angle coefficient and stable
neighborhood; exact screening of V67's frozen-connection heat source;
a probability-explicit native collapse certificate without an estimate
of its spectral event; an exactly defined spatial-flow/constituent source;
its fixed-regulator gauge covariance, time support and unchanged-law
interpretation; the exact homogeneous nonlinear flow, its curvature and
coupled-heat limits; the combined flat first variation; and a nonlinear
influence bound with its cutoff dependence retained.

\emph{Inherited or established background:} V67's fixed-connection
adjoint heat construction and source distinctions; native Wilson transfer
positivity and reflection; the spatial Wilson-flow method itself.
No Gaussian selected-sequence estimate or continuum renormalization
theorem is used to evaluate the interacting law.

\emph{Not asserted:} typicality of the test configurations; a spatial
spectral gap in the native vacuum; collapse of V67's bank under that vacuum;
a physical Hamiltonian gap inferred from a spatial Laplacian; equality
of a flowed-link marginal with a new Wilson law; nonlinear Gaussian
spatial tails; or a reflected variance inferred from a deterministic
curvature limit.

\emph{Still unproved:} a cutoff-uniform, nonzero native reflected Gram for
the specified coupled regulator; its spatial tails and normalization under
the correlated native law; source-directed physical-time contraction with
common dense-core constants; and the interacting four-dimensional
continuum Yang--Mills mass gap.
\end{record}

\begin{firewall}
A small plaquette defect does not make the connection harmless inside a
long spatial heat evolution.  Evolving that connection resolves the stated
deterministic test, not the native continuum problem.  Flow time is not
physical time; the probability law, induced correlations, and physical
transfer remain explicit.  V1--V67 are preserved unchanged.
\end{firewall}
% END F16 V68 APPEND

% BEGIN F16 V69 APPEND -- 2026-09-19
% Append-only continuation of the Post-Fifteenth-Archive V68.
\clearpage
\section[V69: signed spatial-flow sensitivity and native replica localization]{V69: signed spatial-flow sensitivity\\and native replica localization}
\label{f16v69:context}

V68 constructs a spatially flowed source but leaves both its native
reflected covariance and its spatial approximation unpaid.  Its
all-configuration influence bound takes absolute values of the compact
force before exploiting the gradient structure.  This loses a positive
curl form, even at an exactly flat connection.  The present continuation
recovers that signed form at \emph{every} configuration, including center
plaquettes.  It proves an off-diagonal bound for the actual nonlinear
flow differential with a diffusive spatial rate and an explicit integrated
negative-curvature cost.  An exactly growing mode on V68's own background
shows why that cost cannot simply be set to zero.

We then use the bound in a conditional-replica formula under the actual
native slice law.  This gives a source-marked matrix controlling spatial
localization and its effect on reflected pairings.  The matrix is an
explicit expectation, not a proved cutoff-uniform moment.  Thus this
iteration improves and identifies the remaining approximation problem;
it does not report a new native positive-separation lower Gram.

\subsection{Scope, geometry, and the unchanged source}

The spatial torus, unit-quaternion metric, action $E_{\rm sp}$, and flow
$\Phi_{a,t}$ are exactly those of \eqref{f16v68:flow}.  The physical-time
transfer remains $P_a$.  Tangent vectors have inner product
$a^3\sum_e v_e\cdot z_e$ on positively stored spatial links.  The Hessian
endomorphism used throughout is defined with the unscaled product metric.
Using $a^3$-weighted norms on both its input and output changes neither
its operator norm nor its self-adjointness.  We do not redefine the flow
using a differently normalized gradient; its variational equation retains
the factor $a^{-2}$.

This is a signed energy estimate, not a new flow algorithm.  For context,
spatial flow is described by K.~Sakai and S.~Sasaki,
\href{https://arxiv.org/abs/2211.15176}{arXiv:2211.15176}, and exponential
weighting for graph heat bounds is part of the established
Davies--Gaffney method; see F.~Bauer, B.~Hua and S.-T.~Yau,
\href{https://arxiv.org/abs/1402.3457}{arXiv:1402.3457}.
We prove the needed time-dependent compact-link estimates below.
Neither a scalar heat-kernel theorem nor a numerical spectroscopy result
is imported as a theorem about the Wilson vacuum.

For the native localization theorem we choose the already permitted case
$s=0$ of \eqref{f16v68:coupled-source}:
\begin{equation}
 O_i(W)=\frac{\zeta_a^2}{2a}\sum_x f_i(x)|\operatorname{Im}W_{12}(x)|^2,
 \qquad F_i(U)=O_i(\Phi_{a,\tau}U)-\mathbb E_{\pi_a}O_i(\Phi_{a,\tau}U).
 \label{f16v69:source}
\end{equation}
The finite family of real sampled tests has a common finite plaquette
support.  No further frozen-connection heat is needed to define this
source.  A separate derivative estimate for the optional $s>0$ step is
given below; its spatial support is not silently identified with that of
\eqref{f16v69:source}.  All native expectations use the original slice law
$\pi_a$, not a Wilson law at a changed coupling.

\subsection{The exact plaquette Hessian and its signed lower bound}

Fix one oriented plaquette with product $P=w+p$, where $w\in\mathbb R$,
$p\in\operatorname{Im}\mathbb H$, and $w^2+|p|^2=1$.  Vary each stored
link along its product-metric geodesic.  On a reversed occurrence the
velocity is reversed and conjugated.  Transport the resulting four
velocities to the plaquette basepoint, in product order, and call them
$a_1,\ldots,a_4\in\operatorname{Im}\mathbb H$.  Each transport is an
isometry, with its orientation sign retained.  Exactly along these curves,
\begin{equation}
 P(\epsilon)=e^{\epsilon a_1}e^{\epsilon a_2}
              e^{\epsilon a_3}e^{\epsilon a_4}P.
 \label{f16v69:plaquette-curve}
\end{equation}
The $a_j$ in this formula are computed at the reference configuration;
no derivative of a moving transport is omitted.

\begin{theorem}[Signed compact Hessian, with the center sector retained]
\label{f16v69:Hessian}
For $E_P=1-w(P)$,
\begin{equation}
 \boxed{\quad
 \operatorname{Hess}E_P[v,v]
 =w\left|\sum_{j=1}^4a_j\right|^2
       +2p\cdot\sum_{i<j}a_i\times a_j.\quad}
 \label{f16v69:exact-Hessian}
\end{equation}
Define the nonnegative local cost
\begin{equation}
 c(P)=3|\operatorname{Im}P|+4(-w(P))_+,
 \quad \nu_e(W)=\sum_{P\ni e}c(P),
 \quad \nu(W)=\max_e\nu_e(W).
 \label{f16v69:curvature-cost}
\end{equation}
For the full spatial Hessian $\mathsf H_W$,
\begin{equation}
 \mathsf H_W\succeq-\operatorname{diag}(\nu_e(W)I_3)
                \succeq-\nu(W)I,\qquad 0\le\nu(W)\le20.
 \label{f16v69:H-lower}
\end{equation}
If all spatial plaquettes have nonnegative scalar part, then
$\nu(W)\le12\max_P|\operatorname{Im}P|$.  For distinct stored links,
\begin{equation}
 \| (\mathsf H_W)_{ef}\|_{\rm op}\le N_{ef},
 \qquad \sum_{f\ne e}N_{ef}=12,
 \label{f16v69:H-offdiag}
\end{equation}
where $N$ is the exact shared-plaquette incidence matrix of V68.
At a flat connection the Hessian is the nonnegative covariant curl Gram.
\end{theorem}
\begin{proof}
Differentiating \eqref{f16v69:plaquette-curve} twice, and using
$a_i^2=-|a_i|^2$, gives
\[
 P''(0)=\left(-\left|\sum_ia_i\right|^2
                      +\sum_{i<j}[a_i,a_j]\right)P.
\]
For imaginary $b$, $w(bP)=-b\cdot p$, and
$[a_i,a_j]=2a_i\times a_j$.  Taking the negative scalar part proves
\eqref{f16v69:exact-Hessian}.  These are product geodesics, so the
computed second derivative is the Riemannian Hessian, without an
acceleration correction.

The signed curl term is nonnegative when $w\ge0$.  In general,
$w|\sum_i a_i|^2\ge-4(-w)_+\sum_i|a_i|^2$.  Also
\[
 2\left|p\cdot\sum_{i<j}a_i\times a_j\right|
 \le2|p|\sum_{i<j}|a_i||a_j|
 \le3|p|\sum_i|a_i|^2.
\]
Summing over plaquettes proves the first lower bound.  Every spatial
link has four incident plaquettes.  If $w<0$, Cauchy--Schwarz gives
$3\sqrt{1-w^2}+4|w|\le5$; if $w\ge0$, $c(P)\le3$.
This proves the stated constants.

In transported coordinates the diagonal $3\times3$ block of one
plaquette is $wI$.  For $i<j$ its off-diagonal block is
$wI-C_p$, where $C_p z=p\times z$; the opposite block is its transpose.
Since $C_p^2=pp^{\mathsf T}-|p|^2I$, its norm is at most
$\sqrt{w^2+|p|^2}=1$.  Pullback by the oriented isometries preserves
this bound.  Addition over the actual shared plaquettes gives
\eqref{f16v69:H-offdiag}.  At a flat connection, $p=0$, $w=1$, leaving
exactly the sum of squared covariant curls.
\end{proof}

The center term in \eqref{f16v69:curvature-cost} is necessary.
At $P=-I$ the imaginary curvature vanishes but
$\operatorname{Hess}E_P[v,v]=-|\sum_i a_i|^2$.
Taking all $a_i$ equal attains the lower coefficient $-4$ relative to
$\sum_i|a_i|^2$.  An all-configuration estimate using only
$|\operatorname{Im}P|$ would be false.  Conversely, no curvature penalty
is needed at $P=I$; taking absolute values of all force entries there
loses a genuine positive form.

\subsection{A diffusive off-diagonal estimate for the nonlinear flow differential}

Along an actual trajectory $W(t)=\Phi_{a,t}(U)$, let
$\mathsf J(t,t_0)$ be the differential of the flow from $t_0$ to $t$.
Use parallel orthonormal tangent frames separately on each link.
The exact variational equation is
\begin{equation}
 \partial_t\mathsf J=-a^{-2}\mathsf H_{W(t)}\mathsf J,
 \qquad \mathsf J(t_0,t_0)=I.
 \label{f16v69:Jacobi}
\end{equation}
Set
\begin{equation}
 \mathfrak A(t,t_0;U)=a^{-2}\int_{t_0}^t\nu(W(v))\,dv.
 \label{f16v69:amplification}
\end{equation}
Parallel frames do not couple different links and hence preserve the
block support and every scalar spatial weight used next.

\begin{theorem}[Signed weighted flow and spatial separation]
\label{f16v69:weighted-flow}
Let $\psi$ be a real function of link midpoints such that
$|\psi_e-\psi_f|\le\lambda a$ whenever $N_{ef}>0$.
For $h=t-t_0>0$,
\begin{equation}
 \boxed{\quad
 \|e^\psi\mathsf J(t,t_0)e^{-\psi}\|_{2\to2}
 \le\exp\left\{\mathfrak A(t,t_0;U)
               +\frac{12h}{a^2}(\cosh(\lambda a)-1)\right\}.\quad}
 \label{f16v69:weighted-J}
\end{equation}
For link sets $A,B$ separated by physical midpoint distance at least $R$,
\begin{align}
 \|\mathbf1_A\mathsf J(t,t_0)\mathbf1_B\|_{2\to2}
 &\le e^{\mathfrak A(t,t_0;U)-\mathcal I_a(R,h)},
 \label{f16v69:offdiag}\\
 \mathcal I_a(R,h)
 &=\frac Ra\operatorname{arsinh}\frac{Ra}{12h}
 -\frac{12h}{a^2}
       \left(\sqrt{1+\left(\frac{Ra}{12h}\right)^2}-1\right),
 \label{f16v69:rate}\\
 \mathcal I_a(R,h)&\longrightarrow\frac{R^2}{24h}\quad(a\downarrow0).
 \label{f16v69:rate-limit}
\end{align}
The distance is the torus distance when working on a torus.
The bounds are independent of its volume.  The factor $e^{\mathfrak A}$
is not asserted to be bounded under $\pi_a$.
\end{theorem}
\begin{proof}
For a real symmetric block matrix $\mathsf H$, the symmetric part of
$e^\psi\mathsf H e^{-\psi}$ has off-diagonal block
$\cosh(\psi_e-\psi_f)\mathsf H_{ef}$.  The same statement for the
Hermitian part holds after complexification.  Hence for every $z$,
\[
 \operatorname{Re}\langle z,e^\psi\mathsf H_W e^{-\psi}z\rangle
 \ge-\left[\nu(W)+12(\cosh(\lambda a)-1)\right]\|z\|^2.
\]
Indeed the deviation from $\mathsf H_W$ is bounded by
$\sum_{e\ne f}N_{ef}(\cosh(\lambda a)-1)|z_e||z_f|$;
$2|z_e||z_f|\le|z_e|^2+|z_f|^2$ and the row sum twelve finish the estimate.
Apply this bound to the squared norm of a solution of the conjugated
variational equation, and integrate its logarithmic derivative.  This
proves \eqref{f16v69:weighted-J} without positivity of the individual
matrix entries or of the time-ordered propagator.

Midpoints of two links in one plaquette are at distance at most $a$.
Choose the spatial weight
\[\psi_e=\lambda\min\{\operatorname{dist}(e,B),R\}.\]
It is zero on $B$ and equals $\lambda R$ on $A$.
The preceding bound has exponent
\[\mathfrak A-\lambda R+12h(\cosh(\lambda a)-1)/a^2.\]
Its minimizer obeys $\sinh(\lambda a)=Ra/(12h)$, giving
\eqref{f16v69:rate}.  Taylor expansion at zero gives
\eqref{f16v69:rate-limit}.
\end{proof}

This estimate recovers the global cost $e^{20h/a^2}$ only when the
worst-case bound $\nu\le20$ is actually inserted.  At flatness it is a
contraction with a Gaussian spatial rate.  If
$\max_P|\operatorname{Im}W_P(v)|\le a^2K(v)$ and $w(W_P(v))\ge0$
through the stated interval, then
$\mathfrak A\le12\int K(v)\,dv$.  This is a sufficient trajectory
condition, not a theorem that a typical native trajectory satisfies it.
The result concerns a derivative.  A finite difference must integrate
along a family of initial conditions and pay the cost on that family;
it is not enough to inspect only one endpoint trajectory.

\subsection{A genuine growing mode and a polynomial bound on the test family}

The sign cancellation does not make the spatial flow contractive on every
small-plaquette configuration.  We can settle that point exactly without
a separate magnetic-background approximation.

\begin{proposition}[Exact instability of the homogeneous imbalance]
\label{f16v69:imbalance}
In the two-angle homogeneous family of V68, put
$U_1=e^{\alpha\mathbf i}$, $U_2=e^{\beta\mathbf j}$, $U_3=I$.
Linearize the full spatial flow about its equal-angle trajectory
$\alpha=\beta=\theta(t)$, $u(t)=\sin^2\theta(t)$.
The homogeneous tangent $\delta\alpha=-\delta\beta=\eta$ satisfies
\begin{equation}
 \eta'=\frac{4u(t)}{a^2}\eta,
 \qquad
 \frac{\eta(t)}{\eta(t_0)}
 =\left[\frac{u(t_0)(1-u(t))}{u(t)(1-u(t_0))}\right]^{1/2}
 =:\mathcal R(t,t_0)^{1/2}.
 \label{f16v69:imbalance-growth}
\end{equation}
For $u(t_0)<1/\sqrt2$, the differential on the \emph{entire} lattice tangent
space, not just this homogeneous subspace, obeys
\begin{equation}
 \mathcal R(t,t_0)^{1/2}
 \le\|\mathsf J(t,t_0)\|_{2\to2}
 \le e^{\mathfrak A(t,t_0)}
 \le\mathcal R(t,t_0)^{3/2}.
 \label{f16v69:polynomial-J}
\end{equation}
If $u_0\to0$ and $u_0/a^2\to\infty$, the lower factor at fixed $t>0$
grows as $\sqrt{8tu_0/a^2}$.  For fixed $0<t_0<t$, instead,
$\mathcal R(t,t_0)\to t/t_0$.  Thus the same trajectory has a
cutoff-independent signed locality bound after a positive preflow interval,
while its sensitivity back to the rough initial slice grows.
\end{proposition}
\begin{proof}
The literal force calculation of V68 gives
\[
 \alpha'=-4a^{-2}\sin\alpha\cos\alpha\sin^2\beta,\qquad
 \beta'=-4a^{-2}\sin\beta\cos\beta\sin^2\alpha.
\]
At equal angles the antisymmetric eigenvalue of this $2\times2$
linearization is $4u/a^2$.  Along $u'=-8a^{-2}u^2(1-u)$,
\[
 \int_{t_0}^t\frac{4u(v)}{a^2}\,dv
 =\frac12\log\frac{u(t_0)(1-u(t))}{u(t)(1-u(t_0))}.
\]
Changing the two angles is a product-metric tangent with norm proportional
to $|\eta|$, with the same factor at both endpoints.  This gives the
operator lower bound.

Only $12$-plaquettes are nontrivial.  Their imaginary part has magnitude
$2u\sqrt{1-u^2}$, and their scalar part is $1-2u^2>0$.
A direction-$1$ or direction-$2$ link has two such incident plaquettes;
a direction-$3$ link has none.  Consequently
\[
 \nu(W(v))=12u(v)\sqrt{1-u(v)^2}\le12u(v),\qquad
 \mathfrak A(t,t_0)\le\tfrac32\log\mathcal R(t,t_0).
\]
Theorem~\ref{f16v69:weighted-flow} at $\lambda=0$ gives the upper bound
on all tangents.  The asymptotics follow from V68's exact flow brackets.
\end{proof}

This mode is not an assertion of a source-norm failure: the symmetric
quadratic curvature source can have zero first derivative in this
particular imbalance direction.  It proves that replacing the signed
curvature cost by zero in a link-flow stability estimate is invalid.
Neither the homogeneous configurations nor their growing mode have been
assigned a probability under the native vacuum.

\subsection{The optional constituent heat has a paid connection derivative}

The next estimate separates variation of the field from variation of
its connection; freezing the latter would repeat V67's unpaid step.

\begin{proposition}[Duhamel estimate retaining the connection variation]
\label{f16v69:heat-derivative}
Let $d_W$ be the spatial adjoint covariant gradient,
$\mathcal L_W=d_W^*d_W$, and let $W(\epsilon)$ have right-trivialized
velocity $v$.  Then, for $s>0$,
\begin{equation}
 \left\|\left.\frac d{d\epsilon}e^{-s\mathcal L_{W(\epsilon)}}
                      \right|_{\epsilon=0}\right\|_{2\to2}
 \le2\sqrt{\frac{2s}{e}}\,\|\dot d_W\|
 \le\frac{4\sqrt{6s/e}}a\|v\|_\infty.
 \label{f16v69:heat-variation}
\end{equation}
For $\mathcal F_W=\zeta_a a^{-2}\operatorname{Im}W_{12}$ and physical
spatial volume $\mathcal V_a$,
\begin{equation}
 \|\dot{(e^{-s\mathcal L_W}\mathcal F_W)}\|_a
 \le\sqrt8\,\zeta_a a^{-2}\|v\|_{a,\mathrm{links}}
   +4\sqrt{6s/e}\,\zeta_a a^{-3}\sqrt{\mathcal V_a}\,\|v\|_\infty.
 \label{f16v69:complete-heat-derivative}
\end{equation}
This is a global derivative bound, not a spatially localized estimate
for the $s>0$ source under the native law.
\end{proposition}
\begin{proof}
At finite regulator all operators are finite matrices, so Duhamel gives
\[
 \dot e^{-s\mathcal L}
 =-\int_0^s e^{-(s-v)\mathcal L}
           (\dot d^*d+d^*\dot d)e^{-v\mathcal L}\,dv.
\]
The spectral theorem gives
$\|d e^{-v\mathcal L}\|\le(2ev)^{-1/2}$ for $v>0$.
The two terms are therefore bounded by
$\|\dot d\|[(2ev)^{-1/2}+(2e(s-v))^{-1/2}]$.
Integration gives the first inequality.  In the unit-quaternion metric,
$\|\partial_\epsilon\operatorname{Ad}_{W_e}\|\le2|v_e|$.
Each site is the head of three positively oriented links; thus
$\|\dot d\|\le2\sqrt3\|v\|_\infty/a$.

A differentiated plaquette is the sum of four isometrically transported
link velocities.  For the chosen orientation a link occurs in at most
two such plaquettes.  Hence
$\|\dot{\mathcal F}_W\|_a\le\sqrt8\zeta_a a^{-2}\|v\|_{a,\mathrm{links}}$.
Use heat contraction, the product rule, and
$\|\mathcal F_W\|_a\le\zeta_a a^{-2}\sqrt{\mathcal V_a}$.
\end{proof}

The bound grows polynomially in $a^{-1}$ at fixed $s$, rather than as an
unpaid operator exponential.  It does not replace the two terms by an
independent-curvature model, and it does not evaluate their covariance.

\subsection{Native conditional replicas give a local source approximation}

Let $A$ contain every link of the $12$-plaquettes in the tests of
\eqref{f16v69:source}, and let $A_R$ contain all links at midpoint distance
less than $R$ from $A$.  Put
$\mathcal G_R=\sigma(U_e:e\in A_R)$ and
$F_R=\mathbb E_{\pi_a}[F\mid\mathcal G_R]$.
This is a conditional expectation under the native slice law.
It is not a flow with an artificially fixed outer boundary.
For $\lambda>0$, define the explicitly geometric tail
\begin{equation}
 \mathcal S_{a,A}(R,\lambda)
 =a^3\sum_{e\notin A_R}e^{-2\lambda\operatorname{dist}(e,A)},
 \qquad b_a(\lambda)=\frac{12}{a^2}(\cosh(\lambda a)-1).
 \label{f16v69:spatial-tail}
\end{equation}
For fixed physical support and $a\le a_0$, the sum is bounded independently
of the ambient torus by a constant times
$\int_{R-O(a)}^\infty(1+r)^2e^{-2\lambda r}\,dr$; ordinary cubic-cell
counting proves this assertion.  We retain the exact sum in the theorem.

Take two replicas $U,U'$ conditionally independently given the same
$\mathcal G_R$ data, both with the actual conditional law of $\pi_a$.
Join their link values outside $A_R$ by shortest product geodesics
$U^\theta$, $0\le\theta\le1$; inside $A_R$ the path is constant.
The geodesic speeds are at most $\pi$ per link.  Antipodal pairs form a
null set for the smooth positive finite-regulator law; any measurable
choice there gives the same integrals.  Put $W^\theta=\Phi_{a,\tau}U^\theta$
and define the source-gradient Gram at its endpoint by
\begin{equation}
 (\mathsf Q(W))_{ij}
 =\langle\nabla_a O_i(W),\nabla_a O_j(W)\rangle_{a,\mathrm{links}}.
 \label{f16v69:gradient-Gram}
\end{equation}
Here $\nabla_a$ is the gradient for the $a^3$-weighted tangent inner
product.  Its support lies in $A$.  Define the native replica moment
\begin{equation}
 \mathsf M_{a,R}
 =\mathbb E_{\rm rep}\int_0^1
       e^{2\mathfrak A(\tau,0;U^\theta)}\mathsf Q(W^\theta)\,d\theta.
 \label{f16v69:replica-moment}
\end{equation}
The interpolated fields are not postulated to have law $\pi_a$.
Their distribution is the displayed pushforward of the conditional-replica
law, and remains part of this explicit expectation.

\begin{theorem}[Native source-marked spatial localization]
\label{f16v69:replica-localization}
The localization error matrix satisfies
\begin{equation}
 \boxed{\quad
 0\preceq \mathsf E_R:=(\mathbb E_{\pi_a}[(F-F_R)(F-F_R)^*])
 \preceq \frac{\pi^2}{2}\mathcal S_{a,A}(R,\lambda)
                  e^{2\tau b_a(\lambda)}\mathsf M_{a,R}.
 \quad}
 \label{f16v69:replica-bound}
\end{equation}
The result is valid at every finite regulator under the actual Wilson
slice law, including its induced correlations.  The predictor $F_R$ is
gauge invariant and uses only $A_R$ in the original spatial slice.
Every declared exact source relation is preserved.

For a scalar test $f$, a completely explicit but possibly large bound is
\begin{equation}
 \mathsf E_R\le
 4\pi^2\zeta_a^4a^{-8}\|f\|_{2,a}^2
 \mathcal S_{a,A}(R,\lambda)e^{2\tau b_a(\lambda)}
 \mathbb E_{\rm rep}\int_0^1e^{2\mathfrak A(\tau,0;U^\theta)}d\theta.
 \label{f16v69:replica-crude}
\end{equation}
Neither the last exponential moment nor a useful cutoff-uniform bound on
$\mathsf M_{a,R}$ is supplied by this theorem.
\end{theorem}
\begin{proof}
Conditional independence gives the exact matrix identity
\[
 \mathsf E_R=\tfrac12\mathbb E_{\rm rep}
          [(F(U)-F(U'))(F(U)-F(U'))^*].
\]
For a real coefficient vector $c$, differentiate
$c^*O(\Phi_{a,\tau}U^\theta)$ along the product geodesic.
In \eqref{f16v69:weighted-J} take
$\psi_e=-\lambda\operatorname{dist}(e,A)$, which is zero on $A$.
The input tangent vanishes in $A_R$ and has per-link norm at most $\pi$.
Therefore
\[
 \left|\frac d{d\theta}c^*O(W^\theta)\right|
 \le\pi\mathcal S_{a,A}(R,\lambda)^{1/2}
       e^{\mathfrak A(\tau,0;U^\theta)+\tau b_a(\lambda)}
       \|\nabla_a(c^*O)(W^\theta)\|_a.
\]
Integrate in $\theta$, apply Cauchy--Schwarz on $[0,1]$, then take the
replica expectation.  This proves \eqref{f16v69:replica-bound} in every
coefficient direction.  Complex directions follow by complexification.

Gauge transformations preserve the native law and the sigma-field
$\mathcal G_R$; conditional expectation of a gauge-invariant function
is consequently gauge invariant as an $L^2$ class.  Equivalently the
replica and geodesic construction is covariant under the common endpoint
gauge transformations.  Linearity preserves exact source relations.

Finally,
\[
 |dO_f(W)[v]|
 \le\zeta_a^2a^{-1}\sum_x|f(x)|
                       \sum_{e\in\partial P_{12}(x)}|v_e|
 \le\sqrt8\zeta_a^2a^{-4}\|f\|_{2,a}\|v\|_{a,\mathrm{links}}.
\]
The factor eight is four plaquette edges times at most two occurrences
of one stored link in this orientation.  Substitution bounds the gradient
Gram and proves \eqref{f16v69:replica-crude}.
\end{proof}

No independence of the exterior links was used.  In particular, replacing
the replica law in \eqref{f16v69:replica-moment} by product Haar, or assigning
the interpolated configurations the original Wilson distribution, would
change the estimate.  The theorem also makes no relative lower claim for
the entrance or reflected Gram: the error budget is in an absolute native
$L^2$ matrix norm.

\subsection{Reflected pairings and the remaining native moment}

The following consequence specifies precisely how the spatial estimate
enters the original positive-time problem.  It is a Cauchy--Schwarz
application to the new error matrix, not an independent reconstruction
theorem.

\begin{proposition}[Reflected error with its normalization left visible]
\label{f16v69:reflected-error}
Place the sources and their localized predictors at matching positive
times of the stationary native history.  They may be averaged with the
same nonnegative temporal weights of total one, entirely on the positive
side.  Let $D$ be the unsmoothed slice Gram of $F$, and let $H,H_R$ denote
their reflected Grams, also allowing an additional common positive-time
delay.  Then
\begin{equation}
 \|H-H_R\|_{\rm op}\le2\sqrt{\|D\|_{\rm op}\|\mathsf E_R\|_{\rm op}}.
 \label{f16v69:reflection-error}
\end{equation}
If $\mathsf E_R\preceq\epsilon D$, the error conjugated by $D^{-1/2}$
on its exact quotient is at most $2\sqrt\epsilon$.
To conclude a relative error in $H$ instead requires a reflected lower
bound; it does not follow from the entrance bound.
\end{proposition}
\begin{proof}
Conditional expectation contracts the slice $L^2$ norm, so the Gram of
$F_R$ is at most $D$.  Temporal averaging contracts the same bounds by
Jensen and stationarity, both for $F,F_R$ and for their difference.
Expand the reflected difference using one difference factor in each term:
\[
 H-H_R=\mathbb E[(\vartheta(F-F_R))F^*]
                   +\mathbb E[(\vartheta F_R)(F-F_R)^*].
\]
Here the notation includes the indicated time placements.  Reflection
invariance preserves ordinary $L^2$ norms.  Cauchy--Schwarz bounds each
operator by $\sqrt{\|D\|\|\mathsf E_R\|}$.  Congruence by $D^{-1/2}$
gives the quotient assertion.  All reflected objects are computed in
the same native history measure.
\end{proof}

For fixed $\lambda$ the deterministic spatial factor has the finite limit
$e^{2\tau b_a(\lambda)}\to e^{12\tau\lambda^2}$, and its geometric tail
decays with $R$.  The real remaining input is the source-marked moment
\eqref{f16v69:replica-moment}, not the factor $e^{20\tau/a^2}$ from an
absolute force estimate.  A uniform bound on that matrix, together with a
bounded entrance Gram, would give a volume-independent spatial
approximation with an error tending to zero as $R\to\infty$.  The explicit
crude estimate contains $\zeta_a^4a^{-8}$ and an exponential curvature
moment.  It does not prove such a uniform bound.

Nor does V65's near-cancelling-cell estimate bound this new moment.
The cost $\nu$ involves the whole flow trajectory, including center
plaquettes, and its maximum is charged along conditional interpolations.
On a growing spatial torus it may also respond to remote curvature.
One must control that dependence or replace the maximum by a proved local
form estimate.  A rare-event probability without its exponential and
source weighting is insufficient.  A sharp source-specific estimate may
be substantially better than the link-flow norm, as the homogeneous
imbalance already illustrates.

The next dynamical target remains the positive-separation native Gram of
\eqref{f16v69:source}, rather than a new conditional coherent cell.
The present theorem supplies a signed method to approximate this source
spatially and an exact native matrix to estimate.  It does not populate
a nonzero reflected lower bound, and spatial flow time has nowhere been
identified with physical transfer time.

\subsection{Verification, preservation, and result ledger}

The companion diagnostic is
\begin{center}\small
\path{verify_ym_f16_v69_signed_flow.py}.
\end{center}
It distinguishes exact quaternion and incidence identities, floating-point
full-lattice Hessians and weighted propagators, nonlinear flow and imbalance
checks, connection-heat derivatives, and exact finite-probability replica
fixtures.  The latter are algebraic checks, not samples from the
four-dimensional Wilson vacuum.  No numerical result is used as a premise
of an analytic theorem above.  No interval certificate or proof-assistant
verification is claimed.  Removing this appendix and reverting the single
displayed version change recovers this lineage's V68 byte for byte.

\begin{record}[V69 result ledger]
\label{f16v69:ledger}
\emph{Proved here:} the complete all-configuration plaquette Hessian;
its center-sensitive negative-curvature bound and exact off-diagonal
incidence control; a signed time-dependent weighted differential estimate
with a diffusive spatial rate; an exactly amplifying homogeneous mode and
polynomial upper bound on all tangents along the same background;
a Duhamel bound retaining variation of the connection; native
conditional-replica spatial localization for the $s=0$ flowed source,
including its source-gradient matrix and crude constants; and the
resulting absolute reflected-pairing error estimate.

\emph{Inherited or established background:} the spatial Wilson flow,
the fixed native law and time support, V68's homogeneous force and flow
brackets, and the general exponential-weighting method.  No archived
Gaussian covariance theorem is used to evaluate the interacting measure.

\emph{Not asserted:} contractivity on all small-defect configurations;
a typical integrated-curvature bound; a Wilson law for conditional
geodesic interpolations; removal of the center sector; a localized theorem
for the optional $s>0$ step from its global derivative bound; or a
reflected lower Gram from finite-regulator source nonconstancy.

\emph{Still unproved:} a useful cutoff-uniform native replica moment,
spatial approximation with physical normalization, a noncollapsing
positive-separation source family, common dense-core physical-time
contraction, and the interacting continuum Yang--Mills mass gap.
\end{record}

\begin{firewall}
The signed curl form is retained before taking absolute values.
Its curvature cost is proved necessary for the link-flow differential,
not postulated small under the native measure.  The replica law, source
normalization, reflected quotient, and physical clock remain explicit.
V1--V68 are preserved unchanged.
\end{firewall}
% END F16 V69 APPEND

% BEGIN F16 V70 APPEND -- 2026-09-19
% Append-only continuation of the Post-Fifteenth-Archive V69.
\clearpage
\section[V70: native maximum saturation and gauge-completed local propagation]{V70: native maximum saturation\\and gauge-completed local propagation}
\label{f16v70:context}

V69 replaces an absolute force estimate by a signed Hessian bound, but
its amplification still uses a maximum over the entire spatial volume.
We first test whether that maximum can be small under the native law.
It cannot in the thermodynamic order of limits: the finite-energy property
of the Wilson measure forces its spatial translation orbit to encounter
every finite open link pattern.  Spatial Wilson flow is an invertible,
translation-equivariant map at a fixed regulator and preserves this
property.  Consequently the maximum in V69 saturates its worst possible
value at every flow time, almost surely.  The same statement holds along
the conditional-replica interpolations used there.  This is a statement
about an actual native marginal, not about a prescribed coherent boundary.

The replacement is not another assumed small maximum.  For a
gauge-invariant output, we add gauge damping to the \emph{tangent}
equation without changing its derivative.  The completed compact Hessian
has an exact connection-Laplacian decomposition, with a remainder
controlled by local plaquette defects.  A scalar comparison then charges
curvature only along propagating paths.  This gives a source-local replica
bound, a killed-domain continuation identity and a precise occupation-moment
test.  None of them is presented as an evaluation of the native reflected
lower Gram.  In particular, spatially sparse high potentials can still
produce large path weights.

\subsection{Scope, inherited inputs and established methods}

The finite-regulator source is exactly \eqref{f16v69:source}, with $s=0$,
and the spatial flow and original physical transfer are unchanged.  V69's
compact plaquette Hessian and spatial incidence are inherited.  We do not
repeat V65's chessboard estimate, V67's free-source calculation, or V68's
homogeneous evolution.  No archived Gaussian covariance estimate is used.

Gauge damping along flow is established background: L\"uscher and Weisz,
\emph{Perturbative analysis of the gradient flow in non-abelian gauge
theories}, JHEP \textbf{02} (2011), 051,
\href{https://arxiv.org/abs/1101.0963}{arXiv:1101.0963}, \S2.1, explain why
such damping need not change gauge-invariant observables.  Here the needed
statement is proved directly for the finite compact \emph{tangent}
equation, with no gauge slice and no inverse gauge Laplacian.  The
Feynman--Kac and exponential-occupation arguments are also established
methods; the particular Wilson decomposition and constants are derived
below.  These methods do not import a nonperturbative continuum theorem.

All thermodynamic assertions below keep $a>0$ and finite $\gamma$ fixed.
Let $\pi_{a,\gamma}$ be the spatial-slice marginal of an infinite-volume
four-dimensional Wilson Gibbs measure, including local limits of the
stationary Wilson cylinder construction.  Only its finite-energy
inequality is needed.  An interacting continuum measure is not assumed.
Finite-torus differential and source bounds are stated separately, on
nondegenerate spatial tori of side at least four.

\subsection{The global maximum saturates under the native marginal}

Let $\mathcal X=\SU(2)^{\mathcal E(\mathbb Z^3)}$ have the product topology,
and let $T_z$ be spatial translation.  Call a configuration universal when
its translation orbit is dense in $\mathcal X$; this is only a statement
that every finite open link pattern occurs at some translate.

\begin{proposition}[Native finite energy and universal configurations]
\label{f16v70:finite-energy}
For a finite set $S$ of stored spatial links and any measurable
$B\subset\SU(2)^S$,
\begin{equation}
 \pi_{a,\gamma}(U_S\in B\mid U_{S^c})
       \ge e^{-12\gamma|S|}H^S(B)\quad\hbox{almost surely}.
 \label{f16v70:finite-energy-bound}
\end{equation}
In particular $\pi_{a,\gamma}$-almost every configuration is universal.
This conclusion requires neither independent links nor an ergodic Gibbs
state.
\end{proposition}
\begin{proof}
Condition first in the full four-dimensional law on all links outside
$S$.  At most $6|S|$ Wilson plaquettes involve $S$, and the oscillation
of each $\gamma(1-w(U_P))$ is at most $2\gamma$.  The conditional density
relative to normalized product Haar measure is consequently at least
$e^{-12\gamma|S|}$, including all interactions internal to $S$.
The conditional expectation tower proves the same assertion conditional
only on the other spatial-slice links.  The inequality also passes through
local Gibbs limits; no marginal product law is inferred.

For a nonempty open cylinder with support $S$, its Haar mass is positive.
Choose $K$ disjoint spatial translates of $S$.  Conditional on the complete
complement of each translate, its event has probability at least the same
$p=e^{-12\gamma|S|}H^S(B)>0$.  Successively taking expectations gives
\[
 \pi_{a,\gamma}\{\hbox{none of the $K$ patterns occurs}\}\le(1-p)^K.
\]
The same argument after deleting finitely many translates gives infinitely
many occurrences almost surely.  A countable base of finite cylinders,
with rational-radius group balls, suffices for the whole product topology.
Intersecting their probability-one occurrence sets proves density of the
translation orbit.
\end{proof}

\begin{proposition}[Infinite spatial flow preserves every dense orbit]
\label{f16v70:infinite-flow}
At fixed $a>0$, the spatial Wilson flow extends to $\mathcal X$ for every
finite positive or negative flow time.  It is a translation-equivariant
homeomorphism for the product topology, with inverse $\Phi_{a,-t}$.
It maps universal configurations to universal configurations, simultaneously
for every real $t$.
\end{proposition}
\begin{proof}
In quaternion coordinates the vector field is the same finite-range
polynomial as \eqref{f16v68:ambient-flow}.  On a bounded product of ambient
balls its derivative has a uniform finite row sum.  The Banach-space ODE
in $\ell^\infty$ therefore has a unique local solution.  Each component
stays on its unit sphere because its radial derivative vanishes there;
the bounded vector field and derivative extend the solution for all finite
positive and negative times.  Uniqueness gives the inverse and translation
equivariance.

Continuity for the product topology needs more than this $\ell^\infty$
argument.  Iterating the finite-range difference inequality gives V68's
Poisson influence tail, with $|t|$ replacing $t$ for reverse flow.  At fixed
$a,t$ that tail tends to zero with graph distance.  Agreement to sufficient
accuracy on a sufficiently large finite collar thus controls every fixed
output cylinder, uniformly in all more distant links.  This proves product
continuity of both time directions.  Finally
$\{T_z\Phi_{a,t}U\}_z=\Phi_{a,t}\{T_zU\}_z$; a homeomorphism maps a dense
set to a dense set.  This deterministic assertion holds for every $t$ for
any one universal $U$, with no uncountable intersection of measure-one
sets.
\end{proof}

Use V69's exact cost
\[
 c(P)=3|\operatorname{Im}P|+4(-w(P))_+,\qquad
 \nu_e(U)=\sum_{P\ni e}c(P),\qquad
 \nu_\infty(U)=\sup_e\nu_e(U).
\]
There are four spatial plaquettes at each link, so $\nu_\infty\le20$.

\begin{theorem}[Almost-sure saturation of V69's global curvature cost]
\label{f16v70:saturation}
Under the native slice law just specified,
\begin{equation}
 \boxed{\quad
 \nu_\infty(\Phi_{a,t}U)=20\quad\hbox{for every }t\in\mathbb R,
 \qquad
 \mathfrak A_\infty(t,t_0;U)=20(t-t_0)/a^2\quad(t\ge t_0).
 \quad}
 \label{f16v70:max-saturation}
\end{equation}
If two native replicas are conditionally independent given the links in
a finite $A_R$, then every shortest-link-geodesic interpolant $U^\theta$
is universal almost surely, simultaneously for $0\le\theta\le1$.
Consequently the thermodynamic version of V69's matrix factors exactly as
\begin{equation}
 \mathsf M_{a,R}
   =e^{40\tau/a^2}\,
     \mathbb E_{\rm rep}\int_0^1
              \mathsf Q(\Phi_{a,\tau}U^\theta)\,d\theta.
 \label{f16v70:old-moment-factor}
\end{equation}
\end{theorem}
\begin{proof}
The maximum twenty is attainable by a genuine finite link pattern.
Choose one link to be identity, set the connector links of its four
spatial plaquettes to identity, and set each of their distinct opposite
parallel links to $q=-4/5+(3/5)\mathbf i$, with the needed orientation.
Every resulting plaquette is $q$ or $q^{-1}$, so $c(P)=5$.  No instructions
conflict: the four opposite links are private to these four paths.  The
cost at the central link is twenty.  It is a continuous local function,
so its supremum along a dense translation orbit is exactly twenty.
Proposition~\ref{f16v70:infinite-flow} proves the first assertion at all
times, and integration gives the second.

For replicas, take all finite test supports far outside $A_R$.  Their
joint conditional finite-energy lower bound is the product of the two
individual bounds in \eqref{f16v70:finite-energy-bound}; conditioning on
the shared interior does not remove it.  The same disjoint-cylinder
argument gives a dense \emph{joint} translation orbit of $(U,U')$ in
$\mathcal X\times\mathcal X$.  Given a desired finite pattern $V$ and
accuracy $\epsilon$, find a translate where both endpoints are within
$\epsilon/3$ of $V$ on that support.  Every shortest geodesic between those
endpoints stays within $\epsilon$ of $V$, uniformly in $\theta$, by the
triangle inequality.  Thus each interpolant has a dense orbit, and the
argument is uniform in $\theta$.  Antipodal choices are irrelevant to
these near-diagonal cylinders; under the finite-energy replica density,
exact antipodal pairs at any specified link also have probability zero.
Apply the deterministic flow assertion to each interpolant.  Substitution
in V69's definition of $\mathsf M_{a,R}$ gives
\eqref{f16v70:old-moment-factor}.
\end{proof}

This theorem invalidates neither V69's pathwise inequality nor its useful
bound on a prescribed homogeneous background.  It rules out proving a
small \emph{global} maximum by native typicality after the thermodynamic
limit.  It does not prove that the complete source-marked matrix diverges:
its remaining gradient factor may have additional suppression.  It also
does not exchange the thermodynamic and continuum limits or prove the
same equality in a fixed physical finite volume.  The actual source
response can be much smaller than this link-norm envelope.

\subsection{Gauge completion changes the tangent representative, not the source}

Work first on a finite spatial torus.  In right-trivialized coordinates
$v_e\in\operatorname{Im}\mathbb H$, set
\begin{equation}
 (D_W\xi)_{(x,\mu)}
       =\xi_x-\operatorname{Ad}_{W_\mu(x)}\xi_{x+\hat\mu}.
 \label{f16v70:vertical-map}
\end{equation}
This is the infinitesimal local gauge action; $D_W^*$ is its adjoint for
matching unscaled vertex and link metrics.  Matching $a^3$ weights give
the identical matrix.  Let $\mathsf H_W$ be the Riemannian Hessian of
$E_{\rm sp}$, as in V69.  Covariant derivatives below are taken along the
actual spatial-flow trajectory, not along a gauge-fixed probability law.

\begin{theorem}[Exact gauge-damped response of any gauge-invariant output]
\label{f16v70:gauge-completion}
Let $v(t)$ be the native flow variation with initial tangent $v_0$, and
let $\widetilde v(t)$ solve
\begin{equation}
 \nabla_t\widetilde v
  =-a^{-2}(\mathsf H_{W(t)}+D_{W(t)}D_{W(t)}^*)\widetilde v,
 \qquad \widetilde v(0)=v_0.
 \label{f16v70:damped-J}
\end{equation}
Then $\widetilde v(t)-v(t)=D_{W(t)}\xi(t)$ for some vertex field $\xi(t)$.
For every differentiable gauge-invariant output $O$,
\begin{equation}
 \boxed{\quad dO(W(t))[\widetilde v(t)]
                    =dO(W(t))[v(t)].\quad}
 \label{f16v70:source-equivalence}
\end{equation}
There is no inverse of $D_W^*D_W$, no removal of its stabilizer, and no
change to $W(t)$ or to the native law.
\end{theorem}
\begin{proof}
For a fixed vertex field $\xi$, equivariance of the gradient flow implies
\[
 \nabla_t(D_{W(t)}\xi)=-a^{-2}\mathsf H_{W(t)}D_{W(t)}\xi.
\]
Indeed its fundamental gauge vector field commutes with the flow vector
field; the torsion-free connection identifies the two covariant
derivatives, and the derivative of the negative gradient is
$-a^{-2}\mathsf H_W$.  Define
$\xi(0)=0$ and $\xi'=-a^{-2}D_{W(t)}^*\widetilde v(t)$.
The difference $\widetilde v-v-D_W\xi$ then solves the homogeneous native
variational equation with zero initial data, hence is zero.  Gauge
invariance gives $dO(W)D_W=0$, proving
\eqref{f16v70:source-equivalence}.  All equations are finite-dimensional
linear ODEs, and none requires a nondegenerate gauge action.
\end{proof}

\subsection{An exact compact lattice connection-Laplacian decomposition}

For each orientation, regard the link tangent as an adjoint field at its
tail.  Let $\mathsf L_W^{\rm link}$ be the dimensionless covariant
Laplacian on these three fields:
\begin{equation}
 (\mathsf L_W^{\rm link}v)_\mu(x)
  =\sum_{\nu=1}^3\bigl[2v_\mu(x)
   -\operatorname{Ad}_{W_\nu(x)}v_\mu(x+\hat\nu)
   -\operatorname{Ad}_{W_\nu(x-\hat\nu)}^{-1}v_\mu(x-\hat\nu)\bigr].
 \label{f16v70:link-laplacian}
\end{equation}
Unlike the physical Hamiltonian this operator has no physical-time
interpretation.  Its coefficient in the spatial equation is $a^{-2}$.
For each spatial plaquette set
\[
 d_P=|1-W_P|=\sqrt{2(1-w(W_P))},\quad
 r_{ef}=\sum_{P\supset\{e,f\}}d_P\ (e\ne f),\quad
 z_e=\sum_{P\ni e}(1-w(W_P)).
\]

\begin{theorem}[Exact gauge-completed remainder with local Wilson bounds]
\label{f16v70:Weitzenbock}
At every compact link configuration,
\begin{equation}
 \boxed{\quad
 \mathsf H_W+D_WD_W^*
      =\mathsf L_W^{\rm link}+\mathsf R_W,\qquad
 (\mathsf R_W)_{ee}=-z_e I_3,
 \qquad \|(\mathsf R_W)_{ef}\|\le r_{ef}\ (e\ne f).
 \quad}
 \label{f16v70:decomposition}
\end{equation}
The remainder is symmetric and supported on pairs in a common plaquette.
Moreover
\begin{equation}
 \sum_{f\ne e}r_{ef}=3\sum_{P\ni e}d_P\le24,
 \qquad z_e+\sum_{f\ne e}r_{ef}\le32.
 \label{f16v70:local-row}
\end{equation}
At every flat connection the remainder vanishes exactly.  Center
plaquettes are included, not hidden in a small-field exception.
\end{theorem}
\begin{proof}
We give the complete one-plaquette cancellation.  Write its stored links
in product order as $A,B,C,D$, so $P=ABC^{-1}D^{-1}$, and use the same
letters in sans-serif type for their adjoint matrices.  Put
$Z=\operatorname{Ad}_P$ and $M=w(P)I-C_p$, where $C_pv=p\times v$.
The transported velocity maps of V69, from stored right velocities to
the plaquette basepoint, are
\[
 T_1=I,\quad T_2=\mathsf A,\quad
 T_3=-Z\mathsf D,\quad T_4=-Z.
\]
For $i<j$ the Hessian block is $T_i^{\mathsf T}MT_j$.
Quaternion rotation gives the exact identities
\begin{equation}
 MZ=M^{\mathsf T},\qquad Z^{\mathsf T}MZ=M,\qquad
 \|M-I\|=\|M^{\mathsf T}-I\|=d_P.
 \label{f16v70:M-identities}
\end{equation}
They follow also from $C_p^3=-|p|^2C_p$ and
$Z=I+2wC_p+2C_p^2$; no expansion at identity is used.

The two longitudinal neighbors of an edge, and its diagonal contribution
$2I$, come from $D_WD_W^*$ and coincide with the longitudinal part of
\eqref{f16v70:link-laplacian}.  Each plaquette contributes the two
transverse same-orientation bonds of that Laplacian and four mixed
$D_WD_W^*$ blocks.  After these are combined with its Hessian, the six
upper-triangle remainder blocks are exactly
\begin{center}\small
\begin{tabular}{@{}cl@{}}
\toprule
pair & remainder block\\
\midrule
$12$ & $(M-I)\mathsf A$\\
$13$ & $(I-M^{\mathsf T})\mathsf D$\\
$14$ & $I-M^{\mathsf T}$\\
$23$ & $\mathsf A^{\mathsf T}(I-M)Z\mathsf D$\\
$24$ & $\mathsf A^{\mathsf T}(I-M^{\mathsf T})$\\
$34$ & $\mathsf D^{\mathsf T}(M-I)$\\
\bottomrule
\end{tabular}
\end{center}
Every diagonal remainder block is $(w(P)-1)I$; the lower blocks are
transposes.  The relation $Z\mathsf D=\mathsf A\mathsf B\mathsf C^{\mathsf T}$
checks the shared-head $23$ block.  This table accounts for every block
and proves the bound $d_P$, since all other factors are orthogonal.
Summing distinct plaquettes proves \eqref{f16v70:decomposition}.
Each edge meets four spatial plaquettes and three other edges in each,
so \eqref{f16v70:local-row} follows from $d_P\le2$, $1-w(P)\le2$.
\end{proof}

This is a decomposition of the full compact Hessian, not a positive
Hessian assertion.  The sign of its remainder is not discarded before
the longitudinal gauge cancellation.  Applying a componentwise absolute
bound directly to $\mathsf H_W$ would leave positive scalar growth even
at flatness; the completion removes precisely that avoidable loss for
gauge-invariant outputs.

\subsection{A scalar propagation kernel with local occupation costs}

Let $B_{ef}$ be the adjacency of equal-orientation links whose tails differ
by one spatial lattice step.  Each row has sum six.  Along the actual
trajectory define symmetric scalar jump rates and a nonnegative potential
by
\begin{equation}
 \begin{aligned}
 q_{ef}(t)&=a^{-2}(B_{ef}+r_{ef}(W(t)))\quad(e\ne f),\\
 q_{ee}(t)&=-\sum_{f\ne e}q_{ef}(t),\\
 k_e(t)&=a^{-2}\left[z_e(W(t))+\sum_{f\ne e}r_{ef}(W(t))\right].
 \end{aligned}
 \label{f16v70:rates}
\end{equation}
Thus the jump rate is at most $30/a^2$ and $0\le k_e\le32/a^2$.
All jumps join midpoints at distance at most $a$.  The process is
nonexplosive, including on the infinite lattice at fixed regulator.

\begin{theorem}[Occupation-local domination of the gauge-completed tangent]
\label{f16v70:FK}
Let $\mathsf K(t,t_0)$ be the positive scalar propagator of
$\partial_t u=(q(t)+\operatorname{diag}k(t))u$.  The damped variation obeys
\begin{equation}
 \boxed{\quad
 |\widetilde v_e(t)|
       \le\sum_f\mathsf K(t,t_0)_{ef}|\widetilde v_f(t_0)|.
 \quad}
 \label{f16v70:domination}
\end{equation}
For $h=t-t_0$ and nonnegative $g$, it has the representation
\begin{equation}
 (\mathsf K(t,t_0)g)(e)
   =\mathbb E_e\left[
      \exp\left\{\int_0^h k_{X_s}(t-s)\,ds\right\}g(X_h)\right],
 \label{f16v70:FK-formula}
\end{equation}
where the jump process uses the \emph{reversed operator-time} rates
$q(t-s)$.  This process is an auxiliary spatial comparison, not the
Wilson physical-time process and not a new law for the links.
\end{theorem}
\begin{proof}
Pass to right-trivialized tangent frames.  Covariant differentiation then
adds only a skew-adjoint on-link rotation to the ordinary time equation;
it does not affect the derivative of the on-link norm.  The six connection
transports in \eqref{f16v70:link-laplacian} are isometries.  The upper Dini
derivative of $h_e=|\widetilde v_e|$, with the usual regularization at zero,
therefore satisfies
\[
 \partial_t h_e\le a^{-2}\left[-6h_e+\sum_f B_{ef}h_f
                                    +z_eh_e+\sum_{f\ne e}r_{ef}h_f\right].
\]
This is exactly $\partial_t h\le(q+\operatorname{diag}k)h$.
The scalar off-diagonal coefficients are nonnegative, so the positive
linear comparison principle proves \eqref{f16v70:domination}.
On a finite torus it follows by the integrating-factor series; bounded
rates give the same conclusion by exhaustion on the infinite lattice.

Expand the scalar propagator into its time-ordered jump series.  Each
holding interval contributes its escape-rate factor and its local
potential integral, and each jump contributes the corresponding rate.
The latest operator acts on the starting row first, so chronological path
time sees $q(t-s)$, not $q(t_0+s)$.  This proves
\eqref{f16v70:FK-formula}, including its time convention.  Bounded rates
and potentials justify convergence and expectation interchange.
\end{proof}

At flatness $r=z=k=0$, and the kernel is exactly the rate-$a^{-2}$ scalar
nearest-neighbor heat kernel on each orientation copy.  Remote curvature
is charged only by paths that visit it.  The potential bound $32/a^2$ is
not a better \emph{global} constant than V69's twenty; the improvement is
removing the global maximum and retaining gauge-invariant propagation.
No small native occupation moment is inferred from this identity.

\subsection{The source-local native replica bound}

Retain V69's $A$, $A_R$, conditional replicas, output functions $O_i$ and
gradient Gram $\mathsf Q_{ij}(W)=\langle\nabla_aO_i,\nabla_aO_j\rangle_a$.
For each interpolant use the scalar kernel of its own flowed trajectory,
and set
\begin{equation}
 T_R(e;U^\theta)
   =(\mathsf K^{\theta}(\tau,0)\mathbf1_{A_R^c})(e),\qquad
 \mathcal T_R(U^\theta)=a^3\sum_{e\in A}T_R(e;U^\theta)^2.
 \label{f16v70:transport-tail}
\end{equation}
These quantities use \eqref{f16v70:rates} along the same replica
interpolation.  No product probability law is substituted.

\begin{theorem}[Native source localization without a global curvature maximum]
\label{f16v70:native-localization}
For $F_R=\mathbb E_{\pi_a}(F\mid U_{A_R})$, the complete error matrix obeys
\begin{equation}
 \boxed{\quad
 0\preceq\mathsf E_R
 \preceq\frac{\pi^2}{2}\,
   \mathbb E_{\rm rep}\int_0^1
        \mathcal T_R(U^\theta)\,
        \mathsf Q(\Phi_{a,\tau}U^\theta)\,d\theta.
 \quad}
 \label{f16v70:local-replica}
\end{equation}
The statement is simultaneous in every source coefficient direction and
preserves exact source relations.  Gauge invariance, original time support,
and the original native probability law are unchanged.
\end{theorem}
\begin{proof}
For a coefficient vector $c$, the geodesic initial tangent vanishes in
$A_R$ and has norm at most $\pi$ on each exterior link.  Replace its flow
variation by the gauge-damped variation using
Theorem~\ref{f16v70:gauge-completion}.  Equations
\eqref{f16v70:domination}--\eqref{f16v70:transport-tail} give the on-link
bound $|\widetilde v_e(\tau)|\le\pi T_R(e;U^\theta)$.
The output gradient is supported in $A$, so Cauchy--Schwarz for the
$a^3$-weighted sum yields
\[
 \left|\frac d{d\theta}c^*O(\Phi_{a,\tau}U^\theta)\right|^2
 \le\pi^2\mathcal T_R(U^\theta)\,
                           c^*\mathsf Q(\Phi_{a,\tau}U^\theta)c.
\]
Integrate in $\theta$, apply Cauchy--Schwarz on $[0,1]$, then use the exact
conditional-replica identity
$\mathsf E_R=\frac12\mathbb E_{
m rep}[(F(U)-F(U'))(F(U)-F(U'))^*]$.
This proves the matrix inequality.  Gauge transformations rotate tangent
blocks isometrically and leave $d_P,z_e$ invariant, so the construction
respects the gauge action.  Conditional expectation and differentiation
are linear in the source bank; null relations are retained by congruence.
No probability law is assigned to $U^\theta$ other than its replica-induced
law.  Infinite-volume fixed-regulator versions follow by the bounded-rate
kernel construction and local approximation, with bounds allowed to depend
on $a,\tau$.
\end{proof}

For a flat path and $A_R$ defined by midpoint distance, the scalar tail
satisfies the useful benchmark
\begin{equation}
 T_R(e)\le\min\{1,6e^{-J_a(R/\sqrt3,\tau)}\}\quad(e\in A),
 \qquad
 J_a(r,t)=\frac ra\operatorname{arsinh}\frac{ra}{2t}
 -\frac{2t}{a^2}\left(\sqrt{1+(ra/(2t))^2}-1\right).
 \label{f16v70:flat-tail}
\end{equation}
Indeed an endpoint outside $A_R$ is at least distance $R$ from its starting
$e\in A$.  One coordinate of the lifted random walk must have displacement
at least $R/\sqrt3$; its exponential moment is
$\exp[2t(\cosh(\lambda a)-1)/a^2]$.  Chernoff optimization and the six
signed-coordinate union give \eqref{f16v70:flat-tail}.  The same upper bound
holds for torus distance by lifting the walk.  Since $J_a(r,t)\to r^2/(4t)$,
$\mathcal T_R$ has an explicit Gaussian upper tail times $a^3|A|$ in this
benchmark.  This is not a native-flatness assumption, and the curvature
output itself has zero gradient at exact flatness.

V69's reflected-error bound now applies with
\eqref{f16v70:local-replica} in place of its global-max estimate.
In particular $\|H-H_R\|\le2\sqrt{\|D\|\|\mathsf E_R\|}$ for the same
positive-time placements.  This inherited final inequality is not counted
as a new reconstruction result.  A relative estimate still requires a
reflected lower Gram.

\subsection{Killed-domain continuation and an explicit occupation test}

The new path weight should not be bounded by another whole-volume maximum.
For a finite link domain $\mathcal D$, let $\mathsf K_{\mathcal D}(t,s)$ be
the propagator of the restriction of $q+\operatorname{diag}k$ to that
domain, with the full escape rates retained on its diagonal.  Thus paths
are killed on exit, not reflected back into $\mathcal D$.

\begin{proposition}[Exact boundary return and a local exponential-moment test]
\label{f16v70:occupation}
For $u(t)=\mathsf K(t,0)g$, its restriction satisfies
\begin{equation}
 u_{\mathcal D}(t)=\mathsf K_{\mathcal D}(t,0)g_{\mathcal D}
  +\int_0^t\mathsf K_{\mathcal D}(t,s)
       q_{\mathcal D,\mathcal D^c}(s)u_{\mathcal D^c}(s)\,ds.
 \label{f16v70:boundary-return}
\end{equation}
In particular the first term vanishes for $g=\mathbf1_{\mathcal D^c}$.
For the reversed-time process on any one time interval of length at most
$h$, suppose every restart within that interval, at every $e\in\mathcal D$,
has expected potential occupation before exit or the interval end at most
$\eta<1$.  Its killed exponential occupation moment is at most
$(1-\eta)^{-1}$.  Partitioning a longer interval into $m$ such intervals
with the same restart bound gives at most $(1-\eta)^{-m}$.
\end{proposition}
\begin{proof}
Restrict the scalar ODE to the domain; the only cross-domain entries are
the jump rates, not the potential.  Variation of constants proves
\eqref{f16v70:boundary-return}.  Positivity retains every return from the
exterior and prevents silently replacing the full kernel by its killed
restriction.

For the occupation estimate expand the exponential into ordered time
integrals.  By the Markov property, after the first marked time the sum of
the subsequent ordered integrals is bounded using the same restart
hypothesis.  Induction bounds the $n$th ordered integral by $\eta^n$,
including killing.  Summing the geometric series proves the claim.
Iteration at the interval endpoints proves the product bound.  This is
the elementary exponential-occupation argument, with its required
uniformity over restarts explicitly stated.
\end{proof}

The restart hypothesis concerns a local comparison walk along the actual
connection trajectory.  It has not been established under the native
replica law.  Moreover a small spatial density of bad sites cannot supply
it.  The following exact comparison model records the problem quantitatively.

\begin{proposition}[One strong defect can invalidate a density-only path estimate]
\label{f16v70:one-defect}
On $a\mathbb Z^3$ let $q=\Delta/a^2$ be the scalar walk with rate $a^{-2}$
to each neighbor and let $k_x=\beta a^{-2}\mathbf1_{x=0}$, $\beta>6$.
Its Feynman--Kac kernel satisfies
\begin{equation}
 K(t)_{00}\ge e^{(\beta-6)t/a^2}.
 \label{f16v70:defect-growth}
\end{equation}
For a point at graph distance $d\ge1$ from zero and $t\ge da^2$,
\begin{equation}
 K(t)_{x0}\ge
  e^{(\beta-6)t/a^2-\beta d}\frac{d^d}{d!}.
 \label{f16v70:remote-defect}
\end{equation}
Thus even a single defect, of zero asymptotic spatial density, can defeat
a cutoff-uniform occupation bound; at fixed physical separation $d=O(a^{-1})$
its displayed growth still dominates that distance cost.
\end{proposition}
\begin{proof}
For \eqref{f16v70:defect-growth}, retain paths making no jump: their
probability is $e^{-6t/a^2}$ and their weight is $e^{\beta t/a^2}$.
For \eqref{f16v70:remote-defect}, prescribe a shortest path of $d$ jumps
in the first interval of length $da^2$.  The probability of exactly those
ordered directions and no other jumps is $e^{-6d}d^d/d!$.  Then prescribe
no further jump, and retain only the potential accumulated at zero in
the remaining interval.  Multiplying the three factors gives the formula.
This is a scalar comparison example, not a proof that a Wilson observable
has the same instability or that its actual tangent attains its envelope.
\end{proof}

\subsection{The next native comparison and the result ledger}

There is now a definite reason not to pursue V69's global maximum as a
small thermodynamic random variable: it is almost surely twenty for every
finite regulator and flow time.  The replacement instead has three
explicit pieces: exact gauge-equivalent source differentiation, the
local Wilson remainder of \eqref{f16v70:decomposition}, and the joint
native replica average in \eqref{f16v70:local-replica}.  The last average
keeps output gradients correlated with the whole local propagation
history.  No new Wilson measure, Gaussian field or product replica law
has been substituted.

The next quantitative task is a source-weighted estimate of that
\emph{local} occupation/return matrix on a declared physical spatial
region, including the exterior term in \eqref{f16v70:boundary-return}, or
a direct positive-separation native covariance estimate which avoids this
upper-envelope route.  A density of exceptional plaquettes alone cannot
pay the time spent near an amplifying defect.  Conversely failure of a
scalar envelope does not prove failure of the underlying gauge-invariant
source, whose derivative may cancel an unstable gauge or link direction.
This distinction should prevent another circular attempt to infer
noncollapse from a smoothing upper bound.

The companion diagnostic is
\begin{center}\small
\path{verify_ym_f16_v70_local_propagation.py}.
\end{center}
It checks exact quaternion and incidence identities, the complete compact
block decomposition, gauge covariance, native versus damped nonlinear
tangent responses, scalar time ordering and domination, finite-replica
identities, and single-defect kernels.  It is not a Wilson-vacuum simulation,
a proof of a continuum occupation estimate, an interval certificate or
proof-assistant verification.  Removing this delimited appendix and
reverting the one displayed version change recovers V69 byte for byte.

\begin{record}[V70 result ledger]
\label{f16v70:ledger}
\emph{Proved here:} native finite-energy universal-pattern occurrence;
fixed-regulator infinite spatial flow as a translation-equivariant
homeomorphism; exact thermodynamic saturation of V69's global maximum,
also along its native conditional replicas; exact gauge completion of
source derivatives without gauge fixing; the full compact connection-
Laplacian/remainder identity with explicit plaquette blocks; local
Feynman--Kac domination; a source-marked native replica localization
bound without a global maximum; killed-domain continuation with its
exterior return; and explicit single-defect occupation obstructions.

\emph{Inherited or established background:} the Wilson finite-energy
specification and spatial flow, V69's compact Hessian and source geometry,
gauge damping as a general method, scalar Feynman--Kac comparison and the
exponential-occupation argument.  Their specialized forms used above are
proved; no external renormalization or mass-gap conclusion is imported.

\emph{Not asserted:} uniform saturation on fixed physical finite volumes;
divergence of the entire source-gradient-weighted matrix from its scalar
factor alone; a small local native occupation moment; equality of the
comparison walk with physical Wilson time evolution; or source collapse
from a scalar upper-envelope instability.

\emph{Still unproved:} a useful physical-normalization bound on the local
native replica matrix, a noncollapsing positive-separation reflected Gram
for the specified source, common dense-core physical-time contraction,
and the interacting four-dimensional continuum Yang--Mills mass gap.
\end{record}

\begin{firewall}
The thermodynamic maximum is evaluated, not assumed typical.  Gauge
completion changes only the representative tangent seen by a gauge-
invariant output.  Curvature is charged along actual comparison paths,
with their replica-induced environment retained.  This is an upstream
repair of a localization estimate, not a claimed continuum closure.
V1--V69 remain unchanged.
\end{firewall}
% END F16 V70 APPEND

% BEGIN F16 V71 APPEND -- 2026-09-19
% Append-only continuation of the Post-Fifteenth-Archive V70.
\clearpage
\section[V71: gauge-aligned replicas and interpolation-induced curvature]{V71: gauge-aligned replicas\\and interpolation-induced curvature}
\label{f16v71:context}

V70 replaces the thermodynamic curvature maximum by a local propagator.
Its source estimate still joins independently sampled links by shortest
linkwise geodesics.  This interpolation is gauge equivariant under a
\emph{common} transformation, but it does not respect independent choices
of gauge representative for the two endpoints.  We test that earlier
step before asking for another occupation estimate.  Two flat,
gauge-equivalent endpoints can generate a center plaquette halfway along
their linkwise geodesic.  More strongly, under the native conditional-replica
law this mechanism has a coefficient-independent positive probability
whenever the endpoint plaquettes concentrate near identity.  Thus a
native small-plaquette estimate does not transfer even to the midpoint
of that interpolation.

The replacement aligns the second replica on an anchored spanning tree
before interpolating.  It preserves the exact native source difference
and the retained interior.  Gauge coordinates integrate out with normalized
Haar measure; the remaining chord holonomies retain their full induced
law, not independent one-link laws.  V70's propagator can then be charged
to actual holonomy discrepancies instead of arbitrary gauge displacement.
We also retain an exact signed, source-directed response identity.  These
results remove a proved source of artificial curvature, not the true
curvature, winding data, or the missing reflected lower Gram.

\subsection{Scope, inherited ingredients and notation}

The spatial flow, native slice law $\pi_a$, source family
\eqref{f16v69:source}, and physical transfer $P_a$ are unchanged.  Work on
a finite nondegenerate spatial torus when differentiating or disintegrating
its smooth positive density.  The probability statements use only native
local gauge invariance and an explicitly stated endpoint plaquette bound;
they therefore also apply to the corresponding local Gibbs limits.  The
anchored-tree construction is finite; no infinite-tree locality limit is
asserted.

Tree gauge coordinates are established background, not a new gauge-fixing
method.  For a modern account on general graphs see I.~M.~Burbano and
C.~W.~Bauer, \emph{Gauge Loop-String-Hadron Formulation on General Graphs
and Applications to Fully Gauge Fixed Hamiltonian Lattice Gauge Theory},
\href{https://arxiv.org/abs/2409.13812}{arXiv:2409.13812}, Sections III
and C.3.2.  The results below concern a different use: aligning conditional
replicas while preserving their original source difference.  The Haar
factorization and all estimates needed here are proved directly.  No
continuum conclusion is imported from tree gauge coordinates or from
V70's scalar comparison.

Write $d(U,V)=\arccos w(UV^{-1})\in[0,\pi]$ for the round
bi-invariant group distance, and $\rho(P)=d(I,P)$.  Quaternion chord norm
and this distance obey
\begin{equation}
 |U-V|\le d(U,V)\le\frac\pi2|U-V|,
 \qquad \rho(P)^2\le\frac{\pi^2}{2}(1-w(P)).
 \label{f16v71:metric}
\end{equation}
A midpoint of a nonantipodal pair of unit quaternions is
$m(x,y)=(x+y)/|x+y|$.  Reversing a stored edge reverses its path and
conjugates its group value.  All plaquette products below use their
cyclic orientations.

\subsection{Flat endpoints do not imply a flat linkwise interpolation}

\begin{proposition}[A pure-gauge interpolation creates a center plaquette]
\label{f16v71:flat-artifact}
On an actual spatial lattice patch there are two globally flat link
configurations $U,U'$ which agree on any prescribed disjoint interior
region, are related by a gauge transformation equal to identity on that
region, and have unique shortest linkwise geodesics, such that
\begin{equation}
 P(U)=P(U')=I,\qquad P(U^\theta)=e^{-2\pi\theta\mathbf i}
 \quad(0\le\theta\le1)
 \label{f16v71:flat-midpoint}
\end{equation}
on a marked square.  In particular $P(U^{1/2})=-I$.
For every gauge-invariant source and every finite spatial flow time, the
two endpoint source values agree exactly.  The interpolation in
\eqref{f16v71:flat-midpoint} need not have constant source value.
\end{proposition}
\begin{proof}
Put $U_e=I$ everywhere.  Label the marked square's vertices cyclically
as $0,1,2,3$ and assign vertex gauges with angles
\[
 \varphi_0=0,\quad \varphi_1=\pi/2,\quad
 \varphi_2=\pi-\varepsilon,\quad \varphi_3=-\pi/2,
 \qquad 0<\varepsilon<\pi/4.
\]
Let $g_v=e^{\varphi_v\mathbf i}$ at these vertices and $g_v=I$
elsewhere, and set $U'_{xy}=g_xg_y^{-1}$.  Every plaquette telescopes
to identity.  The four principal cyclic edge angles of $U'$ are
\[
 -\pi/2,\quad-\pi/2+\varepsilon,\quad
 -\pi/2-\varepsilon,\quad-\pi/2.
\]
They lie strictly in $(-\pi,\pi)$ and sum to $-2\pi$.  All other
nonidentity edge angles are also strictly inside that interval.  The
shortest linkwise paths multiply these four angles by $\theta$,
proving \eqref{f16v71:flat-midpoint}.  The selected vertices can be placed
outside the retained interior and its incident vertices, so the two
configurations agree there.  A continuous vertex-gauge path from $I$ to
$g$ instead keeps every plaquette flat throughout.  Gauge equivariance
of the spatial flow and gauge invariance of the source give equality of
the endpoint values, for any flow time.

Already at zero flow time the source $|\operatorname{Im}P|^2$ along the
linkwise path is $\sin^2(2\pi\theta)$.  Its signed derivative integrates
to zero, while
\[
 \int_0^1\left|\frac d{d\theta}\sin^2(2\pi\theta)\right|^2d\theta
 =2\pi^2.
\]
Thus squaring the interpolation derivative can lose an exact endpoint
cancellation even for the same quadratic curvature observable.
\end{proof}

This example concerns representatives, not curvature hidden in a flat
connection.  Both endpoints have zero action; their interpolant does not.
Tangent gauge completion in V70 does not change that interpolating
configuration path.  It therefore does not remove this particular loss.

\begin{theorem}[Native midpoint defects need not become rare at weak coupling]
\label{f16v71:native-midpoint}
Let $U,U'$ have the native conditional-replica law given the same retained
link data on $C$.  Choose a plaquette all of whose vertices lie outside
the vertices incident to $C$.  Its midpoint uses the original, unaligned
linkwise geodesics.  Put
\begin{equation}
 \epsilon_0=1/64,\qquad
 m_0=H\{X:d(I,X)<\epsilon_0\}
   =\frac{\epsilon_0-\tfrac12\sin(2\epsilon_0)}\pi,
 \quad p_a=\pi_a\{\rho(P)>\epsilon_0\}.
 \label{f16v71:midpoint-constants}
\end{equation}
Then
\begin{equation}
 \boxed{\quad
 \mathbb P_{\rm rep}\{w(P(U^{1/2}))\le-3/4\}
 \ge m_0^6\,(1-2p_a)_+ .\quad}
 \label{f16v71:native-midpoint-bound}
\end{equation}
In particular the lower limit is at least $m_0^6>0$ along any native
sequence with $p_a\to0$.  For the symmetric periodic Wilson measures
and local limits of V64 one may insert
\begin{equation}
 p_a\le\min\left\{1,
 \frac{B_{\gamma_a}}{\gamma_a(1-\cos\epsilon_0)}\right\},
 \qquad B_\gamma=8+\log\gamma-\tfrac23\log\frac8{3\pi^3}.
 \label{f16v71:midpoint-pressure}
\end{equation}
The assertion is about the replica midpoint before spatial flow.  It is
not a lower bound on a flowed occupation moment or on a physical-time
correlation.
\end{theorem}
\begin{proof}
Condition on the retained data.  Gauge transformations at the three
nonroot vertices of the square leave those data fixed.  The joint replica
law is invariant under independent such transformations of the two
replicas.  We may therefore average its event indicator over six
independent normalized Haar vertex gauges without changing the desired
probability.

For one fixed endpoint configuration, the first three cyclic edges after
this random gauge transformation are independent Haar variables: successively
solving for the three vertex gauges is a triangular product of left and
right Haar translations.  The rooted plaquette product $P$ is unchanged,
since the root gauge is identity.  If $\rho(P)\le\epsilon_0$, forcing these
three edges into balls of radius $\epsilon_0$ about $I$ makes the fourth
edge lie within $4\epsilon_0$ of $I$.  This event has probability $m_0^3$.
For the second replica force its first three edges instead into the balls
about $-\mathbf i$.  The target product $(-\mathbf i)^4$ is identity, so
its fourth edge lies within $4\epsilon_0$ of $-\mathbf i$; the probability
is again $m_0^3$.  This argument uses no assumption on the law of the
plaquette holonomies themselves.

For any of the four pairs on these events,
$|x-I|\le4\epsilon_0$, $|y+\mathbf i|\le4\epsilon_0$.
Normalization of a nonzero vector satisfies
$|v/|v|-v_0/|v_0||\le2|v-v_0|/|v_0|$.  With
$v_0=I-\mathbf i$, each midpoint differs from
$(I-\mathbf i)/\sqrt2$ by at most $8\sqrt2\epsilon_0$.
The sums $x+y$ cannot vanish, so their shortest midpoints are unique.
Telescoping the four unit factors gives
\[
 |P(U^{1/2})+I|\le32\sqrt2\epsilon_0,
 \qquad 1+w(P(U^{1/2}))\le1024\epsilon_0^2=1/4.
\]
Average the two independent gauge events and keep the event that both
endpoint holonomies have distance at most $\epsilon_0$.  Each replica has
marginal $\pi_a$, even though they share interior data.  A union bound
therefore gives its probability at least $(1-2p_a)_+$.  This proves
\eqref{f16v71:native-midpoint-bound}.  Markov's inequality for
$1-w(P)$ and V64's stated native expectation bound give
\eqref{f16v71:midpoint-pressure} in exactly that class of measures.
\end{proof}

The numerical constant is small and no practical sampling claim is made.
The conclusion is qualitative but uniform in the coefficient: local
endpoint flatness in probability does not imply flatness in probability
of independent-representative midpoints.  The gauge-averaging part holds
for any native gauge-invariant slice law; the additional pressure estimate
is not asserted for an arbitrary boundary-conditioned state.  Subsequent
flow can remove a midpoint defect, so this theorem alone neither bounds
nor invalidates V70's integrated local scalar kernel.

\subsection{Anchored tree alignment preserves the exact conditional source difference}

Let the finite spatial graph be connected.  Assume the retained edges $C$
form a connected subgraph on vertices $V_C$; the region may be enlarged
and redeclared when necessary.  Choose a root $o\in V_C$, a spanning tree
$T_C$ of this subgraph, and a spanning tree $T$ of the whole graph extending
$T_C$.  In particular no added tree edge joins two vertices already in
$V_C$.  For a tree path from $o$ to $x$ write its ordered parallel transport
as $A_U(x)$.  A chord $e=(x,y)\notin T$ has rooted holonomy
\begin{equation}
 H_e(U)=A_U(x)U_eA_U(y)^{-1}.
 \label{f16v71:rooted-holonomy}
\end{equation}
All formulas on a reversed tree edge use the inverse group value.

\begin{theorem}[Exact replica alignment and its holonomy kernel]
\label{f16v71:tree}
For replicas agreeing on $C$, there is a unique vertex gauge $g$ with
$g_o=I$ such that $\bar U'=g\cdot U'$ agrees with $U$ on $T$.  It is given
recursively along a tree edge $x\to y$ by
\begin{equation}
 g_y=U_e^{-1}g_xU'_e.
 \label{f16v71:recursion}
\end{equation}
It obeys $g_x=I$ on $V_C$ and $\bar U'_e=U_e$ on every retained edge.
The aligned geodesic speeds satisfy
\begin{equation}
 \boxed{\quad
 \ell_e:=d(U_e,\bar U'_e)=
 \begin{cases}
 0,&e\in T\text{ or }e\in C,\\
 d(H_e(U),H_e(U')),&e\notin T.
 \end{cases}\quad}
 \label{f16v71:speeds}
\end{equation}
All speeds vanish precisely when the endpoints are gauge equivalent by
a gauge equal to identity on $V_C$.  For any gauge-invariant finite source
family $F$, including the spatially flowed family of V69,
\begin{equation}
 \begin{aligned}
 F(U')&=F(\bar U'),\\
 \mathsf E_C&:=\mathbb E[(F-\mathbb E(F\mid U_C))
                      (F-\mathbb E(F\mid U_C))^*]\\
 &=\tfrac12\mathbb E_{\rm rep}
       [(F(U)-F(\bar U'))(F(U)-F(\bar U'))^*].
 \end{aligned}
 \label{f16v71:exact-replica}
\end{equation}
The construction is equivariant under a common gauge transformation of
the replicas.  No minimization, gauge Laplacian inverse, or stabilizer
removal is required.
\end{theorem}
\begin{proof}
The tree recursion is unique and solves
$g_xU'_eg_y^{-1}=U_e$ on each tree edge.  The shared values on $T_C$
force $g_x=I$ throughout $V_C$, so every shared edge remains fixed.
Tree transports of the aligned replica are exactly $A_U(x)$.  Rooted
holonomies transform by conjugation at $o$, and $g_o=I$; hence
\[
 A_U(x)\bar U'_eA_U(y)^{-1}=H_e(U').
\]
Bi-invariance of the metric proves \eqref{f16v71:speeds}.  If all speeds
vanish, the constructed anchored gauge maps $U'$ to $U$.  Conversely, any
such gauge satisfies the same tree recursion and thus is the constructed
one.  Gauge invariance proves equality of the source values, after which
the original conditional-replica variance identity gives
\eqref{f16v71:exact-replica}.  The aligned replicas need not be conditionally
independent; that is not used after the equality of source values.

If both endpoints are transformed by $h$, the recursion produces
$g_x^{\rm new}=h_xg_xh_x^{-1}$.  This verifies equivariance directly.
The tree map is smooth.  Chordwise shortest geodesics are measurable and
unique away from antipodal pairs, a null set under the finite positive
conditional density.  Constant tree edges cause no ambiguity.
\end{proof}

\begin{proposition}[The conditional Haar coordinates and their induced law]
\label{f16v71:Haar}
Fix the first replica $U$ and the shared data on $C$.  Under the conditional
native law of $U'$, the alignment coordinates consist of vertex gauges
$(g_x)_{x\notin V_C}$ and aligned free chords.  The vertex gauges are
independent normalized Haar variables, independent of those chords.
The chord density is the exact original conditional density evaluated
with its tree edges set to $U_T$.  No product chord law or new Wilson
coupling is inferred.
\end{proposition}
\begin{proof}
There are exactly $|V|-|V_C|$ free tree edges.  Recursive Haar translation
changes those variables bijectively to the same number of vertex gauges,
with unit Haar Jacobian.  Conditional on these gauges, changing every
free chord by $U'_e\mapsto g_xU'_eg_y^{-1}$ is a left-right Haar
translation.  Product Haar measure thus factors into the gauge coordinates
and the aligned chord measures.  Since $g$ is identity on $V_C$, native
gauge invariance leaves the conditional density unchanged and makes it
independent of the gauge coordinates.  Their normalized integral is one.
This proves the factorization, including its normalization.  Interactions
among the chords, the shared boundary data and all variables previously
integrated into the slice density remain in that density.  This is not
an assertion that the original slice marginal had a local product action.
\end{proof}

The two replicas may therefore be aligned without a gauge-volume penalty
or a discarded change-of-variables determinant.  A \emph{common} gauge
transformation of both endpoints, unlike this relative alignment, would
not change their original linkwise geodesic distances or cure
Proposition~\ref{f16v71:flat-artifact}.

\subsection{Local propagation charged to aligned holonomy discrepancies}

Let $A$ be the output-gradient support of V69's $O_i$, and take
$C=A_R$ with the connectedness just stated.  For each original replica
pair, form the aligned shortest-link path $Z^\theta$ from $U$ to
$\bar U'$.  Its speeds are $\ell_e$ from \eqref{f16v71:speeds}; in particular
it is constant on $C$.  Let $\mathsf K_T^\theta(\tau,0)$ be V70's local
scalar propagator along the actual spatial trajectory
$\Phi_{a,t}(Z^\theta)$.  Set
\begin{equation}
 L_T(e;U,U',\theta)=(\mathsf K_T^\theta(\tau,0)\ell)(e),\qquad
 \Theta_T=a^3\sum_{e\in A}L_T(e;U,U',\theta)^2.
 \label{f16v71:aligned-tail}
\end{equation}
The rates, local potential, complete escape rates, and reversed operator
time are exactly those of \eqref{f16v70:rates}--\eqref{f16v70:FK-formula}.
In particular the environment here is the \emph{aligned} interpolation,
not the environment of the old geodesic.

\begin{theorem}[Native gauge-aligned source localization]
\label{f16v71:localization}
For the unchanged native source and predictor,
\begin{equation}
 \boxed{\quad
 0\preceq\mathsf E_C\preceq
 \frac12\mathbb E_{\rm rep}\int_0^1
 \Theta_T(U,U',\theta)\,
 \mathsf Q(\Phi_{a,\tau}Z^\theta)\,d\theta .\quad}
 \label{f16v71:aligned-localization}
\end{equation}
This is a simultaneous inequality in every source direction.  All declared
exact source relations are retained.  For an anchored gauge-equivalent
pair the complete integrand is zero, rather than a positive cost for a
pure-gauge change.  On a simply connected box this includes all pairs of
flat endpoints agreeing on $C$.
\end{theorem}
\begin{proof}
For a fixed pair, its alignment gauge is fixed as $\theta$ varies.  The
initial interpolation tangent has on-link norm $\ell_e$.  Gauge completion
of this tangent leaves the output derivative unchanged by
Theorem~\ref{f16v70:gauge-completion}.  V70's scalar domination gives
$|\widetilde v_e(\tau)|\le L_T(e)$, with the actual aligned trajectory
retained.  For every real coefficient vector $c$, Cauchy--Schwarz in the
$a^3$-weighted tangent metric gives
\[
 \left|\frac d{d\theta}c^*O(\Phi_{a,\tau}Z^\theta)\right|^2
 \le\Theta_T\,c^*\mathsf Q(\Phi_{a,\tau}Z^\theta)c.
\]
Integrating and using Cauchy--Schwarz on $[0,1]$, followed by the exact
identity \eqref{f16v71:exact-replica}, proves the matrix assertion.
Complex coefficients follow by complexification.  Gauge covariance of
the construction and linearity in the source preserve the claimed
relations.  An anchored gauge-equivalent pair has $\ell=0$, so its bound
is zero without invoking any bound on its propagator.  On a simply
connected box a flat connection has identity holonomy on every
contractible loop; the tree holonomies are consequently all identity.
Theorem~\ref{f16v71:tree} applies.
\end{proof}

This bound is not claimed smaller than V70's in every curved pair.  The
alignment can lengthen some chord displacements or produce a different
curved interpolation.  Its rigorous gain is exact removal of independent
pure-gauge displacement.  No product estimate of a holonomy moment and
an occupation moment is justified: they remain correlated inside
\eqref{f16v71:aligned-localization}.

\begin{proposition}[Exact signed source response before scalar domination]
\label{f16v71:signed}
Let $\mathsf J_d^\theta$ be the gauge-completed flow differential along
$Z^\theta$, and define its output-directed covectors
\[
 G_i^\theta=(\mathsf J_d^\theta(\tau,0))^*
            \nabla_a O_i(\Phi_{a,\tau}Z^\theta).
\]
Then $G_i^\theta=\nabla_a(O_i\circ\Phi_{a,\tau})(Z^\theta)$ and
$D_{Z^\theta}^*G_i^\theta=0$.  Write $\omega_e^\theta$ for the right
velocity of the aligned geodesic and
\[
 q_i(\theta)=a^3\sum_{e\notin T\cup C}
                  G_{i,e}^\theta\cdot\omega_e^\theta.
\]
For the vector $q=(q_i)_i$, exactly,
\begin{equation}
 \mathsf E_C=\frac12\mathbb E_{\rm rep}
       \left[\left(\int_0^1q(\theta)d\theta\right)
             \left(\int_0^1q(\theta)d\theta\right)^*\right]
 \preceq\frac12\mathbb E_{\rm rep}\int_0^1q(\theta)q(\theta)^*d\theta.
 \label{f16v71:signed-source}
\end{equation}
\end{proposition}
\begin{proof}
V70's source-equivalence identity holds for every initial tangent, so its
adjoint gives equality of the two covectors.  Differentiating gauge
invariance gives their orthogonality to the initial vertical map.
The chain rule and the zero tree and retained-edge velocities give
$q_i(\theta)$.  The fundamental theorem of calculus gives the source
difference, whose sign is immaterial in its Gram.  Use
\eqref{f16v71:exact-replica} and Jensen's inequality.  All expressions
use one original replica law; no sign or cancellation is lost until the
last displayed inequality.
\end{proof}

The exact form \eqref{f16v71:signed-source} is the alternative when an
absolute local kernel remains too large.  It is not a proof that the
actual signed expectation is small.  Nor does horizontality remove
physical curvature fluctuations or imply a norm bound on that expectation.

\subsection{Paid curvature distances, explicit geometry and the winding exception}

\begin{proposition}[Native displacement control for a declared contractible chord]
\label{f16v71:surface}
Suppose the rooted chord loop admits an explicit factorization into $m_e$
conjugated plaquette holonomies or inverses, counting multiplicities.
For the aligned speed,
\begin{equation}
 \ell_e^2\le\min\left\{\pi^2,
 2m_e\sum_{j=1}^{m_e}
 [\rho(P_j(U))^2+\rho(P_j(U'))^2]\right\}.
 \label{f16v71:surface-speed}
\end{equation}
If the relevant native plaquettes obey
$\mathbb E_{\pi_a}(1-w(P_j))\le b_a$, then
\begin{equation}
 \boxed{\quad
 \mathbb E_{\rm rep}\ell_e^2
 \le\min\{\pi^2,\,2\pi^2m_e^2b_a\}.\quad}
 \label{f16v71:native-speed}
\end{equation}
For V64's symmetric native measures one may take
$b_a=B_{\gamma_a}/\gamma_a$.  Replica independence away from the shared
interior is not assumed.
\end{proposition}
\begin{proof}
Bi-invariance and the triangle inequality bound the distance of each
rooted holonomy from identity by the sum of the plaquette distances in
its factorization.  Conjugations and inverses do not change these
distances.  Apply the triangle inequality between the two rooted
holonomies, followed by Cauchy--Schwarz to their $2m_e$ terms.  This
proves \eqref{f16v71:surface-speed}.  The two endpoint marginals are both
$\pi_a$.  Equation~\eqref{f16v71:metric} bounds each expected squared
plaquette distance by $\pi^2 b_a/2$, giving
\eqref{f16v71:native-speed}.  The distance is always at most $\pi$.
\end{proof}

There is a useful explicit geometric example, without concealing an area
factor.  On the open cube with vertices $0,\ldots,L$ in each coordinate,
use the comb tree whose root path first moves in direction 1, then 2,
then 3.  Every direction-3 link is in the tree.  A direction-2 chord at
height $x_3$ bounds a width-one $23$ strip with $m_e=x_3$ plaquettes.
A direction-1 chord bounds the union of a width-one $13$ strip at its
$x_2$ coordinate and a width-one $12$ strip in the bottom plane, with
$m_e=x_2+x_3\le2L$.  These factorizations are exact: peeling one strip
plaquette at a time conjugates its holonomy by the peeled boundary path
and shortens the loop; after the strips are peeled the remaining tree
path retraces itself.  The order of the conjugated factors is retained;
non-Abelian holonomies are not replaced by sums.

Thus for this specified geometry,
\begin{equation}
 \mathbb E_{\rm rep}\ell_e^2
 \le\min\{\pi^2,8\pi^2L^2B_\gamma/\gamma\}.
 \label{f16v71:comb-cost}
\end{equation}
A tree required to extend a larger retained interior must have its own
filling counts checked.  The comb estimate is not assigned automatically
to that different tree.  Even for the comb, a fixed physical width has
$L=O(a^{-1})$; under an exponential $a(\gamma)$ the displayed bound need
not tend to zero.  It proves vanishing displacement for a fixed microscopic
loop, not for an unspecified growing physical support.  Moreover
\eqref{f16v71:native-speed} alone cannot be multiplied by a separate
average of the source-marked propagator in
\eqref{f16v71:aligned-localization}.

\begin{proposition}[The fixed-patch midpoint artifact is removed under the native law]
\label{f16v71:aligned-curvature}
For a marked plaquette $P$, assume every nonconstant aligned boundary
edge has the disk factorization of Proposition~\ref{f16v71:surface}.
Give constant tree or retained edges the count $m_e=0$.  Under its same
native bound $b_a$, the complete aligned interpolation satisfies
\begin{equation}
 \boxed{\quad
 \mathbb E_{\rm rep}\sup_{0\le\theta\le1}\rho(P(Z^\theta))^2
 \le\frac{5\pi^2}{2}\left(1+4\sum_{e\in\partial P}m_e^2\right)b_a.
 \quad}
 \label{f16v71:aligned-curvature-bound}
\end{equation}
Consequently the probability that this path ever has
$w(P(Z^\theta))\le-3/4$ is at most the right side divided by
$[\arccos(-3/4)]^2$, capped at one.  For fixed filling counts and
$b_a\to0$, this probability tends to zero, unlike the unaligned
lower bound \eqref{f16v71:native-midpoint-bound}.
\end{proposition}
\begin{proof}
Telescoping the four unit factors in the bi-invariant metric gives
\[
 \rho(P(Z^\theta))\le\rho(P(U))+\sum_{e\in\partial P}\ell_e
 \quad\hbox{for every }\theta.
\]
The square of a sum of five nonnegative terms is at most five times
the sum of their squares.  Use \eqref{f16v71:metric} for the first
term and \eqref{f16v71:native-speed} for the other four.  This gives
\eqref{f16v71:aligned-curvature-bound} without an independence
assumption or a time discretization of the supremum.  Markov's inequality
gives the event estimate.
\end{proof}

A fixed microscopic comb patch can be attached to the anchored tree
by a single branch, with no retained vertices inside that patch.  Its
root stem only conjugates the disk holonomies and adds no plaquette
filling cost.  Thus the last comparison can be made at the same marked
exterior square as Theorem~\ref{f16v71:native-midpoint}, independently
of the surrounding volume.  This does not prove a bound uniform in
a growing patch, over a positive spatial-flow interval, or after
multiplication by the source-dependent propagation weight.

\begin{proposition}[Topology is not removed by the alignment]
\label{f16v71:winding}
On an $L^3$ torus put $U=I$ and let $U'$ have constant direction-1 links
$e^{\alpha\mathbf i/L}$ and all other links identity, where
$0<\alpha<\pi$.  Both configurations are flat.  No anchored tree
alignment can make all their chord discrepancies zero: a winding loop
has holonomies $I$ and $e^{\alpha\mathbf i}$.  Therefore a bound by
plaquette defects alone is false for general torus chords.
\end{proposition}
\begin{proof}
All links commute and are spatially constant, so every plaquette is
identity.  The ordered product around the direction-1 cycle has the
stated values.  A gauge transformation conjugates a rooted cycle
holonomy and cannot change its trace.  Hence the two configurations
are not gauge equivalent.  The exact kernel statement of
Theorem~\ref{f16v71:tree} proves that some discrepancy survives.  The
noncontractible loop has no plaquette disk factorization of the kind
assumed in Proposition~\ref{f16v71:surface}.
\end{proof}

This example is only a scope check on the alignment, not another analysis
of the archived toron dynamics.  Contractible curvature and winding
holonomies must both be retained when the chosen graph has cycles not
filled by its plaquettes.

\subsection{The native reflected target and the noncircular boundary}

The error matrix and predictor on the left of
\eqref{f16v71:aligned-localization} are exactly those of V69--V70.
Consequently the inherited positive-time comparison still reads
\[
 \|H-H_C\|_{\rm op}\le2\sqrt{\|D\|_{\rm op}\|\mathsf E_C\|_{\rm op}}
\]
for their specified common temporal placements.  Alignment changes the
coupling used to estimate $\mathsf E_C$, not the predictor, the physical
transfer, the source normalization, or the probability measure defining
$D$ and $H$.  A relative bound on $H$ still needs a reflected lower Gram.

The main upstream conclusion is not just the existence of one bad path.
Theorem~\ref{f16v71:native-midpoint} proves that independent local gauges
supply a persistent source of midpoint center defects under the native
replica law even when endpoint plaquettes become flat in probability.
Asking native endpoint estimates to pay those defects would estimate a
gauge-coordinate artifact.  Anchored alignment removes that artifact
exactly, including its entire propagation cost on anchored gauge-equivalent
pairs.  The induced holonomy law, true curved interpolation, filling-size
costs and winding variables are not removed.

The next quantitative target is a source-directed bound on
\eqref{f16v71:signed-source}, or its aligned local domination
\eqref{f16v71:aligned-localization}, for a declared physical region.  A
change to another scalar envelope without estimating its joint native
law would repeat the same interface.  Equally, none of the upper bounds
here substitutes for a direct nonzero positive-separation covariance.
A noncollapsing reflected family, physical-time contraction with common
dense-core constants and the interacting continuum mass gap remain
unproved.

\subsection{Verification, append-only preservation and result ledger}

The companion diagnostic is
\begin{center}\small
\path{verify_ym_f16_v71_gauge_aligned_replicas.py}.
\end{center}
It separates exact quaternion and finite-group fixtures, literal lattice
tree and gauge identities, midpoint robustness, full compact source-flow
checks, signed derivative comparisons and scalar-propagator diagnostics.
Finite-group enumeration is a distinct algebraic test, not a replacement
for the continuous Wilson measure.  No numerical integration is used as
a premise of the proofs.  These checks are neither a native four-dimensional
vacuum simulation nor an outward-rounded interval certificate or
proof-assistant verification.  Removing the delimited appendix and
reverting the single displayed version change recovers V70 byte for byte.

\clearpage
\begin{record}[V71 result ledger]
\label{f16v71:ledger}
\emph{Proved here:} an actual flat-endpoint interpolation with a center
midpoint; a native conditional-replica lower bound for that gauge-generated
midpoint event under local endpoint concentration; anchored replica
alignment with an exact holonomy kernel and Haar-coordinate factorization;
a source-marked native localization inequality charging holonomy speeds;
an exact horizontal signed source-response identity; native contractible
holonomy-displacement estimates with filling counts; a uniform-in-interpolation
native curvature bound removing the fixed-patch artifact; an explicit comb-tree
count and a winding obstruction to plaquette-only control.

\emph{Inherited or established background:} native gauge invariance and
V64's plaquette expectation estimate in its stated class; the unchanged
source and spatial flow; V70's gauge-completed response and local scalar
propagator; the general use of tree gauge coordinates.  Earlier tree or
toron analyses are not claimed as new and the older archive is not
recertified.

\emph{Not asserted:} a Wilson law for unaligned or aligned interpolants;
rarity of their defects from endpoint estimates; independence of the
aligned replicas or their chord holonomies; domination of the old bound
on every curved pair; a cutoff-uniform estimate from a microscopic filling
count; elimination of winding data; or a lower reflected Gram from source
localization.

\emph{Still unproved:} the physical-normalization bound on the joint
aligned native response, a noncollapsing positive-separation reflected
source family, common dense-core physical-time contraction, and the
interacting four-dimensional continuum Yang--Mills mass gap.
\end{record}

\begin{firewall}
The replica coupling is changed only by a symmetry that leaves its source
difference identical.  Its gauge coordinates are paid exactly, while its
induced holonomy law and source-directed propagation remain explicit.
Artificial midpoint curvature is separated from actual native curvature.
V1--V70 remain unchanged.
\end{firewall}
% END F16 V71 APPEND

% BEGIN F16 V72 APPEND -- 2026-09-19
% Append-only continuation of the Post-Fifteenth-Archive V71.
\clearpage
\section[V72: native thin-loop screening and shared-path singlets]{V72: native thin-loop screening\\and shared-path singlets}
\label{f16v72:context}

V71 removes independent gauge displacement from the replica interpolation,
but leaves its long rooted holonomies under the actual native law.  Before
asking their discrepancies to become small on a physical region, we compute
an unconditional obstruction: an unflowed rectangular holonomy tends to
Haar measure when its perimeter in lattice units grows faster than the
Wilson coefficient.  This holds uniformly conditional on the unselected
links.  The resulting aligned chord discrepancy has a nonzero Haar limit,
even when individual plaquettes concentrate near identity.  This is actual
holonomy noise, not V71's pure-gauge midpoint artifact.

The same calculation gives a direct native reflected-covariance estimate.
A bounded or polynomially normalized source depending on one such thin
loop has vanishing positive-separation reflected norm, although its
ordinary variance need not vanish.  We do not transfer this statement to
the spatially flowed source.  In particular, products of loops sharing a
long path have a singlet contribution which must be extracted before any
perimeter estimate.  For two fundamental loops we compute that contribution
and bound its complete native two-time covariance error.  This supplies a
signed source reduction, rather than another unevaluated scalar propagation
moment.

\subsection{Scope, ownership and the native conditional law}

Perimeter-law upper bounds are established results; see B.~Simon and
L.~G.~Yaffe, \emph{Rigorous perimeter law upper bound on Wilson loops},
Phys.\ Lett.\ B \textbf{115} (1982), 145--147,
\href{https://doi.org/10.1016/0370-2693(82)90815-2}{doi:10.1016/0370-2693(82)90815-2}.
We do not claim discovery of perimeter decay.  What is needed here is its
explicit conditional Fourier form, its consequences for the V71 replica
law and the native reflected source quotient, and the shared-path singlet
which survives the same conditional integration.  Those specialized
statements are proved below.  V58's independent-link geometry and V63's
one-link character coefficients are inherited, not counted again as new
methods.  No Gaussian comparison, spatial-flow moment, or continuum clock
calibration is assumed.

Use the isotropic four-dimensional \(\SU(2)\) Wilson law with coefficient
\(\gamma>0\), normalized Haar links and local boundary conditions.  All
selected links are interior and each belongs to at most six plaquettes.
The statements apply on nondegenerate finite tori, under conditional Gibbs
laws with the selected links free, and by the conditional expectation tower
to their slice marginals.  Stationary time statements use the actual
vacuum history and its positive self-adjoint contraction \(P_a\).  They
follow by the fixed-spatial-regulator cylinder limit with boundary data
explicit; no nonlocal Perron endpoint weight is absorbed into a local link
factor.  Bounds for bounded local observables pass to the corresponding
Gibbs limits.

Call a set \(S\) of stored links \emph{plaquette-separated} when no elementary
plaquette contains two of its members.  Conditional on every link outside
\(S\), the selected links are independent with densities
\begin{equation}
 d\nu_{h_e}(x_e)=\frac{e^{h_e\cdot x_e}}{\mathfrak z(|h_e|)}d\sigma(x_e),
 \qquad |h_e|\le6\gamma,\qquad
 \mathfrak z(t)=2I_1(t)/t.
 \label{f16v72:conditional}
\end{equation}
Here \(h_e\) is the actual sum of the complementary Wilson staple vectors,
including their orientations and the coefficient \(\gamma\).  Independence
is asserted only after conditioning; the unselected links retain their
full native distribution.

Let \(D^{(k)}\) be the unitary irreducible representation of dimension
\(d_k=k+1\), with \(\chi_k=\Tr D^{(k)}\).  Write
\begin{equation}
 r_\gamma=\frac{I_2(6\gamma)}{I_1(6\gamma)},\qquad
 q_{\gamma,m}=r_\gamma^m,\qquad
 \mathcal S(z)=\sum_{k=1}^{\infty}(k+1)^2z^k
             =\frac{z(4-3z+z^2)}{(1-z)^3}.
 \label{f16v72:constants}
\end{equation}
The notation \(q\) below means \(q_{\gamma,m}\), not a jump generator.

\subsection{Coefficient-explicit conditional loss of nontrivial representations}

\begin{proposition}[One-link bounds for every representation]
\label{f16v72:one-link}
For \(t>0\) set \(\lambda_k(t)=I_{k+1}(t)/I_1(t)\), with
\(\lambda_0(0)=1\) and \(\lambda_k(0)=0\) for \(k\ge1\).  Then
\begin{equation}
 0\le\lambda_k(t)\le\lambda_1(t)^k,
 \qquad
 r_\gamma\le\exp\left[-\operatorname{arsinh}\frac1{4\gamma}\right]<1.
 \label{f16v72:ratio-upper}
\end{equation}
If \(P=h/|h|\) when \(h\ne0\), the mean of \(D^{(k)}(X)\) under
\eqref{f16v72:conditional} is
\(\lambda_k(|h|)D^{(k)}(P)\).  Inversion or deterministic left and right
multiplication leaves its operator norm equal to \(\lambda_k(|h|)\).
\end{proposition}
\begin{proof}
The last identity is V63's one-link central character convolution, followed
by Haar translation.  It applies to each selected link before any other
variable is marginalized.  V63 also proves monotonicity of \(\lambda_k\) in
its nonnegative argument.  We provide the additional uniform-in-\(k\) upper
bound.  Put \(R_\nu=I_{\nu+1}/I_\nu\).  Its recurrence gives
\[
 R_\nu'=1-R_\nu^2-(2\nu+1)R_\nu/t,\qquad
 R_\nu(t)=t/(2\nu+2)+O(t^3).
\]
For \(\nu>1\), an upward first crossing of \(R_\nu\) through \(R_1\) is
impossible: at a crossing their derivative difference would be
\(-2(\nu-1)R_1/t<0\).  Thus \(R_\nu\le R_1\), and
\(\lambda_k=\prod_{\nu=1}^kR_\nu\le R_1^k\).

For the mean \(r=R_1(t)\) of the scalar coordinate in the tilted \(S^3\)
law, integration by parts in that coordinate gives
\(1-\mathbb E[(x^0)^2]=3r/t\).  Since \(\mathbb E[(x^0)^2]\ge r^2\),
\[
 r^2+3r/t\le1,\qquad
 r\le\sqrt{1+(3/(2t))^2}-3/(2t)
      =e^{-\operatorname{arsinh}(3/(2t))}.
\]
Take \(t=6\gamma\).  At a zero field all nontrivial means vanish and no
choice of its direction is needed.
\end{proof}

A path word below contains each selected stored link exactly once, possibly
with reversed orientation.  Other links may repeat.  In particular a
rooted loop may have a retraced root stem, but its selected links must lie
in the once-traversed part and not in that stem.

\begin{theorem}[Native conditional holonomy approaches Haar]
\label{f16v72:Haar}
Let \(H\) be such a path word with \(m\ge1\) plaquette-separated selected
links, and let \(\Xi\) be the sigma-field of all unselected links.  The
conditional law of \(H\) has density \(p_\Xi\) relative to normalized Haar
measure, and
\begin{align}
 \left\|\mathbb E[D^{(k)}(H)\mid\Xi]\right\|_{\rm op}&\le q^k,
       && k\ge1, \label{f16v72:matrix-Fourier}\\
 \|p_\Xi-1\|_\infty&\le\mathcal S(q)=:\varepsilon_q,
 &\quad \|p_\Xi-1\|_2^2&\le\mathcal S(q^2)=:\eta_q.
 \label{f16v72:density}
\end{align}
The same bounds hold conditional on any sub-sigma-field of \(\Xi\), with
its induced native law.  In particular, for every Haar-centered
\(f\in L^2(H)\),
\begin{equation}
 \left|\mathbb E[f(H)\mid\Xi]\right|^2
       \le\eta_q\|f\|_{L^2(H)}^2.
 \label{f16v72:conditional-source}
\end{equation}
All constants are independent of the exterior and the surrounding volume.
\end{theorem}
\begin{proof}
After conditioning, \(H\) is an ordered product of deterministic group
factors and independent selected variables.  Each selected variable is a
fixed center times a central random variable with density \(k_{|h_e|}\).
Move the centers through the word by conjugating the independent central
variables.  This leaves their joint law unchanged.  Consequently the
conditional law is a translate of
\(\mathop{*}_{e\in S}k_{|h_e|}\).  Equivalently its \(k\)-th matrix mean is
a unitary matrix times \(\prod_e\lambda_k(|h_e|)\), which is at most
\(r_\gamma^{km}\) by Proposition~\ref{f16v72:one-link}.  The same reasoning
handles reversed occurrences.  If any field is zero the product law is
exactly Haar.

Fourier inversion has absolute majorant
\(\sum_{k\ge1}d_k^2q^k\); it converges since \(q<1\).  The trivial
representation is the constant one.  This proves the first density bound
and identifies the actual continuous density by uniqueness of Fourier
coefficients.  Plancherel gives
\[
 \|p_\Xi-1\|_2^2
 =\sum_{k\ge1}d_k\|\widehat p_\Xi(k)\|_{\rm HS}^2
 \le\sum_{k\ge1}d_k^2q^{2k}.
\]
A conditional mixture preserves the Fourier norm bound and both density
bounds, using Jensen for the squared \(L^2\) norm.  Finally integrate
\(f(p_\Xi-1)\) and use Cauchy--Schwarz.  This proof never assigns a product
law to the unselected variables.
\end{proof}

\begin{proposition}[All but four edges of a rectangle can be selected]
\label{f16v72:rectangle}
A nonwrapping spatial rectangle of side lengths \(L_1,L_2\ge2\) in lattice
units has a plaquette-separated subset of its boundary of size
\begin{equation}
 m=2(L_1+L_2)-4=|\partial R|-4.
 \label{f16v72:rectangle-m}
\end{equation}
On a torus assume also that the complementary distances in the two
rectangle directions are at least two lattice steps.  The subset can be
used simultaneously with its time translate by \(n\ge2\) steps: no
plaquette meets both selected sets.
\end{proposition}
\begin{proof}
Select every edge of the two horizontal sides.  On each vertical side
select the edges other than its first and last edge.  Consecutive
collinear edges are never two edges of one elementary plaquette.  The two
parallel sides are separated by at least two steps, also through the
periodic complement.  A selected vertical edge cannot share a corner
plaquette with a horizontal side, precisely because its two endpoint
edges were omitted.  The same separation applies to the two vertical
sides.  Temporal or transverse spatial plaquettes introduce no second
selected edge.  The count is \(2L_1+2(L_2-2)\).  An elementary temporal
plaquette spans only one time step, proving the last assertion.  The
omitted four edges are a geometric choice, not four discarded action
terms; their Wilson factors remain in \(\Xi\).
\end{proof}

For a fixed positive physical rectangle perimeter \(\mathcal P\),
\(m_a a\to\mathcal P\).  Thus \(q_a\to0\) whenever
\(m_a\operatorname{arsinh}(1/(4\gamma_a))\to\infty\); in particular
\(a\gamma_a\to0\) suffices if \(\gamma_a\to\infty\).  Under a separately
specified trajectory
\begin{equation}
 \log(1/a)=c\gamma_a+O(\log\gamma_a),\qquad c>0,
 \label{f16v72:clock}
\end{equation}
\(q_a\) is smaller than every fixed power of \(a\).  This is only an
application under that trajectory, not a derivation of it.  The decay is
in perimeter, not a confining area-law assertion.

\subsection{The actual aligned replica discrepancy does not vanish}

\begin{theorem}[A Haar limit for long aligned chord speeds]
\label{f16v72:replica}
Let \(H,H'\) be the same rooted loop functional of the two native replicas
conditioned on retained data \(C\).  Suppose its selected links avoid \(C\)
and the retraced root stem.  Under the original conditional-replica law,
\begin{equation}
 \left\|\frac{d\,\operatorname{Law}(H^{-1}H'\mid C)}{dH}-1\right\|_\infty
       \le\eta_q=\mathcal S(q^2).
 \label{f16v72:replica-Haar}
\end{equation}
If \(H\) is a fundamental chord holonomy of the anchored tree in V71,
its aligned speed \(\ell=d(H,H')\) therefore satisfies
\begin{equation}
 \boxed{\quad
 \left|\mathbb E_{\rm rep}\ell^2-
             \left(\frac{\pi^2}{3}-\frac12\right)\right|
 \le\eta_q\left(\frac{\pi^2}{3}-\frac12\right).
 \quad}
 \label{f16v72:replica-speed}
\end{equation}
In particular its mean-square displacement tends to
\(\pi^2/3-1/2>0\) when \(q\to0\), including for a physical-size
rectangular chord under \eqref{f16v72:clock}.
\end{theorem}
\begin{proof}
Before the alignment, the two holonomies are independent conditional on
\(C\).  Let \(A_k=\mathbb E[D^{(k)}(H)\mid C]\).  The previous theorem
and the conditional expectation tower give \(\|A_k\|\le q^k\).  Hence
\[
 \mathbb E[D^{(k)}(H^{-1}H')\mid C]=A_k^*A_k,
 \qquad \|A_k^*A_k\|\le q^{2k}.
\]
Fourier inversion proves \eqref{f16v72:replica-Haar}.  V71's exact tree
identity equates the aligned link speed with \(d(H,H')\); no independence
of the aligned links is asserted or needed.  The radial Haar law is
\((2/\pi)\sin^2\theta\,d\theta\), and direct integration gives
\[
 \frac2\pi\int_0^\pi\theta^2\sin^2\theta\,d\theta
       =\frac{\pi^2}{3}-\frac12.
\]
Multiply the density error by \(\theta^2\) and integrate.  A tree containing
all but one boundary edge of a disjoint rectangle, joined to the retained
interior by a root stem, supplies the stated literal example.
\end{proof}

This result is compatible with V71's fixed microscopic filling estimate.
That estimate cannot be extended to physical-size thin chords by assuming
that its growth factor is merely wasteful: the discrepancy has an actual
nonzero limit.  The corresponding endpoint plaquettes may still tend to
identity under V64's native estimate.  Here the nonvanishing quantity is
gauge-invariant holonomy data, unlike the midpoint artifact that V71
removed.  A different tree can redistribute these data among chords, but
this theorem concerns the declared tree and does not assert an optimal
orbit-distance lower bound.

\subsection{A direct native reflected-covariance theorem for a thin-loop bank}

Let \(f=(f_1,\ldots,f_d)\) be real Haar-centered class functions of one loop
holonomy, square-integrable with respect to Haar measure.  They may depend
on the regulator.  Set
\begin{equation}
 G_H=\int f f^*dH,\quad
 F_t=f(H_t)-\mathbb E_{\pi_a}f(H_t),\quad
 G_a=\mathbb E F_0F_0^*,\quad C_{a,n}=\mathbb E F_0F_n^*.
 \label{f16v72:source-bank}
\end{equation}
All coefficient statements are on the exact quotient of \(G_H\).
The loops at different times are translates of the same unflowed spatial
loop, with disjoint selected sets as in Proposition~\ref{f16v72:rectangle}
when their time separation is at least two.

\begin{theorem}[Full native reflected screening, not a conditional-flat calculation]
\label{f16v72:reflected}
Put \(b_q=1-\varepsilon_q-\eta_q\).  Then
\begin{equation}
 b_qG_H\preceq G_a\preceq(1+\varepsilon_q)G_H,
 \qquad
 \boxed{\quad 0\preceq C_{a,n}\preceq\eta_qG_H\quad(n\ge2).\quad}
 \label{f16v72:Gram-screen}
\end{equation}
The source kernel equals the Haar kernel exactly at every finite regulator,
including when the displayed lower coefficient \(b_q\) is not positive.
When \(b_q>0\), the reflected Gram of a source placed \(k\ge1\) steps
beyond the reflection plane obeys
\begin{equation}
 0\preceq G_a^{-1/2}C_{a,2k}G_a^{-1/2}
       \preceq\frac{\eta_q}{b_q}I.
 \label{f16v72:relative-reflection}
\end{equation}
The inverse is taken only on the exact source quotient.  No conditional
tube transfer or auxiliary spatial propagator appears in this estimate.
\end{theorem}
\begin{proof}
For \(g=c^*f\), the density bound gives
\[
 (1-\varepsilon_q)\|g\|_H^2
 \le\mathbb E|g(H)|^2\le(1+\varepsilon_q)\|g\|_H^2,
 \qquad |\mathbb E g(H)|^2\le\eta_q\|g\|_H^2.
\]
Subtracting the mean proves the entrance inequalities.  The finite-regulator
holonomy density is strictly positive: one selected conditional density is
strictly positive and compact convolution and mixture preserve positivity.
Thus a source combination has native zero variance precisely when it is
constant Haar-almost everywhere.  Haar centering makes that constant zero.
This proves the exact kernel statement without a numerical threshold.

For the two-time comparison condition on the complement \(\Xi_{0,n}\) of
the union of the two selected sets.  No plaquette contains two selected
links, even across the two times.  Consequently \(H_0,H_n\) are
conditionally independent and, for \(a_j=\mathbb E[g(H_j)\mid\Xi_{0,n}]\),
\[
 \operatorname{Cov}(g(H_0),g(H_n))=\operatorname{Cov}(a_0,a_n),
 \qquad \mathbb E|a_j|^2\le\eta_q\|g\|_H^2.
\]
Cauchy--Schwarz bounds the absolute covariance by \(\eta_q\|g\|_H^2\).
The stationary native transfer has \(P_a\succeq0\), so the matching
covariance is \(\langle g-\mathbb Eg,P_a^n(g-\mathbb Eg)\rangle\ge0\).
Applying this to every coefficient proves the matrix inequality.  General
\(L^2(H)\) class functions follow by truncation, since the uniform density
upper bound controls the native \(L^2\) limit.  Reflection gives the delay
\(2k\); congruence proves \eqref{f16v72:relative-reflection}.
\end{proof}

The fundamental writer has sharper constants, without summing all
representations.

\begin{proposition}[Explicit fundamental source and physical-time windows]
\label{f16v72:fundamental}
For \(W_t=w(H_t)\), let \(v_a=\operatorname{Var}(W_0)\).  Then
\begin{equation}
 |\mathbb E W_0|\le q,\qquad
 \frac{1-7q^2}{4}\le v_a\le\frac{1+3q^2}{4},\qquad
 0\le\operatorname{Cov}(W_0,W_n)\le q^2\quad(n\ge2).
 \label{f16v72:fundamental-bounds}
\end{equation}
In particular \(v_a\to1/4\) when \(q_a\to0\), while its ordinary-norm
normalized reflected Gram tends to zero.  More generally, for the bank
\eqref{f16v72:source-bank}, a probability-weighted temporal window with
weights \(h_j\ge0\), \(\sum_jh_j=1\), entirely at integer times \(k+j\ge1\),
has ordinary Gram \(D_h\) and reflected Gram \(H_h\) satisfying
\begin{equation}
 D_h\succeq b_q\left(\sum_jh_j^2\right)G_H,
 \qquad 0\preceq H_h\preceq\eta_qG_H.
 \label{f16v72:window}
\end{equation}
Thus the relative reflected bound costs only
\(\eta_q/[b_q\sum_jh_j^2]\), when \(b_q>0\).
\end{proposition}
\begin{proof}
The fundamental conditional mean has magnitude at most \(q\).
The identity \(4w^2=1+\chi_2\) and
\(|\mathbb E\chi_2|\le3\lambda_2(6\gamma)^m\le3q^2\) give the variance
bounds after centering.  In the preceding two-time proof replace the
all-representation mean estimate by \(|a_j|\le q\).  This gives the sharper
covariance bound.  For the window, every native lag covariance is positive
semidefinite; its diagonal contributions alone give the stated lower
ordinary Gram.  Every reflected pair has lag \(2k+i+j\ge2\), so summing
\eqref{f16v72:Gram-screen} with the nonnegative weights gives its upper bound.
\end{proof}

Under \eqref{f16v72:clock}, a fixed physical rectangle and a window with
\((\sum h_j^2)^{-1}\) polynomial in \(a^{-1},\gamma_a\) therefore give
vanishing relative reflection.  If \(\|G_H\|\) is itself only polynomial,
the \emph{unnormalized} reflected Gram also vanishes faster than every
power of \(a\).  This includes any fixed finite-degree character bank with
polynomial coefficients, but the relative theorem does not require a
fixed degree.

For a matching normalized slice source, the actual physical-energy
spectral measure satisfies, for every fixed \(E<\infty\),
\begin{equation}
 \mu_a([0,E])\le e^{2aE}\frac{\eta_{q_a}}{b_{q_a}}\longrightarrow0.
 \label{f16v72:energy-escape}
\end{equation}
Indeed its correlation at two steps is
\(\int e^{-2a\lambda}d\mu_a(\lambda)\), so its low-energy contribution is
at least \(e^{-2aE}\mu_a([0,E])\).  This is source-directed ultraviolet
escape, not a statement about the bottom of the full physical spectrum.

There is an essential normalization qualification.  V67 already shows that
a composite source can have a divergent contact norm and a nonzero
canonically normalized reflected limit.  The preceding result excludes
bounded-norm or polynomially normalized \emph{bare thin-loop} banks on the
stated trajectory; it does not exclude exponentially large perimeter
renormalization, or assert that such a renormalization exists.  Wilson-loop
renormalization is a separate issue; see, for example, M.~Berwein,
N.~Brambilla and A.~Vairo, \emph{Renormalization of Loop Functions in QCD},
\href{https://arxiv.org/abs/1312.6651}{arXiv:1312.6651}.  No renormalization
constant from that literature is imported.  None of these estimates applies
to \eqref{f16v69:source} merely because a flowed observable can be expanded
formally in loop functions.

\subsection{A shared-path singlet survives the same native integration}

The obstruction does not authorize bounding factors of a composite source
separately.  Let \(X\) be a long open path from \(p\) to \(q\), containing
\(m\) plaquette-separated selected links exactly once.  Let \(A,B\) be
retained return paths from \(q\) to \(p\), with no selected occurrence.
The following are genuine gauge-invariant closed-path functions:
\begin{equation}
 Z=w(XA)w(XB),\qquad R=w(AB^{-1}).
 \label{f16v72:fused-writers}
\end{equation}
The returns need not be independent under the native law.

\begin{theorem}[A native signed singlet reduction with a uniform error]
\label{f16v72:singlet}
Conditional on all unselected links \(\Xi\),
\begin{equation}
 \boxed{\quad
 \mathbb E[Z\mid\Xi]=\frac14R+e_\Xi,
 \qquad |e_\Xi|\le e_q:=\frac34\lambda_2(6\gamma)^m
                         \le\frac34q^2.
 \quad}
 \label{f16v72:singlet-bound}
\end{equation}
This retains the full native normalization and is uniform over every
choice of the return-path values.  The constant \(R/4\) is the Haar-singlet
part of the product in its common path variable; it can be nonconstant as
a function of the retained links.  More explicitly, if \(M\) is the
conditional product center of \(X\) and
\(\beta_2=\prod_{e\in S}\lambda_2(|h_e|)\), then
\begin{equation}
 e_\Xi=\beta_2\left[w(MA)w(MB)-\frac14R\right].
 \label{f16v72:exact-singlet-error}
\end{equation}
If some selected field is zero, \(\beta_2=0\) and \(M\) is immaterial.
\end{theorem}
\begin{proof}
For fixed \(A,B\), write \(w(XA)=x\cdot a\), \(w(XB)=x\cdot b\), where
\(a,b\) are unit quaternions and \(a\cdot b=w(AB^{-1})=R\).  Haar second
moments give
\(\int(x\cdot a)(x\cdot b)d\sigma=R/4\).
The difference
\[
 f_2(x)=(x\cdot a)(x\cdot b)-R/4
\]
is a traceless quadratic spherical harmonic on \(S^3\), hence lies entirely
in the Peter--Weyl sector \(k=2\).  This identification follows also from
the scalar-plus-traceless decomposition of symmetric quadratics in four
coordinates: its traceless space has dimension nine, the dimension of that
sector.  Haar fourth moments give exactly
\[
 \|f_2\|_2^2=\frac{2+R^2}{48}\le\frac1{16}.
\]
The conditional common-path density has \(k=2\) projection of \(L^2\) norm
at most \(3\prod_{e\in S}\lambda_2(|h_e|)\), by the same matrix Fourier
calculation and Plancherel as above.  Pair only that projection with
\(f_2\).  Cauchy--Schwarz gives the error at most
\((3/4)\lambda_2(6\gamma)^m\).  Higher sectors cannot contribute, and the
trivial sector is exactly \(R/4\).  A translated central law has second
moment
\[
 \mathbb E[xx^{\mathsf T}\mid\Xi]
      =\frac{1-\beta_2}{4}I_4+\beta_2MM^{\mathsf T};
\]
its longitudinal entry follows from \(4(x^0)^2=1+\chi_2\),
and its three transverse entries are equal by rotational symmetry.
This also proves \eqref{f16v72:exact-singlet-error} directly.
No estimate of the separate means of
\(w(XA)\) and \(w(XB)\) can replace this calculation.
\end{proof}

There is a literal short-loop example: take a return \(A\) avoiding the
selected long path, and put \(B=P_qA\), where \(P_q\) is an elementary
plaquette loop at \(q\), also avoiding the selected links.  Then
\(AB^{-1}=P_q^{-1}\) and \(R=w(P_q)\).  Thus two individually screened
long loops retain a specified elementary Wilson writer in their product.
A plaquette insertion may make the second loop self-intersecting; this is
allowed because only the selected common-path links must occur once.
It does not invalidate the separate simple-rectangle theorem.

\begin{proposition}[The complete two-time fused covariance is controlled]
\label{f16v72:fused-covariance}
Translate the common-path and return-path geometry to times \(0,n\), with
\(n\ge2\), and suppose the selected sets are jointly plaquette-separated.
Under the stationary native history,
\begin{equation}
 \boxed{\quad
 \left|\operatorname{Cov}(Z_0,Z_n)
       -\frac1{16}\operatorname{Cov}(R_0,R_n)\right|
 \le\frac12e_q+e_q^2
 \le\frac38q^2+\frac9{16}q^4.
 \quad}
 \label{f16v72:fused-correlation}
\end{equation}
For any coarser retained sigma-field \(C\subset\Xi\), the conditional
predictor likewise satisfies
\(\mathbb E[Z\mid C]=\frac14\mathbb E[R\mid C]+\mathbb E[e_\Xi\mid C]\),
with the same supremum error \(e_q\).
\end{proposition}
\begin{proof}
Condition on the complement of both selected sets.  The two common-path
variables are independent there, whereas every \(R_t\) is measurable.
The exact law of total covariance therefore gives
\[
 \operatorname{Cov}(Z_0,Z_n)
 =\operatorname{Cov}(R_0/4+e_0,R_n/4+e_n).
\]
The standard deviation of \(R_t/4\) is at most \(1/4\), and that of
\(e_t\) is at most \(e_q\).  The two cross covariances cost at most
\(e_q/4\) each and the error covariance at most \(e_q^2\).
This proves the stated complete covariance bound, with centering under the
actual marginal.  The final predictor statement is the conditional
expectation tower, not a replacement law for the retained paths.
\end{proof}

The surviving sector is therefore explicit, with an evaluated error under
the original interacting measure.  This is still not a positive physical
lower Gram: the retained \(R\) needs its own reflected estimate.  In the
short-loop example it remains a microscopic plaquette, not a constructed
finite-energy continuum field.  The gain is that the nontrivial long-path
sector is removed while the correct retained singlet and its covariance
are left intact, rather than suppressing an entire composite incorrectly.

\subsection{The next source calculation and the result ledger}

The next step should not try to make all raw tree-holonomy speeds small.
Theorem~\ref{f16v72:replica} proves the opposite on a declared growing
native rectangle.  Nor should one infer a surviving physical source from
the Haar-sized ordinary variance of an unflowed loop: the reflected bound
\eqref{f16v72:Gram-screen} is a direct calculation in the full native history.
It does not use a prescribed flat or cancelling exterior.

For the spatially flowed source, coincident path occurrences must instead
be combined into their local representation sectors before integrating
any plaquette-separated link set.  The two-loop calculation above is a
fully paid example of that source-directed operation.  A path expansion of
\eqref{f16v69:source} has increasing lengths, shared links, and coefficients
correlated with its nonlinear flow.  No summable expansion, control of its
retained singlet bank, or positive-separation lower Gram for that bank is
proved here.  Proving one of those native statements would move beyond the
present result; applying \(q^m\) separately to its uncombined factors would
not.  The original physical clock, source normalization and reflection
support remain unchanged.

The companion diagnostic is
\begin{center}\small
\path{verify_ym_f16_v72_thin_loops.py}.
\end{center}
It separates exact quaternion and Haar-moment identities, literal
four-dimensional rectangle incidence, positive compact conditional
integrals, Fourier and replica calculations, finite algebraic covariance
fixtures, and floating-point Bessel diagnostics.  It does not simulate the
full native Wilson vacuum or supply a rigorous numerical enclosure of a
continuum limit.  The analytic proofs do not use its numerical outcomes.
Removing the delimited appendix and reverting the one displayed version
change recovers V71 byte for byte.

\begin{record}[V72 result ledger]
\label{f16v72:ledger}
\emph{Proved here in the stated native setting:} coefficient-explicit
all-representation conditional holonomy bounds; a Haar-density estimate
uniform in the exterior; the all-but-four rectangle selection; a nonzero
Haar limit for declared long aligned chord speeds; simultaneous native
entrance and reflected Gram bounds for a thin-loop class-function bank;
its fundamental constants, temporal-window extension and normalized
physical-energy escape; an exact shared-path singlet with a uniform
nontrivial-sector error; and the complete native two-time retained-covariance
comparison, including its predictor.

\emph{Inherited or established background:} independent-link Wilson
conditioning, one-link character multipliers, classical perimeter-law
upper bounds, Haar representation theory, the anchored-tree identity and
native transfer/reflection positivity.  The perimeter phenomenon itself
is not claimed as a new discovery.  No earlier archive is recertified.

\emph{Not asserted:} a confining area law; failure of the spatially flowed
source; disappearance of every product of long loops; smallness of an
optimal gauge-orbit distance; exclusion or construction of exponential
perimeter renormalization; a physical interpretation of a Haar-sized
contact variance; or a reflected lower bound inferred from the singlet
identity without estimating the retained writer.

\emph{Still unproved:} a controlled physical-size singlet/source expansion
for the nonlinear spatially flowed bank; its noncollapsing positive-separation
native Gram and physical normalization; common dense-core physical-time
contraction; and the interacting four-dimensional continuum Yang--Mills
mass gap.
\end{record}

\begin{firewall}
The unselected native law is never replaced.  Bare holonomy screening is
separated from gauge-coordinate artifacts and from shared-path singlets.
A source-normalized high-energy limit is not a statement about the vacuum
spectral bottom.  V1--V71 remain unchanged.
\end{firewall}
% END F16 V72 APPEND
% BEGIN F16 V73 APPEND -- 2026-09-19
% Append-only continuation of the Post-Fifteenth-Archive V72.
\clearpage
\section[V73: native spectral exclusion for the frozen connection]{V73: native spectral exclusion\\for the frozen connection}
\label{f16v73:context}

V72 shows that shared-path singlets must be retained before screening a
composite.  Rather than introduce an uncontrolled expansion of the nonlinear
flow, we use its exact independent-link conditioning to answer an earlier
native question.  V68 gave a deterministic screening example and a
probability-dependent collapse certificate for V67's \emph{unflowed}
connection heat; it did not estimate the native spectral event.  We now
estimate it.  A small eigenvalue on a microscopic Neumann block forces
several independent chord holonomies to nearly preserve the \emph{same}
adjoint direction.  A two-dimensional sphere net pays for that unknown
direction.  The resulting bound is conditional on the actual exterior,
polynomial in the Wilson coefficient, and has arbitrarily high order at
zero energy when the fixed block is enlarged.

Neumann bracketing then gives two distinct native conclusions: exclusion
of finite spatial energies on a fixed physical torus with high probability,
and suppression of the averaged local heat weight uniformly in the ambient
volume.  The second conclusion handles curvature insertions correlated with
the connection, without replacing their law by an independent one.  It
proves that the specified frozen-connection source vanishes under every
polynomial normalization on the stated clock.  It does not prove collapse
of V68's preflowed source.  No auxiliary spatial spectrum is identified
with the physical Hamiltonian.

\subsection{Scope, notation and the conditional density input}

The measure is the isotropic compact four-dimensional Wilson measure with
coefficient \(\gamma\), or its stationary spatial-slice marginal
\(\pi_{a,\gamma}\).  Selected links are interior and touch at most six
plaquettes.  Bounds conditional on all other four-dimensional links pass
to the slice marginal by the expectation tower.  No nonlocal Perron
endpoint density is put into a one-link factor.  For spatially local heat
expectations we use translation-invariant periodic measures, their
fixed-spatial-regulator stationary time limits, or their spatial local
limits.  The block probability estimate itself needs no translation
invariance and is uniform in exterior data.

Let
\[
 \mathsf L_U=a^2\mathcal L_{a,U},\qquad
 \langle v,\mathsf L_Uv\rangle
 =\sum_{e=(x,y)}|v_x-\operatorname{Ad}_{U_e}v_y|^2.
\]
The sum is over positively stored spatial links, with the unweighted vertex
norm.  Multiplying both norms by \(a^3\) does not change eigenvalues.  For
a connected finite vertex block \(B\), let \(\mathsf L^N_{B,U}\) keep
\emph{only} links whose endpoints are in \(B\), without a boundary
penalty.  This is its Neumann quadratic form.  It is not a conditional
Wilson action or a new transfer operator.

For numerical studies of localization and adjoint spectral scales, see
\href{https://arxiv.org/abs/hep-lat/0504008}{Greensite et al. (2005)} and
\href{https://arxiv.org/abs/hep-lat/0606008}{Greensite et al. (2006)}.
Those studies are not premises of the estimates below; their operators
and observables must not be interchanged with a three-dimensional slice
heat kernel.  The proofs here use compact Haar geometry, finite-dimensional
min--max, and the scalar heat domination already proved in V67.  No claim
of novelty for these general methods, or of verification of the accumulated
archive, is made.

Write
\begin{equation}
 \kappa_\gamma=\frac{3\pi e}{4}(1+6\gamma)^{3/2}.
 \label{f16v73:density-constant}
\end{equation}
For each plaquette-separated selected link, V72's actual conditional law
has density \(e^{h\cdot x}/\mathfrak z(|h|)\), \(|h|\le6\gamma\), and
\begin{equation}
 \left\|\frac{d\nu_h}{dH}\right\|_\infty\le\kappa_\gamma.
 \label{f16v73:one-density}
\end{equation}
Indeed V63's cap estimate is
\(e^{-t}\mathfrak z(t)\ge4(1+t)^{-3/2}/(3\pi e)\).
This is an \emph{upper} density bound.  Its polynomial cost, rather than
a finite-energy bound exponential in \(\gamma\), is essential here.

\subsection{Several native chords cannot preserve an arbitrary common direction cheaply}

Take a rooted spanning tree of \(B\), of maximum root distance \(d\),
and let \(n=|B|\).  Choose \(r\ge2\) non-tree links \(e_1,\ldots,e_r\)
that are plaquette-separated in the \emph{full four-dimensional lattice}.
Let \(A_x\) be the ordered tree transport from the root to \(x\).  Conditional
on all unselected links, the chord holonomies
\[
 H_e=A_xU_eA_y^{-1},\qquad e=(x,y),
\]
are independent, each with density at most \(\kappa_\gamma\) relative to
Haar.  All tree and remaining chord data stay fixed in this conditioning.
Put
\begin{equation}
 K(n,d)=2\sqrt n(1+2\sqrt d),\qquad
 A(r,n,d;\gamma)=64\kappa_\gamma^r K(n,d)^{2r-2}.
 \label{f16v73:block-constants}
\end{equation}

\begin{theorem}[Conditional low-eigenvalue probability with the direction cost paid]
\label{f16v73:block-smallball}
For every \(\epsilon\ge0\), uniformly conditional on the complement of the
selected links,
\begin{equation}
 \boxed{\quad
 \mathbb P\{\lambda_{\min}(\mathsf L^N_{B,U})\le\epsilon\mid\Xi\}
 \le\min\{1,A(r,n,d;\gamma)\epsilon^{r-1}\}.
 \quad}
 \label{f16v73:block-tail}
\end{equation}
The same inequality holds under the full native law and its slice marginal.
There is no independence assumption on the unselected links.
\end{theorem}
\begin{proof}
Choose a real normalized lowest eigenvector \(v\); the matrix is real
symmetric.  If its energy is at most \(\epsilon\), the sum of all squared
edge discrepancies is at most \(\epsilon\).  Transport its values to the
root: \(w_x=\operatorname{Ad}_{A_x}v_x\).  Telescoping along each tree path
and Cauchy--Schwarz give
\[
 |w_x-w_o|\le\sqrt{d\epsilon}.
\]
Some vertex has \(|v_x|\ge n^{-1/2}\).  Therefore
\(|w_o|\ge(2\sqrt n)^{-1}\) if \(\epsilon\le(4nd)^{-1}\).
For a selected chord its edge discrepancy and the two tree discrepancies
imply
\[
 |w_o-\operatorname{Ad}_{H_e}w_o|
 \le(1+2\sqrt d)\sqrt\epsilon.
\]
Thus, with \(z=w_o/|w_o|\in S^2\),
\begin{equation}
 |z-\operatorname{Ad}_{H_e}z|\le K(n,d)\sqrt\epsilon=:\delta
 \quad\hbox{for all selected }e.
 \label{f16v73:common-direction}
\end{equation}
If \(\delta\le1\), the root lower bound used above is valid as well.

A Euclidean \(\delta/2\)-net on the unit two-sphere can be chosen with at
most \(64/\delta^2\) elements.  To prove this count, take a maximal
\(\delta/2\)-separated set.  Chordal caps of radius \(\delta/4\) are
disjoint, and each has area \(\pi\delta^2/16\), out of total area
\(4\pi\).  Maximality gives the claimed net.  If \(z\) satisfies
\eqref{f16v73:common-direction}, one net point \(z_j\) satisfies
\(|z_j-\operatorname{Ad}_{H_e}z_j|\le2\delta\) for every \(e\), since
\(\|I-\operatorname{Ad}_{H_e}\|\le2\).

For fixed \(z_j\), a Haar adjoint rotation sends \(z_j\) uniformly to
\(S^2\).  The Haar probability of its chordal \(2\delta\)-cap is exactly
\(\delta^2\), for \(0\le\delta\le1\).  Conditional independence and
\eqref{f16v73:one-density} therefore bound the probability for all \(r\)
chords at this net point by \(\kappa_\gamma^r\delta^{2r}\).  The union
over the net gives \(64\kappa_\gamma^r\delta^{2r-2}\), proving the
stated coefficient.  The selected variables determine the eigenvector;
the net is precisely what pays that dependence.  If \(\delta>1\), the
right-hand power is already at least one and the trivial bound suffices.
Taking \(\epsilon\downarrow0\) proves the zero-threshold assertion.
Finally use the conditional expectation tower.
\end{proof}

A single rotation always has a fixed adjoint axis.  Ignoring the choice of
a \emph{common} axis would give the incorrect exponent \(r\) in
\eqref{f16v73:block-tail}.  Conversely, the retained tree is not integrated
as independent Haar data.  Shared paths are handled together by parallel
transport; the argument does not screen uncombined scalar factors and does
not discard V72's singlet contributions.

\begin{proposition}[An explicit family of literal cubic blocks]
\label{f16v73:geometry}
For any fixed \(r\ge2\), put
\begin{equation}
 \ell_r=2\lceil\sqrt r\rceil+1,\quad
 n_r=(2\ell_r-1)^3,\quad d_r=3(2\ell_r-2),\quad
 A_{r,\gamma}=64\kappa_\gamma^r K(n_r,d_r)^{2r-2}.
 \label{f16v73:uniform-constants}
\end{equation}
Every rectangular spatial block whose three vertex side lengths lie
between \(\ell_r\) and \(2\ell_r-1\) contains a tree and \(r\)
plaquette-separated chords to which \eqref{f16v73:block-tail} applies with
coefficient \(A_{r,\gamma}\).  These constants do not depend on the
surrounding volume.
\end{proposition}
\begin{proof}
Use coordinates starting at zero.  The tree consists of the first-direction
axis at \(y=z=0\), all second-direction rows at \(z=0\), and all
third-direction columns.  It connects every vertex with no cycle, and its
root distance is at most \(d_r\).  Select \(r\) of the first-direction
links from \((0,2i,2j)\) to \((1,2i,2j)\),
\(1\le i,j\le\lceil\sqrt r\rceil\).  None belongs to the tree.
Two selected parallel links are separated transversely by at least two
steps, so they cannot share a plaquette.  A plaquette involving time
contains no second selected link either.  The vertex number is at most
\(n_r\).  The construction works also on a surrounding nondegenerate
torus: through a periodic seam the extreme selected coordinates remain
at least two steps apart.  Increasing the actual tree-distance and vertex
bounds to \(d_r,n_r\) only enlarges the constant.
\end{proof}

\subsection{Bracketing gives a native spectral estimate rather than a typicality hypothesis}

Let \(N_a\) be the number of spatial vertices of a periodic torus, each
side containing at least \(\ell_r\) vertices.  Partition each side into
intervals of lengths between \(\ell_r\) and \(2\ell_r-1\): start with
intervals of length \(\ell_r\) and add the remainder to the last interval.
Their products give blocks of the preceding proposition.  Remove only
cross-block edges from the \emph{quadratic form}; keep the probability
measure unchanged.  Positivity of each removed edge gives
\begin{equation}
 \mathsf L_U\succeq\bigoplus_B\mathsf L^N_{B,U}.
 \label{f16v73:bracketing}
\end{equation}

\begin{theorem}[Finite-volume exclusion and volume-normalized heat weight]
\label{f16v73:trace}
For the actual native slice law and \(E\ge0\),
\begin{equation}
 \boxed{\quad
 \mathbb P\{\lambda_{\min}(\mathcal L_{a,U})\le E\}
 \le\min\{1,N_a A_{r,\gamma}(a^2E)^{r-1}\}.
 \quad}
 \label{f16v73:finite-volume-spectrum}
\end{equation}
For every \(t>0\), without translation invariance,
\begin{equation}
 \frac1{N_a}\mathbb E\Tr e^{-t\mathcal L_{a,U}}
 \le\Theta_{r,\gamma}(a,t):=
 \min\{3,3\Gamma(r)A_{r,\gamma}(a^2/t)^{r-1}\}.
 \label{f16v73:heat-trace}
\end{equation}
For a translation-invariant native slice law this also gives, at every vertex,
\begin{equation}
 \boxed{\quad
 \mathbb E\operatorname{tr}_{\mathbb R^3}
           e^{-t\mathcal L_{a,U}}(x,x)
 \le\Theta_{r,\gamma}(a,t).
 \quad}
 \label{f16v73:diagonal}
\end{equation}
The local diagonal bound is uniform in the ambient torus and passes to its
translation-invariant spatial Gibbs limits at each fixed regulator.
\end{theorem}
\begin{proof}
If the left operator in \eqref{f16v73:bracketing} has lowest eigenvalue at
most \(a^2E\), some block does also.  The union bound over at most \(N_a\)
blocks and Theorem~\ref{f16v73:block-smallball} prove
\eqref{f16v73:finite-volume-spectrum}; independence of different blocks is
not needed.

Let \(\lambda_B=\lambda_{\min}(\mathsf L^N_{B,U})\).
The tail estimate implies, for \(T>0\),
\[
 \mathbb E e^{-T\lambda_B}
 =\int_0^\infty T e^{-T\epsilon}\mathbb P\{\lambda_B\le\epsilon\}
       d\epsilon
 \le A_{r,\gamma}\Gamma(r)T^{-(r-1)}.
\]
The identity is also obtained directly by Tonelli from
\(e^{-T\lambda}=\int_\lambda^\infty Te^{-T\epsilon}d\epsilon\).
A block has \(3|B|\) eigenvalues, so its expected heat trace is at most
\(3|B|\) times this bound and also at most \(3|B|\).
Min--max in \eqref{f16v73:bracketing} compares the \emph{ordered eigenvalues},
and hence their heat traces.  Summing the block estimates and setting
\(T=t/a^2\) proves \eqref{f16v73:heat-trace}.  We do not assert the false
operator implication \(e^{-t\mathcal L}\preceq e^{-t\mathcal L_N}\).
Translation invariance makes the expectations of the color traces on the
diagonal identical, proving \eqref{f16v73:diagonal}.

At fixed \(a\), the dimensionless Laplacian is bounded by twelve and has
finite range.  Its heat series can be uniformly approximated by polynomials
in that operator; each diagonal polynomial is a bounded local continuous
link observable.  This proves passage through spatial local limits.  Time
cylinder limits are handled by the same conditional-law tower and local
approximation, not by changing the vacuum normalization.
\end{proof}

For fixed \(r\), \(A_{r,\gamma}=C_r(1+6\gamma)^{3r/2}\) with an explicitly
specified, large \(C_r\).  The arbitrary order \(r-1\) is obtained by
choosing a larger \emph{fixed microscopic} block before taking limits.
These constants are not optimized; the bounds are not advertised as useful
numerical estimates at moderate coupling.

\subsection{The earlier collapse certificate is populated at fixed physical volume}

Recall V67's frozen-connection field, writing \(s=r_{\rm heat}^2/2>0\)
to distinguish the physical heat radius from the integer chord count:
\begin{equation}
 \mathcal F_a(U;y)=\zeta_a a^{-2}\operatorname{Im}U_{12}(y),\qquad
 B_s(U)=e^{-s\mathcal L_{a,U}}\mathcal F_a(U),\qquad
 O_f(U)=\frac{a^3}{2}\sum_x f(x)|B_s(U;x)|^2.
 \label{f16v73:source}
\end{equation}
All expectations, including centering \(F_f=O_f-\mathbb E O_f\), use the
original native law.  The initial field and the heat connection may be
arbitrarily correlated.

\begin{proposition}[An explicit native finite-volume collapse bound]
\label{f16v73:finite-collapse}
Under the side-length condition of Theorem~\ref{f16v73:trace}, on a
spatial torus of physical volume \(\mathcal V_a=N_aa^3\), for
\(0<\sigma<2\),
\begin{equation}
 \begin{split}
 \operatorname{Var}_{\pi_a}(O_f)
 \le{}&\frac14\|f\|_\infty^2\mathcal V_a^2\zeta_a^4a^{-8}\,
 \left[e^{-4s a^{-\sigma}}
 +\min\{1,\mathcal V_a A_{r,\gamma}
                       a^{(2-\sigma)(r-1)-3}\}\right].
 \end{split}
 \label{f16v73:finite-collapse-bound}
\end{equation}
In particular \(\sigma=1,r=14\), bounded physical volume and bounded
\(\|f\|_\infty\) give a constant times
\(\zeta_a^4[a^{-8}e^{-4s/a}+(1+6\gamma)^{21}a^2]\).
\end{proposition}
\begin{proof}
Use \eqref{f16v73:finite-volume-spectrum} with \(E=a^{-\sigma}\).
On the complementary event, heat contraction gives
\(\|B_s\|_a^2\le e^{-2sa^{-\sigma}}\mathcal V_a\zeta_a^2a^{-4}\);
on the event use the same bound without the exponential.  Bound \(|O_f|\)
by \(\|f\|_\infty\|B_s\|_a^2/2\), square and average.  This is exactly
the previously conditional certificate with a now-proved probability input.
\end{proof}

This already proves native collapse on a fixed physical volume under the
clock below.  To avoid making the source conclusion depend on that volume
or on a favorable whole-volume spectral bottom, the next argument uses
\eqref{f16v73:diagonal} instead.

\subsection{A local bound for correlated insertions, uniform in surrounding volume}

Let \(K_s^U(x,y)\) be the \(3\times3\) matrix kernel of
\(e^{-s\mathcal L_{a,U}}\), with counting measure in the kernel convention.
V67's exact walk expansion gives
\(\|K_s^U(x,y)\|_{\rm op}\le p_s^a(x,y)\), where \(p_s^a\) is the
scalar rate-\(a^{-2}\) nearest-neighbor kernel.  Self-adjointness and the
semigroup law give, for every fixed configuration,
\begin{equation}
 \sum_y\|K_s^U(x,y)\|_{\rm HS}^2
 =\operatorname{tr}K_{2s}^U(x,x).
 \label{f16v73:HS-row}
\end{equation}
These identities include all repeated path occurrences and their singlets.

For \(R>0\), let \(N_R=(2\lfloor R/a\rfloor+1)^3\), an upper bound on
the number of sites in the torus coordinate cube of radius \(R\).  Set
\begin{align}
 J_a(R,s)&=\frac Ra\operatorname{arsinh}\frac{Ra}{2s}
 -\frac{2s}{a^2}\left(\sqrt{1+(Ra/(2s))^2}-1\right),
     \label{f16v73:walk-rate}\\
 h_a(R,s)&=\min\{1,6e^{-J_a(R,s)}\},\qquad
 \mathcal B_{r,a,\gamma}(R,s)
 =N_R\Theta_{r,\gamma}(a,2s)+h_a(R,s)^2.
     \label{f16v73:local-budget}
\end{align}
On a torus the far-walk probability is zero when this cube covers the
torus, and is otherwise bounded by \(h_a\).  The displayed bound remains
valid in either case.

\begin{theorem}[Averaged local suppression with the connection correlation retained]
\label{f16v73:local-source}
Under the translation-invariant native laws of
Theorem~\ref{f16v73:trace}, any possibly connection-dependent adjoint field
\(\mathcal F_a(U;y)\) with \(|\mathcal F_a(U;y)|\le M_a\) satisfies
\begin{equation}
 \boxed{\quad
 \mathbb E|e^{-s\mathcal L_{a,U}}\mathcal F_a(U;x)|^2
 \le2M_a^2\mathcal B_{r,a,\gamma}(R,s).
 \quad}
 \label{f16v73:local-field}
\end{equation}
For \eqref{f16v73:source}, \(M_a=\zeta_a a^{-2}\).  If
\(\|f\|_{1,a}=a^3\sum_x|f(x)|\), then
\begin{equation}
 \boxed{\quad
 \operatorname{Var}(O_f)\le\mathbb E|O_f|^2
 \le\frac12\zeta_a^4a^{-8}\|f\|_{1,a}^2
                    \mathcal B_{r,a,\gamma}(R,s).
 \quad}
 \label{f16v73:local-variance}
\end{equation}
All constants are independent of the surrounding spatial volume.  In
particular the selected chords need not avoid the curvature insertions.
\end{theorem}
\begin{proof}
Split the kernel sum into the coordinate cube and its complement.  The
near sum is bounded by
\(M_a\sqrt{N_R\operatorname{tr}K_{2s}^U(x,x)}\), by Cauchy--Schwarz and
\eqref{f16v73:HS-row}.  The far sum is at most \(M_ah_a(R,s)\), by scalar
heat domination.  Each signed coordinate displacement has exponential
moment
\[\exp[2s(\cosh(\lambda a)-1)/a^2].\]
Applying this bound in all six directions and optimizing \(\lambda\)
gives the stated tail.  The lifted walk
also bounds torus distance.  Square the two-term bound using
\((b+c)^2\le2b^2+2c^2\), then apply \eqref{f16v73:diagonal}.
Nothing in this pointwise estimate factors the field from its connection.
Scalar domination makes the bounded-field kernel sum absolutely convergent
also in infinite volume; the same argument follows there by exhaustion.

Scalar domination also gives \(|B_s(x)|\le M_a\), hence
\(|O_f|\le M_a^2\|f\|_{1,a}/2\).  Equation
\eqref{f16v73:local-field} gives
\(\mathbb E|O_f|\le M_a^2\|f\|_{1,a}\mathcal B\).
Multiply these two inequalities to bound the second moment.  This proves
\eqref{f16v73:local-variance}, including signed tests.
\end{proof}

For orientation, a fully specified finite-order consequence uses \(r=8\)
and \(R_a=\sqrt{48s\log(1/a)}\).  For all sufficiently small \(a\),
\(aR_a/(2s)\le1\), and the elementary bound
\(J_a(R_a,s)\ge R_a^2/(8s)\) gives \(h_a^2\le36a^{12}\).
Thus, for fixed \(s>0\),
\begin{equation}
 \operatorname{Var}(O_f)
 \le C_s\zeta_a^4\|f\|_{1,a}^2
 \left[(1+6\gamma)^{12}a^3(\log(1/a))^{3/2}+a^4\right].
 \label{f16v73:eight-chord-example}
\end{equation}
The constants in \eqref{f16v73:uniform-constants} determine \(C_s\);
the large block cost remains explicit.
The growing cube in this \emph{estimate} is not a change of the observable.

\subsection{The frozen-connection source vanishes faster than every power}

\begin{theorem}[Native source collapse under polynomial normalization]
\label{f16v73:collapse}
Let \(a\downarrow0\), with spatial side lengths in lattice units tending
to infinity, and suppose
\begin{equation}
 \log(1+\gamma_a)=o(\log(1/a)),\qquad
 \zeta_a\le C a^{-d}(1+\gamma_a)^b
 \label{f16v73:scale-assumptions}
\end{equation}
for fixed finite \(b,d\ge0\).  Fix \(s>0\), and use a finite family of
sampled spatial tests with bounded \(\|f_i\|_{1,a}\).  In the
translation-invariant native laws above, for every \(N>0\),
\begin{equation}
 \boxed{\quad
 \mathbb E|B_s(U;x)|^2=o(a^N),\qquad
 \operatorname{Var}(O_{f_i})=o(a^N),\qquad
 \|D_a\|_{\rm op}=o(a^N),
 \quad}
 \label{f16v73:superpolynomial}
\end{equation}
where \(D_a=\operatorname{Cov}(O_{f_i})_{ij}\).
The same absolute conclusion holds for the reflected Gram of any common
probability-weighted positive-time average and any additional nonnegative
physical-time delay.  Any further polynomial coefficient normalization is
also absorbed by \eqref{f16v73:superpolynomial}.
\end{theorem}
\begin{proof}
For each requested power choose a \emph{fixed} integer \(r\) sufficiently
large.  Its constant obeys
\(A_{r,\gamma_a}=a^{-o(1)}\), by \eqref{f16v73:scale-assumptions}.
Take \(R_a=\sqrt{4sL\log(1/a)}\), where \(L\) is another fixed, sufficiently
large number.  For small \(a\), \(J_a\ge R_a^2/(8s)\), so
\(h_a^2\le36a^L\), while
\(N_{R_a}=O(a^{-3}(\log(1/a))^{3/2})\).  Equations
\eqref{f16v73:local-field}--\eqref{f16v73:local-variance} give respectively
powers at least \(2r-9-2d-o(1)\) and \(2r-13-4d-o(1)\), together with
the arbitrarily high far-tail powers obtained from \(L\).
Choose both exponents strictly larger than \(N\).  The finite bank has
\(\|D_a\|\le\Tr D_a\), proving its claim.  Every sampled linear source
relation remains exact; a coefficient map pulls the covariance back by
congruence, and a polynomial norm only uses a larger requested power.

For time averaging, stationarity and Jensen bound its ordinary covariance
by \(D_a\) in every coefficient direction.  The original native transfer
is a positive contraction, and reflection is norm preserving in the
Euclidean history.  Its reflected Gram is therefore positive semidefinite
and bounded above by \(D_a\), also after an additional delay.  These are
absolute estimates under the same centering and native measure.  No
assertion about a reflected-to-ordinary ratio is needed.
\end{proof}

The separately stated trajectory
\(\log(1/a)=c\gamma_a+O(\log\gamma_a)\), \(c>0\), satisfies
\eqref{f16v73:scale-assumptions}.  The argument does not derive that
trajectory.  It also does not choose \(r=r(a)\): each asymptotic rate is
proved with its own fixed microscopic geometry and its full fixed constant.
An additive modification of \(\mathcal L_{a,U}\), or a faster-than-polynomial
multiplicative normalization, defines a different problem.  Neither is
excluded or constructed by the polynomial conclusion.

\subsection{Zero infinite-volume bottom is compatible with the local theorem}

\begin{proposition}[Rare almost-flat regions keep the spatial spectral bottom at zero]
\label{f16v73:bottom}
At fixed \(a>0\) and finite \(\gamma\), under the infinite-volume native
finite-energy laws in V70,
\begin{equation}
 \inf\operatorname{spec}(\mathcal L_{a,U})=0\quad\hbox{almost surely}.
 \label{f16v73:infinite-bottom}
\end{equation}
This is compatible with every estimate above.  It is not a claim of an
\(L^2\) zero eigenvector.
\end{proposition}
\begin{proof}
V70's proved finite-energy pattern occurrence supplies, for each integer
\(L\), a translated cube and its incident-edge collar on which every link
is within \(1/L\) of identity in group distance.  Take a scalar Dirichlet
sine test in that cube times one fixed unit adjoint vector, extended by zero.
Its scalar gradient energy divided by its norm is at most \(C/L^2\).
Since \(\|\operatorname{Ad}_{U_e}-I\|\le2/L\) on the collar,
\(|v_x-\operatorname{Ad}_{U_e}v_y|^2
\le2|v_x-v_y|^2+8|v_y|^2/L^2\).
Each vertex is an endpoint of boundedly many links, so the covariant
Rayleigh quotient is at most \(C'/(a^2L^2)\).  A countable intersection
over \(L\) has full probability.  Sending \(L\to\infty\), positivity of
the operator proves \eqref{f16v73:infinite-bottom}.
\end{proof}

The finite-volume probability \eqref{f16v73:finite-volume-spectrum} pays
its number of blocks.  The local heat estimate instead averages spectral
\emph{weight}; it never assumes that the infinite-volume bottom is
positive.  This separation is why remote rare regions do not invalidate
the local source conclusion.  None of these spatial spectra is the native
physical spectrum of \((-\log P_a)/a\).

\subsection{What the new native result changes, and what it does not}

V68's deterministic warning is now strengthened to a theorem under the
actual native law for the original frozen-connection heat source.  Its
ordinary and reflected norms vanish with every polynomial normalization on
the stated clock, even when the inserted curvature depends on exactly the
same links as the connection.  A fixed-order singlet example from V72 does
not overturn this result: \eqref{f16v73:HS-row} and the Neumann form include
all repeated paths and test their aggregate propagation rather than
suppressing scalar factors separately.

The proof does \emph{not} apply unchanged after spatial preflow.  The links
\(W=\Phi_{a,\tau}(U)\) have the actual pushforward law, not a Wilson law
with the original coefficient.  Conditional independence and the bound
\(\kappa_\gamma^r\) for \(r\) selected \(W\)-links have not been proved;
V68 already keeps the change-of-variables density and Jacobian explicit.
The current \(s=0,\tau>0\) source in V69--V71 does not apply a frozen
connection heat at all.  Its reflected lower Gram remains unestimated.
It would be circular to use the present vanishing upper bound as evidence
of noncollapse for that different observable.

The next native target remains the positive-separation covariance of the
\emph{preflowed} bank, or a justified renormalized alternative.  This
iteration retires one earlier source candidate in a precisely specified
normalization class.  It supplies neither a new interacting continuum
source nor physical-time contraction.  An actual lower Gram for the
surviving bank, its physical normalization and spatial control, and common
dense-core constants remain separate tasks.

\subsection{Verification, preservation and result ledger}

The companion diagnostic is
\begin{center}\small
\path{verify_ym_f16_v73_native_spectral_exclusion.py}.
\end{center}
It separates exact Haar-cap and incidence fixtures, literal compact
connection matrices, deterministic low-mode/common-axis tests,
Neumann eigenvalue and trace comparisons, scalar walk tails and
correlated-insertion inequalities.  Its random compact configurations
are \emph{diagnostics}, not samples from the four-dimensional Wilson
vacuum.  Analytic probability bounds are proved above, not estimated from
rare-event simulation.  No interval certificate, formal proof-assistant
verification or external archive recertification is claimed.  Removing
this delimited appendix and reverting the displayed version recovers V72
byte for byte.

\clearpage
\begin{record}[V73 result ledger]
\label{f16v73:ledger}
\emph{Proved here:} an exterior-uniform conditional low-eigenvalue bound
with the common-axis entropy paid; explicit cubic selected-chord geometry;
native finite-volume spectral exclusion and volume-normalized heat-trace
bounds; a local averaged heat estimate uniform in surrounding volume;
collapse of the frozen-connection quadratic bank with correlated insertions
and arbitrary polynomial normalization on the stated scaling; the
positive-time reflected consequence; and compatibility with zero
infinite-volume spatial spectral bottom.

\emph{Inherited or established background:} V72's literal independent-link
conditioning, V63's compact cap normalization, V67's adjoint heat and
scalar domination, V70's finite-energy pattern occurrence, and ordinary
min--max and sphere-net arguments.  None is a native physical mass-gap input.

\emph{Not asserted:} a uniform positive spatial spectral bottom in infinite
volume; localization of individual eigenvectors or a sharp Lifshitz exponent;
a false operator-monotonicity statement for the exponential; independence
of the field from its connection; a useful moderate-coupling numerical
constant; a conclusion for exponentially normalized sources; or the same
conditional law for flowed links.

\emph{Still unproved:} a nonzero cutoff-uniform reflected Gram for the
specified preflowed bank, its joint native covariance and physical
normalization, common dense-core physical-time contraction, and the
interacting four-dimensional continuum Yang--Mills mass gap.
\end{record}

\begin{firewall}
The native spectral event is estimated, not declared typical.  Every
conditional normalization and shared tree transport is retained.  The
random insertion remains correlated with the connection.  A spatial
source-collapse theorem is not a theorem about the physical spectral gap.
V1--V72 remain unchanged.
\end{firewall}
% END F16 V73 APPEND
% BEGIN F16 V74 APPEND -- 2026-09-19
% Append-only continuation of the Post-Fifteenth-Archive V73.
\clearpage
\section[V74: transported native densities and the preflow scale]{V74: transported native densities\\and the preflow scale}
\label{f16v74:context}

V73 excludes the frozen-connection heat source but explicitly does not give
its selected-link law after spatial preflow.  The present continuation
computes a replacement for that missing law.  It does not assume that
flowed links are conditionally independent.  A bound on the logarithmic
derivative of the actual pushed density gives a conditional small-cap
estimate for their adjoint directions, and this is enough to repeat the
common-direction test without independence.  The resulting theorem excludes
an additional, logarithmically growing layer of preflow before the fixed
constituent heat.

There is also an opposing native estimate.  The Jacobian produces an exact
entropy identity.  At weak bare coupling it forces large \emph{conditional}
concentration after a positive physical spatial-flow time, although every
individual flowed-link marginal remains exactly Haar.  Thus retaining the
original polynomial one-link density bound at that time is not merely
unjustified: it is incompatible with the native entropy.  The exclusion and
entropy estimates concern different ranges of the flow parameter; neither
is a positive-separation lower Gram for the current $s=0$ source.

\subsection{The actual slice law and the flow Jacobian}

Work first on a finite spatial cubic torus of side at least four.  Its
$M=3N$ stored spatial links and $M$ spatial plaquettes carry normalized
product Haar measure $dH^M$.  Let $\pi_{a,\gamma}=\rho_0dH^M$ be the actual
spatial-slice marginal of the four-dimensional isotropic Wilson law,
including its stationary fixed-spatial-regulator limit.  The density is
smooth and positive.  No Perron endpoint factor is absorbed into a
one-link Wilson density.  All gradients and Laplacians on the configuration
manifold use the product of round unit-quaternion metrics.

Use precisely the spatial flow of V68:
\[
 b(W)=-a^{-2}\nabla E_{\rm sp}(W),\qquad
 E_{\rm sp}(W)=\sum_{P\ {
m spatial}}(1-w(W_P)),\qquad
 W=\Phi_{a,\tau}(U).
\]
Set $u=\tau/a^2$ and
\begin{equation}
 \mu_\tau=(\Phi_{a,\tau})_*\pi_{a,\gamma}=\rho_\tau dH^M,
 \qquad
 \ell(\rho)=\max_e\sup_W|\nabla_e\log\rho(W)|.
 \label{f16v74:law-score}
\end{equation}
The notation $\ell(\rho)$ denotes a local logarithmic-derivative bound,
not a lattice length.  The flow parameter has spatial-length-squared units;
the physical-time transfer remains $P_a$.

The change-of-variables method and the integrated-divergence expression for
a Wilson-flow Jacobian are established background; see M.~L\"uscher,
\emph{Trivializing maps, the Wilson flow and the HMC algorithm},
Commun.\ Math.\ Phys.\ \textbf{293} (2010), 899--919,
\href{https://arxiv.org/abs/0907.5491}{arXiv:0907.5491}, \S3.2.
That reference uses different group and action conventions.  The
$\SU(2)$ spatial coefficients and the conditional estimates below are
derived here, not imported from a trivializing-map assertion.

\begin{theorem}[Exact Jacobian and a volume-independent local score bound]
\label{f16v74:score}
At the stated finite regulator,
\begin{align}
 \log\det D\Phi_{a,\tau}(U)
 &=-\frac{12}{a^2}\int_0^\tau\sum_P w((\Phi_{a,v}U)_P)\,dv,
 \label{f16v74:Jacobian}\\
 \log\rho_\tau(W)
 &=\log\rho_0(\Phi_{a,-\tau}W)
   +\frac{12}{a^2}\int_0^\tau\sum_Pw((\Phi_{a,-v}W)_P)\,dv.
 \label{f16v74:pushed-density}
\end{align}
In particular,
\begin{equation}
 \boxed{\quad
 \ell(\rho_0)\le6\gamma,\qquad
 \ell(\rho_\tau)\le L_{\gamma,u}:=(6\gamma+3)e^{16u}-3.
 \quad}
 \label{f16v74:score-bound}
\end{equation}
The bound is independent of the surrounding spatial volume and temporal
cylinder length.  The formulas retain the original normalization exactly.
\end{theorem}
\begin{proof}
For one spatial link in the full four-dimensional measure, differentiating
the logarithm of its marginal density gives the conditional expectation
of its Wilson score.  There are at most six incident plaquettes and the
norm of the derivative of each normalized trace is at most one.  Thus
$|\nabla_e\log\rho_0|\le6\gamma$.  Differentiation under the compact
integral, followed by the fixed-regulator cylinder limit, proves the same
assertion for the native stationary slice.  Equivalently one may first
use the uniformly bounded local score and then pass its translation
inequalities to that limit.

The round $S^3$ Laplacian sends a coordinate-linear function to minus
three times that function.  A plaquette trace is coordinate-linear in
each of its four distinct links.  Consequently
\begin{equation}
 \Delta E_{\rm sp}=12\sum_Pw(W_P),\qquad
 \operatorname{div}b=-12a^{-2}\sum_Pw(W_P).
 \label{f16v74:divergence}
\end{equation}
Liouville's formula proves \eqref{f16v74:Jacobian}.  Substitution of
$U=\Phi_{a,-\tau}W$ and reversal of the integration parameter give
\eqref{f16v74:pushed-density}, with its positive Jacobian contribution
to the logarithm of the pushed density.

It remains to differentiate this extensive expression without an
extensive estimate.  V69's exact plaquette Hessian has diagonal block
$w(W_P)I_3$ at each incident link and off-diagonal block norm at most
one.  Thus the full Hessian has diagonal norm at most four and
$\|\mathsf H_{ef}\|\le N_{ef}$, where $N$ is the shared-plaquette
matrix with row and column sums twelve.  In separate parallel frames,
the forward or backward differential obeys the entrywise block bound
\begin{equation}
 \|(D\Phi_{a,\pm v})_{ef}\|
 \le \left[e^{v(4I+N)/a^2}\right]_{ef}.
 \label{f16v74:inverse-local-bound}
\end{equation}
This follows by iteration of the linear variational equation; changing
the time direction changes signs, not the displayed norm majorant.
Both row and column sums on the right equal $e^{16v/a^2}$.  Hence for
any scalar function $g$ with $\max_e|\nabla_eg|\le A$,
\[
 \max_e|\nabla_e(g\circ\Phi_{a,-v})|\le A e^{16v/a^2}.
\]
The spatial trace sum has local gradient norm at most four.  Differentiating
\eqref{f16v74:pushed-density} therefore gives
\[
 \ell(\rho_\tau)
 \le6\gamma e^{16u}
       +\frac{48}{a^2}\int_0^\tau e^{16v/a^2}\,dv
 =(6\gamma+3)e^{16u}-3.
\]
Every extensive sum was controlled by a local gradient and a column sum;
no factor proportional to $M$ was omitted.
\end{proof}

The determinant in \eqref{f16v74:Jacobian} is with respect to Haar volume,
which is a constant multiple of the Riemannian volume.  Parallel-frame
changes have unit determinant.  The identity does not imply that
$\rho_\tau$ is a Wilson density at any effective coefficient.

\subsection{Conditional adjoint caps without conditional independence}

\begin{proposition}[Group-density and two-dimensional cap consequences]
\label{f16v74:conditional-caps}
Let $\rho>0$ be a smooth probability density on a finite product of $S^3$
with $\ell(\rho)\le L$.  The same logarithmic-derivative bound holds for
every marginal and for every one-link conditional density, including
conditional densities after other variables have been integrated out.
Writing $c_H=8/(3\pi^3)$ and $C_H=e/c_H$, one has
\begin{equation}
 \left\|\frac{d\mu(W_S\mid W_{S^c})}{dH^S}\right\|_\infty
 \le [C_H(1+L)^3]^{|S|}
 \label{f16v74:joint-density}
\end{equation}
for every selected set $S$.  Independence of $S$ is not required.

For any fixed unit adjoint vectors $z,z'$ and any one selected link,
its conditional adjoint-image density relative to uniform probability on
$S^2$ is at most
\begin{equation}
 \alpha(L):=4e(1+L)^2.
 \label{f16v74:angular-constant}
\end{equation}
In particular, for $0\le\delta\le1$,
\begin{equation}
 \mathbb P\{|\operatorname{Ad}_{W_e}z-z'|\le2\delta
                  \mid W_{e^c}\}
 \le \min\{1,\alpha(L)\delta^2\}.
 \label{f16v74:angular-cap}
\end{equation}
Fixed left and right group factors preserve this bound.
\end{proposition}
\begin{proof}
The logarithmic derivative of a marginal is the conditional expectation
of the original logarithmic derivative.  A conditional normalizer is
constant in the variable being differentiated.  These observations give
the first assertion also for mixed marginal/conditional laws.

For a normalized one-link density $p$ with score at most $L$, choose a
maximum $x_0$.  On its geodesic ball of radius $(1+L)^{-1}$,
$p\ge e^{-1}p(x_0)$.  The Haar mass of that ball is at least
$c_H(1+L)^{-3}$, by the cap estimate already used in V64.
Normalization gives $p(x_0)\le C_H(1+L)^3$.
Factor a joint conditional density by the probability chain rule.
Each factor, with the unused selected variables integrated out, satisfies
this same bound, proving \eqref{f16v74:joint-density}.

The sharper exponent two in \eqref{f16v74:angular-constant} is important.
Disintegrate Haar measure over $W\mapsto\operatorname{Ad}_Wz$.  A rotation
between two image directions at spherical distance $d$ has a lift of
$S^3$ length $d/2$; left multiplication by that lift preserves normalized
Haar measure on their fibers.  Therefore the image density $q$ satisfies
$|\log q(y)-\log q(y')|\le Ld(y,y')/2$.  Around its maximum, take a
\emph{chordal} cap of radius $r=(1+L)^{-1}$.  Its normalized area is
$r^2/4$, and its spherical radius $2\arcsin(r/2)$ is at most $2r$.
On that cap $q\ge e^{-1}\|q\|_\infty$.  This proves
$\|q\|_\infty\le4e(1+L)^2$.  A chordal cap of radius $2\delta$
has normalized area exactly $\delta^2$, proving the last bound.
Haar isometries on the original link and rotations of its image preserve
all these constants.
\end{proof}

In the application $L=L_{\gamma,u}$, these are estimates for the
\emph{actual correlated law of the flowed links}.  Merely replacing the
old Wilson coefficient by a time-dependent number would not prove them.
Nor is the group-density exponent three needed for the following spectral
test: that test concerns adjoint directions and uses exponent two.

\subsection{A flowed native spectral estimate with its density cost included}

Reuse a block, rooted tree, and $r$ distinct non-tree links from V73.  Put
$K_r=K(n_r,d_r)$ with exactly the geometry in
\eqref{f16v73:uniform-constants}, and define
\begin{equation}
 \begin{split}
 A^{\rm fl}_{r,\gamma,u}
   &=64[4e(1+L_{\gamma,u})^2]^rK_r^{2r-2},\\
 \Theta^{\rm fl}_r(a,t;\gamma,u)
   &=\min\{3,3\Gamma(r)A^{\rm fl}_{r,\gamma,u}(a^2/t)^{r-1}\}.
 \end{split}
 \label{f16v74:flowed-constants}
\end{equation}
In particular
\begin{equation}
 A^{\rm fl}_{r,\gamma,u}
 \le C_r(6\gamma+3)^{2r}e^{32ru},
 \qquad C_r=64(4e)^rK_r^{2r-2}.
 \label{f16v74:coefficient-cost}
\end{equation}

\begin{theorem}[Native preflow small-ball and heat-weight bounds]
\label{f16v74:flowed-spectral}
For $W=\Phi_{a,\tau}(U)$ with its exact pushed slice law,
\begin{equation}
 \boxed{\quad
 \mu_\tau\{\lambda_{\min}(\mathsf L^N_{B,W})\le\epsilon\mid\Xi\}
 \le\min\{1,A^{\rm fl}_{r,\gamma,u}\epsilon^{r-1}\}.
 \quad}
 \label{f16v74:flowed-smallball}
\end{equation}
Here $\Xi$ fixes the flowed tree and all other unselected flowed links.
The selected links need not be conditionally independent.  Neumann
bracketing gives
\begin{align}
 \mu_\tau\{\lambda_{\min}(\mathcal L_{a,W})\le E\}
 &\le\min\{1,N_aA^{\rm fl}_{r,\gamma,u}(a^2E)^{r-1}\},
 \label{f16v74:flowed-bottom}\\
 \frac1{N_a}\mathbb E_{\mu_\tau}\Tr e^{-t\mathcal L_{a,W}}
 &\le\Theta^{\rm fl}_r(a,t;\gamma,u).
 \label{f16v74:flowed-trace}
\end{align}
For translation-invariant native laws, the same right side bounds
$\mathbb E_{\mu_\tau}\operatorname{tr}e^{-t\mathcal L_{a,W}}(x,x)$.
The averaged local estimates are uniform in the ambient volume.
\end{theorem}
\begin{proof}
The deterministic first half of V73's block argument is unchanged:
a block eigenvalue at most $\epsilon$ forces a unit root vector $z$
with $|z-\operatorname{Ad}_{H_e(W)}z|\le\delta=K_r\sqrt\epsilon$
for every selected chord.  Its $\delta/2$ sphere net has at most
$64/\delta^2$ points.  At a fixed net point, each chord constraint is a
cap of radius $2\delta$ in the image of its own selected link.
All tree transports are fixed by $\Xi$.

Apply \eqref{f16v74:angular-cap} to one link conditional on \emph{all}
other links.  The other cap indicators are measurable there.  Remove
that indicator by taking conditional expectation, bound it by
$\alpha(L)\delta^2$, and repeat for the other selected links.  The
conditional expectation tower bounds the simultaneous event by
$[\alpha(L)\delta^2]^r$ without factorizing their law.  The sphere net
gives $64\alpha(L)^r\delta^{2r-2}$.  If $\delta>1$ the stated coefficient
already exceeds one, and the trivial bound applies.  This proves
\eqref{f16v74:flowed-smallball}.  In fact any $r$ distinct non-tree chords
would work here; V73's separated construction is retained to avoid changing
its constants and bracketing geometry.

Delete cross-block edges only in the covariant quadratic form, keeping
$\mu_\tau$ unchanged.  The identical min--max and scalar integration
steps of Theorem~\ref{f16v73:trace} apply with
$A_{r,\gamma}$ replaced by $A^{\rm fl}_{r,\gamma,u}$.  This gives
\eqref{f16v74:flowed-bottom}--\eqref{f16v74:flowed-trace}.
Flow equivariance preserves translation invariance, proving the local
statement.  No operator monotonicity of the exponential is used.

For fixed $a,\gamma,\tau$, V70's finite-range ODE approximation makes
$U\mapsto\Phi_{a,\tau}U$ continuous in the local product topology,
uniformly on finite output sets.  Thus finite-torus pushed laws have the
corresponding local limits.  Block small-ball bounds pass first through
continuous approximants to their distribution functions, and then by
monotone limits; the same bounds hold at a closed threshold by approaching
it from above.  Polynomial approximation to the bounded covariant heat
operator proves passage of the local heat bound.  No cutoff-uniform
continuity of the inverse flow is required for this fixed-regulator passage.
\end{proof}

This theorem estimates an auxiliary \emph{spatial} spectrum after preflow.
It neither constructs a new time transfer nor estimates the bottom of
$(-\log P_a)/a$.

\subsection{The additional heat still suppresses logarithmically shallow preflow}

Let $s>0$ be fixed, allow $\tau=\tau_a\ge0$, and define
\begin{equation}
 B_{a,\tau,s}(U)
 =e^{-s\mathcal L_{a,W}}\zeta_a a^{-2}\operatorname{Im}W_{12},
 \qquad W=\Phi_{a,\tau}(U),\qquad
 O_{a,\tau,s,f}=\frac{a^3}{2}\sum_xf(x)|B_{a,\tau,s}(x)|^2.
 \label{f16v74:source}
\end{equation}
This is V68's coupled source, not a new definition of the current
$s=0$ observable.  Its statistics are evaluated under $\pi_{a,\gamma}$,
equivalently under $\mu_\tau$ for $W$.

\begin{proposition}[A local source estimate under the transported law]
\label{f16v74:source-bound}
With $N_R$ and $h_a(R,s)$ exactly as in V73, put
\[
 \mathcal B^{\rm fl}_r
 =N_R\Theta^{\rm fl}_r(a,2s;\gamma,u)+h_a(R,s)^2.
\]
For any possibly $W$-dependent insertion bounded by $M_a$,
\begin{equation}
 \mathbb E|e^{-s\mathcal L_{a,W}}\mathcal F(W;x)|^2
 \le2M_a^2\mathcal B^{\rm fl}_r.
 \label{f16v74:field-bound}
\end{equation}
For \eqref{f16v74:source},
\begin{equation}
 \operatorname{Var}_{\pi_a}(O_{a,\tau,s,f})
 \le\frac12\zeta_a^4a^{-8}\|f\|_{1,a}^2\mathcal B^{\rm fl}_r.
 \label{f16v74:variance-bound}
\end{equation}
\end{proposition}
\begin{proof}
Scalar heat domination and
$\sum_y\|K_s^W(x,y)\|_{\rm HS}^2=\operatorname{tr}K_{2s}^W(x,x)$
are pointwise identities for every compact connection.  V73's near/far
split therefore applies with the newly proved averaged local heat bound.
The curvature insertion can depend on all the same flowed links.
Finally $|B|\le M_a=\zeta_a a^{-2}$, so the supremum of $|O_f|$ times
its first absolute moment gives \eqref{f16v74:variance-bound}.
No conditional independence of the insertion and the connection is used.
\end{proof}

\begin{theorem}[An explicit excluded layer of spatial preflow]
\label{f16v74:shallow-collapse}
Suppose the spatial side lengths in lattice units tend to infinity,
$\log(1+\gamma_a)=o(\log(1/a))$, and
$\zeta_a\le C a^{-d}(1+\gamma_a)^b$ for fixed $b,d\ge0$.
If
\begin{equation}
 \boxed{\qquad
 \limsup_{a\downarrow0}
       \frac{\tau_a}{a^2\log(1/a)}<\frac1{16},
 \qquad}
 \label{f16v74:preflow-window}
\end{equation}
then for every $N>0$, every fixed $s>0$, and each finite family of tests
with bounded sampled $L^1$ norms,
\begin{equation}
 \mathbb E|B_{a,\tau_a,s}(x)|^2=o(a^N),\qquad
 \|D_a\|_{\rm op}=o(a^N),\qquad \|H_a\|_{\rm op}=o(a^N).
 \label{f16v74:collapse}
\end{equation}
Here $D_a$ is the actual centered entrance Gram and $H_a$ any matching
positive-time reflected Gram with common probability weights and an
additional nonnegative delay.  Polynomial coefficient maps are allowed.
\end{theorem}
\begin{proof}
Choose $c<1/16$ strictly above the limsup in
\eqref{f16v74:preflow-window}.  For every fixed chord count $r$,
\eqref{f16v74:coefficient-cost} gives
$A^{\rm fl}_{r,\gamma_a,u_a}\le a^{-32cr-o(1)}$.
Choose $R_a=\sqrt{4sL_0\log(1/a)}$ with $L_0$ a fixed large number.
The V73 walk estimate gives $h_a^2\le36a^{L_0}$ for sufficiently small
$a$, and $N_{R_a}=O(a^{-3}(\log(1/a))^{3/2})$.
The near terms in \eqref{f16v74:field-bound} and
\eqref{f16v74:variance-bound} have exponents at least
\[
 (2-32c)r-9-2d-o(1),\qquad
 (2-32c)r-13-4d-o(1),
\]
respectively.  Their coefficient of $r$ is positive.  Choose a
\emph{fixed} $r$ large enough for the requested $N$, and then a fixed
$L_0$ large enough for the far term.  A finite covariance matrix is bounded
by its trace.  Stationarity, Jensen, and native transfer/reflection
positivity give the stated absolute reflected estimates exactly as in
V73, without changing centering or normalization.
\end{proof}

The number $1/16$ is a sufficient exclusion constant, not a sharp critical
flow time.  The theorem permits $\tau_a/a^2\to\infty$ logarithmically.
For this $s>0$ bank, a uniform nonzero reflected lower norm under the
stated normalization assumptions would necessarily imply
\[
 \liminf_{a\downarrow0}\frac{\tau_a}{a^2\log(1/a)}\ge\frac1{16}:
\]
otherwise a subsequence would satisfy \eqref{f16v74:preflow-window}.
This necessary condition is not sufficient for survival.  Fixed positive
physical $\tau$ is outside the excluded layer, and the current $s=0$
bank is outside the theorem altogether.

\subsection{Native entropy proves that physical preflow changes conditional concentration}

The preceding upper score estimate is not an assertion that its
exponential factor is sharp.  A different calculation proves that an
exponential deterioration is genuinely necessary for the full native
flowed law in the weak-coupling regime.

\begin{theorem}[Exact entropy production and native conditional concentration]
\label{f16v74:entropy}
On the finite spatial torus write
\[
 \mathcal H_\tau=D(\mu_\tau\Vert H^M),\qquad
 \bar w_\tau=M^{-1}\mathbb E_{\mu_\tau}\sum_Pw(W_P).
\]
Then
\begin{equation}
 \boxed{\quad
 \frac{\mathcal H_\tau-\mathcal H_0}{M}
 =\frac{12}{a^2}\int_0^\tau\bar w_v\,dv,
 \qquad
 \bar w_v\ge\bar w_0.
 \quad}
 \label{f16v74:entropy-identity}
\end{equation}
Every individual flowed-link marginal is exactly Haar.  Nevertheless,
with the expectations below taken in the actual pushed law,
\begin{equation}
 \frac1M\sum_e\mathbb E
 D\bigl(\mu_\tau(W_e\mid W_{e^c})\Vert H\bigr)
 \ge\frac{\mathcal H_\tau}{M}
 \ge12u\bar w_0.
 \label{f16v74:conditional-entropy}
\end{equation}
Let $\mathcal C_\tau$ be the largest essential supremum, over links and
conditioned complements, of a one-link conditional density relative to
Haar.  Then
\begin{equation}
 \boxed{\quad
 \mathcal C_\tau\ge e^{12u\bar w_0},\qquad
 1+\ell(\rho_\tau)
 \ge C_H^{-1/3}e^{4u\bar w_0}.
 \quad}
 \label{f16v74:concentration-lower}
\end{equation}
For periodic four-dimensional Wilson measures with all sides at least four,
and their fixed-spatial-regulator stationary time limits,
\begin{equation}
 \bar w_0\ge1-\frac{2B_\gamma}{\gamma},\qquad
 B_\gamma=8+\log\gamma-\tfrac23\log c_H\quad(\gamma\ge1).
 \label{f16v74:initial-spatial-energy}
\end{equation}
Thus at fixed physical $\tau>0$, as $\gamma_a\to\infty$, the right sides
of \eqref{f16v74:concentration-lower} grow as
$\exp[(12-o(1))\tau/a^2]$ and $C_H^{-1/3}\exp[(4-o(1))\tau/a^2]$.
\end{theorem}
\begin{proof}
Change variables in the entropy integral by $W=\Phi_{a,\tau}U$.
Equation~\eqref{f16v74:Jacobian} gives exactly
\[
 \mathcal H_\tau
 =\mathcal H_0-\mathbb E_{\pi_a}\log\det D\Phi_{a,\tau}(U).
\]
This proves the identity.  The pointwise energy law of V68 makes
$E_{\rm sp}(\Phi_{a,v}U)$ nonincreasing, so $\bar w_v\ge\bar w_0$.
Also $\mathcal H_0\ge0$.  The lower bounds remain valid when
$\bar w_0<0$, but are then not useful concentration statements.

The native law is gauge invariant and spatial flow is gauge equivariant.
A gauge transformation at one endpoint of a selected link makes its marginal
left invariant, since that endpoint differs from the other one.  Its
normalized marginal is therefore Haar.  The entropy in the theorem is
thus dependence among links, not concentration of an individual link's
unconditional marginal.

Order the links.  The chain rule for relative entropy gives
\[
 \mathcal H_\tau=\sum_e\mathbb E
 D\bigl(\mu_\tau(W_e\mid\hbox{earlier links})\Vert H\bigr).
\]
Conditioning on all the other links can only increase each expected
conditional relative entropy, by convexity under conditional mixtures.
This proves \eqref{f16v74:conditional-entropy}.  Each earlier-link
conditional density is a mixture of the full-complement conditional
densities, and hence is bounded by $\mathcal C_\tau$.  Thus
$\mathcal H_\tau\le M\log\mathcal C_\tau$.
Proposition~\ref{f16v74:conditional-caps} also gives
$\mathcal C_\tau\le C_H(1+\ell(\rho_\tau))^3$.
Together these prove \eqref{f16v74:concentration-lower}.
If spatial translation and axis-permutation symmetries are present, all
terms on the left of \eqref{f16v74:conditional-entropy} agree, so its
lower bound holds for the expected conditional entropy at each link.

For the last assertion, the product-cap and entropy argument of V64
bounds the total four-dimensional mean defect by $N_PB_\gamma/\gamma$.
It uses only the counts $N_P=6N_v$, $N_L=4N_v$, not equality of side lengths.
All temporal defects are nonnegative; spatial defects account for
$3N_v$ plaquettes.  Time translation gives at each slice a mean spatial
defect at most $2B_\gamma/\gamma$ per spatial plaquette.
This proves \eqref{f16v74:initial-spatial-energy}.  The factor two is
retained rather than assuming a symmetry exchanging unequal time and
space lengths.  Taking the periodic-time limit at fixed spatial regulator
preserves the local expectation bound.
\end{proof}

In particular an exterior-uniform polynomial Haar-density bound for the
flowed links cannot survive a fixed positive physical spatial-flow time on
the usual logarithmic-coupling trajectory.  This is stronger than merely
not knowing a factorization: even the one-link conditional density has
necessarily large concentration.  The averaged conditional entropy shows
that the conclusion is not obtained by searching for a remote worst
plaquette.  It does not identify the correlations with those of a Wilson
measure at an effective coupling, or give a lower reflected covariance of
a centered curvature source.

\subsection{An exact compact diagnostic for the two density scales}

\begin{proposition}[An exact four-link subflow with Haar marginals]
\label{f16v74:one-link-example}
For an isolated cyclic four-link plaquette $Q=X_1X_2X_3X_4$, use
$E=1-w(Q)$ and flow all four links by its negative product-metric gradient
in dimensionless time $u$.  Start from independent Haar links.  If
$C=\cosh(4u)$ and $S=\sinh(4u)$, its exact pushed density is
\begin{equation}
 p_u(Y_1,\ldots,Y_4)=(C-Sw(Y_1Y_2Y_3Y_4))^{-3},\qquad
 \|p_u\|_\infty=e^{12u},\qquad
 \ell(p_u)=3\sinh(4u).
 \label{f16v74:pinned-density}
\end{equation}
Each individual link still has Haar marginal.  The flow Jacobian at its
input is $(C+Sw(Q))^{-3}$.  Consequently the cubic power of $1+L$ in a
general conditional score-to-group-density bound cannot be replaced by
a lower power with a universal constant.
\end{proposition}
\begin{proof}
Write $Q=w+p$.  The four right velocities, transported to its basepoint,
are all $-p$.  Hence $Q'=-4pQ$ and $w'=4(1-w^2)$.
It follows that the output scalar is
$(w+\tanh4u)/(1+w\tanh4u)$.  The four-link divergence is $-12w$,
so integration gives the stated Jacobian.  Inverting the scalar map
gives \eqref{f16v74:pinned-density}.  The product of independent Haar
links is Haar; integrating any one of the other three links normalizes
this density for every fixed value of a selected link.  Thus all four
single-link marginals are Haar.

For each link, $|\nabla_e w(Q)|=\sqrt{1-w(Q)^2}$.  The conditional
logarithmic derivative has norm
$3S\sqrt{1-w^2}/(C-Sw)$, maximized at $w=S/C$ with value $3S$.
The density maximum occurs at $w=1$.  Its conditional normalization is
one with the other three links fixed, so the same density and score
statements hold for every one-link conditional law.
\end{proof}

This is an exactly solved compact plaquette subflow, not the simultaneous
spatial-torus flow or its four-dimensional vacuum.  It independently
checks the Jacobian sign, normalization, conditional correlation and
concentration powers.  The native full-flow lower estimate is
Theorem~\ref{f16v74:entropy}, not an inference from this example.

\subsection{Next native comparison and result ledger}

The formerly unspecified pushed-law input now has two quantitative sides.
For a logarithmically shallow preflow followed by fixed constituent heat,
its conditional angular concentration is still too weak to support the
candidate: Theorem~\ref{f16v74:shallow-collapse} proves native collapse.
At fixed physical preflow, entropy forces strong conditional correlations
and precludes transporting the old polynomial density bound.  Neither
statement says that the present $s=0$ source has a nonzero continuum norm.

For that source the target remains the covariance of
$O(\Phi_{a,\tau}U_t)$ and its positively separated reflected copy under
the same native history.  Working in flowed variables must retain the
joint pushforward of the two slices, not merely their one-slice
$\rho_\tau$.  The determinant above changes each slice's coordinate
density; it does not decouple the slices or bound the connected mixed
source derivative.  The native entropy is an averaged statement about
conditional link dependence, not a physical-time contraction or a source
Gram lower bound.  A useful next estimate must measure the particular
flowed curvature combinations in that joint law, rather than again
postulate a Wilson conditional law or mistake a smoothing upper bound
for source survival.

The companion diagnostic is
\begin{center}\small
\path{verify_ym_f16_v74_transported_density.py}.
\end{center}
It checks exact round-Laplacian and exponent identities, full compact
lattice Hessian traces and incidence, inverse-flow density gradients,
Jacobian evolution, conditional angular laws with dependent links,
finite probability entropy chains, four-link pushed densities, and the
stated scaling windows.  Random compact configurations and small auxiliary
probability models are diagnostics, not Wilson-vacuum samples.  No
numerical outcome is a premise of the analytic proofs.  There is no
outward-rounded interval or proof-assistant certification.  Removing this
delimited appendix and reverting the single displayed version change
recovers V73 byte for byte.

\begin{record}[V74 result ledger]
\label{f16v74:ledger}
\emph{Proved here:} the normalized native spatial-flow Jacobian and a
volume-independent local score bound; conditional group and adjoint-cap
bounds without conditional independence; a flowed-law block spectral
estimate with all density costs paid; local heat and correlated-insertion
bounds under that law; superpolynomial collapse for the specified
logarithmically shallow preflow followed by fixed constituent heat; exact
native entropy production, Haar one-link marginals and quantitatively
necessary conditional concentration after physical preflow; and an exact
four-link check of the density scales.

\emph{Inherited or established background:} spatial Wilson flow and its
finite-regulator invertibility, V69's exact Hessian blocks, V73's
common-direction geometry and min--max/local-heat arguments, V64's
product-cap pressure estimate, the general flow-Jacobian formula and
ordinary relative-entropy identities.  Earlier archive claims are not
recertified.

\emph{Not asserted:} independence of flowed links; a Wilson pushed law at
another coupling; a sharp $1/16$ transition; collapse at fixed physical
preflow or of the $s=0$ bank; a lower physical source norm from entropy;
or control of a two-slice connected covariance from one-slice densities.

\emph{Still unproved:} a nonzero cutoff-uniform native reflected Gram for
the actual positive-physical-preflow bank, its physical normalization and
joint source-directed estimates, common dense-core physical-time
contraction, and the interacting continuum Yang--Mills mass gap.
\end{record}

\begin{firewall}
The Jacobian and native normalization are retained.  A bound for correlated
conditional laws replaces, rather than assumes, independence.  The excluded
preflow layer and the physically flowed high-concentration law are distinct
regimes.  Entropy measures conditional dependence, not the desired spectral
gap.  V1--V73 remain unchanged.
\end{firewall}
% END F16 V74 APPEND
% BEGIN F16 V75 APPEND -- 2026-09-19
% Append-only continuation of the Post-Fifteenth-Archive V74.
\clearpage
\section[V75: two-time invariance and native antipodal source weights]{V75: two-time invariance\\and native antipodal source weights}
\label{f16v75:context}

V74 determines the one-slice transported density but leaves the two-time
source covariance open.  We first separate the part of that density which
can carry temporal information: the two Jacobians cancel exactly from the
joint density relative to its own slice marginals.  The physical transfer
is only unitarily conjugated.  In particular the extensive entropy increase
of V74 cannot itself improve a temporal correlation.

A source-dependent calculation is then possible for the exact four-link
\emph{subflow} used at the end of V74.  This time the initial data are not
independent Haar links: the observable is evaluated under the full native
four-dimensional Wilson history.  Its ordinary and reflected Grams have
different powers of the subflow contraction.  An explicit non-polynomial
renormalization has a finite, strictly positive reflected limit at each
fixed finite regulator, while its ordinary variance diverges.  The latter
is generated by an antipodal layer in the original links.  All native
covariance-reduction errors are bounded uniformly in the surrounding volume
and in integer time separation at least two.

The isolated subflow is a test of a local replacement, \emph{not} a change of
definition of the simultaneous spatial flow in \eqref{f16v69:source}.  Its
failure under polynomial normalization does not prove failure of that
source.  Conversely its fixed-regulator singular reflected limit is not an
interacting continuum limit.  The distinction is needed before using either
entropy or smoothness as evidence of source survival.

\subsection{The joint transported law, rather than its two coordinate densities}

Fix a finite spatial regulator and finite $\gamma>0$.  Write
$\pi=\rho\,dH^M$ for the native stationary slice law and $P$ for its positive,
self-adjoint Markov transfer.  For an integer $n\ge1$ write
\[
 d\Pi_n(U,V)=R_n(U,V)\,d\pi(U)d\pi(V).
\]
At this fixed regulator the densities and the positive-time kernel are
smooth and strictly positive on the compact slice space.  Gauge-invariant
functions are the physical slice algebra.  If a Gauss projection is used
in the kernel representation, every statement below is restricted to that
algebra.  The statements are not about a conditional tube transfer.

Let $\Phi=\Phi_{a,\tau}$ be V68's simultaneous spatial flow at a finite time,
$\mu=\Phi_*\pi$, and $\Pi_n^\Phi=(\Phi,\Phi)_*\Pi_n$.  Its Jacobian is the
one in \eqref{f16v74:Jacobian}.  The change-of-variables method is established
background; see L\"uscher, \emph{Trivializing maps, the Wilson flow and the HMC
algorithm}, \href{https://arxiv.org/abs/0907.5491}{arXiv:0907.5491}.
We use only the finite compact identity proved here.  Spatial-only flow as
an observable construction is discussed by Sakai and Sasaki,
\href{https://arxiv.org/abs/2211.15176}{arXiv:2211.15176}; no numerical or
continuum conclusion from that work is used as an input.

\begin{theorem}[Exact two-time invariants of the full spatial flow]
\label{f16v75:joint}
The density of the joint pushed law relative to its own marginals is
\begin{equation}
 \boxed{\quad
 \frac{d\Pi_n^\Phi}{d(\mu\otimes\mu)}(W,Z)
   =R_n(\Phi^{-1}W,\Phi^{-1}Z).
 \quad}
 \label{f16v75:ratio}
\end{equation}
Define the unitary map $\mathcal U:L^2(\mu)\to L^2(\pi)$ by
$\mathcal U f=f\circ\Phi$.  The Markov transfer in the new slice coordinates
is exactly $\widetilde P=\mathcal U^{-1}P\mathcal U$.  Thus the centered
operator norms, the physical spectrum of $(-\log P)/a$, and
$D(\Pi_n\Vert\pi\otimes\pi)$ are unchanged.  In contrast, their Haar-relative
entropies obey
\begin{equation}
 D(\Pi_n^\Phi\Vert H^{2M})-D(\Pi_n\Vert H^{2M})
       =2\{D(\mu\Vert H^M)-D(\pi\Vert H^M)\}.
 \label{f16v75:joint-entropy}
\end{equation}
With V74's conventions, the right side is
$24a^{-2}\int_0^\tau\mathbb E_\pi\sum_Q w((\Phi_{a,v}U)_Q)\,dv$.
\end{theorem}
\begin{proof}
If $J(U)=\det D\Phi(U)>0$, the joint Haar density at $(W,Z)$ is the old
joint Haar density at $(U,V)$ divided by $J(U)J(V)$.  The two marginal
Haar densities are $\rho(U)/J(U)$ and $\rho(V)/J(V)$.  Dividing gives
\eqref{f16v75:ratio} exactly.  Conditional expectation, or integration of
this identity against a test function, gives the stated conjugation.
It preserves constants and the gauge-invariant subspaces, since the flow
is gauge equivariant.  Spectral calculus gives the Hamiltonian statement;
no physical time rescaling has been made.

Changing variables in the integral of $R_n\log R_n$ proves invariance of
the relative entropy to the product marginals.  Alternatively use the
exact entropy decomposition
\[
 D(\Pi_n\Vert H^{2M})
  =D(\Pi_n\Vert\pi\otimes\pi)+2D(\pi\Vert H^M),
\]
which follows by splitting the logarithm of its density and using
stationarity.  Apply it before and after the map to obtain
\eqref{f16v75:joint-entropy}.  All entropies here are finite at the fixed
compact regulator.  The coefficient twenty-four follows from V74's
one-slice identity, not from a new temporal interaction.
\end{proof}

For a fixed formula $O$ on the \emph{output} links, the relevant source is
$O\circ\Phi$ on the original slice; this function changes with $\tau$.
Unitary equivalence therefore does not make its particular covariance
constant.  It proves instead that any change belongs to this source's
projection on the unchanged transfer, not to two independent density
Jacobians.  No lower Gram follows from \eqref{f16v75:joint-entropy}.

\subsection{The inherited plaquette subflow and its exact Haar moments}

Take a nondegenerate spatial plaquette $Q=X_1X_2X_3X_4$ with cyclic
orientations.  Flow its four links by the negative gradient of $1-w(Q)$,
leaving all other terms out of this \emph{observable's subflow}.  This is
exactly Proposition~\ref{f16v74:one-link-example}, now used with arbitrary
native input links, not a Haar statistical ansatz.  Write $u$ for its
dimensionless time, $\epsilon=e^{-4u}\in(0,1]$, and $w=w(Q)$.  Its scalar
ODE $w'=4(1-w^2)$ gives
\begin{equation}
 A_\epsilon(Q):=|\operatorname{Im}Q(u)|^2
  =\frac{4\epsilon^2(1-w^2)}{
       [1+w+\epsilon^2(1-w)]^2}.
 \label{f16v75:A}
\end{equation}
The value is zero at $Q=\pm I$, and the formula extends smoothly there for
every $\epsilon>0$.  A physical subflow parameter $\tau$ corresponds to
$u=\tau/a^2$.  This is not the simultaneous flow of all spatial plaquettes.

\begin{proposition}[Exact contact moments and an integrable singular limit]
\label{f16v75:Haar}
For normalized Haar measure on $\SU(2)$, set $m_j(\epsilon)=\int A_\epsilon^j
dH$ and $v_H=m_2-m_1^2$.  Exactly,
\begin{align}
 m_1(\epsilon)&=\frac{12\epsilon^2}{(1+\epsilon)^4},&
 m_2(\epsilon)&=\frac{40\epsilon^3}{(1+\epsilon)^6},
 \label{f16v75:moments}\\
 v_H(\epsilon)&=
 \frac{8\epsilon^3(5-8\epsilon+5\epsilon^2)}{(1+\epsilon)^8}
                 \ge\frac{\epsilon^3}{16}.
 \label{f16v75:vH}
\end{align}
Define $g(Q)=4(1-w(Q))/(1+w(Q))$ off $-I$.  Then $g\in L^1(H)$ but
$g\notin L^2(H)$, and
\begin{equation}
 0\le\epsilon^{-2}A_\epsilon\le g,\qquad
 \int g\,dH=12,\qquad
 \int(g-\epsilon^{-2}A_\epsilon)dH
 =:\delta_\epsilon=12[1-(1+\epsilon)^{-4}]\le48\epsilon.
 \label{f16v75:L1}
\end{equation}
No value assigned at the Haar-null point $-I$ affects these statements.
\end{proposition}
\begin{proof}
Put $r=\tan(\theta/2)$ when $w=\cos\theta$.  The radial Haar density and
observable become
\[
 \frac{16}{\pi}\frac{r^2}{(1+r^2)^3}\,dr,
 \qquad A_\epsilon=\frac{4\epsilon^2r^2}{(1+\epsilon^2r^2)^2}.
\]
Consequently
\[
 m_j=\frac{16\,4^j\epsilon^{2j}}{\pi}
       \int_0^\infty
       \frac{r^{2j+2}\,dr}{(1+r^2)^3(1+\epsilon^2r^2)^{2j}}.
\]
For $j=1,2$, partial fractions followed by
$\int_0^\infty(1+r^2)^{-k}dr
=\sqrt\pi\,\Gamma(k-1/2)/(2\Gamma(k))$
give \eqref{f16v75:moments}; equivalently the even rational integral is
$\pi i$ times the sum of its residues at $i$ and $i/\epsilon$.
The coincident-pole case $\epsilon=1$ follows by continuity.  Subtraction
gives $v_H$.  The derivative of $v_H/\epsilon^3$ has the sign of
$-6(8-11\epsilon+5\epsilon^2)$, which is negative.  Its minimum on
$(0,1]$ is its value $1/16$ at one.

In these coordinates $g=4r^2$.  Its Haar first moment is twelve; its
second moment diverges because the radial integrand tends to a positive
constant at infinity.  The pointwise comparison and the exact $L^1$
error follow from the denominator $(1+\epsilon^2r^2)^2$ and the formula
for $m_1$.  Finally $1-(1+\epsilon)^{-4}\le4\epsilon$.
\end{proof}

\subsection{The same subflow observables under the full native history}

Choose $d$ spatial plaquettes $Q_1,\ldots,Q_d$.  For each choose one boundary
link $e_j$ absent from every other writer, with the selected links jointly
plaquette-separated.  Disjoint plaquettes in separated fixed-size boxes
suffice.  The choices and their translates to times $0,n$ remain jointly
separated for $n\ge2$.  All unselected Wilson plaquettes, temporal links,
and boundary data remain in the measure.  We use the original native
stationary history, not the isolated subflow as a sampling algorithm.

Set
\begin{equation}
 K_\gamma=\frac{3\pi e}{4}(1+6\gamma)^{3/2},\qquad
 a_{\gamma,d}=e^{-12\gamma d}.
 \label{f16v75:constants}
\end{equation}
The symbol $K_\gamma$ equals the inherited $\kappa_\gamma$ in
\eqref{f16v73:density-constant}.  V72's conditional factorization gives,
for a single selected link, a translated one-link density bounded between
$e^{-12\gamma}$ and $K_\gamma$: the upper bound is
\eqref{f16v73:one-density}, and the lower bound follows from
$e^{h\cdot x}/\mathfrak z(|h|)\ge e^{-2|h|}$, $|h|\le6\gamma$.
Changing that link to its plaquette product is a Haar isometry.
It follows that the joint one-time law of $(Q_j)$ has density at least
$a_{\gamma,d}$ relative to $H^d$.  Each individual density is at most
$K_\gamma$.  These are consequences of the inherited conditional law,
not independence assertions about the native plaquettes.

Let $A_t=(A_\epsilon(Q_{j,t}))_{j=1}^d$ and define the native matrices
\[
 D_\epsilon=\operatorname{Cov}(A_0),\qquad
 C_{\epsilon,n}=\operatorname{Cov}(A_0,A_n).
\]
They use the same actual centering at the two stationary times.

\begin{theorem}[Explicit native entrance and positive-time Gram bounds]
\label{f16v75:native-Grams}
For every finite regulator and every $\epsilon\in(0,1]$,
\begin{align}
 a_{\gamma,d}v_H(\epsilon)I_d
       &\preceq D_\epsilon
        \preceq dK_\gamma m_2(\epsilon)I_d,
 \label{f16v75:D}\\
 0\preceq C_{\epsilon,n}
       &\preceq\frac d4 K_\gamma^2m_1(\epsilon)^2I_d
        \preceq36dK_\gamma^2\epsilon^4I_d,
       \qquad n\ge2.
 \label{f16v75:C}
\end{align}
In particular, for the reflected placement $n=2k$, $k\ge1$,
\begin{equation}
 \boxed{\quad
 0\preceq D_\epsilon^{-1/2}C_{\epsilon,2k}D_\epsilon^{-1/2}
 \preceq
 \min\{1,576dK_\gamma^2e^{12\gamma d}\epsilon\}\,I_d.
 \quad}
 \label{f16v75:relative}
\end{equation}
A declared redundant coefficient map pulls back all inequalities and
preserves its exact kernel.  No numerical eigenvalue threshold is used.
\end{theorem}
\begin{proof}
For real $c$, variance is the infimum of the second moment after subtracting
a constant.  The joint Haar-density lower bound gives
\[
 \operatorname{Var}_\pi(c^*A)
 \ge a_{\gamma,d}\inf_b\int|c^*A-b|^2dH^d
 =a_{\gamma,d}v_H|c|^2.
\]
The independent factors in this last integral are \emph{Haar reference}
variables, not native variables.  Each native second moment is at most
$K_\gamma m_2$, so the covariance trace gives the upper bound in
\eqref{f16v75:D}.

For the time comparison condition on the complement $\Xi$ of the selected
sets at both times.  These sets are jointly plaquette-separated.  All
components on the two time slices are conditionally independent across
the slices, and every conditional mean $a_{j,t}$ lies in
$[0,K_\gamma m_1]$.  Thus
\[
 \operatorname{Cov}(c^*A_0,c^*A_n)
       =\operatorname{Cov}(c^*a_0,c^*a_n).
\]
The range width of $c^*a_t$ is at most $K_\gamma m_1\|c\|_1$.
The variance of a real variable in an interval of width $L$ is at most
$L^2/4$; Cauchy--Schwarz and $\|c\|_1^2\le d|c|^2$ give the upper
bound.  Matching native lag covariances are nonnegative because $P\succeq0$.
This proves the matrix assertion.  Complex coefficients follow from the
real symmetric matrices by complexification.  Use $v_H\ge\epsilon^3/16$
and congruence for \eqref{f16v75:relative}; $C_{\epsilon,n}\preceq D_\epsilon$
also follows from the native contraction.  Strict positive definiteness
of $D_\epsilon$ proves the exact kernel statement.
\end{proof}

This is a source estimate in the full native law.  Unlike the full spatial
flow, the source formula here depends only on one input plaquette per
component, so an original selected link can still be integrated exactly.
Applying that same conditioning after silently replacing $A_\epsilon(Q)$
by the full $O(\Phi_{a,\tau}U)$ would not be justified.

\subsection{The contact norm is carried by an antipodal layer}

\begin{theorem}[Exact leading native contact weight]
\label{f16v75:antipodal}
At a fixed finite regulator let $p_j$ be the native Haar density of $Q_j$.
It is continuous and strictly positive.  As $\epsilon\downarrow0$,
\begin{equation}
 \boxed{\quad
 \epsilon^{-3}D_\epsilon
    \longrightarrow 40\operatorname{diag}(p_1(-I),\ldots,p_d(-I)),
 \qquad e^{-12\gamma}\le p_j(-I)\le K_\gamma.
 \quad}
 \label{f16v75:contact-limit}
\end{equation}
For any fixed neighborhood $N$ of $-I$,
\begin{equation}
 \frac{\mathbb E[A_\epsilon(Q_j)^2\mathbf1_{Q_j\in N}]}
      {\mathbb E A_\epsilon(Q_j)^2}\longrightarrow1.
 \label{f16v75:layer}
\end{equation}
For one component, an explicit uniform estimate for $0<\epsilon\le1$ is
\begin{equation}
 \left|\epsilon^{-3}\operatorname{Var}(A_\epsilon(Q_j))-40p_j(-I)\right|
 \le\frac{4096}{\pi}K_\gamma(1+\gamma)\epsilon[1+\log(1/\epsilon)]
        +144K_\gamma^2\epsilon.
 \label{f16v75:contact-error}
\end{equation}
The limit is at fixed regulator unless the displayed errors are separately
paid along a varying sequence.
\end{theorem}
\begin{proof}
Gauge invariance makes each $p_j$ a class function.  Put $y=\epsilon r$
in the radial integral for its second moment.  Exactly,
\[
 \epsilon^{-3}\mathbb E A_\epsilon(Q_j)^2
 =\frac{256}{\pi}\int_0^\infty
  \frac{y^6}{(\epsilon^2+y^2)^3(1+y^2)^4}
  p_j(\cos(2\arctan(y/\epsilon)))\,dy.
\]
The notation $p_j(\cos\theta)$ refers to its class value at angle $\theta$.
For each $y>0$ the group argument tends to $-I$.  The integrand is bounded
by $K_\gamma(1+y^2)^{-4}$, which is integrable.  Since
$\int_0^\infty(1+y^2)^{-4}dy=5\pi/32$, dominated convergence gives the
coefficient forty.  The squared mean is $O(\epsilon^4)$.
For distinct plaquettes the two selected links are conditionally
independent, so their product first moment is at most
$K_\gamma^2m_1^2=O(\epsilon^4)$.  This proves \eqref{f16v75:contact-limit}.
The same dominated convergence with the indicator of the complement of
$N$ proves \eqref{f16v75:layer}.

For the error, the individual conditional translated densities have
logarithmic gradient at most $6\gamma$.  Their mixture has
$|\nabla p_j|\le6\gamma K_\gamma$; here the complementary variables are
integrated using their original marginal, independent of the evaluation
argument $Q_j$.  Thus for $y\ge\epsilon$ the difference from $p_j(-I)$
is at most $12\gamma K_\gamma\epsilon/y$.  On $0<y<\epsilon$ use
$2K_\gamma$.  Also
\[
 \int_0^\infty\frac{1-y^6/(\epsilon^2+y^2)^3}{(1+y^2)^4}dy
 \le\int_0^\infty\min\{1,3\epsilon^2/y^2\}\,dy\le4\epsilon.
\]
Splitting the remaining integral at $\epsilon$ and at one bounds it by
$K_\gamma\epsilon[2+12\gamma\log(1/\epsilon)+3\gamma/2]$.
Multiplication by $256/\pi$, and a harmless enlargement to the constant
in \eqref{f16v75:contact-error}, pay these two errors.  Finally
$\epsilon^{-3}(\mathbb E A_\epsilon)^2\le144K_\gamma^2\epsilon$.
\end{proof}

In particular, typical forward concentration near $I$ does not describe
the normalized second moment of this source.  Its leading weight comes
from inputs near the unstable point $-I$, whose width is of order
$\epsilon$.  This assertion is source-weighted and is not an assertion
that an antipodal input is typical under the native law.

\subsection{A singular source has a controlled native reflected limit}

Define $g_j=g(Q_j)$ as in \eqref{f16v75:L1}, centered with its finite
native mean.  These functions are native $L^1$ but not native $L^2$.
For time separation $n\ge2$ the products at distinct times are integrable:
the joint two-time holonomy density is at most $K_\gamma^{2d}$, and the
one-component conditional-mean bound below is sharper.  Define
\[
 H^{\rm sing}_n=\operatorname{Cov}(g(Q_{j,0}),g(Q_{l,n}))_{j,l=1}^d.
\]
This definition uses the actual two-time law, not a product of slice laws.

\begin{theorem}[An explicit reflected renormalization with a paid covariance error]
\label{f16v75:singular}
For every $n\ge2$,
\begin{align}
 0\preceq H^{\rm sing}_n&\preceq36dK_\gamma^2I_d,
 \label{f16v75:sing-bound}\\
 \boxed{\quad
 \|\epsilon^{-4}C_{\epsilon,n}-H^{\rm sing}_n\|_{\rm op}
 \le6dK_\gamma^2\delta_\epsilon
 \le288dK_\gamma^2\epsilon.
 \quad}
 \label{f16v75:sing-error}
\end{align}
The constants are independent of $n\ge2$ and surrounding volume.  At each
fixed finite spatial regulator, $H^{\rm sing}_{2k}\succ0$ for every
$k\ge1$.  More precisely the centered observables
$\epsilon^{-2}A_\epsilon(Q_j)$ converge in the positive-time reflected
seminorm at distance $k$ to the singular family, with error Gram at most
\begin{equation}
 \frac d4K_\gamma^2\delta_\epsilon^2I_d.
 \label{f16v75:OS-error}
\end{equation}
They do not converge in ordinary $L^2$.  Their entrance Grams diverge as
$40\epsilon^{-1}\operatorname{diag}(p_j(-I))$ at fixed regulator.
\end{theorem}
\begin{proof}
Condition on the complement of the selected links at the two times.
From \eqref{f16v75:L1}, each conditional mean of $g_j$ and of
$\epsilon^{-2}A_\epsilon(Q_j)$ belongs to $[0,12K_\gamma]$, while their
nonnegative difference belongs to $[0,K_\gamma\delta_\epsilon]$.
Conditional independence makes the complete two-time covariance the
covariance of these predictors, including the original means.  For a
real coefficient vector $c$, their range widths are respectively at most
$12K_\gamma\|c\|_1$ and $K_\gamma\delta_\epsilon\|c\|_1$.
The interval variance inequality proves the bound $36dK_\gamma^2$.
Expanding the covariance difference into two terms, each with one
predictor difference, gives $6dK_\gamma^2\delta_\epsilon$.
Using two differences gives \eqref{f16v75:OS-error}.  Approximation by
bounded functions, with these integrable bounds, justifies every use of
total covariance.  Positivity follows by taking the limit of the native
positive lag Grams.

For strict positivity at even delay one can use the fixed-regulator
transfer directly.  Its smooth kernel maps $L^1(\pi)$ to smooth bounded
functions.  Thus $P^k g_c$ is well-defined and square integrable, where
$g_c=\sum_jc_j(g_j-\mathbb E g_j)$.  Semigroup composition, or bounded
approximation followed by the preceding estimates, gives
\[
 c^*H^{\rm sing}_{2k}c=\|P^kg_c\|_{L^2(\pi)}^2.
\]
Injectivity here extends to invariant $L^1$ inputs, not just to $L^2$.
Indeed the finite Wilson transfer is a positive smooth magnetic
multiplier, a product of central one-link convolutions with strictly
positive representation multipliers, and a second positive magnetic
multiplier, with the Gauss projection restricted to its invariant domain.
The multiplication factors and the Perron ground-state transform are
invertible on distributions at a fixed regulator.  If a convolution
annihilates an invariant distribution, every Peter--Weyl coefficient is
zero because none of its multipliers vanishes.  The distribution is
therefore zero.  Iterating proves injectivity of $P^k$ on $g_c$.
Finally $g_c$ cannot be constant unless $c=0$: vary its distinct private
selected links independently in the Haar reference measure.  The native
slice density is strictly positive, so the same independence modulo
constants holds almost everywhere natively.  This proves strict
positivity.  The ordinary-norm divergence follows from
\eqref{f16v75:contact-limit} after multiplication by $\epsilon^{-4}$.
\end{proof}

The coefficient $\epsilon^{-2}=e^{8u}$ is part of an explicitly different
source normalization.  This reflected construction is at finite regulator;
no cutoff-uniform lower eigenvalue of $H^{\rm sing}_{2k}$ at
$ka\to\delta>0$ has been established.  The singular source samples a
microscopic plaquette, not a proved physical continuum field.  Conversely,
normalizing by the ordinary variance would lose its entire reflected
limit by \eqref{f16v75:relative}.  The two conclusions use different,
fully specified normalizations.

\subsection{What is excluded on a physical subflow scale}

\begin{proposition}[Polynomially normalized local subflow sources collapse]
\label{f16v75:physical-collapse}
Let $u_a=\tau_a/a^2$, $\epsilon_a=e^{-4u_a}$, and suppose
$u_a/\log(1/a)\to\infty$ and
$\log(1+\gamma_a)=o(u_a)$.  Any finite source family formed from sums of
$A_{\epsilon_a}(Q)$ with total absolute coefficient mass bounded by a
fixed polynomial in $a^{-1}$ and $1+\gamma_a$ has native entrance and
matching reflected Grams $o(a^N)$ for every $N>0$.  The result allows
polynomially many marked plaquettes and overlapping marks.

For a fixed private $d$-plaquette bank, the reflected-to-entrance ratio
also tends to zero whenever
\begin{equation}
 4u_a-12d\gamma_a-2\log K_{\gamma_a}\longrightarrow+\infty.
 \label{f16v75:physical-ratio-condition}
\end{equation}
Both conclusions hold at fixed $\tau_a=\tau>0$ under the separately
specified trajectory $\log(1/a)=c\gamma_a+O(\log\gamma_a)$, $c>0$.
\end{proposition}
\begin{proof}
The marginal upper bound alone gives
$\mathbb E A_{\epsilon}(Q)^2\le40K_\gamma\epsilon^3$ for each mark.
For arbitrary coefficients the $L^2$ triangle inequality consequently
gives
\[
 \operatorname{Var}\left(\sum_Qb_QA_\epsilon(Q)\right)
 \le40K_\gamma\epsilon^3\left(\sum_Q|b_Q|\right)^2.
\]
No independence or disjointness of these marks is needed.  The logarithm
of the right side is $-12u_a+o(u_a)+O(\log(1/a))$, so it is smaller than
$N\log a$ for each fixed $N$.  Matrix traces bound finite-family Grams;
reflection and Cauchy--Schwarz bound matching reflected norms by ordinary
norms.  Common probability-weighted stationary temporal averages preserve
these bounds.  The ratio assertion follows from \eqref{f16v75:relative}.
At fixed positive physical subflow time the stated logarithmic trajectory
makes $u_a$ larger than $\gamma_a$ and $\log(1/a)$ by a diverging factor.
\end{proof}

The spatial sum in this proposition evaluates a separate scalar subflow
function of each input plaquette.  It is not a simultaneous overlapping
plaquette evolution.  For the current full spatial flow, curvature is
transported through a number of links growing in physical units; the
single-private-link conditioning used above does not evaluate that source.
The proposition excludes a local replacement, not V69's current bank.

\subsection{Noncircular endpoint, diagnostics, and result ledger}

The one-slice density route has now been separated from temporal dynamics
at the level of exact native measures.  V74's entropy increase is paid
identically by both coordinate densities; the joint dependence kernel
retains the original physical transfer.  The source-marked subflow example
then evaluates a genuinely different question: its native contact weight
is antipodal, its separated covariance is one extra power smaller, and
its explicitly renormalized reflected remainder is uniformly controlled.
Neither an independent Haar endpoint law nor an isolated temporal transfer
was substituted for the native history in these estimates.

For the simultaneous spatial-flow source the remaining task is still its
positive-separation native Gram.  A proposed reduction to independent local
plaquette flows is excluded by Proposition~\ref{f16v75:physical-collapse}.
A valid comparison must retain the spatially collective part of the source
and the unchanged joint dependence kernel of \eqref{f16v75:ratio}.  It must
also control source-weighted tails rather than infer second moments from
one-slice concentration.  The new singular microscopic bank is not asserted
to solve that physical problem.  No further unmarked Jacobian, entropy,
or conditional-density estimate will by itself populate its missing mixed
source derivative.

The companion diagnostic is
\begin{center}\small
\path{verify_ym_f16_v75_joint_source_weights.py}.
\end{center}
It separates exact rational/Haar identities, finite probability
change-of-variables and entropy checks, literal compact plaquette-flow
identities, density-weighted moment quadrature, and positive convolution
benchmarks.  A benchmark convolution is not the full Wilson vacuum.
Numerical values are not premises of the analytic arguments, and no
outward-rounded interval or proof-assistant verification is claimed.
Removing this delimited appendix and reverting the displayed version
recovers V74 byte for byte.

\clearpage
\begin{record}[V75 result ledger]
\label{f16v75:ledger}
\emph{Proved here:} exact two-time likelihood-ratio and transfer invariance
under the simultaneous spatial flow; its joint entropy balance; exact Haar
moments for the inherited local subflow observable; native all-direction
entrance and reflected bounds for that observable; its antipodal contact
asymptotic with an explicit error; a singular reflected source limit with
uniform two-time covariance and seminorm errors; finite-regulator strict
positivity of that reflected limit; and collapse of polynomially normalized
local-subflow sums at the stated physical subflow scale.

\emph{Inherited or elementary background:} the finite compact Wilson
transfer and conditional-link law, the positive one-link Fourier
multipliers, V74's spatial-flow diffeomorphism/Jacobian and isolated
plaquette scalar ODE, Haar integration, unitary coordinate change and the
entropy chain rule.  These general methods are not claimed as discoveries.
No prior archive is recertified.

\emph{Not asserted:} disappearance of temporal dependence under an
invertible map; equality of isolated and simultaneous spatial flow;
collapse of the current $s=0$ full-flow source; a cutoff-uniform reflected
lower bound for the singular microscopic bank; a mass gap from its fixed
regulator positivity; or interchange of the long-subflow and continuum
limits without the displayed error and scale conditions.

\emph{Still unproved:} a nonzero physical-scale reflected Gram for the
simultaneously flowed native source, its normalization and collective
spatial control, common dense-core physical-time contraction, and the
interacting four-dimensional continuum Yang--Mills mass gap.
\end{record}

\begin{firewall}
The two Jacobians cancel only from the joint dependence ratio, not from an
arbitrary observable integral.  Antipodal source weight is computed under
the native law.  A finite reflected limit after exponential subflow
renormalization is not a continuum physical-source theorem.  The original
physical transfer and all V1--V74 mathematical text remain unchanged.
\end{firewall}
% END F16 V75 APPEND

% BEGIN F16 V76 APPEND -- 2026-09-19
% Append-only continuation of the Post-Fifteenth-Archive V75.
\clearpage
\section[V76: singular-source matching and finite character reconstruction]{V76: singular-source matching\\and finite character reconstruction}
\label{f16v76:context}

V75 produces a nonzero reflected limit at each fixed regulator by an
exponentially large normalization of an isolated plaquette subflow.  It
also proves that this limit is the singular function
\(g(Q)=4(1-w(Q))/(1+w(Q))\), not the simultaneously flowed observable.
Before developing that singular bank as another continuum candidate, we
identify its weak-coupling reflected class under the original native law.
The antipodal divergence of its ordinary second moment is not ignored.
It is estimated at two separated times, where the product is integrable.
The resulting native theorem identifies its thermally normalized reflected
covariance with that of twice the elementary Wilson action density, with an
error uniform in the integer separation and surrounding volume.

A second result gives an explicit finite character reconstruction of the
singular source.  Its reflected error is exponentially small in the
representation cutoff once that cutoff exceeds a specified multiple of
\(\gamma\).  This controls growing physical-size sums of these microscopic
sources without using an unproved flow expansion.  The reconstruction is
not ordinary \(L^2\) convergence.  Neither result supplies the remaining
positive-separation lower Gram for the simultaneous spatial-flow source.

\subsection{Native inputs, scope and the separated-source norm}

We use the original isotropic four-dimensional \(\SU(2)\) Wilson measure
and its stationary cylinder limits.  For estimates using reflection
blocks, take the compatible rectangular dyadic tori and limits specified
in V65, with side lengths at least twice the block lengths used below.
The time side may tend to infinity first.  The finite spatial regulator is
kept until an explicitly uniform estimate is passed to a local limit.
No arbitrary boundary-conditioned state is assigned these probability
bounds.  The character reconstruction, in contrast, only needs the native
single-link conditional law and is valid with the same admissible local
boundary conventions as V72--V75.

The reflection/chessboard and partition-function inputs are precisely
\eqref{f16v65:chessboard-input} and \eqref{f16v65:cap-input}; their general
method is inherited.  For external context see Biskup,
\emph{Reflection Positivity and Phase Transitions in Lattice Spin Models},
\href{https://arxiv.org/abs/math-ph/0610025}{arXiv:math-ph/0610025},
Theorem 5.8.  The character formulas below use the elementary Chebyshev
identity recorded in \href{https://dlmf.nist.gov/18.5.E2}{NIST DLMF,
18.5.2}; the source-specific coefficients and native error bounds are
proved here.  No continuum renormalization theorem is imported.

Write \(d(Q)=1-w(Q)\in[0,2]\), and retain
\begin{equation}
 c_H=\frac8{3\pi^3},\qquad
 B_\gamma=8+\log\gamma-\frac23\log c_H,\qquad
 b_\gamma=384B_\gamma,\qquad
 K_\gamma=\frac{3\pi e}{4}(1+6\gamma)^{3/2}.
 \label{f16v76:constants}
\end{equation}
Here \(K_\gamma\) is V75's conditional-density constant.  Selecting one
boundary link of a plaquette and changing that link to its plaquette
product is a Haar isometry.  Conditional on the other links, its density
is bounded between \(e^{-12\gamma}\) and \(K_\gamma\).  At two time
slices separated by \(n\ge2\), the two selected links belong to no common
plaquette.  Their joint conditional density is a product.  In particular,
\begin{equation}
 \frac{d\operatorname{Law}(Q_0,Q_n)}{d(H\otimes H)}\le K_\gamma^2.
 \label{f16v76:two-density}
\end{equation}
Only these selected links are conditionally independent; the complement
keeps its full interacting native distribution.

For a real one-plaquette function \(h\in L^1(H)\), let
\(h^\circ=h(Q)-\mathbb E_\pi h(Q)\).  At fixed finite spatial regulator
\(P h^\circ\) is well-defined in \(L^2(\pi)\), because the native
positive-time kernel is smooth.  Put
\begin{equation}
 \|h\|_+^2:=\|P h^\circ\|_{L^2(\pi)}^2
 =\operatorname{Cov}(h(Q_0),h(Q_2)).
 \label{f16v76:plus-norm}
\end{equation}
All subsequent uses of singular functions are justified by truncation and
\eqref{f16v76:two-density}.  For \(n\ge2\),
\begin{equation}
 \operatorname{Cov}(h(Q_0),j(Q_n))
 =\langle P h^\circ,P^{n-2}P j^\circ\rangle.
 \label{f16v76:separated-transfer}
\end{equation}
Thus \(\|\cdot\|_+\) controls every larger separation.  It is an
unscaled two-step reflected norm, not an entrance norm and not a choice
of the physical time unit.  The same estimates apply to a reflection
placement \(k\ge1\) and to common probability-weighted positive-time
averages.  In local limits the right-hand pairings and their bounded
approximations define the corresponding reflected forms.

\subsection{Native defect moments and singular tails are paid separately}

\begin{proposition}[One marked plaquette: exponential tail and all fixed moments]
\label{f16v76:tail}
For \(\gamma\ge1\) in the reflection-compatible native laws above and
\(0\le t\le2\),
\begin{equation}
 \pi\{d(Q)\ge t\}\le\min\{1,e^{b_\gamma-\gamma t}\}.
 \label{f16v76:defect-tail}
\end{equation}
For each positive integer \(m\), define
\begin{equation}
 M_{m,\gamma}=\gamma^{-m}
   \sum_{j=0}^m\binom mj b_\gamma^{m-j}j!.
 \label{f16v76:moment-constant}
\end{equation}
Then \(\mathbb E_\pi d(Q)^m\le M_{m,\gamma}\).  In particular
\[
 M_{2,\gamma}=\frac{b_\gamma^2+2b_\gamma+2}{\gamma^2},\qquad
 M_{4,\gamma}=\frac{b_\gamma^4+4b_\gamma^3+12b_\gamma^2+24b_\gamma+24}
                    {\gamma^4}.
\]
The constants are independent of surrounding volume and time-cylinder
length.  They are large sufficient constants, not moderate-coupling
numerical predictions.
\end{proposition}
\begin{proof}
Place the marked plaquette at the central square of a closed block of
lengths \((4,4,2,2)\): its first two coordinates run from 1 to 2 and its
other coordinates are 1.  Every one of its edges is strictly inside the
block.  Distinct alternating reflected copies therefore have disjoint
witness plaquettes.  If the threshold event is disseminated to all
\(N_b\) blocks, the complete Wilson energy is at least \(tN_b\).
The normalized disseminated probability is at most
\(Z^{-1}e^{-\gamma tN_b}\).  The inherited chessboard inequality takes
its \(N_b\)-th root, while the cap bound pays
\(Z^{-1/N_b}\le e^{6\cdot64B_\gamma}\).  This proves
\eqref{f16v76:defect-tail} without a product law for blocks.  Continuous
majorants at nearby thresholds pass the bound to the stated local limits.

The tail of \(Y=\gamma d(Q)\) is dominated by that of
\(b_\gamma+Z_1\), where \(Z_1\) is a unit-rate exponential random
variable.  Integrating the tail, or expanding the moments of that latter
variable, gives \eqref{f16v76:moment-constant}.  This comparison is a scalar
distribution bound, not a claim that a plaquette actually has that law.
\end{proof}

The moment estimate alone cannot be applied to \(g^2\), which is not
integrable at one time.  The next estimate instead uses both actual time
slices before estimating the singular part.

\begin{proposition}[A native reflected bound for an integrable singular tail]
\label{f16v76:singular-tail}
Let \(p>1\), \(h\in L^p(H)\), and put
\(h_B(Q)=h(Q)\mathbf1_{\{d(Q)>1\}}\).  With
\(p_\gamma=\min\{1,e^{b_\gamma-\gamma}\}\),
\begin{equation}
 \boxed{\quad
 \|h_B\|_+\le
 K_\gamma^{1/p}\|h\|_{L^p(H)}
                  p_\gamma^{(p-1)/(2p)}.\quad}
 \label{f16v76:tail-plus}
\end{equation}
The same bound controls the square root of the matching covariance at
any \(n\ge2\).  No one-time \(L^2\) assumption is required.
\end{proposition}
\begin{proof}
Let \(B_0,B_n\) be the two bad events.  Their intersection has probability
at most \(p_\gamma\), not a postulated product probability.  H\"older's
inequality under the full native two-time law and
\eqref{f16v76:two-density} give
\[
 \mathbb E|h(Q_0)h(Q_n)|\mathbf1_{B_0\cap B_n}
 \le K_\gamma^{2/p}\|h\|_p^2p_\gamma^{1-1/p}.
\]
The matching centered covariance is nonnegative by native transfer
positivity and is no larger than this absolute product moment: its
subtracted squared mean is nonnegative.  Bounded truncations justify
positivity for the integrable limit.  Take square roots at \(n=2\);
\eqref{f16v76:separated-transfer} gives the remaining delays as well.
\end{proof}

\subsection{The reflected thermal class is fixed by the identity expansion}

\begin{theorem}[Native first-order matching with a singular remainder]
\label{f16v76:jet}
Let \(h\) be a real class function with \(h(Q)=h_0+c\,d(Q)+r(Q)\), where
\(r\in L^p(H)\) for a fixed \(p>1\), and assume
\(|r(Q)|\le C_r d(Q)^2\) when \(d(Q)\le1\).
Define
\begin{align}
 J_\gamma&=\gamma\sqrt{M_{2,\gamma}},
 \label{f16v76:J}\\
 E_\gamma(h)&=\gamma\left[
 C_r\sqrt{M_{4,\gamma}}
 +K_\gamma^{1/p}\|r\|_p p_\gamma^{(p-1)/(2p)}\right].
 \label{f16v76:E}
\end{align}
Then
\begin{equation}
 \boxed{\quad
 \|\gamma h-c\gamma d\|_+\le E_\gamma(h),\qquad
 \|\gamma d\|_+\le J_\gamma.\quad}
 \label{f16v76:jet-plus}
\end{equation}
For fixed \(h,p,C_r\), \(E_\gamma(h)=O((\log\gamma)^2/\gamma)\),
up to an explicitly exponentially decaying term, and
\(J_\gamma=O(\log\gamma)\).

For marks \(Q_j\), source coefficients \(b_{ij}\), and
\(L_i=\sum_j|b_{ij}|\), compare the centered families
\(X_i=\gamma\sum_j b_{ij}h(Q_j)\) and
\(Y_i=c\gamma\sum_jb_{ij}d(Q_j)\).
For every integer separation \(n\ge2\),
\begin{equation}
 \boxed{\quad
 \|C_X(n)-C_Y(n)\|_{\rm op}
 \le\|L\|_2^2\bigl[2|c|J_\gamma E_\gamma(h)+E_\gamma(h)^2\bigr].\quad}
 \label{f16v76:jet-cov}
\end{equation}
The marks may overlap; their marginal one-plaquette estimates are not
multiplied.  The right side tends to zero for a fixed finite bank and
fixed coefficients.  More generally its displayed coefficient cost must
be paid.  The estimates preserve declared coefficient-map relations.
\end{theorem}
\begin{proof}
Split \(r=r_G+r_B\) at \(d=1\).  The ordinary norm of the centered
\(r_G\) is at most \(C_r\sqrt{\mathbb E d^4}\); ordinary-to-reflected
contraction bounds its \(\|\cdot\|_+\) norm by the same quantity.
Proposition~\ref{f16v76:singular-tail} pays \(r_B\).
The triangle inequality in the reflected seminorm proves
\eqref{f16v76:jet-plus}.  Constants disappear under native centering.
For fixed data, \(b_\gamma=O(\log\gamma)\), and
\(p_\gamma=e^{-\gamma+O(\log\gamma)}\) eventually.  The stated rates
follow without setting the singular tail to zero.

For each component of the marked bank, the norm of its difference after
one native time step is at most \(L_iE_\gamma(h)\); that of its target
is at most \(L_i|c|J_\gamma\).  Form their column synthesis operators
\(V_X,V_Y\) in \(L^2(\pi)\).  Their operator norms are bounded by the
corresponding \(\ell^2\) sums of column bounds.  By
\eqref{f16v76:separated-transfer},
\(C_X(n)=V_X^*P^{n-2}V_X\).  Expand its difference from the expression
with \(V_Y\), retaining the two linear error terms and the quadratic
error.  Since \(P^{n-2}\) is a contraction, this proves
\eqref{f16v76:jet-cov}.  No separate conditional independence for the
whole marked bank is needed.  Matching positive-time averages have the
same bounds by convexity; congruence retains the prescribed relations.
\end{proof}

For example, \(|\operatorname{Im}Q|^2=2d-d^2\) and the elementary Wilson
action \(2d\) have the same fixed-bank thermal reflected class by this
theorem.  This assertion concerns input plaquettes under the original
measure, not plaquettes after a positive physical simultaneous flow.

For the V75 singular source the remainder is
\begin{equation}
 g=2d+r_g,\qquad r_g=\frac{2d^2}{2-d},\qquad
 0\le r_g\le g,\qquad r_g\le2d^2\quad(d\le1).
 \label{f16v76:g-remainder}
\end{equation}
Choose \(p=4/3\), for which
\begin{equation}
 G_p:=\|g\|_{L^p(H)}
 =\left[4^p\frac8\pi
       \mathrm B(p+\tfrac32,\tfrac32-p)\right]^{1/p}<\infty.
 \label{f16v76:Lp}
\end{equation}
Here \(\mathrm B\) is Euler's beta integral, not the Wilson constant
\(B_\gamma\).  To check the formula, \(x=d/2\) has Haar density
\((8/\pi)\sqrt{x(1-x)}\,dx\), and \(g=4x/(1-x)\).
Put
\begin{equation}
 e_\gamma=\gamma\left[
 2\sqrt{M_{4,\gamma}}+K_\gamma^{3/4}G_{4/3}p_\gamma^{1/8}\right].
 \label{f16v76:g-error}
\end{equation}

\begin{theorem}[The renormalized antipodal bank matches the thermal Wilson bank]
\label{f16v76:subflow-match}
Use V75's exact \(A_\epsilon\), \(0<\epsilon\le1\), and native
centering.  For one mark, and for every larger reflected separation,
\begin{equation}
 \boxed{\quad
 \|\gamma\epsilon^{-2}A_\epsilon-2\gamma d\|_+
 \le\Delta_{\gamma,\epsilon}:=
         e_\gamma+24\gamma K_\gamma\epsilon.\quad}
 \label{f16v76:subflow-plus}
\end{equation}
For a marked bank with row absolute masses \(L_i\),
\begin{equation}
 \left\|\gamma^2\epsilon^{-4}C_A(n)
                   -4\gamma^2 C_d(n)\right\|_{\rm op}
 \le\|L\|_2^2
     [4J_\gamma\Delta_{\gamma,\epsilon}
                         +\Delta_{\gamma,\epsilon}^2],\qquad n\ge2.
 \label{f16v76:subflow-cov}
\end{equation}
Thus for a fixed bank the right side tends to zero if \(\gamma\to\infty\)
and \(\gamma K_\gamma\epsilon\log\gamma\to0\).  This includes
\(\epsilon_a=e^{-4\tau/a^2}\), fixed \(\tau>0\), on the separately
specified trajectory \(\log(1/a)=c\gamma_a+O(\log\gamma_a)\), \(c>0\).
No relation between \(n\) and \(a\) is required by the error bound.
\end{theorem}
\begin{proof}
Theorem~\ref{f16v76:jet} and \eqref{f16v76:g-remainder} give
\(\|\gamma g-2\gamma d\|_+\le e_\gamma\).
The nonnegative difference \(g-\epsilon^{-2}A_\epsilon\) has Haar
first moment \(\delta_\epsilon\le48\epsilon\), by V75.  Its native
conditional mean, after one selected-link integration, belongs to
\([0,K_\gamma\delta_\epsilon]\).  Integrate the two separated selected
links.  The variance of each conditional mean is at most
\(K_\gamma^2\delta_\epsilon^2/4\), and Cauchy--Schwarz gives the same
upper bound for its matching covariance.  Its \(\|\cdot\|_+\) norm is
therefore at most \(24K_\gamma\epsilon\).  Add the two errors and
repeat the synthesis-operator argument of
Theorem~\ref{f16v76:jet}, with target coefficient two.
\end{proof}

For fixed coefficients the covariance error from the singular matching
alone is \(O((\log\gamma)^3/\gamma)\).  This is an \emph{absolute}
comparison: a nonzero limit of the elementary thermal bank is not assumed.
If either bank has a finite physical-separation Gram limit, the other has
the same limit.  If the elementary bank has zero limit, the singular
normalization does not create another one.

The ordinary contact statement can be opposite.  For one mark V75's lower
bound gives
\begin{equation}
 \operatorname{Var}_\pi(\gamma\epsilon^{-2}A_\epsilon)
       \ge\frac{\gamma^2}{16}e^{-12\gamma}\epsilon^{-1}.
 \label{f16v76:contact-lower}
\end{equation}
It diverges on the physical subflow trajectory just described.  Thus the
new identification does not erase V75's antipodal contact weight: it
shows that this weight does not distinguish the thermally normalized
\emph{separated} source from the ordinary plaquette-energy bank.  The two
orders of limits are compared only through the explicit error above.

\subsection{An exact character reconstruction without an ordinary-norm limit}

The fixed-bank first-order error is only inverse-polynomial in
\(\gamma\).  It cannot be multiplied by an arbitrarily large physical
spatial sampling factor.  A different, representation-resolved
approximation gives an exponentially small separated-source error.

\begin{proposition}[The singular source has constant alternating character coefficients]
\label{f16v76:characters}
Let \(\chi_k\) be the \(\SU(2)\) character of dimension \(k+1\).
The Haar coefficients of \(g\) are
\begin{equation}
 \int g\chi_0\,dH=12,\qquad
 \int g\chi_k\,dH=16(-1)^k\quad(k\ge1).
 \label{f16v76:coefficients}
\end{equation}
The following Abel representation converges in \(L^1(H)\):
\begin{equation}
 g_\rho=12+16\sum_{k=1}^\infty(-\rho)^k\chi_k
       =\frac{16}{1+2\rho w+\rho^2}-4
       \longrightarrow g\qquad(\rho\uparrow1).
 \label{f16v76:Abel}
\end{equation}
For the finite polynomial
\begin{equation}
 G_K=12+16\sum_{k=1}^K(-1)^k\chi_k,\qquad K\ge1,
 \label{f16v76:GK}
\end{equation}
\(\|G_K\|_\infty\le12+8K(K+3)\) and
\(\operatorname{Var}_H(G_K)=256K\).  In the actual native one-plaquette
law,
\begin{equation}
 256e^{-12\gamma}K\le\operatorname{Var}_\pi(G_K)
                         \le256K_\gamma K.
 \label{f16v76:polynomial-contact}
\end{equation}
In particular no ordinary \(L^2(\pi)\) convergence is claimed at fixed
regulator.
\end{proposition}
\begin{proof}
The recursion \(\chi_{k+1}=2w\chi_k-\chi_{k-1}\), with
\(\chi_0=1,\chi_1=2w\), gives
\(\sum_{k\ge0}z^k\chi_k=(1-2wz+z^2)^{-1}\) for \(|z|<1\).
Take \(z=-\rho\).  When \(\rho\ge1/2\), its denominator is
\((1-\rho)^2+2\rho(1+w)\ge1+w\), while \(g_\rho\ge0\).
Thus \(g_\rho\le16/(1+w)\), an integrable Haar majorant.
The limit is \(8/(1+w)-4=g\).  Orthogonality and dominated convergence
prove \eqref{f16v76:coefficients}.  The supremum bound uses
\(|\chi_k|\le k+1\); orthogonality gives the Haar variance.
For the native lower bound minimize the second moment over constants and
use its density lower bound \(e^{-12\gamma}\).  For the upper bound
subtract the Haar mean twelve and use the density upper bound
\(K_\gamma\).  No assertion of native character orthogonality is used.
\end{proof}

\begin{theorem}[Finite-polynomial reconstruction in the native reflected norm]
\label{f16v76:reconstruction}
For \(K+2>3\gamma\), put
\begin{equation}
 T_{\gamma,K}=
 \frac{16(3\gamma)^{K+1}}{(K+1)!}
       \frac1{1-3\gamma/(K+2)}.
 \label{f16v76:T}
\end{equation}
For every admissible selected-link complement \(\Xi\),
\begin{equation}
 \boxed{\quad
 |\mathbb E[g(Q)-G_K(Q)\mid\Xi]|\le T_{\gamma,K},
 \qquad \|g-G_K\|_+\le T_{\gamma,K}.\quad}
 \label{f16v76:reconstruction-bound}
\end{equation}
For a bank \(\sum_jb_{ij}g(Q_j)\) and its common substitution
\(\sum_jb_{ij}G_K(Q_j)\), with row masses \(L_i\),
\begin{equation}
 \boxed{\quad
 \|C_g(n)-C_{G_K}(n)\|_{\rm op}
 \le\|L\|_2^2[12K_\gamma T_{\gamma,K}+T_{\gamma,K}^2],
 \qquad n\ge2.\quad}
 \label{f16v76:reconstruction-cov}
\end{equation}
Their reflected error Gram has norm at most
\(\|L\|_2^2T_{\gamma,K}^2\).  These bounds are uniform in surrounding
volume, integer separation and the positions or overlaps of the marks.
The coefficients and the cutoff may depend on the regulator.
\end{theorem}
\begin{proof}
The selected conditional holonomy density is a translate of the central
one-link law.  If its field length is \(t\le6\gamma\) and its center
is \(M\), its conditional character mean is
\(\lambda_k(t)\chi_k(M)\).  At \(t=0\) every nontrivial term vanishes.
Using V63's factorial bound,
\[
 16|\lambda_k(t)\chi_k(M)|
 \le16(k+1)\frac{(t/2)^k}{(k+1)!}
 \le16\frac{(3\gamma)^k}{k!}.
\]
Apply \eqref{f16v76:Abel}, pass to its \(L^1\) limit under the bounded
conditional density, and use this summable majorant.  The conditional
expectation of the omitted terms is bounded by
\(16\sum_{k>K}(3\gamma)^k/k!\).  Its successive ratios starting at
\(k=K+1\) are at most \(3\gamma/(K+2)<1\), proving
\eqref{f16v76:T}.  The translated center and the actual field lengths
have not been replaced by coherent values.

For the same mark at times 0 and 2, condition on both selected links'
complement.  The centered covariance of \(g-G_K\) is the covariance of
two conditional means in \([-T_{\gamma,K},T_{\gamma,K}]\).  Its absolute
value is at most \(T_{\gamma,K}^2\).  Native positivity and truncation
give the reflected bound.  Moreover \(\|g\|_+\le6K_\gamma\): its
conditional mean lies in \([0,12K_\gamma]\), as in V75.
For a marked sum use the triangle inequality after applying one native
transfer step.  The synthesis norms of the target and error are at most
\(6K_\gamma\|L\|_2\) and \(T_{\gamma,K}\|L\|_2\), respectively.
The two linear errors and one quadratic error prove
\eqref{f16v76:reconstruction-cov}.  Singular products are integrable by
\eqref{f16v76:two-density}, so each limiting step is legitimate.
\end{proof}

The partial sums need not approximate \(g\) in the ordinary norm, nor is
that norm used to estimate the error.  They approximate its conditional
means uniformly, which is the quantity entering the separated native
pairing.  Coefficients alternate and must remain signed; bounding the
uncombined character terms of the original observable separately in an
ordinary norm would miss this reconstruction.

\begin{proposition}[An explicit cutoff pays polynomially many physical marks]
\label{f16v76:physical-cutoff}
If \(K+1\ge6e\gamma\), then
\begin{equation}
 T_{\gamma,K}\le32\,2^{-(K+1)}.
 \label{f16v76:T-simple}
\end{equation}
For \(0<a<1\), define
\[
 A_a=\max\{1,\|L_a\|_2^2\},\qquad B_a^{\sharp}=12K_{\gamma_a}+1.
\]
Given \(N>0\), choose an integer \(K_a\ge1\) such that
\begin{equation}
 K_a+1\ge
 \max\left\{6e\gamma_a,
 \frac{N\log(1/a)+\log(32A_aB_a^{\sharp})}{\log2}\right\}.
 \label{f16v76:cutoff-choice}
\end{equation}
Then the covariance error in \eqref{f16v76:reconstruction-cov} is at most
\(a^N\), for all \(n\ge2\).  If \(\gamma_a=O(\log(1/a))\) and
\(\|L_a\|_2\) is polynomially bounded in \(a^{-1},1+\gamma_a\), this
requires only \(K_a=O(\log(1/a))\).
\end{proposition}
\begin{proof}
The inequality \(m!\ge(m/e)^m\), with \(m=K+1\), bounds the first
term in \eqref{f16v76:T} by \(16\,2^{-m}\).  Its denominator is at
least \(1-1/(2e)>1/2\).  This proves \eqref{f16v76:T-simple}.
The chosen cutoff makes \(T\le a^N/(A_aB_a^{\sharp})\le1\).
Hence \(\|L_a\|_2^2(12K_\gamma T+T^2)\le a^N\).
The degree assertion follows with all coefficient masses retained.
\end{proof}

This physical-size statement reconstructs sums of the singular
\emph{microscopic plaquette} functions.  It does not identify their limit
with the first thermal jet for such growing sums.  That latter estimate
still costs \(\|L_a\|_2^2O((\log\gamma_a)^3/\gamma_a)\), which can be
large on an exponential physical trajectory.  Increasing the character
cutoff pays the full singular reconstruction but does not prove a
nonzero reflected Gram for the reconstructed polynomial bank.

\subsection{The remaining simultaneous-flow problem and verification}

The exponentially renormalized local-subflow route no longer represents
an unidentified reflected source.  For fixed thermally normalized banks
it matches twice the elementary Wilson action density under the stated
native weak-coupling limits.  For growing sums its singular source has a
finite, quantitatively controlled character reconstruction.  Neither
statement creates temporal dependence that was absent from the native
transfer, and neither uses an independent Haar vacuum.

The current source is instead \(O(\Phi_{a,\tau}U)\), with simultaneous
spatial flow and fixed positive physical \(\tau\).  To apply the present
identity-expansion argument to \emph{flowed} plaquettes one would need
quantitative moments and source-weighted singular tails under their actual
pushed law.  The elementary energy decrease and V74's entropy identity do
not imply those higher moments.  Nor does the original private-link
conditional convolution remain valid after silently replacing \(Q\)
by \((\Phi_{a,\tau}U)_Q\).  The character theorem is therefore not a
controlled expansion of that collective source.

A next useful result must estimate the simultaneous-flow native reflected
Gram itself, or its genuine collective source expansion, rather than
claim survival from V75's fixed-regulator singular positivity.  For the
microscopic branch, the still-uncomputed object can now be taken to be the
ordinary thermal plaquette Gram or the explicit finite character bank,
with the displayed normalization and sampling costs.  A physical lower
Gram, common dense-core physical-time contraction and the interacting
continuum mass gap remain distinct obligations.

The companion diagnostic is
\begin{center}\small
\path{verify_ym_f16_v76_singular_matching.py}.
\end{center}
It separates exact character and moment algebra, literal reflected-block
incidence, analytic bounds tested on compact conditional integrals,
and explicitly labelled positive-convolution or latent-mixture reference
models.  No numerical value is an input to a proof.  These are not native
four-dimensional vacuum samples, interval certificates, or proof-assistant
verification.  The append-only check removes this delimited appendix and
reverses the single displayed version change to recover V75 byte for byte.
The preceding archive is preserved, not independently recertified.

\clearpage
\begin{record}[V76 result ledger]
\label{f16v76:ledger}
\emph{Proved here:} a native marked-plaquette tail and fixed-moment bound
from the inherited reflection input; a separated singular-tail estimate;
all-direction thermal matching from the identity expansion, including an
\(L^p\) non-\(L^2\) remainder; a uniform reflected comparison of V75's
exponentially renormalized subflow bank with the thermal Wilson bank;
exact alternating character coefficients of the singular source; native
finite-character reconstruction with paid factorial tails and divergent
contact norm; and a representation cutoff paying polynomially growing
physical sampling masses.

\emph{Inherited or established background:} V65's closed-block reflection
and cap normalization, V72's native conditional link geometry, V63's
factorial character bound, V75's subflow formula and singular \(L^1\)
error, native transfer positivity and the general use of Abel summation.
No unmarked entropy estimate is differentiated into a source covariance.

\emph{Not asserted:} ordinary \(L^2\) convergence of either reconstruction;
a thermal first-jet theorem for an unpaid growing sampling mass; a Wilson
conditional law for flowed links; a nonzero physical plaquette Gram from
finite-regulator positivity; a continuum construction of the singular
bank; or replacement of simultaneous spatial flow by separate subflows.

\emph{Still unproved:} a nonzero physical-scale native reflected Gram for
the simultaneous spatial-flow source, its collective spatial and source
normalization control, common dense-core physical-time contraction, and
the interacting four-dimensional continuum Yang--Mills mass gap.
\end{record}

\begin{firewall}
Antipodal contact divergence is retained while its separated contribution
is estimated.  The reflected source class is identified under the original
native law, not inferred from one-slice concentration.  Every growing
coefficient mass and representation cutoff is explicit.  The physical
transfer and all V1--V75 mathematical text remain unchanged.
\end{firewall}
% END F16 V76 APPEND

% BEGIN F16 V77 APPEND -- 2026-09-19
% Append-only continuation of the Post-Fifteenth-Archive V76.
\clearpage
\section[V77: derivative-free reconstruction of the simultaneous-flow source]{V77: derivative-free reconstruction\\of the simultaneous-flow source}
\label{f16v77:context}

V76 identifies the singular isolated-subflow bank but does not expand the
simultaneously flowed source.  We return to that original collective
observable.  The new argument does not differentiate its spatial flow,
postulate a conditional Wilson law for flowed links, or suppress shared
path occurrences separately.  Instead it moves a finite positive
polynomial filter from the observable to the \emph{original} conditional
slice density.  A local integration-by-parts cancellation bounds the
resulting density error.  Tensorization costs the number of input links
linearly, not exponentially.  Conditional independence between two
separated slices then converts that error into the native reflected norm.

This gives an explicit finite gauge-invariant polynomial reconstruction of
\(O(\Phi_{a,\tau}U)\), including fixed positive physical \(\tau\), with a
cutoff polynomial in the regulator for polynomially bounded source sizes.
The theorem is an approximation of its actual native pairings, not a
calculation of their limiting value.  The number of polynomial coefficients
can still be exponential in the spatial volume, and a nonzero reflected
lower Gram remains unproved.

\subsection{Original-law scope and a positive finite polynomial filter}

Work on the nondegenerate finite spatial tori of the preceding sections,
with the original isotropic four-dimensional \(\SU(2)\) Wilson coefficient
\(\gamma>0\).  A selected finite set \(S\) of spatial links has size \(M\).
The source is a bounded function of precisely these original input
variables.  For the simultaneous-flow application, \(S\) is the entire
spatial slice.  Conditional on every other link in a finite Euclidean
cylinder, one selected link, with all the other selected links also fixed,
has density
\begin{equation}
 q_h(X)=\frac{e^{h\cdot x}}{\mathfrak z(|h|)},\qquad |h|\le6\gamma,
 \qquad \mathfrak z(t)=2I_1(t)/t.
 \label{f16v77:conditional}
\end{equation}
The joint law on \(S\) is \emph{not} asserted to be a product: it includes
its internal spatial Wilson plaquettes.  We pass the uniform conditional
estimates through the fixed-spatial-regulator stationary-cylinder limit.
A nonlocal Perron endpoint weight is not absorbed into
\eqref{f16v77:conditional}.  All later covariances use the same native
stationary transfer \(P_a\), with \(0\preceq P_a\preceq I\).

Positive polynomial means on spheres are established approximation tools.
The kernel below is the normalized \(S^3\) de la Vall\'ee Poussin kernel;
see R.~Yang, F.~Cao and J.~Xiong, \emph{The strong converse inequality for
de la Vall\'ee Poussin means on the sphere},
\href{https://arxiv.org/abs/1105.4062}{arXiv:1105.4062}, equations (4)--(5).
The use of gauge-invariant representation coefficients is also established;
see J.~M.~Aroca, H.~Fort and R.~Gambini,
\href{https://doi.org/10.1103/PhysRevD.58.045007}{Phys.\ Rev.\ D
\textbf{58} (1998), 045007}.
Neither the kernel nor the existence of a spin-network basis is claimed
as new.  The input-law density estimate, its all-slice tensorization,
and the resulting derivative-free native source bound are proved here.

For an integer \(K\ge1\), put
\begin{equation}
 \begin{aligned}
 J_K(X)&=\frac{(1+w(X))^K}{z_K},\qquad
 z_K=\int(1+w(X))^K\,dH(X),\\
 (A_K f)(U)&=\int J_K(VU^{-1})f(V)\,dH(V).
 \end{aligned}
 \label{f16v77:kernel}
\end{equation}
Haar measure is normalized.  This filter is an operation on functions of
the original links, not another Wilson flow or a physical-time evolution.

\begin{proposition}[Exact kernel coefficients and concentration]
\label{f16v77:kernel-prop}
The operator \(A_K\) preserves constants, positivity and real functions,
is self-adjoint on Haar \(L^2\), and contracts every Haar \(L^p\), including
\(p=1,\infty\).  Its Peter--Weyl multiplier on representation degree \(k\)
is
\begin{equation}
 \alpha_{K,k}=\begin{cases}
 \displaystyle\frac{K!(K+2)!}{(K-k)!(K+k+2)!},&0\le k\le K,\\[1mm]
 0,&k>K.
 \end{cases}
 \label{f16v77:multiplier}
\end{equation}
In particular it is a positive semidefinite finite-rank operator.  If
\(\theta=d(I,X)\in[0,\pi]\), then
\begin{equation}
 \int(1-w)J_K\,dH=\frac3{K+3},\qquad
 \int\theta^2J_K\,dH\le\frac{3\pi^2}{2(K+3)}.
 \label{f16v77:kernel-moment}
\end{equation}
\end{proposition}
\begin{proof}
The kernel is nonnegative, central, inversion symmetric, and has integral
one.  Haar invariance gives all contraction and symmetry statements.
For the coefficients, expand
\[
 (1+\cos\theta)^K=2^{-K}\left\{\binom{2K}{K}
             +2\sum_{j=1}^K\binom{2K}{K-j}\cos(j\theta)\right\}.
\]
Use the radial Haar law \((2/\pi)\sin^2\theta\,d\theta\) and
\(\chi_k(\theta)=\sin((k+1)\theta)/\sin\theta\).  Its unnormalized
character coefficient is
\[
 2^{-K}\left[\binom{2K}{K-k}-\binom{2K}{K-k-2}\right],
\]
with out-of-range binomial coefficients zero.  Divide by the coefficient
at \(k=0\) and by \(k+1\); factorial cancellation gives
\eqref{f16v77:multiplier}.  All retained multipliers are positive.
Alternatively, \((1+w)/2\) under \(J_KdH\) has beta parameters
\(K+3/2,3/2\).  This gives the first moment in
\eqref{f16v77:kernel-moment}.  The inequality
\(\theta^2\le(\pi^2/2)(1-\cos\theta)\), from
\(\sin(\theta/2)\ge\theta/\pi\), gives the second.
\end{proof}

\subsection{A density estimate that never differentiates the source}

\begin{proposition}[A one-link \(L^1\) cancellation under the native field]
\label{f16v77:one-link}
For every \(h\in\mathbb R^4\),
\begin{equation}
 \boxed{\quad
 \|A_Kq_h-q_h\|_{L^1(H)}
 \le\min\left\{2,\frac{3\pi^2|h|}{2(K+3)}\right\}.
 \quad}
 \label{f16v77:one-link-bound}
\end{equation}
The estimate has no exponential factor in \(|h|\).
\end{proposition}
\begin{proof}
Fix a unit imaginary quaternion \(n\) and let \(D_n\) generate left
multiplication by \(e^{tn}\).  For \(l(X)=h\cdot x\),
\[
 D_n^2l=-l,\qquad
 D_n^2q_h=\{(D_nl)^2-l\}q_h.
\]
Integration of the second derivative over Haar measure is zero.  Hence
\(\mathbb E_{q_h}(D_nl)^2=\mathbb E_{q_h}l\ge0\), and
\begin{equation}
 \|D_n^2q_h\|_1
 \le\mathbb E_{q_h}(D_nl)^2+\mathbb E_{q_h}|l|
 \le2|h|.
 \label{f16v77:directional-L1}
\end{equation}
This integrates the squared first derivative before bounding it; the
pointwise estimate of that square would instead introduce \(|h|^2\).

Let \(T_t\) be left translation on functions.  Paired Taylor integration
and its \(L^1\) isometry give
\[
 \left\|\tfrac12(T_\theta+T_{-\theta})q_h-q_h\right\|_1
 \le\frac{\theta^2}{2}\|D_n^2q_h\|_1
 \le |h|\theta^2.
\]
The central kernel averages \(n\) symmetrically at each \(\theta\).
Integrate this inequality against its radial law and use
\eqref{f16v77:kernel-moment}.  Two probability densities differ in
\(L^1\) by at most two.  At \(h=0\) both sides before that cap are zero.
\end{proof}

Let \(A_{K_e}^{(e)}\) act only on input link \(e\), and set
\begin{equation}
 \mathcal A_{\boldsymbol K,S}=\prod_{e\in S}A_{K_e}^{(e)},\qquad
 \varepsilon_{\boldsymbol K,S}
 =\min\left\{2,9\pi^2\gamma\sum_{e\in S}\frac1{K_e+3}\right\}.
 \label{f16v77:product-filter}
\end{equation}
The factors commute because they act on different variables.

\begin{theorem}[Correlated-slice tensorization with only a linear link cost]
\label{f16v77:slice-density}
Let \(p_\Xi\) be the native conditional density of \(U_S\) given its
full complement \(\Xi\).  Then
\begin{equation}
 \boxed{\quad
 \|p_\Xi-\mathcal A_{\boldsymbol K,S}p_\Xi\|_1
 \le\varepsilon_{\boldsymbol K,S}.
 \quad}
 \label{f16v77:density-bound}
\end{equation}
For any bounded real source \(F\) of these variables, with
\(b(F)=\tfrac12(\sup F-\inf F)\),
\begin{equation}
 \left|\mathbb E[F-\mathcal A_{\boldsymbol K,S}F\mid\Xi]\right|
 \le b(F)\varepsilon_{\boldsymbol K,S}.
 \label{f16v77:mean-bound}
\end{equation}
Both bounds are uniform in the exterior and surrounding volume.  They do
not use derivatives, polynomial degree, path length, or flow sensitivity
of \(F\).
\end{theorem}
\begin{proof}
For a fixed \(e\), disintegrate the \emph{original} density as
\(p_\Xi=p_{\Xi,-e}q_{h_e}\).  The marginal factor is independent of
\(U_e\).  Integrating \eqref{f16v77:one-link-bound} against that marginal
therefore gives
\[
 \|(I-A_{K_e}^{(e)})p_\Xi\|_1\le\frac{9\pi^2\gamma}{K_e+3}.
\]
Enumerate \(S\) and use the exact telescoping identity
\[
 I-A_1\cdots A_M=\sum_{j=1}^M A_1\cdots A_{j-1}(I-A_j).
\]
Every preceding factor is an \(L^1\) contraction.  Crucially, each
\((I-A_j)\) in this identity is applied to the original \(p_\Xi\),
not to an already filtered law assigned a new Wilson conditional.
Taking norms and summing proves the uncapped estimate; normalization
of both densities supplies the cap at two.  Self-adjointness moves the
filter from \(F\) to the density.  Subtract the midpoint of the range of
\(F\) before estimating; the density difference has integral zero.
This proves \eqref{f16v77:mean-bound}.
\end{proof}

\subsection{The complete reflected source approximation}

Let \(F=(F_1,\ldots,F_d)\) be a finite real bounded slice family, and let
\(F^{\boldsymbol K}=\mathcal A_{\boldsymbol K,S}F\).  Native centering is
performed \emph{after} either function is chosen.  Put
\begin{equation}
 B_F^2=\sum_i b(F_i)^2,\qquad
 C_F(n)=\operatorname{Cov}_{\pi}(F(U_0),F(U_n)),\qquad
 E=F-F^{\boldsymbol K}.
 \label{f16v77:source-data}
\end{equation}
No native means or covariances are computed under product Haar measure.

\begin{theorem}[Derivative-free native reflected reconstruction]
\label{f16v77:reflected}
For integer separation \(n\ge2\),
\begin{equation}
 0\preceq C_E(n)\preceq
          \varepsilon_{\boldsymbol K,S}^2B_F^2I_d.
          \label{f16v77:error-Gram}
\end{equation}
\begin{equation}
 \boxed{\quad
 \|C_F(n)-C_{F^{\boldsymbol K}}(n)\|_{\rm op}
 \le2\varepsilon_{\boldsymbol K,S}B_F^2.
 \quad}
          \label{f16v77:cov-error}
\end{equation}
The same estimates hold for matching reflected placements at least one
lattice step beyond the reflection plane, additional nonnegative delays,
and common nonnegative temporal weights of total one.  The estimates
preserve declared linear identities by applying the same filter to every
column.  They do not assert that the cutoff introduces no additional
null directions.
\end{theorem}
\begin{proof}
Condition on the complement of \(S\) at times zero and two.  No Wilson
plaquette meets both variable sets: a temporal plaquette spans only one
time step.  The two \emph{whole slices} are conditionally independent,
although links within each slice are not.  For any coefficient vector
\(c\), the conditional error means have magnitude at most
\(\varepsilon\sum_i|c_i|b(F_i)\le\varepsilon B_F|c|\).
Their covariance is the full error covariance, by conditional independence.
Cauchy--Schwarz bounds it by \(\varepsilon^2B_F^2|c|^2\).
The actual positive transfer makes this matching covariance nonnegative.
Equivalently, with centering denoted by a circle,
\[
 \|P_a(c^*E)^\circ\|_2\le\varepsilon B_F|c|.
\]
This one-step vector estimate controls every \(n\ge2\), since
\(C_E(n)=V_E^*P_a^{n-2}V_E\), where
\(V_Ec=P_a(c^*E)^\circ\).

The filter is a positive average and cannot increase the oscillation of
any real function.  Popoviciu's elementary variance bound gives
\(\operatorname{Var}(c^*F)\le(\sum_i|c_i|b(F_i))^2\), and the same is
true for \(c^*F^{\boldsymbol K}\).  Hence the source syntheses after one
native step have norms at most \(B_F\); their difference has norm at most
\(\varepsilon B_F\).  Expand
\[
 V_F^*P_a^{n-2}V_F-V_K^*P_a^{n-2}V_K
 =(V_F-V_K)^*P_a^{n-2}V_F
       +V_K^*P_a^{n-2}(V_F-V_K).
\]
This proves \eqref{f16v77:cov-error}.  More generally a reflected average
has synthesis \(\sum_j h_jP_a^{t_j}F^\circ\), with every \(t_j\ge1\).
Its norm and error obey the same bounds by positivity of the weights and
contraction.  Insert any further power of \(P_a\) between the two syntheses.
Linearity proves the statement about source identities.  Cylinder limits
retain the local conditional proof and all constants.
\end{proof}

\subsection{Application to the actual simultaneous spatial-flow bank}

Use the source of \eqref{f16v69:source}, without modifying its simultaneous
spatial evolution.  Before centering write
\begin{equation}
 O_i^\tau(U)=\frac{\zeta_a^2}{2a}\sum_x f_i(x)
       |\operatorname{Im}(\Phi_{a,\tau}U)_{12}(x)|^2,
 \qquad \|f_i\|_{1,a}=a^3\sum_x|f_i(x)|.
 \label{f16v77:full-flow-source}
\end{equation}
Let \(S\) be all \(M_a=3\mathcal V_a/a^3\) spatial links of the torus,
and apply the common cutoff \(K\) to those \emph{input} variables.

\begin{theorem}[A populated collective polynomial cutoff at physical flow time]
\label{f16v77:physical}
For every finite \(\tau\ge0\), including regulator-dependent values,
define
\(O_i^{\tau,K}=\mathcal A_{K,S}O_i^\tau\).
It is a finite gauge-invariant polynomial in the original slice links.
At every native integer separation \(n\ge2\),
\begin{equation}
 \boxed{\quad
 \|C_{O^\tau}(n)-C_{O^{\tau,K}}(n)\|_{\rm op}
 \le\frac{9\pi^2}{8}\,
 \frac{M_a\gamma\zeta_a^4}{a^8(K+3)}
               \sum_i\|f_i\|_{1,a}^2.
 \quad}
 \label{f16v77:flow-error}
\end{equation}
In particular, for any prescribed \(N>0\), it suffices to choose
\begin{equation}
 K+3\ge\frac{27\pi^2}{8}\,
 \mathcal V_a\gamma\zeta_a^4
       \left(\sum_i\|f_i\|_{1,a}^2\right)a^{-(N+11)},\qquad K\ge1,
 \label{f16v77:physical-cutoff}
\end{equation}
to make the error at most \(a^N\), uniformly in \(n\) and \(\tau\).
The corresponding reflected error-Gram bound is
\eqref{f16v77:error-Gram} with
\[
 B_F^2=\frac{\zeta_a^4}{16a^8}\sum_i\|f_i\|_{1,a}^2
\]
as a valid upper budget.
\end{theorem}
\begin{proof}
Every compact flowed plaquette has squared imaginary part in \([0,1]\).
Even for signed \(f_i\), the range length of \(O_i^\tau\) is at most
\(\zeta_a^2\|f_i\|_{1,a}/(2a^4)\).  Thus its half-oscillation is at most
\(\zeta_a^2\|f_i\|_{1,a}/(4a^4)\).  Substitute these values and
\(\varepsilon\le9\pi^2M_a\gamma/(K+3)\) into
\eqref{f16v77:cov-error}.  This proves \eqref{f16v77:flow-error} and,
using the actual link count, \eqref{f16v77:physical-cutoff}.

Finite polynomial degree follows from
\eqref{f16v77:multiplier}.  For gauge invariance, under
\(U_e\mapsto g_xU_eg_y^{-1}\) change the integrated input link similarly.
The argument of \(J_K\) changes by conjugation at \(x\), to which the
central kernel is invariant.  The original \(O_i^\tau\) is gauge invariant
by the inherited flow equivariance.  Thus its filtered function is gauge
invariant as well.  Only original links on the same slice have been used;
no time-reflection support is crossed.
\end{proof}

For fixed physical \(\mathcal V_a\) and sampled \(L^1\) norms, a
normalization \(\zeta_a=O(a^{-d}(1+\gamma_a)^b)\) and
\(\gamma_a=O(\log(1/a))\) require at most
\[
 K_a=O\!\left(a^{-(N+11+4d)}
                   (\log(1/a))^{1+4b}\right).
\]
There is no \(e^{C\tau/a^2}\) factor.  This is not a sharp cutoff and not a
claim of efficient computation.  Unlike V76's logarithmic cutoff for a
particular microscopic singular function, it controls an arbitrary bounded
collective output.  The flow itself is unchanged.  Haar integration in the
filter defines an approximating \emph{observable}; it does not replace the
native measure by noisy links or a new coupling.  V74's concentrated
\emph{flowed-link} conditional laws are irrelevant to the proof, which
uses only original input-link conditionals.

\subsection{Explicit finite coefficients and preservation of shared-path singlets}

For an orthonormal product Peter--Weyl basis write
\[
 \varphi_\alpha(U)=\prod_{e\in S}\sqrt{k_e+1}\,
        D^{(k_e)}_{m_en_e}(U_e),\qquad
 \widehat F_\alpha=\int F(U)\overline{\varphi_\alpha(U)}\,dH^S(U).
\]
The coefficients integrate the \emph{actual} function
\(O(\Phi_{a,\tau}U)\), not a linearized or separately flowed substitute.

\begin{proposition}[Finite source expansion and its coefficient cost]
\label{f16v77:coefficients}
The filtered function is exactly
\begin{equation}
 \boxed{\quad
 F^{\boldsymbol K}(U)=
 \sum_{0\le k_e\le K_e}
 \left(\prod_e\alpha_{K_e,k_e}\right)
              \widehat F_\alpha\varphi_\alpha(U),
 \quad}
 \label{f16v77:finite-expansion}
\end{equation}
where the sum includes all matrix indices.  The common-cutoff ambient
coefficient count is
\begin{equation}
 \left[\frac{(K+1)(K+2)(2K+3)}6\right]^M.
 \label{f16v77:dimension}
\end{equation}
Gauge invariance restricts these tensors to invariant vertex contractions;
it does not justify deleting their signed coefficients.  There is no
asserted polynomial bound on their number as \(M\to\infty\).
\end{proposition}
\begin{proof}
Each one-link convolution multiplies its matrix coefficients by
\(\alpha_{K_e,k_e}\).  The product operators act diagonally in the
orthonormal tensor basis and kill every index beyond its cutoff.  The
resulting series is finite, giving \eqref{f16v77:finite-expansion} for
bounded \(F\in L^2(H^S)\).  There are \((k+1)^2\) matrix coefficients
at degree \(k\); summing their dimensions gives
\eqref{f16v77:dimension}.  The gauge-equivariance argument above makes the
coefficient tensor invariant under every vertex action, which is precisely
the condition for contraction through invariant tensors.  No probability
interpretation is assigned to the Haar expansion coefficients.
\end{proof}

\begin{proposition}[The common-path singlet is preserved]
\label{f16v77:singlet}
Let \(Z=w(XA)w(XB)\) and \(R=w(AB^{-1})\) be V72's shared-path functions.
Filter only \(m\) input links appearing once in the common path \(X\)
and not in the retained returns \(A,B\).  For the common cutoff,
\begin{equation}
 \boxed{\quad
 \mathcal A_{K,S}Z=\frac14R+
          \left[\frac{K(K-1)}{(K+3)(K+4)}\right]^m
                           \left(Z-\frac14R\right).
 \quad}
 \label{f16v77:filtered-singlet}
\end{equation}
The bracket is zero for \(K=1\).  In contrast a single common-path
fundamental factor has multiplier \([K/(K+3)]^m\).
\end{proposition}
\begin{proof}
In this operator calculation the replacement of one input link is
multiplication by independent central \(J_K\) noise.  Inversion and
conjugation of that noise preserve its law.  Moving all deterministic
factors through the common path leaves its center \(X\) and the central
convolution of \(m\) copies.  The product of its two fundamental factors
is the sum of its Haar singlet \(R/4\) and a degree-two harmonic, as proved
in V72.  Equation \eqref{f16v77:multiplier} supplies multipliers one and
\(\alpha_{K,2}^m\), respectively.  This proves the formula.  A single
fundamental factor uses \(\alpha_{K,1}^m\).  No assertion of independence
under the Wilson law is needed for this identity of filtered functions.
\end{proof}

\subsection{What the reconstruction does not preserve automatically}

\begin{proposition}[Ordinary error and source rank are separate]
\label{f16v77:rank}
The derivative-free reflected bounds do not imply a uniform ordinary
\(L^2\) approximation for all bounded cutoff-dependent sources.  They
preserve every original linear identity but can create additional null
columns at a finite cutoff.
\end{proposition}
\begin{proof}
On one native elementary plaquette take
\(F_K(Q)=\cos((K+3)\theta(Q))\), so \(|F_K|\le1\).  The exact identity
\[
 \cos(l\theta)=\tfrac12(\chi_l-\chi_{l-2}),\qquad l\ge2,
\]
shows that filtering even one boundary link with cutoff \(K\) gives zero:
both representation degrees exceed \(K\).  On the other hand, at any
fixed finite \(\gamma\), the conditional Fourier coefficient bound
\(\lambda_j(6\gamma)\le(3\gamma)^j/(j+1)!\) implies
\(\mathbb E F_K\to0\) and \(\mathbb E F_K^2\to1/2\).
Thus \(\operatorname{Var}_\pi(F_K)\to1/2\), whereas the filtered variance
is zero.  This also gives a genuine new null column.  Its separated
covariance is small, consistently with Theorem~\ref{f16v77:reflected}.
The source changes with \(K\); ordinary convergence for one fixed smooth
finite-regulator function is not denied.
\end{proof}

An actual lower-Gram certificate is now finite but remains to be populated.
If, on a declared coefficient subspace with a stated Euclidean norm,
\(H_K\succeq hI\) and the bound \(2\varepsilon B_F^2\) is less than
\(h\), then
\begin{equation}
 H\succeq(h-2\varepsilon B_F^2)I.
 \label{f16v77:lower-certificate}
\end{equation}
This is Weyl's elementary perturbation bound, not a proved value of \(h\).
For numerical coefficients there is another cost: if a computed polynomial
\(\widetilde F^K\) differs componentwise in supremum norm by at most
\(\delta_i\), put \(\Delta^2=\sum_i\delta_i^2\).  The complete covariance
error is at most
\begin{equation}
 2\varepsilon B_F^2+2B_F\Delta+\Delta^2.
 \label{f16v77:coefficient-error}
\end{equation}
Indeed the exact filtered synthesis has norm at most \(B_F\), and the
centered coefficient-error synthesis has norm at most \(\Delta\).
For expansion \eqref{f16v77:finite-expansion}, one can use the explicit
supremum-error budget
\[
 \delta_i=\sum_\alpha|\delta c_{i\alpha}|\prod_e\sqrt{k_e+1}.
\]
No numerical values for the full collective coefficients or the native
matrix \(H_K\) are supplied here.

\subsection{Noncircular consequence and result ledger}

The missing collective expansion is now available on every finite spatial
regulator with a proved physical-size reflected error.  Its coefficients
are specified by the actual simultaneous flow, and its covariance is
always evaluated under the original Wilson history.  This avoids the
unstable derivative, the concentrated pushed conditional density, and
separate suppression of shared path factors.  None of those three objects
has been assumed small.

This is not yet spatial localization uniformly in an infinite volume: the
link cost \(M\) is the size of the input support, which for the full flow
on a torus is the whole slice.  The finite count
\eqref{f16v77:dimension} can be prohibitive.  The next substantive
calculation is a signed native covariance or lower-Gram estimate for this
\emph{specified collective polynomial bank}, including its retained
singlets and coefficient errors.  Another ordinary polynomial-density
statement or another local-subflow formula would not evaluate that matrix.
No physical source is inferred from finite-rank availability alone.

The companion diagnostic
\begin{center}\small
\path{verify_ym_f16_v77_collective_reconstruction.py}
\end{center}
checks exact kernel coefficients, Haar moments and shared-path sectors;
the directional density cancellation; correlated-slice tensorization;
separated conditional-mean and matrix identities; literal simultaneous
spatial-flow gauge covariance; and the physical cutoff algebra.  Numerical
compact integration and finite-probability examples are diagnostics, not
Wilson-vacuum sampling, interval certificates or formal verification.
Removing this delimited appendix and reversing the displayed title change
recovers V76 byte for byte.

\begin{record}[V77 result ledger]
\label{f16v77:ledger}
\emph{Proved here:} an explicit native \(L^1\) one-link density estimate;
correlated-slice tensorization with a linear link count; derivative-free
all-direction reflected reconstruction of arbitrary bounded slice banks;
a finite gauge-invariant polynomial reconstruction of the actual
simultaneous-flow source at arbitrary finite flow time; an explicit
physical-size cutoff and coefficient count; exact retention of a shared-path
singlet; and error certificates which keep numerical coefficient error,
possible new null directions, and lower-Gram normalization distinct.

\emph{Inherited or established background:} positive de la Vall\'ee Poussin
means, Peter--Weyl decomposition and invariant tensors, original Wilson
one-link conditionals, the literal simultaneous spatial flow, and native
transfer/reflection positivity.  No archive theorem is recertified.

\emph{Not asserted:} an efficiently computable full coefficient bank;
ordinary-norm convergence at a derivative-free rate; exact equality of
source kernels at each cutoff; spatially local thermodynamic reconstruction
without its support cost; a filtered Wilson measure; a sharp representation
cutoff; or positivity of a native lower Gram from polynomial availability.

\emph{Still unproved:} a nonzero fixed-physical-separation native Gram for
the collective bank, its thermodynamic and normalization control,
common dense-core physical-time contraction, and the interacting
four-dimensional continuum Yang--Mills mass gap.
\end{record}

\begin{firewall}
The filter acts on the entire original source before any factors are split.
Its error is charged to the original conditional density, not to the spatial
flow derivative or a substituted flowed-link law.  A finite approximation
with a proved reflected error is not an evaluation of its native covariance.
The physical transfer and all V1--V76 mathematical text remain unchanged.
\end{firewall}
% END F16 V77 APPEND

% BEGIN F16 V78 APPEND -- 2026-09-19
% Append-only continuation of the Post-Fifteenth-Archive V77.
\clearpage
\section[V78: native lower Grams from finite kinetic witnesses]{V78: native lower Grams\\from finite kinetic witnesses}
\label{f16v78:context}

V77 reconstructs the simultaneous-flow source but leaves a native lower
Gram unevaluated.  We now prove a positive lower bound, with all constants
specified, for that actual source on each finite spatial regulator.  The
bound is deliberately distinguished from a cutoff-uniform physical
noncollapse theorem.  Its large volume, flow-time and representation costs
are displayed; they do not tend to a positive continuum constant.

The essential change is to separate a \emph{kinetic witness resolution}
from the final reconstruction cutoff.  Applying a smallest-multiplier
bound to the entire reconstructed bank makes the lower bound deteriorate
whenever that cutoff is increased.  A fixed witness instead tests the
source before the kinetic smoothing.  An explicitly paid Haar tail then
controls this overlap, and increasing the reconstruction cutoff no longer
changes the resulting positive lower constant.  Native centering, the
magnetic multiplier, the actual Perron normalization and the physical
time step are all retained.

\subsection{Scope, inherited transfer and constants}

Work on a nondegenerate finite periodic cubic spatial lattice with side
at least four.  Let $M$ be its number of spatial links and $N_p$ its number
of spatial plaquettes.  For a cubic torus both equal three times the
number of vertices.  Write $dm$ for normalized product Haar measure.
The original magnetic action and the known-scalar-normalized Wilson
transfer are
\begin{equation}
 S(U)=\sum_{P\ {\rm spatial}}(1-w(U_P)),\quad
 b(U)=e^{-\gamma S(U)/2},\quad
 T=M_b C M_b,\quad C=C_\gamma^{\otimes M},\quad
 k_\gamma(X)=\frac{e^{\gamma w(X)}}{\mathfrak z(\gamma)}.
 \label{f16v78:transfer}
\end{equation}
Here $C_\gamma$ is convolution by $k_\gamma$, with eigenvalue
$\lambda_j(\gamma)=I_{j+1}(\gamma)/I_1(\gamma)$ in representation degree
$j$.  The known scalar removed in \eqref{f16v78:transfer} does not change
the vacuum-normalized physical transfer.  On gauge-invariant functions
the gauge average is the identity; it is not dropped on a noninvariant
physical state.  The positive extension $M_bCM_b$ to Haar space has a
unique positive top eigenfunction, which is gauge invariant.

This transfer factorization, positivity, and character convention are
inherited from V51--V52 and f15 V79.  The latter already supplies
finite-rank certificates for the \emph{transfer spectrum}; we do not
repeat those certificates as new results.  For background on the positive
Wilson transfer see M.~Creutz,
\href{https://doi.org/10.1103/PhysRevD.15.1128}{Phys.\ Rev.\ D
\textbf{15} (1977), 1128--1136}, and G.~Kanwar and M.~L.~Wagman,
\href{https://arxiv.org/abs/2103.02602}{arXiv:2103.02602}.
The source-specific lower bounds below are proved directly, not imported
from those works or from the archived transfer-gap certificates.

Let $T\phi=\rho\phi$, $\phi>0$, $\|\phi\|_{L^2(m)}=1$.  The native slice
law and its Markov transfer are
\begin{equation}
 d\pi=\phi^2dm,\qquad
 (Pf)(U)=\rho^{-1}\phi(U)^{-1}T(\phi f)(U),\qquad 0\preceq P\preceq I.
 \label{f16v78:Markov}
\end{equation}
For a real finite family $F$, $G_H(F)$ denotes its \emph{Haar-centered}
Gram, $G_\pi(F)$ its native-centered Gram, and
$C_F(n)=\operatorname{Cov}_\pi(F(U_0),F(U_n))$.  Introduce
\begin{equation}
 R=e^{\gamma N_p+2\gamma M},\qquad
 J=e^{2\gamma(N_p+M)},\qquad
 \alpha=J^{-2}=e^{-4\gamma(N_p+M)}.
 \label{f16v78:constants}
\end{equation}
These large constants are finite-regulator bounds, not pressure-density
or thermodynamic estimates.

\begin{proposition}[Explicit vacuum ratios with the same native top]
\label{f16v78:vacuum}
For the transfer \eqref{f16v78:transfer}, $0<\rho\le1$ and
\begin{equation}
 \frac{\phi_{\max}}{\phi_{\min}}\le R,\quad
 R^{-2}\le\phi^2\le R^2,\quad
 \frac{(b\phi)_{\max}}{(b\phi)_{\min}}\le J,
 \qquad
 R^{-2}G_H(F)\preceq G_\pi(F)\preceq R^2G_H(F).
 \label{f16v78:vacuum-bounds}
\end{equation}
The Gram comparison holds for every square-integrable finite family.
\end{proposition}
\begin{proof}
The magnetic factor lies in $[e^{-\gamma N_p},1]$, and $C$ is a positive
contraction.  Its strictly positive continuous kernel is a product of
normalized one-link kernels.  For any $U,V,Z$ its kernel ratio at
$(U,Z)$ and $(V,Z)$ is at most $e^{2\gamma M}$.  Divide the two Perron
equations after keeping the common positive integrand $b(Z)\phi(Z)$.
This gives $\phi(U)/\phi(V)\le e^{\gamma N_p+2\gamma M}$.
Compactness gives the extrema.  Normalization then implies
$\phi_{\min}\ge R^{-1}$ and $\phi_{\max}\le R$.
Multiplication by $b$ gives the bound $J$.
For any scalar source, minimize
$\int|f-c|^2\phi^2dm$ over constants $c$ and compare it with the same
minimum under Haar.  Apply this in every coefficient direction.  The
actual value of $\rho$ is not replaced; only the valid upper bound
$\rho\le1$ is used later.
\end{proof}

\subsection{An effective lower Gram for a finite polynomial bank}

Let $g=F-\int Fdm$ be a gauge-invariant polynomial family.  Its finitely
many nonconstant Peter--Weyl coefficients form a matrix $A$, with one
column per source.  For a product representation index $\beta$ put
$c_\beta=\prod_e\lambda_{j_e(\beta)}(\gamma)$.  Define the finite matrices
\begin{equation}
 G=A^*A=G_H(F),\qquad E=A^*\operatorname{diag}(c_\beta^{-1})A,
 \qquad
 \sigma=\sup_{v^*Gv>0}\frac{v^*Ev}{v^*Gv}\ge1.
 \label{f16v78:inverse-kinetic}
\end{equation}
All inverses and generalized eigenvalues are taken on the exact quotient
of $G$.  The common cutoff $j_e\le L$ gives the optional bound
\begin{equation}
 \sigma\le\Lambda_L^{-1},\qquad
 \Lambda_L=\lambda_L(\gamma)^M>0.
 \label{f16v78:Lambda}
\end{equation}
The source-specific matrix $E$ can be much smaller than this worst-band
bound.  It is a Haar coefficient calculation, not a native covariance
whose lower eigenvalue has been assumed.

\begin{theorem}[Native polynomial lower Gram with paid degree and delay]
\label{f16v78:polynomial}
For every integer $n\ge1$,
\begin{equation}
 \boxed{\quad
 C_F(n)\succeq(\alpha/\sigma)^nG_\pi(F)
          \succeq R^{-2}(\alpha/\sigma)^nG_H(F).
 \quad}
 \label{f16v78:polynomial-lower}
\end{equation}
Thus every finite-delay native Gram has exactly the Haar source kernel.
No invariance of the polynomial source range under $P$ is assumed.
\end{theorem}
\begin{proof}
Fix a real coefficient vector and write $g$ for its Haar-centered
polynomial, $x=\|g\|_H^2$, $y=\langle g,C^{-1}g\rangle_H\ge x$.
Its native-centered version is $f=g-\mu$, $\mu=\mathbb E_\pi g$.
The constant is an eigenvector of $C$ with eigenvalue one, so
$\|f\|_H^2=x+\mu^2$ and
$\langle f,C^{-1}f\rangle_H=y+\mu^2$.  With $h=b\phi$, spectral
Cauchy--Schwarz gives
\[
 \langle hf,C hf\rangle_H
 \ge\frac{|\langle hf,f\rangle_H|^2}
           {\langle f,C^{-1}f\rangle_H}
 \ge h_{\min}^2\frac{(x+\mu^2)^2}{y+\mu^2}.
\]
The test $f$ belongs to the inverse-kinetic form domain because it is a
finite polynomial plus a constant.  The vector $hf$ need not be a
polynomial; its omitted representations have not been deleted.
Since $G_\pi(f)\le\phi_{\max}^2(x+\mu^2)$,
\[
 \frac{\langle f,Pf\rangle_\pi}{\|f\|_\pi^2}
 \ge\frac{h_{\min}^2}{\rho\phi_{\max}^2}
          \frac{x+\mu^2}{y+\mu^2}
 \ge\alpha\frac{x}{y}\ge\alpha/\sigma.
\]
Here $h_{\min}/\phi_{\max}\ge e^{-\gamma N_p}/R$ and $\rho\le1$;
$y\ge x$ proves the middle fraction inequality.  For a normalized source
vector, its spectral measure for $P$ is a probability on $[0,1]$.
Jensen's inequality for $t\mapsto t^n$ propagates this first-moment lower
bound to every $n\ge1$.  Apply it separately to every source direction;
this is not the generally invalid operation of raising an operator-order
inequality to a power.  The preceding proposition gives the last bound.
Zero Haar variance means a constant source and hence zero native-centered
source, proving both inclusions of the kernel identity.
\end{proof}

For an explicit coefficient-only version of \eqref{f16v78:Lambda}, V63's
ratio barrier gives
\begin{equation}
 \lambda_L(\gamma)\ge\prod_{j=1}^L\frac{\gamma}{\gamma+2j+2}
 \ge\exp[-L(L+3)/\gamma].
 \label{f16v78:Bessel-lower}
\end{equation}
This inherited scalar bound is optional; using the actual Bessel ratios
or a certified enclosure preserves the sharper finite number.

\subsection{A fixed witness avoids a moving-cutoff lower-bound loop}

The previous theorem alone does not solve the reconstruction problem:
$\sigma$ may grow rapidly with the final representation cutoff.  Use
instead a fixed Haar projection $Q_L$ onto product labels $j_e\le L$.
For a general real $L^2(H)$ family with Haar centering, write
\begin{equation}
 g=F-\mathbb E_HF=p+r,\qquad p=Q_Lg,\qquad r=(I-Q_L)g.
 \label{f16v78:witness-split}
\end{equation}
Both terms are Haar-centered, and their Haar cross Gram is zero.
Let $G_p=\int pp^*dm$ and $G_r=\int rr^*dm$.  Suppose
\begin{equation}
 G_r\preceq d^2G_p,\qquad Jd<1,
 \label{f16v78:overlap-condition}
\end{equation}
with the constant $J$ from \eqref{f16v78:constants}.  Define $\sigma_p$
by \eqref{f16v78:inverse-kinetic} for the finite witness $p$.

\begin{theorem}[Native witness bound before kinetic smoothing]
\label{f16v78:witness}
Under \eqref{f16v78:overlap-condition}, put
\begin{equation}
 \theta_L=\frac{\alpha}{\sigma_p}
                   \frac{(1-Jd)^2}{1+d^2}>0.
 \label{f16v78:theta}
\end{equation}
For every integer $n\ge1$,
\begin{equation}
 \boxed{\quad C_F(n)\succeq\theta_L^nG_\pi(F)
                    \succeq R^{-2}\theta_L^nG_H(F).\quad}
 \label{f16v78:witness-lower}
\end{equation}
The tail is paid in the source overlap, before the inverse-kinetic cost.
The condition on $d$ does not grow with the time separation $n$.
\end{theorem}
\begin{proof}
For a coefficient direction put $\mu=\mathbb E_\pi(p+r)$, $t=p-\mu$,
so the native-centered source is $f=t+r$.  Since $r$ is Haar-orthogonal
to both $p$ and constants,
\[
 \|r\|_H\le d\|t\|_H,\qquad
 \|f\|_H^2=\|t\|_H^2+\|r\|_H^2\le(1+d^2)\|t\|_H^2.
\]
The signed overlap obeys
\[
 \langle hf,t\rangle_H
 \ge h_{\min}\|t\|_H^2-h_{\max}\|r\|_H\|t\|_H
 \ge h_{\min}(1-Jd)\|t\|_H^2.
\]
Also
$\langle t,C^{-1}t\rangle_H\le\sigma_p\|t\|_H^2$,
including the constant $-\mu$ because $\sigma_p\ge1$.
Test $C$ against this $t$, exactly as in the preceding theorem, and use
$\|f\|_\pi^2\le\phi_{\max}^2(1+d^2)\|t\|_H^2$.
This proves the first-moment lower quotient $\theta_L$.  Scalar spectral
Jensen and the native--Haar Gram comparison finish the proof.  No inverse
of $C$ is applied to the nonpolynomial remainder $r$.
\end{proof}

For a fixed finite family independent modulo constants, the condition
\eqref{f16v78:overlap-condition} holds for some finite $L$: product
Peter--Weyl projections converge in Haar $L^2$, whereas the least
nonzero eigenvalue of $G_H(F)$ is fixed and positive.  This fact alone is
not an effective estimate on a changing regulator.  The next calculation
supplies explicit values for the actual simultaneous-flow family.

\subsection{An explicit witness for the original simultaneous-flow source}

Use exactly V69's $s=0$ source on the complete finite spatial torus:
\begin{equation}
 O_i(W)=\frac{\zeta_a^2}{2a}\sum_x f_i(x)|\operatorname{Im}W_{12}(x)|^2,
 \qquad F_i(U)=O_i(\Phi_{a,\tau}U),\qquad u=\tau/a^2.
 \label{f16v78:flow-source}
\end{equation}
Native centering is implicit only when a covariance is taken.
Repeated marks are first combined.  For the sampled coefficient matrix
put
\begin{equation}
 \mathcal S_{ij}=\sum_x f_i(x)f_j(x),\qquad
 G_0=\frac{\zeta_a^4}{64a^2}\mathcal S,
 \qquad \zeta_a>0.
 \label{f16v78:sampling}
\end{equation}
This is an exactly known matrix of samples, not a native source Gram.

\begin{proposition}[Haar rank and an effective simultaneous-flow witness tail]
\label{f16v78:flow-tail}
For \eqref{f16v78:flow-source},
\begin{equation}
 G_H(O)=G_0,\qquad G_H(F)\succeq e^{-12N_pu}G_0,
 \qquad \ker G_H(F)=\ker\mathcal S.
 \label{f16v78:flow-Haar}
\end{equation}
For $g=F-\mathbb E_HF$ and its split \eqref{f16v78:witness-split},
\begin{equation}
 G_r\preceq t_LG_H(F),\qquad
 t_L=\frac{128e^{(32+12N_p)u}}{(L+1)(L+3)}.
 \label{f16v78:tail-explicit}
\end{equation}
These assertions hold for signed spatial tests and every finite $\tau$.
\end{proposition}
\begin{proof}
For a Haar plaquette, $\chi_2=4w^2-1$ has mean zero and squared norm
one, so $|\operatorname{Im}U_P|^2=3/4-\chi_2(U_P)/4$ has variance
$1/16$.  Distinct elementary $12$-plaquettes have a boundary link absent
from the other; integrating it makes their $\chi_2$ cross moment zero.
The side-length condition excludes degeneracies.  This proves $G_H(O)=G_0$.

V74's exact spatial-flow Jacobian is
\[
 \log\det D\Phi_{a,\tau}(U)
       =-\frac{12}{a^2}\int_0^\tau\sum_Pw((\Phi_{a,v}U)_P)\,dv.
\]
Consequently the pushforward of \emph{Haar} by this finite diffeomorphism
has density at least $e^{-12N_pu}$.  Minimize the variance over constants
to obtain the second assertion.  This use of the Haar pushforward
estimates a Haar coefficient norm only; the native covariance remains
under $\pi$.  Surjectivity and the first identity give the exact kernel.

The Hessian of the unscaled spatial action has diagonal block norm at
most four and off-diagonal block row sum at most twelve, by V69's exact
plaquette calculation.  Its operator norm is at most sixteen.  The native
flow variational equation thus gives
$\|D\Phi_{a,\tau}\|\le e^{16u}$ in the unscaled product metric.
For a plaquette with transported tangent sum $v_P$,
\[
 |D|\operatorname{Im}W_P|^2[v]|
       =|2w(W_P)\operatorname{Im}W_P\cdot v_P|\le |v_P|.
\]
There are four edges per marked plaquette and at most two such plaquettes
per link.  Cauchy--Schwarz therefore proves, in every test direction,
\[
 \|\nabla O_c(W)\|^2\le\frac{2\zeta_a^4}{a^2}\sum_x f_c(x)^2
                       =128\,c^*G_0c.
\]
After the full flow, integrate the bound
$\|\nabla F_c\|^2\le128e^{32u}c^*G_0c$ under Haar.
The positive product Haar Laplacian has eigenvalue
$\sum_e j_e(j_e+2)$ on product labels.  Every mode omitted by $Q_L$
has eigenvalue at least $(L+1)(L+3)$.  Its spectral energy inequality gives
\[
 \|r_c\|_H^2\le
 \frac{128e^{32u}}{(L+1)(L+3)}c^*G_0c.
\]
Use \eqref{f16v78:flow-Haar} to conclude \eqref{f16v78:tail-explicit}.
The flow is neither split into plaquette subflows nor replaced by its
linearization anywhere in this proof.
\end{proof}

\begin{theorem}[A populated finite-regulator lower Gram for the collective bank]
\label{f16v78:flow-lower}
Choose an integer $L\ge1$ such that
\begin{equation}
 (L+1)(L+3)\ge
 128e^{(32+12N_p)u}(1+16J^2),\qquad
 \Theta=\frac{\alpha\Lambda_L}{5}>0.
 \label{f16v78:explicit-L}
\end{equation}
Then, for every integer $n\ge1$,
\begin{equation}
 \boxed{\quad
 C_F(n)\succeq R^{-2}e^{-12N_pu}\Theta^n
                   \frac{\zeta_a^4}{64a^2}\mathcal S.
 \quad}
 \label{f16v78:actual-lower}
\end{equation}
Every finite-delay native source Gram has kernel exactly
$\ker\mathcal S$.  For common temporal weights $h_j\ge0$, $\sum_jh_j=1$,
at positive integer times $t_j$, and additional delay $n\ge0$, its
reflected Gram is bounded below by the same matrix with $\Theta^n$
replaced by
\begin{equation}
 \Theta^n\left(\sum_jh_j\Theta^{t_j}\right)^2.
 \label{f16v78:window-lower}
\end{equation}
\end{theorem}
\begin{proof}
Since $G_H(F)=G_p+G_r$, the explicit tail implies
$G_r\preceq[t_L/(1-t_L)]G_p$.  The chosen $L$ gives
$d\le1/(4J)$.  Also $\sigma_p\le\Lambda_L^{-1}$.
The witness factor in \eqref{f16v78:theta} is then at least
$\alpha\Lambda_L(3/4)^2/(1+1/16)$, which is larger than
$\Theta$.  Combine the witness theorem with the Haar Gram lower bound.
A sampled zero combination is the zero function; conversely the positive
lower bound excludes every other null direction.  For the reflected
window each matching matrix is $C_F(n+t_i+t_j)$, so summing the
Loewner lower bounds with the nonnegative weights proves
\eqref{f16v78:window-lower}.  The time step $a$ is not changed.
\end{proof}

Unlike an assertion that a finite-regulator source is nonconstant, this
is an explicit quantitative lower matrix depending only on the declared
lattice, coupling, flow time, normalization and sampled tests.  It does
not claim to approximate the magnitude of that matrix sharply.

\subsection{A stable lower bound for the V77 reconstruction itself}

Let $F^K=\mathcal A_{K,S}F$ be exactly V77's positive filter on all original
input links, with the same function $F$ and independent native centering.
The witness resolution $L$ is fixed first.  It is \emph{not} replaced by
$K$ when the latter is increased.

\begin{theorem}[Two resolutions and a terminating finite positivity certificate]
\label{f16v78:stable-K}
Use the $L$ in \eqref{f16v78:explicit-L}.  For every integer $K\ge L$
satisfying
\begin{equation}
 K+3\ge2M L(L+2),
 \label{f16v78:K-witness}
\end{equation}
the reconstructed collective bank obeys, for $n\ge1$,
\begin{equation}
 \boxed{\quad
 C_{F^K}(n)\succeq
 \frac{R^{-2}e^{-12N_pu}}5\Theta^nG_0.
 \quad}
 \label{f16v78:stable-lower}
\end{equation}
Its exact source kernel is $\ker\mathcal S$.  The right side is independent
of the final cutoff $K$ once \eqref{f16v78:K-witness} holds.
\end{theorem}
\begin{proof}
The filter commutes with $Q_L$ and $C$ and contracts Haar $L^2$.
V77's multiplier can be written
\[
 \alpha_{K,j}=\prod_{s=0}^{j-1}\frac{K-s}{K+s+3},\qquad
 1-\alpha_{K,j}\le\frac{j(j+2)}{K+3}.
\]
For the last inequality use $1-\prod(1-x_s)\le\sum x_s$ with
$x_s=(2s+3)/(K+s+3)$.  Hence the product multiplier on the entire
witness range is at least
\[
 \omega=\alpha_{K,L}^{M}\ge1-\frac{ML(L+2)}{K+3}\ge\frac12.
\]
Write $p^K=\mathcal A p$, $r^K=\mathcal A r$.  Then
$G_{p^K}\succeq\omega^2G_p$ and
$G_{r^K}\preceq G_r\preceq(d/\omega)^2G_{p^K}$.
Its witness parameter is therefore at most $1/(2J)$, and its inverse
kinetic quotient is at most $\Lambda_L^{-1}$, independently of its
higher-mode remainder.  The factor
$(1-Jd_K)^2/(1+d_K^2)$ is at least $1/5$.
Apply the witness theorem to $F^K$ to get $C_{F^K}(n)\succeq
R^{-2}\Theta^nG_H(F^K)$.  Finally
\[
 G_H(F^K)\succeq\tfrac14G_p
 \succeq\frac{1}{4(1+d^2)}G_H(F)
 \succeq\frac15G_H(F)
 \succeq\frac15e^{-12N_pu}G_0.
\]
The positive lower matrix and preservation of every original relation
prove the exact kernel assertion.  No claim that a filter preserves the
kernel of every arbitrary bank at every cutoff is made; V77's contrary
example remains valid below the witness threshold.
\end{proof}

On a specified coefficient subspace let $s_*>0$ be the least eigenvalue of
$\mathcal S$ and fix a finite lag $n\ge2$.  Define the completely explicit
positive number
\begin{equation}
 h_n=\frac{R^{-2}e^{-12N_pu}}5\Theta^n
                           \frac{\zeta_a^4}{64a^2}s_*.
 \label{f16v78:h-populated}
\end{equation}
V77's error is at most $18\pi^2M\gamma B_F^2/(K+3)$, where
$B_F^2$ is its half-oscillation budget.  Thus, in addition to
\eqref{f16v78:K-witness}, choosing
\begin{equation}
 K+3\ge36\pi^2M\gamma B_F^2/h_n
 \label{f16v78:finite-certificate}
\end{equation}
makes its reconstruction error at most $h_n/2$.
The finite lower-Gram certificate of V77 is now populated and terminates
at each fixed regulator, lag and positive sampling quotient.  Its lower
constant is not diminished by this last increase of $K$.  Numerical
coefficient errors would have to be added using V77's separate budget;
no full set of such coefficients is claimed computed here.

\subsection{Sharp kinetic test, physical cost, and the remaining task}

The degree and time dependence cannot be omitted even when the magnetic
factor is absent.  In the exactly solvable \emph{reference} $b=1$,
$\phi=1$, $P=C$, a nonconstant Haar-orthonormal bank contained in a single
product sector of multiplier $c_*$ has
\begin{equation}
 C_F(n)=c_*^nG_H(F),\qquad \sigma=c_*^{-1}.
 \label{f16v78:reference-sharp}
\end{equation}
With the exact constants $R=J=\alpha=1$ and no witness tail, the polynomial
theorem is an equality.  For a four-edge plaquette source
$|\operatorname{Im}U_P|^2$, its nonconstant part is $-\chi_2(U_P)/4$;
therefore $c_*=\lambda_2(\gamma)^4$ and its covariance is exactly
$\lambda_2(\gamma)^{4n}/16$ in this reference.  This calculation tests
the source-specific inverse-kinetic quotient, not the four-dimensional
interacting vacuum.  The full magnetic action is present in every native
bound above.

The explicit physical limitations are severe.  Already
$\alpha^n=\exp[-4\gamma(M+N_p)n]$ tends to zero on a fixed positive
physical volume and delay with $M,N_p\asymp a^{-3}$,
$n\asymp a^{-1}$ and $\gamma\to\infty$.  The Haar-flow comparison adds
$\exp[-12N_p\tau/a^2]$.  Moreover the sufficient witness choice satisfies,
up to fixed numerical constants,
\begin{equation}
 \log L=O\!\left((16+6N_p)\tau/a^2+
                         2\gamma(N_p+M)+1\right).
 \label{f16v78:cost}
\end{equation}
For fixed physical $\tau>0$ its leading displayed term is of order
$a^{-5}$ on fixed physical volume.  The Bessel product $\Lambda_L$ can
make the lower constant much smaller still.  The final sufficient $K$
in \eqref{f16v78:finite-certificate} need not be polynomial in $a^{-1}$.
This does not contradict V77's polynomial cutoff for a prescribed
\emph{absolute} error; resolving the present tiny lower bound is a much
more demanding target.

These are limitations of the explicit sufficient lower bounds, not upper
bounds on the true correlations and not a collapse theorem for the
simultaneously flowed source.  We have closed the fixed-regulator logical
gap between a reconstructed bank and a certified positive native Gram.
We have not established a positive cutoff-uniform lower Gram.  In
particular, a lower correlation bound is not an upper contraction bound
and cannot by itself establish a mass gap.

A noncircular next step is now more precise: estimate the source's
\emph{actual weighted low-kinetic overlap} without replacing $b\phi$ by
its whole-volume extrema, or control a reflected Gram directly in its
physical normalization.  The test in Theorem~\ref{f16v78:witness} retains
that overlap before any kinetic inverse is used; its estimates show
exactly where the extensive worst-case loss enters.  Another increase
of the output cutoff or another bare positivity assertion does not remove
that loss.  A useful improvement must also address the growing physical
source support and the limit order, rather than relabel the finite
constants as uniform ones.

\subsection{Verification, preservation and result ledger}

The companion diagnostic is
\begin{center}\small
\path{verify_ym_f16_v78_kinetic_witnesses.py}.
\end{center}
It checks exact Haar and multiplier identities, the centering-sensitive
inverse-kinetic inequality, finite positive convolution/magnetic
sandwiches, noninvariant source ranges, witness tails and two-resolution
bounds, literal lattice source gradients, and the physical cost ledger.
Finite-group and compact-quadrature fixtures are identified separately.
They are not four-dimensional Wilson-vacuum simulations, rigorous
interval certificates or proof-assistant verification.  No numerical
outcome is a premise of the analytic arguments.  Removing this delimited
appendix and reverting the single displayed version change recovers V77
byte for byte; the inherited archive is preserved, not recertified.

\begin{record}[V78 result ledger]
\label{f16v78:ledger}
\emph{Proved here:} a native all-direction polynomial covariance lower
bound in terms of a finite Haar inverse-kinetic matrix; an overlap witness
bound for nonpolynomial sources with native centering paid explicitly;
the full simultaneous-flow sampling kernel and an effective Haar witness
tail; an explicit strictly positive finite-regulator reflected lower
matrix for the actual collective source; a reconstruction lower bound
stable under increasing the final cutoff; and a terminating, though
potentially prohibitive, finite positivity certificate.

\emph{Inherited or established background:} the positive compact Wilson
transfer and its magnetic sandwich, its representation multipliers, V69's
compact Hessian blocks, V74's flow Jacobian, V77's positive input filter and
reflected error, spectral Cauchy--Schwarz and scalar Jensen.  The archived
finite-rank transfer-gap certificates are not claimed again as new results.

\emph{Not asserted:} a useful sharp value of the native Gram; a polynomial
cost for resolving this tiny lower constant; a bound uniform in spatial
volume, continuum cutoff or physical delay; invariance of a polynomial
range under the interacting transfer; a discarded native centering term;
or a physical mass gap from a lower correlation estimate.

\emph{Still unproved:} a nonzero fixed-physical-scale native reflected Gram
with cutoff-uniform normalization for the simultaneous-flow bank,
its thermodynamic control, common dense-core physical-time contraction,
and the interacting four-dimensional continuum Yang--Mills mass gap.
\end{record}

\begin{firewall}
The native lower Gram is bounded, not assumed.  Its finite witness is
chosen before the reconstruction cutoff, and every extensive vacuum,
flow and kinetic cost remains visible.  Finite positivity is separated
from physical noncollapse and from a spectral gap.  The physical transfer
and all V1--V77 mathematical text remain unchanged.
\end{firewall}
% END F16 V78 APPEND

% BEGIN F16 V79 APPEND -- 2026-09-19
% Append-only continuation of the Post-Fifteenth-Archive V78.
\clearpage
\section[V79: exact vacuum cancellation and volume-uniform source bounds]{V79: exact vacuum cancellation\\and volume-uniform source bounds}
\label{f16v79:context}

V78 makes the finite-regulator lower Gram explicit, but replaces the
vacuum weight by extrema over the entire slice.  We remove that operation,
not merely improve its constants.  The exponential quaternion kernel has
an exact coherent-vector realization.  A polynomial creation-operator
test pairs with the source through the Perron equation and yields its
\emph{native} entrance Gram exactly.  Normal ordering expresses the test
norm as a finite sum of source derivatives.  The unknown vacuum weight
and its normalization cancel, with no whole-volume density comparison.
The coherent vectors are an auxiliary realization of the original compact
kernel, not a Gaussian approximation to the Wilson measure.

The resulting bound applies in particular to V77's reconstructed
simultaneous-flow bank.  We also populate it, without a source-count loss,
for arbitrary weighted sums of all spatial plaquette traces.  A two-stage
private-link argument supplies a native entrance lower Gram for this
entire bank; its derivative cost is a configuration-independent polynomial
in $\gamma^{-1}$ times the sampling Gram.  Both constants are independent
of the surrounding volume and of the number of marked plaquettes.
They still do not provide fixed-physical-time noncollapse as $a\to0$.

\subsection{The native transfer and polynomial representatives}

Use precisely the finite-regulator transfer and normalization of
\eqref{f16v78:transfer}--\eqref{f16v78:Markov}.  With $M$ spatial links,
$x(U)\in(S^3)^M\subset\mathbb R^{4M}$, normalized Haar measure $dm$, and
$z_\gamma=2I_1(\gamma)/\gamma$, write
\begin{equation}
 C(U,V)=z_\gamma^{-M}e^{\gamma x(U)\cdot x(V)},\quad
 T=bCb,\quad T\phi=\rho\phi,\quad d\pi=\phi^2dm,\quad h=b\phi.
 \label{f16v79:setup}
\end{equation}
In particular the exact pointwise Perron identity is
\begin{equation}
                    h\,Ch=\rho\phi^2.\label{f16v79:Perron}
\end{equation}
The positive Markov extension $\mathcal P=\rho^{-1}\phi^{-1}T\phi$ acts
on all slice functions and agrees with the physical $P$ on gauge-invariant
functions.  We use the extension only for auxiliary derivative functions;
no Gauss projection is inserted into their intermediate contractions.
It has the same positive top and satisfies $0\preceq\mathcal P\preceq I$.
All stated source correlations use the physical gauge-invariant bank.

The positive Wilson transfer is inherited, not reproved.  For external
context see M.~Creutz, \href{https://doi.org/10.1103/PhysRevD.15.1128}
{Phys.\ Rev.\ D \textbf{15} (1977), 1128--1136}.
The exponential coherent kernel and creation/annihilation calculus are
established tools; see V.~Bargmann,
\href{https://doi.org/10.1002/cpa.3160140303}
{Comm.\ Pure Appl.\ Math.\ \textbf{14} (1961), 187--214}.
The particular compact-kernel/native-source identity needed here is
proved below.  No statement about a free physical vacuum is imported.

Let $F=(F_1,\ldots,F_d)$ be real gauge-invariant polynomials on the slice.
Fix real ambient polynomial representatives $\widetilde F_i(x)$, and put
$f_i=F_i-\mathbb E_\pi F_i$.  The derivatives below are \emph{ambient}
quaternion-coordinate derivatives, not tangent derivatives on $S^3$.
A useful fixed choice is the separate harmonic extension: expand each
link's restriction in homogeneous harmonic polynomials and use their
sum in $\mathbb R^4$.  This choice is linear, equivariant under orthogonal
changes of link coordinates, and sends an exact zero restriction to the
zero ambient polynomial.  An arbitrary other representative is allowed
but can give a larger, representation-dependent derivative bound.
Native centering changes only the constant polynomial term.

For multiindices $\beta\in\mathbb N^{4M}$, set
\begin{align}
 G&=\operatorname{Cov}_\pi(F),\qquad C_1=C_F(1),\nonumber\\
 (\mathsf B_\gamma)_{ij}
  &=\sum_{|\beta|\ge1}\frac{\gamma^{-|\beta|}}{\beta!}
       \mathbb E_\pi[(\partial^\beta\widetilde F_i)
                         (\partial^\beta\widetilde F_j)],
       \label{f16v79:B}\\
 (\mathsf R_\gamma)_{ij}
  &=\sum_{|\beta|\ge1}\frac{\gamma^{-|\beta|}}{\beta!}
       \langle\partial^\beta\widetilde F_i,
       \mathcal P\partial^\beta\widetilde F_j\rangle_\pi.
       \label{f16v79:R}
\end{align}
These sums are finite.  The derivatives in both expressions are
\emph{not centered}.  Positivity and contraction give
$0\preceq\mathsf R_\gamma\preceq\mathsf B_\gamma$.
The multiindex sum equals the sum of fully contracted order-$k$ derivative
tensors with coefficients $(k!\gamma^k)^{-1}$, so both costs respect
gauge rotations of the fixed equivariant representatives.

\subsection{An exact source witness cancels the vacuum weight}

\begin{theorem}[Native polynomial block Gram without vacuum extrema]
\label{f16v79:block}
For the original magnetic sandwich \eqref{f16v79:setup} and any such
finite polynomial bank,
\begin{equation}
 \boxed{\qquad
 \begin{pmatrix}C_1&G\\G&C_1+\mathsf R_\gamma\end{pmatrix}\succeq0,
 \qquad
 \begin{pmatrix}C_1&G\\G&C_1+\mathsf B_\gamma\end{pmatrix}\succeq0.
 \qquad}\label{f16v79:block-Gram}
\end{equation}
No minimum or maximum of $b$, $\phi$, or $b\phi$ occurs.  The actual
$\rho$ cancels rather than being replaced by an upper bound.
\end{theorem}
\begin{proof}
Let $\mathcal F$ be symmetric Fock space over $\mathbb C^{4M}$, with
exponential vectors
$\mathrm e(z)=\bigoplus_{n\ge0}z^{\otimes n}/\sqrt{n!}$.
Their inner product is $e^{\bar z\cdot w}$.  Define
\[
 \Psi(U)=z_\gamma^{-M/2}\mathrm e(\sqrt\gamma\,x(U)),\qquad
 \eta_q=\int q(U)\Psi(U)\,dm(U).
\]
Thus $\langle\Psi(U),\Psi(V)\rangle=C(U,V)$ and
$\langle\eta_q,\eta_r\rangle=\langle q,Cr\rangle_H$ for real functions.
The annihilation operators $a_A$ obey
$a_A\Psi(U)=\sqrt\gamma\,x_A(U)\Psi(U)$ and
$[a_A,a_B^*]=\delta_{AB}$.  All polynomial operator expressions below
are well defined on $\eta_h$: creation-operator norms on an exponential
vector are bounded by a fixed polynomial in $|z|$ times $e^{|z|^2/2}$,
and $|z|^2=\gamma M$ on the compact integration domain.  Finite-degree
normal ordering can first be performed on exponential vectors and then
integrated.  No unbounded-operator domain interchange is being assumed.

Use the two families of auxiliary vectors
\begin{equation}
 V_i=\rho^{-1/2}\eta_{hf_i},\qquad
 W_i=\rho^{-1/2}\widetilde f_i(a^*/\sqrt\gamma)\eta_h.
 \label{f16v79:VW}
\end{equation}
Evaluation against $\Psi(U)$ gives
$\langle\Psi(U),\widetilde f_j(a^*/\sqrt\gamma)\eta_h\rangle
=f_j(U)(Ch)(U)$.  Therefore
\[
 \langle V_i,V_j\rangle=(C_1)_{ij},\qquad
 \langle V_i,W_j\rangle
 =\rho^{-1}\int h f_if_j Ch\,dm
 =\int f_if_j\phi^2dm=G_{ij}.
\]
This is where the Perron identity, with native centering, removes the
unknown weight exactly.  It is not a bound on an unweighted Haar overlap.

For real polynomials $p,q$ the finite normal-ordering identity is
\begin{equation}
 p(a/\sqrt\gamma)q(a^*/\sqrt\gamma)
 =\sum_\beta\frac{\gamma^{-|\beta|}}{\beta!}
 (\partial^\beta q)(a^*/\sqrt\gamma)
 (\partial^\beta p)(a/\sqrt\gamma).
 \label{f16v79:normal-order}
\end{equation}
For one mode it follows by induction from
$a^m(a^*)^n=\sum_k\binom mk\binom nk k!(a^*)^{n-k}a^{m-k}$;
commuting different modes and extending linearly proves the formula.
Also $(\partial^\beta\widetilde f_i)(a/\sqrt\gamma)\eta_h
=\eta_{h\partial^\beta\widetilde f_i}$.  Consequently
\[
 \langle W_i,W_j\rangle
 =\rho^{-1}\sum_\beta\frac{\gamma^{-|\beta|}}{\beta!}
 \langle h\partial^\beta\widetilde f_i,
                 C h\partial^\beta\widetilde f_j\rangle_H
 =(C_1+\mathsf R_\gamma)_{ij}.
\]
The zero derivative supplies $C_1$; the others are uncentered derivative
pairings.  The Gram of $(V,W)$ proves the first matrix inequality.
Since the Markov extension is a positive contraction, replacing
$\mathsf R_\gamma$ by $\mathsf B_\gamma$ adds a positive block and proves
the second.  Every quantity on the left concerns the original native
source, even though the feature space and derivative functions are auxiliary.
\end{proof}

It would be incorrect to center the derivatives in \eqref{f16v79:B}.
For a linear polynomial its derivatives are constants; deleting them
would force $C_1\succeq G$ and hence falsely make a nonconstant source
invariant.  Likewise the auxiliary coherent vectors do not identify the
Wilson state with a bosonic vacuum or add physical field species.

\begin{theorem}[Source-directed delay bounds and the physical-time test]
\label{f16v79:delay}
The same bank satisfies
\begin{equation}
 0\preceq G-C_1\preceq\tfrac12\mathsf R_\gamma
                      \preceq\tfrac12\mathsf B_\gamma.
 \label{f16v79:deficit}
\end{equation}
If $\mathsf B_\gamma\preceq\delta G$ on a declared source quotient, then
\begin{equation}
 \boxed{\quad C_F(n)\succeq q(\delta)^nG,\qquad
 q(\delta)=\frac{2}{\sqrt{\delta^2+4}+\delta}
          =e^{-\operatorname{arsinh}(\delta/2)},\quad n\ge1.\quad}
 \label{f16v79:delay-bound}
\end{equation}
Common positive-time weights have the lower factor
$q(\delta)^n(\sum_j h_j q(\delta)^{t_j})^2$, as in V78.
In particular $\mathsf B_{\gamma_a}\preceq2aE_*G_a$ implies
$C_F(n)\succeq e^{-naE_*}G_a$ and
\begin{equation}
 \mu_{a,c}((E,\infty))\le
 \frac{aE_*}{1-e^{-aE}},\qquad E>0,
 \label{f16v79:spectral-tail}
\end{equation}
for each normalized source's physical spectral measure.
\end{theorem}
\begin{proof}
Test the first block Gram on $(c,-c)$ and use $P\preceq I$ to obtain
\eqref{f16v79:deficit}.  For a nonzero source direction set
$g=c^*Gc$, $x=c^*C_1c$.  Its scalar $2\times2$ Gram gives
$g^2\le x(x+c^*\mathsf B_\gamma c)\le x(x+\delta g)$.
The positive root is $x/g\ge q(\delta)$.  Scalar Jensen in the spectral
probability measure of this direction gives $C_F(n)\succeq q(\delta)^nG$.
No noncommuting matrix square root or power inequality is used.
Sum over temporal weights to get the stated factor.  Finally
$\operatorname{arsinh}(aE_*)\le aE_*$ and
$\langle 1-P\rangle\le aE_*$ prove both physical-time statements;
on physical energies exceeding $E$, $1-P\ge1-e^{-aE}$.
\end{proof}

The condition $\mathsf B_{\gamma_a}=O(aG_a)$ is a sufficient
\emph{source} nonescape test, not a physical gap condition.  It allows
spectral weight near zero.  It also supplies no absolute nonzero limit of
$G_a$ unless its normalization and rank have been established separately.

\subsection{Finite coefficient certificates and exact source relations}

The derivative matrix is a native expectation, but it admits a
coefficient-only upper bound.  Suppose the chosen representatives have
separate degrees at most $L_e$ on a finite support $S$, and set
\begin{equation}
 \mathcal N_{\boldsymbol L}=
 \prod_{e\in S}\frac{(L_e+1)(L_e+2)(2L_e+3)}6,\qquad
 \mathsf B^H_\gamma=
 \sum_{|\beta|\ge1}\frac{\gamma^{-|\beta|}}{\beta!}
                  \int\partial^\beta\widetilde F
                       (\partial^\beta\widetilde F)^*dm.
 \label{f16v79:coefficient-budget}
\end{equation}

\begin{proposition}[A finite coefficient bound with only support costs]
\label{f16v79:coefficients}
Pointwise and under the native law,
\begin{equation}
 \sum_{|\beta|\ge1}\frac{\gamma^{-|\beta|}}{\beta!}
 \partial^\beta\widetilde F(\partial^\beta\widetilde F)^*
 \preceq\mathcal N_{\boldsymbol L}\mathsf B^H_\gamma,
 \qquad
 \mathsf B_\gamma\preceq\mathcal N_{\boldsymbol L}\mathsf B^H_\gamma.
 \label{f16v79:coefficient-upper}
\end{equation}
Together with V65's original-law polynomial bound this gives
\begin{equation}
 G\succeq\mathcal L_{\gamma,\boldsymbol L}^{-1}G_H(F),\quad
 \delta=\mathcal L_{\gamma,\boldsymbol L}\mathcal N_{\boldsymbol L}
 \lambda_{\max}\!\left(G_H(F)^{-1/2}\mathsf B^H_\gamma G_H(F)^{-1/2}\right)
 \label{f16v79:coefficient-delta}
\end{equation}
as a sufficient value in \eqref{f16v79:delay-bound}.  Inverses are on
the exact Haar quotient, using representatives respecting its relations.
No constant depends on untouched surrounding links.
\end{proposition}
\begin{proof}
The diagonal of the reproducing kernel of the separate-degree product
space is the constant $\mathcal N_{\boldsymbol L}$, by orthonormal
matrix coefficients and Haar homogeneity.  Thus $|p(U)|^2\le
\mathcal N_{\boldsymbol L}\|p\|_H^2$ for every polynomial in that space.
Each ambient derivative, restricted to the spheres, belongs to the same
space.  Apply the inequality to every derivative of a source combination
and sum.  This proves the pointwise matrix bound without changing the
measure.  For the first assertion of \eqref{f16v79:coefficient-delta},
apply \eqref{f16v65:leverage} to $F_c-\mathbb E_\pi F_c$ and use
$\|F_c-\mathbb E_\pi F_c\|_H^2\ge\operatorname{Var}_H(F_c)$.
The generalized eigenvalue comparison completes the proof.
\end{proof}

All entries of $\mathsf B^H_\gamma$ are finite Haar polynomial integrals.
For exponents $2r_i$ on one link they are evaluated by
\[
 \int\prod_{i=0}^3x_i^{2r_i}\,d\sigma
 =\frac{\prod_i(2r_i-1)!!}{4\cdot6\cdots(2\sum_i r_i+2)};
\]
odd-coordinate moments vanish, and the empty product is one.
The separate harmonic representative prevents a zero sphere restriction
from leaving a spurious derivative cost.  The simple ambient zero
$\sum_i x_i^2-1$ illustrates why that convention matters.
A common coefficient map then preserves every exact source relation.
The support and degree products in \eqref{f16v79:coefficient-delta}
remain explicit; they need not be small for V77's whole-slice reconstruction.

\subsection{A populated native entrance Gram for all spatial plaquettes}

We can do substantially better for a large concrete bank, without paying
its full support product.  Work on an even spatial torus of side at least
four, and let $\mathcal P$ be any set of distinct unoriented elementary
spatial plaquettes.  Repeated or reversed marks are combined first.
For real coefficients $s_{ip}$ set
\begin{equation}
 F_i=\sum_{p\in\mathcal P}s_{ip}w(U_p),\qquad
 \mathcal S_{ij}=\sum_{p\in\mathcal P}s_{ip}s_{jp},\qquad
 d_\gamma=\frac1{32000(1+6\gamma)^2},\quad
 \kappa_\gamma=\frac43d_\gamma^2.
 \label{f16v79:plaquette-bank}
\end{equation}
This includes signed spatial sampling, any number of plaquettes, and all
three orientations.  The same centered bank is obtained, up to sign,
from $1-w(U_p)$.

\begin{theorem}[Volume- and mark-count-uniform native entrance bound]
\label{f16v79:entrance}
Under the original Wilson stationary history,
\begin{equation}
                 \boxed{\quad G_\pi(F)\succeq
                       \kappa_\gamma\mathcal S.\quad}
 \label{f16v79:entrance-lower}
\end{equation}
The same bound holds in finite Wilson cylinders with the chosen slice
interior.  It passes to the corresponding local limits.  Its constant is
independent of the number, positions, and overlaps of marked plaquettes.
\end{theorem}
\begin{proof}
It suffices to prove the scalar assertion for coefficients $s_p$.
Use the twelve link colors of \eqref{f16v59:12-coloring}:
$\operatorname{col}(x,\mu)=(\mu,(x_\nu\bmod2)_{\nu\ne\mu})$.
No four-dimensional Wilson plaquette contains two links of one color
on the chosen slice.  Conditional on the \emph{full history} outside
that color, the links are independent with fields of norm at most
$6\gamma$.  The inherited one-link bound is
$\operatorname{Cov}(x_e\mid\mathrm{complement})\succeq d_\gamma I_4$.
Because each marked plaquette has at most one free link,
\[
 F=\text{constant}+\sum_{e\in\mathrm{color}}x_e\cdot B_e,\qquad
 B_e=\sum_{p\ni e}s_p\,p_{p,e},
\]
where $p_{p,e}$ is the actual complementary three-link path, oriented
as the stored link.  Conditional variance and then averaging give
$\operatorname{Var}(F)\ge d_\gamma
\sum_{e\in\mathrm{color}}\mathbb E|B_e|^2$.

For each fixed $e=(x,\mu)$, consider the four opposite $\mu$-links
in its four incident spatial plaquettes, at $x\pm\hat\nu$, $\nu\ne\mu$.
They are distinct, none is a connector in any of those four paths, and
no elementary plaquette contains two of them: their relative displacements
are two transverse steps or a sum of two distinct transverse steps, never
one.  The even side-length and nondegeneracy assumptions also handle the
periodic complement.  Conditional on all other history links, these four
private links are independent.  Each path vector $p_{p,e}$ is an
orthogonal linear image of its own private unit quaternion, even when
its orientation is reversed.  Hence
\[
 \operatorname{Cov}(B_e\mid\mathrm{complement})
       \succeq d_\gamma\sum_{p\ni e}s_p^2 I_4,
 \qquad
 \mathbb E|B_e|^2\ge4d_\gamma\sum_{p\ni e}s_p^2.
\]
This second conditional calculation estimates an unconditional moment;
it is not a sequential claim that a marginalized color retains an
independent-link Wilson law.

Average the twelve separate lower bounds, rather than adding twelve
copies of the same variance.  Each marked plaquette has four edges, so
\[
 \operatorname{Var}(F)\ge
 \frac{4d_\gamma^2}{12}\sum_e\sum_{p\ni e}s_p^2
 =\frac43d_\gamma^2\sum_p s_p^2.
\]
All conditionings were in the original finite Euclidean history.  Uniform
constants and bounded local moments pass the result to stationary and
local limits.  Applying the scalar assertion in every coefficient
direction proves the matrix bound.
\end{proof}

\subsection{The complete plaquette derivative budget is local}

\begin{proposition}[All derivative orders for a compact Wilson plaquette]
\label{f16v79:plaquette-derivatives}
For the multilinear quaternion representative of $W=w(U_p)$ and each
specified subset $J$ of $k$ distinct boundary links,
\begin{equation}
 \sum_{A_e\in\{0,1,2,3\},\ e\in J}
        |\partial_{(e,A_e):e\in J}W|^2=4^{k-1},\qquad 1\le k\le4.
 \label{f16v79:derivative-exact}
\end{equation}
It follows that the bank \eqref{f16v79:plaquette-bank} obeys pointwise
and natively
\begin{equation}
 \boxed{\quad \mathsf B_\gamma(F)\preceq4u_\gamma\mathcal S,\qquad
 u_\gamma=\frac{(1+4/\gamma)^4-1}{4}
 =\frac4\gamma+\frac{24}{\gamma^2}+\frac{64}{\gamma^3}
                                      +\frac{64}{\gamma^4}.\quad}
 \label{f16v79:plaquette-budget}
\end{equation}
The factor four is a local spatial-plaquette incidence, not the number
of marks.  Keeping only one plaquette orientation permits replacing it
by two.
\end{proposition}
\begin{proof}
Fix all but one differentiated coordinate.  The scalar part of an ordered
quaternion product is a linear functional of the remaining quaternion
with squared coefficient norm equal to the product of the squared norms
of the other factors.  Every unit quaternion and every quaternion basis
vector has norm one.  Summing that last coordinate therefore gives one;
the other $k-1$ free coordinates each contribute a factor four.
Conjugation of reversed links is orthogonal and does not change the sum.
No link is differentiated twice because $W$ is multilinear.  Summing over
its $\binom4k$ subsets proves the exact one-plaquette cost $u_\gamma$.

For any nonzero multiindex contributing to a sum of plaquettes, choose
one link it differentiates.  At most four spatial plaquettes contain that
link.  Thus
$|\sum_p s_p\partial^\beta W_p|^2\le
4\sum_p s_p^2|\partial^\beta W_p|^2$.
Sum the positive derivative weights and use the single-plaquette identity.
This is a pointwise inequality, so averaging under the native law adds
no density ratio or decorrelation assumption.  With one fixed orientation
a link is in at most two marked plaquettes.
\end{proof}

\begin{theorem}[A populated native lower Gram with no ambient-volume loss]
\label{f16v79:plaquette-lower}
Define
\begin{equation}
 \delta_\gamma=\frac{3u_\gamma}{d_\gamma^2},\qquad
 q_\gamma=\frac2{\sqrt{\delta_\gamma^2+4}+\delta_\gamma}.
 \label{f16v79:q-explicit}
\end{equation}
For every finite bank \eqref{f16v79:plaquette-bank} and integer $n\ge1$,
\begin{equation}
 \boxed{\quad
 C_F(n)\succeq q_\gamma^nG_\pi(F)
          \succeq\frac43d_\gamma^2q_\gamma^n\mathcal S,
 \qquad
 0\preceq G_\pi(F)-C_F(1)\preceq2u_\gamma\mathcal S.
 \quad}\label{f16v79:native-main}
\end{equation}
The constants are independent of the surrounding volume, source count,
and coefficient magnitudes.  At every finite delay the kernel is exactly
$\ker\mathcal S$.  The same reflected-window factor as in
Theorem~\ref{f16v79:delay} applies.
\end{theorem}
\begin{proof}
Combine $\mathsf B_\gamma\preceq4u_\gamma\mathcal S$ with
$G_\pi\succeq(4/3)d_\gamma^2\mathcal S$ to obtain
$\mathsf B_\gamma\preceq\delta_\gamma G_\pi$.
Theorem~\ref{f16v79:delay} proves both inequalities.
A zero sampled coefficient combination is identically constant after
native centering; conversely the positive lower matrix excludes any
other null direction.  The constants survive bounded local limits at
fixed $n$ and $\gamma$ without a lower bound on a whole-volume Perron
function.  No invariance of the plaquette span under $P$ was required.
\end{proof}

For fixed large $\gamma$, $u_\gamma\sim4/\gamma$ and
$d_\gamma\sim(1152000\gamma^2)^{-1}$.  Thus
$q_\gamma\asymp\gamma^{-3}$, with explicitly small sufficient constants.
This replaces V78's extensive $e^{-O(\gamma M)n}$ factor for this bank
by a factor depending only on $\gamma$ and delay.  It does not make that
factor useful at fixed physical delay $n\asymp a^{-1}$: the displayed
lower bound still tends to zero.  That is a limitation of the bound,
not an upper estimate or a collapse theorem for the native plaquette bank.

\subsection{Checks against false improvements and the collective target}

In the exactly soluble kinetic reference $b=1$, $\phi=1$, a coordinate
source on one Haar link has $G=1/4$, $C_1=\lambda_1(\gamma)/4$, and
$\mathsf B_\gamma=1/\gamma$.  The scalar lower quotient is
$q(4/\gamma)$; it has the same leading $\gamma/4$ behavior as
$\lambda_1(\gamma)$ when $\gamma\downarrow0$.  For a four-link plaquette
in that reference,
\[
 G=\tfrac14,\quad C_1=\tfrac14\lambda_1(\gamma)^4,\quad
 C_1+\mathsf R_\gamma=\tfrac14(\lambda_1(\gamma)+4/\gamma)^4.
\]
These are exact reference tests of the full normal-ordering budget, not
substitutes for an interacting covariance.  The same feature proof works
with arbitrary positive magnetic weights; neither reference is its premise.

V77's $F^K=\mathcal A_{K,S}(O\circ\Phi_{a,\tau})$ is a finite invariant
polynomial, so Theorem~\ref{f16v79:block} applies to that \emph{actual}
collective bank.  Its derivative matrix can be formed from its full
coefficients using \eqref{f16v79:coefficient-budget}, or evaluated under
the original slice law.  This is not permission to apply the simple
plaquette formula \eqref{f16v79:plaquette-budget} to a flowed plaquette:
its dependence on original links is different.  Nor does a small reflected
reconstruction error bound its ambient derivative cost.  Increasing $K$
alone can increase that cost and is not a new lower-Gram argument.

For a scalar illustration of the distinction, in the Haar kinetic
reference a high-character source can have an arbitrarily small multiplier
while remaining a smooth polynomial with nonzero variance.  No coefficient-
independent $O(a)$ derivative budget follows from polynomial availability.
Conversely spectator links that are absent from a polynomial have zero
ambient derivatives and introduce no direct factor into
\eqref{f16v79:B}.  The finite coefficient bound depends on its actual
support, and the fully populated plaquette result avoids even a growing
support penalty by retaining its local overlap structure.

The new physical target is quantitative: for a declared collective
reconstruction and physical normalization, control its actual
$\mathsf B_{\gamma_a}$ relative to $aG_a$, or a sharper source-specific
version of the exact block Gram.  The vacuum overlap is no longer an
unestimated whole-volume extremum; it cancels exactly in
\eqref{f16v79:Perron}.  This does not prove the required relative
source estimate or a nonzero limit of $G_a$.  Even if both were available,
low-energy participation would still not exclude weight arbitrarily
near zero.  A physical mass gap requires a separate upper contraction
estimate with common dense-core constants.

\subsection{Verification, preservation and result ledger}

The companion diagnostic is
\begin{center}\small
\path{verify_ym_f16_v79_vacuum_cancellation.py}.
\end{center}
The tests distinguish exact operator and quaternion algebra, literal
lattice colors and private links, finite positive kernels with magnetic
weights, and coefficient estimates.  Finite group and quadrature examples
test the operator identity; they are not Wilson vacuum samples.  The analytic arguments do not use their
numerical outcomes.  They are not interval certificates or proof-assistant
verification.  Removing this delimited appendix and reversing the one
displayed version change recovers V78 byte for byte.

\begin{record}[V79 result ledger]
\label{f16v79:ledger}
\emph{Proved here:} an exact coherent-vector/native-Perron cancellation;
the full polynomial block Gram with every uncentered derivative order;
source-directed one-step deficits and delay lower bounds without vacuum
extrema; a finite coefficient certificate on the actual source support;
a two-stage native entrance lower Gram for arbitrary spatial plaquette
sums; their exact all-order derivative budget and volume-/mark-count-
uniform reflected lower bound; and an explicit sufficient physical-energy
nonescape criterion, without claiming its hypotheses for the flowed bank.

\emph{Inherited or established background:} the compact Wilson magnetic
sandwich, positive extension and native reflection; coherent-vector and
normal-ordering calculus; V58's one-link covariance, V59's twelve-color
geometry, V65's polynomial norm estimate, and V77's full-source
reconstruction.  Their new uses above retain the original normalization;
no earlier archive is independently recertified.

\emph{Not asserted:} a Gaussian Wilson measure; centered derivative costs;
a Gauss projection on the noninvariant intermediate derivatives; a
matrix-power inference from a noncommuting lower bound; the elementary
plaquette derivative budget for a flowed or arbitrarily reconstructed
observable; ordinary derivative convergence from V77's reflected error;
or cutoff-uniform physical noncollapse from the displayed lattice bound.

\emph{Still unproved:} a nonzero physically normalized reflected limit of
the simultaneous-flow bank; its required collective derivative or sharper
weighted-source control; thermodynamic/continuum compatibility for that
bank; common dense-core physical-time upper contraction; and the
interacting four-dimensional continuum Yang--Mills mass gap.
\end{record}

\begin{firewall}
The native vacuum weight cancels by its exact Perron equation, not by
assuming typicality or replacing it by Haar measure.  The auxiliary Fock
space represents the original compact kernel and has no new physical
interpretation.  All derivative orders, local overlaps, source kernels,
and physical-time costs remain explicit.  V1--V78 are preserved unchanged.
\end{firewall}
% END F16 V79 APPEND

% BEGIN F16 V80 APPEND -- 2026-09-19
% Append-only continuation of the Post-Fifteenth-Archive V79.
\clearpage
\section[V80: sphere-constraint reduction and intrinsic source costs]{V80: sphere-constraint reduction\\and intrinsic source costs}
\label{f16v80:context}

V79 cancels the native vacuum weight exactly, but its sufficient matrix
$\mathsf B_\gamma$ uses derivatives of a polynomial away from the product
of unit-quaternion spheres.  That continuation is not determined by the
physical source.  Before asking a reconstructed bank to satisfy
$\mathsf B_{\gamma_a}=O(aG_a)$, we separate removable normal derivatives
from the intrinsic derivatives that every continuation must retain.
The source, its centering, its physical transfer and its native law are
unchanged throughout this operation.

There are two distinct conclusions.  A finite sphere-constraint reduction
monotonically improves the native block-Gram certificate.  An explicit
second-order radially constant extension then removes the leading normal
cost for a plaquette and improves the populated large-coupling bound.
However, mixed tangential second derivatives impose a positive
$\gamma^{-2}$ floor for \emph{every} polynomial extension of a plaquette
bank.  Combined with the already proved native defect moments, that floor
excludes V79's sufficient $O(aG)$ test for a fixed microscopic bank on the
stated exponential clock.  This is an obstruction to that test, not a
proof of ultraviolet escape or a no-go theorem for a continuum field.
The growing simultaneous-flow bank remains a different source.

\subsection{Inputs and the difference between a source and an extension}

Use exactly \eqref{f16v79:setup}--\eqref{f16v79:R}, on
$\mathcal M=(S^3)^M$ with its unscaled product round metric.  In particular,
$P$ is the original gauge-invariant native transfer and $\mathcal P$ its
specified positive extension.  No Gauss projection is inserted into the
noninvariant derivative pairings.  For a real polynomial extension $p$,
write the pointwise cost
\begin{equation}
 \mathcal B_\gamma[p](U)=
  \sum_{|\beta|\ge1}\frac{\gamma^{-|\beta|}}{\beta!}
           |\partial^\beta p(U)|^2,
 \qquad
 \mathsf B_\gamma(p,q)=
  \mathbb E_\pi\sum_{|\beta|\ge1}
  \frac{\gamma^{-|\beta|}}{\beta!}\partial^\beta p\,\partial^\beta q.
 \label{f16v80:cost}
\end{equation}
Vector costs are obtained by polarization.  Constants have zero cost,
but derivatives of a nonconstant source remain \emph{uncentered}.
The matrix $\mathsf R_\gamma$ has the same definition with the uncentered
$\mathcal P$ pairing in place of the equal-time expectation.

The relation between an ambient Hessian and a submanifold Hessian is
established differential geometry; see P.-A.~Absil, R.~Mahony and
J.~Trumpf, \emph{An extrinsic look at the Riemannian Hessian},
\href{https://doi.org/10.1007/978-3-642-40020-9_39}{GSI 2013,
pp.~361--368}.  The product-sphere formulas, compact plaquette constants,
and native implications needed here are proved below.  V79's
coherent-vector block Gram, V69's exact plaquette Hessian and V76's
reflection-compatible defect moments are inherited inputs, not new
claims.  No previously proved derivative cost is silently redefined as
an intrinsic cost.

\subsection{A finite reduction by exact sphere constraints}

Let $p=(p_1,\ldots,p_d)$ extend the physical bank $F$, and let
$q=(q_1,\ldots,q_m)$ be any finite list of real polynomials whose
restrictions to $\mathcal M$ vanish identically.  An explicit supply is
\begin{equation}
 q_{e,\alpha}(x)=(|x_e|^2-1)x^\alpha.
 \label{f16v80:null-polynomials}
\end{equation}
Adding linear combinations of these polynomials changes neither the
source nor any of its native covariances.  Set $G=G_\pi(F)$,
$C_1=C_F(1)$, and form the derivative Gram blocks $B_{pp},B_{pq},B_{qq}$
using \eqref{f16v80:cost}; use $R$ for their transfer-paired counterparts.

\begin{theorem}[Constraint reduction of the native source certificate]
\label{f16v80:short}
Define, with Moore--Penrose inverses on the exact ranges,
\begin{equation}
 B_{p\mid q}=B_{pp}-B_{pq}B_{qq}^{\dagger}B_{qp},
 \qquad
 R_{p\mid q}=R_{pp}-R_{pq}R_{qq}^{\dagger}R_{qp}.
 \label{f16v80:Schur}
\end{equation}
Then
\begin{equation}
 0\preceq R_{p\mid q}\preceq B_{p\mid q},\qquad
 \boxed{\quad
 \begin{pmatrix}C_1&G\\G&C_1+R_{p\mid q}\end{pmatrix}\succeq0,
 \qquad
 \begin{pmatrix}C_1&G\\G&C_1+B_{p\mid q}\end{pmatrix}\succeq0.
 \quad}
 \label{f16v80:short-block}
\end{equation}
Consequently V79's deficit and delay conclusions hold with either
reduced cost in place of the corresponding unreduced cost.
Enlarging the span of $q$ can only decrease either reduced matrix.
For the equal-time cost, $B_{p\mid q}$ is exactly the cost of the
representative family
\begin{equation}
 p_i^{\rm opt}=p_i-\sum_\alpha q_\alpha
                   (B_{qq}^{\dagger}B_{qp})_{\alpha i}.
 \label{f16v80:optimal-extension}
\end{equation}
No source direction is removed by this optimization.
\end{theorem}
\begin{proof}
Apply V79's exact feature construction to $(p,q)$.  The source vectors
associated with $q$ vanish.  Its creation-operator vectors need not
vanish, but are orthogonal to every source vector: their overlap is the
native Gram of a physical source with the zero restriction.  The Gram
of the remaining vectors is therefore
\[
 \begin{pmatrix}
 C_1&G&0\\
 G&C_1+R_{pp}&R_{pq}\\
 0&R_{qp}&R_{qq}
 \end{pmatrix}\succeq0.
\]
Projecting the middle family orthogonally off the last family gives the
first block inequality in \eqref{f16v80:short-block}.  This proof also
pays the range condition when $R_{qq}$ is singular.  Equivalently it is
the Schur complement on that exact range.  Replacing the bottom-right
derivative Gram of the full $(p,q)$ family by its larger $B$ Gram gives
the second inequality by the same argument.

For every coefficient vector $c$, the Schur quadratic form is
\[
 c^*B_{p\mid q}c=\min_v\mathsf B_\gamma(c^*p+v^*q,c^*p+v^*q),
\]
with an analogous formula for $R$.  These nonnegative minima exist in
the finite-dimensional quotient; zero-cost directions have zero cross
pairings.  Since $R\preceq B$ on the full polynomial list, the minima
are ordered.  The minimizing linear map gives
\eqref{f16v80:optimal-extension}.  A larger null space permits more
competitors, proving monotonicity.  All competitors restrict to the
same physical function, so centering, its exact source kernel and its
native correlations remain unchanged.  Testing the reduced block on
$(c,-c)$, or using its scalar determinant as in V79, proves the deficit
and delay assertions.
\end{proof}

A hierarchy can use the null subspace of polynomials of total degree at
most $D$, with $D$ increasing.  These are finite, algebraically specified
spaces.  Their union is the ideal generated by $|x_e|^2-1$:
separate harmonic decomposition writes a polynomial as a finite sum of
products of homogeneous harmonics times polynomials in the radii
$|x_e|^2$.  If its sphere restriction is zero, orthogonality of the
product harmonic basis makes each radial coefficient vanish at
$(1,\ldots,1)$.  Telescoping that coefficient about this point factors
it into a sum of $|x_e|^2-1$ multiples.  This proves the assertion without
assuming a numerical rank threshold.

The hierarchy is finite at every stage, but its native moment matrices
have not been evaluated for the full collective bank.  No convergence
rate or equality of its limit with the exact physical deficit is claimed.
The optimizing matrices for $R$ and $B$ can differ.  A lower obstruction
for $B$ below must not be assigned to $R$.

\subsection{A canonical extension with two radial derivatives removed}

For a homogeneous harmonic polynomial $H_j$ on $\mathbb R^4$, set
\begin{equation}
 (\mathcal E_2 H_j)(x)=H_j(x)
 \left[1-\frac j2(|x|^2-1)
       +\frac{j(j+2)}8(|x|^2-1)^2\right].
 \label{f16v80:E2}
\end{equation}
Use the separate harmonic expansion in every link, apply this formula
to each factor, and extend linearly.  This defines $\mathcal E_2F$ from
the sphere restriction itself.  It respects exact zero relations,
constants and orthogonal changes of link coordinates.  In particular
it preserves the vertex-gauge invariance of the physical bank.  Each
separate degree increases by at most four, not by a new spectator-link
factor.

Write $\Gamma_{ij}=\langle\nabla_{\mathcal M}F_i,
\nabla_{\mathcal M}F_j\rangle$ and
$\mathcal H_{ij}=\langle\operatorname{Hess}_{\mathcal M}F_i,
\operatorname{Hess}_{\mathcal M}F_j\rangle_{\rm HS}$.

\begin{proposition}[Exact intrinsic content of the first two orders]
\label{f16v80:jets}
On $\mathcal M$, the first and second radial derivatives of
$\mathcal E_2F$ vanish separately at every link, and
\begin{equation}
 \boxed{\quad
 \mathsf B_\gamma(\mathcal E_2F)=
 (\gamma^{-1}+\gamma^{-2})\mathbb E_\pi\Gamma
 +\frac1{2\gamma^2}\mathbb E_\pi\mathcal H
 +\mathsf B^{(\ge3)}_\gamma(\mathcal E_2F),
 \qquad \mathsf B^{(\ge3)}_\gamma\succeq0.
 \quad}
 \label{f16v80:intrinsic-two-jet}
\end{equation}
The last term is the full finite sum of derivative orders at least three;
it is not discarded or asserted small for a growing-degree source.
\end{proposition}
\begin{proof}
Along a radial line, $H_j(rx)=r^jH_j(x)$.  The bracket in
\eqref{f16v80:E2} is the quadratic Taylor polynomial of
$s^{-j/2}$ at $s=1$, with $s=r^2$.  The product is
$1+O((r-1)^3)$, proving the two radial cancellations.  They hold as
identities in the other link coordinates.

The ambient gradient is consequently tangent and equals the intrinsic
gradient.  In an orthonormal frame consisting of the unit normal $x_e$
and three tangents at link $e$, the same-link ambient Hessian has blocks
\[
 D^2_{ee}(\mathcal E_2F)=
 \begin{pmatrix}0&-(\nabla_eF)^*\\
                 -\nabla_eF&\operatorname{Hess}_{ee}F\end{pmatrix}.
\]
For example, differentiating $x_e\cdot\nabla_e(\mathcal E_2F)=0$
in a tangential direction gives the displayed mixed normal block;
the radial-radial block is zero by construction.  Mixed Hessians at
distinct links have only tangent-tangent blocks, since each first
radial derivative vanishes identically along the other sphere factors.
Thus, by polarization, the ambient gradient Gram is $\Gamma$ and the
ambient Hessian Gram is $\mathcal H+2\Gamma$.  The multiindex sum at
order two is one half the full ordered Hessian contraction.
Separating those first two orders in \eqref{f16v80:cost} proves the
formula, including its positive higher-order remainder.
\end{proof}

This operation does not replace ambient differentiation by tangent
differentiation inside V79's proof.  It chooses a different, admissible
polynomial and computes its complete ambient cost.  Simply keeping
$\gamma^{-1}\Gamma$ would not be justified by
\eqref{f16v80:intrinsic-two-jet}.

\subsection{All orders of a radially flattened plaquette are explicit}

For one degree-one link factor, \eqref{f16v80:E2} amounts to the map
\begin{equation}
 \mathcal R_2(x)=x\left[1-\frac12(|x|^2-1)
                           +\frac38(|x|^2-1)^2\right].
 \label{f16v80:radial-map}
\end{equation}
For a literal plaquette $W=w(x_1x_2x_3^{-1}x_4^{-1})$, its multilinear
ambient representative uses quaternion conjugation for inverse
occurrences.  Replace every factor by $\mathcal R_2(x_e)$, keeping those
conjugations.  Call the resulting polynomial $\widetilde W^{[2]}$.
It equals $W$ on the spheres and has degree at most five at each of its
four links.

Put $t=\gamma^{-1}$ and define
\begin{equation}
 A(t)=t+t^2+\frac92t^3+36t^4+54t^5,\qquad
 B(t)=1-t+\frac12t^2+39t^3+126t^4.
 \label{f16v80:AB}
\end{equation}
These $A,B$ are scalar polynomials, not the native Gram matrices.

\begin{theorem}[Exact compact plaquette budget after radial correction]
\label{f16v80:exact-budget}
At every compact link configuration, with $y=w(U_P)^2\in[0,1]$,
\begin{equation}
 \boxed{\quad
 \mathcal B_\gamma[\widetilde W^{[2]}]
 =\Pi(t,y):=(B(t)^4-1)y+
                 \frac{(B(t)+4A(t))^4-B(t)^4}{4}.
 \quad}
 \label{f16v80:Pi}
\end{equation}
In particular
\begin{equation}
 \Pi(t,y)=4t(1-y)+(16+8y)t^2+O(t^3),
 \label{f16v80:Pi-leading}
\end{equation}
with an exact polynomial remainder of degree at most twenty.  A uniform
finite bound, valid for $\gamma\ge16$, is
\begin{equation}
 0\le\Pi(\gamma^{-1},y)
       -4\gamma^{-1}(1-y)\le45\gamma^{-2}.
 \label{f16v80:Pi-bound}
\end{equation}
\end{theorem}
\begin{proof}
For a unit vector $x$, rotational equivariance gives the matrix identity
\[
 \sum_\beta\frac{t^{|\beta|}}{\beta!}
 (\partial^\beta\mathcal R_2(x))
 (\partial^\beta\mathcal R_2(x))^*
       =A(t)I_4+B(t)xx^*.
\]
Here the zero derivative is included.  To evaluate every coefficient,
take $x=e_0$, differentiate the displayed degree-five polynomial, and
contract the coordinate derivatives.  For a unit longitudinal or
transverse test vector the respective coefficients are
\begin{center}\small
\begin{tabular}{@{}c|rrrrrr@{}}
\toprule
order $k$&0&1&2&3&4&5\\
\midrule
longitudinal&1&0&$3/2$&$87/2$&162&54\\
transverse&0&1&1&$9/2$&36&54\\
\bottomrule
\end{tabular}
\end{center}
The weights are $1/\beta!$, equivalently $1/k!$ for ordered derivatives.
For instance the second derivative of $a\cdot\mathcal R_2$ at a unit $x$
is
\[
 -a x^*-x a^*-(a\cdot x)I+3(a\cdot x)xx^*,
\]
whose weighted squared norm is $|a|^2+\tfrac12(a\cdot x)^2$.
The remaining two scalar entries at each order follow by differentiating
$(15-10|x|^2+3|x|^4)(a\cdot x)/8$; all higher derivatives vanish.
This proves \eqref{f16v80:AB}.

Contract the four link matrices with the literal quaternion
plaquette tensor twice.  Choosing the $Bxx^*$ term at every link gives
$B^4y$.  Choosing $A I$ at $k\ge1$ of the four links gives
$A^kB^{4-k}4^{k-1}$, by the exact quaternion derivative contraction
of V79.  Sum its $\binom4k$ choices and subtract the zero-order value
$y$.  This proves \eqref{f16v80:Pi}, including reversed orientations.
Its first two coefficients give \eqref{f16v80:Pi-leading}.

Every coefficient at positive derivative order is a sum of squares;
in particular the endpoint polynomials $\Pi(t,0)-4t$ and $\Pi(t,1)$
have nonnegative coefficients and start at order two.  Their quotients
by $t^2$ increase on $t\ge0$.  At $t=1/16$ their exact values are,
respectively,
\[
 \frac{124540631143913751329}{2^{62}}<28,
 \qquad
 \frac{204334005218988033825}{2^{62}}<45.
\]
Since the remainder is affine in $y$, these evaluations give
\eqref{f16v80:Pi-bound} for $0\le t\le1/16$.  This is a finite
polynomial inequality, not a numerical quadrature premise.
\end{proof}

\subsection{An improved populated native bound with local overlap retained}

Return to V79's arbitrary spatial plaquette bank, with distinct marks,
\begin{equation}
 F_i=\sum_{p\in\mathcal P}s_{ip}w(U_p),\qquad
 \mathcal S=\sum_p s_p s_p^*,\qquad
 d_\gamma=[32000(1+6\gamma)^2]^{-1},\quad
 \kappa_\gamma=\tfrac43d_\gamma^2.
 \label{f16v80:bank}
\end{equation}
Let $\chi$ be an upper bound on the number of marked plaquettes through
one link: $\chi=4$ always works, and $\chi=2$ works for one orientation.
Use the linear sum of the extensions $\widetilde W_p^{[2]}$.

\begin{theorem}[Intrinsic first order and a paid all-order native remainder]
\label{f16v80:native-upper}
For $\gamma\ge16$,
\begin{equation}
 \boxed{\quad
 \mathsf B^{[2]}_\gamma(F)
 \preceq\gamma^{-1}\mathbb E_\pi\Gamma(F)
                      +45\chi\gamma^{-2}\mathcal S.
 \quad}
 \label{f16v80:collective-upper}
\end{equation}
In the reflection-compatible periodic laws and limits of V76, put
\begin{equation}
 b_\gamma=384\left(8+\log\gamma
                      -\tfrac23\log\frac8{3\pi^3}\right),
 \qquad v_\gamma=\frac{\chi(8b_\gamma+53)}{\gamma^2}.
 \label{f16v80:native-constants}
\end{equation}
Then $\mathsf B^{[2]}_\gamma\preceq v_\gamma\mathcal S$.  Consequently
\begin{equation}
 \boxed{\quad
 C_F(n)\succeq q(v_\gamma/\kappa_\gamma)^nG_\pi(F)
       \succeq\kappa_\gamma q(v_\gamma/\kappa_\gamma)^n\mathcal S,
 \qquad
 0\preceq G_\pi(F)-C_F(1)\preceq\tfrac12v_\gamma\mathcal S.
 \quad}
 \label{f16v80:native-lower}
\end{equation}
Here $q(\delta)=2/(\sqrt{\delta^2+4}+\delta)$ is V79's exact scalar
factor.  Constants do not depend on the number of marks or the ambient
volume.  The exact finite-delay kernel remains $\ker\mathcal S$.
\end{theorem}
\begin{proof}
The order-one derivatives of the corrected representative are the actual
intrinsic derivatives of the \emph{complete sum}, so retain their Gram
$\Gamma(F)$ without splitting it.  At any higher nonzero multiindex,
choose one differentiated link.  At most $\chi$ marked plaquettes can
contribute.  Cauchy--Schwarz and \eqref{f16v80:Pi-bound}, applied to the
remaining orders of each individual plaquette, give
$45\chi\gamma^{-2}\mathcal S$.  This proves
\eqref{f16v80:collective-upper} pointwise before its native expectation.

For one plaquette and one boundary link,
$|\nabla_e w(U_p)|^2=1-w(U_p)^2$.  The same incidence argument gives
\[
 \Gamma(F)\preceq4\chi\sum_p(1-w(U_p)^2)s_ps_p^*.
\]
Use $1-w^2\le2(1-w)$ and the inherited native moment bound
$\mathbb E(1-w(U_p))\le(b_\gamma+1)/\gamma$.
Thus $\gamma^{-1}\mathbb E\Gamma\preceq
8\chi(b_\gamma+1)\gamma^{-2}\mathcal S$, which yields $v_\gamma$.
V79 supplies $G_\pi(F)\succeq\kappa_\gamma\mathcal S$ and its
block-Gram delay inequality applies to the new extension.  Those two
inputs prove \eqref{f16v80:native-lower} in every source direction.
A sampled zero combination is the same zero physical function before
and after extension; the positive lower bound excludes all other
finite-delay null directions.
\end{proof}

The displayed factor is asymptotic, up to explicit constants, to
$(\gamma^2\log\gamma)^{-1}$ instead of V79's $\gamma^{-3}$.
The large $b_\gamma$ is not a practical moderate-coupling prediction;
the new numerical lower factor is not asserted larger at every coupling.
Both bounds remain valid, so one can retain their better scalar bound.
Equivalently the constraint reduction can include the difference of the
two extensions, and is then no worse than either native derivative
matrix.  None of these lower factors is uniform at physical delay
$n\asymp a^{-1}$.

\subsection{An unavoidable mixed-derivative floor}

The radial correction removes an artificial first-order cost, but cannot
remove the intrinsic mixed response of the plaquette.  The next statement
applies to \emph{every} polynomial extension, including ones of
arbitrarily high degree and regulator-dependent coefficients.

\begin{theorem}[Extension-independent floor for a plaquette bank]
\label{f16v80:floor}
For the bank \eqref{f16v80:bank} on a nondegenerate spatial torus, every
ambient representative family $\widetilde F$ satisfies pointwise
\begin{equation}
 \boxed{\quad
 \sum_{|\beta|\ge1}\frac{\gamma^{-|\beta|}}{\beta!}
    \partial^\beta\widetilde F
    (\partial^\beta\widetilde F)^*
 \succeq\frac6{\gamma^2}\sum_p(2+w(U_p)^2)s_ps_p^*
 \succeq\frac{12}{\gamma^2}\mathcal S.
 \quad}
 \label{f16v80:floor-matrix}
\end{equation}
It follows that the same floor holds for $\mathsf B_\gamma$ and every
finite constraint reduction $B_{p\mid q}$.  No such lower bound for
$\mathsf R_\gamma$ is asserted.
\end{theorem}
\begin{proof}
For two distinct link variables, the ambient second derivative projected
onto their two tangent spaces is determined entirely by the physical
restriction.  Indeed differentiating along separate link geodesics has
no mixed acceleration.  A polynomial vanishing on the product of spheres
therefore has zero such mixed tangent block.

For two distinct edges of a plaquette, V69's exact intrinsic Hessian is,
up to left and right orthogonal coordinate changes and a sign,
$wI_3-C_p$, where $p=\operatorname{Im}U_p$ and $C_pv=p\times v$.
Its squared Frobenius norm is exactly
\[
 \|wI_3-C_p\|_{\rm HS}^2=3w^2+2|p|^2=2+w^2.
\]
Two distinct elementary spatial plaquettes cannot contain the same pair
of distinct boundary links under the stated lattice nondegeneracy.
Hence for each of the six unordered pairs in plaquette $p$, the mixed
block of the complete sampled source is precisely $s_p$ times this
block; no other plaquette coefficient enters it.  The order-two ambient
cost for that pair has weight $\gamma^{-2}$ and dominates the squared
norm of its tangent projection.  The six pair-index sets for different
plaquettes are disjoint.  Summing proves \eqref{f16v80:floor-matrix};
all unused orders and coordinate components are nonnegative.

Integrate under the original native law.  Each candidate in the
minimization defining $B_{p\mid q}$ is an extension of the same physical
bank, so the same lower matrix survives the minimum.  The argument uses
squared equal-time derivatives.  Inserting $\mathcal P$ can suppress
those derivatives; it would be invalid to transfer this floor to $R$.
\end{proof}

\begin{proposition}[The flat plaquette's optimal leading cost is nonzero]
\label{f16v80:flat-sharp}
For a single plaquette at the all-identity links and any ambient
polynomial extension $p$ of its trace,
\begin{equation}
 \mathcal B_\gamma[p](\mathbf1)
 \ge\frac{18}{\gamma^2}+\frac{12}{\gamma(2\gamma+3)}.
 \label{f16v80:flat-floor}
\end{equation}
Moreover
\begin{equation}
 \boxed{\quad
 \lim_{\gamma\to\infty}\gamma^2
       \inf_{p|_{\mathcal M}=W}\mathcal B_\gamma[p](\mathbf1)=24.
 \quad}
 \label{f16v80:flat-optimum}
\end{equation}
The infimum permits different finite-degree extensions at each $\gamma$.
\end{proposition}
\begin{proof}
The intrinsic first derivative is zero at identity.  Write $r_e$ for
any extension's radial first derivative at each of the four links.
Its first-order cost is at least $\gamma^{-1}\sum_e r_e^2$.
The intrinsic same-link Hessian is $-I_3$, so its ambient
 tangent-tangent block is $(r_e-1)I_3$.  This follows directly by
differentiating along a unit sphere geodesic, whose ambient acceleration
is $-|v_e|^2x_e$.  The six mixed blocks contribute $18/\gamma^2$ as
above.  Thus
\[
 \mathcal B_\gamma[p](\mathbf1)\ge
 \frac{18}{\gamma^2}+\sum_{e=1}^4
 \left[\frac{r_e^2}{\gamma}
               +\frac{3(r_e-1)^2}{2\gamma^2}\right].
\]
Minimizing each explicit quadratic gives $3/[\gamma(2\gamma+3)]$.
Other derivatives, including derivatives in spectator links, cannot
subtract this cost.  Conversely \eqref{f16v80:Pi-leading} at $y=1$
gives $\mathcal B_\gamma[\widetilde W^{[2]}](\mathbf1)
=24\gamma^{-2}+O(\gamma^{-3})$.  The lower and upper limits prove
\eqref{f16v80:flat-optimum}.
\end{proof}

In particular a zero intrinsic first gradient at a coherent configuration
does not make the full source budget zero.  The $\gamma^{-2}$ term is
not a removable normal-coordinate artifact.  Conversely the original
multilinear representative had cost $4/\gamma+O(\gamma^{-2})$ there;
that leading normal cost really was removable.

\subsection{A native obstruction to the fixed-bank physical test}

\begin{theorem}[Fixed plaquette banks cannot pass the derivative test]
\label{f16v80:fixed-bank}
In the reflection-compatible native laws above, let
$N=|\mathcal P|<\infty$ and
$D_\gamma=b_\gamma^2+2b_\gamma+2$.  Every ambient extension and every finite
constraint reduction of $B$ satisfy
\begin{equation}
 \boxed{\quad
 \mathsf B_\gamma\succeq
             \frac{12}{N D_\gamma}\,G_\pi(F).
 \quad}
 \label{f16v80:relative-floor}
\end{equation}
Consequently, for nonzero source quotient, the sufficient condition
$\mathsf B_{\gamma_a}\preceq2aE_*G_a$ with a fixed $0<E_*<\infty$ is
impossible whenever
\begin{equation}
          aN_aD_{\gamma_a}\longrightarrow0.
 \label{f16v80:test-exclusion}
\end{equation}
In particular this includes fixed $N$, $\gamma_a\to\infty$ and
$\log(1/a)=c\gamma_a+O(\log\gamma_a)$ with $c>0$.  Arbitrary declared
coefficient rescalings or changes of basis cannot remove the obstruction.
\end{theorem}
\begin{proof}
For a scalar combination use the defects $d_p=1-w(U_p)$ and subtract the
constant sum of its coefficients.  Cauchy--Schwarz and V76's native second
moment give
\[
 \operatorname{Var}_\pi\sum_p s_pw(U_p)
 \le\mathbb E_\pi\left|\sum_p s_pd_p\right|^2
 \le N\frac{D_\gamma}{\gamma^2}\sum_p s_p^2.
\]
Thus $G_\pi(F)\preceq N D_\gamma\gamma^{-2}\mathcal S$.
Combine with \eqref{f16v80:floor-matrix}.  If the stated upper test also
held, any nonzero source direction would require
$aND_\gamma\ge6/E_*$.  This contradicts
\eqref{f16v80:test-exclusion}.  The argument applies in every coefficient
direction, including changing coefficients, and is invariant under
congruence on the exact quotient.  Finally $D_\gamma=O((\log\gamma)^2)$
as $\gamma\to\infty$ in the declared native bound; the exponential
trajectory satisfies \eqref{f16v80:test-exclusion} for fixed $N$.
\end{proof}

This theorem excludes a \emph{sufficient derivative criterion}, not the
source's actual finite-energy participation.  It supplies no upper bound
on $C_F(n)$ and no collapse theorem.  In particular V67's distinction
between contact normalization and a canonical reflected normalization
remains in force.  For a growing physically smeared source,
$N_a$ may be large enough that \eqref{f16v80:test-exclusion} fails;
no conclusion is inferred then.  The exact transfer-paired cost $R$ is
also not subject to the proved equal-time floor.

For V77's collective simultaneous-flow polynomial, the constraint
reduction and \eqref{f16v80:intrinsic-two-jet} apply without replacing
its coefficients by elementary plaquettes.  They retain its exact
source and shared-path contributions.  However, its intrinsic Hessian,
higher derivatives and native expectations still have to be estimated.
A small reflected reconstruction error does not bound any of those
higher costs.  The next useful source calculation is the
constraint-reduced \emph{transfer-paired} matrix for that declared
collective bank, or a direct physical-separation Gram estimate.  The
fixed-microscopic $B=O(aG)$ target should not be pursued by changing
ambient representatives: it is excluded under the stated scaling,
even though the unnecessary leading normal cost has been removed.

\subsection{Verification, preservation and result ledger}

The companion diagnostic is
\begin{center}\small
\path{verify_ym_f16_v80_sphere_constraints.py}.
\end{center}
It separates exact rational radial-jet and polynomial identities,
literal plaquette pair incidence, compact quaternion Hessian tests,
and finite positive-kernel magnetic-sandwich tests of the two Schur
reductions.  The latter are explicitly auxiliary models, not native
four-dimensional Wilson samples.  The analytic proofs do not depend
on their numerical values.  No interval certificate or proof-assistant
verification is claimed.  Removing the delimited appendix and reverting
the single displayed version change recovers V79 byte for byte.

\begin{record}[V80 result ledger]
\label{f16v80:ledger}
\emph{Proved here:} exact finite sphere-constraint reductions of both
native block-Gram costs; monotonicity and unchanged source relations;
a canonical second-order radially constant polynomial extension and its
intrinsic two-jet identity; the complete degree-twenty corrected
plaquette budget; a native local-incidence improvement to the populated
plaquette bound; a representative-independent mixed-derivative floor;
the sharp leading flat-plaquette extension cost; and an explicit native
obstruction to the fixed-bank $B=O(aG)$ criterion.

\emph{Inherited or established background:} V79's exact Perron/coherent
block Gram, positive transfer extension and native entrance bound;
V69's compact plaquette Hessian; V76's reflection-compatible moments;
ordinary harmonic decomposition, submanifold Hessian identities and
finite Gram projection.  These tools are not claimed as new theories.

\emph{Not asserted:} suppression of all higher derivatives by radial
flattening; an evaluated optimal collective moment matrix; convergence
of derivative budgets from reflected polynomial approximation; the
same obstruction for $R$ as for $B$; absence of physical fields from
failure of this sufficient test; a collapse theorem for a growing
spatially flowed bank; or a continuum mass gap from any lower correlation
bound.

\emph{Still unproved:} a nonzero physically normalized reflected limit
for the actual simultaneous-flow source; suitable collective
constraint-reduced or other source-directed control; thermodynamic and
continuum compatibility; common dense-core physical-time upper
contraction; and the interacting four-dimensional continuum
Yang--Mills mass gap.
\end{record}

\begin{firewall}
The physical source is unchanged when its off-sphere continuation is
optimized.  Removable radial cost is separated from unavoidable mixed
intrinsic cost.  The native moment input and physical clock are explicit;
failure of a sufficient derivative criterion is not failure of the
underlying theory.  All V1--V79 mathematical text remains unchanged.
\end{firewall}
% END F16 V80 APPEND

% BEGIN F16 V81 APPEND -- 2026-09-19
% Append-only continuation of the Post-Fifteenth-Archive V80.
\clearpage
\section[V81: transfer-paired floors and the complete feature quotient]{V81: transfer-paired floors\\and the complete feature quotient}
\label{f16v81:context}

V80 excludes a fixed microscopic plaquette bank from the equal-time
$B=O(aG)$ test, but expressly does not assign that floor to the
transfer-paired cost $R$.  We examine the latter cost directly, before
building a larger table of polynomial constraints.  There are two
independent losses: an intrinsic inverse-transfer moment and a residual
in the auxiliary coherent feature space.  A complete projection in that
space removes the second loss exactly.  In contrast, even the union of
all polynomial sphere constraints need not remove it.  A compact
four-link example proves this strict distinction by evaluation at a
complex zero of the vacuum feature function.

A second, native calculation applies the two-stage conditional variance
argument of V79 inside a time separator.  It gives a floor for the actual
innovation, and hence for every transfer-paired constraint cost, including
the complete feature minimum.  Thus changing from $B$ to $R$ does not
rescue the fixed microscopic $O(aG)$ test on the stated trajectory.  This
is not a collapse theorem.  Indeed an explicit compact kinetic reference
has a uniformly bounded source bank with a surviving reflected limit and
bounded mean physical energy, but a divergent minimum inverse-transfer
cost.  The last result constructs the minimum witness of a positively
buffered native source using only forward transfer powers.  Its residual
normalization requirement is left visible rather than inferred from
smoothing.

\subsection{The native feature operator and its exact minimum}

All native statements first keep a finite nondegenerate spatial
regulator.  Use exactly V79's $T=bCb$, $T\phi=\rho\phi$, $d\pi=\phi^2dm$,
$h=b\phi$, exponential vectors $\Psi$, and $\eta_q=\int q\Psi\,dm$.
The full positive Markov extension is denoted by $\mathcal P$; on
invariant sources it is the physical $P$.  Define the bounded operator
\begin{equation}
 \mathcal J:L^2(\pi)\longrightarrow\mathcal F,
 \qquad \mathcal Jf=\rho^{-1/2}\eta_{hf},
 \qquad \mathcal J^*\mathcal J=\mathcal P.
 \label{f16v81:J}
\end{equation}
Here $\mathcal F$ is V79's auxiliary symmetric Fock space.  Let
$\mathcal K=\overline{\operatorname{Ran}\mathcal J}$ and let $\Pi$ be its
orthogonal projection.  This is the compact-restriction feature space,
not a Gauss projection.  No noninvariant derivative is projected onto
physical gauge states during an intermediate contraction.

The kinetic convolution has strictly positive representation
multipliers, and $b,\phi$ are positive and bounded above and below at
this fixed regulator.  Hence $\mathcal P$ is injective.  Its polar
factorization is
\begin{equation}
 \mathcal J=\mathcal U\mathcal P^{1/2},
 \qquad \mathcal U:L^2(\pi)\longrightarrow\mathcal K
 \quad\hbox{unitary onto }\mathcal K.
 \label{f16v81:polar}
\end{equation}
Inverse powers below are unbounded spectral operators, with their form
domains stated.  They are not formed by inverting a finite covariance
matrix.

For an invariant polynomial bank, let $f_i=F_i-\mathbb E_\pi F_i$ and
choose centered ambient representatives $p_i$.  V79's creation witnesses
are
\[
 W_{p_i}=\rho^{-1/2}p_i(a^*/\sqrt\gamma)\eta_h,
 \qquad V_i=\mathcal Jf_i.
\]
Its exact Perron identity gives $\mathcal J^*W_{p_i}=f_i$, not merely the
overlap with the chosen finite bank.  Write $G=(\langle f_i,f_j\rangle)$,
$C_n=(\langle f_i,P^nf_j\rangle)$ and define
\begin{equation}
 (C_{-1})_{ij}=
 \langle P^{-1/2}f_i,P^{-1/2}f_j\rangle_\pi,
 \qquad \mathsf Q_-=C_{-1}-C_1.
 \label{f16v81:inverse-Gram}
\end{equation}
These forms are finite for the polynomial bank, as the next proof shows.
The orthogonal-projection and polar-decomposition methods are standard;
the native identification and the distinction from the polynomial
constraint hierarchy are derived here.  For the explicit special-function
example below, the Bessel connection formula and simplicity of positive
zeros are recorded in NIST DLMF,
\href{https://dlmf.nist.gov/10.27.E6}{10.27.6} and
\href{https://dlmf.nist.gov/10.21.i}{10.21(i)}.  No external continuum or
reconstruction theorem is used.

\begin{theorem}[Complete feature minimum and the exact remaining constraint cost]
\label{f16v81:minimum}
Every polynomial $f_i$ belongs to $\operatorname{Dom}P^{-1/2}$, and
\begin{equation}
 \boxed{\quad
 W_i^{\min}:=\Pi W_{p_i}=\mathcal U P^{-1/2}f_i,
 \qquad \mathcal J^*W_i^{\min}=f_i.
 \quad}
 \label{f16v81:minimum-vector}
\end{equation}
It is the unique minimum-norm solution of this last equation.  Its joint
Gram with the original source feature vectors is
\begin{equation}
 \boxed{\quad
 \begin{pmatrix}C_1&G\\G&C_{-1}\end{pmatrix}\succeq0,
 \qquad \mathsf Q_-=C_{-1}-C_1\succeq0.
 \quad}
 \label{f16v81:minimum-block}
\end{equation}
For a finite list $q$ of sphere-null polynomials, put
$\mathcal N_q=\operatorname{span}\{W_{q_\alpha}\}$ and let $\Pi_q$ be its
orthogonal projection.  With $Z_i=(I-\Pi)W_{p_i}$, V80's exact reduction
satisfies
\begin{equation}
 \boxed{\quad
 R_{p\mid q}=\mathsf Q_-+
       \bigl(\langle (I-\Pi_q)Z_i,(I-\Pi_q)Z_j\rangle\bigr)_{ij}.
 \quad}
 \label{f16v81:finite-residual}
\end{equation}
In particular $R_{p\mid q}\succeq\mathsf Q_-$.  Along all increasing
finite sphere-null polynomial spaces, define
$\mathcal N_\infty=\overline{\bigcup_q\mathcal N_q}\subset\mathcal K^\perp$.
Then the decreasing matrices converge to
\begin{equation}
 R_{p\mid\infty}=\mathsf Q_-+
 \bigl(\langle (I-\Pi_{\mathcal N_\infty})Z_i,
                  (I-\Pi_{\mathcal N_\infty})Z_j\rangle\bigr)_{ij}.
 \label{f16v81:infinite-residual}
\end{equation}
Equality with $\mathsf Q_-$ holds exactly when every $Z_i$ lies in
$\mathcal N_\infty$; it is not automatic from algebraic generation of
the sphere ideal.
\end{theorem}
\begin{proof}
The identity $\mathcal J^*W_p=f$ follows by pairing with an arbitrary
$g\in L^2(\pi)$ and using $hCh=\rho\phi^2$, exactly as in V79.  Initially
one can use bounded $g$ and then pass by continuity.  From
$\mathcal J^*=\mathcal P^{1/2}\mathcal U^*\Pi$ we obtain
$f=\mathcal P^{1/2}\mathcal U^*W_p$.  Thus $f$ belongs to the range of
$\mathcal P^{1/2}$, equivalently the domain of $\mathcal P^{-1/2}$.
An invariant $f$ has its inverse in the invariant reducing subspace,
so the inverse there is the physical $P^{-1/2}$.

Every other solution differs from $\mathcal U P^{-1/2}f$ by a vector
in $\ker\mathcal J^*=\mathcal K^\perp$.  Orthogonality proves the
minimum and its uniqueness.  Taking the Gram of
$(\mathcal U P^{1/2}f_i,\mathcal U P^{-1/2}f_i)_i$ gives
\eqref{f16v81:minimum-block}.  Positivity of $\mathsf Q_-$ follows from
$\lambda^{-1}-\lambda\ge0$ for $0<\lambda\le1$.

For sphere-null $q_\alpha$, the same overlap identity gives
$\mathcal J^*W_{q_\alpha}=0$, hence $\mathcal N_q\subset\mathcal K^\perp$.
V79 gives $\operatorname{Gram}(W_p)=C_1+R_{pp}$, and V80's Schur
complement is precisely subtraction of the projection onto
$\mathcal N_q$.  Its compact-feature component remains $W^{\min}$;
its orthogonal component becomes $(I-\Pi_q)Z$.  Subtracting $C_1$ proves
\eqref{f16v81:finite-residual}, including singular null Gram matrices.
Increasing closed spans have strongly convergent orthogonal projections
onto their closure.  Apply this to the finitely many $Z_i$ to obtain
\eqref{f16v81:infinite-residual}.  The residual Gram vanishes exactly
when its columns vanish.  This is a Hilbert-space closure statement;
the algebraic sphere-ideal statement in V80 does not imply it.
\end{proof}

\begin{proposition}[The minimum is an inverse-transfer moment, not the deficit]
\label{f16v81:spectral}
For the same bank,
\begin{equation}
 \begin{aligned}
 \mathsf Q_-&\succeq G-C_2,\qquad \mathsf Q_-\succeq2(G-C_1),\\
 \mathsf Q_--2(G-C_1)
 &=\bigl(\langle f_i,(P^{-1}+P-2I)f_j\rangle\bigr)_{ij}\succeq0.
 \end{aligned}
 \label{f16v81:floor-deficit}
\end{equation}
The final pairing is understood as a quadratic form.  For a normalized
source direction with physical spectral measure $\mu_{a,c}$ and
$H_a=-a^{-1}\log P_a$,
\begin{equation}
 \boxed{\quad
 \frac{c^*\mathsf Q_-c}{c^*Gc}
       =\int_0^\infty2\sinh(aE)\,d\mu_{a,c}(E).
 \quad}
 \label{f16v81:sinh}
\end{equation}
Thus $\mathsf Q_-=O(aG)$ controls an exponential high-energy moment,
not only participation at bounded physical energies.
\end{proposition}
\begin{proof}
Apply spectral calculus to
$\lambda^{-1}-\lambda\ge1-\lambda^2$ and
$\lambda^{-1}-\lambda-2(1-\lambda)=\lambda^{-1}+\lambda-2\ge0$.
All forms are finite by the theorem.  Substitution $\lambda=e^{-aE}$
gives \eqref{f16v81:sinh}.  It also gives
$\int E\,d\mu_{a,c}\le c^*\mathsf Q_-c/(2a\,c^*Gc)$, but the converse
is not furnished by this inequality.
\end{proof}

\subsection{All polynomial sphere constraints can leave a strict residual}

The following is an exact \emph{compact kinetic reference}: four cyclic
links, $b=1$, $\phi=\rho=1$, and $P=C_\gamma^{\otimes4}$.  It is not the
interacting four-dimensional Wilson vacuum.  The source
$F=w(X_1X_2X_3X_4)$ is invariant under the four vertex gauge rotations,
has Haar mean zero and variance $1/4$.  Put
\[
 z_\gamma=2I_1(\gamma)/\gamma,\qquad
 \lambda=I_2(\gamma)/I_1(\gamma),\qquad c_\gamma=\lambda^4.
\]
Then $PF=c_\gamma F$ and the complete feature minimum is
$Q_-=(c_\gamma^{-1}-c_\gamma)/4$.

\begin{theorem}[A complex-zero obstruction to exhausting the feature quotient]
\label{f16v81:zero}
Let $j>0$ be the first positive zero of $J_1$, and set
\begin{equation}
 d_\gamma^{\rm zero}=
 \frac{4J_2(j)^2 z_\gamma^2}{j^2\lambda^2}
             \exp[-j^2/\gamma-3\gamma]>0.
 \label{f16v81:zero-constant}
\end{equation}
For every real ambient polynomial $p$ whose restriction is $F$, and
after any finite or complete polynomial sphere-constraint reduction,
\begin{equation}
 \boxed{\quad
 R_{p\mid q}\ge
 \frac{c_\gamma^{-1}-c_\gamma}{4}+d_\gamma^{\rm zero}.
 \quad}
 \label{f16v81:strict-residual}
\end{equation}
In particular the polynomial hierarchy does not attain the complete
feature minimum, even when its degree is allowed to grow without bound
at this fixed $\gamma$.
\end{theorem}
\begin{proof}
For a Fock vector $v$ use its entire evaluation
$\widehat v(z)=\langle\mathrm e(\bar z),v\rangle$.  It obeys
$|\widehat v(z)|\le e^{|z|^2/2}\|v\|$ and sends a creation operator
to multiplication by the corresponding coordinate.  In this reference,
\[
 \widehat{\eta_1}(z)=z_\gamma^{-2}
       \prod_{e=1}^4\int_{S^3}e^{\sqrt\gamma z_e\cdot x}\,d\sigma(x),
 \qquad
 \widehat W_p(z)=p(z/\sqrt\gamma)\widehat{\eta_1}(z).
\]
Choose
\begin{equation}
 z_*=(ij e_0/\sqrt\gamma,\sqrt\gamma e_0,
                    \sqrt\gamma e_0,\sqrt\gamma e_0).
 \label{f16v81:zstar}
\end{equation}
The first factor in $\widehat{\eta_1}(z_*)$ is $2J_1(j)/j=0$.
Therefore every polynomial witness $W_p$, every sphere-null witness,
and every norm limit of their linear combinations vanishes at $z_*$.
The last assertion uses continuity of the evaluation functional, not a
formal assertion about divisibility of a limiting power series.

The compact-feature minimum is $W^{\min}=\eta_F/c_\gamma$.
The one-link identities needed to evaluate it are
\[
 \int x\,e^{\gamma x^0}d\sigma=z_\gamma\lambda e_0,
 \qquad
 \int x\,e^{ijx^0}d\sigma=i\frac{2J_2(j)}j e_0.
\]
The second follows by differentiating $2J_1(j)/j$ and using
$(J_1(t)/t)'=-J_2(t)/t$.  Multilinearity of the scalar quaternion product
now gives
\[
 \widehat{\eta_F}(z_*)=
 i\frac{2J_2(j)}j z_\gamma\lambda^3,
 \qquad
 \widehat W^{\min}(z_*)=
 i\frac{2J_2(j)}j\frac{z_\gamma}{\lambda}\ne0.
\]
A positive zero of $J_1$ is simple; its derivative there is $-J_2(j)$,
so the asserted nonvanishing follows.  These statements also follow
from the Bessel equation's uniqueness and its standard positive zeros.

For any admissible witness, including its finite constraint optimum,
write $W=W^{\min}+Z$ with $Z\perp\mathcal K$.  Since
$\widehat W(z_*)=0$, continuity of evaluation gives
\[
 \|Z\|^2\ge
 e^{-|z_*|^2}|\widehat W^{\min}(z_*)|^2=d_\gamma^{\rm zero}.
\]
The same inequality passes to the complete hierarchy, either by norm
limits of its orthogonal projections or by monotonicity of the reduced
costs.  Equation \eqref{f16v81:finite-residual} proves the result.
\end{proof}

This example does not give a quantitative lower bound on the residual
for the interacting vacuum.  It proves that a proposed general equality
between polynomial sphere reduction and the complete feature quotient
would be false.  Algebraically generating every zero restriction is not
the same as generating every invisible auxiliary feature vector.  The
strict gap above may itself become very small with $\gamma$; no uniform
physical obstruction is inferred from its size.

\subsection{A native innovation floor for the transfer-paired cost}

We now return to the full original Wilson history and the elementary
spatial plaquette bank, on the even nondegenerate tori of V79:
\begin{equation}
 F_i=\sum_{p\in\mathcal P}s_{ip}w(U_p),\quad
 \mathcal S=\sum_p s_ps_p^*,\quad
 d_\gamma=[32000(1+6\gamma)^2]^{-1},\quad
 \kappa_\gamma=\tfrac43d_\gamma^2.
 \label{f16v81:bank}
\end{equation}
Repeated marks are combined as before.  Let $\mathcal Z_0$ be the
sigma-field of all history links except the spatial links on slice zero.
In particular it contains the full spatial slice at time one.

\begin{theorem}[The actual time innovation bounds every transfer-paired reduction]
\label{f16v81:native-floor}
At the finite regulator, and uniformly under the stated cylinder limits,
\begin{equation}
 \operatorname{Cov}(F(U_0)\mid\mathcal Z_0)
       \succeq\kappa_\gamma\mathcal S\quad\hbox{almost surely},
 \qquad
 \boxed{\quad G-C_2\succeq\kappa_\gamma\mathcal S.\quad}
 \label{f16v81:innovation}
\end{equation}
Consequently every ambient polynomial extension, every finite constraint
reduction and its decreasing limit satisfy
\begin{equation}
 \boxed{\quad
 R_{p\mid q}\succeq\mathsf Q_-\succeq G-C_2
                   \succeq\kappa_\gamma\mathcal S.
 \quad}
 \label{f16v81:R-native-floor}
\end{equation}
This includes the complete feature minimum, not just the polynomial
hierarchy.
\end{theorem}
\begin{proof}
Condition first on $\mathcal Z_0$.  The remaining slice density keeps all
its spatial plaquettes and both sets of temporal staples.  Its one-link
conditional fields still have norm at most $6\gamma$.  V79's two-stage
argument can therefore be performed inside this conditional law: for
each of its twelve link colors, conditional variance is at least
$d_\gamma\sum_{e\ {\mathrm{in\ color}}}\mathbb E(|B_e|^2\mid\mathcal Z_0)$.
For each fixed $e$, its four private opposite links give
\[
 \mathbb E(|B_e|^2\mid\mathcal Z_0)
       \ge4d_\gamma\sum_{p\ni e}s_p^2.
\]
The color and private-link integrations each use their exact original
conditionals, with $\mathcal Z_0$ held fixed.  No marginalized slice is
assigned a new independent law.  Averaging the twelve lower bounds and
counting the four edges of a plaquette gives $\kappa_\gamma\sum_ps_p^2$
for every scalar coefficient direction.  This is the inherited geometry
used conditionally at a native time separator, not a new variance
factorization assumption.

Conditional covariance decreases on average upon further conditioning.
Since $\sigma(U_1)\subset\mathcal Z_0$,
\[
 \mathbb E\operatorname{Cov}(F(U_0)\mid U_1)
 \succeq\mathbb E\operatorname{Cov}(F(U_0)\mid\mathcal Z_0)
 \succeq\kappa_\gamma\mathcal S.
\]
Reversibility and the native Markov property identify the left side with
$G-C_2$: the conditional mean of the centered source is $Pf(U_1)$.
Uniform conditional constants justify the stationary cylinder passage.
Equation \eqref{f16v81:floor-deficit} and the feature theorem finish the
proof.  Positivity of the inverse form is used at finite spatial regulator;
no inverse-transfer domain is claimed in an unspecified infinite-volume
limit.
\end{proof}

\begin{proposition}[The microscopic transfer-paired $O(aG)$ test also fails]
\label{f16v81:relative}
In the reflection-compatible native laws of V80, put
\[
 N=|\mathcal P|,\qquad
 D_\gamma=b_\gamma^2+2b_\gamma+2,
 \qquad b_\gamma=384\left(8+\log\gamma-\tfrac23\log\frac8{3\pi^3}\right).
\]
Then
\begin{equation}
 R_{p\mid q}\succeq\mathsf Q_-
       \succeq\frac{\kappa_\gamma\gamma^2}{ND_\gamma}\,G.
 \label{f16v81:relative-floor}
\end{equation}
For $\gamma_a\ge1$, no nonzero source quotient can satisfy
$R_{p\mid q}\preceq2aE_*G$, with fixed $0<E_*<\infty$, along a sequence
with
\begin{equation}
                aN_a\gamma_a^2D_{\gamma_a}\longrightarrow0.
 \label{f16v81:R-exclusion}
\end{equation}
The same exclusion holds for $\mathsf Q_-$ itself.  Fixed microscopic
banks on the stated exponential $a(\gamma)$ trajectory are included.
\end{proposition}
\begin{proof}
V80's native second-moment bound is
$G\preceq ND_\gamma\gamma^{-2}\mathcal S$.  Combine it with
\eqref{f16v81:R-native-floor}.  For $\gamma\ge1$,
\[
 \kappa_\gamma\gamma^2
 \ge\frac{4}{3\cdot32000^2\cdot7^4}\,\gamma^{-2}.
\]
Thus the proposed upper bound would force $aN\gamma^2D_\gamma$ to be
at least a fixed positive constant depending on $E_*$.  This contradicts
\eqref{f16v81:R-exclusion}.  The estimates are matrix inequalities and
are invariant under changing coefficients on the declared quotient.
\end{proof}

This is a new obstruction derived from a native innovation, not an
incorrect transfer of V80's equal-time mixed-derivative floor to $R$.
Its condition has an additional factor $\gamma^2$ and is correspondingly
weaker.  It is still not an upper bound on the source correlation, not a
proof of ultraviolet escape, and not an obstruction for arbitrary
growing physical source banks.  A small fraction of fast spectral weight
can obstruct an inverse-moment criterion without destroying a surviving
finite-energy component.

\subsection{A bounded-energy source can fail the minimum witness test}

\begin{proposition}[Uniform source convergence and an inverse-moment divergence]
\label{f16v81:counterexample}
In the four-link compact kinetic reference, independently declare
$0<a<1/2$, $\gamma_a=1/a$, and $K_a=\lceil a^{-2}\rceil$.  Let $Q$ be the
cyclic plaquette and set
\begin{equation}
 f_a=\sqrt{1-a^6}\,\chi_1(Q)+a^3\chi_{K_a}(Q).
 \label{f16v81:two-sector-source}
\end{equation}
These are centered gauge-invariant polynomials, $\|f_a\|_H=1$, uniformly
bounded in supremum norm, and $f_a\to\chi_1(Q)$ uniformly.  For every
fixed $t>0$,
\begin{equation}
 \langle f_a,P_a^{\lceil t/a\rceil}f_a\rangle_H\longrightarrow e^{-6t},
 \qquad
 \int E\,d\mu_a(E)\longrightarrow6.
 \label{f16v81:surviving-reference}
\end{equation}
Nevertheless its complete feature cost diverges:
\begin{equation}
                 \langle f_a,(P_a^{-1}-P_a)f_a\rangle_H\longrightarrow\infty.
 \label{f16v81:inverse-divergence}
\end{equation}
In particular bounded mean physical energy, and even uniform convergence
of these bounded sources, do not imply the inverse-moment test.
\end{proposition}
\begin{proof}
The holonomy $Q$ is Haar under independent Haar links.  Its distinct
characters are orthonormal and satisfy
$P_a\chi_k(Q)=c_{k,a}\chi_k(Q)$, where
$c_{k,a}=\lambda_k(1/a)^4$.  Therefore the variance is one and every
correlation is exactly
\[
 (1-a^6)c_{1,a}^{\lceil t/a\rceil}
                   +a^6c_{K_a,a}^{\lceil t/a\rceil}.
\]
Since $|\chi_k|\le k+1$, the high component has supremum at most
$a^3(K_a+1)=O(a)$.  The fixed-order Bessel expansion, equivalently the
one-link compact Laplace calculation, gives
$\lambda_1(1/a)=1-3a/2+O(a^2)$; hence the first sector's physical energy
tends to six and the displayed correlations have the asserted limit.

The inherited lower barrier
$\lambda_k(\gamma)\ge e^{-k(k+3)/\gamma}$ gives
$-a^{-1}\log c_{K_a,a}\le4K_a(K_a+3)=O(a^{-4})$.
Its contribution to the mean energy is at most $O(a^2)$, proving the
second assertion of \eqref{f16v81:surviving-reference}.
In the other direction, the inherited factorial bound gives
\[
 c_{K_a,a}^{-1}\ge
       \bigl[(K_a+1)!(2a)^{K_a}\bigr]^4.
\]
Stirling's elementary lower bound $m!\ge(m/e)^m$ shows that the logarithm
on the right grows at least as a positive multiple of
$a^{-2}\log(1/a)$.  Multiplication by $a^6$ cannot offset it.  The high
sector contribution $a^6(c_{K_a,a}^{-1}-c_{K_a,a})$ therefore diverges.
All these are exact compact kinetic sectors, not a Gaussian approximation
or an assertion about the interacting four-dimensional source.
\end{proof}

The spacing $a=1/\gamma$ in this reference is explicitly chosen only to
test necessity of the criterion.  It is not the physical scaling proposed
for the Wilson continuum program.  The counterexample's role is logical:
an inverse-transfer certificate can fail although its source has a
nonzero positive-separation limit and a convergent first energy moment.
The preceding native obstruction consequently must not be promoted to a
native source-collapse statement.

\subsection{A complete minimum witness built only from forward native time}

The remaining construction applies to the actual native source, including
the simultaneous spatial-flow function, without ambient differentiability.
Let $f$ be any finite centered $L^2(\pi)$ bank and $C_n=f^*P^nf$ in
synthesis notation, so $C_0=G$.  Choose an integer $k\ge1$ and let
$f^{[k]}=P^kf$.  Its entrance Gram is $G_k=C_{2k}$.  In the reflected
history interpretation this is the vector created by placing the original
slice source $k$ steps beyond the reflection plane, not a new physical
clock and not a claim that $P^kf$ is spatially local.

\begin{theorem}[Forward construction, exact cost and a paid buffer bound]
\label{f16v81:buffer}
For every such bank the minimum witness is explicitly
\begin{equation}
 W^{\min}_{[k]}=\mathcal J P^{k-1}f,
 \qquad \mathcal J^*W^{\min}_{[k]}=P^kf.
 \label{f16v81:forward-witness}
\end{equation}
It realizes
\begin{equation}
 \boxed{\quad
 \begin{pmatrix}C_{2k+1}&C_{2k}\\C_{2k}&C_{2k-1}\end{pmatrix}\succeq0,
 \qquad
 \mathsf Q_{-,k}=C_{2k-1}-C_{2k+1}.
 \quad}
 \label{f16v81:forward-block}
\end{equation}
There is no inverse-transfer computation or polynomial extension in this
construction.  Moreover
\begin{equation}
 0\preceq\mathsf Q_{-,k}\preceq m_kG,\qquad
 m_k=\frac2{2k+1}
       \left(\frac{2k-1}{2k+1}\right)^{(2k-1)/2}
       \le\frac1k.
 \label{f16v81:buffer-cost}
\end{equation}
If, separately, $G_k\succeq r_aG$ with $r_a>0$, then
$\mathsf Q_{-,k}\preceq(m_k/r_a)G_k$.  For
$k=\lceil\sigma/a\rceil$, fixed $\sigma>0$, a uniform $r_a\ge r_*>0$
gives the sufficient relative cost
$\mathsf Q_{-,k}\preceq a(\sigma r_*)^{-1}G_k$.
\end{theorem}
\begin{proof}
By \eqref{f16v81:J},
$\mathcal J^*(\mathcal J P^{k-1}f)=P^kf$.  The proposed witness is in
$\mathcal K$, so it is the minimum-norm solution.  It is well-defined for
all $f\in L^2(\pi)$; equivalently $P^kf$ belongs to
$\operatorname{Dom}P^{-1/2}$ because $P^{k-1/2}$ is bounded.
Take the Gram of $(\mathcal J P^kf,\mathcal J P^{k-1}f)$ to obtain
\eqref{f16v81:forward-block}.  Its bottom-right minus top-left block is
exactly the inverse-transfer minimum for $f^{[k]}$.

The nonnegative scalar multiplier of this difference is
$\lambda^{2k-1}(1-\lambda^2)$.  Its maximum on $[0,1]$ occurs at
$\lambda^2=(2k-1)/(2k+1)$, giving $m_k$.
Since $m_k\le2/(2k+1)\le1/k$, spectral calculus proves
\eqref{f16v81:buffer-cost} in every direction.  A separately supplied
retention bound permits replacing $G$ by $G_k/r_a$; the final statement
uses $k\ge\sigma/a$.  No lower estimate on $G_k$ was used to prove the
absolute inequality.
\end{proof}

The block in \eqref{f16v81:forward-block} is ordinary Hankel positivity
for a positive transfer, here identified with the exact minimum of V79's
feature problem.  We do not claim a new spectral positivity principle.
The new point is that a buffered source has an attained, derivative-free
minimum witness made solely of forward native vectors; it need not be
approximated by the polynomial-null module which the complex-zero example
shows can be incomplete.

The normalization qualification remains essential.  A fixed physical
buffer may remove most of the original contact Gram, as V66--V67
already demonstrate in their distinct examples.  The absolute
$O(aG)$ estimate above is not an $O(aG_k)$ theorem without retention,
and imposing retention of a divergent contact norm is not necessary for
a canonically normalized reflected field.  A direct comparison of
$C_{2k-1}-C_{2k+1}$ with $C_{2k}$ can instead be appropriate.  Neither
comparison has been evaluated at physical scale for the full
simultaneous-flow bank here.

\subsection{Noncircular consequence, verification and result ledger}

For the actual collective source, the two unresolved numerical objects
can now be stated without an off-sphere extension:
\[
 G_{a,k}=C_F(2k),\qquad
 \mathsf Q_{a,k}=C_F(2k-1)-C_F(2k+1),\qquad
                  ka\longrightarrow\sigma>0.
\]
They are three forward covariances of the same native simultaneous-flow
source.  The feature quotient eliminates auxiliary residuals exactly;
it does not evaluate these covariances.  A useful next advance must
bound this signed adjacent-lag difference in the source's physical
normalization and establish a nonzero reflected lower Gram, or directly
estimate the surviving source's spectral distribution.  Repeatedly
increasing the sphere-constraint degree is not a substitute: in a compact
reference it leaves a strictly positive extra cost, and even the ideal
minimum can demand a needlessly strong high-energy moment.

The finite microscopic $R=O(aG)$ route is now excluded under the stated
native scaling, independently of the equal-time argument.  The growing
and time-buffered collective bank is not covered by that exclusion.
A lower source-energy estimate still does not exclude weight at energy
zero.  Common dense-core upper contraction and the interacting continuum
construction remain separate obligations.

The companion diagnostic is
\begin{center}\small
\path{verify_ym_f16_v81_feature_quotient.py}.
\end{center}
It separates exact spectral identities, finite positive-kernel feature
and constraint tests, compact Bessel/character calculations, explicit
four-link gauge identities, conditional-variance fixtures and lattice
incidence, and forward-buffer bounds.  The Bessel evaluations and finite
matrix checks are floating-point diagnostics, not outward-rounded
certificates.  They do not sample the four-dimensional Wilson vacuum or
prove the native physical-scale covariance input.  No numerical outcome
is a premise of the analytic proofs.  Removing this delimited appendix
and reverting the one displayed version change recovers V80 byte for
byte; the inherited archive is preserved, not independently recertified.

\begin{record}[V81 result ledger]
\label{f16v81:ledger}
\emph{Proved here:} the exact complete compact-feature minimum and its
inverse-transfer domain; the orthogonal residual formula for finite and
complete polynomial sphere-constraint reduction; a strict complex-zero
obstruction to exhausting that minimum in a gauge-invariant compact
kinetic reference; a native conditional-innovation floor for every
transfer-paired reduction and the complete minimum; the corresponding
fixed-microscopic $O(aG)$ exclusion; a bounded-source kinetic reference
with convergent mean energy and reflected limit but divergent minimum
inverse moment; and an attained forward-native witness for buffered
sources with its exact cost and explicit absolute bound.

\emph{Inherited or established background:} V79's exact Perron/coherent
construction, V80's finite Schur reduction, V79's conditional plaquette
geometry, V76--V80's declared native moments, positive native transfer,
Hilbert-space polar decomposition, Bessel identities and scalar spectral
calculus.  No Gaussian or kinetic reference replaces the native law.

\emph{Not asserted:} completeness of the polynomial-null witness module;
equality of the minimum with twice the one-step deficit; necessity of an
inverse-exponential moment for physical noncollapse; collapse of the
microscopic or simultaneous-flow source from a failed sufficient test;
a relative buffered estimate without its normalization; or physical
spectral information obtained by choosing a new time unit.

\emph{Still unproved:} a nonzero reflected Gram in physical normalization
for the simultaneous spatial-flow bank; control of its signed forward
covariances and spatial dependence; compatibility of the thermodynamic
and continuum limits; common dense-core upper contraction in physical
time; and the interacting four-dimensional continuum Yang--Mills mass gap.
\end{record}

\begin{firewall}
Polynomial source identities, feature-space orthogonality and spectral
moments are distinct operations.  Native innovation gives the new floor;
kinetic references test necessity and hierarchy completeness only.  The
physical source, its original law, and its time step are never replaced
by an auxiliary quotient.  All V1--V80 mathematical text remains unchanged.
\end{firewall}
% END F16 V81 APPEND

% BEGIN F16 V82 APPEND -- 2026-09-19
% Append-only continuation of the Post-Fifteenth-Archive V81.
\clearpage
\section[V82: positive-time secants and stable buffered reconstruction]{V82: positive-time secants\\and stable buffered reconstruction}
\label{f16v82:context}

V81 identifies the complete buffered feature cost with the signed native
packet $C_{2k-1}-C_{2k+1}$, but its absolute estimate uses the original
contact Gram.  That Gram need not remain finite in the normalization of a
surviving reflected composite field.  We replace contact retention by a
comparison between two \emph{positive-time} Grams.  Operator compression,
not an entrywise matrix-power argument, gives a sharp secant bound for the
entire source quotient.  A bounded earlier positive-time Gram and a
nonzero later one then imply the desired relative $O(a)$ cost without a
condition on the contact norm.  Neither of these two physical-scale native
bounds is assumed to have been proved for the simultaneous-flow bank.

There is also a quantitative reconstruction result.  Estimating the two
adjacent covariances independently would amplify their errors by $a^{-1}$
when forming the physical energy matrix.  We instead propagate their
\emph{common source error} through the positive spectral multiplier of the
signed difference.  At a positive physical buffer this removes that extra
loss.  V77's existing polynomial cutoff therefore controls both the Gram
and the physical signed difference of the actual simultaneous-flow source.
The final native task is reduced to positive-time source information;
it is not replaced by another off-sphere constraint calculation.

\subsection{Source convention and inherited spectral input}

At each finite spatial regulator use the same centered, invariant source
synthesis $V:\mathbb C^r\to L^2(\pi)$ and the original positive injective
Markov transfer $P$ of V81.  Write
\begin{equation}
 C_n=V^*P^nV,\quad H=-a^{-1}\log P,\quad
 \mathcal Q_n=C_{n-1}-C_{n+1}\quad(n\ge1).
 \label{f16v82:data}
\end{equation}
At $n=2k$, $C_n$ is the reflected Gram of the original slice source placed
$k$ steps beyond the reflection plane, and $\mathcal Q_n$ is exactly
V81's attained minimum cost.  All these are forward native covariances.
The spectral interpolation $C(t)=V^*e^{-tH}V$ is used only to prove
estimates; it does not declare fractional lattice slices or alter $a$.

Injectivity makes all $C_n$ have the same exact kernel as $V$.  Work on a
declared complement of that kernel, so the displayed inverse square roots
are finite matrices.  A regulator-dependent common coefficient map is
allowed, with its quotient and normalization fixed before the estimates.
No numerical eigenvalue threshold is used to choose the quotient.

The operator-power inequalities used below are standard instances of
operator Jensen; see F.~Hansen and G.~K.~Pedersen,
\emph{Jensen's Operator Inequality}, Bull.\ London Math.\ Soc.\ \textbf{35}
(2003), 553--564,
\href{https://arxiv.org/abs/math/0204049}{arXiv:math/0204049}.
We give the compression proof, including the unbounded form limit, rather
than import a covariance or continuum theorem.  V81's forward witness and
V77's native reconstruction estimate are inherited.  The earlier scalar
pressure secants and V66's history-window retention do not estimate the
noncommuting positive-time matrices considered here.

\subsection{A positive-time matrix secant controls the complete minimum}

\begin{theorem}[Two positive-time Grams, without contact retention]
\label{f16v82:secant}
Let $n>d\ge1$ be integers, and set
\begin{equation}
 G=C_n,\qquad B=G^{-1/2}C_{n-d}G^{-1/2}\succeq I.
 \label{f16v82:B}
\end{equation}
For every integer $1\le j\le d$,
\begin{equation}
 \boxed{\quad
 C_{n-j}\preceq G^{1/2}B^{j/d}G^{1/2},\qquad
 C_{n+j}\succeq G^{1/2}B^{-j/d}G^{1/2}.
 \quad}
 \label{f16v82:two-powers}
\end{equation}
In particular,
\begin{equation}
 \boxed{\quad
 0\preceq\mathcal Q_n
 \preceq G^{1/2}(B^{1/d}-B^{-1/d})G^{1/2}.
 \quad}
 \label{f16v82:Q-secant}
\end{equation}
Let $D=G^{-1/2}C_{n+d}G^{-1/2}$ and
$A_n=G^{-1/2}V^*He^{-naH}VG^{-1/2}$.  Then
\begin{equation}
 -\frac1{da}\log D\preceq A_n\preceq\frac1{da}\log B.
 \label{f16v82:energy-secants}
\end{equation}
No pair of these compressed matrices is assumed to commute.
\end{theorem}
\begin{proof}
The map $J=P^{n/2}VG^{-1/2}$ is an isometry.  Put $X=P^{-d}\succeq I$
as a positive spectral form.  Its compression is finite because
$X^{1/2}J=P^{(n-d)/2}VG^{-1/2}$ is bounded, and $J^*XJ=B$.
For $s>0$, spectral Cauchy--Schwarz, equivalently the minimum quadratic
form for the inverse, gives
\[
 J^*(X+sI)^{-1}J\succeq(B+sI)^{-1}.
\]
For $0<r<1$ the positive integral representations
\[
 x^r=\frac{\sin\pi r}{\pi}\int_0^\infty
       s^{r-1}\frac{x}{x+s}\,ds,\qquad
 x^{-r}=\frac{\sin\pi r}{\pi}\int_0^\infty
       s^{-r}\frac1{x+s}\,ds
\]
therefore imply
$J^*X^rJ\preceq B^r$ and $J^*X^{-r}J\succeq B^{-r}$.
One can first replace $X$ by $\min(X,L)$; its compressed first moment
converges to $B$, the positive powers converge in form on $\operatorname{Ran}J$,
and the negative powers converge boundedly.  This proves the assertions
also for unbounded $X$.  At $r=1$ use the exact first moment and the
inverse inequality.  Taking $r=j/d$ proves \eqref{f16v82:two-powers}.
Subtract its two valid bounds at $j=1$; $P\preceq I$ gives positivity
of the actual difference.

The same resolvent argument for $\log x$, or its integral representation,
gives $J^*\log XJ\preceq\log B$.  Since $\log X=daH$, this is the upper
energy bound.  Apply convexity of $-\log$ to $Y=P^d$ to get
$-\log(J^*YJ)\preceq J^*(-\log Y)J$, proving the lower bound.  These
forms are finite: $H e^{-(n-d)aH}$ is bounded after any further positive
time, or directly $\log X\le X$ on this compression.  Regularizing
$Y$ away from zero justifies its inverse/logarithmic limit.  There is no
operation of applying a scalar monotone function to a claimed order
between two unrelated covariance matrices.
\end{proof}

\begin{proposition}[Relative cost, future correlations and energy tails]
\label{f16v82:relative}
If $C_{n-d}\preceq R C_n$ with $1\le R<\infty$, then
\begin{equation}
 \boxed{\quad
 \mathcal Q_n\preceq2\sinh\!\left(\frac{\log R}{d}\right)C_n,
 \qquad C_{n+m}\succeq R^{-m/d}C_n\quad(m\ge0).
 \quad}
 \label{f16v82:scalar-secant}
\end{equation}
For every normalized direction in the buffered source, its physical
spectral probability measure $\mu_{n,c}$ obeys
\begin{equation}
 \int e^{daE}\,d\mu_{n,c}(E)\le R,\qquad
 \mu_{n,c}((E,\infty))\le R e^{-daE},\qquad
 \int E\,d\mu_{n,c}(E)\le\frac{\log R}{da}.
 \label{f16v82:exponential-tail}
\end{equation}
In particular if $d_aa\to h>0$ and $R_a\le R_*$, then
$\mathcal Q_{n_a}=O(aC_{n_a})$.  A sufficient input is, on the same
coefficient quotient with a declared metric $S_a\succ0$,
\begin{equation}
 C_{n_a-d_a}\preceq M S_a,\qquad C_{n_a}\succeq m S_a,
 \qquad 0<m\le M<\infty.
 \label{f16v82:anchors}
\end{equation}
This needs no bound or retention fraction for $C_0$.
\end{proposition}
\begin{proof}
The spectrum of $B$ lies in $[1,R]$.  Apply scalar spectral calculus
\emph{within this one matrix} to \eqref{f16v82:Q-secant}.  The exponential
moment in \eqref{f16v82:exponential-tail} is exactly the ratio
$c^*C_{n-d}c/(c^*C_nc)$.  Markov's inequality proves its tail bound, and
scalar Jensen for $e^{daE}$ bounds its mean.  Jensen for $e^{-maE}$ then
gives the second part of \eqref{f16v82:scalar-secant} for every $m$, not
just $m\le d$.  Applying this separately in every coefficient direction
proves the Loewner bound.  Since $1/d=a/(da)$, the claimed $O(a)$ follows.
The two anchors give $R\le M/m$.  For a single exact energy sector all
the secant and scalar cost bounds are equalities, so their constants
cannot be improved using only that ratio.
\end{proof}

The lower energy matrix in \eqref{f16v82:energy-secants} is a lower bound
on a compressed \emph{mean} energy, not on the support of the physical
spectrum.  The upper bounds above ensure nonescape at high energy once
\eqref{f16v82:anchors} is known; they allow spectral mass tending to zero.
Neither anchor is supplied by finite-regulator positivity alone.  For the
actual simultaneous-flow family their physical-scale values remain unpaid.

\subsection{The signed packet has a controlled physical derivative}

\begin{proposition}[An earlier reflected norm pays the central-difference error]
\label{f16v82:derivative}
Let $h=da>a$, $n>d$, and put $E_n=V^*He^{-naH}V$.  Then
\begin{equation}
 \boxed{\quad
 0\preceq\frac{\mathcal Q_n}{2a}-E_n
 \preceq\frac{a^2}{6}
          \left(\frac3{e(h-a)}\right)^3 C_{n-d}.
 \quad}
 \label{f16v82:central-error}
\end{equation}
More generally, for $t>s>0$ and integer $r\ge1$,
\begin{equation}
 0\preceq V^*H^re^{-tH}V
 \preceq\left(\frac r{e(t-s)}\right)^r C(s).
 \label{f16v82:moment-bound}
\end{equation}
These are absolute native matrix bounds at positive separation, not bounds
in terms of the potentially divergent contact norm.
\end{proposition}
\begin{proof}
For $x\ge0$, Taylor integration gives
$0\le\sinh x-x\le x^3e^x/6$.  Multiply by $e^{-naE}/a$ with $x=aE$.
After retaining $e^{-(n-d)aE}$, the remaining multiplier is at most
$a^2E^3e^{-(h-a)E}/6$.  Its supremum is the displayed constant.
The second inequality similarly uses
$\sup_{E\ge0}E^re^{-(t-s)E}=(r/[e(t-s)])^r$.
Both inequalities are scalar spectral comparisons before source
compression, so they remain valid for noncommuting source Grams.
\end{proof}

\begin{theorem}[Positive-time convergence already controls the buffered cost]
\label{f16v82:convergence}
Consider a sequence of finite native regulators and a fixed finite source
coefficient space, with all normalizations included in $V_a$.  Suppose
$C_a(t_0)$ is uniformly bounded for one $t_0>0$.  Suppose the covariance
matrices converge at every real $t$ in some nonempty open interval
$I\subset(t_0,\infty)$.  It suffices instead to have convergence at the
corresponding lattice times $n_a(t)a\to t$.  Then their spectral
interpolations and every derivative converge locally uniformly on
$\{z:\operatorname{Re}z>t_0\}$ to a holomorphic matrix function $C(z)$.
For lattice times $n_aa\to t>t_0$,
\begin{equation}
 \boxed{\quad
 \frac{C_{a,n_a-1}-C_{a,n_a+1}}{2a}\longrightarrow -C'(t).
 \quad}
 \label{f16v82:cost-limit}
\end{equation}
If $C(t)$ is strictly positive on the declared quotient, the normalized
energy matrices converge to
$C(t)^{-1/2}[-C'(t)]C(t)^{-1/2}$.  In particular the relative buffered
$O(a)$ test then has no separate inverse-moment obstruction.
\end{theorem}
\begin{proof}
Factor $C_a(z)=V_{a,t_0}^*e^{-(z-t_0)H_a}V_{a,t_0}$ with
$V_{a,t_0}=e^{-t_0H_a/2}V_a$.  On $\operatorname{Re}z\ge t_0$ this has
norm at most $\|C_a(t_0)\|$.  Cauchy's integral formula on disks within
the open half-plane gives uniform bounds and equicontinuity for every
derivative.  Successive compact exhaustion therefore gives locally
uniform analytic subsequential limits.  Their values on $I$ coincide
by hypothesis; the identity theorem makes all limits identical, proving
full convergence and derivative convergence.  Equicontinuity also
converts convergence on approaching lattice times to convergence at the
real argument.  Alternatively the derivative bounds follow directly
from \eqref{f16v82:moment-bound} on each smaller real interval.

Use \eqref{f16v82:central-error} with a fixed earlier positive separation,
rounded to a lattice anchor strictly between $t_0$ and $t$.  Its norm is
uniformly bounded by monotonicity from $t_0$, and $h-a$ stays positive.
The $O(a^2)$ error and convergence of $-C_a'$ prove
\eqref{f16v82:cost-limit}.  Strict positivity makes inverse square roots
continuous on the fixed quotient.  No convergence of $C_a(0)$, no
ordinary $L^2$ limit of the raw source, and no interacting continuum
measure have been inferred from these conditional hypotheses.
\end{proof}

This theorem is a consequence of positive-transfer spectral calculus,
not a construction of the assumed covariance limit.  Its purpose is to
remove an unnecessary extra regularity task: once a canonically
normalized reflected limit is bounded on an earlier positive-time
neighborhood and nonzero later, the adjacent cost follows.  A positive
value at one isolated time, with no earlier control, is not enough.

\subsection{Common-source reconstruction avoids an inverse-spacing error}

Use exactly V77's positive input filter $F^K=\mathcal A_{K,S}F$ for a
bounded real source bank on $M$ original slice links.  Set
\begin{equation}
 B_F^2=\sum_i\left[\frac{\sup F_i-\inf F_i}{2}\right]^2,
 \qquad \epsilon_K=\min\{2,9\pi^2M\gamma/(K+3)\}.
 \label{f16v82:filter-data}
\end{equation}
Center each family under the same native law.  V77 proves the synthesis
bounds $\|Pf\|,\|Pf^K\|\le B_F$ and
$\|P(f-f^K)\|\le\epsilon_K B_F$.

\begin{theorem}[Stable reconstruction of the complete buffered cost]
\label{f16v82:stable}
For $k\ge2$ put
\begin{equation}
 \mu_k=\frac2{2k-1}
       \left(\frac{2k-3}{2k-1}\right)^{(2k-3)/2}\le\frac1{k-1}.
 \label{f16v82:mu}
\end{equation}
Then the actual, jointly reconstructed packets satisfy
\begin{equation}
 \boxed{\quad
 \|\mathcal Q_{2k}(F)-\mathcal Q_{2k}(F^K)\|
       \le2\epsilon_K B_F^2\mu_k,
 \qquad
 \left\|\frac{\mathcal Q_{2k}(F)-\mathcal Q_{2k}(F^K)}{2a}\right\|
       \le\frac{\epsilon_K B_F^2}{a(k-1)}.
 \quad}
 \label{f16v82:stable-cost}
\end{equation}
Thus $k=\lceil\sigma/a\rceil$ and $a\le\sigma/2$ give a bound
$2\epsilon_KB_F^2/\sigma$ for the physical signed-difference error.
There is no additional factor $a^{-1}$ at fixed positive physical buffer.
\end{theorem}
\begin{proof}
Write $A=Pf$ and $A_K=Pf^K$ for column syntheses in the same native
Hilbert space.  The signed packet is
\[
 \mathcal Q_{2k}(F)=A^*D_kA,\qquad
 D_k=P^{2k-3}(I-P^2)\succeq0.
\]
Its scalar multiplier has maximum $\mu_k$ on $[0,1]$, by the elementary
maximization already used for V81's absolute buffer bound.  Expand the
\emph{one} quadratic-form difference as
$(A-A_K)^*D_kA+A_K^*D_k(A-A_K)$.
The inherited norm bounds give \eqref{f16v82:stable-cost}.  Finally
$a(k-1)\ge\sigma-a\ge\sigma/2$.
The small multiplier belongs to the signed difference before any norm is
taken; it is lost if the two covariances are bounded independently.
\end{proof}

\begin{proposition}[The existing simultaneous-flow cutoff pays both quantities]
\label{f16v82:cutoff}
For the unchanged source
\[
 O_i^\tau(U)=\frac{\zeta_a^2}{2a}\sum_x f_i(x)
             |\operatorname{Im}(\Phi_{a,\tau}U)_{12}(x)|^2,
 \qquad M_a=3\mathcal V_a/a^3,
\]
use V77's common input cutoff and write $F=O^\tau$, $F^K=O^{\tau,K}$.
For $k=\lceil\sigma/a\rceil$, $a\le\sigma/2$, both
\[
 \|C_{2k}(F)-C_{2k}(F^K)\|,\qquad
 \sigma\left\|\frac{\mathcal Q_{2k}(F)-\mathcal Q_{2k}(F^K)}{2a}\right\|
\]
are at most
\begin{equation}
 \boxed{\quad
 \mathcal D_{a,K}=
 \frac{27\pi^2}{8}\,
 \frac{\mathcal V_a\gamma_a\zeta_a^4}{a^{11}(K+3)}
                   \sum_i\|f_i\|_{1,a}^2.
 \quad}
 \label{f16v82:physical-error}
\end{equation}
The estimate holds for every finite simultaneous flow time $\tau$, also
when it depends on $a$.  Taking the same cutoff as V77,
\begin{equation}
 K+3\ge\frac{27\pi^2}{8}\mathcal V_a\gamma_a\zeta_a^4
          \left(\sum_i\|f_i\|_{1,a}^2\right)a^{-(N+11)},
 \label{f16v82:physical-K}
\end{equation}
pays error $a^N$ in both displayed quantities, with no extra power of $a$.
\end{proposition}
\begin{proof}
V77's exact oscillation budget is
$B_F^2\le\zeta_a^4(16a^8)^{-1}\sum_i\|f_i\|_{1,a}^2$.
Its covariance error is at most $2\epsilon_KB_F^2$, and the preceding
theorem gives the same bound after multiplying the energy error by
$\sigma$.  Substitute $\epsilon_K\le9\pi^2M_a\gamma_a/(K+3)$ and the
actual whole-slice link count.  This proof estimates the genuine source
approximation under the original law, not a polynomial approximation to
the flow differential.  The shared-path and source-identity properties
of V77 remain unchanged.
\end{proof}

If an additionally computed polynomial has componentwise supremum errors
$\delta_i$, put $\Delta^2=\sum_i\delta_i^2$.  Relative to the exact filtered
bank, the covariance error costs $2B_F\Delta+\Delta^2$.  The same
common-source argument bounds its physical signed-difference error by
\begin{equation}
 \frac{\mu_k}{2a}(2B_F\Delta+\Delta^2)
 \le\frac1{\sigma}(2B_F\Delta+\Delta^2)
 \quad(a\le\sigma/2).
 \label{f16v82:coefficient-error}
\end{equation}
These are not bounds for independent rounding errors in separately
estimated covariance entries.  Such errors may still be amplified by
$a^{-1}$.  Likewise the coefficient count can remain exponential in the
input-link count.  The new stability theorem is not an efficient
algorithm for evaluating the reconstructed native matrices.

\subsection{A finite macroscopic certificate and exact tests of its scope}

A practical certificate can avoid adjacent-lag subtraction altogether.
Let $n=2k$, choose $d<n$ with $(n-d)a>0$, and suppose Hermitian estimates
$\widehat C_-,\widehat C_0$ have certified absolute errors $\eta_-,\eta_0$
for $C_{n-d},C_n$, respectively, on the fixed coefficient quotient.  These
errors include reconstruction and any numerical integration error.

\begin{proposition}[Two-matrix certificate with explicit uncertainty]
\label{f16v82:certificate}
If
\[
 L_0=\widehat C_0-\eta_0I\succ0,\qquad
 U_-=\widehat C_-+\eta_-I,
 \qquad
 R=\max\{1,\lambda_{\max}(L_0^{-1/2}U_-L_0^{-1/2})\},
\]
then $C_n\succeq L_0$ and
\begin{equation}
 \boxed{\quad
 C_{n-d}\preceq R C_n,\qquad
 \mathcal Q_n\preceq2\sinh((\log R)/d)C_n.
 \quad}
 \label{f16v82:certified-secant}
\end{equation}
All tail and delay conclusions of Proposition~\ref{f16v82:relative}
follow.  No value of $R$ for the interacting physical-scale collective
bank is supplied here.
\end{proposition}
\begin{proof}
The error intervals give $C_n\succeq L_0$ and $C_{n-d}\preceq U_-$.
By the generalized eigenvalue definition $U_-\preceq R L_0\preceq R C_n$.
Apply the secant theorem.  The test is performed on a predetermined
source quotient; failure of $L_0\succ0$ cannot be repaired by silently
deleting a numerically small direction.
\end{proof}

The following compact reference tests distinguish the two anchors from
contact retention and from an isolated-time normalization.  They are not
the Wilson exponential kernel or a replacement for its interacting law.
On each of four cyclic $\SU(2)$ links use the central Poisson kernel
\begin{equation}
 k_r(X)=\frac{1-r^2}{(1-2rw(X)+r^2)^2}
       =\sum_{j\ge0}(j+1)r^j\chi_j(X),\qquad r=e^{-a/4}.
 \label{f16v82:Poisson}
\end{equation}
It is positive, normalized and has Fourier multipliers $r^j$.  The series
follows by differentiating $r/(1-2rw+r^2)$; Haar orthogonality gives its
normalization.  The fourfold convolution is a reversible positive
Markov transfer and has
$P_a\chi_j(Q)=e^{-aj}\chi_j(Q)$ for $Q=X_1X_2X_3X_4$.

\begin{proposition}[Divergent contact norm is compatible with bounded positive-time cost]
\label{f16v82:contact-reference}
Set $L_a=\lceil a^{-2}\rceil$ and use the centered invariant polynomial
$f_a=\sum_{j=1}^{L_a}e^{-aj/2}\chi_j(Q)$.  Exactly,
\begin{equation}
 C_{a,n}=\sum_{j=1}^{L_a}e^{-(n+1)aj}.
 \label{f16v82:reference-C}
\end{equation}
Then $C_{a,0}\sim a^{-1}$, while for $n_aa\to t>0$,
\begin{equation}
 C_{a,n_a}\longrightarrow\frac1{e^t-1},\qquad
 \frac{\mathcal Q_{a,n_a}}{2a}
       \longrightarrow\frac{e^t}{(e^t-1)^2}.
 \label{f16v82:reference-limit}
\end{equation}
In particular $C_{a,n_a}/C_{a,0}\to0$, but both positive-time anchors
and the complete buffered relative $O(a)$ cost are bounded and nonzero.
\end{proposition}
\begin{proof}
Characters of $Q$ are orthonormal under four independent Haar links.
The exact multipliers give \eqref{f16v82:reference-C}; its geometric sum
at $n=0$ has the stated asymptotic.  On any interval $t\ge t_*>0$, its
tail after $L_a$ tends to zero uniformly, giving the first limit.
The signed multiplier is
$\sum_{j\le L_a}e^{-(n_a+1)aj}\sinh(aj)/a$.
For $a<t_*/2$, it is dominated by a summable constant times
$j e^{-t_*j/2}$.  Termwise convergence gives
$\sum_{j\ge1}j e^{-tj}=e^t/(e^t-1)^2$.  The source is a finite
polynomial at each regulator; its unbounded contact norm is not removed
before evaluating the reflected pairings.
\end{proof}

\begin{proposition}[An isolated-time lower Gram does not control the cost or a gap]
\label{f16v82:one-time}
In the same reference take $a=m^{-2}$ and
$f_a=e^{m/2}\chi_m(Q)$.  At the exactly sampled time $n_aa=1$,
$C_{a,n_a}=1$, but
\begin{equation}
 C_{a,n_a-d_a}=e^{d_aa m},\qquad
 \frac{\mathcal Q_{a,n_a}}{2a}
       =\frac{\sinh(am)}a\sim m\longrightarrow\infty.
 \label{f16v82:isolated-time}
\end{equation}
Thus a positive Gram at one time alone does not imply the secant input.
Conversely choose instead $r_a=e^{-a^2/4}$ in \eqref{f16v82:Poisson}
and $f=\chi_1(Q)$.  Both anchors are uniformly bounded and positive, and
the relative signed cost divided by $2a$ tends to zero, while this source's
physical energy is $a\to0$.
\end{proposition}
\begin{proof}
Both claims follow from the exact character eigenvalue.  In the first,
$C_{a,n}=e^{m}e^{-nam}$; the second uses eigenvalue $e^{-a^2}$.
They are distinct explicitly chosen kinetic references.  The first
normalization makes the earlier anchor grow; the second confirms that
high-energy control is not an upper-contraction or mass-gap theorem.
Neither example describes the simultaneous Wilson-flow vacuum.
\end{proof}

\subsection{The remaining native estimate and result ledger}

The two physical-scale quantities in V81 are no longer independent
targets once canonical positive-time bounds are available.  A bounded
Gram at an earlier positive separation and a nonzero Gram at a later
one imply an $O(a)$ complete buffered cost, exponential source-energy
tails and lower future correlations.  Positive-time covariance
convergence similarly forces convergence of the signed derivative.
This avoids retention of a divergent contact norm and an additional
inverse-exponential moment at the original slice.

The simultaneous-flow polynomial reconstruction now approximates that
physical signed difference with the same sufficient cutoff as its Gram.
It is essential that this is one common approximation of the source;
three independently approximated matrix entries have no such automatic
stability.  The two-matrix certificate states exactly how native numerical
errors, if supplied, would enter without a microscopic subtraction.

What remains unproved is a cutoff-uniform instance of
\eqref{f16v82:anchors}, or an equivalent canonical positive-time limit,
for the actual simultaneous spatial-flow source.  V78's finite positive
lower bound can be inserted at fixed regulator, but its severe constants
do not establish such uniformity.  We have not evaluated the interacting
matrices in the certificate.  The next calculation should target these
two positive-time native Grams on the same physical source quotient,
with thermodynamic and spatial-approximation costs retained.  Another
constraint-degree increase, contact-retention requirement, or unmarked
one-slice density estimate does not replace that information.  Even
successful nonescape and noncollapse leave the separate common dense-core
\emph{upper} contraction needed for a mass gap.

The companion script is
\begin{center}\small
\path{verify_ym_f16_v82_positive_time_secants.py}.
\end{center}
It checks noncommuting compression inequalities, signed spectral
multipliers, reconstruction and coefficient-error bounds, finite
positive quaternion magnetic kernels, exact compact Poisson reference
identities, and the append-only recovery.  Floating-point matrix and
special-function checks are diagnostics, not outward-rounded interval
certificates or formal proof verification.  No full four-dimensional
Wilson-vacuum data or physical-scale lower-Gram values are produced by
those tests; the analytic proofs do not depend on their outcomes.

\clearpage
\begin{record}[V82 result ledger]
\label{f16v82:ledger}
\emph{Proved here:} the noncommuting two-positive-time secant estimate
for the complete buffered minimum; its scalar relative, exponential-tail
and future-correlation consequences without contact retention; an explicit
central-difference error and a positive-time convergence implication;
stable common-source reconstruction of the physical signed cost with
V77's existing cutoff; a two-matrix uncertainty certificate; and exact
compact reference tests distinguishing positive-time control from contact
retention, isolated-time positivity and a mass gap.

\emph{Inherited or standard:} the original positive native transfer,
V81's attained forward witness, V77's conditional-density reconstruction,
operator Jensen and resolvent compression, spectral calculus, and elementary
analytic compactness.  No external covariance or continuum theorem is
imported, and the inherited manuscript is preserved rather than recertified.

\emph{Not asserted:} a uniform earlier/later native Gram sandwich for the
collective bank; ordinary-norm convergence of its reconstruction; immunity
to unrelated numerical subtraction errors; an efficient full coefficient
calculation; a shared eigenbasis for source matrices; or a lower physical
spectral edge inferred from a compressed mean energy.

\emph{Still unproved:} nonzero canonical positive-time native Grams for the
simultaneous-flow bank with the required uniform bounds; its spatial and
thermodynamic/continuum compatibility; common dense-core physical-time
upper contraction; and the interacting four-dimensional continuum
Yang--Mills mass gap.
\end{record}

\begin{firewall}
All time arguments and matrices belong to the original native source and
its original transfer.  Positive-time secants do not assume retention of
the contact norm.  The shared reconstruction error is propagated before
subtraction; native lower-Gram data are not fabricated from the resulting
certificate.  All V1--V81 mathematical text remains unchanged.
\end{firewall}
% END F16 V82 APPEND

% BEGIN F16 V83 APPEND -- 2026-09-19
% Append-only continuation of the Post-Fifteenth-Archive V82.
\clearpage
\section[V83: magnetic channel mixing and a native covariance calculation]{V83: magnetic channel mixing\\and a native covariance calculation}
\label{f16v83:context}

V82 reduces the buffered source problem to two positive-time native Grams,
but does not evaluate either of them uniformly on the proposed continuum
trajectory.  Rather than introduce another sufficient condition at that
interface, we compute a covariance of the \emph{actual simultaneous-flow}
source in a controlled regime.  The regime is explicitly different:
throughout the asymptotic calculations below the Wilson coefficient
$\gamma$ tends to zero on a fixed finite spatial lattice.  This is strong
bare coupling, not the $\gamma\to\infty$ trajectory needed by the current
continuum programme.  No continuation between the two regimes is assumed.

The calculation isolates the first gauge-invariant kinetic channel and
keeps its magnetic dressing and the actual Perron vector.  For the squared
curvature source the unflowed Haar projection on that channel vanishes.
Both simultaneous spatial flow and the magnetic vacuum restore it at first
order, with a definite interference term in the native covariance.  The
result supplies a signed, fully normalized source calculation, including
an explicit finite-coefficient error bound.  It also identifies why neither
an undressed Haar projection nor a separately flowed local plaquette gives
a justified replacement for the collective source at physical flow time.

\subsection{Fixed-volume scope and the original transfer}

Take an even cubic spatial torus of side at least six.  This excludes
length-four winding loops and ensures that distinct elementary plaquettes
share at most one stored link.  There are $M$ spatial links and $N_p$
spatial plaquettes.  All scalar products in this section, unless marked
otherwise, use normalized product Haar measure $dm$ on the \emph{full}
spatial slice; native covariances always use the stationary Wilson law.
Let $\mathcal H_{\rm inv}\subset L^2(dm)$ be the vertex-gauge-invariant
space.  Define
\begin{equation}
 W=\sum_p w(U_p),\quad b_\gamma=e^{\gamma W/2},\quad
 T_\gamma=b_\gamma C_\gamma^{\otimes M}b_\gamma,\quad
 T_\gamma\phi_\gamma=\rho_\gamma\phi_\gamma,\quad
 \|\phi_\gamma\|_2=1.
 \label{f16v83:setup}
\end{equation}
Compared with V78's $b=e^{-\gamma\sum_p(1-w(U_p))/2}$, this only removes
the known scalar $e^{-\gamma N_p}$ from the transfer.  Its normalized
physical operator and its native density $d\pi_\gamma=\phi_\gamma^2dm$
are unchanged.  No magnetic interaction is removed.  Restriction to
$\mathcal H_{\rm inv}$ is essential in the kinetic statements below.

Strong-coupling transfer and glueball expansions are established methods;
see G.~M\"unster, \emph{Strong coupling expansions for the mass gap in
lattice gauge theories}, Nucl.\ Phys.\ B \textbf{190} (1981), 439--453,
\href{https://doi.org/10.1016/0550-3213(81)90570-8}
{doi:10.1016/0550-3213(81)90570-8}.  The leading plaquette channel is not
claimed as a new confinement or gap mechanism.  Spatial flow as an
operator construction is likewise established; see Sakai and Sasaki,
\href{https://arxiv.org/abs/2211.15176}{arXiv:2211.15176}.
The simultaneous-source overlap, interference coefficient, and explicit
native comparison used here are derived below.  No published numerical
spectroscopy result or extrapolation supplies a premise of the proofs.

Put $\chi_p=2w(U_p)$, let $\Pi_0=|1\rangle\langle1|$, and let $Q$ be the
Haar orthogonal projection onto the span of all $\chi_p$.  These $N_p$
functions are orthonormal: each has squared norm one, and a link private
to either of two distinct plaquettes makes their cross integral zero.
Write
\begin{equation}
 r_\gamma=I_2(\gamma)/I_1(\gamma),\qquad c_\gamma=r_\gamma^4.
 \label{f16v83:kinetic-scale}
\end{equation}
The symbols $r_\gamma,c_\gamma$ here are not V79's derivative-bound
constants.  In particular $r_\gamma=\gamma/4+O(\gamma^3)$.

\begin{proposition}[The first invariant kinetic sector and its remainder]
\label{f16v83:sector}
On $\mathcal H_{\rm inv}$ the exact kinetic convolution has the decomposition
\begin{equation}
 C_\gamma^{\otimes M}=\Pi_0+c_\gamma Q+R_\gamma^{\rm kin},\qquad
 0\preceq R_\gamma^{\rm kin}\preceq
       r_\gamma^6(I-\Pi_0-Q).
 \label{f16v83:kinetic-split}
\end{equation}
Moreover $QWQ=0$, where the middle $W$ denotes multiplication.
\end{proposition}
\begin{proof}
In a product Peter--Weyl sector with link labels $k_e$, the multiplier is
$\prod_e\lambda_{k_e}(\gamma)$, where
$\lambda_k=I_{k+1}/I_1\le r_\gamma^k$ by V72.  A nontrivial invariant
sector cannot have a leaf in its nonzero-label support: integration of the
gauge action at that vertex would leave a single nontrivial irreducible
representation with no invariant vector.  A support with at most four
units of total label must therefore contain a cycle.  The graph has no
cycles shorter than four, and its length-four cycles are exactly the
elementary plaquettes.  At total label four every edge of such a cycle has
label one; its intertwiner space is one dimensional and yields $\chi_p$.
There are no total-label-five invariant sectors.  Indeed the graph is
bipartite and the central gauge transformation at each black vertex makes
the sum of its incident labels even; adding those sums counts each edge
label once.  Thus all remaining nonconstant sectors have total label at
least six.  Positivity and the multiplier bound prove
\eqref{f16v83:kinetic-split}, including its exact complement.

Finally $\langle\chi_p,W\chi_q\rangle$ is a sum of three-plaquette
fundamental character integrals.  If three plaquettes are distinct, one
of them has an edge absent from the other two.  If some coincide, an odd
remaining plaquette still has a link of odd parity.  Haar invariance under
$U_e\mapsto-U_e$ makes each integral zero.  This proves $QWQ=0$.
\end{proof}

\subsection{An explicit native covariance expansion, uniform in the output flow}

Define the vacuum-subtracted operator on the same invariant Haar space,
\begin{equation}
 L_\gamma=T_\gamma/\rho_\gamma-
                  |\phi_\gamma\rangle\langle\phi_\gamma|\succeq0.
 \label{f16v83:L}
\end{equation}
It vanishes on $\phi_\gamma$.  For bounded invariant slice functions $F_i$,
subtract any chosen constants and put
$B_F^2=\sum_i\|F_i\|_\infty^2$.  The constants need not be native means.
Writing $A_\gamma c=\phi_\gamma\sum_i c_iF_i$, exact native centering gives
\begin{equation}
 C_F(n)=A_\gamma^*L_\gamma^n A_\gamma\quad(n\ge1),\qquad
 H_F=\left(\int\chi_p F_i\,dm\right)_{p,i}^{\!*}
       \left(\int\chi_p F_i\,dm\right)_{p,i}.
 \label{f16v83:source-C}
\end{equation}
In particular $H_F$ is a coefficient matrix, not an assumed native Gram.

\begin{theorem}[Quantified full-native channel calculation]
\label{f16v83:quantified}
Set $x=\gamma N_p$, $e=e^x c_\gamma$, and assume $e\le1/4$.  Put
\begin{align}
 \eta_\gamma&=4(e^x-1)+e^xr_\gamma^2+6e^{2x}c_\gamma,\nonumber\\
 \xi_\gamma&=2(e^{x/2}-1)+2e,\qquad
 E_{\gamma,n}=n\eta_\gamma(1+\eta_\gamma)^{n-1}+2\xi_\gamma.
 \label{f16v83:error-constants}
\end{align}
Then for every integer $n\ge1$,
\begin{equation}
 \boxed{\quad
 \|L_\gamma/c_\gamma-Q\|\le\eta_\gamma,\qquad
 \|c_\gamma^{-n}C_F(n)-H_F\|\le B_F^2 E_{\gamma,n}.
 \quad}
 \label{f16v83:native-expansion}
\end{equation}
The bound holds for the actual simultaneous-flow output at every finite
flow time, since only its supremum bounds enter.  Its constants retain
$N_p$; they are not asserted uniform in spatial volume.
\end{theorem}
\begin{proof}
Let $\beta=\|b_\gamma\|_2^2\ge1$, $s=b_\gamma/\sqrt\beta$, and
$S=|s\rangle\langle s|$.  The inequality $\beta\ge1$ is Jensen with
$\int Wdm=0$.  The exact rank-one splitting is
$T_\gamma=\beta S+E$, with $0\preceq E$ and $\|E\|\le e$.
Thus $\beta\le\rho_\gamma\le\beta+e$.
Projecting the Perron equation off $s$ gives
\[
 \|\Pi_\phi-S\|\le\frac{e/\beta}{1-e/\beta},\qquad
 \|\phi_\gamma-s\|_2\le2e/\beta\le2e.
\]
Here $\langle\phi_\gamma,s\rangle>0$ fixes the sign.  To verify the first
bound, solve the projected equation on $s^\perp$; its inverse has norm at
most $(\rho_\gamma-e)^{-1}$.  The second follows from the elementary
angle formula for normalized vectors and $e/\beta\le1/4$.

Compress $T_\gamma$ to $\phi_\gamma^\perp$.  Replacing that projection
by $I-S$, and $\rho_\gamma$ by $\beta$, costs at most
$6(e/\beta)^2$ in norm: the rank-one term costs the squared projection
distance, the two $E$ projection errors cost at most twice that distance
times $e/\beta$, and the denominator error costs $(e/\beta)^2$.
Their sum is at most $(7e/(3\beta))^2<6(e/\beta)^2$.
The remainder in \eqref{f16v83:kinetic-split} costs at most
$e^xr_\gamma^6/\beta$.  Consequently
\[
 \left\|L_\gamma-\frac{c_\gamma}{\beta}
            (I-S)b_\gamma Qb_\gamma(I-S)\right\|
 \le e^xr_\gamma^6+6e^{2x}c_\gamma^2.
\]
Let $t=e^{x/2}-1$.  Since $1\le\beta\le e^x$, one has
$\|s-1\|_2\le2t$ and
$\|(I-S)b_\gamma Q-Q\|\le3t$.  Multiplying these two factors and
accounting for $\beta^{-1}$ bounds the remaining difference from $Q$
by $3t(2+t)+(e^x-1)=4(e^x-1)$.
This proves the first inequality and $\|\phi_\gamma-1\|_2\le\xi_\gamma$.

For any positive operator $B$ with $\|B-Q\|\le\eta$, the telescoping
identity for $B^n-Q^n$ gives
$\|B^n-Q\|\le n\eta(1+\eta)^{n-1}$.  Both source syntheses,
$A_\gamma$ and $A_0c=\sum_i c_iF_i$, have norm at most $B_F$ and
$\|A_\gamma-A_0\|\le B_F\xi_\gamma$.  Insert these in
\eqref{f16v83:source-C} and use one error factor in each of the two
remaining quadratic-form terms.  This proves the second inequality.
All centering is effected by the exact Perron projection in $L_\gamma$.
\end{proof}

The theorem already computes a normalized native covariance, not just a
finite-rank approximation of the observable.  It does not say the Haar
coefficient matrix is nonzero.  In particular it vanishes for the
unflowed squared-curvature bank.  Resolving that zero requires the next
source and magnetic terms together.

\subsection{The simultaneous spatial flow creates a definite first channel}

Use dimensionless spatial flow time $u=\tau/a^2$ and the full vector field
$X=-\nabla\sum_p(1-w(U_p))=\nabla W$ in the unit-quaternion product metric.
Let $\Phi_u$ denote this full flow and put
\begin{equation}
 A_p=1-w(U_p)^2=\frac34-\frac14\chi_2(U_p),\qquad
 F_i^u=\sum_{p\in\mathcal P}s_{ip}A_p(\Phi_u U),\qquad
 \mathcal S_{ij}=\sum_p s_{ip}s_{jp}.
 \label{f16v83:flow-bank}
\end{equation}
There are $m=|\mathcal P|$ distinct marks, with signed coefficients allowed.
The source in V82 is the fixed-orientation case, multiplied by
$\zeta_a^2/(2a)$.  This scalar and any declared coefficient maps multiply
the resulting Gram by congruence; they do not change the calculation.

\begin{proposition}[An exact collective first-flow coefficient]
\label{f16v83:flow-jet}
For all spatial plaquettes $p,r$,
\begin{equation}
 \int A_p\chi_r\,dm=0,\qquad
 \int (XA_p)\chi_r\,dm=-2\mathbf1_{p=r},\qquad
 \|A_p\circ\Phi_u-A_p-uXA_p\|_\infty\le640u^2.
 \label{f16v83:flow-jet-identities}
\end{equation}
No neighboring plaquette force is omitted.  In particular, for
$0<u\le(640\sqrt m)^{-1}$,
\begin{equation}
 u^2\mathcal S\preceq H_{F^u}\preceq9u^2\mathcal S.
 \label{f16v83:flow-overlap}
\end{equation}
\end{proposition}
\begin{proof}
The first identity follows from the center parity of each link.
Expand $X=\sum_q\nabla w_q\cdot\nabla$.  Since $A_p$ is even under
all individual link-center changes, the integral of
$\chi_r\nabla w_q\cdot\nabla A_p$ vanishes unless $q=r$; distinct
plaquette boundaries have a link of unmatched parity.  For $q=r$,
Haar integration by parts and the product-sphere Laplacian give
\[
 \int \chi_q\nabla w_q\cdot\nabla A_p\,dm
 =\int A_p\{24w_q^2-2|\nabla w_q|^2\}\,dm
 =8\int A_p\chi_2(U_q)\,dm=-2\mathbf1_{p=q}.
\]
We used $\Delta w_q=-12w_q$ and
$|\nabla w_q|^2=4(1-w_q^2)$.  Distinct adjoint plaquette characters
are Haar orthogonal by integration of a private link.

For the quantitative remainder, every link force has norm at most four.
On the four links of $p$, $\|X\|^2\le64$.
The plaquette trace Hessian has operator norm at most four; hence
$\|\operatorname{Hess}A_p\|\le16$, and each link gradient of $A_p$ has
norm at most one.  The full action Hessian has block row norm at most
sixteen, so $|(\nabla_XX)_e|\le64$.
It follows that $|X^2A_p|\le16\cdot64+4\cdot64=1280$ at every compact
configuration.  Taylor integration along the actual flow proves the
remainder.  For a coefficient direction, let
$s_p=\sum_i c_i s_{ip}$.  Projection onto $Q$ is an $L^2$ contraction,
and $\sum_p|s_p|\le\sqrt m(\sum_p|s_p|^2)^{1/2}$.  Thus the projected
remainder has norm at most $640u^2\sqrt m\|s\|_2$, whereas the first
coefficient has norm $2u\|s\|_2$.  This proves both sides of
\eqref{f16v83:flow-overlap}.
\end{proof}

Combining the scalar-direction version of
Theorem~\ref{f16v83:quantified} with this proposition gives a populated,
finite-parameter certificate.  The half-range of $\sum_p s_p A_p(\Phi_uU)$
is at most $\sum_p|s_p|/2$, so whenever $e\le1/4$,
$0<u\le(640\sqrt m)^{-1}$ and $E_{\gamma,n}\le2u^2/m$,
\begin{equation}
 \boxed{\quad
 \frac12u^2c_\gamma^n\mathcal S
 \preceq C_{F^u}(n)
 \preceq\frac{19}{2}u^2c_\gamma^n\mathcal S.
 \quad}
 \label{f16v83:populated-certificate}
\end{equation}
Both sides concern the full native stationary law.  For fixed lattice,
$n$ and such $u>0$, the conditions hold for sufficiently small $\gamma$,
with an explicit test in \eqref{f16v83:error-constants}.
This is not a claim that the bounds remain uniform as $u=\tau/a^2$
diverges at fixed positive physical flow time.

\subsection{Magnetic dressing and flow must be combined before squaring}

A vanishing unflowed Haar projection does not imply absence of the same
channel in the native covariance.  The following joint limit evaluates
its first nonzero term, with the source and the vacuum both changing.

\begin{theorem}[Native magnetic--flow interference in the covariance]
\label{f16v83:mixed-limit}
Fix the spatial torus and a finite coefficient bank in
\eqref{f16v83:flow-bank}.  Let $\gamma\downarrow0$ and $u_\gamma\ge0$
with $u_\gamma/\gamma\to\kappa\in[0,\infty)$.
For every integer sequence $n_\gamma\ge2$ satisfying
$n_\gamma\gamma^2\to0$,
\begin{equation}
 \boxed{\quad
 \gamma^{-2}c_\gamma^{-n_\gamma}
 C_{F^{u_\gamma}}(n_\gamma)
 \longrightarrow
 \left(2\kappa+\frac18\right)^2\mathcal S.
 \quad}
 \label{f16v83:mixed-main}
\end{equation}
In particular the result holds at each fixed integer lag $n\ge2$.
The normalization and the limit retain the full magnetic transfer and
its actual Perron projection.  The statement is not volume-uniform.
\end{theorem}
\begin{proof}
All operators are first restricted to $\mathcal H_{\rm inv}$.
The compact kernel and magnetic multiplier are analytic in $\gamma$
near zero in operator norm.  The isolated eigenvalue one of $\Pi_0$
therefore has an analytic Perron projection and normalized eigenvector:
this follows directly by a resolvent Neumann series on a fixed circle
about one.  The estimates already proved imply that
$L_\gamma=O(\gamma^4)$; consequently its analytic series is divisible
by $c_\gamma=\gamma^4/256+O(\gamma^6)$.
Let $B_\gamma=L_\gamma/c_\gamma$, continued at zero by $Q$.

Here are the needed first derivatives, including the vacuum motion.
Writing $|W\rangle$ for the Haar function $W$, one has
\[
 b_\gamma=I+\tfrac\gamma2M_W+O(\gamma^2),\quad
 \beta=1+O(\gamma^2),\quad
 \phi_\gamma=1+\tfrac\gamma2W+O_{L^2}(\gamma^2).
\]
In the proof of Theorem~\ref{f16v83:quantified}, the difference between
$B_\gamma$ and
$\beta^{-1}(I-S)b_\gamma Qb_\gamma(I-S)$ is $O(\gamma^2)$.
Thus, with $A=(M_W/2)Q-|1\rangle\langle W|$,
\begin{equation}
 B_\gamma=Q+\gamma(A+A^*)+O(\gamma^2),\qquad
 QA=0,
 \label{f16v83:B-first}
\end{equation}
where $QA=0$ uses $QWQ=0$.

Let $E_\gamma$ be the spectral projection of $B_\gamma$ on the circle
$|z-1|=1/2$.  For small $\gamma$ this has rank $N_p$.
Expanding its resolvent gives
$E_\gamma=Q+\gamma(A+A^*)+O(\gamma^2)$.
The canonical isometry from $\operatorname{Ran}Q$ to this range is
\[
 J_\gamma=E_\gamma Q(QE_\gamma Q)^{-1/2}
          =Q+\gamma A+O(\gamma^2).
\]
Its compressed transfer obeys
$J_\gamma^*B_\gamma J_\gamma=I+O(\gamma^2)$.
The absence of a first-order band eigenvalue change follows from the
vanishing $Q$--$Q$ block in \eqref{f16v83:B-first}; no common eigenbasis
or invariant plaquette source span is assumed.

Subtract the fixed Haar constant $3\sum_p s_p/4$ from each source and
write $g_0=-\sum_p s_p\chi_2(U_p)/4$ and
$v_\gamma=\phi_\gamma(F^{u_\gamma}-3\sum_p s_p/4)$.
Since $Qg_0=0$ and $\int g_0dm=0$, the coefficient in the moving band is
\begin{align*}
 J_\gamma^*v_\gamma
 &=u_\gamma QXF^0
   +\gamma\{\tfrac12QWg_0+A^*g_0\}
   +o_{L^2}(\gamma)\\
 &=u_\gamma QXF^0+\gamma QWg_0+o_{L^2}(\gamma)\\
 &=-\left(2u_\gamma+\frac\gamma8\right)
                  \sum_p s_p\chi_p+o_{L^2}(\gamma).
\end{align*}
The first equality uses the uniform flow Taylor remainder above.
The last one also uses the exact Haar identity
\[
 \int\chi_r\chi_q\chi_2(U_p)\,dm=\mathbf1_{p=q=r}.
\]
If the plaquettes do not all coincide, a private link carries either a
single nontrivial representation or unmatched center parity; its integral
is zero.  In the coincident case $\chi_1^2=1+\chi_2$ gives the value one.
The two magnetic halves, from $\phi_\gamma$ and $J_\gamma$, are both
necessary for the coefficient $1/8$.

There is a uniform bound on the remainder of the spectrum, not just an
expansion in a guessed band.  By \eqref{f16v83:kinetic-split},
$T_\gamma$ is the sum of an operator of rank at most $1+N_p$ and a
positive remainder of norm at most $e^xr_\gamma^6$.
Min--max bounds every remaining normalized eigenvalue by
$e^xr_\gamma^6/\rho_\gamma$.
After division by $c_\gamma$, the complement of $E_\gamma$ is bounded
by $e^xr_\gamma^2$ (and includes the zero vacuum eigenvalue).
Its contribution divided by $\gamma^2$ is therefore
$O((e^xr_\gamma^2)^{n_\gamma}/\gamma^2)=o(1)$ for $n_\gamma\ge2$.
Inside the band,
$[J_\gamma^*B_\gamma J_\gamma]^{n_\gamma}=I+o(1)$ because
$n_\gamma\gamma^2\to0$.
The evaluated overlap and Haar orthonormality of the $\chi_p$ now prove
\eqref{f16v83:mixed-main} in every coefficient direction and hence as a
finite matrix.  All constants in these analytic remainders may depend
on the fixed lattice and bank.
\end{proof}

The coefficient is
\[
 \left(2\kappa+\tfrac18\right)^2
       =4\kappa^2+\frac\kappa2+\frac1{64}.
\]
Adding separate squared flow and magnetic effects would lose $\kappa/2$.
At zero flow the native leading covariance is
$\gamma^2c_\gamma^n\mathcal S/64$, not the undressed kinetic answer.
Indeed in the distinct reference $b=1$, $\phi=1$, the unflowed bank is a
sum of adjoint plaquette characters and has exactly
$\lambda_2(\gamma)^{4n}\mathcal S/16$.  Its leading power is $\gamma^{8n}$,
rather than the native $\gamma^{4n+2}$ obtained here for $n\ge2$.
The difference is a concrete source-weighted magnetic effect; it is not
permission to replace either law by the other.

\subsection{What is computed, and why it is not the physical anchor theorem}

The current calculation includes the complete simultaneous vector field,
a native infinite-time Perron limit at fixed spatial regulator, its
magnetic projection, and a source-dependent upper and lower covariance
comparison.  Its main mixed limit is new work on these stated objects,
not a use of V82's secant conclusion as an assumption.  The classical
strong-coupling gap mechanism by itself is not the advance claimed here.

For clarity an elementary exact bound identifies the clock limitation.
Since $T_\gamma=\beta S+E$ with $\|E\|\le e^xc_\gamma$ and
$\rho_\gamma\ge\beta\ge1$, min--max gives on the nonvacuum invariant space
\begin{equation}
 \|L_\gamma\|\le e^xc_\gamma=:s_\gamma.
 \label{f16v83:strong-upper}
\end{equation}
Whenever $s_\gamma<1$, every bounded invariant source satisfies
\begin{equation}
 0\preceq C_F(n+d)\preceq s_\gamma^d C_F(n),\qquad n\ge0,\ d\ge1.
 \label{f16v83:clock-test}
\end{equation}
At $n=0$ use the exact centered Gram, not $L_\gamma^0$ on uncentered
functions.  At larger $n$ the inequality is scalar spectral calculus on
the original transfer.  If along a shrinking-spacing sequence this
particular sufficient bound remains $s_\gamma\le s_*<1$, two positive
times separated by $d_aa\to h>0$ cannot obey V82's cutoff-uniform earlier
upper and later positive lower bounds in one common normalization:
\eqref{f16v83:clock-test} would force the later Gram to zero relative to
the earlier one.  This is the usual wrong-clock issue in a controlled
strong-coupling regime, not a statement about $\gamma\to\infty$.

More importantly, the physical source uses $u_a=\tau/a^2\to\infty$,
whereas the explicit mixed coefficient above uses $u_\gamma=O(\gamma)$
as $\gamma\downarrow0$.  The derivative-free estimate
\eqref{f16v83:native-expansion} remains valid for arbitrary finite $u$,
but then the full coefficient matrix $H_{F^u}$ has not been evaluated.
Its value cannot be obtained by extending the first jet $-2u$ to large
$u$, and the factor $x=\gamma N_p$ prevents a claim of thermodynamic
uniformity for the displayed expansion.  The initial Gaussian/weak-field
and exponential physical trajectory results in other parts of the
manuscript are not spliced into this opposite-coupling limit.

The next estimate in the physical programme remains the native collective
covariance at positive physical flow time and weak bare coupling.  This
section contributes a fully paid test for that calculation: its spectral
channel and its source overlap must be transported \emph{together}, with
the Perron motion included.  Further terms of the present strong-coupling
series, by themselves, would not bridge that missing regime.

\subsection{Verification, preservation and result ledger}

The companion diagnostic is
\begin{center}\small
\path{verify_ym_f16_v83_channel_mixing.py}.
\end{center}
It distinguishes exact Haar and character identities, literal spatial
cycle and plaquette-incidence checks, compact simultaneous-flow derivatives,
finite positive-matrix tests of the rank-one remainder estimate, and a
separately labelled compact one-plaquette magnetic-transfer quadrature.
That last graph checks the same local interference coefficient but is not
a simulation of the three-dimensional spatial torus or four-dimensional
Wilson vacuum.  Numerical cutoffs and resolution comparisons are reported,
not presented as interval enclosures.  The analytic proofs above do not
use fitted constants or numerical outcomes.  Removing this delimited
appendix and reverting the single displayed version change recovers V82
byte for byte.

\clearpage
\begin{record}[V83 result ledger]
\label{f16v83:ledger}
\emph{Proved here:} the specified invariant-channel remainder and explicit
native covariance error; a populated small-coefficient covariance sandwich
for the actual simultaneous-flow source; its exact collective first-flow
coefficient and uniform second-flow remainder; the native magnetic--flow
interference limit with coefficient $(2\kappa+1/8)^2$ and a controlled
nonvacuum spectral complement; its distinction from the undressed kinetic
reference; and the stated strong-coupling clock limitation.

\emph{Inherited or standard:} the original positive Wilson transfer,
compact representation multipliers, simultaneous spatial Wilson flow,
Haar character calculus, elementary isolated-spectral-band perturbation,
and the classical strong-coupling plaquette mechanism.  These do not
supply a weak-coupling continuum covariance theorem.

\emph{Not asserted:} a small-$\gamma$ calculation at $\gamma\to\infty$;
a uniform large-volume or fixed-physical-flow remainder from the first
jet; a gauge-invariant low-degree sector on the unprojected charged
space; a polynomial plaquette span invariant under the magnetic transfer;
independent squared source contributions; or analytic continuation of
the two positive-time bounds to the intended physical trajectory.

\emph{Still unproved:} V82's cutoff-uniform positive-time native Gram
bounds for the physically normalized collective bank, its spatial and
thermodynamic/continuum control, common dense-core physical-time upper
contraction on the intended trajectory, and the interacting
four-dimensional continuum Yang--Mills mass gap.
\end{record}

\begin{firewall}
This is a direct source calculation under the full native magnetic
transfer in its explicitly stated strong-coupling regime.  Magnetic
vacuum motion and simultaneous flow interfere before squaring.  Its
finite-regulator result and clock are not relabelled as a continuum
construction.  All V1--V82 mathematical text remains unchanged.
\end{firewall}
% END F16 V83 APPEND

% BEGIN F16 V84 APPEND -- 2026-09-19
% Append-only continuation of the Post-Fifteenth-Archive V83.
\clearpage
\section[V84: forward-source completion and hidden low-energy weight]{V84: forward-source completion\\and hidden low-energy weight}
\label{f16v84:context}

V83 computes the simultaneous-source covariance at small Wilson coefficient,
not on the intended weak-bare-coupling physical trajectory.  We do not add
further coefficients to that opposite-regime expansion.  Instead we examine
the remaining spectral-coverage question: what does a finite collection of
positive-time source Grams actually certify about a lower energy edge?
A strict contraction on a finite source space need not persist after that
space is completed by forward time translates.  There is an exact native
obstruction, not just a numerical conditioning concern: two smooth invariant
sources of the same finite Wilson transfer can have identical covariances
at any prescribed finite list of times and different lowest visible energies.
The transfer and its full mass gap are the same in that construction.

The constructive result is a forward-source hierarchy with an explicit
bound on the low-energy weight left unresolved at finite depth.  A Chebyshev
polynomial turns a finite matrix inequality into exponential suppression
in the temporal depth, with its numerical and reconstruction errors retained.
V77's genuine simultaneous-flow reconstruction pays those errors at
logarithmically growing depth by a polynomial representation cutoff.
No required inequality for that bank is assumed to have been verified.
The native positive-time lower Grams of V82 remain a separate open input.

\subsection{Scope and the exact native finite-record obstruction}

At each finite spatial regulator retain the original gauge-invariant Wilson
Hilbert space, stationary measure $\pi$, centered positive transfer $P$,
and physical Hamiltonian $H=-a^{-1}\log P$.  On the centered invariant
space $P$ is injective and compact, with infinitely many positive
eigenvalues tending to zero; its eigenfunctions are smooth and bounded.
These properties follow from the inherited positive convolution/magnetic
sandwich: convolution has positive multipliers in every representation,
the invariant space is infinite dimensional, and its compact smooth kernel
maps each nonzero-eigenvalue eigenfunction to a smooth function.  The
Perron eigenvector is removed only by native centering.

Generalized eigenvalue and matrix-polynomial methods for lattice correlators
are established; see Blossier et al.,
\href{https://arxiv.org/abs/0902.1265}{arXiv:0902.1265}, and Fleming,
\href{https://arxiv.org/abs/2309.05111}{arXiv:2309.05111}.
Finite-rank termination is related to the moment flat-extension principle;
see Laurent and Mourrain,
\href{https://arxiv.org/abs/0812.2563}{arXiv:0812.2563}.
The elementary operator proofs needed here are included.  These methods,
Hankel positivity, and Rayleigh--Ritz monotonicity are not claimed as new.
V57's Krylov resolvent estimate and V81's two-vector feature block are not
reproved as source-support theorems.  The additional claims here concern
native finite-record ambiguity, the quantitative unresolved spectral weight,
and a physical-depth reconstruction budget.

\begin{theorem}[Exact finite records can hide a lower native source energy]
\label{f16v84:alias}
Fix any finite native regulator with $\gamma>0$, and any integer $J\ge0$.
There are two real smooth bounded invariant centered sources $f_0,f_1$,
each of native variance one, such that
\begin{equation}
 \langle f_0,P^nf_0\rangle_\pi
 =\langle f_1,P^nf_1\rangle_\pi\quad(0\le n\le J),
 \label{f16v84:alias-moments}
\end{equation}
but the infima of the physical spectral supports of these two source
vectors are different.  The sources and their spectral weights are given
explicitly in terms of finitely many native eigenpairs below.
\end{theorem}
\begin{proof}
Choose distinct nonvacuum eigenvalues
$1>s>x_1>\cdots>x_{J+1}>0$ and real orthonormal invariant eigenfunctions
$\psi_*,\psi_1,\ldots,\psi_{J+1}$.  Compact injectivity on an infinite
dimensional invariant space supplies as many distinct eigenvalues as needed.
Choose $w_i>0$ with $\sum_iw_i=1$ and set
\[
 \ell_i(s)=\prod_{j\ne i}\frac{s-x_j}{x_i-x_j}.
\]
Polynomial interpolation gives
$\sum_i\ell_i(s)x_i^n=s^n$ for $0\le n\le J$, including
$\sum_i\ell_i(s)=1$.  At least one $\ell_i(s)$ is positive.  Choose
\[
 0<\epsilon<\min_{\ell_i(s)>0}\frac{w_i}{\ell_i(s)},\qquad
 w_i'=w_i-\epsilon\ell_i(s)>0.
\]
Then
\begin{equation}
 f_0=\sum_i\sqrt{w_i}\psi_i,\qquad
 f_1=\sqrt\epsilon\,\psi_*+\sum_i\sqrt{w_i'}\psi_i
 \label{f16v84:alias-sources}
\end{equation}
are centered and have norm one.  Their moment difference is
$\epsilon[s^n-\sum_i\ell_i(s)x_i^n]=0$ in the stated range.
Their lowest source energies are $-a^{-1}\log x_1$ and
$-a^{-1}\log s$, respectively.  Each is a finite linear combination of
smooth bounded eigenfunctions.  All statements concern one and the same
native law and transfer, not two hypothetical Wilson measures.
\end{proof}

A finite list of nonconsecutive integer times is covered by taking $J$ to
be its largest member.  The theorem changes the source; it does not assert
that the particular simultaneous-flow formula has this ambiguity once all
of its additional structure is used.  The constructed sources may depend
on the entire finite slice, and their supremum norms are not claimed
uniform in the regulator.  These records show that finite-horizon contraction alone is not a general
certificate of source spectral exhaustion.  Exact termination, treated
below, is a distinct sufficient condition.  The theorem does not change
or disprove the mass gap of the underlying fixed transfer.

\begin{proposition}[An exact compact four-link record with rational weights]
\label{f16v84:compact-alias}
In the compact Poisson reference of V82 choose each link's parameter
$r=2^{-1/4}$, so $P\chi_j(Q)=2^{-j}\chi_j(Q)$ for the cyclic product
$Q=X_1X_2X_3X_4$.  Take
\[
 f_0=\tfrac12(\chi_2+\chi_3+\chi_4+\chi_5),\qquad
 f_1=\sqrt{\tfrac1{1000}}\chi_1+\sqrt{\tfrac{47}{200}}\chi_2
     +\sqrt{\tfrac8{25}}\chi_3+\sqrt{\tfrac{13}{100}}\chi_4
     +\sqrt{\tfrac{157}{500}}\chi_5.
\]
Both source records at $n=0,1,2,3$ are exactly
\begin{equation}
       1,\quad\frac{15}{128},\quad\frac{85}{4096},\quad
                       \frac{585}{131072}.
 \label{f16v84:rational-record}
\end{equation}
Their largest visible transfer multipliers are $1/4$ and $1/2$.
\end{proposition}
\begin{proof}
Haar character orthogonality converts the squared coefficients to spectral
weights.  For $s=1/2$ and $x=(1/4,1/8,1/16,1/32)$ the interpolation
weights in the preceding proof are $(15,-70,120,-64)$.  Subtract
$10^{-3}$ times these weights from $(1/4,1/4,1/4,1/4)$.
The displayed record follows by rational arithmetic.  The reference is a
positive compact gauge-invariant convolution model, not the full Wilson
magnetic law; the preceding theorem supplies the separate native statement.
\end{proof}

\subsection{Forward completion and the direction of the spectral bound}

Return to an arbitrary finite centered native source synthesis
$V:\mathbb C^b\to L^2(\pi)$.  Choose a buffer $k\ge1$ and a block lag
$d\ge1$, and set
\begin{equation}
 U=P^kV,\qquad T=P^d,\qquad
 M_j=U^*T^jU=C_F(2k+jd),\qquad
 \mathbf M(B)=U^*\mathbf1_B(T)U.
 \label{f16v84:moments}
\end{equation}
All moments are covariances of the same original source.  The block lag has
physical duration $da$, not a newly assigned clock.  Work on the declared
quotient of $M_0$; higher moment blocks can have further exact dependencies.
For $r\ge0$ put
\begin{equation}
 \mathsf H_r=(M_{i+j})_{i,j=0}^r,\qquad
 \mathsf K_r=(M_{i+j+1})_{i,j=0}^r,\qquad
 \mathcal X_r=\operatorname{span}\{T^jUc:0\le j\le r\}.
 \label{f16v84:Hankel}
\end{equation}
The names $\mathsf H_r,\mathsf K_r$ denote finite matrices, not the physical
Hamiltonian or an auxiliary spatial heat operator.

\begin{theorem}[Exact upper contraction requires forward-source completion]
\label{f16v84:completion}
Let $\mathcal X=\overline{\bigcup_r\mathcal X_r}$ and define
\[
 q_r=\sup_{z^*\mathsf H_rz>0}
        \frac{z^*\mathsf K_rz}{z^*\mathsf H_rz}.
\]
Then $0\preceq\mathsf K_r\preceq\mathsf H_r$ and
\begin{equation}
       q_r\uparrow q_\infty=\|T|_{\mathcal X}\|.
 \label{f16v84:Ritz}
\end{equation}
For $0\le q\le1$ the following are equivalent:
\begin{equation}
 \mathsf K_r\preceq q\mathsf H_r\ \hbox{for every }r;\qquad
 T|_{\mathcal X}\preceq qI;\qquad
 \mathbf M((q,1])=0.
 \label{f16v84:equivalent}
\end{equation}
In particular $-(da)^{-1}\log q_r$ is an \emph{upper} variational bound
on the lowest energy visible in this cyclic source space, not a lower
mass bound.  One finite-depth strict contraction is not spectral exhaustion.
\end{theorem}
\begin{proof}
Let $X_r(z_0,\ldots,z_r)=\sum_{j=0}^rT^jUz_j$.
Exactly $\mathsf H_r=X_r^*X_r$ and $\mathsf K_r=X_r^*TX_r$.
The Rayleigh quotient is therefore the largest compression value on
$\mathcal X_r$.  These spaces increase, and their union is dense in
$\mathcal X$.  This proves monotonicity and its limit.  Since $T$ is
bounded and self-adjoint, the closed invariant space $\mathcal X$ is
reducing.  A quadratic-form bound on its dense polynomial union extends
to it.  Conversely a spectral-support bound gives all the finite
inequalities.  Finally $\mathbf M((q,1])=0$ means that its spectral
projection annihilates $U$ and hence every $T^jU$, proving the last
implication.  The energy conclusion follows because $-\log x$ decreases.
No common eigenbasis for the compressed source matrices is assumed.
\end{proof}

The native ambiguity theorem can make the entire pair
$(\mathsf H_r,\mathsf K_r)$ identical while changing $q_\infty$:
choose $J\ge2k+(2r+1)d$.  In the explicit compact record
\eqref{f16v84:rational-record}, even the depth-one unbuffered generalized
problem is identical.  The statement does not contradict V83: its
\emph{full-operator} bound $\|L_\gamma\|\le s_\gamma$ gives
$\mathsf K_r\preceq s_\gamma^d\mathsf H_r$ at every depth in its stated
regime.  That is more information than a single compressed eigenvalue.

\subsection{A finite depth gives a quantitative bound on unresolved soft weight}

Denote the Chebyshev polynomial of the first kind by $\mathrm T_r$.
We use its elementary identities
$\mathrm T_r(\cos\theta)=\cos(r\theta)$ and
$\mathrm T_r(\cosh t)=\cosh(rt)$; see
\href{https://dlmf.nist.gov/18.5.E1}{DLMF 18.5.1}.
For $0<q<1$ define
\begin{equation}
 p_{r,q}(x)=\mathrm T_r(2x/q-1)=\sum_{j=0}^r a_jx^j,
 \quad D_r(q)=\sum_j a_j^2,\quad
 A_r(q)=\mathrm T_r(1+2/q).
 \label{f16v84:Cheb}
\end{equation}
All roots of $p_{r,q}$ are in $(0,q)$, so its coefficient signs alternate
and $\sum_j|a_j|=A_r(q)$.  Hence $D_r(q)\le A_r(q)^2$, also at $r=0$.

\begin{theorem}[Matrix soft-weight bound with a paid finite-depth defect]
\label{f16v84:soft}
Suppose $0<q<s\le1$ and
\begin{equation}
 \mathsf K_r\preceq q\mathsf H_r+
              \eta\,(I_{r+1}\otimes M_0),\qquad \eta\ge0.
 \label{f16v84:localizer}
\end{equation}
Then
\begin{equation}
 \boxed{\quad
 \mathbf M([s,1])\preceq
 \min\left\{1,
 \frac{q+\eta D_r(q)}
 {(s-q)\mathrm T_r(2s/q-1)^2}\right\}M_0.
 \quad}
 \label{f16v84:soft-bound}
\end{equation}
One may replace $D_r(q)$ by the explicit larger value $A_r(q)^2$.
If $\eta=0$, the fraction is at most
\begin{equation}
 \frac{4q}{s-q}
       e^{-2r\operatorname{arcosh}(2s/q-1)}.
 \label{f16v84:soft-rate}
\end{equation}
For $T=P^d$, $s=e^{-daE_0}$, the matrix on the left is exactly the
buffered source weight at physical energies $0\le E\le E_0$.
\end{theorem}
\begin{proof}
Insert the coefficient vector $(a_0c,\ldots,a_rc)$ into
\eqref{f16v84:localizer}.  In matrix form it yields
\[
 \int(q-x)p_{r,q}(x)^2\,d\mathbf M(x)
                    \succeq-\eta D_r(q)M_0.
\]
On $[0,q]$ the integrand is at most $q$ because $|p_{r,q}|\le1$.
On $(q,s)$ it is nonpositive.  On $[s,1]$ it is at most
$-(s-q)\mathrm T_r(2s/q-1)^2$, since $\mathrm T_r$ is positive and
nondecreasing beyond one.  Thus the same integral is at most
$qM_0-(s-q)\mathrm T_r(2s/q-1)^2\mathbf M([s,1])$.
Combine the two bounds and use $\mathbf M([s,1])\preceq M_0$.
This is an integral of scalar inequalities against a positive matrix
measure; no entrywise comparison or simultaneous diagonalization is used.
The hyperbolic identity gives $\mathrm T_r(z)\ge
\tfrac12e^{r\operatorname{arcosh}z}$ for $z>1$.
\end{proof}

For a direction carrying a fraction $w>0$ on $[s,1]$, an exact
$\mathsf K_r\preceq q\mathsf H_r$ is impossible once
\begin{equation}
 r>
 \frac{\operatorname{arcosh}\sqrt{\max\{1,q/[w(s-q)]\}}}
      {\operatorname{arcosh}(2s/q-1)}.
 \label{f16v84:detection-depth}
\end{equation}
If $q/[w(s-q)]<1$, depth zero already contradicts the bound.
Thus arbitrarily small spectral weights can require arbitrarily large
depth.  At depth $r$ the latest required covariance is
$C_F(2k+(2r+1)d)$; the time horizon is not suppressed from the cost.
The theorem bounds a fraction, not the existence of an eigenvalue: a
strictly positive bound cannot be reinterpreted as exact absence of soft
spectrum.

\subsection{When a finite termination is genuine}

There is a precise exceptional case in which finite data do exhaust the
cyclic source spectrum.  Put
\[
 \mathsf L_r=(M_{i+j+2})_{i,j=0}^r,\qquad
 \mathsf E_r=\mathsf L_r-\mathsf K_r\mathsf H_r^\dagger\mathsf K_r.
\]
This uses one additional moment.  All pseudoinverses refer to exact ranges.

\begin{proposition}[Exact residual, flatness, and the nonzero-residual distinction]
\label{f16v84:flat}
Let $\Pi_r$ be the orthogonal projection onto $\mathcal X_r$.  Then
\begin{equation}
 \mathsf E_r=X_r^*T(I-\Pi_r)TX_r\succeq0.
 \label{f16v84:residual}
\end{equation}
The following are equivalent:
\[
 \mathsf E_r=0;\qquad T\mathcal X_r\subset\mathcal X_r;\qquad
 \operatorname{rank}\mathsf H_{r+1}
                      =\operatorname{rank}\mathsf H_r.
\]
In this case $\mathcal X=\mathcal X_r$, and its complete transfer spectrum
is the spectrum of
$\mathsf H_r^{\dagger/2}\mathsf K_r\mathsf H_r^{\dagger/2}$
on $\operatorname{Ran}\mathsf H_r$.
A merely small nonzero $\mathsf E_r$ does not give this conclusion.
\end{proposition}
\begin{proof}
The identity $\Pi_r=X_r\mathsf H_r^\dagger X_r^*$ proves
\eqref{f16v84:residual}.  Its vanishing is equivalent to
$(I-\Pi_r)TX_r=0$.  The range of $(X_r,TX_r)$ is precisely
$\mathcal X_{r+1}$, proving the rank equivalence.  Invariance makes all
later translates stay in $\mathcal X_r$, and the isometry
$X_r\mathsf H_r^{\dagger/2}$ identifies the compressed operator with
the full cyclic restriction.  For the last assertion take a scalar
probability measure $(1-\epsilon)\delta_x+\epsilon\delta_s$, $x<s$.
At $r=0$, the residual is $\epsilon(1-\epsilon)(s-x)^2\to0$, while
its largest visible multiplier is always $s$.  These two weights can
also be put on two actual native invariant eigenfunctions.
\end{proof}

Exact flatness is a structural identity, not a numerical rank threshold.
Small nonzero residuals can still give useful \emph{weight} bounds.  For
example, on the exact quotient put
$A=\mathsf H_r^{-1/2}\mathsf K_r\mathsf H_r^{-1/2}$ and
$\beta^2=\|\mathsf H_r^{-1/2}\mathsf E_r\mathsf H_r^{-1/2}\|$.
For $s>q_r$ the elementary Sylvester estimate gives
\begin{equation}
       \mathbf M([s,1])\preceq
       \min\{1,\beta^2/(s-q_r)^2\}M_0.
 \label{f16v84:residual-weight}
\end{equation}
Indeed for the isometry $J=X_r\mathsf H_r^{-1/2}$ write
$TJ=JA+R$, where $\|R\|=\beta$.  With $E_s=\mathbf1_{[s,1]}(T)$,
\[
 E_sJ=\int_0^\infty e^{-tT|_{E_s}}E_sR e^{tA}\,dt,
 \qquad\|E_sJ\|\le\beta/(s-q_r).
\]
Since $\operatorname{Ran}U\subset\operatorname{Ran}J$, this proves
\eqref{f16v84:residual-weight}.  The same formulas on a singular moment
block use its exact range.  This shorting identity is not a new Feshbach
principle; it identifies exactly when a finite source calculation has
actually terminated.

\subsection{Uncertainty and the actual simultaneous-source reconstruction}

Suppose Hermitian estimates $\widehat M_j$ have certified errors
$\|\widehat M_j-M_j\|\le\delta$ for $0\le j\le2r+1$, in a fixed
coefficient metric.  Set the corresponding block matrices
$\widehat{\mathsf H}_r,\widehat{\mathsf K}_r$.  Suppose
\begin{equation}
 \widehat{\mathsf K}_r-q\widehat{\mathsf H}_r\preceq\omega I,
 \quad \omega\ge0,\qquad
 m=\lambda_{\min}(\widehat M_0)-\delta>0.
 \label{f16v84:observed}
\end{equation}

\begin{proposition}[A defect certificate without moment-matrix inversion]
\label{f16v84:noise}
Under \eqref{f16v84:observed}, the hypothesis
\eqref{f16v84:localizer} holds with
\begin{equation}
 \boxed{\quad
 \eta=\frac{\omega+(r+1)(1+q)\delta}{m}.
 \quad}
 \label{f16v84:noise-eta}
\end{equation}
Thus \eqref{f16v84:soft-bound} is valid with all errors included.  A
predetermined source quotient is used; failure of $m>0$ does not permit
silently deleting a small source direction.
\end{proposition}
\begin{proof}
For a block vector $z=(z_i)$ the perturbation of either Hankel matrix
has quadratic form at most
$\delta(\sum_i|z_i|)^2\le(r+1)\delta\sum_i|z_i|^2$ in absolute value.
Hence $\mathsf K_r-q\mathsf H_r\preceq
[\omega+(r+1)(1+q)\delta]I$.
Also $M_0\succeq mI$, so $I_{r+1}\otimes M_0\succeq mI$.
This proves the assertion without an ill-conditioned Hankel inverse.
\end{proof}

For the original simultaneous spatial-flow bank use exactly
\[
 O_i^\tau(U)=\frac{\zeta_a^2}{2a}\sum_x f_i(x)
   |\operatorname{Im}(\Phi_{a,\tau}U)_{12}(x)|^2,
 \qquad O^{\tau,K}=\mathcal A_{K,S}O^\tau.
\]
V77 supplies, simultaneously for every integer separation at least two,
\begin{equation}
 \|C_{O^\tau}(n)-C_{O^{\tau,K}}(n)\|\le
 \mathcal D_{a,K}:=\frac{27\pi^2}{8}
 \frac{\mathcal V_a\gamma_a\zeta_a^4}{a^{11}(K+3)}
                    \sum_i\|f_i\|_{1,a}^2.
 \label{f16v84:source-error}
\end{equation}
It applies to every entry used above because $k\ge1$.  Thus this value,
plus numerical covariance and coefficient errors when present, may be
used for $\delta$.  The common filter retains shared-path singlets and
acts on original inputs, not on a substituted flowed-link law.
The cost of forming temporal blocks is $(r+1)$, while their Chebyshev
combination carries $D_r(q)$; neither factor is hidden.  No native
covariance matrices or their numerical error enclosures are claimed
computed by this substitution.

\begin{theorem}[Logarithmic temporal depth with a polynomial reconstruction budget]
\label{f16v84:continuum-test}
Fix $h,\sigma,c,\Delta>0$, put $q=e^{-h\Delta}$, and let
$k_aa\to\sigma$, $d_aa\to h$, $r_a=\lfloor c\log(1/a)\rfloor$.
Assume in a fixed physical coefficient metric that the buffered Grams
satisfy $mI\preceq M_{a,0}\preceq MI$ with fixed $0<m\le M$.
Suppose estimates obey \eqref{f16v84:observed} with a uniformly positive
certified lower value and $\delta_a,\omega_a\le a^N$, where
\begin{equation}
                 N>2c\operatorname{arcosh}(1+2/q).
 \label{f16v84:precision-condition}
\end{equation}
Then, for every $0\le E_0<\Delta$,
\begin{equation}
 \left\|M_{a,0}^{-1/2}U_a^*
      \mathbf1_{[0,E_0]}(H_a)U_aM_{a,0}^{-1/2}\right\|
                          \longrightarrow0.
 \label{f16v84:low-energy-limit}
\end{equation}
The approximation part of this test is paid for the actual simultaneous
flow by choosing
\begin{equation}
 K_a+3\ge\frac{27\pi^2}{8}\mathcal V_a\gamma_a\zeta_a^4
       \left(\sum_i\|f_i\|_{1,a}^2\right)a^{-(N+11)}.
 \label{f16v84:physical-K}
\end{equation}
Other errors require their own allocated portions of the $a^N$ budget.
The latest physical time is $2k_aa+(2r_a+1)d_aa=O(\log(1/a))$.
\end{theorem}
\begin{proof}
The preceding proposition gives $\eta_a=O((r_a+1)a^N)$.
Since $A_r(q)=\cosh[r\operatorname{arcosh}(1+2/q)]$,
\eqref{f16v84:precision-condition} implies
$\eta_aA_{r_a}(q)^2\to0$.
For $s_a=e^{-d_aaE_0}$, one has $s_a\to e^{-hE_0}>q$.
The denominator in \eqref{f16v84:soft-bound} grows exponentially in
$r_a$, proving \eqref{f16v84:low-energy-limit}.  At $E_0=0$ use $s_a=1$.
The reconstruction statement is direct substitution in
\eqref{f16v84:source-error}; its constant is independent of the largest
integer time and of finite $\tau$.  The physical time horizon follows
from the definitions.  There is no microscopic difference quotient and
no additional inverse-spacing subtraction cost.
\end{proof}

For bounded physical volume and test norms, polynomial normalization, and
$\gamma_a=O(\log(1/a))$, the displayed $K_a$ is polynomial in $a^{-1}$.
The number of collective representation coefficients can still be
exponential in the number of input links; this is not an efficient native
integration algorithm.  More importantly, neither the approximate
upper-contraction inequality in \eqref{f16v84:observed} nor its positive
lower $m$ is supplied by the reconstruction theorem.  The theorem is a
quantitative sufficient test, not a populated weak-coupling instance.

If, separately, an interacting limiting Hilbert space and compatible
positive-time source limits are constructed, with convergence of their
buffered spectral measures, these bounds exclude spectral weight below
$\Delta$ on their cyclic source spaces.  To conclude a gap on the complete
vacuum-orthogonal Hilbert space one needs a collection of such spaces with
dense union and the \emph{same} $\Delta$.  A spectral projection annihilating
a dense union is zero, which proves that final implication.  V82's earlier
positive-time bounds can supply high-energy tightness when established;
neither that construction nor dense source coverage follows from the
moment inequalities alone.  Approximate flatness, an isolated compressed
energy, or a positive finite Gram cannot replace these hypotheses.

\subsection{Next native calculation and result ledger}

The next all-coupling target now has a definite form: on the declared
simultaneous-flow source quotient, estimate the first growing-depth signed
matrix $\mathsf K_r-q\mathsf H_r$ which is not paid by an existing
full-operator bound.  A positive direction is an explicit witness against
the proposed spectral edge; a valid upper inequality gives the quantified
soft-weight bound, not automatic exact absence.  Source-specific structure
or genuine termination can provide additional information, whereas another
bare finite-source contraction cannot.  The two positive-time lower/upper
anchors themselves remain unproved at the desired weak-coupling physical
scale.  This iteration neither extrapolates V83 nor asserts a new native
value for those matrices.

The companion diagnostic is
\begin{center}\small
\path{verify_ym_f16_v84_forward_completion.py}.
\end{center}
The companion separates exact rational matching and compact characters,
finite magnetic references, noncommuting matrix measures, Chebyshev and
noise bounds, and exact/near-flat residuals.  These are diagnostics, not
Wilson-vacuum samples, interval certificates or proof-assistant verification.
Removing the appendix and reversing the title change recovers V83 byte for
byte; the inherited archive is preserved, not recertified.

\begin{record}[V84 result ledger]
\label{f16v84:ledger}
\emph{Proved here on the stated objects:} exact finite-record ambiguity
between smooth invariant sources of the same native finite Wilson transfer;
a compact rational realization; the source-support interpretation and
direction of the forward moment hierarchy; a matrix low-energy-weight bound
with explicit Chebyshev and error costs; exact termination and its residual
criterion; an uncertainty certificate without Hankel inversion; and a
logarithmic-depth, polynomial-reconstruction sufficient test for the
unchanged simultaneous-flow bank.

\emph{Inherited or standard:} the native positive transfer and exact source
quotients; V77's common-source reconstruction and V82's positive-time
interpretation; moment/Rayleigh--Ritz methods, polynomial interpolation,
Chebyshev identities, and finite-rank shorting.  The general methods are
not asserted to be new spectral theory.

\emph{Not asserted:} that the constructed alias sources are the current
flowed bank; a lower mass bound from a finite Ritz energy; exact zero
spectral weight from a small upper bound or a numerical rank threshold;
independence or commutation of covariance matrices; efficient evaluation
of collective coefficients; or a physical gap from nonzero Gram data.

\emph{Still unproved:} the required positive-time native Gram bounds and
growing-depth upper inequalities for the physical simultaneous-flow bank;
its spatial and thermodynamic/continuum compatibility; dense limiting
source coverage with common constants; and the interacting four-dimensional
continuum Yang--Mills mass gap.
\end{record}

\begin{firewall}
The native law, source normalization and physical clock are unchanged.
Finite moments bound unresolved low-energy weight; spectral exhaustion
needs the stated completion or exact termination.  No physical-scale
covariance value is inferred.  V1--V83 remain unchanged.
\end{firewall}
% END F16 V84 APPEND

% BEGIN F16 V85 APPEND -- 2026-09-19
% Append-only continuation of the Post-Fifteenth-Archive V84.
\clearpage
\section[V85: local score witnesses and exact contact removal]{V85: local score witnesses\\and exact contact removal}
\label{f16v85:context}

V84 distinguishes finite covariance information from spectral exhaustion.
Before estimating another temporal localizer, we return to a native quantity
whose current bound loses two powers of the Wilson coefficient: the entrance
Gram of an arbitrarily large plaquette bank.  A pair of link derivatives
isolates one elementary plaquette.  Moving those derivatives onto the full
Wilson density produces a local, mean-zero witness.  Its covariance with a
distant witness is exactly zero by integration by parts, without any mixing
assumption on the measure.  Score moments, rather than worst-case scores,
then give a volume-uniform lower bound of order $(1+\gamma)^{-2}$.

The same witnesses identify what this stronger lower bound actually pays.
They are contact observables: after a positive-time placement they are null
in the reflected quotient.  A specified local source correction removes a
positive part of the ordinary variance while preserving every sufficiently
separated covariance exactly.  Thus a stronger native entrance bound is not
misidentified as a new propagating component.  All results concern the
original Wilson law; none replaces it by a product or Gaussian measure.

\subsection{Native setup, scope, and the local adjoint}

Use a four-dimensional hypercubic Wilson lattice with periodic side lengths
at least four, or a finite open patch with the marked links interior.  Each
stored link meets at most six elementary plaquettes.  There are no
length-two periodic identifications.  The conditional law of any specified
free link set, with all other links fixed, has density
\begin{equation}
 d\mu=Z^{-1}\exp\!\left(\gamma\sum_q W_q\right)dm,
 \qquad W_q=w(U_q),\qquad \gamma\ge0,
 \label{f16v85:law}
\end{equation}
where the sum includes every plaquette meeting the free set, and $dm$ is
normalized product Haar measure.  Additive constant Wilson defects cancel
from this density.  Every selected pair below consists of free links.
The general bounds require no translation symmetry.  Uniform local bounds
pass to the original stationary-cylinder and Gibbs limits at fixed
regulator.  A nonlocal Perron endpoint density is not absorbed into
\eqref{f16v85:law}; condition in the finite Euclidean history first.

Haar integration by parts and lattice Schwinger--Dyson identities are
established tools; for a finite-$N$ lattice treatment see V.~Kazakov and
Z.~Zheng, \emph{Bootstrap for Finite N Lattice Yang--Mills Theory},
\href{https://arxiv.org/abs/2404.16925}{arXiv:2404.16925}, Section 2.3.
We use only the elementary compact integration identity proved below, not
a numerical bootstrap bound or a continuum loop-equation closure.  V79's
private-link entrance theorem and V80's compact mixed-derivative geometry
are comparison inputs; their method is not repeated as the new argument.

Let $D_e^a$, $a=1,2,3$, generate $U_e\mapsto e^{t\mathbf e_a}U_e$, with
the unit-quaternion metric.  They are divergence-free under Haar measure.
Put
\begin{equation}
 s_e^a=D_e^a\log\mu=\gamma\sum_{q\ni e}D_e^aW_q,
 \qquad (D_e^a)^*=-D_e^a-s_e^a.
 \label{f16v85:adjoint}
\end{equation}
The adjoint is taken in the actual conditional $L^2(\mu)$ space, not in
Haar space.  Derivatives on distinct links commute, and so do their
weighted adjoints: both statements follow either by direct differentiation
or by conjugating $-D_e^a$ with the strictly positive density.  No
commutation of different generators on the same link will be used.

\begin{proposition}[Native score moments with no density ratio]
\label{f16v85:score-moments}
Uniformly in every admissible exterior and every free link,
\begin{equation}
 \mathbb E_\mu|s_e|^2\le18\gamma,
 \qquad \mathbb E_\mu|s_e|^4\le540\gamma^2.
 \label{f16v85:score-bounds}
\end{equation}
\end{proposition}
\begin{proof}
Condition on all links but $e$.  Its log density, up to a constant, is
$l(x)=h\cdot x$ on $S^3$, with $|h|\le6\gamma$.
The score norm is $|\nabla l|$, independent of the chosen orthonormal
left frame.  On the unit sphere,
$\Delta l=-3l$ and $\operatorname{Hess}l=-l\,g$.
Integrate $\operatorname{div}(e^l\nabla l)$ and
$\operatorname{div}(e^l|\nabla l|^2\nabla l)$ over the compact sphere.
There is no boundary term, and exactly
\[
 \mathbb E|\nabla l|^2=3\mathbb E l,
 \qquad
 \mathbb E|\nabla l|^4=5\mathbb E[l|\nabla l|^2].
\]
Since $l\le|h|$, these imply $3|h|$ and $15|h|^2$, respectively.
Average over the unchanged complement.  The case $h=0$ is included.
In particular the fourth moment is $O(\gamma^2)$, not the
$O(\gamma^4)$ obtained by bounding the score pointwise.
\end{proof}

\subsection{A local double-adjoint witness isolates a plaquette}

For each distinct unoriented elementary plaquette $p$, choose once a pair
of opposite boundary links $e_p,f_p$.  Reversed occurrences in the product
are differentiated with their actual signs and conjugations.  No other
elementary plaquette contains both of these stored links.  Set
\begin{equation}
 h_p^{ab}=D_{e_p}^aD_{f_p}^bW_p,
 \qquad
 Y_p=\sum_{a,b=1}^3(D_{e_p}^a)^*(D_{f_p}^b)^*h_p^{ab}.
 \label{f16v85:witness}
\end{equation}
The symbol $Y_p$ denotes a local function of the full Euclidean history;
it is not in general a function of one spatial slice.

\begin{theorem}[Exact local duality, gauge invariance, and witness size]
\label{f16v85:duality}
The witnesses are gauge invariant and mean zero.  For every smooth test
function $A$,
\begin{equation}
 \mathbb E_\mu[A Y_p]
 =\mathbb E_\mu\sum_{a,b}
      (D_{e_p}^aD_{f_p}^b A)h_p^{ab}.
 \label{f16v85:duality-identity}
\end{equation}
In particular, for distinct elementary plaquettes $q$,
\begin{equation}
 \mathbb E_\mu[W_qY_p]
   =\mathbf1_{p=q}k_p,
 \qquad k_p=\mathbb E_\mu(2+W_p^2)\in[2,3].
 \label{f16v85:biorthogonal}
\end{equation}
For every exterior,
\begin{equation}
                         \|Y_p\|_{L^2(\mu)}\le40(1+\gamma).
 \label{f16v85:Y-size}
\end{equation}
\end{theorem}
\begin{proof}
Two applications of \eqref{f16v85:adjoint} prove
\eqref{f16v85:duality-identity}.  Taking $A=1$ gives zero mean.
If $q\ne p$, one of the chosen links is absent from $q$, giving zero.
The two tangent-coordinate bases can be transported to the plaquette
basepoint by orthogonal maps.  The mixed block is then $wI-C_v$ or its
transpose, up to orthogonal factors, where $U_p=w+v$ and $C_vz=v\times z$.
Since $C_v^2=vv^{\mathsf T}-|v|^2I$,
\begin{equation}
 \sum_{a,b}(h_p^{ab})^2=3w^2+2|v|^2=2+w^2,
 \qquad \|h_p\|_{\rm op}\le1.
 \label{f16v85:mixed-block}
\end{equation}
This also follows directly by differentiating the ordered quaternion
product; no small-field expansion is used.  Gauge transformations rotate
the two frames separately, and every index in \eqref{f16v85:witness}
is contracted.  The density is gauge invariant, so its weighted adjoints
transform in the same way.  This proves gauge invariance of $Y_p$.

For its norm, expansion of the two adjoints gives the useful exact formula
\begin{equation}
 \begin{split}
 Y_p={}&9W_p+\gamma(2+W_p^2)\\
 &-3\{s_{e_p}\cdot\nabla_{e_p}W_p+
                    s_{f_p}\cdot\nabla_{f_p}W_p\}
   +s_{e_p}^{\mathsf T}h_p s_{f_p}.
 \end{split}
 \label{f16v85:Y-expanded}
\end{equation}
Indeed $(D_e^a)^2W_p=-W_p$, so the sum of fourth derivatives is $9W_p$.
The mixed density derivative is
$D_{e_p}^as_{f_p}^b=\gamma h_p^{ab}$: their one shared plaquette is $p$.
The remaining two first-score terms sum to the displayed factors $-3$.
For one link $|\nabla_eW_p|^2=1-W_p^2\le1$.
H\"older and \eqref{f16v85:score-bounds} therefore imply
\[
 \|Y_p\|_2\le
 9+3\gamma+6\sqrt{18\gamma}+\sqrt{540}\,\gamma
 =9+(3+6\sqrt{15})\gamma+18\sqrt{2\gamma}.
\]
Using $2\sqrt\gamma\le1+\gamma$ bounds the last line by
$40(1+\gamma)$.  All scores refer to the full conditional Wilson action,
including temporal and neighboring spatial plaquettes.
\end{proof}

Let $\mathcal B_p$ be the union of links in all plaquettes incident to
$e_p$ or $f_p$.  Formula \eqref{f16v85:Y-expanded} shows that $Y_p$ depends
only on $\mathcal B_p$, restricted to the free variables when conditioning.

\begin{proposition}[Exact finite-range witness covariance]
\label{f16v85:sparsity}
If at least one of $e_p,f_p$ is absent from $\mathcal B_q$, then
$\mathbb E_\mu[Y_pY_q]=0$.  For a finite collection of marks, let
$\mathsf B=(\mathbb E_\mu Y_pY_q)_{p,q}$ and
$\mathsf K=\operatorname{diag}(k_p)$.  Then
\begin{equation}
 0\preceq\mathsf B\preceq L_\gamma I,
 \qquad L_\gamma=182400(1+\gamma)^2,
 \qquad 2I\preceq\mathsf K\preceq3I.
 \label{f16v85:B-bound}
\end{equation}
The constants are independent of the number of marks and the surrounding
spatial or temporal volume.
\end{proposition}
\begin{proof}
Apply \eqref{f16v85:duality-identity} to $A=Y_q$.  The derivative in a
missing variable is zero.  This is orthogonality by integration by parts,
not independence or exponential decay of the Wilson measure.
For a fixed link $e$, at most $1+6\cdot3=19$ stored links share a
plaquette with $e$, including $e$ itself.  If $e\in\mathcal B_q$, at
least one chosen link of $q$ belongs to that set.  Each such link is in
at most six elementary plaquettes, so there are at most $19\cdot6=114$
possible $q$.  A potentially nonzero covariance must satisfy the missing-
variable test in both directions, giving a symmetric relation with row
count at most $114$, including the diagonal.  By Cauchy--Schwarz every
entry has magnitude at most $1600(1+\gamma)^2$.  The symmetric row bound
therefore gives \eqref{f16v85:B-bound}.  Conditioning can only delete
variable dependencies; it does not enlarge this count.
\end{proof}

\subsection{An all-coupling native thermal-order covariance floor}

For any finite real sampling matrix $S=(s_{ip})$, combine repeated marks
and reversed orientations first, and put
\begin{equation}
 F_i=\sum_p s_{ip}W_p,\qquad
 \mathcal S=SS^{\mathsf T},\qquad
 \kappa_\gamma^\sharp=\frac1{45600(1+\gamma)^2}.
 \label{f16v85:constants}
\end{equation}
All four-dimensional plaquette orientations and arbitrarily many time
layers are allowed in this theorem.

\begin{theorem}[Native covariance lower bound at the thermal order]
\label{f16v85:entrance}
Under \eqref{f16v85:law}, with the chosen pairs free,
\begin{equation}
               \boxed{\quad
               \operatorname{Cov}_\mu(F)
                   \succeq\kappa_\gamma^\sharp\mathcal S.\quad}
 \label{f16v85:main}
\end{equation}
This holds uniformly conditional on every admissible exterior.  It passes
to stationary-cylinder and local Gibbs limits at fixed $\gamma$.
Its exact coefficient kernel is $\ker\mathcal S$ when the source marks
satisfy these hypotheses.  Complex coefficients follow by complexification.
\end{theorem}
\begin{proof}
For scalar coefficients $s=(s_p)$, set $F_s=\sum_p s_pW_p$ and
$Y_s=\sum_p s_pY_p$.  The exact overlap and norm bound are
\[
 \mathbb E[(F_s-\mathbb EF_s)Y_s]
       =\sum_p k_ps_p^2\ge2|s|^2,
 \qquad \mathbb E Y_s^2\le L_\gamma|s|^2.
\]
Cauchy--Schwarz gives
$\operatorname{Var}(F_s)\ge4|s|^2/L_\gamma$, which is
\eqref{f16v85:main}.  Apply it to $s=S^{\mathsf T}c$ for every $c$.
A sampled zero combination is the zero function; the positive lower bound
excludes all other null directions.  No source-count factor is introduced.
At fixed $\gamma$ the local witness functions are bounded.  Finite-history
identities and bounds therefore pass to the stated local limits, or follow
directly from the conditional Gibbs specification.  This does not assert
uniform convergence in $\gamma$ or construct a continuum measure.
\end{proof}

The order $(1+\gamma)^{-2}$ improves V79's
$\kappa_\gamma=(4/3)d_\gamma^2=O(\gamma^{-4})$ for the same spatial
bank.  More importantly it applies to a complete space--time sampled bank,
without a temporal-duration factor.  The proof neither assumes a spectral
gap nor derives a long-distance covariance lower bound from this fact.

\begin{proposition}[The exponent is sharp on an actual conditional Wilson cell]
\label{f16v85:sharp}
Free two opposite links $X,Y$ of one elementary plaquette and set every
other link to identity.  Their actual conditional four-dimensional density
is proportional to
\[
 \exp\{\gamma[w(XY^{-1})+5w(X)+5w(Y)]\}\,dH(X)dH(Y).
\]
For $W=w(XY^{-1})$,
\begin{equation}
                  \gamma^2\operatorname{Var}(W)\longrightarrow\frac6{49}.
 \label{f16v85:sharp-limit}
\end{equation}
Thus a uniform conditional lower bound with a better power than
$\gamma^{-2}$ is impossible in this class.
\end{proposition}
\begin{proof}
The complementary five plaquettes of each free link contain no other free
link; their traces are $w(X)$ and $w(Y)$.  The unique minimum of the
nonnegative energy is $X=Y=I$.  In northern charts put the imaginary parts
of $X,Y$ equal to $z/\sqrt\gamma,u/\sqrt\gamma$.  The scaled energy tends to
\[
 \tfrac12(5|z|^2+5|u|^2+|z-u|^2).
\]
Outside fixed charts the pinning energy is bounded away from zero, whereas
a ball of radius $\gamma^{-1/2}$ gives a partition lower bound
$c\gamma^{-3}$.  Inside the charts the pinning bounds the integrand by an
integrable Gaussian times every fixed polynomial needed for fourth moments.
Dominated convergence therefore applies to the normalized compact law, as
in V63; no Gaussian law is substituted at finite $\gamma$.
For each color the limiting precision is
$\left(\begin{smallmatrix}6&-1\\-1&6\end{smallmatrix}\right)$,
so $z-u$ has variance $2/7$ per component.  Also
$\gamma(1-W)\to|z-u|^2/2$ in second moments.  Its limiting variance is
$(3/2)(2/7)^2=6/49$.
\end{proof}

\subsection{The witnesses are exact contact terms in the reflected theory}

Return to the stationary native history and mark spatial plaquettes on one
slice.  A witness $Y_{p,t}$ uses only links in the plaquettes incident to its
chosen spatial pair: its temporal support lies between slices $t-1$ and
$t+1$.  For any observable independent of at least one chosen link,
\eqref{f16v85:duality-identity} remains zero, by approximation if necessary.
Consequently
\begin{equation}
 \mathbb E[Y_{p,0}A(U_n)]=0\quad(n\ne0),\qquad
 \mathbb E[Y_{p,0}Y_{q,n}]=0\quad(|n|\ge2),
 \label{f16v85:time-contact}
\end{equation}
where $A$ is any bounded spatial-slice observable, including a simultaneous-
flow function of that slice.  No spatial restriction on $A$ is required.

\begin{theorem}[Exact local contact removal without a change of separated covariance]
\label{f16v85:removal}
For a spatial bank $F_t=S W_t$, take witnesses on that same set of marks,
center $F$ natively, and set
\begin{equation}
 \lambda_\gamma=\frac2{L_\gamma},\qquad
 F_t^{\rm sub}=F_t^\circ-\lambda_\gamma S Y_t.
 \label{f16v85:sub-source}
\end{equation}
Then
\begin{equation}
 \boxed{\begin{aligned}
 0\preceq G_{\rm sub}&\preceq G_F-
                       \kappa_\gamma^\sharp\mathcal S,\\
 \operatorname{Cov}(F_0^{\rm sub},F_n^{\rm sub})
                    &=C_F(n)\qquad(|n|\ge2).
 \end{aligned}}
 \label{f16v85:subtraction}
\end{equation}
The correction is gauge invariant, local in a one-step temporal collar,
and changes neither the probability measure nor any sufficiently separated
source pairing.  Placed at $t\ge2$, each $Y_{p,t}$ is null in the
Osterwalder--Schrader reflected quotient.  In particular $F^{\rm sub}$ and
$F$ have the same reflected class at these placements.
\end{theorem}
\begin{proof}
Let $\mathsf K=\operatorname{diag}(k_p)$ and $\mathsf B=\operatorname{Cov}Y$.
The exact equal-time covariance is
\[
 G_{\rm sub}=G_F-
 S(2\lambda_\gamma\mathsf K-\lambda_\gamma^2\mathsf B)S^{\mathsf T}.
\]
Using $\mathsf K\succeq2I$ and $\mathsf B\preceq L_\gamma I$, the
subtracted matrix is at least
$(4\lambda_\gamma-\lambda_\gamma^2L_\gamma)\mathcal S
=(4/L_\gamma)\mathcal S$.
Nonnegativity is automatic because $G_{\rm sub}$ is a covariance.
Equation \eqref{f16v85:time-contact} kills both linear correction terms
and the quadratic one at every separation at least two.

For reflected nullity, put the marked spatial pair at $t\ge2$, so its
entire collar is on the positive side.  Reflecting it places the two
integration-by-parts variables strictly on the negative side.  Every
positive-time cylinder observable is independent of those variables.
Thus its pairing with the reflected witness is zero by
\eqref{f16v85:duality-identity}.  Taking the observable to be the witness
itself gives zero reflected norm.  The same identity handles any finite
linear combination and common positive-time placement.  The extension by
continuity uses the original reflection-positive form, not an altered law.
\end{proof}

These nonzero Euclidean witnesses are not nonzero slice vectors propagated
by the positive injective physical transfer: they use temporal plaquettes
and hence adjacent slices.  Their exact reflected nullity therefore does
not contradict that injectivity.  It also explains why a large ordinary
variance need not be evidence of physical source noncollapse.
The displayed local correction removes at least the stated variance; it
is not claimed to remove every divergent contact contribution.

For completeness the optimal finite linear score subtraction has equal-time
Gram
\begin{equation}
 G_F-\mathsf L\mathsf B^{-1}\mathsf L^*,
 \qquad \mathsf L_{ip}=\mathbb E[F_i^\circ Y_p],
 \label{f16v85:optimal}
\end{equation}
for any bounded smooth gauge-invariant slice bank.  Here $\mathsf B$ is
strictly positive for distinct marks: a zero witness combination would
contradict \eqref{f16v85:biorthogonal}.  Exact redundant marks use its
range instead.  The coefficients in \eqref{f16v85:optimal} may connect
all chosen marks, so the explicit scalar choice in
\eqref{f16v85:sub-source}, not this inverse, is the local construction.
For a simultaneous-flow bank $\mathsf L$ is an actual mixed-derivative
expectation by \eqref{f16v85:duality-identity}; its value is not evaluated
here.  Every such subtraction still preserves the separated pairings.

\subsection{A native contact divergence and sharper microscopic consequences}

\begin{theorem}[Thermally normalized space--time contact divergence]
\label{f16v85:contact-divergence}
Let $\psi_i$ be fixed continuous compactly supported space--time tests,
sampled on a sequence of tori or cylinders large enough to contain their
support away from boundaries.  Use one fixed spatial plaquette orientation
and set, with the original Wilson parameter $\gamma_a\to\infty$,
\[
 \mathcal E_{a,i}=\gamma_a\sum_z\psi_i(az)
                    [1-W_z-\mathbb E(1-W_z)].
\]
Then in the quadratic-form sense on the finite coefficient space,
\begin{equation}
 \boxed{\quad
 \liminf_{a\downarrow0}a^4\operatorname{Cov}(\mathcal E_a)
 \succeq\frac1{45600}
       \left(\int_{\mathbb R^4}\psi_i(x)\psi_j(x)\,d^4x\right)_{ij}.
 \quad}
 \label{f16v85:contact-limit}
\end{equation}
No Gaussian covariance comparison or relation between $a$ and $\gamma_a$
is assumed.  For tests supported a fixed positive distance beyond a time
reflection plane, the score correction above removes at least this order
of ordinary covariance while leaving the reflected class unchanged.
\end{theorem}
\begin{proof}
Apply \eqref{f16v85:main} to all marked space--time plaquettes at once.
Its constant is independent of the time duration.  Multiply by
$a^4\gamma_a^2$ and use
$\gamma_a^2/(1+\gamma_a)^2\to1$ and the Riemann-sum identity for the test
Gram.  In every coefficient direction this proves
\eqref{f16v85:contact-limit}.  The same equal-time algebra used in
\eqref{f16v85:subtraction} applies to the whole space--time mark set, because
\eqref{f16v85:B-bound} was proved in four dimensions.  Its removed covariance
has the same lower bound.  At positive physical support distance, the
one-lattice-step collars eventually stay strictly on the positive side,
so the reflected-null proof applies term by term.  Neither statement
constructs a finite reflected limit or an ordinary random-field limit.
\end{proof}

\begin{proposition}[Sharper native innovation and delay bounds]
\label{f16v85:temporal}
For a spatial plaquette bank under the original stationary transfer,
\begin{equation}
                  G_F-C_F(2)\succeq
                            \kappa_\gamma^\sharp\mathcal S.
 \label{f16v85:innovation}
\end{equation}
The complete feature cost of V81 and all its larger polynomial witness
costs dominate the same matrix.  In the reflection-compatible laws where
V80's corrected-extension estimate holds, put, for $\gamma\ge16$,
\[
 v_\gamma=\frac{\chi(8b_\gamma+53)}{\gamma^2},\qquad
 \delta_\gamma^\sharp=\frac{v_\gamma}{\kappa_\gamma^\sharp},\qquad
 q(\delta)=\frac2{\sqrt{\delta^2+4}+\delta}.
\]
Here $\chi\le4$ is the spatial mark incidence and $b_\gamma=384B_\gamma$
is exactly the inherited defect-tail constant.  Then
\begin{equation}
 C_F(n)\succeq q(\delta_\gamma^\sharp)^nG_F
        \succeq\kappa_\gamma^\sharp q(\delta_\gamma^\sharp)^n\mathcal S,
 \qquad n\ge1.
 \label{f16v85:delay}
\end{equation}
The sufficient one-step factor now has order $(\log\gamma)^{-1}$.
\end{proposition}
\begin{proof}
Condition on the complete history except the spatial links on slice zero.
The conditional version of \eqref{f16v85:main} gives the lower variance
matrix for every exterior.  This conditioning is finer than conditioning
only on slice one.  The conditional-variance decomposition and native
reversibility therefore give
\[
 G_F-C_F(2)=\mathbb E\operatorname{Cov}(F(U_0)\mid U_1)
 \succeq\kappa_\gamma^\sharp\mathcal S.
\]
V81's spectral inequality $C_{-1}-C_1\succeq G-C_2$ gives its feature-cost
consequence; no equal-time derivative floor is assigned to a different
operator.  Finally use V80's proved
$\mathsf B_\gamma^{[2]}\preceq v_\gamma\mathcal S$ and the new entrance
bound in V79's block-Gram delay theorem.  Since
$b_\gamma=O(\log\gamma)$, $\delta_\gamma^\sharp=O(\log\gamma)$ with
comparable large-$\gamma$ constants, proving the stated order.
\end{proof}

The improvement replaces V80's sufficient $(\gamma^2\log\gamma)^{-1}$
factor for this bank; its constants remain deliberately conservative.
At a fixed physical delay $n\asymp a^{-1}$ the displayed lower factor
still tends to zero.  This is not a source-collapse theorem and not a mass
gap.  If $G_F\preceq ND_\gamma\gamma^{-2}\mathcal S$ is the inherited
microscopic moment bound, \eqref{f16v85:innovation} also improves V81's
necessary lower ratio to
$\gamma^2/[45600(1+\gamma)^2ND_\gamma]$.
Thus its sufficient $R\preceq2aE_*G$ test is impossible when
$aN_aD_{\gamma_a}\to0$; this is a criterion obstruction, not absence of
finite-energy source weight.

\subsection{The remaining reflected problem and result ledger}

We have evaluated a new native matrix, not another unspecified certificate:
the full space--time plaquette covariance has a uniform thermal-order floor.
Its proof simultaneously constructs local contact corrections that remove
a positive part of that floor without changing any separated covariance.
For plaquette data used in V84 with buffer $k\ge1$, all moments
$C_F(2k+jd)$ are consequently unchanged by the score substitution.
For literal positive-time representatives use $k\ge2$ so the collars do
not cross the plane.  The entire forward moment hierarchy, not merely one
compressed eigenvalue, is then the same.  Contact optimization cannot turn
an unproved growing-depth upper inequality into a proved one.

For the actual simultaneously flowed source, the same local score functions
can be paired with its original-input mixed derivatives, and their source
corrections are still reflected-null.  The elementary isolation identity
\eqref{f16v85:biorthogonal} does not hold with $W_q$ silently replaced by a
flowed plaquette.  That collective derivative pairing and its surviving
positive-separation Gram remain unestimated.  A useful next calculation
must control those noncontact correlations or V84's native signed localizer,
not infer a physical state from the stronger entrance floor.  No action,
physical clock, stationary law, or simultaneous-flow map has changed.

The companion diagnostic is
\begin{center}\small
\path{verify_ym_f16_v85_local_scores.py}.
\end{center}
It separates exact Haar-series identities, literal four-dimensional link
incidence, compact quaternion derivative and gauge checks, conditional
compact quadrature, and finite matrix contact-subtraction fixtures.
These are not full Wilson-vacuum simulations, outward-rounded interval
certificates, or proof-assistant verification.  The proofs above do not use
numerical outcomes.  The predecessor is preserved, not recertified.

\clearpage
\begin{record}[V85 result ledger]
\label{f16v85:ledger}
\emph{Proved here:} conditional second and fourth score moments; local
weighted double-adjoint witnesses with exact plaquette duality; finite-range
witness covariance without a native mixing assumption; the explicit
$(1+\gamma)^{-2}$ covariance floor for arbitrary space--time plaquette
banks; its sharp exponent in a literal conditional two-link cell; local
source corrections with a quantitative variance removal and identical
separated correlations; exact reflected nullity with support retained;
a native canonical contact-divergence lower bound; and stronger native
innovation and microscopic delay consequences.

\emph{Inherited or standard:} compact Haar integration by parts and its
Schwinger--Dyson interpretation; the Wilson conditional specification,
unit-quaternion derivative convention, native transfer and reflection framework,
V80's corrected-extension upper budget and defect moments, V81's complete
feature ordering, and V84's source-completion distinction.  No external
continuum or bootstrap result is imported as a hypothesis.

\emph{Not asserted:} independence of distant native plaquettes; an optimal
numerical constant; removal of every contact divergence; a slice-vector
interpretation of the temporally extended score witness; a local inverse
of its covariance matrix; the plaquette isolation identity for flowed
outputs; or a reflected lower Gram inferred from a contact covariance.

\emph{Still unproved:} nonzero, physically normalized positive-time native
Grams for the simultaneous-flow bank and their spatial/thermodynamic control;
V84's growing-depth upper inequalities with common dense-core constants;
a compatible interacting continuum construction; and its Yang--Mills
mass gap.
\end{record}

\begin{firewall}
The new lower bound is native and all-coupling, but the score modes producing
it are contact terms in the reflected theory.  Their nullity is proved with
their actual time support, not by declaring their ordinary variance zero.
Improved contact normalization leaves the original positive-time spectral
question unchanged.  All V1--V84 mathematical text remains unchanged.
\end{firewall}
% END F16 V85 APPEND


% BEGIN F16 V86 APPEND -- 2026-09-19
% Append-only continuation of the Post-Fifteenth-Archive V85.
\clearpage
\section[V86: an invisible parity sector and a separating source algebra]{V86: an invisible parity sector\\and a separating source algebra}
\label{f16v86:context}

V85 improves a native contact bound without changing the positive-time
source class.  V84 separately requires dense source coverage before a
common cyclic-space estimate can imply a full spectral gap.  Before
optimizing more contact representatives, we check that coverage for the
actual family being studied.  There is an exact obstruction: linear sums
of individually flowed elementary-plaquette class functions, even with
arbitrary spatial tests, flow times and forward physical-time translates,
have no zero-momentum, odd-parity component.  A specified contractible
eight-link loop supplies a nonzero native vector in that missing sector.
A local score calculation gives its bank a thermal-order covariance floor
uniform in volume, and an explicit finite-delay lower bound.

The constructive repair is an algebra of sources, not a further time
completion of the same linear family.  We give finite rooted-loop
invariants which separate every compact $\SU(2)$ gauge orbit, including
reducible configurations.  Products of these invariants after the actual
simultaneous spatial flow are dense at each finite regulator.  This is a
finite-regulator coverage statement: it does not construct a nonzero
physical-scale limit of the odd source or give common continuum spectral
constants.  The earlier source-restricted theorems remain unchanged.

\subsection{Native symmetry and the precise linear family}

Work on a periodic cubic spatial lattice $\Lambda$ of even side length
$L\ge8$, at spacing $a>0$ and finite $\gamma>0$.  Write $N_v=|\Lambda|$.
Use the original invariant Wilson Hilbert space $L^2(\pi)$, its centered
subspace, and its positive injective Markov transfer $P$.  The unique
positive Perron function makes $\pi$ invariant under spatial translations
and inversion.  Indeed the magnetic sandwich and its kernel commute with
these transformations, and the transformed positive normalized Perron
function must be the same function.  No boundary-conditioned law is
assigned these global symmetries.  Conditional covariance bounds below
are stated separately and do not require symmetry of the exterior.

Let $\mathsf T_z$ be the unitary translation of observables and let
$\mathsf I$ implement spatial inversion.  On positively stored links the
latter configuration map is
\begin{equation}
 (\iota U)_\mu(x)=U_\mu(-x-\hat\mu)^{-1},\qquad
 \Pi_0=N_v^{-1}\sum_z\mathsf T_z,\qquad
 \Pi_{0,-}=\tfrac12(I-\mathsf I)\Pi_0.
 \label{f16v86:symmetry}
\end{equation}
Thus $\Pi_{0,-}$ is the orthogonal zero-momentum, odd-parity projection.
Both projections commute with $P$.  The simultaneous spatial flow
$\Phi_{a,\tau}$ commutes with these configuration symmetries, because its
spatial action and product metric are invariant and its finite-dimensional
ODE has unique solutions.  In particular this assertion concerns the full
flow, not an isolated plaquette evolution.

For context, organizing lattice spectra by spatial representations and
parity is standard; see Athenodorou and Teper,
\href{https://doi.org/10.1007/JHEP11(2020)172}{JHEP \textbf{11} (2020), 172}.
The density of Wilson-loop algebras is also established, for example by
T.~L\'evy, \emph{Wilson loops in the light of spin networks},
\href{https://arxiv.org/abs/math-ph/0306059}{arXiv:math-ph/0306059},
J.~Geom.~Phys.\ \textbf{52} (2004), 382--397.
Neither symmetry projection nor the general density principle is claimed
as new.  The particular missing native sector, its literal loop and local
covariance constants, and the finite separating construction used here are
proved below.  No spectrum from those references is imported.

Consider the closed linear span $\mathcal C_{\rm pl}$ generated by all
native-centered functions
\begin{equation}
 P^n\left[\sum_x c_x\,
 g\bigl(w((\Phi_{a,\tau}U)_{\mu\nu}(x))\bigr)\right]^\circ,
 \qquad n\ge0,\quad \tau\ge0,
 \label{f16v86:pl-family}
\end{equation}
with arbitrary real coefficients, orientations and bounded real functions
$g$.  The current squared-imaginary-curvature bank is included by
$g(t)=1-t^2$.  This deliberately large \emph{linear} family does not include
arbitrary products of its members.  Fractional spectral powers of $P$,
when used only as Hilbert-space operators, give the same symmetry conclusion.

\begin{proposition}[An exact absent sector of every elementary-plaquette flow bank]
\label{f16v86:absent}
The whole linear space \eqref{f16v86:pl-family} satisfies
\begin{equation}
                 \boxed{\quad \Pi_{0,-}\mathcal C_{\rm pl}=0.\quad}
 \label{f16v86:absent-sector}
\end{equation}
The conclusion persists after common positive-time buffering, temporal
averaging, and V77's common input-link polynomial filter.  Reflected-null
score corrections do not change the conclusion for the physical classes.
\end{proposition}
\begin{proof}
Inversion sends a $\mu\nu$ plaquette to the same unoriented shape translated
to $-x-\hat\mu-\hat\nu$.  A reversed loop has the same real $\SU(2)$
trace.  The spatial sum of its flowed class function is therefore even
under inversion.  Applying $\Pi_0$ to a weighted sum in
\eqref{f16v86:pl-family} gives that even sum multiplied by
$N_v^{-1}\sum_xc_x$.  Centering changes only a constant.  Since $P$ and
$\Pi_0$ commute with inversion, the odd projection vanishes before and
after every forward translate, and hence on the closed span.

The common input filter commutes with translations and inversion: its
central kernel is invariant under inversion, and Haar measure is unchanged
by reversing and permuting stored links.  It cannot create the missing
component.  Positive-time averages preserve the same linear identities.
A reflected-null correction leaves the represented Hilbert vector itself
unchanged, by V85, so it cannot add an odd physical class.
\end{proof}

\subsection{A literal contractible loop outside that source space}

Use the ordered spatial step word
\begin{equation}
 \Gamma=(+1,+2,+1,+3,-1,-1,-2,-3),\qquad
 \Gamma^-=-\Gamma
 \label{f16v86:loop}
\end{equation}
with the signs in $-\Gamma$ changed without reversing their order.
The successive vertices of $\Gamma$, rooted at zero, are
\[
 (0,0,0),(1,0,0),(1,1,0),(2,1,0),(2,1,1),
 (1,1,1),(0,1,1),(0,0,1),(0,0,0).
\]
It has eight distinct stored edges, lies in a $2\times1\times1$ box and
is contractible.  Let $W_{+,x}=w(U_{x+\Gamma})$ and
$W_{-,x}=w(U_{x+\Gamma^-})$.  Define
\begin{equation}
 Z_\Lambda(U)=\frac1{\sqrt{2N_v}}
                   \sum_{x\in\Lambda}(W_{+,x}-W_{-,x}).
 \label{f16v86:odd-source}
\end{equation}
No continuum spin label is assigned to this source; its claimed quantum
numbers here are only zero lattice momentum and odd spatial parity.

\begin{proposition}[Distinct shapes, exact Haar norm and native orthogonality]
\label{f16v86:odd}
All $2N_v$ unoriented stored-edge sets in \eqref{f16v86:odd-source} are
distinct.  The source is real, smooth, gauge invariant, translation
invariant, parity odd and nonzero.  Exactly,
\begin{equation}
 \mathbb E_H Z_\Lambda=0,\qquad
 \|Z_\Lambda\|_H^2=\tfrac14,\qquad
 \mathbb E_\pi Z_\Lambda=0.
 \label{f16v86:Haar-odd}
\end{equation}
For every $A\in\mathcal C_{\rm pl}$ and $m\ge0$,
\begin{equation}
 \boxed{\quad \langle Z_\Lambda,P^mA\rangle_\pi=0,\qquad
       \langle Z_\Lambda,P^mZ_\Lambda\rangle_\pi>0.\quad}
 \label{f16v86:orthogonal}
\end{equation}
The same absence applies to the nonzero odd source
$Z_\Lambda\circ\Phi_{a,\tau_0}$ at every finite $\tau_0\ge0$.
\end{proposition}
\begin{proof}
We give an edge marker which also supplies the quantitative estimate below.
Choose the positive direction-$3$ edge at $(0,0,0)$ and the positive
direction-$1$ edge at $(1,1,0)$ in $\Gamma$.  They belong to no common
elementary four-dimensional plaquette.  The set of differences
``direction-$1$ tail minus direction-$3$ tail'' in $\Gamma$ is
\begin{equation}
 \begin{split}
 \mathcal D=\{&(0,0,0),(1,1,0),(1,1,1),(0,1,1),\\
             &(-2,-1,0),(-1,0,0),(-1,0,1),(-2,0,1)\}.
 \end{split}
 \label{f16v86:displacements}
\end{equation}
Every entry occurs once.  In the inverted shape the set becomes
$-\mathcal D-\hat1+\hat3$.  The chosen difference $(1,1,0)$ would then
require $(-2,-1,1)\in\mathcal D$, which is false.  Thus its two actual
stored edges occur together in exactly one translated shape from either
family.  Inverting the statement gives the same uniqueness for each
negative shape.  Periodic side length at least eight prevents a distinct
listed displacement from becoming equal modulo the periods.  This also
proves that no two marked edge sets coincide.

Each simple loop has Haar mean zero and second moment $1/4$: integrating
one of its once-occurring links makes its holonomy Haar.  Distinct sets
have a link present in just one of the two traces, so their Haar cross
moment is zero.  Summing proves \eqref{f16v86:Haar-odd}.  Translations
permute each sum, and inversion exchanges the sums.  Native symmetry
therefore makes its mean zero.  The positive native density has full
support, so its nonzero continuous restriction has positive native norm.
It lies in $\operatorname{Ran}\Pi_{0,-}$; the previous proposition and
commutation of $P$ give the orthogonality.  Positive injectivity of $P$
gives the strictly positive quadratic form at every finite $m$.
Finally, the finite-regulator flow is a surjective gauge-equivariant
symmetry-equivariant diffeomorphism.  Composition preserves nonconstancy
and the stated symmetry sector; the same argument applies.
\end{proof}

This proves failure of dense coverage for the stated linear elementary
bank, not failure of any earlier theorem restricted to its own cyclic
space.  All of its V84 forward moments can be controlled while the odd
sector remains completely untested.  The source \eqref{f16v86:odd-source}
is also neutral under each flat $\mathbb Z_2$ cut transformation of a
spatial torus: a contractible loop has zero winding.  Thus merely restricting
to that neutral sector does not remove this example.  No nonzero limiting
odd-parity vector is inferred from finite-regulator nonconstancy.

\subsection{A populated native Gram for the missing shape family}

The preceding positivity can be quantified without whole-volume density
extrema.  Let $\mathcal L$ be any finite collection of translates of the
two shapes, also allowing marks on different time slices.  Use the two
markers above and their inverted/translational images for each $\ell$.
They are assumed free when a conditional law is used.  Put
\[
 F_i=\sum_{\ell\in\mathcal L}s_{i\ell}W_\ell,\qquad
 \mathcal S_{ij}=\sum_\ell s_{i\ell}s_{j\ell},\qquad
 \kappa_\gamma^{\rm odd}=\frac1{18400(1+\gamma)^2}.
\]

\begin{theorem}[Thermal-order covariance for the explicit eight-link bank]
\label{f16v86:entrance}
Under the original Wilson law, and uniformly conditional on every exterior
with those marker pairs free,
\begin{equation}
 \boxed{\quad
 \operatorname{Cov}_\mu(F)\succeq
             \kappa_\gamma^{\rm odd}\mathcal S.\quad}
 \label{f16v86:native-entrance}
\end{equation}
The constant is independent of the number of marks and surrounding spatial
or temporal volume.  The coefficient kernel is exactly $\ker\mathcal S$.
In particular
\begin{equation}
 \operatorname{Var}_\pi Z_\Lambda\ge\kappa_\gamma^{\rm odd},\qquad
 \liminf_{\gamma\to\infty}
       \operatorname{Var}_\pi(\gamma Z_\Lambda)\ge\frac1{18400},
 \label{f16v86:odd-floor}
\end{equation}
including changing finite spatial volumes in this inequality.
\end{theorem}
\begin{proof}
Retain V85's weighted adjoints and score moments under the full native
conditional action.  For a marked loop with chosen links $e,f$ define
\[
 h_\ell^{ab}=D_e^aD_f^bW_\ell,\qquad
 Y_\ell=\sum_{a,b}(D_e^a)^*(D_f^b)^*h_\ell^{ab}.
\]
The trace is multilinear in each of its eight distinct quaternion links.
Transporting the two tangent frames through its other unit factors gives
\[
 \sum_{a,b}(h_\ell^{ab})^2=2+W_\ell^2,\quad
 \|h_\ell\|\le1,\quad |\nabla_eW_\ell|\le1.
\]
These identities do not require $e,f$ to border one plaquette.  The marker
uniqueness in \eqref{f16v86:displacements} gives the exact overlap
\begin{equation}
 \mathbb E[W_{\ell'}Y_\ell]
 =\mathbf1_{\ell=\ell'}k_\ell,\qquad
 k_\ell=\mathbb E(2+W_\ell^2)\in[2,3].
 \label{f16v86:dual}
\end{equation}
The witnesses are gauge invariant and mean zero by the same two compact
integrations by parts.  Because $e,f$ meet no action plaquette,
$D_e^as_f^b=0$, so their exact expansion is now
\begin{equation}
 Y_\ell=9W_\ell
 -3(s_e\cdot\nabla_eW_\ell+s_f\cdot\nabla_fW_\ell)
 +s_e^{\mathsf T}h_\ell s_f.
 \label{f16v86:score}
\end{equation}
In particular there is no $\gamma(2+W_\ell^2)$ term.  V85's score bounds
imply
$\|Y_\ell\|_2\le9+6\sqrt{18\gamma}+\sqrt{540}\gamma
\le40(1+\gamma)$, uniformly in the exterior.

The support of $Y_\ell$ is contained in the eight loop edges and the two
full plaquette stars of its markers.  If either marker of $\ell$ is absent
from that support for $Y_{\ell'}$, their covariance is exactly zero by
integration by parts.  At a fixed stored edge at most eight translated
shapes contain that edge: each shape has four direction-$1$, two
direction-$2$ and two direction-$3$ edges, and there are two shapes.
There are at most nineteen edges sharing an action plaquette with a given
edge, counting the edge itself.  Each such edge can be a chosen marker
of at most two shape instances, since the marker directions are distinct.
Consequently at most $8+19\cdot2=46$ witness supports contain a prescribed
marker.  This bounds every covariance row, also with marks on several
time slices.  Cauchy--Schwarz therefore gives
\[
 0\preceq\operatorname{Cov}(Y)
       \preceq73600(1+\gamma)^2I.
\]
For scalar coefficients $c_\ell$, the overlap with $\sum c_\ell Y_\ell$
is at least $2\sum c_\ell^2$.  Cauchy--Schwarz yields
$\operatorname{Var}(\sum c_\ell W_\ell)
\ge[18400(1+\gamma)^2]^{-1}\sum c_\ell^2$.
Apply this to every row combination of the sampling matrix.  For
$Z_\Lambda$ the squared coefficient sum is one.  Finite-history identities
and uniform local bounds pass to the same stationary and local limits as
V85.  No native decorrelation assumption occurs.
\end{proof}

\begin{proposition}[Finite-delay odd-sector lower bound and its contact limitation]
\label{f16v86:delay}
For a spatial bank of these loops set
\begin{equation}
 v_\gamma^{(8)}=\frac{(1+4/\gamma)^8-1}{4},\quad
 \delta_\gamma^{\rm odd}=147200(1+\gamma)^2v_\gamma^{(8)},\quad
 q_\gamma^{\rm odd}=\frac2{\sqrt{(\delta_\gamma^{\rm odd})^2+4}
                                      +\delta_\gamma^{\rm odd}}.
 \label{f16v86:delay-constants}
\end{equation}
Then the original stationary transfer satisfies
\begin{equation}
 \boxed{\quad C_F(n)\succeq(q_\gamma^{\rm odd})^nG_\pi(F)
 \succeq\kappa_\gamma^{\rm odd}(q_\gamma^{\rm odd})^n\mathcal S,
 \qquad n\ge1.\quad}
 \label{f16v86:delay-bound}
\end{equation}
Also $G_\pi(F)-C_F(2)\succeq\kappa_\gamma^{\rm odd}\mathcal S$.
The thermal entrance floor is not a positive-time continuum lower bound.
\end{proposition}
\begin{proof}
For a specified set of $j$ differentiated links of a simple unit-quaternion
loop, the sum of squared ambient coordinate derivatives of its trace is
$4^{j-1}$.  Sum one coordinate first, using norm multiplicativity, and then
the remaining $j-1$ coordinates; each has four unit basis choices.  No
link is differentiated twice.  Thus its full V79 derivative cost is
$v_\gamma^{(8)}$.  A nonzero derivative of the sampled sum contains a
specified differentiated edge, which occurs in at most eight marked loops.
The pointwise Cauchy--Schwarz bound gives
$\mathsf B_\gamma\preceq8v_\gamma^{(8)}\mathcal S$.
Combine with \eqref{f16v86:native-entrance} and apply V79's exact native
block-Gram theorem, without any vacuum extrema.  This yields
\eqref{f16v86:delay-constants}--\eqref{f16v86:delay-bound}.
For the innovation, apply the conditional entrance theorem with the whole
history except slice zero fixed.  This conditioning is finer than the
next spatial slice, so the variance decomposition gives $G-C_2$ as in V85.
The witnesses have a one-step temporal collar and remain reflected-null
at positive placements; their support proof does not use loop length four.
\end{proof}

For example $q_\gamma^{\rm odd}\sim(1177600\gamma)^{-1}$ as
$\gamma\to\infty$.  The lower bound at $n\asymp a^{-1}$ still vanishes
on the contemplated physical trajectory.  The coefficient-normalized shape
sum has supremum at most $\sqrt{2N_v}$; no volume-uniform supremum or
limiting odd state is asserted.  The new estimate proves that the absent
finite-regulator sector is native and quantitatively nonzero, not that its
unknown physical-scale mass is smaller than the scalar-sector mass.

\subsection{A finite separating algebra, rather than more translates of one shape}

We now specify enough invariant data to avoid that coverage defect at a
finite regulator.  Let $\mathcal G=(\mathcal V,\mathcal E)$ be a finite
connected spatial graph, choose a root and a spanning tree, and put
$b=|\mathcal E|-|\mathcal V|+1$.  The $b$ non-tree edges define rooted
fundamental holonomies $H_1,\ldots,H_b$.  Tree transport sets the tree
links to identity; the residual root gauge transformation conjugates all
$H_i$ simultaneously.  This is a parametrization of gauge orbits, not a
replacement of the native probability density by independent holonomies.

Define the following real functions, in the displayed order of indices:
\begin{equation}
 t_i=w(H_i),\qquad t_{ij}=w(H_iH_j)\ (i<j),\qquad
 t_{ijk}=w(H_iH_jH_k)\ (i<j<k).
 \label{f16v86:generators}
\end{equation}
Each is a genuine closed-path trace, possibly with a retraced root stem.

\begin{theorem}[Single, pair and triple traces separate every compact gauge orbit]
\label{f16v86:separate}
The finite family \eqref{f16v86:generators}, of size
$b+\binom b2+\binom b3$, separates all gauge orbits on $\mathcal G$.
Its unital real polynomial algebra is uniformly dense in continuous
gauge-invariant functions, and hence dense in the native invariant
$L^2$ space.  Reducible configurations and central holonomies are included.
\end{theorem}
\begin{proof}
Write $H_i=a_i+v_i$ with $v_i\in\mathbb R^3$.  The singles give $a_i$;
unit norm gives $|v_i|^2=1-a_i^2$.  The pairs give every off-diagonal entry
of the vector Gram matrix through
\[
 v_i\cdot v_j=a_ia_j-t_{ij}.
\]
Quaternion multiplication gives
\begin{equation}
 \det(v_i,v_j,v_k)=a_ia_ja_k-a_kv_i\cdot v_j
 -a_iv_j\cdot v_k-a_jv_i\cdot v_k-t_{ijk}.
 \label{f16v86:triple}
\end{equation}
Thus all oriented triple volumes are known as well.  Two tuples with the
same Gram matrix are related by an orthogonal transformation of their real
spans.  If the span is three dimensional, one nonzero oriented triple fixes
its determinant to $+1$.  If the span has dimension at most two, the map
can be extended to $\mathbb R^3$ with determinant $+1$ by choosing the
orientation of the orthogonal complement.  Every such rotation is an
adjoint $\SU(2)$ transformation.  The holonomy tuples are therefore
simultaneously conjugate, and the original links are gauge equivalent.
This proof uses no inverse of a vanishing Gram minor.

The gauge quotient is compact Hausdorff.  The real algebra generated by
\eqref{f16v86:generators} contains constants and separates its points;
the real Stone--Weierstrass theorem gives uniform density.  Continuous
functions are dense for the finite native Borel probability measure on
that quotient, proving the $L^2$ assertion.  The known density theorem is
therefore realized here by specified finite generators, not by assuming
that the original linear plaquette bank is dense.
\end{proof}

The triple information cannot in general be omitted.  For $0<\theta<\pi/2$
compare the three holonomies
\[
 (\cos\theta+\mathbf i\sin\theta,
  \cos\theta+\mathbf j\sin\theta,
  \cos\theta+\mathbf k\sin\theta)
\]
with the tuple obtained by changing the last imaginary sign.  Their single
and pair traces agree, while their ordered triple traces are
$\cos^3\theta-\sin^3\theta$ and
$\cos^3\theta+\sin^3\theta$.  They are not simultaneously conjugate.
Set tree links to identity and the three chord links to these values to
obtain literal finite-graph configurations.  This is not a claim that
higher derivatives of all flowed plaquette functions have the same data;
it establishes precisely what the stated pair-only generator list misses.

\begin{proposition}[Actual spatial flow and forward buffering preserve finite coverage]
\label{f16v86:flow-density}
Fix any finite $\tau_0\ge0$.  Compose the generators
\eqref{f16v86:generators} with the complete simultaneous spatial flow on
the finite lattice.  The polynomial algebra of those composed functions
is uniformly dense in continuous gauge-invariant slice functions.
After native centering and any fixed integer buffer $k\ge1$, the span of
its images under $P^k$ is dense in the vacuum-orthogonal Hilbert space.
\end{proposition}
\begin{proof}
The flow is a gauge-equivariant homeomorphism with a gauge-equivariant
inverse, and so descends to a homeomorphism of the compact orbit space.
Its composed generators still separate all orbits.  The preceding proof
then gives density without asserting any uniform Jacobian bound.
Centering is a bounded map onto the centered space.  A bounded injective
self-adjoint $P^k$ has dense range there, since the orthogonal complement
of its range is $\ker P^k=0$.  Approximation before applying the bounded
operator proves the last statement.  No inverse physical-time evolution
is applied to a source.
\end{proof}

Products are indispensable to this algebra assertion.  It is not a claim
that the finite list of generators, or their forward time translates,
alone forms a dense linear set.  The rooted paths may be long on a growing
lattice, and the number of monomials increases with degree and cycle rank.
No uniform thermodynamic or physical-source approximation rate follows.

\subsection{An explicit symmetry-correct first enlargement and its costs}

The smallest enlargement supplied by the present calculation is the odd
shape itself, or its actual simultaneously flowed version
\begin{equation}
 Z_\Lambda^{\tau_0}=Z_\Lambda\circ\Phi_{a,\tau_0}.
 \label{f16v86:flowed-odd}
\end{equation}
It stays exactly orthogonal to the zero-momentum even bank at every native
time lag and has positive finite-regulator reflected norm.  Its physical
normalization and positive-separation limit remain to be estimated.
For the unflowed shape, V77's common filter obeys the exact identity
\begin{equation}
 \mathcal A_{K,S}Z_\Lambda
       =\left(\frac K{K+3}\right)^8 Z_\Lambda,\qquad K\ge1,
 \label{f16v86:bare-filter}
\end{equation}
because each loop uses eight distinct fundamental links once.  The filter
therefore retains this newly added direction, whereas it could never have
created it from the old even zero-momentum family.

For the genuine flowed odd source, its half-oscillation is at most
$\sqrt{2N_v}$ independently of finite $\tau_0$.  Applying V77 to all
$M=3N_v$ original input links gives, for $n\ge2$,
\begin{equation}
 \left|C_{Z^{\tau_0}}(n)-C_{\mathcal A_KZ^{\tau_0}}(n)\right|
       \le \frac{108\pi^2N_v^2\gamma}{K+3}.
 \label{f16v86:odd-reconstruction}
\end{equation}
A scalar source normalization multiplies this budget by its square.
The filtered function is still odd and translation invariant; the native
law and the flow map are unchanged.  At fixed physical volume
$N_v=\mathcal V/a^3$, the sufficient unscaled cutoff for absolute error
$a^N$ is $K+3\ge108\pi^2\mathcal V^2\gamma a^{-(N+6)}$.
This is a direct application of the inherited reconstruction theorem,
not a new evaluation of the odd reflected covariance.

More generally enumerate polynomials of increasing total degree in
\eqref{f16v86:generators} after one common finite flow time.  Symmetry
projections are finite averages on these functions.  For instance
$\Pi_{0,\pm}=\frac12(I\pm\mathsf I)\Pi_0$ gives separate even and odd
banks.  Their cross Grams vanish at all native time separations, including
every V84 Hankel and localizer entry.  Exact dependencies are taken on
function-space quotients; a numerical rank threshold is not a proof of
coverage or of spectral exhaustion.  Other lattice symmetry sectors must
also be retained if the full invariant Hilbert space is the target.

On a torus, flipping signs of direction-$j$ links crossing one fixed
periodic cut leaves every elementary plaquette unchanged.  It is an
isometry commuting with the transfer and the spatial flow.  These three
commuting involutions define the flat-center group $\mathbb Z_2^3$.
Every contractible loop, including \eqref{f16v86:loop}, is neutral.  The
full rooted generator algebra includes winding loops and covers all these
sectors; averaging that algebra over the center group is dense in the
neutral sector alone.  Thus a deliberate neutral-sector restriction can be
implemented exactly, but cannot exclude the missing odd example above.
No assertion is made about the continuum energy of finite-torus winding
states.

\subsection{Physical scope and the next source packet}

The new native calculation is the missing-sector Gram
\eqref{f16v86:native-entrance}--\eqref{f16v86:delay-bound}, together with
its exact orthogonality to the entire old linear family.  The separating
algebra specifies a finite-regulator route to actual source coverage.
It does not convert an estimate on one sector into an estimate on the
others, or prove a common continuum constant across its growing products.

The next declared physical packet should therefore include both an even
simultaneous-flow source and \eqref{f16v86:flowed-odd}, with their actual
positive-time Grams and their V84 forward localizers treated in their
separate exact sectors.  A nonzero physically normalized odd limit has not
been proved; neither has a theorem eliminating that sector from the
limiting local vacuum Hilbert space.  If a larger algebra is used, its
products, spatial support, normalization, and limit order require their
own estimates.  Finite-regulator uniform density after an invertible flow
is not uniform density of limiting physical sources.  The dense-core step
must eventually be justified in the constructed limiting Hilbert space,
with common upper-contraction constants.  No earlier positive-time native
covariance problem, and no continuum mass gap, is solved merely by adding
a symmetry label.

The companion diagnostic is
\begin{center}\small
\path{verify_ym_f16_v86_parity_coverage.py}.
\end{center}
It separates exact lattice-marker incidence, quaternion and trace algebra,
finite symmetry identities, compact conditional-score quadrature, and
finite positive-matrix spectral checks.  Full spatial-flow tests use the
literal simultaneous force, not a plaquette subflow.  These are not full
Wilson-vacuum samples, interval certificates or proof-assistant
verification.  No numerical output is a premise of the proofs.  Removing
this delimited appendix and reversing the one displayed version change
recovers V85 byte for byte.

\begin{record}[V86 result ledger]
\label{f16v86:ledger}
\emph{Proved here:} exact absence of zero-momentum odd parity from the
closed linear elementary-plaquette flow family; a contractible eight-link
native source outside it with exact Haar norm and nonzero finite-delay
pairings; a unique nonadjacent-link marker geometry and a volume-/mark-count-
uniform thermal-order native Gram for that shape family; its explicit
finite-delay lower bound; a single/pair/triple rooted-trace orbit separator
with reducible configurations included; finite-regulator dense flowed
source algebras and dense forward-buffered images; a retained odd-channel
reconstruction budget; and exact neutral-sector and parity bookkeeping.

\emph{Inherited or standard:} native transfer and spatial symmetry,
finite-regulator simultaneous-flow invertibility, V85's weighted-adjoint
score moments, V79's compact derivative-cost theorem, V77's input filter,
V84's cyclic versus dense-source distinction, symmetry selection rules,
Wilson-loop algebra density, and Stone--Weierstrass.  The preceding
archive is preserved, not independently recertified.

\emph{Not asserted:} that the linear symmetry obstruction applies to all
products of the old observables; a continuum spin assignment to the odd
shape; a missing physical pole or a new mass value; a continuum odd-source
limit from its thermal contact floor; a volume-uniform inverse flow;
spatial locality of the rooted separator on growing graphs; or a common
spectral constant from finite-regulator algebra density.

\emph{Still unproved:} physically normalized positive-time even and odd
native Grams with compatible continuum and thermodynamic limits; their
required growing-depth upper inequalities; dense limiting source coverage
with a common physical contraction constant; and the interacting
four-dimensional continuum Yang--Mills mass gap.
\end{record}

\begin{firewall}
The missing sector is proved under the same native law, not introduced by
changing the transfer.  Its thermal variance and finite-delay positivity
are kept separate from physical noncollapse.  Source multiplication and
gauge-orbit separation repair finite coverage, not the unestimated dynamics.
All V1--V85 mathematical text remains unchanged.
\end{firewall}
% END F16 V86 APPEND

% BEGIN F16 V87 APPEND -- 2026-09-19
% Append-only continuation of the Post-Fifteenth-Archive V86.
\clearpage
\section[V87: cubic pseudoscalar sources and their curvature-derivative limit]{V87: cubic pseudoscalar sources\\and their curvature-derivative limit}
\label{f16v87:context}

V86 supplies an odd-parity source but deliberately does not assign it a
continuum spin or a physical normalization.  Before estimating that source
as another scalar energy density, we identify its rotationally scalar
component and its first nonzero smooth-field coefficient.  The component
does not vanish.  Its coefficient is a magnetic curvature paired with a
covariant curl, not the dimension-four magnetic energy.  A three-link
score witness gives this complete rotated family a native covariance
floor.  The different power of the cutoff required by its smooth-field
normalization is then explicit.

Three statements below concern different objects.  The symmetry and score
bounds hold under the original compact Wilson law.  The small-loop limit
holds for links sampled from a declared smooth connection, with a uniform
remainder on bounded smooth families.  The final positive-time covariance
is computed in a separately defined free-curvature reference.  The exact
second variation of the simultaneous lattice-flow source connects the last
two calculations; it does not identify either reference with the native
vacuum.  In particular no smoothness or Gaussian law is assigned to native
flowed links.

\subsection{The full cubic projection and its literal marker geometry}

Keep V86's even periodic spatial lattice of side $L\ge8$, with $N_v$ vertices,
and its original stationary Wilson transfer $P$ and slice law $\pi$.  Let
$\mathcal O$ be the 24 proper signed permutation matrices in three dimensions,
and let $\mathcal O_h=\mathcal O\cup(-\mathcal O)$.  Use exactly V86's word
\[
 \Gamma=(1,2,1,3,-1,-1,-2,-3).
\]
For $Q\in\mathcal O_h$ let $Q\Gamma$ mean the transformed directed steps,
not a reversal of their order.  Put
\begin{equation}
 Z_{\rm ps}=\frac1{\sqrt{48N_v}}
       \sum_{x\in\Lambda}\sum_{Q\in\mathcal O_h}
                 \det(Q)\,w(U_{x+Q\Gamma}).
 \label{f16v87:Z}
\end{equation}
This is $\sqrt{24}$ times the proper-rotation average of V86's $Z_\Lambda$.
The label ``ps'' refers to the scalar representation of the proper cubic
group with odd spatial parity, usually denoted $A_1^-$.  A finite-lattice
representation is not by itself a unique continuum angular momentum.

Spatial symmetry projection and spatial gradient flow are established
methods; see Sakai and Sasaki,
\href{https://arxiv.org/abs/2211.15176}{arXiv:2211.15176}.
For the smooth expansion, covariant Taylor expansion in radial gauge is
standard; see Ivanov and Kharuk,
\href{https://arxiv.org/abs/1906.04019}{arXiv:1906.04019}, \S\S3.2--3.3.
The particular loop moments, cubic coefficient, native three-marker bound,
and reference covariance are derived here.  No spectroscopy or interacting
continuum theorem is imported from those sources.

\begin{proposition}[A nonzero cubic pseudoscalar and a unique three-link marker]
\label{f16v87:geometry}
All $48N_v$ stored unoriented edge sets in \eqref{f16v87:Z} are distinct.
For the untransformed word, the three stored edges
\begin{equation}
 e=((0,0,0),1),\qquad f=((0,0,1),2),\qquad g=((1,1,1),1)
 \label{f16v87:markers}
\end{equation}
are pairwise plaquette-separated in four dimensions.  They occur together
in exactly one translate of one shape $Q\Gamma$.  Their transformed images
have the same uniqueness property.  Consequently
\begin{equation}
 \mathbb E_H Z_{\rm ps}=0,\quad \|Z_{\rm ps}\|_H^2=\tfrac14,
 \quad \mathbb E_\pi Z_{\rm ps}=0,
 \quad \langle Z_{\rm ps},P^n Z_{\rm ps}\rangle_\pi>0\quad(n\ge0).
 \label{f16v87:rank}
\end{equation}
The source is gauge and translation invariant, invariant under every proper
cubic rotation, and odd under inversion.  These properties and finite-delay
strict positivity persist after any one finite simultaneous spatial flow.
\end{proposition}
\begin{proof}
The eight vertices are listed in V86.  The markers in
\eqref{f16v87:markers} are among its eight edges.  Two parallel markers
have transverse displacement in two coordinates; the perpendicular pairs
cannot be sides of one elementary plaquette.  Thus none of the three pairs
shares an action plaquette, including a temporal one.

For completeness, the following finite table makes the shape-isolation
check reproducible over the integers.  A signed permutation is specified
by its images $(a,b,c)$ of directions $(1,2,3)$, so its word is
$(a,b,a,c,-a,-a,-b,-c)$.  These are all its translated instances containing
the first two markers; only the indicated row also contains $g$.
\begin{center}\small
\begin{tabular}{@{}rrrlc@{}}
\toprule
$a$&$b$&$c$&translation&contains $g$\\
\midrule
1&-2&-3&$(-1,1,1)$&no\\
1&2&3&$(0,0,0)$&yes\\
2&-1&3&$(1,-1,0)$&no\\
2&1&-3&$(0,0,1)$&no\\
-1&-2&-3&$(1,1,1)$&no\\
1&-3&-2&$(0,1,1)$&no\\
-2&-1&3&$(1,1,0)$&no\\
2&3&-1&$(1,0,0)$&no\\
3&-1&2&$(1,0,-1)$&no\\
-3&1&2&$(0,0,1)$&no\\
3&2&1&$(0,0,0)$&no\\
-3&-2&1&$(0,1,2)$&no\\
\bottomrule
\end{tabular}
\end{center}
To obtain the table, map each of the eight stored edges to $e$, which fixes
the translation for each signed permutation, and retain exactly the cases
containing $f$.  Substitution of the eight tails gives the displayed list.
The third-edge test then proves uniqueness.  All relevant relative
coordinates have span less than eight, so no new equality is created on
the stated torus.  Transforming this argument proves uniqueness for every
shape and also excludes any duplicate stored-edge set.

Each loop has one once-occurring Haar link, so its trace has mean zero and
second moment $1/4$.  Distinct edge sets have a link in their symmetric
difference; integrating that link kills the cross moment.  The squared
coefficient sum in \eqref{f16v87:Z} is one, proving its Haar norm.  Proper
rotations permute the positive and negative terms separately; inversion
exchanges them.  Native symmetries give mean zero.  Full support of the
positive slice density gives nonzero native variance, and positive
injectivity of $P$ gives every finite-delay assertion.  The simultaneous
flow is a symmetry-equivariant, gauge-equivariant diffeomorphism at the
fixed regulator, so the same argument applies to its composition.  No
uniform inverse-flow estimate is used.
\end{proof}

\subsection{A native covariance bound for all rotated shapes}

Let $\mathcal L$ contain distinct translates of these 48 shapes, allowing
marks on several time slices.  For real coefficient columns set
\[
 F_i=\sum_{\ell\in\mathcal L}s_{i\ell}W_\ell,
 \qquad \mathcal S_{ij}=\sum_{\ell\in\mathcal L}s_{i\ell}s_{j\ell}.
\]
Choose the transformed marker triple for each loop.  The following result
uses the original Wilson action, not a conditionally independent loop law.

\begin{theorem}[A volume-uniform native cubic-family Gram]
\label{f16v87:native}
Under the native finite Wilson law, and uniformly conditional on admissible
exteriors with the chosen marker triples free,
\begin{equation}
 \boxed{\quad
 \operatorname{Cov}_\mu(F)\succeq
          \frac{1}{568620(1+\gamma)^3}\mathcal S.\quad}
 \label{f16v87:native-Gram}
\end{equation}
The bound passes to the same stationary-cylinder and local limits as V85.
It is independent of the number of marks and surrounding space--time
volume.  Its coefficient kernel is exactly $\ker\mathcal S$.  In
particular $\operatorname{Var}_\pi Z_{\rm ps}\ge[568620(1+\gamma)^3]^{-1}$.
The power three is a proved sufficient bound for this marker construction,
not an optimality assertion about the native variance.
\end{theorem}
\begin{proof}
Use the unit left generators $D_e^a$, native scores $s_e^a=D_e^a\log\mu$
and weighted adjoints $(D_e^a)^*=-D_e^a-s_e^a$ of V85.  Its original
one-link identity gives $\mathbb E(|s_e|^2\mid e^c)\le18\gamma$.
For a marker triple define
\[
 h_\ell^{abc}=D_e^aD_f^bD_g^cW_\ell,\qquad
 Y_\ell=\sum_{a,b,c}(D_e^a)^*(D_f^b)^*(D_g^c)^*h_\ell^{abc}.
\]
Three integrations by parts and marker uniqueness give
\begin{equation}
 \mathbb E[W_{\ell'}Y_\ell]
 =\mathbf1_{\ell=\ell'}\mathbb E(7-W_\ell^2),
 \qquad 6\le\mathbb E(7-W_\ell^2)\le7.
 \label{f16v87:dual}
\end{equation}
Here is the exact derivative contraction.  Transport all marked tangent
bases through the intervening unit links.  Each remains an orthonormal
imaginary quaternion basis.  If $S_j(P)$ is the sum of the squared scalar
parts after $j$ such independent insertions, then $S_0=w(P)^2$ and
\[
 S_1=1-w^2,\qquad S_2=2+w^2,\qquad S_3=7-w^2.
\]
The recursion follows from $\sum_{a=1}^3w(\mathbf e_aP)^2=1-w(P)^2$:
if $S_j=A+Bw^2$, then $S_{j+1}=3A+B-Bw^2$.
This proves the factor in \eqref{f16v87:dual} without a weak-field expansion.

No action plaquette contains two of $e,f,g$.  Conditional on their full
complement their link laws, and hence their local scores, are independent.
Cross-link derivatives of these three scores vanish.  Since
$\Delta_eW_\ell=-3W_\ell$, the full witness is exactly
\[
 Y_\ell=(3-s_e\cdot D_e)(3-s_f\cdot D_f)
                         (3-s_g\cdot D_g)W_\ell.
\]
Every contracted derivative of the multilinear trace has magnitude at most
the product of its tangent lengths.  Conditional score independence and
then the expectation tower give
\begin{equation}
 \|Y_\ell\|_2\le(3+\sqrt{18\gamma})^3,
 \qquad \|Y_\ell\|_2^2\le19683(1+\gamma)^3.
 \label{f16v87:witness-norm}
\end{equation}
The second inequality is $(3+\sqrt{18\gamma})^2\le27(1+\gamma)$.

The support of $Y_\ell$ lies in its eight edges and the three full
four-dimensional plaquette stars of its markers.  If a marker of $\ell$
is absent from the support of $Y_{\ell'}$, their covariance is exactly
zero by integration by parts.  A specified spatial edge belongs to at most
$48\cdot8/3=128$ translated loop instances, and is a marker in at most
$48\cdot3/3=48$ instances.  At most 19 stored edges, including the edge
itself, share an action plaquette with it.  Hence at most
$128+19\cdot48=1040$ witness supports contain a specified marker.  This
count includes all time layers.  By the row-sum bound and
\eqref{f16v87:witness-norm},
\[
 \operatorname{Cov}(Y)\preceq20470320(1+\gamma)^3I.
\]
The overlap of $\sum c_\ell W_\ell$ with $\sum c_\ell Y_\ell$ is at
least $6\sum c_\ell^2$.  Cauchy--Schwarz gives
$36/[20470320(1+\gamma)^3]=[568620(1+\gamma)^3]^{-1}$.
Apply this in every source direction.  Exact zero coefficient combinations
vanish as functions, and the lower bound excludes every other kernel.
All identities were conditional compact integrations under the original
law; there is no sequential assignment of a new Wilson marginal.
\end{proof}

For a spatial bank one may also populate V79's delay theorem.  A specified
edge occurs in at most 128 of these shapes, so its complete ambient
derivative matrix obeys
\[
 \mathsf B_\gamma\preceq128v_\gamma^{(8)}\mathcal S,
 \qquad v_\gamma^{(8)}=((1+4/\gamma)^8-1)/4.
\]
Thus, with $\delta=128\cdot568620(1+\gamma)^3v_\gamma^{(8)}$,
\begin{equation}
 C_F(n)\succeq q(\delta)^nG_\pi(F),\qquad
 q(\delta)=\frac2{\sqrt{\delta^2+4}+\delta}.
 \label{f16v87:delay}
\end{equation}
This direct application retains the exact source quotient and unchanged
physical time.  Its sufficient factor has order $\gamma^{-2}$ at large
$\gamma$; it is not a useful fixed-physical-time lower constant.

\subsection{The signed geometric moments of the eight-link contour}

The smooth calculation uses a $C^7$ connection $A_i(x)\in\operatorname{Im}
\mathbb H$ on a fixed periodic spatial continuum box.  Set
\[
 F_{ij}=\partial_iA_j-\partial_jA_i+[A_i,A_j],\quad
 B_i=\tfrac12\epsilon_{ijk}F_{jk},\quad D_jB_i=\partial_jB_i+[A_j,B_i].
\]
Use the quaternion inner product $v\cdot w=-\tfrac12\Tr(vw)$ on imaginary
quaternions.  A link is exact ordered parallel transport along its straight
segment, with the product convention whose small $ij$ plaquette is
$1+a^2F_{ij}+O(a^3)$.  Denote the trace around the signed-scale contour
$x+a\Gamma$ by $W_\Gamma(a,x)$; negative $a$ gives the inverted contour,
not its reversed traversal.

If $r_0,\ldots,r_8=r_0=0$ are the eight contour vertices, define fan moments
\begin{equation}
 \mathcal A=\frac12\sum_jr_j\times r_{j+1},\qquad
 \mathcal M=\frac16\sum_j(r_j\times r_{j+1})(r_j+r_{j+1})^{\mathsf T}.
 \label{f16v87:moments}
\end{equation}
The triangles may be regarded as an oriented surface chain; no assumption
that the contour is planar is made.  Direct rational evaluation gives
\begin{equation}
 \mathcal A=\begin{pmatrix}1\\-2\\1\end{pmatrix},\qquad
 \mathcal M=\begin{pmatrix}
 2/3&1/2&1/2\\-2&-4/3&-1\\1/2&1/2&2/3
 \end{pmatrix}.
 \label{f16v87:explicit-moments}
\end{equation}
In particular
\begin{equation}
 \sum_{i,j,l}\epsilon_{ijl}\mathcal A_i\mathcal M_{lj}=-1,
 \qquad
 \frac1{24}\sum_{R\in\mathcal O}
 (R\mathcal A)_i(R\mathcal M R^{\mathsf T})_{lj}
 =-\frac16\epsilon_{ijl}.
 \label{f16v87:cubic-tensor}
\end{equation}
For the second identity, a rank-three invariant tensor of the proper cubic
group must have three distinct indices: the half turns remove the other
components.  Right-angle rotations make the surviving components fully
antisymmetric.  Contracting with $\epsilon_{ijl}$ fixes the coefficient.
This also proves that the pseudoscalar moment cannot be a translated proper
rotation or a reversal of its opposite handed shape: the displayed scalar
is unchanged by translations and reversal, but changes sign under inversion.

\begin{theorem}[The leading odd field is a covariant curvature derivative]
\label{f16v87:smooth}
Uniformly on bounded $C^7$ families of the specified connections,
\begin{equation}
 \begin{split}
 W_\Gamma(a,x)-W_\Gamma(-a,x)
 ={}&-2a^5\sum_{i,l,j}\mathcal A_i\mathcal M_{lj}
                         B_i(x)\cdot D_jB_l(x)+O(a^7).
 \end{split}
 \label{f16v87:odd-expansion}
\end{equation}
Consequently the proper-rotation average
\[
 D_a(x)=\frac1{24}\sum_{R\in\mathcal O}
       [W_{R\Gamma}(a,x)-W_{-R\Gamma}(a,x)]
\]
satisfies
\begin{equation}
 \boxed{\quad
 D_a(x)=\frac{a^5}{3}\,\mathcal H_A(x)+O(a^7),\qquad
 \mathcal H_A=\sum_{i,j,l}\epsilon_{ijl}B_i\cdot D_jB_l
                      =B\cdot(D\times B).\quad}
 \label{f16v87:helicity}
\end{equation}
The remainder is deterministic and its constant depends on the stated
smooth bounds.  This theorem assigns no such bounds to the native law.
\end{theorem}
\begin{proof}
Choose radial gauge locally about $x$.  The radial transport formula, which
follows by differentiating transport along rays, is
\[
 A_i(x+y)=\tfrac12y_jF_{ji}(x)
             +\tfrac13y_jy_kD_kF_{ji}(x)+O(|y|^3).
\]
At the basepoint $A=0$, so covariant first derivatives agree with ordinary
ones.  Integrating the first two terms along $a\Gamma$ gives
\[
 \log U_{a\Gamma}
 =a^2\mathcal A_iB_i+a^3\mathcal M_{ij}D_jB_i+O(a^4).
\]
For the cubic term this is ordinary Stokes on the fan for the linearized
curvature.  The Bianchi identity $D_iB_i=0$ ensures that the first-order
curvature is closed; equivalently one can integrate the displayed radial
polynomial on each edge.  Commutators between line integrals start at
order $a^4$ and do not change either coefficient.

The logarithm is an imaginary quaternion near identity and
$w(e^X)=1-|X|^2/2+O(|X|^4)$.  Its order-$a^5$ trace coefficient is
$-\mathcal A_i\mathcal M_{lj}B_i\cdot D_jB_l$.  There are no trace
terms below order four.  Taylor-expand the smooth signed-scale trace
through order six; subtraction at $a$ and $-a$ cancels every even
coefficient, leaving an $O(a^7)$ remainder.  Smooth dependence of the
finite path ODE makes that remainder uniform under the stated $C^7$
bounds.  Traces are gauge invariant, so the local gauge calculation
proves \eqref{f16v87:odd-expansion} in any gauge.  Apply
\eqref{f16v87:cubic-tensor} for \eqref{f16v87:helicity}.
\end{proof}

The scalar $\mathcal H_A$ is gauge invariant, invariant under proper
spatial rotations and odd under inversion.  It has one more spatial
derivative than $B^2$; with $A$ of inverse-length dimension it has
dimension five.  It is not generally a spatial divergence.  On a periodic
box, the commuting connection
\[
 A_1=\alpha\cos(kx_3)\mathbf k,\quad
 A_2=\alpha\sin(kx_3)\mathbf k,\quad A_3=0
\]
has $\mathcal H_A=-\alpha^2k^3$ everywhere, for any compatible nonzero
wave number $k$.  Its spatial integral is nonzero.  This supplies an
explicit smooth test rather than an inferred continuum glueball state.

\subsection{Physical sampling and an exact simultaneous-flow second variation}

Let $\mathcal V=N_va^3$ be the fixed physical spatial volume.  The
smooth-field normalization determined by \eqref{f16v87:helicity} is
\begin{equation}
 \boxed{\quad
 \Xi_a=\frac{\sqrt3}{2}a^{-7/2}Z_{\rm ps}
 =\frac1{8a^2\sqrt{\mathcal V}}
    \sum_{x,R\in\mathcal O}
       [w(U_{x+R\Gamma})-w(U_{x-R\Gamma})].\quad}
 \label{f16v87:Xi}
\end{equation}
For links of a fixed smooth periodic connection,
\begin{equation}
 \Xi_a\longrightarrow\mathcal V^{-1/2}\int\mathcal H_A(x)\,d^3x.
 \label{f16v87:smooth-sum}
\end{equation}
Indeed $\Xi_a=(3/\sqrt{\mathcal V})a^3\sum_xD_a(x)/a^5$; the local
$O(a^2)$ error and the ordinary Riemann-sum limit prove the assertion.
A possible coupling-dependent multiplicative renormalization is a separate
choice and is not fixed by this geometric computation.

The native statement can be very different at zero flow.  From
\eqref{f16v87:native-Gram},
\begin{equation}
 \operatorname{Var}_\pi(\Xi_a)
 \ge\frac1{758160\,a^7(1+\gamma)^3}.
 \label{f16v87:contact}
\end{equation}
For example this lower bound diverges when $\gamma=O(\log(1/a))$.
This is a native contact lower bound, not a contradiction to
\eqref{f16v87:smooth-sum}: typicality of uniformly smooth input
connections was never assumed.  The bound is not assigned to $\Xi_a$
after fixed positive physical spatial flow.

Define the actual collective source
$\Xi_a^\tau=\Xi_a\circ\Phi_{a,\tau}$, with $\tau>0$ and unchanged
native centering (its odd symmetry makes the mean zero).  On oriented
link cochains let $d$ be the dimensionless lattice curl, $d^*$ its Haar-flat
adjoint, and let $\ell_Cv$ be the signed sum of link tangents around a loop.

\begin{proposition}[The exact quadratic part of the full flowed pseudoscalar]
\label{f16v87:flow-jet}
At identity the first derivative of $\Xi_a^\tau$ vanishes.  With
$v_\tau=e^{-(\tau/a^2)d^*d}v$, its exact second variation is
\begin{equation}
 \frac12D^2\Xi_a^\tau(\mathbf1)[v,v]
 =-\frac1{16a^2\sqrt{\mathcal V}}
   \sum_{x,R\in\mathcal O}
   \left[|\ell_{x+R\Gamma}v_\tau|^2
                    -|\ell_{x-R\Gamma}v_\tau|^2\right].
 \label{f16v87:exact-jet}
\end{equation}
If $v_e$ is the exact line integral of a fixed trigonometric-polynomial
connection $A$ along $e$, then as $a\downarrow0$ the right side tends to
\begin{equation}
 \mathcal V^{-1/2}\int
 (e^{\tau\Delta}\nabla\times A)\cdot
       \nabla\times(e^{\tau\Delta}\nabla\times A)\,d^3x.
 \label{f16v87:jet-limit}
\end{equation}
This is a first nonzero variation of the actual native source formula,
not its covariance under the interacting law.
\end{proposition}
\begin{proof}
The linearized full spatial Wilson flow is
$v'=-a^{-2}d^*dv$, by the exact Hessian at flatness.  A loop trace has
first derivative zero and second derivative $-|\ell_Cv|^2$ there.
The chain rule therefore has no term involving the second derivative of
the flow: it would multiply a zero first derivative of the output.
This proves \eqref{f16v87:exact-jet}.  Closed loop currents annihilate
exact lattice gradients.  On the complementary cochains the cubic Hodge
identity makes $d^*d$ the scalar lattice Laplacian, whose multiplier is
$4\sum_j\sin^2(ap_j/2)$.  Thus, on each fixed Fourier mode, its flowed
closed-loop current has multiplier tending to $e^{-\tau|p|^2}$.
The finite set of Fourier modes permits termwise passage to the limit.
The quadratic part of \eqref{f16v87:helicity} is obtained by replacing
$B$ with $\nabla\times A$ and $D$ with $\nabla$.  This yields
\eqref{f16v87:jet-limit}.  No uniform nonlinear flow approximation on
native rough connections follows from this finite-mode calculation.
\end{proof}

\subsection{A computed positive-time reference for the odd component}

To test that the newly identified quadratic component can have nonzero
reflected norm, specify an independent reference, with three color copies
of a transverse real Gaussian magnetic field.  Its Fourier covariance is
\begin{equation}
 \mathbb E[B_i^a(t,p)B_j^b(s,q)]
 =(2\pi)^3\delta(p+q)\delta_{ab}\frac{|p|}{2}
 (\delta_{ij}-\widehat p_i\widehat p_j)e^{-|p||t-s|}.
 \label{f16v87:free-law}
\end{equation}
Set $B_\tau=e^{\tau\Delta}B$, $\tau>0$, and
$h_\tau=\sum_a B_\tau^a\cdot(\nabla\times B_\tau^a)$.
Its expectation is zero by parity.  A periodic-box definition uses
$\mathcal V^{-1/2}\int h_\tau$, excludes the zero magnetic mode, and
then takes the ordinary infinite-volume limit of its covariance.
This law is not the Wilson measure, a conditional Wilson cell, or a
claimed perturbative evaluation of its full interacting expectation.

\begin{theorem}[Odd reference covariance and spectral density]
\label{f16v87:free}
The zero-momentum covariance density of the specified reference is
\begin{equation}
 \boxed{\quad
 \mathcal C_\tau(t)
 =\frac3{2\pi^2}\int_0^\infty p^6e^{-4\tau p^2-2|t|p}\,dp
 =\frac3{256\pi^2}\int_0^\infty E^6e^{-\tau E^2-|t|E}\,dE.\quad}
 \label{f16v87:free-C}
\end{equation}
In particular
\begin{equation}
 \mathcal C_\tau(0)=\frac{45}{4096\pi^{3/2}\tau^{7/2}},\qquad
 0<\mathcal C_\tau(t)\le\mathcal C_\tau(0)\quad(t\ge0).
 \label{f16v87:free-contact}
\end{equation}
Every prescribed pair of positive times has bounded earlier and strictly
positive later Grams.  The spectral weight has a Gaussian ultraviolet
tail but is nonzero arbitrarily close to zero.  In fact
\begin{equation}
 \mathcal C_\tau(t)\sim\frac{135}{16\pi^2}\,t^{-7}
                         \quad(t\longrightarrow\infty).
 \label{f16v87:free-tail}
\end{equation}
\end{theorem}
\begin{proof}
At finite Fourier cutoff Wick's identity for a centered real quadratic
form gives $2\Tr(KCKC)$ per color, with
$K(p)=i[p\times]$ and
$C(p)=(|p|/2)\Pi_T(p)e^{-2\tau p^2-|t||p|}$.
On the transverse plane $K^2=p^2\Pi_T$ and $\Tr\Pi_T=2$.
Thus the two Wick contractions together give $p^4e^{-4\tau p^2-2|t|p}$
per color.  Summing three colors and integrating $d^3p/(2\pi)^3$
gives the first formula.  A periodic box gives the same integrand in its
Fourier Riemann sum; Gaussian decay and its harmless behavior at zero
justify the infinite-volume and ultraviolet-cutoff limits.  The change
of variables $E=2p$ gives the positive spectral density in
\eqref{f16v87:free-C}.  The Gaussian sixth moment proves
\eqref{f16v87:free-contact}.  All positive-time derivatives follow by
integrable differentiation.  Finally set $y=tE$ and use dominated
convergence with $y^6e^{-y}$; since $\Gamma(7)=720$, the coefficient is
$3\cdot720/(256\pi^2)=135/(16\pi^2)$.
\end{proof}

For an explicit later lower bound one may retain just
$E\in[1/(2\sqrt\tau),1/\sqrt\tau]$ in \eqref{f16v87:free-C}:
\[
 \mathcal C_\tau(t)\ge
 \frac3{32768\pi^2\tau^{7/2}}e^{-1-t/\sqrt\tau}.
\]
This populates V82's two-positive-time comparison \emph{in this reference}.
It does not supply that comparison under the original Wilson history.
The finite continuum covariance also does not imply a mass gap: its exact
spectral density has support all the way down to zero.

\subsection{A symmetry-preserving native reconstruction and the remaining estimate}

At every finite simultaneous flow time, the source in \eqref{f16v87:Xi}
has half of its range bounded by $6\sqrt{\mathcal V}a^{-5}$.
Apply V77's common input-link filter to all $3\mathcal V/a^3$ original
spatial links.  Its native covariance error therefore obeys
\begin{equation}
 \boxed{\quad
 |C_{\Xi^\tau}(n)-C_{\mathcal A_K\Xi^\tau}(n)|
 \le\frac{1944\pi^2\mathcal V^2\gamma}
                   {a^{13}(K+3)},\qquad n\ge2.\quad}
 \label{f16v87:reconstruction}
\end{equation}
The proof is direct substitution in V77's $2\epsilon_KB_F^2$ bound:
$B_F^2\le36\mathcal V a^{-10}$ and
$\epsilon_K\le27\pi^2\mathcal V\gamma a^{-3}/(K+3)$.
All source values and covariances are still under the original law.  An
additional scalar normalization $\lambda_a$ multiplies this error by
$|\lambda_a|^2$.  For fixed physical volume and polynomial normalizations,
a sufficient cutoff for error $a^N$ remains polynomial, now with the
explicit power $a^{-(N+13)}$.  There is no claim of an efficiently sized
coefficient bank.

The common filter preserves $A_1^-$, and the unflowed function has exactly
the multiplier $[K/(K+3)]^8$.  Native even--odd cross covariances vanish at
every lag.  The positive-time Gram and V84 localizers of the added channel
must nevertheless be estimated separately from the even channel.  The
contact floor \eqref{f16v87:contact}, the smooth-field limit, and the free
reference are not three estimates of the same probability law.

The new upstream identification is precise: the rotationally scalar part
of V86's actual loop probes a dimension-five magnetic derivative at its
first smooth order.  Its physical sampling cannot be copied from the
scalar magnetic-energy bank.  The full nonlinear simultaneous-flow source
has the exact quadratic part \eqref{f16v87:exact-jet}; its higher variations
and their native correlations have not been bounded here.  The next
unpaid object is its original-law positive-separation covariance with
this declared normalization (and any explicitly chosen coupling factor),
not another symmetry or finite-regulator density statement.  A controlled
native comparison must in particular retain the evolving connection,
source-weighted rough configurations, and the distinction between a
formal smooth expansion and a limit of quantum expectations.

\subsection{Verification, preservation and result ledger}

The companion diagnostic is
\begin{center}\small
\path{verify_ym_f16_v87_pseudoscalar.py}.
\end{center}
It separates exact rational contour moments and cubic tensors, literal
marker and support enumeration, quaternion derivative contractions,
conditional score tests, smooth compact-holonomy checks, the full
simultaneous-flow quadratic identity, and the explicit Gaussian reference.
Random compact configurations are not Wilson-vacuum samples.  Numerical
outputs are diagnostics, not premises of the analytic proofs, interval
certificates, or proof-assistant verification.  Removing this delimited
appendix and reversing the one displayed version change recovers V86
byte for byte.

\begin{record}[V87 result ledger]
\label{f16v87:ledger}
\emph{Proved here on the specified objects:} the complete cubic pseudoscalar
projection and three-link marker isolation; a volume-uniform native
covariance floor; exact contour moments and the dimension-five coefficient
with its smooth-family remainder; the associated sampling and contact
bound; the exact simultaneous-flow quadratic part; a computed positive-time
free covariance with gapless spectral density; and the correctly normalized
native reconstruction budget.

\emph{Inherited or standard:} the loop and native symmetries from V86;
invertibility of the finite-regulator flow; V85's score method; V79's
delay inequality; V77's original-input reconstruction; radial-gauge
covariant Taylor expansion; cubic projection; and Gaussian moment pairing.
These methods do not recertify the previous archive or establish a native
Gaussian limit.

\emph{Not asserted:} a pure continuum spin from a cubic label; optimality of
the three-marker coupling exponent; a uniformly smooth native input or
flowed connection; substitution of the free covariance for the interacting
one; a physical reflected lower Gram from the contact bound or a classical
nonzero field value; or a mass gap from a positive finite reference Gram.

\emph{Still unproved:} canonically normalized positive-time native Grams
for the even and the new cubic-odd simultaneously flowed sources; their
spatial and thermodynamic/continuum control; the growing-depth upper
inequalities and dense limiting coverage with a common physical constant;
and the interacting four-dimensional continuum Yang--Mills mass gap.
\end{record}

\begin{firewall}
The original source is projected and normalized explicitly, not replaced by
a reference field.  Native, deterministic smooth, and Gaussian-reference
claims keep their separate hypotheses and clocks.  Source coverage does
not imply source noncollapse or spectral contraction.  All V1--V86
mathematical text remains unchanged.
\end{firewall}
% END F16 V87 APPEND


% BEGIN F16 V88 APPEND -- 2026-09-19
% Append-only continuation of the Post-Fifteenth-Archive V87.
\clearpage
\section[V88: nonlinear ultraviolet drift in the spatial-flow reference]{V88: nonlinear ultraviolet drift\\in the spatial-flow reference}
\label{f16v88:context}

V87 computes a finite positive-time covariance for the quadratic part of
its odd source.  It does not control the higher variations of the
simultaneous flow under a rough input law.  We test that missing step,
rather than compute another smooth contour coefficient.  The test uses
\emph{the same free magnetic covariance as V87}, completed to a transverse
Gaussian connection, and the full nonlinear spatial Yang--Mills equation.
At third order in the input amplitude its first Gaussian-chaos component
has a logarithmic ultraviolet divergence.  Both the iterated quadratic
term and the cubic term are included: they contribute three and minus two
copies, respectively, of the surviving coefficient.

For the gauge-invariant curvature-derivative source this gives a logarithmic
sixth-order covariance at strictly positive physical time separation.
It is not a contact-only term.  A subtraction in the third-order reference
jet makes that jet converge, but is not prescribed as a counterterm for the
native Wilson measure.  The original lattice source and law remain
unchanged.  The conclusions are a calculated obstruction to an unrenormalized
Gaussian-input perturbative comparison, not a theorem that the native
simultaneous-flow source fails to exist.

\subsection{Reference law, conventions, and the nonlinear equation}

Use the spatial torus of side $2\pi$, volume $V=(2\pi)^3$, and Fourier
normalization
\[A(x)=V^{-1/2}\sum_p\widehat A(p)e^{ip\cdot x}.\]
Let $T^a$, $1\le a\le3$, be the imaginary quaternion basis, with
$[T^a,T^b]=2\epsilon_{abc}T^c$.  Write $f^{abc}=2\epsilon_{abc}$ and
$C_A=8$, so $\sum_{b,c}f^{abc}f^{dbc}=C_A\delta_{ad}$.
Set $\Pi(p)=I-pp^{\mathsf T}/|p|^2$ for $p\ne0$.

At physical times $t,s$ take a jointly real Gaussian connection with zero
spatial zero mode and covariance
\begin{equation}
 \mathbb E[\widehat A_i^a(t,p)\widehat A_j^b(s,q)]
 =\mathbf1_{q=-p}\delta_{ab}\,
       \frac{\Pi_{ij}(p)}{2|p|}e^{-|p||t-s|}.
 \label{f16v88:input}
\end{equation}
The cutoff $A_\Lambda$ retains $0<|p|\le\Lambda$.  The curl of this
connection has exactly V87's magnetic covariance.  This is a specified
Gaussian reference, not a gauge fixing or a marginal of the native Wilson
law.  Its covariance is $|p|^{-1}$, not the $|p|^{-2}$ covariance of the
three-dimensional Euclidean Gaussian free field.

For each smooth cutoff input, evolve each physical slice independently in
spatial flow time $v$ by the DeTurck form
\begin{equation}
 \partial_v A_i=\Delta A_i+
 \sum_j\{2[A_j,\partial_jA_i]-[A_j,\partial_iA_j]
                         +[A_j,[A_j,A_i]]\}.
 \label{f16v88:flow}
\end{equation}
The DeTurck term is a gauge motion for smooth solutions; gauge-invariant
outputs agree with those of the ordinary spatial Yang--Mills flow for the
same smooth initial connection.  This equation and its gauge convention
are standard; see Chevyrev, \emph{Norm inflation for a non-linear heat
equation with Gaussian initial conditions},
\href{https://doi.org/10.1007/s40072-023-00317-6}{Stoch.
PDE: Anal.
Comp.
\textbf{12} (2024), 1745--1768}, Example 2.2.  The rough-data results of
Cao and Chatterjee,
\href{https://arxiv.org/abs/2111.10652}{arXiv:2111.10652}, include the 3D
Gaussian free field; they are not imported for the different covariance
\eqref{f16v88:input}.  No norm-inflation theorem from either source is used
below.

For an input $gA_\Lambda$, define the amplitude jets at $g=0$ by
\begin{equation}
 A(v;gA_\Lambda)=gX_\Lambda(v)+g^2Y_\Lambda(v)
                         +g^3Z_\Lambda(v)+O_\Lambda(g^4).
 \label{f16v88:jets}
\end{equation}
This is a finite-cutoff derivative identity; its remainder is not uniform
in $\Lambda$.  Let $Q$ be the symmetric polarization of the quadratic
part of \eqref{f16v88:flow}, and let
$R_i(X)=\sum_j[X_j,[X_j,X_i]]$.  Exactly,
\begin{align}
 X(v)&=e^{v\Delta}A_\Lambda,\\
 Y(v)&=\int_0^v e^{(v-r)\Delta}Q(X(r),X(r))\,dr,\\
 Z(v)&=\int_0^v e^{(v-r)\Delta}
             \{2Q(X(r),Y(r))+R(X(r))\}\,dr.
 \label{f16v88:Duhamel}
\end{align}
The cubic source term and the second use of the quadratic term are
separate summands of one jet, not alternative approximations.  Fourier
projection is imposed on the \emph{initial field only}; intermediate
products retain their generated modes.

\subsection{The complete first-chaos coefficient is logarithmic}

Let $\mathcal P_1$ denote orthogonal projection onto the first Gaussian
chaos of the initial variables.  Define the explicit divergent scalar
\begin{equation}
 d_\Lambda=\frac{C_A}{6V}
               \sum_{0<|q|\le\Lambda}|q|^{-3}
 =\frac{C_A}{12\pi^2}\log\Lambda+O(1).
 \label{f16v88:d}
\end{equation}
The equality is an asymptotic of a lattice sum, not a fitted coefficient.

\begin{theorem}[Full third-order drift, including the quadratic return]
\label{f16v88:drift}
For each fixed $v>0$ and every Sobolev index $m$, the fields
\begin{equation}
 X_\Lambda(v),\qquad Y_\Lambda(v),\qquad
 Z_\Lambda^{\rm sub}(v):=Z_\Lambda(v)-d_\Lambda X_\Lambda(v)
 \label{f16v88:subtracted}
\end{equation}
converge in $L^p(\Omega;H^m)$, for every finite $p\ge1$, as
$\Lambda\to\infty$.  Convergence is uniform when $v$ ranges over a compact
subset of $(0,\infty)$.  In particular
\begin{equation}
 \boxed{\quad
 \mathcal P_1 Z_\Lambda(v)=d_\Lambda X_\Lambda(v)+O_{v,m,L^p}(1).
 \quad}
 \label{f16v88:full-drift}
\end{equation}
The two contributions in \eqref{f16v88:Duhamel} obey separately
\begin{equation}
 \mathcal P_1 Z^{Q}_\Lambda=3d_\Lambda X_\Lambda+O(1),\qquad
 \mathcal P_1 Z^{R}_\Lambda=-2d_\Lambda X_\Lambda+O(1).
 \label{f16v88:three-minus-two}
\end{equation}
Here the bounded remainders converge in the same spaces.  For $C_A=8$,
the surviving logarithmic coefficient is $2/(3\pi^2)$.
\end{theorem}
\begin{proof}
We provide the momentum contraction and the summability argument.  The
symmetrized quadratic vertex, with input momenta $q,r$ and spatial input
indices $j,k$, is $if^{abc}S_{i;jk}(q,r)$, where
\begin{equation}
 S_{i;jk}(q,r)=r_j\delta_{ik}-q_k\delta_{ij}
                          +\tfrac12(q_i-r_i)\delta_{jk}.
 \label{f16v88:vertex}
\end{equation}
Symmetrization uses the antisymmetry of $f^{abc}$; it does not replace the
non-Abelian bracket by a scalar product.  The first-chaos part of $Y$ is
zero, and its mean is zero, by the color contraction.

Fix an external nonzero momentum $p$ and a unit polarization $e\perp p$.
After factoring out $e^{-v|p|^2}$, the coefficient of the corresponding
initial color component in $\mathcal P_1Z$ is exactly
\begin{equation}
 \frac{C_A}{V}\sum_{0<|q|\le\Lambda}\frac1{2|q|}
 \left[4\mathcal T(p,q,e)I_v(A,B)
                 -\mathcal U(q,e)J_v(2|q|^2)\right],
 \label{f16v88:exact-coefficient}
\end{equation}
where
\begin{align}
 A&=2|q|^2-2p\cdot q,\qquad B=2p\cdot q,\\
 J_v(z)&=\int_0^v e^{-sz}\,ds,\\
 I_v(A,B)&=\int_0^v e^{-sA}\int_0^s e^{-rB}\,dr\,ds,\\
 \mathcal T(p,q,e)&=p^{\mathsf T}\Pi(q)p
 +\tfrac12(|q|^2+p\cdot q)e^{\mathsf T}\Pi(q)e
 +2(q\cdot e)^2,\\
 \mathcal U(q,e)&=1+(q\cdot e)^2/|q|^2.
 \label{f16v88:kernel-data}
\end{align}
All removable zero denominators are interpreted by these integrals.
For example $I_v(A,B)=[J_v(A)-J_v(A+B)]/B$ if $B\ne0$.

For the quadratic return there are two positions of $Y$ and two ways to
contract its internal first-chaos field with the other $X$, giving the
factor four.  Its color contraction is
$f^{abc}f^{cbe}=-C_A\delta_{ae}$; the two Fourier derivatives supply the
second minus sign.  The spatial contraction before evaluation is
\[
 \sum_{i,j,k,l,m}e_i S_{i;jk}(q,p-q)\Pi_{jl}(q)
                         S_{k;lm}(-q,p)e_m=\mathcal T(p,q,e).
\]
For the cubic term the two nonzero Wick contractions give
$-C_A[(\operatorname{tr}\Pi)I-\Pi]$, yielding $-\mathcal U$.
The third possible contraction is zero in color.  These computations
prove \eqref{f16v88:exact-coefficient}; its off-diagonal version follows
from the same two-vertex tensor contraction.

For $|q|>4(1+|p|)$, replace $I_v(A,B)$ by
$[A(A+B)]^{-1}$ and $J_v(2|q|^2)$ by $(2|q|^2)^{-1}$.
The errors decay exponentially in $v|q|^2$.  The change from $p$ to zero
in the leading coefficients is $O_{v,p}(|q|^{-4})$ after the factor
$(2|q|)^{-1}$ is included; this is summable in three dimensions.
At $p=0$ the two leading matrix contractions are
\begin{equation}
 \frac{I+3\widehat q\widehat q^{\mathsf T}}{4|q|^3},
 \qquad
 -\frac{I+\widehat q\widehat q^{\mathsf T}}{4|q|^3}.
 \label{f16v88:leading-tensors}
\end{equation}
The cutoff ball is invariant under signed permutations, so exactly
\[
 \sum_{0<|q|\le\Lambda}\frac{q_iq_j}{|q|^5}
 =\frac{\delta_{ij}}3\sum_{0<|q|\le\Lambda}|q|^{-3}.
\]
The first tensor sums to $3d_\Lambda I$, the second to $-2d_\Lambda I$.
Their sum is $d_\Lambda I$.  At nonzero $p$ it acts on the initial
transverse vector.  Comparing $|x|^{-3}$ on unit cubes outside a fixed
ball with its radial integral gives $4\pi\log\Lambda+O(1)$; the cube
errors are summable because its gradient is $O(|x|^{-4})$.
This proves the coefficient and fixed-mode convergence.

Here are details that also justify the Gaussian-chaos and spatial limits.
Write $\langle q\rangle=1+|q|$.  After removing first-chaos contractions,
the second and third jets have respectively second and third homogeneous
Wick kernels.  Heat integration and the vertex bound
$|S(q,r)|\le C(|q|+|r|)$ give, for a fixed output $p$, the bounds
\begin{align*}
 |K_2(p;q,r)|&\le
 \frac{C_{v,p}}{\langle q\rangle+\langle r\rangle},
                   &&q+r=p,\\
 |K_3(p;q_1,q_2,q_3)|&\le C_{v,p}
 \sum_{\mathrm{cyc}}
 \frac1{(\langle q_1\rangle+\langle p-q_1\rangle)
                      (\langle q_2\rangle+\langle q_3\rangle)},
                   &&q_1+q_2+q_3=p.
\end{align*}
The cubic vertex without derivatives is included in the second bound.
To see the bounds, outside $|q|\le4(1+|p|)$ a quadratic time integral
pays $c|q|^2$, while its numerator costs at most $C|q|$.
For a nested vertex apply this argument first to the inner pair and then
to its sum and the remaining input.  On the complementary bounded-momentum
regions use the interval length in place of a denominator.  The positive
heat exponents along a tree are at least $v|p|^2/3$ in total, by
Cauchy--Schwarz at each split.  Reserving half of that damping shows that
$C_{v,p}$ can be taken as a fixed polynomial in $\langle p\rangle$ times
$C_v e^{-v|p|^2/6}$, uniformly on positive compact $v$-intervals.

For clarity the potentially delicate inner high--high sum is paid by
\begin{equation}
 \sum_q\frac1{\langle q\rangle\langle r-q\rangle
              (\langle q\rangle+\langle r-q\rangle)^2}
       \le\frac{C}{\langle r\rangle}.
 \label{f16v88:convolution-bound}
\end{equation}
Split into $|q|\le|r|/2$, $|r-q|\le|r|/2$, their remaining region of
size $O(|r|)$, and dyadic outer shells.  Each contributes at most the
right side, also for bounded $r$.  Wick's isometry and
\eqref{f16v88:convolution-bound} bound the squared third-chaos kernel
by a constant times
$\sum_{q_1}\langle q_1\rangle^{-3}\langle p-q_1\rangle^{-1}<\infty$.
The second-chaos bound uses the same convolution directly.  All cutoff
kernels converge pointwise, so dominated summation proves convergence of
these chaoses.  The first-chaos remainder was already summable above,
with polynomial output cost and the reserved heat damping.  Thus every
$H^m$ norm is controlled.  Gaussian polynomial moment bounds, obtained
by the finite Wick pairing formula in each fixed chaos, upgrade $L^2$
convergence to every finite $L^p$.  The estimates are uniform on compact
positive flow intervals, completing the proof.
\end{proof}

The subtraction in \eqref{f16v88:subtracted} is a statement about the
\emph{third derivative of this reference map}.  It is not yet a
renormalized nonlinear solution.  Removing only the cubic tadpole would
leave $3d_\Lambda X$; deleting both terms without computing their return
would also give the wrong answer.  Since the divergent coefficient
multiplies $X(v)$ independently of positive $v$, it is a boundary-amplitude
term at this order.  A constant bulk mass insertion gives instead
$-m_\Lambda^2 vX(v)$ and cannot cancel it at two distinct flow times with
one $v$-independent coefficient.

\subsection{A positive-separation divergence of the actual nonlinear reference source}

Use the smooth gauge-invariant functional identified in V87,
\begin{equation}
 \mathcal H(A)=V^{-1/2}\int_{\mathbb T^3}
 B_i(A)\cdot\epsilon_{ijl}D_j(A)B_l(A)\,dx,
 \quad B_i(A)=\epsilon_{ijk}\partial_jA_k+\tfrac12\epsilon_{ijk}[A_j,A_k].
 \label{f16v88:source}
\end{equation}
The flow is applied before evaluating this full functional, not just its
quadratic part.  At each cutoff write
\begin{equation}
 \mathcal H(A(v;gA_\Lambda))
 =g^2H_{2,\Lambda}+g^3H_{3,\Lambda}+g^4H_{4,\Lambda}+O_\Lambda(g^5).
 \label{f16v88:source-jets}
\end{equation}
Spatial parity of the reference law makes the mean of the full source and
of each jet zero.  At physical separation $t\ge0$ define
\begin{align}
 C_{4,\Lambda}(t)&=\mathbb E[H_{2,\Lambda}(0)H_{2,\Lambda}(t)],\\
 C_{6,\Lambda}(t)&=\mathbb E[H_{3,\Lambda}(0)H_{3,\Lambda}(t)]
 +\mathbb E[H_{2,\Lambda}(0)H_{4,\Lambda}(t)]
 +\mathbb E[H_{4,\Lambda}(0)H_{2,\Lambda}(t)].
 \label{f16v88:cov-jets}
\end{align}
These are Taylor coefficients, not an assertion of a cutoff-uniform
convergent series at nonzero $g$.

\begin{theorem}[The first nonlinear covariance coefficient is not ultraviolet finite]
\label{f16v88:covariance}
At fixed $v>0$,
\begin{equation}
 \boxed{\quad
 C_{4,\Lambda}(t)\longrightarrow C_{4}(t)
 =\frac3V\sum_{p\ne0}|p|^4e^{-4v|p|^2-2t|p|}>0,
 \quad}
 \label{f16v88:C4}
\end{equation}
and
\begin{equation}
 \boxed{\quad
 C_{6,\Lambda}(t)-4d_\Lambda C_{4,\Lambda}(t)
                 \text{ has a finite limit},\qquad
 \frac{C_{6,\Lambda}(t)}{\log\Lambda}
        \longrightarrow\frac{C_A}{3\pi^2}C_4(t).
 \quad}
 \label{f16v88:C6}
\end{equation}
The statements hold also at every strictly positive physical separation,
so the divergent coefficient is not supported only at contact.  For the
quaternion convention the coefficient is $8/(3\pi^2)$ times $C_4(t)$.
\end{theorem}
\begin{proof}
Separate the functional \eqref{f16v88:source} into its homogeneous parts
$\mathcal H^{[j]}$ in the connection; its first nonzero part is quadratic.
Let $\mathcal B$ be the symmetric bilinear polarization of that quadratic
part.  Substitution of the flow jets gives exactly
\begin{align*}
 H_2&=\mathcal B(X,X),\\
 H_3&=2\mathcal B(X,Y)+\mathcal H^{[3]}(X),\\
 H_4&=2\mathcal B(X,Z)+\mathcal B(Y,Y)
             +D\mathcal H^{[3]}(X)[Y]+\mathcal H^{[4]}(X).
\end{align*}
All commutator terms in the curvature and the covariant curl are retained
in the last two terms.  The drift theorem gives
\begin{equation}
 H_{4,\Lambda}=2d_\Lambda H_{2,\Lambda}
                         +H_{4,\Lambda}^{\rm sub},
 \label{f16v88:source-subtraction}
\end{equation}
where $H_2,H_3,H_4^{\rm sub}$ converge in every finite $L^p$.
Indeed they are finite products of derivatives of the convergent smooth
positive-flow fields $X,Y,Z^{\rm sub}$; take a sufficiently large Sobolev
index and use its product bound, followed by the moment bounds in the
preceding proof.  This establishes convergence without omitting any fourth
order source term.  The two physical slices have the fixed joint Gaussian
law \eqref{f16v88:input}; Cauchy--Schwarz justifies every mixed limit.
Inserting \eqref{f16v88:source-subtraction} in \eqref{f16v88:cov-jets}
proves its first assertion and the factor four.  V87's two Wick contractions,
or their Fourier sum on this torus, give \eqref{f16v88:C4}.  Its Gaussian
tail makes $d_\Lambda(C_{4,\Lambda}-C_4)\to0$.
Finally use \eqref{f16v88:d}.  At $t>0$ the positive factor $C_4(t)$
remains, proving that the divergence is not a pure contact distribution.
\end{proof}

The free infinite-volume expression inherited from V87 is
\[
 C_4(t)=\frac3{256\pi^2}\int_0^\infty
        E^6e^{-vE^2-tE}\,dE
\]
when its Fourier sum is passed separately to infinite volume.  The new
logarithmic coefficient was proved at fixed spatial volume.  No uniform
bound on the finite part as the volume grows is inferred by interchanging
these two limits.

\subsection{What is and is not changed on a bare-coupling trajectory}

\begin{proposition}[Vanishing bare amplitude alone does not pay the computed correction]
\label{f16v88:scaling}
Let $g_\Lambda>0$ tend to zero.  Normalize the reference source by
$g^{-2}$, so its covariance has formal
coefficients $C_{4,\Lambda}+g^2C_{6,\Lambda}+\cdots$.
For any positive $t,v$ the computed order-$g^2$ correction tends to zero
relative to $C_4(t)$ only if
\begin{equation}
                         g_\Lambda^2\log\Lambda\longrightarrow0,
 \label{f16v88:small-parameter}
\end{equation}
provided it is this coefficient, with no compensating input or source
counterterm, that is being used.  If
$g_\Lambda^2\log\Lambda\to b\in(0,\infty)$, its limit is
$(C_A b/(3\pi^2))C_4(t)$, not zero.  If that product diverges the
coefficient contribution diverges.
\end{proposition}
\begin{proof}
Equation \eqref{f16v88:C6} has a bounded finite remainder, and
$C_4(t)>0$.  Multiplication by $g_\Lambda^2\to0$ proves the statements.
They concern this coefficient, not the sum of an unproved uniform series.
\end{proof}

For example, taking $\Lambda\asymp a^{-1}$,
$g_a^2\asymp\gamma_a^{-1}$ and
$\log(1/a)=c\gamma_a+O(\log\gamma_a)$ leaves the logarithmic expansion
parameter of order one.  This is a scale comparison under declared
hypotheses, not a derivation of the native running coupling.  The higher
coefficients may contain further logarithms; no conclusion about their
resummation or the fixed-nonzero-amplitude limit is proved here.

The third-order subtraction admits a concrete reference interpretation.
Replacing the initial amplitude by
$gA_\Lambda-g^3d_\Lambda A_\Lambda$ changes its third flow jet to
$Z^{\rm sub}$ and its fourth source jet to
$H_4-2d_\Lambda H_2$.  Equivalently a formal source factor
$1-2g^2d_\Lambda$ removes the displayed sixth-order covariance logarithm.
Neither prescription is assigned to the original Wilson source: the
native input correlations, field normalization, and interaction corrections
have not been calculated in this reference.  Wick ordering just the local
cubic monomial is not this complete prescription.

There is also no contradiction to V87's exact lattice quadratic jet or
its smooth $a\to0$ contour expansion.  At each fixed Fourier cutoff the
reference input is smooth and its small-amplitude jets are well defined.
The new divergence occurs when that cutoff is then removed in a higher
jet.  Bounds on bounded $C^7$ families are not uniform on these Gaussian
cutoff fields.  One-time smoothness of a deterministic output and finiteness
of its quadratic Gaussian covariance therefore do not establish a
nonlinear native covariance comparison.

\subsection{The next native matching problem and verification}

The positive-separation divergence identifies a previously unpaid
boundary contribution in the proposed comparison with the free reference.
For the native even and odd simultaneous-flow sources, the next comparison
must transport the input law and the source normalization at the same
perturbative order, or bound the native positive-time covariance directly.
The transverse Gaussian law \eqref{f16v88:input} is not a substitute for
the Wilson law at that order.  In particular the coefficient
$C_A/(3\pi^2)$ above is \emph{not} asserted to be a native anomalous
dimension or a Wilson beta-function coefficient.

The native loop source, its positive physical-time placement, its cubic
sector, and V77's finite reconstruction are unmodified.  Their earlier
finite-regulator estimates remain valid.  No new native positive-time
lower Gram, growing-depth upper localizer, or common dense-core contraction
is supplied by this reference computation.  It prevents inferring any of
those conclusions by stopping at the finite quadratic covariance.

The companion diagnostic is
\begin{center}\small
\path{verify_ym_f16_v88_nonlinear_reference.py}.
\end{center}
It separates exact color and vertex contractions, Duhamel identities,
lattice momentum sums, cubic-symmetry averages, Gaussian-chaos algebra,
reference covariance sums and cutoff asymptotics.  No numerical result is
a premise of a proof.  The Gaussian and finite algebra fixtures are not
Wilson-vacuum samples, interval certificates, or proof-assistant
verification.  Removing this appendix and reverting the displayed version
recovers V87 byte for byte; the earlier archive is preserved, not recertified.

\clearpage
\begin{record}[V88 result ledger]
\label{f16v88:ledger}
\emph{Proved here for the specified reference:} the complete first-chaos
third-order spatial-flow coefficient; its three-minus-two cancellation
and explicit logarithmic lattice sum; convergence of the subtracted
third-order jet at positive flow time; the full sixth-order pseudoscalar
covariance logarithm, including the higher source terms; its persistence
at positive physical separation; and the explicit boundary/source
subtractions and expansion-parameter comparison at this order.

\emph{Inherited or standard:} V87's loop normalization and odd functional;
its quadratic field and magnetic covariance; the smooth DeTurck equation;
Duhamel expansion; Wick pairing and fixed-chaos moment bounds.  Results
for 3D Gaussian-free-field inputs are not used for the present law.

\emph{Not asserted:} divergence or convergence of the full nonlinear flow
at fixed nonzero amplitude; a native Wilson ultraviolet counterterm;
equality of native and Gaussian higher correlations; interchange of
thermodynamic and ultraviolet limits; a uniform perturbative remainder
on the physical trajectory; or a mass gap from either a reference
subtraction or a positive quadratic covariance.

\emph{Still unproved:} the canonically normalized native positive-time
Grams of the simultaneous-flow banks, their coupled source and input
renormalization or direct nonperturbative comparison, spatial and
thermodynamic/continuum compatibility, dense limiting source coverage
with common upper-contraction constants, and the interacting continuum
Yang--Mills mass gap.
\end{record}

\begin{firewall}
A positive-flow Gaussian quadratic covariance does not pay the nonlinear
ultraviolet boundary contribution.  The logarithm is computed in its
stated reference, with both flow terms and the full source expansion.
It is not silently promoted to a theorem about the native Wilson measure.
All V1--V87 mathematical text and the original physical clock remain
unchanged.
\end{firewall}
% END F16 V88 APPEND

% BEGIN F16 V89 APPEND -- 2026-09-19
% Append-only continuation of the Post-Fifteenth-Archive V88.
\clearpage
\section[V89: joint-law matching and finite nonlinear covariance ratios]{V89: joint-law matching\\and finite nonlinear covariance ratios}
\label{f16v89:context}

V88 finds a common logarithm in a positive-separation covariance, but does
not calculate the native input correction at the same order.  There are two
questions to separate before interpreting that logarithm as a spectral
obstruction.  A common multiplicative divergence cancels from a normalized
temporal ratio at its first perturbative order.  Conversely, agreement of
\emph{all} one-slice laws and all linear two-time covariances does not fix
that ratio for a composite source.  We prove the cancellation and construct
an explicit stationary, reflection-positive family that proves the second
assertion.  Its change in the first nonlinear source coefficient is
computed, not left as an unspecified higher cumulant.

The construction is a reference comparison, not another Wilson law or a
new physical-time choice.  It also yields a convergent, positive finite-jet
reference observable with an explicit amplitude remainder.  That observable
is not identified with the full nonlinear flow at nonzero amplitude.
The last calculation states exactly which normalized joint-law insertions
would be needed for a native matching calculation.

\subsection{Scope and the inherited subtracted jets}

Use V88's torus of volume $V=(2\pi)^3$, transverse Gaussian slice measure
$\nu$, quaternion convention $C_A=8$, and original Fourier cutoff on the
\emph{input} only.  Physical separation is denoted by $t$; $\tau_i>0$
are spatial flow times.  For a finite bank of these times write
\begin{equation}
 X_{i,\Lambda}=H_{2,\Lambda}^{\tau_i},\qquad
 Y_{i,\Lambda}=H_{3,\Lambda}^{\tau_i},\qquad
 Z_{i,\Lambda}=H_{4,\Lambda}^{\tau_i}
                         -2d_\Lambda H_{2,\Lambda}^{\tau_i},
 \qquad d_\Lambda=\frac{C_A}{6V}
                \sum_{0<|p|\le\Lambda}|p|^{-3}.
 \label{f16v89:jets}
\end{equation}
Thus $X,Y,Z$ here are \emph{source} jets, not V88's connection jets.
They have zero slice mean by spatial parity.  Under global field sign
reversal, $X,Z$ are even and $Y$ is odd.

The $L^p(\nu)$ convergence of \eqref{f16v89:jets}, for every finite $p$,
is inherited from Theorem~\ref{f16v88:drift} and its source substitution.
It is not independently recertified below.  Assertions about the exact
alternative input processes and their quadratic-source covariances do not
need that theorem.  Only assertions using $Y,Z$ in the cutoff limit use it.
In particular, no native Wilson Gaussian limit is among our hypotheses.

The Mehler/Hermite diagonalization used below is standard; for context see
J.~Teuwen, \emph{On the integral kernels of derivatives of the
Ornstein--Uhlenbeck semigroup},
\href{https://arxiv.org/abs/1509.05702}{arXiv:1509.05702}, and the generating
identity \href{https://dlmf.nist.gov/18.12.E16}{DLMF 18.12.16}.
We give the transition kernel and its elementary spectral proof.  The
refresh operation itself and Gaussian chaos calculus are not claimed as
new methods.

\subsection{The common logarithm drops out of normalized spectral data}

Let $P_t^0$ denote the Gaussian reference time semigroup, with each momentum
frequency $|p|$.  With column-synthesis notation in $L^2(\nu)$ set
\begin{align}
 A_\Lambda(t)&=X_\Lambda^*P_t^0X_\Lambda,\\
 B_\Lambda(t)&=Y_\Lambda^*P_t^0Y_\Lambda
       +X_\Lambda^*P_t^0Z_\Lambda+Z_\Lambda^*P_t^0X_\Lambda.
 \label{f16v89:AB}
\end{align}
For the amplitude-normalized \emph{full finite-cutoff} nonlinear source,
its covariance Taylor coefficients are
\begin{equation}
 A_\Lambda(t)+g^2\{4d_\Lambda A_\Lambda(t)+B_\Lambda(t)\}
                       +O_\Lambda(g^4).
 \label{f16v89:formal-C}
\end{equation}
This is a statement about derivatives at $g=0$, not a uniform remainder
for the full nonlinear flow.  Parity under field sign makes the covariance
even in $g$.  The matrices $A_\Lambda,B_\Lambda$ converge uniformly in
$t\ge0$, by the inherited $L^2$ convergence and semigroup contraction.

\begin{proposition}[Finite normalized coefficients despite the common logarithm]
\label{f16v89:ratio}
For a scalar bank with $A(t),A(t_0)>0$, the coefficient of $g^2$ in the
logarithm of the covariance ratio at $t,t_0$ has the finite limit
\begin{equation}
 \boxed{\quad
 \frac{B_\Lambda(t)}{A_\Lambda(t)}
       -\frac{B_\Lambda(t_0)}{A_\Lambda(t_0)}
 \longrightarrow
 \frac{B(t)}{A(t)}-\frac{B(t_0)}{A(t_0)}.\quad}
 \label{f16v89:ratio-limit}
\end{equation}
The term $4d_\Lambda$ cancels exactly at each cutoff.  For a finite matrix
bank the same cancellation holds in the first perturbation of any
positive-leading-Gram generalized eigenvalue problem at two times.  If
$A(t_1)v=\lambda A(t_0)v$, $v^*A(t_0)v=1$, and this eigenvalue is simple,
its limiting first perturbation is
\begin{equation}
                 v^*\{B(t_1)-\lambda B(t_0)\}v.
 \label{f16v89:pencil}
\end{equation}
At a fixed cutoff a multiple eigenvalue uses the analogous matrix on its
leading eigenspace.  The normalized matrix first jets converge also through
limiting degeneracies; convergence of separately differentiated ordered
eigenvalues across a closing degeneracy is not asserted.  The scalar
cancellation also applies to any fixed V84 forward moment pencil on a
declared positive leading quotient.
\end{proposition}
\begin{proof}
Differentiate the scalar logarithm at $g^2=0$.  For the pencil, differentiate
its equation and multiply by the normalized leading eigenvector.  The
common term gives $4d_\Lambda v^*(A_1-\lambda A_0)v=0$.
More generally factor the common scalar $1+4d_\Lambda g^2$ from both
members before taking their first jets.  This leaves the finite perturbations
$B_0,B_1$.  Hermitian perturbation on a positive metric gives the eigenspace formula
at each cutoff.  Whitening is a smooth matrix operation on a uniformly
positive leading metric, so its first jet converges.  Individual sorted
eigenvalue derivatives require the stated separation condition.  The same
scalar factor multiplies every covariance in a fixed temporal pencil;
no commutation of its matrix entries is required.
\end{proof}

For distinct positive flow times $\tau_i$, the leading matrix is explicitly
\begin{equation}
 A_{ij}(t)=\frac3V\sum_{p\ne0}|p|^4
       e^{-2(\tau_i+\tau_j)|p|^2-2t|p|}.
 \label{f16v89:A-matrix}
\end{equation}
It is strictly positive: a zero direction would make
$\sum_i c_i e^{-2\tau_i|p|^2}=0$ at every momentum.  Along $p=(n,0,0)$,
successively divide by the slowest exponential as $n\to\infty$ to force
every $c_i=0$.  This is positivity in the stated reference, not native data.
The proposition does not justify substituting a trajectory with
$g^2\log\Lambda=O(1)$ into a truncated expansion.  Higher logarithms and
nonuniform full-flow remainders are still unpaid.

\subsection{Different reflection-positive inputs with identical two-point data}

Group each pair $\{p,-p\}$ of nonzero spatial momenta into one real block:
its cosine and sine coefficients, both transverse polarizations, and all
three colors.  After the fixed covariance scaling, its measure is a standard
finite-dimensional Gaussian $\nu_p$.  Let $N_p$ be its Hermite number
operator, let $E_p$ average that \emph{entire} block, and put $\omega_p=|p|$.
For $0\le\alpha<1$ define
\begin{equation}
 L_\alpha=\sum_{\{p,-p\}}\omega_p
       \{(1-\alpha)N_p+\alpha(I-E_p)\},
 \qquad P_t^\alpha=e^{-tL_\alpha}.
 \label{f16v89:input-generator}
\end{equation}
A finite cutoff keeps only its indicated blocks.  The coefficients of the
connection itself retain exactly the normalization in
\eqref{f16v88:input}.

\begin{theorem}[Exact joint-law nonuniqueness beyond the covariance]
\label{f16v89:input-law}
The operators \eqref{f16v89:input-generator} define stationary reversible
Markov input processes with the same Gaussian law $\nu$ on every slice.
Their Euclidean histories are reflection positive.  For all $\alpha$ they
have \emph{exactly} V88's connection covariance at all physical times,
all momenta and all cutoff values:
\begin{equation}
 \mathbb E_\alpha[
 \widehat A_i^a(t,p)\widehat A_j^b(s,q)]
 =\mathbf1_{q=-p}\delta_{ab}\frac{\Pi_{ij}(p)}{2|p|}
                                            e^{-|p||t-s|}.
 \label{f16v89:same-two-point}
\end{equation}
For standardized coordinates $x(0),x(t)$ of one block of frequency $\omega$,
its mixed fourth cumulant is instead
\begin{equation}
 \boxed{\quad
 \operatorname{cum}_\alpha(x(0),x(0),x(t),x(t))
 =2\{e^{-(2-\alpha)\omega t}-e^{-2\omega t}\}.\quad}
 \label{f16v89:cumulant}
\end{equation}
It is strictly positive when $\alpha,t>0$.  The input is then not jointly
Gaussian in time, although its complete one-time distribution is Gaussian.
\end{theorem}
\begin{proof}
On one standardized block let
\[
 (\mathcal M_s f)(x)=\int f(e^{-s}x+\sqrt{1-e^{-2s}}y)\,d\nu_p(y).
\]
This is a positive, constant-preserving, reversible Markov kernel.  The
Hermite generating identity proves that its multiplier on degree $m$ is
$e^{-ms}$.  Since $E_p\mathcal M_s=E_p$, the one-block kernel is exactly
\begin{equation}
 P_{p,t}^\alpha
 =e^{-\alpha\omega t}\mathcal M_{(1-\alpha)\omega t}
                       +(1-e^{-\alpha\omega t})E_p.
 \label{f16v89:transition}
\end{equation}
It can be realized as Ornstein--Uhlenbeck evolution at rate
$(1-\alpha)\omega$, together with independent refreshes to $\nu_p$ at
rate $\alpha\omega$.  Its degree-$m$ multiplier, for $m\ge1$, is
$e^{-\omega[(1-\alpha)m+\alpha]t}$; constants have multiplier one.
For $m=1$ this is $e^{-\omega t}$, independently of $\alpha$.
Independent products over blocks prove \eqref{f16v89:same-two-point}.

The degree-two polynomial $x^2-1$ has variance two and multiplier
$e^{-(2-\alpha)\omega t}$.  Consequently,
\[\mathbb E[x(0)^2x(t)^2]=1+2e^{-(2-\alpha)\omega t}.\]
Subtract the three pairings determined by the unchanged two-point law to
obtain \eqref{f16v89:cumulant}.

Finite-block processes are compatible under deletion of blocks, giving
an infinite product stationary process.  Its time marginals are supported
on $H^{-1-\epsilon}$ for every $\epsilon>0$, since the expected squared
norm is a constant times $\sum_{p\ne0}\langle p\rangle^{-2-2\epsilon}/|p|$.
Equivalently define the semigroup first on cylinder functions and extend
its positive $L^2(\nu)$ contractions by density.  Given the time-zero
configuration, the stationary past and future are conditionally independent
and have reflected identical laws.  For a positive-time cylinder function
$F$, the reflection pairing is
$\mathbb E|\mathbb E(F\mid A(0))|^2\ge0$.
This proves reflection positivity and its $L^2$ extension.

Block isotropy preserves color rotations and the rotations between sine
and cosine modes induced by translations.  Frequencies depend only on
$|p|$, so cubic rotations and spatial inversion are also preserved.  No
local gauge invariance is asserted for a transverse reference measure;
just as in V88, the observable itself is gauge invariant before this input
is chosen.  These laws are not the native Wilson histories.
\end{proof}

Every measurable single-slice output, including a finite-cutoff full
simultaneous spatial-flow observable whenever defined, consequently has
\emph{the same one-time law} for all $\alpha$.  All such one-time data and
\eqref{f16v89:same-two-point} still do not determine its separated
covariance.  This is stronger input agreement than matching only the
covariance at coincident time.

\subsection{A computed composite covariance and its first input correction}

\begin{theorem}[The odd quadratic source distinguishes these inputs exactly]
\label{f16v89:quadratic}
For the bank $X_i=H_2^{\tau_i}$, the full covariance in
\eqref{f16v89:input-generator} is
\begin{equation}
 \boxed{\quad
 A_{ij}^{\alpha}(t)=\frac3V\sum_{p\ne0}|p|^4
 e^{-2(\tau_i+\tau_j)|p|^2-(2-\alpha)t|p|}
 =A_{ij}((1-\alpha/2)t).\quad}
 \label{f16v89:quadratic-C}
\end{equation}
The finite-cutoff version has the same exact formula with $|p|\le\Lambda$.
All sums and their $t$-derivatives converge at positive flow times.
On this torus the lowest visible source energy is $2-\alpha$, whereas
the linear input two-point energy at $|p|=1$ is always one.
For one flow time, its separately taken infinite-volume covariance-density
limit is
\begin{equation}
 A^\alpha(t)=\frac{3}{2\pi^2(2-\alpha)^7}
 \int_0^\infty E^6
       e^{-4\tau E^2/(2-\alpha)^2-tE}\,dE.
 \label{f16v89:infinite-C}
\end{equation}
The limiting spectral density is gapless, with
$A^\alpha(t)\sim1080[\pi^2(2-\alpha)^7]^{-1}t^{-7}$.  This is a limit of
the normalized source covariances; no local infinite-volume gauge-action
realization of the refresh process is inferred.
\end{theorem}
\begin{proof}
The spatial integral defining $H_2$ pairs only opposite momenta.  Its
contribution from one block is a centered quadratic polynomial: the trace
of the helicity operator on the transverse polarizations vanishes.
Contributions from distinct blocks are orthogonal under the unchanged
slice measure.  Each is in Hermite degree two of that block, so
\eqref{f16v89:transition} multiplies it by
$e^{-(2-\alpha)|p|t}$.  Its equal-time weight is exactly the weight
$3|p|^4e^{-2(\tau_i+\tau_j)|p|^2}/V$ already evaluated in V87 by the two
\emph{equal-time} Wick contractions.  No two-time Wick factorization is
used for $\alpha>0$.  This proves \eqref{f16v89:quadratic-C}.
The weight at $|p|=1$ is nonzero.  Gaussian spatial damping justifies the
Fourier and derivative limits.  The radial momentum integral and
$E=(2-\alpha)|p|$ give \eqref{f16v89:infinite-C}; rescaling $E$ by $t$
and using $\Gamma(7)=720$ gives its tail.
\end{proof}

The equality with a rescaled argument in \eqref{f16v89:quadratic-C} is
\emph{not} a change of physical clock.  The same clock and every linear
correlation remain fixed.  Only composite-sector relaxation has changed.
In particular a source normalization determined from one-time data cannot
repair these different time dependences.

Now let the input deformation enter at the order left unpaid in V88:
\begin{equation}
                 \alpha=\beta g^2,\qquad \beta\ge0.
 \label{f16v89:weak-input}
\end{equation}
For each finite cutoff take Taylor coefficients at $g=0$ of the covariance
of the \emph{same full nonlinear} source as in V88, under this input law.

\begin{theorem}[A two-point-invisible correction at the first nonlinear order]
\label{f16v89:matching}
Its coefficients obey
\begin{equation}
 \boxed{\quad
 C_{4,\Lambda}^{\beta}(t)=A_\Lambda(t),\qquad
 C_{6,\Lambda}^{\beta}(t)
 =4d_\Lambda A_\Lambda(t)+B_\Lambda(t)+\beta D_\Lambda(t),\quad}
 \label{f16v89:C6}
\end{equation}
where the new input term is completely explicit:
\begin{equation}
 D_{ij,\Lambda}(t)
 =\frac{3t}{V}\sum_{0<|p|\le\Lambda}|p|^5
       e^{-2(\tau_i+\tau_j)|p|^2-2t|p|}
 =-\frac t2 A'_{ij,\Lambda}(t).
 \label{f16v89:D}
\end{equation}
It has a finite cutoff limit, despite being invisible to \emph{all}
linear two-time and all one-time input laws.  For a scalar bank the limiting
normalized log-ratio coefficient is shifted by
\begin{equation}
          \beta\left\{\frac{D(t)}{A(t)}-
                              \frac{D(t_0)}{A(t_0)}\right\}.
 \label{f16v89:ratio-shift}
\end{equation}
No scalar source factor $1+z_\Lambda g^2$ can remove this shift at every
positive separation.
\end{theorem}
\begin{proof}
On product Hermite polynomials set
\[
 L_0=\sum_p\omega_pN_p,\qquad
 R=\sum_p\omega_p\{N_p-(I-E_p)\}.
\]
These commute, $0\preceq R\preceq L_0$, and $L_\alpha=L_0-\alpha R$.
Thus $\partial_\alpha P_t^\alpha|_0=tP_t^0R$.
At covariance order $g^6$, deformation of the input acts only on its
order-$g^4$ pair $H_2H_2$.  All other order-$g^6$ pairs use $P_t^0$;
odd pairs vanish by field-sign symmetry.  On a degree-two block $R$
has eigenvalue $|p|$, giving \eqref{f16v89:D}.  V88 supplies the common
$4d_\Lambda A_\Lambda$ and the finite $B_\Lambda$.
The inherited $L^2$ convergence and the absolutely convergent heat sum
justify the limits.

A scalar source factor adds $2z_\Lambda A_\Lambda$ to its sixth coefficient,
so cancels from every normalized ratio.  It cannot remove
\eqref{f16v89:ratio-shift} at all times: on the fixed torus
$D(t)/A(t)=t\langle |p|\rangle_t$ tends to zero as $t\downarrow0$ and
is asymptotic to $t$ as $t\to\infty$.  It is not constant on any open
positive-time interval, by analyticity.  The claim does not require every
chosen pair of times to distinguish every possible perturbation.
\end{proof}

There is an explicit normalized density derivative behind this correction.
For one standardized block, let $j_r(x,y)$ be its joint Gaussian density
relative to $\nu_p\otimes\nu_p$, with correlation $r=e^{-\omega t}$.
For $t>0$, the exact joint density is
\[
 j_{\alpha,t}=e^{-\alpha\omega t}j_{e^{-(1-\alpha)\omega t}}
                                      +1-e^{-\alpha\omega t}.
\]
Its derivative divided by $j_r$ is
\begin{equation}
 \ell_{2,p}(x,y)=\beta\omega t
       \{j_r^{-1}-1+r\partial_r\log j_r\}.
 \label{f16v89:score}
\end{equation}
At a finite cutoff the joint score is the sum of these block scores.
It integrates to zero, annihilates the derivative of every one-time
expectation and every linear two-time covariance, but pairs with the
quadratic-source product to give $\beta D_\Lambda(t)$.
Thus the missing input insertion is not an unnormalized probability
perturbation or an artificially changed two-point function.

\subsection{A convergent positive finite-jet model, with an explicit remainder}

The following construction is deliberately distinguished from the full
nonlinear solution.  It retains the inherited subtracted jets as actual
random variables on the common Gaussian slice space.  Put
\begin{equation}
 S_\Lambda(g)=X_\Lambda+gY_\Lambda+g^2Z_\Lambda,
 \qquad S(g)=X+gY+g^2Z.
 \label{f16v89:finite-source}
\end{equation}
Evaluate them in the genuine input process $P_t^{\beta g^2}$.

\begin{theorem}[A positive subtracted reference and a uniform finite-jet error]
\label{f16v89:positive-model}
For $0\le\beta g^2\le1/2$, the covariances of
\eqref{f16v89:finite-source} converge uniformly for $t\ge0$ and bounded
$g$ to those of $S(g)$.  The limiting process of observables is stationary
and reflection positive.  Define synthesis norms $M_2=\|X\|$,
$M_3=\|Y\|$, $M_4=\|Z\|$ in $L^2(\nu)$.  Then
\begin{equation}
 \left\|\operatorname{Cov}(S(g,0),S(g,t))
     -A(t)-g^2\{B(t)+\beta D(t)\}\right\|
 \le g^4\left\{
 \frac{8\beta^2}{e^2}M_2^2+
 \frac{2\beta}{e}(M_3^2+2M_2M_4)+M_4^2\right\}.
 \label{f16v89:remainder}
\end{equation}
The same estimate holds at finite cutoff with its own synthesis norms;
those norms are uniformly bounded by the inherited convergence.
\end{theorem}
\begin{proof}
All processes have the same slice measure.  Hence the $L^2$ norm of the
source difference is at most the corresponding sum of its three jet
errors, uniformly in the allowed input parameter.  Semigroup contraction
bounds covariance errors by the product of this norm with the bounded
source norms.  This proves uniform covariance convergence.  Reflection
positivity passes through these $L^2$ limits of same-slice observables;
no truncation of a covariance series is used to define the positive model.

By joint spectral calculus for $R,L_0$, and $0\le R\le L_0$,
\[
 \|P_t^\alpha-P_t^0\|\le\frac{\alpha}{e(1-\alpha)},\qquad
 \|P_t^\alpha-P_t^0-\alpha tP_t^0R\|
                   \le\frac{2\alpha^2}{e^2(1-\alpha)^2}.
\]
For example differentiate twice and use
$\sup_{x\ge0}x^2e^{-(1-\alpha)x}=4/[e^2(1-\alpha)^2]$;
Taylor integration supplies the factor $1/2$.
Sign symmetry eliminates the cross pairs of $Y$ with $X,Z$.  The exact
remaining covariance is
\[
 X^*P_t^\alpha X+g^2
 (Y^*P_t^\alpha Y+X^*P_t^\alpha Z+Z^*P_t^\alpha X)
                              +g^4Z^*P_t^\alpha Z.
\]
Insert the two operator bounds, set $\alpha=\beta g^2\le1/2$, and use
\eqref{f16v89:D}.  This proves \eqref{f16v89:remainder}.
\end{proof}

This is a uniform remainder for the \emph{specified finite-jet source},
not for $g^{-2}\mathcal H(A(\tau;gA_\Lambda))$.  In particular it does not
sum higher nonlinear flow orders on a running-coupling trajectory.  It
provides a legitimate positive reference completion of the computed orders,
while exhibiting the remaining dependence on the chosen joint input law.

\subsection{Normalized joint-law insertions required for native matching}

The exact input example can be related to the usual perturbative ledger
without claiming that the Wilson law has been matched to it.  At fixed
cutoff suppose a joint two-time law has the normalized expansion
\begin{equation}
 d\mu_g=\frac{e^{gL_1+g^2L_2+o(g^2)}}
                  {\mathbb E_0e^{gL_1+g^2L_2+o(g^2)}}\,d\mu_0.
 \label{f16v89:law-expansion}
\end{equation}
Assume sufficient weighted integrability to differentiate all source
products below; a merely pointwise logarithmic expansion is not sufficient.
Suppose the laws preserve spatial parity, so the odd source has zero mean,
and the expansion respects field sign with $g\mapsto-g$.  At the two times
write
\[
 Q_4=H_2(0)H_2(t),\quad
 Q_5=H_2(0)H_3(t)+H_3(0)H_2(t),
\]
\[
 Q_6=H_3(0)H_3(t)+H_2(0)H_4(t)+H_4(0)H_2(t),
\]
with outer products for a vector bank.

\begin{proposition}[The full normalized order-six input contribution]
\label{f16v89:law-ledger}
Under the stated finite-cutoff hypotheses, the sixth covariance coefficient
is
\begin{equation}
 \boxed{\quad
 \mathbb E_0Q_6+
 \operatorname{Cov}_0(Q_5,L_1)+
 \operatorname{Cov}_0(Q_4,L_2)+
 \tfrac12\operatorname{cum}_0(Q_4,L_1,L_1).\quad}
 \label{f16v89:full-input}
\end{equation}
The last three terms, denoted by $\mathcal I_\Lambda(t)$, retain the joint-law
normalization.  A scalar source correction $1+z_\Lambda g^2$ cancels V88's
logarithm at this order only if, jointly over the required time packet,
\begin{equation}
 4d_\Lambda A_\Lambda(t)+\mathcal I_\Lambda(t)
                          +2z_\Lambda A_\Lambda(t)=O(1).
 \label{f16v89:native-target}
\end{equation}
This is not a claim that such a native expansion, remainder, or choice of
$z_\Lambda$ has been established.
\end{proposition}
\begin{proof}
Expand the normalized density through order $g^2$.  Its first expectation
derivative is covariance with $L_1$; its second Taylor coefficient is
covariance with $L_2$ plus half the joint cumulant with two $L_1$ insertions.
Multiplying by $g^4Q_4+g^5Q_5+g^6Q_6$ gives
\eqref{f16v89:full-input}.  The odd coefficient vanishes by the stipulated
field-sign symmetry.  Spatial parity eliminates source means before the
covariance is formed.  Multiplying each endpoint by $1+z_\Lambda g^2$
adds $2z_\Lambda A_\Lambda$; use V88's evaluated flow contribution.
For the explicit family \eqref{f16v89:weak-input}, $L_1=0$ and
$L_2=\sum_p\ell_{2,p}$, and the sole nonzero input term is exactly
$\beta D_\Lambda(t)$ by \eqref{f16v89:score}.  Thus the formula has a
populated, normalized example, not only a symbolic insertion rule.
\end{proof}

The native Wilson two-time law and its possible gauge-coordinate
representation are not supplied by \eqref{f16v89:law-expansion}.  In
particular, neither its first two scores nor their source-weighted tails
are computed here.  The explicit examples prove that matching complete
one-time laws and every input two-point function cannot replace those
joint insertions.  The ratio cancellation proves at the same time that
an absolute common logarithm cannot by itself be read as a divergence of
the first normalized spectral coefficient.

The next native target is a common positive-time even/odd source packet
with its joint input correction or a direct nonperturbative covariance
bound.  Its original lattice source, simultaneous flow, physical clock,
and V77 reconstruction remain unchanged.  A log-free reference ratio
is not a nonzero native reflected Gram or a growing-depth upper bound.
Even the exactly soluble alternative references above remain gapless in
infinite volume.

\clearpage
\subsection{Verification, preservation and result ledger}

The companion diagnostic is
\begin{center}\small
\path{verify_ym_f16_v89_joint_matching.py}.
\end{center}
It separates exact algebra, reference quadrature and spectral checks.
These are not Wilson samples, interval certificates or proof-assistant
verification.  The unchanged V88 diagnostic was also run; this does not
recertify its analytic jet theorem.  Removing this appendix and reversing
the displayed version recovers V88 byte for byte.

\begin{record}[V89 result ledger]
\label{f16v89:ledger}
\emph{Proved independently:} the reversible reflection-positive input
family; identical complete Gaussian slices and linear time covariances,
but distinct fourth cumulants, composite covariances and visible energies;
the explicit joint score and normalized input-insertion formula.

\emph{Using the stated V88 jet input:} first-coefficient logarithmic
cancellation in normalized matrix data; the additional finite correction
$\beta D$; and a convergent positive finite-jet model with a uniform remainder.
The full nonlinear map at nonzero amplitude is not identified with that model.

\emph{Not asserted:} native input matching or counterterms, a uniform
full-flow expansion at $g^2\log\Lambda=O(1)$, a local gauge-action law for
the refresh process, or a mass gap from its finite-torus spectrum.

\emph{Still unproved:} native physical-scale source Grams, coupled input
and source control, compatible thermodynamic and continuum limits, dense
limiting coverage with common upper contraction, and the interacting
continuum Yang--Mills mass gap.
\end{record}

\begin{firewall}
The cancellation concerns Taylor coefficients.  The input examples are
positive reference laws, not Wilson histories.  All V1--V88 mathematical
text is preserved unchanged.  Neither the native physical clock nor the
lattice source is redefined.
\end{firewall}
% END F16 V89 APPEND
\end{document}
